%% major-theorems.tex -- a standalone wrapper around the generated list of
%% MAJOR THEOREMS (reference/major-theorems.tex).
%%
%% A MAJOR THEOREM is a result that appears in the user's notes as a Theorem,
%% Proposition, Lemma or Corollary.  The body of this document is GENERATED at
%% every library load by
%%
%%     (write-major-theorems-tex! "reference/major-theorems.sexp"
%%                                "reference/major-theorems.tex")
%%
%% from the hand-written table reference/major-theorems.sexp; edit the table,
%% not the generated file, and not this wrapper's body.
%%
%% Build:  pdflatex major-theorems.tex     (twice, for the labels)
%%
%% The preamble below is the manual's, cut down to what the generated fragment
%% needs.  It is COPIED rather than \input from manual.tex on purpose: the
%% manual is an amsbook whose preamble carries the index, hyperref, microtype,
%% the verbatim environments and two Kindle builds, none of which this document
%% has any use for.

\documentclass[11pt]{article}
\usepackage[T1]{fontenc}
\usepackage[utf8]{inputenc}
\usepackage{amsmath,amssymb,amsthm}
\usepackage[margin=1in]{geometry}
%% graphicx: \resizebox, which the generated \vnbstmt uses to keep a wide
%% statement inside the text width.
\usepackage{graphicx}
%% longtable: the Deviations table, which runs over a page boundary.
\usepackage{array}
\usepackage{longtable}
\usepackage[hidelinks]{hyperref}

%% The manual's theorem environments (docs/manual.tex), with the counter
%% attached to `section' rather than to `chapter' -- this is an article, and
%% one section per document of the notes is the structure the generator emits.
\newtheorem{thm}{Theorem}[section]
\newtheorem{prop}[thm]{Proposition}
\newtheorem{cor}[thm]{Corollary}
\newtheorem{lem}[thm]{Lemma}
\newtheorem{defn}[thm]{Definition}

\setlength{\parindent}{0pt}
\setlength{\parskip}{2pt}

%% The Tactics / Lemmas paragraphs are lists of \texttt names, and a typewriter
%% token neither hyphenates nor breaks at its hyphens.  \sloppy lets TeX stretch
%% the interword glue instead of running into the margin; \emergencystretch
%% gives it a last pass for the few lines \sloppy alone cannot set.  Without
%% both, every long lemma list overfills by 20--150 pt.
\setlength{\emergencystretch}{3em}
\sloppy

\title{VNB: the proven theorems}
\author{}
\date{}

\begin{document}
\maketitle

\noindent
Every result that the user's notes state as a Theorem, Proposition, Lemma or
Corollary, in the order of the notes, with the statement the VNB library makes
about it and the tactics, lemmas, oracles and bill of its proof.  Where the
library's statement does not say what the notes say, the difference is recorded
under \textbf{Deviation from the notes}; the final section tabulates those and
the results that are not proven.

\medskip
\noindent
This document is generated: the statements and proof data are read from the
library at load time, and the mapping of the notes onto it is the table
\texttt{reference/major-theorems.sexp}.

\tableofcontents

% all-theorems.tex -- GENERATED by (write-all-theorems-tex! ...) at load.
% Every proven theorem, grouped by the file that proves it, in load order.
% Do not hand-edit.
% \vnbstmt -- a displayed statement that shrinks to the text width only
% when it would overflow it, so short statements are set at normal size and
% a wide one never runs into the margin.  Needs graphicx.
%
% The leading \leavevmode is not decoration.  amsthm defers the theorem head
% (`Theorem 1.1 (Title).') into \everypar, so it is injected at the start of
% the first paragraph of the body; without \leavevmode the opening \par is a
% no-op on an empty paragraph, the head lands in the SAME paragraph as the
% box below, and head-width + \linewidth overfills every statement by 300 pt.
% \leavevmode starts a paragraph for the head to attach to, and the \par then
% closes it before the box.
\providecommand{\vnbstmt}[1]{%
  \leavevmode\par\nobreak\smallskip\noindent
  \resizebox{\ifdim\width>\linewidth\linewidth\else\width\fi}{!}{$\displaystyle #1$}%
  \par\smallskip}
\providecommand{\vnbmissing}[1]{\textit{No theorem \texttt{#1} is installed in the library.}}
% \vnbprf -- one paragraph of the proof block.  Set RAGGED RIGHT: these are
% comma-separated lists of \texttt identifiers, which neither hyphenate nor
% break at their hyphens, so justified setting either overfills the measure or
% (under \sloppy) buys the fit with very loose interword glue.
\providecommand{\vnbprf}[1]{{\raggedright #1\par}}
% \vnbpart -- one PART of a statement that the library gives under several names:
% its own paragraph, the name as a run-in label (the user, 2026-09-21).
\providecommand{\vnbpart}[1]{\par\medskip{\raggedright\noindent\textup{\textbf{Part} #1.}\par}\nopagebreak\smallskip}
\providecommand{\vnbproofpart}[1]{\par\medskip\noindent\textbf{Proof of} #1\textbf{.}\par\nopagebreak}
\providecommand{\vnbentry}[2]{\par\medskip\noindent\textbf{#1}\par\nopagebreak#2\par}
\providecommand{\vnbfile}[2]{\section*{\href{file:../#2}{\texttt{#1}}}\addcontentsline{toc}{section}{\texttt{#1}}}
\noindent 3601 theorems in 497 files, in load order.  Each entry is the statement as the library holds it, with the bill when it is not \texttt{modulo 0}; the section is the file that proves it (its path is a link).  Open the proof's graph from Emacs with \texttt{M-x vnb-graph-browse-theorem}.

\tableofcontents

\vnbfile{theorem-library/ord-well-ordered-proof.scm}{theorem-library/ord-well-ordered-proof.scm}
\noindent{\small ord-well-ordered-proof.scm -- the well-ordering of ORD, PROVEN.}\par
\vnbentry{\texttt{ord-\allowbreak{}well-\allowbreak{}ordered}}{\vnbstmt{\begin{array}{@{}l@{}} \forall l. \\ \quad (\forall x \in l.\; (x \in \mathrm{Ord}) \wedge \exists w.\; (w \in l)) \Rightarrow  \\ \quad \quad \exists m \in l.\; \forall k \in l.\; \operatorname{ord\text{-}le}(m, k) \end{array}}}

\vnbfile{theorem-library/equality-basics.scm}{theorem-library/equality-basics.scm}
\noindent{\small equality-basics.scm -- eq-sym, eq-trans, neq-sym, PROVEN.}\par
\vnbentry{\texttt{eq-\allowbreak{}sym}}{\vnbstmt{\forall a, b.\; ((a = b) \Rightarrow (b = a))}}
\vnbentry{\texttt{eq-\allowbreak{}trans}}{\vnbstmt{\forall a, b, c.\; (((a = b) \wedge (b = c)) \Rightarrow (a = c))}}
\vnbentry{\texttt{neq-\allowbreak{}sym}}{\vnbstmt{\forall a, b.\; (\neg (a = b) \Rightarrow \neg (b = a))}}

\vnbfile{theorem-library/factorial-in-nn.scm}{theorem-library/factorial-in-nn.scm}
\noindent{\small factorial-in-nn.scm -- n! is a natural number, PROVEN. Batch 27-A follow-on, 2026-09-24.}\par
\vnbentry{\texttt{factorial-\allowbreak{}in-\allowbreak{}nn}}{\vnbstmt{\forall n \in \mathbb{N}.\; (\operatorname{factorial}(n) \in \mathbb{N})}}

\vnbfile{theorem-library/subseq-apply.scm}{theorem-library/subseq-apply.scm}
\noindent{\small subseq-apply.scm -- the VALUE of a subsequence at one TYPED index.}\par
\vnbentry{\texttt{subseq-\allowbreak{}apply}}{\vnbstmt{\begin{array}{@{}l@{}} \forall f, i.; \forall k \in \mathbb{N}. \\ \quad (\operatorname{subseq}(f, i)(k) \simeq f(i(k))) \end{array}}}

\vnbfile{theorem-library/rake-compose-typing.scm}{theorem-library/rake-compose-typing.scm}
\noindent{\small rake-compose-typing.scm -- the seven RAN / composition typing supports of structure-library/compose-typing.scm, PROVEN modulo 0.}\par
\vnbentry{\texttt{range-\allowbreak{}membership}}{\vnbstmt{\begin{array}{@{}l@{}} \forall f.; \forall a \in \operatorname{dom}(f). \\ \quad (f(a) \in \operatorname{ran}(f)) \end{array}}}
\vnbentry{\texttt{fun-\allowbreak{}range-\allowbreak{}membership}}{\vnbstmt{\begin{array}{@{}l@{}} \forall x, f, a. \\ \quad ((f \in \mathrm{Fun}(x)) \wedge (a \in x)) \Rightarrow  \\ \quad \quad (f(a) \in \operatorname{ran}(f)) \end{array}}}
\vnbentry{\texttt{ran-\allowbreak{}subset-\allowbreak{}codomain}}{\vnbstmt{\forall x, y.; \forall f \in (x \to y).\; (\operatorname{ran}(f) \subseteq y)}}
\vnbentry{\texttt{compose-\allowbreak{}type-\allowbreak{}2}}{\vnbstmt{\begin{array}{@{}l@{}} \forall x, f, g, a. \\ \quad ((f \in \mathrm{Fun}(x)) \wedge ((g \in \mathrm{Fun}(\operatorname{ran}(f))) \wedge (a \in x))) \Rightarrow  \\ \quad \quad (g(f(a)) \in \operatorname{ran}(g)) \end{array}}}
\vnbentry{\texttt{compose-\allowbreak{}type-\allowbreak{}3}}{\vnbstmt{\begin{array}{@{}l@{}} \forall x, f, g, h, a. \\ \quad (f \in \mathrm{Fun}(x)) \wedge \\ \quad (g \in \mathrm{Fun}(\operatorname{ran}(f))) \wedge \\ \quad (h \in \mathrm{Fun}(\operatorname{ran}(g))) \wedge \\ \quad (a \in x) \Rightarrow \\ \quad \quad (h(g(f(a))) \in \operatorname{ran}(h)) \end{array}}}
\vnbentry{\texttt{compose-\allowbreak{}type-\allowbreak{}4}}{\vnbstmt{\begin{array}{@{}l@{}} \forall x, f, g, h, k, a. \\ \quad (f \in \mathrm{Fun}(x)) \wedge \\ \quad (g \in \mathrm{Fun}(\operatorname{ran}(f))) \wedge \\ \quad (h \in \mathrm{Fun}(\operatorname{ran}(g))) \wedge \\ \quad (k \in \mathrm{Fun}(\operatorname{ran}(h))) \wedge \\ \quad (a \in x) \Rightarrow \\ \quad \quad (k(h(g(f(a)))) \in \operatorname{ran}(k)) \end{array}}}
\vnbentry{\texttt{compose-\allowbreak{}type-\allowbreak{}5}}{\vnbstmt{\begin{array}{@{}l@{}} \forall x, f, g, h, k, m, a. \\ \quad (f \in \mathrm{Fun}(x)) \wedge \\ \quad (g \in \mathrm{Fun}(\operatorname{ran}(f))) \wedge \\ \quad (h \in \mathrm{Fun}(\operatorname{ran}(g))) \wedge \\ \quad (k \in \mathrm{Fun}(\operatorname{ran}(h))) \wedge \\ \quad (m \in \mathrm{Fun}(\operatorname{ran}(k))) \wedge \\ \quad (a \in x) \Rightarrow \\ \quad \quad (m(k(h(g(f(a))))) \in \operatorname{ran}(m)) \end{array}}}

\vnbfile{theorem-library/co-eq-lt-trans.scm}{theorem-library/co-eq-lt-trans.scm}
\noindent{\small co-eq-lt-trans.scm -- co-eq-lt-trans and co-lt-eq-trans, PROVEN.}\par
\vnbentry{\texttt{co-\allowbreak{}eq-\allowbreak{}lt-\allowbreak{}trans}}{\vnbstmt{\forall a, b, c.\; (((a = b) \wedge (b < c)) \Rightarrow (a < c))}}
\vnbentry{\texttt{co-\allowbreak{}lt-\allowbreak{}eq-\allowbreak{}trans}}{\vnbstmt{\forall a, b, c.\; (((a < b) \wedge (b = c)) \Rightarrow (a < c))}}

\vnbfile{theorem-library/interval-basics.scm}{theorem-library/interval-basics.scm}
\noindent{\small interval-basics.scm -- INTERVAL's read-offs, PROVEN from its definition.}\par
\vnbentry{\texttt{interval-\allowbreak{}unfold}}{\vnbstmt{\begin{array}{@{}l@{}} \forall a, b. \\ \quad (\{i \in \mathbb{N} : ((a \leq i) \wedge (i \leq b))\} \simeq \operatorname{interval}(a, b)) \end{array}}}
\vnbentry{\texttt{interval-\allowbreak{}in-\allowbreak{}set}}{\vnbstmt{\forall a, b.\; (\operatorname{interval}(a, b) \in \mathbf{Set})}}
\vnbentry{\texttt{interval-\allowbreak{}elt-\allowbreak{}in-\allowbreak{}nn}}{\vnbstmt{\forall a, b.; \forall i \in \operatorname{interval}(a, b).\; (i \in \mathbb{N})}}
\vnbentry{\texttt{interval-\allowbreak{}lo}}{\vnbstmt{\forall a, b.; \forall i \in \operatorname{interval}(a, b).\; (a \leq i)}}
\vnbentry{\texttt{interval-\allowbreak{}hi}}{\vnbstmt{\forall a, b.; \forall i \in \operatorname{interval}(a, b).\; (i \leq b)}}

\vnbfile{theorem-library/interval-mem-intro.scm}{theorem-library/interval-mem-intro.scm}
\noindent{\small interval-mem-intro.scm -- membership INTRODUCTION for INTERVAL, proven.}\par
\vnbentry{\texttt{interval-\allowbreak{}mem-\allowbreak{}intro}}{\vnbstmt{\begin{array}{@{}l@{}} \forall a, b.; \forall i \in \mathbb{N}. \\ \quad ((a \leq i) \wedge (i \leq b)) \Rightarrow  \\ \quad \quad (i \in \operatorname{interval}(a, b)) \end{array}}}

\vnbfile{theorem-library/ord-segment-nn-succ-proof.scm}{theorem-library/ord-segment-nn-succ-proof.scm}
\noindent{\small ord-segment-nn-succ-proof.scm -- the NN-flavoured ORD-SEGMENT successor characterisation, proven from three ordinal primitives.}\par
\vnbentry{\texttt{ord-\allowbreak{}segment-\allowbreak{}nn-\allowbreak{}succ}}{\vnbstmt{\begin{array}{@{}l@{}} \forall n \in \mathbb{N}, k. \\ \quad ((k \in \mathrm{OS}((n + 1))) \Leftrightarrow ((k \in \mathrm{OS}(n)) \vee (k = n))) \end{array}}}

\vnbfile{theorem-library/ord-segment-nn-subset-proof.scm}{theorem-library/ord-segment-nn-subset-proof.scm}
\noindent{\small ord-segment-nn-subset-proof.scm -- a member of ORD-SEGMENT(m), m in NN, is a natural number.  PROVEN, retiring the `informal' support of the same name (theorem-library/ord-segment-nn-subset.scm:9).}\par
\vnbentry{\texttt{ord-\allowbreak{}segment-\allowbreak{}zero-\allowbreak{}no-\allowbreak{}members}}{\vnbstmt{\forall k.\; \neg (k \in \mathrm{OS}(0))}}
\vnbentry{\texttt{ord-\allowbreak{}segment-\allowbreak{}nn-\allowbreak{}subset}}{\vnbstmt{\forall m \in \mathbb{N}, k \in \mathrm{OS}(m).\; (k \in \mathbb{N})}}

\vnbfile{theorem-library/inf-subsets-is-set.scm}{theorem-library/inf-subsets-is-set.scm}
\noindent{\small inf-subsets-is-set.scm -- INF-SUBSETS(A) is a set when A is, proven.}\par
\vnbentry{\texttt{inf-\allowbreak{}subsets-\allowbreak{}is-\allowbreak{}set}}{\vnbstmt{\forall a \in \mathbf{Set}.\; (\mathrm{Inf}(a) \in \mathbf{Set})}}

\vnbfile{theorem-library/bdd-metric-carrier.scm}{theorem-library/bdd-metric-carrier.scm}
\noindent{\small bdd-metric-carrier.scm -- BDD-METRIC keeps the point set, proven.}\par
\vnbentry{\texttt{bdd-\allowbreak{}metric-\allowbreak{}carrier}}{\vnbstmt{\begin{array}{@{}l@{}} \forall s. \\ \quad (\operatorname{pts}(\operatorname{bdd\text{-}metric}(s)) \simeq \operatorname{pts}(s)) \end{array}}}

\vnbfile{theorem-library/poly-membership.scm}{theorem-library/poly-membership.scm}
\noindent{\small poly-membership.scm -- SUPP / FINSUPP / POLY read-offs, PROVEN.}\par
\vnbentry{\texttt{supp-\allowbreak{}unfold}}{\vnbstmt{\begin{array}{@{}l@{}} \forall a, m, f. \\ \quad (\operatorname{supp}(a, m, f) \simeq \{x \in \operatorname{carr}(m) : \neg (f(x) = \operatorname{zero}(a))\}) \end{array}}}
\vnbentry{\texttt{supp-\allowbreak{}membership}}{\vnbstmt{\begin{array}{@{}l@{}} \forall a, m, f, x. \\ \quad ((x \in \operatorname{supp}(a, m, f)) \Leftrightarrow ((x \in \operatorname{carr}(m)) \wedge \neg (f(x) = \operatorname{zero}(a)))) \end{array}}}
\vnbentry{\texttt{supp-\allowbreak{}in-\allowbreak{}set}}{\vnbstmt{\begin{array}{@{}l@{}} \forall a, m, f. \\ \quad (\operatorname{carr}(m) \in \mathbf{Set}) \Rightarrow  \\ \quad \quad (\operatorname{supp}(a, m, f) \in \mathbf{Set}) \end{array}}}
\vnbentry{\texttt{finsupp-\allowbreak{}unfold}}{\vnbstmt{\begin{array}{@{}l@{}} \forall a, m. \\ \quad (\operatorname{finsupp}(a, m) \simeq \{f \in (\operatorname{carr}(m) \to \operatorname{carr}(a)) : (|\operatorname{supp}(a, m, f)| \in \mathbb{N})\}) \end{array}}}
\vnbentry{\texttt{finsupp-\allowbreak{}membership}}{\vnbstmt{\begin{array}{@{}l@{}} \forall a, m, f. \\ \quad ((f \in \operatorname{finsupp}(a, m)) \Leftrightarrow ((f \in (\operatorname{carr}(m) \to \operatorname{carr}(a))) \wedge (|\operatorname{supp}(a, m, f)| \in \mathbb{N}))) \end{array}}}
\vnbentry{\texttt{poly-\allowbreak{}carrier}}{\vnbstmt{\begin{array}{@{}l@{}} \forall a. \\ \quad (\operatorname{carr}(\operatorname{poly}(a)) \simeq \operatorname{finsupp}(a, \operatorname{nn\text{-}add\text{-}monoid})) \end{array}}}
\vnbentry{\texttt{poly-\allowbreak{}membership}}{\vnbstmt{\begin{array}{@{}l@{}} \forall a, f. \\ \quad ((f \in \operatorname{carr}(\operatorname{poly}(a))) \Leftrightarrow ((f \in \operatorname{carr}(a)^{\mathbb{N}}) \wedge (|\operatorname{supp}(a, \operatorname{nn\text{-}add\text{-}monoid}, f)| \in \mathbb{N}))) \end{array}}}

\vnbfile{theorem-library/normed-field-ring-view.scm}{theorem-library/normed-field-ring-view.scm}
\noindent{\small normed-field-ring-view.scm -- THE SCALAR BRIDGE, PROVEN.}\par
\vnbentry{\texttt{normed-\allowbreak{}field-\allowbreak{}ring-\allowbreak{}view-\allowbreak{}carr}}{\vnbstmt{\begin{array}{@{}l@{}} \forall f. \\ \quad (\operatorname{carr}(\operatorname{normed\text{-}field\text{-}as\text{-}commutative\text{-}ring}(f)) \simeq \operatorname{carr}(f)) \end{array}}}
\vnbentry{\texttt{normed-\allowbreak{}field-\allowbreak{}ring-\allowbreak{}view-\allowbreak{}add}}{\vnbstmt{\begin{array}{@{}l@{}} \forall f. \\ \quad (\operatorname{add}(\operatorname{normed\text{-}field\text{-}as\text{-}commutative\text{-}ring}(f)) \simeq \operatorname{add}(f)) \end{array}}}
\vnbentry{\texttt{normed-\allowbreak{}field-\allowbreak{}ring-\allowbreak{}view-\allowbreak{}mul}}{\vnbstmt{\begin{array}{@{}l@{}} \forall f. \\ \quad (\operatorname{mul}(\operatorname{normed\text{-}field\text{-}as\text{-}commutative\text{-}ring}(f)) \simeq \operatorname{mul}(f)) \end{array}}}
\vnbentry{\texttt{normed-\allowbreak{}field-\allowbreak{}ring-\allowbreak{}view-\allowbreak{}neg}}{\vnbstmt{\begin{array}{@{}l@{}} \forall f. \\ \quad (\operatorname{neg}(\operatorname{normed\text{-}field\text{-}as\text{-}commutative\text{-}ring}(f)) \simeq \operatorname{neg}(f)) \end{array}}}
\vnbentry{\texttt{normed-\allowbreak{}field-\allowbreak{}ring-\allowbreak{}view-\allowbreak{}zero}}{\vnbstmt{\begin{array}{@{}l@{}} \forall f. \\ \quad (\operatorname{zero}(\operatorname{normed\text{-}field\text{-}as\text{-}commutative\text{-}ring}(f)) \simeq \operatorname{zero}(f)) \end{array}}}
\vnbentry{\texttt{normed-\allowbreak{}field-\allowbreak{}ring-\allowbreak{}view-\allowbreak{}one}}{\vnbstmt{\begin{array}{@{}l@{}} \forall f. \\ \quad (\operatorname{one}(\operatorname{normed\text{-}field\text{-}as\text{-}commutative\text{-}ring}(f)) \simeq \operatorname{one}(f)) \end{array}}}
\vnbentry{\texttt{rr-\allowbreak{}scalar-\allowbreak{}ring-\allowbreak{}carr}}{\vnbstmt{(\operatorname{carr}(\operatorname{normed\text{-}field\text{-}as\text{-}commutative\text{-}ring}(\operatorname{rr\text{-}normed\text{-}field})) \simeq \mathbb{R})}}
\vnbentry{\texttt{rr-\allowbreak{}scalar-\allowbreak{}ring-\allowbreak{}add}}{\vnbstmt{(\operatorname{add}(\operatorname{normed\text{-}field\text{-}as\text{-}commutative\text{-}ring}(\operatorname{rr\text{-}normed\text{-}field})) \simeq (\lambda \langle \operatorname{x\_}, \operatorname{y\_} \rangle \in (\mathbb{R} \times \mathbb{R}).\; (\operatorname{x\_} + \operatorname{y\_})))}}
\vnbentry{\texttt{rr-\allowbreak{}scalar-\allowbreak{}ring-\allowbreak{}mul}}{\vnbstmt{(\operatorname{mul}(\operatorname{normed\text{-}field\text{-}as\text{-}commutative\text{-}ring}(\operatorname{rr\text{-}normed\text{-}field})) \simeq (\lambda \langle \operatorname{x\_}, \operatorname{y\_} \rangle \in (\mathbb{R} \times \mathbb{R}).\; (\operatorname{x\_} \cdot \operatorname{y\_})))}}
\vnbentry{\texttt{rr-\allowbreak{}scalar-\allowbreak{}ring-\allowbreak{}neg}}{\vnbstmt{(\operatorname{neg}(\operatorname{normed\text{-}field\text{-}as\text{-}commutative\text{-}ring}(\operatorname{rr\text{-}normed\text{-}field})) \simeq (\lambda x \in \mathbb{R}.\; -x))}}
\vnbentry{\texttt{rr-\allowbreak{}scalar-\allowbreak{}ring-\allowbreak{}zero}}{\vnbstmt{(\operatorname{zero}(\operatorname{normed\text{-}field\text{-}as\text{-}commutative\text{-}ring}(\operatorname{rr\text{-}normed\text{-}field})) \simeq 0)}}
\vnbentry{\texttt{rr-\allowbreak{}scalar-\allowbreak{}ring-\allowbreak{}one}}{\vnbstmt{(\operatorname{one}(\operatorname{normed\text{-}field\text{-}as\text{-}commutative\text{-}ring}(\operatorname{rr\text{-}normed\text{-}field})) \simeq 1)}}

\vnbfile{theorem-library/ag-view-read-offs.scm}{theorem-library/ag-view-read-offs.scm}
\noindent{\small ag-view-read-offs.scm -- THE ADDITIVE-GROUP VIEW READ-OFFS, PROVEN.}\par
\vnbentry{\texttt{ras-\allowbreak{}carr}}{\vnbstmt{\begin{array}{@{}l@{}} \forall a. \\ \quad \operatorname{is\text{-}ring}(a) \Rightarrow  \\ \quad \quad (\operatorname{carr}(\operatorname{ring\text{-}additive\text{-}ag}(a)) = \operatorname{carr}(a)) \end{array}}}
\vnbentry{\texttt{ras-\allowbreak{}op}}{\vnbstmt{\begin{array}{@{}l@{}} \forall a. \\ \quad \operatorname{is\text{-}ring}(a) \Rightarrow  \\ \quad \quad (\operatorname{opr}(\operatorname{ring\text{-}additive\text{-}ag}(a)) = \operatorname{add}(a)) \end{array}}}
\vnbentry{\texttt{ras-\allowbreak{}id}}{\vnbstmt{\begin{array}{@{}l@{}} \forall a. \\ \quad \operatorname{is\text{-}ring}(a) \Rightarrow  \\ \quad \quad (\operatorname{iden}(\operatorname{ring\text{-}additive\text{-}ag}(a)) = \operatorname{zero}(a)) \end{array}}}
\vnbentry{\texttt{mvag-\allowbreak{}carr}}{\vnbstmt{\begin{array}{@{}l@{}} \forall d. \\ \quad \operatorname{is\text{-}module}(d) \Rightarrow  \\ \quad \quad (\operatorname{carr}(\operatorname{module\text{-}vector\text{-}ag}(d)) = \operatorname{vec}(d)) \end{array}}}
\vnbentry{\texttt{mvag-\allowbreak{}op}}{\vnbstmt{\begin{array}{@{}l@{}} \forall d. \\ \quad \operatorname{is\text{-}module}(d) \Rightarrow  \\ \quad \quad (\operatorname{opr}(\operatorname{module\text{-}vector\text{-}ag}(d)) = \operatorname{vadd}(d)) \end{array}}}
\vnbentry{\texttt{mvag-\allowbreak{}id}}{\vnbstmt{\begin{array}{@{}l@{}} \forall d. \\ \quad \operatorname{is\text{-}module}(d) \Rightarrow  \\ \quad \quad (\operatorname{iden}(\operatorname{module\text{-}vector\text{-}ag}(d)) = \operatorname{vzero}(d)) \end{array}}}

\vnbfile{theorem-library/binary-minus-laws.scm}{theorem-library/binary-minus-laws.scm}
\noindent{\small binary-minus-laws.scm -- what follows from `binary-minus-def'.}\par
\vnbentry{\texttt{rr-\allowbreak{}sub-\allowbreak{}in-\allowbreak{}rr}}{\vnbstmt{\begin{array}{@{}l@{}} \forall u, v. \\ \quad ((u \in \mathbb{R}) \wedge (v \in \mathbb{R})) \Rightarrow  \\ \quad \quad ((u - v) \in \mathbb{R}) \end{array}}}
\vnbentry{\texttt{zz-\allowbreak{}sub-\allowbreak{}in-\allowbreak{}zz}}{\vnbstmt{\begin{array}{@{}l@{}} \forall u, v. \\ \quad ((u \in \mathbb{Z}) \wedge (v \in \mathbb{Z})) \Rightarrow  \\ \quad \quad ((u - v) \in \mathbb{Z}) \end{array}}}
\vnbentry{\texttt{rr-\allowbreak{}add-\allowbreak{}in-\allowbreak{}rr}}{\vnbstmt{\begin{array}{@{}l@{}} \forall a, b. \\ \quad ((a \in \mathbb{R}) \wedge (b \in \mathbb{R})) \Rightarrow  \\ \quad \quad ((a + b) \in \mathbb{R}) \end{array}}}
\vnbentry{\texttt{rr-\allowbreak{}mul-\allowbreak{}in-\allowbreak{}rr}}{\vnbstmt{\begin{array}{@{}l@{}} \forall a, b. \\ \quad ((a \in \mathbb{R}) \wedge (b \in \mathbb{R})) \Rightarrow  \\ \quad \quad ((a \cdot b) \in \mathbb{R}) \end{array}}}

\vnbfile{theorem-library/fun-apply-type-proof.scm}{theorem-library/fun-apply-type-proof.scm}
\noindent{\small fun-apply-type-proof.scm -- f : A -$>$ B and x in A give f(x) in B, PROVEN.}\par
\vnbentry{\texttt{fun-\allowbreak{}apply-\allowbreak{}type-\allowbreak{}c}}{\vnbstmt{\begin{array}{@{}l@{}} \forall f, a, b, x. \\ \quad (((f \in (a \to b)) \wedge (x \in a)) \Rightarrow (f(x) \in b)) \end{array}}}

\vnbfile{theorem-library/image-subset-codomain.scm}{theorem-library/image-subset-codomain.scm}
\noindent{\small image-subset-codomain.scm -- the image of phi : dm -$>$ cod lies in cod.}\par
\vnbentry{\texttt{image-\allowbreak{}subset-\allowbreak{}codomain}}{\vnbstmt{\begin{array}{@{}l@{}} \forall m, d.; \forall i \in (m \to d), w \in \operatorname{image}(i, m). \\ \quad (w \in d) \end{array}}}

\vnbfile{theorem-library/rake-dc-on-nn.scm}{theorem-library/rake-dc-on-nn.scm}
\noindent{\small rake-dc-on-nn.scm -- dependent choice on NN, PROVEN.}\par
\vnbentry{\texttt{dc-\allowbreak{}iter-\allowbreak{}in}}{\vnbstmt{\begin{array}{@{}l@{}} \forall n \in \mathbb{N}, s.; \forall a \in s, x. \\ \quad \forall k \in \mathbb{N}, u \in s.\; \exists y \in s.\; (y \in x(k, u)) \Rightarrow  \\ \quad \quad (\operatorname{dc\text{-}iter}(s, a, x, n) \in s) \end{array}}}
\vnbentry{\texttt{dc-\allowbreak{}iter-\allowbreak{}step}}{\vnbstmt{\begin{array}{@{}l@{}} \forall n \in \mathbb{N}, s.; \forall a \in s, x. \\ \quad \forall k \in \mathbb{N}, u \in s.\; \exists y \in s.\; (y \in x(k, u)) \Rightarrow  \\ \quad \quad \operatorname{dc\text{-}iter}(s, a, x, (n + 1)) \\ \quad \quad \in\; x(n, \operatorname{dc\text{-}iter}(s, a, x, n)) \end{array}}}
\vnbentry{\texttt{dc-\allowbreak{}on-\allowbreak{}nn-\allowbreak{}pred}}{\vnbstmt{\begin{array}{@{}l@{}} \forall x \in \mathbf{Set}, a \in x, t. \\ \quad \forall k \in \mathbb{N}, u \in x.\; \exists y \in x.\; (y \in t(k, u)) \Rightarrow  \\ \quad \quad \exists f \in (\mathbb{N} \to x). \\ \quad \quad \quad ((f(0) = a) \wedge \forall k \in \mathbb{N}.\; (f((k + 1)) \in t(k, f(k)))) \end{array}}}
\vnbentry{\texttt{dc-\allowbreak{}on-\allowbreak{}nn}}{\vnbstmt{\begin{array}{@{}l@{}} \forall x \in \mathbf{Set}, a \in x, r \in \mathbf{Set}. \\ \quad \forall k \in \mathbb{N}, u \in x.\; \exists y \in x.\; (\langle k, u, y \rangle \in r) \Rightarrow  \\ \quad \quad \exists f \in (\mathbb{N} \to x). \\ \quad \quad \quad ((f(0) = a) \wedge \forall k \in \mathbb{N}.\; (\langle k, f(k), f((k + 1)) \rangle \in r)) \end{array}}}

\vnbfile{theorem-library/compose-apply-proof.scm}{theorem-library/compose-apply-proof.scm}
\noindent{\small compose-apply-proof.scm -- (COMPOSE f g)(x) = f(g(x)), PROVEN modulo 0.}\par
\vnbentry{\texttt{compose-\allowbreak{}apply}}{\vnbstmt{\begin{array}{@{}l@{}} \forall a, b, c.; \forall f \in (b \to c), g \in (a \to b), x \in a. \\ \quad (\operatorname{compose}(f, g)(x) = f(g(x))) \end{array}}}
\vnbentry{\texttt{compose-\allowbreak{}type}}{\vnbstmt{\begin{array}{@{}l@{}} \forall a, b, c, f, g. \\ \quad ((a \in \mathbf{Set}) \wedge ((g \in (a \to b)) \wedge (f \in (b \to c)))) \Rightarrow  \\ \quad \quad (\operatorname{compose}(f, g) \in (a \to c)) \end{array}}}

\vnbfile{theorem-library/id-fun-laws.scm}{theorem-library/id-fun-laws.scm}
\noindent{\small id-fun-laws.scm -- the identity map and the composition laws over FUN.}\par
\vnbentry{\texttt{id-\allowbreak{}fun-\allowbreak{}type}}{\vnbstmt{\forall a \in \mathbf{Set}.\; (\operatorname{id\text{-}fun}(a) \in (a \to a))}}
\vnbentry{\texttt{id-\allowbreak{}fun-\allowbreak{}apply}}{\vnbstmt{\forall a.; \forall x \in a.\; (\operatorname{id\text{-}fun}(a)(x) = x)}}
\vnbentry{\texttt{compose-\allowbreak{}id-\allowbreak{}left}}{\vnbstmt{\begin{array}{@{}l@{}} \forall a, b, f. \\ \quad ((a \in \mathbf{Set}) \wedge ((b \in \mathbf{Set}) \wedge (f \in (a \to b)))) \Rightarrow  \\ \quad \quad (\operatorname{compose}(\operatorname{id\text{-}fun}(b), f) = f) \end{array}}}
\vnbentry{\texttt{compose-\allowbreak{}id-\allowbreak{}right}}{\vnbstmt{\begin{array}{@{}l@{}} \forall a, b, f. \\ \quad ((a \in \mathbf{Set}) \wedge (f \in (a \to b))) \Rightarrow  \\ \quad \quad (\operatorname{compose}(f, \operatorname{id\text{-}fun}(a)) = f) \end{array}}}
\vnbentry{\texttt{compose-\allowbreak{}assoc}}{\vnbstmt{\begin{array}{@{}l@{}} \forall a, b, c, u, f, g, h. \\ \quad (a \in \mathbf{Set}) \wedge \\ \quad (b \in \mathbf{Set}) \wedge \\ \quad (f \in (a \to b)) \wedge \\ \quad (g \in (b \to c)) \wedge \\ \quad (h \in (c \to u)) \Rightarrow \\ \quad \quad \operatorname{compose}(h, \operatorname{compose}(g, f)) \\ \quad \quad =\; \operatorname{compose}(\operatorname{compose}(h, g), f) \end{array}}}

\vnbfile{theorem-library/pair-tuple-sethood.scm}{theorem-library/pair-tuple-sethood.scm}
\noindent{\small pair-tuple-sethood.scm -- a pair [u,v] is a set, PROVEN modulo 0.}\par
\vnbentry{\texttt{pair-\allowbreak{}in-\allowbreak{}cartesian}}{\vnbstmt{\begin{array}{@{}l@{}} \forall x, y, u, v. \\ \quad ((u \in x) \wedge (v \in y)) \Rightarrow  \\ \quad \quad (\langle u, v \rangle \in (x \times y)) \end{array}}}
\vnbentry{\texttt{pair-\allowbreak{}tuple-\allowbreak{}is-\allowbreak{}set}}{\vnbstmt{\begin{array}{@{}l@{}} \forall x, y, u, v. \\ \quad ((u \in x) \wedge (v \in y)) \Rightarrow  \\ \quad \quad (\langle u, v \rangle \in \mathbf{Set}) \end{array}}}

\vnbfile{theorem-library/rake-setoid.scm}{theorem-library/rake-setoid.scm}
\noindent{\small rake-setoid.scm -- BATCH F of the 2026-09-17 rake: setoids and carriers. Seven asserted leaves of the PSS, PROVEN.  Every statement here is its support's statement UNCHANGED.}\par
\vnbentry{\texttt{class-\allowbreak{}unfold}}{\vnbstmt{\begin{array}{@{}l@{}} \forall s, a. \\ \quad (\operatorname{class}(s, a) \simeq \{b \in \operatorname{pts}(s) : \operatorname{related}(s, a, b)\}) \end{array}}}
\vnbentry{\texttt{proj-\allowbreak{}unfold}}{\vnbstmt{\begin{array}{@{}l@{}} \forall s. \\ \quad (\operatorname{proj}(s) \simeq (\lambda a \in \operatorname{pts}(s).\; \operatorname{class}(s, a))) \end{array}}}
\vnbentry{\texttt{quotient-\allowbreak{}unfold}}{\vnbstmt{\begin{array}{@{}l@{}} \forall s. \\ \quad (\operatorname{quotient}(s) \simeq \operatorname{image}(\operatorname{proj}(s), \operatorname{pts}(s))) \end{array}}}
\vnbentry{\texttt{respects-\allowbreak{}unfold}}{\vnbstmt{\begin{array}{@{}l@{}} \forall s, f. \\ \quad (\operatorname{respects}(s, f) \Leftrightarrow \forall a, b \in \operatorname{pts}(s).\; (\operatorname{related}(s, a, b) \Rightarrow (f(a) = f(b)))) \end{array}}}
\vnbentry{\texttt{respects2-\allowbreak{}unfold}}{\vnbstmt{\begin{array}{@{}l@{}} \forall s, f. \\ \quad (\operatorname{respects2}(s, f) \Leftrightarrow \forall a, b, c, d.\; (((a \in \operatorname{pts}(s)) \wedge ((b \in \operatorname{pts}(s)) \wedge ((c \in \operatorname{pts}(s)) \wedge ((d \in \operatorname{pts}(s)) \wedge (\operatorname{related}(s, a, c) \wedge \operatorname{related}(s, b, d)))))) \Rightarrow (f(a, b) = f(c, d)))) \end{array}}}
\vnbentry{\texttt{is-\allowbreak{}vector-\allowbreak{}space-\allowbreak{}unfold}}{\vnbstmt{\begin{array}{@{}l@{}} \forall m. \\ \quad (\operatorname{is\text{-}vector\text{-}space}(m) \Leftrightarrow (\operatorname{is\text{-}module}(m) \wedge \operatorname{is\text{-}field\text{-}ring}(\operatorname{scal}(m)))) \end{array}}}
\vnbentry{\texttt{class-\allowbreak{}is-\allowbreak{}set}}{\vnbstmt{\begin{array}{@{}l@{}} \forall s.; \forall a \in \operatorname{pts}(s). \\ \quad \operatorname{is\text{-}setoid}(s) \Rightarrow  \\ \quad \quad (\operatorname{class}(s, a) \in \mathbf{Set}) \end{array}}}
\vnbentry{\texttt{quotient-\allowbreak{}is-\allowbreak{}set}}{\vnbstmt{\begin{array}{@{}l@{}} \forall s. \\ \quad \operatorname{is\text{-}setoid}(s) \Rightarrow  \\ \quad \quad (\operatorname{quotient}(s) \in \mathbf{Set}) \end{array}}}
\vnbentry{\texttt{proj-\allowbreak{}in-\allowbreak{}fun}}{\vnbstmt{\begin{array}{@{}l@{}} \forall s. \\ \quad \operatorname{is\text{-}setoid}(s) \Rightarrow  \\ \quad \quad (\operatorname{proj}(s) \in (\operatorname{pts}(s) \to \operatorname{quotient}(s))) \end{array}}}
\vnbentry{\texttt{descend-\allowbreak{}in-\allowbreak{}fun}}{\vnbstmt{\begin{array}{@{}l@{}} \forall s, z.; \forall f \in (\operatorname{pts}(s) \to z). \\ \quad (\operatorname{is\text{-}setoid}(s) \wedge \operatorname{respects}(s, f)) \Rightarrow  \\ \quad \quad (\operatorname{descend}(s, f) \in (\operatorname{quotient}(s) \to z)) \end{array}}}
\vnbentry{\texttt{descend2-\allowbreak{}in-\allowbreak{}fun}}{\vnbstmt{\begin{array}{@{}l@{}} \forall s, z.; \forall f \in ((\operatorname{pts}(s) \times \operatorname{pts}(s)) \to z). \\ \quad (\operatorname{is\text{-}setoid}(s) \wedge \operatorname{respects2}(s, f)) \Rightarrow  \\ \quad \quad \operatorname{descend2}(s, f) \\ \quad \quad \in\; ((\operatorname{quotient}(s) \times \operatorname{quotient}(s)) \to z) \end{array}}}
\vnbentry{\texttt{vec-\allowbreak{}is-\allowbreak{}set}}{\vnbstmt{\begin{array}{@{}l@{}} \forall m. \\ \quad \operatorname{is\text{-}normed\text{-}vector\text{-}space}(m) \Rightarrow  \\ \quad \quad (\operatorname{vec}(m) \in \mathbf{Set}) \end{array}}}
\vnbentry{\texttt{vspace-\allowbreak{}vec-\allowbreak{}is-\allowbreak{}set}}{\vnbstmt{\begin{array}{@{}l@{}} \forall m. \\ \quad \operatorname{is\text{-}vector\text{-}space}(m) \Rightarrow  \\ \quad \quad (\operatorname{vec}(m) \in \mathbf{Set}) \end{array}}}

\vnbfile{theorem-library/nn-order-ord.scm}{theorem-library/nn-order-ord.scm}
\noindent{\small nn-order-ord.scm -- the elementary NN order block, PROVEN.  Seven facts that were `well-known' supports in structure-library/order-lemmas.scm are theorems here, every one billing `modulo 0'.}\par
\vnbentry{\texttt{nn-\allowbreak{}one-\allowbreak{}in}}{\vnbstmt{(1 \in \mathbb{N})}}
\vnbentry{\texttt{nn-\allowbreak{}zero-\allowbreak{}le}}{\vnbstmt{\forall n \in \mathbb{N}.\; (0 \leq n)}}
\vnbentry{\texttt{nn-\allowbreak{}le-\allowbreak{}succ}}{\vnbstmt{\forall k \in \mathbb{N}.\; (k \leq (k + 1))}}
\vnbentry{\texttt{nn-\allowbreak{}le-\allowbreak{}imp-\allowbreak{}neq-\allowbreak{}succ}}{\vnbstmt{\forall k, j \in \mathbb{N}.\; ((j \leq k) \Rightarrow \neg (j = (k + 1)))}}
\vnbentry{\texttt{nn-\allowbreak{}succ-\allowbreak{}mono}}{\vnbstmt{\forall a, b \in \mathbb{N}.\; ((a \leq b) \Rightarrow ((a + 1) \leq (b + 1)))}}
\vnbentry{\texttt{nn-\allowbreak{}one-\allowbreak{}le-\allowbreak{}succ}}{\vnbstmt{\forall n \in \mathbb{N}.\; (1 \leq (n + 1))}}
\vnbentry{\texttt{nn-\allowbreak{}le-\allowbreak{}succ-\allowbreak{}cases}}{\vnbstmt{\begin{array}{@{}l@{}} \forall k, j \in \mathbb{N}. \\ \quad ((j \leq (k + 1)) \Rightarrow ((j \leq k) \vee (j = (k + 1)))) \end{array}}}

\vnbfile{theorem-library/nn-not-le-succ-le.scm}{theorem-library/nn-not-le-succ-le.scm}
\noindent{\small nn-not-le-succ-le.scm -- NN is totally ordered and discrete, PROVEN.}\par
\vnbentry{\texttt{nn-\allowbreak{}not-\allowbreak{}le-\allowbreak{}succ-\allowbreak{}le}}{\vnbstmt{\forall m, n \in \mathbb{N}.\; (\neg (m \leq n) \Rightarrow ((n + 1) \leq m))}}

\vnbfile{theorem-library/nn-order-basics.scm}{theorem-library/nn-order-basics.scm}
\noindent{\small nn-order-basics.scm -- the elementary NN order facts, PROVEN.}\par
\vnbentry{\texttt{nn-\allowbreak{}in-\allowbreak{}rr}}{\vnbstmt{\forall k \in \mathbb{N}.\; (k \in \mathbb{R})}}
\vnbentry{\texttt{rr-\allowbreak{}le-\allowbreak{}trans}}{\vnbstmt{\begin{array}{@{}l@{}} \forall x, y, z \in \mathbb{R}. \\ \quad (((x \leq y) \wedge (y \leq z)) \Rightarrow (x \leq z)) \end{array}}}
\vnbentry{\texttt{rr-\allowbreak{}le-\allowbreak{}trans-\allowbreak{}c}}{\vnbstmt{\begin{array}{@{}l@{}} \forall x, y, z \in \mathbb{R}. \\ \quad (((x \leq y) \wedge (y \leq z)) \Rightarrow (x \leq z)) \end{array}}}
\vnbentry{\texttt{nn-\allowbreak{}le-\allowbreak{}trans-\allowbreak{}guarded}}{\vnbstmt{\begin{array}{@{}l@{}} \forall a, b, c \in \mathbb{N}. \\ \quad (((a \leq b) \wedge (b \leq c)) \Rightarrow (a \leq c)) \end{array}}}
\vnbentry{\texttt{nn-\allowbreak{}le-\allowbreak{}refl}}{\vnbstmt{\forall k \in \mathbb{N}.\; (k \leq k)}}
\vnbentry{\texttt{nn-\allowbreak{}le-\allowbreak{}add-\allowbreak{}right}}{\vnbstmt{\forall n, m \in \mathbb{N}.\; (m \leq (m + n))}}
\vnbentry{\texttt{nn-\allowbreak{}pair-\allowbreak{}upper-\allowbreak{}bound}}{\vnbstmt{\begin{array}{@{}l@{}} \forall a, b. \\ \quad ((a \in \mathbb{N}) \wedge (b \in \mathbb{N})) \Rightarrow  \\ \quad \quad \exists c \in \mathbb{N}.\; ((a \leq c) \wedge (b \leq c)) \end{array}}}
\vnbentry{\texttt{zz-\allowbreak{}in-\allowbreak{}rr}}{\vnbstmt{\forall k \in \mathbb{Z}.\; (k \in \mathbb{R})}}

\vnbfile{theorem-library/mat-basics.scm}{theorem-library/mat-basics.scm}
\noindent{\small mat-basics.scm -- MAT's read-offs, PROVEN from its definition.}\par
\vnbentry{\texttt{mat-\allowbreak{}unfold}}{\vnbstmt{\begin{array}{@{}l@{}} \forall m, n, x. \\ \quad (\{p \in \operatorname{matrix}(x) : ((n \in \mathbb{N}) \wedge (((\operatorname{length}(p) = 0) \Rightarrow (m = 0)) \wedge (\neg (\operatorname{length}(p) = 0) \Rightarrow (\operatorname{size}(p) = \langle m, n \rangle))))\} \simeq \operatorname{mat}(m, n, x)) \end{array}}}
\vnbentry{\texttt{size-\allowbreak{}unfold}}{\vnbstmt{\begin{array}{@{}l@{}} \forall m. \\ \quad (\langle \operatorname{length}(m), \left(\text{if }(\operatorname{length}(m) = 0)\text{ then }0\text{ else }\operatorname{length}(\operatorname{nth}(1, m))\right) \rangle \simeq \operatorname{size}(m)) \end{array}}}
\vnbentry{\texttt{size-\allowbreak{}nonempty}}{\vnbstmt{\begin{array}{@{}l@{}} \forall m. \\ \quad \neg (\operatorname{length}(m) = 0) \Rightarrow  \\ \quad \quad (\langle \operatorname{length}(m), \operatorname{length}(\operatorname{nth}(1, m)) \rangle \simeq \operatorname{size}(m)) \end{array}}}
\vnbentry{\texttt{nth1-\allowbreak{}pair}}{\vnbstmt{\forall a, b.\; (\operatorname{nth}(1, \langle a, b \rangle) \simeq a)}}
\vnbentry{\texttt{mat-\allowbreak{}in-\allowbreak{}matrix}}{\vnbstmt{\begin{array}{@{}l@{}} \forall m, n, x.; \forall q \in \operatorname{mat}(m, n, x). \\ \quad (q \in \operatorname{matrix}(x)) \end{array}}}
\vnbentry{\texttt{mat-\allowbreak{}cols-\allowbreak{}in-\allowbreak{}nn}}{\vnbstmt{\begin{array}{@{}l@{}} \forall m, n, x.; \forall q \in \operatorname{mat}(m, n, x). \\ \quad (n \in \mathbb{N}) \end{array}}}
\vnbentry{\texttt{mat-\allowbreak{}length}}{\vnbstmt{\begin{array}{@{}l@{}} \forall m, n, x.; \forall q \in \operatorname{mat}(m, n, x). \\ \quad (\operatorname{length}(q) = m) \end{array}}}
\vnbentry{\texttt{mat-\allowbreak{}rows-\allowbreak{}in-\allowbreak{}nn}}{\vnbstmt{\begin{array}{@{}l@{}} \forall m, n, x.; \forall q \in \operatorname{mat}(m, n, x). \\ \quad (m \in \mathbb{N}) \end{array}}}
\vnbentry{\texttt{mat-\allowbreak{}size}}{\vnbstmt{\begin{array}{@{}l@{}} \forall m, n, x.; \forall q \in \operatorname{mat}(m, n, x). \\ \quad ((1 \leq m) \Rightarrow (\operatorname{size}(q) = \langle m, n \rangle)) \end{array}}}
\vnbentry{\texttt{mat-\allowbreak{}cols}}{\vnbstmt{\begin{array}{@{}l@{}} \forall m, n, x.; \forall q \in \operatorname{mat}(m, n, x). \\ \quad ((1 \leq m) \Rightarrow (\operatorname{length}(\operatorname{nth}(1, q)) = n)) \end{array}}}
\vnbentry{\texttt{mat-\allowbreak{}size-\allowbreak{}rows}}{\vnbstmt{\begin{array}{@{}l@{}} \forall m, n, x.; \forall q \in \operatorname{mat}(m, n, x). \\ \quad (\operatorname{nth}(1, \operatorname{size}(q)) = m) \end{array}}}
\vnbentry{\texttt{mat-\allowbreak{}size-\allowbreak{}cols}}{\vnbstmt{\begin{array}{@{}l@{}} \forall m, n, x.; \forall q \in \operatorname{mat}(m, n, x). \\ \quad ((1 \leq m) \Rightarrow (\operatorname{nth}(2, \operatorname{size}(q)) = n)) \end{array}}}
\vnbentry{\texttt{mat-\allowbreak{}intro}}{\vnbstmt{\begin{array}{@{}l@{}} \forall m, n, x.; \forall q \in \operatorname{matrix}(x). \\ \quad (n \in \mathbb{N}) \wedge \\ \quad (1 \leq m) \wedge \\ \quad (\operatorname{length}(q) = m) \wedge \\ \quad (\operatorname{length}(\operatorname{nth}(1, q)) = n) \Rightarrow \\ \quad \quad (q \in \operatorname{mat}(m, n, x)) \end{array}}}
\vnbentry{\texttt{mat-\allowbreak{}intro-\allowbreak{}empty}}{\vnbstmt{\begin{array}{@{}l@{}} \forall n, x.; \forall q \in \operatorname{matrix}(x). \\ \quad ((n \in \mathbb{N}) \wedge (\operatorname{length}(q) = 0)) \Rightarrow  \\ \quad \quad (q \in \operatorname{mat}(0, n, x)) \end{array}}}
\vnbentry{\texttt{nil-\allowbreak{}in-\allowbreak{}mat}}{\vnbstmt{\begin{array}{@{}l@{}} \forall n \in \mathbb{N}, x. \\ \quad (\langle  \rangle \in \operatorname{mat}(0, n, x)) \end{array}}}
\vnbentry{\texttt{mat-\allowbreak{}size-\allowbreak{}cols-\allowbreak{}in-\allowbreak{}nn}}{\vnbstmt{\begin{array}{@{}l@{}} \forall m, n, x.; \forall q \in \operatorname{mat}(m, n, x). \\ \quad (\operatorname{nth}(2, \operatorname{size}(q)) \in \mathbb{N}) \end{array}}}
\vnbentry{\texttt{mat-\allowbreak{}size-\allowbreak{}cols-\allowbreak{}pos}}{\vnbstmt{\begin{array}{@{}l@{}} \forall m, n, x.; \forall q \in \operatorname{mat}(m, n, x). \\ \quad (1 \leq \operatorname{nth}(2, \operatorname{size}(q))) \Rightarrow  \\ \quad \quad (\operatorname{nth}(2, \operatorname{size}(q)) = n) \end{array}}}
\vnbentry{\texttt{mat-\allowbreak{}empty-\allowbreak{}any-\allowbreak{}cols}}{\vnbstmt{\begin{array}{@{}l@{}} \forall c, x, q, k. \\ \quad ((q \in \operatorname{mat}(0, c, x)) \wedge (k \in \mathbb{N})) \Rightarrow  \\ \quad \quad (q \in \operatorname{mat}(0, k, x)) \end{array}}}
\vnbentry{\texttt{interval-\allowbreak{}1-\allowbreak{}0-\allowbreak{}empty}}{\vnbstmt{(\operatorname{interval}(1, 0) = \emptyset)}}

\vnbfile{theorem-library/nn-pos-of-nonzero.scm}{theorem-library/nn-pos-of-nonzero.scm}
\noindent{\small nn-pos-of-nonzero.scm -- a nonzero natural is strictly positive.}\par
\vnbentry{\texttt{nn-\allowbreak{}pos-\allowbreak{}of-\allowbreak{}nonzero}}{\vnbstmt{\forall k \in \mathbb{N}.\; (\neg (k = 0) \Rightarrow (0 < k))}}

\vnbfile{theorem-library/interval-membership.scm}{theorem-library/interval-membership.scm}
\noindent{\small interval-membership.scm -- INTERVAL's membership IFF and the [1,1] instance, both PROVEN, both by citation of the read-offs already in the tree.}\par
\vnbentry{\texttt{interval-\allowbreak{}membership}}{\vnbstmt{\begin{array}{@{}l@{}} \forall a, b, i. \\ \quad ((i \in \operatorname{interval}(a, b)) \Leftrightarrow ((i \in \mathbb{N}) \wedge ((a \leq i) \wedge (i \leq b)))) \end{array}}}
\vnbentry{\texttt{one-\allowbreak{}in-\allowbreak{}interval-\allowbreak{}1}}{\vnbstmt{(1 \in \operatorname{interval}(1, 1))}}

\vnbfile{theorem-library/nn-not-lt-le-proof.scm}{theorem-library/nn-not-lt-le-proof.scm}
\noindent{\small nn-not-lt-le-proof.scm -- NOT (k $<$ j) =$>$ j $<$= k on NN, `modulo 0'.}\par
\vnbentry{\texttt{nn-\allowbreak{}not-\allowbreak{}lt-\allowbreak{}le}}{\vnbstmt{\forall k, j \in \mathbb{N}.\; (\neg (k < j) \Rightarrow (j \leq k))}}

\vnbfile{theorem-library/entry-in-carrier.scm}{theorem-library/entry-in-carrier.scm}
\noindent{\small entry-in-carrier.scm -- the entries of an m-by-n matrix over X lie in X.}\par
\vnbentry{\texttt{entry-\allowbreak{}in-\allowbreak{}carrier}}{\vnbstmt{\begin{array}{@{}l@{}} \forall m, n, x, p, i, j. \\ \quad (p \in \operatorname{mat}(m, n, x)) \wedge \\ \quad (i \in \operatorname{interval}(1, m)) \wedge \\ \quad (j \in \operatorname{interval}(1, n)) \Rightarrow \\ \quad \quad (\operatorname{entry}(p, i, j) \in x) \end{array}}}

\vnbfile{theorem-library/bt-shims.scm}{theorem-library/bt-shims.scm}
\noindent{\small bt-shims.scm -- the five "Bernstein typing" shims of theorem-library/binomial.scm, PROVEN.}\par
\vnbentry{\texttt{bt-\allowbreak{}nn-\allowbreak{}in-\allowbreak{}zz}}{\vnbstmt{\forall n \in \mathbb{N}.\; (n \in \mathbb{Z})}}
\vnbentry{\texttt{bt-\allowbreak{}succ-\allowbreak{}in-\allowbreak{}nn}}{\vnbstmt{\forall n \in \mathbb{N}.\; ((n + 1) \in \mathbb{N})}}
\vnbentry{\texttt{bt-\allowbreak{}neg1-\allowbreak{}in-\allowbreak{}zz}}{\vnbstmt{((0 - 1) \in \mathbb{Z})}}
\vnbentry{\texttt{bt-\allowbreak{}lt-\allowbreak{}succ}}{\vnbstmt{\forall n \in \mathbb{N}.\; (n < (n + 1))}}
\vnbentry{\texttt{bt-\allowbreak{}neg1-\allowbreak{}neg}}{\vnbstmt{((0 - 1) < 0)}}

\vnbfile{theorem-library/rr-order-basics.scm}{theorem-library/rr-order-basics.scm}
\noindent{\small rr-order-basics.scm -- the elementary RR order facts, PROVEN.}\par
\vnbentry{\texttt{rr-\allowbreak{}lt-\allowbreak{}trans}}{\vnbstmt{\forall x, y, z \in \mathbb{R}.\; (((x < y) \wedge (y < z)) \Rightarrow (x < z))}}
\vnbentry{\texttt{rr-\allowbreak{}lt-\allowbreak{}le-\allowbreak{}trans}}{\vnbstmt{\begin{array}{@{}l@{}} \forall x, y, z \in \mathbb{R}. \\ \quad (((x < y) \wedge (y \leq z)) \Rightarrow (x < z)) \end{array}}}
\vnbentry{\texttt{rr-\allowbreak{}le-\allowbreak{}lt-\allowbreak{}trans}}{\vnbstmt{\begin{array}{@{}l@{}} \forall x, y, z \in \mathbb{R}. \\ \quad (((x \leq y) \wedge (y < z)) \Rightarrow (x < z)) \end{array}}}
\vnbentry{\texttt{rr-\allowbreak{}lt-\allowbreak{}implies-\allowbreak{}le}}{\vnbstmt{\forall x, y \in \mathbb{R}.\; ((x < y) \Rightarrow (x \leq y))}}
\vnbentry{\texttt{rr-\allowbreak{}le-\allowbreak{}diff-\allowbreak{}nonpos}}{\vnbstmt{\forall u, v \in \mathbb{R}.\; ((u \leq v) \Rightarrow ((u - v) \leq 0))}}
\vnbentry{\texttt{rr-\allowbreak{}lt-\allowbreak{}diff-\allowbreak{}pos}}{\vnbstmt{\forall u, v \in \mathbb{R}.\; ((u < v) \Rightarrow (0 < (v - u)))}}
\vnbentry{\texttt{rr-\allowbreak{}lt-\allowbreak{}diff-\allowbreak{}neg}}{\vnbstmt{\forall u, v \in \mathbb{R}.\; ((u < v) \Rightarrow ((u - v) < 0))}}
\vnbentry{\texttt{rr-\allowbreak{}diff-\allowbreak{}zero-\allowbreak{}eq}}{\vnbstmt{\forall u, v \in \mathbb{R}.\; ((0 = (u - v)) \Rightarrow (u = v))}}
\vnbentry{\texttt{rr-\allowbreak{}le-\allowbreak{}neg}}{\vnbstmt{\forall u, v \in \mathbb{R}.\; ((u \leq v) \Rightarrow (-v \leq -u))}}
\vnbentry{\texttt{rr-\allowbreak{}neg-\allowbreak{}eq-\allowbreak{}zero}}{\vnbstmt{\forall u \in \mathbb{R}.\; ((-u = 0) \Rightarrow (u = 0))}}
\vnbentry{\texttt{rr-\allowbreak{}sub-\allowbreak{}ne-\allowbreak{}zero}}{\vnbstmt{\forall u, v \in \mathbb{R}.\; (\neg (u = v) \Rightarrow \neg ((u - v) = 0))}}
\vnbentry{\texttt{rr-\allowbreak{}pos-\allowbreak{}ne-\allowbreak{}zero}}{\vnbstmt{\forall c \in \mathbb{R}.\; ((0 < c) \Rightarrow \neg (c = 0))}}
\vnbentry{\texttt{rr-\allowbreak{}lt-\allowbreak{}irrefl}}{\vnbstmt{\forall a \in \mathbb{R}.\; \neg (a < a)}}
\vnbentry{\texttt{rr-\allowbreak{}le-\allowbreak{}cases}}{\vnbstmt{\forall u, v \in \mathbb{R}.\; ((u \leq v) \Rightarrow ((u < v) \vee (u = v)))}}
\vnbentry{\texttt{rr-\allowbreak{}le-\allowbreak{}ne-\allowbreak{}lt}}{\vnbstmt{\begin{array}{@{}l@{}} \forall u, v \in \mathbb{R}. \\ \quad (((u \leq v) \wedge \neg (u = v)) \Rightarrow (u < v)) \end{array}}}
\vnbentry{\texttt{rr-\allowbreak{}le-\allowbreak{}add}}{\vnbstmt{\begin{array}{@{}l@{}} \forall x, y, u, v. \\ \quad (x \in \mathbb{R}) \wedge \\ \quad (y \in \mathbb{R}) \wedge \\ \quad (u \in \mathbb{R}) \wedge \\ \quad (v \in \mathbb{R}) \wedge \\ \quad (x \leq y) \wedge \\ \quad (u \leq v) \Rightarrow \\ \quad \quad ((x + u) \leq (y + v)) \end{array}}}
\vnbentry{\texttt{rr-\allowbreak{}lt-\allowbreak{}add}}{\vnbstmt{\begin{array}{@{}l@{}} \forall x, y, u, v. \\ \quad (x \in \mathbb{R}) \wedge \\ \quad (y \in \mathbb{R}) \wedge \\ \quad (u \in \mathbb{R}) \wedge \\ \quad (v \in \mathbb{R}) \wedge \\ \quad (x < y) \wedge \\ \quad (u \leq v) \Rightarrow \\ \quad \quad ((x + u) < (y + v)) \end{array}}}
\vnbentry{\texttt{rr-\allowbreak{}add-\allowbreak{}nonneg}}{\vnbstmt{\begin{array}{@{}l@{}} \forall x, y. \\ \quad ((x \in \mathbb{R}) \wedge ((y \in \mathbb{R}) \wedge ((0 \leq x) \wedge (0 \leq y)))) \Rightarrow  \\ \quad \quad (0 \leq (x + y)) \end{array}}}
\vnbentry{\texttt{rr-\allowbreak{}le-\allowbreak{}from-\allowbreak{}diff-\allowbreak{}nonneg}}{\vnbstmt{\begin{array}{@{}l@{}} \forall x, y. \\ \quad ((x \in \mathbb{R}) \wedge ((y \in \mathbb{R}) \wedge (0 \leq (y - x)))) \Rightarrow  \\ \quad \quad (x \leq y) \end{array}}}
\vnbentry{\texttt{rr-\allowbreak{}double-\allowbreak{}nonneg}}{\vnbstmt{\forall x \in \mathbb{R}.\; ((0 \leq (x + x)) \Rightarrow (0 \leq x))}}
\vnbentry{\texttt{rr-\allowbreak{}no-\allowbreak{}zero-\allowbreak{}divisors}}{\vnbstmt{\begin{array}{@{}l@{}} \forall a, b. \\ \quad ((a \in \mathbb{R}) \wedge ((b \in \mathbb{R}) \wedge ((a \cdot b) = 0))) \Rightarrow  \\ \quad \quad ((a = 0) \vee (b = 0)) \end{array}}}
\vnbentry{\texttt{rr-\allowbreak{}cancel-\allowbreak{}mul-\allowbreak{}right}}{\vnbstmt{\begin{array}{@{}l@{}} \forall u, v, c. \\ \quad (u \in \mathbb{R}) \wedge \\ \quad (v \in \mathbb{R}) \wedge \\ \quad (c \in \mathbb{R}) \wedge \\ \quad \neg (c = 0) \wedge \\ \quad ((u \cdot c) = (v \cdot c)) \Rightarrow \\ \quad \quad (u = v) \end{array}}}
\vnbentry{\texttt{rr-\allowbreak{}lt-\allowbreak{}scale-\allowbreak{}pos}}{\vnbstmt{\begin{array}{@{}l@{}} \forall c, x, y. \\ \quad (c \in \mathbb{R}) \wedge \\ \quad (x \in \mathbb{R}) \wedge \\ \quad (y \in \mathbb{R}) \wedge \\ \quad (0 < c) \wedge \\ \quad (x < y) \Rightarrow \\ \quad \quad ((c \cdot x) < (c \cdot y)) \end{array}}}
\vnbentry{\texttt{rr-\allowbreak{}le-\allowbreak{}scale-\allowbreak{}nonneg}}{\vnbstmt{\begin{array}{@{}l@{}} \forall c, x, y. \\ \quad (c \in \mathbb{R}) \wedge \\ \quad (x \in \mathbb{R}) \wedge \\ \quad (y \in \mathbb{R}) \wedge \\ \quad (0 \leq c) \wedge \\ \quad (x \leq y) \Rightarrow \\ \quad \quad ((c \cdot x) \leq (c \cdot y)) \end{array}}}
\vnbentry{\texttt{rr-\allowbreak{}sq-\allowbreak{}nonneg}}{\vnbstmt{\forall x \in \mathbb{R}.\; (0 \leq (x \cdot x))}}
\vnbentry{\texttt{rr-\allowbreak{}lt-\allowbreak{}trichotomy}}{\vnbstmt{\forall u, v \in \mathbb{R}.\; ((u < v) \vee ((u = v) \vee (v < u)))}}
\vnbentry{\texttt{rr-\allowbreak{}upper-\allowbreak{}of-\allowbreak{}two}}{\vnbstmt{\begin{array}{@{}l@{}} \forall u, v. \\ \quad ((u \in \mathbb{R}) \wedge (v \in \mathbb{R})) \Rightarrow  \\ \quad \quad \exists w \in \mathbb{R}.\; ((u \leq w) \wedge (v \leq w)) \end{array}}}
\vnbentry{\texttt{rr-\allowbreak{}upper-\allowbreak{}of-\allowbreak{}two-\allowbreak{}below}}{\vnbstmt{\begin{array}{@{}l@{}} \forall u, v, m. \\ \quad (u \in \mathbb{R}) \wedge \\ \quad (v \in \mathbb{R}) \wedge \\ \quad (m \in \mathbb{R}) \wedge \\ \quad (u < m) \wedge \\ \quad (v < m) \Rightarrow \\ \quad \quad \exists w \in \mathbb{R}.\; ((u \leq w) \wedge ((v \leq w) \wedge (w < m))) \end{array}}}
\vnbentry{\texttt{rr-\allowbreak{}min-\allowbreak{}pos}}{\vnbstmt{\begin{array}{@{}l@{}} \forall u, v. \\ \quad ((u \in \mathbb{R}) \wedge ((v \in \mathbb{R}) \wedge ((0 < u) \wedge (0 < v)))) \Rightarrow  \\ \quad \quad \exists w \in \mathbb{R}.\; ((0 < w) \wedge ((w \leq u) \wedge (w \leq v))) \end{array}}}
\vnbentry{\texttt{rr-\allowbreak{}prod-\allowbreak{}le-\allowbreak{}prod}}{\vnbstmt{\begin{array}{@{}l@{}} \forall a, b, u, v. \\ \quad (a \in \mathbb{R}) \wedge \\ \quad (b \in \mathbb{R}) \wedge \\ \quad (u \in \mathbb{R}) \wedge \\ \quad (v \in \mathbb{R}) \wedge \\ \quad (0 \leq a) \wedge \\ \quad (a \leq b) \wedge \\ \quad (0 \leq u) \wedge \\ \quad (u \leq v) \Rightarrow \\ \quad \quad ((a \cdot u) \leq (b \cdot v)) \end{array}}}
\vnbentry{\texttt{rr-\allowbreak{}nonneg-\allowbreak{}cancel-\allowbreak{}pos}}{\vnbstmt{\begin{array}{@{}l@{}} \forall c, d. \\ \quad ((c \in \mathbb{R}) \wedge ((d \in \mathbb{R}) \wedge ((0 < c) \wedge (0 \leq (c \cdot d))))) \Rightarrow  \\ \quad \quad (0 \leq d) \end{array}}}
\vnbentry{\texttt{rr-\allowbreak{}le-\allowbreak{}total}}{\vnbstmt{\forall u, v \in \mathbb{R}.\; ((u \leq v) \vee (v \leq u))}}
\vnbentry{\texttt{rr-\allowbreak{}lt-\allowbreak{}not-\allowbreak{}le}}{\vnbstmt{\forall a, b \in \mathbb{R}.\; ((a < b) \Rightarrow \neg (b \leq a))}}
\vnbentry{\texttt{rr-\allowbreak{}not-\allowbreak{}le-\allowbreak{}lt}}{\vnbstmt{\forall a, b \in \mathbb{R}.\; (\neg (a \leq b) \Rightarrow (b < a))}}

\vnbfile{theorem-library/co-order-trans-guarded.scm}{theorem-library/co-order-trans-guarded.scm}
\noindent{\small co-order-trans-guarded.scm -- the six `co-' order composition lemmas, PROVEN.  Four of them GUARDED on RR; two of them exactly as stated.}\par
\vnbentry{\texttt{co-\allowbreak{}lt-\allowbreak{}trans}}{\vnbstmt{\forall a, b, c \in \mathbb{R}.\; (((a < b) \wedge (b < c)) \Rightarrow (a < c))}}
\vnbentry{\texttt{co-\allowbreak{}le-\allowbreak{}lt-\allowbreak{}trans}}{\vnbstmt{\begin{array}{@{}l@{}} \forall a, b, c \in \mathbb{R}. \\ \quad (((a \leq b) \wedge (b < c)) \Rightarrow (a < c)) \end{array}}}
\vnbentry{\texttt{co-\allowbreak{}lt-\allowbreak{}le-\allowbreak{}trans}}{\vnbstmt{\begin{array}{@{}l@{}} \forall a, b, c \in \mathbb{R}. \\ \quad (((a < b) \wedge (b \leq c)) \Rightarrow (a < c)) \end{array}}}
\vnbentry{\texttt{co-\allowbreak{}le-\allowbreak{}eq-\allowbreak{}trans}}{\vnbstmt{\forall a, b, c.\; (((a \leq b) \wedge (b = c)) \Rightarrow (a \leq c))}}
\vnbentry{\texttt{co-\allowbreak{}eq-\allowbreak{}le-\allowbreak{}trans}}{\vnbstmt{\forall a, b, c.\; (((a = b) \wedge (b \leq c)) \Rightarrow (a \leq c))}}

\vnbfile{theorem-library/pos-rr-bridges.scm}{theorem-library/pos-rr-bridges.scm}
\noindent{\small pos-rr-bridges.scm -- the two projections of POS-RR, and one more positivity fact, PROVEN modulo 0.}\par
\vnbentry{\texttt{rr-\allowbreak{}pos-\allowbreak{}rr-\allowbreak{}in-\allowbreak{}rr}}{\vnbstmt{\forall t.\; (\operatorname{pos\text{-}rr}(t) \Rightarrow (t \in \mathbb{R}))}}
\vnbentry{\texttt{rr-\allowbreak{}lt-\allowbreak{}of-\allowbreak{}pos-\allowbreak{}rr}}{\vnbstmt{\forall t.\; (\operatorname{pos\text{-}rr}(t) \Rightarrow (0 < t))}}
\vnbentry{\texttt{rr-\allowbreak{}one-\allowbreak{}plus-\allowbreak{}nonneg-\allowbreak{}pos}}{\vnbstmt{\forall d \in \mathbb{R}.\; ((0 \leq d) \Rightarrow (0 < (1 + d)))}}

\vnbfile{theorem-library/rr-abs-basics.scm}{theorem-library/rr-abs-basics.scm}
\noindent{\small rr-abs-basics.scm -- absolute value on RR, PROVEN from its DEFINITION.}\par
\vnbentry{\texttt{rr-\allowbreak{}abs-\allowbreak{}of-\allowbreak{}nonneg}}{\vnbstmt{\forall x \in \mathbb{R}.\; ((0 \leq x) \Rightarrow (\lvert x \rvert = x))}}
\vnbentry{\texttt{rr-\allowbreak{}abs-\allowbreak{}of-\allowbreak{}nonpos}}{\vnbstmt{\forall x \in \mathbb{R}.\; ((x \leq 0) \Rightarrow (\lvert x \rvert = -x))}}
\vnbentry{\texttt{rr-\allowbreak{}abs-\allowbreak{}of-\allowbreak{}neg}}{\vnbstmt{\forall x \in \mathbb{R}.\; ((x < 0) \Rightarrow (\lvert x \rvert = -x))}}
\vnbentry{\texttt{rr-\allowbreak{}abs-\allowbreak{}cases}}{\vnbstmt{\begin{array}{@{}l@{}} \forall x \in \mathbb{R}. \\ \quad (((0 \leq x) \wedge (\lvert x \rvert = x)) \vee (\neg (0 \leq x) \wedge (\lvert x \rvert = -x))) \end{array}}}
\vnbentry{\texttt{rr-\allowbreak{}abs-\allowbreak{}closed}}{\vnbstmt{\forall a \in \mathbb{R}.\; (\lvert a \rvert \in \mathbb{R})}}
\vnbentry{\texttt{rr-\allowbreak{}abs-\allowbreak{}nonneg}}{\vnbstmt{\forall a \in \mathbb{R}.\; (0 \leq \lvert a \rvert)}}
\vnbentry{\texttt{rr-\allowbreak{}abs-\allowbreak{}zero}}{\vnbstmt{\forall a \in \mathbb{R}.\; ((\lvert a \rvert = 0) \Leftrightarrow (a = 0))}}
\vnbentry{\texttt{rr-\allowbreak{}le-\allowbreak{}abs}}{\vnbstmt{\forall x \in \mathbb{R}.\; (x \leq \lvert x \rvert)}}
\vnbentry{\texttt{rr-\allowbreak{}neg-\allowbreak{}abs-\allowbreak{}le}}{\vnbstmt{\forall x \in \mathbb{R}.\; (-\lvert x \rvert \leq x)}}
\vnbentry{\texttt{rr-\allowbreak{}abs-\allowbreak{}bound}}{\vnbstmt{\begin{array}{@{}l@{}} \forall x, c \in \mathbb{R}. \\ \quad ((\lvert x \rvert \leq c) \Leftrightarrow ((-c \leq x) \wedge (x \leq c))) \end{array}}}
\vnbentry{\texttt{rr-\allowbreak{}abs-\allowbreak{}triangle}}{\vnbstmt{\begin{array}{@{}l@{}} \forall a, b. \\ \quad ((a \in \mathbb{R}) \wedge (b \in \mathbb{R})) \Rightarrow  \\ \quad \quad (\lvert (a + b) \rvert \leq (\lvert a \rvert + \lvert b \rvert)) \end{array}}}
\vnbentry{\texttt{rr-\allowbreak{}abs-\allowbreak{}triangle-\allowbreak{}c}}{\vnbstmt{\begin{array}{@{}l@{}} \forall a, b. \\ \quad ((a \in \mathbb{R}) \wedge (b \in \mathbb{R})) \Rightarrow  \\ \quad \quad (\lvert (a + b) \rvert \leq (\lvert a \rvert + \lvert b \rvert)) \end{array}}}
\vnbentry{\texttt{rr-\allowbreak{}abs-\allowbreak{}neg}}{\vnbstmt{\forall a \in \mathbb{R}.\; (\lvert -a \rvert = \lvert a \rvert)}}
\vnbentry{\texttt{rr-\allowbreak{}abs-\allowbreak{}sub-\allowbreak{}sym}}{\vnbstmt{\begin{array}{@{}l@{}} \forall u, v. \\ \quad ((u \in \mathbb{R}) \wedge (v \in \mathbb{R})) \Rightarrow  \\ \quad \quad (\lvert (u - v) \rvert = \lvert (v - u) \rvert) \end{array}}}
\vnbentry{\texttt{rr-\allowbreak{}abs-\allowbreak{}mult}}{\vnbstmt{\begin{array}{@{}l@{}} \forall a, b. \\ \quad ((a \in \mathbb{R}) \wedge (b \in \mathbb{R})) \Rightarrow  \\ \quad \quad (\lvert (a \cdot b) \rvert = (\lvert a \rvert \cdot \lvert b \rvert)) \end{array}}}
\vnbentry{\texttt{rr-\allowbreak{}abs-\allowbreak{}reverse-\allowbreak{}triangle}}{\vnbstmt{\begin{array}{@{}l@{}} \forall x, y \in \mathbb{R}. \\ \quad (\lvert (\lvert x \rvert - \lvert y \rvert) \rvert \leq \lvert (x - y) \rvert) \end{array}}}
\vnbentry{\texttt{rr-\allowbreak{}abs-\allowbreak{}sum-\allowbreak{}bound}}{\vnbstmt{\begin{array}{@{}l@{}} \forall u, v, p, q, h, e. \\ \quad (u \in \mathbb{R}) \wedge \\ \quad (v \in \mathbb{R}) \wedge \\ \quad (p \in \mathbb{R}) \wedge \\ \quad (q \in \mathbb{R}) \wedge \\ \quad (h \in \mathbb{R}) \wedge \\ \quad (e \in \mathbb{R}) \wedge \\ \quad (\lvert (u - p) \rvert \leq h) \wedge \\ \quad (\lvert (v - q) \rvert \leq h) \wedge \\ \quad ((h + h) = e) \Rightarrow \\ \quad \quad (\lvert ((u + v) - (p + q)) \rvert \leq e) \end{array}}}
\vnbentry{\texttt{rr-\allowbreak{}abs-\allowbreak{}prod-\allowbreak{}bound}}{\vnbstmt{\begin{array}{@{}l@{}} \forall u, v, p, q, m, k, t, e. \\ \quad (u \in \mathbb{R}) \wedge \\ \quad (v \in \mathbb{R}) \wedge \\ \quad (p \in \mathbb{R}) \wedge \\ \quad (q \in \mathbb{R}) \wedge \\ \quad (m \in \mathbb{R}) \wedge \\ \quad (k \in \mathbb{R}) \wedge \\ \quad (t \in \mathbb{R}) \wedge \\ \quad (e \in \mathbb{R}) \wedge \\ \quad (\lvert u \rvert \leq m) \wedge \\ \quad (\lvert q \rvert \leq k) \wedge \\ \quad (\lvert (u - p) \rvert \leq t) \wedge \\ \quad (\lvert (v - q) \rvert \leq t) \wedge \\ \quad (((m + k) \cdot t) \leq e) \Rightarrow \\ \quad \quad (\lvert ((u \cdot v) - (p \cdot q)) \rvert \leq e) \end{array}}}

\vnbfile{theorem-library/rr-max-basics.scm}{theorem-library/rr-max-basics.scm}
\noindent{\small rr-max-basics.scm -- the usable facts about MAX on RR, PROVEN from its defining equation.}\par
\vnbentry{\texttt{rr-\allowbreak{}max-\allowbreak{}closed}}{\vnbstmt{\forall x, y \in \mathbb{R}.\; (\operatorname{max}(x, y) \in \mathbb{R})}}
\vnbentry{\texttt{rr-\allowbreak{}max-\allowbreak{}cases}}{\vnbstmt{\begin{array}{@{}l@{}} \forall x, y \in \mathbb{R}. \\ \quad ((\operatorname{max}(x, y) = x) \vee (\operatorname{max}(x, y) = y)) \end{array}}}
\vnbentry{\texttt{rr-\allowbreak{}le-\allowbreak{}max-\allowbreak{}left}}{\vnbstmt{\forall x, y \in \mathbb{R}.\; (x \leq \operatorname{max}(x, y))}}
\vnbentry{\texttt{rr-\allowbreak{}le-\allowbreak{}max-\allowbreak{}right}}{\vnbstmt{\forall x, y \in \mathbb{R}.\; (y \leq \operatorname{max}(x, y))}}
\vnbentry{\texttt{rr-\allowbreak{}max-\allowbreak{}le}}{\vnbstmt{\begin{array}{@{}l@{}} \forall x, y, z \in \mathbb{R}. \\ \quad (((x \leq z) \wedge (y \leq z)) \Rightarrow (\operatorname{max}(x, y) \leq z)) \end{array}}}

\vnbfile{theorem-library/rr-min-basics.scm}{theorem-library/rr-min-basics.scm}
\noindent{\small rr-min-basics.scm -- the usable facts about MIN on RR, PROVEN from its defining equation, and the NN closure of both MIN and MAX.}\par
\vnbentry{\texttt{rr-\allowbreak{}min-\allowbreak{}closed}}{\vnbstmt{\forall x, y \in \mathbb{R}.\; (\operatorname{min}(x, y) \in \mathbb{R})}}
\vnbentry{\texttt{rr-\allowbreak{}min-\allowbreak{}cases}}{\vnbstmt{\begin{array}{@{}l@{}} \forall x, y \in \mathbb{R}. \\ \quad ((\operatorname{min}(x, y) = x) \vee (\operatorname{min}(x, y) = y)) \end{array}}}
\vnbentry{\texttt{rr-\allowbreak{}min-\allowbreak{}le-\allowbreak{}left}}{\vnbstmt{\forall x, y \in \mathbb{R}.\; (\operatorname{min}(x, y) \leq x)}}
\vnbentry{\texttt{rr-\allowbreak{}min-\allowbreak{}le-\allowbreak{}right}}{\vnbstmt{\forall x, y \in \mathbb{R}.\; (\operatorname{min}(x, y) \leq y)}}
\vnbentry{\texttt{rr-\allowbreak{}le-\allowbreak{}min}}{\vnbstmt{\begin{array}{@{}l@{}} \forall x, y, z \in \mathbb{R}. \\ \quad (((z \leq x) \wedge (z \leq y)) \Rightarrow (z \leq \operatorname{min}(x, y))) \end{array}}}
\vnbentry{\texttt{nn-\allowbreak{}max-\allowbreak{}closed}}{\vnbstmt{\forall x, y \in \mathbb{N}.\; (\operatorname{max}(x, y) \in \mathbb{N})}}
\vnbentry{\texttt{nn-\allowbreak{}min-\allowbreak{}closed}}{\vnbstmt{\forall x, y \in \mathbb{N}.\; (\operatorname{min}(x, y) \in \mathbb{N})}}
\vnbentry{\texttt{rr-\allowbreak{}min-\allowbreak{}one-\allowbreak{}subadditive}}{\vnbstmt{\begin{array}{@{}l@{}} \forall x, y \in \mathbb{R}. \\ \quad ((0 \leq x) \wedge (0 \leq y)) \Rightarrow  \\ \quad \quad (\operatorname{min}(1, (x + y)) \leq (\operatorname{min}(1, x) + \operatorname{min}(1, y))) \end{array}}}

\vnbfile{theorem-library/cc-real-imag.scm}{theorem-library/cc-real-imag.scm}
\noindent{\small cc-real-imag.scm -- the real and imaginary parts of a complex number.}\par
\vnbentry{\texttt{cc-\allowbreak{}re-\allowbreak{}im-\allowbreak{}of}}{\vnbstmt{\begin{array}{@{}l@{}} \forall z \in \mathbb{C}, a, b \in \mathbb{R}. \\ \quad (z = (a + (b \cdot +i))) \Rightarrow  \\ \quad \quad ((\operatorname{real\text{-}part}(z) = a) \wedge (\operatorname{imag\text{-}part}(z) = b)) \end{array}}}
\vnbentry{\texttt{real-\allowbreak{}part-\allowbreak{}in-\allowbreak{}rr}}{\vnbstmt{\forall z \in \mathbb{C}.\; (\operatorname{real\text{-}part}(z) \in \mathbb{R})}}
\vnbentry{\texttt{imag-\allowbreak{}part-\allowbreak{}in-\allowbreak{}rr}}{\vnbstmt{\forall z \in \mathbb{C}.\; (\operatorname{imag\text{-}part}(z) \in \mathbb{R})}}
\vnbentry{\texttt{cc-\allowbreak{}re-\allowbreak{}im-\allowbreak{}decompose}}{\vnbstmt{\begin{array}{@{}l@{}} \forall z \in \mathbb{C}. \\ \quad (z = (\operatorname{real\text{-}part}(z) + (+i \cdot \operatorname{imag\text{-}part}(z)))) \end{array}}}
\vnbentry{\texttt{cc-\allowbreak{}plus-\allowbreak{}conj-\allowbreak{}is-\allowbreak{}2re}}{\vnbstmt{\begin{array}{@{}l@{}} \forall z \in \mathbb{C}. \\ \quad ((z + \operatorname{conjugate}(z)) = (2 \cdot \operatorname{real\text{-}part}(z))) \end{array}}}
\vnbentry{\texttt{cc-\allowbreak{}re-\allowbreak{}sq-\allowbreak{}le-\allowbreak{}mod-\allowbreak{}sq}}{\vnbstmt{\begin{array}{@{}l@{}} \forall z \in \mathbb{C}. \\ \quad (\operatorname{real\text{-}part}(z) \cdot \operatorname{real\text{-}part}(z)) \\ \quad \leq\; (z \cdot \operatorname{conjugate}(z)) \end{array}}}

\vnbfile{theorem-library/rr-recip-order.scm}{theorem-library/rr-recip-order.scm}
\noindent{\small rr-recip-order.scm -- the L1 rung of the "attach the loose ends" program: the reciprocal's SIGN, and the two field facts it stands on.}\par
\vnbentry{\texttt{rr-\allowbreak{}mul-\allowbreak{}zero}}{\vnbstmt{\forall a \in \mathbb{R}.\; ((a \cdot 0) = 0)}}
\vnbentry{\texttt{rr-\allowbreak{}zero-\allowbreak{}lt-\allowbreak{}one}}{\vnbstmt{(0 < 1)}}
\vnbentry{\texttt{rr-\allowbreak{}mul-\allowbreak{}pos}}{\vnbstmt{\begin{array}{@{}l@{}} \forall a, b \in \mathbb{R}. \\ \quad (((0 < a) \wedge (0 < b)) \Rightarrow (0 < (a \cdot b))) \end{array}}}
\vnbentry{\texttt{rr-\allowbreak{}recip-\allowbreak{}pos}}{\vnbstmt{\forall a \in \mathbb{R}.\; ((0 < a) \Rightarrow (0 < \operatorname{recip}(a)))}}

\vnbfile{theorem-library/pos-rr-of-lt.scm}{theorem-library/pos-rr-of-lt.scm}
\noindent{\small pos-rr-of-lt.scm -- 0 $<$ x IMPLIES POS-RR(x), the bridge the tree did not have, and `nn-recip-succ-pos' retired as a four-line consequence of it.}\par
\vnbentry{\texttt{rr-\allowbreak{}pos-\allowbreak{}rr-\allowbreak{}of-\allowbreak{}lt}}{\vnbstmt{\begin{array}{@{}l@{}} \forall x \in \mathbb{R}. \\ \quad ((0 < x) \Rightarrow \operatorname{pos\text{-}rr}(x)) \end{array}}}
\vnbentry{\texttt{nn-\allowbreak{}recip-\allowbreak{}succ-\allowbreak{}pos}}{\vnbstmt{\begin{array}{@{}l@{}} \forall n \in \mathbb{N}. \\ \quad \operatorname{pos\text{-}rr}(\operatorname{recip}((n + 1))) \end{array}}}

\vnbfile{theorem-library/rr-halving.scm}{theorem-library/rr-halving.scm}
\noindent{\small rr-halving.scm -- every positive real splits into two equal positive halves, PROVEN.}\par
\vnbentry{\texttt{rr-\allowbreak{}pos-\allowbreak{}halvable}}{\vnbstmt{\begin{array}{@{}l@{}} \forall s. \\ \quad \operatorname{pos\text{-}rr}(s) \Rightarrow  \\ \quad \quad \exists d.\; (\operatorname{pos\text{-}rr}(d) \wedge ((d + d) = s)) \end{array}}}

\vnbfile{theorem-library/rr-order-bundle.scm}{theorem-library/rr-order-bundle.scm}
\noindent{\small rr-order-bundle.scm -- four elementary real-order facts, PROVEN modulo 0.}\par
\vnbentry{\texttt{rr-\allowbreak{}prod-\allowbreak{}nonpos-\allowbreak{}pos}}{\vnbstmt{\begin{array}{@{}l@{}} \forall u, v \in \mathbb{R}. \\ \quad (((0 < v) \wedge ((u \cdot v) \leq 0)) \Rightarrow (u \leq 0)) \end{array}}}
\vnbentry{\texttt{rr-\allowbreak{}prod-\allowbreak{}nonpos-\allowbreak{}neg}}{\vnbstmt{\begin{array}{@{}l@{}} \forall u, v \in \mathbb{R}. \\ \quad (((v < 0) \wedge ((u \cdot v) \leq 0)) \Rightarrow (0 \leq u)) \end{array}}}
\vnbentry{\texttt{rr-\allowbreak{}midpoint-\allowbreak{}between}}{\vnbstmt{\begin{array}{@{}l@{}} \forall u, v \in \mathbb{R}. \\ \quad ((u < v) \Rightarrow \exists w \in \mathbb{R}.\; ((u < w) \wedge (w < v))) \end{array}}}
\vnbentry{\texttt{rr-\allowbreak{}le-\allowbreak{}scale-\allowbreak{}nonneg-\allowbreak{}right}}{\vnbstmt{\begin{array}{@{}l@{}} \forall x, y, c \in \mathbb{R}. \\ \quad (((0 \leq c) \wedge (x \leq y)) \Rightarrow ((x \cdot c) \leq (y \cdot c))) \end{array}}}

\vnbfile{theorem-library/rr-le-all-pos.scm}{theorem-library/rr-le-all-pos.scm}
\noindent{\small rr-le-all-pos.scm -- a real below every positive real is non-positive, PROVEN.}\par
\vnbentry{\texttt{rr-\allowbreak{}le-\allowbreak{}all-\allowbreak{}pos-\allowbreak{}nonpos}}{\vnbstmt{\begin{array}{@{}l@{}} \forall x \in \mathbb{R}. \\ \quad \forall s.\; (\operatorname{pos\text{-}rr}(s) \Rightarrow (x \leq s)) \Rightarrow  \\ \quad \quad (x \leq 0) \end{array}}}

\vnbfile{theorem-library/subset-lemmas.scm}{theorem-library/subset-lemmas.scm}
\noindent{\small subset-lemmas.scm -- the four inclusion facts, PROVEN.}\par
\vnbentry{\texttt{subset-\allowbreak{}mem-\allowbreak{}fwd}}{\vnbstmt{\forall a, b, x.\; (((a \subseteq b) \wedge (x \in a)) \Rightarrow (x \in b))}}
\vnbentry{\texttt{subset-\allowbreak{}trans}}{\vnbstmt{\begin{array}{@{}l@{}} \forall a, b, c. \\ \quad (((a \subseteq b) \wedge (b \subseteq c)) \Rightarrow (a \subseteq c)) \end{array}}}
\vnbentry{\texttt{subclass-\allowbreak{}of-\allowbreak{}set-\allowbreak{}is-\allowbreak{}set}}{\vnbstmt{\begin{array}{@{}l@{}} \forall a.; \forall b \in \mathbf{Set}. \\ \quad ((a \subseteq b) \Rightarrow (a \in \mathbf{Set})) \end{array}}}

\vnbfile{theorem-library/fun-image-set.scm}{theorem-library/fun-image-set.scm}
\noindent{\small fun-image-set.scm -- the image of ANY class under a function into a SET is a set, proven with no citation of `image-set' (replacement).}\par
\vnbentry{\texttt{fun-\allowbreak{}image-\allowbreak{}set}}{\vnbstmt{\begin{array}{@{}l@{}} \forall a, b, f, s. \\ \quad ((f \in (a \to b)) \wedge (b \in \mathbf{Set})) \Rightarrow  \\ \quad \quad (\operatorname{image}(f, s) \in \mathbf{Set}) \end{array}}}

\vnbfile{theorem-library/rake-analysis2.scm}{theorem-library/rake-analysis2.scm}
\noindent{\small rake-analysis2.scm -- rake batch L (2026-09-17): five asserted leaves of the analysis / sets area, plus the general lemma that unblocked one of them.}\par
\vnbentry{\texttt{fun-\allowbreak{}domain-\allowbreak{}in-\allowbreak{}set}}{\vnbstmt{\forall a, b.; \forall f \in (a \to b).\; (a \in \mathbf{Set})}}
\vnbentry{\texttt{bijection-\allowbreak{}compose}}{\vnbstmt{\begin{array}{@{}l@{}} \forall x, y, z, i, a. \\ \quad ((i \in \mathrm{Bij}(x, y)) \wedge (a \in \mathrm{Bij}(y, z))) \Rightarrow  \\ \quad \quad ((\lambda b \in x.\; a(i(b))) \in \mathrm{Bij}(x, z)) \end{array}}}
\vnbentry{\texttt{subseq-\allowbreak{}is-\allowbreak{}fun}}{\vnbstmt{\begin{array}{@{}l@{}} \forall s, f, i. \\ \quad (f \in (\mathbb{N} \to \operatorname{pts}(s))) \wedge \\ \quad \operatorname{strictly\text{-}mono\text{-}nn}(i) \Rightarrow \\ \quad \quad (\operatorname{subseq}(f, i) \in (\mathbb{N} \to \operatorname{pts}(s))) \end{array}}}
\vnbentry{\texttt{metric-\allowbreak{}hom-\allowbreak{}is-\allowbreak{}continuous}}{\vnbstmt{\begin{array}{@{}l@{}} \forall s, t, f. \\ \quad \operatorname{is\text{-}hom\text{-}metric\text{-}space}(s, t, f) \Rightarrow  \\ \quad \quad \operatorname{is\text{-}continuous}(s, t, f) \end{array}}}
\vnbentry{\texttt{le-\allowbreak{}bound-\allowbreak{}mono}}{\vnbstmt{\begin{array}{@{}l@{}} \forall p, a, b, n. \\ \quad (p \in \mathbb{R}) \wedge \\ \quad (a \in \mathbb{R}) \wedge \\ \quad (b \in \mathbb{R}) \wedge \\ \quad (n \in \mathbb{R}) \wedge \\ \quad (0 \leq n) \wedge \\ \quad (p \leq (a \cdot n)) \wedge \\ \quad (a \leq b) \Rightarrow \\ \quad \quad (p \leq (b \cdot n)) \end{array}}}
\vnbentry{\texttt{centre-\allowbreak{}set-\allowbreak{}contains-\allowbreak{}choice}}{\vnbstmt{\begin{array}{@{}l@{}} \forall s, r, f.; \forall u \in f. \\ \quad \exists c.\; (c \in \operatorname{centres}(s, u, r)) \Rightarrow  \\ \quad \quad (\varepsilon(\operatorname{centres}(s, u, r)) \in \operatorname{centre\text{-}set}(s, r, f)) \end{array}}}

\vnbfile{theorem-library/rake-inverse-bij.scm}{theorem-library/rake-inverse-bij.scm}
\noindent{\small theorem-library/rake-inverse-bij.scm -- rake batch N (2026-09-17), part 1 of 2. --------------------------------------------------------------------------- NOTE, 2026-09-20 (batch 9-B).  This file was written while cardinality was AXIOMATISED under the name CARD and the defined constant was its companion}\par
\vnbentry{\texttt{inverse-\allowbreak{}bij-\allowbreak{}right}}{\vnbstmt{\forall m, d.; \forall i \in \mathrm{Bij}(m, d), y \in d.\; (i(i^{-1}(y)) = y)}}
\vnbentry{\texttt{inverse-\allowbreak{}bij-\allowbreak{}left}}{\vnbstmt{\forall m, d.; \forall i \in \mathrm{Bij}(m, d), x \in m.\; (i^{-1}(i(x)) = x)}}
\vnbentry{\texttt{inverse-\allowbreak{}bij-\allowbreak{}in-\allowbreak{}fun}}{\vnbstmt{\forall x, y.; \forall i \in \mathrm{Bij}(x, y).\; (i^{-1} \in (y \to x))}}
\vnbentry{\texttt{inverse-\allowbreak{}bij-\allowbreak{}is-\allowbreak{}bijection}}{\vnbstmt{\begin{array}{@{}l@{}} \forall x, y.; \forall i \in \mathrm{Bij}(x, y). \\ \quad (i^{-1} \in \mathrm{Bij}(y, x)) \end{array}}}
\vnbentry{\texttt{ord-\allowbreak{}segment-\allowbreak{}self}}{\vnbstmt{\forall n \in \mathbb{N}.\; \neg (n \in \mathrm{OS}(n))}}

\vnbfile{theorem-library/nn-parity-proof.scm}{theorem-library/nn-parity-proof.scm}
\noindent{\small nn-parity-proof.scm -- Peano arithmetic on NN: from the two recursion equations (structure-library/nn-arith.scm) to parity and even-square.}\par
\vnbentry{\texttt{nn-\allowbreak{}one-\allowbreak{}is-\allowbreak{}succ-\allowbreak{}zero}}{\vnbstmt{(1 = (0 + 1))}}
\vnbentry{\texttt{nn-\allowbreak{}succ-\allowbreak{}inj}}{\vnbstmt{\forall a, b \in \mathbb{N}.\; (((a + 1) = (b + 1)) \Rightarrow (a = b))}}
\vnbentry{\texttt{nn-\allowbreak{}succ-\allowbreak{}nonzero}}{\vnbstmt{\forall n \in \mathbb{N}.\; \neg ((n + 1) = 0)}}
\vnbentry{\texttt{nn-\allowbreak{}succ-\allowbreak{}plus-\allowbreak{}one}}{\vnbstmt{\forall n \in \mathbb{N}.\; ((n + 1) = (n + 1))}}
\vnbentry{\texttt{nn-\allowbreak{}mul-\allowbreak{}succ}}{\vnbstmt{\forall a, b \in \mathbb{N}.\; ((a \cdot (b + 1)) = ((a \cdot b) + a))}}
\vnbentry{\texttt{nn-\allowbreak{}two-\allowbreak{}is-\allowbreak{}succ-\allowbreak{}succ-\allowbreak{}zero}}{\vnbstmt{(2 = ((0 + 1) + 1))}}
\vnbentry{\texttt{nn-\allowbreak{}plus-\allowbreak{}two}}{\vnbstmt{\forall x \in \mathbb{N}.\; ((x + 2) = ((x + 1) + 1))}}
\vnbentry{\texttt{nn-\allowbreak{}two-\allowbreak{}mul-\allowbreak{}succ}}{\vnbstmt{\forall k \in \mathbb{N}.\; ((2 \cdot (k + 1)) = (((2 \cdot k) + 1) + 1))}}
\vnbentry{\texttt{nn-\allowbreak{}nonzero-\allowbreak{}is-\allowbreak{}succ}}{\vnbstmt{\begin{array}{@{}l@{}} \forall n \in \mathbb{N}. \\ \quad (\neg (n = 0) \Rightarrow \exists q \in \mathbb{N}.\; (n = (q + 1))) \end{array}}}
\vnbentry{\texttt{nn-\allowbreak{}parity}}{\vnbstmt{\begin{array}{@{}l@{}} \forall n \in \mathbb{N}. \exists k \in \mathbb{N}. \\ \quad ((n = (2 \cdot k)) \vee (n = ((2 \cdot k) + 1))) \end{array}}}
\vnbentry{\texttt{nn-\allowbreak{}zero-\allowbreak{}or-\allowbreak{}succ}}{\vnbstmt{\begin{array}{@{}l@{}} \forall y \in \mathbb{N}. \\ \quad ((y = 0) \vee \exists q \in \mathbb{N}.\; (y = (q + 1))) \end{array}}}
\vnbentry{\texttt{nn-\allowbreak{}parity-\allowbreak{}exclusive}}{\vnbstmt{\forall x, y \in \mathbb{N}.\; \neg ((2 \cdot x) = ((2 \cdot y) + 1))}}
\vnbentry{\texttt{nn-\allowbreak{}even-\allowbreak{}square}}{\vnbstmt{\begin{array}{@{}l@{}} \forall p \in \mathbb{N}. \\ \quad \exists m \in \mathbb{N}.\; ((p \cdot p) = (2 \cdot m)) \Rightarrow  \\ \quad \quad \exists k \in \mathbb{N}.\; (p = (2 \cdot k)) \end{array}}}

\vnbfile{theorem-library/nn-pred.scm}{theorem-library/nn-pred.scm}
\noindent{\small nn-pred.scm -- the predecessor on NN, DEFINED by definite description.}\par
\vnbentry{\texttt{pred-\allowbreak{}property}}{\vnbstmt{\begin{array}{@{}l@{}} \forall n \in \mathbb{N}. \\ \quad \neg (n = 0) \Rightarrow  \\ \quad \quad ((\operatorname{pred}(n) \in \mathbb{N}) \wedge ((\operatorname{pred}(n) + 1) = n)) \end{array}}}
\vnbentry{\texttt{pred-\allowbreak{}in-\allowbreak{}nn}}{\vnbstmt{\begin{array}{@{}l@{}} \forall n \in \mathbb{N}. \\ \quad (\neg (n = 0) \Rightarrow (\operatorname{pred}(n) \in \mathbb{N})) \end{array}}}
\vnbentry{\texttt{pred-\allowbreak{}succ}}{\vnbstmt{\begin{array}{@{}l@{}} \forall n \in \mathbb{N}. \\ \quad (\neg (n = 0) \Rightarrow ((\operatorname{pred}(n) + 1) = n)) \end{array}}}
\vnbentry{\texttt{pred-\allowbreak{}of-\allowbreak{}succ}}{\vnbstmt{\forall m \in \mathbb{N}.\; (\operatorname{pred}((m + 1)) = m)}}

\vnbfile{theorem-library/ord-segment-arith.scm}{theorem-library/ord-segment-arith.scm}
\noindent{\small ord-segment-arith.scm -- the two bridges between segment membership and arithmetic on NN.}\par
\vnbentry{\texttt{seg-\allowbreak{}mem-\allowbreak{}lt}}{\vnbstmt{\forall n, j \in \mathbb{N}.\; ((j \in \mathrm{OS}(n)) \Leftrightarrow (j < n))}}
\vnbentry{\texttt{seg-\allowbreak{}mem-\allowbreak{}succ-\allowbreak{}le}}{\vnbstmt{\begin{array}{@{}l@{}} \forall n, j \in \mathbb{N}. \\ \quad ((j \in \mathrm{OS}((n + 1))) \Leftrightarrow (j \leq n)) \end{array}}}

\vnbfile{theorem-library/finite-surgery.scm}{theorem-library/finite-surgery.scm}
\noindent{\small finite-surgery.scm -- the finite-surgery kit, first member: COLLAPSE-AT.}\par
\vnbentry{\texttt{nn-\allowbreak{}lt-\allowbreak{}succ-\allowbreak{}le}}{\vnbstmt{\forall k, j \in \mathbb{N}.\; ((k < j) \Rightarrow ((k + 1) \leq j))}}
\vnbentry{\texttt{collapse-\allowbreak{}at-\allowbreak{}lo}}{\vnbstmt{\begin{array}{@{}l@{}} \forall k, j \in \mathbb{N}. \\ \quad (\neg (k < j) \Rightarrow (\operatorname{collapse\text{-}at}(k)(j) = j)) \end{array}}}
\vnbentry{\texttt{collapse-\allowbreak{}at-\allowbreak{}hi}}{\vnbstmt{\begin{array}{@{}l@{}} \forall k, j \in \mathbb{N}. \\ \quad (k < j) \Rightarrow  \\ \quad \quad (\operatorname{collapse\text{-}at}(k)(j) = \operatorname{pred}(j)) \end{array}}}
\vnbentry{\texttt{collapse-\allowbreak{}at-\allowbreak{}in-\allowbreak{}nn}}{\vnbstmt{\begin{array}{@{}l@{}} \forall k, j \in \mathbb{N}. \\ \quad (\operatorname{collapse\text{-}at}(k)(j) \in \mathbb{N}) \end{array}}}
\vnbentry{\texttt{collapse-\allowbreak{}at-\allowbreak{}lt}}{\vnbstmt{\begin{array}{@{}l@{}} \forall n, k, j. \\ \quad (n \in \mathbb{N}) \wedge \\ \quad (k \in \mathbb{N}) \wedge \\ \quad (j \in \mathbb{N}) \wedge \\ \quad (k \leq n) \wedge \\ \quad (j \leq n) \wedge \\ \quad \neg (j = k) \Rightarrow \\ \quad \quad (\operatorname{collapse\text{-}at}(k)(j) < n) \end{array}}}
\vnbentry{\texttt{collapse-\allowbreak{}at-\allowbreak{}inj}}{\vnbstmt{\begin{array}{@{}l@{}} \forall k, a, b. \\ \quad (k \in \mathbb{N}) \wedge \\ \quad (a \in \mathbb{N}) \wedge \\ \quad (b \in \mathbb{N}) \wedge \\ \quad \neg (a = k) \wedge \\ \quad \neg (b = k) \wedge \\ \quad (\operatorname{collapse\text{-}at}(k)(a) = \operatorname{collapse\text{-}at}(k)(b)) \Rightarrow \\ \quad \quad (a = b) \end{array}}}
\vnbentry{\texttt{collapse-\allowbreak{}at-\allowbreak{}in-\allowbreak{}seg}}{\vnbstmt{\begin{array}{@{}l@{}} \forall n, k, j. \\ \quad (n \in \mathbb{N}) \wedge \\ \quad (k \in \mathbb{N}) \wedge \\ \quad (j \in \mathbb{N}) \wedge \\ \quad (k \in \mathrm{OS}((n + 1))) \wedge \\ \quad (j \in \mathrm{OS}((n + 1))) \wedge \\ \quad \neg (j = k) \Rightarrow \\ \quad \quad (\operatorname{collapse\text{-}at}(k)(j) \in \mathrm{OS}(n)) \end{array}}}

\vnbfile{theorem-library/pigeonhole-segments.scm}{theorem-library/pigeonhole-segments.scm}
\noindent{\small pigeonhole-segments.scm -- Track A of the CARD plan: finite pigeonhole for ordinal segments, the one piece of real mathematics the plan needs.}\par
\vnbentry{\texttt{fun-\allowbreak{}into-\allowbreak{}seg-\allowbreak{}zero-\allowbreak{}absurd}}{\vnbstmt{\begin{array}{@{}l@{}} \forall a, f, x. \\ \quad (((f \in (a \to \mathrm{OS}(0))) \wedge (x \in a)) \Rightarrow \mathit{falsity}) \end{array}}}
\vnbentry{\texttt{no-\allowbreak{}injection-\allowbreak{}seg-\allowbreak{}1-\allowbreak{}into-\allowbreak{}seg-\allowbreak{}0}}{\vnbstmt{\begin{array}{@{}l@{}} \forall f. \\ \quad \neg (f \in \operatorname{injection}(\mathrm{OS}((0 + 1)), \mathrm{OS}(0))) \end{array}}}
\vnbentry{\texttt{pigeonhole-\allowbreak{}step}}{\vnbstmt{\begin{array}{@{}l@{}} \forall n \in \mathbb{N}, g. \\ \quad \forall f.\; \neg (f \in \operatorname{injection}(\mathrm{OS}((n + 1)), \mathrm{OS}(n))) \Rightarrow  \\ \quad \quad \neg (g \in \operatorname{injection}(\mathrm{OS}(((n + 1) + 1)), \mathrm{OS}((n + 1)))) \end{array}}}
\vnbentry{\texttt{pigeonhole-\allowbreak{}segments}}{\vnbstmt{\begin{array}{@{}l@{}} \forall n \in \mathbb{N}, f. \\ \quad \neg (f \in \operatorname{injection}(\mathrm{OS}((n + 1)), \mathrm{OS}(n))) \end{array}}}
\vnbentry{\texttt{pigeonhole-\allowbreak{}segments-\allowbreak{}gen}}{\vnbstmt{\begin{array}{@{}l@{}} \forall m, n, f. \\ \quad ((m \in \mathbb{N}) \wedge ((n \in \mathbb{N}) \wedge (n < m))) \Rightarrow  \\ \quad \quad \neg (f \in \operatorname{injection}(\mathrm{OS}(m), \mathrm{OS}(n))) \end{array}}}

\vnbfile{theorem-library/enum-append.scm}{theorem-library/enum-append.scm}
\noindent{\small theorem-library/enum-append.scm -- enum-append-is-bijection, moved out of theorem-library/finsum-insert.scm on 2026-09-20 (batch 9-B, CARD := CARD-STAR).}\par
\vnbentry{\texttt{enum-\allowbreak{}append-\allowbreak{}is-\allowbreak{}bijection}}{\vnbstmt{\begin{array}{@{}l@{}} \forall n \in \mathbb{N}, g, t \in \mathbf{Set}, m \in \mathrm{Bij}(\mathrm{OS}(n), g). \\ \quad \neg (t \in g) \Rightarrow  \\ \quad \quad (\lambda i \in \mathrm{OS}((n + 1)).\; \\ \quad \quad \quad \left(\text{if }(i \in \mathrm{OS}(n))\text{ then }m(i)\text{ else }t\right)) \\ \quad \quad \in\; \mathrm{Bij}(\mathrm{OS}((n + 1)), (g \cup \{t\})) \end{array}}}

\vnbfile{theorem-library/union-laws.scm}{theorem-library/union-laws.scm}
\noindent{\small theorem-library/union-laws.scm -- union-empty-right and union-assoc, moved out of theorem-library/fin-subsets.scm on 2026-09-20 (batch 9-B, CARD := CARD-STAR).}\par
\vnbentry{\texttt{union-\allowbreak{}empty-\allowbreak{}right}}{\vnbstmt{\forall a.\; ((a \cup \emptyset) = a)}}
\vnbentry{\texttt{union-\allowbreak{}assoc}}{\vnbstmt{\forall a, b, c.\; (((a \cup b) \cup c) = (a \cup (b \cup c)))}}

\vnbfile{structure-library/bijection-derived.scm}{structure-library/bijection-derived.scm}
\noindent{\small bijection-derived.scm -- the three BIJECTION projection lemmas, PROVEN.}\par
\vnbentry{\texttt{bijection-\allowbreak{}in-\allowbreak{}fun}}{\vnbstmt{\forall x, y.; \forall i \in \mathrm{Bij}(x, y).\; (i \in (x \to y))}}
\vnbentry{\texttt{bijection-\allowbreak{}injective}}{\vnbstmt{\begin{array}{@{}l@{}} \forall x, y.; \forall i \in \mathrm{Bij}(x, y), a, b \in x. \\ \quad ((i(a) = i(b)) \Rightarrow (a = b)) \end{array}}}
\vnbentry{\texttt{bijection-\allowbreak{}surjective}}{\vnbstmt{\begin{array}{@{}l@{}} \forall x, y.; \forall i \in \mathrm{Bij}(x, y), w \in y. \exists z \in x. \\ \quad (i(z) = w) \end{array}}}

\vnbfile{theorem-library/bijection-identity-proof.scm}{theorem-library/bijection-identity-proof.scm}
\noindent{\small bijection-identity-proof.scm -- the identity lambda on a SET is a bijection.}\par
\vnbentry{\texttt{bijection-\allowbreak{}identity}}{\vnbstmt{\forall x \in \mathbf{Set}.\; ((\lambda a \in x.\; a) \in \mathrm{Bij}(x, x))}}

\vnbfile{theorem-library/card-defined.scm}{theorem-library/card-defined.scm}
\noindent{\small card-defined.scm -- CARD, cardinality DEFINED rather than axiomatised.}\par
\vnbentry{\texttt{cd-\allowbreak{}body-\allowbreak{}unique}}{\vnbstmt{\begin{array}{@{}l@{}} \forall a, l, e. \\ \quad (l \in \mathrm{Ord}) \wedge \\ \quad \exists i.\; (i \in \mathrm{Bij}(a, \mathrm{OS}(l))) \wedge \\ \quad \forall b. \\ \quad \quad \operatorname{ord\text{-}lt}(b, l) \Rightarrow  \\ \quad \quad \quad \neg \exists c.\; (c \in \mathrm{Bij}(a, \mathrm{OS}(b))) \wedge \\ \quad (e \in \mathrm{Ord}) \wedge \\ \quad \exists i.\; (i \in \mathrm{Bij}(a, \mathrm{OS}(e))) \wedge \\ \quad \forall b. \\ \quad \quad \operatorname{ord\text{-}lt}(b, e) \Rightarrow  \\ \quad \quad \quad \neg \exists c.\; (c \in \mathrm{Bij}(a, \mathrm{OS}(b))) \Rightarrow \\ \quad \quad (l = e) \end{array}}}
\vnbentry{\texttt{cd-\allowbreak{}seg-\allowbreak{}body}}{\vnbstmt{\begin{array}{@{}l@{}} \forall n \in \mathbb{N}. \\ \quad (n \in \mathrm{Ord}) \wedge \\ \quad \exists i.\; (i \in \mathrm{Bij}(\mathrm{OS}(n), \mathrm{OS}(n))) \wedge \\ \quad \forall a. \\ \quad \quad \operatorname{ord\text{-}lt}(a, n) \Rightarrow  \\ \quad \quad \quad \neg \exists b.\; (b \in \mathrm{Bij}(\mathrm{OS}(n), \mathrm{OS}(a))) \end{array}}}
\vnbentry{\texttt{card-\allowbreak{}segment}}{\vnbstmt{\forall n \in \mathbb{N}.\; (|\mathrm{OS}(n)| = n)}}

\vnbfile{theorem-library/card-finite.scm}{theorem-library/card-finite.scm}
\noindent{\small card-finite.scm -- the finite layer of the DEFINED cardinal, CARD. --------------------------------------------------------------------------- NOTE, 2026-09-20 (batch 9-B).  This file was written while cardinality was AXIOMATISED under the name CARD and the defined constant was its companion}\par
\vnbentry{\texttt{card-\allowbreak{}from-\allowbreak{}body}}{\vnbstmt{\begin{array}{@{}l@{}} \forall a.; \forall l \in \mathrm{Ord}. \\ \quad \exists i.\; (i \in \mathrm{Bij}(a, \mathrm{OS}(l))) \wedge \\ \quad \forall b. \\ \quad \quad \operatorname{ord\text{-}lt}(b, l) \Rightarrow  \\ \quad \quad \quad \neg \exists c.\; (c \in \mathrm{Bij}(a, \mathrm{OS}(b))) \Rightarrow \\ \quad \quad (|a| = l) \end{array}}}
\vnbentry{\texttt{card-\allowbreak{}bij}}{\vnbstmt{\begin{array}{@{}l@{}} \forall n \in \mathbb{N}, a, h, t. \\ \quad ((h \in \mathrm{Bij}(a, \mathrm{OS}(n))) \wedge (t \in \mathrm{Bij}(\mathrm{OS}(n), a))) \Rightarrow  \\ \quad \quad (|a| = n) \end{array}}}
\vnbentry{\texttt{seg0-\allowbreak{}empty}}{\vnbstmt{(\emptyset = \mathrm{OS}(0))}}
\vnbentry{\texttt{card-\allowbreak{}empty}}{\vnbstmt{(|\emptyset| = 0)}}

\vnbfile{theorem-library/ord-no-injection.scm}{theorem-library/ord-no-injection.scm}
\noindent{\small ord-no-injection.scm -- NO SET RECEIVES AN INJECTIVE CLASS FUNCTION FROM ORD.}\par
\vnbentry{\texttt{ord-\allowbreak{}no-\allowbreak{}injection-\allowbreak{}into-\allowbreak{}set}}{\vnbstmt{\begin{array}{@{}l@{}} \forall d, i. \\ \quad (d \in \mathbf{Set}) \wedge \\ \quad \forall a \in \mathrm{Ord}.\; (i(a) \in d) \wedge \\ \quad \forall a_{1}, a_{2} \in \mathrm{Ord}. \\ \quad \quad ((i(a_{1}) = i(a_{2})) \Rightarrow (a_{1} = a_{2})) \Rightarrow \\ \quad \quad \mathit{falsity} \end{array}}}

\vnbfile{theorem-library/zen-step.scm}{theorem-library/zen-step.scm}
\noindent{\small zen-step.scm -- ZEN's three recursion equations, collapsed into ONE.}\par
\vnbentry{\texttt{ord-\allowbreak{}lt-\allowbreak{}succ-\allowbreak{}iff-\allowbreak{}le}}{\vnbstmt{\begin{array}{@{}l@{}} \forall l, z \in \mathrm{Ord}. \\ \quad (\operatorname{ord\text{-}lt}(z, (l + 1)) \Leftrightarrow \operatorname{ord\text{-}le}(z, l)) \end{array}}}
\vnbentry{\texttt{zen-\allowbreak{}step}}{\vnbstmt{\begin{array}{@{}l@{}} \forall a \in \mathrm{Ord}, x. \\ \quad \operatorname{zen}(x, a) \\ \quad =\; \varepsilon(\{y \in x : \forall c.\; (\operatorname{ord\text{-}lt}(c, a) \Rightarrow \neg (\operatorname{zen}(x, c) = y))\}) \end{array}}}
\vnbentry{\texttt{zen-\allowbreak{}hits}}{\vnbstmt{\begin{array}{@{}l@{}} \forall a \in \mathrm{Ord}, s. \\ \quad \exists x. \\ \quad \quad x \\ \quad \quad \in\; \{y \in s : \forall c.\; (\operatorname{ord\text{-}lt}(c, a) \Rightarrow \neg (\operatorname{zen}(s, c) = y))\} \Rightarrow \\ \quad \quad \operatorname{zen}(s, a) \\ \quad \quad \in\; \{y \in s : \forall c.\; (\operatorname{ord\text{-}lt}(c, a) \Rightarrow \neg (\operatorname{zen}(s, c) = y))\} \end{array}}}
\vnbentry{\texttt{zen-\allowbreak{}exhausts}}{\vnbstmt{\begin{array}{@{}l@{}} \forall s \in \mathbf{Set}. \exists a \in \mathrm{Ord}. \\ \quad \neg \exists x.\; (x \in \{y \in s : \forall c.\; (\operatorname{ord\text{-}lt}(c, a) \Rightarrow \neg (\operatorname{zen}(s, c) = y))\}) \end{array}}}

\vnbfile{theorem-library/rake-zermelo.scm}{theorem-library/rake-zermelo.scm}
\noindent{\small rake-zermelo.scm -- CARD-free Zermelo: every set bijects with an ordinal --------------------------------------------------------------------------- NOTE, 2026-09-20 (batch 9-B).  This file was written while cardinality was AXIOMATISED under the name CARD and the defined constant was its companion}\par
\vnbentry{\texttt{zermelo-\allowbreak{}least-\allowbreak{}ordinal}}{\vnbstmt{\begin{array}{@{}l@{}} \forall s \in \mathbf{Set}. \exists l \in \mathrm{Ord}. \\ \quad \neg \exists x.\; (x \in \{y \in s : \forall c.\; (\operatorname{ord\text{-}lt}(c, l) \Rightarrow \neg (\operatorname{zen}(s, c) = y))\}) \wedge \\ \quad \forall z. \\ \quad \quad \operatorname{ord\text{-}lt}(z, l) \Rightarrow  \\ \quad \quad \quad \exists x. \\ \quad \quad \quad \quad x \\ \quad \quad \quad \quad \in\; \{y \in s : \forall c.\; (\operatorname{ord\text{-}lt}(c, z) \Rightarrow \neg (\operatorname{zen}(s, c) = y))\} \end{array}}}
\vnbentry{\texttt{zen-\allowbreak{}enum-\allowbreak{}in-\allowbreak{}fun}}{\vnbstmt{\begin{array}{@{}l@{}} \forall l \in \mathrm{Ord}, s \in \mathbf{Set}. \\ \quad \forall z. \\ \quad \quad \operatorname{ord\text{-}lt}(z, l) \Rightarrow  \\ \quad \quad \quad \exists x. \\ \quad \quad \quad \quad x \\ \quad \quad \quad \quad \in\; \{y \in s : \forall c.\; (\operatorname{ord\text{-}lt}(c, z) \Rightarrow \neg (\operatorname{zen}(s, c) = y))\} \Rightarrow \\ \quad \quad ((\lambda a \in \mathrm{OS}(l).\; \operatorname{zen}(s, a)) \in (\mathrm{OS}(l) \to s)) \end{array}}}
\vnbentry{\texttt{zen-\allowbreak{}enum-\allowbreak{}injective}}{\vnbstmt{\begin{array}{@{}l@{}} \forall l \in \mathrm{Ord}, s \in \mathbf{Set}, a_{1}, a_{2} \in \mathrm{OS}(l). \\ \quad \forall z. \\ \quad \quad \operatorname{ord\text{-}lt}(z, l) \Rightarrow  \\ \quad \quad \quad \exists x. \\ \quad \quad \quad \quad x \\ \quad \quad \quad \quad \in\; \{y \in s : \forall c.\; (\operatorname{ord\text{-}lt}(c, z) \Rightarrow \neg (\operatorname{zen}(s, c) = y))\} \wedge \\ \quad (\lambda a \in \mathrm{OS}(l).\; \operatorname{zen}(s, a))(a_{1}) \\ \quad =\; (\lambda a \in \mathrm{OS}(l).\; \operatorname{zen}(s, a))(a_{2}) \Rightarrow \\ \quad \quad (a_{1} = a_{2}) \end{array}}}
\vnbentry{\texttt{zen-\allowbreak{}enum-\allowbreak{}surjective}}{\vnbstmt{\begin{array}{@{}l@{}} \forall l \in \mathrm{Ord}, s \in \mathbf{Set}, w \in s. \\ \quad \neg \exists x.\; (x \in \{y \in s : \forall c.\; (\operatorname{ord\text{-}lt}(c, l) \Rightarrow \neg (\operatorname{zen}(s, c) = y))\}) \Rightarrow \\ \quad \quad \exists z \in \mathrm{OS}(l). \\ \quad \quad \quad ((\lambda a \in \mathrm{OS}(l).\; \operatorname{zen}(s, a))(z) = w) \end{array}}}
\vnbentry{\texttt{zermelo-\allowbreak{}bijection}}{\vnbstmt{\begin{array}{@{}l@{}} \forall s \in \mathbf{Set}. \exists l \in \mathrm{Ord}, i. \\ \quad (i \in \mathrm{Bij}(\mathrm{OS}(l), s)) \end{array}}}

\vnbfile{theorem-library/rake-ord-pigeonhole.scm}{theorem-library/rake-ord-pigeonhole.scm}
\noindent{\small rake-ord-pigeonhole.scm -- rake batch 5c, assignment 5c-U (2026-09-18). --------------------------------------------------------------------------- NOTE, 2026-09-20 (batch 9-B).  This file was written while cardinality was AXIOMATISED under the name CARD and the defined constant was its companion}\par
\vnbentry{\texttt{ord-\allowbreak{}segment-\allowbreak{}pigeonhole-\allowbreak{}guarded}}{\vnbstmt{\begin{array}{@{}l@{}} \forall l, e \in \mathbb{N}, f. \\ \quad \operatorname{ord\text{-}lt}(e, l) \Rightarrow  \\ \quad \quad \neg (f \in \operatorname{injection}(\mathrm{OS}(l), \mathrm{OS}(e))) \end{array}}}
\vnbentry{\texttt{card-\allowbreak{}body-\allowbreak{}exists}}{\vnbstmt{\begin{array}{@{}l@{}} \forall s \in \mathbf{Set}. \exists l \in \mathrm{Ord}. \\ \quad \exists i.\; (i \in \mathrm{Bij}(s, \mathrm{OS}(l))) \wedge \\ \quad \forall a. \\ \quad \quad \operatorname{ord\text{-}lt}(a, l) \Rightarrow  \\ \quad \quad \quad \neg \exists b.\; (b \in \mathrm{Bij}(s, \mathrm{OS}(a))) \end{array}}}
\vnbentry{\texttt{card-\allowbreak{}zermelo}}{\vnbstmt{\begin{array}{@{}l@{}} \forall s \in \mathbf{Set}. \\ \quad ((|s| \in \mathrm{Ord}) \wedge \exists i.\; (i \in \mathrm{Bij}(s, \mathrm{OS}(|s|)))) \end{array}}}
\vnbentry{\texttt{well-\allowbreak{}ordering-\allowbreak{}principle}}{\vnbstmt{\begin{array}{@{}l@{}} \forall s \in \mathbf{Set}. \exists i. \\ \quad (i \in \mathrm{Bij}(\mathrm{OS}(|s|), s)) \end{array}}}

\vnbfile{theorem-library/rake-card-star-laws.scm}{theorem-library/rake-card-star-laws.scm}
\noindent{\small rake-card-star-laws.scm -- rake batch 5c, assignment 5c-X (2026-09-18). --------------------------------------------------------------------------- NOTE, 2026-09-20 (batch 9-B).  This file was written while cardinality was AXIOMATISED under the name CARD and the defined constant was its companion}\par
\vnbentry{\texttt{card-\allowbreak{}insert-\allowbreak{}curried}}{\vnbstmt{\begin{array}{@{}l@{}} \forall a, x \in \mathbf{Set}. \\ \quad ((|a| \in \mathbb{N}) \wedge \neg (x \in a)) \Rightarrow  \\ \quad \quad (|(a \cup \{x\})| = (|a| + 1)) \end{array}}}
\vnbentry{\texttt{seg-\allowbreak{}restrict-\allowbreak{}is-\allowbreak{}bijection}}{\vnbstmt{\begin{array}{@{}l@{}} \forall n \in \mathbb{N}, b.; \forall g \in \mathrm{Bij}(\mathrm{OS}((n + 1)), b). \\ \quad (\lambda i \in \mathrm{OS}(n).\; g(i)) \\ \quad \in\; \mathrm{Bij}(\mathrm{OS}(n), \operatorname{image}(g, \mathrm{OS}(n))) \end{array}}}
\vnbentry{\texttt{card-\allowbreak{}seg-\allowbreak{}image}}{\vnbstmt{\begin{array}{@{}l@{}} \forall n \in \mathbb{N}, b.; \forall g \in \mathrm{Bij}(\mathrm{OS}((n + 1)), b). \\ \quad (|\operatorname{image}(g, \mathrm{OS}(n))| = n) \end{array}}}
\vnbentry{\texttt{seg-\allowbreak{}image-\allowbreak{}subset}}{\vnbstmt{\begin{array}{@{}l@{}} \forall n \in \mathbb{N}, b.; \forall g \in \mathrm{Bij}(\mathrm{OS}((n + 1)), b), w \in \operatorname{image}(g, \mathrm{OS}(n)). \\ \quad (w \in b) \end{array}}}
\vnbentry{\texttt{seg-\allowbreak{}image-\allowbreak{}omits-\allowbreak{}top}}{\vnbstmt{\begin{array}{@{}l@{}} \forall n \in \mathbb{N}, b.; \forall g \in \mathrm{Bij}(\mathrm{OS}((n + 1)), b). \\ \quad \neg (g(n) \in \operatorname{image}(g, \mathrm{OS}(n))) \end{array}}}
\vnbentry{\texttt{seg-\allowbreak{}image-\allowbreak{}insert}}{\vnbstmt{\begin{array}{@{}l@{}} \forall n \in \mathbb{N}, b.; \forall g \in \mathrm{Bij}(\mathrm{OS}((n + 1)), b). \\ \quad (b = (\operatorname{image}(g, \mathrm{OS}(n)) \cup \{g(n)\})) \end{array}}}
\vnbentry{\texttt{card-\allowbreak{}zero-\allowbreak{}is-\allowbreak{}empty}}{\vnbstmt{\forall b \in \mathbf{Set}.\; ((|b| = 0) \Rightarrow (b = \emptyset))}}
\vnbentry{\texttt{intersection-\allowbreak{}empty-\allowbreak{}transfer}}{\vnbstmt{\begin{array}{@{}l@{}} \forall a \in \mathbf{Set}, b, c. \\ \quad (((a \cap b) = \emptyset) \wedge \forall w \in c.\; (w \in b)) \Rightarrow  \\ \quad \quad ((a \cap c) = \emptyset) \end{array}}}
\vnbentry{\texttt{card-\allowbreak{}union-\allowbreak{}disjoint-\allowbreak{}ind}}{\vnbstmt{\begin{array}{@{}l@{}} \forall n \in \mathbb{N}, a, b \in \mathbf{Set}. \\ \quad ((|a| \in \mathbb{N}) \wedge ((|b| = n) \wedge ((a \cap b) = \emptyset))) \Rightarrow  \\ \quad \quad (|(a \cup b)| = (|a| + n)) \end{array}}}
\vnbentry{\texttt{card-\allowbreak{}union-\allowbreak{}disjoint-\allowbreak{}curried}}{\vnbstmt{\begin{array}{@{}l@{}} \forall a, b \in \mathbf{Set}. \\ \quad ((|a| \in \mathbb{N}) \wedge ((|b| \in \mathbb{N}) \wedge ((a \cap b) = \emptyset))) \Rightarrow  \\ \quad \quad (|(a \cup b)| = (|a| + |b|)) \end{array}}}
\vnbentry{\texttt{card-\allowbreak{}finite-\allowbreak{}induction-\allowbreak{}aux}}{\vnbstmt{\begin{array}{@{}l@{}} \forall n \in \mathbb{N}, c.; \forall s \in \mathbf{Set}. \\ \quad (\emptyset \in c) \wedge \\ \quad \forall s, x \in \mathbf{Set}. \\ \quad \quad (((|s| \in \mathbb{N}) \wedge (s \in c)) \wedge \neg (x \in s)) \Rightarrow  \\ \quad \quad \quad ((s \cup \{x\}) \in c) \wedge \\ \quad (|s| = n) \Rightarrow \\ \quad \quad (s \in c) \end{array}}}
\vnbentry{\texttt{finite-\allowbreak{}set-\allowbreak{}induction}}{\vnbstmt{\begin{array}{@{}l@{}} \forall c.; \forall s \in \mathbf{Set}. \\ \quad (\emptyset \in c) \wedge \\ \quad \forall s, x \in \mathbf{Set}. \\ \quad \quad (((|s| \in \mathbb{N}) \wedge (s \in c)) \wedge \neg (x \in s)) \Rightarrow  \\ \quad \quad \quad ((s \cup \{x\}) \in c) \wedge \\ \quad (|s| \in \mathbb{N}) \Rightarrow \\ \quad \quad (s \in c) \end{array}}}
\vnbentry{\texttt{injection-\allowbreak{}image-\allowbreak{}is-\allowbreak{}bijection}}{\vnbstmt{\begin{array}{@{}l@{}} \forall m \in \mathbf{Set}, d.; \forall h \in \operatorname{injection}(m, d). \\ \quad ((\lambda z \in m.\; h(z)) \in \mathrm{Bij}(m, \operatorname{image}(h, m))) \end{array}}}
\vnbentry{\texttt{card-\allowbreak{}image-\allowbreak{}injection-\allowbreak{}curried}}{\vnbstmt{\begin{array}{@{}l@{}} \forall m \in \mathbf{Set}, d.; \forall h \in \operatorname{injection}(m, d). \\ \quad ((|m| \in \mathbb{N}) \Rightarrow (|\operatorname{image}(h, m)| = |m|)) \end{array}}}

\vnbfile{theorem-library/card-laws.scm}{theorem-library/card-laws.scm}
\noindent{\small theorem-library/card-laws.scm -- the former CARD axioms whose counterpart has a different SHAPE, restated character for character as the axioms stood and proven (2026-09-20, batch 9-B: CARD := CARD-STAR, the user's decision).}\par
\vnbentry{\texttt{card-\allowbreak{}in-\allowbreak{}ord}}{\vnbstmt{\forall a \in \mathbf{Set}.\; (|a| \in \mathrm{Ord})}}
\vnbentry{\texttt{card-\allowbreak{}insert}}{\vnbstmt{\begin{array}{@{}l@{}} \forall a, x \in \mathbf{Set}. \\ \quad ((|a| \in \mathbb{N}) \wedge \neg (x \in a)) \Rightarrow  \\ \quad \quad (|(a \cup \{x\})| = (|a| + 1)) \end{array}}}
\vnbentry{\texttt{card-\allowbreak{}finite-\allowbreak{}bij}}{\vnbstmt{\begin{array}{@{}l@{}} \forall a \in \mathbf{Set}. \\ \quad (|a| \in \mathbb{N}) \Rightarrow  \\ \quad \quad \exists i.\; (i \in \mathrm{Bij}(\mathrm{OS}(|a|), a)) \end{array}}}
\vnbentry{\texttt{card-\allowbreak{}union-\allowbreak{}disjoint}}{\vnbstmt{\begin{array}{@{}l@{}} \forall a, b \in \mathbf{Set}. \\ \quad ((|a| \in \mathbb{N}) \wedge ((|b| \in \mathbb{N}) \wedge ((a \cap b) = \emptyset))) \Rightarrow  \\ \quad \quad (|(a \cup b)| = (|a| + |b|)) \end{array}}}
\vnbentry{\texttt{card-\allowbreak{}image-\allowbreak{}injection}}{\vnbstmt{\begin{array}{@{}l@{}} \forall m, d.; \forall i \in \operatorname{injection}(m, d). \\ \quad ((m \in \mathbf{Set}) \wedge (|m| \in \mathbb{N})) \Rightarrow  \\ \quad \quad (|\operatorname{image}(i, m)| = |m|) \end{array}}}

\vnbfile{theorem-library/card-singleton-proof.scm}{theorem-library/card-singleton-proof.scm}
\noindent{\small card-singleton-proof.scm -- the cardinality of a singleton is 1, proven.}\par
\vnbentry{\texttt{card-\allowbreak{}singleton}}{\vnbstmt{\forall x \in \mathbf{Set}.\; (|\{x\}| = (0 + 1))}}

\vnbfile{theorem-library/rake-fin-enum.scm}{theorem-library/rake-fin-enum.scm}
\noindent{\small theorem-library/rake-fin-enum.scm -- fin-enum-is-bijection, moved out of theorem-library/rake-inverse-bij.scm on 2026-09-20 (batch 9-B, CARD := CARD-STAR).}\par
\vnbentry{\texttt{fin-\allowbreak{}enum-\allowbreak{}is-\allowbreak{}bijection}}{\vnbstmt{\begin{array}{@{}l@{}} \forall s \in \mathbf{Set}. \\ \quad (|s| \in \mathbb{N}) \Rightarrow  \\ \quad \quad (\operatorname{fin\text{-}enum}(s) \in \mathrm{Bij}(\mathrm{OS}(|s|), s)) \end{array}}}

\vnbfile{theorem-library/difference-laws.scm}{theorem-library/difference-laws.scm}
\noindent{\small difference-laws.scm -- the two DIFFERENCE laws, PROVEN.}\par
\vnbentry{\texttt{difference-\allowbreak{}membership}}{\vnbstmt{\begin{array}{@{}l@{}} \forall u, b, x. \\ \quad ((x \in \operatorname{difference}(u, b)) \Leftrightarrow ((x \in u) \wedge \neg (x \in b))) \end{array}}}
\vnbentry{\texttt{difference-\allowbreak{}set}}{\vnbstmt{\begin{array}{@{}l@{}} \forall x \in \mathbf{Set}, b. \\ \quad (\operatorname{difference}(x, b) \in \mathbf{Set}) \end{array}}}

\vnbfile{theorem-library/intersection-of-laws.scm}{theorem-library/intersection-of-laws.scm}
\noindent{\small intersection-of-laws.scm -- the laws of INTERSECTION-OF, PROVEN.}\par
\vnbentry{\texttt{intersection-\allowbreak{}of-\allowbreak{}unfold}}{\vnbstmt{\begin{array}{@{}l@{}} \forall c. \\ \quad (\operatorname{intersection\text{-}of}(c) \simeq \{v \in \operatorname{big\text{-}union}(w, c, w) : \forall u \in c.\; (v \in u)\}) \end{array}}}
\vnbentry{\texttt{intersection-\allowbreak{}of-\allowbreak{}membership}}{\vnbstmt{\begin{array}{@{}l@{}} \forall c, x. \\ \quad ((x \in \operatorname{intersection\text{-}of}(c)) \Leftrightarrow (\exists u.\; (u \in c) \wedge \forall u \in c.\; (x \in u))) \end{array}}}
\vnbentry{\texttt{intersection-\allowbreak{}of-\allowbreak{}subset-\allowbreak{}member}}{\vnbstmt{\begin{array}{@{}l@{}} \forall c.; \forall a \in c. \\ \quad (\operatorname{intersection\text{-}of}(c) \subseteq a) \end{array}}}
\vnbentry{\texttt{intersection-\allowbreak{}of-\allowbreak{}empty-\allowbreak{}family}}{\vnbstmt{\begin{array}{@{}l@{}} \forall c. \\ \quad \neg \exists u.\; (u \in c) \Rightarrow  \\ \quad \quad (\operatorname{intersection\text{-}of}(c) = \emptyset) \end{array}}}
\vnbentry{\texttt{intersection-\allowbreak{}of-\allowbreak{}is-\allowbreak{}set}}{\vnbstmt{\forall c.\; (\operatorname{intersection\text{-}of}(c) \in \mathbf{Set})}}
\vnbentry{\texttt{intersection-\allowbreak{}of-\allowbreak{}antitone}}{\vnbstmt{\begin{array}{@{}l@{}} \forall c, d. \\ \quad ((c \subseteq d) \wedge \exists u.\; (u \in c)) \Rightarrow  \\ \quad \quad (\operatorname{intersection\text{-}of}(d) \subseteq \operatorname{intersection\text{-}of}(c)) \end{array}}}
\vnbentry{\texttt{complement-\allowbreak{}in-\allowbreak{}big-\allowbreak{}union}}{\vnbstmt{\begin{array}{@{}l@{}} \forall y, c, x. \\ \quad ((x \in \operatorname{complement\text{-}in}(y, \operatorname{big\text{-}union}(w, c, w))) \Leftrightarrow ((x \in y) \wedge \forall u \in c.\; \neg (x \in u))) \end{array}}}
\vnbentry{\texttt{complement-\allowbreak{}in-\allowbreak{}intersection-\allowbreak{}of}}{\vnbstmt{\begin{array}{@{}l@{}} \forall y, c. \\ \quad \exists u.\; (u \in c) \Rightarrow  \\ \quad \quad \operatorname{complement\text{-}in}(y, \operatorname{intersection\text{-}of}(c)) \\ \quad \quad =\; \operatorname{big\text{-}union}(w, c, \operatorname{complement\text{-}in}(y, w)) \end{array}}}

\vnbfile{theorem-library/rake-complement-laws.scm}{theorem-library/rake-complement-laws.scm}
\noindent{\small rake-complement-laws.scm -- two set-theoretic bricks the compactness duality needs, PROVEN (2026-09-20, batch 11-E; helper prefix `rcl-').}\par
\vnbentry{\texttt{complement-\allowbreak{}in-\allowbreak{}double}}{\vnbstmt{\begin{array}{@{}l@{}} \forall y, a. \\ \quad (a \subseteq y) \Rightarrow  \\ \quad \quad (\operatorname{complement\text{-}in}(y, \operatorname{complement\text{-}in}(y, a)) = a) \end{array}}}
\vnbentry{\texttt{family-\allowbreak{}of-\allowbreak{}subsets-\allowbreak{}is-\allowbreak{}set}}{\vnbstmt{\begin{array}{@{}l@{}} \forall y, c. \\ \quad ((y \in \mathbf{Set}) \wedge \forall u \in c.\; (u \subseteq y)) \Rightarrow  \\ \quad \quad (c \in \mathbf{Set}) \end{array}}}

\vnbfile{theorem-library/rake-fun-codomain.scm}{theorem-library/rake-fun-codomain.scm}
\noindent{\small theorem-library/rake-fun-codomain.scm -- codomain widening for FUN, PROVEN.}\par
\vnbentry{\texttt{fun-\allowbreak{}codomain-\allowbreak{}superset}}{\vnbstmt{\begin{array}{@{}l@{}} \forall f, a, b, c. \\ \quad (((f \in (a \to b)) \wedge (b \subseteq c)) \Rightarrow (f \in (a \to c))) \end{array}}}

\vnbfile{theorem-library/rake-rr-order-leaves.scm}{theorem-library/rake-rr-order-leaves.scm}
\noindent{\small theorem-library/rake-rr-order-leaves.scm -- the two asserted order leaves of `deriv-pos-strictly-increasing', PROVEN.}\par
\vnbentry{\texttt{rr-\allowbreak{}lt-\allowbreak{}from-\allowbreak{}diff-\allowbreak{}pos}}{\vnbstmt{\forall x, y \in \mathbb{R}.\; ((0 < (y - x)) \Rightarrow (x < y))}}

\vnbfile{theorem-library/rake-infinity-points.scm}{theorem-library/rake-infinity-points.scm}
\noindent{\small theorem-library/rake-infinity-points.scm -- the five DERIVABLE axioms of structure-library/extended-reals.scm, proven.  Rake "infinity points".}\par
\vnbentry{\texttt{rr-\allowbreak{}subset-\allowbreak{}rr-\allowbreak{}star}}{\vnbstmt{(\mathbb{R} \subseteq \operatorname{rr\text{-}star})}}
\vnbentry{\texttt{pos-\allowbreak{}inf-\allowbreak{}in-\allowbreak{}rr-\allowbreak{}star}}{\vnbstmt{(\operatorname{pos\text{-}inf} \in \operatorname{rr\text{-}star})}}
\vnbentry{\texttt{neg-\allowbreak{}inf-\allowbreak{}in-\allowbreak{}rr-\allowbreak{}star}}{\vnbstmt{(\operatorname{neg\text{-}inf} \in \operatorname{rr\text{-}star})}}
\vnbentry{\texttt{pos-\allowbreak{}inf-\allowbreak{}not-\allowbreak{}in-\allowbreak{}rr}}{\vnbstmt{\neg (\operatorname{pos\text{-}inf} \in \mathbb{R})}}
\vnbentry{\texttt{pos-\allowbreak{}inf-\allowbreak{}neq-\allowbreak{}neg-\allowbreak{}inf}}{\vnbstmt{\neg (\operatorname{pos\text{-}inf} = \operatorname{neg\text{-}inf})}}

\vnbfile{theorem-library/rake-rr-pos-star.scm}{theorem-library/rake-rr-pos-star.scm}
\noindent{\small theorem-library/rake-rr-pos-star.scm -- the six CARRIER facts of RR-POS-STAR = [0, +inf], proven.  Rake batch 5c-K.}\par
\vnbentry{\texttt{zero-\allowbreak{}in-\allowbreak{}rr-\allowbreak{}pos-\allowbreak{}star}}{\vnbstmt{(0 \in \operatorname{rr\text{-}pos\text{-}star})}}
\vnbentry{\texttt{pos-\allowbreak{}inf-\allowbreak{}in-\allowbreak{}rr-\allowbreak{}pos-\allowbreak{}star}}{\vnbstmt{(\operatorname{pos\text{-}inf} \in \operatorname{rr\text{-}pos\text{-}star})}}
\vnbentry{\texttt{neg-\allowbreak{}inf-\allowbreak{}not-\allowbreak{}in-\allowbreak{}rr-\allowbreak{}pos-\allowbreak{}star}}{\vnbstmt{\neg (\operatorname{neg\text{-}inf} \in \operatorname{rr\text{-}pos\text{-}star})}}
\vnbentry{\texttt{rr-\allowbreak{}pos-\allowbreak{}star-\allowbreak{}nonneg}}{\vnbstmt{\forall x \in \operatorname{rr\text{-}pos\text{-}star}.\; (0 \leq x)}}
\vnbentry{\texttt{rr-\allowbreak{}pos-\allowbreak{}star-\allowbreak{}below-\allowbreak{}pos-\allowbreak{}inf}}{\vnbstmt{\begin{array}{@{}l@{}} \forall x \in \operatorname{rr\text{-}pos\text{-}star}. \\ \quad (x \leq \operatorname{pos\text{-}inf}) \end{array}}}
\vnbentry{\texttt{rr-\allowbreak{}pos-\allowbreak{}star-\allowbreak{}subset-\allowbreak{}rr-\allowbreak{}star}}{\vnbstmt{(\operatorname{rr\text{-}pos\text{-}star} \subseteq \operatorname{rr\text{-}star})}}

\vnbfile{theorem-library/rake-eplus-defined.scm}{theorem-library/rake-eplus-defined.scm}
\noindent{\small theorem-library/rake-eplus-defined.scm -- extended addition on [0,+inf] DEFINED, and its laws proven.  Rake batch 5c-V.}\par
\vnbentry{\texttt{apply-\allowbreak{}congruence-\allowbreak{}2}}{\vnbstmt{\forall f, g, u, v.\; ((f \simeq g) \Rightarrow (f(u, v) \simeq g(u, v)))}}
\vnbentry{\texttt{rr-\allowbreak{}pos-\allowbreak{}star-\allowbreak{}is-\allowbreak{}set}}{\vnbstmt{(\operatorname{rr\text{-}pos\text{-}star} \in \mathbf{Set})}}
\vnbentry{\texttt{eplus-\allowbreak{}in-\allowbreak{}fun}}{\vnbstmt{\begin{array}{@{}l@{}} \mathit{eplus} \\ \in\; ((\operatorname{rr\text{-}pos\text{-}star} \times \operatorname{rr\text{-}pos\text{-}star}) \to \operatorname{rr\text{-}pos\text{-}star}) \end{array}}}
\vnbentry{\texttt{eplus-\allowbreak{}pos-\allowbreak{}inf-\allowbreak{}left}}{\vnbstmt{\begin{array}{@{}l@{}} \forall y \in \operatorname{rr\text{-}pos\text{-}star}. \\ \quad (\operatorname{eplus}(\operatorname{pos\text{-}inf}, y) = \operatorname{pos\text{-}inf}) \end{array}}}
\vnbentry{\texttt{eplus-\allowbreak{}pos-\allowbreak{}inf-\allowbreak{}right}}{\vnbstmt{\begin{array}{@{}l@{}} \forall x \in \operatorname{rr\text{-}pos\text{-}star}. \\ \quad (\operatorname{eplus}(x, \operatorname{pos\text{-}inf}) = \operatorname{pos\text{-}inf}) \end{array}}}
\vnbentry{\texttt{eplus-\allowbreak{}real-\allowbreak{}defined}}{\vnbstmt{\begin{array}{@{}l@{}} \forall x, y. \\ \quad ((x \in \mathbb{R}) \wedge ((0 \leq x) \wedge ((y \in \mathbb{R}) \wedge (0 \leq y)))) \Rightarrow  \\ \quad \quad (\operatorname{eplus}(x, y) = (x + y)) \end{array}}}

\vnbfile{theorem-library/rake-etimes-defined.scm}{theorem-library/rake-etimes-defined.scm}
\noindent{\small theorem-library/rake-etimes-defined.scm -- extended MULTIPLICATION on [0,+inf] DEFINED, and its laws proven.  Rake batch 5c-Y.}\par
\vnbentry{\texttt{etimes-\allowbreak{}in-\allowbreak{}fun}}{\vnbstmt{\begin{array}{@{}l@{}} \mathit{etimes} \\ \in\; ((\operatorname{rr\text{-}pos\text{-}star} \times \operatorname{rr\text{-}pos\text{-}star}) \to \operatorname{rr\text{-}pos\text{-}star}) \end{array}}}
\vnbentry{\texttt{etimes-\allowbreak{}zero-\allowbreak{}pos-\allowbreak{}inf}}{\vnbstmt{(\operatorname{etimes}(0, \operatorname{pos\text{-}inf}) = 0)}}
\vnbentry{\texttt{etimes-\allowbreak{}pos-\allowbreak{}inf-\allowbreak{}zero}}{\vnbstmt{(\operatorname{etimes}(\operatorname{pos\text{-}inf}, 0) = 0)}}
\vnbentry{\texttt{etimes-\allowbreak{}pos-\allowbreak{}inf-\allowbreak{}left}}{\vnbstmt{\begin{array}{@{}l@{}} \forall y \in \operatorname{rr\text{-}pos\text{-}star}. \\ \quad \neg (y = 0) \Rightarrow  \\ \quad \quad (\operatorname{etimes}(\operatorname{pos\text{-}inf}, y) = \operatorname{pos\text{-}inf}) \end{array}}}
\vnbentry{\texttt{etimes-\allowbreak{}pos-\allowbreak{}inf-\allowbreak{}right}}{\vnbstmt{\begin{array}{@{}l@{}} \forall x \in \operatorname{rr\text{-}pos\text{-}star}. \\ \quad \neg (x = 0) \Rightarrow  \\ \quad \quad (\operatorname{etimes}(x, \operatorname{pos\text{-}inf}) = \operatorname{pos\text{-}inf}) \end{array}}}
\vnbentry{\texttt{etimes-\allowbreak{}real-\allowbreak{}defined}}{\vnbstmt{\begin{array}{@{}l@{}} \forall x, y. \\ \quad ((x \in \mathbb{R}) \wedge ((0 \leq x) \wedge ((y \in \mathbb{R}) \wedge (0 \leq y)))) \Rightarrow  \\ \quad \quad (\operatorname{etimes}(x, y) = (x \cdot y)) \end{array}}}

\vnbfile{theorem-library/rr-sup-approx.scm}{theorem-library/rr-sup-approx.scm}
\noindent{\small rr-sup-approx.scm -- the SUP APPROXIMATION lemma, PROVEN.}\par
\vnbentry{\texttt{rr-\allowbreak{}sup-\allowbreak{}approx}}{\vnbstmt{\begin{array}{@{}l@{}} \forall s, d. \\ \quad (s \subseteq \mathbb{R}) \wedge \\ \quad \exists x.\; (x \in s) \wedge \\ \quad \operatorname{rr\text{-}bounded\text{-}above}(s) \wedge \\ \quad (d \in \mathbb{R}) \wedge \\ \quad (0 < d) \Rightarrow \\ \quad \quad \exists w \in s.\; ((\operatorname{sup}(s) - d) < w) \end{array}}}

\vnbfile{theorem-library/nn-unbounded-in-rr.scm}{theorem-library/nn-unbounded-in-rr.scm}
\noindent{\small nn-unbounded-in-rr.scm -- the ARCHIMEDEAN PROPERTY, PROVEN.}\par
\vnbentry{\texttt{nn-\allowbreak{}unbounded-\allowbreak{}in-\allowbreak{}rr}}{\vnbstmt{\forall x \in \mathbb{R}.\; \exists n \in \mathbb{N}.\; (x < n)}}

\vnbfile{theorem-library/nn-recip-succ-small.scm}{theorem-library/nn-recip-succ-small.scm}
\noindent{\small nn-recip-succ-small.scm -- 1/(n+1) gets below any positive eps, PROVEN.}\par
\vnbentry{\texttt{nn-\allowbreak{}recip-\allowbreak{}succ-\allowbreak{}small}}{\vnbstmt{\begin{array}{@{}l@{}} \forall s. \\ \quad \operatorname{pos\text{-}rr}(s) \Rightarrow  \\ \quad \quad \exists n \in \mathbb{N}.\; (\operatorname{recip}((n + 1)) < s) \end{array}}}

\vnbfile{theorem-library/discrete-space.scm}{theorem-library/discrete-space.scm}
\noindent{\small discrete-space.scm -- THE FIRST WITNESS OF A MEASURABLE SPACE, and the first outright inhabitant of TOP-SPACE.}\par
\vnbentry{\texttt{subset-\allowbreak{}refl}}{\vnbstmt{\forall a.\; (a \subseteq a)}}
\vnbentry{\texttt{power-\allowbreak{}mem-\allowbreak{}sethood}}{\vnbstmt{\forall a.; \forall b \in \operatorname{power}(a).\; (b \in \mathbf{Set})}}
\vnbentry{\texttt{power-\allowbreak{}mem-\allowbreak{}in}}{\vnbstmt{\begin{array}{@{}l@{}} \forall a, b, z. \\ \quad (((b \in \operatorname{power}(a)) \wedge (z \in b)) \Rightarrow (z \in a)) \end{array}}}
\vnbentry{\texttt{power-\allowbreak{}mem-\allowbreak{}intro}}{\vnbstmt{\begin{array}{@{}l@{}} \forall a.; \forall b \in \mathbf{Set}. \\ \quad (\forall z \in b.\; (z \in a) \Rightarrow (b \in \operatorname{power}(a))) \end{array}}}
\vnbentry{\texttt{power-\allowbreak{}whole-\allowbreak{}in}}{\vnbstmt{\forall a \in \mathbf{Set}.\; (a \in \operatorname{power}(a))}}
\vnbentry{\texttt{power-\allowbreak{}empty-\allowbreak{}in}}{\vnbstmt{\forall a.\; (\emptyset \in \operatorname{power}(a))}}
\vnbentry{\texttt{power-\allowbreak{}inter-\allowbreak{}closed}}{\vnbstmt{\begin{array}{@{}l@{}} \forall a, u, v. \\ \quad ((u \in \operatorname{power}(a)) \wedge (v \in \operatorname{power}(a))) \Rightarrow  \\ \quad \quad ((u \cap v) \in \operatorname{power}(a)) \end{array}}}
\vnbentry{\texttt{power-\allowbreak{}complement-\allowbreak{}closed}}{\vnbstmt{\begin{array}{@{}l@{}} \forall a \in \mathbf{Set}, b. \\ \quad (\operatorname{complement\text{-}in}(a, b) \in \operatorname{power}(a)) \end{array}}}
\vnbentry{\texttt{power-\allowbreak{}big-\allowbreak{}union-\allowbreak{}closed}}{\vnbstmt{\begin{array}{@{}l@{}} \forall a, m. \\ \quad ((a \in \mathbf{Set}) \wedge (m \subseteq \operatorname{power}(a))) \Rightarrow  \\ \quad \quad (\operatorname{big\text{-}union}(u, m, u) \in \operatorname{power}(a)) \end{array}}}
\vnbentry{\texttt{power-\allowbreak{}fun-\allowbreak{}union-\allowbreak{}closed}}{\vnbstmt{\begin{array}{@{}l@{}} \forall a.; \forall f \in (\mathbb{N} \to \operatorname{power}(a)). \\ \quad (\operatorname{big\text{-}union}(n, \mathbb{N}, f(n)) \in \operatorname{power}(a)) \end{array}}}
\vnbentry{\texttt{ds-\allowbreak{}pts}}{\vnbstmt{\begin{array}{@{}l@{}} \forall a. \\ \quad (\operatorname{pts}(\operatorname{discrete\text{-}space}(a)) \simeq a) \end{array}}}
\vnbentry{\texttt{ds-\allowbreak{}opens}}{\vnbstmt{\begin{array}{@{}l@{}} \forall a. \\ \quad (\operatorname{opens}(\operatorname{discrete\text{-}space}(a)) \simeq \operatorname{power}(a)) \end{array}}}
\vnbentry{\texttt{ds-\allowbreak{}sigma}}{\vnbstmt{\begin{array}{@{}l@{}} \forall a. \\ \quad (\sigma(\operatorname{discrete\text{-}space}(a)) \simeq \operatorname{power}(a)) \end{array}}}
\vnbentry{\texttt{discrete-\allowbreak{}space-\allowbreak{}is-\allowbreak{}top-\allowbreak{}space}}{\vnbstmt{\begin{array}{@{}l@{}} \forall a \in \mathbf{Set}. \\ \quad \operatorname{is\text{-}top\text{-}space}(\operatorname{discrete\text{-}space}(a)) \end{array}}}
\vnbentry{\texttt{discrete-\allowbreak{}space-\allowbreak{}is-\allowbreak{}measurable-\allowbreak{}space}}{\vnbstmt{\begin{array}{@{}l@{}} \forall a \in \mathbf{Set}. \\ \quad \operatorname{is\text{-}measurable\text{-}space}(\operatorname{discrete\text{-}space}(a)) \end{array}}}

\vnbfile{theorem-library/rake-measure.scm}{theorem-library/rake-measure.scm}
\noindent{\small rake-measure.scm -- rake batch 5, batch W: the light end of measure theory.}\par
\vnbentry{\texttt{sigma-\allowbreak{}algebra-\allowbreak{}ambient-\allowbreak{}set}}{\vnbstmt{\begin{array}{@{}l@{}} \forall a, b. \\ \quad \operatorname{is\text{-}sigma\text{-}algebra}(a, b) \Rightarrow  \\ \quad \quad (a \in \mathbf{Set}) \end{array}}}
\vnbentry{\texttt{sigma-\allowbreak{}algebra-\allowbreak{}whole-\allowbreak{}in}}{\vnbstmt{\begin{array}{@{}l@{}} \forall a, b. \\ \quad (\operatorname{is\text{-}sigma\text{-}algebra}(a, b) \Rightarrow (a \in b)) \end{array}}}
\vnbentry{\texttt{sigma-\allowbreak{}algebra-\allowbreak{}complement-\allowbreak{}closed}}{\vnbstmt{\begin{array}{@{}l@{}} \forall a, b, \alpha. \\ \quad (\operatorname{is\text{-}sigma\text{-}algebra}(a, b) \wedge (\alpha \in b)) \Rightarrow  \\ \quad \quad (\operatorname{complement\text{-}in}(a, \alpha) \in b) \end{array}}}
\vnbentry{\texttt{sigma-\allowbreak{}algebra-\allowbreak{}union-\allowbreak{}closed}}{\vnbstmt{\begin{array}{@{}l@{}} \forall a, b, f. \\ \quad (\operatorname{is\text{-}sigma\text{-}algebra}(a, b) \wedge (f \in (\mathbb{N} \to b))) \Rightarrow  \\ \quad \quad (\operatorname{big\text{-}union}(n, \mathbb{N}, f(n)) \in b) \end{array}}}
\vnbentry{\texttt{sigma-\allowbreak{}generated-\allowbreak{}in-\allowbreak{}power}}{\vnbstmt{\begin{array}{@{}l@{}} \forall a, e.; \forall \alpha \in \operatorname{sigma\text{-}generated}(a, e). \\ \quad (\alpha \in \operatorname{power}(a)) \end{array}}}
\vnbentry{\texttt{sigma-\allowbreak{}generated-\allowbreak{}in-\allowbreak{}every}}{\vnbstmt{\begin{array}{@{}l@{}} \forall a, e, \alpha, b. \\ \quad (\alpha \in \operatorname{sigma\text{-}generated}(a, e)) \wedge \\ \quad \operatorname{is\text{-}sigma\text{-}algebra}(a, b) \wedge \\ \quad (e \subseteq b) \Rightarrow \\ \quad \quad (\alpha \in b) \end{array}}}
\vnbentry{\texttt{sigma-\allowbreak{}generated-\allowbreak{}intro}}{\vnbstmt{\begin{array}{@{}l@{}} \forall a, e.; \forall \alpha \in \operatorname{power}(a). \\ \quad \forall b. \\ \quad \quad (\operatorname{is\text{-}sigma\text{-}algebra}(a, b) \wedge (e \subseteq b)) \Rightarrow  \\ \quad \quad \quad (\alpha \in b) \Rightarrow \\ \quad \quad (\alpha \in \operatorname{sigma\text{-}generated}(a, e)) \end{array}}}
\vnbentry{\texttt{sigma-\allowbreak{}generated-\allowbreak{}is-\allowbreak{}sigma-\allowbreak{}algebra}}{\vnbstmt{\begin{array}{@{}l@{}} \forall a, e. \\ \quad ((a \in \mathbf{Set}) \wedge (e \subseteq \operatorname{power}(a))) \Rightarrow  \\ \quad \quad \operatorname{is\text{-}sigma\text{-}algebra}(a, \operatorname{sigma\text{-}generated}(a, e)) \end{array}}}
\vnbentry{\texttt{measurable-\allowbreak{}map-\allowbreak{}source-\allowbreak{}algebra}}{\vnbstmt{\begin{array}{@{}l@{}} \forall a, b, a_{2}, c, f. \\ \quad \operatorname{is\text{-}measurable\text{-}map}(a, b, a_{2}, c, f) \Rightarrow  \\ \quad \quad \operatorname{is\text{-}sigma\text{-}algebra}(a, b) \end{array}}}
\vnbentry{\texttt{measurable-\allowbreak{}map-\allowbreak{}target-\allowbreak{}algebra}}{\vnbstmt{\begin{array}{@{}l@{}} \forall a, b, a_{2}, c, f. \\ \quad \operatorname{is\text{-}measurable\text{-}map}(a, b, a_{2}, c, f) \Rightarrow  \\ \quad \quad \operatorname{is\text{-}sigma\text{-}algebra}(a_{2}, c) \end{array}}}
\vnbentry{\texttt{measurable-\allowbreak{}map-\allowbreak{}in-\allowbreak{}fun}}{\vnbstmt{\begin{array}{@{}l@{}} \forall a, b, a_{2}, c, f. \\ \quad \operatorname{is\text{-}measurable\text{-}map}(a, b, a_{2}, c, f) \Rightarrow  \\ \quad \quad (f \in (a \to a_{2})) \end{array}}}
\vnbentry{\texttt{measurable-\allowbreak{}map-\allowbreak{}preimage}}{\vnbstmt{\begin{array}{@{}l@{}} \forall a, b, a_{2}, c, f, \beta. \\ \quad (\operatorname{is\text{-}measurable\text{-}map}(a, b, a_{2}, c, f) \wedge (\beta \in c)) \Rightarrow  \\ \quad \quad (\{x \in a : (f(x) \in \beta)\} \in b) \end{array}}}
\vnbentry{\texttt{measurable-\allowbreak{}map-\allowbreak{}intro}}{\vnbstmt{\begin{array}{@{}l@{}} \forall a, b, a_{2}, c, f. \\ \quad \operatorname{is\text{-}sigma\text{-}algebra}(a, b) \wedge \\ \quad \operatorname{is\text{-}sigma\text{-}algebra}(a_{2}, c) \wedge \\ \quad (f \in (a \to a_{2})) \wedge \\ \quad \forall \beta \in c.\; (\{x \in a : (f(x) \in \beta)\} \in b) \Rightarrow \\ \quad \quad \operatorname{is\text{-}measurable\text{-}map}(a, b, a_{2}, c, f) \end{array}}}
\vnbentry{\texttt{measurable-\allowbreak{}map-\allowbreak{}compose}}{\vnbstmt{\begin{array}{@{}l@{}} \forall a, b, a_{2}, c, a_{3}, d, f, g. \\ \quad \operatorname{is\text{-}measurable\text{-}map}(a, b, a_{2}, c, f) \wedge \\ \quad \operatorname{is\text{-}measurable\text{-}map}(a_{2}, c, a_{3}, d, g) \Rightarrow \\ \quad \quad \operatorname{is\text{-}measurable\text{-}map}(a, b, a_{3}, d, \operatorname{compose}(g, f)) \end{array}}}

\vnbfile{theorem-library/rake-measure2.scm}{theorem-library/rake-measure2.scm}
\noindent{\small rake-measure2.scm -- rake batch 5b-F: the indicator function, and the four SET-LEVEL bricks sigma-algebra.scm's header lists as "NEXT (not stated)".}\par
\vnbentry{\texttt{sigma-\allowbreak{}algebra-\allowbreak{}empty-\allowbreak{}in}}{\vnbstmt{\begin{array}{@{}l@{}} \forall a, b. \\ \quad \operatorname{is\text{-}sigma\text{-}algebra}(a, b) \Rightarrow  \\ \quad \quad (\emptyset \in b) \end{array}}}
\vnbentry{\texttt{sigma-\allowbreak{}algebra-\allowbreak{}union-\allowbreak{}closed-\allowbreak{}2}}{\vnbstmt{\begin{array}{@{}l@{}} \forall a, b, \alpha, \beta. \\ \quad (\operatorname{is\text{-}sigma\text{-}algebra}(a, b) \wedge ((\alpha \in b) \wedge (\beta \in b))) \Rightarrow  \\ \quad \quad ((\alpha \cup \beta) \in b) \end{array}}}
\vnbentry{\texttt{sigma-\allowbreak{}algebra-\allowbreak{}inter-\allowbreak{}closed-\allowbreak{}2}}{\vnbstmt{\begin{array}{@{}l@{}} \forall a, b, \alpha, \beta. \\ \quad (\operatorname{is\text{-}sigma\text{-}algebra}(a, b) \wedge ((\alpha \in b) \wedge (\beta \in b))) \Rightarrow  \\ \quad \quad ((\alpha \cap \beta) \in b) \end{array}}}
\vnbentry{\texttt{sigma-\allowbreak{}algebra-\allowbreak{}difference-\allowbreak{}closed}}{\vnbstmt{\begin{array}{@{}l@{}} \forall a, b, \alpha, \beta. \\ \quad (\operatorname{is\text{-}sigma\text{-}algebra}(a, b) \wedge ((\alpha \in b) \wedge (\beta \in b))) \Rightarrow  \\ \quad \quad (\operatorname{complement\text{-}in}(\alpha, \beta) \in b) \end{array}}}
\vnbentry{\texttt{measurable-\allowbreak{}fn-\allowbreak{}indicator}}{\vnbstmt{\begin{array}{@{}l@{}} \forall a, b, \alpha. \\ \quad (\operatorname{is\text{-}sigma\text{-}algebra}(a, b) \wedge (\alpha \in b)) \Rightarrow  \\ \quad \quad \operatorname{is\text{-}measurable\text{-}fn}(a, b, \operatorname{indicator}(a, \alpha)) \end{array}}}

\vnbfile{structure-library/subtype-laws.scm}{structure-library/subtype-laws.scm}
\noindent{\small subtype-laws.scm -- the trivial subtype-subsumption laws, PROVEN.}\par
\vnbentry{\texttt{abelian-\allowbreak{}group-\allowbreak{}is-\allowbreak{}group}}{\vnbstmt{\begin{array}{@{}l@{}} \forall s. \\ \quad \operatorname{is\text{-}abelian\text{-}group}(s) \Rightarrow  \\ \quad \quad \operatorname{is\text{-}group}(s) \end{array}}}
\vnbentry{\texttt{commutative-\allowbreak{}ring-\allowbreak{}is-\allowbreak{}ring}}{\vnbstmt{\begin{array}{@{}l@{}} \forall s. \\ \quad \operatorname{is\text{-}commutative\text{-}ring}(s) \Rightarrow  \\ \quad \quad \operatorname{is\text{-}ring}(s) \end{array}}}
\vnbentry{\texttt{integral-\allowbreak{}domain-\allowbreak{}is-\allowbreak{}commutative-\allowbreak{}ring}}{\vnbstmt{\begin{array}{@{}l@{}} \forall s. \\ \quad \operatorname{is\text{-}integral\text{-}domain}(s) \Rightarrow  \\ \quad \quad \operatorname{is\text{-}commutative\text{-}ring}(s) \end{array}}}
\vnbentry{\texttt{euclidean-\allowbreak{}ring-\allowbreak{}is-\allowbreak{}integral-\allowbreak{}domain}}{\vnbstmt{\begin{array}{@{}l@{}} \forall s. \\ \quad \operatorname{is\text{-}euclidean\text{-}ring}(s) \Rightarrow  \\ \quad \quad \operatorname{is\text{-}integral\text{-}domain}(s) \end{array}}}
\vnbentry{\texttt{comm-\allowbreak{}monoid-\allowbreak{}is-\allowbreak{}monoid}}{\vnbstmt{\begin{array}{@{}l@{}} \forall s. \\ \quad \operatorname{is\text{-}comm\text{-}monoid}(s) \Rightarrow  \\ \quad \quad \operatorname{is\text{-}monoid}(s) \end{array}}}
\vnbentry{\texttt{bijection-\allowbreak{}is-\allowbreak{}injection}}{\vnbstmt{\begin{array}{@{}l@{}} \forall x, y.; \forall i \in \mathrm{Bij}(x, y). \\ \quad (i \in \operatorname{injection}(x, y)) \end{array}}}
\vnbentry{\texttt{abelian-\allowbreak{}group-\allowbreak{}opr-\allowbreak{}comm}}{\vnbstmt{\begin{array}{@{}l@{}} \forall s.; \forall a, b \in \operatorname{carr}(s). \\ \quad \operatorname{is\text{-}abelian\text{-}group}(s) \Rightarrow  \\ \quad \quad (\operatorname{opr}(s)(a, b) = \operatorname{opr}(s)(b, a)) \end{array}}}
\vnbentry{\texttt{group-\allowbreak{}assoc}}{\vnbstmt{\begin{array}{@{}l@{}} \forall s.; \forall a, b, c \in \operatorname{carr}(s). \\ \quad \operatorname{is\text{-}group}(s) \Rightarrow  \\ \quad \quad \operatorname{opr}(s)(\operatorname{opr}(s)(a, b), c) \\ \quad \quad =\; \operatorname{opr}(s)(a, \operatorname{opr}(s)(b, c)) \end{array}}}
\vnbentry{\texttt{group-\allowbreak{}left-\allowbreak{}id}}{\vnbstmt{\begin{array}{@{}l@{}} \forall s.; \forall a \in \operatorname{carr}(s). \\ \quad \operatorname{is\text{-}group}(s) \Rightarrow  \\ \quad \quad (\operatorname{opr}(s)(\operatorname{iden}(s), a) = a) \end{array}}}
\vnbentry{\texttt{group-\allowbreak{}left-\allowbreak{}inv}}{\vnbstmt{\begin{array}{@{}l@{}} \forall s.; \forall a \in \operatorname{carr}(s). \\ \quad \operatorname{is\text{-}group}(s) \Rightarrow  \\ \quad \quad (\operatorname{opr}(s)(\operatorname{inv}(s)(a), a) = \operatorname{iden}(s)) \end{array}}}
\vnbentry{\texttt{group-\allowbreak{}identity-\allowbreak{}in}}{\vnbstmt{\begin{array}{@{}l@{}} \forall s. \\ \quad \operatorname{is\text{-}group}(s) \Rightarrow  \\ \quad \quad (\operatorname{iden}(s) \in \operatorname{carr}(s)) \end{array}}}

\vnbfile{theorem-library/op-typing.scm}{theorem-library/op-typing.scm}
\noindent{\small op-typing.scm -- the codomain typing of a structure operation, in APPLIED form, PROVEN.  One driver, seven theorems.}\par
\vnbentry{\texttt{ring-\allowbreak{}add-\allowbreak{}closed}}{\vnbstmt{\begin{array}{@{}l@{}} \forall s.; \forall a, b \in \operatorname{carr}(s). \\ \quad \operatorname{is\text{-}ring}(s) \Rightarrow  \\ \quad \quad (\operatorname{add}(s)(a, b) \in \operatorname{carr}(s)) \end{array}}}
\vnbentry{\texttt{ring-\allowbreak{}carrier-\allowbreak{}closed-\allowbreak{}mul}}{\vnbstmt{\begin{array}{@{}l@{}} \forall r.; \forall a, b \in \operatorname{carr}(r). \\ \quad \operatorname{is\text{-}ring}(r) \Rightarrow  \\ \quad \quad (\operatorname{mul}(r)(a, b) \in \operatorname{carr}(r)) \end{array}}}
\vnbentry{\texttt{ring-\allowbreak{}neg-\allowbreak{}in-\allowbreak{}carr}}{\vnbstmt{\begin{array}{@{}l@{}} \forall s.; \forall a \in \operatorname{carr}(s). \\ \quad \operatorname{is\text{-}ring}(s) \Rightarrow  \\ \quad \quad (\operatorname{neg}(s)(a) \in \operatorname{carr}(s)) \end{array}}}
\vnbentry{\texttt{bt-\allowbreak{}add-\allowbreak{}in-\allowbreak{}carr}}{\vnbstmt{\begin{array}{@{}l@{}} \forall r.; \forall a, b \in \operatorname{carr}(r). \\ \quad \operatorname{is\text{-}commutative\text{-}ring}(r) \Rightarrow  \\ \quad \quad (\operatorname{add}(r)(a, b) \in \operatorname{carr}(r)) \end{array}}}
\vnbentry{\texttt{metric-\allowbreak{}dist-\allowbreak{}real}}{\vnbstmt{\begin{array}{@{}l@{}} \forall s.; \forall x, y \in \operatorname{pts}(s). \\ \quad \operatorname{is\text{-}metric\text{-}space}(s) \Rightarrow  \\ \quad \quad (\operatorname{dist}(s)(x, y) \in \mathbb{R}) \end{array}}}
\vnbentry{\texttt{vnrm-\allowbreak{}real}}{\vnbstmt{\begin{array}{@{}l@{}} \forall m, w. \\ \quad (\operatorname{is\text{-}normed\text{-}vector\text{-}space}(m) \wedge (w \in \operatorname{vec}(m))) \Rightarrow  \\ \quad \quad (\operatorname{vnrm}(m)(w) \in \mathbb{R}) \end{array}}}
\vnbentry{\texttt{nvs-\allowbreak{}vadd-\allowbreak{}in-\allowbreak{}vec}}{\vnbstmt{\begin{array}{@{}l@{}} \forall m.; \forall x, y \in \operatorname{vec}(m). \\ \quad \operatorname{is\text{-}normed\text{-}vector\text{-}space}(m) \Rightarrow  \\ \quad \quad (\operatorname{vadd}(m)(x, y) \in \operatorname{vec}(m)) \end{array}}}

\vnbfile{theorem-library/finsum-type-proof.scm}{theorem-library/finsum-type-proof.scm}
\noindent{\small theorem-library/finsum-type-proof.scm -- the FINSUM typing floor, PROVEN.}\par
\vnbentry{\texttt{group-\allowbreak{}carrier-\allowbreak{}closed-\allowbreak{}opr}}{\vnbstmt{\begin{array}{@{}l@{}} \forall s.; \forall a, b \in \operatorname{carr}(s). \\ \quad \operatorname{is\text{-}group}(s) \Rightarrow  \\ \quad \quad (\operatorname{opr}(s)(a, b) \in \operatorname{carr}(s)) \end{array}}}
\vnbentry{\texttt{sum-\allowbreak{}ag-\allowbreak{}type-\allowbreak{}ind}}{\vnbstmt{\begin{array}{@{}l@{}} \forall n \in \mathbb{N}, g.; \forall f \in (\mathbb{N} \to \operatorname{carr}(g)). \\ \quad \operatorname{is\text{-}abelian\text{-}group}(g) \Rightarrow  \\ \quad \quad (\operatorname{sum\text{-}ag}(g, f, n) \in \operatorname{carr}(g)) \end{array}}}
\vnbentry{\texttt{sum-\allowbreak{}ag-\allowbreak{}type}}{\vnbstmt{\begin{array}{@{}l@{}} \forall g.; \forall f \in (\mathbb{N} \to \operatorname{carr}(g)), n \in \mathbb{N}. \\ \quad \operatorname{is\text{-}abelian\text{-}group}(g) \Rightarrow  \\ \quad \quad (\operatorname{sum\text{-}ag}(g, f, n) \in \operatorname{carr}(g)) \end{array}}}
\vnbentry{\texttt{enum-\allowbreak{}fam-\allowbreak{}in-\allowbreak{}fun}}{\vnbstmt{\begin{array}{@{}l@{}} \forall n \in \mathbb{N}, g, s.; \forall i \in (\mathrm{OS}(n) \to s), f \in (s \to \operatorname{carr}(g)). \\ \quad \operatorname{is\text{-}group}(g) \Rightarrow  \\ \quad \quad (\operatorname{enum\text{-}fam}(g, f, i, n) \in (\mathbb{N} \to \operatorname{carr}(g))) \end{array}}}
\vnbentry{\texttt{finsum-\allowbreak{}type}}{\vnbstmt{\begin{array}{@{}l@{}} \forall g.; \forall s \in \mathbf{Set}, f \in (s \to \operatorname{carr}(g)). \\ \quad (\operatorname{is\text{-}abelian\text{-}group}(g) \wedge (|s| \in \mathbb{N})) \Rightarrow  \\ \quad \quad (\operatorname{finsum}(g, f, s) \in \operatorname{carr}(g)) \end{array}}}
\vnbentry{\texttt{sum-\allowbreak{}ag-\allowbreak{}all-\allowbreak{}id-\allowbreak{}ind}}{\vnbstmt{\begin{array}{@{}l@{}} \forall n \in \mathbb{N}, g, a. \\ \quad \operatorname{is\text{-}abelian\text{-}group}(g) \wedge \\ \quad \forall i \in \mathrm{OS}(n).\; (a(i) = \operatorname{iden}(g)) \Rightarrow \\ \quad \quad (\operatorname{sum\text{-}ag}(g, a, n) = \operatorname{iden}(g)) \end{array}}}
\vnbentry{\texttt{finsum-\allowbreak{}all-\allowbreak{}id}}{\vnbstmt{\begin{array}{@{}l@{}} \forall g.; \forall s \in \mathbf{Set}, f \in (s \to \operatorname{carr}(g)). \\ \quad \operatorname{is\text{-}abelian\text{-}group}(g) \wedge \\ \quad (|s| \in \mathbb{N}) \wedge \\ \quad \forall z \in s.\; (f(z) = \operatorname{iden}(g)) \Rightarrow \\ \quad \quad (\operatorname{finsum}(g, f, s) = \operatorname{iden}(g)) \end{array}}}

\vnbfile{theorem-library/finsum-single-support.scm}{theorem-library/finsum-single-support.scm}
\noindent{\small finsum-single-support.scm -- a finite sum whose summand vanishes off a single index equals its value there.  Both theorems `modulo 0'.}\par
\vnbentry{\texttt{sum-\allowbreak{}ag-\allowbreak{}single-\allowbreak{}ind}}{\vnbstmt{\begin{array}{@{}l@{}} \forall n \in \mathbb{N}, g, a.; \forall p \in \mathrm{OS}(n). \\ \quad \operatorname{is\text{-}abelian\text{-}group}(g) \wedge \\ \quad (a(p) \in \operatorname{carr}(g)) \wedge \\ \quad \forall i \in \mathrm{OS}(n). \\ \quad \quad (\neg (i = p) \Rightarrow (a(i) = \operatorname{iden}(g))) \Rightarrow \\ \quad \quad (\operatorname{sum\text{-}ag}(g, a, n) = a(p)) \end{array}}}
\vnbentry{\texttt{finsum-\allowbreak{}single-\allowbreak{}support}}{\vnbstmt{\begin{array}{@{}l@{}} \forall g.; \forall s \in \mathbf{Set}, f \in (s \to \operatorname{carr}(g)), a \in s. \\ \quad \operatorname{is\text{-}abelian\text{-}group}(g) \wedge \\ \quad (|s| \in \mathbb{N}) \wedge \\ \quad \forall j \in s.\; (\neg (j = a) \Rightarrow (f(j) = \operatorname{iden}(g))) \Rightarrow \\ \quad \quad (\operatorname{finsum}(g, f, s) = f(a)) \end{array}}}

\vnbfile{structure-library/integral-domain-laws.scm}{structure-library/integral-domain-laws.scm}
\noindent{\small integral-domain-laws.scm -- integral-domain-cancel-zero, PROVEN from the definition.}\par
\vnbentry{\texttt{integral-\allowbreak{}domain-\allowbreak{}cancel-\allowbreak{}zero}}{\vnbstmt{\begin{array}{@{}l@{}} \forall s.; \forall a, b \in \operatorname{carr}(s). \\ \quad \operatorname{is\text{-}integral\text{-}domain}(s) \wedge \\ \quad (\operatorname{mul}(s)(a, b) = \operatorname{zero}(s)) \wedge \\ \quad \neg (b = \operatorname{zero}(s)) \Rightarrow \\ \quad \quad (a = \operatorname{zero}(s)) \end{array}}}
\vnbentry{\texttt{integral-\allowbreak{}domain-\allowbreak{}nontrivial}}{\vnbstmt{\begin{array}{@{}l@{}} \forall s. \\ \quad \operatorname{is\text{-}integral\text{-}domain}(s) \Rightarrow  \\ \quad \quad \neg (\operatorname{one}(s) = \operatorname{zero}(s)) \end{array}}}

\vnbfile{theorem-library/lambda-slot-apply.scm}{theorem-library/lambda-slot-apply.scm}
\noindent{\small lambda-slot-apply.scm -- taking an operation SLOT down to the surface.}\par
\vnbentry{\texttt{lam-\allowbreak{}slot-\allowbreak{}add-\allowbreak{}apply}}{\vnbstmt{\begin{array}{@{}l@{}} \forall m, a, b. \\ \quad ((a \in m) \wedge (b \in m)) \Rightarrow  \\ \quad \quad ((\lambda \langle \operatorname{x\_}, \operatorname{y\_} \rangle \in (m \times m).\; (\operatorname{x\_} + \operatorname{y\_}))(a, b) \simeq (a + b)) \end{array}}}
\vnbentry{\texttt{lam-\allowbreak{}slot-\allowbreak{}mul-\allowbreak{}apply}}{\vnbstmt{\begin{array}{@{}l@{}} \forall m, a, b. \\ \quad ((a \in m) \wedge (b \in m)) \Rightarrow  \\ \quad \quad ((\lambda \langle \operatorname{x\_}, \operatorname{y\_} \rangle \in (m \times m).\; (\operatorname{x\_} \cdot \operatorname{y\_}))(a, b) \simeq (a \cdot b)) \end{array}}}
\vnbentry{\texttt{lam-\allowbreak{}slot-\allowbreak{}neg-\allowbreak{}apply}}{\vnbstmt{\forall m.; \forall a \in m.\; ((\lambda x \in m.\; -x)(a) \simeq -a)}}

\vnbfile{theorem-library/zz-ring-is-ring.scm}{theorem-library/zz-ring-is-ring.scm}
\noindent{\small zz-ring-is-ring.scm -- IS-RING(ZZ-RING), PROVED.}\par
\vnbentry{\texttt{zz-\allowbreak{}is-\allowbreak{}ring}}{\vnbstmt{\operatorname{is\text{-}ring}(\operatorname{zz\text{-}ring})}}
\vnbentry{\texttt{qq-\allowbreak{}is-\allowbreak{}ring}}{\vnbstmt{\operatorname{is\text{-}ring}(\operatorname{qq\text{-}ring})}}

\vnbfile{theorem-library/rr-is-normed-field.scm}{theorem-library/rr-is-normed-field.scm}
\noindent{\small rr-is-normed-field.scm -- IS-NORMED-FIELD(RR-NORMED-FIELD), PROVED.}\par
\vnbentry{\texttt{rr-\allowbreak{}is-\allowbreak{}normed-\allowbreak{}field}}{\vnbstmt{\operatorname{is\text{-}normed\text{-}field}(\operatorname{rr\text{-}normed\text{-}field})}}

\vnbfile{theorem-library/rr-nvs-exemplification.scm}{theorem-library/rr-nvs-exemplification.scm}
\noindent{\small rr-nvs-exemplification.scm -- THE FIRST WITNESS OF A NORMED VECTOR SPACE.}\par
\vnbentry{\texttt{rr-\allowbreak{}nvs-\allowbreak{}is-\allowbreak{}normed-\allowbreak{}vector-\allowbreak{}space}}{\vnbstmt{\operatorname{is\text{-}normed\text{-}vector\text{-}space}(\operatorname{rr\text{-}nvs})}}

\vnbfile{theorem-library/zz-integral-domain.scm}{theorem-library/zz-integral-domain.scm}
\noindent{\small zz-integral-domain.scm -- IS-COMMUTATIVE-RING(ZZ-RING) and IS-INTEGRAL-DOMAIN(ZZ-RING), PROVED.}\par
\vnbentry{\texttt{zz-\allowbreak{}is-\allowbreak{}commutative-\allowbreak{}ring}}{\vnbstmt{\operatorname{is\text{-}commutative\text{-}ring}(\operatorname{zz\text{-}ring})}}
\vnbentry{\texttt{zz-\allowbreak{}no-\allowbreak{}zero-\allowbreak{}divisors}}{\vnbstmt{\begin{array}{@{}l@{}} \forall a, b \in \mathbb{Z}. \\ \quad (((a \cdot b) = 0) \Rightarrow ((a = 0) \vee (b = 0))) \end{array}}}
\vnbentry{\texttt{zz-\allowbreak{}is-\allowbreak{}integral-\allowbreak{}domain}}{\vnbstmt{\operatorname{is\text{-}integral\text{-}domain}(\operatorname{zz\text{-}ring})}}

\vnbfile{theorem-library/nn-add-monoid.scm}{theorem-library/nn-add-monoid.scm}
\noindent{\small nn-add-monoid.scm -- IS-COMM-MONOID(NN-ADD-MONOID), PROVED.}\par
\vnbentry{\texttt{nn-\allowbreak{}add-\allowbreak{}monoid-\allowbreak{}is-\allowbreak{}comm-\allowbreak{}monoid}}{\vnbstmt{\operatorname{is\text{-}comm\text{-}monoid}(\operatorname{nn\text{-}add\text{-}monoid})}}

\vnbfile{theorem-library/cancellation.scm}{theorem-library/cancellation.scm}
\noindent{\small cancellation.scm -- cancellation, proved ONCE where it lives, then carried.}\par
\vnbentry{\texttt{group-\allowbreak{}inv-\allowbreak{}in}}{\vnbstmt{\begin{array}{@{}l@{}} \forall s.; \forall a \in \operatorname{carr}(s). \\ \quad \operatorname{is\text{-}group}(s) \Rightarrow  \\ \quad \quad (\operatorname{inv}(s)(a) \in \operatorname{carr}(s)) \end{array}}}
\vnbentry{\texttt{group-\allowbreak{}cancel-\allowbreak{}left}}{\vnbstmt{\begin{array}{@{}l@{}} \forall s.; \forall a, b, c \in \operatorname{carr}(s). \\ \quad \operatorname{is\text{-}group}(s) \wedge \\ \quad (\operatorname{opr}(s)(c, a) = \operatorname{opr}(s)(c, b)) \Rightarrow \\ \quad \quad (a = b) \end{array}}}
\vnbentry{\texttt{ag-\allowbreak{}cancel-\allowbreak{}right}}{\vnbstmt{\begin{array}{@{}l@{}} \forall s.; \forall a, b, c \in \operatorname{carr}(s). \\ \quad \operatorname{is\text{-}abelian\text{-}group}(s) \wedge \\ \quad (\operatorname{opr}(s)(a, c) = \operatorname{opr}(s)(b, c)) \Rightarrow \\ \quad \quad (a = b) \end{array}}}
\vnbentry{\texttt{abelian-\allowbreak{}group-\allowbreak{}idempotent-\allowbreak{}is-\allowbreak{}id}}{\vnbstmt{\begin{array}{@{}l@{}} \forall s.; \forall a \in \operatorname{carr}(s). \\ \quad (\operatorname{is\text{-}abelian\text{-}group}(s) \wedge (\operatorname{opr}(s)(a, a) = a)) \Rightarrow  \\ \quad \quad (a = \operatorname{iden}(s)) \end{array}}}
\vnbentry{\texttt{zz-\allowbreak{}add-\allowbreak{}cancel}}{\vnbstmt{\forall a, b, c \in \mathbb{Z}.\; (((a + c) = (b + c)) \Rightarrow (a = b))}}
\vnbentry{\texttt{qq-\allowbreak{}add-\allowbreak{}cancel}}{\vnbstmt{\forall a, b, c \in \mathbb{Q}.\; (((a + c) = (b + c)) \Rightarrow (a = b))}}
\vnbentry{\texttt{nn-\allowbreak{}add-\allowbreak{}cancel}}{\vnbstmt{\forall a, b, c \in \mathbb{N}.\; (((a + c) = (b + c)) \Rightarrow (a = b))}}

\vnbfile{theorem-library/ring-zero-one-power.scm}{theorem-library/ring-zero-one-power.scm}
\noindent{\small ring-zero-one-power.scm -- four ring facts that were asserted, PROVEN.}\par
\vnbentry{\texttt{ring-\allowbreak{}one-\allowbreak{}in}}{\vnbstmt{\begin{array}{@{}l@{}} \forall s. \\ \quad \operatorname{is\text{-}ring}(s) \Rightarrow  \\ \quad \quad (\operatorname{one}(s) \in \operatorname{carr}(s)) \end{array}}}
\vnbentry{\texttt{ring-\allowbreak{}mul-\allowbreak{}zero-\allowbreak{}left}}{\vnbstmt{\begin{array}{@{}l@{}} \forall s.; \forall a \in \operatorname{carr}(s). \\ \quad \operatorname{is\text{-}ring}(s) \Rightarrow  \\ \quad \quad (\operatorname{mul}(s)(\operatorname{zero}(s), a) = \operatorname{zero}(s)) \end{array}}}
\vnbentry{\texttt{ring-\allowbreak{}mul-\allowbreak{}zero-\allowbreak{}right}}{\vnbstmt{\begin{array}{@{}l@{}} \forall s.; \forall a \in \operatorname{carr}(s). \\ \quad \operatorname{is\text{-}ring}(s) \Rightarrow  \\ \quad \quad (\operatorname{mul}(s)(a, \operatorname{zero}(s)) = \operatorname{zero}(s)) \end{array}}}
\vnbentry{\texttt{ring-\allowbreak{}power-\allowbreak{}zero}}{\vnbstmt{\begin{array}{@{}l@{}} \forall r, x. \\ \quad (\operatorname{ring\text{-}power}(r, x, 0) \simeq \operatorname{one}(r)) \end{array}}}

\vnbfile{theorem-library/nn-integral.scm}{theorem-library/nn-integral.scm}
\noindent{\small nn-integral.scm -- the NN arithmetic support layer.}\par
\vnbentry{\texttt{zz-\allowbreak{}mul-\allowbreak{}cancel-\allowbreak{}zero}}{\vnbstmt{\begin{array}{@{}l@{}} \forall a, b \in \mathbb{Z}. \\ \quad ((((a \cdot b) = 0) \wedge \neg (b = 0)) \Rightarrow (a = 0)) \end{array}}}
\vnbentry{\texttt{nn-\allowbreak{}mul-\allowbreak{}nonzero}}{\vnbstmt{\begin{array}{@{}l@{}} \forall a, b \in \mathbb{N}. \\ \quad ((\neg (a = 0) \wedge \neg (b = 0)) \Rightarrow \neg ((a \cdot b) = 0)) \end{array}}}
\vnbentry{\texttt{nn-\allowbreak{}mul-\allowbreak{}cancel}}{\vnbstmt{\begin{array}{@{}l@{}} \forall a, b, c \in \mathbb{N}. \\ \quad ((\neg (c = 0) \wedge ((a \cdot c) = (b \cdot c))) \Rightarrow (a = b)) \end{array}}}
\vnbentry{\texttt{nn-\allowbreak{}2-\allowbreak{}cancel}}{\vnbstmt{\forall x, y \in \mathbb{N}.\; (((2 \cdot x) = (2 \cdot y)) \Rightarrow (x = y))}}
\vnbentry{\texttt{nn-\allowbreak{}lt-\allowbreak{}double}}{\vnbstmt{\forall k \in \mathbb{N}.\; (\neg (k = 0) \Rightarrow (k < (2 \cdot k)))}}

\vnbfile{theorem-library/comb-kk-laws.scm}{theorem-library/comb-kk-laws.scm}
\noindent{\small comb-kk-laws.scm -- the six COMB-KK / Bernstein plumbing supports of theorem-library/binomial.scm, PROVEN.}\par
\vnbentry{\texttt{bt-\allowbreak{}succ-\allowbreak{}minus-\allowbreak{}1}}{\vnbstmt{\forall n \in \mathbb{N}.\; (((n + 1) - 1) = n)}}
\vnbentry{\texttt{sum-\allowbreak{}in-\allowbreak{}carr-\allowbreak{}zz-\allowbreak{}ind}}{\vnbstmt{\begin{array}{@{}l@{}} \forall n \in \mathbb{N}, r.; \forall g \in (\mathbb{Z} \to \operatorname{carr}(r)). \\ \quad \operatorname{is\text{-}commutative\text{-}ring}(r) \Rightarrow  \\ \quad \quad (\operatorname{sum}(r, g, n) \in \operatorname{carr}(r)) \end{array}}}
\vnbentry{\texttt{bt-\allowbreak{}sum-\allowbreak{}in-\allowbreak{}carr-\allowbreak{}zz}}{\vnbstmt{\begin{array}{@{}l@{}} \forall r.; \forall g \in (\mathbb{Z} \to \operatorname{carr}(r)), n \in \mathbb{N}. \\ \quad \operatorname{is\text{-}commutative\text{-}ring}(r) \Rightarrow  \\ \quad \quad (\operatorname{sum}(r, g, n) \in \operatorname{carr}(r)) \end{array}}}
\vnbentry{\texttt{comb-\allowbreak{}kk-\allowbreak{}0-\allowbreak{}0}}{\vnbstmt{\begin{array}{@{}l@{}} \forall r.; \forall x, y \in \operatorname{carr}(r). \\ \quad \operatorname{is\text{-}commutative\text{-}ring}(r) \Rightarrow  \\ \quad \quad (\operatorname{comb\text{-}kk}(r, x, y, 0)(0) = \operatorname{one}(r)) \end{array}}}
\vnbentry{\texttt{comb-\allowbreak{}kk-\allowbreak{}in-\allowbreak{}fun-\allowbreak{}ind}}{\vnbstmt{\begin{array}{@{}l@{}} \forall m \in \mathbb{N}, r.; \forall x, y \in \operatorname{carr}(r). \\ \quad \operatorname{is\text{-}commutative\text{-}ring}(r) \Rightarrow  \\ \quad \quad (\operatorname{comb\text{-}kk}(r, x, y, m) \in (\mathbb{Z} \to \operatorname{carr}(r))) \end{array}}}
\vnbentry{\texttt{comb-\allowbreak{}kk-\allowbreak{}in-\allowbreak{}fun}}{\vnbstmt{\begin{array}{@{}l@{}} \forall r.; \forall x, y \in \operatorname{carr}(r), m \in \mathbb{N}. \\ \quad \operatorname{is\text{-}commutative\text{-}ring}(r) \Rightarrow  \\ \quad \quad (\operatorname{comb\text{-}kk}(r, x, y, m) \in (\mathbb{Z} \to \operatorname{carr}(r))) \end{array}}}
\vnbentry{\texttt{comb-\allowbreak{}kk-\allowbreak{}null-\allowbreak{}ind}}{\vnbstmt{\begin{array}{@{}l@{}} \forall m \in \mathbb{N}, r.; \forall x, y \in \operatorname{carr}(r), k \in \mathbb{Z}. \\ \quad (\operatorname{is\text{-}commutative\text{-}ring}(r) \wedge (k < 0)) \Rightarrow  \\ \quad \quad (\operatorname{comb\text{-}kk}(r, x, y, m)(k) = \operatorname{zero}(r)) \end{array}}}
\vnbentry{\texttt{comb-\allowbreak{}kk-\allowbreak{}null}}{\vnbstmt{\begin{array}{@{}l@{}} \forall r.; \forall x, y \in \operatorname{carr}(r), m \in \mathbb{N}, k \in \mathbb{Z}. \\ \quad (\operatorname{is\text{-}commutative\text{-}ring}(r) \wedge (k < 0)) \Rightarrow  \\ \quad \quad (\operatorname{comb\text{-}kk}(r, x, y, m)(k) = \operatorname{zero}(r)) \end{array}}}
\vnbentry{\texttt{comb-\allowbreak{}kk-\allowbreak{}above-\allowbreak{}ind}}{\vnbstmt{\begin{array}{@{}l@{}} \forall m \in \mathbb{N}, r.; \forall x, y \in \operatorname{carr}(r), k \in \mathbb{Z}. \\ \quad (\operatorname{is\text{-}commutative\text{-}ring}(r) \wedge (m < k)) \Rightarrow  \\ \quad \quad (\operatorname{comb\text{-}kk}(r, x, y, m)(k) = \operatorname{zero}(r)) \end{array}}}
\vnbentry{\texttt{comb-\allowbreak{}kk-\allowbreak{}above}}{\vnbstmt{\begin{array}{@{}l@{}} \forall r.; \forall x, y \in \operatorname{carr}(r), m \in \mathbb{N}, k \in \mathbb{Z}. \\ \quad (\operatorname{is\text{-}commutative\text{-}ring}(r) \wedge (m < k)) \Rightarrow  \\ \quad \quad (\operatorname{comb\text{-}kk}(r, x, y, m)(k) = \operatorname{zero}(r)) \end{array}}}

\vnbfile{theorem-library/nn-mod3-proof.scm}{theorem-library/nn-mod3-proof.scm}
\noindent{\small nn-mod3-proof.scm -- mod-3 arithmetic on NN, the mirror of nn-parity-proof.scm.}\par
\vnbentry{\texttt{nn-\allowbreak{}three-\allowbreak{}is-\allowbreak{}succ3-\allowbreak{}zero}}{\vnbstmt{(3 = (((0 + 1) + 1) + 1))}}
\vnbentry{\texttt{nn-\allowbreak{}plus-\allowbreak{}three}}{\vnbstmt{\forall x \in \mathbb{N}.\; ((x + 3) = (((x + 1) + 1) + 1))}}
\vnbentry{\texttt{nn-\allowbreak{}three-\allowbreak{}mul-\allowbreak{}succ}}{\vnbstmt{\forall k \in \mathbb{N}.\; ((3 \cdot (k + 1)) = ((((3 \cdot k) + 1) + 1) + 1))}}
\vnbentry{\texttt{nn-\allowbreak{}trichotomy-\allowbreak{}3}}{\vnbstmt{\begin{array}{@{}l@{}} \forall n \in \mathbb{N}. \exists k \in \mathbb{N}. \\ \quad ((n = (3 \cdot k)) \vee ((n = ((3 \cdot k) + 1)) \vee (n = (((3 \cdot k) + 1) + 1)))) \end{array}}}
\vnbentry{\texttt{nn-\allowbreak{}3-\allowbreak{}not-\allowbreak{}succ}}{\vnbstmt{\forall x, y \in \mathbb{N}.\; \neg ((3 \cdot x) = ((3 \cdot y) + 1))}}
\vnbentry{\texttt{nn-\allowbreak{}3-\allowbreak{}div-\allowbreak{}square}}{\vnbstmt{\begin{array}{@{}l@{}} \forall p \in \mathbb{N}. \\ \quad \exists m \in \mathbb{N}.\; ((p \cdot p) = (3 \cdot m)) \Rightarrow  \\ \quad \quad \exists k \in \mathbb{N}.\; (p = (3 \cdot k)) \end{array}}}
\vnbentry{\texttt{nn-\allowbreak{}3-\allowbreak{}cancel}}{\vnbstmt{\forall x, y \in \mathbb{N}.\; (((3 \cdot x) = (3 \cdot y)) \Rightarrow (x = y))}}
\vnbentry{\texttt{nn-\allowbreak{}lt-\allowbreak{}triple}}{\vnbstmt{\forall k \in \mathbb{N}.\; (\neg (k = 0) \Rightarrow (k < (3 \cdot k)))}}

\vnbfile{theorem-library/nn-order-via-rr.scm}{theorem-library/nn-order-via-rr.scm}
\noindent{\small nn-order-via-rr.scm -- elementary NN order facts, PROVEN by taking the contradiction in RR.}\par
\vnbentry{\texttt{nn-\allowbreak{}not-\allowbreak{}le-\allowbreak{}zero-\allowbreak{}pos}}{\vnbstmt{\forall j \in \mathbb{N}.\; ((1 \leq j) \Rightarrow \neg (j \leq 0))}}
\vnbentry{\texttt{nn-\allowbreak{}succ-\allowbreak{}le-\allowbreak{}antisym}}{\vnbstmt{\forall a, b \in \mathbb{N}.\; (((a + 1) \leq b) \Rightarrow \neg (b \leq a))}}

\vnbfile{theorem-library/rake-intervals.scm}{theorem-library/rake-intervals.scm}
\noindent{\small rake-intervals.scm -- the interval / monus-predecessor supports of structure-library/order-lemmas.scm, PROVEN (batch K, 2026-09-17).}\par
\vnbentry{\texttt{nn-\allowbreak{}succ-\allowbreak{}le-\allowbreak{}cancel}}{\vnbstmt{\forall a, b \in \mathbb{N}.\; (((a + 1) \leq (b + 1)) \Rightarrow (a \leq b))}}
\vnbentry{\texttt{nn-\allowbreak{}minus-\allowbreak{}succ-\allowbreak{}1}}{\vnbstmt{\forall z \in \mathbb{N}.\; (\operatorname{nn\text{-}minus}((z + 1), 1) = z)}}
\vnbentry{\texttt{one-\allowbreak{}in-\allowbreak{}interval}}{\vnbstmt{\forall n \in \mathbb{N}.\; (1 \in \operatorname{interval}(1, (n + 1)))}}
\vnbentry{\texttt{succ-\allowbreak{}not-\allowbreak{}one}}{\vnbstmt{\forall q.; \forall z \in \operatorname{interval}(1, q).\; \neg ((z + 1) = 1)}}
\vnbentry{\texttt{nn-\allowbreak{}minus-\allowbreak{}1-\allowbreak{}inj}}{\vnbstmt{\begin{array}{@{}l@{}} \forall i, j. \\ \quad (i \in \mathbb{N}) \wedge \\ \quad (j \in \mathbb{N}) \wedge \\ \quad (1 \leq i) \wedge \\ \quad (1 \leq j) \wedge \\ \quad \neg (i = j) \Rightarrow \\ \quad \quad \neg (\operatorname{nn\text{-}minus}(i, 1) = \operatorname{nn\text{-}minus}(j, 1)) \end{array}}}
\vnbentry{\texttt{succ-\allowbreak{}in-\allowbreak{}interval}}{\vnbstmt{\begin{array}{@{}l@{}} \forall q \in \mathbb{N}, z \in \operatorname{interval}(1, q). \\ \quad ((z + 1) \in \operatorname{interval}(1, (q + 1))) \end{array}}}
\vnbentry{\texttt{pred-\allowbreak{}in-\allowbreak{}interval}}{\vnbstmt{\begin{array}{@{}l@{}} \forall p \in \mathbb{N}, i \in \operatorname{interval}(1, (p + 1)). \\ \quad \neg (i = 1) \Rightarrow  \\ \quad \quad (\operatorname{nn\text{-}minus}(i, 1) \in \operatorname{interval}(1, p)) \end{array}}}

\vnbfile{theorem-library/zz-order.scm}{theorem-library/zz-order.scm}
\noindent{\small zz-order.scm -- the ORDER THEORY OF ZZ, proven.}\par
\vnbentry{\texttt{zz-\allowbreak{}abs-\allowbreak{}in-\allowbreak{}nn}}{\vnbstmt{\forall a \in \mathbb{Z}.\; (\lvert a \rvert \in \mathbb{N})}}
\vnbentry{\texttt{zz-\allowbreak{}abs-\allowbreak{}cases}}{\vnbstmt{\forall a \in \mathbb{Z}.\; ((a = \lvert a \rvert) \vee (a = -\lvert a \rvert))}}
\vnbentry{\texttt{zz-\allowbreak{}trichotomy}}{\vnbstmt{\forall a \in \mathbb{Z}.\; ((a < 0) \vee ((a = 0) \vee (0 < a)))}}
\vnbentry{\texttt{zz-\allowbreak{}discrete}}{\vnbstmt{\forall a \in \mathbb{Z}.\; ((0 < a) \Rightarrow (1 \leq a))}}
\vnbentry{\texttt{zz-\allowbreak{}lt-\allowbreak{}succ-\allowbreak{}le}}{\vnbstmt{\forall a, b \in \mathbb{Z}.\; ((a < b) \Rightarrow ((a + 1) \leq b))}}
\vnbentry{\texttt{zz-\allowbreak{}nonneg-\allowbreak{}in-\allowbreak{}nn}}{\vnbstmt{\forall a \in \mathbb{Z}.\; ((0 \leq a) \Rightarrow (a \in \mathbb{N}))}}

\vnbfile{theorem-library/zz-division.scm}{theorem-library/zz-division.scm}
\noindent{\small zz-division.scm -- DIVISION WITH REMAINDER, and ZZ as a Euclidean ring.}\par
\vnbentry{\texttt{nn-\allowbreak{}division}}{\vnbstmt{\begin{array}{@{}l@{}} \forall a, b \in \mathbb{N}. \\ \quad \neg (b = 0) \Rightarrow  \\ \quad \quad \exists q, r \in \mathbb{N}.\; ((a = ((q \cdot b) + r)) \wedge ((r + 1) \leq b)) \end{array}}}
\vnbentry{\texttt{zz-\allowbreak{}division}}{\vnbstmt{\begin{array}{@{}l@{}} \forall a, b \in \mathbb{Z}. \\ \quad \neg (b = 0) \Rightarrow  \\ \quad \quad \exists q, r \in \mathbb{Z}. \\ \quad \quad \quad ((a = ((q \cdot b) + r)) \wedge ((r = 0) \vee ((\lvert r \rvert + 1) \leq \lvert b \rvert))) \end{array}}}
\vnbentry{\texttt{zz-\allowbreak{}is-\allowbreak{}euclidean-\allowbreak{}ring}}{\vnbstmt{\operatorname{is\text{-}euclidean\text{-}ring}(\operatorname{zz\text{-}ring})}}

\vnbfile{theorem-library/zz-parity-proof.scm}{theorem-library/zz-parity-proof.scm}
\noindent{\small zz-parity-proof.scm -- parity on the INTEGERS, lifted from NN.}\par
\vnbentry{\texttt{zz-\allowbreak{}parity}}{\vnbstmt{\forall a \in \mathbb{Z}.\; (\operatorname{even}(a) \vee \operatorname{odd}(a))}}
\vnbentry{\texttt{nn-\allowbreak{}2p-\allowbreak{}neq-\allowbreak{}1}}{\vnbstmt{\forall p \in \mathbb{N}.\; \neg ((2 \cdot p) = 1)}}
\vnbentry{\texttt{nn-\allowbreak{}2p-\allowbreak{}neq-\allowbreak{}neg1}}{\vnbstmt{\forall p \in \mathbb{N}.\; \neg ((2 \cdot p) = -1)}}
\vnbentry{\texttt{zz-\allowbreak{}2m-\allowbreak{}neq-\allowbreak{}1}}{\vnbstmt{\forall m \in \mathbb{Z}.\; \neg ((2 \cdot m) = 1)}}
\vnbentry{\texttt{zz-\allowbreak{}not-\allowbreak{}even-\allowbreak{}and-\allowbreak{}odd}}{\vnbstmt{\begin{array}{@{}l@{}} \forall x \in \mathbb{Z}. \\ \quad \neg (\operatorname{even}(x) \wedge \operatorname{odd}(x)) \end{array}}}
\vnbentry{\texttt{zz-\allowbreak{}even-\allowbreak{}iff-\allowbreak{}not-\allowbreak{}odd}}{\vnbstmt{\begin{array}{@{}l@{}} \forall x \in \mathbb{Z}. \\ \quad (\operatorname{even}(x) \Leftrightarrow \neg \operatorname{odd}(x)) \end{array}}}
\vnbentry{\texttt{zz-\allowbreak{}even-\allowbreak{}mul}}{\vnbstmt{\begin{array}{@{}l@{}} \forall x, y \in \mathbb{Z}. \\ \quad (\operatorname{even}(x) \Rightarrow \operatorname{even}((x \cdot y))) \end{array}}}
\vnbentry{\texttt{zz-\allowbreak{}odd-\allowbreak{}mul-\allowbreak{}odd}}{\vnbstmt{\begin{array}{@{}l@{}} \forall x, y \in \mathbb{Z}. \\ \quad (\operatorname{odd}(x) \wedge \operatorname{odd}(y)) \Rightarrow  \\ \quad \quad \operatorname{odd}((x \cdot y)) \end{array}}}
\vnbentry{\texttt{zz-\allowbreak{}odd-\allowbreak{}sq-\allowbreak{}odd}}{\vnbstmt{\begin{array}{@{}l@{}} \forall x \in \mathbb{Z}. \\ \quad (\operatorname{odd}(x) \Rightarrow \operatorname{odd}((x \cdot x))) \end{array}}}

\vnbfile{theorem-library/nn-pos-is-succ.scm}{theorem-library/nn-pos-is-succ.scm}
\noindent{\small nn-pos-is-succ.scm -- a natural at least 1 is a successor.}\par
\vnbentry{\texttt{nn-\allowbreak{}pos-\allowbreak{}is-\allowbreak{}succ}}{\vnbstmt{\begin{array}{@{}l@{}} \forall n \in \mathbb{N}. \\ \quad ((1 \leq n) \Rightarrow \exists q \in \mathbb{N}.\; (n = (q + 1))) \end{array}}}

\vnbfile{theorem-library/nn-order-proof.scm}{theorem-library/nn-order-proof.scm}
\noindent{\small nn-order-proof.scm -- three elementary facts about $<$= and + on NN, PROVEN.}\par
\vnbentry{\texttt{nn-\allowbreak{}le-\allowbreak{}zero-\allowbreak{}is-\allowbreak{}zero}}{\vnbstmt{\forall a \in \mathbb{N}.\; ((a \leq 0) \Rightarrow (a = 0))}}
\vnbentry{\texttt{nn-\allowbreak{}add-\allowbreak{}le-\allowbreak{}mono}}{\vnbstmt{\begin{array}{@{}l@{}} \forall c, a, b \in \mathbb{N}. \\ \quad ((b \leq c) \Rightarrow ((a + b) \leq (a + c))) \end{array}}}
\vnbentry{\texttt{nn-\allowbreak{}le-\allowbreak{}add-\allowbreak{}left}}{\vnbstmt{\begin{array}{@{}l@{}} \forall a, b. \\ \quad (((a \in \mathbb{N}) \wedge (b \in \mathbb{N})) \Rightarrow (b \leq (a + b))) \end{array}}}
\vnbentry{\texttt{nn-\allowbreak{}le-\allowbreak{}antisym}}{\vnbstmt{\begin{array}{@{}l@{}} \forall a, b \in \mathbb{N}. \\ \quad (((a \leq b) \wedge (b \leq a)) \Rightarrow (a = b)) \end{array}}}

\vnbfile{theorem-library/rake-algebra.scm}{theorem-library/rake-algebra.scm}
\noindent{\small rake-algebra.scm -- BATCH D of the 2026-09-17 rake: seven asserted leaves of the PSS, PROVEN.  Every statement here is its support's statement UNCHANGED.}\par
\vnbentry{\texttt{ring-\allowbreak{}carr-\allowbreak{}in-\allowbreak{}set}}{\vnbstmt{\begin{array}{@{}l@{}} \forall r. \\ \quad \operatorname{is\text{-}ring}(r) \Rightarrow  \\ \quad \quad (\operatorname{carr}(r) \in \mathbf{Set}) \end{array}}}
\vnbentry{\texttt{ideal-\allowbreak{}elt-\allowbreak{}in-\allowbreak{}carrier}}{\vnbstmt{\begin{array}{@{}l@{}} \forall s, i, x. \\ \quad (\operatorname{is\text{-}ideal}(s, i) \wedge (x \in i)) \Rightarrow  \\ \quad \quad (x \in \operatorname{carr}(s)) \end{array}}}
\vnbentry{\texttt{principal-\allowbreak{}ideal-\allowbreak{}in-\allowbreak{}ideal}}{\vnbstmt{\begin{array}{@{}l@{}} \forall s, i, b.; \forall x \in \operatorname{principal\text{-}ideal}(s, b). \\ \quad ((\operatorname{is\text{-}ideal}(s, i) \wedge (b \in i)) \Rightarrow (x \in i)) \end{array}}}
\vnbentry{\texttt{diagonal-\allowbreak{}off-\allowbreak{}entry}}{\vnbstmt{\begin{array}{@{}l@{}} \forall a, m, n, d.; \forall i \in \operatorname{interval}(1, m), j \in \operatorname{interval}(1, n). \\ \quad (\operatorname{is\text{-}diagonal}(a, m, n, d) \wedge \neg (i = j)) \Rightarrow  \\ \quad \quad (\operatorname{entry}(d, i, j) = \operatorname{zero}(a)) \end{array}}}
\vnbentry{\texttt{ring-\allowbreak{}add-\allowbreak{}right-\allowbreak{}id}}{\vnbstmt{\begin{array}{@{}l@{}} \forall s.; \forall a \in \operatorname{carr}(s). \\ \quad \operatorname{is\text{-}ring}(s) \Rightarrow  \\ \quad \quad (\operatorname{add}(s)(a, \operatorname{zero}(s)) = a) \end{array}}}
\vnbentry{\texttt{span-\allowbreak{}add-\allowbreak{}one-\allowbreak{}unfold}}{\vnbstmt{\begin{array}{@{}l@{}} \forall m, s, v. \\ \quad (\operatorname{span\text{-}add\text{-}one}(m, s, v) \simeq \{y \in \operatorname{vec}(m) : \exists x \in s, r \in \mathbb{R}.\; (y = \operatorname{vadd}(m)(x, \operatorname{act}(m)(r, v)))\}) \end{array}}}
\vnbentry{\texttt{span-\allowbreak{}add-\allowbreak{}one-\allowbreak{}membership}}{\vnbstmt{\begin{array}{@{}l@{}} \forall m, s, v, w. \\ \quad ((w \in \operatorname{span\text{-}add\text{-}one}(m, s, v)) \Leftrightarrow ((w \in \operatorname{vec}(m)) \wedge \exists y \in s, r \in \mathbb{R}.\; (w = \operatorname{vadd}(m)(y, \operatorname{act}(m)(r, v))))) \end{array}}}
\vnbentry{\texttt{strictly-\allowbreak{}mono-\allowbreak{}ge-\allowbreak{}id-\allowbreak{}ind}}{\vnbstmt{\begin{array}{@{}l@{}} \forall k \in \mathbb{N}, i. \\ \quad (\operatorname{strictly\text{-}mono\text{-}nn}(i) \Rightarrow (k \leq i(k))) \end{array}}}
\vnbentry{\texttt{strictly-\allowbreak{}mono-\allowbreak{}ge-\allowbreak{}id}}{\vnbstmt{\begin{array}{@{}l@{}} \forall i.; \forall k \in \mathbb{N}. \\ \quad (\operatorname{strictly\text{-}mono\text{-}nn}(i) \Rightarrow (k \leq i(k))) \end{array}}}

\vnbfile{theorem-library/rake-algebra2.scm}{theorem-library/rake-algebra2.scm}
\noindent{\small rake-algebra2.scm -- BATCH I of the 2026-09-17 rake: the algebra two-liners and the four unwarranted algebra axioms, PROVEN.  Every statement is its site's statement UNCHANGED (extracted from the site, not retyped).}\par
\vnbentry{\texttt{abelian-\allowbreak{}group-\allowbreak{}assoc}}{\vnbstmt{\begin{array}{@{}l@{}} \forall s.; \forall a, b, c \in \operatorname{carr}(s). \\ \quad \operatorname{is\text{-}abelian\text{-}group}(s) \Rightarrow  \\ \quad \quad \operatorname{opr}(s)(\operatorname{opr}(s)(a, b), c) \\ \quad \quad =\; \operatorname{opr}(s)(a, \operatorname{opr}(s)(b, c)) \end{array}}}
\vnbentry{\texttt{abelian-\allowbreak{}group-\allowbreak{}right-\allowbreak{}id}}{\vnbstmt{\begin{array}{@{}l@{}} \forall s.; \forall a \in \operatorname{carr}(s). \\ \quad \operatorname{is\text{-}abelian\text{-}group}(s) \Rightarrow  \\ \quad \quad (\operatorname{opr}(s)(a, \operatorname{iden}(s)) = a) \end{array}}}
\vnbentry{\texttt{abelian-\allowbreak{}group-\allowbreak{}inverse-\allowbreak{}unique}}{\vnbstmt{\begin{array}{@{}l@{}} \forall s.; \forall a, b \in \operatorname{carr}(s). \\ \quad \operatorname{is\text{-}abelian\text{-}group}(s) \wedge \\ \quad (\operatorname{opr}(s)(a, b) = \operatorname{iden}(s)) \Rightarrow \\ \quad \quad (b = \operatorname{inv}(s)(a)) \end{array}}}
\vnbentry{\texttt{ring-\allowbreak{}add-\allowbreak{}right-\allowbreak{}inv}}{\vnbstmt{\begin{array}{@{}l@{}} \forall s.; \forall a \in \operatorname{carr}(s). \\ \quad \operatorname{is\text{-}ring}(s) \Rightarrow  \\ \quad \quad (\operatorname{add}(s)(a, \operatorname{neg}(s)(a)) = \operatorname{zero}(s)) \end{array}}}
\vnbentry{\texttt{ring-\allowbreak{}add-\allowbreak{}inverse-\allowbreak{}unique}}{\vnbstmt{\begin{array}{@{}l@{}} \forall s.; \forall a, b \in \operatorname{carr}(s). \\ \quad (\operatorname{is\text{-}ring}(s) \wedge (\operatorname{add}(s)(a, b) = \operatorname{zero}(s))) \Rightarrow  \\ \quad \quad (b = \operatorname{neg}(s)(a)) \end{array}}}
\vnbentry{\texttt{ring-\allowbreak{}neg-\allowbreak{}neg}}{\vnbstmt{\begin{array}{@{}l@{}} \forall s.; \forall r \in \operatorname{carr}(s). \\ \quad \operatorname{is\text{-}ring}(s) \Rightarrow  \\ \quad \quad (\operatorname{neg}(s)(\operatorname{neg}(s)(r)) = r) \end{array}}}
\vnbentry{\texttt{ring-\allowbreak{}neg-\allowbreak{}mul-\allowbreak{}left}}{\vnbstmt{\begin{array}{@{}l@{}} \forall s.; \forall a, b \in \operatorname{carr}(s). \\ \quad \operatorname{is\text{-}ring}(s) \Rightarrow  \\ \quad \quad \operatorname{mul}(s)(\operatorname{neg}(s)(a), b) \\ \quad \quad =\; \operatorname{neg}(s)(\operatorname{mul}(s)(a, b)) \end{array}}}
\vnbentry{\texttt{matrix-\allowbreak{}sethood}}{\vnbstmt{\forall s \in \mathbf{Set}.\; (\operatorname{matrix}(s) \in \mathbf{Set})}}
\vnbentry{\texttt{monoid-\allowbreak{}identity-\allowbreak{}in}}{\vnbstmt{\begin{array}{@{}l@{}} \forall m. \\ \quad \operatorname{is\text{-}monoid}(m) \Rightarrow  \\ \quad \quad (\operatorname{iden}(m) \in \operatorname{carr}(m)) \end{array}}}
\vnbentry{\texttt{monoid-\allowbreak{}carrier-\allowbreak{}closed-\allowbreak{}opr}}{\vnbstmt{\begin{array}{@{}l@{}} \forall m, a, b. \\ \quad \operatorname{is\text{-}monoid}(m) \wedge \\ \quad (a \in \operatorname{carr}(m)) \wedge \\ \quad (b \in \operatorname{carr}(m)) \Rightarrow \\ \quad \quad (\operatorname{opr}(m)(a, b) \in \operatorname{carr}(m)) \end{array}}}
\vnbentry{\texttt{mpow-\allowbreak{}type-\allowbreak{}ind}}{\vnbstmt{\begin{array}{@{}l@{}} \forall n \in \mathbb{N}, m.; \forall x \in \operatorname{carr}(m). \\ \quad \operatorname{is\text{-}monoid}(m) \Rightarrow  \\ \quad \quad (\operatorname{mpow}(m, x, n) \in \operatorname{carr}(m)) \end{array}}}
\vnbentry{\texttt{mpow-\allowbreak{}type}}{\vnbstmt{\begin{array}{@{}l@{}} \forall m.; \forall x \in \operatorname{carr}(m), n \in \mathbb{N}. \\ \quad \operatorname{is\text{-}monoid}(m) \Rightarrow  \\ \quad \quad (\operatorname{mpow}(m, x, n) \in \operatorname{carr}(m)) \end{array}}}

\vnbfile{theorem-library/rake-monoid-laws.scm}{theorem-library/rake-monoid-laws.scm}
\noindent{\small rake-monoid-laws.scm -- BATCH 5c-P of the 2026-09-18 rake: the three MONOID shape projections, PROVEN.}\par
\vnbentry{\texttt{monoid-\allowbreak{}assoc}}{\vnbstmt{\begin{array}{@{}l@{}} \forall s.; \forall a, b, c \in \operatorname{carr}(s). \\ \quad \operatorname{is\text{-}monoid}(s) \Rightarrow  \\ \quad \quad \operatorname{opr}(s)(\operatorname{opr}(s)(a, b), c) \\ \quad \quad =\; \operatorname{opr}(s)(a, \operatorname{opr}(s)(b, c)) \end{array}}}
\vnbentry{\texttt{monoid-\allowbreak{}left-\allowbreak{}id}}{\vnbstmt{\begin{array}{@{}l@{}} \forall s.; \forall a \in \operatorname{carr}(s). \\ \quad \operatorname{is\text{-}monoid}(s) \Rightarrow  \\ \quad \quad (\operatorname{opr}(s)(\operatorname{iden}(s), a) = a) \end{array}}}
\vnbentry{\texttt{monoid-\allowbreak{}right-\allowbreak{}id}}{\vnbstmt{\begin{array}{@{}l@{}} \forall s.; \forall a \in \operatorname{carr}(s). \\ \quad \operatorname{is\text{-}monoid}(s) \Rightarrow  \\ \quad \quad (\operatorname{opr}(s)(a, \operatorname{iden}(s)) = a) \end{array}}}

\vnbfile{theorem-library/rake-zz-act.scm}{theorem-library/rake-zz-act.scm}
\noindent{\small rake-zz-act.scm -- BATCH 6 of the rake: ZZ-ACT made a DEFINITION, and the five billed ZZ-ACT laws proven from it, each `modulo 0'.}\par
\vnbentry{\texttt{ag-\allowbreak{}as-\allowbreak{}monoid-\allowbreak{}carr}}{\vnbstmt{\begin{array}{@{}l@{}} \forall g. \\ \quad \operatorname{is\text{-}abelian\text{-}group}(g) \Rightarrow  \\ \quad \quad (\operatorname{carr}(g) \simeq \operatorname{carr}(\operatorname{abelian\text{-}group\text{-}as\text{-}monoid}(g))) \end{array}}}
\vnbentry{\texttt{ag-\allowbreak{}as-\allowbreak{}monoid-\allowbreak{}iden}}{\vnbstmt{\begin{array}{@{}l@{}} \forall g. \\ \quad \operatorname{is\text{-}abelian\text{-}group}(g) \Rightarrow  \\ \quad \quad (\operatorname{iden}(g) \simeq \operatorname{iden}(\operatorname{abelian\text{-}group\text{-}as\text{-}monoid}(g))) \end{array}}}
\vnbentry{\texttt{ag-\allowbreak{}as-\allowbreak{}monoid-\allowbreak{}opr}}{\vnbstmt{\begin{array}{@{}l@{}} \forall g. \\ \quad \operatorname{is\text{-}abelian\text{-}group}(g) \Rightarrow  \\ \quad \quad (\operatorname{opr}(g) \simeq \operatorname{opr}(\operatorname{abelian\text{-}group\text{-}as\text{-}monoid}(g))) \end{array}}}
\vnbentry{\texttt{abelian-\allowbreak{}group-\allowbreak{}right-\allowbreak{}inv}}{\vnbstmt{\begin{array}{@{}l@{}} \forall s.; \forall a \in \operatorname{carr}(s). \\ \quad \operatorname{is\text{-}abelian\text{-}group}(s) \Rightarrow  \\ \quad \quad (\operatorname{opr}(s)(a, \operatorname{inv}(s)(a)) = \operatorname{iden}(s)) \end{array}}}
\vnbentry{\texttt{abelian-\allowbreak{}group-\allowbreak{}inv-\allowbreak{}iden}}{\vnbstmt{\begin{array}{@{}l@{}} \forall g. \\ \quad \operatorname{is\text{-}abelian\text{-}group}(g) \Rightarrow  \\ \quad \quad (\operatorname{inv}(g)(\operatorname{iden}(g)) = \operatorname{iden}(g)) \end{array}}}
\vnbentry{\texttt{abelian-\allowbreak{}group-\allowbreak{}inv-\allowbreak{}inv}}{\vnbstmt{\begin{array}{@{}l@{}} \forall s.; \forall a \in \operatorname{carr}(s). \\ \quad \operatorname{is\text{-}abelian\text{-}group}(s) \Rightarrow  \\ \quad \quad (\operatorname{inv}(s)(\operatorname{inv}(s)(a)) = a) \end{array}}}
\vnbentry{\texttt{abelian-\allowbreak{}group-\allowbreak{}inv-\allowbreak{}opr}}{\vnbstmt{\begin{array}{@{}l@{}} \forall s.; \forall a, b \in \operatorname{carr}(s). \\ \quad \operatorname{is\text{-}abelian\text{-}group}(s) \Rightarrow  \\ \quad \quad \operatorname{inv}(s)(\operatorname{opr}(s)(a, b)) \\ \quad \quad =\; \operatorname{opr}(s)(\operatorname{inv}(s)(a), \operatorname{inv}(s)(b)) \end{array}}}
\vnbentry{\texttt{mpow-\allowbreak{}one}}{\vnbstmt{\begin{array}{@{}l@{}} \forall m.; \forall x \in \operatorname{carr}(m). \\ \quad \operatorname{is\text{-}monoid}(m) \Rightarrow  \\ \quad \quad (\operatorname{mpow}(m, x, 1) = x) \end{array}}}
\vnbentry{\texttt{mpow-\allowbreak{}add-\allowbreak{}ind}}{\vnbstmt{\begin{array}{@{}l@{}} \forall j \in \mathbb{N}, m.; \forall x \in \operatorname{carr}(m), k \in \mathbb{N}. \\ \quad \operatorname{is\text{-}monoid}(m) \Rightarrow  \\ \quad \quad \operatorname{mpow}(m, x, (j + k)) \\ \quad \quad =\; \operatorname{opr}(m)(\operatorname{mpow}(m, x, j), \operatorname{mpow}(m, x, k)) \end{array}}}
\vnbentry{\texttt{mpow-\allowbreak{}add}}{\vnbstmt{\begin{array}{@{}l@{}} \forall m.; \forall x \in \operatorname{carr}(m), j, k \in \mathbb{N}. \\ \quad \operatorname{is\text{-}monoid}(m) \Rightarrow  \\ \quad \quad \operatorname{mpow}(m, x, (j + k)) \\ \quad \quad =\; \operatorname{opr}(m)(\operatorname{mpow}(m, x, j), \operatorname{mpow}(m, x, k)) \end{array}}}
\vnbentry{\texttt{ag-\allowbreak{}mpow-\allowbreak{}diff-\allowbreak{}nonneg}}{\vnbstmt{\begin{array}{@{}l@{}} \forall g.; \forall d, n, m \in \mathbb{N}, a \in \operatorname{carr}(g). \\ \quad (\operatorname{is\text{-}abelian\text{-}group}(g) \wedge ((d + n) = m)) \Rightarrow  \\ \quad \quad \operatorname{opr}(g)(\operatorname{mpow}(\operatorname{abelian\text{-}group\text{-}as\text{-}monoid}(g), a, m), \\ \quad \quad \quad \operatorname{inv}(g)(\operatorname{mpow}(\operatorname{abelian\text{-}group\text{-}as\text{-}monoid}(g), a, n))) \\ \quad \quad =\; \operatorname{mpow}(\operatorname{abelian\text{-}group\text{-}as\text{-}monoid}(g), a, d) \end{array}}}
\vnbentry{\texttt{ag-\allowbreak{}mpow-\allowbreak{}diff-\allowbreak{}neg}}{\vnbstmt{\begin{array}{@{}l@{}} \forall g.; \forall p, m, n \in \mathbb{N}, a \in \operatorname{carr}(g). \\ \quad (\operatorname{is\text{-}abelian\text{-}group}(g) \wedge ((p + m) = n)) \Rightarrow  \\ \quad \quad \operatorname{opr}(g)(\operatorname{mpow}(\operatorname{abelian\text{-}group\text{-}as\text{-}monoid}(g), a, m), \\ \quad \quad \quad \operatorname{inv}(g)(\operatorname{mpow}(\operatorname{abelian\text{-}group\text{-}as\text{-}monoid}(g), a, n))) \\ \quad \quad =\; \operatorname{inv}(g)(\operatorname{mpow}(\operatorname{abelian\text{-}group\text{-}as\text{-}monoid}(g), a, p)) \end{array}}}
\vnbentry{\texttt{zz-\allowbreak{}act-\allowbreak{}nonneg}}{\vnbstmt{\begin{array}{@{}l@{}} \forall g.; \forall k \in \mathbb{N}, a \in \operatorname{carr}(g). \\ \quad \operatorname{is\text{-}abelian\text{-}group}(g) \Rightarrow  \\ \quad \quad \operatorname{zz\text{-}act}(g, k, a) \\ \quad \quad =\; \operatorname{mpow}(\operatorname{abelian\text{-}group\text{-}as\text{-}monoid}(g), a, k) \end{array}}}
\vnbentry{\texttt{zz-\allowbreak{}act-\allowbreak{}neg}}{\vnbstmt{\begin{array}{@{}l@{}} \forall g.; \forall k \in \mathbb{N}, a \in \operatorname{carr}(g). \\ \quad \operatorname{is\text{-}abelian\text{-}group}(g) \Rightarrow  \\ \quad \quad \operatorname{zz\text{-}act}(g, -k, a) \\ \quad \quad =\; \operatorname{inv}(g)(\operatorname{mpow}(\operatorname{abelian\text{-}group\text{-}as\text{-}monoid}(g), a, k)) \end{array}}}
\vnbentry{\texttt{zz-\allowbreak{}act-\allowbreak{}zero}}{\vnbstmt{\begin{array}{@{}l@{}} \forall g.; \forall a \in \operatorname{carr}(g). \\ \quad \operatorname{is\text{-}abelian\text{-}group}(g) \Rightarrow  \\ \quad \quad (\operatorname{zz\text{-}act}(g, 0, a) = \operatorname{iden}(g)) \end{array}}}
\vnbentry{\texttt{zz-\allowbreak{}act-\allowbreak{}one}}{\vnbstmt{\begin{array}{@{}l@{}} \forall g.; \forall a \in \operatorname{carr}(g). \\ \quad \operatorname{is\text{-}abelian\text{-}group}(g) \Rightarrow  \\ \quad \quad (\operatorname{zz\text{-}act}(g, 1, a) = a) \end{array}}}
\vnbentry{\texttt{zz-\allowbreak{}act-\allowbreak{}type}}{\vnbstmt{\begin{array}{@{}l@{}} \forall g.; \forall k \in \mathbb{Z}, a \in \operatorname{carr}(g). \\ \quad \operatorname{is\text{-}abelian\text{-}group}(g) \Rightarrow  \\ \quad \quad (\operatorname{zz\text{-}act}(g, k, a) \in \operatorname{carr}(g)) \end{array}}}
\vnbentry{\texttt{zz-\allowbreak{}act-\allowbreak{}neg-\allowbreak{}sign}}{\vnbstmt{\begin{array}{@{}l@{}} \forall g.; \forall k \in \mathbb{Z}, a \in \operatorname{carr}(g). \\ \quad \operatorname{is\text{-}abelian\text{-}group}(g) \Rightarrow  \\ \quad \quad \operatorname{zz\text{-}act}(g, -k, a) \\ \quad \quad =\; \operatorname{inv}(g)(\operatorname{zz\text{-}act}(g, k, a)) \end{array}}}
\vnbentry{\texttt{zz-\allowbreak{}act-\allowbreak{}add}}{\vnbstmt{\begin{array}{@{}l@{}} \forall g.; \forall j, k \in \mathbb{Z}, a \in \operatorname{carr}(g). \\ \quad \operatorname{is\text{-}abelian\text{-}group}(g) \Rightarrow  \\ \quad \quad \operatorname{zz\text{-}act}(g, (j + k), a) \\ \quad \quad =\; \operatorname{opr}(g)(\operatorname{zz\text{-}act}(g, j, a), \operatorname{zz\text{-}act}(g, k, a)) \end{array}}}

\vnbfile{theorem-library/rake-ringoid-additive.scm}{theorem-library/rake-ringoid-additive.scm}
\noindent{\small rake-ringoid-additive.scm -- the additive group of a RINGOID, proven.}\par
\vnbentry{\texttt{r6b-\allowbreak{}radd-\allowbreak{}type}}{\vnbstmt{\begin{array}{@{}l@{}} \forall r, a, b. \\ \quad \operatorname{is\text{-}ringoid}(r) \wedge \\ \quad (a \in \operatorname{carr}(r)) \wedge \\ \quad (b \in \operatorname{carr}(r)) \Rightarrow \\ \quad \quad (\operatorname{add}(r)(a, b) \in \operatorname{carr}(r)) \end{array}}}
\vnbentry{\texttt{r6b-\allowbreak{}rneg-\allowbreak{}type}}{\vnbstmt{\begin{array}{@{}l@{}} \forall r.; \forall a \in \operatorname{carr}(r). \\ \quad \operatorname{is\text{-}ringoid}(r) \Rightarrow  \\ \quad \quad (\operatorname{neg}(r)(a) \in \operatorname{carr}(r)) \end{array}}}
\vnbentry{\texttt{r6b-\allowbreak{}radd-\allowbreak{}right-\allowbreak{}id}}{\vnbstmt{\begin{array}{@{}l@{}} \forall r.; \forall a \in \operatorname{carr}(r). \\ \quad \operatorname{is\text{-}ringoid}(r) \Rightarrow  \\ \quad \quad (\operatorname{add}(r)(a, \operatorname{zero}(r)) = a) \end{array}}}
\vnbentry{\texttt{r6b-\allowbreak{}radd-\allowbreak{}interchange}}{\vnbstmt{\begin{array}{@{}l@{}} \forall r, a, b, c, d. \\ \quad \operatorname{is\text{-}ringoid}(r) \wedge \\ \quad (a \in \operatorname{carr}(r)) \wedge \\ \quad (b \in \operatorname{carr}(r)) \wedge \\ \quad (c \in \operatorname{carr}(r)) \wedge \\ \quad (d \in \operatorname{carr}(r)) \Rightarrow \\ \quad \quad \operatorname{add}(r)(\operatorname{add}(r)(a, b), \operatorname{add}(r)(c, d)) \\ \quad \quad =\; \operatorname{add}(r)(\operatorname{add}(r)(a, c), \operatorname{add}(r)(b, d)) \end{array}}}
\vnbentry{\texttt{r6b-\allowbreak{}rinv-\allowbreak{}unique}}{\vnbstmt{\begin{array}{@{}l@{}} \forall r, a, b. \\ \quad \operatorname{is\text{-}ringoid}(r) \wedge \\ \quad (a \in \operatorname{carr}(r)) \wedge \\ \quad (b \in \operatorname{carr}(r)) \wedge \\ \quad (\operatorname{add}(r)(a, b) = \operatorname{zero}(r)) \Rightarrow \\ \quad \quad (b = \operatorname{neg}(r)(a)) \end{array}}}
\vnbentry{\texttt{r6b-\allowbreak{}rneg-\allowbreak{}add}}{\vnbstmt{\begin{array}{@{}l@{}} \forall r, a, b. \\ \quad \operatorname{is\text{-}ringoid}(r) \wedge \\ \quad (a \in \operatorname{carr}(r)) \wedge \\ \quad (b \in \operatorname{carr}(r)) \Rightarrow \\ \quad \quad \operatorname{neg}(r)(\operatorname{add}(r)(a, b)) \\ \quad \quad =\; \operatorname{add}(r)(\operatorname{neg}(r)(a), \operatorname{neg}(r)(b)) \end{array}}}
\vnbentry{\texttt{r6b-\allowbreak{}rneg-\allowbreak{}neg}}{\vnbstmt{\begin{array}{@{}l@{}} \forall r.; \forall a \in \operatorname{carr}(r). \\ \quad \operatorname{is\text{-}ringoid}(r) \Rightarrow  \\ \quad \quad (\operatorname{neg}(r)(\operatorname{neg}(r)(a)) = a) \end{array}}}
\vnbentry{\texttt{r6b-\allowbreak{}rdiff-\allowbreak{}add}}{\vnbstmt{\begin{array}{@{}l@{}} \forall r, a, b, u, v. \\ \quad \operatorname{is\text{-}ringoid}(r) \wedge \\ \quad (a \in \operatorname{carr}(r)) \wedge \\ \quad (b \in \operatorname{carr}(r)) \wedge \\ \quad (u \in \operatorname{carr}(r)) \wedge \\ \quad (v \in \operatorname{carr}(r)) \Rightarrow \\ \quad \quad \operatorname{add}(r)(\operatorname{add}(r)(a, b), \operatorname{neg}(r)(\operatorname{add}(r)(u, v))) \\ \quad \quad =\; \operatorname{add}(r)(\operatorname{add}(r)(a, \operatorname{neg}(r)(u)), \operatorname{add}(r)(b, \operatorname{neg}(r)(v))) \end{array}}}
\vnbentry{\texttt{ringoid-\allowbreak{}neg-\allowbreak{}diff}}{\vnbstmt{\begin{array}{@{}l@{}} \forall r, a, b. \\ \quad \operatorname{is\text{-}ringoid}(r) \wedge \\ \quad (a \in \operatorname{carr}(r)) \wedge \\ \quad (b \in \operatorname{carr}(r)) \Rightarrow \\ \quad \quad \operatorname{neg}(r)(\operatorname{add}(r)(a, \operatorname{neg}(r)(b))) \\ \quad \quad =\; \operatorname{add}(r)(b, \operatorname{neg}(r)(a)) \end{array}}}
\vnbentry{\texttt{ringoid-\allowbreak{}diff-\allowbreak{}telescope}}{\vnbstmt{\begin{array}{@{}l@{}} \forall r, a, b, c. \\ \quad \operatorname{is\text{-}ringoid}(r) \wedge \\ \quad (a \in \operatorname{carr}(r)) \wedge \\ \quad (b \in \operatorname{carr}(r)) \wedge \\ \quad (c \in \operatorname{carr}(r)) \Rightarrow \\ \quad \quad \operatorname{add}(r)(\operatorname{add}(r)(a, \operatorname{neg}(r)(b)), \operatorname{add}(r)(b, \operatorname{neg}(r)(c))) \\ \quad \quad =\; \operatorname{add}(r)(a, \operatorname{neg}(r)(c)) \end{array}}}

\vnbfile{theorem-library/rake-finsum-typing.scm}{theorem-library/rake-finsum-typing.scm}
\noindent{\small rake-finsum-typing.scm -- the finsum / monoid-power typing leaves of the rake (batch G), PROVEN, together with the guarded form of the one statement in the batch that is NOT provable as stated.}\par
\vnbentry{\texttt{comm-\allowbreak{}monoid-\allowbreak{}carrier-\allowbreak{}closed-\allowbreak{}opr}}{\vnbstmt{\begin{array}{@{}l@{}} \forall s.; \forall a, b \in \operatorname{carr}(s). \\ \quad \operatorname{is\text{-}comm\text{-}monoid}(s) \Rightarrow  \\ \quad \quad (\operatorname{opr}(s)(a, b) \in \operatorname{carr}(s)) \end{array}}}
\vnbentry{\texttt{comm-\allowbreak{}monoid-\allowbreak{}identity-\allowbreak{}in-\allowbreak{}carr}}{\vnbstmt{\begin{array}{@{}l@{}} \forall s. \\ \quad \operatorname{is\text{-}comm\text{-}monoid}(s) \Rightarrow  \\ \quad \quad (\operatorname{iden}(s) \in \operatorname{carr}(s)) \end{array}}}
\vnbentry{\texttt{sum-\allowbreak{}ag-\allowbreak{}comm-\allowbreak{}monoid-\allowbreak{}type-\allowbreak{}ind}}{\vnbstmt{\begin{array}{@{}l@{}} \forall n \in \mathbb{N}, m.; \forall f \in (\mathbb{N} \to \operatorname{carr}(m)). \\ \quad \operatorname{is\text{-}comm\text{-}monoid}(m) \Rightarrow  \\ \quad \quad (\operatorname{sum\text{-}ag}(m, f, n) \in \operatorname{carr}(m)) \end{array}}}
\vnbentry{\texttt{enum-\allowbreak{}fam-\allowbreak{}comm-\allowbreak{}monoid-\allowbreak{}in-\allowbreak{}fun}}{\vnbstmt{\begin{array}{@{}l@{}} \forall n \in \mathbb{N}, m, s.; \forall i \in (\mathrm{OS}(n) \to s), f \in (s \to \operatorname{carr}(m)). \\ \quad \operatorname{is\text{-}comm\text{-}monoid}(m) \Rightarrow  \\ \quad \quad (\operatorname{enum\text{-}fam}(m, f, i, n) \in (\mathbb{N} \to \operatorname{carr}(m))) \end{array}}}
\vnbentry{\texttt{finsum-\allowbreak{}comm-\allowbreak{}monoid-\allowbreak{}type}}{\vnbstmt{\begin{array}{@{}l@{}} \forall m.; \forall s \in \mathbf{Set}, f \in (s \to \operatorname{carr}(m)). \\ \quad (\operatorname{is\text{-}comm\text{-}monoid}(m) \wedge (|s| \in \mathbb{N})) \Rightarrow  \\ \quad \quad (\operatorname{finsum}(m, f, s) \in \operatorname{carr}(m)) \end{array}}}
\vnbentry{\texttt{crmcm-\allowbreak{}carr}}{\vnbstmt{\begin{array}{@{}l@{}} \forall r. \\ \quad \operatorname{is\text{-}commutative\text{-}ring}(r) \Rightarrow  \\ \quad \quad \operatorname{carr}(\operatorname{commutative\text{-}ring\text{-}multiplicative\text{-}cm}(r)) \\ \quad \quad =\; \operatorname{carr}(r) \end{array}}}
\vnbentry{\texttt{prod-\allowbreak{}ring-\allowbreak{}type}}{\vnbstmt{\begin{array}{@{}l@{}} \forall r.; \forall x \in \mathbf{Set}, f \in (x \to \operatorname{carr}(r)). \\ \quad (\operatorname{is\text{-}commutative\text{-}ring}(r) \wedge (|x| \in \mathbb{N})) \Rightarrow  \\ \quad \quad (\operatorname{prod\text{-}ring}(r, f, x) \in \operatorname{carr}(r)) \end{array}}}
\vnbentry{\texttt{mpow-\allowbreak{}comm-\allowbreak{}monoid-\allowbreak{}type-\allowbreak{}ind}}{\vnbstmt{\begin{array}{@{}l@{}} \forall n \in \mathbb{N}, m.; \forall x \in \operatorname{carr}(m). \\ \quad \operatorname{is\text{-}comm\text{-}monoid}(m) \Rightarrow  \\ \quad \quad (\operatorname{mpow}(m, x, n) \in \operatorname{carr}(m)) \end{array}}}
\vnbentry{\texttt{ring-\allowbreak{}power-\allowbreak{}type}}{\vnbstmt{\begin{array}{@{}l@{}} \forall r.; \forall x \in \operatorname{carr}(r), n \in \mathbb{N}. \\ \quad \operatorname{is\text{-}commutative\text{-}ring}(r) \Rightarrow  \\ \quad \quad (\operatorname{ring\text{-}power}(r, x, n) \in \operatorname{carr}(r)) \end{array}}}
\vnbentry{\texttt{sum-\allowbreak{}ag-\allowbreak{}in-\allowbreak{}subset-\allowbreak{}ind}}{\vnbstmt{\begin{array}{@{}l@{}} \forall n \in \mathbb{N}, g, m, a. \\ \quad \operatorname{is\text{-}abelian\text{-}group}(g) \wedge \\ \quad (\operatorname{iden}(g) \in m) \wedge \\ \quad \forall x, y \in m.\; (\operatorname{opr}(g)(x, y) \in m) \wedge \\ \quad \forall i \in \mathrm{OS}(n).\; (a(i) \in m) \Rightarrow \\ \quad \quad (\operatorname{sum\text{-}ag}(g, a, n) \in m) \end{array}}}
\vnbentry{\texttt{finsum-\allowbreak{}in-\allowbreak{}subset}}{\vnbstmt{\begin{array}{@{}l@{}} \forall g, m.; \forall s \in \mathbf{Set}, f \in (s \to \operatorname{carr}(g)). \\ \quad \operatorname{is\text{-}abelian\text{-}group}(g) \wedge \\ \quad (\operatorname{iden}(g) \in m) \wedge \\ \quad \forall x, y \in m.\; (\operatorname{opr}(g)(x, y) \in m) \wedge \\ \quad (|s| \in \mathbb{N}) \wedge \\ \quad \forall z \in s.\; (f(z) \in m) \Rightarrow \\ \quad \quad (\operatorname{finsum}(g, f, s) \in m) \end{array}}}
\vnbentry{\texttt{submodule-\allowbreak{}finsum-\allowbreak{}closed}}{\vnbstmt{\begin{array}{@{}l@{}} \forall d, m.; \forall s \in \mathbf{Set}, f \in (s \to \operatorname{carr}(\operatorname{module\text{-}vector\text{-}ag}(d))). \\ \quad \operatorname{is\text{-}module}(d) \wedge \\ \quad \operatorname{is\text{-}submodule}(d, m) \wedge \\ \quad (|s| \in \mathbb{N}) \wedge \\ \quad \forall z \in s.\; (f(z) \in m) \Rightarrow \\ \quad \quad (\operatorname{finsum}(\operatorname{module\text{-}vector\text{-}ag}(d), f, s) \in m) \end{array}}}
\vnbentry{\texttt{funcomp-\allowbreak{}succ-\allowbreak{}type}}{\vnbstmt{\begin{array}{@{}l@{}} \forall x.; \forall q \in \mathbb{N}, f \in (\operatorname{interval}(1, (q + 1)) \to x). \\ \quad (\lambda z \in \operatorname{interval}(1, q).\; f((z + 1))) \\ \quad \in\; (\operatorname{interval}(1, q) \to x) \end{array}}}

\vnbfile{theorem-library/rake-finsum-welldef.scm}{theorem-library/rake-finsum-welldef.scm}
\noindent{\small theorem-library/rake-finsum-welldef.scm -- rake batch N (2026-09-17), part 2 of 2 --------------------------------------------------------------------------- NOTE, 2026-09-20 (batch 9-B).  This file was written while cardinality was AXIOMATISED under the name CARD and the defined constant was its companion}\par
\vnbentry{\texttt{opr-\allowbreak{}rearrange-\allowbreak{}laws}}{\vnbstmt{\begin{array}{@{}l@{}} \forall s.; \forall a, b, c \in \operatorname{carr}(s). \\ \quad (\operatorname{iden}(s) \in \operatorname{carr}(s)) \wedge \\ \quad \forall \alpha, \beta \in \operatorname{carr}(s). \\ \quad \quad (\operatorname{opr}(s)(\alpha, \beta) \in \operatorname{carr}(s)) \wedge \\ \quad \forall \alpha, \beta, \gamma \in \operatorname{carr}(s). \\ \quad \quad \operatorname{opr}(s)(\operatorname{opr}(s)(\alpha, \beta), \gamma) \\ \quad \quad =\; \operatorname{opr}(s)(\alpha, \operatorname{opr}(s)(\beta, \gamma)) \wedge \\ \quad \forall \alpha, \beta \in \operatorname{carr}(s). \\ \quad \quad (\operatorname{opr}(s)(\alpha, \beta) = \operatorname{opr}(s)(\beta, \alpha)) \Rightarrow \\ \quad \quad \operatorname{opr}(s)(\operatorname{opr}(s)(a, b), c) \\ \quad \quad =\; \operatorname{opr}(s)(\operatorname{opr}(s)(a, c), b) \end{array}}}
\vnbentry{\texttt{sum-\allowbreak{}ag-\allowbreak{}type-\allowbreak{}laws}}{\vnbstmt{\begin{array}{@{}l@{}} \forall n \in \mathbb{N}, s.; \forall g \in (\mathbb{N} \to \operatorname{carr}(s)). \\ \quad (\operatorname{iden}(s) \in \operatorname{carr}(s)) \wedge \\ \quad \forall a, b \in \operatorname{carr}(s). \\ \quad \quad (\operatorname{opr}(s)(a, b) \in \operatorname{carr}(s)) \wedge \\ \quad \forall a, b, c \in \operatorname{carr}(s). \\ \quad \quad \operatorname{opr}(s)(\operatorname{opr}(s)(a, b), c) \\ \quad \quad =\; \operatorname{opr}(s)(a, \operatorname{opr}(s)(b, c)) \wedge \\ \quad \forall a, b \in \operatorname{carr}(s). \\ \quad \quad (\operatorname{opr}(s)(a, b) = \operatorname{opr}(s)(b, a)) \Rightarrow \\ \quad \quad (\operatorname{sum\text{-}ag}(s, g, n) \in \operatorname{carr}(s)) \end{array}}}
\vnbentry{\texttt{sum-\allowbreak{}ag-\allowbreak{}segment-\allowbreak{}congruence}}{\vnbstmt{\begin{array}{@{}l@{}} \forall n \in \mathbb{N}, a, g, h. \\ \quad \forall i \in \mathrm{OS}(n).\; (g(i) \simeq h(i)) \Rightarrow  \\ \quad \quad (\operatorname{sum\text{-}ag}(a, g, n) \simeq \operatorname{sum\text{-}ag}(a, h, n)) \end{array}}}
\vnbentry{\texttt{sum-\allowbreak{}ag-\allowbreak{}replace-\allowbreak{}entry-\allowbreak{}laws}}{\vnbstmt{\begin{array}{@{}l@{}} \forall n \in \mathbb{N}, s.; \forall g, p \in (\mathbb{N} \to \operatorname{carr}(s)), k \in \mathrm{OS}(n). \\ \quad (\operatorname{iden}(s) \in \operatorname{carr}(s)) \wedge \\ \quad \forall a, b \in \operatorname{carr}(s). \\ \quad \quad (\operatorname{opr}(s)(a, b) \in \operatorname{carr}(s)) \wedge \\ \quad \forall a, b, c \in \operatorname{carr}(s). \\ \quad \quad \operatorname{opr}(s)(\operatorname{opr}(s)(a, b), c) \\ \quad \quad =\; \operatorname{opr}(s)(a, \operatorname{opr}(s)(b, c)) \wedge \\ \quad \forall a, b \in \operatorname{carr}(s). \\ \quad \quad (\operatorname{opr}(s)(a, b) = \operatorname{opr}(s)(b, a)) \wedge \\ \quad \forall i \in \mathrm{OS}(n).\; (\neg (i = k) \Rightarrow (g(i) \simeq p(i))) \Rightarrow \\ \quad \quad \operatorname{opr}(s)(\operatorname{sum\text{-}ag}(s, g, n), p(k)) \\ \quad \quad =\; \operatorname{opr}(s)(\operatorname{sum\text{-}ag}(s, p, n), g(k)) \end{array}}}
\vnbentry{\texttt{seg-\allowbreak{}perm-\allowbreak{}restrict-\allowbreak{}is-\allowbreak{}bijection}}{\vnbstmt{\begin{array}{@{}l@{}} \forall n \in \mathbb{N}, i \in \mathrm{Bij}(\mathrm{OS}((n + 1)), \mathrm{OS}((n + 1))), k. \\ \quad (i(n) = k) \Rightarrow  \\ \quad \quad (\lambda a \in \mathrm{OS}(n).\; \left(\text{if }(i(a) = n)\text{ then }k\text{ else }i(a)\right)) \\ \quad \quad \in\; \mathrm{Bij}(\mathrm{OS}(n), \mathrm{OS}(n)) \end{array}}}
\vnbentry{\texttt{sum-\allowbreak{}ag-\allowbreak{}permutation-\allowbreak{}invariance-\allowbreak{}laws}}{\vnbstmt{\begin{array}{@{}l@{}} \forall n \in \mathbb{N}, s.; \forall g, h \in (\mathbb{N} \to \operatorname{carr}(s)), a \in \mathrm{Bij}(\mathrm{OS}(n), \mathrm{OS}(n)). \\ \quad (\operatorname{iden}(s) \in \operatorname{carr}(s)) \wedge \\ \quad \forall \alpha, b \in \operatorname{carr}(s). \\ \quad \quad (\operatorname{opr}(s)(\alpha, b) \in \operatorname{carr}(s)) \wedge \\ \quad \forall \alpha, b, c \in \operatorname{carr}(s). \\ \quad \quad \operatorname{opr}(s)(\operatorname{opr}(s)(\alpha, b), c) \\ \quad \quad =\; \operatorname{opr}(s)(\alpha, \operatorname{opr}(s)(b, c)) \wedge \\ \quad \forall \alpha, b \in \operatorname{carr}(s). \\ \quad \quad (\operatorname{opr}(s)(\alpha, b) = \operatorname{opr}(s)(b, \alpha)) \wedge \\ \quad \forall i \in \mathrm{OS}(n).\; (h(i) = g(a(i))) \Rightarrow \\ \quad \quad (\operatorname{sum\text{-}ag}(s, g, n) = \operatorname{sum\text{-}ag}(s, h, n)) \end{array}}}
\vnbentry{\texttt{sum-\allowbreak{}ag-\allowbreak{}permutation-\allowbreak{}invariance}}{\vnbstmt{\begin{array}{@{}l@{}} \forall n \in \mathbb{N}, a.; \forall g, h \in (\mathbb{N} \to \operatorname{carr}(a)), b \in \mathrm{Bij}(\mathrm{OS}(n), \mathrm{OS}(n)). \\ \quad \operatorname{is\text{-}abelian\text{-}group}(a) \wedge \\ \quad \forall i \in \mathrm{OS}(n).\; (h(i) = g(b(i))) \Rightarrow \\ \quad \quad (\operatorname{sum\text{-}ag}(a, g, n) = \operatorname{sum\text{-}ag}(a, h, n)) \end{array}}}
\vnbentry{\texttt{finsum-\allowbreak{}comm-\allowbreak{}monoid-\allowbreak{}permutation-\allowbreak{}invariance}}{\vnbstmt{\begin{array}{@{}l@{}} \forall n \in \mathbb{N}, m.; \forall g, h \in (\mathbb{N} \to \operatorname{carr}(m)), a \in \mathrm{Bij}(\mathrm{OS}(n), \mathrm{OS}(n)). \\ \quad \operatorname{is\text{-}comm\text{-}monoid}(m) \wedge \\ \quad \forall i \in \mathrm{OS}(n).\; (h(i) = g(a(i))) \Rightarrow \\ \quad \quad (\operatorname{sum\text{-}ag}(m, g, n) = \operatorname{sum\text{-}ag}(m, h, n)) \end{array}}}
\vnbentry{\texttt{finsum-\allowbreak{}well-\allowbreak{}defined}}{\vnbstmt{\begin{array}{@{}l@{}} \forall s \in \mathbf{Set}, g.; \forall f \in (s \to \operatorname{carr}(g)), m \in \mathrm{Bij}(\mathrm{OS}(|s|), s). \\ \quad ((|s| \in \mathbb{N}) \wedge \operatorname{is\text{-}abelian\text{-}group}(g)) \Rightarrow  \\ \quad \quad \operatorname{finsum}(g, f, s) \\ \quad \quad =\; \operatorname{sum\text{-}ag}(g, \operatorname{enum\text{-}fam}(g, f, m, |s|), |s|) \end{array}}}
\vnbentry{\texttt{finsum-\allowbreak{}comm-\allowbreak{}monoid-\allowbreak{}well-\allowbreak{}defined}}{\vnbstmt{\begin{array}{@{}l@{}} \forall s \in \mathbf{Set}, m.; \forall f \in (s \to \operatorname{carr}(m)), a \in \mathrm{Bij}(\mathrm{OS}(|s|), s). \\ \quad ((|s| \in \mathbb{N}) \wedge \operatorname{is\text{-}comm\text{-}monoid}(m)) \Rightarrow  \\ \quad \quad \operatorname{finsum}(m, f, s) \\ \quad \quad =\; \operatorname{sum\text{-}ag}(m, \operatorname{enum\text{-}fam}(m, f, a, |s|), |s|) \end{array}}}

\vnbfile{theorem-library/finsum-insert.scm}{theorem-library/finsum-insert.scm}
\noindent{\small theorem-library/finsum-insert.scm -- the insertion recurrence for FINSUM. --------------------------------------------------------------------------- NOTE, 2026-09-20 (batch 9-B).  This file was written while cardinality was AXIOMATISED under the name CARD and the defined constant was its companion}\par
\vnbentry{\texttt{finsum-\allowbreak{}insert-\allowbreak{}from-\allowbreak{}enum}}{\vnbstmt{\begin{array}{@{}l@{}} \forall m, f, x, k, a. \\ \quad (|x| \in \mathbb{N}) \wedge \\ \quad (|(x \cup \{k\})| = (|x| + 1)) \wedge \\ \quad (a(|x|) \simeq k) \wedge \\ \quad \forall i \in \mathrm{OS}(|x|). \\ \quad \quad (a(i) \simeq \operatorname{fin\text{-}enum}(x)(i)) \wedge \\ \quad \operatorname{finsum}(m, f, (x \cup \{k\})) \\ \quad =\; \operatorname{sum\text{-}ag}(m, \operatorname{enum\text{-}fam}(m, f, a, |(x \cup \{k\})|), \\ \quad \quad |(x \cup \{k\})|) \Rightarrow \\ \quad \quad \operatorname{finsum}(m, f, (x \cup \{k\})) \\ \quad \quad =\; \operatorname{opr}(m)(\operatorname{finsum}(m, f, x), f(k)) \end{array}}}
\vnbentry{\texttt{finsum-\allowbreak{}insert}}{\vnbstmt{\begin{array}{@{}l@{}} \forall m.; \forall x, k \in \mathbf{Set}, f \in ((x \cup \{k\}) \to \operatorname{carr}(m)). \\ \quad (\operatorname{is\text{-}comm\text{-}monoid}(m) \wedge ((|x| \in \mathbb{N}) \wedge \neg (k \in x))) \Rightarrow  \\ \quad \quad \operatorname{finsum}(m, f, (x \cup \{k\})) \\ \quad \quad =\; \operatorname{opr}(m)(\operatorname{finsum}(m, f, x), f(k)) \end{array}}}
\vnbentry{\texttt{finsum-\allowbreak{}insert-\allowbreak{}ag}}{\vnbstmt{\begin{array}{@{}l@{}} \forall g.; \forall x, k \in \mathbf{Set}, f \in ((x \cup \{k\}) \to \operatorname{carr}(g)). \\ \quad \operatorname{is\text{-}abelian\text{-}group}(g) \wedge \\ \quad (|x| \in \mathbb{N}) \wedge \\ \quad \neg (k \in x) \Rightarrow \\ \quad \quad \operatorname{finsum}(g, f, (x \cup \{k\})) \\ \quad \quad =\; \operatorname{opr}(g)(\operatorname{finsum}(g, f, x), f(k)) \end{array}}}
\vnbentry{\texttt{finsum-\allowbreak{}empty}}{\vnbstmt{\begin{array}{@{}l@{}} \forall g, f. \\ \quad (\operatorname{finsum}(g, f, \emptyset) \simeq \operatorname{iden}(g)) \end{array}}}

\vnbfile{theorem-library/rake-extended-order.scm}{theorem-library/rake-extended-order.scm}
\noindent{\small theorem-library/rake-extended-order.scm -- the order and addition theory of RR-POS-STAR = [0, +inf], built on the user's new axiom `pos-inf-above-reals' (structure-library/extended-reals.scm, stamped `primitive', 2026-09-18).}\par
\vnbentry{\texttt{rr-\allowbreak{}pos-\allowbreak{}star-\allowbreak{}le-\allowbreak{}refl}}{\vnbstmt{\forall x \in \operatorname{rr\text{-}pos\text{-}star}.\; (x \leq x)}}
\vnbentry{\texttt{rr-\allowbreak{}pos-\allowbreak{}star-\allowbreak{}le-\allowbreak{}real-\allowbreak{}in-\allowbreak{}rr}}{\vnbstmt{\begin{array}{@{}l@{}} \forall x, y. \\ \quad (((x \in \operatorname{rr\text{-}pos\text{-}star}) \wedge (y \in \mathbb{R})) \wedge (x \leq y)) \Rightarrow  \\ \quad \quad (x \in \mathbb{R}) \end{array}}}
\vnbentry{\texttt{rr-\allowbreak{}pos-\allowbreak{}star-\allowbreak{}le-\allowbreak{}trans}}{\vnbstmt{\begin{array}{@{}l@{}} \forall x, y, z. \\ \quad (x \in \operatorname{rr\text{-}pos\text{-}star}) \wedge \\ \quad (y \in \operatorname{rr\text{-}pos\text{-}star}) \wedge \\ \quad (z \in \operatorname{rr\text{-}pos\text{-}star}) \wedge \\ \quad (x \leq y) \wedge \\ \quad (y \leq z) \Rightarrow \\ \quad \quad (x \leq z) \end{array}}}
\vnbentry{\texttt{rr-\allowbreak{}pos-\allowbreak{}star-\allowbreak{}le-\allowbreak{}antisymm}}{\vnbstmt{\begin{array}{@{}l@{}} \forall x, y. \\ \quad (x \in \operatorname{rr\text{-}pos\text{-}star}) \wedge \\ \quad (y \in \operatorname{rr\text{-}pos\text{-}star}) \wedge \\ \quad (x \leq y) \wedge \\ \quad (y \leq x) \Rightarrow \\ \quad \quad (x = y) \end{array}}}
\vnbentry{\texttt{eplus-\allowbreak{}closed}}{\vnbstmt{\begin{array}{@{}l@{}} \forall x, y. \\ \quad (x \in \operatorname{rr\text{-}pos\text{-}star}) \wedge \\ \quad (y \in \operatorname{rr\text{-}pos\text{-}star}) \Rightarrow \\ \quad \quad (\operatorname{eplus}(x, y) \in \operatorname{rr\text{-}pos\text{-}star}) \end{array}}}
\vnbentry{\texttt{eplus-\allowbreak{}zero-\allowbreak{}left}}{\vnbstmt{\begin{array}{@{}l@{}} \forall x \in \operatorname{rr\text{-}pos\text{-}star}. \\ \quad (\operatorname{eplus}(0, x) = x) \end{array}}}
\vnbentry{\texttt{eplus-\allowbreak{}zero-\allowbreak{}right}}{\vnbstmt{\begin{array}{@{}l@{}} \forall x \in \operatorname{rr\text{-}pos\text{-}star}. \\ \quad (\operatorname{eplus}(x, 0) = x) \end{array}}}
\vnbentry{\texttt{eplus-\allowbreak{}comm}}{\vnbstmt{\begin{array}{@{}l@{}} \forall x, y. \\ \quad (x \in \operatorname{rr\text{-}pos\text{-}star}) \wedge \\ \quad (y \in \operatorname{rr\text{-}pos\text{-}star}) \Rightarrow \\ \quad \quad (\operatorname{eplus}(x, y) = \operatorname{eplus}(y, x)) \end{array}}}
\vnbentry{\texttt{eplus-\allowbreak{}assoc}}{\vnbstmt{\begin{array}{@{}l@{}} \forall x, y, z. \\ \quad (x \in \operatorname{rr\text{-}pos\text{-}star}) \wedge \\ \quad (y \in \operatorname{rr\text{-}pos\text{-}star}) \wedge \\ \quad (z \in \operatorname{rr\text{-}pos\text{-}star}) \Rightarrow \\ \quad \quad \operatorname{eplus}(\operatorname{eplus}(x, y), z) \\ \quad \quad =\; \operatorname{eplus}(x, \operatorname{eplus}(y, z)) \end{array}}}
\vnbentry{\texttt{eplus-\allowbreak{}le-\allowbreak{}left}}{\vnbstmt{\begin{array}{@{}l@{}} \forall x, y. \\ \quad (x \in \operatorname{rr\text{-}pos\text{-}star}) \wedge \\ \quad (y \in \operatorname{rr\text{-}pos\text{-}star}) \Rightarrow \\ \quad \quad (x \leq \operatorname{eplus}(x, y)) \end{array}}}
\vnbentry{\texttt{eplus-\allowbreak{}mono}}{\vnbstmt{\begin{array}{@{}l@{}} \forall x, y, u, v. \\ \quad (x \in \operatorname{rr\text{-}pos\text{-}star}) \wedge \\ \quad (y \in \operatorname{rr\text{-}pos\text{-}star}) \wedge \\ \quad (u \in \operatorname{rr\text{-}pos\text{-}star}) \wedge \\ \quad (v \in \operatorname{rr\text{-}pos\text{-}star}) \wedge \\ \quad (x \leq u) \wedge \\ \quad (y \leq v) \Rightarrow \\ \quad \quad (\operatorname{eplus}(x, y) \leq \operatorname{eplus}(u, v)) \end{array}}}

\vnbfile{theorem-library/rake-esup-defined.scm}{theorem-library/rake-esup-defined.scm}
\noindent{\small theorem-library/rake-esup-defined.scm -- ESUP DEFINED as a definite description, and its four characterising axioms PROVEN from the definition. Assignment "ESUP and ESUM defined" (the user's decision of 2026-09-19). Helper prefix `esd-'.}\par
\vnbentry{\texttt{esup-\allowbreak{}exists}}{\vnbstmt{\begin{array}{@{}l@{}} \forall s. \\ \quad (s \subseteq \operatorname{rr\text{-}pos\text{-}star}) \Rightarrow  \\ \quad \quad \exists \mathit{bb} \in \operatorname{rr\text{-}pos\text{-}star}. \\ \quad \quad \quad \forall x \in s.\; (x \leq \mathit{bb}) \wedge \\ \quad \quad \quad \forall b \in \operatorname{rr\text{-}pos\text{-}star}. \\ \quad \quad \quad \quad (\forall x \in s.\; (x \leq b) \Rightarrow (\mathit{bb} \leq b)) \end{array}}}
\vnbentry{\texttt{esup-\allowbreak{}unique}}{\vnbstmt{\begin{array}{@{}l@{}} \forall s, b_{1}, b_{2}. \\ \quad (b_{1} \in \operatorname{rr\text{-}pos\text{-}star}) \wedge \\ \quad \forall x \in s.\; (x \leq b_{1}) \wedge \\ \quad \forall b \in \operatorname{rr\text{-}pos\text{-}star}. \\ \quad \quad (\forall x \in s.\; (x \leq b) \Rightarrow (b_{1} \leq b)) \wedge \\ \quad (b_{2} \in \operatorname{rr\text{-}pos\text{-}star}) \wedge \\ \quad \forall x \in s.\; (x \leq b_{2}) \wedge \\ \quad \forall b \in \operatorname{rr\text{-}pos\text{-}star}. \\ \quad \quad (\forall x \in s.\; (x \leq b) \Rightarrow (b_{2} \leq b)) \Rightarrow \\ \quad \quad (b_{1} = b_{2}) \end{array}}}
\vnbentry{\texttt{esup-\allowbreak{}prop}}{\vnbstmt{\begin{array}{@{}l@{}} \forall s. \\ \quad (s \subseteq \operatorname{rr\text{-}pos\text{-}star}) \Rightarrow  \\ \quad \quad (\operatorname{esup}(s) \in \operatorname{rr\text{-}pos\text{-}star}) \wedge \\ \quad \quad \forall x \in s.\; (x \leq \operatorname{esup}(s)) \wedge \\ \quad \quad \forall b \in \operatorname{rr\text{-}pos\text{-}star}. \\ \quad \quad \quad (\forall x \in s.\; (x \leq b) \Rightarrow (\operatorname{esup}(s) \leq b)) \end{array}}}
\vnbentry{\texttt{esup-\allowbreak{}in}}{\vnbstmt{\begin{array}{@{}l@{}} \forall s. \\ \quad (s \subseteq \operatorname{rr\text{-}pos\text{-}star}) \Rightarrow  \\ \quad \quad (\operatorname{esup}(s) \in \operatorname{rr\text{-}pos\text{-}star}) \end{array}}}
\vnbentry{\texttt{esup-\allowbreak{}upper}}{\vnbstmt{\begin{array}{@{}l@{}} \forall s.; \forall x \in s. \\ \quad (s \subseteq \operatorname{rr\text{-}pos\text{-}star}) \Rightarrow  \\ \quad \quad (x \leq \operatorname{esup}(s)) \end{array}}}
\vnbentry{\texttt{esup-\allowbreak{}least}}{\vnbstmt{\begin{array}{@{}l@{}} \forall s.; \forall b \in \operatorname{rr\text{-}pos\text{-}star}. \\ \quad ((s \subseteq \operatorname{rr\text{-}pos\text{-}star}) \wedge \forall x \in s.\; (x \leq b)) \Rightarrow  \\ \quad \quad (\operatorname{esup}(s) \leq b) \end{array}}}
\vnbentry{\texttt{esup-\allowbreak{}empty}}{\vnbstmt{(\operatorname{esup}(\emptyset) = 0)}}

\vnbfile{theorem-library/rake-rr-pos-star-monoid.scm}{theorem-library/rake-rr-pos-star-monoid.scm}
\noindent{\small theorem-library/rake-rr-pos-star-monoid.scm -- IS-COMM-MONOID(RR-POS-STAR-ADD-MONOID), PROVEN.  Rake batch 6-C.  Helper prefix `rpm-'.}\par
\vnbentry{\texttt{rr-\allowbreak{}pos-\allowbreak{}star-\allowbreak{}is-\allowbreak{}comm-\allowbreak{}monoid}}{\vnbstmt{\operatorname{is\text{-}comm\text{-}monoid}(\operatorname{rr\text{-}pos\text{-}star\text{-}add\text{-}monoid})}}

\vnbfile{theorem-library/field-ring-view.scm}{theorem-library/field-ring-view.scm}
\noindent{\small field-ring-view.scm -- EVERY FIELD IS A FIELD-RING.}\par
\vnbentry{\texttt{singleton-\allowbreak{}unfold}}{\vnbstmt{\begin{array}{@{}l@{}} \forall y. \\ \quad (\operatorname{singleton}(y) \simeq \operatorname{make\text{-}set}(\langle y \rangle)) \end{array}}}
\vnbentry{\texttt{singleton-\allowbreak{}membership}}{\vnbstmt{\begin{array}{@{}l@{}} \forall y, x. \\ \quad ((x \in \operatorname{singleton}(y)) \Leftrightarrow ((x \in \mathbf{Set}) \wedge (x = y))) \end{array}}}
\vnbentry{\texttt{field-\allowbreak{}id-\allowbreak{}carr}}{\vnbstmt{\begin{array}{@{}l@{}} \forall f. \\ \quad (\operatorname{carr}(\operatorname{field\text{-}as\text{-}integral\text{-}domain}(f)) \simeq \operatorname{carr}(f)) \end{array}}}
\vnbentry{\texttt{field-\allowbreak{}id-\allowbreak{}mul}}{\vnbstmt{\begin{array}{@{}l@{}} \forall f. \\ \quad (\operatorname{mul}(\operatorname{field\text{-}as\text{-}integral\text{-}domain}(f)) \simeq \operatorname{mul}(f)) \end{array}}}
\vnbentry{\texttt{field-\allowbreak{}id-\allowbreak{}one}}{\vnbstmt{\begin{array}{@{}l@{}} \forall f. \\ \quad (\operatorname{one}(\operatorname{field\text{-}as\text{-}integral\text{-}domain}(f)) \simeq \operatorname{one}(f)) \end{array}}}
\vnbentry{\texttt{field-\allowbreak{}id-\allowbreak{}zero}}{\vnbstmt{\begin{array}{@{}l@{}} \forall f. \\ \quad (\operatorname{zero}(\operatorname{field\text{-}as\text{-}integral\text{-}domain}(f)) \simeq \operatorname{zero}(f)) \end{array}}}
\vnbentry{\texttt{field-\allowbreak{}id-\allowbreak{}add}}{\vnbstmt{\begin{array}{@{}l@{}} \forall f. \\ \quad (\operatorname{add}(\operatorname{field\text{-}as\text{-}integral\text{-}domain}(f)) \simeq \operatorname{add}(f)) \end{array}}}
\vnbentry{\texttt{field-\allowbreak{}id-\allowbreak{}neg}}{\vnbstmt{\begin{array}{@{}l@{}} \forall f. \\ \quad (\operatorname{neg}(\operatorname{field\text{-}as\text{-}integral\text{-}domain}(f)) \simeq \operatorname{neg}(f)) \end{array}}}
\vnbentry{\texttt{field-\allowbreak{}is-\allowbreak{}field-\allowbreak{}ring}}{\vnbstmt{\begin{array}{@{}l@{}} \forall s. \\ \quad \operatorname{is\text{-}field}(s) \Rightarrow  \\ \quad \quad \operatorname{is\text{-}field\text{-}ring}(\operatorname{field\text{-}as\text{-}integral\text{-}domain}(s)) \end{array}}}

\vnbfile{theorem-library/hom-laws.scm}{theorem-library/hom-laws.scm}
\noindent{\small hom-laws.scm -- the models of a structure form a CATEGORY: identity arrows, composition, and the hom-set as a set.}\par
\vnbentry{\texttt{hom-\allowbreak{}semigroup-\allowbreak{}member-\allowbreak{}iff}}{\vnbstmt{\begin{array}{@{}l@{}} \forall a, b, f. \\ \quad ((f \in \operatorname{hom\text{-}semigroup}(a, b)) \Leftrightarrow ((f \in (\operatorname{carr}(a) \to \operatorname{carr}(b))) \wedge \operatorname{is\text{-}hom\text{-}semigroup}(a, b, f))) \end{array}}}
\vnbentry{\texttt{hom-\allowbreak{}semigroup-\allowbreak{}id}}{\vnbstmt{\begin{array}{@{}l@{}} \forall a. \\ \quad \operatorname{is\text{-}semigroup}(a) \Rightarrow  \\ \quad \quad \operatorname{is\text{-}hom\text{-}semigroup}(a, a, \operatorname{id\text{-}fun}(\operatorname{carr}(a))) \end{array}}}
\vnbentry{\texttt{hom-\allowbreak{}semigroup-\allowbreak{}compose}}{\vnbstmt{\begin{array}{@{}l@{}} \forall a, b, c, f, g. \\ \quad \operatorname{is\text{-}hom\text{-}semigroup}(a, b, f) \wedge \\ \quad \operatorname{is\text{-}hom\text{-}semigroup}(b, c, g) \Rightarrow \\ \quad \quad \operatorname{is\text{-}hom\text{-}semigroup}(a, c, \operatorname{compose}(g, f)) \end{array}}}
\vnbentry{\texttt{hom-\allowbreak{}semigroup-\allowbreak{}in-\allowbreak{}set}}{\vnbstmt{\begin{array}{@{}l@{}} \forall a, b. \\ \quad (\operatorname{is\text{-}semigroup}(a) \wedge \operatorname{is\text{-}semigroup}(b)) \Rightarrow  \\ \quad \quad (\operatorname{hom\text{-}semigroup}(a, b) \in \mathbf{Set}) \end{array}}}
\vnbentry{\texttt{hom-\allowbreak{}monoid-\allowbreak{}member-\allowbreak{}iff}}{\vnbstmt{\begin{array}{@{}l@{}} \forall a, b, f. \\ \quad ((f \in \operatorname{hom\text{-}monoid}(a, b)) \Leftrightarrow ((f \in (\operatorname{carr}(a) \to \operatorname{carr}(b))) \wedge \operatorname{is\text{-}hom\text{-}monoid}(a, b, f))) \end{array}}}
\vnbentry{\texttt{hom-\allowbreak{}monoid-\allowbreak{}id}}{\vnbstmt{\begin{array}{@{}l@{}} \forall a. \\ \quad \operatorname{is\text{-}monoid}(a) \Rightarrow  \\ \quad \quad \operatorname{is\text{-}hom\text{-}monoid}(a, a, \operatorname{id\text{-}fun}(\operatorname{carr}(a))) \end{array}}}
\vnbentry{\texttt{hom-\allowbreak{}monoid-\allowbreak{}compose}}{\vnbstmt{\begin{array}{@{}l@{}} \forall a, b, c, f, g. \\ \quad \operatorname{is\text{-}hom\text{-}monoid}(a, b, f) \wedge \\ \quad \operatorname{is\text{-}hom\text{-}monoid}(b, c, g) \Rightarrow \\ \quad \quad \operatorname{is\text{-}hom\text{-}monoid}(a, c, \operatorname{compose}(g, f)) \end{array}}}
\vnbentry{\texttt{hom-\allowbreak{}monoid-\allowbreak{}in-\allowbreak{}set}}{\vnbstmt{\begin{array}{@{}l@{}} \forall a, b. \\ \quad (\operatorname{is\text{-}monoid}(a) \wedge \operatorname{is\text{-}monoid}(b)) \Rightarrow  \\ \quad \quad (\operatorname{hom\text{-}monoid}(a, b) \in \mathbf{Set}) \end{array}}}
\vnbentry{\texttt{hom-\allowbreak{}comm-\allowbreak{}monoid-\allowbreak{}member-\allowbreak{}iff}}{\vnbstmt{\begin{array}{@{}l@{}} \forall a, b, f. \\ \quad ((f \in \operatorname{hom\text{-}comm\text{-}monoid}(a, b)) \Leftrightarrow ((f \in (\operatorname{carr}(a) \to \operatorname{carr}(b))) \wedge \operatorname{is\text{-}hom\text{-}comm\text{-}monoid}(a, b, f))) \end{array}}}
\vnbentry{\texttt{hom-\allowbreak{}comm-\allowbreak{}monoid-\allowbreak{}id}}{\vnbstmt{\begin{array}{@{}l@{}} \forall a. \\ \quad \operatorname{is\text{-}comm\text{-}monoid}(a) \Rightarrow  \\ \quad \quad \operatorname{is\text{-}hom\text{-}comm\text{-}monoid}(a, a, \operatorname{id\text{-}fun}(\operatorname{carr}(a))) \end{array}}}
\vnbentry{\texttt{hom-\allowbreak{}comm-\allowbreak{}monoid-\allowbreak{}compose}}{\vnbstmt{\begin{array}{@{}l@{}} \forall a, b, c, f, g. \\ \quad \operatorname{is\text{-}hom\text{-}comm\text{-}monoid}(a, b, f) \wedge \\ \quad \operatorname{is\text{-}hom\text{-}comm\text{-}monoid}(b, c, g) \Rightarrow \\ \quad \quad \operatorname{is\text{-}hom\text{-}comm\text{-}monoid}(a, c, \operatorname{compose}(g, f)) \end{array}}}
\vnbentry{\texttt{hom-\allowbreak{}comm-\allowbreak{}monoid-\allowbreak{}in-\allowbreak{}set}}{\vnbstmt{\begin{array}{@{}l@{}} \forall a, b. \\ \quad \operatorname{is\text{-}comm\text{-}monoid}(a) \wedge \\ \quad \operatorname{is\text{-}comm\text{-}monoid}(b) \Rightarrow \\ \quad \quad (\operatorname{hom\text{-}comm\text{-}monoid}(a, b) \in \mathbf{Set}) \end{array}}}
\vnbentry{\texttt{hom-\allowbreak{}group-\allowbreak{}member-\allowbreak{}iff}}{\vnbstmt{\begin{array}{@{}l@{}} \forall a, b, f. \\ \quad ((f \in \operatorname{hom\text{-}group}(a, b)) \Leftrightarrow ((f \in (\operatorname{carr}(a) \to \operatorname{carr}(b))) \wedge \operatorname{is\text{-}hom\text{-}group}(a, b, f))) \end{array}}}
\vnbentry{\texttt{hom-\allowbreak{}group-\allowbreak{}id}}{\vnbstmt{\begin{array}{@{}l@{}} \forall a. \\ \quad \operatorname{is\text{-}group}(a) \Rightarrow  \\ \quad \quad \operatorname{is\text{-}hom\text{-}group}(a, a, \operatorname{id\text{-}fun}(\operatorname{carr}(a))) \end{array}}}
\vnbentry{\texttt{hom-\allowbreak{}group-\allowbreak{}compose}}{\vnbstmt{\begin{array}{@{}l@{}} \forall a, b, c, f, g. \\ \quad \operatorname{is\text{-}hom\text{-}group}(a, b, f) \wedge \\ \quad \operatorname{is\text{-}hom\text{-}group}(b, c, g) \Rightarrow \\ \quad \quad \operatorname{is\text{-}hom\text{-}group}(a, c, \operatorname{compose}(g, f)) \end{array}}}
\vnbentry{\texttt{hom-\allowbreak{}group-\allowbreak{}in-\allowbreak{}set}}{\vnbstmt{\begin{array}{@{}l@{}} \forall a, b. \\ \quad (\operatorname{is\text{-}group}(a) \wedge \operatorname{is\text{-}group}(b)) \Rightarrow  \\ \quad \quad (\operatorname{hom\text{-}group}(a, b) \in \mathbf{Set}) \end{array}}}
\vnbentry{\texttt{hom-\allowbreak{}abelian-\allowbreak{}group-\allowbreak{}member-\allowbreak{}iff}}{\vnbstmt{\begin{array}{@{}l@{}} \forall a, b, f. \\ \quad ((f \in \operatorname{hom\text{-}abelian\text{-}group}(a, b)) \Leftrightarrow ((f \in (\operatorname{carr}(a) \to \operatorname{carr}(b))) \wedge \operatorname{is\text{-}hom\text{-}abelian\text{-}group}(a, b, f))) \end{array}}}
\vnbentry{\texttt{hom-\allowbreak{}abelian-\allowbreak{}group-\allowbreak{}id}}{\vnbstmt{\begin{array}{@{}l@{}} \forall a. \\ \quad \operatorname{is\text{-}abelian\text{-}group}(a) \Rightarrow  \\ \quad \quad \operatorname{is\text{-}hom\text{-}abelian\text{-}group}(a, a, \operatorname{id\text{-}fun}(\operatorname{carr}(a))) \end{array}}}
\vnbentry{\texttt{hom-\allowbreak{}abelian-\allowbreak{}group-\allowbreak{}compose}}{\vnbstmt{\begin{array}{@{}l@{}} \forall a, b, c, f, g. \\ \quad \operatorname{is\text{-}hom\text{-}abelian\text{-}group}(a, b, f) \wedge \\ \quad \operatorname{is\text{-}hom\text{-}abelian\text{-}group}(b, c, g) \Rightarrow \\ \quad \quad \operatorname{is\text{-}hom\text{-}abelian\text{-}group}(a, c, \operatorname{compose}(g, f)) \end{array}}}
\vnbentry{\texttt{hom-\allowbreak{}abelian-\allowbreak{}group-\allowbreak{}in-\allowbreak{}set}}{\vnbstmt{\begin{array}{@{}l@{}} \forall a, b. \\ \quad \operatorname{is\text{-}abelian\text{-}group}(a) \wedge \\ \quad \operatorname{is\text{-}abelian\text{-}group}(b) \Rightarrow \\ \quad \quad (\operatorname{hom\text{-}abelian\text{-}group}(a, b) \in \mathbf{Set}) \end{array}}}
\vnbentry{\texttt{hom-\allowbreak{}ring-\allowbreak{}member-\allowbreak{}iff}}{\vnbstmt{\begin{array}{@{}l@{}} \forall a, b, f. \\ \quad ((f \in \operatorname{hom\text{-}ring}(a, b)) \Leftrightarrow ((f \in (\operatorname{carr}(a) \to \operatorname{carr}(b))) \wedge \operatorname{is\text{-}hom\text{-}ring}(a, b, f))) \end{array}}}
\vnbentry{\texttt{hom-\allowbreak{}ring-\allowbreak{}id}}{\vnbstmt{\begin{array}{@{}l@{}} \forall a. \\ \quad \operatorname{is\text{-}ring}(a) \Rightarrow  \\ \quad \quad \operatorname{is\text{-}hom\text{-}ring}(a, a, \operatorname{id\text{-}fun}(\operatorname{carr}(a))) \end{array}}}
\vnbentry{\texttt{hom-\allowbreak{}ring-\allowbreak{}compose}}{\vnbstmt{\begin{array}{@{}l@{}} \forall a, b, c, f, g. \\ \quad \operatorname{is\text{-}hom\text{-}ring}(a, b, f) \wedge \\ \quad \operatorname{is\text{-}hom\text{-}ring}(b, c, g) \Rightarrow \\ \quad \quad \operatorname{is\text{-}hom\text{-}ring}(a, c, \operatorname{compose}(g, f)) \end{array}}}
\vnbentry{\texttt{hom-\allowbreak{}ring-\allowbreak{}in-\allowbreak{}set}}{\vnbstmt{\begin{array}{@{}l@{}} \forall a, b. \\ \quad (\operatorname{is\text{-}ring}(a) \wedge \operatorname{is\text{-}ring}(b)) \Rightarrow  \\ \quad \quad (\operatorname{hom\text{-}ring}(a, b) \in \mathbf{Set}) \end{array}}}
\vnbentry{\texttt{hom-\allowbreak{}commutative-\allowbreak{}ring-\allowbreak{}member-\allowbreak{}iff}}{\vnbstmt{\begin{array}{@{}l@{}} \forall a, b, f. \\ \quad ((f \in \operatorname{hom\text{-}commutative\text{-}ring}(a, b)) \Leftrightarrow ((f \in (\operatorname{carr}(a) \to \operatorname{carr}(b))) \wedge \operatorname{is\text{-}hom\text{-}commutative\text{-}ring}(a, b, f))) \end{array}}}
\vnbentry{\texttt{hom-\allowbreak{}commutative-\allowbreak{}ring-\allowbreak{}id}}{\vnbstmt{\begin{array}{@{}l@{}} \forall a. \\ \quad \operatorname{is\text{-}commutative\text{-}ring}(a) \Rightarrow  \\ \quad \quad \operatorname{is\text{-}hom\text{-}commutative\text{-}ring}(a, a, \operatorname{id\text{-}fun}(\operatorname{carr}(a))) \end{array}}}
\vnbentry{\texttt{hom-\allowbreak{}commutative-\allowbreak{}ring-\allowbreak{}compose}}{\vnbstmt{\begin{array}{@{}l@{}} \forall a, b, c, f, g. \\ \quad \operatorname{is\text{-}hom\text{-}commutative\text{-}ring}(a, b, f) \wedge \\ \quad \operatorname{is\text{-}hom\text{-}commutative\text{-}ring}(b, c, g) \Rightarrow \\ \quad \quad \operatorname{is\text{-}hom\text{-}commutative\text{-}ring}(a, c, \operatorname{compose}(g, f)) \end{array}}}
\vnbentry{\texttt{hom-\allowbreak{}commutative-\allowbreak{}ring-\allowbreak{}in-\allowbreak{}set}}{\vnbstmt{\begin{array}{@{}l@{}} \forall a, b. \\ \quad \operatorname{is\text{-}commutative\text{-}ring}(a) \wedge \\ \quad \operatorname{is\text{-}commutative\text{-}ring}(b) \Rightarrow \\ \quad \quad (\operatorname{hom\text{-}commutative\text{-}ring}(a, b) \in \mathbf{Set}) \end{array}}}
\vnbentry{\texttt{hom-\allowbreak{}integral-\allowbreak{}domain-\allowbreak{}member-\allowbreak{}iff}}{\vnbstmt{\begin{array}{@{}l@{}} \forall a, b, f. \\ \quad ((f \in \operatorname{hom\text{-}integral\text{-}domain}(a, b)) \Leftrightarrow ((f \in (\operatorname{carr}(a) \to \operatorname{carr}(b))) \wedge \operatorname{is\text{-}hom\text{-}integral\text{-}domain}(a, b, f))) \end{array}}}
\vnbentry{\texttt{hom-\allowbreak{}integral-\allowbreak{}domain-\allowbreak{}id}}{\vnbstmt{\begin{array}{@{}l@{}} \forall a. \\ \quad \operatorname{is\text{-}integral\text{-}domain}(a) \Rightarrow  \\ \quad \quad \operatorname{is\text{-}hom\text{-}integral\text{-}domain}(a, a, \operatorname{id\text{-}fun}(\operatorname{carr}(a))) \end{array}}}
\vnbentry{\texttt{hom-\allowbreak{}integral-\allowbreak{}domain-\allowbreak{}compose}}{\vnbstmt{\begin{array}{@{}l@{}} \forall a, b, c, f, g. \\ \quad \operatorname{is\text{-}hom\text{-}integral\text{-}domain}(a, b, f) \wedge \\ \quad \operatorname{is\text{-}hom\text{-}integral\text{-}domain}(b, c, g) \Rightarrow \\ \quad \quad \operatorname{is\text{-}hom\text{-}integral\text{-}domain}(a, c, \operatorname{compose}(g, f)) \end{array}}}
\vnbentry{\texttt{hom-\allowbreak{}integral-\allowbreak{}domain-\allowbreak{}in-\allowbreak{}set}}{\vnbstmt{\begin{array}{@{}l@{}} \forall a, b. \\ \quad \operatorname{is\text{-}integral\text{-}domain}(a) \wedge \\ \quad \operatorname{is\text{-}integral\text{-}domain}(b) \Rightarrow \\ \quad \quad (\operatorname{hom\text{-}integral\text{-}domain}(a, b) \in \mathbf{Set}) \end{array}}}
\vnbentry{\texttt{hom-\allowbreak{}euclidean-\allowbreak{}ring-\allowbreak{}member-\allowbreak{}iff}}{\vnbstmt{\begin{array}{@{}l@{}} \forall a, b, f. \\ \quad ((f \in \operatorname{hom\text{-}euclidean\text{-}ring}(a, b)) \Leftrightarrow ((f \in (\operatorname{carr}(a) \to \operatorname{carr}(b))) \wedge \operatorname{is\text{-}hom\text{-}euclidean\text{-}ring}(a, b, f))) \end{array}}}
\vnbentry{\texttt{hom-\allowbreak{}euclidean-\allowbreak{}ring-\allowbreak{}id}}{\vnbstmt{\begin{array}{@{}l@{}} \forall a. \\ \quad \operatorname{is\text{-}euclidean\text{-}ring}(a) \Rightarrow  \\ \quad \quad \operatorname{is\text{-}hom\text{-}euclidean\text{-}ring}(a, a, \operatorname{id\text{-}fun}(\operatorname{carr}(a))) \end{array}}}
\vnbentry{\texttt{hom-\allowbreak{}euclidean-\allowbreak{}ring-\allowbreak{}compose}}{\vnbstmt{\begin{array}{@{}l@{}} \forall a, b, c, f, g. \\ \quad \operatorname{is\text{-}hom\text{-}euclidean\text{-}ring}(a, b, f) \wedge \\ \quad \operatorname{is\text{-}hom\text{-}euclidean\text{-}ring}(b, c, g) \Rightarrow \\ \quad \quad \operatorname{is\text{-}hom\text{-}euclidean\text{-}ring}(a, c, \operatorname{compose}(g, f)) \end{array}}}
\vnbentry{\texttt{hom-\allowbreak{}euclidean-\allowbreak{}ring-\allowbreak{}in-\allowbreak{}set}}{\vnbstmt{\begin{array}{@{}l@{}} \forall a, b. \\ \quad \operatorname{is\text{-}euclidean\text{-}ring}(a) \wedge \\ \quad \operatorname{is\text{-}euclidean\text{-}ring}(b) \Rightarrow \\ \quad \quad (\operatorname{hom\text{-}euclidean\text{-}ring}(a, b) \in \mathbf{Set}) \end{array}}}
\vnbentry{\texttt{hom-\allowbreak{}pid-\allowbreak{}member-\allowbreak{}iff}}{\vnbstmt{\begin{array}{@{}l@{}} \forall a, b, f. \\ \quad ((f \in \operatorname{hom\text{-}pid}(a, b)) \Leftrightarrow ((f \in (\operatorname{carr}(a) \to \operatorname{carr}(b))) \wedge \operatorname{is\text{-}hom\text{-}pid}(a, b, f))) \end{array}}}
\vnbentry{\texttt{hom-\allowbreak{}pid-\allowbreak{}id}}{\vnbstmt{\begin{array}{@{}l@{}} \forall a. \\ \quad \operatorname{is\text{-}pid}(a) \Rightarrow  \\ \quad \quad \operatorname{is\text{-}hom\text{-}pid}(a, a, \operatorname{id\text{-}fun}(\operatorname{carr}(a))) \end{array}}}
\vnbentry{\texttt{hom-\allowbreak{}pid-\allowbreak{}compose}}{\vnbstmt{\begin{array}{@{}l@{}} \forall a, b, c, f, g. \\ \quad \operatorname{is\text{-}hom\text{-}pid}(a, b, f) \wedge \\ \quad \operatorname{is\text{-}hom\text{-}pid}(b, c, g) \Rightarrow \\ \quad \quad \operatorname{is\text{-}hom\text{-}pid}(a, c, \operatorname{compose}(g, f)) \end{array}}}
\vnbentry{\texttt{hom-\allowbreak{}pid-\allowbreak{}in-\allowbreak{}set}}{\vnbstmt{\begin{array}{@{}l@{}} \forall a, b. \\ \quad (\operatorname{is\text{-}pid}(a) \wedge \operatorname{is\text{-}pid}(b)) \Rightarrow  \\ \quad \quad (\operatorname{hom\text{-}pid}(a, b) \in \mathbf{Set}) \end{array}}}
\vnbentry{\texttt{hom-\allowbreak{}field-\allowbreak{}ring-\allowbreak{}member-\allowbreak{}iff}}{\vnbstmt{\begin{array}{@{}l@{}} \forall a, b, f. \\ \quad ((f \in \operatorname{hom\text{-}field\text{-}ring}(a, b)) \Leftrightarrow ((f \in (\operatorname{carr}(a) \to \operatorname{carr}(b))) \wedge \operatorname{is\text{-}hom\text{-}field\text{-}ring}(a, b, f))) \end{array}}}
\vnbentry{\texttt{hom-\allowbreak{}field-\allowbreak{}ring-\allowbreak{}id}}{\vnbstmt{\begin{array}{@{}l@{}} \forall a. \\ \quad \operatorname{is\text{-}field\text{-}ring}(a) \Rightarrow  \\ \quad \quad \operatorname{is\text{-}hom\text{-}field\text{-}ring}(a, a, \operatorname{id\text{-}fun}(\operatorname{carr}(a))) \end{array}}}
\vnbentry{\texttt{hom-\allowbreak{}field-\allowbreak{}ring-\allowbreak{}compose}}{\vnbstmt{\begin{array}{@{}l@{}} \forall a, b, c, f, g. \\ \quad \operatorname{is\text{-}hom\text{-}field\text{-}ring}(a, b, f) \wedge \\ \quad \operatorname{is\text{-}hom\text{-}field\text{-}ring}(b, c, g) \Rightarrow \\ \quad \quad \operatorname{is\text{-}hom\text{-}field\text{-}ring}(a, c, \operatorname{compose}(g, f)) \end{array}}}
\vnbentry{\texttt{hom-\allowbreak{}field-\allowbreak{}ring-\allowbreak{}in-\allowbreak{}set}}{\vnbstmt{\begin{array}{@{}l@{}} \forall a, b. \\ \quad (\operatorname{is\text{-}field\text{-}ring}(a) \wedge \operatorname{is\text{-}field\text{-}ring}(b)) \Rightarrow  \\ \quad \quad (\operatorname{hom\text{-}field\text{-}ring}(a, b) \in \mathbf{Set}) \end{array}}}
\vnbentry{\texttt{hom-\allowbreak{}normed-\allowbreak{}ag-\allowbreak{}member-\allowbreak{}iff}}{\vnbstmt{\begin{array}{@{}l@{}} \forall a, b, f. \\ \quad ((f \in \operatorname{hom\text{-}normed\text{-}ag}(a, b)) \Leftrightarrow ((f \in (\operatorname{carr}(a) \to \operatorname{carr}(b))) \wedge \operatorname{is\text{-}hom\text{-}normed\text{-}ag}(a, b, f))) \end{array}}}
\vnbentry{\texttt{hom-\allowbreak{}normed-\allowbreak{}ag-\allowbreak{}id}}{\vnbstmt{\begin{array}{@{}l@{}} \forall a. \\ \quad \operatorname{is\text{-}normed\text{-}ag}(a) \Rightarrow  \\ \quad \quad \operatorname{is\text{-}hom\text{-}normed\text{-}ag}(a, a, \operatorname{id\text{-}fun}(\operatorname{carr}(a))) \end{array}}}
\vnbentry{\texttt{hom-\allowbreak{}normed-\allowbreak{}ag-\allowbreak{}compose}}{\vnbstmt{\begin{array}{@{}l@{}} \forall a, b, c, f, g. \\ \quad \operatorname{is\text{-}hom\text{-}normed\text{-}ag}(a, b, f) \wedge \\ \quad \operatorname{is\text{-}hom\text{-}normed\text{-}ag}(b, c, g) \Rightarrow \\ \quad \quad \operatorname{is\text{-}hom\text{-}normed\text{-}ag}(a, c, \operatorname{compose}(g, f)) \end{array}}}
\vnbentry{\texttt{hom-\allowbreak{}normed-\allowbreak{}ag-\allowbreak{}in-\allowbreak{}set}}{\vnbstmt{\begin{array}{@{}l@{}} \forall a, b. \\ \quad (\operatorname{is\text{-}normed\text{-}ag}(a) \wedge \operatorname{is\text{-}normed\text{-}ag}(b)) \Rightarrow  \\ \quad \quad (\operatorname{hom\text{-}normed\text{-}ag}(a, b) \in \mathbf{Set}) \end{array}}}
\vnbentry{\texttt{hom-\allowbreak{}normed-\allowbreak{}field-\allowbreak{}member-\allowbreak{}iff}}{\vnbstmt{\begin{array}{@{}l@{}} \forall a, b, f. \\ \quad ((f \in \operatorname{hom\text{-}normed\text{-}field}(a, b)) \Leftrightarrow ((f \in (\operatorname{carr}(a) \to \operatorname{carr}(b))) \wedge \operatorname{is\text{-}hom\text{-}normed\text{-}field}(a, b, f))) \end{array}}}
\vnbentry{\texttt{hom-\allowbreak{}normed-\allowbreak{}field-\allowbreak{}id}}{\vnbstmt{\begin{array}{@{}l@{}} \forall a. \\ \quad \operatorname{is\text{-}normed\text{-}field}(a) \Rightarrow  \\ \quad \quad \operatorname{is\text{-}hom\text{-}normed\text{-}field}(a, a, \operatorname{id\text{-}fun}(\operatorname{carr}(a))) \end{array}}}
\vnbentry{\texttt{hom-\allowbreak{}normed-\allowbreak{}field-\allowbreak{}compose}}{\vnbstmt{\begin{array}{@{}l@{}} \forall a, b, c, f, g. \\ \quad \operatorname{is\text{-}hom\text{-}normed\text{-}field}(a, b, f) \wedge \\ \quad \operatorname{is\text{-}hom\text{-}normed\text{-}field}(b, c, g) \Rightarrow \\ \quad \quad \operatorname{is\text{-}hom\text{-}normed\text{-}field}(a, c, \operatorname{compose}(g, f)) \end{array}}}
\vnbentry{\texttt{hom-\allowbreak{}normed-\allowbreak{}field-\allowbreak{}in-\allowbreak{}set}}{\vnbstmt{\begin{array}{@{}l@{}} \forall a, b. \\ \quad \operatorname{is\text{-}normed\text{-}field}(a) \wedge \\ \quad \operatorname{is\text{-}normed\text{-}field}(b) \Rightarrow \\ \quad \quad (\operatorname{hom\text{-}normed\text{-}field}(a, b) \in \mathbf{Set}) \end{array}}}
\vnbentry{\texttt{hom-\allowbreak{}module-\allowbreak{}member-\allowbreak{}iff}}{\vnbstmt{\begin{array}{@{}l@{}} \forall a, b, f. \\ \quad ((f \in \operatorname{hom\text{-}module}(a, b)) \Leftrightarrow ((f \in (\operatorname{vec}(a) \to \operatorname{vec}(b))) \wedge \operatorname{is\text{-}hom\text{-}module}(a, b, f))) \end{array}}}
\vnbentry{\texttt{hom-\allowbreak{}module-\allowbreak{}id}}{\vnbstmt{\begin{array}{@{}l@{}} \forall a. \\ \quad \operatorname{is\text{-}module}(a) \Rightarrow  \\ \quad \quad \operatorname{is\text{-}hom\text{-}module}(a, a, \operatorname{id\text{-}fun}(\operatorname{vec}(a))) \end{array}}}
\vnbentry{\texttt{hom-\allowbreak{}module-\allowbreak{}compose}}{\vnbstmt{\begin{array}{@{}l@{}} \forall a, b, c, f, g. \\ \quad \operatorname{is\text{-}hom\text{-}module}(a, b, f) \wedge \\ \quad \operatorname{is\text{-}hom\text{-}module}(b, c, g) \Rightarrow \\ \quad \quad \operatorname{is\text{-}hom\text{-}module}(a, c, \operatorname{compose}(g, f)) \end{array}}}
\vnbentry{\texttt{hom-\allowbreak{}module-\allowbreak{}in-\allowbreak{}set}}{\vnbstmt{\begin{array}{@{}l@{}} \forall a, b. \\ \quad (\operatorname{is\text{-}module}(a) \wedge \operatorname{is\text{-}module}(b)) \Rightarrow  \\ \quad \quad (\operatorname{hom\text{-}module}(a, b) \in \mathbf{Set}) \end{array}}}
\vnbentry{\texttt{hom-\allowbreak{}vector-\allowbreak{}space-\allowbreak{}member-\allowbreak{}iff}}{\vnbstmt{\begin{array}{@{}l@{}} \forall a, b, f. \\ \quad ((f \in \operatorname{hom\text{-}vector\text{-}space}(a, b)) \Leftrightarrow ((f \in (\operatorname{vec}(a) \to \operatorname{vec}(b))) \wedge \operatorname{is\text{-}hom\text{-}vector\text{-}space}(a, b, f))) \end{array}}}
\vnbentry{\texttt{hom-\allowbreak{}vector-\allowbreak{}space-\allowbreak{}id}}{\vnbstmt{\begin{array}{@{}l@{}} \forall a. \\ \quad \operatorname{is\text{-}vector\text{-}space}(a) \Rightarrow  \\ \quad \quad \operatorname{is\text{-}hom\text{-}vector\text{-}space}(a, a, \operatorname{id\text{-}fun}(\operatorname{vec}(a))) \end{array}}}
\vnbentry{\texttt{hom-\allowbreak{}vector-\allowbreak{}space-\allowbreak{}compose}}{\vnbstmt{\begin{array}{@{}l@{}} \forall a, b, c, f, g. \\ \quad \operatorname{is\text{-}hom\text{-}vector\text{-}space}(a, b, f) \wedge \\ \quad \operatorname{is\text{-}hom\text{-}vector\text{-}space}(b, c, g) \Rightarrow \\ \quad \quad \operatorname{is\text{-}hom\text{-}vector\text{-}space}(a, c, \operatorname{compose}(g, f)) \end{array}}}
\vnbentry{\texttt{hom-\allowbreak{}vector-\allowbreak{}space-\allowbreak{}in-\allowbreak{}set}}{\vnbstmt{\begin{array}{@{}l@{}} \forall a, b. \\ \quad \operatorname{is\text{-}vector\text{-}space}(a) \wedge \\ \quad \operatorname{is\text{-}vector\text{-}space}(b) \Rightarrow \\ \quad \quad (\operatorname{hom\text{-}vector\text{-}space}(a, b) \in \mathbf{Set}) \end{array}}}
\vnbentry{\texttt{hom-\allowbreak{}normed-\allowbreak{}vector-\allowbreak{}space-\allowbreak{}member-\allowbreak{}iff}}{\vnbstmt{\begin{array}{@{}l@{}} \forall a, b, f. \\ \quad ((f \in \operatorname{hom\text{-}normed\text{-}vector\text{-}space}(a, b)) \Leftrightarrow ((f \in (\operatorname{vec}(a) \to \operatorname{vec}(b))) \wedge \operatorname{is\text{-}hom\text{-}normed\text{-}vector\text{-}space}(a, b, f))) \end{array}}}
\vnbentry{\texttt{hom-\allowbreak{}normed-\allowbreak{}vector-\allowbreak{}space-\allowbreak{}id}}{\vnbstmt{\begin{array}{@{}l@{}} \forall a. \\ \quad \operatorname{is\text{-}normed\text{-}vector\text{-}space}(a) \Rightarrow  \\ \quad \quad \operatorname{is\text{-}hom\text{-}normed\text{-}vector\text{-}space}(a, a, \operatorname{id\text{-}fun}(\operatorname{vec}(a))) \end{array}}}
\vnbentry{\texttt{hom-\allowbreak{}normed-\allowbreak{}vector-\allowbreak{}space-\allowbreak{}compose}}{\vnbstmt{\begin{array}{@{}l@{}} \forall a, b, c, f, g. \\ \quad \operatorname{is\text{-}hom\text{-}normed\text{-}vector\text{-}space}(a, b, f) \wedge \\ \quad \operatorname{is\text{-}hom\text{-}normed\text{-}vector\text{-}space}(b, c, g) \Rightarrow \\ \quad \quad \operatorname{is\text{-}hom\text{-}normed\text{-}vector\text{-}space}(a, c, \operatorname{compose}(g, f)) \end{array}}}
\vnbentry{\texttt{hom-\allowbreak{}normed-\allowbreak{}vector-\allowbreak{}space-\allowbreak{}in-\allowbreak{}set}}{\vnbstmt{\begin{array}{@{}l@{}} \forall a, b. \\ \quad \operatorname{is\text{-}normed\text{-}vector\text{-}space}(a) \wedge \\ \quad \operatorname{is\text{-}normed\text{-}vector\text{-}space}(b) \Rightarrow \\ \quad \quad (\operatorname{hom\text{-}normed\text{-}vector\text{-}space}(a, b) \in \mathbf{Set}) \end{array}}}
\vnbentry{\texttt{hom-\allowbreak{}complex-\allowbreak{}inner-\allowbreak{}product-\allowbreak{}space-\allowbreak{}member-\allowbreak{}iff}}{\vnbstmt{\begin{array}{@{}l@{}} \forall a, b, f. \\ \quad ((f \in \operatorname{hom\text{-}complex\text{-}inner\text{-}product\text{-}space}(a, b)) \Leftrightarrow ((f \in (\operatorname{vec}(a) \to \operatorname{vec}(b))) \wedge \operatorname{is\text{-}hom\text{-}complex\text{-}inner\text{-}product\text{-}space}(a, b, f))) \end{array}}}
\vnbentry{\texttt{hom-\allowbreak{}complex-\allowbreak{}inner-\allowbreak{}product-\allowbreak{}space-\allowbreak{}id}}{\vnbstmt{\begin{array}{@{}l@{}} \forall a. \\ \quad \operatorname{is\text{-}complex\text{-}inner\text{-}product\text{-}space}(a) \Rightarrow  \\ \quad \quad \operatorname{is\text{-}hom\text{-}complex\text{-}inner\text{-}product\text{-}space}(a, a, \operatorname{id\text{-}fun}(\operatorname{vec}(a))) \end{array}}}
\vnbentry{\texttt{hom-\allowbreak{}complex-\allowbreak{}inner-\allowbreak{}product-\allowbreak{}space-\allowbreak{}compose}}{\vnbstmt{\begin{array}{@{}l@{}} \forall a, b, c, f, g. \\ \quad \operatorname{is\text{-}hom\text{-}complex\text{-}inner\text{-}product\text{-}space}(a, b, f) \wedge \\ \quad \operatorname{is\text{-}hom\text{-}complex\text{-}inner\text{-}product\text{-}space}(b, c, g) \Rightarrow \\ \quad \quad \operatorname{is\text{-}hom\text{-}complex\text{-}inner\text{-}product\text{-}space}(a, c, \operatorname{compose}(g, f)) \end{array}}}
\vnbentry{\texttt{hom-\allowbreak{}complex-\allowbreak{}inner-\allowbreak{}product-\allowbreak{}space-\allowbreak{}in-\allowbreak{}set}}{\vnbstmt{\begin{array}{@{}l@{}} \forall a, b. \\ \quad \operatorname{is\text{-}complex\text{-}inner\text{-}product\text{-}space}(a) \wedge \\ \quad \operatorname{is\text{-}complex\text{-}inner\text{-}product\text{-}space}(b) \Rightarrow \\ \quad \quad (\operatorname{hom\text{-}complex\text{-}inner\text{-}product\text{-}space}(a, b) \in \mathbf{Set}) \end{array}}}
\vnbentry{\texttt{hom-\allowbreak{}metric-\allowbreak{}space-\allowbreak{}member-\allowbreak{}iff}}{\vnbstmt{\begin{array}{@{}l@{}} \forall a, b, f. \\ \quad ((f \in \operatorname{hom\text{-}metric\text{-}space}(a, b)) \Leftrightarrow ((f \in (\operatorname{pts}(a) \to \operatorname{pts}(b))) \wedge \operatorname{is\text{-}hom\text{-}metric\text{-}space}(a, b, f))) \end{array}}}
\vnbentry{\texttt{hom-\allowbreak{}metric-\allowbreak{}space-\allowbreak{}id}}{\vnbstmt{\begin{array}{@{}l@{}} \forall a. \\ \quad \operatorname{is\text{-}metric\text{-}space}(a) \Rightarrow  \\ \quad \quad \operatorname{is\text{-}hom\text{-}metric\text{-}space}(a, a, \operatorname{id\text{-}fun}(\operatorname{pts}(a))) \end{array}}}
\vnbentry{\texttt{hom-\allowbreak{}metric-\allowbreak{}space-\allowbreak{}compose}}{\vnbstmt{\begin{array}{@{}l@{}} \forall a, b, c, f, g. \\ \quad \operatorname{is\text{-}hom\text{-}metric\text{-}space}(a, b, f) \wedge \\ \quad \operatorname{is\text{-}hom\text{-}metric\text{-}space}(b, c, g) \Rightarrow \\ \quad \quad \operatorname{is\text{-}hom\text{-}metric\text{-}space}(a, c, \operatorname{compose}(g, f)) \end{array}}}
\vnbentry{\texttt{hom-\allowbreak{}metric-\allowbreak{}space-\allowbreak{}in-\allowbreak{}set}}{\vnbstmt{\begin{array}{@{}l@{}} \forall a, b. \\ \quad \operatorname{is\text{-}metric\text{-}space}(a) \wedge \\ \quad \operatorname{is\text{-}metric\text{-}space}(b) \Rightarrow \\ \quad \quad (\operatorname{hom\text{-}metric\text{-}space}(a, b) \in \mathbf{Set}) \end{array}}}
\vnbentry{\texttt{hom-\allowbreak{}pseudometric-\allowbreak{}space-\allowbreak{}member-\allowbreak{}iff}}{\vnbstmt{\begin{array}{@{}l@{}} \forall a, b, f. \\ \quad ((f \in \operatorname{hom\text{-}pseudometric\text{-}space}(a, b)) \Leftrightarrow ((f \in (\operatorname{pts}(a) \to \operatorname{pts}(b))) \wedge \operatorname{is\text{-}hom\text{-}pseudometric\text{-}space}(a, b, f))) \end{array}}}
\vnbentry{\texttt{hom-\allowbreak{}pseudometric-\allowbreak{}space-\allowbreak{}id}}{\vnbstmt{\begin{array}{@{}l@{}} \forall a. \\ \quad \operatorname{is\text{-}pseudometric\text{-}space}(a) \Rightarrow  \\ \quad \quad \operatorname{is\text{-}hom\text{-}pseudometric\text{-}space}(a, a, \operatorname{id\text{-}fun}(\operatorname{pts}(a))) \end{array}}}
\vnbentry{\texttt{hom-\allowbreak{}pseudometric-\allowbreak{}space-\allowbreak{}compose}}{\vnbstmt{\begin{array}{@{}l@{}} \forall a, b, c, f, g. \\ \quad \operatorname{is\text{-}hom\text{-}pseudometric\text{-}space}(a, b, f) \wedge \\ \quad \operatorname{is\text{-}hom\text{-}pseudometric\text{-}space}(b, c, g) \Rightarrow \\ \quad \quad \operatorname{is\text{-}hom\text{-}pseudometric\text{-}space}(a, c, \operatorname{compose}(g, f)) \end{array}}}
\vnbentry{\texttt{hom-\allowbreak{}pseudometric-\allowbreak{}space-\allowbreak{}in-\allowbreak{}set}}{\vnbstmt{\begin{array}{@{}l@{}} \forall a, b. \\ \quad \operatorname{is\text{-}pseudometric\text{-}space}(a) \wedge \\ \quad \operatorname{is\text{-}pseudometric\text{-}space}(b) \Rightarrow \\ \quad \quad (\operatorname{hom\text{-}pseudometric\text{-}space}(a, b) \in \mathbf{Set}) \end{array}}}
\vnbentry{\texttt{hom-\allowbreak{}c-\allowbreak{}metric-\allowbreak{}space-\allowbreak{}member-\allowbreak{}iff}}{\vnbstmt{\begin{array}{@{}l@{}} \forall a, b, f. \\ \quad ((f \in \operatorname{hom\text{-}c\text{-}metric\text{-}space}(a, b)) \Leftrightarrow ((f \in (\operatorname{pts}(a) \to \operatorname{pts}(b))) \wedge \operatorname{is\text{-}hom\text{-}c\text{-}metric\text{-}space}(a, b, f))) \end{array}}}
\vnbentry{\texttt{hom-\allowbreak{}c-\allowbreak{}metric-\allowbreak{}space-\allowbreak{}id}}{\vnbstmt{\begin{array}{@{}l@{}} \forall a. \\ \quad \operatorname{is\text{-}c\text{-}metric\text{-}space}(a) \Rightarrow  \\ \quad \quad \operatorname{is\text{-}hom\text{-}c\text{-}metric\text{-}space}(a, a, \operatorname{id\text{-}fun}(\operatorname{pts}(a))) \end{array}}}
\vnbentry{\texttt{hom-\allowbreak{}c-\allowbreak{}metric-\allowbreak{}space-\allowbreak{}compose}}{\vnbstmt{\begin{array}{@{}l@{}} \forall a, b, c, f, g. \\ \quad \operatorname{is\text{-}hom\text{-}c\text{-}metric\text{-}space}(a, b, f) \wedge \\ \quad \operatorname{is\text{-}hom\text{-}c\text{-}metric\text{-}space}(b, c, g) \Rightarrow \\ \quad \quad \operatorname{is\text{-}hom\text{-}c\text{-}metric\text{-}space}(a, c, \operatorname{compose}(g, f)) \end{array}}}
\vnbentry{\texttt{hom-\allowbreak{}c-\allowbreak{}metric-\allowbreak{}space-\allowbreak{}in-\allowbreak{}set}}{\vnbstmt{\begin{array}{@{}l@{}} \forall a, b. \\ \quad \operatorname{is\text{-}c\text{-}metric\text{-}space}(a) \wedge \\ \quad \operatorname{is\text{-}c\text{-}metric\text{-}space}(b) \Rightarrow \\ \quad \quad (\operatorname{hom\text{-}c\text{-}metric\text{-}space}(a, b) \in \mathbf{Set}) \end{array}}}
\vnbentry{\texttt{hom-\allowbreak{}field-\allowbreak{}non-\allowbreak{}zero}}{\vnbstmt{\begin{array}{@{}l@{}} \forall a, b, f.; \forall x \in \operatorname{non\text{-}zero}(a). \\ \quad \operatorname{is\text{-}hom\text{-}field}(a, b, f) \Rightarrow  \\ \quad \quad (f(x) \in \operatorname{non\text{-}zero}(b)) \end{array}}}
\vnbentry{\texttt{hom-\allowbreak{}field-\allowbreak{}member-\allowbreak{}iff}}{\vnbstmt{\begin{array}{@{}l@{}} \forall a, b, f. \\ \quad ((f \in \operatorname{hom\text{-}field}(a, b)) \Leftrightarrow ((f \in (\operatorname{carr}(a) \to \operatorname{carr}(b))) \wedge \operatorname{is\text{-}hom\text{-}field}(a, b, f))) \end{array}}}
\vnbentry{\texttt{hom-\allowbreak{}field-\allowbreak{}id}}{\vnbstmt{\begin{array}{@{}l@{}} \forall a. \\ \quad \operatorname{is\text{-}field}(a) \Rightarrow  \\ \quad \quad \operatorname{is\text{-}hom\text{-}field}(a, a, \operatorname{id\text{-}fun}(\operatorname{carr}(a))) \end{array}}}
\vnbentry{\texttt{hom-\allowbreak{}field-\allowbreak{}compose}}{\vnbstmt{\begin{array}{@{}l@{}} \forall a, b, c, f, g. \\ \quad \operatorname{is\text{-}hom\text{-}field}(a, b, f) \wedge \\ \quad \operatorname{is\text{-}hom\text{-}field}(b, c, g) \Rightarrow \\ \quad \quad \operatorname{is\text{-}hom\text{-}field}(a, c, \operatorname{compose}(g, f)) \end{array}}}
\vnbentry{\texttt{hom-\allowbreak{}field-\allowbreak{}in-\allowbreak{}set}}{\vnbstmt{\begin{array}{@{}l@{}} \forall a, b. \\ \quad (\operatorname{is\text{-}field}(a) \wedge \operatorname{is\text{-}field}(b)) \Rightarrow  \\ \quad \quad (\operatorname{hom\text{-}field}(a, b) \in \mathbf{Set}) \end{array}}}
\vnbentry{\texttt{hom-\allowbreak{}top-\allowbreak{}space-\allowbreak{}member-\allowbreak{}iff}}{\vnbstmt{\begin{array}{@{}l@{}} \forall a, b, f. \\ \quad ((f \in \operatorname{hom\text{-}top\text{-}space}(a, b)) \Leftrightarrow ((f \in (\operatorname{pts}(a) \to \operatorname{pts}(b))) \wedge \operatorname{is\text{-}hom\text{-}top\text{-}space}(a, b, f))) \end{array}}}
\vnbentry{\texttt{hom-\allowbreak{}top-\allowbreak{}space-\allowbreak{}id}}{\vnbstmt{\begin{array}{@{}l@{}} \forall s. \\ \quad \operatorname{is\text{-}top\text{-}space}(s) \Rightarrow  \\ \quad \quad \operatorname{is\text{-}hom\text{-}top\text{-}space}(s, s, \operatorname{id\text{-}fun}(\operatorname{pts}(s))) \end{array}}}
\vnbentry{\texttt{hom-\allowbreak{}top-\allowbreak{}space-\allowbreak{}compose}}{\vnbstmt{\begin{array}{@{}l@{}} \forall s, t, w, f, g. \\ \quad \operatorname{is\text{-}hom\text{-}top\text{-}space}(s, t, f) \wedge \\ \quad \operatorname{is\text{-}hom\text{-}top\text{-}space}(t, w, g) \Rightarrow \\ \quad \quad \operatorname{is\text{-}hom\text{-}top\text{-}space}(s, w, \operatorname{compose}(g, f)) \end{array}}}
\vnbentry{\texttt{hom-\allowbreak{}top-\allowbreak{}space-\allowbreak{}in-\allowbreak{}set}}{\vnbstmt{\begin{array}{@{}l@{}} \forall a, b. \\ \quad (\operatorname{is\text{-}top\text{-}space}(a) \wedge \operatorname{is\text{-}top\text{-}space}(b)) \Rightarrow  \\ \quad \quad (\operatorname{hom\text{-}top\text{-}space}(a, b) \in \mathbf{Set}) \end{array}}}
\vnbentry{\texttt{hom-\allowbreak{}metrizable-\allowbreak{}top-\allowbreak{}space-\allowbreak{}member-\allowbreak{}iff}}{\vnbstmt{\begin{array}{@{}l@{}} \forall a, b, f. \\ \quad ((f \in \operatorname{hom\text{-}metrizable\text{-}top\text{-}space}(a, b)) \Leftrightarrow ((f \in (\operatorname{pts}(a) \to \operatorname{pts}(b))) \wedge \operatorname{is\text{-}hom\text{-}metrizable\text{-}top\text{-}space}(a, b, f))) \end{array}}}
\vnbentry{\texttt{hom-\allowbreak{}metrizable-\allowbreak{}top-\allowbreak{}space-\allowbreak{}id}}{\vnbstmt{\begin{array}{@{}l@{}} \forall s. \\ \quad \operatorname{is\text{-}metrizable\text{-}top\text{-}space}(s) \Rightarrow  \\ \quad \quad \operatorname{is\text{-}hom\text{-}metrizable\text{-}top\text{-}space}(s, s, \operatorname{id\text{-}fun}(\operatorname{pts}(s))) \end{array}}}
\vnbentry{\texttt{hom-\allowbreak{}metrizable-\allowbreak{}top-\allowbreak{}space-\allowbreak{}compose}}{\vnbstmt{\begin{array}{@{}l@{}} \forall s, t, w, f, g. \\ \quad \operatorname{is\text{-}hom\text{-}metrizable\text{-}top\text{-}space}(s, t, f) \wedge \\ \quad \operatorname{is\text{-}hom\text{-}metrizable\text{-}top\text{-}space}(t, w, g) \Rightarrow \\ \quad \quad \operatorname{is\text{-}hom\text{-}metrizable\text{-}top\text{-}space}(s, w, \operatorname{compose}(g, f)) \end{array}}}
\vnbentry{\texttt{hom-\allowbreak{}metrizable-\allowbreak{}top-\allowbreak{}space-\allowbreak{}in-\allowbreak{}set}}{\vnbstmt{\begin{array}{@{}l@{}} \forall a, b. \\ \quad \operatorname{is\text{-}metrizable\text{-}top\text{-}space}(a) \wedge \\ \quad \operatorname{is\text{-}metrizable\text{-}top\text{-}space}(b) \Rightarrow \\ \quad \quad (\operatorname{hom\text{-}metrizable\text{-}top\text{-}space}(a, b) \in \mathbf{Set}) \end{array}}}
\vnbentry{\texttt{hom-\allowbreak{}ringoid-\allowbreak{}member-\allowbreak{}iff}}{\vnbstmt{\begin{array}{@{}l@{}} \forall a, b, f. \\ \quad ((f \in \operatorname{hom\text{-}ringoid}(a, b)) \Leftrightarrow ((f \in (\operatorname{carr}(a) \to \operatorname{carr}(b))) \wedge \operatorname{is\text{-}hom\text{-}ringoid}(a, b, f))) \end{array}}}
\vnbentry{\texttt{hom-\allowbreak{}ringoid-\allowbreak{}id}}{\vnbstmt{\begin{array}{@{}l@{}} \forall a. \\ \quad \operatorname{is\text{-}ringoid}(a) \Rightarrow  \\ \quad \quad \operatorname{is\text{-}hom\text{-}ringoid}(a, a, \operatorname{id\text{-}fun}(\operatorname{carr}(a))) \end{array}}}
\vnbentry{\texttt{hom-\allowbreak{}ringoid-\allowbreak{}compose}}{\vnbstmt{\begin{array}{@{}l@{}} \forall a, b, c, f, g. \\ \quad \operatorname{is\text{-}hom\text{-}ringoid}(a, b, f) \wedge \\ \quad \operatorname{is\text{-}hom\text{-}ringoid}(b, c, g) \Rightarrow \\ \quad \quad \operatorname{is\text{-}hom\text{-}ringoid}(a, c, \operatorname{compose}(g, f)) \end{array}}}
\vnbentry{\texttt{hom-\allowbreak{}ringoid-\allowbreak{}in-\allowbreak{}set}}{\vnbstmt{\begin{array}{@{}l@{}} \forall a, b. \\ \quad (\operatorname{is\text{-}ringoid}(a) \wedge \operatorname{is\text{-}ringoid}(b)) \Rightarrow  \\ \quad \quad (\operatorname{hom\text{-}ringoid}(a, b) \in \mathbf{Set}) \end{array}}}

\vnbfile{theorem-library/hom-kinds.scm}{theorem-library/hom-kinds.scm}
\noindent{\small hom-kinds.scm -- the category laws for the arrows the generic driver of hom-laws.scm does not reach: the FAMILY slot kinds, and the RELATION kind (SETOID, since batch 40; two carriers in batch 39).}\par
\vnbentry{\texttt{hom-\allowbreak{}measurable-\allowbreak{}space-\allowbreak{}member-\allowbreak{}iff}}{\vnbstmt{\begin{array}{@{}l@{}} \forall a, b, f. \\ \quad ((f \in \operatorname{hom\text{-}measurable\text{-}space}(a, b)) \Leftrightarrow ((f \in (\operatorname{pts}(a) \to \operatorname{pts}(b))) \wedge \operatorname{is\text{-}hom\text{-}measurable\text{-}space}(a, b, f))) \end{array}}}
\vnbentry{\texttt{hom-\allowbreak{}measurable-\allowbreak{}space-\allowbreak{}id}}{\vnbstmt{\begin{array}{@{}l@{}} \forall a. \\ \quad \operatorname{is\text{-}measurable\text{-}space}(a) \Rightarrow  \\ \quad \quad \operatorname{is\text{-}hom\text{-}measurable\text{-}space}(a, a, \operatorname{id\text{-}fun}(\operatorname{pts}(a))) \end{array}}}
\vnbentry{\texttt{hom-\allowbreak{}measurable-\allowbreak{}space-\allowbreak{}compose}}{\vnbstmt{\begin{array}{@{}l@{}} \forall a, b, c, f, g. \\ \quad \operatorname{is\text{-}hom\text{-}measurable\text{-}space}(a, b, f) \wedge \\ \quad \operatorname{is\text{-}hom\text{-}measurable\text{-}space}(b, c, g) \Rightarrow \\ \quad \quad \operatorname{is\text{-}hom\text{-}measurable\text{-}space}(a, c, \operatorname{compose}(g, f)) \end{array}}}
\vnbentry{\texttt{hom-\allowbreak{}measurable-\allowbreak{}space-\allowbreak{}in-\allowbreak{}set}}{\vnbstmt{\begin{array}{@{}l@{}} \forall a, b. \\ \quad \operatorname{is\text{-}measurable\text{-}space}(a) \wedge \\ \quad \operatorname{is\text{-}measurable\text{-}space}(b) \Rightarrow \\ \quad \quad (\operatorname{hom\text{-}measurable\text{-}space}(a, b) \in \mathbf{Set}) \end{array}}}
\vnbentry{\texttt{hom-\allowbreak{}measure-\allowbreak{}space-\allowbreak{}member-\allowbreak{}iff}}{\vnbstmt{\begin{array}{@{}l@{}} \forall a, b, f. \\ \quad ((f \in \operatorname{hom\text{-}measure\text{-}space}(a, b)) \Leftrightarrow ((f \in (\operatorname{pts}(a) \to \operatorname{pts}(b))) \wedge \operatorname{is\text{-}hom\text{-}measure\text{-}space}(a, b, f))) \end{array}}}
\vnbentry{\texttt{hom-\allowbreak{}measure-\allowbreak{}space-\allowbreak{}id}}{\vnbstmt{\begin{array}{@{}l@{}} \forall a. \\ \quad \operatorname{is\text{-}measure\text{-}space}(a) \Rightarrow  \\ \quad \quad \operatorname{is\text{-}hom\text{-}measure\text{-}space}(a, a, \operatorname{id\text{-}fun}(\operatorname{pts}(a))) \end{array}}}
\vnbentry{\texttt{hom-\allowbreak{}measure-\allowbreak{}space-\allowbreak{}compose}}{\vnbstmt{\begin{array}{@{}l@{}} \forall a, b, c, f, g. \\ \quad \operatorname{is\text{-}hom\text{-}measure\text{-}space}(a, b, f) \wedge \\ \quad \operatorname{is\text{-}hom\text{-}measure\text{-}space}(b, c, g) \Rightarrow \\ \quad \quad \operatorname{is\text{-}hom\text{-}measure\text{-}space}(a, c, \operatorname{compose}(g, f)) \end{array}}}
\vnbentry{\texttt{hom-\allowbreak{}measure-\allowbreak{}space-\allowbreak{}in-\allowbreak{}set}}{\vnbstmt{\begin{array}{@{}l@{}} \forall a, b. \\ \quad \operatorname{is\text{-}measure\text{-}space}(a) \wedge \\ \quad \operatorname{is\text{-}measure\text{-}space}(b) \Rightarrow \\ \quad \quad (\operatorname{hom\text{-}measure\text{-}space}(a, b) \in \mathbf{Set}) \end{array}}}
\vnbentry{\texttt{hom-\allowbreak{}setoid-\allowbreak{}member-\allowbreak{}iff}}{\vnbstmt{\begin{array}{@{}l@{}} \forall a, b, f. \\ \quad ((f \in \operatorname{hom\text{-}setoid}(a, b)) \Leftrightarrow ((f \in (\operatorname{pts}(a) \to \operatorname{pts}(b))) \wedge \operatorname{is\text{-}hom\text{-}setoid}(a, b, f))) \end{array}}}
\vnbentry{\texttt{hom-\allowbreak{}setoid-\allowbreak{}id}}{\vnbstmt{\begin{array}{@{}l@{}} \forall a. \\ \quad \operatorname{is\text{-}setoid}(a) \Rightarrow  \\ \quad \quad \operatorname{is\text{-}hom\text{-}setoid}(a, a, \operatorname{id\text{-}fun}(\operatorname{pts}(a))) \end{array}}}
\vnbentry{\texttt{hom-\allowbreak{}setoid-\allowbreak{}compose}}{\vnbstmt{\begin{array}{@{}l@{}} \forall a, b, c, f, g. \\ \quad \operatorname{is\text{-}hom\text{-}setoid}(a, b, f) \wedge \\ \quad \operatorname{is\text{-}hom\text{-}setoid}(b, c, g) \Rightarrow \\ \quad \quad \operatorname{is\text{-}hom\text{-}setoid}(a, c, \operatorname{compose}(g, f)) \end{array}}}
\vnbentry{\texttt{hom-\allowbreak{}setoid-\allowbreak{}in-\allowbreak{}set}}{\vnbstmt{\begin{array}{@{}l@{}} \forall a, b. \\ \quad (\operatorname{is\text{-}setoid}(a) \wedge \operatorname{is\text{-}setoid}(b)) \Rightarrow  \\ \quad \quad (\operatorname{hom\text{-}setoid}(a, b) \in \mathbf{Set}) \end{array}}}
\vnbentry{\texttt{measure-\allowbreak{}space-\allowbreak{}as-\allowbreak{}measurable-\allowbreak{}space-\allowbreak{}functorial}}{\vnbstmt{\begin{array}{@{}l@{}} \forall a, b, f. \\ \quad \operatorname{is\text{-}hom\text{-}measure\text{-}space}(a, b, f) \Rightarrow  \\ \quad \quad \operatorname{is\text{-}hom\text{-}measurable\text{-}space}(\operatorname{measure\text{-}space\text{-}as\text{-}measurable\text{-}space}(a), \\ \quad \quad \quad \operatorname{measure\text{-}space\text{-}as\text{-}measurable\text{-}space}(b), f) \end{array}}}

\vnbfile{theorem-library/hom-functors.scm}{theorem-library/hom-functors.scm}
\noindent{\small hom-functors.scm -- the Hom functors: Hom(a, -) covariant, Hom(-, b) contravariant, as maps on hom-sets.}\par
\vnbentry{\texttt{hom-\allowbreak{}semigroup-\allowbreak{}post-\allowbreak{}type}}{\vnbstmt{\begin{array}{@{}l@{}} \forall a, b, c, h. \\ \quad (\operatorname{is\text{-}semigroup}(a) \wedge (h \in \operatorname{hom\text{-}semigroup}(b, c))) \Rightarrow  \\ \quad \quad (\lambda f \in \operatorname{hom\text{-}semigroup}(a, b).\; \operatorname{compose}(h, f)) \\ \quad \quad \in\; (\operatorname{hom\text{-}semigroup}(a, b) \to \operatorname{hom\text{-}semigroup}(a, c)) \end{array}}}
\vnbentry{\texttt{hom-\allowbreak{}semigroup-\allowbreak{}pre-\allowbreak{}type}}{\vnbstmt{\begin{array}{@{}l@{}} \forall a, b, c, k. \\ \quad (\operatorname{is\text{-}semigroup}(b) \wedge (k \in \operatorname{hom\text{-}semigroup}(c, a))) \Rightarrow  \\ \quad \quad (\lambda f \in \operatorname{hom\text{-}semigroup}(a, b).\; \operatorname{compose}(f, k)) \\ \quad \quad \in\; (\operatorname{hom\text{-}semigroup}(a, b) \to \operatorname{hom\text{-}semigroup}(c, b)) \end{array}}}
\vnbentry{\texttt{hom-\allowbreak{}monoid-\allowbreak{}post-\allowbreak{}type}}{\vnbstmt{\begin{array}{@{}l@{}} \forall a, b, c, h. \\ \quad (\operatorname{is\text{-}monoid}(a) \wedge (h \in \operatorname{hom\text{-}monoid}(b, c))) \Rightarrow  \\ \quad \quad (\lambda f \in \operatorname{hom\text{-}monoid}(a, b).\; \operatorname{compose}(h, f)) \\ \quad \quad \in\; (\operatorname{hom\text{-}monoid}(a, b) \to \operatorname{hom\text{-}monoid}(a, c)) \end{array}}}
\vnbentry{\texttt{hom-\allowbreak{}monoid-\allowbreak{}pre-\allowbreak{}type}}{\vnbstmt{\begin{array}{@{}l@{}} \forall a, b, c, k. \\ \quad (\operatorname{is\text{-}monoid}(b) \wedge (k \in \operatorname{hom\text{-}monoid}(c, a))) \Rightarrow  \\ \quad \quad (\lambda f \in \operatorname{hom\text{-}monoid}(a, b).\; \operatorname{compose}(f, k)) \\ \quad \quad \in\; (\operatorname{hom\text{-}monoid}(a, b) \to \operatorname{hom\text{-}monoid}(c, b)) \end{array}}}
\vnbentry{\texttt{hom-\allowbreak{}comm-\allowbreak{}monoid-\allowbreak{}post-\allowbreak{}type}}{\vnbstmt{\begin{array}{@{}l@{}} \forall a, b, c, h. \\ \quad \operatorname{is\text{-}comm\text{-}monoid}(a) \wedge \\ \quad (h \in \operatorname{hom\text{-}comm\text{-}monoid}(b, c)) \Rightarrow \\ \quad \quad (\lambda f \in \operatorname{hom\text{-}comm\text{-}monoid}(a, b).\; \operatorname{compose}(h, f)) \\ \quad \quad \in\; (\operatorname{hom\text{-}comm\text{-}monoid}(a, b) \to \operatorname{hom\text{-}comm\text{-}monoid}(a, c)) \end{array}}}
\vnbentry{\texttt{hom-\allowbreak{}comm-\allowbreak{}monoid-\allowbreak{}pre-\allowbreak{}type}}{\vnbstmt{\begin{array}{@{}l@{}} \forall a, b, c, k. \\ \quad \operatorname{is\text{-}comm\text{-}monoid}(b) \wedge \\ \quad (k \in \operatorname{hom\text{-}comm\text{-}monoid}(c, a)) \Rightarrow \\ \quad \quad (\lambda f \in \operatorname{hom\text{-}comm\text{-}monoid}(a, b).\; \operatorname{compose}(f, k)) \\ \quad \quad \in\; (\operatorname{hom\text{-}comm\text{-}monoid}(a, b) \to \operatorname{hom\text{-}comm\text{-}monoid}(c, b)) \end{array}}}
\vnbentry{\texttt{hom-\allowbreak{}group-\allowbreak{}post-\allowbreak{}type}}{\vnbstmt{\begin{array}{@{}l@{}} \forall a, b, c, h. \\ \quad (\operatorname{is\text{-}group}(a) \wedge (h \in \operatorname{hom\text{-}group}(b, c))) \Rightarrow  \\ \quad \quad (\lambda f \in \operatorname{hom\text{-}group}(a, b).\; \operatorname{compose}(h, f)) \\ \quad \quad \in\; (\operatorname{hom\text{-}group}(a, b) \to \operatorname{hom\text{-}group}(a, c)) \end{array}}}
\vnbentry{\texttt{hom-\allowbreak{}group-\allowbreak{}pre-\allowbreak{}type}}{\vnbstmt{\begin{array}{@{}l@{}} \forall a, b, c, k. \\ \quad (\operatorname{is\text{-}group}(b) \wedge (k \in \operatorname{hom\text{-}group}(c, a))) \Rightarrow  \\ \quad \quad (\lambda f \in \operatorname{hom\text{-}group}(a, b).\; \operatorname{compose}(f, k)) \\ \quad \quad \in\; (\operatorname{hom\text{-}group}(a, b) \to \operatorname{hom\text{-}group}(c, b)) \end{array}}}
\vnbentry{\texttt{hom-\allowbreak{}abelian-\allowbreak{}group-\allowbreak{}post-\allowbreak{}type}}{\vnbstmt{\begin{array}{@{}l@{}} \forall a, b, c, h. \\ \quad \operatorname{is\text{-}abelian\text{-}group}(a) \wedge \\ \quad (h \in \operatorname{hom\text{-}abelian\text{-}group}(b, c)) \Rightarrow \\ \quad \quad (\lambda f \in \operatorname{hom\text{-}abelian\text{-}group}(a, b).\; \operatorname{compose}(h, f)) \\ \quad \quad \in\; (\operatorname{hom\text{-}abelian\text{-}group}(a, b) \to \operatorname{hom\text{-}abelian\text{-}group}(a, c)) \end{array}}}
\vnbentry{\texttt{hom-\allowbreak{}abelian-\allowbreak{}group-\allowbreak{}pre-\allowbreak{}type}}{\vnbstmt{\begin{array}{@{}l@{}} \forall a, b, c, k. \\ \quad \operatorname{is\text{-}abelian\text{-}group}(b) \wedge \\ \quad (k \in \operatorname{hom\text{-}abelian\text{-}group}(c, a)) \Rightarrow \\ \quad \quad (\lambda f \in \operatorname{hom\text{-}abelian\text{-}group}(a, b).\; \operatorname{compose}(f, k)) \\ \quad \quad \in\; (\operatorname{hom\text{-}abelian\text{-}group}(a, b) \to \operatorname{hom\text{-}abelian\text{-}group}(c, b)) \end{array}}}
\vnbentry{\texttt{hom-\allowbreak{}ring-\allowbreak{}post-\allowbreak{}type}}{\vnbstmt{\begin{array}{@{}l@{}} \forall a, b, c, h. \\ \quad (\operatorname{is\text{-}ring}(a) \wedge (h \in \operatorname{hom\text{-}ring}(b, c))) \Rightarrow  \\ \quad \quad (\lambda f \in \operatorname{hom\text{-}ring}(a, b).\; \operatorname{compose}(h, f)) \\ \quad \quad \in\; (\operatorname{hom\text{-}ring}(a, b) \to \operatorname{hom\text{-}ring}(a, c)) \end{array}}}
\vnbentry{\texttt{hom-\allowbreak{}ring-\allowbreak{}pre-\allowbreak{}type}}{\vnbstmt{\begin{array}{@{}l@{}} \forall a, b, c, k. \\ \quad (\operatorname{is\text{-}ring}(b) \wedge (k \in \operatorname{hom\text{-}ring}(c, a))) \Rightarrow  \\ \quad \quad (\lambda f \in \operatorname{hom\text{-}ring}(a, b).\; \operatorname{compose}(f, k)) \\ \quad \quad \in\; (\operatorname{hom\text{-}ring}(a, b) \to \operatorname{hom\text{-}ring}(c, b)) \end{array}}}
\vnbentry{\texttt{hom-\allowbreak{}commutative-\allowbreak{}ring-\allowbreak{}post-\allowbreak{}type}}{\vnbstmt{\begin{array}{@{}l@{}} \forall a, b, c, h. \\ \quad \operatorname{is\text{-}commutative\text{-}ring}(a) \wedge \\ \quad (h \in \operatorname{hom\text{-}commutative\text{-}ring}(b, c)) \Rightarrow \\ \quad \quad (\lambda f \in \operatorname{hom\text{-}commutative\text{-}ring}(a, b).\; \\ \quad \quad \quad \operatorname{compose}(h, f)) \\ \quad \quad \in\; (\operatorname{hom\text{-}commutative\text{-}ring}(a, b) \to \operatorname{hom\text{-}commutative\text{-}ring}(a, c)) \end{array}}}
\vnbentry{\texttt{hom-\allowbreak{}commutative-\allowbreak{}ring-\allowbreak{}pre-\allowbreak{}type}}{\vnbstmt{\begin{array}{@{}l@{}} \forall a, b, c, k. \\ \quad \operatorname{is\text{-}commutative\text{-}ring}(b) \wedge \\ \quad (k \in \operatorname{hom\text{-}commutative\text{-}ring}(c, a)) \Rightarrow \\ \quad \quad (\lambda f \in \operatorname{hom\text{-}commutative\text{-}ring}(a, b).\; \\ \quad \quad \quad \operatorname{compose}(f, k)) \\ \quad \quad \in\; (\operatorname{hom\text{-}commutative\text{-}ring}(a, b) \to \operatorname{hom\text{-}commutative\text{-}ring}(c, b)) \end{array}}}
\vnbentry{\texttt{hom-\allowbreak{}integral-\allowbreak{}domain-\allowbreak{}post-\allowbreak{}type}}{\vnbstmt{\begin{array}{@{}l@{}} \forall a, b, c, h. \\ \quad \operatorname{is\text{-}integral\text{-}domain}(a) \wedge \\ \quad (h \in \operatorname{hom\text{-}integral\text{-}domain}(b, c)) \Rightarrow \\ \quad \quad (\lambda f \in \operatorname{hom\text{-}integral\text{-}domain}(a, b).\; \\ \quad \quad \quad \operatorname{compose}(h, f)) \\ \quad \quad \in\; (\operatorname{hom\text{-}integral\text{-}domain}(a, b) \to \operatorname{hom\text{-}integral\text{-}domain}(a, c)) \end{array}}}
\vnbentry{\texttt{hom-\allowbreak{}integral-\allowbreak{}domain-\allowbreak{}pre-\allowbreak{}type}}{\vnbstmt{\begin{array}{@{}l@{}} \forall a, b, c, k. \\ \quad \operatorname{is\text{-}integral\text{-}domain}(b) \wedge \\ \quad (k \in \operatorname{hom\text{-}integral\text{-}domain}(c, a)) \Rightarrow \\ \quad \quad (\lambda f \in \operatorname{hom\text{-}integral\text{-}domain}(a, b).\; \\ \quad \quad \quad \operatorname{compose}(f, k)) \\ \quad \quad \in\; (\operatorname{hom\text{-}integral\text{-}domain}(a, b) \to \operatorname{hom\text{-}integral\text{-}domain}(c, b)) \end{array}}}
\vnbentry{\texttt{hom-\allowbreak{}euclidean-\allowbreak{}ring-\allowbreak{}post-\allowbreak{}type}}{\vnbstmt{\begin{array}{@{}l@{}} \forall a, b, c, h. \\ \quad \operatorname{is\text{-}euclidean\text{-}ring}(a) \wedge \\ \quad (h \in \operatorname{hom\text{-}euclidean\text{-}ring}(b, c)) \Rightarrow \\ \quad \quad (\lambda f \in \operatorname{hom\text{-}euclidean\text{-}ring}(a, b).\; \\ \quad \quad \quad \operatorname{compose}(h, f)) \\ \quad \quad \in\; (\operatorname{hom\text{-}euclidean\text{-}ring}(a, b) \to \operatorname{hom\text{-}euclidean\text{-}ring}(a, c)) \end{array}}}
\vnbentry{\texttt{hom-\allowbreak{}euclidean-\allowbreak{}ring-\allowbreak{}pre-\allowbreak{}type}}{\vnbstmt{\begin{array}{@{}l@{}} \forall a, b, c, k. \\ \quad \operatorname{is\text{-}euclidean\text{-}ring}(b) \wedge \\ \quad (k \in \operatorname{hom\text{-}euclidean\text{-}ring}(c, a)) \Rightarrow \\ \quad \quad (\lambda f \in \operatorname{hom\text{-}euclidean\text{-}ring}(a, b).\; \\ \quad \quad \quad \operatorname{compose}(f, k)) \\ \quad \quad \in\; (\operatorname{hom\text{-}euclidean\text{-}ring}(a, b) \to \operatorname{hom\text{-}euclidean\text{-}ring}(c, b)) \end{array}}}
\vnbentry{\texttt{hom-\allowbreak{}pid-\allowbreak{}post-\allowbreak{}type}}{\vnbstmt{\begin{array}{@{}l@{}} \forall a, b, c, h. \\ \quad (\operatorname{is\text{-}pid}(a) \wedge (h \in \operatorname{hom\text{-}pid}(b, c))) \Rightarrow  \\ \quad \quad (\lambda f \in \operatorname{hom\text{-}pid}(a, b).\; \operatorname{compose}(h, f)) \\ \quad \quad \in\; (\operatorname{hom\text{-}pid}(a, b) \to \operatorname{hom\text{-}pid}(a, c)) \end{array}}}
\vnbentry{\texttt{hom-\allowbreak{}pid-\allowbreak{}pre-\allowbreak{}type}}{\vnbstmt{\begin{array}{@{}l@{}} \forall a, b, c, k. \\ \quad (\operatorname{is\text{-}pid}(b) \wedge (k \in \operatorname{hom\text{-}pid}(c, a))) \Rightarrow  \\ \quad \quad (\lambda f \in \operatorname{hom\text{-}pid}(a, b).\; \operatorname{compose}(f, k)) \\ \quad \quad \in\; (\operatorname{hom\text{-}pid}(a, b) \to \operatorname{hom\text{-}pid}(c, b)) \end{array}}}
\vnbentry{\texttt{hom-\allowbreak{}field-\allowbreak{}ring-\allowbreak{}post-\allowbreak{}type}}{\vnbstmt{\begin{array}{@{}l@{}} \forall a, b, c, h. \\ \quad \operatorname{is\text{-}field\text{-}ring}(a) \wedge \\ \quad (h \in \operatorname{hom\text{-}field\text{-}ring}(b, c)) \Rightarrow \\ \quad \quad (\lambda f \in \operatorname{hom\text{-}field\text{-}ring}(a, b).\; \operatorname{compose}(h, f)) \\ \quad \quad \in\; (\operatorname{hom\text{-}field\text{-}ring}(a, b) \to \operatorname{hom\text{-}field\text{-}ring}(a, c)) \end{array}}}
\vnbentry{\texttt{hom-\allowbreak{}field-\allowbreak{}ring-\allowbreak{}pre-\allowbreak{}type}}{\vnbstmt{\begin{array}{@{}l@{}} \forall a, b, c, k. \\ \quad \operatorname{is\text{-}field\text{-}ring}(b) \wedge \\ \quad (k \in \operatorname{hom\text{-}field\text{-}ring}(c, a)) \Rightarrow \\ \quad \quad (\lambda f \in \operatorname{hom\text{-}field\text{-}ring}(a, b).\; \operatorname{compose}(f, k)) \\ \quad \quad \in\; (\operatorname{hom\text{-}field\text{-}ring}(a, b) \to \operatorname{hom\text{-}field\text{-}ring}(c, b)) \end{array}}}
\vnbentry{\texttt{hom-\allowbreak{}field-\allowbreak{}post-\allowbreak{}type}}{\vnbstmt{\begin{array}{@{}l@{}} \forall a, b, c, h. \\ \quad (\operatorname{is\text{-}field}(a) \wedge (h \in \operatorname{hom\text{-}field}(b, c))) \Rightarrow  \\ \quad \quad (\lambda f \in \operatorname{hom\text{-}field}(a, b).\; \operatorname{compose}(h, f)) \\ \quad \quad \in\; (\operatorname{hom\text{-}field}(a, b) \to \operatorname{hom\text{-}field}(a, c)) \end{array}}}
\vnbentry{\texttt{hom-\allowbreak{}field-\allowbreak{}pre-\allowbreak{}type}}{\vnbstmt{\begin{array}{@{}l@{}} \forall a, b, c, k. \\ \quad (\operatorname{is\text{-}field}(b) \wedge (k \in \operatorname{hom\text{-}field}(c, a))) \Rightarrow  \\ \quad \quad (\lambda f \in \operatorname{hom\text{-}field}(a, b).\; \operatorname{compose}(f, k)) \\ \quad \quad \in\; (\operatorname{hom\text{-}field}(a, b) \to \operatorname{hom\text{-}field}(c, b)) \end{array}}}
\vnbentry{\texttt{hom-\allowbreak{}normed-\allowbreak{}ag-\allowbreak{}post-\allowbreak{}type}}{\vnbstmt{\begin{array}{@{}l@{}} \forall a, b, c, h. \\ \quad \operatorname{is\text{-}normed\text{-}ag}(a) \wedge \\ \quad (h \in \operatorname{hom\text{-}normed\text{-}ag}(b, c)) \Rightarrow \\ \quad \quad (\lambda f \in \operatorname{hom\text{-}normed\text{-}ag}(a, b).\; \operatorname{compose}(h, f)) \\ \quad \quad \in\; (\operatorname{hom\text{-}normed\text{-}ag}(a, b) \to \operatorname{hom\text{-}normed\text{-}ag}(a, c)) \end{array}}}
\vnbentry{\texttt{hom-\allowbreak{}normed-\allowbreak{}ag-\allowbreak{}pre-\allowbreak{}type}}{\vnbstmt{\begin{array}{@{}l@{}} \forall a, b, c, k. \\ \quad \operatorname{is\text{-}normed\text{-}ag}(b) \wedge \\ \quad (k \in \operatorname{hom\text{-}normed\text{-}ag}(c, a)) \Rightarrow \\ \quad \quad (\lambda f \in \operatorname{hom\text{-}normed\text{-}ag}(a, b).\; \operatorname{compose}(f, k)) \\ \quad \quad \in\; (\operatorname{hom\text{-}normed\text{-}ag}(a, b) \to \operatorname{hom\text{-}normed\text{-}ag}(c, b)) \end{array}}}
\vnbentry{\texttt{hom-\allowbreak{}normed-\allowbreak{}field-\allowbreak{}post-\allowbreak{}type}}{\vnbstmt{\begin{array}{@{}l@{}} \forall a, b, c, h. \\ \quad \operatorname{is\text{-}normed\text{-}field}(a) \wedge \\ \quad (h \in \operatorname{hom\text{-}normed\text{-}field}(b, c)) \Rightarrow \\ \quad \quad (\lambda f \in \operatorname{hom\text{-}normed\text{-}field}(a, b).\; \operatorname{compose}(h, f)) \\ \quad \quad \in\; (\operatorname{hom\text{-}normed\text{-}field}(a, b) \to \operatorname{hom\text{-}normed\text{-}field}(a, c)) \end{array}}}
\vnbentry{\texttt{hom-\allowbreak{}normed-\allowbreak{}field-\allowbreak{}pre-\allowbreak{}type}}{\vnbstmt{\begin{array}{@{}l@{}} \forall a, b, c, k. \\ \quad \operatorname{is\text{-}normed\text{-}field}(b) \wedge \\ \quad (k \in \operatorname{hom\text{-}normed\text{-}field}(c, a)) \Rightarrow \\ \quad \quad (\lambda f \in \operatorname{hom\text{-}normed\text{-}field}(a, b).\; \operatorname{compose}(f, k)) \\ \quad \quad \in\; (\operatorname{hom\text{-}normed\text{-}field}(a, b) \to \operatorname{hom\text{-}normed\text{-}field}(c, b)) \end{array}}}
\vnbentry{\texttt{hom-\allowbreak{}module-\allowbreak{}post-\allowbreak{}type}}{\vnbstmt{\begin{array}{@{}l@{}} \forall a, b, c, h. \\ \quad (\operatorname{is\text{-}module}(a) \wedge (h \in \operatorname{hom\text{-}module}(b, c))) \Rightarrow  \\ \quad \quad (\lambda f \in \operatorname{hom\text{-}module}(a, b).\; \operatorname{compose}(h, f)) \\ \quad \quad \in\; (\operatorname{hom\text{-}module}(a, b) \to \operatorname{hom\text{-}module}(a, c)) \end{array}}}
\vnbentry{\texttt{hom-\allowbreak{}module-\allowbreak{}pre-\allowbreak{}type}}{\vnbstmt{\begin{array}{@{}l@{}} \forall a, b, c, k. \\ \quad (\operatorname{is\text{-}module}(b) \wedge (k \in \operatorname{hom\text{-}module}(c, a))) \Rightarrow  \\ \quad \quad (\lambda f \in \operatorname{hom\text{-}module}(a, b).\; \operatorname{compose}(f, k)) \\ \quad \quad \in\; (\operatorname{hom\text{-}module}(a, b) \to \operatorname{hom\text{-}module}(c, b)) \end{array}}}
\vnbentry{\texttt{hom-\allowbreak{}vector-\allowbreak{}space-\allowbreak{}post-\allowbreak{}type}}{\vnbstmt{\begin{array}{@{}l@{}} \forall a, b, c, h. \\ \quad \operatorname{is\text{-}vector\text{-}space}(a) \wedge \\ \quad (h \in \operatorname{hom\text{-}vector\text{-}space}(b, c)) \Rightarrow \\ \quad \quad (\lambda f \in \operatorname{hom\text{-}vector\text{-}space}(a, b).\; \operatorname{compose}(h, f)) \\ \quad \quad \in\; (\operatorname{hom\text{-}vector\text{-}space}(a, b) \to \operatorname{hom\text{-}vector\text{-}space}(a, c)) \end{array}}}
\vnbentry{\texttt{hom-\allowbreak{}vector-\allowbreak{}space-\allowbreak{}pre-\allowbreak{}type}}{\vnbstmt{\begin{array}{@{}l@{}} \forall a, b, c, k. \\ \quad \operatorname{is\text{-}vector\text{-}space}(b) \wedge \\ \quad (k \in \operatorname{hom\text{-}vector\text{-}space}(c, a)) \Rightarrow \\ \quad \quad (\lambda f \in \operatorname{hom\text{-}vector\text{-}space}(a, b).\; \operatorname{compose}(f, k)) \\ \quad \quad \in\; (\operatorname{hom\text{-}vector\text{-}space}(a, b) \to \operatorname{hom\text{-}vector\text{-}space}(c, b)) \end{array}}}
\vnbentry{\texttt{hom-\allowbreak{}normed-\allowbreak{}vector-\allowbreak{}space-\allowbreak{}post-\allowbreak{}type}}{\vnbstmt{\begin{array}{@{}l@{}} \forall a, b, c, h. \\ \quad \operatorname{is\text{-}normed\text{-}vector\text{-}space}(a) \wedge \\ \quad (h \in \operatorname{hom\text{-}normed\text{-}vector\text{-}space}(b, c)) \Rightarrow \\ \quad \quad (\lambda f \in \operatorname{hom\text{-}normed\text{-}vector\text{-}space}(a, b).\; \\ \quad \quad \quad \operatorname{compose}(h, f)) \\ \quad \quad \in\; (\operatorname{hom\text{-}normed\text{-}vector\text{-}space}(a, b) \to \operatorname{hom\text{-}normed\text{-}vector\text{-}space}(a, c)) \end{array}}}
\vnbentry{\texttt{hom-\allowbreak{}normed-\allowbreak{}vector-\allowbreak{}space-\allowbreak{}pre-\allowbreak{}type}}{\vnbstmt{\begin{array}{@{}l@{}} \forall a, b, c, k. \\ \quad \operatorname{is\text{-}normed\text{-}vector\text{-}space}(b) \wedge \\ \quad (k \in \operatorname{hom\text{-}normed\text{-}vector\text{-}space}(c, a)) \Rightarrow \\ \quad \quad (\lambda f \in \operatorname{hom\text{-}normed\text{-}vector\text{-}space}(a, b).\; \\ \quad \quad \quad \operatorname{compose}(f, k)) \\ \quad \quad \in\; (\operatorname{hom\text{-}normed\text{-}vector\text{-}space}(a, b) \to \operatorname{hom\text{-}normed\text{-}vector\text{-}space}(c, b)) \end{array}}}
\vnbentry{\texttt{hom-\allowbreak{}complex-\allowbreak{}inner-\allowbreak{}product-\allowbreak{}space-\allowbreak{}post-\allowbreak{}type}}{\vnbstmt{\begin{array}{@{}l@{}} \forall a, b, c, h. \\ \quad \operatorname{is\text{-}complex\text{-}inner\text{-}product\text{-}space}(a) \wedge \\ \quad (h \in \operatorname{hom\text{-}complex\text{-}inner\text{-}product\text{-}space}(b, c)) \Rightarrow \\ \quad \quad (\lambda f \in \operatorname{hom\text{-}complex\text{-}inner\text{-}product\text{-}space}(a, b).\; \\ \quad \quad \quad \operatorname{compose}(h, f)) \\ \quad \quad \in\; (\operatorname{hom\text{-}complex\text{-}inner\text{-}product\text{-}space}(a, b) \to \operatorname{hom\text{-}complex\text{-}inner\text{-}product\text{-}space}(a, c)) \end{array}}}
\vnbentry{\texttt{hom-\allowbreak{}complex-\allowbreak{}inner-\allowbreak{}product-\allowbreak{}space-\allowbreak{}pre-\allowbreak{}type}}{\vnbstmt{\begin{array}{@{}l@{}} \forall a, b, c, k. \\ \quad \operatorname{is\text{-}complex\text{-}inner\text{-}product\text{-}space}(b) \wedge \\ \quad (k \in \operatorname{hom\text{-}complex\text{-}inner\text{-}product\text{-}space}(c, a)) \Rightarrow \\ \quad \quad (\lambda f \in \operatorname{hom\text{-}complex\text{-}inner\text{-}product\text{-}space}(a, b).\; \\ \quad \quad \quad \operatorname{compose}(f, k)) \\ \quad \quad \in\; (\operatorname{hom\text{-}complex\text{-}inner\text{-}product\text{-}space}(a, b) \to \operatorname{hom\text{-}complex\text{-}inner\text{-}product\text{-}space}(c, b)) \end{array}}}
\vnbentry{\texttt{hom-\allowbreak{}metric-\allowbreak{}space-\allowbreak{}post-\allowbreak{}type}}{\vnbstmt{\begin{array}{@{}l@{}} \forall a, b, c, h. \\ \quad \operatorname{is\text{-}metric\text{-}space}(a) \wedge \\ \quad (h \in \operatorname{hom\text{-}metric\text{-}space}(b, c)) \Rightarrow \\ \quad \quad (\lambda f \in \operatorname{hom\text{-}metric\text{-}space}(a, b).\; \operatorname{compose}(h, f)) \\ \quad \quad \in\; (\operatorname{hom\text{-}metric\text{-}space}(a, b) \to \operatorname{hom\text{-}metric\text{-}space}(a, c)) \end{array}}}
\vnbentry{\texttt{hom-\allowbreak{}metric-\allowbreak{}space-\allowbreak{}pre-\allowbreak{}type}}{\vnbstmt{\begin{array}{@{}l@{}} \forall a, b, c, k. \\ \quad \operatorname{is\text{-}metric\text{-}space}(b) \wedge \\ \quad (k \in \operatorname{hom\text{-}metric\text{-}space}(c, a)) \Rightarrow \\ \quad \quad (\lambda f \in \operatorname{hom\text{-}metric\text{-}space}(a, b).\; \operatorname{compose}(f, k)) \\ \quad \quad \in\; (\operatorname{hom\text{-}metric\text{-}space}(a, b) \to \operatorname{hom\text{-}metric\text{-}space}(c, b)) \end{array}}}
\vnbentry{\texttt{hom-\allowbreak{}pseudometric-\allowbreak{}space-\allowbreak{}post-\allowbreak{}type}}{\vnbstmt{\begin{array}{@{}l@{}} \forall a, b, c, h. \\ \quad \operatorname{is\text{-}pseudometric\text{-}space}(a) \wedge \\ \quad (h \in \operatorname{hom\text{-}pseudometric\text{-}space}(b, c)) \Rightarrow \\ \quad \quad (\lambda f \in \operatorname{hom\text{-}pseudometric\text{-}space}(a, b).\; \\ \quad \quad \quad \operatorname{compose}(h, f)) \\ \quad \quad \in\; (\operatorname{hom\text{-}pseudometric\text{-}space}(a, b) \to \operatorname{hom\text{-}pseudometric\text{-}space}(a, c)) \end{array}}}
\vnbentry{\texttt{hom-\allowbreak{}pseudometric-\allowbreak{}space-\allowbreak{}pre-\allowbreak{}type}}{\vnbstmt{\begin{array}{@{}l@{}} \forall a, b, c, k. \\ \quad \operatorname{is\text{-}pseudometric\text{-}space}(b) \wedge \\ \quad (k \in \operatorname{hom\text{-}pseudometric\text{-}space}(c, a)) \Rightarrow \\ \quad \quad (\lambda f \in \operatorname{hom\text{-}pseudometric\text{-}space}(a, b).\; \\ \quad \quad \quad \operatorname{compose}(f, k)) \\ \quad \quad \in\; (\operatorname{hom\text{-}pseudometric\text{-}space}(a, b) \to \operatorname{hom\text{-}pseudometric\text{-}space}(c, b)) \end{array}}}
\vnbentry{\texttt{hom-\allowbreak{}c-\allowbreak{}metric-\allowbreak{}space-\allowbreak{}post-\allowbreak{}type}}{\vnbstmt{\begin{array}{@{}l@{}} \forall a, b, c, h. \\ \quad \operatorname{is\text{-}c\text{-}metric\text{-}space}(a) \wedge \\ \quad (h \in \operatorname{hom\text{-}c\text{-}metric\text{-}space}(b, c)) \Rightarrow \\ \quad \quad (\lambda f \in \operatorname{hom\text{-}c\text{-}metric\text{-}space}(a, b).\; \\ \quad \quad \quad \operatorname{compose}(h, f)) \\ \quad \quad \in\; (\operatorname{hom\text{-}c\text{-}metric\text{-}space}(a, b) \to \operatorname{hom\text{-}c\text{-}metric\text{-}space}(a, c)) \end{array}}}
\vnbentry{\texttt{hom-\allowbreak{}c-\allowbreak{}metric-\allowbreak{}space-\allowbreak{}pre-\allowbreak{}type}}{\vnbstmt{\begin{array}{@{}l@{}} \forall a, b, c, k. \\ \quad \operatorname{is\text{-}c\text{-}metric\text{-}space}(b) \wedge \\ \quad (k \in \operatorname{hom\text{-}c\text{-}metric\text{-}space}(c, a)) \Rightarrow \\ \quad \quad (\lambda f \in \operatorname{hom\text{-}c\text{-}metric\text{-}space}(a, b).\; \\ \quad \quad \quad \operatorname{compose}(f, k)) \\ \quad \quad \in\; (\operatorname{hom\text{-}c\text{-}metric\text{-}space}(a, b) \to \operatorname{hom\text{-}c\text{-}metric\text{-}space}(c, b)) \end{array}}}
\vnbentry{\texttt{hom-\allowbreak{}measurable-\allowbreak{}space-\allowbreak{}post-\allowbreak{}type}}{\vnbstmt{\begin{array}{@{}l@{}} \forall a, b, c, h. \\ \quad \operatorname{is\text{-}measurable\text{-}space}(a) \wedge \\ \quad (h \in \operatorname{hom\text{-}measurable\text{-}space}(b, c)) \Rightarrow \\ \quad \quad (\lambda f \in \operatorname{hom\text{-}measurable\text{-}space}(a, b).\; \\ \quad \quad \quad \operatorname{compose}(h, f)) \\ \quad \quad \in\; (\operatorname{hom\text{-}measurable\text{-}space}(a, b) \to \operatorname{hom\text{-}measurable\text{-}space}(a, c)) \end{array}}}
\vnbentry{\texttt{hom-\allowbreak{}measurable-\allowbreak{}space-\allowbreak{}pre-\allowbreak{}type}}{\vnbstmt{\begin{array}{@{}l@{}} \forall a, b, c, k. \\ \quad \operatorname{is\text{-}measurable\text{-}space}(b) \wedge \\ \quad (k \in \operatorname{hom\text{-}measurable\text{-}space}(c, a)) \Rightarrow \\ \quad \quad (\lambda f \in \operatorname{hom\text{-}measurable\text{-}space}(a, b).\; \\ \quad \quad \quad \operatorname{compose}(f, k)) \\ \quad \quad \in\; (\operatorname{hom\text{-}measurable\text{-}space}(a, b) \to \operatorname{hom\text{-}measurable\text{-}space}(c, b)) \end{array}}}
\vnbentry{\texttt{hom-\allowbreak{}measure-\allowbreak{}space-\allowbreak{}post-\allowbreak{}type}}{\vnbstmt{\begin{array}{@{}l@{}} \forall a, b, c, h. \\ \quad \operatorname{is\text{-}measure\text{-}space}(a) \wedge \\ \quad (h \in \operatorname{hom\text{-}measure\text{-}space}(b, c)) \Rightarrow \\ \quad \quad (\lambda f \in \operatorname{hom\text{-}measure\text{-}space}(a, b).\; \operatorname{compose}(h, f)) \\ \quad \quad \in\; (\operatorname{hom\text{-}measure\text{-}space}(a, b) \to \operatorname{hom\text{-}measure\text{-}space}(a, c)) \end{array}}}
\vnbentry{\texttt{hom-\allowbreak{}measure-\allowbreak{}space-\allowbreak{}pre-\allowbreak{}type}}{\vnbstmt{\begin{array}{@{}l@{}} \forall a, b, c, k. \\ \quad \operatorname{is\text{-}measure\text{-}space}(b) \wedge \\ \quad (k \in \operatorname{hom\text{-}measure\text{-}space}(c, a)) \Rightarrow \\ \quad \quad (\lambda f \in \operatorname{hom\text{-}measure\text{-}space}(a, b).\; \operatorname{compose}(f, k)) \\ \quad \quad \in\; (\operatorname{hom\text{-}measure\text{-}space}(a, b) \to \operatorname{hom\text{-}measure\text{-}space}(c, b)) \end{array}}}
\vnbentry{\texttt{hom-\allowbreak{}top-\allowbreak{}space-\allowbreak{}post-\allowbreak{}type}}{\vnbstmt{\begin{array}{@{}l@{}} \forall a, b, c, h. \\ \quad \operatorname{is\text{-}top\text{-}space}(a) \wedge \\ \quad (h \in \operatorname{hom\text{-}top\text{-}space}(b, c)) \Rightarrow \\ \quad \quad (\lambda f \in \operatorname{hom\text{-}top\text{-}space}(a, b).\; \operatorname{compose}(h, f)) \\ \quad \quad \in\; (\operatorname{hom\text{-}top\text{-}space}(a, b) \to \operatorname{hom\text{-}top\text{-}space}(a, c)) \end{array}}}
\vnbentry{\texttt{hom-\allowbreak{}top-\allowbreak{}space-\allowbreak{}pre-\allowbreak{}type}}{\vnbstmt{\begin{array}{@{}l@{}} \forall a, b, c, k. \\ \quad \operatorname{is\text{-}top\text{-}space}(b) \wedge \\ \quad (k \in \operatorname{hom\text{-}top\text{-}space}(c, a)) \Rightarrow \\ \quad \quad (\lambda f \in \operatorname{hom\text{-}top\text{-}space}(a, b).\; \operatorname{compose}(f, k)) \\ \quad \quad \in\; (\operatorname{hom\text{-}top\text{-}space}(a, b) \to \operatorname{hom\text{-}top\text{-}space}(c, b)) \end{array}}}
\vnbentry{\texttt{hom-\allowbreak{}metrizable-\allowbreak{}top-\allowbreak{}space-\allowbreak{}post-\allowbreak{}type}}{\vnbstmt{\begin{array}{@{}l@{}} \forall a, b, c, h. \\ \quad \operatorname{is\text{-}metrizable\text{-}top\text{-}space}(a) \wedge \\ \quad (h \in \operatorname{hom\text{-}metrizable\text{-}top\text{-}space}(b, c)) \Rightarrow \\ \quad \quad (\lambda f \in \operatorname{hom\text{-}metrizable\text{-}top\text{-}space}(a, b).\; \\ \quad \quad \quad \operatorname{compose}(h, f)) \\ \quad \quad \in\; (\operatorname{hom\text{-}metrizable\text{-}top\text{-}space}(a, b) \to \operatorname{hom\text{-}metrizable\text{-}top\text{-}space}(a, c)) \end{array}}}
\vnbentry{\texttt{hom-\allowbreak{}metrizable-\allowbreak{}top-\allowbreak{}space-\allowbreak{}pre-\allowbreak{}type}}{\vnbstmt{\begin{array}{@{}l@{}} \forall a, b, c, k. \\ \quad \operatorname{is\text{-}metrizable\text{-}top\text{-}space}(b) \wedge \\ \quad (k \in \operatorname{hom\text{-}metrizable\text{-}top\text{-}space}(c, a)) \Rightarrow \\ \quad \quad (\lambda f \in \operatorname{hom\text{-}metrizable\text{-}top\text{-}space}(a, b).\; \\ \quad \quad \quad \operatorname{compose}(f, k)) \\ \quad \quad \in\; (\operatorname{hom\text{-}metrizable\text{-}top\text{-}space}(a, b) \to \operatorname{hom\text{-}metrizable\text{-}top\text{-}space}(c, b)) \end{array}}}
\vnbentry{\texttt{hom-\allowbreak{}ringoid-\allowbreak{}post-\allowbreak{}type}}{\vnbstmt{\begin{array}{@{}l@{}} \forall a, b, c, h. \\ \quad (\operatorname{is\text{-}ringoid}(a) \wedge (h \in \operatorname{hom\text{-}ringoid}(b, c))) \Rightarrow  \\ \quad \quad (\lambda f \in \operatorname{hom\text{-}ringoid}(a, b).\; \operatorname{compose}(h, f)) \\ \quad \quad \in\; (\operatorname{hom\text{-}ringoid}(a, b) \to \operatorname{hom\text{-}ringoid}(a, c)) \end{array}}}
\vnbentry{\texttt{hom-\allowbreak{}ringoid-\allowbreak{}pre-\allowbreak{}type}}{\vnbstmt{\begin{array}{@{}l@{}} \forall a, b, c, k. \\ \quad (\operatorname{is\text{-}ringoid}(b) \wedge (k \in \operatorname{hom\text{-}ringoid}(c, a))) \Rightarrow  \\ \quad \quad (\lambda f \in \operatorname{hom\text{-}ringoid}(a, b).\; \operatorname{compose}(f, k)) \\ \quad \quad \in\; (\operatorname{hom\text{-}ringoid}(a, b) \to \operatorname{hom\text{-}ringoid}(c, b)) \end{array}}}
\vnbentry{\texttt{hom-\allowbreak{}setoid-\allowbreak{}post-\allowbreak{}type}}{\vnbstmt{\begin{array}{@{}l@{}} \forall a, b, c, h. \\ \quad (\operatorname{is\text{-}setoid}(a) \wedge (h \in \operatorname{hom\text{-}setoid}(b, c))) \Rightarrow  \\ \quad \quad (\lambda f \in \operatorname{hom\text{-}setoid}(a, b).\; \operatorname{compose}(h, f)) \\ \quad \quad \in\; (\operatorname{hom\text{-}setoid}(a, b) \to \operatorname{hom\text{-}setoid}(a, c)) \end{array}}}
\vnbentry{\texttt{hom-\allowbreak{}setoid-\allowbreak{}pre-\allowbreak{}type}}{\vnbstmt{\begin{array}{@{}l@{}} \forall a, b, c, k. \\ \quad (\operatorname{is\text{-}setoid}(b) \wedge (k \in \operatorname{hom\text{-}setoid}(c, a))) \Rightarrow  \\ \quad \quad (\lambda f \in \operatorname{hom\text{-}setoid}(a, b).\; \operatorname{compose}(f, k)) \\ \quad \quad \in\; (\operatorname{hom\text{-}setoid}(a, b) \to \operatorname{hom\text{-}setoid}(c, b)) \end{array}}}
\vnbentry{\texttt{post-\allowbreak{}id}}{\vnbstmt{\begin{array}{@{}l@{}} \forall a, b, s, f. \\ \quad ((a \in \mathbf{Set}) \wedge ((b \in \mathbf{Set}) \wedge ((f \in s) \wedge (f \in (a \to b))))) \Rightarrow  \\ \quad \quad (\lambda \mathit{hmpf} \in s.\; \operatorname{compose}(\operatorname{id\text{-}fun}(b), \mathit{hmpf}))(f) \\ \quad \quad =\; f \end{array}}}
\vnbentry{\texttt{pre-\allowbreak{}id}}{\vnbstmt{\begin{array}{@{}l@{}} \forall a, b, s, f. \\ \quad ((a \in \mathbf{Set}) \wedge ((f \in s) \wedge (f \in (a \to b)))) \Rightarrow  \\ \quad \quad (\lambda \mathit{hmpf} \in s.\; \operatorname{compose}(\mathit{hmpf}, \operatorname{id\text{-}fun}(a)))(f) \\ \quad \quad =\; f \end{array}}}
\vnbentry{\texttt{post-\allowbreak{}compose}}{\vnbstmt{\begin{array}{@{}l@{}} \forall a, b, c, u, s, d, f, h, k. \\ \quad (a \in \mathbf{Set}) \wedge \\ \quad (b \in \mathbf{Set}) \wedge \\ \quad (f \in (a \to b)) \wedge \\ \quad (h \in (b \to c)) \wedge \\ \quad (k \in (c \to u)) \wedge \\ \quad (f \in s) \wedge \\ \quad (\operatorname{compose}(h, f) \in d) \Rightarrow \\ \quad \quad (\lambda \mathit{hmpf} \in s.\; \operatorname{compose}(\operatorname{compose}(k, h), \mathit{hmpf}))(f) \\ \quad \quad =\; (\lambda \mathit{hmpf} \in d.\; \operatorname{compose}(k, \mathit{hmpf}))((\lambda \mathit{hmpf} \in s.\; \operatorname{compose}(h, \mathit{hmpf}))(f)) \end{array}}}
\vnbentry{\texttt{pre-\allowbreak{}compose}}{\vnbstmt{\begin{array}{@{}l@{}} \forall a, b, c, u, s, d, f, h, k. \\ \quad (a \in \mathbf{Set}) \wedge \\ \quad (b \in \mathbf{Set}) \wedge \\ \quad (k \in (a \to b)) \wedge \\ \quad (h \in (b \to c)) \wedge \\ \quad (f \in (c \to u)) \wedge \\ \quad (f \in s) \wedge \\ \quad (\operatorname{compose}(f, h) \in d) \Rightarrow \\ \quad \quad (\lambda \mathit{hmpf} \in s.\; \operatorname{compose}(\mathit{hmpf}, \operatorname{compose}(h, k)))(f) \\ \quad \quad =\; (\lambda \mathit{hmpf} \in d.\; \operatorname{compose}(\mathit{hmpf}, k))((\lambda \mathit{hmpf} \in s.\; \operatorname{compose}(\mathit{hmpf}, h))(f)) \end{array}}}

\vnbfile{theorem-library/hom-views.scm}{theorem-library/hom-views.scm}
\noindent{\small hom-views.scm -- a view's action on arrows.}\par
\vnbentry{\texttt{ringoid-\allowbreak{}as-\allowbreak{}ring-\allowbreak{}functorial}}{\vnbstmt{\begin{array}{@{}l@{}} \forall a, b, f. \\ \quad \operatorname{is\text{-}hom\text{-}ringoid}(a, b, f) \Rightarrow  \\ \quad \quad \operatorname{is\text{-}hom\text{-}ring}(\operatorname{ringoid\text{-}as\text{-}ring}(a), \operatorname{ringoid\text{-}as\text{-}ring}(b), f) \end{array}}}

\vnbfile{theorem-library/qq-field-is-field.scm}{theorem-library/qq-field-is-field.scm}
\noindent{\small qq-field-is-field.scm -- IS-FIELD(QQ-FIELD), PROVED.}\par
\vnbentry{\texttt{qq-\allowbreak{}recip-\allowbreak{}nonzero}}{\vnbstmt{\begin{array}{@{}l@{}} \forall a \in \mathbb{Q}. \\ \quad (\neg (a = 0) \Rightarrow \neg (\operatorname{recip}(a) = 0)) \end{array}}}
\vnbentry{\texttt{qq-\allowbreak{}field-\allowbreak{}mul-\allowbreak{}inv-\allowbreak{}type}}{\vnbstmt{\begin{array}{@{}l@{}} (\lambda x \in \operatorname{difference}(\mathbb{Q}, \operatorname{singleton}(0)).\; \\ \quad \operatorname{recip}(x)) \\ \in\; (\operatorname{difference}(\mathbb{Q}, \operatorname{singleton}(0)) \to \operatorname{difference}(\mathbb{Q}, \operatorname{singleton}(0))) \end{array}}}
\vnbentry{\texttt{qq-\allowbreak{}field-\allowbreak{}is-\allowbreak{}field}}{\vnbstmt{\operatorname{is\text{-}field}(\operatorname{qq\text{-}field})}}

\vnbfile{theorem-library/qq-vs-exemplification.scm}{theorem-library/qq-vs-exemplification.scm}
\noindent{\small qq-vs-exemplification.scm -- THE FIRST WITNESS OF A VECTOR SPACE.}\par
\vnbentry{\texttt{qq-\allowbreak{}line-\allowbreak{}is-\allowbreak{}module}}{\vnbstmt{\operatorname{is\text{-}module}(\operatorname{qq\text{-}line})}}
\vnbentry{\texttt{qq-\allowbreak{}line-\allowbreak{}is-\allowbreak{}vector-\allowbreak{}space}}{\vnbstmt{\operatorname{is\text{-}vector\text{-}space}(\operatorname{qq\text{-}line})}}

\vnbfile{theorem-library/makeset-basics.scm}{theorem-library/makeset-basics.scm}
\noindent{\small makeset-basics.scm -- reading membership out of a literal brace set.}\par
\vnbentry{\texttt{ms-\allowbreak{}eq-\allowbreak{}symm}}{\vnbstmt{\forall x, y.\; ((x = y) \Leftrightarrow (y = x))}}
\vnbentry{\texttt{ms-\allowbreak{}index-\allowbreak{}2}}{\vnbstmt{\begin{array}{@{}l@{}} \forall i \in \mathbb{N}. \\ \quad (((1 \leq i) \wedge (i \leq 2)) \Rightarrow ((i = 1) \vee (i = 2))) \end{array}}}
\vnbentry{\texttt{ms-\allowbreak{}nth-\allowbreak{}2}}{\vnbstmt{\begin{array}{@{}l@{}} \forall a, b.; \forall i \in \mathbb{N}. \\ \quad ((1 \leq i) \wedge (i \leq 2)) \Rightarrow  \\ \quad \quad (\operatorname{nth}(i, \langle a, b \rangle) = a) \vee \\ \quad \quad (\operatorname{nth}(i, \langle a, b \rangle) = b) \end{array}}}
\vnbentry{\texttt{makeset2-\allowbreak{}membership}}{\vnbstmt{\begin{array}{@{}l@{}} \forall x, a, b. \\ \quad ((x \in \operatorname{make\text{-}set}(\langle a, b \rangle)) \Leftrightarrow ((x \in \mathbf{Set}) \wedge ((x = a) \vee (x = b)))) \end{array}}}
\vnbentry{\texttt{makeset2-\allowbreak{}split}}{\vnbstmt{\begin{array}{@{}l@{}} \forall a, b \in \mathbf{Set}. \\ \quad (\operatorname{make\text{-}set}(\langle a, b \rangle) = (\{a\} \cup \{b\})) \end{array}}}
\vnbentry{\texttt{card-\allowbreak{}pair}}{\vnbstmt{\begin{array}{@{}l@{}} \forall a, b \in \mathbf{Set}. \\ \quad \neg (a = b) \Rightarrow  \\ \quad \quad (|\operatorname{make\text{-}set}(\langle a, b \rangle)| = 2) \end{array}}}

\vnbfile{theorem-library/empty-cardinality.scm}{theorem-library/empty-cardinality.scm}
\noindent{\small empty-cardinality.scm -- the two BACKWARD rungs for "this set is empty, so its cardinal is 0".}\par
\vnbentry{\texttt{card-\allowbreak{}empty-\allowbreak{}le}}{\vnbstmt{\forall x.\; ((x = \emptyset) \Rightarrow (|x| \leq 0))}}
\vnbentry{\texttt{makeset-\allowbreak{}of-\allowbreak{}empty-\allowbreak{}tuple}}{\vnbstmt{\begin{array}{@{}l@{}} \forall l. \\ \quad (l = \langle  \rangle) \Rightarrow  \\ \quad \quad (\operatorname{make\text{-}set}(l) = \emptyset) \end{array}}}

\vnbfile{theorem-library/tuples-induction.scm}{theorem-library/tuples-induction.scm}
\noindent{\small tuples-induction.scm -- list induction, DERIVED.}\par
\vnbentry{\texttt{tuples-\allowbreak{}induction}}{\vnbstmt{\begin{array}{@{}l@{}} \forall a, c.; \forall l \in \operatorname{tuples}(a). \\ \quad (\langle  \rangle \in c) \wedge \\ \quad \forall x, m. \\ \quad \quad ((x \in a) \wedge ((m \in \operatorname{tuples}(a)) \wedge (m \in c))) \Rightarrow  \\ \quad \quad \quad (\operatorname{cons}(x, m) \in c) \Rightarrow \\ \quad \quad (l \in c) \end{array}}}

\vnbfile{theorem-library/tuple-extensionality.scm}{theorem-library/tuple-extensionality.scm}
\noindent{\small tuple-extensionality.scm -- two tuples with the same length and the same entry at every index are equal; and, on top of it, matrix-entry-extensionality.}\par
\vnbentry{\texttt{nth2-\allowbreak{}pair}}{\vnbstmt{\forall a, b.\; (\operatorname{nth}(2, \langle a, b \rangle) \simeq b)}}
\vnbentry{\texttt{entry-\allowbreak{}unfold}}{\vnbstmt{\begin{array}{@{}l@{}} \forall m, i, j. \\ \quad (\operatorname{entry}(m, i, j) \simeq \operatorname{nth}(j, \operatorname{nth}(i, m))) \end{array}}}
\vnbentry{\texttt{tuple-\allowbreak{}extensionality}}{\vnbstmt{\begin{array}{@{}l@{}} \forall a, l, m. \\ \quad (l \in \operatorname{tuples}(a)) \wedge \\ \quad (m \in \operatorname{tuples}(a)) \wedge \\ \quad (\operatorname{length}(l) = \operatorname{length}(m)) \wedge \\ \quad \forall k \in \mathbb{N}. \\ \quad \quad ((1 \leq k) \wedge (k \leq \operatorname{length}(l))) \Rightarrow  \\ \quad \quad \quad (\operatorname{nth}(k, l) = \operatorname{nth}(k, m)) \Rightarrow \\ \quad \quad (l = m) \end{array}}}
\vnbentry{\texttt{matrix-\allowbreak{}entry-\allowbreak{}extensionality}}{\vnbstmt{\begin{array}{@{}l@{}} \forall m, n, x, p, q. \\ \quad (p \in \operatorname{mat}(m, n, x)) \wedge \\ \quad (q \in \operatorname{mat}(m, n, x)) \wedge \\ \quad \forall i \in \operatorname{interval}(1, m), j \in \operatorname{interval}(1, n). \\ \quad \quad (\operatorname{entry}(p, i, j) = \operatorname{entry}(q, i, j)) \Rightarrow \\ \quad \quad (p = q) \end{array}}}

\vnbfile{theorem-library/tuple-tabulation.scm}{theorem-library/tuple-tabulation.scm}
\noindent{\small tuple-tabulation.scm -- TABULATION: the missing generation principle for tuples and matrices, and the honest replacement for `matof-exists'.}\par
\vnbentry{\texttt{tab-\allowbreak{}succ-\allowbreak{}add}}{\vnbstmt{\forall a, b \in \mathbb{N}.\; (((a + 1) + b) = ((a + b) + 1))}}
\vnbentry{\texttt{tuple-\allowbreak{}tabulation-\allowbreak{}offset}}{\vnbstmt{\begin{array}{@{}l@{}} \forall n \in \mathbb{N}, a, f.; \forall k \in \mathbb{N}. \\ \quad \forall i \in \mathbb{N}.\; (((1 \leq i) \wedge (i \leq n)) \Rightarrow (f((k + i)) \in a)) \Rightarrow  \\ \quad \quad \exists l \in \operatorname{tuples}(a). \\ \quad \quad \quad (\operatorname{length}(l) = n) \wedge \\ \quad \quad \quad \forall i \in \mathbb{N}. \\ \quad \quad \quad \quad ((1 \leq i) \wedge (i \leq n)) \Rightarrow  \\ \quad \quad \quad \quad \quad (\operatorname{nth}(i, l) = f((k + i))) \end{array}}}
\vnbentry{\texttt{row-\allowbreak{}tabulation-\allowbreak{}offset}}{\vnbstmt{\begin{array}{@{}l@{}} \forall n \in \mathbb{N}, a, g, c.; \forall k \in \mathbb{N}. \\ \quad \forall j \in \mathbb{N}.\; (((1 \leq j) \wedge (j \leq n)) \Rightarrow (g(c, (k + j)) \in a)) \Rightarrow  \\ \quad \quad \exists l \in \operatorname{tuples}(a). \\ \quad \quad \quad (\operatorname{length}(l) = n) \wedge \\ \quad \quad \quad \forall j \in \mathbb{N}. \\ \quad \quad \quad \quad ((1 \leq j) \wedge (j \leq n)) \Rightarrow  \\ \quad \quad \quad \quad \quad (\operatorname{nth}(j, l) = g(c, (k + j))) \end{array}}}
\vnbentry{\texttt{tuple-\allowbreak{}tabulation}}{\vnbstmt{\begin{array}{@{}l@{}} \forall n \in \mathbb{N}, a, f. \\ \quad \forall i \in \mathbb{N}.\; (((1 \leq i) \wedge (i \leq n)) \Rightarrow (f(i) \in a)) \Rightarrow  \\ \quad \quad \exists l \in \operatorname{tuples}(a). \\ \quad \quad \quad (\operatorname{length}(l) = n) \wedge \\ \quad \quad \quad \forall i \in \mathbb{N}. \\ \quad \quad \quad \quad (((1 \leq i) \wedge (i \leq n)) \Rightarrow (\operatorname{nth}(i, l) = f(i))) \end{array}}}
\vnbentry{\texttt{row-\allowbreak{}tabulation}}{\vnbstmt{\begin{array}{@{}l@{}} \forall n \in \mathbb{N}, a, g, c. \\ \quad \forall j \in \mathbb{N}.\; (((1 \leq j) \wedge (j \leq n)) \Rightarrow (g(c, j) \in a)) \Rightarrow  \\ \quad \quad \exists l \in \operatorname{tuples}(a). \\ \quad \quad \quad (\operatorname{length}(l) = n) \wedge \\ \quad \quad \quad \forall j \in \mathbb{N}. \\ \quad \quad \quad \quad ((1 \leq j) \wedge (j \leq n)) \Rightarrow  \\ \quad \quad \quad \quad \quad (\operatorname{nth}(j, l) = g(c, j)) \end{array}}}
\vnbentry{\texttt{matrix-\allowbreak{}tabulation}}{\vnbstmt{\begin{array}{@{}l@{}} \forall m \in \mathbb{N}, a, g, n.; \forall k \in \mathbb{N}. \\ \quad (n \in \mathbb{N}) \wedge \\ \quad \forall i, j \in \mathbb{N}. \\ \quad \quad (((1 \leq i) \wedge (i \leq m)) \wedge ((1 \leq j) \wedge (j \leq n))) \Rightarrow  \\ \quad \quad \quad (g((k + i), j) \in a) \Rightarrow \\ \quad \quad \exists p \in \operatorname{tuples}(\operatorname{tuples}(a)). \\ \quad \quad \quad (\operatorname{length}(p) = m) \wedge \\ \quad \quad \quad \forall i \in \mathbb{N}. \\ \quad \quad \quad \quad ((1 \leq i) \wedge (i \leq m)) \Rightarrow  \\ \quad \quad \quad \quad \quad (\operatorname{length}(\operatorname{nth}(i, p)) = n) \wedge \\ \quad \quad \quad \forall i, j \in \mathbb{N}. \\ \quad \quad \quad \quad (((1 \leq i) \wedge (i \leq m)) \wedge ((1 \leq j) \wedge (j \leq n))) \Rightarrow  \\ \quad \quad \quad \quad \quad (\operatorname{nth}(j, \operatorname{nth}(i, p)) = g((k + i), j)) \end{array}}}
\vnbentry{\texttt{mat-\allowbreak{}tabulation-\allowbreak{}exists}}{\vnbstmt{\begin{array}{@{}l@{}} \forall m, n, x, g. \\ \quad (m \in \mathbb{N}) \wedge \\ \quad (n \in \mathbb{N}) \wedge \\ \quad (1 \leq m) \wedge \\ \quad \forall i \in \operatorname{interval}(1, m), j \in \operatorname{interval}(1, n). \\ \quad \quad (g(i, j) \in x) \Rightarrow \\ \quad \quad \exists p \in \operatorname{mat}(m, n, x). \forall i \in \operatorname{interval}(1, m), j \in \operatorname{interval}(1, n). \\ \quad \quad \quad (\operatorname{entry}(p, i, j) = g(i, j)) \end{array}}}
\vnbentry{\texttt{matof-\allowbreak{}exists-\allowbreak{}image}}{\vnbstmt{\begin{array}{@{}l@{}} \forall m, n, g. \\ \quad (m \in \mathbb{N}) \wedge \\ \quad (n \in \mathbb{N}) \wedge \\ \quad (1 \leq m) \wedge \\ \quad \forall i \in \operatorname{interval}(1, m), j \in \operatorname{interval}(1, n). \\ \quad \quad (g(i, j) \in \mathbf{Set}) \Rightarrow \\ \quad \quad \exists p \in \operatorname{mat}(m, n, \operatorname{image}(g, (\operatorname{interval}(1, m) \times \operatorname{interval}(1, n)))). \forall i \in \operatorname{interval}(1, m), j \in \operatorname{interval}(1, n). \\ \quad \quad \quad (\operatorname{entry}(p, i, j) = g(i, j)) \end{array}}}
\vnbentry{\texttt{entry-\allowbreak{}of-\allowbreak{}matof-\allowbreak{}guarded}}{\vnbstmt{\begin{array}{@{}l@{}} \forall m, n, g.; \forall u \in \operatorname{interval}(1, m), v \in \operatorname{interval}(1, n). \\ \quad (m \in \mathbb{N}) \wedge \\ \quad (n \in \mathbb{N}) \wedge \\ \quad (1 \leq m) \wedge \\ \quad \forall i \in \operatorname{interval}(1, m), j \in \operatorname{interval}(1, n). \\ \quad \quad (g(i, j) \in \mathbf{Set}) \Rightarrow \\ \quad \quad (\operatorname{entry}(\operatorname{matof}(m, n, g), u, v) = g(u, v)) \end{array}}}
\vnbentry{\texttt{matof-\allowbreak{}exists}}{\vnbstmt{\begin{array}{@{}l@{}} \forall m, n, g. \\ \quad (m \in \mathbb{N}) \wedge \\ \quad (n \in \mathbb{N}) \wedge \\ \quad \forall a \in \operatorname{interval}(1, m), b \in \operatorname{interval}(1, n). \\ \quad \quad (g(a, b) \in \mathbf{Set}) \Rightarrow \\ \quad \quad \exists p \in \operatorname{mat}(m, n, \operatorname{image}(g, (\operatorname{interval}(1, m) \times \operatorname{interval}(1, n)))). \forall i \in \operatorname{interval}(1, m), j \in \operatorname{interval}(1, n). \\ \quad \quad \quad (\operatorname{entry}(p, i, j) = g(i, j)) \end{array}}}
\vnbentry{\texttt{entry-\allowbreak{}of-\allowbreak{}matof}}{\vnbstmt{\begin{array}{@{}l@{}} \forall m, n, g, i, j. \\ \quad (m \in \mathbb{N}) \wedge \\ \quad (n \in \mathbb{N}) \wedge \\ \quad \forall a \in \operatorname{interval}(1, m), b \in \operatorname{interval}(1, n). \\ \quad \quad (g(a, b) \in \mathbf{Set}) \wedge \\ \quad (i \in \operatorname{interval}(1, m)) \wedge \\ \quad (j \in \operatorname{interval}(1, n)) \Rightarrow \\ \quad \quad (\operatorname{entry}(\operatorname{matof}(m, n, g), i, j) = g(i, j)) \end{array}}}

\vnbfile{theorem-library/rake-border-entry.scm}{theorem-library/rake-border-entry.scm}
\noindent{\small rake-border-entry.scm -- the BORDER block entry read-offs, PROVEN.}\par
\vnbentry{\texttt{border-\allowbreak{}entries-\allowbreak{}defined}}{\vnbstmt{\begin{array}{@{}l@{}} \forall a, b, m, p, q.; \forall c \in \operatorname{interval}(1, (p + 1)), d \in \operatorname{interval}(1, (q + 1)). \\ \quad \operatorname{is\text{-}ring}(a) \wedge \\ \quad (p \in \mathbb{N}) \wedge \\ \quad (q \in \mathbb{N}) \wedge \\ \quad (b \in \operatorname{carr}(a)) \wedge \\ \quad (m \in \operatorname{mat}(p, q, \operatorname{carr}(a))) \Rightarrow \\ \quad \quad (\lambda \langle i, j \rangle \in (\operatorname{interval}(1, (p + 1)) \times \operatorname{interval}(1, (q + 1))).\; \left(\text{if }(i = 1)\text{ then }\left(\text{if }(j = 1)\text{ then }b\text{ else }\operatorname{zero}(a)\right)\text{ else }\left(\text{if }(j = 1)\text{ then }\operatorname{zero}(a)\text{ else }\operatorname{entry}(m, \operatorname{nn\text{-}minus}(i, 1), \operatorname{nn\text{-}minus}(j, 1))\right)\right))(c, d) \\ \quad \quad \in\; \mathbf{Set} \end{array}}}
\vnbentry{\texttt{border-\allowbreak{}entry-\allowbreak{}block2}}{\vnbstmt{\begin{array}{@{}l@{}} \forall a, b, m, p, q, i, j. \\ \quad \operatorname{is\text{-}ring}(a) \wedge \\ \quad (p \in \mathbb{N}) \wedge \\ \quad (q \in \mathbb{N}) \wedge \\ \quad (b \in \operatorname{carr}(a)) \wedge \\ \quad (m \in \operatorname{mat}(p, q, \operatorname{carr}(a))) \wedge \\ \quad (i \in \operatorname{interval}(1, (p + 1))) \wedge \\ \quad \neg (i = 1) \wedge \\ \quad (j \in \operatorname{interval}(1, (q + 1))) \wedge \\ \quad \neg (j = 1) \Rightarrow \\ \quad \quad \operatorname{entry}(\operatorname{border}(a, b, m, p, q), i, j) \\ \quad \quad =\; \operatorname{entry}(m, \operatorname{nn\text{-}minus}(i, 1), \operatorname{nn\text{-}minus}(j, 1)) \end{array}}}
\vnbentry{\texttt{border-\allowbreak{}entry-\allowbreak{}block}}{\vnbstmt{\begin{array}{@{}l@{}} \forall a, b, m, p, q, i, j. \\ \quad \operatorname{is\text{-}ring}(a) \wedge \\ \quad (p \in \mathbb{N}) \wedge \\ \quad (q \in \mathbb{N}) \wedge \\ \quad (b \in \operatorname{carr}(a)) \wedge \\ \quad (m \in \operatorname{mat}(p, q, \operatorname{carr}(a))) \wedge \\ \quad (i \in \operatorname{interval}(1, p)) \wedge \\ \quad (j \in \operatorname{interval}(1, q)) \Rightarrow \\ \quad \quad \operatorname{entry}(\operatorname{border}(a, b, m, p, q), (i + 1), (j + 1)) \\ \quad \quad =\; \operatorname{entry}(m, i, j) \end{array}}}

\vnbfile{theorem-library/rake-border-siblings.scm}{theorem-library/rake-border-siblings.scm}
\noindent{\small rake-border-siblings.scm -- the three remaining BORDER entry read-offs, PROVEN: border-entry-11, border-entry-1j, border-entry-i1.}\par
\vnbentry{\texttt{border-\allowbreak{}entry-\allowbreak{}11}}{\vnbstmt{\begin{array}{@{}l@{}} \forall a, b, m, p, q. \\ \quad \operatorname{is\text{-}ring}(a) \wedge \\ \quad (p \in \mathbb{N}) \wedge \\ \quad (q \in \mathbb{N}) \wedge \\ \quad (b \in \operatorname{carr}(a)) \wedge \\ \quad (m \in \operatorname{mat}(p, q, \operatorname{carr}(a))) \Rightarrow \\ \quad \quad (\operatorname{entry}(\operatorname{border}(a, b, m, p, q), 1, 1) = b) \end{array}}}
\vnbentry{\texttt{border-\allowbreak{}entry-\allowbreak{}1j}}{\vnbstmt{\begin{array}{@{}l@{}} \forall a, b, m, p, q, j. \\ \quad \operatorname{is\text{-}ring}(a) \wedge \\ \quad (p \in \mathbb{N}) \wedge \\ \quad (q \in \mathbb{N}) \wedge \\ \quad (b \in \operatorname{carr}(a)) \wedge \\ \quad (m \in \operatorname{mat}(p, q, \operatorname{carr}(a))) \wedge \\ \quad (j \in \operatorname{interval}(1, (q + 1))) \wedge \\ \quad \neg (j = 1) \Rightarrow \\ \quad \quad (\operatorname{entry}(\operatorname{border}(a, b, m, p, q), 1, j) = \operatorname{zero}(a)) \end{array}}}
\vnbentry{\texttt{border-\allowbreak{}entry-\allowbreak{}i1}}{\vnbstmt{\begin{array}{@{}l@{}} \forall a, b, m, p, q, i. \\ \quad \operatorname{is\text{-}ring}(a) \wedge \\ \quad (p \in \mathbb{N}) \wedge \\ \quad (q \in \mathbb{N}) \wedge \\ \quad (b \in \operatorname{carr}(a)) \wedge \\ \quad (m \in \operatorname{mat}(p, q, \operatorname{carr}(a))) \wedge \\ \quad (i \in \operatorname{interval}(1, (p + 1))) \wedge \\ \quad \neg (i = 1) \Rightarrow \\ \quad \quad (\operatorname{entry}(\operatorname{border}(a, b, m, p, q), i, 1) = \operatorname{zero}(a)) \end{array}}}
\vnbentry{\texttt{succ-\allowbreak{}nn-\allowbreak{}minus-\allowbreak{}1}}{\vnbstmt{\begin{array}{@{}l@{}} \forall i \in \mathbb{N}. \\ \quad ((1 \leq i) \Rightarrow ((\operatorname{nn\text{-}minus}(i, 1) + 1) = i)) \end{array}}}

\vnbfile{theorem-library/matof-in-mat.scm}{theorem-library/matof-in-mat.scm}
\noindent{\small matof-in-mat.scm -- the tabulated matrix MATOF(m,n,g) is an m-by-n matrix over X when every value g(i,j) on the index box [1,m] x [1,n] lies in X.}\par
\vnbentry{\texttt{tuples-\allowbreak{}mono}}{\vnbstmt{\begin{array}{@{}l@{}} \forall a, b.; \forall l \in \operatorname{tuples}(a). \\ \quad (\forall w \in a.\; (w \in b) \Rightarrow (l \in \operatorname{tuples}(b))) \end{array}}}
\vnbentry{\texttt{matrix-\allowbreak{}mono}}{\vnbstmt{\begin{array}{@{}l@{}} \forall a, b.; \forall m \in \operatorname{matrix}(a). \\ \quad (\forall w \in a.\; (w \in b) \Rightarrow (m \in \operatorname{matrix}(b))) \end{array}}}
\vnbentry{\texttt{mat-\allowbreak{}mono}}{\vnbstmt{\begin{array}{@{}l@{}} \forall m, n, a, b.; \forall q \in \operatorname{mat}(m, n, a). \\ \quad (\forall w \in a.\; (w \in b) \Rightarrow (q \in \operatorname{mat}(m, n, b))) \end{array}}}
\vnbentry{\texttt{matof-\allowbreak{}in-\allowbreak{}mat}}{\vnbstmt{\begin{array}{@{}l@{}} \forall m, n, x, g. \\ \quad (m \in \mathbb{N}) \wedge \\ \quad (n \in \mathbb{N}) \wedge \\ \quad \forall i \in \operatorname{interval}(1, m), j \in \operatorname{interval}(1, n). \\ \quad \quad (g(i, j) \in x) \Rightarrow \\ \quad \quad (\operatorname{matof}(m, n, g) \in \operatorname{mat}(m, n, x)) \end{array}}}

\vnbfile{theorem-library/makeset-card-bound.scm}{theorem-library/makeset-card-bound.scm}
\noindent{\small makeset-card-bound.scm -- a tuple of length n has at most n distinct entries.}\par
\vnbentry{\texttt{union-\allowbreak{}comm}}{\vnbstmt{\forall a, b.\; ((a \cup b) = (b \cup a))}}
\vnbentry{\texttt{union-\allowbreak{}singleton-\allowbreak{}absorb}}{\vnbstmt{\forall s.; \forall y \in s.\; ((\{y\} \cup s) = s)}}
\vnbentry{\texttt{card-\allowbreak{}union-\allowbreak{}singleton-\allowbreak{}bound}}{\vnbstmt{\begin{array}{@{}l@{}} \forall s, y \in \mathbf{Set}. \\ \quad (|s| \in \mathbb{N}) \Rightarrow  \\ \quad \quad ((|(\{y\} \cup s)| \in \mathbb{N}) \wedge (|(\{y\} \cup s)| \leq (|s| + 1))) \end{array}}}
\vnbentry{\texttt{makeset-\allowbreak{}card-\allowbreak{}bound-\allowbreak{}strong}}{\vnbstmt{\begin{array}{@{}l@{}} \forall n \in \mathbb{N}, a.; \forall l \in \operatorname{tuples}(a). \\ \quad (\operatorname{length}(l) = n) \Rightarrow  \\ \quad \quad (\operatorname{make\text{-}set}(l) \in \mathbf{Set}) \wedge \\ \quad \quad (|\operatorname{make\text{-}set}(l)| \in \mathbb{N}) \wedge \\ \quad \quad (|\operatorname{make\text{-}set}(l)| \leq n) \end{array}}}
\vnbentry{\texttt{makeset-\allowbreak{}card-\allowbreak{}bound}}{\vnbstmt{\begin{array}{@{}l@{}} \forall n \in \mathbb{N}, a.; \forall l \in \operatorname{tuples}(a). \\ \quad (\operatorname{length}(l) = n) \Rightarrow  \\ \quad \quad (|\operatorname{make\text{-}set}(l)| \leq n) \end{array}}}

\vnbfile{theorem-library/card-subset-nn.scm}{theorem-library/card-subset-nn.scm}
\noindent{\small card-subset-nn.scm -- a subset of a finite set is finite.}\par
\vnbentry{\texttt{card-\allowbreak{}subset-\allowbreak{}nn}}{\vnbstmt{\begin{array}{@{}l@{}} \forall x, s \in \mathbf{Set}. \\ \quad ((|x| \in \mathbb{N}) \wedge \forall z \in s.\; (z \in x)) \Rightarrow  \\ \quad \quad (|s| \in \mathbb{N}) \end{array}}}

\vnbfile{theorem-library/card-image-finite.scm}{theorem-library/card-image-finite.scm}
\noindent{\small card-image-finite.scm -- the image of a finite set is finite.}\par
\vnbentry{\texttt{card-\allowbreak{}image-\allowbreak{}finite}}{\vnbstmt{\begin{array}{@{}l@{}} \forall i.; \forall x \in \mathbf{Set}. \\ \quad ((|x| \in \mathbb{N}) \Rightarrow (|\operatorname{image}(i, x)| \in \mathbb{N})) \end{array}}}
\vnbentry{\texttt{centre-\allowbreak{}set-\allowbreak{}finite-\allowbreak{}guarded}}{\vnbstmt{\begin{array}{@{}l@{}} \forall s, r.; \forall f \in \mathbf{Set}. \\ \quad (|f| \in \mathbb{N}) \Rightarrow  \\ \quad \quad (|\operatorname{centre\text{-}set}(s, r, f)| \in \mathbb{N}) \end{array}}}

\vnbfile{theorem-library/interval-card-in-nn.scm}{theorem-library/interval-card-in-nn.scm}
\noindent{\small interval-card-in-nn.scm -- an interval with a natural upper bound is finite.}\par
\vnbentry{\texttt{interval-\allowbreak{}card-\allowbreak{}in-\allowbreak{}nn}}{\vnbstmt{\begin{array}{@{}l@{}} \forall a.; \forall b \in \mathbb{N}. \\ \quad (|\operatorname{interval}(a, b)| \in \mathbb{N}) \end{array}}}

\vnbfile{theorem-library/rake-finsum-laws.scm}{theorem-library/rake-finsum-laws.scm}
\noindent{\small theorem-library/rake-finsum-laws.scm -- the ADDITIVE algebra of FINSUM, proven.  Rake batch J (2026-09-17), part 1 of 2; part 2 is theorem-library/rake-finsum-laws2.scm (the three distribution laws and the interval back-peel), which loads AFTER this file and cites two of its}\par
\vnbentry{\texttt{enum-\allowbreak{}fam-\allowbreak{}value}}{\vnbstmt{\begin{array}{@{}l@{}} \forall n \in \mathbb{N}, g, u, a.; \forall i \in \mathrm{OS}(n). \\ \quad (\operatorname{enum\text{-}fam}(g, u, a, n)(i) \simeq u(a(i))) \end{array}}}
\vnbentry{\texttt{sum-\allowbreak{}ag-\allowbreak{}type-\allowbreak{}ptwise}}{\vnbstmt{\begin{array}{@{}l@{}} \forall n \in \mathbb{N}, g, a. \\ \quad \operatorname{is\text{-}abelian\text{-}group}(g) \wedge \\ \quad \forall i \in \mathrm{OS}(n).\; (a(i) \in \operatorname{carr}(g)) \Rightarrow \\ \quad \quad (\operatorname{sum\text{-}ag}(g, a, n) \in \operatorname{carr}(g)) \end{array}}}
\vnbentry{\texttt{abelian-\allowbreak{}group-\allowbreak{}opr-\allowbreak{}interchange}}{\vnbstmt{\begin{array}{@{}l@{}} \forall s.; \forall a, b, c, d \in \operatorname{carr}(s). \\ \quad \operatorname{is\text{-}abelian\text{-}group}(s) \Rightarrow  \\ \quad \quad \operatorname{opr}(s)(\operatorname{opr}(s)(a, b), \operatorname{opr}(s)(c, d)) \\ \quad \quad =\; \operatorname{opr}(s)(\operatorname{opr}(s)(a, c), \operatorname{opr}(s)(b, d)) \end{array}}}
\vnbentry{\texttt{finsum-\allowbreak{}congruence-\allowbreak{}q}}{\vnbstmt{\begin{array}{@{}l@{}} \forall a.; \forall s \in \mathbf{Set}, f, g. \\ \quad \operatorname{is\text{-}abelian\text{-}group}(a) \wedge \\ \quad (|s| \in \mathbb{N}) \wedge \\ \quad \forall z \in s.\; (f(z) = g(z)) \Rightarrow \\ \quad \quad (\operatorname{finsum}(a, f, s) \simeq \operatorname{finsum}(a, g, s)) \end{array}}}
\vnbentry{\texttt{finsum-\allowbreak{}congruence-\allowbreak{}guarded}}{\vnbstmt{\begin{array}{@{}l@{}} \forall a.; \forall s \in \mathbf{Set}, f \in (s \to \operatorname{carr}(a)), g. \\ \quad \operatorname{is\text{-}abelian\text{-}group}(a) \wedge \\ \quad (|s| \in \mathbb{N}) \wedge \\ \quad \forall z \in s.\; (f(z) = g(z)) \Rightarrow \\ \quad \quad (\operatorname{finsum}(a, f, s) = \operatorname{finsum}(a, g, s)) \end{array}}}
\vnbentry{\texttt{finsum-\allowbreak{}type-\allowbreak{}ptwise}}{\vnbstmt{\begin{array}{@{}l@{}} \forall g.; \forall s \in \mathbf{Set}, f. \\ \quad \operatorname{is\text{-}abelian\text{-}group}(g) \wedge \\ \quad (|s| \in \mathbb{N}) \wedge \\ \quad \forall z \in s.\; (f(z) \in \operatorname{carr}(g)) \Rightarrow \\ \quad \quad (\operatorname{finsum}(g, f, s) \in \operatorname{carr}(g)) \end{array}}}
\vnbentry{\texttt{finsum-\allowbreak{}congruence}}{\vnbstmt{\begin{array}{@{}l@{}} \forall a.; \forall s \in \mathbf{Set}, f, g. \\ \quad \operatorname{is\text{-}abelian\text{-}group}(a) \wedge \\ \quad (|s| \in \mathbb{N}) \wedge \\ \quad \forall z \in s.\; (f(z) = g(z)) \wedge \\ \quad \forall z \in s.\; (f(z) \in \operatorname{carr}(a)) \Rightarrow \\ \quad \quad (\operatorname{finsum}(a, f, s) = \operatorname{finsum}(a, g, s)) \end{array}}}
\vnbentry{\texttt{sum-\allowbreak{}ag-\allowbreak{}add-\allowbreak{}ind}}{\vnbstmt{\begin{array}{@{}l@{}} \forall n \in \mathbb{N}, g, a, b, c. \\ \quad \operatorname{is\text{-}abelian\text{-}group}(g) \wedge \\ \quad \forall i \in \mathrm{OS}(n).\; (a(i) \in \operatorname{carr}(g)) \wedge \\ \quad \forall i \in \mathrm{OS}(n).\; (b(i) \in \operatorname{carr}(g)) \wedge \\ \quad \forall i \in \mathrm{OS}(n).\; (c(i) \simeq \operatorname{opr}(g)(a(i), b(i))) \Rightarrow \\ \quad \quad \operatorname{sum\text{-}ag}(g, c, n) \\ \quad \quad =\; \operatorname{opr}(g)(\operatorname{sum\text{-}ag}(g, a, n), \operatorname{sum\text{-}ag}(g, b, n)) \end{array}}}
\vnbentry{\texttt{finsum-\allowbreak{}add-\allowbreak{}ag}}{\vnbstmt{\begin{array}{@{}l@{}} \forall g.; \forall s \in \mathbf{Set}, f, h \in (s \to \operatorname{carr}(g)). \\ \quad (\operatorname{is\text{-}abelian\text{-}group}(g) \wedge (|s| \in \mathbb{N})) \Rightarrow  \\ \quad \quad \operatorname{finsum}(g, (\lambda z \in s.\; \operatorname{opr}(g)(f(z), h(z))), s) \\ \quad \quad =\; \operatorname{opr}(g)(\operatorname{finsum}(g, f, s), \operatorname{finsum}(g, h, s)) \end{array}}}
\vnbentry{\texttt{sum-\allowbreak{}ag-\allowbreak{}two-\allowbreak{}support-\allowbreak{}ind}}{\vnbstmt{\begin{array}{@{}l@{}} \forall n \in \mathbb{N}, g, a, p, q. \\ \quad \operatorname{is\text{-}abelian\text{-}group}(g) \wedge \\ \quad (p \in \mathrm{OS}(n)) \wedge \\ \quad (q \in \mathrm{OS}(n)) \wedge \\ \quad \neg (p = q) \wedge \\ \quad (a(p) \in \operatorname{carr}(g)) \wedge \\ \quad (a(q) \in \operatorname{carr}(g)) \wedge \\ \quad \forall i \in \mathrm{OS}(n). \\ \quad \quad ((\neg (i = p) \wedge \neg (i = q)) \Rightarrow (a(i) = \operatorname{iden}(g))) \Rightarrow \\ \quad \quad (\operatorname{sum\text{-}ag}(g, a, n) = \operatorname{opr}(g)(a(p), a(q))) \end{array}}}
\vnbentry{\texttt{finsum-\allowbreak{}two-\allowbreak{}support}}{\vnbstmt{\begin{array}{@{}l@{}} \forall g.; \forall s \in \mathbf{Set}, f \in (s \to \operatorname{carr}(g)), a, b \in s. \\ \quad \operatorname{is\text{-}abelian\text{-}group}(g) \wedge \\ \quad (|s| \in \mathbb{N}) \wedge \\ \quad \neg (a = b) \wedge \\ \quad \forall j \in s. \\ \quad \quad ((\neg (j = a) \wedge \neg (j = b)) \Rightarrow (f(j) = \operatorname{iden}(g))) \Rightarrow \\ \quad \quad (\operatorname{finsum}(g, f, s) = \operatorname{opr}(g)(f(a), f(b))) \end{array}}}

\vnbfile{theorem-library/rake-finsum-laws2.scm}{theorem-library/rake-finsum-laws2.scm}
\noindent{\small theorem-library/rake-finsum-laws2.scm -- rake batch J, part 2: the three DISTRIBUTION laws of FINSUM (a ring element on either side, a module scalar) and the interval back-peel.  Part 1 is theorem-library/rake-finsum-laws.scm, which must load FIRST: this file cites its `enum-fam-value' and}\par
\vnbentry{\texttt{sum-\allowbreak{}ag-\allowbreak{}ring-\allowbreak{}left-\allowbreak{}ind}}{\vnbstmt{\begin{array}{@{}l@{}} \forall n \in \mathbb{N}, g.; \forall r \in \operatorname{carr}(g), a, b. \\ \quad \operatorname{is\text{-}ring}(g) \wedge \\ \quad \forall i \in \mathrm{OS}(n).\; (a(i) \in \operatorname{carr}(g)) \wedge \\ \quad \forall i \in \mathrm{OS}(n).\; (b(i) \simeq \operatorname{mul}(g)(r, a(i))) \Rightarrow \\ \quad \quad \operatorname{sum\text{-}ag}(\operatorname{ring\text{-}additive\text{-}ag}(g), b, n) \\ \quad \quad =\; \operatorname{mul}(g)(r, \operatorname{sum\text{-}ag}(\operatorname{ring\text{-}additive\text{-}ag}(g), a, n)) \end{array}}}
\vnbentry{\texttt{finsum-\allowbreak{}ring-\allowbreak{}distrib-\allowbreak{}left-\allowbreak{}gen}}{\vnbstmt{\begin{array}{@{}l@{}} \forall g.; \forall r \in \operatorname{carr}(g), s \in \mathbf{Set}, f \in (s \to \operatorname{carr}(g)). \\ \quad (\operatorname{is\text{-}ring}(g) \wedge (|s| \in \mathbb{N})) \Rightarrow  \\ \quad \quad \operatorname{mul}(g)(r, \operatorname{finsum}(\operatorname{ring\text{-}additive\text{-}ag}(g), f, s)) \\ \quad \quad =\; \operatorname{finsum}(\operatorname{ring\text{-}additive\text{-}ag}(g), (\lambda z \in s.\; \operatorname{mul}(g)(r, f(z))), s) \end{array}}}
\vnbentry{\texttt{sum-\allowbreak{}ag-\allowbreak{}ring-\allowbreak{}right-\allowbreak{}ind}}{\vnbstmt{\begin{array}{@{}l@{}} \forall n \in \mathbb{N}, g.; \forall r \in \operatorname{carr}(g), a, b. \\ \quad \operatorname{is\text{-}ring}(g) \wedge \\ \quad \forall i \in \mathrm{OS}(n).\; (a(i) \in \operatorname{carr}(g)) \wedge \\ \quad \forall i \in \mathrm{OS}(n).\; (b(i) \simeq \operatorname{mul}(g)(a(i), r)) \Rightarrow \\ \quad \quad \operatorname{sum\text{-}ag}(\operatorname{ring\text{-}additive\text{-}ag}(g), b, n) \\ \quad \quad =\; \operatorname{mul}(g)(\operatorname{sum\text{-}ag}(\operatorname{ring\text{-}additive\text{-}ag}(g), a, n), r) \end{array}}}
\vnbentry{\texttt{finsum-\allowbreak{}ring-\allowbreak{}distrib-\allowbreak{}right-\allowbreak{}gen}}{\vnbstmt{\begin{array}{@{}l@{}} \forall g.; \forall r \in \operatorname{carr}(g), s \in \mathbf{Set}, f \in (s \to \operatorname{carr}(g)). \\ \quad (\operatorname{is\text{-}ring}(g) \wedge (|s| \in \mathbb{N})) \Rightarrow  \\ \quad \quad \operatorname{mul}(g)(\operatorname{finsum}(\operatorname{ring\text{-}additive\text{-}ag}(g), f, s), r) \\ \quad \quad =\; \operatorname{finsum}(\operatorname{ring\text{-}additive\text{-}ag}(g), (\lambda z \in s.\; \operatorname{mul}(g)(f(z), r)), s) \end{array}}}
\vnbentry{\texttt{module-\allowbreak{}act-\allowbreak{}zero-\allowbreak{}vec}}{\vnbstmt{\begin{array}{@{}l@{}} \forall d.; \forall r \in \operatorname{carr}(\operatorname{scal}(d)). \\ \quad \operatorname{is\text{-}module}(d) \Rightarrow  \\ \quad \quad (\operatorname{act}(d)(r, \operatorname{vzero}(d)) = \operatorname{vzero}(d)) \end{array}}}
\vnbentry{\texttt{sum-\allowbreak{}ag-\allowbreak{}act-\allowbreak{}ind}}{\vnbstmt{\begin{array}{@{}l@{}} \forall n \in \mathbb{N}, d.; \forall r \in \operatorname{carr}(\operatorname{scal}(d)), a, b. \\ \quad \operatorname{is\text{-}module}(d) \wedge \\ \quad \forall i \in \mathrm{OS}(n). \\ \quad \quad (a(i) \in \operatorname{carr}(\operatorname{module\text{-}vector\text{-}ag}(d))) \wedge \\ \quad \forall i \in \mathrm{OS}(n).\; (b(i) \simeq \operatorname{act}(d)(r, a(i))) \Rightarrow \\ \quad \quad \operatorname{sum\text{-}ag}(\operatorname{module\text{-}vector\text{-}ag}(d), b, n) \\ \quad \quad =\; \operatorname{act}(d)(r, \operatorname{sum\text{-}ag}(\operatorname{module\text{-}vector\text{-}ag}(d), a, n)) \end{array}}}
\vnbentry{\texttt{finsum-\allowbreak{}act-\allowbreak{}distrib-\allowbreak{}gen}}{\vnbstmt{\begin{array}{@{}l@{}} \forall d.; \forall r \in \operatorname{carr}(\operatorname{scal}(d)), s \in \mathbf{Set}, f \in (s \to \operatorname{carr}(\operatorname{module\text{-}vector\text{-}ag}(d))). \\ \quad (\operatorname{is\text{-}module}(d) \wedge (|s| \in \mathbb{N})) \Rightarrow  \\ \quad \quad \operatorname{act}(d)(r, \operatorname{finsum}(\operatorname{module\text{-}vector\text{-}ag}(d), f, s)) \\ \quad \quad =\; \operatorname{finsum}(\operatorname{module\text{-}vector\text{-}ag}(d), (\lambda z \in s.\; \operatorname{act}(d)(r, f(z))), s) \end{array}}}
\vnbentry{\texttt{interval-\allowbreak{}succ-\allowbreak{}insert}}{\vnbstmt{\begin{array}{@{}l@{}} \forall n \in \mathbb{N}. \\ \quad (\operatorname{interval}(1, (n + 1)) = (\operatorname{interval}(1, n) \cup \{(n + 1)\})) \end{array}}}
\vnbentry{\texttt{finsum-\allowbreak{}interval-\allowbreak{}peel}}{\vnbstmt{\begin{array}{@{}l@{}} \forall g.; \forall n \in \mathbb{N}, f \in (\operatorname{interval}(1, (n + 1)) \to \operatorname{carr}(g)). \\ \quad \operatorname{is\text{-}abelian\text{-}group}(g) \Rightarrow  \\ \quad \quad \operatorname{finsum}(g, f, \operatorname{interval}(1, (n + 1))) \\ \quad \quad =\; \operatorname{opr}(g)(\operatorname{finsum}(g, f, \operatorname{interval}(1, n)), f((n + 1))) \end{array}}}

\vnbfile{theorem-library/rake-finsum-core.scm}{theorem-library/rake-finsum-core.scm}
\noindent{\small theorem-library/rake-finsum-core.scm -- the FINSUM CORE, proven. Rake batch M (2026-09-17).  Nine of the batch's ten leaves; the tenth (finsum-embed) is BLOCKED and its obstacle is stated at the end of this header.}\par
\vnbentry{\texttt{cra-\allowbreak{}is-\allowbreak{}rag}}{\vnbstmt{\begin{array}{@{}l@{}} \forall g. \\ \quad (\operatorname{commutative\text{-}ring\text{-}additive\text{-}ag}(g) \simeq \operatorname{ring\text{-}additive\text{-}ag}(g)) \end{array}}}
\vnbentry{\texttt{finsum-\allowbreak{}ring-\allowbreak{}distrib-\allowbreak{}left}}{\vnbstmt{\begin{array}{@{}l@{}} \forall g.; \forall r \in \operatorname{carr}(g), s \in \mathbf{Set}, f \in (s \to \operatorname{carr}(g)). \\ \quad (\operatorname{is\text{-}commutative\text{-}ring}(g) \wedge (|s| \in \mathbb{N})) \Rightarrow  \\ \quad \quad \operatorname{mul}(g)(r, \operatorname{finsum}(\operatorname{commutative\text{-}ring\text{-}additive\text{-}ag}(g), f, s)) \\ \quad \quad =\; \operatorname{finsum}(\operatorname{commutative\text{-}ring\text{-}additive\text{-}ag}(g), (\lambda z \in s.\; \operatorname{mul}(g)(r, f(z))), \\ \quad \quad \quad s) \end{array}}}
\vnbentry{\texttt{finsum-\allowbreak{}ring-\allowbreak{}distrib-\allowbreak{}right}}{\vnbstmt{\begin{array}{@{}l@{}} \forall g.; \forall r \in \operatorname{carr}(g), s \in \mathbf{Set}, f \in (s \to \operatorname{carr}(g)). \\ \quad (\operatorname{is\text{-}commutative\text{-}ring}(g) \wedge (|s| \in \mathbb{N})) \Rightarrow  \\ \quad \quad \operatorname{mul}(g)(\operatorname{finsum}(\operatorname{commutative\text{-}ring\text{-}additive\text{-}ag}(g), f, s), r) \\ \quad \quad =\; \operatorname{finsum}(\operatorname{commutative\text{-}ring\text{-}additive\text{-}ag}(g), (\lambda z \in s.\; \operatorname{mul}(g)(f(z), r)), \\ \quad \quad \quad s) \end{array}}}
\vnbentry{\texttt{ord-\allowbreak{}segment-\allowbreak{}insert}}{\vnbstmt{\forall n \in \mathbb{N}.\; (\mathrm{OS}((n + 1)) = (\mathrm{OS}(n) \cup \{n\}))}}
\vnbentry{\texttt{finsum-\allowbreak{}singleton}}{\vnbstmt{\begin{array}{@{}l@{}} \forall g.; \forall x \in \mathbf{Set}, f \in (\{x\} \to \operatorname{carr}(g)). \\ \quad \operatorname{is\text{-}abelian\text{-}group}(g) \Rightarrow  \\ \quad \quad (\operatorname{finsum}(g, f, \{x\}) = f(x)) \end{array}}}
\vnbentry{\texttt{finsum-\allowbreak{}ord-\allowbreak{}peel}}{\vnbstmt{\begin{array}{@{}l@{}} \forall g.; \forall n \in \mathbb{N}, f \in (\mathrm{OS}((n + 1)) \to \operatorname{carr}(g)). \\ \quad \operatorname{is\text{-}abelian\text{-}group}(g) \Rightarrow  \\ \quad \quad \operatorname{finsum}(g, f, \mathrm{OS}((n + 1))) \\ \quad \quad =\; \operatorname{opr}(g)(\operatorname{finsum}(g, f, \mathrm{OS}(n)), f(n)) \end{array}}}
\vnbentry{\texttt{comm-\allowbreak{}monoid-\allowbreak{}opr-\allowbreak{}comm}}{\vnbstmt{\begin{array}{@{}l@{}} \forall s.; \forall a, b \in \operatorname{carr}(s). \\ \quad \operatorname{is\text{-}comm\text{-}monoid}(s) \Rightarrow  \\ \quad \quad (\operatorname{opr}(s)(a, b) = \operatorname{opr}(s)(b, a)) \end{array}}}
\vnbentry{\texttt{comm-\allowbreak{}monoid-\allowbreak{}assoc}}{\vnbstmt{\begin{array}{@{}l@{}} \forall s.; \forall a, b, c \in \operatorname{carr}(s). \\ \quad \operatorname{is\text{-}comm\text{-}monoid}(s) \Rightarrow  \\ \quad \quad \operatorname{opr}(s)(\operatorname{opr}(s)(a, b), c) \\ \quad \quad =\; \operatorname{opr}(s)(a, \operatorname{opr}(s)(b, c)) \end{array}}}
\vnbentry{\texttt{comm-\allowbreak{}monoid-\allowbreak{}left-\allowbreak{}id}}{\vnbstmt{\begin{array}{@{}l@{}} \forall s.; \forall a \in \operatorname{carr}(s). \\ \quad \operatorname{is\text{-}comm\text{-}monoid}(s) \Rightarrow  \\ \quad \quad (\operatorname{opr}(s)(\operatorname{iden}(s), a) = a) \end{array}}}
\vnbentry{\texttt{comm-\allowbreak{}monoid-\allowbreak{}opr-\allowbreak{}interchange}}{\vnbstmt{\begin{array}{@{}l@{}} \forall s.; \forall a, b, c, d \in \operatorname{carr}(s). \\ \quad \operatorname{is\text{-}comm\text{-}monoid}(s) \Rightarrow  \\ \quad \quad \operatorname{opr}(s)(\operatorname{opr}(s)(a, b), \operatorname{opr}(s)(c, d)) \\ \quad \quad =\; \operatorname{opr}(s)(\operatorname{opr}(s)(a, c), \operatorname{opr}(s)(b, d)) \end{array}}}
\vnbentry{\texttt{sum-\allowbreak{}ag-\allowbreak{}comm-\allowbreak{}monoid-\allowbreak{}type-\allowbreak{}ptwise}}{\vnbstmt{\begin{array}{@{}l@{}} \forall n \in \mathbb{N}, m, g. \\ \quad \operatorname{is\text{-}comm\text{-}monoid}(m) \wedge \\ \quad \forall i \in \mathrm{OS}(n).\; (g(i) \in \operatorname{carr}(m)) \Rightarrow \\ \quad \quad (\operatorname{sum\text{-}ag}(m, g, n) \in \operatorname{carr}(m)) \end{array}}}
\vnbentry{\texttt{sum-\allowbreak{}ag-\allowbreak{}comm-\allowbreak{}monoid-\allowbreak{}add-\allowbreak{}ind}}{\vnbstmt{\begin{array}{@{}l@{}} \forall n \in \mathbb{N}, m, a, b, c. \\ \quad \operatorname{is\text{-}comm\text{-}monoid}(m) \wedge \\ \quad \forall i \in \mathrm{OS}(n).\; (a(i) \in \operatorname{carr}(m)) \wedge \\ \quad \forall i \in \mathrm{OS}(n).\; (b(i) \in \operatorname{carr}(m)) \wedge \\ \quad \forall i \in \mathrm{OS}(n).\; (c(i) \simeq \operatorname{opr}(m)(a(i), b(i))) \Rightarrow \\ \quad \quad \operatorname{sum\text{-}ag}(m, c, n) \\ \quad \quad =\; \operatorname{opr}(m)(\operatorname{sum\text{-}ag}(m, a, n), \operatorname{sum\text{-}ag}(m, b, n)) \end{array}}}
\vnbentry{\texttt{finsum-\allowbreak{}add}}{\vnbstmt{\begin{array}{@{}l@{}} \forall m.; \forall s \in \mathbf{Set}, f, h \in (s \to \operatorname{carr}(m)). \\ \quad (\operatorname{is\text{-}comm\text{-}monoid}(m) \wedge (|s| \in \mathbb{N})) \Rightarrow  \\ \quad \quad \operatorname{finsum}(m, (\lambda z \in s.\; \operatorname{opr}(m)(f(z), h(z))), s) \\ \quad \quad =\; \operatorname{opr}(m)(\operatorname{finsum}(m, f, s), \operatorname{finsum}(m, h, s)) \end{array}}}
\vnbentry{\texttt{sum-\allowbreak{}ag-\allowbreak{}act-\allowbreak{}collect-\allowbreak{}ind}}{\vnbstmt{\begin{array}{@{}l@{}} \forall n \in \mathbb{N}, d.; \forall x \in \operatorname{vec}(d), c, v. \\ \quad \operatorname{is\text{-}module}(d) \wedge \\ \quad \forall i \in \mathrm{OS}(n). \\ \quad \quad (c(i) \in \operatorname{carr}(\operatorname{scal}(d))) \wedge \\ \quad \forall i \in \mathrm{OS}(n).\; (v(i) \simeq \operatorname{act}(d)(c(i), x)) \Rightarrow \\ \quad \quad \operatorname{sum\text{-}ag}(\operatorname{module\text{-}vector\text{-}ag}(d), v, n) \\ \quad \quad =\; \operatorname{act}(d)(\operatorname{sum\text{-}ag}(\operatorname{ring\text{-}additive\text{-}ag}(\operatorname{scal}(d)), c, n), \\ \quad \quad \quad x) \end{array}}}
\vnbentry{\texttt{finsum-\allowbreak{}act-\allowbreak{}collect-\allowbreak{}gen}}{\vnbstmt{\begin{array}{@{}l@{}} \forall d.; \forall x \in \operatorname{vec}(d), s \in \mathbf{Set}, c \in (s \to \operatorname{carr}(\operatorname{scal}(d))). \\ \quad (\operatorname{is\text{-}module}(d) \wedge (|s| \in \mathbb{N})) \Rightarrow  \\ \quad \quad \operatorname{act}(d)(\operatorname{finsum}(\operatorname{ring\text{-}additive\text{-}ag}(\operatorname{scal}(d)), c, s), x) \\ \quad \quad =\; \operatorname{finsum}(\operatorname{module\text{-}vector\text{-}ag}(d), (\lambda z \in s.\; \operatorname{act}(d)(c(z), x)), s) \end{array}}}
\vnbentry{\texttt{sum-\allowbreak{}ag-\allowbreak{}fubini-\allowbreak{}ind}}{\vnbstmt{\begin{array}{@{}l@{}} \forall m \in \mathbb{N}, g.; \forall n \in \mathbb{N}, f, \mathit{gg}, s, \mathit{cs}. \\ \quad \operatorname{is\text{-}abelian\text{-}group}(g) \wedge \\ \quad \forall i \in \mathrm{OS}(m), j \in \mathrm{OS}(n). \\ \quad \quad (f(i)(j) \in \operatorname{carr}(g)) \wedge \\ \quad \forall i \in \mathrm{OS}(m), j \in \mathrm{OS}(n). \\ \quad \quad (\operatorname{gg}(j)(i) \simeq f(i)(j)) \wedge \\ \quad \forall i \in \mathrm{OS}(m). \\ \quad \quad (s(i) \simeq \operatorname{sum\text{-}ag}(g, f(i), n)) \wedge \\ \quad \forall j \in \mathrm{OS}(n). \\ \quad \quad (\operatorname{cs}(j) \simeq \operatorname{sum\text{-}ag}(g, \operatorname{gg}(j), m)) \Rightarrow \\ \quad \quad (\operatorname{sum\text{-}ag}(g, s, m) = \operatorname{sum\text{-}ag}(g, \mathit{cs}, n)) \end{array}}}
\vnbentry{\texttt{finsum-\allowbreak{}fubini}}{\vnbstmt{\begin{array}{@{}l@{}} \forall g.; \forall x, y \in \mathbf{Set}, f \in ((x \times y) \to \operatorname{carr}(g)). \\ \quad \operatorname{is\text{-}abelian\text{-}group}(g) \wedge \\ \quad (|x| \in \mathbb{N}) \wedge \\ \quad (|y| \in \mathbb{N}) \Rightarrow \\ \quad \quad \operatorname{finsum}(g, (\lambda i \in x.\; \operatorname{finsum}(g, (\lambda j \in y.\; f(\langle i, j \rangle)), y)), \\ \quad \quad \quad x) \\ \quad \quad =\; \operatorname{finsum}(g, (\lambda j \in y.\; \operatorname{finsum}(g, (\lambda i \in x.\; f(\langle i, j \rangle)), x)), \\ \quad \quad \quad y) \end{array}}}
\vnbentry{\texttt{ras-\allowbreak{}inv}}{\vnbstmt{\begin{array}{@{}l@{}} \forall a. \\ \quad \operatorname{is\text{-}ring}(a) \Rightarrow  \\ \quad \quad (\operatorname{inv}(\operatorname{ring\text{-}additive\text{-}ag}(a)) = \operatorname{neg}(a)) \end{array}}}
\vnbentry{\texttt{commutative-\allowbreak{}ring-\allowbreak{}mul-\allowbreak{}comm}}{\vnbstmt{\begin{array}{@{}l@{}} \forall s.; \forall a, b \in \operatorname{carr}(s). \\ \quad \operatorname{is\text{-}commutative\text{-}ring}(s) \Rightarrow  \\ \quad \quad (\operatorname{mul}(s)(a, b) = \operatorname{mul}(s)(b, a)) \end{array}}}
\vnbentry{\texttt{ring-\allowbreak{}scalar-\allowbreak{}zz-\allowbreak{}nn}}{\vnbstmt{\begin{array}{@{}l@{}} \forall k \in \mathbb{N}, g.; \forall r, a \in \operatorname{carr}(g). \\ \quad \operatorname{is\text{-}commutative\text{-}ring}(g) \Rightarrow  \\ \quad \quad \operatorname{mul}(g)(r, \operatorname{zz\text{-}act}(\operatorname{ring\text{-}additive\text{-}ag}(g), k, a)) \\ \quad \quad =\; \operatorname{zz\text{-}act}(\operatorname{ring\text{-}additive\text{-}ag}(g), k, \operatorname{mul}(g)(r, a)) \end{array}}}
\vnbentry{\texttt{finsum-\allowbreak{}ring-\allowbreak{}scalar-\allowbreak{}zz}}{\vnbstmt{\begin{array}{@{}l@{}} \forall g.; \forall r \in \operatorname{carr}(g), c \in \mathbb{Z}, a \in \operatorname{carr}(g). \\ \quad \operatorname{is\text{-}commutative\text{-}ring}(g) \Rightarrow  \\ \quad \quad \operatorname{mul}(g)(r, \\ \quad \quad \quad \operatorname{zz\text{-}act}(\operatorname{commutative\text{-}ring\text{-}additive\text{-}ag}(g), c, a)) \\ \quad \quad =\; \operatorname{zz\text{-}act}(\operatorname{commutative\text{-}ring\text{-}additive\text{-}ag}(g), c, \operatorname{mul}(g)(r, a)) \end{array}}}

\vnbfile{theorem-library/rake-finsum-cm-ptwise.scm}{theorem-library/rake-finsum-cm-ptwise.scm}
\noindent{\small theorem-library/rake-finsum-cm-ptwise.scm -- finsum-comm-monoid-type-ptwise, the COMMUTATIVE-MONOID twin of `finsum-type-ptwise'.  Rake batch 5c-K.}\par
\vnbentry{\texttt{finsum-\allowbreak{}comm-\allowbreak{}monoid-\allowbreak{}type-\allowbreak{}ptwise}}{\vnbstmt{\begin{array}{@{}l@{}} \forall m.; \forall s \in \mathbf{Set}, f. \\ \quad \operatorname{is\text{-}comm\text{-}monoid}(m) \wedge \\ \quad (|s| \in \mathbb{N}) \wedge \\ \quad \forall z \in s.\; (f(z) \in \operatorname{carr}(m)) \Rightarrow \\ \quad \quad (\operatorname{finsum}(m, f, s) \in \operatorname{carr}(m)) \end{array}}}

\vnbfile{theorem-library/rake-esum-defined.scm}{theorem-library/rake-esum-defined.scm}
\noindent{\small theorem-library/rake-esum-defined.scm -- ESUM DEFINED as the ESUP of the set of finite partial sums, and its three characterising axioms PROVEN from the definition.  Rake batch 6-C, the ESUP half of which is theorem-library/rake-esup-defined.scm (agent 6-B).  Helper prefix `esm-'.}\par
\vnbentry{\texttt{esum-\allowbreak{}partial-\allowbreak{}sum-\allowbreak{}type}}{\vnbstmt{\begin{array}{@{}l@{}} \forall f, s. \\ \quad (f \in (\operatorname{dom}(f) \to \operatorname{rr\text{-}pos\text{-}star})) \wedge \\ \quad (s \in \mathbf{Set}) \wedge \\ \quad (|s| \in \mathbb{N}) \wedge \\ \quad (s \subseteq \operatorname{dom}(f)) \Rightarrow \\ \quad \quad \operatorname{finsum}(\operatorname{rr\text{-}pos\text{-}star\text{-}add\text{-}monoid}, f, s) \\ \quad \quad \in\; \operatorname{rr\text{-}pos\text{-}star} \end{array}}}
\vnbentry{\texttt{esum-\allowbreak{}in}}{\vnbstmt{\begin{array}{@{}l@{}} \forall f \in (\operatorname{dom}(f) \to \operatorname{rr\text{-}pos\text{-}star}). \\ \quad (\operatorname{esum}(f) \in \operatorname{rr\text{-}pos\text{-}star}) \end{array}}}
\vnbentry{\texttt{esum-\allowbreak{}upper}}{\vnbstmt{\begin{array}{@{}l@{}} \forall f \in (\operatorname{dom}(f) \to \operatorname{rr\text{-}pos\text{-}star}), s \in \mathbf{Set}. \\ \quad ((|s| \in \mathbb{N}) \wedge (s \subseteq \operatorname{dom}(f))) \Rightarrow  \\ \quad \quad \operatorname{finsum}(\operatorname{rr\text{-}pos\text{-}star\text{-}add\text{-}monoid}, f, s) \\ \quad \quad \leq\; \operatorname{esum}(f) \end{array}}}
\vnbentry{\texttt{esum-\allowbreak{}least}}{\vnbstmt{\begin{array}{@{}l@{}} \forall f \in (\operatorname{dom}(f) \to \operatorname{rr\text{-}pos\text{-}star}), b \in \operatorname{rr\text{-}pos\text{-}star}. \\ \quad \forall s \in \mathbf{Set}. \\ \quad \quad ((|s| \in \mathbb{N}) \wedge (s \subseteq \operatorname{dom}(f))) \Rightarrow  \\ \quad \quad \quad \operatorname{finsum}(\operatorname{rr\text{-}pos\text{-}star\text{-}add\text{-}monoid}, f, s) \\ \quad \quad \quad \leq\; b \Rightarrow \\ \quad \quad (\operatorname{esum}(f) \leq b) \end{array}}}

\vnbfile{theorem-library/rake-esum-finite.scm}{theorem-library/rake-esum-finite.scm}
\noindent{\small rake-esum-finite.scm -- ESUM: the finiteness dichotomy.}\par
\vnbentry{\texttt{esum-\allowbreak{}bounded-\allowbreak{}implies-\allowbreak{}finite}}{\vnbstmt{\begin{array}{@{}l@{}} \forall f \in (\operatorname{dom}(f) \to \operatorname{rr\text{-}pos\text{-}star}). \\ \quad \exists m \in \mathbb{R}. \forall s \in \mathbf{Set}. \\ \quad \quad ((|s| \in \mathbb{N}) \wedge (s \subseteq \operatorname{dom}(f))) \Rightarrow  \\ \quad \quad \quad \operatorname{finsum}(\operatorname{rr\text{-}pos\text{-}star\text{-}add\text{-}monoid}, f, s) \\ \quad \quad \quad \leq\; m \Rightarrow \\ \quad \quad (\operatorname{esum}(f) \in \mathbb{R}) \end{array}}}
\vnbentry{\texttt{esum-\allowbreak{}finite-\allowbreak{}iff-\allowbreak{}bounded}}{\vnbstmt{\begin{array}{@{}l@{}} \forall f \in (\operatorname{dom}(f) \to \operatorname{rr\text{-}pos\text{-}star}). \\ \quad ((\operatorname{esum}(f) \in \mathbb{R}) \Leftrightarrow \exists m \in \mathbb{R}.\; \forall s \in \mathbf{Set}.\; (((|s| \in \mathbb{N}) \wedge (s \subseteq \operatorname{dom}(f))) \Rightarrow (\operatorname{finsum}(\operatorname{rr\text{-}pos\text{-}star\text{-}add\text{-}monoid}, f, s) \leq m))) \end{array}}}

\vnbfile{theorem-library/rake-finsum-fubini-c.scm}{theorem-library/rake-finsum-fubini-c.scm}
\noindent{\small rake-finsum-fubini-c.scm -- finsum-fubini-c, moved out of rake-finsum-laws.scm on 2026-09-17 (rake batch M) because it cites finsum-fubini, now PROVEN in rake-finsum-core.scm, which loads after rake-finsum-laws.  Window: after rake-finsum-core, before finsum-fiber (its earliest citer).  Helper prefix rkj-}\par
\vnbentry{\texttt{finsum-\allowbreak{}fubini-\allowbreak{}c}}{\vnbstmt{\begin{array}{@{}l@{}} \forall g.; \forall x, y \in \mathbf{Set}, f \in ((x \times y) \to \operatorname{carr}(g)). \\ \quad \operatorname{is\text{-}abelian\text{-}group}(g) \wedge \\ \quad (|x| \in \mathbb{N}) \wedge \\ \quad (|y| \in \mathbb{N}) \Rightarrow \\ \quad \quad \operatorname{finsum}(g, (\lambda i \in x.\; \operatorname{finsum}(g, (\lambda j \in y.\; f(\langle i, j \rangle)), y)), \\ \quad \quad \quad x) \\ \quad \quad =\; \operatorname{finsum}(g, (\lambda j \in y.\; \operatorname{finsum}(g, (\lambda i \in x.\; f(\langle i, j \rangle)), x)), \\ \quad \quad \quad y) \end{array}}}

\vnbfile{theorem-library/mat-typing-bundle.scm}{theorem-library/mat-typing-bundle.scm}
\noindent{\small mat-typing-bundle.scm -- three MATOF typings, PROVEN from the definitions.}\par
\vnbentry{\texttt{mat-\allowbreak{}colcount-\allowbreak{}transfer}}{\vnbstmt{\begin{array}{@{}l@{}} \forall n, k, y, q, m, x, r. \\ \quad (q \in \operatorname{mat}(n, k, y)) \wedge \\ \quad ((n = 0) \Rightarrow ((m = 0) \vee (k = 0))) \wedge \\ \quad (r \in \operatorname{mat}(m, \operatorname{nth}(2, \operatorname{size}(q)), x)) \Rightarrow \\ \quad \quad (r \in \operatorname{mat}(m, k, x)) \end{array}}}
\vnbentry{\texttt{identmat-\allowbreak{}type}}{\vnbstmt{\begin{array}{@{}l@{}} \forall a.; \forall n \in \mathbb{N}. \\ \quad \operatorname{is\text{-}ring}(a) \Rightarrow  \\ \quad \quad (\operatorname{identmat}(a, n) \in \operatorname{mat}(n, n, \operatorname{carr}(a))) \end{array}}}
\vnbentry{\texttt{elem-\allowbreak{}g-\allowbreak{}type}}{\vnbstmt{\begin{array}{@{}l@{}} \forall a, n, r, k, l. \\ \quad (\operatorname{is\text{-}ring}(a) \wedge ((n \in \mathbb{N}) \wedge (r \in \operatorname{carr}(a)))) \Rightarrow  \\ \quad \quad (\operatorname{elem\text{-}g}(a, n, r, k, l) \in \operatorname{mat}(n, n, \operatorname{carr}(a))) \end{array}}}
\vnbentry{\texttt{zeromat-\allowbreak{}type}}{\vnbstmt{\begin{array}{@{}l@{}} \forall a, m, n. \\ \quad (\operatorname{is\text{-}ring}(a) \wedge ((m \in \mathbb{N}) \wedge (n \in \mathbb{N}))) \Rightarrow  \\ \quad \quad (\operatorname{zeromat}(a, m, n) \in \operatorname{mat}(m, n, \operatorname{carr}(a))) \end{array}}}
\vnbentry{\texttt{unitrow-\allowbreak{}type}}{\vnbstmt{\begin{array}{@{}l@{}} \forall a, n, i. \\ \quad (\operatorname{is\text{-}ring}(a) \wedge (n \in \mathbb{N})) \Rightarrow  \\ \quad \quad (\operatorname{unitrow}(a, n, i) \in \operatorname{mat}(1, n, \operatorname{carr}(a))) \end{array}}}
\vnbentry{\texttt{elem-\allowbreak{}h-\allowbreak{}type}}{\vnbstmt{\begin{array}{@{}l@{}} \forall a, n, r, k. \\ \quad (\operatorname{is\text{-}ring}(a) \wedge ((n \in \mathbb{N}) \wedge (r \in \operatorname{carr}(a)))) \Rightarrow  \\ \quad \quad (\operatorname{elem\text{-}h}(a, n, r, k) \in \operatorname{mat}(n, n, \operatorname{carr}(a))) \end{array}}}
\vnbentry{\texttt{matmul-\allowbreak{}type}}{\vnbstmt{\begin{array}{@{}l@{}} \forall a, m, n, k, p, q. \\ \quad \operatorname{is\text{-}ring}(a) \wedge \\ \quad (p \in \operatorname{mat}(m, n, \operatorname{carr}(a))) \wedge \\ \quad (q \in \operatorname{mat}(n, k, \operatorname{carr}(a))) \wedge \\ \quad ((n = 0) \Rightarrow ((m = 0) \vee (k = 0))) \Rightarrow \\ \quad \quad (\operatorname{matmul}(a, p, q) \in \operatorname{mat}(m, k, \operatorname{carr}(a))) \end{array}}}

\vnbfile{theorem-library/rake-mat-typing.scm}{theorem-library/rake-mat-typing.scm}
\noindent{\small rake-mat-typing.scm -- the ELEVEN matrix typing / sethood supports, PROVEN.}\par
\vnbentry{\texttt{mat-\allowbreak{}is-\allowbreak{}set}}{\vnbstmt{\begin{array}{@{}l@{}} \forall x \in \mathbf{Set}, m, n. \\ \quad (\operatorname{mat}(m, n, x) \in \mathbf{Set}) \end{array}}}
\vnbentry{\texttt{block-\allowbreak{}type}}{\vnbstmt{\begin{array}{@{}l@{}} \forall m, n, x, p, k, l. \\ \quad (k \in \mathbb{N}) \wedge \\ \quad (l \in \mathbb{N}) \wedge \\ \quad (p \in \operatorname{mat}(m, n, x)) \wedge \\ \quad (k \leq m) \wedge \\ \quad (l \leq n) \Rightarrow \\ \quad \quad (\operatorname{block}(p, k, l) \in \operatorname{mat}(k, l, x)) \end{array}}}
\vnbentry{\texttt{snoc-\allowbreak{}col-\allowbreak{}type}}{\vnbstmt{\begin{array}{@{}l@{}} \forall x, n, w, v. \\ \quad ((n \in \mathbb{N}) \wedge ((w \in \operatorname{mat}(n, 1, x)) \wedge (v \in x))) \Rightarrow  \\ \quad \quad (\operatorname{snoc\text{-}col}(w, n, v) \in \operatorname{mat}((n + 1), 1, x)) \end{array}}}
\vnbentry{\texttt{snoc-\allowbreak{}row-\allowbreak{}type}}{\vnbstmt{\begin{array}{@{}l@{}} \forall x, n, c, r. \\ \quad ((n \in \mathbb{N}) \wedge ((c \in \operatorname{mat}(1, n, x)) \wedge (r \in x))) \Rightarrow  \\ \quad \quad (\operatorname{snoc\text{-}row}(c, n, r) \in \operatorname{mat}(1, (n + 1), x)) \end{array}}}
\vnbentry{\texttt{matadd-\allowbreak{}type}}{\vnbstmt{\begin{array}{@{}l@{}} \forall a, m, n, p, q. \\ \quad \operatorname{is\text{-}ring}(a) \wedge \\ \quad (p \in \operatorname{mat}(m, n, \operatorname{carr}(a))) \wedge \\ \quad (q \in \operatorname{mat}(m, n, \operatorname{carr}(a))) \Rightarrow \\ \quad \quad (\operatorname{matadd}(a, p, q) \in \operatorname{mat}(m, n, \operatorname{carr}(a))) \end{array}}}
\vnbentry{\texttt{matscale-\allowbreak{}type}}{\vnbstmt{\begin{array}{@{}l@{}} \forall a, m, n, r, p. \\ \quad \operatorname{is\text{-}ring}(a) \wedge \\ \quad (r \in \operatorname{carr}(a)) \wedge \\ \quad (p \in \operatorname{mat}(m, n, \operatorname{carr}(a))) \Rightarrow \\ \quad \quad (\operatorname{matscale}(a, r, p) \in \operatorname{mat}(m, n, \operatorname{carr}(a))) \end{array}}}
\vnbentry{\texttt{matneg-\allowbreak{}type}}{\vnbstmt{\begin{array}{@{}l@{}} \forall a, m, n.; \forall p \in \operatorname{mat}(m, n, \operatorname{carr}(a)). \\ \quad \operatorname{is\text{-}ring}(a) \Rightarrow  \\ \quad \quad (\operatorname{matneg}(a, p) \in \operatorname{mat}(m, n, \operatorname{carr}(a))) \end{array}}}
\vnbentry{\texttt{submat-\allowbreak{}type}}{\vnbstmt{\begin{array}{@{}l@{}} \forall a, p, q, s. \\ \quad (p \in \mathbb{N}) \wedge \\ \quad (q \in \mathbb{N}) \wedge \\ \quad (s \in \operatorname{mat}((p + 1), (q + 1), \operatorname{carr}(a))) \Rightarrow \\ \quad \quad (\operatorname{submat}(s, p, q) \in \operatorname{mat}(p, q, \operatorname{carr}(a))) \end{array}}}
\vnbentry{\texttt{border-\allowbreak{}type}}{\vnbstmt{\begin{array}{@{}l@{}} \forall a, b, m, p, q. \\ \quad \operatorname{is\text{-}ring}(a) \wedge \\ \quad (p \in \mathbb{N}) \wedge \\ \quad (q \in \mathbb{N}) \wedge \\ \quad (b \in \operatorname{carr}(a)) \wedge \\ \quad (m \in \operatorname{mat}(p, q, \operatorname{carr}(a))) \Rightarrow \\ \quad \quad \operatorname{border}(a, b, m, p, q) \\ \quad \quad \in\; \operatorname{mat}((p + 1), (q + 1), \operatorname{carr}(a)) \end{array}}}
\vnbentry{\texttt{minor-\allowbreak{}type}}{\vnbstmt{\begin{array}{@{}l@{}} \forall r, s, p, q, n. \\ \quad \operatorname{is\text{-}ring}(r) \wedge \\ \quad (p \in \mathbb{N}) \wedge \\ \quad (q \in \mathbb{N}) \wedge \\ \quad (n \in \mathbb{N}) \wedge \\ \quad (s \in \operatorname{mat}((n + 1), (n + 1), \operatorname{carr}(r))) \Rightarrow \\ \quad \quad (\operatorname{minor}(s, p, q, n) \in \operatorname{mat}(n, n, \operatorname{carr}(r))) \end{array}}}
\vnbentry{\texttt{rmm-\allowbreak{}carr}}{\vnbstmt{\begin{array}{@{}l@{}} \forall r. \\ \quad \operatorname{is\text{-}ring}(r) \Rightarrow  \\ \quad \quad \operatorname{carr}(\operatorname{ring\text{-}multiplicative\text{-}monoid}(r)) \\ \quad \quad =\; \operatorname{carr}(r) \end{array}}}
\vnbentry{\texttt{det-\allowbreak{}in-\allowbreak{}carrier-\allowbreak{}ind}}{\vnbstmt{\begin{array}{@{}l@{}} \forall r.; \forall n \in \mathbb{N}, a \in \operatorname{mat}(n, n, \operatorname{carr}(r)). \\ \quad \operatorname{is\text{-}ring}(r) \Rightarrow  \\ \quad \quad (\operatorname{det}(r, n, a) \in \operatorname{carr}(r)) \end{array}}}
\vnbentry{\texttt{det-\allowbreak{}in-\allowbreak{}carrier}}{\vnbstmt{\begin{array}{@{}l@{}} \forall r, n, a. \\ \quad \operatorname{is\text{-}ring}(r) \wedge \\ \quad (n \in \mathbb{N}) \wedge \\ \quad (a \in \operatorname{mat}(n, n, \operatorname{carr}(r))) \Rightarrow \\ \quad \quad (\operatorname{det}(r, n, a) \in \operatorname{carr}(r)) \end{array}}}

\vnbfile{theorem-library/rake-det-small.scm}{theorem-library/rake-det-small.scm}
\noindent{\small rake-det-small.scm -- the small determinants, PROVEN (rake batch 5b-C, 2026-09-18).}\par
\vnbentry{\texttt{interval-\allowbreak{}1-\allowbreak{}1}}{\vnbstmt{(\operatorname{interval}(1, 1) = \{1\})}}
\vnbentry{\texttt{rmm-\allowbreak{}op}}{\vnbstmt{\begin{array}{@{}l@{}} \forall r. \\ \quad \operatorname{is\text{-}ring}(r) \Rightarrow  \\ \quad \quad \operatorname{opr}(\operatorname{ring\text{-}multiplicative\text{-}monoid}(r)) \\ \quad \quad =\; \operatorname{mul}(r) \end{array}}}
\vnbentry{\texttt{rmm-\allowbreak{}iden}}{\vnbstmt{\begin{array}{@{}l@{}} \forall r. \\ \quad \operatorname{is\text{-}ring}(r) \Rightarrow  \\ \quad \quad \operatorname{iden}(\operatorname{ring\text{-}multiplicative\text{-}monoid}(r)) \\ \quad \quad =\; \operatorname{one}(r) \end{array}}}
\vnbentry{\texttt{ring-\allowbreak{}mpow-\allowbreak{}one}}{\vnbstmt{\begin{array}{@{}l@{}} \forall r.; \forall x \in \operatorname{carr}(r). \\ \quad \operatorname{is\text{-}ring}(r) \Rightarrow  \\ \quad \quad (\operatorname{mpow}(\operatorname{ring\text{-}multiplicative\text{-}monoid}(r), x, 1) = x) \end{array}}}
\vnbentry{\texttt{minor-\allowbreak{}entry}}{\vnbstmt{\begin{array}{@{}l@{}} \forall r, s, p, q, n, u, v. \\ \quad \operatorname{is\text{-}ring}(r) \wedge \\ \quad (n \in \mathbb{N}) \wedge \\ \quad (s \in \operatorname{mat}((n + 1), (n + 1), \operatorname{carr}(r))) \wedge \\ \quad (u \in \operatorname{interval}(1, n)) \wedge \\ \quad (v \in \operatorname{interval}(1, n)) \Rightarrow \\ \quad \quad \operatorname{entry}(\operatorname{minor}(s, p, q, n), u, v) \\ \quad \quad =\; \operatorname{entry}(s, \left(\text{if }(u < p)\text{ then }u\text{ else }(u + 1)\right), \left(\text{if }(v < q)\text{ then }v\text{ else }(v + 1)\right)) \end{array}}}
\vnbentry{\texttt{det-\allowbreak{}1x1}}{\vnbstmt{\begin{array}{@{}l@{}} \forall r, a. \\ \quad (\operatorname{is\text{-}ring}(r) \wedge (a \in \operatorname{mat}(1, 1, \operatorname{carr}(r)))) \Rightarrow  \\ \quad \quad (\operatorname{det}(r, 1, a) = \operatorname{entry}(a, 1, 1)) \end{array}}}
\vnbentry{\texttt{det-\allowbreak{}2x2}}{\vnbstmt{\begin{array}{@{}l@{}} \forall r, a. \\ \quad (\operatorname{is\text{-}ring}(r) \wedge (a \in \operatorname{mat}(2, 2, \operatorname{carr}(r)))) \Rightarrow  \\ \quad \quad \operatorname{det}(r, 2, a) \\ \quad \quad =\; \operatorname{add}(r)(\operatorname{mul}(r)(\operatorname{entry}(a, 1, 1), \operatorname{entry}(a, 2, 2)), \\ \quad \quad \quad \operatorname{neg}(r)(\operatorname{mul}(r)(\operatorname{entry}(a, 1, 2), \operatorname{entry}(a, 2, 1)))) \end{array}}}
\vnbentry{\texttt{identmat-\allowbreak{}minor}}{\vnbstmt{\begin{array}{@{}l@{}} \forall r, n. \\ \quad (\operatorname{is\text{-}ring}(r) \wedge (n \in \mathbb{N})) \Rightarrow  \\ \quad \quad (\operatorname{minor}(\operatorname{identmat}(r, (n + 1)), 1, 1, n) = \operatorname{identmat}(r, n)) \end{array}}}
\vnbentry{\texttt{mat-\allowbreak{}ring-\allowbreak{}one}}{\vnbstmt{\begin{array}{@{}l@{}} \forall a, n. \\ \quad (\operatorname{is\text{-}ring}(a) \wedge (n \in \mathbb{N})) \Rightarrow  \\ \quad \quad (\operatorname{one}(\operatorname{mat\text{-}ring}(a, n)) = \operatorname{identmat}(a, n)) \end{array}}}
\vnbentry{\texttt{det-\allowbreak{}identity-\allowbreak{}ind}}{\vnbstmt{\begin{array}{@{}l@{}} \forall r.; \forall n \in \mathbb{N}. \\ \quad \operatorname{is\text{-}ring}(r) \Rightarrow  \\ \quad \quad (\operatorname{det}(r, n, \operatorname{identmat}(r, n)) = \operatorname{one}(r)) \end{array}}}
\vnbentry{\texttt{det-\allowbreak{}identity}}{\vnbstmt{\begin{array}{@{}l@{}} \forall r, n. \\ \quad (\operatorname{is\text{-}ring}(r) \wedge (n \in \mathbb{N})) \Rightarrow  \\ \quad \quad \operatorname{det}(r, n, \operatorname{one}(\operatorname{mat\text{-}ring}(r, n))) \\ \quad \quad =\; \operatorname{one}(r) \end{array}}}

\vnbfile{theorem-library/rake-mat-ring-readoffs.scm}{theorem-library/rake-mat-ring-readoffs.scm}
\noindent{\small rake-mat-ring-readoffs.scm -- the MAT-RING slot read-offs and the two curried-operation FUN typings, PROVEN (rake, 2026-09-19).}\par
\vnbentry{\texttt{mat-\allowbreak{}ring-\allowbreak{}carr}}{\vnbstmt{\begin{array}{@{}l@{}} \forall a, n. \\ \quad (\operatorname{is\text{-}ring}(a) \wedge (n \in \mathbb{N})) \Rightarrow  \\ \quad \quad \operatorname{carr}(\operatorname{mat\text{-}ring}(a, n)) \\ \quad \quad =\; \operatorname{mat}(n, n, \operatorname{carr}(a)) \end{array}}}
\vnbentry{\texttt{mat-\allowbreak{}ring-\allowbreak{}add}}{\vnbstmt{\begin{array}{@{}l@{}} \forall a, n. \\ \quad (\operatorname{is\text{-}ring}(a) \wedge (n \in \mathbb{N})) \Rightarrow  \\ \quad \quad \operatorname{add}(\operatorname{mat\text{-}ring}(a, n)) \\ \quad \quad =\; (\lambda \langle p, q \rangle \in (\operatorname{mat}(n, n, \operatorname{carr}(a)) \times \operatorname{mat}(n, n, \operatorname{carr}(a))).\; \\ \quad \quad \quad \operatorname{matadd}(a, p, q)) \end{array}}}
\vnbentry{\texttt{mat-\allowbreak{}ring-\allowbreak{}mul}}{\vnbstmt{\begin{array}{@{}l@{}} \forall a, n. \\ \quad (\operatorname{is\text{-}ring}(a) \wedge (n \in \mathbb{N})) \Rightarrow  \\ \quad \quad \operatorname{mul}(\operatorname{mat\text{-}ring}(a, n)) \\ \quad \quad =\; (\lambda \langle p, q \rangle \in (\operatorname{mat}(n, n, \operatorname{carr}(a)) \times \operatorname{mat}(n, n, \operatorname{carr}(a))).\; \\ \quad \quad \quad \operatorname{matmul}(a, p, q)) \end{array}}}
\vnbentry{\texttt{mat-\allowbreak{}ring-\allowbreak{}neg}}{\vnbstmt{\begin{array}{@{}l@{}} \forall a, n. \\ \quad (\operatorname{is\text{-}ring}(a) \wedge (n \in \mathbb{N})) \Rightarrow  \\ \quad \quad \operatorname{neg}(\operatorname{mat\text{-}ring}(a, n)) \\ \quad \quad =\; (\lambda p \in \operatorname{mat}(n, n, \operatorname{carr}(a)).\; \operatorname{matneg}(a, p)) \end{array}}}
\vnbentry{\texttt{mat-\allowbreak{}ring-\allowbreak{}zero}}{\vnbstmt{\begin{array}{@{}l@{}} \forall a, n. \\ \quad (\operatorname{is\text{-}ring}(a) \wedge (n \in \mathbb{N})) \Rightarrow  \\ \quad \quad (\operatorname{zero}(\operatorname{mat\text{-}ring}(a, n)) = \operatorname{zeromat}(a, n, n)) \end{array}}}
\vnbentry{\texttt{mat-\allowbreak{}ring-\allowbreak{}add-\allowbreak{}fun}}{\vnbstmt{\begin{array}{@{}l@{}} \forall a.; \forall n \in \mathbb{N}. \\ \quad \operatorname{is\text{-}ring}(a) \Rightarrow  \\ \quad \quad (\lambda \langle p, q \rangle \in (\operatorname{mat}(n, n, \operatorname{carr}(a)) \times \operatorname{mat}(n, n, \operatorname{carr}(a))).\; \\ \quad \quad \quad \operatorname{matadd}(a, p, q)) \\ \quad \quad \in\; ((\operatorname{mat}(n, n, \operatorname{carr}(a)) \times \operatorname{mat}(n, n, \operatorname{carr}(a))) \to \operatorname{mat}(n, n, \operatorname{carr}(a))) \end{array}}}
\vnbentry{\texttt{mat-\allowbreak{}ring-\allowbreak{}mul-\allowbreak{}fun}}{\vnbstmt{\begin{array}{@{}l@{}} \forall a.; \forall n \in \mathbb{N}. \\ \quad \operatorname{is\text{-}ring}(a) \Rightarrow  \\ \quad \quad (\lambda \langle p, q \rangle \in (\operatorname{mat}(n, n, \operatorname{carr}(a)) \times \operatorname{mat}(n, n, \operatorname{carr}(a))).\; \\ \quad \quad \quad \operatorname{matmul}(a, p, q)) \\ \quad \quad \in\; ((\operatorname{mat}(n, n, \operatorname{carr}(a)) \times \operatorname{mat}(n, n, \operatorname{carr}(a))) \to \operatorname{mat}(n, n, \operatorname{carr}(a))) \end{array}}}

\vnbfile{theorem-library/matmul-entry-proof.scm}{theorem-library/matmul-entry-proof.scm}
\noindent{\small matmul-entry-proof.scm -- the matrix-product entry formula, PROVEN.}\par
\vnbentry{\texttt{matmul-\allowbreak{}entry}}{\vnbstmt{\begin{array}{@{}l@{}} \forall a, m, n, k, p, q.; \forall i \in \operatorname{interval}(1, m), c \in \operatorname{interval}(1, k). \\ \quad \operatorname{is\text{-}ring}(a) \wedge \\ \quad (p \in \operatorname{mat}(m, n, \operatorname{carr}(a))) \wedge \\ \quad (q \in \operatorname{mat}(n, k, \operatorname{carr}(a))) \wedge \\ \quad (1 \leq n) \Rightarrow \\ \quad \quad \operatorname{entry}(\operatorname{matmul}(a, p, q), i, c) \\ \quad \quad =\; \operatorname{finsum}(\operatorname{ring\text{-}additive\text{-}ag}(a), \\ \quad \quad \quad (\lambda j \in \operatorname{interval}(1, n).\; \\ \quad \quad \quad \quad \operatorname{mul}(a)(\operatorname{entry}(p, i, j), \operatorname{entry}(q, j, c))), \\ \quad \quad \quad \operatorname{interval}(1, n)) \end{array}}}

\vnbfile{theorem-library/matunit-matact-type.scm}{theorem-library/matunit-matact-type.scm}
\noindent{\small matunit-matact-type.scm -- two MATOF typings, PROVEN from the definitions.}\par
\vnbentry{\texttt{matunit-\allowbreak{}type}}{\vnbstmt{\begin{array}{@{}l@{}} \forall a, n, k, l. \\ \quad (\operatorname{is\text{-}ring}(a) \wedge (n \in \mathbb{N})) \Rightarrow  \\ \quad \quad (\operatorname{matunit}(a, n, k, l) \in \operatorname{mat}(n, n, \operatorname{carr}(a))) \end{array}}}
\vnbentry{\texttt{matact-\allowbreak{}type}}{\vnbstmt{\begin{array}{@{}l@{}} \forall d, m, n, q, p, u. \\ \quad \operatorname{is\text{-}module}(d) \wedge \\ \quad (p \in \operatorname{mat}(m, n, \operatorname{carr}(\operatorname{scal}(d)))) \wedge \\ \quad (u \in \operatorname{mat}(n, q, \operatorname{vec}(d))) \wedge \\ \quad ((n = 0) \Rightarrow ((m = 0) \vee (q = 0))) \Rightarrow \\ \quad \quad (\operatorname{matact}(d, p, u) \in \operatorname{mat}(m, q, \operatorname{vec}(d))) \end{array}}}

\vnbfile{theorem-library/nn-finite-subset-bounded.scm}{theorem-library/nn-finite-subset-bounded.scm}
\noindent{\small nn-finite-subset-bounded.scm -- a finite subset of NN is bounded.}\par
\vnbentry{\texttt{nn-\allowbreak{}finite-\allowbreak{}subset-\allowbreak{}bounded}}{\vnbstmt{\begin{array}{@{}l@{}} \forall t. \\ \quad ((t \subseteq \mathbb{N}) \wedge (|t| \in \mathbb{N})) \Rightarrow  \\ \quad \quad \exists n \in \mathbb{N}. \forall y \in \mathbb{N}. \\ \quad \quad \quad ((n \leq y) \Rightarrow \neg (y \in t)) \end{array}}}

\vnbfile{theorem-library/card-inequalities.scm}{theorem-library/card-inequalities.scm}
\noindent{\small card-inequalities.scm -- the two missing CARD inequalities. --------------------------------------------------------------------------- NOTE, 2026-09-20 (batch 9-B).  This file was written while cardinality was AXIOMATISED under the name CARD and the defined constant was its companion}\par
\vnbentry{\texttt{complement-\allowbreak{}in-\allowbreak{}subset}}{\vnbstmt{\forall a, b.\; (\operatorname{complement\text{-}in}(b, a) \subseteq b)}}
\vnbentry{\texttt{union-\allowbreak{}complement-\allowbreak{}in}}{\vnbstmt{\begin{array}{@{}l@{}} \forall a, b. \\ \quad (a \subseteq b) \Rightarrow  \\ \quad \quad (b = (a \cup \operatorname{complement\text{-}in}(b, a))) \end{array}}}
\vnbentry{\texttt{union-\allowbreak{}as-\allowbreak{}disjoint}}{\vnbstmt{\begin{array}{@{}l@{}} \forall a, b. \\ \quad ((a \cup b) = (a \cup \operatorname{complement\text{-}in}(b, a))) \end{array}}}
\vnbentry{\texttt{intersection-\allowbreak{}complement-\allowbreak{}in-\allowbreak{}empty}}{\vnbstmt{\forall a, b.\; ((a \cap \operatorname{complement\text{-}in}(b, a)) = \emptyset)}}
\vnbentry{\texttt{card-\allowbreak{}subset-\allowbreak{}mono}}{\vnbstmt{\begin{array}{@{}l@{}} \forall b \in \mathbf{Set}, a. \\ \quad (((|b| \in \mathbb{N}) \wedge (a \subseteq b)) \Rightarrow (|a| \leq |b|)) \end{array}}}
\vnbentry{\texttt{card-\allowbreak{}union-\allowbreak{}bound}}{\vnbstmt{\begin{array}{@{}l@{}} \forall a, b \in \mathbf{Set}. \\ \quad ((|a| \in \mathbb{N}) \wedge (|b| \in \mathbb{N})) \Rightarrow  \\ \quad \quad (|(a \cup b)| \leq (|a| + |b|)) \end{array}}}
\vnbentry{\texttt{card-\allowbreak{}union-\allowbreak{}nn}}{\vnbstmt{\begin{array}{@{}l@{}} \forall a, b \in \mathbf{Set}. \\ \quad ((|a| \in \mathbb{N}) \wedge (|b| \in \mathbb{N})) \Rightarrow  \\ \quad \quad (|(a \cup b)| \in \mathbb{N}) \end{array}}}

\vnbfile{theorem-library/rake-combinatorics.scm}{theorem-library/rake-combinatorics.scm}
\noindent{\small rake-combinatorics.scm -- BATCH P of the 2026-09-17 rake: NN combinatorics, the ring powers, the powerset insert-split, and the two CHOOSE boundary values.  Seventeen asserted supports, PROVEN; every statement is copied LITERALLY from its site.}\par
\vnbentry{\texttt{succ-\allowbreak{}nn-\allowbreak{}ord}}{\vnbstmt{\forall n \in \mathbb{N}.\; ((n + 1) = (n + 1))}}
\vnbentry{\texttt{union-\allowbreak{}empty-\allowbreak{}left}}{\vnbstmt{\forall a \in \mathbf{Set}.\; ((\emptyset \cup a) = a)}}
\vnbentry{\texttt{ord-\allowbreak{}segment-\allowbreak{}trans}}{\vnbstmt{\begin{array}{@{}l@{}} \forall m \in \mathbb{N}, k \in \mathrm{OS}(m), i \in \mathrm{OS}(k). \\ \quad (i \in \mathrm{OS}(m)) \end{array}}}
\vnbentry{\texttt{power-\allowbreak{}zero-\allowbreak{}base}}{\vnbstmt{\forall n \in \mathbb{N}.\; (\operatorname{power}(0, (n + 1)) = 0)}}
\vnbentry{\texttt{crmcm-\allowbreak{}op}}{\vnbstmt{\begin{array}{@{}l@{}} \forall r. \\ \quad \operatorname{is\text{-}commutative\text{-}ring}(r) \Rightarrow  \\ \quad \quad \operatorname{opr}(\operatorname{commutative\text{-}ring\text{-}multiplicative\text{-}cm}(r)) \\ \quad \quad =\; \operatorname{mul}(r) \end{array}}}
\vnbentry{\texttt{crmcm-\allowbreak{}iden}}{\vnbstmt{\begin{array}{@{}l@{}} \forall r. \\ \quad \operatorname{is\text{-}commutative\text{-}ring}(r) \Rightarrow  \\ \quad \quad \operatorname{iden}(\operatorname{commutative\text{-}ring\text{-}multiplicative\text{-}cm}(r)) \\ \quad \quad =\; \operatorname{one}(r) \end{array}}}
\vnbentry{\texttt{ring-\allowbreak{}power-\allowbreak{}succ-\allowbreak{}left}}{\vnbstmt{\begin{array}{@{}l@{}} \forall r.; \forall x \in \operatorname{carr}(r), n \in \mathbb{N}. \\ \quad \operatorname{is\text{-}commutative\text{-}ring}(r) \Rightarrow  \\ \quad \quad \operatorname{ring\text{-}power}(r, x, (n + 1)) \\ \quad \quad =\; \operatorname{mul}(r)(x, \operatorname{ring\text{-}power}(r, x, n)) \end{array}}}
\vnbentry{\texttt{ring-\allowbreak{}power-\allowbreak{}one}}{\vnbstmt{\begin{array}{@{}l@{}} \forall r.; \forall x \in \operatorname{carr}(r). \\ \quad \operatorname{is\text{-}commutative\text{-}ring}(r) \Rightarrow  \\ \quad \quad (\operatorname{ring\text{-}power}(r, x, 1) = x) \end{array}}}
\vnbentry{\texttt{ring-\allowbreak{}power-\allowbreak{}succ}}{\vnbstmt{\begin{array}{@{}l@{}} \forall r.; \forall x \in \operatorname{carr}(r), n \in \mathbb{N}. \\ \quad \operatorname{is\text{-}commutative\text{-}ring}(r) \Rightarrow  \\ \quad \quad \operatorname{ring\text{-}power}(r, x, (n + 1)) \\ \quad \quad =\; \operatorname{mul}(r)(\operatorname{ring\text{-}power}(r, x, n), x) \end{array}}}
\vnbentry{\texttt{ring-\allowbreak{}mul-\allowbreak{}interchange}}{\vnbstmt{\begin{array}{@{}l@{}} \forall r.; \forall a, b, c, d \in \operatorname{carr}(r). \\ \quad \operatorname{is\text{-}commutative\text{-}ring}(r) \Rightarrow  \\ \quad \quad \operatorname{mul}(r)(\operatorname{mul}(r)(a, b), \operatorname{mul}(r)(c, d)) \\ \quad \quad =\; \operatorname{mul}(r)(\operatorname{mul}(r)(a, c), \operatorname{mul}(r)(b, d)) \end{array}}}
\vnbentry{\texttt{ring-\allowbreak{}power-\allowbreak{}add-\allowbreak{}ind}}{\vnbstmt{\begin{array}{@{}l@{}} \forall k \in \mathbb{N}, r.; \forall x \in \operatorname{carr}(r), j \in \mathbb{N}. \\ \quad \operatorname{is\text{-}commutative\text{-}ring}(r) \Rightarrow  \\ \quad \quad \operatorname{ring\text{-}power}(r, x, (j + k)) \\ \quad \quad =\; \operatorname{mul}(r)(\operatorname{ring\text{-}power}(r, x, j), \operatorname{ring\text{-}power}(r, x, k)) \end{array}}}
\vnbentry{\texttt{ring-\allowbreak{}power-\allowbreak{}add}}{\vnbstmt{\begin{array}{@{}l@{}} \forall r.; \forall x \in \operatorname{carr}(r), j, k \in \mathbb{N}. \\ \quad \operatorname{is\text{-}commutative\text{-}ring}(r) \Rightarrow  \\ \quad \quad \operatorname{ring\text{-}power}(r, x, (j + k)) \\ \quad \quad =\; \operatorname{mul}(r)(\operatorname{ring\text{-}power}(r, x, j), \operatorname{ring\text{-}power}(r, x, k)) \end{array}}}
\vnbentry{\texttt{ring-\allowbreak{}power-\allowbreak{}mult-\allowbreak{}ind}}{\vnbstmt{\begin{array}{@{}l@{}} \forall n \in \mathbb{N}, r.; \forall x, y \in \operatorname{carr}(r). \\ \quad \operatorname{is\text{-}commutative\text{-}ring}(r) \Rightarrow  \\ \quad \quad \operatorname{ring\text{-}power}(r, \operatorname{mul}(r)(x, y), n) \\ \quad \quad =\; \operatorname{mul}(r)(\operatorname{ring\text{-}power}(r, x, n), \operatorname{ring\text{-}power}(r, y, n)) \end{array}}}
\vnbentry{\texttt{ring-\allowbreak{}power-\allowbreak{}mult}}{\vnbstmt{\begin{array}{@{}l@{}} \forall r.; \forall x, y \in \operatorname{carr}(r), n \in \mathbb{N}. \\ \quad \operatorname{is\text{-}commutative\text{-}ring}(r) \Rightarrow  \\ \quad \quad \operatorname{ring\text{-}power}(r, \operatorname{mul}(r)(x, y), n) \\ \quad \quad =\; \operatorname{mul}(r)(\operatorname{ring\text{-}power}(r, x, n), \operatorname{ring\text{-}power}(r, y, n)) \end{array}}}
\vnbentry{\texttt{prod-\allowbreak{}ring-\allowbreak{}empty}}{\vnbstmt{\begin{array}{@{}l@{}} \forall r, f. \\ \quad \operatorname{is\text{-}commutative\text{-}ring}(r) \Rightarrow  \\ \quad \quad (\operatorname{prod\text{-}ring}(r, f, \emptyset) = \operatorname{one}(r)) \end{array}}}
\vnbentry{\texttt{prod-\allowbreak{}ring-\allowbreak{}insert}}{\vnbstmt{\begin{array}{@{}l@{}} \forall r.; \forall x, k \in \mathbf{Set}, f \in ((x \cup \{k\}) \to \operatorname{carr}(r)). \\ \quad \operatorname{is\text{-}commutative\text{-}ring}(r) \wedge \\ \quad (|x| \in \mathbb{N}) \wedge \\ \quad \neg (k \in x) \Rightarrow \\ \quad \quad \operatorname{prod\text{-}ring}(r, f, (x \cup \{k\})) \\ \quad \quad =\; \operatorname{mul}(r)(\operatorname{prod\text{-}ring}(r, f, x), f(k)) \end{array}}}
\vnbentry{\texttt{prod-\allowbreak{}ring-\allowbreak{}singleton}}{\vnbstmt{\begin{array}{@{}l@{}} \forall r.; \forall x \in \mathbf{Set}, f \in (\{x\} \to \operatorname{carr}(r)). \\ \quad \operatorname{is\text{-}commutative\text{-}ring}(r) \Rightarrow  \\ \quad \quad (\operatorname{prod\text{-}ring}(r, f, \{x\}) = f(x)) \end{array}}}
\vnbentry{\texttt{power-\allowbreak{}insert-\allowbreak{}disjoint}}{\vnbstmt{\begin{array}{@{}l@{}} \forall x, k \in \mathbf{Set}. \\ \quad \neg (k \in x) \Rightarrow  \\ \quad \quad (\operatorname{power}(x) \cap \operatorname{image}((\lambda s \in \operatorname{power}(x).\; (s \cup \{k\})), \operatorname{power}(x))) \\ \quad \quad =\; \emptyset \end{array}}}
\vnbentry{\texttt{choose-\allowbreak{}set-\allowbreak{}unfold}}{\vnbstmt{\begin{array}{@{}l@{}} \forall n, m. \\ \quad (\operatorname{choose\text{-}set}(n, m) \simeq \{a \in \operatorname{power}(\mathrm{OS}(n)) : (|a| = m)\}) \end{array}}}
\vnbentry{\texttt{choose-\allowbreak{}0-\allowbreak{}succ}}{\vnbstmt{\forall k \in \mathbb{N}.\; (\operatorname{choose}(0, (k + 1)) = 0)}}
\vnbentry{\texttt{choose-\allowbreak{}n-\allowbreak{}0}}{\vnbstmt{\forall n \in \mathbb{N}.\; (\operatorname{choose}(n, 0) = (0 + 1))}}
\vnbentry{\texttt{power-\allowbreak{}insert-\allowbreak{}cover}}{\vnbstmt{\begin{array}{@{}l@{}} \forall x, k \in \mathbf{Set}. \\ \quad \neg (k \in x) \Rightarrow  \\ \quad \quad \operatorname{power}((x \cup \{k\})) \\ \quad \quad =\; (\operatorname{power}(x) \cup \operatorname{image}((\lambda s \in \operatorname{power}(x).\; (s \cup \{k\})), \operatorname{power}(x))) \end{array}}}
\vnbentry{\texttt{subset-\allowbreak{}of-\allowbreak{}empty-\allowbreak{}is-\allowbreak{}empty}}{\vnbstmt{\begin{array}{@{}l@{}} \forall w \in \mathbf{Set}. \\ \quad (\forall z \in w.\; (z \in \emptyset) \Rightarrow (w = \emptyset)) \end{array}}}
\vnbentry{\texttt{power-\allowbreak{}of-\allowbreak{}empty}}{\vnbstmt{(\operatorname{power}(\emptyset) = \{\emptyset\})}}
\vnbentry{\texttt{card-\allowbreak{}power-\allowbreak{}nn}}{\vnbstmt{\begin{array}{@{}l@{}} \forall x \in \mathbf{Set}. \\ \quad ((|x| \in \mathbb{N}) \Rightarrow (|\operatorname{power}(x)| \in \mathbb{N})) \end{array}}}

\vnbfile{theorem-library/rake-combinatorics2.scm}{theorem-library/rake-combinatorics2.scm}
\noindent{\small rake-combinatorics2.scm -- BATCH U of the 2026-09-18 rake (batch 5, the off-bill supports): finite combinatorics.  Every statement is copied LITERALLY from its support site; the sites are named beside each proof. Six results, all `proven modulo 0'.}\par
\vnbentry{\texttt{interval-\allowbreak{}card}}{\vnbstmt{\forall n \in \mathbb{N}.\; (|\operatorname{interval}(1, n)| = n)}}
\vnbentry{\texttt{choose-\allowbreak{}in-\allowbreak{}nn}}{\vnbstmt{\forall n, m \in \mathbb{N}.\; (\operatorname{choose}(n, m) \in \mathbb{N})}}
\vnbentry{\texttt{injection-\allowbreak{}from-\allowbreak{}empty}}{\vnbstmt{\begin{array}{@{}l@{}} \forall c \in \mathbf{Set}. \\ \quad (|\operatorname{injection}(\emptyset, c)| = (0 + 1)) \end{array}}}
\vnbentry{\texttt{permutations-\allowbreak{}zero}}{\vnbstmt{(|\operatorname{permutations}(0)| = (0 + 1))}}
\vnbentry{\texttt{fibrewise-\allowbreak{}finite}}{\vnbstmt{\begin{array}{@{}l@{}} \forall p.; \forall g, t \in \mathbf{Set}. \\ \quad (|g| \in \mathbb{N}) \wedge \\ \quad \forall y \in t.\; (p(y) \in g) \wedge \\ \quad \forall c \in g.\; (|\{x \in t : (p(x) = c)\}| \in \mathbb{N}) \Rightarrow \\ \quad \quad (|t| \in \mathbb{N}) \end{array}}}
\vnbentry{\texttt{pigeonhole-\allowbreak{}infinite}}{\vnbstmt{\begin{array}{@{}l@{}} \forall s, f \in \mathbf{Set}, p \in (s \to f). \\ \quad (\neg (|s| \in \mathbb{N}) \wedge (|f| \in \mathbb{N})) \Rightarrow  \\ \quad \quad \exists c \in f.\; \neg (|\{x \in s : (p(x) = c)\}| \in \mathbb{N}) \end{array}}}

\vnbfile{theorem-library/rake-choose-succ.scm}{theorem-library/rake-choose-succ.scm}
\noindent{\small rake-choose-succ.scm -- PASCAL'S RULE for the binomial coefficient, and the CHOOSE-SET lemmas it is assembled from.  Rake batch 5b, assignment 5b-G.}\par
\vnbentry{\texttt{choose-\allowbreak{}card-\allowbreak{}unfold}}{\vnbstmt{\begin{array}{@{}l@{}} \forall n, m. \\ \quad (|\operatorname{choose\text{-}set}(n, m)| \simeq \operatorname{choose}(n, m)) \end{array}}}
\vnbentry{\texttt{choose-\allowbreak{}set-\allowbreak{}mem-\allowbreak{}iff}}{\vnbstmt{\begin{array}{@{}l@{}} \forall n, m, a. \\ \quad ((a \in \operatorname{choose\text{-}set}(n, m)) \Leftrightarrow ((a \in \operatorname{power}(\mathrm{OS}(n))) \wedge (|a| = m))) \end{array}}}
\vnbentry{\texttt{choose-\allowbreak{}set-\allowbreak{}avoids-\allowbreak{}point}}{\vnbstmt{\begin{array}{@{}l@{}} \forall n \in \mathbb{N}, m.; \forall a \in \operatorname{choose\text{-}set}(n, m). \\ \quad \neg (n \in a) \end{array}}}
\vnbentry{\texttt{choose-\allowbreak{}set-\allowbreak{}widen}}{\vnbstmt{\begin{array}{@{}l@{}} \forall n \in \mathbb{N}, m.; \forall a \in \operatorname{choose\text{-}set}(n, m). \\ \quad (a \in \operatorname{choose\text{-}set}((n + 1), m)) \end{array}}}
\vnbentry{\texttt{choose-\allowbreak{}set-\allowbreak{}insert-\allowbreak{}in}}{\vnbstmt{\begin{array}{@{}l@{}} \forall n, k \in \mathbb{N}, s \in \operatorname{choose\text{-}set}(n, k). \\ \quad ((s \cup \{n\}) \in \operatorname{choose\text{-}set}((n + 1), (k + 1))) \end{array}}}
\vnbentry{\texttt{choose-\allowbreak{}set-\allowbreak{}remove-\allowbreak{}restores}}{\vnbstmt{\begin{array}{@{}l@{}} \forall n, a \in \mathbf{Set}. \\ \quad ((n \in a) \Rightarrow ((\operatorname{difference}(a, \{n\}) \cup \{n\}) = a)) \end{array}}}
\vnbentry{\texttt{choose-\allowbreak{}set-\allowbreak{}remove-\allowbreak{}in}}{\vnbstmt{\begin{array}{@{}l@{}} \forall n, k \in \mathbb{N}, a \in \operatorname{choose\text{-}set}((n + 1), (k + 1)). \\ \quad (n \in a) \Rightarrow  \\ \quad \quad (\operatorname{difference}(a, \{n\}) \in \operatorname{choose\text{-}set}(n, k)) \end{array}}}
\vnbentry{\texttt{choose-\allowbreak{}set-\allowbreak{}avoid-\allowbreak{}in}}{\vnbstmt{\begin{array}{@{}l@{}} \forall n \in \mathbb{N}, m.; \forall a \in \operatorname{choose\text{-}set}((n + 1), m). \\ \quad (\neg (n \in a) \Rightarrow (a \in \operatorname{choose\text{-}set}(n, m))) \end{array}}}
\vnbentry{\texttt{choose-\allowbreak{}set-\allowbreak{}insert-\allowbreak{}injective}}{\vnbstmt{\begin{array}{@{}l@{}} \forall n \in \mathbb{N}, m.; \forall s, t \in \operatorname{choose\text{-}set}(n, m). \\ \quad (((s \cup \{n\}) = (t \cup \{n\})) \Rightarrow (s = t)) \end{array}}}
\vnbentry{\texttt{choose-\allowbreak{}set-\allowbreak{}split}}{\vnbstmt{\begin{array}{@{}l@{}} \forall n, k \in \mathbb{N}. \\ \quad \operatorname{choose\text{-}set}((n + 1), (k + 1)) \\ \quad =\; (\operatorname{choose\text{-}set}(n, (k + 1)) \cup \operatorname{image}((\lambda s \in \operatorname{choose\text{-}set}(n, k).\; (s \cup \{n\})), \operatorname{choose\text{-}set}(n, k))) \end{array}}}
\vnbentry{\texttt{choose-\allowbreak{}set-\allowbreak{}split-\allowbreak{}disjoint}}{\vnbstmt{\begin{array}{@{}l@{}} \forall n, k \in \mathbb{N}. \\ \quad (\operatorname{choose\text{-}set}(n, (k + 1)) \cap \operatorname{image}((\lambda s \in \operatorname{choose\text{-}set}(n, k).\; (s \cup \{n\})), \operatorname{choose\text{-}set}(n, k))) \\ \quad =\; \emptyset \end{array}}}
\vnbentry{\texttt{choose-\allowbreak{}succ}}{\vnbstmt{\begin{array}{@{}l@{}} \forall n, k \in \mathbb{N}. \\ \quad \operatorname{choose}((n + 1), (k + 1)) \\ \quad =\; (\operatorname{choose}(n, k) + \operatorname{choose}(n, (k + 1))) \end{array}}}

\vnbfile{theorem-library/nn-infinite.scm}{theorem-library/nn-infinite.scm}
\noindent{\small nn-infinite.scm -- NN is infinite, and the two leaves that hang on it.}\par
\vnbentry{\texttt{inf-\allowbreak{}subsets-\allowbreak{}unfold}}{\vnbstmt{\begin{array}{@{}l@{}} \forall a. \\ \quad (\mathrm{Inf}(a) \simeq \{s \in \operatorname{power}(a) : \neg (|s| \in \mathbb{N})\}) \end{array}}}
\vnbentry{\texttt{nn-\allowbreak{}not-\allowbreak{}finite}}{\vnbstmt{\neg (|\mathbb{N}| \in \mathbb{N})}}
\vnbentry{\texttt{nn-\allowbreak{}in-\allowbreak{}inf-\allowbreak{}subsets}}{\vnbstmt{(\mathbb{N} \in \mathrm{Inf}(\mathbb{N}))}}
\vnbentry{\texttt{inf-\allowbreak{}subset-\allowbreak{}nn-\allowbreak{}unbounded}}{\vnbstmt{\begin{array}{@{}l@{}} \forall t \in \mathrm{Inf}(\mathbb{N}), u \in \mathbb{N}. \exists y \in \mathbb{N}. \\ \quad ((y \in t) \wedge (u < y)) \end{array}}}

\vnbfile{theorem-library/fin-subsets.scm}{theorem-library/fin-subsets.scm}
\noindent{\small fin-subsets.scm -- FIN-SUBSETS(a), the finite subsets of a, and the --------------------------------------------------------------------------- NOTE, 2026-09-20 (batch 9-B).  This file was written while cardinality was AXIOMATISED under the name CARD and the defined constant was its companion}\par
\vnbentry{\texttt{fin-\allowbreak{}subsets-\allowbreak{}membership}}{\vnbstmt{\begin{array}{@{}l@{}} \forall a, s. \\ \quad ((s \in \operatorname{fin\text{-}subsets}(a)) \Leftrightarrow ((s \in \mathbf{Set}) \wedge (\forall z \in s.\; (z \in a) \wedge (|s| \in \mathbb{N})))) \end{array}}}
\vnbentry{\texttt{fin-\allowbreak{}subsets-\allowbreak{}is-\allowbreak{}set}}{\vnbstmt{\begin{array}{@{}l@{}} \forall a \in \mathbf{Set}. \\ \quad (\operatorname{fin\text{-}subsets}(a) \in \mathbf{Set}) \end{array}}}
\vnbentry{\texttt{fin-\allowbreak{}subsets-\allowbreak{}has-\allowbreak{}empty}}{\vnbstmt{\forall a.\; (\emptyset \in \operatorname{fin\text{-}subsets}(a))}}
\vnbentry{\texttt{fin-\allowbreak{}subsets-\allowbreak{}union-\allowbreak{}closed}}{\vnbstmt{\begin{array}{@{}l@{}} \forall a, s, u. \\ \quad ((s \in \operatorname{fin\text{-}subsets}(a)) \wedge (u \in \operatorname{fin\text{-}subsets}(a))) \Rightarrow  \\ \quad \quad ((s \cup u) \in \operatorname{fin\text{-}subsets}(a)) \end{array}}}

\vnbfile{theorem-library/rake-finsum-union.scm}{theorem-library/rake-finsum-union.scm}
\noindent{\small theorem-library/rake-finsum-union.scm -- the two set/card bricks that three rake batches built inline, the POINTWISE finsum recurrences they make possible, the DISJOINT-UNION law for FINSUM, and `finsum-embed'. Rake batch 5c, assignment 5c-M.  Helper prefix `r7m-'; every helper is}\par
\vnbentry{\texttt{card-\allowbreak{}insert-\allowbreak{}converse}}{\vnbstmt{\begin{array}{@{}l@{}} \forall a, x \in \mathbf{Set}, k \in \mathbb{N}. \\ \quad ((\neg (x \in a) \wedge (|(a \cup \{x\})| = (k + 1))) \Rightarrow (|a| = k)) \end{array}}}
\vnbentry{\texttt{remove-\allowbreak{}restore}}{\vnbstmt{\begin{array}{@{}l@{}} \forall a \in \mathbf{Set}, x \in a. \\ \quad ((\operatorname{difference}(a, \{x\}) \cup \{x\}) = a) \end{array}}}
\vnbentry{\texttt{finsum-\allowbreak{}insert-\allowbreak{}ptwise}}{\vnbstmt{\begin{array}{@{}l@{}} \forall g.; \forall s, k \in \mathbf{Set}, f. \\ \quad \operatorname{is\text{-}abelian\text{-}group}(g) \wedge \\ \quad (|s| \in \mathbb{N}) \wedge \\ \quad \neg (k \in s) \wedge \\ \quad \forall z \in (s \cup \{k\}).\; (f(z) \in \operatorname{carr}(g)) \Rightarrow \\ \quad \quad \operatorname{finsum}(g, f, (s \cup \{k\})) \\ \quad \quad =\; \operatorname{opr}(g)(\operatorname{finsum}(g, f, s), f(k)) \end{array}}}
\vnbentry{\texttt{finsum-\allowbreak{}union-\allowbreak{}disjoint}}{\vnbstmt{\begin{array}{@{}l@{}} \forall g.; \forall s, t \in \mathbf{Set}, f. \\ \quad \operatorname{is\text{-}abelian\text{-}group}(g) \wedge \\ \quad (|s| \in \mathbb{N}) \wedge \\ \quad (|t| \in \mathbb{N}) \wedge \\ \quad \forall z \in t.\; \neg (z \in s) \wedge \\ \quad \forall z \in (s \cup t).\; (f(z) \in \operatorname{carr}(g)) \Rightarrow \\ \quad \quad \operatorname{finsum}(g, f, (s \cup t)) \\ \quad \quad =\; \operatorname{opr}(g)(\operatorname{finsum}(g, f, s), \operatorname{finsum}(g, f, t)) \end{array}}}
\vnbentry{\texttt{finsum-\allowbreak{}union-\allowbreak{}disjoint-\allowbreak{}fun}}{\vnbstmt{\begin{array}{@{}l@{}} \forall g.; \forall s, t \in \mathbf{Set}, f \in ((s \cup t) \to \operatorname{carr}(g)). \\ \quad \operatorname{is\text{-}abelian\text{-}group}(g) \wedge \\ \quad (|s| \in \mathbb{N}) \wedge \\ \quad (|t| \in \mathbb{N}) \wedge \\ \quad ((s \cap t) = \emptyset) \Rightarrow \\ \quad \quad \operatorname{finsum}(g, f, (s \cup t)) \\ \quad \quad =\; \operatorname{opr}(g)(\operatorname{finsum}(g, f, s), \operatorname{finsum}(g, f, t)) \end{array}}}
\vnbentry{\texttt{finsum-\allowbreak{}all-\allowbreak{}id-\allowbreak{}ptwise}}{\vnbstmt{\begin{array}{@{}l@{}} \forall g.; \forall s \in \mathbf{Set}, f. \\ \quad \operatorname{is\text{-}abelian\text{-}group}(g) \wedge \\ \quad (|s| \in \mathbb{N}) \wedge \\ \quad \forall z \in s.\; (f(z) = \operatorname{iden}(g)) \Rightarrow \\ \quad \quad (\operatorname{finsum}(g, f, s) = \operatorname{iden}(g)) \end{array}}}
\vnbentry{\texttt{finsum-\allowbreak{}embed}}{\vnbstmt{\begin{array}{@{}l@{}} \forall g.; \forall s, a \in \mathbf{Set}, f \in (a \to \operatorname{carr}(g)). \\ \quad \operatorname{is\text{-}abelian\text{-}group}(g) \wedge \\ \quad (|s| \in \mathbb{N}) \wedge \\ \quad (|a| \in \mathbb{N}) \wedge \\ \quad (s \subseteq a) \wedge \\ \quad \forall z \in a.\; (\neg (z \in s) \Rightarrow (f(z) = \operatorname{iden}(g))) \Rightarrow \\ \quad \quad (\operatorname{finsum}(g, f, a) = \operatorname{finsum}(g, f, s)) \end{array}}}

\vnbfile{theorem-library/rake-finsum-cm-union.scm}{theorem-library/rake-finsum-cm-union.scm}
\noindent{\small theorem-library/rake-finsum-cm-union.scm -- the COMMUTATIVE-MONOID mirror of the FINSUM congruence / insertion / disjoint-union block, plus the two set identities that a peeling induction over POWER(X) needs. Rake batch 5c, assignment 5c-T.  Helper prefix `r7t-'; every helper is}\par
\vnbentry{\texttt{finsum-\allowbreak{}cm-\allowbreak{}congruence-\allowbreak{}q}}{\vnbstmt{\begin{array}{@{}l@{}} \forall m.; \forall s \in \mathbf{Set}, f, g. \\ \quad \operatorname{is\text{-}comm\text{-}monoid}(m) \wedge \\ \quad (|s| \in \mathbb{N}) \wedge \\ \quad \forall z \in s.\; (f(z) = g(z)) \Rightarrow \\ \quad \quad (\operatorname{finsum}(m, f, s) \simeq \operatorname{finsum}(m, g, s)) \end{array}}}
\vnbentry{\texttt{finsum-\allowbreak{}cm-\allowbreak{}insert-\allowbreak{}ptwise}}{\vnbstmt{\begin{array}{@{}l@{}} \forall m.; \forall s, k \in \mathbf{Set}, f. \\ \quad \operatorname{is\text{-}comm\text{-}monoid}(m) \wedge \\ \quad (|s| \in \mathbb{N}) \wedge \\ \quad \neg (k \in s) \wedge \\ \quad \forall z \in (s \cup \{k\}).\; (f(z) \in \operatorname{carr}(m)) \Rightarrow \\ \quad \quad \operatorname{finsum}(m, f, (s \cup \{k\})) \\ \quad \quad =\; \operatorname{opr}(m)(\operatorname{finsum}(m, f, s), f(k)) \end{array}}}
\vnbentry{\texttt{finsum-\allowbreak{}cm-\allowbreak{}union-\allowbreak{}disjoint}}{\vnbstmt{\begin{array}{@{}l@{}} \forall m.; \forall s, t \in \mathbf{Set}, f. \\ \quad \operatorname{is\text{-}comm\text{-}monoid}(m) \wedge \\ \quad (|s| \in \mathbb{N}) \wedge \\ \quad (|t| \in \mathbb{N}) \wedge \\ \quad \forall z \in t.\; \neg (z \in s) \wedge \\ \quad \forall z \in (s \cup t).\; (f(z) \in \operatorname{carr}(m)) \Rightarrow \\ \quad \quad \operatorname{finsum}(m, f, (s \cup t)) \\ \quad \quad =\; \operatorname{opr}(m)(\operatorname{finsum}(m, f, s), \operatorname{finsum}(m, f, t)) \end{array}}}
\vnbentry{\texttt{insert-\allowbreak{}difference-\allowbreak{}outside}}{\vnbstmt{\begin{array}{@{}l@{}} \forall x, s.; \forall k \in \mathbf{Set}. \\ \quad \neg (k \in s) \Rightarrow  \\ \quad \quad (\operatorname{difference}((x \cup \{k\}), s) = (\operatorname{difference}(x, s) \cup \{k\})) \end{array}}}
\vnbentry{\texttt{insert-\allowbreak{}difference-\allowbreak{}cancel}}{\vnbstmt{\begin{array}{@{}l@{}} \forall x, s.; \forall k \in \mathbf{Set}. \\ \quad \neg (k \in x) \Rightarrow  \\ \quad \quad (\operatorname{difference}((x \cup \{k\}), (s \cup \{k\})) = \operatorname{difference}(x, s)) \end{array}}}

\vnbfile{theorem-library/rake-sum-set-defined.scm}{theorem-library/rake-sum-set-defined.scm}
\noindent{\small theorem-library/rake-sum-set-defined.scm -- SUM-SET DEFINED, and its four characterising laws PROVEN from the definition.  Rake batch 5c, assignment 5c-S (the user's decision (B) of 2026-09-18).  Helper prefix `r7s-'.}\par
\vnbentry{\texttt{sum-\allowbreak{}set-\allowbreak{}empty}}{\vnbstmt{\begin{array}{@{}l@{}} \forall r, f. \\ \quad (\operatorname{sum\text{-}set}(r, \emptyset, f) \simeq \operatorname{zero}(r)) \end{array}}}
\vnbentry{\texttt{finsum-\allowbreak{}singleton-\allowbreak{}ptwise}}{\vnbstmt{\begin{array}{@{}l@{}} \forall g.; \forall x \in \mathbf{Set}, f. \\ \quad (\operatorname{is\text{-}abelian\text{-}group}(g) \wedge (f(x) \in \operatorname{carr}(g))) \Rightarrow  \\ \quad \quad (\operatorname{finsum}(g, f, \{x\}) = f(x)) \end{array}}}
\vnbentry{\texttt{sum-\allowbreak{}set-\allowbreak{}singleton-\allowbreak{}defined}}{\vnbstmt{\begin{array}{@{}l@{}} \forall r.; \forall x \in \mathbf{Set}, f. \\ \quad (\operatorname{is\text{-}ring}(r) \wedge (f(x) \in \operatorname{carr}(r))) \Rightarrow  \\ \quad \quad (\operatorname{sum\text{-}set}(r, \{x\}, f) = f(x)) \end{array}}}
\vnbentry{\texttt{sum-\allowbreak{}set-\allowbreak{}type-\allowbreak{}defined}}{\vnbstmt{\begin{array}{@{}l@{}} \forall r, x, s, f. \\ \quad \operatorname{is\text{-}ring}(r) \wedge \\ \quad (x \in \mathbf{Set}) \wedge \\ \quad (s \subseteq x) \wedge \\ \quad (f \in (x \to \operatorname{carr}(r))) \wedge \\ \quad (|s| \in \mathbb{N}) \Rightarrow \\ \quad \quad (\operatorname{sum\text{-}set}(r, s, f) \in \operatorname{carr}(r)) \end{array}}}
\vnbentry{\texttt{sum-\allowbreak{}set-\allowbreak{}disjoint-\allowbreak{}union-\allowbreak{}defined}}{\vnbstmt{\begin{array}{@{}l@{}} \forall r, x, a, b, f. \\ \quad \operatorname{is\text{-}ring}(r) \wedge \\ \quad (x \in \mathbf{Set}) \wedge \\ \quad (a \in \mathbf{Set}) \wedge \\ \quad (b \in \mathbf{Set}) \wedge \\ \quad ((a \cap b) = \emptyset) \wedge \\ \quad ((a \cup b) \subseteq x) \wedge \\ \quad (f \in (x \to \operatorname{carr}(r))) \wedge \\ \quad (|a| \in \mathbb{N}) \wedge \\ \quad (|b| \in \mathbb{N}) \Rightarrow \\ \quad \quad \operatorname{sum\text{-}set}(r, (a \cup b), f) \\ \quad \quad =\; \operatorname{add}(r)(\operatorname{sum\text{-}set}(r, a, f), \operatorname{sum\text{-}set}(r, b, f)) \end{array}}}

\vnbfile{theorem-library/rake-prod-set-defined.scm}{theorem-library/rake-prod-set-defined.scm}
\noindent{\small theorem-library/rake-prod-set-defined.scm -- PROD-SET DEFINED, and its four characterising laws PROVEN from the definition.  Rake batch 5c, assignment 5c-P (the user's decision of 2026-09-19).  Helper prefix `psd-'.}\par
\vnbentry{\texttt{prod-\allowbreak{}set-\allowbreak{}empty}}{\vnbstmt{\begin{array}{@{}l@{}} \forall m, f. \\ \quad (\operatorname{prod\text{-}set}(m, \emptyset, f) \simeq \operatorname{iden}(m)) \end{array}}}
\vnbentry{\texttt{finsum-\allowbreak{}cm-\allowbreak{}singleton-\allowbreak{}ptwise}}{\vnbstmt{\begin{array}{@{}l@{}} \forall m.; \forall x \in \mathbf{Set}, f. \\ \quad (\operatorname{is\text{-}comm\text{-}monoid}(m) \wedge (f(x) \in \operatorname{carr}(m))) \Rightarrow  \\ \quad \quad (\operatorname{finsum}(m, f, \{x\}) = f(x)) \end{array}}}
\vnbentry{\texttt{prod-\allowbreak{}set-\allowbreak{}singleton-\allowbreak{}defined}}{\vnbstmt{\begin{array}{@{}l@{}} \forall m.; \forall x \in \mathbf{Set}, f. \\ \quad (\operatorname{is\text{-}comm\text{-}monoid}(m) \wedge (f(x) \in \operatorname{carr}(m))) \Rightarrow  \\ \quad \quad (\operatorname{prod\text{-}set}(m, \{x\}, f) = f(x)) \end{array}}}
\vnbentry{\texttt{prod-\allowbreak{}set-\allowbreak{}type-\allowbreak{}defined}}{\vnbstmt{\begin{array}{@{}l@{}} \forall m, x, s, f. \\ \quad \operatorname{is\text{-}comm\text{-}monoid}(m) \wedge \\ \quad (x \in \mathbf{Set}) \wedge \\ \quad (s \subseteq x) \wedge \\ \quad (f \in (x \to \operatorname{carr}(m))) \wedge \\ \quad (|s| \in \mathbb{N}) \Rightarrow \\ \quad \quad (\operatorname{prod\text{-}set}(m, s, f) \in \operatorname{carr}(m)) \end{array}}}
\vnbentry{\texttt{prod-\allowbreak{}set-\allowbreak{}disjoint-\allowbreak{}union-\allowbreak{}defined}}{\vnbstmt{\begin{array}{@{}l@{}} \forall m, a, b, f. \\ \quad \operatorname{is\text{-}comm\text{-}monoid}(m) \wedge \\ \quad (a \in \mathbf{Set}) \wedge \\ \quad (b \in \mathbf{Set}) \wedge \\ \quad ((a \cap b) = \emptyset) \wedge \\ \quad (f \in ((a \cup b) \to \operatorname{carr}(m))) \wedge \\ \quad (|a| \in \mathbb{N}) \wedge \\ \quad (|b| \in \mathbb{N}) \Rightarrow \\ \quad \quad \operatorname{prod\text{-}set}(m, (a \cup b), f) \\ \quad \quad =\; \operatorname{opr}(m)(\operatorname{prod\text{-}set}(m, a, f), \operatorname{prod\text{-}set}(m, b, f)) \end{array}}}

\vnbfile{theorem-library/rake-sum-set-scalar.scm}{theorem-library/rake-sum-set-scalar.scm}
\noindent{\small rake-sum-set-scalar.scm -- the two scalar pull-outs for SUM-SET, PROVEN.}\par
\vnbentry{\texttt{sum-\allowbreak{}set-\allowbreak{}left-\allowbreak{}scalar}}{\vnbstmt{\begin{array}{@{}l@{}} \forall s.; \forall a \in \operatorname{carr}(s), x \in \mathbf{Set}, f \in (x \to \operatorname{carr}(s)). \\ \quad (\operatorname{is\text{-}ring}(s) \wedge (|x| \in \mathbb{N})) \Rightarrow  \\ \quad \quad \operatorname{mul}(s)(a, \operatorname{sum\text{-}set}(s, x, f)) \\ \quad \quad =\; \operatorname{sum\text{-}set}(s, x, (\lambda z \in x.\; \operatorname{mul}(s)(a, f(z)))) \end{array}}}
\vnbentry{\texttt{sum-\allowbreak{}set-\allowbreak{}right-\allowbreak{}scalar}}{\vnbstmt{\begin{array}{@{}l@{}} \forall s.; \forall b \in \operatorname{carr}(s), x \in \mathbf{Set}, f \in (x \to \operatorname{carr}(s)). \\ \quad (\operatorname{is\text{-}ring}(s) \wedge (|x| \in \mathbb{N})) \Rightarrow  \\ \quad \quad \operatorname{mul}(s)(\operatorname{sum\text{-}set}(s, x, f), b) \\ \quad \quad =\; \operatorname{sum\text{-}set}(s, x, (\lambda z \in x.\; \operatorname{mul}(s)(f(z), b))) \end{array}}}

\vnbfile{theorem-library/fin-subset-monoid.scm}{theorem-library/fin-subset-monoid.scm}
\noindent{\small fin-subset-monoid.scm -- FIN-SUBSETS(a) is a commutative monoid under union.}\par
\vnbentry{\texttt{fsm-\allowbreak{}carr}}{\vnbstmt{\begin{array}{@{}l@{}} \forall a. \\ \quad (\operatorname{carr}(\operatorname{fin\text{-}subset\text{-}monoid}(a)) \simeq \operatorname{fin\text{-}subsets}(a)) \end{array}}}
\vnbentry{\texttt{fsm-\allowbreak{}opr}}{\vnbstmt{\begin{array}{@{}l@{}} \forall a. \\ \quad (\operatorname{opr}(\operatorname{fin\text{-}subset\text{-}monoid}(a)) \simeq (\lambda \langle \operatorname{x\_}, \operatorname{y\_} \rangle \in (\operatorname{fin\text{-}subsets}(a) \times \operatorname{fin\text{-}subsets}(a)).\; (\operatorname{x\_} \cup \operatorname{y\_}))) \end{array}}}
\vnbentry{\texttt{fsm-\allowbreak{}iden}}{\vnbstmt{\begin{array}{@{}l@{}} \forall a. \\ \quad (\operatorname{iden}(\operatorname{fin\text{-}subset\text{-}monoid}(a)) \simeq \emptyset) \end{array}}}
\vnbentry{\texttt{fin-\allowbreak{}subset-\allowbreak{}monoid-\allowbreak{}is-\allowbreak{}comm-\allowbreak{}monoid}}{\vnbstmt{\begin{array}{@{}l@{}} \forall a \in \mathbf{Set}. \\ \quad \operatorname{is\text{-}comm\text{-}monoid}(\operatorname{fin\text{-}subset\text{-}monoid}(a)) \end{array}}}

\vnbfile{theorem-library/nn-pairing.scm}{theorem-library/nn-pairing.scm}
\noindent{\small nn-pairing.scm -- the Cantor pairing NN x NN -$>$ NN, and the arithmetic it needs.  Step one of the missing re-indexing mechanism.}\par
\vnbentry{\texttt{trinum-\allowbreak{}type}}{\vnbstmt{\forall n \in \mathbb{N}.\; (\operatorname{trinum}(n) \in \mathbb{N})}}
\vnbentry{\texttt{trinum-\allowbreak{}mono}}{\vnbstmt{\begin{array}{@{}l@{}} \forall b, a \in \mathbb{N}. \\ \quad (a \leq b) \Rightarrow  \\ \quad \quad (\operatorname{trinum}(a) \leq \operatorname{trinum}(b)) \end{array}}}
\vnbentry{\texttt{nnpair-\allowbreak{}onto}}{\vnbstmt{\begin{array}{@{}l@{}} \forall n \in \mathbb{N}. \exists i, j \in \mathbb{N}. \\ \quad (\operatorname{nnpair}(i, j) = n) \end{array}}}
\vnbentry{\texttt{nnpair-\allowbreak{}diag-\allowbreak{}bound}}{\vnbstmt{\begin{array}{@{}l@{}} \forall s, j. \\ \quad ((s \in \mathbb{N}) \wedge ((j \in \mathbb{N}) \wedge (j \leq s))) \Rightarrow  \\ \quad \quad (((\operatorname{trinum}(s) + j) + 1) \leq \operatorname{trinum}((s + 1))) \end{array}}}
\vnbentry{\texttt{nnpair-\allowbreak{}cross}}{\vnbstmt{\begin{array}{@{}l@{}} \forall s, t, j, k. \\ \quad (s \in \mathbb{N}) \wedge \\ \quad (t \in \mathbb{N}) \wedge \\ \quad (j \in \mathbb{N}) \wedge \\ \quad (k \in \mathbb{N}) \wedge \\ \quad (j \leq s) \wedge \\ \quad ((s + 1) \leq t) \Rightarrow \\ \quad \quad (((\operatorname{trinum}(s) + j) + 1) \leq (\operatorname{trinum}(t) + k)) \end{array}}}
\vnbentry{\texttt{nnpair-\allowbreak{}inj}}{\vnbstmt{\begin{array}{@{}l@{}} \forall i, j, \mathit{ii}, \mathit{jj}. \\ \quad (i \in \mathbb{N}) \wedge \\ \quad (j \in \mathbb{N}) \wedge \\ \quad (\mathit{ii} \in \mathbb{N}) \wedge \\ \quad (\mathit{jj} \in \mathbb{N}) \wedge \\ \quad (\operatorname{nnpair}(i, j) = \operatorname{nnpair}(\mathit{ii}, \mathit{jj})) \Rightarrow \\ \quad \quad ((i = \mathit{ii}) \wedge (j = \mathit{jj})) \end{array}}}
\vnbentry{\texttt{nnpair-\allowbreak{}type}}{\vnbstmt{\begin{array}{@{}l@{}} \forall i, j. \\ \quad ((i \in \mathbb{N}) \wedge (j \in \mathbb{N})) \Rightarrow  \\ \quad \quad (\operatorname{nnpair}(i, j) \in \mathbb{N}) \end{array}}}
\vnbentry{\texttt{nnfst-\allowbreak{}nnpair}}{\vnbstmt{\begin{array}{@{}l@{}} \forall i, j. \\ \quad ((i \in \mathbb{N}) \wedge (j \in \mathbb{N})) \Rightarrow  \\ \quad \quad (\operatorname{nnfst}(\operatorname{nnpair}(i, j)) = i) \end{array}}}
\vnbentry{\texttt{nnsnd-\allowbreak{}nnpair}}{\vnbstmt{\begin{array}{@{}l@{}} \forall i, j. \\ \quad ((i \in \mathbb{N}) \wedge (j \in \mathbb{N})) \Rightarrow  \\ \quad \quad (\operatorname{nnsnd}(\operatorname{nnpair}(i, j)) = j) \end{array}}}
\vnbentry{\texttt{nnfst-\allowbreak{}type}}{\vnbstmt{\forall n \in \mathbb{N}.\; (\operatorname{nnfst}(n) \in \mathbb{N})}}
\vnbentry{\texttt{nnsnd-\allowbreak{}type}}{\vnbstmt{\forall n \in \mathbb{N}.\; (\operatorname{nnsnd}(n) \in \mathbb{N})}}
\vnbentry{\texttt{nn-\allowbreak{}flatten}}{\vnbstmt{\begin{array}{@{}l@{}} \forall a.; \forall h \in (\mathbb{N} \to (\mathbb{N} \to a)). \exists e \in (\mathbb{N} \to a). \forall i, j. \\ \quad ((i \in \mathbb{N}) \wedge (j \in \mathbb{N})) \Rightarrow  \\ \quad \quad \exists m \in \mathbb{N}.\; (e(m) = h(i)(j)) \end{array}}}

\vnbfile{theorem-library/ball-cover-lemmas.scm}{theorem-library/ball-cover-lemmas.scm}
\noindent{\small ball-cover-lemmas.scm -- the membership plumbing under finite-ball-subcover-r-net, PROVEN (2026-09-14).}\par
\vnbentry{\texttt{centres-\allowbreak{}unfold}}{\vnbstmt{\begin{array}{@{}l@{}} \forall s, b, r. \\ \quad (\operatorname{centres}(s, b, r) \simeq \{c \in \operatorname{pts}(s) : (\operatorname{ball}(s, c, r) = b)\}) \end{array}}}
\vnbentry{\texttt{ball-\allowbreak{}cover-\allowbreak{}unfold}}{\vnbstmt{\begin{array}{@{}l@{}} \forall s, r. \\ \quad (\operatorname{ball\text{-}cover}(s, r) \simeq \operatorname{image}((\lambda c \in \operatorname{pts}(s).\; \operatorname{ball}(s, c, r)), \operatorname{pts}(s))) \end{array}}}
\vnbentry{\texttt{centre-\allowbreak{}set-\allowbreak{}unfold}}{\vnbstmt{\begin{array}{@{}l@{}} \forall s, r, f. \\ \quad (\operatorname{centre\text{-}set}(s, r, f) \simeq \operatorname{image}((\lambda b \in f.\; \varepsilon(\operatorname{centres}(s, b, r))), f)) \end{array}}}
\vnbentry{\texttt{ball-\allowbreak{}sep-\allowbreak{}unfold}}{\vnbstmt{\begin{array}{@{}l@{}} \forall s, c, r. \\ \quad (\operatorname{ball}(s, c, r) \simeq \{y \in \operatorname{pts}(s) : ((\operatorname{dist}(s)(c, y) \leq r) \wedge \neg (\operatorname{dist}(s)(c, y) = r))\}) \end{array}}}
\vnbentry{\texttt{centres-\allowbreak{}mem-\allowbreak{}build}}{\vnbstmt{\begin{array}{@{}l@{}} \forall s, b, r.; \forall c \in \operatorname{pts}(s). \\ \quad (\operatorname{ball}(s, c, r) = b) \Rightarrow  \\ \quad \quad (c \in \operatorname{centres}(s, b, r)) \end{array}}}
\vnbentry{\texttt{centres-\allowbreak{}in-\allowbreak{}carrier}}{\vnbstmt{\begin{array}{@{}l@{}} \forall s, b, r.; \forall c \in \operatorname{centres}(s, b, r). \\ \quad (c \in \operatorname{pts}(s)) \end{array}}}
\vnbentry{\texttt{centres-\allowbreak{}ball-\allowbreak{}eq}}{\vnbstmt{\begin{array}{@{}l@{}} \forall s, b, r.; \forall c \in \operatorname{centres}(s, b, r). \\ \quad (\operatorname{ball}(s, c, r) = b) \end{array}}}
\vnbentry{\texttt{ball-\allowbreak{}cover-\allowbreak{}mem-\allowbreak{}fwd}}{\vnbstmt{\begin{array}{@{}l@{}} \forall s, r.; \forall u \in \operatorname{ball\text{-}cover}(s, r). \exists c \in \operatorname{pts}(s). \\ \quad (\operatorname{ball}(s, c, r) = u) \end{array}}}
\vnbentry{\texttt{open-\allowbreak{}cover-\allowbreak{}covers-\allowbreak{}point}}{\vnbstmt{\begin{array}{@{}l@{}} \forall s, f.; \forall p \in \operatorname{pts}(s). \\ \quad \operatorname{is\text{-}open\text{-}cover}(s, f) \Rightarrow  \\ \quad \quad \exists u \in f.\; (p \in u) \end{array}}}
\vnbentry{\texttt{ball-\allowbreak{}point-\allowbreak{}le}}{\vnbstmt{\begin{array}{@{}l@{}} \forall s, c, r, w.; \forall p \in w. \\ \quad (\operatorname{ball}(s, c, r) = w) \Rightarrow  \\ \quad \quad (\operatorname{dist}(s)(c, p) \leq r) \end{array}}}
\vnbentry{\texttt{ball-\allowbreak{}point-\allowbreak{}ne}}{\vnbstmt{\begin{array}{@{}l@{}} \forall s, c, r, w.; \forall p \in w. \\ \quad (\operatorname{ball}(s, c, r) = w) \Rightarrow  \\ \quad \quad \neg (\operatorname{dist}(s)(c, p) = r) \end{array}}}

\vnbfile{calculus/finite-ball-subcover-proof.scm}{calculus/finite-ball-subcover-proof.scm}
\noindent{\small finite-ball-subcover-proof.scm -- machine proof of the centre-extraction lemma  finite-ball-subcover-r-net  (calculus.pdf Prop 3.12, the step that closes compact =$>$ totally bounded), with the choice of centres made EXPLICIT via the global Hilbert epsilon.  Run:}\par
\vnbentry{\texttt{chosen-\allowbreak{}centre-\allowbreak{}is-\allowbreak{}centre}}{\vnbstmt{\begin{array}{@{}l@{}} \forall s, r.; \forall u \in \operatorname{ball\text{-}cover}(s, r). \\ \quad (\varepsilon(\operatorname{centres}(s, u, r)) \in \operatorname{pts}(s)) \wedge \\ \quad (\operatorname{ball}(s, \varepsilon(\operatorname{centres}(s, u, r)), r) = u) \end{array}}}
\vnbentry{\texttt{finite-\allowbreak{}ball-\allowbreak{}subcover-\allowbreak{}r-\allowbreak{}net}}{\vnbstmt{\begin{array}{@{}l@{}} \forall s.; \forall r \in \mathbb{R}. \\ \quad \operatorname{is\text{-}metric\text{-}space}(s) \wedge \\ \quad (0 \leq r) \wedge \\ \quad \neg (0 = r) \wedge \\ \quad \exists f. \\ \quad \quad (f \subseteq \operatorname{ball\text{-}cover}(s, r)) \wedge \\ \quad \quad (|f| \in \mathbb{N}) \wedge \\ \quad \quad \operatorname{is\text{-}open\text{-}cover}(s, f) \Rightarrow \\ \quad \quad \exists n. \\ \quad \quad \quad ((|n| \in \mathbb{N}) \wedge \operatorname{is\text{-}r\text{-}net}(s, n, \operatorname{pts}(s), r)) \end{array}}}

\vnbfile{theorem-library/rake-algebra3.scm}{theorem-library/rake-algebra3.scm}
\noindent{\small theorem-library/rake-algebra3.scm -- rake batch 5, batch V (algebra leftovers).}\par
\vnbentry{\texttt{qq-\allowbreak{}nonneg-\allowbreak{}is-\allowbreak{}fraction}}{\vnbstmt{\begin{array}{@{}l@{}} \forall q \in \mathbb{Q}. \\ \quad (0 \leq q) \Rightarrow  \\ \quad \quad \exists a, b \in \mathbb{N}.\; (\neg (b = 0) \wedge ((q \cdot b) = a)) \end{array}}}
\vnbentry{\texttt{card-\allowbreak{}bijection-\allowbreak{}eq}}{\vnbstmt{\begin{array}{@{}l@{}} \forall m, d, i. \\ \quad ((m \in \mathbf{Set}) \wedge ((|m| \in \mathbb{N}) \wedge (i \in \mathrm{Bij}(m, d)))) \Rightarrow  \\ \quad \quad (|d| = |m|) \end{array}}}
\vnbentry{\texttt{finsum-\allowbreak{}reindex}}{\vnbstmt{\begin{array}{@{}l@{}} \forall m.; \forall s, t \in \mathbf{Set}, i \in \mathrm{Bij}(t, s), f \in (s \to \operatorname{carr}(m)). \\ \quad \operatorname{is\text{-}comm\text{-}monoid}(m) \wedge \\ \quad (|s| \in \mathbb{N}) \wedge \\ \quad (|t| \in \mathbb{N}) \Rightarrow \\ \quad \quad (\operatorname{finsum}(m, f, s) = \operatorname{finsum}(m, (\lambda z \in t.\; f(i(z))), t)) \end{array}}}
\vnbentry{\texttt{finsum-\allowbreak{}reindex-\allowbreak{}ag}}{\vnbstmt{\begin{array}{@{}l@{}} \forall g.; \forall s, t \in \mathbf{Set}, i \in \mathrm{Bij}(t, s), f \in (s \to \operatorname{carr}(g)). \\ \quad \operatorname{is\text{-}abelian\text{-}group}(g) \wedge \\ \quad (|s| \in \mathbb{N}) \wedge \\ \quad (|t| \in \mathbb{N}) \Rightarrow \\ \quad \quad (\operatorname{finsum}(g, f, s) = \operatorname{finsum}(g, (\lambda z \in t.\; f(i(z))), t)) \end{array}}}

\vnbfile{structure-library/metric-laws.scm}{structure-library/metric-laws.scm}
\noindent{\small metric-laws.scm -- the five metric laws, PROVEN from the definition.}\par
\vnbentry{\texttt{metric-\allowbreak{}pos}}{\vnbstmt{\begin{array}{@{}l@{}} \forall s.; \forall x, y \in \operatorname{pts}(s). \\ \quad \operatorname{is\text{-}metric\text{-}space}(s) \Rightarrow  \\ \quad \quad (0 \leq \operatorname{dist}(s)(x, y)) \end{array}}}
\vnbentry{\texttt{metric-\allowbreak{}self-\allowbreak{}zero}}{\vnbstmt{\begin{array}{@{}l@{}} \forall s.; \forall x \in \operatorname{pts}(s). \\ \quad \operatorname{is\text{-}metric\text{-}space}(s) \Rightarrow  \\ \quad \quad (\operatorname{dist}(s)(x, x) = 0) \end{array}}}
\vnbentry{\texttt{metric-\allowbreak{}zero-\allowbreak{}eq}}{\vnbstmt{\begin{array}{@{}l@{}} \forall s.; \forall x, y \in \operatorname{pts}(s). \\ \quad (\operatorname{is\text{-}metric\text{-}space}(s) \wedge (\operatorname{dist}(s)(x, y) = 0)) \Rightarrow  \\ \quad \quad (x = y) \end{array}}}
\vnbentry{\texttt{metric-\allowbreak{}sym}}{\vnbstmt{\begin{array}{@{}l@{}} \forall s.; \forall x, y \in \operatorname{pts}(s). \\ \quad \operatorname{is\text{-}metric\text{-}space}(s) \Rightarrow  \\ \quad \quad (\operatorname{dist}(s)(x, y) = \operatorname{dist}(s)(y, x)) \end{array}}}
\vnbentry{\texttt{metric-\allowbreak{}triangle}}{\vnbstmt{\begin{array}{@{}l@{}} \forall s.; \forall x, y, z \in \operatorname{pts}(s). \\ \quad \operatorname{is\text{-}metric\text{-}space}(s) \Rightarrow  \\ \quad \quad (\operatorname{dist}(s)(x, z) \leq (\operatorname{dist}(s)(x, y) + \operatorname{dist}(s)(y, z))) \end{array}}}

\vnbfile{theorem-library/ball-is-open.scm}{theorem-library/ball-is-open.scm}
\noindent{\small ball-is-open.scm -- an open ball is an open set.}\par
\vnbentry{\texttt{ball-\allowbreak{}is-\allowbreak{}open}}{\vnbstmt{\begin{array}{@{}l@{}} \forall s.; \forall x \in \operatorname{pts}(s), r \in \mathbb{R}. \\ \quad (\operatorname{is\text{-}metric\text{-}space}(s) \wedge ((0 \leq r) \wedge \neg (0 = r))) \Rightarrow  \\ \quad \quad \operatorname{is\text{-}open}(s, \operatorname{ball}(s, x, r)) \end{array}}}

\vnbfile{theorem-library/rake-compact-complete.scm}{theorem-library/rake-compact-complete.scm}
\noindent{\small rake-compact-complete.scm -- two of the three asserted leaves on the bill of compact-implies-complete (calculus/compact-complete-proof.scm), PROVEN.}\par
\vnbentry{\texttt{cauchy-\allowbreak{}seq-\allowbreak{}is-\allowbreak{}fun}}{\vnbstmt{\begin{array}{@{}l@{}} \forall s, f. \\ \quad \operatorname{is\text{-}cauchy\text{-}seq}(s, f) \Rightarrow  \\ \quad \quad (f \in (\mathbb{N} \to \operatorname{pts}(s))) \end{array}}}
\vnbentry{\texttt{cauchy-\allowbreak{}cluster-\allowbreak{}converges}}{\vnbstmt{\begin{array}{@{}l@{}} \forall s, f. \\ \quad \operatorname{is\text{-}cauchy\text{-}seq}(s, f) \wedge \\ \quad \exists x.\; \operatorname{cluster\text{-}point}(s, f, x) \Rightarrow \\ \quad \quad \operatorname{converges}(s, f) \end{array}}}

\vnbfile{theorem-library/rake-compact-cluster.scm}{theorem-library/rake-compact-cluster.scm}
\noindent{\small rake-compact-cluster.scm -- a compact metric space is sequentially compact.}\par
\vnbentry{\texttt{compact-\allowbreak{}seq-\allowbreak{}has-\allowbreak{}cluster}}{\vnbstmt{\begin{array}{@{}l@{}} \forall s, f. \\ \quad (\operatorname{is\text{-}compact}(s) \wedge (f \in (\mathbb{N} \to \operatorname{pts}(s)))) \Rightarrow  \\ \quad \quad \exists x.\; \operatorname{cluster\text{-}point}(s, f, x) \end{array}}}

\vnbfile{calculus/compact-complete-proof.scm}{calculus/compact-complete-proof.scm}
\noindent{\small compact-complete-proof.scm -- machine proof of  compact =$>$ complete (calculus.pdf Prop 3.12, the other half of (1)=$>$(4); the totally-bounded half is calculus/compact-tb-proof.scm).  Run: ./prover calculus/compact-complete-proof.scm}\par
\vnbentry{\texttt{compact-\allowbreak{}implies-\allowbreak{}complete}}{\vnbstmt{\begin{array}{@{}l@{}} \forall s. \\ \quad \operatorname{is\text{-}compact}(s) \Rightarrow  \\ \quad \quad \operatorname{is\text{-}complete}(s) \end{array}}}

\vnbfile{theorem-library/rake-balls.scm}{theorem-library/rake-balls.scm}
\noindent{\small rake-balls.scm -- the open-ball read-offs of structure-library/metric-topology.scm, PROVEN.  Five facts that stood there as `support' + `warrant! ... 'proof' -- a tier that claims a machine-checked proof exists -- and had none:}\par
\vnbentry{\texttt{ball-\allowbreak{}membership}}{\vnbstmt{\begin{array}{@{}l@{}} \forall s, x, r, y. \\ \quad ((y \in \operatorname{ball}(s, x, r)) \Leftrightarrow ((y \in \operatorname{pts}(s)) \wedge ((\operatorname{dist}(s)(x, y) \leq r) \wedge \neg (\operatorname{dist}(s)(x, y) = r)))) \end{array}}}
\vnbentry{\texttt{ball-\allowbreak{}mem-\allowbreak{}from-\allowbreak{}le}}{\vnbstmt{\begin{array}{@{}l@{}} \forall s.; \forall x \in \operatorname{pts}(s), y, d, r. \\ \quad \operatorname{is\text{-}metric\text{-}space}(s) \wedge \\ \quad (y \in \operatorname{pts}(s)) \wedge \\ \quad (d \in \mathbb{R}) \wedge \\ \quad (r \in \mathbb{R}) \wedge \\ \quad (\operatorname{dist}(s)(x, y) \leq d) \wedge \\ \quad (d < r) \Rightarrow \\ \quad \quad (y \in \operatorname{ball}(s, x, r)) \end{array}}}
\vnbentry{\texttt{ball-\allowbreak{}is-\allowbreak{}set}}{\vnbstmt{\begin{array}{@{}l@{}} \forall s, x, r. \\ \quad \operatorname{is\text{-}metric\text{-}space}(s) \Rightarrow  \\ \quad \quad (\operatorname{ball}(s, x, r) \in \mathbf{Set}) \end{array}}}
\vnbentry{\texttt{ball-\allowbreak{}center-\allowbreak{}in}}{\vnbstmt{\begin{array}{@{}l@{}} \forall s.; \forall x \in \operatorname{pts}(s), r \in \mathbb{R}. \\ \quad (\operatorname{is\text{-}metric\text{-}space}(s) \wedge ((0 \leq r) \wedge \neg (0 = r))) \Rightarrow  \\ \quad \quad (x \in \operatorname{ball}(s, x, r)) \end{array}}}
\vnbentry{\texttt{ball-\allowbreak{}2r-\allowbreak{}triangle}}{\vnbstmt{\begin{array}{@{}l@{}} \forall s.; \forall x \in \operatorname{pts}(s), r \in \mathbb{R}, y, z \in \operatorname{ball}(s, x, r). \\ \quad (\operatorname{is\text{-}metric\text{-}space}(s) \wedge ((0 \leq r) \wedge \neg (0 = r))) \Rightarrow  \\ \quad \quad ((\operatorname{dist}(s)(y, z) \leq (r + r)) \wedge \neg (\operatorname{dist}(s)(y, z) = (r + r))) \end{array}}}
\vnbentry{\texttt{cauchy-\allowbreak{}block-\allowbreak{}estimate}}{\vnbstmt{\begin{array}{@{}l@{}} \forall s.; \forall c \in \operatorname{pts}(s), r, u.; \forall y, z \in u, d.; \forall a \in \mathbb{R}. \\ \quad \operatorname{is\text{-}metric\text{-}space}(s) \wedge \\ \quad \operatorname{pos\text{-}rr}(r) \wedge \\ \quad (u = \operatorname{ball}(s, c, r)) \wedge \\ \quad \operatorname{pos\text{-}rr}(d) \wedge \\ \quad (r \leq d) \wedge \\ \quad ((d + d) = a) \Rightarrow \\ \quad \quad (\operatorname{dist}(s)(y, z) \leq a) \end{array}}}

\vnbfile{theorem-library/metric-subspace-laws.scm}{theorem-library/metric-subspace-laws.scm}
\noindent{\small metric-subspace-laws.scm -- the laws of SUBSPACE-MS and RESTRICT (structure-library/metric-subspace.scm).  Everything here is PROVEN; the definition file asserts nothing.}\par
\vnbentry{\texttt{subspace-\allowbreak{}pts}}{\vnbstmt{\begin{array}{@{}l@{}} \forall s, a. \\ \quad (\operatorname{pts}(\operatorname{subspace\text{-}ms}(s, a)) \simeq a) \end{array}}}
\vnbentry{\texttt{subspace-\allowbreak{}dist-\allowbreak{}slot}}{\vnbstmt{\begin{array}{@{}l@{}} \forall s, a. \\ \quad (\operatorname{dist}(\operatorname{subspace\text{-}ms}(s, a)) \simeq (\lambda \langle \operatorname{sbu\_}, \operatorname{sbv\_} \rangle \in (a \times a).\; \operatorname{dist}(s)(\operatorname{sbu\_}, \operatorname{sbv\_}))) \end{array}}}
\vnbentry{\texttt{subspace-\allowbreak{}dist}}{\vnbstmt{\begin{array}{@{}l@{}} \forall s, a, u, v. \\ \quad ((u \in a) \wedge (v \in a)) \Rightarrow  \\ \quad \quad (\operatorname{dist}(\operatorname{subspace\text{-}ms}(s, a))(u, v) \simeq \operatorname{dist}(s)(u, v)) \end{array}}}
\vnbentry{\texttt{subspace-\allowbreak{}is-\allowbreak{}metric-\allowbreak{}space}}{\vnbstmt{\begin{array}{@{}l@{}} \forall s, a. \\ \quad (\operatorname{is\text{-}metric\text{-}space}(s) \wedge (a \subseteq \operatorname{pts}(s))) \Rightarrow  \\ \quad \quad \operatorname{is\text{-}metric\text{-}space}(\operatorname{subspace\text{-}ms}(s, a)) \end{array}}}
\vnbentry{\texttt{restrict-\allowbreak{}apply}}{\vnbstmt{\forall f, a.; \forall u \in a.\; (\operatorname{restrict}(f, a)(u) \simeq f(u))}}
\vnbentry{\texttt{restrict-\allowbreak{}in-\allowbreak{}fun}}{\vnbstmt{\begin{array}{@{}l@{}} \forall f, m, d, a. \\ \quad ((f \in (m \to d)) \wedge (a \subseteq m)) \Rightarrow  \\ \quad \quad (\operatorname{restrict}(f, a) \in (a \to d)) \end{array}}}
\vnbentry{\texttt{subspace-\allowbreak{}whole}}{\vnbstmt{\begin{array}{@{}l@{}} \forall s. \\ \quad (\operatorname{pts}(\operatorname{subspace\text{-}ms}(s, \operatorname{pts}(s))) \simeq \operatorname{pts}(s)) \wedge \\ \quad \forall u, v. \\ \quad \quad ((u \in \operatorname{pts}(s)) \wedge (v \in \operatorname{pts}(s))) \Rightarrow  \\ \quad \quad \quad (\operatorname{dist}(\operatorname{subspace\text{-}ms}(s, \operatorname{pts}(s)))(u, v) \simeq \operatorname{dist}(s)(u, v)) \end{array}}}
\vnbentry{\texttt{subspace-\allowbreak{}ball-\allowbreak{}membership}}{\vnbstmt{\begin{array}{@{}l@{}} \forall s, a.; \forall c \in a, r, y. \\ \quad (\operatorname{is\text{-}metric\text{-}space}(s) \wedge (a \subseteq \operatorname{pts}(s))) \Rightarrow  \\ \quad \quad ((y \in \operatorname{ball}(\operatorname{subspace\text{-}ms}(s, a), c, r)) \Leftrightarrow ((y \in \operatorname{ball}(s, c, r)) \wedge (y \in a))) \end{array}}}
\vnbentry{\texttt{subspace-\allowbreak{}ball}}{\vnbstmt{\begin{array}{@{}l@{}} \forall s, a.; \forall c \in a, r. \\ \quad (\operatorname{is\text{-}metric\text{-}space}(s) \wedge (a \subseteq \operatorname{pts}(s))) \Rightarrow  \\ \quad \quad \operatorname{ball}(\operatorname{subspace\text{-}ms}(s, a), c, r) \\ \quad \quad =\; (\operatorname{ball}(s, c, r) \cap a) \end{array}}}
\vnbentry{\texttt{subspace-\allowbreak{}open-\allowbreak{}trace}}{\vnbstmt{\begin{array}{@{}l@{}} \forall s, a, u. \\ \quad \operatorname{is\text{-}metric\text{-}space}(s) \wedge \\ \quad (a \subseteq \operatorname{pts}(s)) \wedge \\ \quad \operatorname{is\text{-}open}(s, u) \Rightarrow \\ \quad \quad \operatorname{is\text{-}open}(\operatorname{subspace\text{-}ms}(s, a), (u \cap a)) \end{array}}}
\vnbentry{\texttt{restrict-\allowbreak{}continuous-\allowbreak{}at}}{\vnbstmt{\begin{array}{@{}l@{}} \forall s, a, t, f.; \forall \mathit{pa} \in a. \\ \quad \operatorname{is\text{-}metric\text{-}space}(s) \wedge \\ \quad (a \subseteq \operatorname{pts}(s)) \wedge \\ \quad \operatorname{is\text{-}continuous\text{-}at}(s, t, f, \mathit{pa}) \Rightarrow \\ \quad \quad \operatorname{is\text{-}continuous\text{-}at}(\operatorname{subspace\text{-}ms}(s, a), t, \operatorname{restrict}(f, a), \mathit{pa}) \end{array}}}
\vnbentry{\texttt{subspace-\allowbreak{}converges-\allowbreak{}to}}{\vnbstmt{\begin{array}{@{}l@{}} \forall s, a.; \forall g \in (\mathbb{N} \to a), l \in a. \\ \quad (\operatorname{is\text{-}metric\text{-}space}(s) \wedge (a \subseteq \operatorname{pts}(s))) \Rightarrow  \\ \quad \quad (\operatorname{converges\text{-}to}(\operatorname{subspace\text{-}ms}(s, a), g, l) \Leftrightarrow \operatorname{converges\text{-}to}(s, g, l)) \end{array}}}
\vnbentry{\texttt{subspace-\allowbreak{}open-\allowbreak{}is-\allowbreak{}trace}}{\vnbstmt{\begin{array}{@{}l@{}} \forall s, a, \mathit{vs}. \\ \quad \operatorname{is\text{-}metric\text{-}space}(s) \wedge \\ \quad (a \subseteq \operatorname{pts}(s)) \wedge \\ \quad \operatorname{is\text{-}open}(\operatorname{subspace\text{-}ms}(s, a), \mathit{vs}) \Rightarrow \\ \quad \quad \exists \mathit{us}. \\ \quad \quad \quad (\operatorname{is\text{-}open}(s, \mathit{us}) \wedge (\mathit{vs} = (\mathit{us} \cap a))) \end{array}}}
\vnbentry{\texttt{closed-\allowbreak{}limit-\allowbreak{}in}}{\vnbstmt{\begin{array}{@{}l@{}} \forall s, a.; \forall f \in (\mathbb{N} \to a), l. \\ \quad \operatorname{is\text{-}metric\text{-}space}(s) \wedge \\ \quad \operatorname{is\text{-}closed}(s, a) \wedge \\ \quad \operatorname{converges\text{-}to}(s, f, l) \Rightarrow \\ \quad \quad (l \in a) \end{array}}}
\vnbentry{\texttt{subspace-\allowbreak{}is-\allowbreak{}cauchy-\allowbreak{}seq}}{\vnbstmt{\begin{array}{@{}l@{}} \forall s, a.; \forall f \in (\mathbb{N} \to a). \\ \quad (\operatorname{is\text{-}metric\text{-}space}(s) \wedge (a \subseteq \operatorname{pts}(s))) \Rightarrow  \\ \quad \quad (\operatorname{is\text{-}cauchy\text{-}seq}(\operatorname{subspace\text{-}ms}(s, a), f) \Leftrightarrow \operatorname{is\text{-}cauchy\text{-}seq}(s, f)) \end{array}}}
\vnbentry{\texttt{subspace-\allowbreak{}complete}}{\vnbstmt{\begin{array}{@{}l@{}} \forall s, a. \\ \quad (\operatorname{is\text{-}complete}(s) \wedge \operatorname{is\text{-}closed}(s, a)) \Rightarrow  \\ \quad \quad \operatorname{is\text{-}complete}(\operatorname{subspace\text{-}ms}(s, a)) \end{array}}}

\vnbfile{theorem-library/compact-subspace.scm}{theorem-library/compact-subspace.scm}
\noindent{\small compact-subspace.scm -- a CLOSED SUBSET of a COMPACT space is COMPACT as a subspace (structure-library/metric-subspace.scm, theorem-library/metric-subspace-laws.scm).}\par
\vnbentry{\texttt{open-\allowbreak{}in-\allowbreak{}power}}{\vnbstmt{\begin{array}{@{}l@{}} \forall s, u. \\ \quad (\operatorname{is\text{-}metric\text{-}space}(s) \wedge \operatorname{is\text{-}open}(s, u)) \Rightarrow  \\ \quad \quad (u \in \operatorname{power}(\operatorname{pts}(s))) \end{array}}}
\vnbentry{\texttt{closed-\allowbreak{}subset-\allowbreak{}of-\allowbreak{}compact-\allowbreak{}is-\allowbreak{}compact}}{\vnbstmt{\begin{array}{@{}l@{}} \forall s, a. \\ \quad (\operatorname{is\text{-}compact}(s) \wedge \operatorname{is\text{-}closed}(s, a)) \Rightarrow  \\ \quad \quad \operatorname{is\text{-}compact}(\operatorname{subspace\text{-}ms}(s, a)) \end{array}}}

\vnbfile{theorem-library/rake-analysis-typing.scm}{theorem-library/rake-analysis-typing.scm}
\noindent{\small rake-analysis-typing.scm -- BATCH H of the leaf rake (2026-09-17): the four ANALYSIS typing leaves.  Two are proven as stated; two are UNDERDETERMINED as stated -- the same defect that made `nth-deriv-v-in-vec' false -- and are proven here with the guard their scalar twins already carry.}\par
\vnbentry{\texttt{rkt-\allowbreak{}completion-\allowbreak{}pts}}{\vnbstmt{\begin{array}{@{}l@{}} \forall m. \\ \quad (\operatorname{pts}(\operatorname{completion}(m)) \simeq \operatorname{quotient}(\operatorname{cauchy\text{-}setoid}(m))) \end{array}}}
\vnbentry{\texttt{rkt-\allowbreak{}cauchy-\allowbreak{}setoid-\allowbreak{}pts}}{\vnbstmt{\begin{array}{@{}l@{}} \forall m. \\ \quad (\operatorname{pts}(\operatorname{cauchy\text{-}setoid}(m)) \simeq \operatorname{cseq}(m)) \end{array}}}
\vnbentry{\texttt{rkt-\allowbreak{}cseq-\allowbreak{}is-\allowbreak{}set}}{\vnbstmt{\begin{array}{@{}l@{}} \forall m. \\ \quad \operatorname{is\text{-}metric\text{-}space}(m) \Rightarrow  \\ \quad \quad (\operatorname{cseq}(m) \in \mathbf{Set}) \end{array}}}
\vnbentry{\texttt{rkt-\allowbreak{}class-\allowbreak{}is-\allowbreak{}set}}{\vnbstmt{\begin{array}{@{}l@{}} \forall m, w. \\ \quad \operatorname{is\text{-}metric\text{-}space}(m) \Rightarrow  \\ \quad \quad (\operatorname{class}(\operatorname{cauchy\text{-}setoid}(m), w) \in \mathbf{Set}) \end{array}}}
\vnbentry{\texttt{rkt-\allowbreak{}embed-\allowbreak{}seq-\allowbreak{}in-\allowbreak{}fun}}{\vnbstmt{\begin{array}{@{}l@{}} \forall m, u. \\ \quad (\operatorname{is\text{-}metric\text{-}space}(m) \wedge (u \in \operatorname{pts}(m))) \Rightarrow  \\ \quad \quad (\operatorname{embed\text{-}seq}(m, u) \in (\mathbb{N} \to \operatorname{pts}(m))) \end{array}}}
\vnbentry{\texttt{rkt-\allowbreak{}embed-\allowbreak{}seq-\allowbreak{}cauchy}}{\vnbstmt{\begin{array}{@{}l@{}} \forall m, u. \\ \quad (\operatorname{is\text{-}metric\text{-}space}(m) \wedge (u \in \operatorname{pts}(m))) \Rightarrow  \\ \quad \quad \operatorname{is\text{-}cauchy\text{-}seq}(m, \operatorname{embed\text{-}seq}(m, u)) \end{array}}}
\vnbentry{\texttt{rkt-\allowbreak{}embed-\allowbreak{}seq-\allowbreak{}in-\allowbreak{}cseq}}{\vnbstmt{\begin{array}{@{}l@{}} \forall m, u. \\ \quad (\operatorname{is\text{-}metric\text{-}space}(m) \wedge (u \in \operatorname{pts}(m))) \Rightarrow  \\ \quad \quad (\operatorname{embed\text{-}seq}(m, u) \in \operatorname{cseq}(m)) \end{array}}}
\vnbentry{\texttt{embed-\allowbreak{}in-\allowbreak{}fun}}{\vnbstmt{\begin{array}{@{}l@{}} \forall m. \\ \quad \operatorname{is\text{-}metric\text{-}space}(m) \Rightarrow  \\ \quad \quad \operatorname{embed}(m) \\ \quad \quad \in\; (\operatorname{pts}(m) \to \operatorname{pts}(\operatorname{completion}(m))) \end{array}}}

\vnbfile{theorem-library/rake-cont-preimage.scm}{theorem-library/rake-cont-preimage.scm}
\noindent{\small rake-cont-preimage.scm -- continuous-implies-open-preimage, PROVEN.}\par
\vnbentry{\texttt{continuous-\allowbreak{}implies-\allowbreak{}open-\allowbreak{}preimage}}{\vnbstmt{\begin{array}{@{}l@{}} \forall s, t, f, v. \\ \quad (\operatorname{is\text{-}continuous}(s, t, f) \wedge \operatorname{is\text{-}open}(t, v)) \Rightarrow  \\ \quad \quad \operatorname{is\text{-}open}(s, \operatorname{preimage}(s, f, v)) \end{array}}}

\vnbfile{theorem-library/pseudometric-laws.scm}{theorem-library/pseudometric-laws.scm}
\noindent{\small pseudometric-laws.scm -- the FOUR PSEUDOMETRIC LAWS, proven by projecting the `is-pseudometric' property, exactly as structure-library/metric-laws.scm projects `is-metric'.}\par
\vnbentry{\texttt{pseudometric-\allowbreak{}self-\allowbreak{}zero}}{\vnbstmt{\begin{array}{@{}l@{}} \forall t, r.; \forall u \in r. \\ \quad (\operatorname{is\text{-}pseudometric}(t, r) \Rightarrow (t(u, u) = 0)) \end{array}}}
\vnbentry{\texttt{pseudometric-\allowbreak{}pos}}{\vnbstmt{\begin{array}{@{}l@{}} \forall t, r.; \forall u, v \in r. \\ \quad (\operatorname{is\text{-}pseudometric}(t, r) \Rightarrow (0 \leq t(u, v))) \end{array}}}
\vnbentry{\texttt{pseudometric-\allowbreak{}sym}}{\vnbstmt{\begin{array}{@{}l@{}} \forall t, r.; \forall u, v \in r. \\ \quad (\operatorname{is\text{-}pseudometric}(t, r) \Rightarrow (t(u, v) = t(v, u))) \end{array}}}
\vnbentry{\texttt{pseudometric-\allowbreak{}triangle}}{\vnbstmt{\begin{array}{@{}l@{}} \forall t, r.; \forall u, v, w \in r. \\ \quad \operatorname{is\text{-}pseudometric}(t, r) \Rightarrow  \\ \quad \quad (t(u, w) \leq (t(u, v) + t(v, w))) \end{array}}}

\vnbfile{theorem-library/bdd-metric-distance.scm}{theorem-library/bdd-metric-distance.scm}
\noindent{\small bdd-metric-distance.scm -- the BDD-METRIC distance formula, proven.}\par
\vnbentry{\texttt{bdd-\allowbreak{}metric-\allowbreak{}distance}}{\vnbstmt{\begin{array}{@{}l@{}} \forall s.; \forall x, y \in \operatorname{pts}(s). \\ \quad \operatorname{is\text{-}metric\text{-}space}(s) \Rightarrow  \\ \quad \quad \operatorname{dist}(\operatorname{bdd\text{-}metric}(s))(x, y) \\ \quad \quad =\; /(\operatorname{dist}(s)(x, y), (1 + \operatorname{dist}(s)(x, y))) \end{array}}}

\vnbfile{theorem-library/bdd-metric-basics.scm}{theorem-library/bdd-metric-basics.scm}
\noindent{\small bdd-metric-basics.scm -- the t/(1+t) scalar family and the two BDD-METRIC statements that read off it, PROVEN.}\par
\vnbentry{\texttt{bdd-\allowbreak{}fn-\allowbreak{}nonneg}}{\vnbstmt{\forall t \in \mathbb{R}.\; ((0 \leq t) \Rightarrow (0 \leq /(t, (1 + t))))}}
\vnbentry{\texttt{bdd-\allowbreak{}fn-\allowbreak{}lt-\allowbreak{}one}}{\vnbstmt{\forall t \in \mathbb{R}.\; ((0 \leq t) \Rightarrow (/(t, (1 + t)) < 1))}}
\vnbentry{\texttt{bdd-\allowbreak{}fn-\allowbreak{}le-\allowbreak{}arg}}{\vnbstmt{\forall t \in \mathbb{R}.\; ((0 \leq t) \Rightarrow (/(t, (1 + t)) \leq t))}}
\vnbentry{\texttt{bdd-\allowbreak{}fn-\allowbreak{}zero}}{\vnbstmt{\begin{array}{@{}l@{}} \forall t \in \mathbb{R}. \\ \quad (((0 \leq t) \wedge (t = 0)) \Rightarrow (/(t, (1 + t)) = 0)) \end{array}}}
\vnbentry{\texttt{bdd-\allowbreak{}fn-\allowbreak{}zero-\allowbreak{}eq}}{\vnbstmt{\begin{array}{@{}l@{}} \forall t \in \mathbb{R}. \\ \quad (((0 \leq t) \wedge (/(t, (1 + t)) = 0)) \Rightarrow (t = 0)) \end{array}}}
\vnbentry{\texttt{bdd-\allowbreak{}fn-\allowbreak{}mono}}{\vnbstmt{\begin{array}{@{}l@{}} \forall s, t \in \mathbb{R}. \\ \quad (((0 \leq s) \wedge (s \leq t)) \Rightarrow (/(s, (1 + s)) \leq /(t, (1 + t)))) \end{array}}}
\vnbentry{\texttt{bdd-\allowbreak{}fn-\allowbreak{}subadd}}{\vnbstmt{\begin{array}{@{}l@{}} \forall a, b \in \mathbb{R}. \\ \quad ((0 \leq a) \wedge (0 \leq b)) \Rightarrow  \\ \quad \quad (/((a + b), (1 + (a + b))) \leq (/(a, (1 + a)) + /(b, (1 + b)))) \end{array}}}
\vnbentry{\texttt{bdd-\allowbreak{}metric-\allowbreak{}bounded}}{\vnbstmt{\begin{array}{@{}l@{}} \forall s.; \forall x, y \in \operatorname{pts}(s). \\ \quad \operatorname{is\text{-}metric\text{-}space}(s) \Rightarrow  \\ \quad \quad (\operatorname{dist}(\operatorname{bdd\text{-}metric}(s))(x, y) < 1) \end{array}}}
\vnbentry{\texttt{bdd-\allowbreak{}metric-\allowbreak{}is-\allowbreak{}metric-\allowbreak{}space}}{\vnbstmt{\begin{array}{@{}l@{}} \forall s. \\ \quad \operatorname{is\text{-}metric\text{-}space}(s) \Rightarrow  \\ \quad \quad \operatorname{is\text{-}metric\text{-}space}(\operatorname{bdd\text{-}metric}(s)) \end{array}}}

\vnbfile{theorem-library/rake-subseq-leaves.scm}{theorem-library/rake-subseq-leaves.scm}
\noindent{\small rake-subseq-leaves.scm -- the remaining asserted leaves of the subsequence / diagonalization cluster, PROVEN.}\par
\vnbentry{\texttt{nn-\allowbreak{}step-\allowbreak{}mono-\allowbreak{}ptwise}}{\vnbstmt{\begin{array}{@{}l@{}} \forall n \in \mathbb{N}, g.; \forall m \in \mathbb{N}. \\ \quad \forall k \in \mathbb{N}.\; (g(k) \in \mathbb{N}) \wedge \\ \quad \forall k \in \mathbb{N}.\; (g(k) < g((k + 1))) \wedge \\ \quad (m < n) \Rightarrow \\ \quad \quad (g(m) < g(n)) \end{array}}}
\vnbentry{\texttt{nn-\allowbreak{}step-\allowbreak{}strictly-\allowbreak{}mono}}{\vnbstmt{\begin{array}{@{}l@{}} \forall g \in (\mathbb{N} \to \mathbb{N}). \\ \quad \forall k \in \mathbb{N}.\; (g(k) < g((k + 1))) \Rightarrow  \\ \quad \quad \operatorname{strictly\text{-}mono\text{-}nn}(g) \end{array}}}
\vnbentry{\texttt{subsequence-\allowbreak{}capture}}{\vnbstmt{\begin{array}{@{}l@{}} \forall s \in \mathrm{Inf}(\mathbb{N}). \exists f \in (\mathbb{N} \to s). \forall m, n \in \mathbb{N}. \\ \quad ((m < n) \Rightarrow (f(m) < f(n))) \end{array}}}
\vnbentry{\texttt{nn-\allowbreak{}recip-\allowbreak{}succ-\allowbreak{}small-\allowbreak{}at}}{\vnbstmt{\begin{array}{@{}l@{}} \forall s.; \forall k \in \mathbb{N}. \\ \quad (\operatorname{pos\text{-}rr}(s) \wedge (\operatorname{recip}(s) < (k + 1))) \Rightarrow  \\ \quad \quad (\operatorname{recip}((k + 1)) < s) \end{array}}}
\vnbentry{\texttt{null-\allowbreak{}rr-\allowbreak{}seq-\allowbreak{}exists}}{\vnbstmt{\exists d.\; \operatorname{null\text{-}rr\text{-}seq}(d)}}

\vnbfile{theorem-library/rake-dc-consumers.scm}{theorem-library/rake-dc-consumers.scm}
\noindent{\small rake-dc-consumers.scm -- the consumers of recursive choice (rake batch 8-C).}\par
\vnbentry{\texttt{cauchy-\allowbreak{}rapid-\allowbreak{}subsequence}}{\vnbstmt{\begin{array}{@{}l@{}} \forall s, f, d. \\ \quad \operatorname{is\text{-}cauchy\text{-}seq}(s, f) \wedge \\ \quad (d \in (\mathbb{N} \to \mathbb{R})) \wedge \\ \quad \forall k \in \mathbb{N}.\; \operatorname{pos\text{-}rr}(d(k)) \Rightarrow \\ \quad \quad \exists i \in (\mathbb{N} \to \mathbb{N}). \\ \quad \quad \quad \forall m, n \in \mathbb{N}.\; ((m < n) \Rightarrow (i(m) < i(n))) \wedge \\ \quad \quad \quad \forall k \in \mathbb{N}. \\ \quad \quad \quad \quad (\operatorname{dist}(s)(f(i(k)), f(i((k + 1)))) \leq d(k)) \end{array}}}

\vnbfile{theorem-library/rake-compact-iff-seq-compact.scm}{theorem-library/rake-compact-iff-seq-compact.scm}
\noindent{\small rake-compact-iff-seq-compact.scm -- the FORWARD half of Prop 3.12 (1)=$>$(3'): a compact metric space is sequentially compact.}\par
\vnbentry{\texttt{cluster-\allowbreak{}point-\allowbreak{}has-\allowbreak{}convergent-\allowbreak{}subseq}}{\vnbstmt{\begin{array}{@{}l@{}} \forall s, f, x. \\ \quad \operatorname{cluster\text{-}point}(s, f, x) \Rightarrow  \\ \quad \quad \exists i. \\ \quad \quad \quad \operatorname{strictly\text{-}mono\text{-}nn}(i) \wedge \\ \quad \quad \quad \operatorname{converges\text{-}to}(s, \operatorname{subseq}(f, i), x) \end{array}}}
\vnbentry{\texttt{compact-\allowbreak{}implies-\allowbreak{}seq-\allowbreak{}compact}}{\vnbstmt{\begin{array}{@{}l@{}} \forall s. \\ \quad (\operatorname{is\text{-}metric\text{-}space}(s) \wedge \operatorname{is\text{-}compact}(s)) \Rightarrow  \\ \quad \quad \operatorname{seq\text{-}compact}(s) \end{array}}}

\vnbfile{theorem-library/rake-seq-compact-tb.scm}{theorem-library/rake-seq-compact-tb.scm}
\noindent{\small rake-seq-compact-tb.scm -- a sequentially compact metric space is totally bounded (the (C1) leg of Prop 3.12's missing direction; see the closing block of theorem-library/rake-compact-iff-seq-compact.scm).}\par
\vnbentry{\texttt{seq-\allowbreak{}compact-\allowbreak{}implies-\allowbreak{}totally-\allowbreak{}bounded}}{\vnbstmt{\begin{array}{@{}l@{}} \forall s. \\ \quad (\operatorname{is\text{-}metric\text{-}space}(s) \wedge \operatorname{seq\text{-}compact}(s)) \Rightarrow  \\ \quad \quad \operatorname{totally\text{-}bounded}(s) \end{array}}}

\vnbfile{theorem-library/rake-lebesgue-number.scm}{theorem-library/rake-lebesgue-number.scm}
\noindent{\small rake-lebesgue-number.scm -- the Lebesgue number lemma, and the CONVERSE half of Prop 3.12 (1)$<$=$>$(3'): a sequentially compact metric space is compact.}\par
\vnbentry{\texttt{lebesgue-\allowbreak{}number}}{\vnbstmt{\begin{array}{@{}l@{}} \forall s, v. \\ \quad \operatorname{is\text{-}metric\text{-}space}(s) \wedge \\ \quad \operatorname{seq\text{-}compact}(s) \wedge \\ \quad \operatorname{is\text{-}open\text{-}cover}(s, v) \Rightarrow \\ \quad \quad \exists l. \\ \quad \quad \quad \operatorname{pos\text{-}rr}(l) \wedge \\ \quad \quad \quad \forall p \in \operatorname{pts}(s). \exists u \in v. \\ \quad \quad \quad \quad (\operatorname{ball}(s, p, l) \subseteq u) \end{array}}}
\vnbentry{\texttt{seq-\allowbreak{}compact-\allowbreak{}implies-\allowbreak{}compact}}{\vnbstmt{\begin{array}{@{}l@{}} \forall s. \\ \quad (\operatorname{is\text{-}metric\text{-}space}(s) \wedge \operatorname{seq\text{-}compact}(s)) \Rightarrow  \\ \quad \quad \operatorname{is\text{-}compact}(s) \end{array}}}
\vnbentry{\texttt{compact-\allowbreak{}iff-\allowbreak{}seq-\allowbreak{}compact}}{\vnbstmt{\begin{array}{@{}l@{}} \forall s. \\ \quad \operatorname{is\text{-}metric\text{-}space}(s) \Rightarrow  \\ \quad \quad (\operatorname{is\text{-}compact}(s) \Leftrightarrow \operatorname{seq\text{-}compact}(s)) \end{array}}}

\vnbfile{theorem-library/rake-nn-enum.scm}{theorem-library/rake-nn-enum.scm}
\noindent{\small theorem-library/rake-nn-enum.scm ==================================================================== The chosen enumeration NN-ENUM and the two strictly-monotone bricks the refinement tower (convergence-block-tower) wants.}\par
\vnbentry{\texttt{nn-\allowbreak{}enum-\allowbreak{}spec}}{\vnbstmt{\begin{array}{@{}l@{}} \forall s \in \mathrm{Inf}(\mathbb{N}). \\ \quad (\operatorname{nn\text{-}enum}(s) \in (\mathbb{N} \to s)) \wedge \\ \quad \forall m, n \in \mathbb{N}. \\ \quad \quad (m < n) \Rightarrow  \\ \quad \quad \quad (\operatorname{nn\text{-}enum}(s)(m) < \operatorname{nn\text{-}enum}(s)(n)) \end{array}}}
\vnbentry{\texttt{strictly-\allowbreak{}mono-\allowbreak{}image-\allowbreak{}infinite}}{\vnbstmt{\begin{array}{@{}l@{}} \forall i, a. \\ \quad \operatorname{strictly\text{-}mono\text{-}nn}(i) \wedge \\ \quad (a \subseteq \mathbb{N}) \wedge \\ \quad \forall k \in \mathbb{N}.\; (i(k) \in a) \Rightarrow \\ \quad \quad (\{m \in a : \exists j \in \mathbb{N}.\; (m = i(j))\} \in \mathrm{Inf}(\mathbb{N})) \end{array}}}
\vnbentry{\texttt{strictly-\allowbreak{}mono-\allowbreak{}le-\allowbreak{}reflect}}{\vnbstmt{\begin{array}{@{}l@{}} \forall i.; \forall a, j \in \mathbb{N}. \\ \quad (\operatorname{strictly\text{-}mono\text{-}nn}(i) \wedge (i(a) \leq i(j))) \Rightarrow  \\ \quad \quad (a \leq j) \end{array}}}

\vnbfile{theorem-library/rake-tb-leaves.scm}{theorem-library/rake-tb-leaves.scm}
\noindent{\small rake-tb-leaves.scm -- the totally-bounded / pigeonhole leaves of the Cauchy-subsequence and separability arcs.  Two theorems, both modulo 0:}\par
\vnbentry{\texttt{cover-\allowbreak{}block-\allowbreak{}step}}{\vnbstmt{\begin{array}{@{}l@{}} \forall v \in \mathbf{Set}, f \in (\mathbb{N} \to v), c.; \forall j \in \mathrm{Inf}(\mathbb{N}). \\ \quad \operatorname{is\text{-}finite\text{-}cover}(c, v) \Rightarrow  \\ \quad \quad \exists a \in \mathrm{Inf}(\mathbb{N}). \\ \quad \quad \quad ((a \subseteq j) \wedge \exists u \in c.\; \forall i \in a.\; (f(i) \in u)) \end{array}}}
\vnbentry{\texttt{dist-\allowbreak{}le-\allowbreak{}implies-\allowbreak{}in-\allowbreak{}carrier}}{\vnbstmt{\begin{array}{@{}l@{}} \forall s, c, p, r. \\ \quad (\operatorname{is\text{-}metric\text{-}space}(s) \wedge (\operatorname{dist}(s)(c, p) \leq r)) \Rightarrow  \\ \quad \quad (c \in \operatorname{pts}(s)) \end{array}}}

\vnbfile{theorem-library/rake-tb-leaves-2.scm}{theorem-library/rake-tb-leaves-2.scm}
\noindent{\small rake-tb-leaves-2.scm -- the two totally-bounded leaves that the IS-R-NET repair of 2026-09-19 unblocked.}\par
\vnbentry{\texttt{tb-\allowbreak{}scale-\allowbreak{}dense-\allowbreak{}seq}}{\vnbstmt{\begin{array}{@{}l@{}} \forall s.; \forall a \in \operatorname{pts}(s), r. \\ \quad (\operatorname{totally\text{-}bounded}(s) \wedge \operatorname{pos\text{-}rr}(r)) \Rightarrow  \\ \quad \quad \exists g \in (\mathbb{N} \to \operatorname{pts}(s)). \forall p \in \operatorname{pts}(s). \exists j \in \mathbb{N}. \\ \quad \quad \quad (\operatorname{dist}(s)(p, g(j)) < r) \end{array}}}
\vnbentry{\texttt{tb-\allowbreak{}rad-\allowbreak{}ball-\allowbreak{}cover}}{\vnbstmt{\begin{array}{@{}l@{}} \forall s, d. \\ \quad \operatorname{totally\text{-}bounded}(s) \wedge \\ \quad \forall k \in \mathbb{N}.\; \operatorname{pos\text{-}rr}(d(k)) \Rightarrow \\ \quad \quad (\operatorname{pts}(s) \in \mathbf{Set}) \wedge \\ \quad \quad \exists v. \\ \quad \quad \quad \forall k \in \mathbb{N}. \\ \quad \quad \quad \quad \operatorname{is\text{-}finite\text{-}cover}(v(k), \operatorname{pts}(s)) \wedge \\ \quad \quad \quad \forall k \in \mathbb{N}, u \in v(k). \exists c \in \operatorname{pts}(s). \\ \quad \quad \quad \quad (u = \operatorname{ball}(s, c, d(k))) \end{array}}}

\vnbfile{theorem-library/nn-nested-subset-chain-proof.scm}{theorem-library/nn-nested-subset-chain-proof.scm}
\noindent{\small nn-nested-subset-chain-proof.scm -- a descending NN-indexed family is a chain.}\par
\vnbentry{\texttt{nn-\allowbreak{}nested-\allowbreak{}subset-\allowbreak{}chain-\allowbreak{}j}}{\vnbstmt{\begin{array}{@{}l@{}} \forall m.; \forall j, k \in \mathbb{N}. \\ \quad (\forall k \in \mathbb{N}.\; (m((k + 1)) \subseteq m(k)) \wedge (k \leq j)) \Rightarrow  \\ \quad \quad (m(j) \subseteq m(k)) \end{array}}}
\vnbentry{\texttt{nn-\allowbreak{}nested-\allowbreak{}subset-\allowbreak{}chain}}{\vnbstmt{\begin{array}{@{}l@{}} \forall t.; \forall k, j \in \mathbb{N}. \\ \quad (\forall k \in \mathbb{N}.\; (t((k + 1)) \subseteq t(k)) \wedge (k \leq j)) \Rightarrow  \\ \quad \quad (t(j) \subseteq t(k)) \end{array}}}

\vnbfile{theorem-library/diagonalization.scm}{theorem-library/diagonalization.scm}
\noindent{\small theorem-library/diagonalization.scm}\par
\vnbentry{\texttt{diagonalization}}{\vnbstmt{\begin{array}{@{}l@{}} \forall s \in (\mathbb{N} \to \mathrm{Inf}(\mathbb{N})). \\ \quad \forall k \in \mathbb{N}.\; (s((k + 1)) \subseteq s(k)) \Rightarrow  \\ \quad \quad \exists f. \\ \quad \quad \quad \operatorname{strictly\text{-}mono\text{-}nn}(f) \wedge \\ \quad \quad \quad \forall k, j \in \mathbb{N}.\; ((k \leq j) \Rightarrow (f(j) \in s(k))) \end{array}}}

\vnbfile{theorem-library/block-family-combinatorial-proof.scm}{theorem-library/block-family-combinatorial-proof.scm}
\noindent{\small theorem-library/block-family-combinatorial-proof.scm}\par
\vnbentry{\texttt{block-\allowbreak{}family-\allowbreak{}combinatorial}}{\vnbstmt{\begin{array}{@{}l@{}} \forall v \in \mathbf{Set}, f \in (\mathbb{N} \to v), a. \\ \quad \forall k \in \mathbb{N}.\; \operatorname{is\text{-}finite\text{-}cover}(a(k), v) \Rightarrow  \\ \quad \quad \exists b \in (\mathbb{N} \to \mathrm{Inf}(\mathbb{N})). \\ \quad \quad \quad \forall k \in \mathbb{N}.\; (b((k + 1)) \subseteq b(k)) \wedge \\ \quad \quad \quad \forall k \in \mathbb{N}. \exists u \in a(k). \forall i \in b(k). \\ \quad \quad \quad \quad (f(i) \in u) \end{array}}}

\vnbfile{theorem-library/cauchy-subseq-proof.scm}{theorem-library/cauchy-subseq-proof.scm}
\noindent{\small theorem-library/cauchy-subseq-proof.scm ==================================================================== totally-bounded-has-cauchy-subsequence -- PROVEN (was asserted).}\par
\vnbentry{\texttt{totally-\allowbreak{}bounded-\allowbreak{}has-\allowbreak{}cauchy-\allowbreak{}subsequence}}{\vnbstmt{\begin{array}{@{}l@{}} \forall s.; \forall f \in (\mathbb{N} \to \operatorname{pts}(s)). \\ \quad \operatorname{totally\text{-}bounded}(s) \Rightarrow  \\ \quad \quad \exists i. \\ \quad \quad \quad \operatorname{strictly\text{-}mono\text{-}nn}(i) \wedge \\ \quad \quad \quad \operatorname{is\text{-}cauchy\text{-}seq}(s, \operatorname{subseq}(f, i)) \end{array}}}

\vnbfile{theorem-library/rr-ms-dist.scm}{theorem-library/rr-ms-dist.scm}
\noindent{\small rr-ms-dist.scm -- the distance of RR-MS, on the surface.}\par
\vnbentry{\texttt{rr-\allowbreak{}ms-\allowbreak{}dist}}{\vnbstmt{\begin{array}{@{}l@{}} \forall u, v. \\ \quad ((u \in \mathbb{R}) \wedge (v \in \mathbb{R})) \Rightarrow  \\ \quad \quad (\operatorname{dist}(\operatorname{rr\text{-}ms})(u, v) \simeq \lvert (u - v) \rvert) \end{array}}}

\vnbfile{theorem-library/rr-metric-space-proof.scm}{theorem-library/rr-metric-space-proof.scm}
\noindent{\small rr-metric-space-proof.scm -- IS-METRIC-SPACE(RR-MS), PROVEN.}\par
\vnbentry{\texttt{rr-\allowbreak{}is-\allowbreak{}metric-\allowbreak{}space}}{\vnbstmt{\operatorname{is\text{-}metric\text{-}space}(\operatorname{rr\text{-}ms})}}

\vnbfile{theorem-library/trunc-metric-proof.scm}{theorem-library/trunc-metric-proof.scm}
\noindent{\small trunc-metric-proof.scm -- the truncated metric d'(x,y) = min(1, d(x,y)) is a metric, is bounded by 1, and is UNIFORMLY equivalent to d.}\par
\vnbentry{\texttt{rr-\allowbreak{}min-\allowbreak{}one-\allowbreak{}mono}}{\vnbstmt{\begin{array}{@{}l@{}} \forall a, b \in \mathbb{R}. \\ \quad (a \leq b) \Rightarrow  \\ \quad \quad (\operatorname{min}(1, a) \leq \operatorname{min}(1, b)) \end{array}}}
\vnbentry{\texttt{trunc-\allowbreak{}metric-\allowbreak{}dist}}{\vnbstmt{\begin{array}{@{}l@{}} \forall s, u, v. \\ \quad ((u \in \operatorname{pts}(s)) \wedge (v \in \operatorname{pts}(s))) \Rightarrow  \\ \quad \quad (\operatorname{dist}(\operatorname{trunc\text{-}metric}(s))(u, v) \simeq \operatorname{min}(1, \operatorname{dist}(s)(u, v))) \end{array}}}
\vnbentry{\texttt{trunc-\allowbreak{}metric-\allowbreak{}is-\allowbreak{}metric-\allowbreak{}space}}{\vnbstmt{\begin{array}{@{}l@{}} \forall s. \\ \quad \operatorname{is\text{-}metric\text{-}space}(s) \Rightarrow  \\ \quad \quad \operatorname{is\text{-}metric\text{-}space}(\operatorname{trunc\text{-}metric}(s)) \end{array}}}
\vnbentry{\texttt{trunc-\allowbreak{}metric-\allowbreak{}bounded}}{\vnbstmt{\begin{array}{@{}l@{}} \forall s, u, v. \\ \quad \operatorname{is\text{-}metric\text{-}space}(s) \wedge \\ \quad (u \in \operatorname{pts}(s)) \wedge \\ \quad (v \in \operatorname{pts}(s)) \Rightarrow \\ \quad \quad (\operatorname{dist}(\operatorname{trunc\text{-}metric}(s))(u, v) \leq 1) \end{array}}}
\vnbentry{\texttt{trunc-\allowbreak{}metric-\allowbreak{}id-\allowbreak{}uniformly-\allowbreak{}continuous}}{\vnbstmt{\begin{array}{@{}l@{}} \forall s. \\ \quad \operatorname{is\text{-}metric\text{-}space}(s) \Rightarrow  \\ \quad \quad \operatorname{is\text{-}uniformly\text{-}continuous}(s, \operatorname{trunc\text{-}metric}(s), (\lambda a \in \operatorname{pts}(s).\; a)) \end{array}}}
\vnbentry{\texttt{trunc-\allowbreak{}metric-\allowbreak{}id-\allowbreak{}uniformly-\allowbreak{}continuous-\allowbreak{}back}}{\vnbstmt{\begin{array}{@{}l@{}} \forall s. \\ \quad \operatorname{is\text{-}metric\text{-}space}(s) \Rightarrow  \\ \quad \quad \operatorname{is\text{-}uniformly\text{-}continuous}(\operatorname{trunc\text{-}metric}(s), s, (\lambda a \in \operatorname{pts}(s).\; a)) \end{array}}}

\vnbfile{theorem-library/rr-complete-proof.scm}{theorem-library/rr-complete-proof.scm}
\noindent{\small rr-complete-proof.scm -- RR is a COMPLETE METRIC SPACE.  Proven 2026-08-01/02.}\par
\vnbentry{\texttt{rr-\allowbreak{}complete}}{\vnbstmt{\operatorname{is\text{-}complete}(\operatorname{rr\text{-}ms})}}

\vnbfile{theorem-library/metric-limit-unique.scm}{theorem-library/metric-limit-unique.scm}
\noindent{\small metric-limit-unique.scm -- A SEQUENCE IN A METRIC SPACE HAS AT MOST ONE LIMIT, proven once for an arbitrary metric space, with the RR and CC cases taken as instances.}\par
\vnbentry{\texttt{metric-\allowbreak{}limit-\allowbreak{}unique}}{\vnbstmt{\begin{array}{@{}l@{}} \forall s, f, \ell_{1}, \ell_{2}. \\ \quad \operatorname{converges\text{-}to}(s, f, \ell_{1}) \wedge \\ \quad \operatorname{converges\text{-}to}(s, f, \ell_{2}) \Rightarrow \\ \quad \quad (\ell_{1} = \ell_{2}) \end{array}}}
\vnbentry{\texttt{rr-\allowbreak{}limit-\allowbreak{}unique}}{\vnbstmt{\begin{array}{@{}l@{}} \forall f \in (\mathbb{N} \to \mathbb{R}), \ell_{1}, \ell_{2} \in \mathbb{R}. \\ \quad \operatorname{converges\text{-}to}(\operatorname{rr\text{-}ms}, f, \ell_{1}) \wedge \\ \quad \operatorname{converges\text{-}to}(\operatorname{rr\text{-}ms}, f, \ell_{2}) \Rightarrow \\ \quad \quad (\ell_{1} = \ell_{2}) \end{array}}}
\vnbentry{\texttt{cc-\allowbreak{}limit-\allowbreak{}unique}}{\vnbstmt{\begin{array}{@{}l@{}} \forall f \in (\mathbb{N} \to \mathbb{C}), \ell_{1}, \ell_{2} \in \mathbb{C}. \\ \quad \operatorname{converges\text{-}to}(\operatorname{cc\text{-}ms}, f, \ell_{1}) \wedge \\ \quad \operatorname{converges\text{-}to}(\operatorname{cc\text{-}ms}, f, \ell_{2}) \Rightarrow \\ \quad \quad (\ell_{1} = \ell_{2}) \end{array}}}

\vnbfile{theorem-library/seq-limit-core.scm}{theorem-library/seq-limit-core.scm}
\noindent{\small theorem-library/seq-limit-core.scm -- THE LIMIT OF A REAL SEQUENCE AS A TERM (SEQ-LIMIT), its three characterising theorems, the TAIL bound rr-limit-tail-abs-le, and rr-cauchy-converges.}\par
\vnbentry{\texttt{seq-\allowbreak{}limit-\allowbreak{}converges-\allowbreak{}to}}{\vnbstmt{\begin{array}{@{}l@{}} \forall f. \\ \quad \operatorname{converges}(\operatorname{rr\text{-}ms}, f) \Rightarrow  \\ \quad \quad \operatorname{converges\text{-}to}(\operatorname{rr\text{-}ms}, f, \operatorname{seq\text{-}limit}(f)) \end{array}}}
\vnbentry{\texttt{seq-\allowbreak{}limit-\allowbreak{}in-\allowbreak{}rr}}{\vnbstmt{\forall f.\; (\operatorname{seq\text{-}limit}(f) \in \mathbb{R})}}
\vnbentry{\texttt{seq-\allowbreak{}limit-\allowbreak{}value}}{\vnbstmt{\begin{array}{@{}l@{}} \forall f, v. \\ \quad \operatorname{converges\text{-}to}(\operatorname{rr\text{-}ms}, f, v) \Rightarrow  \\ \quad \quad (\operatorname{seq\text{-}limit}(f) = v) \end{array}}}
\vnbentry{\texttt{rr-\allowbreak{}limit-\allowbreak{}tail-\allowbreak{}abs-\allowbreak{}le}}{\vnbstmt{\begin{array}{@{}l@{}} \forall f \in (\mathbb{N} \to \mathbb{R}), v, a, c \in \mathbb{R}, n \in \mathbb{N}. \\ \quad \operatorname{converges\text{-}to}(\operatorname{rr\text{-}ms}, f, v) \wedge \\ \quad \forall k \in \mathbb{N}. \\ \quad \quad ((n \leq k) \Rightarrow (\lvert (f(k) - a) \rvert \leq c)) \Rightarrow \\ \quad \quad (\lvert (v - a) \rvert \leq c) \end{array}}}
\vnbentry{\texttt{rr-\allowbreak{}cauchy-\allowbreak{}converges}}{\vnbstmt{\begin{array}{@{}l@{}} \forall f \in (\mathbb{N} \to \mathbb{R}). \\ \quad \forall s. \\ \quad \quad \operatorname{pos\text{-}rr}(s) \Rightarrow  \\ \quad \quad \quad \exists n \in \mathbb{N}. \forall m, p \in \mathbb{N}. \\ \quad \quad \quad \quad ((n \leq m) \wedge (n \leq p)) \Rightarrow  \\ \quad \quad \quad \quad \quad (\lvert (f(m) - f(p)) \rvert \leq s) \Rightarrow \\ \quad \quad \operatorname{converges}(\operatorname{rr\text{-}ms}, f) \end{array}}}

\vnbfile{theorem-library/ascoli-analytic-cores.scm}{theorem-library/ascoli-analytic-cores.scm}
\noindent{\small ascoli-analytic-cores.scm -- the analytic content of Ascoli-Arzelà, machine-proved. bridge, MACHINE-PROVED, and the one rung it still stands on.}\par
\vnbentry{\texttt{equicont-\allowbreak{}dense-\allowbreak{}conv-\allowbreak{}ptwise-\allowbreak{}cauchy}}{\vnbstmt{\begin{array}{@{}l@{}} \forall s.; \forall m \in (\mathbb{N} \to (\operatorname{pts}(s) \to \mathbb{R})). \\ \quad \operatorname{is\text{-}equicontinuous}(s, \operatorname{rr\text{-}ms}, m) \wedge \\ \quad \exists q. \\ \quad \quad (\operatorname{is\text{-}dense\text{-}seq}(s, q) \wedge \operatorname{converges\text{-}on}(s, m, q)) \Rightarrow \\ \quad \quad \operatorname{is\text{-}ptwise\text{-}cauchy}(s, m) \end{array}}}
\vnbentry{\texttt{uniform-\allowbreak{}limit-\allowbreak{}continuous}}{\vnbstmt{\begin{array}{@{}l@{}} \forall s, m, g. \\ \quad (m \in (\mathbb{N} \to (\operatorname{pts}(s) \to \mathbb{R}))) \wedge \\ \quad (g \in (\operatorname{pts}(s) \to \mathbb{R})) \wedge \\ \quad \forall k \in \mathbb{N}. \\ \quad \quad \operatorname{is\text{-}continuous}(s, \operatorname{rr\text{-}ms}, m(k)) \wedge \\ \quad \operatorname{converges\text{-}uniformly}(s, m, g) \Rightarrow \\ \quad \quad \operatorname{is\text{-}continuous}(s, \operatorname{rr\text{-}ms}, g) \end{array}}}

\vnbfile{theorem-library/rake-open-sets.scm}{theorem-library/rake-open-sets.scm}
\noindent{\small rake-open-sets.scm -- the metric topology axioms, PROVEN.}\par
\vnbentry{\texttt{empty-\allowbreak{}is-\allowbreak{}open}}{\vnbstmt{\begin{array}{@{}l@{}} \forall s. \\ \quad \operatorname{is\text{-}metric\text{-}space}(s) \Rightarrow  \\ \quad \quad \operatorname{is\text{-}open}(s, \emptyset) \end{array}}}
\vnbentry{\texttt{carrier-\allowbreak{}is-\allowbreak{}open}}{\vnbstmt{\begin{array}{@{}l@{}} \forall s. \\ \quad \operatorname{is\text{-}metric\text{-}space}(s) \Rightarrow  \\ \quad \quad \operatorname{is\text{-}open}(s, \operatorname{pts}(s)) \end{array}}}
\vnbentry{\texttt{union-\allowbreak{}of-\allowbreak{}opens-\allowbreak{}open}}{\vnbstmt{\begin{array}{@{}l@{}} \forall s, a, g. \\ \quad \operatorname{is\text{-}metric\text{-}space}(s) \wedge \\ \quad \forall i \in a.\; \operatorname{is\text{-}open}(s, g(i)) \Rightarrow \\ \quad \quad \operatorname{is\text{-}open}(s, \operatorname{big\text{-}union}(i, a, g(i))) \end{array}}}
\vnbentry{\texttt{inter-\allowbreak{}of-\allowbreak{}opens-\allowbreak{}open}}{\vnbstmt{\begin{array}{@{}l@{}} \forall s, u, w. \\ \quad \operatorname{is\text{-}metric\text{-}space}(s) \wedge \\ \quad \operatorname{is\text{-}open}(s, u) \wedge \\ \quad \operatorname{is\text{-}open}(s, w) \Rightarrow \\ \quad \quad \operatorname{is\text{-}open}(s, (u \cap w)) \end{array}}}
\vnbentry{\texttt{preimage-\allowbreak{}complement}}{\vnbstmt{\begin{array}{@{}l@{}} \forall s, t, f, u. \\ \quad \operatorname{is\text{-}metric\text{-}space}(s) \wedge \\ \quad \operatorname{is\text{-}metric\text{-}space}(t) \wedge \\ \quad (f \in (\operatorname{pts}(s) \to \operatorname{pts}(t))) \Rightarrow \\ \quad \quad \operatorname{preimage}(s, f, \operatorname{complement\text{-}in}(\operatorname{pts}(t), u)) \\ \quad \quad =\; \operatorname{complement\text{-}in}(\operatorname{pts}(s), \operatorname{preimage}(s, f, u)) \end{array}}}
\vnbentry{\texttt{open-\allowbreak{}preimage-\allowbreak{}implies-\allowbreak{}continuous}}{\vnbstmt{\begin{array}{@{}l@{}} \forall s, t, f. \\ \quad \operatorname{is\text{-}metric\text{-}space}(s) \wedge \\ \quad \operatorname{is\text{-}metric\text{-}space}(t) \wedge \\ \quad (f \in (\operatorname{pts}(s) \to \operatorname{pts}(t))) \wedge \\ \quad \forall v. \\ \quad \quad \operatorname{is\text{-}open}(t, v) \Rightarrow  \\ \quad \quad \quad \operatorname{is\text{-}open}(s, \operatorname{preimage}(s, f, v)) \Rightarrow \\ \quad \quad \operatorname{is\text{-}continuous}(s, t, f) \end{array}}}
\vnbentry{\texttt{continuous-\allowbreak{}implies-\allowbreak{}closed-\allowbreak{}preimage}}{\vnbstmt{\begin{array}{@{}l@{}} \forall s, t, f, a. \\ \quad (\operatorname{is\text{-}continuous}(s, t, f) \wedge \operatorname{is\text{-}closed}(t, a)) \Rightarrow  \\ \quad \quad \operatorname{is\text{-}closed}(s, \operatorname{preimage}(s, f, a)) \end{array}}}
\vnbentry{\texttt{closed-\allowbreak{}preimage-\allowbreak{}implies-\allowbreak{}continuous}}{\vnbstmt{\begin{array}{@{}l@{}} \forall s, t, f. \\ \quad \operatorname{is\text{-}metric\text{-}space}(s) \wedge \\ \quad \operatorname{is\text{-}metric\text{-}space}(t) \wedge \\ \quad (f \in (\operatorname{pts}(s) \to \operatorname{pts}(t))) \wedge \\ \quad \forall a. \\ \quad \quad \operatorname{is\text{-}closed}(t, a) \Rightarrow  \\ \quad \quad \quad \operatorname{is\text{-}closed}(s, \operatorname{preimage}(s, f, a)) \Rightarrow \\ \quad \quad \operatorname{is\text{-}continuous}(s, t, f) \end{array}}}

\vnbfile{theorem-library/ptwise-cauchy-unif.scm}{theorem-library/ptwise-cauchy-unif.scm}
\noindent{\small ptwise-cauchy-unif.scm -- pointwise Cauchy + equicontinuous on a COMPACT space  =$>$  uniformly Cauchy  (the user's notes, Prop 3.33, Cauchy spelling).}\par
\vnbentry{\texttt{ptwise-\allowbreak{}cauchy-\allowbreak{}compact-\allowbreak{}equicont-\allowbreak{}unif}}{\vnbstmt{\begin{array}{@{}l@{}} \forall s, m. \\ \quad \operatorname{is\text{-}compact}(s) \wedge \\ \quad (m \in (\mathbb{N} \to (\operatorname{pts}(s) \to \mathbb{R}))) \wedge \\ \quad \operatorname{is\text{-}equicontinuous}(s, \operatorname{rr\text{-}ms}, m) \wedge \\ \quad \operatorname{is\text{-}ptwise\text{-}cauchy}(s, m) \Rightarrow \\ \quad \quad \operatorname{is\text{-}unif\text{-}cauchy}(s, m) \end{array}}}

\vnbfile{theorem-library/ascoli-assembly.scm}{theorem-library/ascoli-assembly.scm}
\noindent{\small theorem-library/ascoli-assembly.scm -- equicont-dense-conv-implies-unif-cauchy.}\par
\vnbentry{\texttt{equicont-\allowbreak{}dense-\allowbreak{}conv-\allowbreak{}implies-\allowbreak{}unif-\allowbreak{}cauchy}}{\vnbstmt{\begin{array}{@{}l@{}} \forall s, m. \\ \quad \operatorname{is\text{-}compact}(s) \wedge \\ \quad (m \in (\mathbb{N} \to (\operatorname{pts}(s) \to \mathbb{R}))) \wedge \\ \quad \operatorname{is\text{-}equicontinuous}(s, \operatorname{rr\text{-}ms}, m) \wedge \\ \quad \exists q. \\ \quad \quad (\operatorname{is\text{-}dense\text{-}seq}(s, q) \wedge \operatorname{converges\text{-}on}(s, m, q)) \Rightarrow \\ \quad \quad \operatorname{is\text{-}unif\text{-}cauchy}(s, m) \end{array}}}

\vnbfile{theorem-library/unif-cauchy-limit.scm}{theorem-library/unif-cauchy-limit.scm}
\noindent{\small unif-cauchy-limit.scm -- unif-cauchy-has-uniform-limit, PROVEN (2026-09-14).}\par
\vnbentry{\texttt{unif-\allowbreak{}cauchy-\allowbreak{}has-\allowbreak{}uniform-\allowbreak{}limit}}{\vnbstmt{\begin{array}{@{}l@{}} \forall s.; \forall m \in (\mathbb{N} \to (\operatorname{pts}(s) \to \mathbb{R})). \\ \quad \operatorname{is\text{-}unif\text{-}cauchy}(s, m) \Rightarrow  \\ \quad \quad \exists g \in (\operatorname{pts}(s) \to \mathbb{R}). \\ \quad \quad \quad \operatorname{converges\text{-}uniformly}(s, m, g) \end{array}}}

\vnbfile{theorem-library/ascoli-bridge.scm}{theorem-library/ascoli-bridge.scm}
\noindent{\small ascoli-bridge.scm -- the dense-to-uniform bridge for Ascoli-Arzelà.}\par
\vnbentry{\texttt{unif-\allowbreak{}cauchy-\allowbreak{}cont-\allowbreak{}implies-\allowbreak{}uniform-\allowbreak{}limit}}{\vnbstmt{\begin{array}{@{}l@{}} \forall s.; \forall m \in (\mathbb{N} \to (\operatorname{pts}(s) \to \mathbb{R})). \\ \quad \forall k \in \mathbb{N}. \\ \quad \quad \operatorname{is\text{-}continuous}(s, \operatorname{rr\text{-}ms}, m(k)) \wedge \\ \quad \operatorname{is\text{-}unif\text{-}cauchy}(s, m) \Rightarrow \\ \quad \quad \exists g \in (\operatorname{pts}(s) \to \mathbb{R}). \\ \quad \quad \quad \operatorname{is\text{-}continuous}(s, \operatorname{rr\text{-}ms}, g) \wedge \\ \quad \quad \quad \operatorname{converges\text{-}uniformly}(s, m, g) \end{array}}}
\vnbentry{\texttt{ascoli-\allowbreak{}dense-\allowbreak{}bridge}}{\vnbstmt{\begin{array}{@{}l@{}} \forall s, m. \\ \quad \operatorname{is\text{-}compact}(s) \wedge \\ \quad (m \in (\mathbb{N} \to (\operatorname{pts}(s) \to \mathbb{R}))) \wedge \\ \quad \forall k \in \mathbb{N}. \\ \quad \quad \operatorname{is\text{-}continuous}(s, \operatorname{rr\text{-}ms}, m(k)) \wedge \\ \quad \operatorname{is\text{-}equicontinuous}(s, \operatorname{rr\text{-}ms}, m) \wedge \\ \quad \exists q. \\ \quad \quad (\operatorname{is\text{-}dense\text{-}seq}(s, q) \wedge \operatorname{converges\text{-}on}(s, m, q)) \Rightarrow \\ \quad \quad \exists g \in (\operatorname{pts}(s) \to \mathbb{R}). \\ \quad \quad \quad \operatorname{is\text{-}continuous}(s, \operatorname{rr\text{-}ms}, g) \wedge \\ \quad \quad \quad \operatorname{converges\text{-}uniformly}(s, m, g) \end{array}}}

\vnbfile{theorem-library/injective-star.scm}{theorem-library/injective-star.scm}
\noindent{\small injective-star.scm -- the bridge from the set-function vocabulary to the class-function one.}\par
\vnbentry{\texttt{injection-\allowbreak{}is-\allowbreak{}injective-\allowbreak{}star}}{\vnbstmt{\begin{array}{@{}l@{}} \forall x, y.; \forall f \in \operatorname{injection}(x, y). \\ \quad \operatorname{injective\text{-}star}(f) \end{array}}}

\vnbfile{theorem-library/zorn-proof.scm}{theorem-library/zorn-proof.scm}
\noindent{\small zorn-proof.scm -- proving ZORN'S LEMMA from the well-ordering principle, via the GREEDY-KEPT-CHAIN construction (avoids the cardinality/Hartogs comparison lemma the base lacks).  Multi-session build; this file grows as lemmas land. Started 2026-07-22.}\par
\vnbentry{\texttt{zorn-\allowbreak{}empty-\allowbreak{}is-\allowbreak{}chain}}{\vnbstmt{\forall d \in \mathbf{Set}, l.\; \operatorname{is\text{-}chain}(d, l, \emptyset)}}
\vnbentry{\texttt{zorn-\allowbreak{}chain-\allowbreak{}plus-\allowbreak{}keepset}}{\vnbstmt{\begin{array}{@{}l@{}} \forall i, d, l, t, a. \\ \quad (\operatorname{is\text{-}partial\text{-}order}(d, l) \wedge \operatorname{is\text{-}chain}(d, l, t)) \Rightarrow  \\ \quad \quad \operatorname{is\text{-}chain}(d, l, (t \cup \operatorname{keep\text{-}set}(i, d, l, t, a))) \end{array}}}
\vnbentry{\texttt{zorn-\allowbreak{}zkept-\allowbreak{}monotone}}{\vnbstmt{\begin{array}{@{}l@{}} \forall b \in \mathrm{Ord}, i, d, l.; \forall a \in \mathrm{Ord}. \\ \quad \operatorname{ord\text{-}le}(a, b) \Rightarrow  \\ \quad \quad (\operatorname{zkept}(i, d, l, a) \subseteq \operatorname{zkept}(i, d, l, b)) \end{array}}}
\vnbentry{\texttt{zorn-\allowbreak{}zkept-\allowbreak{}limit-\allowbreak{}is-\allowbreak{}chain}}{\vnbstmt{\begin{array}{@{}l@{}} \forall i, d, l, m. \\ \quad \operatorname{is\text{-}partial\text{-}order}(d, l) \wedge \\ \quad \operatorname{limit\text{-}ord}(m) \wedge \\ \quad \forall a \in \mathrm{OS}(m). \\ \quad \quad \operatorname{is\text{-}chain}(d, l, \operatorname{zkept}(i, d, l, a)) \Rightarrow \\ \quad \quad \operatorname{is\text{-}chain}(d, l, \operatorname{big\text{-}union}(a_{2}, \mathrm{OS}(m), \operatorname{zkept}(i, d, l, a_{2}))) \end{array}}}
\vnbentry{\texttt{zorn-\allowbreak{}zkept-\allowbreak{}is-\allowbreak{}chain}}{\vnbstmt{\begin{array}{@{}l@{}} \forall a \in \mathrm{Ord}, i, d, l. \\ \quad \operatorname{is\text{-}partial\text{-}order}(d, l) \Rightarrow  \\ \quad \quad \operatorname{is\text{-}chain}(d, l, \operatorname{zkept}(i, d, l, a)) \end{array}}}
\vnbentry{\texttt{greedy-\allowbreak{}chain-\allowbreak{}is-\allowbreak{}chain}}{\vnbstmt{\begin{array}{@{}l@{}} \forall i, d, l. \\ \quad \operatorname{is\text{-}partial\text{-}order}(d, l) \Rightarrow  \\ \quad \quad \operatorname{is\text{-}chain}(d, l, \operatorname{greedy\text{-}chain}(i, d, l)) \end{array}}}

\vnbfile{theorem-library/zorn-route-two.scm}{theorem-library/zorn-route-two.scm}
\noindent{\small zorn-route-two.scm -- ZORN'S LEMMA, PROVED.}\par
\vnbentry{\texttt{strictly-\allowbreak{}below-\allowbreak{}irreflexive}}{\vnbstmt{\forall l, x.\; \neg \operatorname{is\text{-}strictly\text{-}below}(l, x, x)}}
\vnbentry{\texttt{strictly-\allowbreak{}below-\allowbreak{}transitive}}{\vnbstmt{\begin{array}{@{}l@{}} \forall d, l, x, y, z. \\ \quad \operatorname{is\text{-}partial\text{-}order}(d, l) \wedge \\ \quad (x \in d) \wedge \\ \quad (y \in d) \wedge \\ \quad (z \in d) \wedge \\ \quad \operatorname{is\text{-}strictly\text{-}below}(l, x, y) \wedge \\ \quad \operatorname{is\text{-}strictly\text{-}below}(l, y, z) \Rightarrow \\ \quad \quad \operatorname{is\text{-}strictly\text{-}below}(l, x, z) \end{array}}}
\vnbentry{\texttt{strictly-\allowbreak{}below-\allowbreak{}not-\allowbreak{}equal}}{\vnbstmt{\begin{array}{@{}l@{}} \forall l, x, y. \\ \quad (x = y) \Rightarrow  \\ \quad \quad \neg \operatorname{is\text{-}strictly\text{-}below}(l, x, y) \end{array}}}
\vnbentry{\texttt{no-\allowbreak{}maximal-\allowbreak{}strictly-\allowbreak{}above}}{\vnbstmt{\begin{array}{@{}l@{}} \forall d, l.; \forall x \in d. \\ \quad \neg \exists \mathit{mx}.\; \operatorname{is\text{-}maximal}(d, l, \mathit{mx}) \Rightarrow  \\ \quad \quad \exists y \in d.\; \operatorname{is\text{-}strictly\text{-}below}(l, x, y) \end{array}}}
\vnbentry{\texttt{po-\allowbreak{}reflexive}}{\vnbstmt{\begin{array}{@{}l@{}} \forall d, l, x. \\ \quad (\operatorname{is\text{-}partial\text{-}order}(d, l) \wedge (x \in d)) \Rightarrow  \\ \quad \quad (\langle x, x \rangle \in l) \end{array}}}
\vnbentry{\texttt{po-\allowbreak{}transitive}}{\vnbstmt{\begin{array}{@{}l@{}} \forall d, l, x, y, z. \\ \quad \operatorname{is\text{-}partial\text{-}order}(d, l) \wedge \\ \quad (x \in d) \wedge \\ \quad (y \in d) \wedge \\ \quad (z \in d) \wedge \\ \quad (\langle x, y \rangle \in l) \wedge \\ \quad (\langle y, z \rangle \in l) \Rightarrow \\ \quad \quad (\langle x, z \rangle \in l) \end{array}}}
\vnbentry{\texttt{ord-\allowbreak{}lt-\allowbreak{}succ-\allowbreak{}cases}}{\vnbstmt{\begin{array}{@{}l@{}} \forall a, b. \\ \quad ((a \in \mathrm{Ord}) \wedge ((b \in \mathrm{Ord}) \wedge \operatorname{ord\text{-}lt}(b, (a + 1)))) \Rightarrow  \\ \quad \quad (\operatorname{ord\text{-}lt}(b, a) \vee (b = a)) \end{array}}}
\vnbentry{\texttt{ord-\allowbreak{}succ-\allowbreak{}lt-\allowbreak{}limit}}{\vnbstmt{\begin{array}{@{}l@{}} \forall m, b. \\ \quad \operatorname{limit\text{-}ord}(m) \wedge \\ \quad (b \in \mathrm{Ord}) \wedge \\ \quad \operatorname{ord\text{-}lt}(b, m) \Rightarrow \\ \quad \quad \operatorname{ord\text{-}lt}((b + 1), m) \end{array}}}
\vnbentry{\texttt{zup-\allowbreak{}tower}}{\vnbstmt{\begin{array}{@{}l@{}} \forall a \in \mathrm{Ord}, d, l. \\ \quad \operatorname{is\text{-}partial\text{-}order}(d, l) \wedge \\ \quad \exists w.\; (w \in d) \wedge \\ \quad \forall h. \\ \quad \quad \operatorname{is\text{-}chain}(d, l, h) \Rightarrow  \\ \quad \quad \quad \exists u.\; \operatorname{is\text{-}upper\text{-}bound}(d, l, h, u) \wedge \\ \quad \forall x \in d. \exists y \in d. \\ \quad \quad \operatorname{is\text{-}strictly\text{-}below}(l, x, y) \Rightarrow \\ \quad \quad (\operatorname{zup}(d, l, a) \in d) \wedge \\ \quad \quad \forall b. \\ \quad \quad \quad \operatorname{ord\text{-}lt}(b, a) \Rightarrow  \\ \quad \quad \quad \quad (\operatorname{zup}(d, l, b) \in d) \wedge \\ \quad \quad \quad \quad \operatorname{is\text{-}strictly\text{-}below}(l, \operatorname{zup}(d, l, b), \operatorname{zup}(d, l, a)) \end{array}}}
\vnbentry{\texttt{zup-\allowbreak{}injective}}{\vnbstmt{\begin{array}{@{}l@{}} \forall d, l, a_{1}, a_{2}. \\ \quad \operatorname{is\text{-}partial\text{-}order}(d, l) \wedge \\ \quad \exists w.\; (w \in d) \wedge \\ \quad \forall h. \\ \quad \quad \operatorname{is\text{-}chain}(d, l, h) \Rightarrow  \\ \quad \quad \quad \exists u.\; \operatorname{is\text{-}upper\text{-}bound}(d, l, h, u) \wedge \\ \quad \forall x \in d. \exists y \in d. \\ \quad \quad \operatorname{is\text{-}strictly\text{-}below}(l, x, y) \wedge \\ \quad (a_{1} \in \mathrm{Ord}) \wedge \\ \quad (a_{2} \in \mathrm{Ord}) \wedge \\ \quad (\operatorname{zup}(d, l, a_{1}) = \operatorname{zup}(d, l, a_{2})) \Rightarrow \\ \quad \quad (a_{1} = a_{2}) \end{array}}}
\vnbentry{\texttt{zorn-\allowbreak{}lemma}}{\vnbstmt{\begin{array}{@{}l@{}} \forall d, l. \\ \quad \operatorname{is\text{-}partial\text{-}order}(d, l) \wedge \\ \quad \exists w.\; (w \in d) \wedge \\ \quad \forall h. \\ \quad \quad \operatorname{is\text{-}chain}(d, l, h) \Rightarrow  \\ \quad \quad \quad \exists u.\; \operatorname{is\text{-}upper\text{-}bound}(d, l, h, u) \Rightarrow \\ \quad \quad \exists x.\; \operatorname{is\text{-}maximal}(d, l, x) \end{array}}}

\vnbfile{theorem-library/metric-closure-laws.scm}{theorem-library/metric-closure-laws.scm}
\noindent{\small metric-closure-laws.scm -- CLOSURE, INTERIOR and the CLOSED BALL.}\par
\vnbentry{\texttt{closed-\allowbreak{}ball-\allowbreak{}sep-\allowbreak{}unfold}}{\vnbstmt{\begin{array}{@{}l@{}} \forall s, c, r. \\ \quad (\operatorname{closed\text{-}ball}(s, c, r) \simeq \{y \in \operatorname{pts}(s) : (\operatorname{dist}(s)(c, y) \leq r)\}) \end{array}}}
\vnbentry{\texttt{closure-\allowbreak{}sep-\allowbreak{}unfold}}{\vnbstmt{\begin{array}{@{}l@{}} \forall s, a. \\ \quad (\operatorname{closure}(s, a) \simeq \{x \in \operatorname{pts}(s) : \forall u.\; ((\operatorname{is\text{-}open}(s, u) \wedge (x \in u)) \Rightarrow \exists y \in u.\; (y \in a))\}) \end{array}}}
\vnbentry{\texttt{interior-\allowbreak{}sep-\allowbreak{}unfold}}{\vnbstmt{\begin{array}{@{}l@{}} \forall s, a. \\ \quad (\operatorname{interior}(s, a) \simeq \{x \in \operatorname{pts}(s) : \exists u.\; (\operatorname{is\text{-}open}(s, u) \wedge ((x \in u) \wedge (u \subseteq a)))\}) \end{array}}}
\vnbentry{\texttt{closed-\allowbreak{}ball-\allowbreak{}membership}}{\vnbstmt{\begin{array}{@{}l@{}} \forall s, c, r, y. \\ \quad ((y \in \operatorname{closed\text{-}ball}(s, c, r)) \Leftrightarrow ((y \in \operatorname{pts}(s)) \wedge (\operatorname{dist}(s)(c, y) \leq r))) \end{array}}}
\vnbentry{\texttt{closure-\allowbreak{}membership}}{\vnbstmt{\begin{array}{@{}l@{}} \forall s, a, x. \\ \quad ((x \in \operatorname{closure}(s, a)) \Leftrightarrow ((x \in \operatorname{pts}(s)) \wedge \forall u.\; ((\operatorname{is\text{-}open}(s, u) \wedge (x \in u)) \Rightarrow \exists y \in u.\; (y \in a)))) \end{array}}}
\vnbentry{\texttt{interior-\allowbreak{}membership}}{\vnbstmt{\begin{array}{@{}l@{}} \forall s, a, x. \\ \quad ((x \in \operatorname{interior}(s, a)) \Leftrightarrow ((x \in \operatorname{pts}(s)) \wedge \exists u.\; (\operatorname{is\text{-}open}(s, u) \wedge ((x \in u) \wedge (u \subseteq a))))) \end{array}}}
\vnbentry{\texttt{closed-\allowbreak{}ball-\allowbreak{}is-\allowbreak{}set}}{\vnbstmt{\begin{array}{@{}l@{}} \forall s, c, r. \\ \quad \operatorname{is\text{-}metric\text{-}space}(s) \Rightarrow  \\ \quad \quad (\operatorname{closed\text{-}ball}(s, c, r) \in \mathbf{Set}) \end{array}}}
\vnbentry{\texttt{closure-\allowbreak{}in-\allowbreak{}carrier}}{\vnbstmt{\forall s, a.\; (\operatorname{closure}(s, a) \subseteq \operatorname{pts}(s))}}
\vnbentry{\texttt{interior-\allowbreak{}in-\allowbreak{}carrier}}{\vnbstmt{\forall s, a.\; (\operatorname{interior}(s, a) \subseteq \operatorname{pts}(s))}}
\vnbentry{\texttt{interior-\allowbreak{}in-\allowbreak{}set}}{\vnbstmt{\forall s, a.\; (\operatorname{interior}(s, a) \subseteq a)}}
\vnbentry{\texttt{closure-\allowbreak{}contains}}{\vnbstmt{\begin{array}{@{}l@{}} \forall s, a. \\ \quad (a \subseteq \operatorname{pts}(s)) \Rightarrow  \\ \quad \quad (a \subseteq \operatorname{closure}(s, a)) \end{array}}}
\vnbentry{\texttt{ball-\allowbreak{}in-\allowbreak{}closed-\allowbreak{}ball}}{\vnbstmt{\begin{array}{@{}l@{}} \forall s, c, r. \\ \quad (\operatorname{ball}(s, c, r) \subseteq \operatorname{closed\text{-}ball}(s, c, r)) \end{array}}}
\vnbentry{\texttt{closed-\allowbreak{}ball-\allowbreak{}in-\allowbreak{}ball}}{\vnbstmt{\begin{array}{@{}l@{}} \forall s.; \forall c \in \operatorname{pts}(s), r, q \in \mathbb{R}. \\ \quad (\operatorname{is\text{-}metric\text{-}space}(s) \wedge (r < q)) \Rightarrow  \\ \quad \quad (\operatorname{closed\text{-}ball}(s, c, r) \subseteq \operatorname{ball}(s, c, q)) \end{array}}}
\vnbentry{\texttt{interior-\allowbreak{}mem-\allowbreak{}intro}}{\vnbstmt{\begin{array}{@{}l@{}} \forall s, a, u, x. \\ \quad (\operatorname{is\text{-}open}(s, u) \wedge ((u \subseteq a) \wedge (x \in u))) \Rightarrow  \\ \quad \quad (x \in \operatorname{interior}(s, a)) \end{array}}}
\vnbentry{\texttt{interior-\allowbreak{}is-\allowbreak{}open}}{\vnbstmt{\begin{array}{@{}l@{}} \forall s, a. \\ \quad \operatorname{is\text{-}metric\text{-}space}(s) \Rightarrow  \\ \quad \quad \operatorname{is\text{-}open}(s, \operatorname{interior}(s, a)) \end{array}}}
\vnbentry{\texttt{closure-\allowbreak{}is-\allowbreak{}closed}}{\vnbstmt{\begin{array}{@{}l@{}} \forall s, a. \\ \quad \operatorname{is\text{-}metric\text{-}space}(s) \Rightarrow  \\ \quad \quad \operatorname{is\text{-}closed}(s, \operatorname{closure}(s, a)) \end{array}}}
\vnbentry{\texttt{closed-\allowbreak{}ball-\allowbreak{}is-\allowbreak{}closed}}{\vnbstmt{\begin{array}{@{}l@{}} \forall s.; \forall c \in \operatorname{pts}(s), r \in \mathbb{R}. \\ \quad \operatorname{is\text{-}metric\text{-}space}(s) \Rightarrow  \\ \quad \quad \operatorname{is\text{-}closed}(s, \operatorname{closed\text{-}ball}(s, c, r)) \end{array}}}
\vnbentry{\texttt{closure-\allowbreak{}in-\allowbreak{}closed}}{\vnbstmt{\begin{array}{@{}l@{}} \forall s, a, t. \\ \quad (\operatorname{is\text{-}closed}(s, t) \wedge (a \subseteq t)) \Rightarrow  \\ \quad \quad (\operatorname{closure}(s, a) \subseteq t) \end{array}}}
\vnbentry{\texttt{closure-\allowbreak{}of-\allowbreak{}ball-\allowbreak{}in-\allowbreak{}closed-\allowbreak{}ball}}{\vnbstmt{\begin{array}{@{}l@{}} \forall s.; \forall c \in \operatorname{pts}(s), r \in \mathbb{R}. \\ \quad \operatorname{is\text{-}metric\text{-}space}(s) \Rightarrow  \\ \quad \quad \operatorname{closure}(s, \operatorname{ball}(s, c, r)) \\ \quad \quad \subseteq\; \operatorname{closed\text{-}ball}(s, c, r) \end{array}}}
\vnbentry{\texttt{closed-\allowbreak{}set-\allowbreak{}limit-\allowbreak{}in}}{\vnbstmt{\begin{array}{@{}l@{}} \forall s, a, f, l. \\ \quad \operatorname{is\text{-}closed}(s, a) \wedge \\ \quad \operatorname{converges\text{-}to}(s, f, l) \wedge \\ \quad \forall n \in \mathbb{N}.\; (f(n) \in a) \Rightarrow \\ \quad \quad (l \in a) \end{array}}}
\vnbentry{\texttt{closure-\allowbreak{}is-\allowbreak{}set}}{\vnbstmt{\begin{array}{@{}l@{}} \forall s, a. \\ \quad \operatorname{is\text{-}metric\text{-}space}(s) \Rightarrow  \\ \quad \quad (\operatorname{closure}(s, a) \in \mathbf{Set}) \end{array}}}
\vnbentry{\texttt{interior-\allowbreak{}is-\allowbreak{}set}}{\vnbstmt{\begin{array}{@{}l@{}} \forall s, a. \\ \quad \operatorname{is\text{-}metric\text{-}space}(s) \Rightarrow  \\ \quad \quad (\operatorname{interior}(s, a) \in \mathbf{Set}) \end{array}}}
\vnbentry{\texttt{nowhere-\allowbreak{}dense-\allowbreak{}ball-\allowbreak{}point}}{\vnbstmt{\begin{array}{@{}l@{}} \forall s, a.; \forall c \in \operatorname{pts}(s), r. \\ \quad (\operatorname{is\text{-}nowhere\text{-}dense}(s, a) \wedge \operatorname{pos\text{-}rr}(r)) \Rightarrow  \\ \quad \quad \exists x \in \operatorname{ball}(s, c, r). \\ \quad \quad \quad \neg (x \in \operatorname{closure}(s, a)) \end{array}}}
\vnbentry{\texttt{nowhere-\allowbreak{}dense-\allowbreak{}closed-\allowbreak{}ball}}{\vnbstmt{\begin{array}{@{}l@{}} \forall s, a.; \forall c \in \operatorname{pts}(s), r, d. \\ \quad \operatorname{is\text{-}nowhere\text{-}dense}(s, a) \wedge \\ \quad \operatorname{pos\text{-}rr}(r) \wedge \\ \quad \operatorname{pos\text{-}rr}(d) \Rightarrow \\ \quad \quad \exists x \in \operatorname{pts}(s), q. \\ \quad \quad \quad \operatorname{pos\text{-}rr}(q) \wedge \\ \quad \quad \quad (q \leq d) \wedge \\ \quad \quad \quad (\operatorname{closed\text{-}ball}(s, x, q) \subseteq \operatorname{ball}(s, c, r)) \wedge \\ \quad \quad \quad \forall z \in \operatorname{closed\text{-}ball}(s, x, q). \\ \quad \quad \quad \quad \neg (z \in \operatorname{closure}(s, a)) \end{array}}}

\vnbfile{theorem-library/rake-baire.scm}{theorem-library/rake-baire.scm}
\noindent{\small rake-baire.scm -- towards the BAIRE CATEGORY theorem (batch 12-F, 2026-09-20).}\par
\vnbentry{\texttt{nested-\allowbreak{}family-\allowbreak{}monotone}}{\vnbstmt{\begin{array}{@{}l@{}} \forall m \in \mathbb{N}, w.; \forall g \in (\mathbb{N} \to w), n \in \mathbb{N}. \\ \quad (\forall n \in \mathbb{N}.\; (g((n + 1)) \subseteq g(n)) \wedge (n \leq m)) \Rightarrow  \\ \quad \quad (g(m) \subseteq g(n)) \end{array}}}

\vnbfile{theorem-library/subseq-convergence-proof.scm}{theorem-library/subseq-convergence-proof.scm}
\noindent{\small theorem-library/subseq-convergence-proof.scm ==================================================================== subseq-of-convergent -- PROVEN.}\par
\vnbentry{\texttt{subseq-\allowbreak{}of-\allowbreak{}convergent}}{\vnbstmt{\begin{array}{@{}l@{}} \forall s, f, i, l. \\ \quad \operatorname{converges\text{-}to}(s, f, l) \wedge \\ \quad \operatorname{strictly\text{-}mono\text{-}nn}(i) \Rightarrow \\ \quad \quad \operatorname{converges\text{-}to}(s, \operatorname{subseq}(f, i), l) \end{array}}}

\vnbfile{theorem-library/subsqn-basics.scm}{theorem-library/subsqn-basics.scm}
\noindent{\small subsqn-basics.scm -- the SUBSQN relation and what is proven about it.}\par
\vnbentry{\texttt{is-\allowbreak{}subsequence-\allowbreak{}iff-\allowbreak{}typed-\allowbreak{}subsqn}}{\vnbstmt{\begin{array}{@{}l@{}} \forall s, y, f. \\ \quad (\operatorname{is\text{-}subsequence}(s, y, f) \Leftrightarrow ((f \in (\mathbb{N} \to \operatorname{pts}(s))) \wedge \operatorname{subsqn}(y, f))) \end{array}}}
\vnbentry{\texttt{subsqn-\allowbreak{}preserves-\allowbreak{}sqn}}{\vnbstmt{\begin{array}{@{}l@{}} \forall a, f, y. \\ \quad ((f \in a^{\mathbb{N}}) \wedge \operatorname{subsqn}(y, f)) \Rightarrow  \\ \quad \quad (y \in a^{\mathbb{N}}) \end{array}}}

\vnbfile{theorem-library/subsequence-principle.scm}{theorem-library/subsequence-principle.scm}
\noindent{\small theorem-library/subsequence-principle.scm ====================================================================== subsequence-principle -- PROVEN.}\par
\vnbentry{\texttt{subsequence-\allowbreak{}principle}}{\vnbstmt{\begin{array}{@{}l@{}} \forall s, f, l. \\ \quad \operatorname{is\text{-}metric\text{-}space}(s) \wedge \\ \quad (f \in (\mathbb{N} \to \operatorname{pts}(s))) \wedge \\ \quad (l \in \operatorname{pts}(s)) \wedge \\ \quad \forall i. \\ \quad \quad \operatorname{strictly\text{-}mono\text{-}nn}(i) \Rightarrow  \\ \quad \quad \quad \exists a. \\ \quad \quad \quad \quad \operatorname{strictly\text{-}mono\text{-}nn}(a) \wedge \\ \quad \quad \quad \quad \operatorname{converges\text{-}to}(s, \operatorname{subseq}(\operatorname{subseq}(f, i), a), l) \Rightarrow \\ \quad \quad \operatorname{converges\text{-}to}(s, f, l) \end{array}}}

\vnbfile{theorem-library/rake-offbill-combinatorial.scm}{theorem-library/rake-offbill-combinatorial.scm}
\noindent{\small theorem-library/rake-offbill-combinatorial.scm ==================================================================== Batch 8-D: the off-bill `combinatorial' PSS entries that batch 6/7 made cheap.  Three theorems, in dependency order:}\par
\vnbentry{\texttt{tb-\allowbreak{}block-\allowbreak{}step}}{\vnbstmt{\begin{array}{@{}l@{}} \forall s.; \forall f \in (\mathbb{N} \to \operatorname{pts}(s)), r.; \forall j \in \mathrm{Inf}(\mathbb{N}). \\ \quad (\operatorname{totally\text{-}bounded}(s) \wedge \operatorname{pos\text{-}rr}(r)) \Rightarrow  \\ \quad \quad \exists a \in \mathrm{Inf}(\mathbb{N}). \\ \quad \quad \quad (a \subseteq j) \wedge \\ \quad \quad \quad \exists c \in \operatorname{pts}(s). \forall i \in a. \\ \quad \quad \quad \quad (f(i) \in \operatorname{ball}(s, c, r)) \end{array}}}
\vnbentry{\texttt{tb-\allowbreak{}has-\allowbreak{}eps-\allowbreak{}cauchy-\allowbreak{}subseq}}{\vnbstmt{\begin{array}{@{}l@{}} \forall s.; \forall f \in (\mathbb{N} \to \operatorname{pts}(s)), a. \\ \quad (\operatorname{totally\text{-}bounded}(s) \wedge \operatorname{pos\text{-}rr}(a)) \Rightarrow  \\ \quad \quad \exists i. \\ \quad \quad \quad \operatorname{strictly\text{-}mono\text{-}nn}(i) \wedge \\ \quad \quad \quad \operatorname{is\text{-}eps\text{-}cauchy\text{-}seq}(s, a, \operatorname{subseq}(f, i)) \end{array}}}
\vnbentry{\texttt{block-\allowbreak{}family}}{\vnbstmt{\begin{array}{@{}l@{}} \forall s.; \forall f \in (\mathbb{N} \to \operatorname{pts}(s)), d. \\ \quad (\operatorname{totally\text{-}bounded}(s) \wedge \operatorname{null\text{-}rr\text{-}seq}(d)) \Rightarrow  \\ \quad \quad \exists a \in (\mathbb{N} \to \mathrm{Inf}(\mathbb{N})). \\ \quad \quad \quad \forall k \in \mathbb{N}.\; (a((k + 1)) \subseteq a(k)) \wedge \\ \quad \quad \quad \forall k \in \mathbb{N}. \exists c \in \operatorname{pts}(s). \forall i \in a(k). \\ \quad \quad \quad \quad (f(i) \in \operatorname{ball}(s, c, d(k))) \end{array}}}

\vnbfile{theorem-library/poly-tail-zero.scm}{theorem-library/poly-tail-zero.scm}
\noindent{\small poly-tail-zero.scm -- "a polynomial is a sequence that is eventually zero".}\par
\vnbentry{\texttt{poly-\allowbreak{}tail-\allowbreak{}zero-\allowbreak{}fwd}}{\vnbstmt{\begin{array}{@{}l@{}} \forall a.; \forall f \in \operatorname{carr}(\operatorname{poly}(a)). \exists n \in \mathbb{N}. \forall k \in \mathbb{N}. \\ \quad ((n \leq k) \Rightarrow (f(k) = \operatorname{zero}(a))) \end{array}}}
\vnbentry{\texttt{poly-\allowbreak{}tail-\allowbreak{}zero-\allowbreak{}bwd}}{\vnbstmt{\begin{array}{@{}l@{}} \forall a.; \forall f \in \operatorname{carr}(a)^{\mathbb{N}}. \\ \quad \exists n \in \mathbb{N}. \forall k \in \mathbb{N}. \\ \quad \quad ((n \leq k) \Rightarrow (f(k) = \operatorname{zero}(a))) \Rightarrow \\ \quad \quad (f \in \operatorname{carr}(\operatorname{poly}(a))) \end{array}}}
\vnbentry{\texttt{poly-\allowbreak{}tail-\allowbreak{}zero}}{\vnbstmt{\begin{array}{@{}l@{}} \forall a, f. \\ \quad ((f \in \operatorname{carr}(\operatorname{poly}(a))) \Leftrightarrow ((f \in \operatorname{carr}(a)^{\mathbb{N}}) \wedge \exists n \in \mathbb{N}.\; \forall k \in \mathbb{N}.\; ((n \leq k) \Rightarrow (f(k) = \operatorname{zero}(a))))) \end{array}}}

\vnbfile{theorem-library/finsum-fiber.scm}{theorem-library/finsum-fiber.scm}
\noindent{\small finsum-fiber.scm -- FIBERED FUBINI: group the terms of a finite sum by the fibers of an arbitrary map, and the CARTESIAN projection brick it needs.}\par
\vnbentry{\texttt{cartesian-\allowbreak{}nth}}{\vnbstmt{\begin{array}{@{}l@{}} \forall x, a, b. \\ \quad (x \in (a \times b)) \Rightarrow  \\ \quad \quad ((\operatorname{nth}(1, x) \in a) \wedge (\operatorname{nth}(2, x) \in b)) \end{array}}}
\vnbentry{\texttt{finsum-\allowbreak{}fiber}}{\vnbstmt{\begin{array}{@{}l@{}} \forall g.; \forall s, t \in \mathbf{Set}, i \in (s \to t), f \in (s \to \operatorname{carr}(g)). \\ \quad \operatorname{is\text{-}abelian\text{-}group}(g) \wedge \\ \quad (|s| \in \mathbb{N}) \wedge \\ \quad (|t| \in \mathbb{N}) \Rightarrow \\ \quad \quad \operatorname{finsum}(g, f, s) \\ \quad \quad =\; \operatorname{finsum}(g, (\lambda a \in t.\; \operatorname{finsum}(g, f, \{b \in s : (i(b) = a)\})), t) \end{array}}}

\vnbfile{theorem-library/rake-baire-2.scm}{theorem-library/rake-baire-2.scm}
\noindent{\small rake-baire-2.scm -- the BAIRE CATEGORY theorem (batch 13-C, 2026-09-20).}\par
\vnbentry{\texttt{closed-\allowbreak{}ball-\allowbreak{}in-\allowbreak{}carrier}}{\vnbstmt{\begin{array}{@{}l@{}} \forall s, v, \mathit{qv}. \\ \quad (\operatorname{closed\text{-}ball}(s, v, \mathit{qv}) \subseteq \operatorname{pts}(s)) \end{array}}}
\vnbentry{\texttt{closed-\allowbreak{}ball-\allowbreak{}center-\allowbreak{}in}}{\vnbstmt{\begin{array}{@{}l@{}} \forall s.; \forall v \in \operatorname{pts}(s), \mathit{qv} \in \mathbb{R}. \\ \quad (\operatorname{is\text{-}metric\text{-}space}(s) \wedge (0 \leq \mathit{qv})) \Rightarrow  \\ \quad \quad (v \in \operatorname{closed\text{-}ball}(s, v, \mathit{qv})) \end{array}}}
\vnbentry{\texttt{closed-\allowbreak{}ball-\allowbreak{}in-\allowbreak{}power}}{\vnbstmt{\begin{array}{@{}l@{}} \forall s, v, \mathit{qv}. \\ \quad \operatorname{is\text{-}metric\text{-}space}(s) \Rightarrow  \\ \quad \quad \operatorname{closed\text{-}ball}(s, v, \mathit{qv}) \\ \quad \quad \in\; \operatorname{power}(\operatorname{pts}(s)) \end{array}}}
\vnbentry{\texttt{closed-\allowbreak{}set-\allowbreak{}limit-\allowbreak{}in-\allowbreak{}tail}}{\vnbstmt{\begin{array}{@{}l@{}} \forall s, v, \mathit{fv}, \mathit{lv}, b. \\ \quad \operatorname{is\text{-}closed}(s, v) \wedge \\ \quad \operatorname{converges\text{-}to}(s, \mathit{fv}, \mathit{lv}) \wedge \\ \quad (b \in \mathbb{N}) \wedge \\ \quad \forall n \in \mathbb{N}. \\ \quad \quad ((b \leq n) \Rightarrow (\operatorname{fv}(n) \in v)) \Rightarrow \\ \quad \quad (\mathit{lv} \in v) \end{array}}}
\vnbentry{\texttt{nested-\allowbreak{}closed-\allowbreak{}balls-\allowbreak{}point}}{\vnbstmt{\begin{array}{@{}l@{}} \forall s.; \forall r \in (\mathbb{N} \to \operatorname{pts}(s)), \mathit{rds} \in (\mathbb{N} \to \mathbb{R}). \\ \quad \operatorname{is\text{-}complete}(s) \wedge \\ \quad \forall n \in \mathbb{N}.\; (0 \leq \operatorname{rds}(n)) \wedge \\ \quad \forall n \in \mathbb{N}. \\ \quad \quad \operatorname{closed\text{-}ball}(s, r((n + 1)), \operatorname{rds}((n + 1))) \\ \quad \quad \subseteq\; \operatorname{closed\text{-}ball}(s, r(n), \operatorname{rds}(n)) \wedge \\ \quad \forall n \in \mathbb{N}. \\ \quad \quad (\operatorname{rds}((n + 1)) \leq \operatorname{recip}((n + 1))) \Rightarrow \\ \quad \quad \exists v \in \operatorname{pts}(s). \forall n \in \mathbb{N}. \\ \quad \quad \quad (v \in \operatorname{closed\text{-}ball}(s, r(n), \operatorname{rds}(n))) \end{array}}}
\vnbentry{\texttt{baire-\allowbreak{}category}}{\vnbstmt{\begin{array}{@{}l@{}} \forall s. \\ \quad (\operatorname{is\text{-}complete}(s) \wedge \exists x.\; (x \in \operatorname{pts}(s))) \Rightarrow  \\ \quad \quad \operatorname{is\text{-}nonmeager}(s, \operatorname{pts}(s)) \end{array}}}

\vnbfile{theorem-library/lam-fun-bricks.scm}{theorem-library/lam-fun-bricks.scm}
\noindent{\small lam-fun-bricks.scm -- the lambda FUN-typing bricks of the matrix arc, PROVEN.}\par
\vnbentry{\texttt{matprod-\allowbreak{}summand-\allowbreak{}type}}{\vnbstmt{\begin{array}{@{}l@{}} \forall a, m, n, b, p, q, i, c. \\ \quad \operatorname{is\text{-}ring}(a) \wedge \\ \quad (p \in \operatorname{mat}(m, n, \operatorname{carr}(a))) \wedge \\ \quad (q \in \operatorname{mat}(n, b, \operatorname{carr}(a))) \wedge \\ \quad (i \in \operatorname{interval}(1, m)) \wedge \\ \quad (c \in \operatorname{interval}(1, b)) \Rightarrow \\ \quad \quad (\lambda j \in \operatorname{interval}(1, n).\; \\ \quad \quad \quad \operatorname{mul}(a)(\operatorname{entry}(p, i, j), \operatorname{entry}(q, j, c))) \\ \quad \quad \in\; (\operatorname{interval}(1, n) \to \operatorname{carr}(\operatorname{ring\text{-}additive\text{-}ag}(a))) \end{array}}}
\vnbentry{\texttt{matmul-\allowbreak{}assoc-\allowbreak{}summand-\allowbreak{}type}}{\vnbstmt{\begin{array}{@{}l@{}} \forall a, m, n, k, l, p, q, r, w, b. \\ \quad \operatorname{is\text{-}ring}(a) \wedge \\ \quad (p \in \operatorname{mat}(m, n, \operatorname{carr}(a))) \wedge \\ \quad (q \in \operatorname{mat}(n, k, \operatorname{carr}(a))) \wedge \\ \quad (r \in \operatorname{mat}(k, l, \operatorname{carr}(a))) \wedge \\ \quad (w \in \operatorname{interval}(1, m)) \wedge \\ \quad (b \in \operatorname{interval}(1, l)) \Rightarrow \\ \quad \quad (\lambda z \in (\operatorname{interval}(1, k) \times \operatorname{interval}(1, n)).\; \\ \quad \quad \quad \operatorname{mul}(a)(\operatorname{mul}(a)(\operatorname{entry}(p, w, \operatorname{nth}(2, z)), \operatorname{entry}(q, \operatorname{nth}(2, z), \operatorname{nth}(1, z))), \\ \quad \quad \quad \quad \operatorname{entry}(r, \operatorname{nth}(1, z), b))) \\ \quad \quad \in\; ((\operatorname{interval}(1, k) \times \operatorname{interval}(1, n)) \to \operatorname{carr}(\operatorname{ring\text{-}additive\text{-}ag}(a))) \end{array}}}
\vnbentry{\texttt{matact-\allowbreak{}summand-\allowbreak{}type}}{\vnbstmt{\begin{array}{@{}l@{}} \forall d, m, n, q, p, u, i, c. \\ \quad \operatorname{is\text{-}module}(d) \wedge \\ \quad (p \in \operatorname{mat}(m, n, \operatorname{carr}(\operatorname{scal}(d)))) \wedge \\ \quad (u \in \operatorname{mat}(n, q, \operatorname{vec}(d))) \wedge \\ \quad (i \in \operatorname{interval}(1, m)) \wedge \\ \quad (c \in \operatorname{interval}(1, q)) \Rightarrow \\ \quad \quad (\lambda j \in \operatorname{interval}(1, n).\; \\ \quad \quad \quad \operatorname{act}(d)(\operatorname{entry}(p, i, j), \operatorname{entry}(u, j, c))) \\ \quad \quad \in\; (\operatorname{interval}(1, n) \to \operatorname{carr}(\operatorname{module\text{-}vector\text{-}ag}(d))) \end{array}}}
\vnbentry{\texttt{matunit-\allowbreak{}summand-\allowbreak{}type}}{\vnbstmt{\begin{array}{@{}l@{}} \forall a, m, n, p, k, l, i, c. \\ \quad \operatorname{is\text{-}ring}(a) \wedge \\ \quad (p \in \operatorname{mat}(m, n, \operatorname{carr}(a))) \wedge \\ \quad (i \in \operatorname{interval}(1, m)) \wedge \\ \quad (c \in \operatorname{interval}(1, n)) \Rightarrow \\ \quad \quad (\lambda j \in \operatorname{interval}(1, n).\; \\ \quad \quad \quad \operatorname{mul}(a)(\operatorname{entry}(p, i, j), \operatorname{entry}(\operatorname{matunit}(a, n, k, l), j, c))) \\ \quad \quad \in\; (\operatorname{interval}(1, n) \to \operatorname{carr}(\operatorname{ring\text{-}additive\text{-}ag}(a))) \end{array}}}
\vnbentry{\texttt{tel-\allowbreak{}outf-\allowbreak{}type}}{\vnbstmt{\begin{array}{@{}l@{}} \forall a, m, n, k, l, p, q, r, w, b. \\ \quad \operatorname{is\text{-}ring}(a) \wedge \\ \quad (p \in \operatorname{mat}(m, n, \operatorname{carr}(a))) \wedge \\ \quad (q \in \operatorname{mat}(n, k, \operatorname{carr}(a))) \wedge \\ \quad (r \in \operatorname{mat}(k, l, \operatorname{carr}(a))) \wedge \\ \quad (w \in \operatorname{interval}(1, m)) \wedge \\ \quad (b \in \operatorname{interval}(1, l)) \wedge \\ \quad (1 \leq n) \Rightarrow \\ \quad \quad (\lambda j \in \operatorname{interval}(1, k).\; \\ \quad \quad \quad \operatorname{mul}(a)(\operatorname{entry}(\operatorname{matmul}(a, p, q), w, j), \operatorname{entry}(r, j, b))) \\ \quad \quad \in\; (\operatorname{interval}(1, k) \to \operatorname{carr}(\operatorname{ring\text{-}additive\text{-}ag}(a))) \end{array}}}
\vnbentry{\texttt{tel-\allowbreak{}tout-\allowbreak{}type}}{\vnbstmt{\begin{array}{@{}l@{}} \forall a, m, n, k, l, p, q, r, w, b. \\ \quad \operatorname{is\text{-}ring}(a) \wedge \\ \quad (p \in \operatorname{mat}(m, n, \operatorname{carr}(a))) \wedge \\ \quad (q \in \operatorname{mat}(n, k, \operatorname{carr}(a))) \wedge \\ \quad (r \in \operatorname{mat}(k, l, \operatorname{carr}(a))) \wedge \\ \quad (w \in \operatorname{interval}(1, m)) \wedge \\ \quad (b \in \operatorname{interval}(1, l)) \Rightarrow \\ \quad \quad (\lambda c \in \operatorname{interval}(1, k).\; \\ \quad \quad \quad \operatorname{finsum}(\operatorname{ring\text{-}additive\text{-}ag}(a), \\ \quad \quad \quad \quad (\lambda j \in \operatorname{interval}(1, n).\; \\ \quad \quad \quad \quad \quad (\lambda z \in (\operatorname{interval}(1, k) \times \operatorname{interval}(1, n)).\; \operatorname{mul}(a)(\operatorname{mul}(a)(\operatorname{entry}(p, w, \operatorname{nth}(2, z)), \operatorname{entry}(q, \operatorname{nth}(2, z), \operatorname{nth}(1, z))), \operatorname{entry}(r, \operatorname{nth}(1, z), b)))(\langle c, j \rangle)), \\ \quad \quad \quad \quad \operatorname{interval}(1, n))) \\ \quad \quad \in\; (\operatorname{interval}(1, k) \to \operatorname{carr}(\operatorname{ring\text{-}additive\text{-}ag}(a))) \end{array}}}
\vnbentry{\texttt{tel-\allowbreak{}inf-\allowbreak{}type}}{\vnbstmt{\begin{array}{@{}l@{}} \forall a, m, n, k, l, p, q, r, w, b, x. \\ \quad \operatorname{is\text{-}ring}(a) \wedge \\ \quad (p \in \operatorname{mat}(m, n, \operatorname{carr}(a))) \wedge \\ \quad (q \in \operatorname{mat}(n, k, \operatorname{carr}(a))) \wedge \\ \quad (r \in \operatorname{mat}(k, l, \operatorname{carr}(a))) \wedge \\ \quad (w \in \operatorname{interval}(1, m)) \wedge \\ \quad (b \in \operatorname{interval}(1, l)) \wedge \\ \quad (x \in \operatorname{interval}(1, k)) \Rightarrow \\ \quad \quad (\lambda j \in \operatorname{interval}(1, n).\; \\ \quad \quad \quad \operatorname{mul}(a)(\operatorname{entry}(p, w, j), \operatorname{entry}(q, j, x))) \\ \quad \quad \in\; (\operatorname{interval}(1, n) \to \operatorname{carr}(a)) \end{array}}}
\vnbentry{\texttt{tel-\allowbreak{}dist-\allowbreak{}type}}{\vnbstmt{\begin{array}{@{}l@{}} \forall a, m, n, k, l, p, q, r, w, b, x. \\ \quad \operatorname{is\text{-}ring}(a) \wedge \\ \quad (p \in \operatorname{mat}(m, n, \operatorname{carr}(a))) \wedge \\ \quad (q \in \operatorname{mat}(n, k, \operatorname{carr}(a))) \wedge \\ \quad (r \in \operatorname{mat}(k, l, \operatorname{carr}(a))) \wedge \\ \quad (w \in \operatorname{interval}(1, m)) \wedge \\ \quad (b \in \operatorname{interval}(1, l)) \wedge \\ \quad (x \in \operatorname{interval}(1, k)) \Rightarrow \\ \quad \quad (\lambda z \in \operatorname{interval}(1, n).\; \\ \quad \quad \quad \operatorname{mul}(a)((\lambda j \in \operatorname{interval}(1, n).\; \operatorname{mul}(a)(\operatorname{entry}(p, w, j), \operatorname{entry}(q, j, x)))(z), \\ \quad \quad \quad \quad \operatorname{entry}(r, x, b))) \\ \quad \quad \in\; (\operatorname{interval}(1, n) \to \operatorname{carr}(\operatorname{ring\text{-}additive\text{-}ag}(a))) \end{array}}}
\vnbentry{\texttt{tel-\allowbreak{}red-\allowbreak{}type}}{\vnbstmt{\begin{array}{@{}l@{}} \forall a, m, n, k, l, p, q, r, w, b, x. \\ \quad \operatorname{is\text{-}ring}(a) \wedge \\ \quad (p \in \operatorname{mat}(m, n, \operatorname{carr}(a))) \wedge \\ \quad (q \in \operatorname{mat}(n, k, \operatorname{carr}(a))) \wedge \\ \quad (r \in \operatorname{mat}(k, l, \operatorname{carr}(a))) \wedge \\ \quad (w \in \operatorname{interval}(1, m)) \wedge \\ \quad (b \in \operatorname{interval}(1, l)) \wedge \\ \quad (x \in \operatorname{interval}(1, k)) \Rightarrow \\ \quad \quad (\lambda j \in \operatorname{interval}(1, n).\; \\ \quad \quad \quad \operatorname{mul}(a)(\operatorname{mul}(a)(\operatorname{entry}(p, w, j), \operatorname{entry}(q, j, x)), \operatorname{entry}(r, x, b))) \\ \quad \quad \in\; (\operatorname{interval}(1, n) \to \operatorname{carr}(\operatorname{ring\text{-}additive\text{-}ag}(a))) \end{array}}}
\vnbentry{\texttt{ter-\allowbreak{}outf-\allowbreak{}type}}{\vnbstmt{\begin{array}{@{}l@{}} \forall a, m, n, k, l, p, q, r, w, b. \\ \quad \operatorname{is\text{-}ring}(a) \wedge \\ \quad (p \in \operatorname{mat}(m, n, \operatorname{carr}(a))) \wedge \\ \quad (q \in \operatorname{mat}(n, k, \operatorname{carr}(a))) \wedge \\ \quad (r \in \operatorname{mat}(k, l, \operatorname{carr}(a))) \wedge \\ \quad (w \in \operatorname{interval}(1, m)) \wedge \\ \quad (b \in \operatorname{interval}(1, l)) \wedge \\ \quad (1 \leq k) \Rightarrow \\ \quad \quad (\lambda j \in \operatorname{interval}(1, n).\; \\ \quad \quad \quad \operatorname{mul}(a)(\operatorname{entry}(p, w, j), \operatorname{entry}(\operatorname{matmul}(a, q, r), j, b))) \\ \quad \quad \in\; (\operatorname{interval}(1, n) \to \operatorname{carr}(\operatorname{ring\text{-}additive\text{-}ag}(a))) \end{array}}}
\vnbentry{\texttt{ter-\allowbreak{}tout-\allowbreak{}type}}{\vnbstmt{\begin{array}{@{}l@{}} \forall a, m, n, k, l, p, q, r, w, b. \\ \quad \operatorname{is\text{-}ring}(a) \wedge \\ \quad (p \in \operatorname{mat}(m, n, \operatorname{carr}(a))) \wedge \\ \quad (q \in \operatorname{mat}(n, k, \operatorname{carr}(a))) \wedge \\ \quad (r \in \operatorname{mat}(k, l, \operatorname{carr}(a))) \wedge \\ \quad (w \in \operatorname{interval}(1, m)) \wedge \\ \quad (b \in \operatorname{interval}(1, l)) \Rightarrow \\ \quad \quad (\lambda j \in \operatorname{interval}(1, n).\; \\ \quad \quad \quad \operatorname{finsum}(\operatorname{ring\text{-}additive\text{-}ag}(a), \\ \quad \quad \quad \quad (\lambda c \in \operatorname{interval}(1, k).\; \\ \quad \quad \quad \quad \quad (\lambda z \in (\operatorname{interval}(1, k) \times \operatorname{interval}(1, n)).\; \operatorname{mul}(a)(\operatorname{mul}(a)(\operatorname{entry}(p, w, \operatorname{nth}(2, z)), \operatorname{entry}(q, \operatorname{nth}(2, z), \operatorname{nth}(1, z))), \operatorname{entry}(r, \operatorname{nth}(1, z), b)))(\langle c, j \rangle)), \\ \quad \quad \quad \quad \operatorname{interval}(1, k))) \\ \quad \quad \in\; (\operatorname{interval}(1, n) \to \operatorname{carr}(\operatorname{ring\text{-}additive\text{-}ag}(a))) \end{array}}}
\vnbentry{\texttt{ter-\allowbreak{}gj-\allowbreak{}type}}{\vnbstmt{\begin{array}{@{}l@{}} \forall a, m, n, k, l, p, q, r, w, b, x. \\ \quad \operatorname{is\text{-}ring}(a) \wedge \\ \quad (p \in \operatorname{mat}(m, n, \operatorname{carr}(a))) \wedge \\ \quad (q \in \operatorname{mat}(n, k, \operatorname{carr}(a))) \wedge \\ \quad (r \in \operatorname{mat}(k, l, \operatorname{carr}(a))) \wedge \\ \quad (w \in \operatorname{interval}(1, m)) \wedge \\ \quad (b \in \operatorname{interval}(1, l)) \wedge \\ \quad (x \in \operatorname{interval}(1, n)) \Rightarrow \\ \quad \quad (\lambda j \in \operatorname{interval}(1, k).\; \\ \quad \quad \quad \operatorname{mul}(a)(\operatorname{entry}(q, x, j), \operatorname{entry}(r, j, b))) \\ \quad \quad \in\; (\operatorname{interval}(1, k) \to \operatorname{carr}(a)) \end{array}}}
\vnbentry{\texttt{ter-\allowbreak{}dist-\allowbreak{}type}}{\vnbstmt{\begin{array}{@{}l@{}} \forall a, m, n, k, l, p, q, r, w, b, x. \\ \quad \operatorname{is\text{-}ring}(a) \wedge \\ \quad (p \in \operatorname{mat}(m, n, \operatorname{carr}(a))) \wedge \\ \quad (q \in \operatorname{mat}(n, k, \operatorname{carr}(a))) \wedge \\ \quad (r \in \operatorname{mat}(k, l, \operatorname{carr}(a))) \wedge \\ \quad (w \in \operatorname{interval}(1, m)) \wedge \\ \quad (b \in \operatorname{interval}(1, l)) \wedge \\ \quad (x \in \operatorname{interval}(1, n)) \Rightarrow \\ \quad \quad (\lambda z \in \operatorname{interval}(1, k).\; \\ \quad \quad \quad \operatorname{mul}(a)(\operatorname{entry}(p, w, x), \\ \quad \quad \quad \quad (\lambda j \in \operatorname{interval}(1, k).\; \operatorname{mul}(a)(\operatorname{entry}(q, x, j), \operatorname{entry}(r, j, b)))(z))) \\ \quad \quad \in\; (\operatorname{interval}(1, k) \to \operatorname{carr}(\operatorname{ring\text{-}additive\text{-}ag}(a))) \end{array}}}
\vnbentry{\texttt{ter-\allowbreak{}red-\allowbreak{}type}}{\vnbstmt{\begin{array}{@{}l@{}} \forall a, m, n, k, l, p, q, r, w, b, x. \\ \quad \operatorname{is\text{-}ring}(a) \wedge \\ \quad (p \in \operatorname{mat}(m, n, \operatorname{carr}(a))) \wedge \\ \quad (q \in \operatorname{mat}(n, k, \operatorname{carr}(a))) \wedge \\ \quad (r \in \operatorname{mat}(k, l, \operatorname{carr}(a))) \wedge \\ \quad (w \in \operatorname{interval}(1, m)) \wedge \\ \quad (b \in \operatorname{interval}(1, l)) \wedge \\ \quad (x \in \operatorname{interval}(1, n)) \Rightarrow \\ \quad \quad (\lambda c \in \operatorname{interval}(1, k).\; \\ \quad \quad \quad \operatorname{mul}(a)(\operatorname{mul}(a)(\operatorname{entry}(p, w, x), \operatorname{entry}(q, x, c)), \operatorname{entry}(r, c, b))) \\ \quad \quad \in\; (\operatorname{interval}(1, k) \to \operatorname{carr}(\operatorname{ring\text{-}additive\text{-}ag}(a))) \end{array}}}
\vnbentry{\texttt{mal-\allowbreak{}outf-\allowbreak{}type}}{\vnbstmt{\begin{array}{@{}l@{}} \forall d, m, n, k, l, p, q, u, w, a. \\ \quad \operatorname{is\text{-}module}(d) \wedge \\ \quad (p \in \operatorname{mat}(m, n, \operatorname{carr}(\operatorname{scal}(d)))) \wedge \\ \quad (q \in \operatorname{mat}(n, k, \operatorname{carr}(\operatorname{scal}(d)))) \wedge \\ \quad (u \in \operatorname{mat}(k, l, \operatorname{vec}(d))) \wedge \\ \quad (w \in \operatorname{interval}(1, m)) \wedge \\ \quad (a \in \operatorname{interval}(1, l)) \wedge \\ \quad (1 \leq n) \Rightarrow \\ \quad \quad (\lambda j \in \operatorname{interval}(1, k).\; \\ \quad \quad \quad \operatorname{act}(d)(\operatorname{entry}(\operatorname{matmul}(\operatorname{scal}(d), p, q), w, j), \operatorname{entry}(u, j, a))) \\ \quad \quad \in\; (\operatorname{interval}(1, k) \to \operatorname{carr}(\operatorname{module\text{-}vector\text{-}ag}(d))) \end{array}}}
\vnbentry{\texttt{mal-\allowbreak{}tout-\allowbreak{}type}}{\vnbstmt{\begin{array}{@{}l@{}} \forall d, m, n, k, l, p, q, u, w, a. \\ \quad \operatorname{is\text{-}module}(d) \wedge \\ \quad (p \in \operatorname{mat}(m, n, \operatorname{carr}(\operatorname{scal}(d)))) \wedge \\ \quad (q \in \operatorname{mat}(n, k, \operatorname{carr}(\operatorname{scal}(d)))) \wedge \\ \quad (u \in \operatorname{mat}(k, l, \operatorname{vec}(d))) \wedge \\ \quad (w \in \operatorname{interval}(1, m)) \wedge \\ \quad (a \in \operatorname{interval}(1, l)) \Rightarrow \\ \quad \quad (\lambda c \in \operatorname{interval}(1, k).\; \\ \quad \quad \quad \operatorname{finsum}(\operatorname{module\text{-}vector\text{-}ag}(d), \\ \quad \quad \quad \quad (\lambda j \in \operatorname{interval}(1, n).\; \\ \quad \quad \quad \quad \quad (\lambda z \in (\operatorname{interval}(1, k) \times \operatorname{interval}(1, n)).\; \operatorname{act}(d)(\operatorname{mul}(\operatorname{scal}(d))(\operatorname{entry}(p, w, \operatorname{nth}(2, z)), \operatorname{entry}(q, \operatorname{nth}(2, z), \operatorname{nth}(1, z))), \operatorname{entry}(u, \operatorname{nth}(1, z), a)))(\langle c, j \rangle)), \\ \quad \quad \quad \quad \operatorname{interval}(1, n))) \\ \quad \quad \in\; (\operatorname{interval}(1, k) \to \operatorname{carr}(\operatorname{module\text{-}vector\text{-}ag}(d))) \end{array}}}
\vnbentry{\texttt{mal-\allowbreak{}inf-\allowbreak{}type}}{\vnbstmt{\begin{array}{@{}l@{}} \forall d, m, n, k, l, p, q, u, w, a, x. \\ \quad \operatorname{is\text{-}module}(d) \wedge \\ \quad (p \in \operatorname{mat}(m, n, \operatorname{carr}(\operatorname{scal}(d)))) \wedge \\ \quad (q \in \operatorname{mat}(n, k, \operatorname{carr}(\operatorname{scal}(d)))) \wedge \\ \quad (u \in \operatorname{mat}(k, l, \operatorname{vec}(d))) \wedge \\ \quad (w \in \operatorname{interval}(1, m)) \wedge \\ \quad (a \in \operatorname{interval}(1, l)) \wedge \\ \quad (x \in \operatorname{interval}(1, k)) \Rightarrow \\ \quad \quad (\lambda j \in \operatorname{interval}(1, n).\; \\ \quad \quad \quad \operatorname{mul}(\operatorname{scal}(d))(\operatorname{entry}(p, w, j), \operatorname{entry}(q, j, x))) \\ \quad \quad \in\; (\operatorname{interval}(1, n) \to \operatorname{carr}(\operatorname{scal}(d))) \end{array}}}
\vnbentry{\texttt{mal-\allowbreak{}dist-\allowbreak{}type}}{\vnbstmt{\begin{array}{@{}l@{}} \forall d, m, n, k, l, p, q, u, w, a, x. \\ \quad \operatorname{is\text{-}module}(d) \wedge \\ \quad (p \in \operatorname{mat}(m, n, \operatorname{carr}(\operatorname{scal}(d)))) \wedge \\ \quad (q \in \operatorname{mat}(n, k, \operatorname{carr}(\operatorname{scal}(d)))) \wedge \\ \quad (u \in \operatorname{mat}(k, l, \operatorname{vec}(d))) \wedge \\ \quad (w \in \operatorname{interval}(1, m)) \wedge \\ \quad (a \in \operatorname{interval}(1, l)) \wedge \\ \quad (x \in \operatorname{interval}(1, k)) \Rightarrow \\ \quad \quad (\lambda z \in \operatorname{interval}(1, n).\; \\ \quad \quad \quad \operatorname{act}(d)((\lambda j \in \operatorname{interval}(1, n).\; \operatorname{mul}(\operatorname{scal}(d))(\operatorname{entry}(p, w, j), \operatorname{entry}(q, j, x)))(z), \\ \quad \quad \quad \quad \operatorname{entry}(u, x, a))) \\ \quad \quad \in\; (\operatorname{interval}(1, n) \to \operatorname{carr}(\operatorname{module\text{-}vector\text{-}ag}(d))) \end{array}}}
\vnbentry{\texttt{mal-\allowbreak{}red-\allowbreak{}type}}{\vnbstmt{\begin{array}{@{}l@{}} \forall d, m, n, k, l, p, q, u, w, a, x. \\ \quad \operatorname{is\text{-}module}(d) \wedge \\ \quad (p \in \operatorname{mat}(m, n, \operatorname{carr}(\operatorname{scal}(d)))) \wedge \\ \quad (q \in \operatorname{mat}(n, k, \operatorname{carr}(\operatorname{scal}(d)))) \wedge \\ \quad (u \in \operatorname{mat}(k, l, \operatorname{vec}(d))) \wedge \\ \quad (w \in \operatorname{interval}(1, m)) \wedge \\ \quad (a \in \operatorname{interval}(1, l)) \wedge \\ \quad (x \in \operatorname{interval}(1, k)) \Rightarrow \\ \quad \quad (\lambda j \in \operatorname{interval}(1, n).\; \\ \quad \quad \quad \operatorname{act}(d)(\operatorname{mul}(\operatorname{scal}(d))(\operatorname{entry}(p, w, j), \operatorname{entry}(q, j, x)), \\ \quad \quad \quad \quad \operatorname{entry}(u, x, a))) \\ \quad \quad \in\; (\operatorname{interval}(1, n) \to \operatorname{carr}(\operatorname{module\text{-}vector\text{-}ag}(d))) \end{array}}}
\vnbentry{\texttt{mar-\allowbreak{}outf-\allowbreak{}type}}{\vnbstmt{\begin{array}{@{}l@{}} \forall d, m, n, k, l, p, q, u, w, a. \\ \quad \operatorname{is\text{-}module}(d) \wedge \\ \quad (p \in \operatorname{mat}(m, n, \operatorname{carr}(\operatorname{scal}(d)))) \wedge \\ \quad (q \in \operatorname{mat}(n, k, \operatorname{carr}(\operatorname{scal}(d)))) \wedge \\ \quad (u \in \operatorname{mat}(k, l, \operatorname{vec}(d))) \wedge \\ \quad (w \in \operatorname{interval}(1, m)) \wedge \\ \quad (a \in \operatorname{interval}(1, l)) \wedge \\ \quad (1 \leq k) \Rightarrow \\ \quad \quad (\lambda j \in \operatorname{interval}(1, n).\; \\ \quad \quad \quad \operatorname{act}(d)(\operatorname{entry}(p, w, j), \operatorname{entry}(\operatorname{matact}(d, q, u), j, a))) \\ \quad \quad \in\; (\operatorname{interval}(1, n) \to \operatorname{carr}(\operatorname{module\text{-}vector\text{-}ag}(d))) \end{array}}}
\vnbentry{\texttt{mar-\allowbreak{}tout-\allowbreak{}type}}{\vnbstmt{\begin{array}{@{}l@{}} \forall d, m, n, k, l, p, q, u, w, a. \\ \quad \operatorname{is\text{-}module}(d) \wedge \\ \quad (p \in \operatorname{mat}(m, n, \operatorname{carr}(\operatorname{scal}(d)))) \wedge \\ \quad (q \in \operatorname{mat}(n, k, \operatorname{carr}(\operatorname{scal}(d)))) \wedge \\ \quad (u \in \operatorname{mat}(k, l, \operatorname{vec}(d))) \wedge \\ \quad (w \in \operatorname{interval}(1, m)) \wedge \\ \quad (a \in \operatorname{interval}(1, l)) \Rightarrow \\ \quad \quad (\lambda j \in \operatorname{interval}(1, n).\; \\ \quad \quad \quad \operatorname{finsum}(\operatorname{module\text{-}vector\text{-}ag}(d), \\ \quad \quad \quad \quad (\lambda c \in \operatorname{interval}(1, k).\; \\ \quad \quad \quad \quad \quad (\lambda z \in (\operatorname{interval}(1, k) \times \operatorname{interval}(1, n)).\; \operatorname{act}(d)(\operatorname{mul}(\operatorname{scal}(d))(\operatorname{entry}(p, w, \operatorname{nth}(2, z)), \operatorname{entry}(q, \operatorname{nth}(2, z), \operatorname{nth}(1, z))), \operatorname{entry}(u, \operatorname{nth}(1, z), a)))(\langle c, j \rangle)), \\ \quad \quad \quad \quad \operatorname{interval}(1, k))) \\ \quad \quad \in\; (\operatorname{interval}(1, n) \to \operatorname{carr}(\operatorname{module\text{-}vector\text{-}ag}(d))) \end{array}}}
\vnbentry{\texttt{mar-\allowbreak{}gj-\allowbreak{}type}}{\vnbstmt{\begin{array}{@{}l@{}} \forall d, m, n, k, l, p, q, u, w, a, x. \\ \quad \operatorname{is\text{-}module}(d) \wedge \\ \quad (p \in \operatorname{mat}(m, n, \operatorname{carr}(\operatorname{scal}(d)))) \wedge \\ \quad (q \in \operatorname{mat}(n, k, \operatorname{carr}(\operatorname{scal}(d)))) \wedge \\ \quad (u \in \operatorname{mat}(k, l, \operatorname{vec}(d))) \wedge \\ \quad (w \in \operatorname{interval}(1, m)) \wedge \\ \quad (a \in \operatorname{interval}(1, l)) \wedge \\ \quad (x \in \operatorname{interval}(1, n)) \Rightarrow \\ \quad \quad (\lambda j \in \operatorname{interval}(1, k).\; \\ \quad \quad \quad \operatorname{act}(d)(\operatorname{entry}(q, x, j), \operatorname{entry}(u, j, a))) \\ \quad \quad \in\; (\operatorname{interval}(1, k) \to \operatorname{carr}(\operatorname{module\text{-}vector\text{-}ag}(d))) \end{array}}}
\vnbentry{\texttt{mar-\allowbreak{}dist-\allowbreak{}type}}{\vnbstmt{\begin{array}{@{}l@{}} \forall d, m, n, k, l, p, q, u, w, a, x. \\ \quad \operatorname{is\text{-}module}(d) \wedge \\ \quad (p \in \operatorname{mat}(m, n, \operatorname{carr}(\operatorname{scal}(d)))) \wedge \\ \quad (q \in \operatorname{mat}(n, k, \operatorname{carr}(\operatorname{scal}(d)))) \wedge \\ \quad (u \in \operatorname{mat}(k, l, \operatorname{vec}(d))) \wedge \\ \quad (w \in \operatorname{interval}(1, m)) \wedge \\ \quad (a \in \operatorname{interval}(1, l)) \wedge \\ \quad (x \in \operatorname{interval}(1, n)) \Rightarrow \\ \quad \quad (\lambda z \in \operatorname{interval}(1, k).\; \\ \quad \quad \quad \operatorname{act}(d)(\operatorname{entry}(p, w, x), \\ \quad \quad \quad \quad (\lambda j \in \operatorname{interval}(1, k).\; \operatorname{act}(d)(\operatorname{entry}(q, x, j), \operatorname{entry}(u, j, a)))(z))) \\ \quad \quad \in\; (\operatorname{interval}(1, k) \to \operatorname{carr}(\operatorname{module\text{-}vector\text{-}ag}(d))) \end{array}}}
\vnbentry{\texttt{mar-\allowbreak{}red-\allowbreak{}type}}{\vnbstmt{\begin{array}{@{}l@{}} \forall d, m, n, k, l, p, q, u, w, a, x. \\ \quad \operatorname{is\text{-}module}(d) \wedge \\ \quad (p \in \operatorname{mat}(m, n, \operatorname{carr}(\operatorname{scal}(d)))) \wedge \\ \quad (q \in \operatorname{mat}(n, k, \operatorname{carr}(\operatorname{scal}(d)))) \wedge \\ \quad (u \in \operatorname{mat}(k, l, \operatorname{vec}(d))) \wedge \\ \quad (w \in \operatorname{interval}(1, m)) \wedge \\ \quad (a \in \operatorname{interval}(1, l)) \wedge \\ \quad (x \in \operatorname{interval}(1, n)) \Rightarrow \\ \quad \quad (\lambda c \in \operatorname{interval}(1, k).\; \\ \quad \quad \quad \operatorname{act}(d)(\operatorname{mul}(\operatorname{scal}(d))(\operatorname{entry}(p, w, x), \operatorname{entry}(q, x, c)), \\ \quad \quad \quad \quad \operatorname{entry}(u, c, a))) \\ \quad \quad \in\; (\operatorname{interval}(1, k) \to \operatorname{carr}(\operatorname{module\text{-}vector\text{-}ag}(d))) \end{array}}}
\vnbentry{\texttt{matact-\allowbreak{}assoc-\allowbreak{}summand-\allowbreak{}type}}{\vnbstmt{\begin{array}{@{}l@{}} \forall d, m, n, k, l, p, q, u, w, a. \\ \quad \operatorname{is\text{-}module}(d) \wedge \\ \quad (p \in \operatorname{mat}(m, n, \operatorname{carr}(\operatorname{scal}(d)))) \wedge \\ \quad (q \in \operatorname{mat}(n, k, \operatorname{carr}(\operatorname{scal}(d)))) \wedge \\ \quad (u \in \operatorname{mat}(k, l, \operatorname{vec}(d))) \wedge \\ \quad (w \in \operatorname{interval}(1, m)) \wedge \\ \quad (a \in \operatorname{interval}(1, l)) \Rightarrow \\ \quad \quad (\lambda z \in (\operatorname{interval}(1, k) \times \operatorname{interval}(1, n)).\; \\ \quad \quad \quad \operatorname{act}(d)(\operatorname{mul}(\operatorname{scal}(d))(\operatorname{entry}(p, w, \operatorname{nth}(2, z)), \operatorname{entry}(q, \operatorname{nth}(2, z), \operatorname{nth}(1, z))), \\ \quad \quad \quad \quad \operatorname{entry}(u, \operatorname{nth}(1, z), a))) \\ \quad \quad \in\; ((\operatorname{interval}(1, k) \times \operatorname{interval}(1, n)) \to \operatorname{carr}(\operatorname{module\text{-}vector\text{-}ag}(d))) \end{array}}}
\vnbentry{\texttt{mra-\allowbreak{}combined-\allowbreak{}summand-\allowbreak{}type}}{\vnbstmt{\begin{array}{@{}l@{}} \forall d, n, a, b, u. \\ \quad \operatorname{is\text{-}module}(d) \wedge \\ \quad (a \in \operatorname{mat}(1, n, \operatorname{carr}(\operatorname{scal}(d)))) \wedge \\ \quad (b \in \operatorname{mat}(1, n, \operatorname{carr}(\operatorname{scal}(d)))) \wedge \\ \quad (u \in \operatorname{mat}(n, 1, \operatorname{vec}(d))) \Rightarrow \\ \quad \quad (\lambda z \in \operatorname{interval}(1, n).\; \\ \quad \quad \quad \operatorname{opr}(\operatorname{module\text{-}vector\text{-}ag}(d))((\lambda j \in \operatorname{interval}(1, n).\; \operatorname{act}(d)(\operatorname{entry}(a, 1, j), \operatorname{entry}(u, j, 1)))(z), \\ \quad \quad \quad \quad (\lambda j \in \operatorname{interval}(1, n).\; \operatorname{act}(d)(\operatorname{entry}(b, 1, j), \operatorname{entry}(u, j, 1)))(z))) \\ \quad \quad \in\; (\operatorname{interval}(1, n) \to \operatorname{carr}(\operatorname{module\text{-}vector\text{-}ag}(d))) \end{array}}}
\vnbentry{\texttt{mrs-\allowbreak{}scaled-\allowbreak{}summand-\allowbreak{}type}}{\vnbstmt{\begin{array}{@{}l@{}} \forall d, n, r, c, u. \\ \quad \operatorname{is\text{-}module}(d) \wedge \\ \quad (r \in \operatorname{carr}(\operatorname{scal}(d))) \wedge \\ \quad (c \in \operatorname{mat}(1, n, \operatorname{carr}(\operatorname{scal}(d)))) \wedge \\ \quad (u \in \operatorname{mat}(n, 1, \operatorname{vec}(d))) \Rightarrow \\ \quad \quad (\lambda z \in \operatorname{interval}(1, n).\; \\ \quad \quad \quad \operatorname{act}(d)(r, \\ \quad \quad \quad \quad (\lambda j \in \operatorname{interval}(1, n).\; \operatorname{act}(d)(\operatorname{entry}(c, 1, j), \operatorname{entry}(u, j, 1)))(z))) \\ \quad \quad \in\; (\operatorname{interval}(1, n) \to \operatorname{carr}(\operatorname{module\text{-}vector\text{-}ag}(d))) \end{array}}}

\vnbfile{theorem-library/rake-identmat.scm}{theorem-library/rake-identmat.scm}
\noindent{\small rake-identmat.scm -- the two IDENTMAT identity laws and the module twin of matmul-entry, PROVEN (batch K, 2026-09-17).}\par
\vnbentry{\texttt{identmat-\allowbreak{}left-\allowbreak{}identity}}{\vnbstmt{\begin{array}{@{}l@{}} \forall a, m, n.; \forall p \in \operatorname{mat}(m, n, \operatorname{carr}(a)). \\ \quad \operatorname{is\text{-}ring}(a) \Rightarrow  \\ \quad \quad (\operatorname{matmul}(a, \operatorname{identmat}(a, m), p) = p) \end{array}}}
\vnbentry{\texttt{identmat-\allowbreak{}right-\allowbreak{}identity}}{\vnbstmt{\begin{array}{@{}l@{}} \forall a, m, n.; \forall p \in \operatorname{mat}(m, n, \operatorname{carr}(a)). \\ \quad \operatorname{is\text{-}ring}(a) \Rightarrow  \\ \quad \quad (\operatorname{matmul}(a, p, \operatorname{identmat}(a, n)) = p) \end{array}}}
\vnbentry{\texttt{matact-\allowbreak{}entry}}{\vnbstmt{\begin{array}{@{}l@{}} \forall d, m, n, q, p, u.; \forall i \in \operatorname{interval}(1, m), c \in \operatorname{interval}(1, q). \\ \quad \operatorname{is\text{-}module}(d) \wedge \\ \quad (p \in \operatorname{mat}(m, n, \operatorname{carr}(\operatorname{scal}(d)))) \wedge \\ \quad (u \in \operatorname{mat}(n, q, \operatorname{vec}(d))) \wedge \\ \quad (1 \leq n) \Rightarrow \\ \quad \quad \operatorname{entry}(\operatorname{matact}(d, p, u), i, c) \\ \quad \quad =\; \operatorname{finsum}(\operatorname{module\text{-}vector\text{-}ag}(d), \\ \quad \quad \quad (\lambda j \in \operatorname{interval}(1, n).\; \\ \quad \quad \quad \quad \operatorname{act}(d)(\operatorname{entry}(p, i, j), \operatorname{entry}(u, j, c))), \\ \quad \quad \quad \operatorname{interval}(1, n)) \end{array}}}

\vnbfile{theorem-library/rake-generates-coeff.scm}{theorem-library/rake-generates-coeff.scm}
\noindent{\small rake-generates-coeff.scm -- generates-coeff-matrix, proven.}\par
\vnbentry{\texttt{generates-\allowbreak{}coeff-\allowbreak{}matrix}}{\vnbstmt{\begin{array}{@{}l@{}} \forall d, n, m, u, v. \\ \quad \operatorname{is\text{-}module}(d) \wedge \\ \quad (u \in \operatorname{mat}(n, 1, \operatorname{vec}(d))) \wedge \\ \quad (v \in \operatorname{mat}(m, 1, \operatorname{vec}(d))) \wedge \\ \quad ((n = 0) \Rightarrow (m = 0)) \wedge \\ \quad \operatorname{generates}(d, n, u) \Rightarrow \\ \quad \quad \exists a \in \operatorname{mat}(m, n, \operatorname{carr}(\operatorname{scal}(d))). \\ \quad \quad \quad (v = \operatorname{matact}(d, a, u)) \end{array}}}

\vnbfile{theorem-library/monalg-laws.scm}{theorem-library/monalg-laws.scm}
\noindent{\small monalg-laws.scm -- the monoid-algebra laws that go through without reindexing, PROVEN.}\par
\vnbentry{\texttt{monoid-\allowbreak{}carrier-\allowbreak{}is-\allowbreak{}set}}{\vnbstmt{\begin{array}{@{}l@{}} \forall m. \\ \quad \operatorname{is\text{-}monoid}(m) \Rightarrow  \\ \quad \quad (\operatorname{carr}(m) \in \mathbf{Set}) \end{array}}}
\vnbentry{\texttt{monalg-\allowbreak{}add-\allowbreak{}apply}}{\vnbstmt{\begin{array}{@{}l@{}} \forall a, m, f, g.; \forall x \in \operatorname{carr}(m). \\ \quad (\operatorname{monalg\text{-}add}(a, m, f, g)(x) \simeq \operatorname{add}(a)(f(x), g(x))) \end{array}}}
\vnbentry{\texttt{monalg-\allowbreak{}add-\allowbreak{}fun}}{\vnbstmt{\begin{array}{@{}l@{}} \forall a, m, f, g. \\ \quad \operatorname{is\text{-}ring}(a) \wedge \\ \quad \operatorname{is\text{-}monoid}(m) \wedge \\ \quad (f \in \operatorname{finsupp}(a, m)) \wedge \\ \quad (g \in \operatorname{finsupp}(a, m)) \Rightarrow \\ \quad \quad (\operatorname{monalg\text{-}add}(a, m, f, g) \in \operatorname{finsupp}(a, m)) \end{array}}}

\vnbfile{theorem-library/poly-zero.scm}{theorem-library/poly-zero.scm}
\noindent{\small poly-zero.scm -- the ZERO of a monoid algebra, pointwise, and the EXTENSIONALITY BRIDGE for polynomials.  Nothing here is asserted; every statement is `proven modulo 0'.}\par
\vnbentry{\texttt{monalg-\allowbreak{}zero-\allowbreak{}apply}}{\vnbstmt{\begin{array}{@{}l@{}} \forall a, m.; \forall x \in \operatorname{carr}(m). \\ \quad (\operatorname{monalg\text{-}zero}(a, m)(x) \simeq \operatorname{zero}(a)) \end{array}}}
\vnbentry{\texttt{monalg-\allowbreak{}zero-\allowbreak{}fun}}{\vnbstmt{\begin{array}{@{}l@{}} \forall a, m. \\ \quad (\operatorname{carr}(m) \in \mathbf{Set}) \wedge \\ \quad (\operatorname{zero}(a) \in \operatorname{carr}(a)) \Rightarrow \\ \quad \quad (\operatorname{monalg\text{-}zero}(a, m) \in (\operatorname{carr}(m) \to \operatorname{carr}(a))) \end{array}}}
\vnbentry{\texttt{poly-\allowbreak{}zero-\allowbreak{}unfold}}{\vnbstmt{\begin{array}{@{}l@{}} \forall a. \\ \quad (\operatorname{zero}(\operatorname{poly}(a)) \simeq \operatorname{monalg\text{-}zero}(a, \operatorname{nn\text{-}add\text{-}monoid})) \end{array}}}
\vnbentry{\texttt{poly-\allowbreak{}zero-\allowbreak{}apply}}{\vnbstmt{\begin{array}{@{}l@{}} \forall a, k. \\ \quad (\operatorname{is\text{-}ring}(a) \wedge (k \in \mathbb{N})) \Rightarrow  \\ \quad \quad (\operatorname{zero}(\operatorname{poly}(a))(k) = \operatorname{zero}(a)) \end{array}}}
\vnbentry{\texttt{poly-\allowbreak{}zero-\allowbreak{}fun}}{\vnbstmt{\begin{array}{@{}l@{}} \forall a. \\ \quad \operatorname{is\text{-}ring}(a) \Rightarrow  \\ \quad \quad (\operatorname{zero}(\operatorname{poly}(a)) \in (\mathbb{N} \to \operatorname{carr}(a))) \end{array}}}
\vnbentry{\texttt{poly-\allowbreak{}zero-\allowbreak{}iff}}{\vnbstmt{\begin{array}{@{}l@{}} \forall a, p. \\ \quad (\operatorname{is\text{-}ring}(a) \wedge (p \in \operatorname{carr}(\operatorname{poly}(a)))) \Rightarrow  \\ \quad \quad ((p = \operatorname{zero}(\operatorname{poly}(a))) \Leftrightarrow \forall k \in \mathbb{N}.\; (p(k) = \operatorname{zero}(a))) \end{array}}}

\vnbfile{theorem-library/monalg-is-ring.scm}{theorem-library/monalg-is-ring.scm}
\noindent{\small monalg-is-ring.scm -- the monoid algebra A[M] is a ring (structure-library/polynomial.scm's umbrella seed), PROVEN, with every additive law, both distributive laws, both unit laws and the typing of the convolution proven on the way.  The one law still chained to its seed is the associativity of convolution, ...}\par
\vnbentry{\texttt{monalg-\allowbreak{}carr}}{\vnbstmt{\begin{array}{@{}l@{}} \forall a, m. \\ \quad (\operatorname{carr}(\operatorname{monalg}(a, m)) \simeq \operatorname{finsupp}(a, m)) \end{array}}}
\vnbentry{\texttt{monalg-\allowbreak{}add-\allowbreak{}op}}{\vnbstmt{\begin{array}{@{}l@{}} \forall a, m. \\ \quad (\operatorname{add}(\operatorname{monalg}(a, m)) \simeq (\lambda \langle f, g \rangle \in (\operatorname{finsupp}(a, m) \times \operatorname{finsupp}(a, m)).\; \operatorname{monalg\text{-}add}(a, m, f, g))) \end{array}}}
\vnbentry{\texttt{monalg-\allowbreak{}mul-\allowbreak{}op}}{\vnbstmt{\begin{array}{@{}l@{}} \forall a, m. \\ \quad (\operatorname{mul}(\operatorname{monalg}(a, m)) \simeq (\lambda \langle f, g \rangle \in (\operatorname{finsupp}(a, m) \times \operatorname{finsupp}(a, m)).\; \operatorname{monalg\text{-}mul}(a, m, f, g))) \end{array}}}
\vnbentry{\texttt{monalg-\allowbreak{}neg-\allowbreak{}op}}{\vnbstmt{\begin{array}{@{}l@{}} \forall a, m. \\ \quad (\operatorname{neg}(\operatorname{monalg}(a, m)) \simeq (\lambda f \in \operatorname{finsupp}(a, m).\; \operatorname{monalg\text{-}neg}(a, m, f))) \end{array}}}
\vnbentry{\texttt{monalg-\allowbreak{}zero-\allowbreak{}slot}}{\vnbstmt{\begin{array}{@{}l@{}} \forall a, m. \\ \quad (\operatorname{zero}(\operatorname{monalg}(a, m)) \simeq \operatorname{monalg\text{-}zero}(a, m)) \end{array}}}
\vnbentry{\texttt{monalg-\allowbreak{}one-\allowbreak{}slot}}{\vnbstmt{\begin{array}{@{}l@{}} \forall a, m. \\ \quad (\operatorname{one}(\operatorname{monalg}(a, m)) \simeq \operatorname{monalg\text{-}one}(a, m)) \end{array}}}
\vnbentry{\texttt{finsupp-\allowbreak{}in-\allowbreak{}set}}{\vnbstmt{\begin{array}{@{}l@{}} \forall a, m. \\ \quad (\operatorname{is\text{-}ring}(a) \wedge \operatorname{is\text{-}monoid}(m)) \Rightarrow  \\ \quad \quad (\operatorname{finsupp}(a, m) \in \mathbf{Set}) \end{array}}}
\vnbentry{\texttt{monalg-\allowbreak{}ext}}{\vnbstmt{\begin{array}{@{}l@{}} \forall a, m, f, g. \\ \quad (f \in \operatorname{finsupp}(a, m)) \wedge \\ \quad (g \in \operatorname{finsupp}(a, m)) \wedge \\ \quad \forall x \in \operatorname{carr}(m).\; (f(x) = g(x)) \Rightarrow \\ \quad \quad (f = g) \end{array}}}
\vnbentry{\texttt{monalg-\allowbreak{}neg-\allowbreak{}apply}}{\vnbstmt{\begin{array}{@{}l@{}} \forall a, m, f.; \forall x \in \operatorname{carr}(m). \\ \quad (\operatorname{monalg\text{-}neg}(a, m, f)(x) \simeq \operatorname{neg}(a)(f(x))) \end{array}}}
\vnbentry{\texttt{monalg-\allowbreak{}one-\allowbreak{}apply}}{\vnbstmt{\begin{array}{@{}l@{}} \forall a, m.; \forall x \in \operatorname{carr}(m). \\ \quad (\operatorname{monalg\text{-}one}(a, m)(x) \simeq \left(\text{if }(x = \operatorname{iden}(m))\text{ then }\operatorname{one}(a)\text{ else }\operatorname{zero}(a)\right)) \end{array}}}
\vnbentry{\texttt{ring-\allowbreak{}neg-\allowbreak{}zero}}{\vnbstmt{\begin{array}{@{}l@{}} \forall a. \\ \quad \operatorname{is\text{-}ring}(a) \Rightarrow  \\ \quad \quad (\operatorname{neg}(a)(\operatorname{zero}(a)) = \operatorname{zero}(a)) \end{array}}}
\vnbentry{\texttt{monalg-\allowbreak{}neg-\allowbreak{}fun}}{\vnbstmt{\begin{array}{@{}l@{}} \forall a, m, f. \\ \quad \operatorname{is\text{-}ring}(a) \wedge \\ \quad \operatorname{is\text{-}monoid}(m) \wedge \\ \quad (f \in \operatorname{finsupp}(a, m)) \Rightarrow \\ \quad \quad (\operatorname{monalg\text{-}neg}(a, m, f) \in \operatorname{finsupp}(a, m)) \end{array}}}
\vnbentry{\texttt{monalg-\allowbreak{}zero-\allowbreak{}in}}{\vnbstmt{\begin{array}{@{}l@{}} \forall a, m. \\ \quad (\operatorname{is\text{-}ring}(a) \wedge \operatorname{is\text{-}monoid}(m)) \Rightarrow  \\ \quad \quad (\operatorname{monalg\text{-}zero}(a, m) \in \operatorname{finsupp}(a, m)) \end{array}}}
\vnbentry{\texttt{monalg-\allowbreak{}one-\allowbreak{}in}}{\vnbstmt{\begin{array}{@{}l@{}} \forall a, m. \\ \quad (\operatorname{is\text{-}ring}(a) \wedge \operatorname{is\text{-}monoid}(m)) \Rightarrow  \\ \quad \quad (\operatorname{monalg\text{-}one}(a, m) \in \operatorname{finsupp}(a, m)) \end{array}}}
\vnbentry{\texttt{monalg-\allowbreak{}add-\allowbreak{}assoc}}{\vnbstmt{\begin{array}{@{}l@{}} \forall a, m, f, g, h. \\ \quad \operatorname{is\text{-}ring}(a) \wedge \\ \quad \operatorname{is\text{-}monoid}(m) \wedge \\ \quad (f \in \operatorname{finsupp}(a, m)) \wedge \\ \quad (g \in \operatorname{finsupp}(a, m)) \wedge \\ \quad (h \in \operatorname{finsupp}(a, m)) \Rightarrow \\ \quad \quad \operatorname{monalg\text{-}add}(a, m, \operatorname{monalg\text{-}add}(a, m, f, g), h) \\ \quad \quad =\; \operatorname{monalg\text{-}add}(a, m, f, \operatorname{monalg\text{-}add}(a, m, g, h)) \end{array}}}
\vnbentry{\texttt{monalg-\allowbreak{}add-\allowbreak{}comm}}{\vnbstmt{\begin{array}{@{}l@{}} \forall a, m, f, g. \\ \quad \operatorname{is\text{-}ring}(a) \wedge \\ \quad \operatorname{is\text{-}monoid}(m) \wedge \\ \quad (f \in \operatorname{finsupp}(a, m)) \wedge \\ \quad (g \in \operatorname{finsupp}(a, m)) \Rightarrow \\ \quad \quad (\operatorname{monalg\text{-}add}(a, m, f, g) = \operatorname{monalg\text{-}add}(a, m, g, f)) \end{array}}}
\vnbentry{\texttt{monalg-\allowbreak{}add-\allowbreak{}zero-\allowbreak{}left}}{\vnbstmt{\begin{array}{@{}l@{}} \forall a, m, f. \\ \quad \operatorname{is\text{-}ring}(a) \wedge \\ \quad \operatorname{is\text{-}monoid}(m) \wedge \\ \quad (f \in \operatorname{finsupp}(a, m)) \Rightarrow \\ \quad \quad (\operatorname{monalg\text{-}add}(a, m, \operatorname{monalg\text{-}zero}(a, m), f) = f) \end{array}}}
\vnbentry{\texttt{monalg-\allowbreak{}add-\allowbreak{}zero-\allowbreak{}right}}{\vnbstmt{\begin{array}{@{}l@{}} \forall a, m, f. \\ \quad \operatorname{is\text{-}ring}(a) \wedge \\ \quad \operatorname{is\text{-}monoid}(m) \wedge \\ \quad (f \in \operatorname{finsupp}(a, m)) \Rightarrow \\ \quad \quad (\operatorname{monalg\text{-}add}(a, m, f, \operatorname{monalg\text{-}zero}(a, m)) = f) \end{array}}}
\vnbentry{\texttt{monalg-\allowbreak{}add-\allowbreak{}neg-\allowbreak{}left}}{\vnbstmt{\begin{array}{@{}l@{}} \forall a, m, f. \\ \quad \operatorname{is\text{-}ring}(a) \wedge \\ \quad \operatorname{is\text{-}monoid}(m) \wedge \\ \quad (f \in \operatorname{finsupp}(a, m)) \Rightarrow \\ \quad \quad \operatorname{monalg\text{-}add}(a, m, \operatorname{monalg\text{-}neg}(a, m, f), f) \\ \quad \quad =\; \operatorname{monalg\text{-}zero}(a, m) \end{array}}}
\vnbentry{\texttt{monalg-\allowbreak{}add-\allowbreak{}neg-\allowbreak{}right}}{\vnbstmt{\begin{array}{@{}l@{}} \forall a, m, f. \\ \quad \operatorname{is\text{-}ring}(a) \wedge \\ \quad \operatorname{is\text{-}monoid}(m) \wedge \\ \quad (f \in \operatorname{finsupp}(a, m)) \Rightarrow \\ \quad \quad \operatorname{monalg\text{-}add}(a, m, f, \operatorname{monalg\text{-}neg}(a, m, f)) \\ \quad \quad =\; \operatorname{monalg\text{-}zero}(a, m) \end{array}}}
\vnbentry{\texttt{card-\allowbreak{}cartesian-\allowbreak{}nn}}{\vnbstmt{\begin{array}{@{}l@{}} \forall x, y. \\ \quad (x \in \mathbf{Set}) \wedge \\ \quad (|x| \in \mathbb{N}) \wedge \\ \quad (y \in \mathbf{Set}) \wedge \\ \quad (|y| \in \mathbb{N}) \Rightarrow \\ \quad \quad (|(x \times y)| \in \mathbb{N}) \end{array}}}
\vnbentry{\texttt{raag-\allowbreak{}carr}}{\vnbstmt{\begin{array}{@{}l@{}} \forall a. \\ \quad (\operatorname{carr}(\operatorname{ring\text{-}additive\text{-}ag}(a)) \simeq \operatorname{carr}(a)) \end{array}}}
\vnbentry{\texttt{raag-\allowbreak{}opr}}{\vnbstmt{\begin{array}{@{}l@{}} \forall a. \\ \quad (\operatorname{opr}(\operatorname{ring\text{-}additive\text{-}ag}(a)) \simeq \operatorname{add}(a)) \end{array}}}
\vnbentry{\texttt{raag-\allowbreak{}iden}}{\vnbstmt{\begin{array}{@{}l@{}} \forall a. \\ \quad (\operatorname{iden}(\operatorname{ring\text{-}additive\text{-}ag}(a)) \simeq \operatorname{zero}(a)) \end{array}}}
\vnbentry{\texttt{monalg-\allowbreak{}mul-\allowbreak{}apply}}{\vnbstmt{\begin{array}{@{}l@{}} \forall a, m, f, g.; \forall x \in \operatorname{carr}(m). \\ \quad (\operatorname{monalg\text{-}mul}(a, m, f, g)(x) \simeq \operatorname{finsum}(\operatorname{ring\text{-}additive\text{-}ag}(a), (\lambda p \in \{p \in (\operatorname{supp}(a, m, f) \times \operatorname{supp}(a, m, g)) : (\operatorname{opr}(m)(\operatorname{nth}(1, p), \operatorname{nth}(2, p)) = x)\}.\; \operatorname{mul}(a)(f(\operatorname{nth}(1, p)), g(\operatorname{nth}(2, p)))), \{p \in (\operatorname{supp}(a, m, f) \times \operatorname{supp}(a, m, g)) : (\operatorname{opr}(m)(\operatorname{nth}(1, p), \operatorname{nth}(2, p)) = x)\})) \end{array}}}
\vnbentry{\texttt{monalg-\allowbreak{}mul-\allowbreak{}value-\allowbreak{}in-\allowbreak{}carr}}{\vnbstmt{\begin{array}{@{}l@{}} \forall a, m, f, g, x. \\ \quad \operatorname{is\text{-}ring}(a) \wedge \\ \quad \operatorname{is\text{-}monoid}(m) \wedge \\ \quad (f \in \operatorname{finsupp}(a, m)) \wedge \\ \quad (g \in \operatorname{finsupp}(a, m)) \wedge \\ \quad (x \in \operatorname{carr}(m)) \Rightarrow \\ \quad \quad (\operatorname{monalg\text{-}mul}(a, m, f, g)(x) \in \operatorname{carr}(a)) \end{array}}}
\vnbentry{\texttt{monalg-\allowbreak{}mul-\allowbreak{}fun}}{\vnbstmt{\begin{array}{@{}l@{}} \forall a, m, f, g. \\ \quad \operatorname{is\text{-}ring}(a) \wedge \\ \quad \operatorname{is\text{-}monoid}(m) \wedge \\ \quad (f \in \operatorname{finsupp}(a, m)) \wedge \\ \quad (g \in \operatorname{finsupp}(a, m)) \Rightarrow \\ \quad \quad (\operatorname{monalg\text{-}mul}(a, m, f, g) \in \operatorname{finsupp}(a, m)) \end{array}}}
\vnbentry{\texttt{cartesian-\allowbreak{}pair-\allowbreak{}eq}}{\vnbstmt{\begin{array}{@{}l@{}} \forall x, a, b. \\ \quad (x \in (a \times b)) \Rightarrow  \\ \quad \quad (x = \langle \operatorname{nth}(1, x), \operatorname{nth}(2, x) \rangle) \end{array}}}
\vnbentry{\texttt{supp-\allowbreak{}add-\allowbreak{}subset}}{\vnbstmt{\begin{array}{@{}l@{}} \forall a, m, g, h, z. \\ \quad \operatorname{is\text{-}ring}(a) \wedge \\ \quad (g \in \operatorname{finsupp}(a, m)) \wedge \\ \quad (h \in \operatorname{finsupp}(a, m)) \wedge \\ \quad (z \in \operatorname{supp}(a, m, \operatorname{monalg\text{-}add}(a, m, g, h))) \Rightarrow \\ \quad \quad (z \in (\operatorname{supp}(a, m, g) \cup \operatorname{supp}(a, m, h))) \end{array}}}
\vnbentry{\texttt{monalg-\allowbreak{}distrib-\allowbreak{}left-\allowbreak{}at}}{\vnbstmt{\begin{array}{@{}l@{}} \forall a, m, f, g, h, x. \\ \quad \operatorname{is\text{-}ring}(a) \wedge \\ \quad \operatorname{is\text{-}monoid}(m) \wedge \\ \quad (f \in \operatorname{finsupp}(a, m)) \wedge \\ \quad (g \in \operatorname{finsupp}(a, m)) \wedge \\ \quad (h \in \operatorname{finsupp}(a, m)) \wedge \\ \quad (x \in \operatorname{carr}(m)) \Rightarrow \\ \quad \quad \operatorname{finsum}(\operatorname{ring\text{-}additive\text{-}ag}(a), \\ \quad \quad \quad (\lambda p \in \{p \in (\operatorname{supp}(a, m, f) \times \operatorname{supp}(a, m, \operatorname{monalg\text{-}add}(a, m, g, h))) : (\operatorname{opr}(m)(\operatorname{nth}(1, p), \operatorname{nth}(2, p)) = x)\}.\; \\ \quad \quad \quad \quad \operatorname{mul}(a)(f(\operatorname{nth}(1, p)), \operatorname{monalg\text{-}add}(a, m, g, h)(\operatorname{nth}(2, p)))), \\ \quad \quad \quad \{p \in (\operatorname{supp}(a, m, f) \times \operatorname{supp}(a, m, \operatorname{monalg\text{-}add}(a, m, g, h))) : (\operatorname{opr}(m)(\operatorname{nth}(1, p), \operatorname{nth}(2, p)) = x)\}) \\ \quad \quad =\; \operatorname{add}(a)(\operatorname{finsum}(\operatorname{ring\text{-}additive\text{-}ag}(a), \\ \quad \quad \quad \quad (\lambda p \in \{p \in (\operatorname{supp}(a, m, f) \times \operatorname{supp}(a, m, g)) : (\operatorname{opr}(m)(\operatorname{nth}(1, p), \operatorname{nth}(2, p)) = x)\}.\; \\ \quad \quad \quad \quad \quad \operatorname{mul}(a)(f(\operatorname{nth}(1, p)), g(\operatorname{nth}(2, p)))), \\ \quad \quad \quad \quad \{p \in (\operatorname{supp}(a, m, f) \times \operatorname{supp}(a, m, g)) : (\operatorname{opr}(m)(\operatorname{nth}(1, p), \operatorname{nth}(2, p)) = x)\}), \\ \quad \quad \quad \operatorname{finsum}(\operatorname{ring\text{-}additive\text{-}ag}(a), \\ \quad \quad \quad \quad (\lambda p \in \{p \in (\operatorname{supp}(a, m, f) \times \operatorname{supp}(a, m, h)) : (\operatorname{opr}(m)(\operatorname{nth}(1, p), \operatorname{nth}(2, p)) = x)\}.\; \\ \quad \quad \quad \quad \quad \operatorname{mul}(a)(f(\operatorname{nth}(1, p)), h(\operatorname{nth}(2, p)))), \\ \quad \quad \quad \quad \{p \in (\operatorname{supp}(a, m, f) \times \operatorname{supp}(a, m, h)) : (\operatorname{opr}(m)(\operatorname{nth}(1, p), \operatorname{nth}(2, p)) = x)\})) \end{array}}}
\vnbentry{\texttt{monalg-\allowbreak{}distrib-\allowbreak{}right-\allowbreak{}at}}{\vnbstmt{\begin{array}{@{}l@{}} \forall a, m, f, g, h, x. \\ \quad \operatorname{is\text{-}ring}(a) \wedge \\ \quad \operatorname{is\text{-}monoid}(m) \wedge \\ \quad (f \in \operatorname{finsupp}(a, m)) \wedge \\ \quad (g \in \operatorname{finsupp}(a, m)) \wedge \\ \quad (h \in \operatorname{finsupp}(a, m)) \wedge \\ \quad (x \in \operatorname{carr}(m)) \Rightarrow \\ \quad \quad \operatorname{finsum}(\operatorname{ring\text{-}additive\text{-}ag}(a), \\ \quad \quad \quad (\lambda p \in \{p \in (\operatorname{supp}(a, m, \operatorname{monalg\text{-}add}(a, m, f, g)) \times \operatorname{supp}(a, m, h)) : (\operatorname{opr}(m)(\operatorname{nth}(1, p), \operatorname{nth}(2, p)) = x)\}.\; \\ \quad \quad \quad \quad \operatorname{mul}(a)(\operatorname{monalg\text{-}add}(a, m, f, g)(\operatorname{nth}(1, p)), h(\operatorname{nth}(2, p)))), \\ \quad \quad \quad \{p \in (\operatorname{supp}(a, m, \operatorname{monalg\text{-}add}(a, m, f, g)) \times \operatorname{supp}(a, m, h)) : (\operatorname{opr}(m)(\operatorname{nth}(1, p), \operatorname{nth}(2, p)) = x)\}) \\ \quad \quad =\; \operatorname{add}(a)(\operatorname{finsum}(\operatorname{ring\text{-}additive\text{-}ag}(a), \\ \quad \quad \quad \quad (\lambda p \in \{p \in (\operatorname{supp}(a, m, f) \times \operatorname{supp}(a, m, h)) : (\operatorname{opr}(m)(\operatorname{nth}(1, p), \operatorname{nth}(2, p)) = x)\}.\; \\ \quad \quad \quad \quad \quad \operatorname{mul}(a)(f(\operatorname{nth}(1, p)), h(\operatorname{nth}(2, p)))), \\ \quad \quad \quad \quad \{p \in (\operatorname{supp}(a, m, f) \times \operatorname{supp}(a, m, h)) : (\operatorname{opr}(m)(\operatorname{nth}(1, p), \operatorname{nth}(2, p)) = x)\}), \\ \quad \quad \quad \operatorname{finsum}(\operatorname{ring\text{-}additive\text{-}ag}(a), \\ \quad \quad \quad \quad (\lambda p \in \{p \in (\operatorname{supp}(a, m, g) \times \operatorname{supp}(a, m, h)) : (\operatorname{opr}(m)(\operatorname{nth}(1, p), \operatorname{nth}(2, p)) = x)\}.\; \\ \quad \quad \quad \quad \quad \operatorname{mul}(a)(g(\operatorname{nth}(1, p)), h(\operatorname{nth}(2, p)))), \\ \quad \quad \quad \quad \{p \in (\operatorname{supp}(a, m, g) \times \operatorname{supp}(a, m, h)) : (\operatorname{opr}(m)(\operatorname{nth}(1, p), \operatorname{nth}(2, p)) = x)\})) \end{array}}}
\vnbentry{\texttt{monalg-\allowbreak{}one-\allowbreak{}supp-\allowbreak{}subset}}{\vnbstmt{\begin{array}{@{}l@{}} \forall a, m, z. \\ \quad \operatorname{is\text{-}ring}(a) \wedge \\ \quad \operatorname{is\text{-}monoid}(m) \wedge \\ \quad (z \in \operatorname{supp}(a, m, \operatorname{monalg\text{-}one}(a, m))) \Rightarrow \\ \quad \quad (z \in \{\operatorname{iden}(m)\}) \end{array}}}
\vnbentry{\texttt{monalg-\allowbreak{}one-\allowbreak{}left-\allowbreak{}at}}{\vnbstmt{\begin{array}{@{}l@{}} \forall a, m, f, x. \\ \quad \operatorname{is\text{-}ring}(a) \wedge \\ \quad \operatorname{is\text{-}monoid}(m) \wedge \\ \quad (f \in \operatorname{finsupp}(a, m)) \wedge \\ \quad (x \in \operatorname{carr}(m)) \Rightarrow \\ \quad \quad \operatorname{finsum}(\operatorname{ring\text{-}additive\text{-}ag}(a), \\ \quad \quad \quad (\lambda p \in \{p \in (\operatorname{supp}(a, m, \operatorname{monalg\text{-}one}(a, m)) \times \operatorname{supp}(a, m, f)) : (\operatorname{opr}(m)(\operatorname{nth}(1, p), \operatorname{nth}(2, p)) = x)\}.\; \\ \quad \quad \quad \quad \operatorname{mul}(a)(\operatorname{monalg\text{-}one}(a, m)(\operatorname{nth}(1, p)), f(\operatorname{nth}(2, p)))), \\ \quad \quad \quad \{p \in (\operatorname{supp}(a, m, \operatorname{monalg\text{-}one}(a, m)) \times \operatorname{supp}(a, m, f)) : (\operatorname{opr}(m)(\operatorname{nth}(1, p), \operatorname{nth}(2, p)) = x)\}) \\ \quad \quad =\; f(x) \end{array}}}
\vnbentry{\texttt{monalg-\allowbreak{}one-\allowbreak{}right-\allowbreak{}at}}{\vnbstmt{\begin{array}{@{}l@{}} \forall a, m, f, x. \\ \quad \operatorname{is\text{-}ring}(a) \wedge \\ \quad \operatorname{is\text{-}monoid}(m) \wedge \\ \quad (f \in \operatorname{finsupp}(a, m)) \wedge \\ \quad (x \in \operatorname{carr}(m)) \Rightarrow \\ \quad \quad \operatorname{finsum}(\operatorname{ring\text{-}additive\text{-}ag}(a), \\ \quad \quad \quad (\lambda p \in \{p \in (\operatorname{supp}(a, m, f) \times \operatorname{supp}(a, m, \operatorname{monalg\text{-}one}(a, m))) : (\operatorname{opr}(m)(\operatorname{nth}(1, p), \operatorname{nth}(2, p)) = x)\}.\; \\ \quad \quad \quad \quad \operatorname{mul}(a)(f(\operatorname{nth}(1, p)), \operatorname{monalg\text{-}one}(a, m)(\operatorname{nth}(2, p)))), \\ \quad \quad \quad \{p \in (\operatorname{supp}(a, m, f) \times \operatorname{supp}(a, m, \operatorname{monalg\text{-}one}(a, m))) : (\operatorname{opr}(m)(\operatorname{nth}(1, p), \operatorname{nth}(2, p)) = x)\}) \\ \quad \quad =\; f(x) \end{array}}}
\vnbentry{\texttt{monalg-\allowbreak{}distrib-\allowbreak{}left}}{\vnbstmt{\begin{array}{@{}l@{}} \forall a, m, f, g, h. \\ \quad \operatorname{is\text{-}ring}(a) \wedge \\ \quad \operatorname{is\text{-}monoid}(m) \wedge \\ \quad (f \in \operatorname{finsupp}(a, m)) \wedge \\ \quad (g \in \operatorname{finsupp}(a, m)) \wedge \\ \quad (h \in \operatorname{finsupp}(a, m)) \Rightarrow \\ \quad \quad \operatorname{monalg\text{-}mul}(a, m, f, \operatorname{monalg\text{-}add}(a, m, g, h)) \\ \quad \quad =\; \operatorname{monalg\text{-}add}(a, m, \operatorname{monalg\text{-}mul}(a, m, f, g), \operatorname{monalg\text{-}mul}(a, m, f, h)) \end{array}}}
\vnbentry{\texttt{monalg-\allowbreak{}distrib-\allowbreak{}right}}{\vnbstmt{\begin{array}{@{}l@{}} \forall a, m, f, g, h. \\ \quad \operatorname{is\text{-}ring}(a) \wedge \\ \quad \operatorname{is\text{-}monoid}(m) \wedge \\ \quad (f \in \operatorname{finsupp}(a, m)) \wedge \\ \quad (g \in \operatorname{finsupp}(a, m)) \wedge \\ \quad (h \in \operatorname{finsupp}(a, m)) \Rightarrow \\ \quad \quad \operatorname{monalg\text{-}mul}(a, m, \operatorname{monalg\text{-}add}(a, m, f, g), h) \\ \quad \quad =\; \operatorname{monalg\text{-}add}(a, m, \operatorname{monalg\text{-}mul}(a, m, f, h), \operatorname{monalg\text{-}mul}(a, m, g, h)) \end{array}}}
\vnbentry{\texttt{monalg-\allowbreak{}one-\allowbreak{}left-\allowbreak{}fn}}{\vnbstmt{\begin{array}{@{}l@{}} \forall a, m, f. \\ \quad \operatorname{is\text{-}ring}(a) \wedge \\ \quad \operatorname{is\text{-}monoid}(m) \wedge \\ \quad (f \in \operatorname{finsupp}(a, m)) \Rightarrow \\ \quad \quad (\operatorname{monalg\text{-}mul}(a, m, \operatorname{monalg\text{-}one}(a, m), f) = f) \end{array}}}
\vnbentry{\texttt{monalg-\allowbreak{}one-\allowbreak{}right-\allowbreak{}fn}}{\vnbstmt{\begin{array}{@{}l@{}} \forall a, m, f. \\ \quad \operatorname{is\text{-}ring}(a) \wedge \\ \quad \operatorname{is\text{-}monoid}(m) \wedge \\ \quad (f \in \operatorname{finsupp}(a, m)) \Rightarrow \\ \quad \quad (\operatorname{monalg\text{-}mul}(a, m, f, \operatorname{monalg\text{-}one}(a, m)) = f) \end{array}}}
\vnbentry{\texttt{monalg-\allowbreak{}one-\allowbreak{}left}}{\vnbstmt{\begin{array}{@{}l@{}} \forall a, m, f, x. \\ \quad \operatorname{is\text{-}ring}(a) \wedge \\ \quad \operatorname{is\text{-}monoid}(m) \wedge \\ \quad (f \in \operatorname{finsupp}(a, m)) \wedge \\ \quad (x \in \operatorname{carr}(m)) \Rightarrow \\ \quad \quad (\operatorname{monalg\text{-}mul}(a, m, \operatorname{monalg\text{-}one}(a, m), f)(x) = f(x)) \end{array}}}
\vnbentry{\texttt{bijection-\allowbreak{}from-\allowbreak{}inverse}}{\vnbstmt{\begin{array}{@{}l@{}} \forall x, y, i, a. \\ \quad (i \in (x \to y)) \wedge \\ \quad (a \in (y \to x)) \wedge \\ \quad \forall u \in x.\; (a(i(u)) = u) \wedge \\ \quad \forall v \in y.\; (i(a(v)) = v) \Rightarrow \\ \quad \quad (i \in \mathrm{Bij}(x, y)) \end{array}}}
\vnbentry{\texttt{lambda-\allowbreak{}compose-\allowbreak{}value}}{\vnbstmt{\begin{array}{@{}l@{}} \forall f, h, m.; \forall t \in m. \\ \quad ((\lambda z \in m.\; f(h(z)))(t) \simeq f(h(t))) \end{array}}}
\vnbentry{\texttt{finsum-\allowbreak{}reindex-\allowbreak{}inverse}}{\vnbstmt{\begin{array}{@{}l@{}} \forall g.; \forall s, t \in \mathbf{Set}, i \in (t \to s), a \in (s \to t), f \in (s \to \operatorname{carr}(g)). \\ \quad \operatorname{is\text{-}abelian\text{-}group}(g) \wedge \\ \quad (|s| \in \mathbb{N}) \wedge \\ \quad (|t| \in \mathbb{N}) \wedge \\ \quad \forall u \in t.\; (a(i(u)) = u) \wedge \\ \quad \forall v \in s.\; (i(a(v)) = v) \Rightarrow \\ \quad \quad (\operatorname{finsum}(g, f, s) = \operatorname{finsum}(g, (\lambda z \in t.\; f(i(z))), t)) \end{array}}}
\vnbentry{\texttt{finsum-\allowbreak{}reindex-\allowbreak{}inverse-\allowbreak{}ptwise}}{\vnbstmt{\begin{array}{@{}l@{}} \forall g.; \forall s, t \in \mathbf{Set}, i \in (t \to s), a \in (s \to t), f. \\ \quad \operatorname{is\text{-}abelian\text{-}group}(g) \wedge \\ \quad (|s| \in \mathbb{N}) \wedge \\ \quad (|t| \in \mathbb{N}) \wedge \\ \quad \forall u \in t.\; (a(i(u)) = u) \wedge \\ \quad \forall v \in s.\; (i(a(v)) = v) \wedge \\ \quad \forall b \in s.\; (f(b) \in \operatorname{carr}(g)) \Rightarrow \\ \quad \quad (\operatorname{finsum}(g, f, s) = \operatorname{finsum}(g, (\lambda z \in t.\; f(i(z))), t)) \end{array}}}
\vnbentry{\texttt{finsum-\allowbreak{}fiber-\allowbreak{}value}}{\vnbstmt{\begin{array}{@{}l@{}} \forall g, f, h, s, t.; \forall a \in t. \\ \quad ((\lambda b \in t.\; \operatorname{finsum}(g, f, \{c \in s : (h(c) = b)\}))(a) \simeq \operatorname{finsum}(g, f, \{c \in s : (h(c) = a)\})) \end{array}}}
\vnbentry{\texttt{finsum-\allowbreak{}fiber-\allowbreak{}slice}}{\vnbstmt{\begin{array}{@{}l@{}} \forall g.; \forall s, k \in \mathbf{Set}, h, p, f, l, e. \\ \quad \operatorname{is\text{-}abelian\text{-}group}(g) \wedge \\ \quad (|s| \in \mathbb{N}) \wedge \\ \quad (|k| \in \mathbb{N}) \wedge \\ \quad (l \in (k \to s)) \wedge \\ \quad \forall a \in k.\; (h(l(a)) = p) \wedge \\ \quad \forall w \in s.\; ((h(w) = p) \Rightarrow (e(w) \in k)) \wedge \\ \quad \forall a \in k.\; (e(l(a)) = a) \wedge \\ \quad \forall w \in s.\; ((h(w) = p) \Rightarrow (l(e(w)) = w)) \wedge \\ \quad \forall b \in s.\; (f(b) \in \operatorname{carr}(g)) \Rightarrow \\ \quad \quad \operatorname{finsum}(g, f, \{c \in s : (h(c) = p)\}) \\ \quad \quad =\; \operatorname{finsum}(g, (\lambda z \in k.\; f(l(z))), k) \end{array}}}
\vnbentry{\texttt{monalg-\allowbreak{}mul-\allowbreak{}supp-\allowbreak{}in-\allowbreak{}image}}{\vnbstmt{\begin{array}{@{}l@{}} \forall a, m, f, g, z. \\ \quad \operatorname{is\text{-}ring}(a) \wedge \\ \quad \operatorname{is\text{-}monoid}(m) \wedge \\ \quad (f \in \operatorname{finsupp}(a, m)) \wedge \\ \quad (g \in \operatorname{finsupp}(a, m)) \wedge \\ \quad (z \in \operatorname{supp}(a, m, \operatorname{monalg\text{-}mul}(a, m, f, g))) \Rightarrow \\ \quad \quad z \\ \quad \quad \in\; \operatorname{image}((\lambda q \in (\operatorname{supp}(a, m, f) \times \operatorname{supp}(a, m, g)).\; \\ \quad \quad \quad \quad \operatorname{opr}(m)(\operatorname{nth}(1, q), \operatorname{nth}(2, q))), \\ \quad \quad \quad (\operatorname{supp}(a, m, f) \times \operatorname{supp}(a, m, g))) \end{array}}}
\vnbentry{\texttt{monalg-\allowbreak{}mul-\allowbreak{}assoc-\allowbreak{}left-\allowbreak{}core}}{\vnbstmt{\begin{array}{@{}l@{}} \forall a, m, f, g, h, x. \\ \quad \operatorname{is\text{-}ring}(a) \wedge \\ \quad \operatorname{is\text{-}monoid}(m) \wedge \\ \quad (f \in \operatorname{finsupp}(a, m)) \wedge \\ \quad (g \in \operatorname{finsupp}(a, m)) \wedge \\ \quad (h \in \operatorname{finsupp}(a, m)) \wedge \\ \quad (x \in \operatorname{carr}(m)) \Rightarrow \\ \quad \quad \operatorname{finsum}(\operatorname{ring\text{-}additive\text{-}ag}(a), \\ \quad \quad \quad (\lambda p \in \{p \in (\operatorname{image}((\lambda q \in (\operatorname{supp}(a, m, f) \times \operatorname{supp}(a, m, g)).\; \operatorname{opr}(m)(\operatorname{nth}(1, q), \operatorname{nth}(2, q))), (\operatorname{supp}(a, m, f) \times \operatorname{supp}(a, m, g))) \times \operatorname{supp}(a, m, h)) : (\operatorname{opr}(m)(\operatorname{nth}(1, p), \operatorname{nth}(2, p)) = x)\}.\; \\ \quad \quad \quad \quad \operatorname{mul}(a)(\operatorname{monalg\text{-}mul}(a, m, f, g)(\operatorname{nth}(1, p)), h(\operatorname{nth}(2, p)))), \\ \quad \quad \quad \{p \in (\operatorname{image}((\lambda q \in (\operatorname{supp}(a, m, f) \times \operatorname{supp}(a, m, g)).\; \operatorname{opr}(m)(\operatorname{nth}(1, q), \operatorname{nth}(2, q))), (\operatorname{supp}(a, m, f) \times \operatorname{supp}(a, m, g))) \times \operatorname{supp}(a, m, h)) : (\operatorname{opr}(m)(\operatorname{nth}(1, p), \operatorname{nth}(2, p)) = x)\}) \\ \quad \quad =\; \operatorname{finsum}(\operatorname{ring\text{-}additive\text{-}ag}(a), \\ \quad \quad \quad (\lambda w \in \{w \in ((\operatorname{supp}(a, m, f) \times \operatorname{supp}(a, m, g)) \times \operatorname{supp}(a, m, h)) : (\operatorname{opr}(m)(\operatorname{opr}(m)(\operatorname{nth}(1, \operatorname{nth}(1, w)), \operatorname{nth}(2, \operatorname{nth}(1, w))), \operatorname{nth}(2, w)) = x)\}.\; \\ \quad \quad \quad \quad \operatorname{mul}(a)(\operatorname{mul}(a)(f(\operatorname{nth}(1, \operatorname{nth}(1, w))), g(\operatorname{nth}(2, \operatorname{nth}(1, w)))), \\ \quad \quad \quad \quad \quad h(\operatorname{nth}(2, w)))), \\ \quad \quad \quad \{w \in ((\operatorname{supp}(a, m, f) \times \operatorname{supp}(a, m, g)) \times \operatorname{supp}(a, m, h)) : (\operatorname{opr}(m)(\operatorname{opr}(m)(\operatorname{nth}(1, \operatorname{nth}(1, w)), \operatorname{nth}(2, \operatorname{nth}(1, w))), \operatorname{nth}(2, w)) = x)\}) \end{array}}}
\vnbentry{\texttt{monalg-\allowbreak{}mul-\allowbreak{}assoc-\allowbreak{}right-\allowbreak{}core}}{\vnbstmt{\begin{array}{@{}l@{}} \forall a, m, f, g, h, x. \\ \quad \operatorname{is\text{-}ring}(a) \wedge \\ \quad \operatorname{is\text{-}monoid}(m) \wedge \\ \quad (f \in \operatorname{finsupp}(a, m)) \wedge \\ \quad (g \in \operatorname{finsupp}(a, m)) \wedge \\ \quad (h \in \operatorname{finsupp}(a, m)) \wedge \\ \quad (x \in \operatorname{carr}(m)) \Rightarrow \\ \quad \quad \operatorname{finsum}(\operatorname{ring\text{-}additive\text{-}ag}(a), \\ \quad \quad \quad (\lambda p \in \{p \in (\operatorname{supp}(a, m, f) \times \operatorname{image}((\lambda q \in (\operatorname{supp}(a, m, g) \times \operatorname{supp}(a, m, h)).\; \operatorname{opr}(m)(\operatorname{nth}(1, q), \operatorname{nth}(2, q))), (\operatorname{supp}(a, m, g) \times \operatorname{supp}(a, m, h)))) : (\operatorname{opr}(m)(\operatorname{nth}(1, p), \operatorname{nth}(2, p)) = x)\}.\; \\ \quad \quad \quad \quad \operatorname{mul}(a)(f(\operatorname{nth}(1, p)), \operatorname{monalg\text{-}mul}(a, m, g, h)(\operatorname{nth}(2, p)))), \\ \quad \quad \quad \{p \in (\operatorname{supp}(a, m, f) \times \operatorname{image}((\lambda q \in (\operatorname{supp}(a, m, g) \times \operatorname{supp}(a, m, h)).\; \operatorname{opr}(m)(\operatorname{nth}(1, q), \operatorname{nth}(2, q))), (\operatorname{supp}(a, m, g) \times \operatorname{supp}(a, m, h)))) : (\operatorname{opr}(m)(\operatorname{nth}(1, p), \operatorname{nth}(2, p)) = x)\}) \\ \quad \quad =\; \operatorname{finsum}(\operatorname{ring\text{-}additive\text{-}ag}(a), \\ \quad \quad \quad (\lambda w \in \{w \in ((\operatorname{supp}(a, m, f) \times \operatorname{supp}(a, m, g)) \times \operatorname{supp}(a, m, h)) : (\operatorname{opr}(m)(\operatorname{opr}(m)(\operatorname{nth}(1, \operatorname{nth}(1, w)), \operatorname{nth}(2, \operatorname{nth}(1, w))), \operatorname{nth}(2, w)) = x)\}.\; \\ \quad \quad \quad \quad \operatorname{mul}(a)(f(\operatorname{nth}(1, \operatorname{nth}(1, w))), \\ \quad \quad \quad \quad \quad \operatorname{mul}(a)(g(\operatorname{nth}(2, \operatorname{nth}(1, w))), h(\operatorname{nth}(2, w))))), \\ \quad \quad \quad \{w \in ((\operatorname{supp}(a, m, f) \times \operatorname{supp}(a, m, g)) \times \operatorname{supp}(a, m, h)) : (\operatorname{opr}(m)(\operatorname{opr}(m)(\operatorname{nth}(1, \operatorname{nth}(1, w)), \operatorname{nth}(2, \operatorname{nth}(1, w))), \operatorname{nth}(2, w)) = x)\}) \end{array}}}
\vnbentry{\texttt{monalg-\allowbreak{}mul-\allowbreak{}assoc-\allowbreak{}at}}{\vnbstmt{\begin{array}{@{}l@{}} \forall a, m, f, g, h, x. \\ \quad \operatorname{is\text{-}ring}(a) \wedge \\ \quad \operatorname{is\text{-}monoid}(m) \wedge \\ \quad (f \in \operatorname{finsupp}(a, m)) \wedge \\ \quad (g \in \operatorname{finsupp}(a, m)) \wedge \\ \quad (h \in \operatorname{finsupp}(a, m)) \wedge \\ \quad (x \in \operatorname{carr}(m)) \Rightarrow \\ \quad \quad \operatorname{monalg\text{-}mul}(a, m, \operatorname{monalg\text{-}mul}(a, m, f, g), h)(x) \\ \quad \quad =\; \operatorname{monalg\text{-}mul}(a, m, f, \operatorname{monalg\text{-}mul}(a, m, g, h))(x) \end{array}}}
\vnbentry{\texttt{monalg-\allowbreak{}mul-\allowbreak{}assoc}}{\vnbstmt{\begin{array}{@{}l@{}} \forall a, m, f, g, h. \\ \quad \operatorname{is\text{-}ring}(a) \wedge \\ \quad \operatorname{is\text{-}monoid}(m) \wedge \\ \quad (f \in \operatorname{finsupp}(a, m)) \wedge \\ \quad (g \in \operatorname{finsupp}(a, m)) \wedge \\ \quad (h \in \operatorname{finsupp}(a, m)) \Rightarrow \\ \quad \quad \operatorname{monalg\text{-}mul}(a, m, \operatorname{monalg\text{-}mul}(a, m, f, g), h) \\ \quad \quad =\; \operatorname{monalg\text{-}mul}(a, m, f, \operatorname{monalg\text{-}mul}(a, m, g, h)) \end{array}}}
\vnbentry{\texttt{monalg-\allowbreak{}is-\allowbreak{}ring}}{\vnbstmt{\begin{array}{@{}l@{}} \forall a, m. \\ \quad (\operatorname{is\text{-}ring}(a) \wedge \operatorname{is\text{-}monoid}(m)) \Rightarrow  \\ \quad \quad \operatorname{is\text{-}ring}(\operatorname{monalg}(a, m)) \end{array}}}

\vnbfile{theorem-library/rake-prod-of-sums.scm}{theorem-library/rake-prod-of-sums.scm}
\noindent{\small rake-prod-of-sums.scm -- the product-of-sums expansion, proven.}\par
\vnbentry{\texttt{difference-\allowbreak{}empty-\allowbreak{}left}}{\vnbstmt{\forall s.\; (\operatorname{difference}(\emptyset, s) = \emptyset)}}
\vnbentry{\texttt{insert-\allowbreak{}remove}}{\vnbstmt{\begin{array}{@{}l@{}} \forall s.; \forall k \in \mathbf{Set}. \\ \quad \neg (k \in s) \Rightarrow  \\ \quad \quad (\operatorname{difference}((s \cup \{k\}), \{k\}) = s) \end{array}}}
\vnbentry{\texttt{prod-\allowbreak{}ring-\allowbreak{}type-\allowbreak{}ptwise}}{\vnbstmt{\begin{array}{@{}l@{}} \forall r.; \forall x \in \mathbf{Set}, f. \\ \quad \operatorname{is\text{-}commutative\text{-}ring}(r) \wedge \\ \quad (|x| \in \mathbb{N}) \wedge \\ \quad \forall z \in x.\; (f(z) \in \operatorname{carr}(r)) \Rightarrow \\ \quad \quad (\operatorname{prod\text{-}ring}(r, f, x) \in \operatorname{carr}(r)) \end{array}}}
\vnbentry{\texttt{prod-\allowbreak{}ring-\allowbreak{}insert-\allowbreak{}ptwise}}{\vnbstmt{\begin{array}{@{}l@{}} \forall r.; \forall x, k \in \mathbf{Set}, f. \\ \quad \operatorname{is\text{-}commutative\text{-}ring}(r) \wedge \\ \quad (|x| \in \mathbb{N}) \wedge \\ \quad \neg (k \in x) \wedge \\ \quad \forall z \in (x \cup \{k\}).\; (f(z) \in \operatorname{carr}(r)) \Rightarrow \\ \quad \quad \operatorname{prod\text{-}ring}(r, f, (x \cup \{k\})) \\ \quad \quad =\; \operatorname{mul}(r)(\operatorname{prod\text{-}ring}(r, f, x), f(k)) \end{array}}}
\vnbentry{\texttt{prod-\allowbreak{}ring-\allowbreak{}congruence-\allowbreak{}q}}{\vnbstmt{\begin{array}{@{}l@{}} \forall r.; \forall x \in \mathbf{Set}, f, g. \\ \quad \operatorname{is\text{-}commutative\text{-}ring}(r) \wedge \\ \quad (|x| \in \mathbb{N}) \wedge \\ \quad \forall z \in x.\; (f(z) = g(z)) \Rightarrow \\ \quad \quad (\operatorname{prod\text{-}ring}(r, f, x) \simeq \operatorname{prod\text{-}ring}(r, g, x)) \end{array}}}
\vnbentry{\texttt{cra-\allowbreak{}carr}}{\vnbstmt{\begin{array}{@{}l@{}} \forall r. \\ \quad \operatorname{is\text{-}commutative\text{-}ring}(r) \Rightarrow  \\ \quad \quad \operatorname{carr}(\operatorname{commutative\text{-}ring\text{-}additive\text{-}ag}(r)) \\ \quad \quad =\; \operatorname{carr}(r) \end{array}}}
\vnbentry{\texttt{cra-\allowbreak{}op}}{\vnbstmt{\begin{array}{@{}l@{}} \forall r. \\ \quad \operatorname{is\text{-}commutative\text{-}ring}(r) \Rightarrow  \\ \quad \quad \operatorname{opr}(\operatorname{commutative\text{-}ring\text{-}additive\text{-}ag}(r)) \\ \quad \quad =\; \operatorname{add}(r) \end{array}}}
\vnbentry{\texttt{difference-\allowbreak{}in-\allowbreak{}power}}{\vnbstmt{\begin{array}{@{}l@{}} \forall w \in \mathbf{Set}, s. \\ \quad (\operatorname{difference}(w, s) \in \operatorname{power}(w)) \end{array}}}
\vnbentry{\texttt{prod-\allowbreak{}ring-\allowbreak{}type-\allowbreak{}subset}}{\vnbstmt{\begin{array}{@{}l@{}} \forall r.; \forall w \in \mathbf{Set}, f.; \forall s \in \operatorname{power}(w). \\ \quad \operatorname{is\text{-}commutative\text{-}ring}(r) \wedge \\ \quad (|w| \in \mathbb{N}) \wedge \\ \quad \forall z \in w.\; (f(z) \in \operatorname{carr}(r)) \Rightarrow \\ \quad \quad (\operatorname{prod\text{-}ring}(r, f, s) \in \operatorname{carr}(r)) \end{array}}}
\vnbentry{\texttt{prod-\allowbreak{}of-\allowbreak{}sums-\allowbreak{}summand-\allowbreak{}type}}{\vnbstmt{\begin{array}{@{}l@{}} \forall r.; \forall w \in \mathbf{Set}, a, b.; \forall s \in \operatorname{power}(w). \\ \quad \operatorname{is\text{-}commutative\text{-}ring}(r) \wedge \\ \quad (|w| \in \mathbb{N}) \wedge \\ \quad \forall z \in w.\; (a(z) \in \operatorname{carr}(r)) \wedge \\ \quad \forall z \in w.\; (b(z) \in \operatorname{carr}(r)) \Rightarrow \\ \quad \quad \operatorname{mul}(r)(\operatorname{prod\text{-}ring}(r, a, s), \operatorname{prod\text{-}ring}(r, b, \operatorname{difference}(w, s))) \\ \quad \quad \in\; \operatorname{carr}(r) \end{array}}}
\vnbentry{\texttt{power-\allowbreak{}insert-\allowbreak{}injective}}{\vnbstmt{\begin{array}{@{}l@{}} \forall x, k \in \mathbf{Set}. \\ \quad \neg (k \in x) \Rightarrow  \\ \quad \quad (\lambda s \in \operatorname{power}(x).\; (s \cup \{k\})) \\ \quad \quad \in\; \operatorname{injection}(\operatorname{power}(x), \operatorname{power}((x \cup \{k\}))) \end{array}}}
\vnbentry{\texttt{lambda-\allowbreak{}eta-\allowbreak{}value}}{\vnbstmt{\forall f, m.; \forall t \in m.\; ((\lambda z \in m.\; f(z))(t) \simeq f(t))}}
\vnbentry{\texttt{prod-\allowbreak{}of-\allowbreak{}sums-\allowbreak{}ptwise}}{\vnbstmt{\begin{array}{@{}l@{}} \forall r.; \forall w \in \mathbf{Set}, a, b. \\ \quad \operatorname{is\text{-}commutative\text{-}ring}(r) \wedge \\ \quad (|w| \in \mathbb{N}) \wedge \\ \quad \forall z \in w.\; (a(z) \in \operatorname{carr}(r)) \wedge \\ \quad \forall z \in w.\; (b(z) \in \operatorname{carr}(r)) \Rightarrow \\ \quad \quad \operatorname{prod\text{-}ring}(r, (\lambda v \in w.\; \operatorname{add}(r)(a(v), b(v))), w) \\ \quad \quad =\; \operatorname{finsum}(\operatorname{commutative\text{-}ring\text{-}additive\text{-}ag}(r), \\ \quad \quad \quad (\lambda \mathit{sv} \in \operatorname{power}(w).\; \\ \quad \quad \quad \quad \operatorname{mul}(r)(\operatorname{prod\text{-}ring}(r, a, \mathit{sv}), \operatorname{prod\text{-}ring}(r, b, \operatorname{difference}(w, \mathit{sv})))), \\ \quad \quad \quad \operatorname{power}(w)) \end{array}}}
\vnbentry{\texttt{prod-\allowbreak{}of-\allowbreak{}sums-\allowbreak{}expansion}}{\vnbstmt{\begin{array}{@{}l@{}} \forall r.; \forall x \in \mathbf{Set}, a, b \in (x \to \operatorname{carr}(r)). \\ \quad (\operatorname{is\text{-}commutative\text{-}ring}(r) \wedge (|x| \in \mathbb{N})) \Rightarrow  \\ \quad \quad \operatorname{prod\text{-}ring}(r, (\lambda k \in x.\; \operatorname{add}(r)(a(k), b(k))), x) \\ \quad \quad =\; \operatorname{finsum}(\operatorname{commutative\text{-}ring\text{-}additive\text{-}ag}(r), \\ \quad \quad \quad (\lambda s \in \operatorname{power}(x).\; \\ \quad \quad \quad \quad \operatorname{mul}(r)(\operatorname{prod\text{-}ring}(r, a, s), \operatorname{prod\text{-}ring}(r, b, \operatorname{difference}(x, s)))), \\ \quad \quad \quad \operatorname{power}(x)) \end{array}}}

\vnbfile{theorem-library/rake-smith-leaves.scm}{theorem-library/rake-smith-leaves.scm}
\noindent{\small rake-smith-leaves.scm -- finsum-interval-shift, PROVEN.}\par
\vnbentry{\texttt{finsum-\allowbreak{}interval-\allowbreak{}shift}}{\vnbstmt{\begin{array}{@{}l@{}} \forall g.; \forall n \in \mathbb{N}, f \in (\operatorname{interval}(1, (n + 1)) \to \operatorname{carr}(g)). \\ \quad \operatorname{is\text{-}abelian\text{-}group}(g) \Rightarrow  \\ \quad \quad \operatorname{finsum}(g, f, \operatorname{interval}(1, (n + 1))) \\ \quad \quad =\; \operatorname{opr}(g)(f(1), \\ \quad \quad \quad \operatorname{finsum}(g, (\lambda z \in \operatorname{interval}(1, n).\; f((z + 1))), \operatorname{interval}(1, n))) \end{array}}}

\vnbfile{theorem-library/rake-monalg-comm.scm}{theorem-library/rake-monalg-comm.scm}
\noindent{\small theorem-library/rake-monalg-comm.scm -- rake batch 5b, assignment 5b-A.}\par
\vnbentry{\texttt{monalg-\allowbreak{}index-\allowbreak{}swap-\allowbreak{}fun}}{\vnbstmt{\begin{array}{@{}l@{}} \forall a, m, f, g, x. \\ \quad \operatorname{is\text{-}comm\text{-}monoid}(m) \Rightarrow  \\ \quad \quad (\lambda q \in \{p \in (\operatorname{supp}(a, m, f) \times \operatorname{supp}(a, m, g)) : (\operatorname{opr}(m)(\operatorname{nth}(1, p), \operatorname{nth}(2, p)) = x)\}.\; \\ \quad \quad \quad \langle \operatorname{nth}(2, q), \operatorname{nth}(1, q) \rangle) \\ \quad \quad \in\; (\{p \in (\operatorname{supp}(a, m, f) \times \operatorname{supp}(a, m, g)) : (\operatorname{opr}(m)(\operatorname{nth}(1, p), \operatorname{nth}(2, p)) = x)\} \to \{p \in (\operatorname{supp}(a, m, g) \times \operatorname{supp}(a, m, f)) : (\operatorname{opr}(m)(\operatorname{nth}(1, p), \operatorname{nth}(2, p)) = x)\}) \end{array}}}
\vnbentry{\texttt{monalg-\allowbreak{}index-\allowbreak{}swap-\allowbreak{}inverse}}{\vnbstmt{\begin{array}{@{}l@{}} \forall a, m, f, g, x.; \forall u \in \{p \in (\operatorname{supp}(a, m, f) \times \operatorname{supp}(a, m, g)) : (\operatorname{opr}(m)(\operatorname{nth}(1, p), \operatorname{nth}(2, p)) = x)\}. \\ \quad \operatorname{is\text{-}comm\text{-}monoid}(m) \Rightarrow  \\ \quad \quad (\lambda q \in \{p \in (\operatorname{supp}(a, m, g) \times \operatorname{supp}(a, m, f)) : (\operatorname{opr}(m)(\operatorname{nth}(1, p), \operatorname{nth}(2, p)) = x)\}.\; \langle \operatorname{nth}(2, q), \operatorname{nth}(1, q) \rangle)((\lambda q \in \{p \in (\operatorname{supp}(a, m, f) \times \operatorname{supp}(a, m, g)) : (\operatorname{opr}(m)(\operatorname{nth}(1, p), \operatorname{nth}(2, p)) = x)\}.\; \langle \operatorname{nth}(2, q), \operatorname{nth}(1, q) \rangle)(u)) \\ \quad \quad =\; u \end{array}}}
\vnbentry{\texttt{monalg-\allowbreak{}index-\allowbreak{}swap-\allowbreak{}bijection}}{\vnbstmt{\begin{array}{@{}l@{}} \forall a, m, f, g, x. \\ \quad \operatorname{is\text{-}comm\text{-}monoid}(m) \Rightarrow  \\ \quad \quad (\lambda q \in \{p \in (\operatorname{supp}(a, m, f) \times \operatorname{supp}(a, m, g)) : (\operatorname{opr}(m)(\operatorname{nth}(1, p), \operatorname{nth}(2, p)) = x)\}.\; \\ \quad \quad \quad \langle \operatorname{nth}(2, q), \operatorname{nth}(1, q) \rangle) \\ \quad \quad \in\; \mathrm{Bij}(\{p \in (\operatorname{supp}(a, m, f) \times \operatorname{supp}(a, m, g)) : (\operatorname{opr}(m)(\operatorname{nth}(1, p), \operatorname{nth}(2, p)) = x)\}, \{p \in (\operatorname{supp}(a, m, g) \times \operatorname{supp}(a, m, f)) : (\operatorname{opr}(m)(\operatorname{nth}(1, p), \operatorname{nth}(2, p)) = x)\}) \end{array}}}
\vnbentry{\texttt{monalg-\allowbreak{}mul-\allowbreak{}comm-\allowbreak{}at}}{\vnbstmt{\begin{array}{@{}l@{}} \forall a, m, f, g, x. \\ \quad \operatorname{is\text{-}commutative\text{-}ring}(a) \wedge \\ \quad \operatorname{is\text{-}comm\text{-}monoid}(m) \wedge \\ \quad (f \in \operatorname{finsupp}(a, m)) \wedge \\ \quad (g \in \operatorname{finsupp}(a, m)) \wedge \\ \quad (x \in \operatorname{carr}(m)) \Rightarrow \\ \quad \quad (\operatorname{monalg\text{-}mul}(a, m, f, g)(x) = \operatorname{monalg\text{-}mul}(a, m, g, f)(x)) \end{array}}}
\vnbentry{\texttt{monalg-\allowbreak{}mul-\allowbreak{}comm}}{\vnbstmt{\begin{array}{@{}l@{}} \forall a, m, f, g. \\ \quad \operatorname{is\text{-}commutative\text{-}ring}(a) \wedge \\ \quad \operatorname{is\text{-}comm\text{-}monoid}(m) \wedge \\ \quad (f \in \operatorname{finsupp}(a, m)) \wedge \\ \quad (g \in \operatorname{finsupp}(a, m)) \Rightarrow \\ \quad \quad (\operatorname{monalg\text{-}mul}(a, m, f, g) = \operatorname{monalg\text{-}mul}(a, m, g, f)) \end{array}}}
\vnbentry{\texttt{monalg-\allowbreak{}comm}}{\vnbstmt{\begin{array}{@{}l@{}} \forall a, m. \\ \quad \operatorname{is\text{-}commutative\text{-}ring}(a) \wedge \\ \quad \operatorname{is\text{-}comm\text{-}monoid}(m) \Rightarrow \\ \quad \quad \operatorname{is\text{-}commutative\text{-}ring}(\operatorname{monalg}(a, m)) \end{array}}}

\vnbfile{theorem-library/coord-block-estimate-proof.scm}{theorem-library/coord-block-estimate-proof.scm}
\noindent{\small theorem-library/coord-block-estimate-proof.scm ==================================================================== coord-block-estimate -- PROVEN.}\par
\vnbentry{\texttt{coord-\allowbreak{}block-\allowbreak{}estimate}}{\vnbstmt{\begin{array}{@{}l@{}} \forall s, g, b, p, a, c. \\ \quad \operatorname{converges\text{-}along}(s, g, b, p) \wedge \\ \quad \operatorname{strictly\text{-}mono\text{-}nn}(a) \wedge \\ \quad (c \in \mathbb{N}) \wedge \\ \quad \forall j \in \mathbb{N}.\; ((c \leq j) \Rightarrow (a(j) \in b)) \Rightarrow \\ \quad \quad \operatorname{converges\text{-}to}(s, \operatorname{subseq}(g, a), p) \end{array}}}

\vnbfile{theorem-library/module-zero-act.scm}{theorem-library/module-zero-act.scm}
\noindent{\small module-zero-act.scm -- 0 . x = 0\_V : the ring zero annihilates every vector.}\par
\vnbentry{\texttt{module-\allowbreak{}zero-\allowbreak{}act}}{\vnbstmt{\begin{array}{@{}l@{}} \forall s.; \forall x \in \operatorname{vec}(s). \\ \quad \operatorname{is\text{-}module}(s) \Rightarrow  \\ \quad \quad (\operatorname{act}(s)(\operatorname{zero}(\operatorname{scal}(s)), x) = \operatorname{vzero}(s)) \end{array}}}

\vnbfile{theorem-library/nn-least-element.scm}{theorem-library/nn-least-element.scm}
\noindent{\small theorem-library/nn-least-element.scm =================================================================== The well-ordering of NN, PROVEN from the well-ordering of ORD.}\par
\vnbentry{\texttt{nn-\allowbreak{}least-\allowbreak{}element}}{\vnbstmt{\begin{array}{@{}l@{}} \forall t. \\ \quad ((t \subseteq \mathbb{N}) \wedge \exists n.\; (n \in t)) \Rightarrow  \\ \quad \quad \exists m \in t.\; \forall k \in t.\; (m \leq k) \end{array}}}

\vnbfile{theorem-library/poly-degree-laws.scm}{theorem-library/poly-degree-laws.scm}
\noindent{\small poly-degree-laws.scm -- DEG / LEADCOEF / MONOMIAL, PROVEN.}\par
\vnbentry{\texttt{deg-\allowbreak{}well-\allowbreak{}defined}}{\vnbstmt{\begin{array}{@{}l@{}} \forall a.; \forall p \in \operatorname{carr}(\operatorname{poly}(a)). \\ \quad (\operatorname{deg}(a, p) \in \mathbb{N}) \wedge \\ \quad \operatorname{deg\text{-}bound}(a, p, \operatorname{deg}(a, p)) \wedge \\ \quad \forall j \in \mathbb{N}. \\ \quad \quad \operatorname{deg\text{-}bound}(a, p, j) \Rightarrow  \\ \quad \quad \quad (\operatorname{deg}(a, p) \leq j) \end{array}}}
\vnbentry{\texttt{deg-\allowbreak{}bound-\allowbreak{}apply}}{\vnbstmt{\begin{array}{@{}l@{}} \forall a, p, n, k. \\ \quad (\operatorname{deg\text{-}bound}(a, p, n) \wedge ((k \in \mathbb{N}) \wedge (n < k))) \Rightarrow  \\ \quad \quad (p(k) = \operatorname{zero}(a)) \end{array}}}
\vnbentry{\texttt{deg-\allowbreak{}unfold}}{\vnbstmt{\forall a, p.\; (\operatorname{deg}(a, p) \simeq \iota(n))}}
\vnbentry{\texttt{leadcoef-\allowbreak{}unfold}}{\vnbstmt{\begin{array}{@{}l@{}} \forall a, p. \\ \quad (\operatorname{leadcoef}(a, p) \simeq p(\operatorname{deg}(a, p))) \end{array}}}
\vnbentry{\texttt{nn-\allowbreak{}succ-\allowbreak{}not-\allowbreak{}le}}{\vnbstmt{\forall n \in \mathbb{N}.\; \neg ((n + 1) \leq n)}}
\vnbentry{\texttt{leadcoef-\allowbreak{}nonzero-\allowbreak{}ptwise}}{\vnbstmt{\begin{array}{@{}l@{}} \forall a.; \forall p \in \operatorname{carr}(\operatorname{poly}(a)). \\ \quad \exists k \in \mathbb{N}.\; \neg (p(k) = \operatorname{zero}(a)) \Rightarrow  \\ \quad \quad \neg (\operatorname{leadcoef}(a, p) = \operatorname{zero}(a)) \end{array}}}
\vnbentry{\texttt{leadcoef-\allowbreak{}nonzero}}{\vnbstmt{\begin{array}{@{}l@{}} \forall a, p. \\ \quad \operatorname{is\text{-}ring}(a) \wedge \\ \quad (p \in \operatorname{carr}(\operatorname{poly}(a))) \wedge \\ \quad \neg (p = \operatorname{zero}(\operatorname{poly}(a))) \Rightarrow \\ \quad \quad \neg (\operatorname{leadcoef}(a, p) = \operatorname{zero}(a)) \end{array}}}
\vnbentry{\texttt{monomial-\allowbreak{}at}}{\vnbstmt{\begin{array}{@{}l@{}} \forall a, c, n. \\ \quad ((c \in \operatorname{carr}(a)) \wedge (n \in \mathbb{N})) \Rightarrow  \\ \quad \quad (\operatorname{monomial}(a, c, n)(n) = c) \end{array}}}
\vnbentry{\texttt{monomial-\allowbreak{}off}}{\vnbstmt{\begin{array}{@{}l@{}} \forall a, c, n, k. \\ \quad ((c \in \operatorname{carr}(a)) \wedge ((k \in \mathbb{N}) \wedge \neg (k = n))) \Rightarrow  \\ \quad \quad (\operatorname{monomial}(a, c, n)(k) = \operatorname{zero}(a)) \end{array}}}
\vnbentry{\texttt{monomial-\allowbreak{}in-\allowbreak{}poly}}{\vnbstmt{\begin{array}{@{}l@{}} \forall a, c, n. \\ \quad (\operatorname{is\text{-}ring}(a) \wedge ((c \in \operatorname{carr}(a)) \wedge (n \in \mathbb{N}))) \Rightarrow  \\ \quad \quad (\operatorname{monomial}(a, c, n) \in \operatorname{carr}(\operatorname{poly}(a))) \end{array}}}
\vnbentry{\texttt{monomial-\allowbreak{}support}}{\vnbstmt{\begin{array}{@{}l@{}} \forall a, c, n, k. \\ \quad \operatorname{is\text{-}ring}(a) \wedge \\ \quad (c \in \operatorname{carr}(a)) \wedge \\ \quad \neg (c = \operatorname{zero}(a)) \wedge \\ \quad (n \in \mathbb{N}) \Rightarrow \\ \quad \quad ((k \in \operatorname{supp}(a, \operatorname{nn\text{-}add\text{-}monoid}, \operatorname{monomial}(a, c, n))) \Leftrightarrow ((k \in \mathbb{N}) \wedge (k = n))) \end{array}}}
\vnbentry{\texttt{deg-\allowbreak{}is}}{\vnbstmt{\begin{array}{@{}l@{}} \forall a, p, m. \\ \quad (p \in \operatorname{carr}(\operatorname{poly}(a))) \wedge \\ \quad (m \in \mathbb{N}) \wedge \\ \quad \operatorname{deg\text{-}bound}(a, p, m) \wedge \\ \quad \forall j \in \mathbb{N}. \\ \quad \quad (\operatorname{deg\text{-}bound}(a, p, j) \Rightarrow (m \leq j)) \Rightarrow \\ \quad \quad (\operatorname{deg}(a, p) = m) \end{array}}}
\vnbentry{\texttt{monomial-\allowbreak{}deg}}{\vnbstmt{\begin{array}{@{}l@{}} \forall a, c, n. \\ \quad \operatorname{is\text{-}ring}(a) \wedge \\ \quad (c \in \operatorname{carr}(a)) \wedge \\ \quad \neg (c = \operatorname{zero}(a)) \wedge \\ \quad (n \in \mathbb{N}) \Rightarrow \\ \quad \quad (\operatorname{deg}(a, \operatorname{monomial}(a, c, n)) = n) \end{array}}}
\vnbentry{\texttt{monomial-\allowbreak{}leadcoef}}{\vnbstmt{\begin{array}{@{}l@{}} \forall a, c, n. \\ \quad \operatorname{is\text{-}ring}(a) \wedge \\ \quad (c \in \operatorname{carr}(a)) \wedge \\ \quad \neg (c = \operatorname{zero}(a)) \wedge \\ \quad (n \in \mathbb{N}) \Rightarrow \\ \quad \quad (\operatorname{leadcoef}(a, \operatorname{monomial}(a, c, n)) = c) \end{array}}}
\vnbentry{\texttt{deg-\allowbreak{}add-\allowbreak{}le}}{\vnbstmt{\begin{array}{@{}l@{}} \forall a, f, g, n. \\ \quad \operatorname{is\text{-}ring}(a) \wedge \\ \quad (f \in \operatorname{carr}(\operatorname{poly}(a))) \wedge \\ \quad (g \in \operatorname{carr}(\operatorname{poly}(a))) \wedge \\ \quad (n \in \mathbb{N}) \wedge \\ \quad (\operatorname{deg}(a, f) \leq n) \wedge \\ \quad (\operatorname{deg}(a, g) \leq n) \Rightarrow \\ \quad \quad \operatorname{deg}(a, \operatorname{monalg\text{-}add}(a, \operatorname{nn\text{-}add\text{-}monoid}, f, g)) \\ \quad \quad \leq\; n \end{array}}}

\vnbfile{theorem-library/sqrt2-proof.scm}{theorem-library/sqrt2-proof.scm}
\noindent{\small sqrt2-proof.scm -- sqrt(2) is irrational, stated on the naturals:}\par
\vnbentry{\texttt{sqrt2-\allowbreak{}irrational}}{\vnbstmt{\begin{array}{@{}l@{}} \forall p, q \in \mathbb{N}. \\ \quad (\neg (q = 0) \Rightarrow \neg ((p \cdot p) = (2 \cdot (q \cdot q)))) \end{array}}}

\vnbfile{theorem-library/sqrt3-proof.scm}{theorem-library/sqrt3-proof.scm}
\noindent{\small sqrt3-proof.scm -- sqrt(3) is irrational, on the naturals: forall p,q in NN.  q /= 0  =$>$  p*p /= 3*(q*q).}\par
\vnbentry{\texttt{sqrt3-\allowbreak{}irrational}}{\vnbstmt{\begin{array}{@{}l@{}} \forall p, q \in \mathbb{N}. \\ \quad (\neg (q = 0) \Rightarrow \neg ((p \cdot p) = (3 \cdot (q \cdot q)))) \end{array}}}

\vnbfile{theorem-library/gauge-is-degree.scm}{theorem-library/gauge-is-degree.scm}
\noindent{\small gauge-is-degree.scm -- the named Euclidean gauge GAUGE(s) is a degree function, PROVEN, together with the four supports its warrant cites.}\par
\vnbentry{\texttt{euclidean-\allowbreak{}gauges-\allowbreak{}unfold}}{\vnbstmt{\begin{array}{@{}l@{}} \forall s. \\ \quad (\operatorname{euclidean\text{-}gauges}(s) \simeq \{g \in (\operatorname{carr}(s) \to \mathbb{N}) : \operatorname{has\text{-}div\text{-}remainder}(s, g)\}) \end{array}}}
\vnbentry{\texttt{euclidean-\allowbreak{}ring-\allowbreak{}has-\allowbreak{}gauge}}{\vnbstmt{\begin{array}{@{}l@{}} \forall s. \\ \quad \operatorname{is\text{-}euclidean\text{-}ring}(s) \Rightarrow  \\ \quad \quad \exists g \in (\operatorname{carr}(s) \to \mathbb{N}). \\ \quad \quad \quad \operatorname{has\text{-}div\text{-}remainder}(s, g) \end{array}}}
\vnbentry{\texttt{gauges-\allowbreak{}mem-\allowbreak{}build}}{\vnbstmt{\begin{array}{@{}l@{}} \forall s.; \forall g \in (\operatorname{carr}(s) \to \mathbb{N}). \\ \quad \operatorname{has\text{-}div\text{-}remainder}(s, g) \Rightarrow  \\ \quad \quad (g \in \operatorname{euclidean\text{-}gauges}(s)) \end{array}}}
\vnbentry{\texttt{gauges-\allowbreak{}in-\allowbreak{}fun}}{\vnbstmt{\begin{array}{@{}l@{}} \forall s.; \forall g \in \operatorname{euclidean\text{-}gauges}(s). \\ \quad (g \in (\operatorname{carr}(s) \to \mathbb{N})) \end{array}}}
\vnbentry{\texttt{gauges-\allowbreak{}spec}}{\vnbstmt{\begin{array}{@{}l@{}} \forall s.; \forall g \in \operatorname{euclidean\text{-}gauges}(s). \\ \quad \operatorname{has\text{-}div\text{-}remainder}(s, g) \end{array}}}
\vnbentry{\texttt{gauge-\allowbreak{}is-\allowbreak{}degree}}{\vnbstmt{\begin{array}{@{}l@{}} \forall s. \\ \quad \operatorname{is\text{-}euclidean\text{-}ring}(s) \Rightarrow  \\ \quad \quad (\operatorname{gauge}(s) \in (\operatorname{carr}(s) \to \mathbb{N})) \wedge \\ \quad \quad \operatorname{has\text{-}div\text{-}remainder}(s, \operatorname{gauge}(s)) \end{array}}}

\vnbfile{theorem-library/euclidean-ideal-generator-proof.scm}{theorem-library/euclidean-ideal-generator-proof.scm}
\noindent{\small euclidean-ideal-generator-proof.scm -- every ideal of a Euclidean ring has a generator: some b in I with I subset (b).}\par
\vnbentry{\texttt{euclidean-\allowbreak{}ideal-\allowbreak{}has-\allowbreak{}generator}}{\vnbstmt{\begin{array}{@{}l@{}} \forall s, i. \\ \quad (\operatorname{is\text{-}euclidean\text{-}ring}(s) \wedge \operatorname{is\text{-}ideal}(s, i)) \Rightarrow  \\ \quad \quad \exists b \in i. \forall a \in i. \\ \quad \quad \quad (a \in \operatorname{principal\text{-}ideal}(s, b)) \end{array}}}

\vnbfile{theorem-library/zz-bezout-proof.scm}{theorem-library/zz-bezout-proof.scm}
\noindent{\small zz-bezout-proof.scm -- BEZOUT on the integers, from the Euclidean-ring core.}\par
\vnbentry{\texttt{zz-\allowbreak{}bezout-\allowbreak{}set-\allowbreak{}is-\allowbreak{}ideal}}{\vnbstmt{\begin{array}{@{}l@{}} \forall a, b. \\ \quad ((a \in \mathbb{Z}) \wedge (b \in \mathbb{Z})) \Rightarrow  \\ \quad \quad \operatorname{is\text{-}ideal}(\operatorname{zz\text{-}ring}, \operatorname{zz\text{-}bezout\text{-}set}(a, b)) \end{array}}}
\vnbentry{\texttt{zz-\allowbreak{}bezout}}{\vnbstmt{\begin{array}{@{}l@{}} \forall a, b. \\ \quad ((a \in \mathbb{Z}) \wedge (b \in \mathbb{Z})) \Rightarrow  \\ \quad \quad \exists d \in \operatorname{zz\text{-}bezout\text{-}set}(a, b). \\ \quad \quad \quad (\operatorname{zz\text{-}divides}(d, a) \wedge \operatorname{zz\text{-}divides}(d, b)) \end{array}}}

\vnbfile{theorem-library/continuity-basics.scm}{theorem-library/continuity-basics.scm}
\noindent{\small continuity-basics.scm -- the constant and identity maps are continuous, PROVEN, with their FUN typings.}\par
\vnbentry{\texttt{const-\allowbreak{}lam-\allowbreak{}in-\allowbreak{}fun}}{\vnbstmt{\begin{array}{@{}l@{}} \forall c \in \mathbb{R}. \\ \quad ((\lambda x \in \mathbb{R}.\; c) \in (\mathbb{R} \to \mathbb{R})) \end{array}}}
\vnbentry{\texttt{ident-\allowbreak{}lam-\allowbreak{}in-\allowbreak{}fun}}{\vnbstmt{((\lambda x \in \mathbb{R}.\; x) \in (\mathbb{R} \to \mathbb{R}))}}
\vnbentry{\texttt{const-\allowbreak{}continuous-\allowbreak{}at}}{\vnbstmt{\begin{array}{@{}l@{}} \forall c, a \in \mathbb{R}. \\ \quad \operatorname{is\text{-}continuous\text{-}at}(\operatorname{rr\text{-}ms}, \operatorname{rr\text{-}ms}, (\lambda x \in \mathbb{R}.\; c), a) \end{array}}}
\vnbentry{\texttt{identity-\allowbreak{}continuous-\allowbreak{}at}}{\vnbstmt{\begin{array}{@{}l@{}} \forall a \in \mathbb{R}. \\ \quad \operatorname{is\text{-}continuous\text{-}at}(\operatorname{rr\text{-}ms}, \operatorname{rr\text{-}ms}, (\lambda x \in \mathbb{R}.\; x), a) \end{array}}}

\vnbfile{theorem-library/neg-continuous.scm}{theorem-library/neg-continuous.scm}
\noindent{\small neg-continuous.scm -- NEGATION preserves membership in FUN(RR,RR) and continuity, PROVEN.}\par
\vnbentry{\texttt{neg-\allowbreak{}fun-\allowbreak{}in-\allowbreak{}fun}}{\vnbstmt{\begin{array}{@{}l@{}} \forall f \in (\mathbb{R} \to \mathbb{R}). \\ \quad ((\lambda z \in \mathbb{R}.\; -f(z)) \in (\mathbb{R} \to \mathbb{R})) \end{array}}}
\vnbentry{\texttt{neg-\allowbreak{}continuous-\allowbreak{}at}}{\vnbstmt{\begin{array}{@{}l@{}} \forall f, x. \\ \quad (f \in (\mathbb{R} \to \mathbb{R})) \wedge \\ \quad (x \in \mathbb{R}) \wedge \\ \quad \operatorname{is\text{-}continuous\text{-}at}(\operatorname{rr\text{-}ms}, \operatorname{rr\text{-}ms}, f, x) \Rightarrow \\ \quad \quad \operatorname{is\text{-}continuous\text{-}at}(\operatorname{rr\text{-}ms}, \operatorname{rr\text{-}ms}, (\lambda z \in \mathbb{R}.\; -f(z)), x) \end{array}}}

\vnbfile{theorem-library/continuity-sum.scm}{theorem-library/continuity-sum.scm}
\noindent{\small continuity-sum.scm -- the pointwise SUM of two maps continuous at a point is continuous there, PROVEN, with its FUN typing.}\par
\vnbentry{\texttt{sum-\allowbreak{}lam-\allowbreak{}in-\allowbreak{}fun}}{\vnbstmt{\begin{array}{@{}l@{}} \forall g, h. \\ \quad ((g \in (\mathbb{R} \to \mathbb{R})) \wedge (h \in (\mathbb{R} \to \mathbb{R}))) \Rightarrow  \\ \quad \quad ((\lambda x \in \mathbb{R}.\; (g(x) + h(x))) \in (\mathbb{R} \to \mathbb{R})) \end{array}}}
\vnbentry{\texttt{sum-\allowbreak{}continuous-\allowbreak{}at}}{\vnbstmt{\begin{array}{@{}l@{}} \forall g, h, a. \\ \quad \operatorname{is\text{-}continuous\text{-}at}(\operatorname{rr\text{-}ms}, \operatorname{rr\text{-}ms}, g, a) \wedge \\ \quad \operatorname{is\text{-}continuous\text{-}at}(\operatorname{rr\text{-}ms}, \operatorname{rr\text{-}ms}, h, a) \Rightarrow \\ \quad \quad \operatorname{is\text{-}continuous\text{-}at}(\operatorname{rr\text{-}ms}, \operatorname{rr\text{-}ms}, (\lambda x \in \mathbb{R}.\; (g(x) + h(x))), a) \end{array}}}

\vnbfile{theorem-library/continuity-product.scm}{theorem-library/continuity-product.scm}
\noindent{\small continuity-product.scm -- the pointwise PRODUCT of two maps continuous at a point is continuous there, PROVEN, with its FUN typing.}\par
\vnbentry{\texttt{prod-\allowbreak{}lam-\allowbreak{}in-\allowbreak{}fun}}{\vnbstmt{\begin{array}{@{}l@{}} \forall g, h. \\ \quad ((g \in (\mathbb{R} \to \mathbb{R})) \wedge (h \in (\mathbb{R} \to \mathbb{R}))) \Rightarrow  \\ \quad \quad ((\lambda x \in \mathbb{R}.\; (g(x) \cdot h(x))) \in (\mathbb{R} \to \mathbb{R})) \end{array}}}
\vnbentry{\texttt{product-\allowbreak{}continuous-\allowbreak{}at}}{\vnbstmt{\begin{array}{@{}l@{}} \forall g, h, a. \\ \quad \operatorname{is\text{-}continuous\text{-}at}(\operatorname{rr\text{-}ms}, \operatorname{rr\text{-}ms}, g, a) \wedge \\ \quad \operatorname{is\text{-}continuous\text{-}at}(\operatorname{rr\text{-}ms}, \operatorname{rr\text{-}ms}, h, a) \Rightarrow \\ \quad \quad \operatorname{is\text{-}continuous\text{-}at}(\operatorname{rr\text{-}ms}, \operatorname{rr\text{-}ms}, (\lambda x \in \mathbb{R}.\; (g(x) \cdot h(x))), a) \end{array}}}

\vnbfile{theorem-library/continuity-transfer.scm}{theorem-library/continuity-transfer.scm}
\noindent{\small continuity-transfer.scm -- continuity reads only the VALUES: a map that agrees pointwise with a map continuous at a is itself continuous at a. PROVEN.}\par
\vnbentry{\texttt{cont-\allowbreak{}transfer-\allowbreak{}ptwise-\allowbreak{}eq}}{\vnbstmt{\begin{array}{@{}l@{}} \forall f, g, a. \\ \quad (f \in (\mathbb{R} \to \mathbb{R})) \wedge \\ \quad \operatorname{is\text{-}continuous\text{-}at}(\operatorname{rr\text{-}ms}, \operatorname{rr\text{-}ms}, g, a) \wedge \\ \quad \forall x \in \mathbb{R}.\; (f(x) = g(x)) \Rightarrow \\ \quad \quad \operatorname{is\text{-}continuous\text{-}at}(\operatorname{rr\text{-}ms}, \operatorname{rr\text{-}ms}, f, a) \end{array}}}

\vnbfile{theorem-library/continuity-sub.scm}{theorem-library/continuity-sub.scm}
\noindent{\small continuity-sub.scm -- the pointwise DIFFERENCE of two maps continuous at a point is continuous there, PROVEN, with its FUN typing.}\par
\vnbentry{\texttt{sub-\allowbreak{}lam-\allowbreak{}in-\allowbreak{}fun}}{\vnbstmt{\begin{array}{@{}l@{}} \forall g, h. \\ \quad ((g \in (\mathbb{R} \to \mathbb{R})) \wedge (h \in (\mathbb{R} \to \mathbb{R}))) \Rightarrow  \\ \quad \quad ((\lambda x \in \mathbb{R}.\; (g(x) - h(x))) \in (\mathbb{R} \to \mathbb{R})) \end{array}}}
\vnbentry{\texttt{sub-\allowbreak{}continuous-\allowbreak{}at}}{\vnbstmt{\begin{array}{@{}l@{}} \forall g, h, a. \\ \quad \operatorname{is\text{-}continuous\text{-}at}(\operatorname{rr\text{-}ms}, \operatorname{rr\text{-}ms}, g, a) \wedge \\ \quad \operatorname{is\text{-}continuous\text{-}at}(\operatorname{rr\text{-}ms}, \operatorname{rr\text{-}ms}, h, a) \Rightarrow \\ \quad \quad \operatorname{is\text{-}continuous\text{-}at}(\operatorname{rr\text{-}ms}, \operatorname{rr\text{-}ms}, (\lambda x \in \mathbb{R}.\; (g(x) - h(x))), a) \end{array}}}

\vnbfile{theorem-library/continuity-scale.scm}{theorem-library/continuity-scale.scm}
\noindent{\small continuity-scale.scm -- the SCALAR MULTIPLE of a real map: its FUN typing and its continuity, both PROVEN `modulo 0'.}\par
\vnbentry{\texttt{scale-\allowbreak{}lam-\allowbreak{}in-\allowbreak{}fun}}{\vnbstmt{\begin{array}{@{}l@{}} \forall c, f. \\ \quad ((c \in \mathbb{R}) \wedge (f \in (\mathbb{R} \to \mathbb{R}))) \Rightarrow  \\ \quad \quad ((\lambda x \in \mathbb{R}.\; (c \cdot f(x))) \in (\mathbb{R} \to \mathbb{R})) \end{array}}}
\vnbentry{\texttt{scale-\allowbreak{}continuous-\allowbreak{}at}}{\vnbstmt{\begin{array}{@{}l@{}} \forall c \in \mathbb{R}, f, a. \\ \quad \operatorname{is\text{-}continuous\text{-}at}(\operatorname{rr\text{-}ms}, \operatorname{rr\text{-}ms}, f, a) \Rightarrow \\ \quad \quad \operatorname{is\text{-}continuous\text{-}at}(\operatorname{rr\text{-}ms}, \operatorname{rr\text{-}ms}, (\lambda x \in \mathbb{R}.\; (c \cdot f(x))), a) \end{array}}}

\vnbfile{theorem-library/continuity-compose.scm}{theorem-library/continuity-compose.scm}
\noindent{\small continuity-compose.scm -- COMPOSITION OF CONTINUOUS MAPS, PROVEN.}\par
\vnbentry{\texttt{compose-\allowbreak{}continuous-\allowbreak{}at}}{\vnbstmt{\begin{array}{@{}l@{}} \forall g, f, a. \\ \quad \operatorname{is\text{-}continuous\text{-}at}(\operatorname{rr\text{-}ms}, \operatorname{rr\text{-}ms}, f, a) \wedge \\ \quad \operatorname{is\text{-}continuous\text{-}at}(\operatorname{rr\text{-}ms}, \operatorname{rr\text{-}ms}, g, f(a)) \Rightarrow \\ \quad \quad \operatorname{is\text{-}continuous\text{-}at}(\operatorname{rr\text{-}ms}, \operatorname{rr\text{-}ms}, \operatorname{compose}(g, f), a) \end{array}}}

\vnbfile{theorem-library/continuity-recip.scm}{theorem-library/continuity-recip.scm}
\noindent{\small continuity-recip.scm -- THE RECIPROCAL OF A NON-VANISHING CONTINUOUS MAP, PROVEN `modulo 0'.}\par
\vnbentry{\texttt{recip-\allowbreak{}lam-\allowbreak{}in-\allowbreak{}fun}}{\vnbstmt{\begin{array}{@{}l@{}} \forall h \in (\mathbb{R} \to \mathbb{R}). \\ \quad \forall x \in \mathbb{R}.\; \neg (h(x) = 0) \Rightarrow  \\ \quad \quad ((\lambda x \in \mathbb{R}.\; \operatorname{recip}(h(x))) \in (\mathbb{R} \to \mathbb{R})) \end{array}}}
\vnbentry{\texttt{recip-\allowbreak{}continuous-\allowbreak{}at}}{\vnbstmt{\begin{array}{@{}l@{}} \forall h, p. \\ \quad (h \in (\mathbb{R} \to \mathbb{R})) \wedge \\ \quad \forall x \in \mathbb{R}.\; \neg (h(x) = 0) \wedge \\ \quad \operatorname{is\text{-}continuous\text{-}at}(\operatorname{rr\text{-}ms}, \operatorname{rr\text{-}ms}, h, p) \Rightarrow \\ \quad \quad \operatorname{is\text{-}continuous\text{-}at}(\operatorname{rr\text{-}ms}, \operatorname{rr\text{-}ms}, (\lambda x \in \mathbb{R}.\; \operatorname{recip}(h(x))), p) \end{array}}}

\vnbfile{theorem-library/cont-agree-off-pt.scm}{theorem-library/cont-agree-off-pt.scm}
\noindent{\small cont-agree-off-pt.scm -- two functions continuous at a point and agreeing AWAY from it agree AT it, PROVEN.}\par
\vnbentry{\texttt{cont-\allowbreak{}agree-\allowbreak{}off-\allowbreak{}pt}}{\vnbstmt{\begin{array}{@{}l@{}} \forall a, b.; \forall t \in \mathbb{R}. \\ \quad \operatorname{is\text{-}continuous\text{-}at}(\operatorname{rr\text{-}ms}, \operatorname{rr\text{-}ms}, a, t) \wedge \\ \quad \operatorname{is\text{-}continuous\text{-}at}(\operatorname{rr\text{-}ms}, \operatorname{rr\text{-}ms}, b, t) \wedge \\ \quad \forall x \in \mathbb{R}.\; (\neg (x = t) \Rightarrow (a(x) = b(x))) \Rightarrow \\ \quad \quad (a(t) = b(t)) \end{array}}}

\vnbfile{theorem-library/continuous-one-sided-sign.scm}{theorem-library/continuous-one-sided-sign.scm}
\noindent{\small continuous-one-sided-sign.scm -- sign preservation under continuity, from ONE side, PROVEN.  Two theorems, one driver parameterised by the side:}\par
\vnbentry{\texttt{continuous-\allowbreak{}nonneg-\allowbreak{}left}}{\vnbstmt{\begin{array}{@{}l@{}} \forall g, a, h. \\ \quad (g \in (\mathbb{R} \to \mathbb{R})) \wedge \\ \quad (a \in \mathbb{R}) \wedge \\ \quad (h \in \mathbb{R}) \wedge \\ \quad (a < h) \wedge \\ \quad \operatorname{is\text{-}continuous\text{-}at}(\operatorname{rr\text{-}ms}, \operatorname{rr\text{-}ms}, g, h) \wedge \\ \quad \forall x \in \mathbb{R}.\; (((a < x) \wedge (x < h)) \Rightarrow (0 \leq g(x))) \Rightarrow \\ \quad \quad (0 \leq g(h)) \end{array}}}
\vnbentry{\texttt{continuous-\allowbreak{}nonpos-\allowbreak{}right}}{\vnbstmt{\begin{array}{@{}l@{}} \forall g, h, b. \\ \quad (g \in (\mathbb{R} \to \mathbb{R})) \wedge \\ \quad (h \in \mathbb{R}) \wedge \\ \quad (b \in \mathbb{R}) \wedge \\ \quad (h < b) \wedge \\ \quad \operatorname{is\text{-}continuous\text{-}at}(\operatorname{rr\text{-}ms}, \operatorname{rr\text{-}ms}, g, h) \wedge \\ \quad \forall x \in \mathbb{R}.\; (((h < x) \wedge (x < b)) \Rightarrow (g(x) \leq 0)) \Rightarrow \\ \quad \quad (g(h) \leq 0) \end{array}}}

\vnbfile{theorem-library/differentiation.scm}{theorem-library/differentiation.scm}
\noindent{\small differentiation.scm -- Chapter 2 (Differentiation) of calculus.pdf, founded on the CARATHEODORY / o(h) ALGEBRAIC definition (user's choice 2026-06-27).}\par
\vnbentry{\texttt{diff-\allowbreak{}implies-\allowbreak{}continuous}}{\vnbstmt{\begin{array}{@{}l@{}} \forall f, a, l. \\ \quad \operatorname{is\text{-}diff\text{-}at}(f, a, l) \Rightarrow  \\ \quad \quad \operatorname{is\text{-}continuous\text{-}at}(\operatorname{rr\text{-}ms}, \operatorname{rr\text{-}ms}, f, a) \end{array}}}
\vnbentry{\texttt{derivative-\allowbreak{}unique}}{\vnbstmt{\begin{array}{@{}l@{}} \forall f, a, l, m. \\ \quad \operatorname{is\text{-}diff\text{-}at}(f, a, l) \wedge \\ \quad \operatorname{is\text{-}diff\text{-}at}(f, a, m) \Rightarrow \\ \quad \quad (l = m) \end{array}}}
\vnbentry{\texttt{deriv-\allowbreak{}const}}{\vnbstmt{\begin{array}{@{}l@{}} \forall c, a. \\ \quad ((c \in \mathbb{R}) \wedge (a \in \mathbb{R})) \Rightarrow  \\ \quad \quad \operatorname{is\text{-}diff\text{-}at}((\lambda x \in \mathbb{R}.\; c), a, 0) \end{array}}}
\vnbentry{\texttt{deriv-\allowbreak{}identity}}{\vnbstmt{\begin{array}{@{}l@{}} \forall a \in \mathbb{R}. \\ \quad \operatorname{is\text{-}diff\text{-}at}((\lambda x \in \mathbb{R}.\; x), a, 1) \end{array}}}

\vnbfile{theorem-library/diff-transfer.scm}{theorem-library/diff-transfer.scm}
\noindent{\small diff-transfer.scm -- differentiability reads only the VALUES: a map that agrees pointwise with a map differentiable at a is itself differentiable at a, with the same derivative.  PROVEN, `modulo 0'.}\par
\vnbentry{\texttt{diff-\allowbreak{}transfer-\allowbreak{}ptwise-\allowbreak{}eq}}{\vnbstmt{\begin{array}{@{}l@{}} \forall f, g, a, l. \\ \quad (f \in (\mathbb{R} \to \mathbb{R})) \wedge \\ \quad \operatorname{is\text{-}diff\text{-}at}(g, a, l) \wedge \\ \quad \forall x \in \mathbb{R}.\; (f(x) \simeq g(x)) \Rightarrow \\ \quad \quad \operatorname{is\text{-}diff\text{-}at}(f, a, l) \end{array}}}

\vnbfile{theorem-library/deriv-sum-product.scm}{theorem-library/deriv-sum-product.scm}
\noindent{\small deriv-sum-product.scm -- calculus.pdf Prop 2.5 (the SUM rule) and Prop 2.6 (the PRODUCT rule), PROVEN, both `modulo 0'.}\par
\vnbentry{\texttt{deriv-\allowbreak{}sum}}{\vnbstmt{\begin{array}{@{}l@{}} \forall f, g, a, l, m. \\ \quad \operatorname{is\text{-}diff\text{-}at}(f, a, l) \wedge \\ \quad \operatorname{is\text{-}diff\text{-}at}(g, a, m) \Rightarrow \\ \quad \quad \operatorname{is\text{-}diff\text{-}at}((\lambda x \in \mathbb{R}.\; (f(x) + g(x))), a, (l + m)) \end{array}}}
\vnbentry{\texttt{deriv-\allowbreak{}product}}{\vnbstmt{\begin{array}{@{}l@{}} \forall f, g, a, l, m. \\ \quad \operatorname{is\text{-}diff\text{-}at}(f, a, l) \wedge \\ \quad \operatorname{is\text{-}diff\text{-}at}(g, a, m) \Rightarrow \\ \quad \quad \operatorname{is\text{-}diff\text{-}at}((\lambda x \in \mathbb{R}.\; (f(x) \cdot g(x))), a, ((l \cdot g(a)) + (f(a) \cdot m))) \end{array}}}

\vnbfile{theorem-library/chain-rule.scm}{theorem-library/chain-rule.scm}
\noindent{\small chain-rule.scm -- THE CHAIN RULE (calculus.pdf Prop 2.8), PROVEN.}\par
\vnbentry{\texttt{deriv-\allowbreak{}chain}}{\vnbstmt{\begin{array}{@{}l@{}} \forall f, g, a, l, m. \\ \quad \operatorname{is\text{-}diff\text{-}at}(f, a, l) \wedge \\ \quad \operatorname{is\text{-}diff\text{-}at}(g, f(a), m) \Rightarrow \\ \quad \quad \operatorname{is\text{-}diff\text{-}at}(\operatorname{compose}(g, f), a, (m \cdot l)) \end{array}}}
\vnbentry{\texttt{deriv-\allowbreak{}of-\allowbreak{}is-\allowbreak{}diff-\allowbreak{}at}}{\vnbstmt{\begin{array}{@{}l@{}} \forall f, a, v. \\ \quad \operatorname{is\text{-}diff\text{-}at}(f, a, v) \Rightarrow  \\ \quad \quad (\operatorname{deriv}(f, a) = v) \end{array}}}
\vnbentry{\texttt{deriv-\allowbreak{}chain-\allowbreak{}value}}{\vnbstmt{\begin{array}{@{}l@{}} \forall f, g, a, l, m. \\ \quad \operatorname{is\text{-}diff\text{-}at}(f, a, l) \wedge \\ \quad \operatorname{is\text{-}diff\text{-}at}(g, f(a), m) \Rightarrow \\ \quad \quad \operatorname{deriv}(\operatorname{compose}(g, f), a) \\ \quad \quad =\; (\operatorname{deriv}(g, f(a)) \cdot \operatorname{deriv}(f, a)) \end{array}}}

\vnbfile{theorem-library/inverse-function.scm}{theorem-library/inverse-function.scm}
\noindent{\small inverse-function.scm -- THE INVERSE FUNCTION THEOREM ON THE LINE, Caratheodory form, PROVEN:}\par
\vnbentry{\texttt{rr-\allowbreak{}recip-\allowbreak{}solve}}{\vnbstmt{\begin{array}{@{}l@{}} \forall c, u, v. \\ \quad (c \in \mathbb{R}) \wedge \\ \quad (u \in \mathbb{R}) \wedge \\ \quad (v \in \mathbb{R}) \wedge \\ \quad \neg (c = 0) \wedge \\ \quad (u = (c \cdot v)) \Rightarrow \\ \quad \quad (v = (\operatorname{recip}(c) \cdot u)) \end{array}}}
\vnbentry{\texttt{deriv-\allowbreak{}inverse}}{\vnbstmt{\begin{array}{@{}l@{}} \forall f, g, a, l. \\ \quad \operatorname{is\text{-}diff\text{-}at}(f, a, l) \wedge \\ \quad \neg (l = 0) \wedge \\ \quad (g \in (\mathbb{R} \to \mathbb{R})) \wedge \\ \quad \forall x \in \mathbb{R}.\; (g(f(x)) = x) \wedge \\ \quad \forall y \in \mathbb{R}.\; (f(g(y)) = y) \wedge \\ \quad \operatorname{is\text{-}continuous\text{-}at}(\operatorname{rr\text{-}ms}, \operatorname{rr\text{-}ms}, g, f(a)) \Rightarrow \\ \quad \quad \operatorname{is\text{-}diff\text{-}at}(g, f(a), \operatorname{recip}(l)) \end{array}}}

\vnbfile{theorem-library/higher-derivatives.scm}{theorem-library/higher-derivatives.scm}
\noindent{\small higher-derivatives.scm -- calculus.pdf Def 2.2: the n-th derivative f\textasciicircum{}(n), by induction on n, built on DERIV / IS-DIFF-AT (differentiation.scm).}\par
\vnbentry{\texttt{nth-\allowbreak{}deriv-\allowbreak{}one}}{\vnbstmt{\begin{array}{@{}l@{}} \forall f. \\ \quad (\operatorname{nth\text{-}deriv}(f, 1) \simeq (\lambda x \in \mathbb{R}.\; \operatorname{deriv}(f, x))) \end{array}}}

\vnbfile{theorem-library/triple-entry-proof.scm}{theorem-library/triple-entry-proof.scm}
\noindent{\small triple-entry-proof.scm -- the two matrix-product entry-expansion lemmas that matmul-assoc rests on, PROVEN (were warranted in matrix.scm) from the (B) finite-sum bricks (finsum-additive.scm): finsum-congruence, and general-ring distribution of a scalar into a finite sum (finsum-ring-distrib-left/right-gen).}\par
\vnbentry{\texttt{triple-\allowbreak{}entry-\allowbreak{}left}}{\vnbstmt{\begin{array}{@{}l@{}} \forall a, m, n, k, l, p, q, r, w, b. \\ \quad \operatorname{is\text{-}ring}(a) \wedge \\ \quad (p \in \operatorname{mat}(m, n, \operatorname{carr}(a))) \wedge \\ \quad (q \in \operatorname{mat}(n, k, \operatorname{carr}(a))) \wedge \\ \quad (r \in \operatorname{mat}(k, l, \operatorname{carr}(a))) \wedge \\ \quad (1 \leq n) \wedge \\ \quad (1 \leq k) \wedge \\ \quad (w \in \operatorname{interval}(1, m)) \wedge \\ \quad (b \in \operatorname{interval}(1, l)) \Rightarrow \\ \quad \quad \operatorname{entry}(\operatorname{matmul}(a, \operatorname{matmul}(a, p, q), r), w, b) \\ \quad \quad =\; \operatorname{finsum}(\operatorname{ring\text{-}additive\text{-}ag}(a), \\ \quad \quad \quad (\lambda c \in \operatorname{interval}(1, k).\; \\ \quad \quad \quad \quad \operatorname{finsum}(\operatorname{ring\text{-}additive\text{-}ag}(a), \\ \quad \quad \quad \quad \quad (\lambda j \in \operatorname{interval}(1, n).\; \\ \quad \quad \quad \quad \quad \quad (\lambda z \in (\operatorname{interval}(1, k) \times \operatorname{interval}(1, n)).\; \operatorname{mul}(a)(\operatorname{mul}(a)(\operatorname{entry}(p, w, \operatorname{nth}(2, z)), \operatorname{entry}(q, \operatorname{nth}(2, z), \operatorname{nth}(1, z))), \operatorname{entry}(r, \operatorname{nth}(1, z), b)))(\langle c, j \rangle)), \\ \quad \quad \quad \quad \quad \operatorname{interval}(1, n))), \\ \quad \quad \quad \operatorname{interval}(1, k)) \end{array}}}
\vnbentry{\texttt{triple-\allowbreak{}entry-\allowbreak{}right}}{\vnbstmt{\begin{array}{@{}l@{}} \forall a, m, n, k, l, p, q, r, w, b. \\ \quad \operatorname{is\text{-}ring}(a) \wedge \\ \quad (p \in \operatorname{mat}(m, n, \operatorname{carr}(a))) \wedge \\ \quad (q \in \operatorname{mat}(n, k, \operatorname{carr}(a))) \wedge \\ \quad (r \in \operatorname{mat}(k, l, \operatorname{carr}(a))) \wedge \\ \quad (1 \leq n) \wedge \\ \quad (1 \leq k) \wedge \\ \quad (w \in \operatorname{interval}(1, m)) \wedge \\ \quad (b \in \operatorname{interval}(1, l)) \Rightarrow \\ \quad \quad \operatorname{entry}(\operatorname{matmul}(a, p, \operatorname{matmul}(a, q, r)), w, b) \\ \quad \quad =\; \operatorname{finsum}(\operatorname{ring\text{-}additive\text{-}ag}(a), \\ \quad \quad \quad (\lambda j \in \operatorname{interval}(1, n).\; \\ \quad \quad \quad \quad \operatorname{finsum}(\operatorname{ring\text{-}additive\text{-}ag}(a), \\ \quad \quad \quad \quad \quad (\lambda c \in \operatorname{interval}(1, k).\; \\ \quad \quad \quad \quad \quad \quad (\lambda z \in (\operatorname{interval}(1, k) \times \operatorname{interval}(1, n)).\; \operatorname{mul}(a)(\operatorname{mul}(a)(\operatorname{entry}(p, w, \operatorname{nth}(2, z)), \operatorname{entry}(q, \operatorname{nth}(2, z), \operatorname{nth}(1, z))), \operatorname{entry}(r, \operatorname{nth}(1, z), b)))(\langle c, j \rangle)), \\ \quad \quad \quad \quad \quad \operatorname{interval}(1, k))), \\ \quad \quad \quad \operatorname{interval}(1, n)) \end{array}}}

\vnbfile{theorem-library/matmul-assoc-proof.scm}{theorem-library/matmul-assoc-proof.scm}
\noindent{\small matmul-assoc-proof.scm -- matrix multiplication is associative, via Fubini.}\par
\vnbentry{\texttt{matmul-\allowbreak{}assoc}}{\vnbstmt{\begin{array}{@{}l@{}} \forall a, m, n, k, l, p, q, r. \\ \quad \operatorname{is\text{-}ring}(a) \wedge \\ \quad (p \in \operatorname{mat}(m, n, \operatorname{carr}(a))) \wedge \\ \quad (q \in \operatorname{mat}(n, k, \operatorname{carr}(a))) \wedge \\ \quad (r \in \operatorname{mat}(k, l, \operatorname{carr}(a))) \wedge \\ \quad ((n = 0) \Rightarrow ((m = 0) \vee ((k = 0) \wedge (l = 0)))) \wedge \\ \quad ((k = 0) \Rightarrow ((l = 0) \vee ((n = 0) \wedge (m = 0)))) \Rightarrow \\ \quad \quad \operatorname{matmul}(a, \operatorname{matmul}(a, p, q), r) \\ \quad \quad =\; \operatorname{matmul}(a, p, \operatorname{matmul}(a, q, r)) \end{array}}}

\vnbfile{theorem-library/elem-entry-readoffs.scm}{theorem-library/elem-entry-readoffs.scm}
\noindent{\small elem-entry-readoffs.scm -- the elementary-matrix ENTRY read-offs, PROVEN.}\par
\vnbentry{\texttt{matunit-\allowbreak{}entries-\allowbreak{}defined}}{\vnbstmt{\begin{array}{@{}l@{}} \forall a, n, k, l.; \forall b, c \in \operatorname{interval}(1, n). \\ \quad (\operatorname{is\text{-}ring}(a) \wedge (n \in \mathbb{N})) \Rightarrow  \\ \quad \quad (\lambda \langle i, j \rangle \in (\operatorname{interval}(1, n) \times \operatorname{interval}(1, n)).\; \left(\text{if }((i = k) \wedge (j = l))\text{ then }\operatorname{one}(a)\text{ else }\operatorname{zero}(a)\right))(b, c) \\ \quad \quad \in\; \mathbf{Set} \end{array}}}
\vnbentry{\texttt{elem-\allowbreak{}f-\allowbreak{}entries-\allowbreak{}defined}}{\vnbstmt{\begin{array}{@{}l@{}} \forall a, n, k, l.; \forall b, c \in \operatorname{interval}(1, n). \\ \quad (\operatorname{is\text{-}ring}(a) \wedge (n \in \mathbb{N})) \Rightarrow  \\ \quad \quad (\lambda \langle i, j \rangle \in (\operatorname{interval}(1, n) \times \operatorname{interval}(1, n)).\; \left(\text{if }(((i = k) \wedge (j = l)) \vee (((i = l) \wedge (j = k)) \vee ((i = j) \wedge (\neg (i = k) \wedge \neg (i = l)))))\text{ then }\operatorname{one}(a)\text{ else }\operatorname{zero}(a)\right))(b, c) \\ \quad \quad \in\; \mathbf{Set} \end{array}}}
\vnbentry{\texttt{elem-\allowbreak{}g-\allowbreak{}entries-\allowbreak{}defined}}{\vnbstmt{\begin{array}{@{}l@{}} \forall a, n, r, k, l.; \forall b, c \in \operatorname{interval}(1, n). \\ \quad (\operatorname{is\text{-}ring}(a) \wedge ((n \in \mathbb{N}) \wedge (r \in \operatorname{carr}(a)))) \Rightarrow  \\ \quad \quad (\lambda \langle i, j \rangle \in (\operatorname{interval}(1, n) \times \operatorname{interval}(1, n)).\; \left(\text{if }(i = j)\text{ then }\operatorname{one}(a)\text{ else }\left(\text{if }((i = k) \wedge (j = l))\text{ then }r\text{ else }\operatorname{zero}(a)\right)\right))(b, c) \\ \quad \quad \in\; \mathbf{Set} \end{array}}}
\vnbentry{\texttt{elem-\allowbreak{}h-\allowbreak{}entries-\allowbreak{}defined}}{\vnbstmt{\begin{array}{@{}l@{}} \forall a, n, r, k.; \forall b, c \in \operatorname{interval}(1, n). \\ \quad (\operatorname{is\text{-}ring}(a) \wedge ((n \in \mathbb{N}) \wedge (r \in \operatorname{carr}(a)))) \Rightarrow  \\ \quad \quad (\lambda \langle i, j \rangle \in (\operatorname{interval}(1, n) \times \operatorname{interval}(1, n)).\; \left(\text{if }(i = j)\text{ then }\left(\text{if }(i = k)\text{ then }r\text{ else }\operatorname{one}(a)\right)\text{ else }\operatorname{zero}(a)\right))(b, c) \\ \quad \quad \in\; \mathbf{Set} \end{array}}}
\vnbentry{\texttt{entry-\allowbreak{}of-\allowbreak{}matunit}}{\vnbstmt{\begin{array}{@{}l@{}} \forall a, n, k, l, i, j. \\ \quad \operatorname{is\text{-}ring}(a) \wedge \\ \quad (n \in \mathbb{N}) \wedge \\ \quad (i \in \operatorname{interval}(1, n)) \wedge \\ \quad (j \in \operatorname{interval}(1, n)) \Rightarrow \\ \quad \quad \operatorname{entry}(\operatorname{matunit}(a, n, k, l), i, j) \\ \quad \quad =\; \left(\text{if }((i = k) \wedge (j = l))\text{ then }\operatorname{one}(a)\text{ else }\operatorname{zero}(a)\right) \end{array}}}
\vnbentry{\texttt{matunit-\allowbreak{}entry-\allowbreak{}off-\allowbreak{}row}}{\vnbstmt{\begin{array}{@{}l@{}} \forall a, n, k, l, j, c. \\ \quad \operatorname{is\text{-}ring}(a) \wedge \\ \quad (n \in \mathbb{N}) \wedge \\ \quad (j \in \operatorname{interval}(1, n)) \wedge \\ \quad (c \in \operatorname{interval}(1, n)) \wedge \\ \quad \neg (j = k) \Rightarrow \\ \quad \quad (\operatorname{entry}(\operatorname{matunit}(a, n, k, l), j, c) = \operatorname{zero}(a)) \end{array}}}
\vnbentry{\texttt{matunit-\allowbreak{}entry-\allowbreak{}k-\allowbreak{}row}}{\vnbstmt{\begin{array}{@{}l@{}} \forall a, n, k, l, c. \\ \quad \operatorname{is\text{-}ring}(a) \wedge \\ \quad (n \in \mathbb{N}) \wedge \\ \quad (k \in \operatorname{interval}(1, n)) \wedge \\ \quad (c \in \operatorname{interval}(1, n)) \Rightarrow \\ \quad \quad \operatorname{entry}(\operatorname{matunit}(a, n, k, l), k, c) \\ \quad \quad =\; \left(\text{if }(c = l)\text{ then }\operatorname{one}(a)\text{ else }\operatorname{zero}(a)\right) \end{array}}}
\vnbentry{\texttt{entry-\allowbreak{}of-\allowbreak{}elem-\allowbreak{}f}}{\vnbstmt{\begin{array}{@{}l@{}} \forall a, n, k, l, i, j. \\ \quad \operatorname{is\text{-}ring}(a) \wedge \\ \quad (n \in \mathbb{N}) \wedge \\ \quad (i \in \operatorname{interval}(1, n)) \wedge \\ \quad (j \in \operatorname{interval}(1, n)) \Rightarrow \\ \quad \quad \operatorname{entry}(\operatorname{elem\text{-}f}(a, n, k, l), i, j) \\ \quad \quad =\; \left(\text{if }(((i = k) \wedge (j = l)) \vee (((i = l) \wedge (j = k)) \vee ((i = j) \wedge (\neg (i = k) \wedge \neg (i = l)))))\text{ then }\operatorname{one}(a)\text{ else }\operatorname{zero}(a)\right) \end{array}}}
\vnbentry{\texttt{entry-\allowbreak{}of-\allowbreak{}elem-\allowbreak{}g}}{\vnbstmt{\begin{array}{@{}l@{}} \forall a, n, r, k, l, i, j. \\ \quad \operatorname{is\text{-}ring}(a) \wedge \\ \quad (n \in \mathbb{N}) \wedge \\ \quad (r \in \operatorname{carr}(a)) \wedge \\ \quad (i \in \operatorname{interval}(1, n)) \wedge \\ \quad (j \in \operatorname{interval}(1, n)) \Rightarrow \\ \quad \quad \operatorname{entry}(\operatorname{elem\text{-}g}(a, n, r, k, l), i, j) \\ \quad \quad =\; \left(\text{if }(i = j)\text{ then }\operatorname{one}(a)\text{ else }\left(\text{if }((i = k) \wedge (j = l))\text{ then }r\text{ else }\operatorname{zero}(a)\right)\right) \end{array}}}
\vnbentry{\texttt{entry-\allowbreak{}of-\allowbreak{}elem-\allowbreak{}h}}{\vnbstmt{\begin{array}{@{}l@{}} \forall a, n, r, k, i, j. \\ \quad \operatorname{is\text{-}ring}(a) \wedge \\ \quad (n \in \mathbb{N}) \wedge \\ \quad (r \in \operatorname{carr}(a)) \wedge \\ \quad (i \in \operatorname{interval}(1, n)) \wedge \\ \quad (j \in \operatorname{interval}(1, n)) \Rightarrow \\ \quad \quad \operatorname{entry}(\operatorname{elem\text{-}h}(a, n, r, k), i, j) \\ \quad \quad =\; \left(\text{if }(i = j)\text{ then }\left(\text{if }(i = k)\text{ then }r\text{ else }\operatorname{one}(a)\right)\text{ else }\operatorname{zero}(a)\right) \end{array}}}
\vnbentry{\texttt{elem-\allowbreak{}g-\allowbreak{}entry-\allowbreak{}off-\allowbreak{}l-\allowbreak{}off-\allowbreak{}diag}}{\vnbstmt{\begin{array}{@{}l@{}} \forall a, n, r, k, l, j, c. \\ \quad \operatorname{is\text{-}ring}(a) \wedge \\ \quad (n \in \mathbb{N}) \wedge \\ \quad (r \in \operatorname{carr}(a)) \wedge \\ \quad (j \in \operatorname{interval}(1, n)) \wedge \\ \quad (c \in \operatorname{interval}(1, n)) \wedge \\ \quad \neg (c = l) \wedge \\ \quad \neg (j = c) \Rightarrow \\ \quad \quad (\operatorname{entry}(\operatorname{elem\text{-}g}(a, n, r, k, l), j, c) = \operatorname{zero}(a)) \end{array}}}
\vnbentry{\texttt{elem-\allowbreak{}g-\allowbreak{}entry-\allowbreak{}off-\allowbreak{}l-\allowbreak{}diag}}{\vnbstmt{\begin{array}{@{}l@{}} \forall a, n, r, k, l, c. \\ \quad \operatorname{is\text{-}ring}(a) \wedge \\ \quad (n \in \mathbb{N}) \wedge \\ \quad (r \in \operatorname{carr}(a)) \wedge \\ \quad (c \in \operatorname{interval}(1, n)) \wedge \\ \quad \neg (c = l) \Rightarrow \\ \quad \quad (\operatorname{entry}(\operatorname{elem\text{-}g}(a, n, r, k, l), c, c) = \operatorname{one}(a)) \end{array}}}
\vnbentry{\texttt{elem-\allowbreak{}g-\allowbreak{}entry-\allowbreak{}l-\allowbreak{}vanish}}{\vnbstmt{\begin{array}{@{}l@{}} \forall a, n, r, k, l, j, c. \\ \quad \operatorname{is\text{-}ring}(a) \wedge \\ \quad (n \in \mathbb{N}) \wedge \\ \quad (r \in \operatorname{carr}(a)) \wedge \\ \quad (j \in \operatorname{interval}(1, n)) \wedge \\ \quad (c \in \operatorname{interval}(1, n)) \wedge \\ \quad (c = l) \wedge \\ \quad \neg (j = l) \wedge \\ \quad \neg (j = k) \Rightarrow \\ \quad \quad (\operatorname{entry}(\operatorname{elem\text{-}g}(a, n, r, k, l), j, c) = \operatorname{zero}(a)) \end{array}}}
\vnbentry{\texttt{elem-\allowbreak{}g-\allowbreak{}entry-\allowbreak{}l-\allowbreak{}at-\allowbreak{}l}}{\vnbstmt{\begin{array}{@{}l@{}} \forall a, n, r, k, l, c. \\ \quad \operatorname{is\text{-}ring}(a) \wedge \\ \quad (n \in \mathbb{N}) \wedge \\ \quad (r \in \operatorname{carr}(a)) \wedge \\ \quad (c \in \operatorname{interval}(1, n)) \wedge \\ \quad (c = l) \Rightarrow \\ \quad \quad (\operatorname{entry}(\operatorname{elem\text{-}g}(a, n, r, k, l), l, c) = \operatorname{one}(a)) \end{array}}}
\vnbentry{\texttt{elem-\allowbreak{}f-\allowbreak{}ck-\allowbreak{}off}}{\vnbstmt{\begin{array}{@{}l@{}} \forall a, n, k, l, j, c. \\ \quad \operatorname{is\text{-}ring}(a) \wedge \\ \quad (n \in \mathbb{N}) \wedge \\ \quad (j \in \operatorname{interval}(1, n)) \wedge \\ \quad (c \in \operatorname{interval}(1, n)) \wedge \\ \quad (c = k) \wedge \\ \quad \neg (k = l) \wedge \\ \quad \neg (j = l) \Rightarrow \\ \quad \quad (\operatorname{entry}(\operatorname{elem\text{-}f}(a, n, k, l), j, c) = \operatorname{zero}(a)) \end{array}}}
\vnbentry{\texttt{elem-\allowbreak{}f-\allowbreak{}cl-\allowbreak{}off}}{\vnbstmt{\begin{array}{@{}l@{}} \forall a, n, k, l, j, c. \\ \quad \operatorname{is\text{-}ring}(a) \wedge \\ \quad (n \in \mathbb{N}) \wedge \\ \quad (j \in \operatorname{interval}(1, n)) \wedge \\ \quad (c \in \operatorname{interval}(1, n)) \wedge \\ \quad (c = l) \wedge \\ \quad \neg (k = l) \wedge \\ \quad \neg (j = k) \Rightarrow \\ \quad \quad (\operatorname{entry}(\operatorname{elem\text{-}f}(a, n, k, l), j, c) = \operatorname{zero}(a)) \end{array}}}
\vnbentry{\texttt{elem-\allowbreak{}f-\allowbreak{}co-\allowbreak{}off}}{\vnbstmt{\begin{array}{@{}l@{}} \forall a, n, k, l, j, c. \\ \quad \operatorname{is\text{-}ring}(a) \wedge \\ \quad (n \in \mathbb{N}) \wedge \\ \quad (j \in \operatorname{interval}(1, n)) \wedge \\ \quad (c \in \operatorname{interval}(1, n)) \wedge \\ \quad \neg (c = k) \wedge \\ \quad \neg (c = l) \wedge \\ \quad \neg (j = c) \Rightarrow \\ \quad \quad (\operatorname{entry}(\operatorname{elem\text{-}f}(a, n, k, l), j, c) = \operatorname{zero}(a)) \end{array}}}
\vnbentry{\texttt{elem-\allowbreak{}f-\allowbreak{}co-\allowbreak{}at}}{\vnbstmt{\begin{array}{@{}l@{}} \forall a, n, k, l, c. \\ \quad \operatorname{is\text{-}ring}(a) \wedge \\ \quad (n \in \mathbb{N}) \wedge \\ \quad (c \in \operatorname{interval}(1, n)) \wedge \\ \quad \neg (c = k) \wedge \\ \quad \neg (c = l) \Rightarrow \\ \quad \quad (\operatorname{entry}(\operatorname{elem\text{-}f}(a, n, k, l), c, c) = \operatorname{one}(a)) \end{array}}}
\vnbentry{\texttt{elem-\allowbreak{}h-\allowbreak{}entry-\allowbreak{}off-\allowbreak{}diag}}{\vnbstmt{\begin{array}{@{}l@{}} \forall a, n, r, k, j, c. \\ \quad \operatorname{is\text{-}ring}(a) \wedge \\ \quad (n \in \mathbb{N}) \wedge \\ \quad (r \in \operatorname{carr}(a)) \wedge \\ \quad (j \in \operatorname{interval}(1, n)) \wedge \\ \quad (c \in \operatorname{interval}(1, n)) \wedge \\ \quad \neg (j = c) \Rightarrow \\ \quad \quad (\operatorname{entry}(\operatorname{elem\text{-}h}(a, n, r, k), j, c) = \operatorname{zero}(a)) \end{array}}}
\vnbentry{\texttt{elem-\allowbreak{}h-\allowbreak{}entry-\allowbreak{}diag}}{\vnbstmt{\begin{array}{@{}l@{}} \forall a, n, r, k, c. \\ \quad \operatorname{is\text{-}ring}(a) \wedge \\ \quad (n \in \mathbb{N}) \wedge \\ \quad (r \in \operatorname{carr}(a)) \wedge \\ \quad (c \in \operatorname{interval}(1, n)) \Rightarrow \\ \quad \quad \operatorname{entry}(\operatorname{elem\text{-}h}(a, n, r, k), c, c) \\ \quad \quad =\; \left(\text{if }(c = k)\text{ then }r\text{ else }\operatorname{one}(a)\right) \end{array}}}
\vnbentry{\texttt{elem-\allowbreak{}f-\allowbreak{}rk-\allowbreak{}off}}{\vnbstmt{\begin{array}{@{}l@{}} \forall a, n, k, l, i, j. \\ \quad \operatorname{is\text{-}ring}(a) \wedge \\ \quad (n \in \mathbb{N}) \wedge \\ \quad (i \in \operatorname{interval}(1, n)) \wedge \\ \quad (j \in \operatorname{interval}(1, n)) \wedge \\ \quad (i = k) \wedge \\ \quad \neg (k = l) \wedge \\ \quad \neg (j = l) \Rightarrow \\ \quad \quad (\operatorname{entry}(\operatorname{elem\text{-}f}(a, n, k, l), i, j) = \operatorname{zero}(a)) \end{array}}}
\vnbentry{\texttt{elem-\allowbreak{}f-\allowbreak{}rl-\allowbreak{}off}}{\vnbstmt{\begin{array}{@{}l@{}} \forall a, n, k, l, i, j. \\ \quad \operatorname{is\text{-}ring}(a) \wedge \\ \quad (n \in \mathbb{N}) \wedge \\ \quad (i \in \operatorname{interval}(1, n)) \wedge \\ \quad (j \in \operatorname{interval}(1, n)) \wedge \\ \quad (i = l) \wedge \\ \quad \neg (k = l) \wedge \\ \quad \neg (j = k) \Rightarrow \\ \quad \quad (\operatorname{entry}(\operatorname{elem\text{-}f}(a, n, k, l), i, j) = \operatorname{zero}(a)) \end{array}}}
\vnbentry{\texttt{elem-\allowbreak{}f-\allowbreak{}ro-\allowbreak{}off}}{\vnbstmt{\begin{array}{@{}l@{}} \forall a, n, k, l, i, j. \\ \quad \operatorname{is\text{-}ring}(a) \wedge \\ \quad (n \in \mathbb{N}) \wedge \\ \quad (i \in \operatorname{interval}(1, n)) \wedge \\ \quad (j \in \operatorname{interval}(1, n)) \wedge \\ \quad \neg (i = k) \wedge \\ \quad \neg (i = l) \wedge \\ \quad \neg (j = i) \Rightarrow \\ \quad \quad (\operatorname{entry}(\operatorname{elem\text{-}f}(a, n, k, l), i, j) = \operatorname{zero}(a)) \end{array}}}
\vnbentry{\texttt{elem-\allowbreak{}g-\allowbreak{}rk-\allowbreak{}vanish}}{\vnbstmt{\begin{array}{@{}l@{}} \forall a, n, r, k, l, i, j. \\ \quad \operatorname{is\text{-}ring}(a) \wedge \\ \quad (n \in \mathbb{N}) \wedge \\ \quad (r \in \operatorname{carr}(a)) \wedge \\ \quad (i \in \operatorname{interval}(1, n)) \wedge \\ \quad (j \in \operatorname{interval}(1, n)) \wedge \\ \quad (i = k) \wedge \\ \quad \neg (j = k) \wedge \\ \quad \neg (j = l) \Rightarrow \\ \quad \quad (\operatorname{entry}(\operatorname{elem\text{-}g}(a, n, r, k, l), i, j) = \operatorname{zero}(a)) \end{array}}}
\vnbentry{\texttt{elem-\allowbreak{}g-\allowbreak{}ro-\allowbreak{}off}}{\vnbstmt{\begin{array}{@{}l@{}} \forall a, n, r, k, l, i, j. \\ \quad \operatorname{is\text{-}ring}(a) \wedge \\ \quad (n \in \mathbb{N}) \wedge \\ \quad (r \in \operatorname{carr}(a)) \wedge \\ \quad (i \in \operatorname{interval}(1, n)) \wedge \\ \quad (j \in \operatorname{interval}(1, n)) \wedge \\ \quad \neg (i = k) \wedge \\ \quad \neg (j = i) \Rightarrow \\ \quad \quad (\operatorname{entry}(\operatorname{elem\text{-}g}(a, n, r, k, l), i, j) = \operatorname{zero}(a)) \end{array}}}
\vnbentry{\texttt{elem-\allowbreak{}h-\allowbreak{}ro-\allowbreak{}off}}{\vnbstmt{\begin{array}{@{}l@{}} \forall a, n, r, k, i, j. \\ \quad \operatorname{is\text{-}ring}(a) \wedge \\ \quad (n \in \mathbb{N}) \wedge \\ \quad (r \in \operatorname{carr}(a)) \wedge \\ \quad (i \in \operatorname{interval}(1, n)) \wedge \\ \quad (j \in \operatorname{interval}(1, n)) \wedge \\ \quad \neg (j = i) \Rightarrow \\ \quad \quad (\operatorname{entry}(\operatorname{elem\text{-}h}(a, n, r, k), i, j) = \operatorname{zero}(a)) \end{array}}}
\vnbentry{\texttt{elem-\allowbreak{}f-\allowbreak{}ck-\allowbreak{}at}}{\vnbstmt{\begin{array}{@{}l@{}} \forall a, n, k, l, c. \\ \quad \operatorname{is\text{-}ring}(a) \wedge \\ \quad (n \in \mathbb{N}) \wedge \\ \quad (l \in \operatorname{interval}(1, n)) \wedge \\ \quad (c \in \operatorname{interval}(1, n)) \wedge \\ \quad (c = k) \Rightarrow \\ \quad \quad (\operatorname{entry}(\operatorname{elem\text{-}f}(a, n, k, l), l, c) = \operatorname{one}(a)) \end{array}}}
\vnbentry{\texttt{elem-\allowbreak{}f-\allowbreak{}cl-\allowbreak{}at}}{\vnbstmt{\begin{array}{@{}l@{}} \forall a, n, k, l, c. \\ \quad \operatorname{is\text{-}ring}(a) \wedge \\ \quad (n \in \mathbb{N}) \wedge \\ \quad (k \in \operatorname{interval}(1, n)) \wedge \\ \quad (c \in \operatorname{interval}(1, n)) \wedge \\ \quad (c = l) \Rightarrow \\ \quad \quad (\operatorname{entry}(\operatorname{elem\text{-}f}(a, n, k, l), k, c) = \operatorname{one}(a)) \end{array}}}
\vnbentry{\texttt{elem-\allowbreak{}f-\allowbreak{}col-\allowbreak{}k}}{\vnbstmt{\begin{array}{@{}l@{}} \forall a, n, k, l, i. \\ \quad \operatorname{is\text{-}ring}(a) \wedge \\ \quad (n \in \mathbb{N}) \wedge \\ \quad (i \in \operatorname{interval}(1, n)) \wedge \\ \quad (k \in \operatorname{interval}(1, n)) \wedge \\ \quad \neg (k = l) \Rightarrow \\ \quad \quad \operatorname{entry}(\operatorname{elem\text{-}f}(a, n, k, l), i, k) \\ \quad \quad =\; \left(\text{if }(i = l)\text{ then }\operatorname{one}(a)\text{ else }\operatorname{zero}(a)\right) \end{array}}}
\vnbentry{\texttt{elem-\allowbreak{}f-\allowbreak{}col-\allowbreak{}l}}{\vnbstmt{\begin{array}{@{}l@{}} \forall a, n, k, l, i. \\ \quad \operatorname{is\text{-}ring}(a) \wedge \\ \quad (n \in \mathbb{N}) \wedge \\ \quad (i \in \operatorname{interval}(1, n)) \wedge \\ \quad (l \in \operatorname{interval}(1, n)) \wedge \\ \quad \neg (k = l) \Rightarrow \\ \quad \quad \operatorname{entry}(\operatorname{elem\text{-}f}(a, n, k, l), i, l) \\ \quad \quad =\; \left(\text{if }(i = k)\text{ then }\operatorname{one}(a)\text{ else }\operatorname{zero}(a)\right) \end{array}}}
\vnbentry{\texttt{elem-\allowbreak{}f-\allowbreak{}col-\allowbreak{}other}}{\vnbstmt{\begin{array}{@{}l@{}} \forall a, n, k, l, i, c. \\ \quad \operatorname{is\text{-}ring}(a) \wedge \\ \quad (n \in \mathbb{N}) \wedge \\ \quad (i \in \operatorname{interval}(1, n)) \wedge \\ \quad (c \in \operatorname{interval}(1, n)) \wedge \\ \quad \neg (c = k) \wedge \\ \quad \neg (c = l) \Rightarrow \\ \quad \quad \operatorname{entry}(\operatorname{elem\text{-}f}(a, n, k, l), i, c) \\ \quad \quad =\; \left(\text{if }(i = c)\text{ then }\operatorname{one}(a)\text{ else }\operatorname{zero}(a)\right) \end{array}}}
\vnbentry{\texttt{elem-\allowbreak{}f-\allowbreak{}rk-\allowbreak{}at}}{\vnbstmt{\begin{array}{@{}l@{}} \forall a, n, k, l, i. \\ \quad \operatorname{is\text{-}ring}(a) \wedge \\ \quad (n \in \mathbb{N}) \wedge \\ \quad (i \in \operatorname{interval}(1, n)) \wedge \\ \quad (l \in \operatorname{interval}(1, n)) \wedge \\ \quad (i = k) \wedge \\ \quad \neg (k = l) \Rightarrow \\ \quad \quad (\operatorname{entry}(\operatorname{elem\text{-}f}(a, n, k, l), i, l) = \operatorname{one}(a)) \end{array}}}
\vnbentry{\texttt{elem-\allowbreak{}f-\allowbreak{}rl-\allowbreak{}at}}{\vnbstmt{\begin{array}{@{}l@{}} \forall a, n, k, l, i. \\ \quad \operatorname{is\text{-}ring}(a) \wedge \\ \quad (n \in \mathbb{N}) \wedge \\ \quad (i \in \operatorname{interval}(1, n)) \wedge \\ \quad (k \in \operatorname{interval}(1, n)) \wedge \\ \quad (i = l) \wedge \\ \quad \neg (k = l) \Rightarrow \\ \quad \quad (\operatorname{entry}(\operatorname{elem\text{-}f}(a, n, k, l), i, k) = \operatorname{one}(a)) \end{array}}}
\vnbentry{\texttt{elem-\allowbreak{}g-\allowbreak{}entry-\allowbreak{}l-\allowbreak{}at-\allowbreak{}k}}{\vnbstmt{\begin{array}{@{}l@{}} \forall a, n, r, k, l, c. \\ \quad \operatorname{is\text{-}ring}(a) \wedge \\ \quad (n \in \mathbb{N}) \wedge \\ \quad (r \in \operatorname{carr}(a)) \wedge \\ \quad (k \in \operatorname{interval}(1, n)) \wedge \\ \quad (c \in \operatorname{interval}(1, n)) \wedge \\ \quad (c = l) \wedge \\ \quad \neg (k = l) \Rightarrow \\ \quad \quad (\operatorname{entry}(\operatorname{elem\text{-}g}(a, n, r, k, l), k, c) = r) \end{array}}}
\vnbentry{\texttt{elem-\allowbreak{}g-\allowbreak{}col-\allowbreak{}k}}{\vnbstmt{\begin{array}{@{}l@{}} \forall a, n, r, k, l, i. \\ \quad \operatorname{is\text{-}ring}(a) \wedge \\ \quad (n \in \mathbb{N}) \wedge \\ \quad (r \in \operatorname{carr}(a)) \wedge \\ \quad (i \in \operatorname{interval}(1, n)) \wedge \\ \quad (k \in \operatorname{interval}(1, n)) \wedge \\ \quad \neg (k = l) \Rightarrow \\ \quad \quad \operatorname{entry}(\operatorname{elem\text{-}g}(a, n, r, k, l), i, k) \\ \quad \quad =\; \left(\text{if }(i = k)\text{ then }\operatorname{one}(a)\text{ else }\operatorname{zero}(a)\right) \end{array}}}
\vnbentry{\texttt{elem-\allowbreak{}g-\allowbreak{}col-\allowbreak{}l}}{\vnbstmt{\begin{array}{@{}l@{}} \forall a, n, r, k, l, i. \\ \quad \operatorname{is\text{-}ring}(a) \wedge \\ \quad (n \in \mathbb{N}) \wedge \\ \quad (r \in \operatorname{carr}(a)) \wedge \\ \quad (i \in \operatorname{interval}(1, n)) \wedge \\ \quad (l \in \operatorname{interval}(1, n)) \Rightarrow \\ \quad \quad \operatorname{entry}(\operatorname{elem\text{-}g}(a, n, r, k, l), i, l) \\ \quad \quad =\; \left(\text{if }(i = l)\text{ then }\operatorname{one}(a)\text{ else }\left(\text{if }(i = k)\text{ then }r\text{ else }\operatorname{zero}(a)\right)\right) \end{array}}}
\vnbentry{\texttt{elem-\allowbreak{}g-\allowbreak{}col-\allowbreak{}other}}{\vnbstmt{\begin{array}{@{}l@{}} \forall a, n, r, k, l, i, c. \\ \quad \operatorname{is\text{-}ring}(a) \wedge \\ \quad (n \in \mathbb{N}) \wedge \\ \quad (r \in \operatorname{carr}(a)) \wedge \\ \quad (i \in \operatorname{interval}(1, n)) \wedge \\ \quad (c \in \operatorname{interval}(1, n)) \wedge \\ \quad \neg (c = l) \Rightarrow \\ \quad \quad \operatorname{entry}(\operatorname{elem\text{-}g}(a, n, r, k, l), i, c) \\ \quad \quad =\; \left(\text{if }(i = c)\text{ then }\operatorname{one}(a)\text{ else }\operatorname{zero}(a)\right) \end{array}}}
\vnbentry{\texttt{elem-\allowbreak{}g-\allowbreak{}rk-\allowbreak{}at-\allowbreak{}k}}{\vnbstmt{\begin{array}{@{}l@{}} \forall a, n, r, k, l, i. \\ \quad \operatorname{is\text{-}ring}(a) \wedge \\ \quad (n \in \mathbb{N}) \wedge \\ \quad (r \in \operatorname{carr}(a)) \wedge \\ \quad (i \in \operatorname{interval}(1, n)) \wedge \\ \quad (k \in \operatorname{interval}(1, n)) \wedge \\ \quad (i = k) \Rightarrow \\ \quad \quad (\operatorname{entry}(\operatorname{elem\text{-}g}(a, n, r, k, l), i, k) = \operatorname{one}(a)) \end{array}}}
\vnbentry{\texttt{elem-\allowbreak{}g-\allowbreak{}rk-\allowbreak{}at-\allowbreak{}l}}{\vnbstmt{\begin{array}{@{}l@{}} \forall a, n, r, k, l, i. \\ \quad \operatorname{is\text{-}ring}(a) \wedge \\ \quad (n \in \mathbb{N}) \wedge \\ \quad (r \in \operatorname{carr}(a)) \wedge \\ \quad (i \in \operatorname{interval}(1, n)) \wedge \\ \quad (l \in \operatorname{interval}(1, n)) \wedge \\ \quad (i = k) \wedge \\ \quad \neg (k = l) \Rightarrow \\ \quad \quad (\operatorname{entry}(\operatorname{elem\text{-}g}(a, n, r, k, l), i, l) = r) \end{array}}}
\vnbentry{\texttt{elem-\allowbreak{}g-\allowbreak{}ro-\allowbreak{}at}}{\vnbstmt{\begin{array}{@{}l@{}} \forall a, n, r, k, l, i. \\ \quad \operatorname{is\text{-}ring}(a) \wedge \\ \quad (n \in \mathbb{N}) \wedge \\ \quad (r \in \operatorname{carr}(a)) \wedge \\ \quad (i \in \operatorname{interval}(1, n)) \wedge \\ \quad \neg (i = k) \Rightarrow \\ \quad \quad (\operatorname{entry}(\operatorname{elem\text{-}g}(a, n, r, k, l), i, i) = \operatorname{one}(a)) \end{array}}}
\vnbentry{\texttt{identmat-\allowbreak{}entries-\allowbreak{}defined}}{\vnbstmt{\begin{array}{@{}l@{}} \forall a, n.; \forall b, c \in \operatorname{interval}(1, n). \\ \quad (\operatorname{is\text{-}ring}(a) \wedge (n \in \mathbb{N})) \Rightarrow  \\ \quad \quad (\lambda \langle i, j \rangle \in (\operatorname{interval}(1, n) \times \operatorname{interval}(1, n)).\; \left(\text{if }(i = j)\text{ then }\operatorname{one}(a)\text{ else }\operatorname{zero}(a)\right))(b, c) \\ \quad \quad \in\; \mathbf{Set} \end{array}}}
\vnbentry{\texttt{unitrow-\allowbreak{}entries-\allowbreak{}defined}}{\vnbstmt{\begin{array}{@{}l@{}} \forall a, n, i.; \forall b \in \operatorname{interval}(1, 1), j \in \operatorname{interval}(1, n). \\ \quad (\operatorname{is\text{-}ring}(a) \wedge (n \in \mathbb{N})) \Rightarrow  \\ \quad \quad (\lambda \langle \mathit{rw}, \mathit{cl} \rangle \in (\operatorname{interval}(1, 1) \times \operatorname{interval}(1, n)).\; \left(\text{if }(\mathit{cl} = i)\text{ then }\operatorname{one}(a)\text{ else }\operatorname{zero}(a)\right))(b, j) \\ \quad \quad \in\; \mathbf{Set} \end{array}}}
\vnbentry{\texttt{zeromat-\allowbreak{}entries-\allowbreak{}defined}}{\vnbstmt{\begin{array}{@{}l@{}} \forall a, m, n.; \forall b \in \operatorname{interval}(1, m), c \in \operatorname{interval}(1, n). \\ \quad (\operatorname{is\text{-}ring}(a) \wedge ((m \in \mathbb{N}) \wedge (n \in \mathbb{N}))) \Rightarrow  \\ \quad \quad (\lambda \langle i, j \rangle \in (\operatorname{interval}(1, m) \times \operatorname{interval}(1, n)).\; \operatorname{zero}(a))(b, c) \\ \quad \quad \in\; \mathbf{Set} \end{array}}}
\vnbentry{\texttt{entry-\allowbreak{}of-\allowbreak{}identmat}}{\vnbstmt{\begin{array}{@{}l@{}} \forall a, n, i, j. \\ \quad \operatorname{is\text{-}ring}(a) \wedge \\ \quad (n \in \mathbb{N}) \wedge \\ \quad (i \in \operatorname{interval}(1, n)) \wedge \\ \quad (j \in \operatorname{interval}(1, n)) \Rightarrow \\ \quad \quad \operatorname{entry}(\operatorname{identmat}(a, n), i, j) \\ \quad \quad =\; \left(\text{if }(i = j)\text{ then }\operatorname{one}(a)\text{ else }\operatorname{zero}(a)\right) \end{array}}}
\vnbentry{\texttt{identmat-\allowbreak{}entry-\allowbreak{}diag}}{\vnbstmt{\begin{array}{@{}l@{}} \forall a, n, i. \\ \quad \operatorname{is\text{-}ring}(a) \wedge \\ \quad (n \in \mathbb{N}) \wedge \\ \quad (i \in \operatorname{interval}(1, n)) \Rightarrow \\ \quad \quad (\operatorname{entry}(\operatorname{identmat}(a, n), i, i) = \operatorname{one}(a)) \end{array}}}
\vnbentry{\texttt{identmat-\allowbreak{}entry-\allowbreak{}off}}{\vnbstmt{\begin{array}{@{}l@{}} \forall a, n, i, j. \\ \quad \operatorname{is\text{-}ring}(a) \wedge \\ \quad (n \in \mathbb{N}) \wedge \\ \quad (i \in \operatorname{interval}(1, n)) \wedge \\ \quad (j \in \operatorname{interval}(1, n)) \wedge \\ \quad \neg (i = j) \Rightarrow \\ \quad \quad (\operatorname{entry}(\operatorname{identmat}(a, n), i, j) = \operatorname{zero}(a)) \end{array}}}
\vnbentry{\texttt{unitrow-\allowbreak{}entry-\allowbreak{}at}}{\vnbstmt{\begin{array}{@{}l@{}} \forall a, n, i. \\ \quad \operatorname{is\text{-}ring}(a) \wedge \\ \quad (n \in \mathbb{N}) \wedge \\ \quad (i \in \operatorname{interval}(1, n)) \Rightarrow \\ \quad \quad (\operatorname{entry}(\operatorname{unitrow}(a, n, i), 1, i) = \operatorname{one}(a)) \end{array}}}
\vnbentry{\texttt{unitrow-\allowbreak{}entry-\allowbreak{}off}}{\vnbstmt{\begin{array}{@{}l@{}} \forall a, n, i, j. \\ \quad \operatorname{is\text{-}ring}(a) \wedge \\ \quad (n \in \mathbb{N}) \wedge \\ \quad (i \in \operatorname{interval}(1, n)) \wedge \\ \quad (j \in \operatorname{interval}(1, n)) \wedge \\ \quad \neg (j = i) \Rightarrow \\ \quad \quad (\operatorname{entry}(\operatorname{unitrow}(a, n, i), 1, j) = \operatorname{zero}(a)) \end{array}}}
\vnbentry{\texttt{entry-\allowbreak{}of-\allowbreak{}zeromat}}{\vnbstmt{\begin{array}{@{}l@{}} \forall a, m, n, i, j. \\ \quad \operatorname{is\text{-}ring}(a) \wedge \\ \quad (m \in \mathbb{N}) \wedge \\ \quad (n \in \mathbb{N}) \wedge \\ \quad (i \in \operatorname{interval}(1, m)) \wedge \\ \quad (j \in \operatorname{interval}(1, n)) \Rightarrow \\ \quad \quad (\operatorname{entry}(\operatorname{zeromat}(a, m, n), i, j) = \operatorname{zero}(a)) \end{array}}}
\vnbentry{\texttt{elem-\allowbreak{}f-\allowbreak{}type}}{\vnbstmt{\begin{array}{@{}l@{}} \forall a, n, k, l. \\ \quad (\operatorname{is\text{-}ring}(a) \wedge (n \in \mathbb{N})) \Rightarrow  \\ \quad \quad (\operatorname{elem\text{-}f}(a, n, k, l) \in \operatorname{mat}(n, n, \operatorname{carr}(a))) \end{array}}}
\vnbentry{\texttt{elem-\allowbreak{}f-\allowbreak{}symmetric}}{\vnbstmt{\begin{array}{@{}l@{}} \forall a, n, k, l. \\ \quad (\operatorname{is\text{-}ring}(a) \wedge (n \in \mathbb{N})) \Rightarrow  \\ \quad \quad (\operatorname{elem\text{-}f}(a, n, k, l) = \operatorname{elem\text{-}f}(a, n, l, k)) \end{array}}}
\vnbentry{\texttt{entry-\allowbreak{}of-\allowbreak{}block}}{\vnbstmt{\begin{array}{@{}l@{}} \forall p, k, l, i, j, m, n, x. \\ \quad (k \in \mathbb{N}) \wedge \\ \quad (l \in \mathbb{N}) \wedge \\ \quad (p \in \operatorname{mat}(m, n, x)) \wedge \\ \quad (k \leq m) \wedge \\ \quad (l \leq n) \wedge \\ \quad (i \in \operatorname{interval}(1, k)) \wedge \\ \quad (j \in \operatorname{interval}(1, l)) \Rightarrow \\ \quad \quad (\operatorname{entry}(\operatorname{block}(p, k, l), i, j) = \operatorname{entry}(p, i, j)) \end{array}}}
\vnbentry{\texttt{entry-\allowbreak{}of-\allowbreak{}snoc-\allowbreak{}col}}{\vnbstmt{\begin{array}{@{}l@{}} \forall w, n, v, i, x. \\ \quad (n \in \mathbb{N}) \wedge \\ \quad (w \in \operatorname{mat}(n, 1, x)) \wedge \\ \quad (v \in x) \wedge \\ \quad (i \in \operatorname{interval}(1, n)) \Rightarrow \\ \quad \quad \operatorname{entry}(\operatorname{snoc\text{-}col}(w, n, v), i, 1) \\ \quad \quad =\; \operatorname{entry}(w, i, 1) \end{array}}}
\vnbentry{\texttt{snoc-\allowbreak{}col-\allowbreak{}last}}{\vnbstmt{\begin{array}{@{}l@{}} \forall w, n, v, x. \\ \quad ((n \in \mathbb{N}) \wedge ((w \in \operatorname{mat}(n, 1, x)) \wedge (v \in x))) \Rightarrow  \\ \quad \quad (\operatorname{entry}(\operatorname{snoc\text{-}col}(w, n, v), (n + 1), 1) = v) \end{array}}}
\vnbentry{\texttt{entry-\allowbreak{}of-\allowbreak{}snoc-\allowbreak{}row}}{\vnbstmt{\begin{array}{@{}l@{}} \forall c, n, r, j, x. \\ \quad (n \in \mathbb{N}) \wedge \\ \quad (c \in \operatorname{mat}(1, n, x)) \wedge \\ \quad (r \in x) \wedge \\ \quad (j \in \operatorname{interval}(1, n)) \Rightarrow \\ \quad \quad \operatorname{entry}(\operatorname{snoc\text{-}row}(c, n, r), 1, j) \\ \quad \quad =\; \operatorname{entry}(c, 1, j) \end{array}}}
\vnbentry{\texttt{snoc-\allowbreak{}row-\allowbreak{}last}}{\vnbstmt{\begin{array}{@{}l@{}} \forall c, n, r, x. \\ \quad ((n \in \mathbb{N}) \wedge ((c \in \operatorname{mat}(1, n, x)) \wedge (r \in x))) \Rightarrow  \\ \quad \quad (\operatorname{entry}(\operatorname{snoc\text{-}row}(c, n, r), 1, (n + 1)) = r) \end{array}}}
\vnbentry{\texttt{matadd-\allowbreak{}entry}}{\vnbstmt{\begin{array}{@{}l@{}} \forall a, m, n, p, q.; \forall i \in \operatorname{interval}(1, m), j \in \operatorname{interval}(1, n). \\ \quad \operatorname{is\text{-}ring}(a) \wedge \\ \quad (p \in \operatorname{mat}(m, n, \operatorname{carr}(a))) \wedge \\ \quad (q \in \operatorname{mat}(m, n, \operatorname{carr}(a))) \Rightarrow \\ \quad \quad \operatorname{entry}(\operatorname{matadd}(a, p, q), i, j) \\ \quad \quad =\; \operatorname{add}(a)(\operatorname{entry}(p, i, j), \operatorname{entry}(q, i, j)) \end{array}}}
\vnbentry{\texttt{matscale-\allowbreak{}entry}}{\vnbstmt{\begin{array}{@{}l@{}} \forall a, m, n, r, p.; \forall i \in \operatorname{interval}(1, m), j \in \operatorname{interval}(1, n). \\ \quad \operatorname{is\text{-}ring}(a) \wedge \\ \quad (r \in \operatorname{carr}(a)) \wedge \\ \quad (p \in \operatorname{mat}(m, n, \operatorname{carr}(a))) \Rightarrow \\ \quad \quad \operatorname{entry}(\operatorname{matscale}(a, r, p), i, j) \\ \quad \quad =\; \operatorname{mul}(a)(r, \operatorname{entry}(p, i, j)) \end{array}}}
\vnbentry{\texttt{entry-\allowbreak{}of-\allowbreak{}submat}}{\vnbstmt{\begin{array}{@{}l@{}} \forall s, p, q, i, j, x. \\ \quad (p \in \mathbb{N}) \wedge \\ \quad (q \in \mathbb{N}) \wedge \\ \quad (s \in \operatorname{mat}((p + 1), (q + 1), x)) \wedge \\ \quad (i \in \operatorname{interval}(1, p)) \wedge \\ \quad (j \in \operatorname{interval}(1, q)) \Rightarrow \\ \quad \quad \operatorname{entry}(\operatorname{submat}(s, p, q), i, j) \\ \quad \quad =\; \operatorname{entry}(s, (i + 1), (j + 1)) \end{array}}}

\vnbfile{theorem-library/rake-matrix-laws.scm}{theorem-library/rake-matrix-laws.scm}
\noindent{\small rake-matrix-laws.scm -- the eight matrix-algebra laws of `mat-ring-is-ring', PROVEN by ONE entrywise driver.}\par
\vnbentry{\texttt{matneg-\allowbreak{}entry}}{\vnbstmt{\begin{array}{@{}l@{}} \forall a, m, n, p.; \forall i \in \operatorname{interval}(1, m), j \in \operatorname{interval}(1, n). \\ \quad (\operatorname{is\text{-}ring}(a) \wedge (p \in \operatorname{mat}(m, n, \operatorname{carr}(a)))) \Rightarrow  \\ \quad \quad \operatorname{entry}(\operatorname{matneg}(a, p), i, j) \\ \quad \quad =\; \operatorname{neg}(a)(\operatorname{entry}(p, i, j)) \end{array}}}
\vnbentry{\texttt{matadd-\allowbreak{}comm}}{\vnbstmt{\begin{array}{@{}l@{}} \forall a, m, n, p, q. \\ \quad \operatorname{is\text{-}ring}(a) \wedge \\ \quad (p \in \operatorname{mat}(m, n, \operatorname{carr}(a))) \wedge \\ \quad (q \in \operatorname{mat}(m, n, \operatorname{carr}(a))) \Rightarrow \\ \quad \quad (\operatorname{matadd}(a, p, q) = \operatorname{matadd}(a, q, p)) \end{array}}}
\vnbentry{\texttt{matadd-\allowbreak{}assoc}}{\vnbstmt{\begin{array}{@{}l@{}} \forall a, m, n, p, q, r. \\ \quad \operatorname{is\text{-}ring}(a) \wedge \\ \quad (p \in \operatorname{mat}(m, n, \operatorname{carr}(a))) \wedge \\ \quad (q \in \operatorname{mat}(m, n, \operatorname{carr}(a))) \wedge \\ \quad (r \in \operatorname{mat}(m, n, \operatorname{carr}(a))) \Rightarrow \\ \quad \quad \operatorname{matadd}(a, \operatorname{matadd}(a, p, q), r) \\ \quad \quad =\; \operatorname{matadd}(a, p, \operatorname{matadd}(a, q, r)) \end{array}}}
\vnbentry{\texttt{matadd-\allowbreak{}zero-\allowbreak{}left}}{\vnbstmt{\begin{array}{@{}l@{}} \forall a, m, n.; \forall p \in \operatorname{mat}(m, n, \operatorname{carr}(a)). \\ \quad \operatorname{is\text{-}ring}(a) \Rightarrow  \\ \quad \quad (\operatorname{matadd}(a, \operatorname{zeromat}(a, m, n), p) = p) \end{array}}}
\vnbentry{\texttt{matadd-\allowbreak{}zero-\allowbreak{}right}}{\vnbstmt{\begin{array}{@{}l@{}} \forall a, m, n.; \forall p \in \operatorname{mat}(m, n, \operatorname{carr}(a)). \\ \quad \operatorname{is\text{-}ring}(a) \Rightarrow  \\ \quad \quad (\operatorname{matadd}(a, p, \operatorname{zeromat}(a, m, n)) = p) \end{array}}}
\vnbentry{\texttt{matadd-\allowbreak{}neg-\allowbreak{}left}}{\vnbstmt{\begin{array}{@{}l@{}} \forall a, m, n.; \forall p \in \operatorname{mat}(m, n, \operatorname{carr}(a)). \\ \quad \operatorname{is\text{-}ring}(a) \Rightarrow  \\ \quad \quad (\operatorname{matadd}(a, \operatorname{matneg}(a, p), p) = \operatorname{zeromat}(a, m, n)) \end{array}}}
\vnbentry{\texttt{matadd-\allowbreak{}neg-\allowbreak{}right}}{\vnbstmt{\begin{array}{@{}l@{}} \forall a, m, n.; \forall p \in \operatorname{mat}(m, n, \operatorname{carr}(a)). \\ \quad \operatorname{is\text{-}ring}(a) \Rightarrow  \\ \quad \quad (\operatorname{matadd}(a, p, \operatorname{matneg}(a, p)) = \operatorname{zeromat}(a, m, n)) \end{array}}}
\vnbentry{\texttt{matmul-\allowbreak{}left-\allowbreak{}dist-\allowbreak{}guarded}}{\vnbstmt{\begin{array}{@{}l@{}} \forall a, m, n, k, p, q, r. \\ \quad \operatorname{is\text{-}ring}(a) \wedge \\ \quad (p \in \operatorname{mat}(m, n, \operatorname{carr}(a))) \wedge \\ \quad (q \in \operatorname{mat}(n, k, \operatorname{carr}(a))) \wedge \\ \quad (r \in \operatorname{mat}(n, k, \operatorname{carr}(a))) \wedge \\ \quad ((n = 0) \Rightarrow ((m = 0) \vee (k = 0))) \Rightarrow \\ \quad \quad \operatorname{matmul}(a, p, \operatorname{matadd}(a, q, r)) \\ \quad \quad =\; \operatorname{matadd}(a, \operatorname{matmul}(a, p, q), \operatorname{matmul}(a, p, r)) \end{array}}}
\vnbentry{\texttt{matmul-\allowbreak{}right-\allowbreak{}dist-\allowbreak{}guarded}}{\vnbstmt{\begin{array}{@{}l@{}} \forall a, m, n, k, p, q, r. \\ \quad \operatorname{is\text{-}ring}(a) \wedge \\ \quad (p \in \operatorname{mat}(m, n, \operatorname{carr}(a))) \wedge \\ \quad (q \in \operatorname{mat}(m, n, \operatorname{carr}(a))) \wedge \\ \quad (r \in \operatorname{mat}(n, k, \operatorname{carr}(a))) \wedge \\ \quad ((n = 0) \Rightarrow ((m = 0) \vee (k = 0))) \Rightarrow \\ \quad \quad \operatorname{matmul}(a, \operatorname{matadd}(a, p, q), r) \\ \quad \quad =\; \operatorname{matadd}(a, \operatorname{matmul}(a, p, r), \operatorname{matmul}(a, q, r)) \end{array}}}

\vnbfile{theorem-library/mat-ring-proof.scm}{theorem-library/mat-ring-proof.scm}
\noindent{\small mat-ring-proof.scm -- the n-by-n matrices over a ring form a ring.}\par
\vnbentry{\texttt{mat-\allowbreak{}ring-\allowbreak{}is-\allowbreak{}ring}}{\vnbstmt{\begin{array}{@{}l@{}} \forall a.; \forall n \in \mathbb{N}. \\ \quad \operatorname{is\text{-}ring}(a) \Rightarrow  \\ \quad \quad \operatorname{is\text{-}ring}(\operatorname{mat\text{-}ring}(a, n)) \end{array}}}

\vnbfile{theorem-library/matunit-shift-proof.scm}{theorem-library/matunit-shift-proof.scm}
\noindent{\small matunit-shift-proof.scm -- Lemma 3.3 (algebraic-numbers.pdf ch.3): the column-shift formula for right-multiplication by a matrix unit.}\par
\vnbentry{\texttt{matunit-\allowbreak{}col-\allowbreak{}shift}}{\vnbstmt{\begin{array}{@{}l@{}} \forall a, m, n, p, k, l.; \forall i \in \operatorname{interval}(1, m), c \in \operatorname{interval}(1, n). \\ \quad \operatorname{is\text{-}ring}(a) \wedge \\ \quad (p \in \operatorname{mat}(m, n, \operatorname{carr}(a))) \wedge \\ \quad (k \in \operatorname{interval}(1, n)) \wedge \\ \quad (l \in \operatorname{interval}(1, n)) \Rightarrow \\ \quad \quad \operatorname{entry}(\operatorname{matmul}(a, p, \operatorname{matunit}(a, n, k, l)), i, c) \\ \quad \quad =\; \left(\text{if }(c = l)\text{ then }\operatorname{entry}(p, i, k)\text{ else }\operatorname{zero}(a)\right) \end{array}}}

\vnbfile{theorem-library/elem-actions-proof.scm}{theorem-library/elem-actions-proof.scm}
\noindent{\small elem-actions-proof.scm -- Prop 3.5 (algebraic-numbers.pdf ch.3): the action of the elementary column matrices on a matrix by right-multiplication.  Each is the matmul-entry expansion + a finsum collapse (single-support for H, and for the off-l column of G; two-support for G's l-column), closed by an}\par
\vnbentry{\texttt{elem-\allowbreak{}h-\allowbreak{}action}}{\vnbstmt{\begin{array}{@{}l@{}} \forall a, m, n, p, r, k.; \forall i \in \operatorname{interval}(1, m), c \in \operatorname{interval}(1, n). \\ \quad \operatorname{is\text{-}ring}(a) \wedge \\ \quad (p \in \operatorname{mat}(m, n, \operatorname{carr}(a))) \wedge \\ \quad (r \in \operatorname{carr}(a)) \wedge \\ \quad (k \in \operatorname{interval}(1, n)) \Rightarrow \\ \quad \quad \operatorname{entry}(\operatorname{matmul}(a, p, \operatorname{elem\text{-}h}(a, n, r, k)), i, c) \\ \quad \quad =\; \left(\text{if }(c = k)\text{ then }\operatorname{mul}(a)(\operatorname{entry}(p, i, k), r)\text{ else }\operatorname{entry}(p, i, c)\right) \end{array}}}
\vnbentry{\texttt{elem-\allowbreak{}g-\allowbreak{}action}}{\vnbstmt{\begin{array}{@{}l@{}} \forall a, m, n, p, r, k, l.; \forall i \in \operatorname{interval}(1, m), c \in \operatorname{interval}(1, n). \\ \quad \operatorname{is\text{-}ring}(a) \wedge \\ \quad (p \in \operatorname{mat}(m, n, \operatorname{carr}(a))) \wedge \\ \quad (r \in \operatorname{carr}(a)) \wedge \\ \quad (k \in \operatorname{interval}(1, n)) \wedge \\ \quad (l \in \operatorname{interval}(1, n)) \wedge \\ \quad \neg (k = l) \Rightarrow \\ \quad \quad \operatorname{entry}(\operatorname{matmul}(a, p, \operatorname{elem\text{-}g}(a, n, r, k, l)), i, c) \\ \quad \quad =\; \left(\text{if }(c = l)\text{ then }\operatorname{add}(a)(\operatorname{entry}(p, i, l), \operatorname{mul}(a)(\operatorname{entry}(p, i, k), r))\text{ else }\operatorname{entry}(p, i, c)\right) \end{array}}}
\vnbentry{\texttt{elem-\allowbreak{}f-\allowbreak{}action}}{\vnbstmt{\begin{array}{@{}l@{}} \forall a, m, n, p, k, l.; \forall i \in \operatorname{interval}(1, m), c \in \operatorname{interval}(1, n). \\ \quad \operatorname{is\text{-}ring}(a) \wedge \\ \quad (p \in \operatorname{mat}(m, n, \operatorname{carr}(a))) \wedge \\ \quad (k \in \operatorname{interval}(1, n)) \wedge \\ \quad (l \in \operatorname{interval}(1, n)) \wedge \\ \quad \neg (k = l) \Rightarrow \\ \quad \quad \operatorname{entry}(\operatorname{matmul}(a, p, \operatorname{elem\text{-}f}(a, n, k, l)), i, c) \\ \quad \quad =\; \left(\text{if }(c = k)\text{ then }\operatorname{entry}(p, i, l)\text{ else }\left(\text{if }(c = l)\text{ then }\operatorname{entry}(p, i, k)\text{ else }\operatorname{entry}(p, i, c)\right)\right) \end{array}}}

\vnbfile{theorem-library/elem-inverses-proof.scm}{theorem-library/elem-inverses-proof.scm}
\noindent{\small elem-inverses-proof.scm -- Cor 3.6 (algebraic-numbers.pdf ch.3): the elementary column matrices are invertible, with elementary inverses.  Each reduces, via matrix-entry-extensionality, to an entry identity whose LHS is the PROVEN Prop 3.5 action (with P = the would-be inverse) and whose RHS is entry-of-identmat;}\par
\vnbentry{\texttt{elem-\allowbreak{}h-\allowbreak{}inverse}}{\vnbstmt{\begin{array}{@{}l@{}} \forall a, n, r, s, k. \\ \quad \operatorname{is\text{-}ring}(a) \wedge \\ \quad (n \in \mathbb{N}) \wedge \\ \quad (r \in \operatorname{carr}(a)) \wedge \\ \quad (s \in \operatorname{carr}(a)) \wedge \\ \quad (\operatorname{mul}(a)(r, s) = \operatorname{one}(a)) \wedge \\ \quad (k \in \operatorname{interval}(1, n)) \Rightarrow \\ \quad \quad \operatorname{matmul}(a, \operatorname{elem\text{-}h}(a, n, r, k), \operatorname{elem\text{-}h}(a, n, s, k)) \\ \quad \quad =\; \operatorname{identmat}(a, n) \end{array}}}
\vnbentry{\texttt{elem-\allowbreak{}g-\allowbreak{}inverse}}{\vnbstmt{\begin{array}{@{}l@{}} \forall a, n, r, k, l. \\ \quad \operatorname{is\text{-}ring}(a) \wedge \\ \quad (n \in \mathbb{N}) \wedge \\ \quad (r \in \operatorname{carr}(a)) \wedge \\ \quad (k \in \operatorname{interval}(1, n)) \wedge \\ \quad (l \in \operatorname{interval}(1, n)) \wedge \\ \quad \neg (k = l) \Rightarrow \\ \quad \quad \operatorname{matmul}(a, \operatorname{elem\text{-}g}(a, n, r, k, l), \operatorname{elem\text{-}g}(a, n, \operatorname{neg}(a)(r), k, l)) \\ \quad \quad =\; \operatorname{identmat}(a, n) \end{array}}}
\vnbentry{\texttt{elem-\allowbreak{}f-\allowbreak{}inverse}}{\vnbstmt{\begin{array}{@{}l@{}} \forall a, n, k, l. \\ \quad \operatorname{is\text{-}ring}(a) \wedge \\ \quad (n \in \mathbb{N}) \wedge \\ \quad (k \in \operatorname{interval}(1, n)) \wedge \\ \quad (l \in \operatorname{interval}(1, n)) \wedge \\ \quad \neg (k = l) \Rightarrow \\ \quad \quad \operatorname{matmul}(a, \operatorname{elem\text{-}f}(a, n, k, l), \operatorname{elem\text{-}f}(a, n, l, k)) \\ \quad \quad =\; \operatorname{identmat}(a, n) \end{array}}}

\vnbfile{theorem-library/elem-row-actions-proof.scm}{theorem-library/elem-row-actions-proof.scm}
\noindent{\small elem-row-actions-proof.scm -- Prop 3.29 (row form, algebraic-numbers.pdf ch.3): the action of the elementary matrices by LEFT-multiplication (row operations). Each mirrors theorem-library/elem-actions-proof.scm with the elementary matrix as the LEFT factor of matmul-entry: (E.P)\_\{ic\} = FINSUM\_j E\_\{ij\}.P\_\{jc\}, so a}\par
\vnbentry{\texttt{elem-\allowbreak{}h-\allowbreak{}row-\allowbreak{}action}}{\vnbstmt{\begin{array}{@{}l@{}} \forall a, m, n, p, r, k.; \forall i \in \operatorname{interval}(1, m), c \in \operatorname{interval}(1, n). \\ \quad \operatorname{is\text{-}ring}(a) \wedge \\ \quad (p \in \operatorname{mat}(m, n, \operatorname{carr}(a))) \wedge \\ \quad (r \in \operatorname{carr}(a)) \wedge \\ \quad (k \in \operatorname{interval}(1, m)) \Rightarrow \\ \quad \quad \operatorname{entry}(\operatorname{matmul}(a, \operatorname{elem\text{-}h}(a, m, r, k), p), i, c) \\ \quad \quad =\; \left(\text{if }(i = k)\text{ then }\operatorname{mul}(a)(r, \operatorname{entry}(p, k, c))\text{ else }\operatorname{entry}(p, i, c)\right) \end{array}}}
\vnbentry{\texttt{elem-\allowbreak{}g-\allowbreak{}row-\allowbreak{}action}}{\vnbstmt{\begin{array}{@{}l@{}} \forall a, m, n, p, r, k, l.; \forall i \in \operatorname{interval}(1, m), c \in \operatorname{interval}(1, n). \\ \quad \operatorname{is\text{-}ring}(a) \wedge \\ \quad (p \in \operatorname{mat}(m, n, \operatorname{carr}(a))) \wedge \\ \quad (r \in \operatorname{carr}(a)) \wedge \\ \quad (k \in \operatorname{interval}(1, m)) \wedge \\ \quad (l \in \operatorname{interval}(1, m)) \wedge \\ \quad \neg (k = l) \Rightarrow \\ \quad \quad \operatorname{entry}(\operatorname{matmul}(a, \operatorname{elem\text{-}g}(a, m, r, k, l), p), i, c) \\ \quad \quad =\; \left(\text{if }(i = k)\text{ then }\operatorname{add}(a)(\operatorname{entry}(p, k, c), \operatorname{mul}(a)(r, \operatorname{entry}(p, l, c)))\text{ else }\operatorname{entry}(p, i, c)\right) \end{array}}}
\vnbentry{\texttt{elem-\allowbreak{}f-\allowbreak{}row-\allowbreak{}action}}{\vnbstmt{\begin{array}{@{}l@{}} \forall a, m, n, p, k, l.; \forall i \in \operatorname{interval}(1, m), c \in \operatorname{interval}(1, n). \\ \quad \operatorname{is\text{-}ring}(a) \wedge \\ \quad (p \in \operatorname{mat}(m, n, \operatorname{carr}(a))) \wedge \\ \quad (k \in \operatorname{interval}(1, m)) \wedge \\ \quad (l \in \operatorname{interval}(1, m)) \wedge \\ \quad \neg (k = l) \Rightarrow \\ \quad \quad \operatorname{entry}(\operatorname{matmul}(a, \operatorname{elem\text{-}f}(a, m, k, l), p), i, c) \\ \quad \quad =\; \left(\text{if }(i = k)\text{ then }\operatorname{entry}(p, l, c)\text{ else }\left(\text{if }(i = l)\text{ then }\operatorname{entry}(p, k, c)\text{ else }\operatorname{entry}(p, i, c)\right)\right) \end{array}}}

\vnbfile{theorem-library/det-rows.scm}{theorem-library/det-rows.scm}
\noindent{\small det-rows.scm -- the determinant is LINEAR in each row and ALTERNATING; the interchange of two rows, the elementary row operations, and the expansion along any row (batch 25-C, 2026-09-24).  Linear-algebra 3.18 (row clauses), the row form of 3.19, and 3.22 (row half).  All over a COMMUTATIVE ring.}\par
\vnbentry{\texttt{dtr-\allowbreak{}ring-\allowbreak{}lin1}}{\vnbstmt{\begin{array}{@{}l@{}} \forall r, s, a, b, c, d. \\ \quad \operatorname{is\text{-}commutative\text{-}ring}(r) \wedge \\ \quad (s \in \operatorname{carr}(r)) \wedge \\ \quad (a \in \operatorname{carr}(r)) \wedge \\ \quad (b \in \operatorname{carr}(r)) \wedge \\ \quad (c \in \operatorname{carr}(r)) \wedge \\ \quad (d \in \operatorname{carr}(r)) \Rightarrow \\ \quad \quad \operatorname{mul}(r)(s, \operatorname{mul}(r)(\operatorname{add}(r)(\operatorname{mul}(r)(c, a), b), d)) \\ \quad \quad =\; \operatorname{add}(r)(\operatorname{mul}(r)(c, \operatorname{mul}(r)(s, \operatorname{mul}(r)(a, d))), \operatorname{mul}(r)(s, \operatorname{mul}(r)(b, d))) \end{array}}}
\vnbentry{\texttt{dtr-\allowbreak{}ring-\allowbreak{}lin2}}{\vnbstmt{\begin{array}{@{}l@{}} \forall r, s, a, c, x, y. \\ \quad \operatorname{is\text{-}commutative\text{-}ring}(r) \wedge \\ \quad (s \in \operatorname{carr}(r)) \wedge \\ \quad (a \in \operatorname{carr}(r)) \wedge \\ \quad (c \in \operatorname{carr}(r)) \wedge \\ \quad (x \in \operatorname{carr}(r)) \wedge \\ \quad (y \in \operatorname{carr}(r)) \Rightarrow \\ \quad \quad \operatorname{mul}(r)(s, \operatorname{mul}(r)(a, \operatorname{add}(r)(\operatorname{mul}(r)(c, x), y))) \\ \quad \quad =\; \operatorname{add}(r)(\operatorname{mul}(r)(c, \operatorname{mul}(r)(s, \operatorname{mul}(r)(a, x))), \operatorname{mul}(r)(s, \operatorname{mul}(r)(a, y))) \end{array}}}
\vnbentry{\texttt{dtr-\allowbreak{}finsum-\allowbreak{}lin}}{\vnbstmt{\begin{array}{@{}l@{}} \forall r, s, c, f, g, h. \\ \quad \operatorname{is\text{-}ring}(r) \wedge \\ \quad (s \in \mathbf{Set}) \wedge \\ \quad (|s| \in \mathbb{N}) \wedge \\ \quad (c \in \operatorname{carr}(r)) \wedge \\ \quad \forall z \in s.\; (f(z) \in \operatorname{carr}(r)) \wedge \\ \quad \forall z \in s.\; (g(z) \in \operatorname{carr}(r)) \wedge \\ \quad \forall z \in s. \\ \quad \quad (h(z) = \operatorname{add}(r)(\operatorname{mul}(r)(c, f(z)), g(z))) \Rightarrow \\ \quad \quad \operatorname{finsum}(\operatorname{ring\text{-}additive\text{-}ag}(r), h, s) \\ \quad \quad =\; \operatorname{add}(r)(\operatorname{mul}(r)(c, \operatorname{finsum}(\operatorname{ring\text{-}additive\text{-}ag}(r), f, s)), \\ \quad \quad \quad \operatorname{finsum}(\operatorname{ring\text{-}additive\text{-}ag}(r), g, s)) \end{array}}}
\vnbentry{\texttt{dtr-\allowbreak{}succ-\allowbreak{}in}}{\vnbstmt{\begin{array}{@{}l@{}} \forall n, u. \\ \quad ((n \in \mathbb{N}) \wedge (u \in \operatorname{interval}(1, n))) \Rightarrow  \\ \quad \quad ((u + 1) \in \operatorname{interval}(1, (n + 1))) \end{array}}}
\vnbentry{\texttt{dtr-\allowbreak{}skip-\allowbreak{}in}}{\vnbstmt{\begin{array}{@{}l@{}} \forall n, j, v. \\ \quad ((n \in \mathbb{N}) \wedge (v \in \operatorname{interval}(1, n))) \Rightarrow  \\ \quad \quad \left(\text{if }(v < j)\text{ then }v\text{ else }(v + 1)\right) \\ \quad \quad \in\; \operatorname{interval}(1, (n + 1)) \end{array}}}
\vnbentry{\texttt{dtr-\allowbreak{}minor1-\allowbreak{}entry}}{\vnbstmt{\begin{array}{@{}l@{}} \forall r, n, s, j, u, v. \\ \quad \operatorname{is\text{-}ring}(r) \wedge \\ \quad (n \in \mathbb{N}) \wedge \\ \quad (s \in \operatorname{mat}((n + 1), (n + 1), \operatorname{carr}(r))) \wedge \\ \quad (u \in \operatorname{interval}(1, n)) \wedge \\ \quad (v \in \operatorname{interval}(1, n)) \Rightarrow \\ \quad \quad \operatorname{entry}(\operatorname{minor}(s, 1, j, n), u, v) \\ \quad \quad =\; \operatorname{entry}(s, (u + 1), \left(\text{if }(v < j)\text{ then }v\text{ else }(v + 1)\right)) \end{array}}}
\vnbentry{\texttt{dtr-\allowbreak{}sign-\allowbreak{}type}}{\vnbstmt{\begin{array}{@{}l@{}} \forall r, j. \\ \quad (\operatorname{is\text{-}ring}(r) \wedge (j \in \mathbb{N})) \Rightarrow  \\ \quad \quad \operatorname{mpow}(\operatorname{ring\text{-}multiplicative\text{-}monoid}(r), \operatorname{neg}(r)(\operatorname{one}(r)), j) \\ \quad \quad \in\; \operatorname{carr}(r) \end{array}}}
\vnbentry{\texttt{dtr-\allowbreak{}minor1-\allowbreak{}congr}}{\vnbstmt{\begin{array}{@{}l@{}} \forall r, n, x, y, p, j. \\ \quad \operatorname{is\text{-}ring}(r) \wedge \\ \quad (n \in \mathbb{N}) \wedge \\ \quad (x \in \operatorname{mat}((n + 1), (n + 1), \operatorname{carr}(r))) \wedge \\ \quad (y \in \operatorname{mat}((n + 1), (n + 1), \operatorname{carr}(r))) \wedge \\ \quad (j \in \mathbb{N}) \wedge \\ \quad (p = 1) \wedge \\ \quad \forall i, k \in \operatorname{interval}(1, (n + 1)). \\ \quad \quad \neg (i = p) \Rightarrow  \\ \quad \quad \quad (\operatorname{entry}(x, i, k) = \operatorname{entry}(y, i, k)) \Rightarrow \\ \quad \quad (\operatorname{minor}(x, 1, j, n) = \operatorname{minor}(y, 1, j, n)) \end{array}}}
\vnbentry{\texttt{det-\allowbreak{}row-\allowbreak{}linear-\allowbreak{}ind}}{\vnbstmt{\begin{array}{@{}l@{}} \forall n \in \mathbb{N}, r, a, b, c, p, \mathit{dsc}. \\ \quad \operatorname{is\text{-}commutative\text{-}ring}(r) \wedge \\ \quad (a \in \operatorname{mat}(n, n, \operatorname{carr}(r))) \wedge \\ \quad (b \in \operatorname{mat}(n, n, \operatorname{carr}(r))) \wedge \\ \quad (c \in \operatorname{mat}(n, n, \operatorname{carr}(r))) \wedge \\ \quad (p \in \operatorname{interval}(1, n)) \wedge \\ \quad (\mathit{dsc} \in \operatorname{carr}(r)) \wedge \\ \quad \forall i, k \in \operatorname{interval}(1, n). \\ \quad \quad \neg (i = p) \Rightarrow  \\ \quad \quad \quad (\operatorname{entry}(b, i, k) = \operatorname{entry}(a, i, k)) \wedge \\ \quad \forall i, k \in \operatorname{interval}(1, n). \\ \quad \quad \neg (i = p) \Rightarrow  \\ \quad \quad \quad (\operatorname{entry}(c, i, k) = \operatorname{entry}(a, i, k)) \wedge \\ \quad \forall k \in \operatorname{interval}(1, n). \\ \quad \quad \operatorname{entry}(c, p, k) \\ \quad \quad =\; \operatorname{add}(r)(\operatorname{mul}(r)(\mathit{dsc}, \operatorname{entry}(a, p, k)), \operatorname{entry}(b, p, k)) \Rightarrow \\ \quad \quad \operatorname{det}(r, n, c) \\ \quad \quad =\; \operatorname{add}(r)(\operatorname{mul}(r)(\mathit{dsc}, \operatorname{det}(r, n, a)), \operatorname{det}(r, n, b)) \end{array}}}
\vnbentry{\texttt{det-\allowbreak{}row-\allowbreak{}linear}}{\vnbstmt{\begin{array}{@{}l@{}} \forall r, n, a, b, c, p, d. \\ \quad \operatorname{is\text{-}commutative\text{-}ring}(r) \wedge \\ \quad (n \in \mathbb{N}) \wedge \\ \quad (a \in \operatorname{mat}(n, n, \operatorname{carr}(r))) \wedge \\ \quad (b \in \operatorname{mat}(n, n, \operatorname{carr}(r))) \wedge \\ \quad (c \in \operatorname{mat}(n, n, \operatorname{carr}(r))) \wedge \\ \quad (p \in \operatorname{interval}(1, n)) \wedge \\ \quad (d \in \operatorname{carr}(r)) \wedge \\ \quad \forall i, k \in \operatorname{interval}(1, n). \\ \quad \quad \neg (i = p) \Rightarrow  \\ \quad \quad \quad (\operatorname{entry}(b, i, k) = \operatorname{entry}(a, i, k)) \wedge \\ \quad \forall i, k \in \operatorname{interval}(1, n). \\ \quad \quad \neg (i = p) \Rightarrow  \\ \quad \quad \quad (\operatorname{entry}(c, i, k) = \operatorname{entry}(a, i, k)) \wedge \\ \quad \forall k \in \operatorname{interval}(1, n). \\ \quad \quad \operatorname{entry}(c, p, k) \\ \quad \quad =\; \operatorname{add}(r)(\operatorname{mul}(r)(d, \operatorname{entry}(a, p, k)), \operatorname{entry}(b, p, k)) \Rightarrow \\ \quad \quad \operatorname{det}(r, n, c) \\ \quad \quad =\; \operatorname{add}(r)(\operatorname{mul}(r)(d, \operatorname{det}(r, n, a)), \operatorname{det}(r, n, b)) \end{array}}}
\vnbentry{\texttt{dtr-\allowbreak{}sign-\allowbreak{}succ}}{\vnbstmt{\begin{array}{@{}l@{}} \forall r, x. \\ \quad (\operatorname{is\text{-}ring}(r) \wedge (x \in \mathbb{N})) \Rightarrow  \\ \quad \quad \operatorname{mpow}(\operatorname{ring\text{-}multiplicative\text{-}monoid}(r), \operatorname{neg}(r)(\operatorname{one}(r)), (x + 1)) \\ \quad \quad =\; \operatorname{mul}(r)(\operatorname{neg}(r)(\operatorname{one}(r)), \\ \quad \quad \quad \operatorname{mpow}(\operatorname{ring\text{-}multiplicative\text{-}monoid}(r), \operatorname{neg}(r)(\operatorname{one}(r)), x)) \end{array}}}
\vnbentry{\texttt{dtr-\allowbreak{}minor-\allowbreak{}minor}}{\vnbstmt{\begin{array}{@{}l@{}} \forall r, n, a, j, k. \\ \quad \operatorname{is\text{-}ring}(r) \wedge \\ \quad (n \in \mathbb{N}) \wedge \\ \quad (a \in \operatorname{mat}(((n + 1) + 1), ((n + 1) + 1), \operatorname{carr}(r))) \wedge \\ \quad (j \in \mathbb{N}) \wedge \\ \quad (k \in \mathbb{N}) \wedge \\ \quad (1 \leq j) \wedge \\ \quad (j \leq k) \wedge \\ \quad (k \leq (n + 1)) \Rightarrow \\ \quad \quad \operatorname{minor}(\operatorname{minor}(a, 1, j, (n + 1)), 1, k, n) \\ \quad \quad =\; \operatorname{minor}(\operatorname{minor}(a, 1, (k + 1), (n + 1)), 1, j, n) \end{array}}}
\vnbentry{\texttt{dtr-\allowbreak{}ring-\allowbreak{}zero-\allowbreak{}add}}{\vnbstmt{\begin{array}{@{}l@{}} \forall r, y, z. \\ \quad \operatorname{is\text{-}commutative\text{-}ring}(r) \wedge \\ \quad (y \in \operatorname{carr}(r)) \wedge \\ \quad (z \in \operatorname{carr}(r)) \Rightarrow \\ \quad \quad \operatorname{add}(r)(\operatorname{add}(r)(\operatorname{zero}(r), y), z) \\ \quad \quad =\; \operatorname{add}(r)(y, z) \end{array}}}
\vnbentry{\texttt{dtr-\allowbreak{}dsum-\allowbreak{}anti}}{\vnbstmt{\begin{array}{@{}l@{}} \forall n \in \mathbb{N}, r, f. \\ \quad \operatorname{is\text{-}commutative\text{-}ring}(r) \wedge \\ \quad \forall j \in \operatorname{interval}(1, (n + 1)), k \in \operatorname{interval}(1, n). \\ \quad \quad (f(j, k) \in \operatorname{carr}(r)) \wedge \\ \quad \forall j, k \in \mathbb{N}. \\ \quad \quad ((1 \leq j) \wedge ((j \leq k) \wedge (k \leq n))) \Rightarrow  \\ \quad \quad \quad (f((k + 1), j) = \operatorname{neg}(r)(f(j, k))) \Rightarrow \\ \quad \quad \operatorname{finsum}(\operatorname{ring\text{-}additive\text{-}ag}(r), \\ \quad \quad \quad (\lambda a \in \operatorname{interval}(1, (n + 1)).\; \\ \quad \quad \quad \quad \operatorname{finsum}(\operatorname{ring\text{-}additive\text{-}ag}(r), (\lambda b \in \operatorname{interval}(1, n).\; f(a, b)), \\ \quad \quad \quad \quad \quad \operatorname{interval}(1, n))), \\ \quad \quad \quad \operatorname{interval}(1, (n + 1))) \\ \quad \quad =\; \operatorname{zero}(r) \end{array}}}
\vnbentry{\texttt{dtr-\allowbreak{}ring-\allowbreak{}anti}}{\vnbstmt{\begin{array}{@{}l@{}} \forall r, s, t, a, b, d. \\ \quad \operatorname{is\text{-}commutative\text{-}ring}(r) \wedge \\ \quad (s \in \operatorname{carr}(r)) \wedge \\ \quad (t \in \operatorname{carr}(r)) \wedge \\ \quad (a \in \operatorname{carr}(r)) \wedge \\ \quad (b \in \operatorname{carr}(r)) \wedge \\ \quad (d \in \operatorname{carr}(r)) \Rightarrow \\ \quad \quad \operatorname{mul}(r)(\operatorname{mul}(r)(\operatorname{neg}(r)(\operatorname{one}(r)), t), \operatorname{mul}(r)(b, \operatorname{mul}(r)(s, \operatorname{mul}(r)(a, d)))) \\ \quad \quad =\; \operatorname{neg}(r)(\operatorname{mul}(r)(s, \operatorname{mul}(r)(a, \operatorname{mul}(r)(t, \operatorname{mul}(r)(b, d))))) \end{array}}}
\vnbentry{\texttt{dtr-\allowbreak{}det-\allowbreak{}rows12}}{\vnbstmt{\begin{array}{@{}l@{}} \forall n, r, a. \\ \quad (n \in \mathbb{N}) \wedge \\ \quad \operatorname{is\text{-}commutative\text{-}ring}(r) \wedge \\ \quad (a \in \operatorname{mat}(((n + 1) + 1), ((n + 1) + 1), \operatorname{carr}(r))) \wedge \\ \quad \forall k \in \operatorname{interval}(1, ((n + 1) + 1)). \\ \quad \quad (\operatorname{entry}(a, 1, k) = \operatorname{entry}(a, 2, k)) \Rightarrow \\ \quad \quad (\operatorname{det}(r, ((n + 1) + 1), a) = \operatorname{zero}(r)) \end{array}}}
\vnbentry{\texttt{dtr-\allowbreak{}det-\allowbreak{}rows12b}}{\vnbstmt{\begin{array}{@{}l@{}} \forall m, n, r, a. \\ \quad (n \in \mathbb{N}) \wedge \\ \quad (m = ((n + 1) + 1)) \wedge \\ \quad \operatorname{is\text{-}commutative\text{-}ring}(r) \wedge \\ \quad (a \in \operatorname{mat}(m, m, \operatorname{carr}(r))) \wedge \\ \quad \forall k \in \operatorname{interval}(1, m). \\ \quad \quad (\operatorname{entry}(a, 1, k) = \operatorname{entry}(a, 2, k)) \Rightarrow \\ \quad \quad (\operatorname{det}(r, m, a) = \operatorname{zero}(r)) \end{array}}}
\vnbentry{\texttt{dtr-\allowbreak{}alt-\allowbreak{}low}}{\vnbstmt{\begin{array}{@{}l@{}} \forall n \in \mathbb{N}, r, a, p, q. \\ \quad \forall r, a, p, q. \\ \quad \quad \operatorname{is\text{-}commutative\text{-}ring}(r) \wedge \\ \quad \quad (a \in \operatorname{mat}(n, n, \operatorname{carr}(r))) \wedge \\ \quad \quad (p \in \operatorname{interval}(1, n)) \wedge \\ \quad \quad (q \in \operatorname{interval}(1, n)) \wedge \\ \quad \quad \neg (p = q) \wedge \\ \quad \quad \forall k \in \operatorname{interval}(1, n). \\ \quad \quad \quad (\operatorname{entry}(a, p, k) = \operatorname{entry}(a, q, k)) \Rightarrow \\ \quad \quad \quad (\operatorname{det}(r, n, a) = \operatorname{zero}(r)) \wedge \\ \quad \operatorname{is\text{-}commutative\text{-}ring}(r) \wedge \\ \quad (a \in \operatorname{mat}((n + 1), (n + 1), \operatorname{carr}(r))) \wedge \\ \quad (p \in \operatorname{interval}(1, (n + 1))) \wedge \\ \quad (q \in \operatorname{interval}(1, (n + 1))) \wedge \\ \quad \neg (p = q) \wedge \\ \quad \neg (p = 1) \wedge \\ \quad \neg (q = 1) \wedge \\ \quad \forall k \in \operatorname{interval}(1, (n + 1)). \\ \quad \quad (\operatorname{entry}(a, p, k) = \operatorname{entry}(a, q, k)) \Rightarrow \\ \quad \quad (\operatorname{det}(r, (n + 1), a) = \operatorname{zero}(r)) \end{array}}}
\vnbentry{\texttt{dtr-\allowbreak{}ring-\allowbreak{}final}}{\vnbstmt{\begin{array}{@{}l@{}} \forall r, t, u_{3}, u_{4}, u_{5}, u_{6}, u_{7}, u_{8}. \\ \quad \operatorname{is\text{-}commutative\text{-}ring}(r) \wedge \\ \quad (t \in \operatorname{carr}(r)) \wedge \\ \quad (u_{3} \in \operatorname{carr}(r)) \wedge \\ \quad (u_{4} \in \operatorname{carr}(r)) \wedge \\ \quad (u_{5} \in \operatorname{carr}(r)) \wedge \\ \quad (u_{6} \in \operatorname{carr}(r)) \wedge \\ \quad (u_{7} \in \operatorname{carr}(r)) \wedge \\ \quad (u_{8} \in \operatorname{carr}(r)) \wedge \\ \quad (u_{3} = \operatorname{add}(r)(\operatorname{mul}(r)(\operatorname{one}(r), u_{4}), u_{5})) \wedge \\ \quad (u_{4} = \operatorname{add}(r)(\operatorname{mul}(r)(\operatorname{one}(r), u_{6}), t)) \wedge \\ \quad (u_{5} = \operatorname{add}(r)(\operatorname{mul}(r)(\operatorname{one}(r), u_{7}), u_{8})) \wedge \\ \quad (u_{3} = \operatorname{zero}(r)) \wedge \\ \quad (u_{6} = \operatorname{zero}(r)) \wedge \\ \quad (u_{7} = \operatorname{zero}(r)) \wedge \\ \quad (u_{8} = \operatorname{zero}(r)) \Rightarrow \\ \quad \quad (t = \operatorname{zero}(r)) \end{array}}}
\vnbentry{\texttt{dtr-\allowbreak{}alt-\allowbreak{}row1}}{\vnbstmt{\begin{array}{@{}l@{}} \forall n \in \mathbb{N}, r, a, q. \\ \quad \forall r, a, p, q. \\ \quad \quad \operatorname{is\text{-}commutative\text{-}ring}(r) \wedge \\ \quad \quad (a \in \operatorname{mat}(n, n, \operatorname{carr}(r))) \wedge \\ \quad \quad (p \in \operatorname{interval}(1, n)) \wedge \\ \quad \quad (q \in \operatorname{interval}(1, n)) \wedge \\ \quad \quad \neg (p = q) \wedge \\ \quad \quad \forall k \in \operatorname{interval}(1, n). \\ \quad \quad \quad (\operatorname{entry}(a, p, k) = \operatorname{entry}(a, q, k)) \Rightarrow \\ \quad \quad \quad (\operatorname{det}(r, n, a) = \operatorname{zero}(r)) \wedge \\ \quad \operatorname{is\text{-}commutative\text{-}ring}(r) \wedge \\ \quad (a \in \operatorname{mat}((n + 1), (n + 1), \operatorname{carr}(r))) \wedge \\ \quad (q \in \operatorname{interval}(1, (n + 1))) \wedge \\ \quad \neg (q = 1) \wedge \\ \quad \forall k \in \operatorname{interval}(1, (n + 1)). \\ \quad \quad (\operatorname{entry}(a, 1, k) = \operatorname{entry}(a, q, k)) \Rightarrow \\ \quad \quad (\operatorname{det}(r, (n + 1), a) = \operatorname{zero}(r)) \end{array}}}
\vnbentry{\texttt{det-\allowbreak{}alternating-\allowbreak{}ind}}{\vnbstmt{\begin{array}{@{}l@{}} \forall n \in \mathbb{N}, r, a, p, q. \\ \quad \operatorname{is\text{-}commutative\text{-}ring}(r) \wedge \\ \quad (a \in \operatorname{mat}(n, n, \operatorname{carr}(r))) \wedge \\ \quad (p \in \operatorname{interval}(1, n)) \wedge \\ \quad (q \in \operatorname{interval}(1, n)) \wedge \\ \quad \neg (p = q) \wedge \\ \quad \forall k \in \operatorname{interval}(1, n). \\ \quad \quad (\operatorname{entry}(a, p, k) = \operatorname{entry}(a, q, k)) \Rightarrow \\ \quad \quad (\operatorname{det}(r, n, a) = \operatorname{zero}(r)) \end{array}}}
\vnbentry{\texttt{det-\allowbreak{}alternating-\allowbreak{}rows}}{\vnbstmt{\begin{array}{@{}l@{}} \forall r, n, a, p, q. \\ \quad \operatorname{is\text{-}commutative\text{-}ring}(r) \wedge \\ \quad (n \in \mathbb{N}) \wedge \\ \quad (a \in \operatorname{mat}(n, n, \operatorname{carr}(r))) \wedge \\ \quad (p \in \operatorname{interval}(1, n)) \wedge \\ \quad (q \in \operatorname{interval}(1, n)) \wedge \\ \quad \neg (p = q) \wedge \\ \quad \forall k \in \operatorname{interval}(1, n). \\ \quad \quad (\operatorname{entry}(a, p, k) = \operatorname{entry}(a, q, k)) \Rightarrow \\ \quad \quad (\operatorname{det}(r, n, a) = \operatorname{zero}(r)) \end{array}}}
\vnbentry{\texttt{dtr-\allowbreak{}ring-\allowbreak{}swap}}{\vnbstmt{\begin{array}{@{}l@{}} \forall r, a, b, u_{1}, u_{2}, u_{3}, u_{4}, u_{5}. \\ \quad \operatorname{is\text{-}commutative\text{-}ring}(r) \wedge \\ \quad (a \in \operatorname{carr}(r)) \wedge \\ \quad (b \in \operatorname{carr}(r)) \wedge \\ \quad (u_{1} \in \operatorname{carr}(r)) \wedge \\ \quad (u_{2} \in \operatorname{carr}(r)) \wedge \\ \quad (u_{3} \in \operatorname{carr}(r)) \wedge \\ \quad (u_{4} \in \operatorname{carr}(r)) \wedge \\ \quad (u_{5} \in \operatorname{carr}(r)) \wedge \\ \quad (u_{1} = \operatorname{add}(r)(\operatorname{mul}(r)(\operatorname{one}(r), u_{2}), u_{3})) \wedge \\ \quad (u_{2} = \operatorname{add}(r)(\operatorname{mul}(r)(\operatorname{one}(r), u_{4}), a)) \wedge \\ \quad (u_{3} = \operatorname{add}(r)(\operatorname{mul}(r)(\operatorname{one}(r), b), u_{5})) \wedge \\ \quad (u_{1} = \operatorname{zero}(r)) \wedge \\ \quad (u_{4} = \operatorname{zero}(r)) \wedge \\ \quad (u_{5} = \operatorname{zero}(r)) \Rightarrow \\ \quad \quad (b = \operatorname{neg}(r)(a)) \end{array}}}
\vnbentry{\texttt{dtr-\allowbreak{}swap-\allowbreak{}s}}{\vnbstmt{\begin{array}{@{}l@{}} \forall n, r, a, b, p, q. \\ \quad (n \in \mathbb{N}) \wedge \\ \quad \operatorname{is\text{-}commutative\text{-}ring}(r) \wedge \\ \quad (a \in \operatorname{mat}((n + 1), (n + 1), \operatorname{carr}(r))) \wedge \\ \quad (b \in \operatorname{mat}((n + 1), (n + 1), \operatorname{carr}(r))) \wedge \\ \quad (p \in \operatorname{interval}(1, (n + 1))) \wedge \\ \quad (q \in \operatorname{interval}(1, (n + 1))) \wedge \\ \quad \neg (p = q) \wedge \\ \quad \forall i, k \in \operatorname{interval}(1, (n + 1)). \\ \quad \quad (\neg (i = p) \wedge \neg (i = q)) \Rightarrow  \\ \quad \quad \quad (\operatorname{entry}(b, i, k) = \operatorname{entry}(a, i, k)) \wedge \\ \quad \forall k \in \operatorname{interval}(1, (n + 1)). \\ \quad \quad (\operatorname{entry}(b, p, k) = \operatorname{entry}(a, q, k)) \wedge \\ \quad \forall k \in \operatorname{interval}(1, (n + 1)). \\ \quad \quad (\operatorname{entry}(b, q, k) = \operatorname{entry}(a, p, k)) \Rightarrow \\ \quad \quad (\operatorname{det}(r, (n + 1), b) = \operatorname{neg}(r)(\operatorname{det}(r, (n + 1), a))) \end{array}}}
\vnbentry{\texttt{det-\allowbreak{}swap-\allowbreak{}rows}}{\vnbstmt{\begin{array}{@{}l@{}} \forall r, n, a, b, p, q. \\ \quad \operatorname{is\text{-}commutative\text{-}ring}(r) \wedge \\ \quad (n \in \mathbb{N}) \wedge \\ \quad (a \in \operatorname{mat}(n, n, \operatorname{carr}(r))) \wedge \\ \quad (b \in \operatorname{mat}(n, n, \operatorname{carr}(r))) \wedge \\ \quad (p \in \operatorname{interval}(1, n)) \wedge \\ \quad (q \in \operatorname{interval}(1, n)) \wedge \\ \quad \neg (p = q) \wedge \\ \quad \forall i, k \in \operatorname{interval}(1, n). \\ \quad \quad (\neg (i = p) \wedge \neg (i = q)) \Rightarrow  \\ \quad \quad \quad (\operatorname{entry}(b, i, k) = \operatorname{entry}(a, i, k)) \wedge \\ \quad \forall k \in \operatorname{interval}(1, n). \\ \quad \quad (\operatorname{entry}(b, p, k) = \operatorname{entry}(a, q, k)) \wedge \\ \quad \forall k \in \operatorname{interval}(1, n). \\ \quad \quad (\operatorname{entry}(b, q, k) = \operatorname{entry}(a, p, k)) \Rightarrow \\ \quad \quad (\operatorname{det}(r, n, b) = \operatorname{neg}(r)(\operatorname{det}(r, n, a))) \end{array}}}
\vnbentry{\texttt{dtr-\allowbreak{}ring-\allowbreak{}negself}}{\vnbstmt{\begin{array}{@{}l@{}} \forall r.; \forall u \in \operatorname{carr}(r). \\ \quad \operatorname{is\text{-}commutative\text{-}ring}(r) \wedge \\ \quad u \\ \quad =\; \operatorname{add}(r)(\operatorname{mul}(r)(\operatorname{neg}(r)(\operatorname{one}(r)), u), u) \Rightarrow \\ \quad \quad (u = \operatorname{zero}(r)) \end{array}}}
\vnbentry{\texttt{dtr-\allowbreak{}zero-\allowbreak{}row-\allowbreak{}s}}{\vnbstmt{\begin{array}{@{}l@{}} \forall n, r, a, k. \\ \quad (n \in \mathbb{N}) \wedge \\ \quad \operatorname{is\text{-}commutative\text{-}ring}(r) \wedge \\ \quad (a \in \operatorname{mat}((n + 1), (n + 1), \operatorname{carr}(r))) \wedge \\ \quad (k \in \operatorname{interval}(1, (n + 1))) \wedge \\ \quad \forall \mathit{drk} \in \operatorname{interval}(1, (n + 1)). \\ \quad \quad (\operatorname{entry}(a, k, \mathit{drk}) = \operatorname{zero}(r)) \Rightarrow \\ \quad \quad (\operatorname{det}(r, (n + 1), a) = \operatorname{zero}(r)) \end{array}}}
\vnbentry{\texttt{dtr-\allowbreak{}elem-\allowbreak{}h-\allowbreak{}s}}{\vnbstmt{\begin{array}{@{}l@{}} \forall n, r, a, \mathit{dsr}, k. \\ \quad (n \in \mathbb{N}) \wedge \\ \quad \operatorname{is\text{-}commutative\text{-}ring}(r) \wedge \\ \quad (a \in \operatorname{mat}((n + 1), (n + 1), \operatorname{carr}(r))) \wedge \\ \quad (\mathit{dsr} \in \operatorname{carr}(r)) \wedge \\ \quad (k \in \operatorname{interval}(1, (n + 1))) \Rightarrow \\ \quad \quad \operatorname{det}(r, (n + 1), \operatorname{matmul}(r, \operatorname{elem\text{-}h}(r, (n + 1), \mathit{dsr}, k), a)) \\ \quad \quad =\; \operatorname{mul}(r)(\mathit{dsr}, \operatorname{det}(r, (n + 1), a)) \end{array}}}
\vnbentry{\texttt{dtr-\allowbreak{}elem-\allowbreak{}g-\allowbreak{}s}}{\vnbstmt{\begin{array}{@{}l@{}} \forall n, r, a, \mathit{dsr}, k, l. \\ \quad (n \in \mathbb{N}) \wedge \\ \quad \operatorname{is\text{-}commutative\text{-}ring}(r) \wedge \\ \quad (a \in \operatorname{mat}((n + 1), (n + 1), \operatorname{carr}(r))) \wedge \\ \quad (\mathit{dsr} \in \operatorname{carr}(r)) \wedge \\ \quad (k \in \operatorname{interval}(1, (n + 1))) \wedge \\ \quad (l \in \operatorname{interval}(1, (n + 1))) \wedge \\ \quad \neg (k = l) \Rightarrow \\ \quad \quad \operatorname{det}(r, (n + 1), \operatorname{matmul}(r, \operatorname{elem\text{-}g}(r, (n + 1), \mathit{dsr}, k, l), a)) \\ \quad \quad =\; \operatorname{det}(r, (n + 1), a) \end{array}}}
\vnbentry{\texttt{dtr-\allowbreak{}elem-\allowbreak{}f-\allowbreak{}s}}{\vnbstmt{\begin{array}{@{}l@{}} \forall n, r, a, k, l. \\ \quad (n \in \mathbb{N}) \wedge \\ \quad \operatorname{is\text{-}commutative\text{-}ring}(r) \wedge \\ \quad (a \in \operatorname{mat}((n + 1), (n + 1), \operatorname{carr}(r))) \wedge \\ \quad (k \in \operatorname{interval}(1, (n + 1))) \wedge \\ \quad (l \in \operatorname{interval}(1, (n + 1))) \wedge \\ \quad \neg (k = l) \Rightarrow \\ \quad \quad \operatorname{det}(r, (n + 1), \operatorname{matmul}(r, \operatorname{elem\text{-}f}(r, (n + 1), k, l), a)) \\ \quad \quad =\; \operatorname{neg}(r)(\operatorname{det}(r, (n + 1), a)) \end{array}}}
\vnbentry{\texttt{det-\allowbreak{}elem-\allowbreak{}h-\allowbreak{}row}}{\vnbstmt{\begin{array}{@{}l@{}} \forall r, n, a, c, k. \\ \quad \operatorname{is\text{-}commutative\text{-}ring}(r) \wedge \\ \quad (n \in \mathbb{N}) \wedge \\ \quad (a \in \operatorname{mat}(n, n, \operatorname{carr}(r))) \wedge \\ \quad (c \in \operatorname{carr}(r)) \wedge \\ \quad (k \in \operatorname{interval}(1, n)) \Rightarrow \\ \quad \quad \operatorname{det}(r, n, \operatorname{matmul}(r, \operatorname{elem\text{-}h}(r, n, c, k), a)) \\ \quad \quad =\; \operatorname{mul}(r)(c, \operatorname{det}(r, n, a)) \end{array}}}
\vnbentry{\texttt{det-\allowbreak{}elem-\allowbreak{}g-\allowbreak{}row}}{\vnbstmt{\begin{array}{@{}l@{}} \forall r, n, a, c, k, l. \\ \quad \operatorname{is\text{-}commutative\text{-}ring}(r) \wedge \\ \quad (n \in \mathbb{N}) \wedge \\ \quad (a \in \operatorname{mat}(n, n, \operatorname{carr}(r))) \wedge \\ \quad (c \in \operatorname{carr}(r)) \wedge \\ \quad (k \in \operatorname{interval}(1, n)) \wedge \\ \quad (l \in \operatorname{interval}(1, n)) \wedge \\ \quad \neg (k = l) \Rightarrow \\ \quad \quad \operatorname{det}(r, n, \operatorname{matmul}(r, \operatorname{elem\text{-}g}(r, n, c, k, l), a)) \\ \quad \quad =\; \operatorname{det}(r, n, a) \end{array}}}
\vnbentry{\texttt{det-\allowbreak{}elem-\allowbreak{}f-\allowbreak{}row}}{\vnbstmt{\begin{array}{@{}l@{}} \forall r, n, a, k, l. \\ \quad \operatorname{is\text{-}commutative\text{-}ring}(r) \wedge \\ \quad (n \in \mathbb{N}) \wedge \\ \quad (a \in \operatorname{mat}(n, n, \operatorname{carr}(r))) \wedge \\ \quad (k \in \operatorname{interval}(1, n)) \wedge \\ \quad (l \in \operatorname{interval}(1, n)) \wedge \\ \quad \neg (k = l) \Rightarrow \\ \quad \quad \operatorname{det}(r, n, \operatorname{matmul}(r, \operatorname{elem\text{-}f}(r, n, k, l), a)) \\ \quad \quad =\; \operatorname{neg}(r)(\operatorname{det}(r, n, a)) \end{array}}}
\vnbentry{\texttt{dtr-\allowbreak{}ring-\allowbreak{}rx}}{\vnbstmt{\begin{array}{@{}l@{}} \forall r, d, e, t. \\ \quad \operatorname{is\text{-}commutative\text{-}ring}(r) \wedge \\ \quad (d \in \operatorname{carr}(r)) \wedge \\ \quad (e \in \operatorname{carr}(r)) \wedge \\ \quad (t \in \operatorname{carr}(r)) \wedge \\ \quad (e = \operatorname{neg}(r)(d)) \wedge \\ \quad (e = t) \Rightarrow \\ \quad \quad d \\ \quad \quad =\; \operatorname{add}(r)(\operatorname{mul}(r)(\operatorname{neg}(r)(\operatorname{one}(r)), t), \operatorname{zero}(r)) \end{array}}}
\vnbentry{\texttt{dtr-\allowbreak{}ring-\allowbreak{}rx2}}{\vnbstmt{\begin{array}{@{}l@{}} \forall r, s, a, d. \\ \quad \operatorname{is\text{-}commutative\text{-}ring}(r) \wedge \\ \quad (s \in \operatorname{carr}(r)) \wedge \\ \quad (a \in \operatorname{carr}(r)) \wedge \\ \quad (d \in \operatorname{carr}(r)) \Rightarrow \\ \quad \quad \operatorname{mul}(r)(\operatorname{mul}(r)(\operatorname{neg}(r)(\operatorname{one}(r)), s), \operatorname{mul}(r)(a, d)) \\ \quad \quad =\; \operatorname{add}(r)(\operatorname{mul}(r)(\operatorname{neg}(r)(\operatorname{one}(r)), \operatorname{mul}(r)(s, \operatorname{mul}(r)(a, d))), \\ \quad \quad \quad \operatorname{zero}(r)) \end{array}}}
\vnbentry{\texttt{det-\allowbreak{}expand-\allowbreak{}row-\allowbreak{}ind}}{\vnbstmt{\begin{array}{@{}l@{}} \forall i \in \mathbb{N}, n, r, a. \\ \quad (n \in \mathbb{N}) \wedge \\ \quad \operatorname{is\text{-}commutative\text{-}ring}(r) \wedge \\ \quad (a \in \operatorname{mat}((n + 1), (n + 1), \operatorname{carr}(r))) \wedge \\ \quad (1 \leq i) \wedge \\ \quad (i \leq (n + 1)) \Rightarrow \\ \quad \quad \operatorname{det}(r, (n + 1), a) \\ \quad \quad =\; \operatorname{finsum}(\operatorname{ring\text{-}additive\text{-}ag}(r), \\ \quad \quad \quad (\lambda j \in \operatorname{interval}(1, (n + 1)).\; \\ \quad \quad \quad \quad \operatorname{mul}(r)(\operatorname{mpow}(\operatorname{ring\text{-}multiplicative\text{-}monoid}(r), \operatorname{neg}(r)(\operatorname{one}(r)), (i + j)), \\ \quad \quad \quad \quad \quad \operatorname{mul}(r)(\operatorname{entry}(a, i, j), \operatorname{det}(r, n, \operatorname{minor}(a, i, j, n))))), \\ \quad \quad \quad \operatorname{interval}(1, (n + 1))) \end{array}}}
\vnbentry{\texttt{det-\allowbreak{}expand-\allowbreak{}row}}{\vnbstmt{\begin{array}{@{}l@{}} \forall r, n, a, i. \\ \quad \operatorname{is\text{-}commutative\text{-}ring}(r) \wedge \\ \quad (n \in \mathbb{N}) \wedge \\ \quad (a \in \operatorname{mat}((n + 1), (n + 1), \operatorname{carr}(r))) \wedge \\ \quad (i \in \operatorname{interval}(1, (n + 1))) \Rightarrow \\ \quad \quad \operatorname{det}(r, (n + 1), a) \\ \quad \quad =\; \operatorname{finsum}(\operatorname{ring\text{-}additive\text{-}ag}(r), \\ \quad \quad \quad (\lambda j \in \operatorname{interval}(1, (n + 1)).\; \\ \quad \quad \quad \quad \operatorname{mul}(r)(\operatorname{mpow}(\operatorname{ring\text{-}multiplicative\text{-}monoid}(r), \operatorname{neg}(r)(\operatorname{one}(r)), (i + j)), \\ \quad \quad \quad \quad \quad \operatorname{mul}(r)(\operatorname{entry}(a, i, j), \operatorname{det}(r, n, \operatorname{minor}(a, i, j, n))))), \\ \quad \quad \quad \operatorname{interval}(1, (n + 1))) \end{array}}}

\vnbfile{theorem-library/det-mul.scm}{theorem-library/det-mul.scm}
\noindent{\small det-mul.scm -- det(E A) = det(E) det(A) for an ELEMENTARY matrix E (batch 25-C, 2026-09-24).  The multiplicativity of the determinant as far as the tree carries it; see the report (scratchpad/triage/BATCH25-REPORTS.md, 25-C) for what stands between this and linear-algebra 3.20 / 3.27.}\par
\vnbentry{\texttt{det-\allowbreak{}elem-\allowbreak{}h}}{\vnbstmt{\begin{array}{@{}l@{}} \forall r, n, c, k. \\ \quad \operatorname{is\text{-}commutative\text{-}ring}(r) \wedge \\ \quad (n \in \mathbb{N}) \wedge \\ \quad (c \in \operatorname{carr}(r)) \wedge \\ \quad (k \in \operatorname{interval}(1, n)) \Rightarrow \\ \quad \quad (\operatorname{det}(r, n, \operatorname{elem\text{-}h}(r, n, c, k)) = c) \end{array}}}
\vnbentry{\texttt{det-\allowbreak{}elem-\allowbreak{}g}}{\vnbstmt{\begin{array}{@{}l@{}} \forall r, n, c, k, l. \\ \quad \operatorname{is\text{-}commutative\text{-}ring}(r) \wedge \\ \quad (n \in \mathbb{N}) \wedge \\ \quad (c \in \operatorname{carr}(r)) \wedge \\ \quad (k \in \operatorname{interval}(1, n)) \wedge \\ \quad (l \in \operatorname{interval}(1, n)) \wedge \\ \quad \neg (k = l) \Rightarrow \\ \quad \quad (\operatorname{det}(r, n, \operatorname{elem\text{-}g}(r, n, c, k, l)) = \operatorname{one}(r)) \end{array}}}
\vnbentry{\texttt{det-\allowbreak{}elem-\allowbreak{}f}}{\vnbstmt{\begin{array}{@{}l@{}} \forall r, n, k, l. \\ \quad \operatorname{is\text{-}commutative\text{-}ring}(r) \wedge \\ \quad (n \in \mathbb{N}) \wedge \\ \quad (k \in \operatorname{interval}(1, n)) \wedge \\ \quad (l \in \operatorname{interval}(1, n)) \wedge \\ \quad \neg (k = l) \Rightarrow \\ \quad \quad \operatorname{det}(r, n, \operatorname{elem\text{-}f}(r, n, k, l)) \\ \quad \quad =\; \operatorname{neg}(r)(\operatorname{one}(r)) \end{array}}}
\vnbentry{\texttt{det-\allowbreak{}mul-\allowbreak{}elem-\allowbreak{}h}}{\vnbstmt{\begin{array}{@{}l@{}} \forall r, n, a, c, k. \\ \quad \operatorname{is\text{-}commutative\text{-}ring}(r) \wedge \\ \quad (n \in \mathbb{N}) \wedge \\ \quad (a \in \operatorname{mat}(n, n, \operatorname{carr}(r))) \wedge \\ \quad (c \in \operatorname{carr}(r)) \wedge \\ \quad (k \in \operatorname{interval}(1, n)) \Rightarrow \\ \quad \quad \operatorname{det}(r, n, \operatorname{matmul}(r, \operatorname{elem\text{-}h}(r, n, c, k), a)) \\ \quad \quad =\; \operatorname{mul}(r)(\operatorname{det}(r, n, \operatorname{elem\text{-}h}(r, n, c, k)), \operatorname{det}(r, n, a)) \end{array}}}
\vnbentry{\texttt{det-\allowbreak{}mul-\allowbreak{}elem-\allowbreak{}g}}{\vnbstmt{\begin{array}{@{}l@{}} \forall r, n, a, c, k, l. \\ \quad \operatorname{is\text{-}commutative\text{-}ring}(r) \wedge \\ \quad (n \in \mathbb{N}) \wedge \\ \quad (a \in \operatorname{mat}(n, n, \operatorname{carr}(r))) \wedge \\ \quad (c \in \operatorname{carr}(r)) \wedge \\ \quad (k \in \operatorname{interval}(1, n)) \wedge \\ \quad (l \in \operatorname{interval}(1, n)) \wedge \\ \quad \neg (k = l) \Rightarrow \\ \quad \quad \operatorname{det}(r, n, \operatorname{matmul}(r, \operatorname{elem\text{-}g}(r, n, c, k, l), a)) \\ \quad \quad =\; \operatorname{mul}(r)(\operatorname{det}(r, n, \operatorname{elem\text{-}g}(r, n, c, k, l)), \operatorname{det}(r, n, a)) \end{array}}}
\vnbentry{\texttt{det-\allowbreak{}mul-\allowbreak{}elem-\allowbreak{}f}}{\vnbstmt{\begin{array}{@{}l@{}} \forall r, n, a, k, l. \\ \quad \operatorname{is\text{-}commutative\text{-}ring}(r) \wedge \\ \quad (n \in \mathbb{N}) \wedge \\ \quad (a \in \operatorname{mat}(n, n, \operatorname{carr}(r))) \wedge \\ \quad (k \in \operatorname{interval}(1, n)) \wedge \\ \quad (l \in \operatorname{interval}(1, n)) \wedge \\ \quad \neg (k = l) \Rightarrow \\ \quad \quad \operatorname{det}(r, n, \operatorname{matmul}(r, \operatorname{elem\text{-}f}(r, n, k, l), a)) \\ \quad \quad =\; \operatorname{mul}(r)(\operatorname{det}(r, n, \operatorname{elem\text{-}f}(r, n, k, l)), \operatorname{det}(r, n, a)) \end{array}}}

\vnbfile{theorem-library/euclidean-division-proof.scm}{theorem-library/euclidean-division-proof.scm}
\noindent{\small euclidean-division-proof.scm -- Brick B2 (LA Phase B, algebraic-numbers.pdf ch.3): division-with-remainder in usable form over any euclidean ring.  For a euclidean ring s, a in CARR, b in CARR, b/=0, there are q,r with a = q.b + r and (r = 0 or deg(r) $<$ deg(b)), where deg = GAUGE(s) (the named degree).  This}\par
\vnbentry{\texttt{euclidean-\allowbreak{}division}}{\vnbstmt{\begin{array}{@{}l@{}} \forall s.; \forall a, b \in \operatorname{carr}(s). \\ \quad (\operatorname{is\text{-}euclidean\text{-}ring}(s) \wedge \neg (b = \operatorname{zero}(s))) \Rightarrow  \\ \quad \quad \exists q, r \in \operatorname{carr}(s). \\ \quad \quad \quad (a = \operatorname{add}(s)(\operatorname{mul}(s)(q, b), r)) \wedge \\ \quad \quad \quad (r = \operatorname{zero}(s)) \vee \\ \quad \quad \quad ((\operatorname{gauge}(s)(r) + 1) \leq \operatorname{gauge}(s)(b)) \end{array}}}

\vnbfile{theorem-library/pivot-col-reduce-proof.scm}{theorem-library/pivot-col-reduce-proof.scm}
\noindent{\small pivot-col-reduce-proof.scm -- Brick B2b (LA Phase B, algebraic-numbers.pdf ch.3): ONE elementary column operation reduces a first-row entry against the pivot.  Over a euclidean ring A, for P with nonzero pivot P\_\{1,1\} and target column j/=1, there are q,r with (P . G[-q,1,j])\_\{1,j\} = r and (r=0 or}\par
\vnbentry{\texttt{pivot-\allowbreak{}col-\allowbreak{}reduce}}{\vnbstmt{\begin{array}{@{}l@{}} \forall a, m, n, p, j. \\ \quad \operatorname{is\text{-}euclidean\text{-}ring}(a) \wedge \\ \quad (p \in \operatorname{mat}(m, n, \operatorname{carr}(a))) \wedge \\ \quad (1 \in \operatorname{interval}(1, m)) \wedge \\ \quad (1 \in \operatorname{interval}(1, n)) \wedge \\ \quad (j \in \operatorname{interval}(1, n)) \wedge \\ \quad \neg (1 = j) \wedge \\ \quad \neg (\operatorname{entry}(p, 1, 1) = \operatorname{zero}(a)) \Rightarrow \\ \quad \quad \exists q, r \in \operatorname{carr}(a). \\ \quad \quad \quad \operatorname{entry}(\operatorname{matmul}(a, p, \operatorname{elem\text{-}g}(a, n, \operatorname{neg}(a)(q), 1, j)), 1, j) \\ \quad \quad \quad =\; r \wedge \\ \quad \quad \quad (r = \operatorname{zero}(a)) \vee \\ \quad \quad \quad ((\operatorname{gauge}(a)(r) + 1) \leq \operatorname{gauge}(a)(\operatorname{entry}(p, 1, 1))) \end{array}}}

\vnbfile{theorem-library/pivot-row-reduce-proof.scm}{theorem-library/pivot-row-reduce-proof.scm}
\noindent{\small pivot-row-reduce-proof.scm -- Brick B3a (LA Phase B, algebraic-numbers.pdf ch.3): the ROW-reduction step, mirror of pivot-col-reduce.  Over a euclidean ring A, for P with nonzero pivot P\_\{1,1\} and target row i/=1, there are q,r with (ELEM-G[-q,i,1] . P)\_\{i,1\} = r and (r=0 or deg(r) $<$ deg(P\_\{1,1\})) -- one}\par
\vnbentry{\texttt{pivot-\allowbreak{}row-\allowbreak{}reduce}}{\vnbstmt{\begin{array}{@{}l@{}} \forall a, m, n, p, i. \\ \quad \operatorname{is\text{-}euclidean\text{-}ring}(a) \wedge \\ \quad (p \in \operatorname{mat}(m, n, \operatorname{carr}(a))) \wedge \\ \quad (1 \in \operatorname{interval}(1, m)) \wedge \\ \quad (1 \in \operatorname{interval}(1, n)) \wedge \\ \quad (i \in \operatorname{interval}(1, m)) \wedge \\ \quad \neg (i = 1) \wedge \\ \quad \neg (\operatorname{entry}(p, 1, 1) = \operatorname{zero}(a)) \Rightarrow \\ \quad \quad \exists q, r \in \operatorname{carr}(a). \\ \quad \quad \quad \operatorname{entry}(\operatorname{matmul}(a, \operatorname{elem\text{-}g}(a, m, \operatorname{neg}(a)(q), i, 1), p), i, 1) \\ \quad \quad \quad =\; r \wedge \\ \quad \quad \quad (r = \operatorname{zero}(a)) \vee \\ \quad \quad \quad ((\operatorname{gauge}(a)(r) + 1) \leq \operatorname{gauge}(a)(\operatorname{entry}(p, 1, 1))) \end{array}}}

\vnbfile{theorem-library/mat-equiv-proof.scm}{theorem-library/mat-equiv-proof.scm}
\noindent{\small mat-equiv-proof.scm -- Brick B3 infrastructure: basic laws of \textasciitilde{} (mat-equiv.scm). invertible-mat-is-mat, identmat-invertible, mat-equiv-refl, mat-equiv-right-mult, mat-equiv-left-mult, elem-f-invertible.  \textasciitilde{} is reflexive and each elementary swap (row/col) yields an equivalent matrix.  elem-g-invertible + transitivity + ...}\par
\vnbentry{\texttt{invertible-\allowbreak{}mat-\allowbreak{}is-\allowbreak{}mat}}{\vnbstmt{\begin{array}{@{}l@{}} \forall a, n, u. \\ \quad \operatorname{is\text{-}invertible\text{-}mat}(a, n, u) \Rightarrow  \\ \quad \quad (u \in \operatorname{mat}(n, n, \operatorname{carr}(a))) \end{array}}}
\vnbentry{\texttt{identmat-\allowbreak{}invertible}}{\vnbstmt{\begin{array}{@{}l@{}} \forall a.; \forall n \in \mathbb{N}. \\ \quad \operatorname{is\text{-}ring}(a) \Rightarrow  \\ \quad \quad \operatorname{is\text{-}invertible\text{-}mat}(a, n, \operatorname{identmat}(a, n)) \end{array}}}
\vnbentry{\texttt{mat-\allowbreak{}equiv-\allowbreak{}refl}}{\vnbstmt{\begin{array}{@{}l@{}} \forall a, m, n.; \forall c \in \operatorname{mat}(m, n, \operatorname{carr}(a)). \\ \quad \operatorname{is\text{-}ring}(a) \Rightarrow  \\ \quad \quad \operatorname{mat\text{-}equiv}(a, m, n, c, c) \end{array}}}
\vnbentry{\texttt{mat-\allowbreak{}equiv-\allowbreak{}right-\allowbreak{}mult}}{\vnbstmt{\begin{array}{@{}l@{}} \forall a, m, n, c, v. \\ \quad \operatorname{is\text{-}ring}(a) \wedge \\ \quad (c \in \operatorname{mat}(m, n, \operatorname{carr}(a))) \wedge \\ \quad \operatorname{is\text{-}invertible\text{-}mat}(a, n, v) \Rightarrow \\ \quad \quad \operatorname{mat\text{-}equiv}(a, m, n, c, \operatorname{matmul}(a, c, v)) \end{array}}}
\vnbentry{\texttt{mat-\allowbreak{}equiv-\allowbreak{}left-\allowbreak{}mult}}{\vnbstmt{\begin{array}{@{}l@{}} \forall a, m, n, c, u. \\ \quad \operatorname{is\text{-}ring}(a) \wedge \\ \quad (c \in \operatorname{mat}(m, n, \operatorname{carr}(a))) \wedge \\ \quad \operatorname{is\text{-}invertible\text{-}mat}(a, m, u) \Rightarrow \\ \quad \quad \operatorname{mat\text{-}equiv}(a, m, n, c, \operatorname{matmul}(a, u, c)) \end{array}}}
\vnbentry{\texttt{elem-\allowbreak{}f-\allowbreak{}invertible}}{\vnbstmt{\begin{array}{@{}l@{}} \forall a, n, k, l. \\ \quad \operatorname{is\text{-}ring}(a) \wedge \\ \quad (n \in \mathbb{N}) \wedge \\ \quad (k \in \operatorname{interval}(1, n)) \wedge \\ \quad (l \in \operatorname{interval}(1, n)) \wedge \\ \quad \neg (k = l) \Rightarrow \\ \quad \quad \operatorname{is\text{-}invertible\text{-}mat}(a, n, \operatorname{elem\text{-}f}(a, n, k, l)) \end{array}}}
\vnbentry{\texttt{elem-\allowbreak{}g-\allowbreak{}param-\allowbreak{}cong}}{\vnbstmt{\begin{array}{@{}l@{}} \forall a, n, r, s, k, l. \\ \quad \operatorname{is\text{-}ring}(a) \wedge \\ \quad (n \in \mathbb{N}) \wedge \\ \quad (s \in \operatorname{carr}(a)) \wedge \\ \quad (r = s) \Rightarrow \\ \quad \quad (\operatorname{elem\text{-}g}(a, n, r, k, l) = \operatorname{elem\text{-}g}(a, n, s, k, l)) \end{array}}}
\vnbentry{\texttt{elem-\allowbreak{}g-\allowbreak{}invertible}}{\vnbstmt{\begin{array}{@{}l@{}} \forall a, n, r, k, l. \\ \quad \operatorname{is\text{-}ring}(a) \wedge \\ \quad (n \in \mathbb{N}) \wedge \\ \quad (r \in \operatorname{carr}(a)) \wedge \\ \quad (k \in \operatorname{interval}(1, n)) \wedge \\ \quad (l \in \operatorname{interval}(1, n)) \wedge \\ \quad \neg (k = l) \Rightarrow \\ \quad \quad \operatorname{is\text{-}invertible\text{-}mat}(a, n, \operatorname{elem\text{-}g}(a, n, r, k, l)) \end{array}}}
\vnbentry{\texttt{product-\allowbreak{}of-\allowbreak{}invertibles-\allowbreak{}is-\allowbreak{}invertible}}{\vnbstmt{\begin{array}{@{}l@{}} \forall a, n, b, c. \\ \quad \operatorname{is\text{-}ring}(a) \wedge \\ \quad \operatorname{is\text{-}invertible\text{-}mat}(a, n, b) \wedge \\ \quad \operatorname{is\text{-}invertible\text{-}mat}(a, n, c) \Rightarrow \\ \quad \quad \operatorname{is\text{-}invertible\text{-}mat}(a, n, \operatorname{matmul}(a, b, c)) \end{array}}}
\vnbentry{\texttt{mat-\allowbreak{}equiv-\allowbreak{}trans}}{\vnbstmt{\begin{array}{@{}l@{}} \forall a, m, n, c, d, e. \\ \quad \operatorname{is\text{-}ring}(a) \wedge \\ \quad (c \in \operatorname{mat}(m, n, \operatorname{carr}(a))) \wedge \\ \quad \operatorname{mat\text{-}equiv}(a, m, n, c, d) \wedge \\ \quad \operatorname{mat\text{-}equiv}(a, m, n, d, e) \Rightarrow \\ \quad \quad \operatorname{mat\text{-}equiv}(a, m, n, c, e) \end{array}}}

\vnbfile{theorem-library/min-degree-entry-proof.scm}{theorem-library/min-degree-entry-proof.scm}
\noindent{\small min-degree-entry-proof.scm -- a matrix with a nonzero entry has a nonzero entry of MINIMAL Euclidean degree.  The "mu is achieved" fact that starts the Smith reduction (Prop 3.36).}\par
\vnbentry{\texttt{min-\allowbreak{}degree-\allowbreak{}entry}}{\vnbstmt{\begin{array}{@{}l@{}} \forall a, m, n.; \forall p \in \operatorname{mat}(m, n, \operatorname{carr}(a)). \\ \quad \operatorname{is\text{-}euclidean\text{-}ring}(a) \wedge \\ \quad \exists b \in \operatorname{interval}(1, m), c \in \operatorname{interval}(1, n). \\ \quad \quad \neg (\operatorname{entry}(p, b, c) = \operatorname{zero}(a)) \Rightarrow \\ \quad \quad \exists s \in \operatorname{interval}(1, m), d \in \operatorname{interval}(1, n). \\ \quad \quad \quad \neg (\operatorname{entry}(p, s, d) = \operatorname{zero}(a)) \wedge \\ \quad \quad \quad \forall i, j. \\ \quad \quad \quad \quad (i \in \operatorname{interval}(1, m)) \wedge \\ \quad \quad \quad \quad (j \in \operatorname{interval}(1, n)) \wedge \\ \quad \quad \quad \quad \neg (\operatorname{entry}(p, i, j) = \operatorname{zero}(a)) \Rightarrow \\ \quad \quad \quad \quad \quad \operatorname{gauge}(a)(\operatorname{entry}(p, s, d)) \\ \quad \quad \quad \quad \quad \leq\; \operatorname{gauge}(a)(\operatorname{entry}(p, i, j)) \end{array}}}

\vnbfile{theorem-library/smith-proof.scm}{theorem-library/smith-proof.scm}
\noindent{\small smith-proof.scm -- Brick B3 (LA Phase B): the Smith normal-form reduction, built on the \textasciitilde{} equivalence infrastructure (mat-equiv-proof.scm). equiv-mul-both  -- A \textasciitilde{} U.A.V for U,V invertible (sandwiching preserves \textasciitilde{}); swap-to-corner  -- bring any entry (i0,j0) to position (1,1) preserving \textasciitilde{}}\par
\vnbentry{\texttt{equiv-\allowbreak{}mul-\allowbreak{}both}}{\vnbstmt{\begin{array}{@{}l@{}} \forall r, m, n, a, u, v. \\ \quad \operatorname{is\text{-}ring}(r) \wedge \\ \quad (a \in \operatorname{mat}(m, n, \operatorname{carr}(r))) \wedge \\ \quad \operatorname{is\text{-}invertible\text{-}mat}(r, m, u) \wedge \\ \quad \operatorname{is\text{-}invertible\text{-}mat}(r, n, v) \Rightarrow \\ \quad \quad \operatorname{mat\text{-}equiv}(r, m, n, a, \operatorname{matmul}(r, \operatorname{matmul}(r, u, a), v)) \end{array}}}
\vnbentry{\texttt{swap-\allowbreak{}to-\allowbreak{}corner}}{\vnbstmt{\begin{array}{@{}l@{}} \forall r, m, n, a, c, d. \\ \quad \operatorname{is\text{-}ring}(r) \wedge \\ \quad (a \in \operatorname{mat}(m, n, \operatorname{carr}(r))) \wedge \\ \quad (1 \in \operatorname{interval}(1, m)) \wedge \\ \quad (1 \in \operatorname{interval}(1, n)) \wedge \\ \quad (c \in \operatorname{interval}(1, m)) \wedge \\ \quad (d \in \operatorname{interval}(1, n)) \wedge \\ \quad \neg (1 = c) \wedge \\ \quad \neg (1 = d) \Rightarrow \\ \quad \quad \exists b. \\ \quad \quad \quad \operatorname{mat\text{-}equiv}(r, m, n, a, b) \wedge \\ \quad \quad \quad (\operatorname{entry}(b, 1, 1) = \operatorname{entry}(a, c, d)) \end{array}}}
\vnbentry{\texttt{swap-\allowbreak{}to-\allowbreak{}corner-\allowbreak{}gen}}{\vnbstmt{\begin{array}{@{}l@{}} \forall r, m, n, a, c, d. \\ \quad \operatorname{is\text{-}ring}(r) \wedge \\ \quad (a \in \operatorname{mat}(m, n, \operatorname{carr}(r))) \wedge \\ \quad (1 \in \operatorname{interval}(1, m)) \wedge \\ \quad (1 \in \operatorname{interval}(1, n)) \wedge \\ \quad (c \in \operatorname{interval}(1, m)) \wedge \\ \quad (d \in \operatorname{interval}(1, n)) \Rightarrow \\ \quad \quad \exists b. \\ \quad \quad \quad \operatorname{mat\text{-}equiv}(r, m, n, a, b) \wedge \\ \quad \quad \quad (\operatorname{entry}(b, 1, 1) = \operatorname{entry}(a, c, d)) \end{array}}}
\vnbentry{\texttt{place-\allowbreak{}min-\allowbreak{}pivot}}{\vnbstmt{\begin{array}{@{}l@{}} \forall a, m, n.; \forall p \in \operatorname{mat}(m, n, \operatorname{carr}(a)). \\ \quad \operatorname{is\text{-}euclidean\text{-}ring}(a) \wedge \\ \quad (1 \in \operatorname{interval}(1, m)) \wedge \\ \quad (1 \in \operatorname{interval}(1, n)) \wedge \\ \quad \exists c \in \operatorname{interval}(1, m), d \in \operatorname{interval}(1, n). \\ \quad \quad \neg (\operatorname{entry}(p, c, d) = \operatorname{zero}(a)) \Rightarrow \\ \quad \quad \exists b. \\ \quad \quad \quad \operatorname{mat\text{-}equiv}(a, m, n, p, b) \wedge \\ \quad \quad \quad \neg (\operatorname{entry}(b, 1, 1) = \operatorname{zero}(a)) \wedge \\ \quad \quad \quad \forall i, j. \\ \quad \quad \quad \quad (i \in \operatorname{interval}(1, m)) \wedge \\ \quad \quad \quad \quad (j \in \operatorname{interval}(1, n)) \wedge \\ \quad \quad \quad \quad \neg (\operatorname{entry}(p, i, j) = \operatorname{zero}(a)) \Rightarrow \\ \quad \quad \quad \quad \quad \operatorname{gauge}(a)(\operatorname{entry}(b, 1, 1)) \\ \quad \quad \quad \quad \quad \leq\; \operatorname{gauge}(a)(\operatorname{entry}(p, i, j)) \end{array}}}

\vnbfile{theorem-library/class-min-pivot-proof.scm}{theorem-library/class-min-pivot-proof.scm}
\noindent{\small class-min-pivot-proof.scm -- the minimal-degree pivot over the whole \textasciitilde{}-equivalence class, and the typing lemma it needs.}\par
\vnbentry{\texttt{mat-\allowbreak{}equiv-\allowbreak{}target-\allowbreak{}is-\allowbreak{}mat}}{\vnbstmt{\begin{array}{@{}l@{}} \forall a, m, n, c, d. \\ \quad \operatorname{is\text{-}ring}(a) \wedge \\ \quad (c \in \operatorname{mat}(m, n, \operatorname{carr}(a))) \wedge \\ \quad \operatorname{mat\text{-}equiv}(a, m, n, c, d) \Rightarrow \\ \quad \quad (d \in \operatorname{mat}(m, n, \operatorname{carr}(a))) \end{array}}}
\vnbentry{\texttt{class-\allowbreak{}min-\allowbreak{}pivot}}{\vnbstmt{\begin{array}{@{}l@{}} \forall a, m, n.; \forall p \in \operatorname{mat}(m, n, \operatorname{carr}(a)). \\ \quad \operatorname{is\text{-}euclidean\text{-}ring}(a) \wedge \\ \quad (1 \in \operatorname{interval}(1, m)) \wedge \\ \quad (1 \in \operatorname{interval}(1, n)) \wedge \\ \quad \exists d \in \operatorname{interval}(1, m), f \in \operatorname{interval}(1, n). \\ \quad \quad \neg (\operatorname{entry}(p, d, f) = \operatorname{zero}(a)) \Rightarrow \\ \quad \quad \exists b. \\ \quad \quad \quad \operatorname{mat\text{-}equiv}(a, m, n, p, b) \wedge \\ \quad \quad \quad \neg (\operatorname{entry}(b, 1, 1) = \operatorname{zero}(a)) \wedge \\ \quad \quad \quad \forall c, i, j. \\ \quad \quad \quad \quad \operatorname{mat\text{-}equiv}(a, m, n, p, c) \wedge \\ \quad \quad \quad \quad (i \in \operatorname{interval}(1, m)) \wedge \\ \quad \quad \quad \quad (j \in \operatorname{interval}(1, n)) \wedge \\ \quad \quad \quad \quad \neg (\operatorname{entry}(c, i, j) = \operatorname{zero}(a)) \Rightarrow \\ \quad \quad \quad \quad \quad \operatorname{gauge}(a)(\operatorname{entry}(b, 1, 1)) \\ \quad \quad \quad \quad \quad \leq\; \operatorname{gauge}(a)(\operatorname{entry}(c, i, j)) \end{array}}}

\vnbfile{theorem-library/smith-clear-proof.scm}{theorem-library/smith-clear-proof.scm}
\noindent{\small smith-clear-proof.scm -- Brick B3 (LA Phase B): the DESCENT step of the Smith reduction.  pivot-clears-col: over a euclidean ring, if P's pivot P\_\{1,1\} is nonzero and of MINIMAL degree over the whole equivalence class of P (class-min-pivot supplies such a P), then a single elementary COLUMN op zeroes}\par
\vnbentry{\texttt{pivot-\allowbreak{}clears-\allowbreak{}col}}{\vnbstmt{\begin{array}{@{}l@{}} \forall a, m, n, p, j. \\ \quad \operatorname{is\text{-}euclidean\text{-}ring}(a) \wedge \\ \quad (p \in \operatorname{mat}(m, n, \operatorname{carr}(a))) \wedge \\ \quad (1 \in \operatorname{interval}(1, m)) \wedge \\ \quad (1 \in \operatorname{interval}(1, n)) \wedge \\ \quad (j \in \operatorname{interval}(1, n)) \wedge \\ \quad \neg (1 = j) \wedge \\ \quad \neg (\operatorname{entry}(p, 1, 1) = \operatorname{zero}(a)) \wedge \\ \quad \forall c, i, b. \\ \quad \quad \operatorname{mat\text{-}equiv}(a, m, n, p, c) \wedge \\ \quad \quad (i \in \operatorname{interval}(1, m)) \wedge \\ \quad \quad (b \in \operatorname{interval}(1, n)) \wedge \\ \quad \quad \neg (\operatorname{entry}(c, i, b) = \operatorname{zero}(a)) \Rightarrow \\ \quad \quad \quad \operatorname{gauge}(a)(\operatorname{entry}(p, 1, 1)) \\ \quad \quad \quad \leq\; \operatorname{gauge}(a)(\operatorname{entry}(c, i, b)) \Rightarrow \\ \quad \quad \exists q. \\ \quad \quad \quad \operatorname{mat\text{-}equiv}(a, m, n, p, q) \wedge \\ \quad \quad \quad (\operatorname{entry}(q, 1, 1) = \operatorname{entry}(p, 1, 1)) \wedge \\ \quad \quad \quad (\operatorname{entry}(q, 1, j) = \operatorname{zero}(a)) \wedge \\ \quad \quad \quad \forall i, c. \\ \quad \quad \quad \quad (i \in \operatorname{interval}(1, m)) \wedge \\ \quad \quad \quad \quad (c \in \operatorname{interval}(1, n)) \wedge \\ \quad \quad \quad \quad \neg (c = j) \Rightarrow \\ \quad \quad \quad \quad \quad (\operatorname{entry}(q, i, c) = \operatorname{entry}(p, i, c)) \end{array}}}

\vnbfile{theorem-library/clear-first-row-proof.scm}{theorem-library/clear-first-row-proof.scm}
\noindent{\small clear-first-row-proof.scm -- Brick B3 (LA Phase B): clear the whole first row. Iterates pivot-clears-col across columns 2..n by NN-induction on the number of cleared columns, then instantiates at n.  With a class-minimal pivot at (1,1), the result is a matrix \textasciitilde{} P, same pivot, whose first row is zero off the pivot.}\par
\vnbentry{\texttt{classmin-\allowbreak{}transport}}{\vnbstmt{\begin{array}{@{}l@{}} \forall a, m, n, p, q, c, i, j. \\ \quad \operatorname{is\text{-}ring}(a) \wedge \\ \quad (p \in \operatorname{mat}(m, n, \operatorname{carr}(a))) \wedge \\ \quad \operatorname{mat\text{-}equiv}(a, m, n, p, q) \wedge \\ \quad (\operatorname{entry}(q, 1, 1) = \operatorname{entry}(p, 1, 1)) \wedge \\ \quad \forall c, i, j. \\ \quad \quad \operatorname{mat\text{-}equiv}(a, m, n, p, c) \wedge \\ \quad \quad (i \in \operatorname{interval}(1, m)) \wedge \\ \quad \quad (j \in \operatorname{interval}(1, n)) \wedge \\ \quad \quad \neg (\operatorname{entry}(c, i, j) = \operatorname{zero}(a)) \Rightarrow \\ \quad \quad \quad \operatorname{gauge}(a)(\operatorname{entry}(p, 1, 1)) \\ \quad \quad \quad \leq\; \operatorname{gauge}(a)(\operatorname{entry}(c, i, j)) \wedge \\ \quad \operatorname{mat\text{-}equiv}(a, m, n, q, c) \wedge \\ \quad (i \in \operatorname{interval}(1, m)) \wedge \\ \quad (j \in \operatorname{interval}(1, n)) \wedge \\ \quad \neg (\operatorname{entry}(c, i, j) = \operatorname{zero}(a)) \Rightarrow \\ \quad \quad \operatorname{gauge}(a)(\operatorname{entry}(q, 1, 1)) \\ \quad \quad \leq\; \operatorname{gauge}(a)(\operatorname{entry}(c, i, j)) \end{array}}}
\vnbentry{\texttt{clear-\allowbreak{}row-\allowbreak{}upto}}{\vnbstmt{\begin{array}{@{}l@{}} \forall a, m, n.; \forall p \in \operatorname{mat}(m, n, \operatorname{carr}(a)), k \in \mathbb{N}. \\ \quad \operatorname{is\text{-}euclidean\text{-}ring}(a) \wedge \\ \quad (1 \in \operatorname{interval}(1, m)) \wedge \\ \quad (1 \in \operatorname{interval}(1, n)) \wedge \\ \quad \neg (\operatorname{entry}(p, 1, 1) = \operatorname{zero}(a)) \wedge \\ \quad \forall c, i, b. \\ \quad \quad \operatorname{mat\text{-}equiv}(a, m, n, p, c) \wedge \\ \quad \quad (i \in \operatorname{interval}(1, m)) \wedge \\ \quad \quad (b \in \operatorname{interval}(1, n)) \wedge \\ \quad \quad \neg (\operatorname{entry}(c, i, b) = \operatorname{zero}(a)) \Rightarrow \\ \quad \quad \quad \operatorname{gauge}(a)(\operatorname{entry}(p, 1, 1)) \\ \quad \quad \quad \leq\; \operatorname{gauge}(a)(\operatorname{entry}(c, i, b)) \Rightarrow \\ \quad \quad \exists q. \\ \quad \quad \quad \operatorname{mat\text{-}equiv}(a, m, n, p, q) \wedge \\ \quad \quad \quad (\operatorname{entry}(q, 1, 1) = \operatorname{entry}(p, 1, 1)) \wedge \\ \quad \quad \quad \forall j \in \operatorname{interval}(1, n). \\ \quad \quad \quad \quad (\neg (j = 1) \wedge (j \leq k)) \Rightarrow  \\ \quad \quad \quad \quad \quad (\operatorname{entry}(q, 1, j) = \operatorname{zero}(a)) \end{array}}}
\vnbentry{\texttt{clear-\allowbreak{}first-\allowbreak{}row}}{\vnbstmt{\begin{array}{@{}l@{}} \forall a, m, n.; \forall p \in \operatorname{mat}(m, n, \operatorname{carr}(a)). \\ \quad \operatorname{is\text{-}euclidean\text{-}ring}(a) \wedge \\ \quad (m \in \mathbb{N}) \wedge \\ \quad (n \in \mathbb{N}) \wedge \\ \quad (1 \in \operatorname{interval}(1, m)) \wedge \\ \quad (1 \in \operatorname{interval}(1, n)) \wedge \\ \quad \neg (\operatorname{entry}(p, 1, 1) = \operatorname{zero}(a)) \wedge \\ \quad \forall c, i, b. \\ \quad \quad \operatorname{mat\text{-}equiv}(a, m, n, p, c) \wedge \\ \quad \quad (i \in \operatorname{interval}(1, m)) \wedge \\ \quad \quad (b \in \operatorname{interval}(1, n)) \wedge \\ \quad \quad \neg (\operatorname{entry}(c, i, b) = \operatorname{zero}(a)) \Rightarrow \\ \quad \quad \quad \operatorname{gauge}(a)(\operatorname{entry}(p, 1, 1)) \\ \quad \quad \quad \leq\; \operatorname{gauge}(a)(\operatorname{entry}(c, i, b)) \Rightarrow \\ \quad \quad \exists q. \\ \quad \quad \quad \operatorname{mat\text{-}equiv}(a, m, n, p, q) \wedge \\ \quad \quad \quad (\operatorname{entry}(q, 1, 1) = \operatorname{entry}(p, 1, 1)) \wedge \\ \quad \quad \quad \forall j \in \operatorname{interval}(1, n). \\ \quad \quad \quad \quad \neg (j = 1) \Rightarrow  \\ \quad \quad \quad \quad \quad (\operatorname{entry}(q, 1, j) = \operatorname{zero}(a)) \end{array}}}

\vnbfile{theorem-library/clear-first-col-proof.scm}{theorem-library/clear-first-col-proof.scm}
\noindent{\small clear-first-col-proof.scm -- Brick B3 (LA Phase B): clear the whole first column. The ROW-op mirror of clear-first-row: left-multiply by elementary row matrices. pivot-clears-row -- one row op zeroes an off-pivot column-1 entry (i,1); the euclidean remainder is forced to vanish by class-minimality}\par
\vnbentry{\texttt{pivot-\allowbreak{}clears-\allowbreak{}row}}{\vnbstmt{\begin{array}{@{}l@{}} \forall a, m, n, p, i. \\ \quad \operatorname{is\text{-}euclidean\text{-}ring}(a) \wedge \\ \quad (p \in \operatorname{mat}(m, n, \operatorname{carr}(a))) \wedge \\ \quad (1 \in \operatorname{interval}(1, m)) \wedge \\ \quad (1 \in \operatorname{interval}(1, n)) \wedge \\ \quad (i \in \operatorname{interval}(1, m)) \wedge \\ \quad \neg (i = 1) \wedge \\ \quad \neg (\operatorname{entry}(p, 1, 1) = \operatorname{zero}(a)) \wedge \\ \quad \forall c, b, j. \\ \quad \quad \operatorname{mat\text{-}equiv}(a, m, n, p, c) \wedge \\ \quad \quad (b \in \operatorname{interval}(1, m)) \wedge \\ \quad \quad (j \in \operatorname{interval}(1, n)) \wedge \\ \quad \quad \neg (\operatorname{entry}(c, b, j) = \operatorname{zero}(a)) \Rightarrow \\ \quad \quad \quad \operatorname{gauge}(a)(\operatorname{entry}(p, 1, 1)) \\ \quad \quad \quad \leq\; \operatorname{gauge}(a)(\operatorname{entry}(c, b, j)) \Rightarrow \\ \quad \quad \exists q. \\ \quad \quad \quad \operatorname{mat\text{-}equiv}(a, m, n, p, q) \wedge \\ \quad \quad \quad (\operatorname{entry}(q, 1, 1) = \operatorname{entry}(p, 1, 1)) \wedge \\ \quad \quad \quad (\operatorname{entry}(q, i, 1) = \operatorname{zero}(a)) \wedge \\ \quad \quad \quad \forall w, c. \\ \quad \quad \quad \quad (w \in \operatorname{interval}(1, m)) \wedge \\ \quad \quad \quad \quad (c \in \operatorname{interval}(1, n)) \wedge \\ \quad \quad \quad \quad \neg (w = i) \Rightarrow \\ \quad \quad \quad \quad \quad (\operatorname{entry}(q, w, c) = \operatorname{entry}(p, w, c)) \end{array}}}
\vnbentry{\texttt{clear-\allowbreak{}col-\allowbreak{}upto}}{\vnbstmt{\begin{array}{@{}l@{}} \forall a, m, n.; \forall p \in \operatorname{mat}(m, n, \operatorname{carr}(a)), k \in \mathbb{N}. \\ \quad \operatorname{is\text{-}euclidean\text{-}ring}(a) \wedge \\ \quad (1 \in \operatorname{interval}(1, m)) \wedge \\ \quad (1 \in \operatorname{interval}(1, n)) \wedge \\ \quad \neg (\operatorname{entry}(p, 1, 1) = \operatorname{zero}(a)) \wedge \\ \quad \forall c, b, j. \\ \quad \quad \operatorname{mat\text{-}equiv}(a, m, n, p, c) \wedge \\ \quad \quad (b \in \operatorname{interval}(1, m)) \wedge \\ \quad \quad (j \in \operatorname{interval}(1, n)) \wedge \\ \quad \quad \neg (\operatorname{entry}(c, b, j) = \operatorname{zero}(a)) \Rightarrow \\ \quad \quad \quad \operatorname{gauge}(a)(\operatorname{entry}(p, 1, 1)) \\ \quad \quad \quad \leq\; \operatorname{gauge}(a)(\operatorname{entry}(c, b, j)) \Rightarrow \\ \quad \quad \exists q. \\ \quad \quad \quad \operatorname{mat\text{-}equiv}(a, m, n, p, q) \wedge \\ \quad \quad \quad (\operatorname{entry}(q, 1, 1) = \operatorname{entry}(p, 1, 1)) \wedge \\ \quad \quad \quad \forall i \in \operatorname{interval}(1, m). \\ \quad \quad \quad \quad (\neg (i = 1) \wedge (i \leq k)) \Rightarrow  \\ \quad \quad \quad \quad \quad (\operatorname{entry}(q, i, 1) = \operatorname{zero}(a)) \wedge \\ \quad \quad \quad \forall c \in \operatorname{interval}(1, n). \\ \quad \quad \quad \quad (\operatorname{entry}(q, 1, c) = \operatorname{entry}(p, 1, c)) \end{array}}}
\vnbentry{\texttt{clear-\allowbreak{}first-\allowbreak{}col}}{\vnbstmt{\begin{array}{@{}l@{}} \forall a, m, n.; \forall p \in \operatorname{mat}(m, n, \operatorname{carr}(a)). \\ \quad \operatorname{is\text{-}euclidean\text{-}ring}(a) \wedge \\ \quad (m \in \mathbb{N}) \wedge \\ \quad (n \in \mathbb{N}) \wedge \\ \quad (1 \in \operatorname{interval}(1, m)) \wedge \\ \quad (1 \in \operatorname{interval}(1, n)) \wedge \\ \quad \neg (\operatorname{entry}(p, 1, 1) = \operatorname{zero}(a)) \wedge \\ \quad \forall c, b, j. \\ \quad \quad \operatorname{mat\text{-}equiv}(a, m, n, p, c) \wedge \\ \quad \quad (b \in \operatorname{interval}(1, m)) \wedge \\ \quad \quad (j \in \operatorname{interval}(1, n)) \wedge \\ \quad \quad \neg (\operatorname{entry}(c, b, j) = \operatorname{zero}(a)) \Rightarrow \\ \quad \quad \quad \operatorname{gauge}(a)(\operatorname{entry}(p, 1, 1)) \\ \quad \quad \quad \leq\; \operatorname{gauge}(a)(\operatorname{entry}(c, b, j)) \Rightarrow \\ \quad \quad \exists q. \\ \quad \quad \quad \operatorname{mat\text{-}equiv}(a, m, n, p, q) \wedge \\ \quad \quad \quad (\operatorname{entry}(q, 1, 1) = \operatorname{entry}(p, 1, 1)) \wedge \\ \quad \quad \quad \forall i \in \operatorname{interval}(1, m). \\ \quad \quad \quad \quad \neg (i = 1) \Rightarrow  \\ \quad \quad \quad \quad \quad (\operatorname{entry}(q, i, 1) = \operatorname{zero}(a)) \wedge \\ \quad \quad \quad \forall c \in \operatorname{interval}(1, n). \\ \quad \quad \quad \quad (\operatorname{entry}(q, 1, c) = \operatorname{entry}(p, 1, c)) \end{array}}}

\vnbfile{theorem-library/interval-widen.scm}{theorem-library/interval-widen.scm}
\noindent{\small interval-widen.scm -- an interval membership survives widening the top end.}\par
\vnbentry{\texttt{interval-\allowbreak{}widen}}{\vnbstmt{\begin{array}{@{}l@{}} \forall b, c \in \mathbb{N}, a.; \forall i \in \operatorname{interval}(a, b). \\ \quad ((b \leq c) \Rightarrow (i \in \operatorname{interval}(a, c))) \end{array}}}

\vnbfile{theorem-library/matact-summand-type-le-proof.scm}{theorem-library/matact-summand-type-le-proof.scm}
\noindent{\small matact-summand-type-le-proof.scm -- matact-summand-type-le, GUARDED and PROVEN.}\par
\vnbentry{\texttt{matact-\allowbreak{}summand-\allowbreak{}type-\allowbreak{}le-\allowbreak{}guarded}}{\vnbstmt{\begin{array}{@{}l@{}} \forall d, m, n, q, p, u, i, c, k. \\ \quad \operatorname{is\text{-}module}(d) \wedge \\ \quad (p \in \operatorname{mat}(m, n, \operatorname{carr}(\operatorname{scal}(d)))) \wedge \\ \quad (u \in \operatorname{mat}(n, q, \operatorname{vec}(d))) \wedge \\ \quad (i \in \operatorname{interval}(1, m)) \wedge \\ \quad (c \in \operatorname{interval}(1, q)) \wedge \\ \quad (k \in \mathbb{N}) \wedge \\ \quad (k \leq n) \Rightarrow \\ \quad \quad (\lambda j \in \operatorname{interval}(1, k).\; \\ \quad \quad \quad \operatorname{act}(d)(\operatorname{entry}(p, i, j), \operatorname{entry}(u, j, c))) \\ \quad \quad \in\; (\operatorname{interval}(1, k) \to \operatorname{carr}(\operatorname{module\text{-}vector\text{-}ag}(d))) \end{array}}}

\vnbfile{theorem-library/border-mult-proof.scm}{theorem-library/border-mult-proof.scm}
\noindent{\small border-mult-proof.scm -- Brick B3 (LA Phase B): block-matrix multiplication. BORDER(A,b,X,p,q) . BORDER(A,d,Y,q,r) = BORDER(A, b*d, X.Y, p, r)   [trust:none] The [[b,0],[0,X]] . [[d,0],[0,Y]] = [[bd,0],[0,XY]] direct-sum law -- the crux that lets the Smith recursion border a block equivalence back to the full}\par
\vnbentry{\texttt{border-\allowbreak{}mult}}{\vnbstmt{\begin{array}{@{}l@{}} \forall a, p, q, r, x, y, b, d. \\ \quad \operatorname{is\text{-}ring}(a) \wedge \\ \quad (p \in \mathbb{N}) \wedge \\ \quad (q \in \mathbb{N}) \wedge \\ \quad (r \in \mathbb{N}) \wedge \\ \quad (x \in \operatorname{mat}(p, q, \operatorname{carr}(a))) \wedge \\ \quad (y \in \operatorname{mat}(q, r, \operatorname{carr}(a))) \wedge \\ \quad (b \in \operatorname{carr}(a)) \wedge \\ \quad (d \in \operatorname{carr}(a)) \wedge \\ \quad ((q = 0) \Rightarrow ((p = 0) \vee (r = 0))) \Rightarrow \\ \quad \quad \operatorname{matmul}(a, \operatorname{border}(a, b, x, p, q), \operatorname{border}(a, d, y, q, r)) \\ \quad \quad =\; \operatorname{border}(a, \operatorname{mul}(a)(b, d), \operatorname{matmul}(a, x, y), p, r) \end{array}}}

\vnbfile{theorem-library/border-assembly-proof.scm}{theorem-library/border-assembly-proof.scm}
\vnbentry{\texttt{border-\allowbreak{}identity}}{\vnbstmt{\begin{array}{@{}l@{}} \forall a.; \forall p \in \mathbb{N}. \\ \quad \operatorname{is\text{-}ring}(a) \Rightarrow  \\ \quad \quad \operatorname{border}(a, \operatorname{one}(a), \operatorname{identmat}(a, p), p, p) \\ \quad \quad =\; \operatorname{identmat}(a, (p + 1)) \end{array}}}
\vnbentry{\texttt{border-\allowbreak{}is-\allowbreak{}diagonal}}{\vnbstmt{\begin{array}{@{}l@{}} \forall a, b, e, p, q. \\ \quad \operatorname{is\text{-}ring}(a) \wedge \\ \quad (p \in \mathbb{N}) \wedge \\ \quad (q \in \mathbb{N}) \wedge \\ \quad \operatorname{is\text{-}diagonal}(a, p, q, e) \wedge \\ \quad (b \in \operatorname{carr}(a)) \wedge \\ \quad (e \in \operatorname{mat}(p, q, \operatorname{carr}(a))) \Rightarrow \\ \quad \quad \operatorname{is\text{-}diagonal}(a, (p + 1), (q + 1), \operatorname{border}(a, b, e, p, q)) \end{array}}}
\vnbentry{\texttt{border-\allowbreak{}invertible}}{\vnbstmt{\begin{array}{@{}l@{}} \forall a, n, u. \\ \quad \operatorname{is\text{-}ring}(a) \wedge \\ \quad (n \in \mathbb{N}) \wedge \\ \quad \operatorname{is\text{-}invertible\text{-}mat}(a, n, u) \Rightarrow \\ \quad \quad \operatorname{is\text{-}invertible\text{-}mat}(a, (n + 1), \operatorname{border}(a, \operatorname{one}(a), u, n, n)) \end{array}}}
\vnbentry{\texttt{bordering}}{\vnbstmt{\begin{array}{@{}l@{}} \forall a, p, q, b, w, g, e. \\ \quad \operatorname{is\text{-}ring}(a) \wedge \\ \quad (p \in \mathbb{N}) \wedge \\ \quad (q \in \mathbb{N}) \wedge \\ \quad (b \in \operatorname{carr}(a)) \wedge \\ \quad (w \in \operatorname{mat}(p, q, \operatorname{carr}(a))) \wedge \\ \quad (g \in \operatorname{mat}((p + 1), (q + 1), \operatorname{carr}(a))) \wedge \\ \quad (g = \operatorname{border}(a, b, w, p, q)) \wedge \\ \quad \operatorname{mat\text{-}equiv}(a, p, q, w, e) \Rightarrow \\ \quad \quad \operatorname{mat\text{-}equiv}(a, (p + 1), (q + 1), g, \operatorname{border}(a, b, e, p, q)) \end{array}}}

\vnbfile{theorem-library/bordered-eq-border-proof.scm}{theorem-library/bordered-eq-border-proof.scm}
\noindent{\small bordered-eq-border -- a matrix C whose first row and first column are cleared off the pivot IS literally BORDER(C\_\{1,1\}, SUBMAT(C)): C = [[C11,0],[0,C']]. The bridge from "cleared cross" (clear-first-row + clear-first-col) to the BORDER block form, so the Smith recursion can border a sub-block equivalence}\par
\vnbentry{\texttt{bordered-\allowbreak{}eq-\allowbreak{}border}}{\vnbstmt{\begin{array}{@{}l@{}} \forall a, p, q, c. \\ \quad \operatorname{is\text{-}ring}(a) \wedge \\ \quad (p \in \mathbb{N}) \wedge \\ \quad (q \in \mathbb{N}) \wedge \\ \quad (c \in \operatorname{mat}((p + 1), (q + 1), \operatorname{carr}(a))) \wedge \\ \quad \forall j \in \operatorname{interval}(1, (q + 1)). \\ \quad \quad \neg (j = 1) \Rightarrow  \\ \quad \quad \quad (\operatorname{entry}(c, 1, j) = \operatorname{zero}(a)) \wedge \\ \quad \forall i \in \operatorname{interval}(1, (p + 1)). \\ \quad \quad \neg (i = 1) \Rightarrow  \\ \quad \quad \quad (\operatorname{entry}(c, i, 1) = \operatorname{zero}(a)) \Rightarrow \\ \quad \quad (c = \operatorname{border}(a, \operatorname{entry}(c, 1, 1), \operatorname{submat}(c, p, q), p, q)) \end{array}}}

\vnbfile{theorem-library/det-alternating-form.scm}{theorem-library/det-alternating-form.scm}
\noindent{\small det-alternating-form.scm -- Hoffman-Kunze 5.3 Theorem 2 (batch 26-B, 2026-09-24): an ALTERNATING n-LINEAR function D on the n x n matrices over a commutative ring R satisfies  D(A) = det(A) D(I)  for every A.  Proven WITHOUT permutations.}\par
\vnbentry{\texttt{daf-\allowbreak{}ring-\allowbreak{}rdist}}{\vnbstmt{\begin{array}{@{}l@{}} \forall r, s, x, y, z. \\ \quad \operatorname{is\text{-}commutative\text{-}ring}(r) \wedge \\ \quad (s \in \operatorname{carr}(r)) \wedge \\ \quad (x \in \operatorname{carr}(r)) \wedge \\ \quad (y \in \operatorname{carr}(r)) \wedge \\ \quad (z \in \operatorname{carr}(r)) \Rightarrow \\ \quad \quad \operatorname{mul}(r)(\operatorname{add}(r)(\operatorname{mul}(r)(s, x), y), z) \\ \quad \quad =\; \operatorname{add}(r)(\operatorname{mul}(r)(s, \operatorname{mul}(r)(x, z)), \operatorname{mul}(r)(y, z)) \end{array}}}
\vnbentry{\texttt{daf-\allowbreak{}ring-\allowbreak{}cancel0}}{\vnbstmt{\begin{array}{@{}l@{}} \forall r, s, x, y. \\ \quad \operatorname{is\text{-}commutative\text{-}ring}(r) \wedge \\ \quad (s \in \operatorname{carr}(r)) \wedge \\ \quad (x \in \operatorname{carr}(r)) \wedge \\ \quad (y \in \operatorname{carr}(r)) \wedge \\ \quad (x = \operatorname{add}(r)(\operatorname{mul}(r)(s, \operatorname{zero}(r)), y)) \Rightarrow \\ \quad \quad (y = x) \end{array}}}
\vnbentry{\texttt{daf-\allowbreak{}det-\allowbreak{}fun}}{\vnbstmt{\begin{array}{@{}l@{}} \forall r, n. \\ \quad (\operatorname{is\text{-}commutative\text{-}ring}(r) \wedge (n \in \mathbb{N})) \Rightarrow  \\ \quad \quad (\lambda l \in \operatorname{mat}(n, n, \operatorname{carr}(r)).\; \operatorname{det}(r, n, l)) \\ \quad \quad \in\; (\operatorname{mat}(n, n, \operatorname{carr}(r)) \to \operatorname{carr}(r)) \end{array}}}
\vnbentry{\texttt{daf-\allowbreak{}det-\allowbreak{}lin}}{\vnbstmt{\begin{array}{@{}l@{}} \forall r, n, a, b, c, p, s. \\ \quad \operatorname{is\text{-}commutative\text{-}ring}(r) \wedge \\ \quad (n \in \mathbb{N}) \wedge \\ \quad (a \in \operatorname{mat}(n, n, \operatorname{carr}(r))) \wedge \\ \quad (b \in \operatorname{mat}(n, n, \operatorname{carr}(r))) \wedge \\ \quad (c \in \operatorname{mat}(n, n, \operatorname{carr}(r))) \wedge \\ \quad (p \in \operatorname{interval}(1, n)) \wedge \\ \quad (s \in \operatorname{carr}(r)) \wedge \\ \quad \forall i, k \in \operatorname{interval}(1, n). \\ \quad \quad \neg (i = p) \Rightarrow  \\ \quad \quad \quad (\operatorname{entry}(b, i, k) = \operatorname{entry}(a, i, k)) \wedge \\ \quad \forall i, k \in \operatorname{interval}(1, n). \\ \quad \quad \neg (i = p) \Rightarrow  \\ \quad \quad \quad (\operatorname{entry}(c, i, k) = \operatorname{entry}(a, i, k)) \wedge \\ \quad \forall k \in \operatorname{interval}(1, n). \\ \quad \quad \operatorname{entry}(c, p, k) \\ \quad \quad =\; \operatorname{add}(r)(\operatorname{mul}(r)(s, \operatorname{entry}(a, p, k)), \operatorname{entry}(b, p, k)) \Rightarrow \\ \quad \quad (\lambda l \in \operatorname{mat}(n, n, \operatorname{carr}(r)).\; \operatorname{det}(r, n, l))(c) \\ \quad \quad =\; \operatorname{add}(r)(\operatorname{mul}(r)(s, (\lambda l \in \operatorname{mat}(n, n, \operatorname{carr}(r)).\; \operatorname{det}(r, n, l))(a)), \\ \quad \quad \quad (\lambda l \in \operatorname{mat}(n, n, \operatorname{carr}(r)).\; \operatorname{det}(r, n, l))(b)) \end{array}}}
\vnbentry{\texttt{daf-\allowbreak{}det-\allowbreak{}alt}}{\vnbstmt{\begin{array}{@{}l@{}} \forall r, n, a, p, q. \\ \quad \operatorname{is\text{-}commutative\text{-}ring}(r) \wedge \\ \quad (n \in \mathbb{N}) \wedge \\ \quad (a \in \operatorname{mat}(n, n, \operatorname{carr}(r))) \wedge \\ \quad (p \in \operatorname{interval}(1, n)) \wedge \\ \quad (q \in \operatorname{interval}(1, n)) \wedge \\ \quad \neg (p = q) \wedge \\ \quad \forall k \in \operatorname{interval}(1, n). \\ \quad \quad (\operatorname{entry}(a, p, k) = \operatorname{entry}(a, q, k)) \Rightarrow \\ \quad \quad (\lambda l \in \operatorname{mat}(n, n, \operatorname{carr}(r)).\; \operatorname{det}(r, n, l))(a) \\ \quad \quad =\; \operatorname{zero}(r) \end{array}}}
\vnbentry{\texttt{daf-\allowbreak{}mm-\allowbreak{}row-\allowbreak{}eq}}{\vnbstmt{\begin{array}{@{}l@{}} \forall r, n, x, y, q, i, j, k. \\ \quad \operatorname{is\text{-}commutative\text{-}ring}(r) \wedge \\ \quad (n \in \mathbb{N}) \wedge \\ \quad (x \in \operatorname{mat}(n, n, \operatorname{carr}(r))) \wedge \\ \quad (y \in \operatorname{mat}(n, n, \operatorname{carr}(r))) \wedge \\ \quad (q \in \operatorname{mat}(n, n, \operatorname{carr}(r))) \wedge \\ \quad (i \in \operatorname{interval}(1, n)) \wedge \\ \quad (j \in \operatorname{interval}(1, n)) \wedge \\ \quad (k \in \operatorname{interval}(1, n)) \wedge \\ \quad \forall \mathit{drk} \in \operatorname{interval}(1, n). \\ \quad \quad (\operatorname{entry}(x, i, \mathit{drk}) = \operatorname{entry}(y, j, \mathit{drk})) \Rightarrow \\ \quad \quad \operatorname{entry}(\operatorname{matmul}(r, x, q), i, k) \\ \quad \quad =\; \operatorname{entry}(\operatorname{matmul}(r, y, q), j, k) \end{array}}}
\vnbentry{\texttt{daf-\allowbreak{}mm-\allowbreak{}row-\allowbreak{}lin}}{\vnbstmt{\begin{array}{@{}l@{}} \forall r, n, a, b, c, q, p, s, k. \\ \quad \operatorname{is\text{-}commutative\text{-}ring}(r) \wedge \\ \quad (n \in \mathbb{N}) \wedge \\ \quad (a \in \operatorname{mat}(n, n, \operatorname{carr}(r))) \wedge \\ \quad (b \in \operatorname{mat}(n, n, \operatorname{carr}(r))) \wedge \\ \quad (c \in \operatorname{mat}(n, n, \operatorname{carr}(r))) \wedge \\ \quad (q \in \operatorname{mat}(n, n, \operatorname{carr}(r))) \wedge \\ \quad (p \in \operatorname{interval}(1, n)) \wedge \\ \quad (s \in \operatorname{carr}(r)) \wedge \\ \quad (k \in \operatorname{interval}(1, n)) \wedge \\ \quad \forall \mathit{drk} \in \operatorname{interval}(1, n). \\ \quad \quad \operatorname{entry}(c, p, \mathit{drk}) \\ \quad \quad =\; \operatorname{add}(r)(\operatorname{mul}(r)(s, \operatorname{entry}(a, p, \mathit{drk})), \operatorname{entry}(b, p, \mathit{drk})) \Rightarrow \\ \quad \quad \operatorname{entry}(\operatorname{matmul}(r, c, q), p, k) \\ \quad \quad =\; \operatorname{add}(r)(\operatorname{mul}(r)(s, \operatorname{entry}(\operatorname{matmul}(r, a, q), p, k)), \\ \quad \quad \quad \operatorname{entry}(\operatorname{matmul}(r, b, q), p, k)) \end{array}}}
\vnbentry{\texttt{daf-\allowbreak{}rm-\allowbreak{}fun}}{\vnbstmt{\begin{array}{@{}l@{}} \forall r, n, d, q. \\ \quad \operatorname{is\text{-}commutative\text{-}ring}(r) \wedge \\ \quad (n \in \mathbb{N}) \wedge \\ \quad (d \in (\operatorname{mat}(n, n, \operatorname{carr}(r)) \to \operatorname{carr}(r))) \wedge \\ \quad (q \in \operatorname{mat}(n, n, \operatorname{carr}(r))) \Rightarrow \\ \quad \quad (\lambda l \in \operatorname{mat}(n, n, \operatorname{carr}(r)).\; \\ \quad \quad \quad d(\operatorname{matmul}(r, l, q))) \\ \quad \quad \in\; (\operatorname{mat}(n, n, \operatorname{carr}(r)) \to \operatorname{carr}(r)) \end{array}}}
\vnbentry{\texttt{daf-\allowbreak{}rm-\allowbreak{}lin}}{\vnbstmt{\begin{array}{@{}l@{}} \forall r, n, d, q, a, b, c, p, s. \\ \quad \operatorname{is\text{-}commutative\text{-}ring}(r) \wedge \\ \quad (n \in \mathbb{N}) \wedge \\ \quad (d \in (\operatorname{mat}(n, n, \operatorname{carr}(r)) \to \operatorname{carr}(r))) \wedge \\ \quad \forall a, b, c, p, s. \\ \quad \quad (a \in \operatorname{mat}(n, n, \operatorname{carr}(r))) \wedge \\ \quad \quad (b \in \operatorname{mat}(n, n, \operatorname{carr}(r))) \wedge \\ \quad \quad (c \in \operatorname{mat}(n, n, \operatorname{carr}(r))) \wedge \\ \quad \quad (p \in \operatorname{interval}(1, n)) \wedge \\ \quad \quad (s \in \operatorname{carr}(r)) \wedge \\ \quad \quad \forall i, k \in \operatorname{interval}(1, n). \\ \quad \quad \quad \neg (i = p) \Rightarrow  \\ \quad \quad \quad \quad (\operatorname{entry}(b, i, k) = \operatorname{entry}(a, i, k)) \wedge \\ \quad \quad \forall i, k \in \operatorname{interval}(1, n). \\ \quad \quad \quad \neg (i = p) \Rightarrow  \\ \quad \quad \quad \quad (\operatorname{entry}(c, i, k) = \operatorname{entry}(a, i, k)) \wedge \\ \quad \quad \forall k \in \operatorname{interval}(1, n). \\ \quad \quad \quad \operatorname{entry}(c, p, k) \\ \quad \quad \quad =\; \operatorname{add}(r)(\operatorname{mul}(r)(s, \operatorname{entry}(a, p, k)), \operatorname{entry}(b, p, k)) \Rightarrow \\ \quad \quad \quad (d(c) = \operatorname{add}(r)(\operatorname{mul}(r)(s, d(a)), d(b))) \wedge \\ \quad (q \in \operatorname{mat}(n, n, \operatorname{carr}(r))) \wedge \\ \quad (a \in \operatorname{mat}(n, n, \operatorname{carr}(r))) \wedge \\ \quad (b \in \operatorname{mat}(n, n, \operatorname{carr}(r))) \wedge \\ \quad (c \in \operatorname{mat}(n, n, \operatorname{carr}(r))) \wedge \\ \quad (p \in \operatorname{interval}(1, n)) \wedge \\ \quad (s \in \operatorname{carr}(r)) \wedge \\ \quad \forall i, k \in \operatorname{interval}(1, n). \\ \quad \quad \neg (i = p) \Rightarrow  \\ \quad \quad \quad (\operatorname{entry}(b, i, k) = \operatorname{entry}(a, i, k)) \wedge \\ \quad \forall i, k \in \operatorname{interval}(1, n). \\ \quad \quad \neg (i = p) \Rightarrow  \\ \quad \quad \quad (\operatorname{entry}(c, i, k) = \operatorname{entry}(a, i, k)) \wedge \\ \quad \forall k \in \operatorname{interval}(1, n). \\ \quad \quad \operatorname{entry}(c, p, k) \\ \quad \quad =\; \operatorname{add}(r)(\operatorname{mul}(r)(s, \operatorname{entry}(a, p, k)), \operatorname{entry}(b, p, k)) \Rightarrow \\ \quad \quad (\lambda l \in \operatorname{mat}(n, n, \operatorname{carr}(r)).\; d(\operatorname{matmul}(r, l, q)))(c) \\ \quad \quad =\; \operatorname{add}(r)(\operatorname{mul}(r)(s, \\ \quad \quad \quad \quad (\lambda l \in \operatorname{mat}(n, n, \operatorname{carr}(r)).\; d(\operatorname{matmul}(r, l, q)))(a)), \\ \quad \quad \quad (\lambda l \in \operatorname{mat}(n, n, \operatorname{carr}(r)).\; d(\operatorname{matmul}(r, l, q)))(b)) \end{array}}}
\vnbentry{\texttt{daf-\allowbreak{}rm-\allowbreak{}alt}}{\vnbstmt{\begin{array}{@{}l@{}} \forall r, n, d, q, a, p, \mathit{dfq}. \\ \quad \operatorname{is\text{-}commutative\text{-}ring}(r) \wedge \\ \quad (n \in \mathbb{N}) \wedge \\ \quad (d \in (\operatorname{mat}(n, n, \operatorname{carr}(r)) \to \operatorname{carr}(r))) \wedge \\ \quad \forall a, p, \mathit{dfq}. \\ \quad \quad (a \in \operatorname{mat}(n, n, \operatorname{carr}(r))) \wedge \\ \quad \quad (p \in \operatorname{interval}(1, n)) \wedge \\ \quad \quad (\mathit{dfq} \in \operatorname{interval}(1, n)) \wedge \\ \quad \quad \neg (p = \mathit{dfq}) \wedge \\ \quad \quad \forall k \in \operatorname{interval}(1, n). \\ \quad \quad \quad (\operatorname{entry}(a, p, k) = \operatorname{entry}(a, \mathit{dfq}, k)) \Rightarrow \\ \quad \quad \quad (d(a) = \operatorname{zero}(r)) \wedge \\ \quad (q \in \operatorname{mat}(n, n, \operatorname{carr}(r))) \wedge \\ \quad (a \in \operatorname{mat}(n, n, \operatorname{carr}(r))) \wedge \\ \quad (p \in \operatorname{interval}(1, n)) \wedge \\ \quad (\mathit{dfq} \in \operatorname{interval}(1, n)) \wedge \\ \quad \neg (p = \mathit{dfq}) \wedge \\ \quad \forall k \in \operatorname{interval}(1, n). \\ \quad \quad (\operatorname{entry}(a, p, k) = \operatorname{entry}(a, \mathit{dfq}, k)) \Rightarrow \\ \quad \quad (\lambda l \in \operatorname{mat}(n, n, \operatorname{carr}(r)).\; d(\operatorname{matmul}(r, l, q)))(a) \\ \quad \quad =\; \operatorname{zero}(r) \end{array}}}
\vnbentry{\texttt{daf-\allowbreak{}det-\allowbreak{}rm-\allowbreak{}fun}}{\vnbstmt{\begin{array}{@{}l@{}} \forall r, n, q. \\ \quad \operatorname{is\text{-}commutative\text{-}ring}(r) \wedge \\ \quad (n \in \mathbb{N}) \wedge \\ \quad (q \in \operatorname{mat}(n, n, \operatorname{carr}(r))) \Rightarrow \\ \quad \quad (\lambda l \in \operatorname{mat}(n, n, \operatorname{carr}(r)).\; \\ \quad \quad \quad \operatorname{det}(r, n, \operatorname{matmul}(r, l, q))) \\ \quad \quad \in\; (\operatorname{mat}(n, n, \operatorname{carr}(r)) \to \operatorname{carr}(r)) \end{array}}}
\vnbentry{\texttt{daf-\allowbreak{}det-\allowbreak{}rm-\allowbreak{}lin}}{\vnbstmt{\begin{array}{@{}l@{}} \forall r, n, q, a, b, c, p, s. \\ \quad \operatorname{is\text{-}commutative\text{-}ring}(r) \wedge \\ \quad (n \in \mathbb{N}) \wedge \\ \quad (q \in \operatorname{mat}(n, n, \operatorname{carr}(r))) \wedge \\ \quad (a \in \operatorname{mat}(n, n, \operatorname{carr}(r))) \wedge \\ \quad (b \in \operatorname{mat}(n, n, \operatorname{carr}(r))) \wedge \\ \quad (c \in \operatorname{mat}(n, n, \operatorname{carr}(r))) \wedge \\ \quad (p \in \operatorname{interval}(1, n)) \wedge \\ \quad (s \in \operatorname{carr}(r)) \wedge \\ \quad \forall i, k \in \operatorname{interval}(1, n). \\ \quad \quad \neg (i = p) \Rightarrow  \\ \quad \quad \quad (\operatorname{entry}(b, i, k) = \operatorname{entry}(a, i, k)) \wedge \\ \quad \forall i, k \in \operatorname{interval}(1, n). \\ \quad \quad \neg (i = p) \Rightarrow  \\ \quad \quad \quad (\operatorname{entry}(c, i, k) = \operatorname{entry}(a, i, k)) \wedge \\ \quad \forall k \in \operatorname{interval}(1, n). \\ \quad \quad \operatorname{entry}(c, p, k) \\ \quad \quad =\; \operatorname{add}(r)(\operatorname{mul}(r)(s, \operatorname{entry}(a, p, k)), \operatorname{entry}(b, p, k)) \Rightarrow \\ \quad \quad (\lambda l \in \operatorname{mat}(n, n, \operatorname{carr}(r)).\; \operatorname{det}(r, n, \operatorname{matmul}(r, l, q)))(c) \\ \quad \quad =\; \operatorname{add}(r)(\operatorname{mul}(r)(s, \\ \quad \quad \quad \quad (\lambda l \in \operatorname{mat}(n, n, \operatorname{carr}(r)).\; \operatorname{det}(r, n, \operatorname{matmul}(r, l, q)))(a)), \\ \quad \quad \quad (\lambda l \in \operatorname{mat}(n, n, \operatorname{carr}(r)).\; \operatorname{det}(r, n, \operatorname{matmul}(r, l, q)))(b)) \end{array}}}
\vnbentry{\texttt{daf-\allowbreak{}det-\allowbreak{}rm-\allowbreak{}alt}}{\vnbstmt{\begin{array}{@{}l@{}} \forall r, n, q, a, p, \mathit{dfq}. \\ \quad \operatorname{is\text{-}commutative\text{-}ring}(r) \wedge \\ \quad (n \in \mathbb{N}) \wedge \\ \quad (q \in \operatorname{mat}(n, n, \operatorname{carr}(r))) \wedge \\ \quad (a \in \operatorname{mat}(n, n, \operatorname{carr}(r))) \wedge \\ \quad (p \in \operatorname{interval}(1, n)) \wedge \\ \quad (\mathit{dfq} \in \operatorname{interval}(1, n)) \wedge \\ \quad \neg (p = \mathit{dfq}) \wedge \\ \quad \forall k \in \operatorname{interval}(1, n). \\ \quad \quad (\operatorname{entry}(a, p, k) = \operatorname{entry}(a, \mathit{dfq}, k)) \Rightarrow \\ \quad \quad (\lambda l \in \operatorname{mat}(n, n, \operatorname{carr}(r)).\; \operatorname{det}(r, n, \operatorname{matmul}(r, l, q)))(a) \\ \quad \quad =\; \operatorname{zero}(r) \end{array}}}
\vnbentry{\texttt{daf-\allowbreak{}zero-\allowbreak{}row}}{\vnbstmt{\begin{array}{@{}l@{}} \forall r, n, d, a, k. \\ \quad \operatorname{is\text{-}commutative\text{-}ring}(r) \wedge \\ \quad (n \in \mathbb{N}) \wedge \\ \quad (d \in (\operatorname{mat}(n, n, \operatorname{carr}(r)) \to \operatorname{carr}(r))) \wedge \\ \quad \forall \mathit{dfa}, b, c, p, s. \\ \quad \quad (\mathit{dfa} \in \operatorname{mat}(n, n, \operatorname{carr}(r))) \wedge \\ \quad \quad (b \in \operatorname{mat}(n, n, \operatorname{carr}(r))) \wedge \\ \quad \quad (c \in \operatorname{mat}(n, n, \operatorname{carr}(r))) \wedge \\ \quad \quad (p \in \operatorname{interval}(1, n)) \wedge \\ \quad \quad (s \in \operatorname{carr}(r)) \wedge \\ \quad \quad \forall i, \mathit{drk} \in \operatorname{interval}(1, n). \\ \quad \quad \quad \neg (i = p) \Rightarrow  \\ \quad \quad \quad \quad (\operatorname{entry}(b, i, \mathit{drk}) = \operatorname{entry}(\mathit{dfa}, i, \mathit{drk})) \wedge \\ \quad \quad \forall i, \mathit{drk} \in \operatorname{interval}(1, n). \\ \quad \quad \quad \neg (i = p) \Rightarrow  \\ \quad \quad \quad \quad (\operatorname{entry}(c, i, \mathit{drk}) = \operatorname{entry}(\mathit{dfa}, i, \mathit{drk})) \wedge \\ \quad \quad \forall \mathit{drk} \in \operatorname{interval}(1, n). \\ \quad \quad \quad \operatorname{entry}(c, p, \mathit{drk}) \\ \quad \quad \quad =\; \operatorname{add}(r)(\operatorname{mul}(r)(s, \operatorname{entry}(\mathit{dfa}, p, \mathit{drk})), \operatorname{entry}(b, p, \mathit{drk})) \Rightarrow \\ \quad \quad \quad (d(c) = \operatorname{add}(r)(\operatorname{mul}(r)(s, d(\mathit{dfa})), d(b))) \wedge \\ \quad (a \in \operatorname{mat}(n, n, \operatorname{carr}(r))) \wedge \\ \quad (k \in \operatorname{interval}(1, n)) \wedge \\ \quad \forall \mathit{drk} \in \operatorname{interval}(1, n). \\ \quad \quad (\operatorname{entry}(a, k, \mathit{drk}) = \operatorname{zero}(r)) \Rightarrow \\ \quad \quad (d(a) = \operatorname{zero}(r)) \end{array}}}
\vnbentry{\texttt{daf-\allowbreak{}bd-\allowbreak{}fun}}{\vnbstmt{\begin{array}{@{}l@{}} \forall r, n, d. \\ \quad \operatorname{is\text{-}commutative\text{-}ring}(r) \wedge \\ \quad (n \in \mathbb{N}) \wedge \\ \quad (d \in (\operatorname{mat}((n + 1), (n + 1), \operatorname{carr}(r)) \to \operatorname{carr}(r))) \Rightarrow \\ \quad \quad (\lambda l \in \operatorname{mat}(n, n, \operatorname{carr}(r)).\; \\ \quad \quad \quad d(\operatorname{border}(r, \operatorname{one}(r), l, n, n))) \\ \quad \quad \in\; (\operatorname{mat}(n, n, \operatorname{carr}(r)) \to \operatorname{carr}(r)) \end{array}}}
\vnbentry{\texttt{daf-\allowbreak{}bd-\allowbreak{}lin}}{\vnbstmt{\begin{array}{@{}l@{}} \forall r, n, d, a, b, c, p, s. \\ \quad \operatorname{is\text{-}commutative\text{-}ring}(r) \wedge \\ \quad (n \in \mathbb{N}) \wedge \\ \quad (d \in (\operatorname{mat}((n + 1), (n + 1), \operatorname{carr}(r)) \to \operatorname{carr}(r))) \wedge \\ \quad \forall a, b, c, p, s. \\ \quad \quad (a \in \operatorname{mat}((n + 1), (n + 1), \operatorname{carr}(r))) \wedge \\ \quad \quad (b \in \operatorname{mat}((n + 1), (n + 1), \operatorname{carr}(r))) \wedge \\ \quad \quad (c \in \operatorname{mat}((n + 1), (n + 1), \operatorname{carr}(r))) \wedge \\ \quad \quad (p \in \operatorname{interval}(1, (n + 1))) \wedge \\ \quad \quad (s \in \operatorname{carr}(r)) \wedge \\ \quad \quad \forall i, k \in \operatorname{interval}(1, (n + 1)). \\ \quad \quad \quad \neg (i = p) \Rightarrow  \\ \quad \quad \quad \quad (\operatorname{entry}(b, i, k) = \operatorname{entry}(a, i, k)) \wedge \\ \quad \quad \forall i, k \in \operatorname{interval}(1, (n + 1)). \\ \quad \quad \quad \neg (i = p) \Rightarrow  \\ \quad \quad \quad \quad (\operatorname{entry}(c, i, k) = \operatorname{entry}(a, i, k)) \wedge \\ \quad \quad \forall k \in \operatorname{interval}(1, (n + 1)). \\ \quad \quad \quad \operatorname{entry}(c, p, k) \\ \quad \quad \quad =\; \operatorname{add}(r)(\operatorname{mul}(r)(s, \operatorname{entry}(a, p, k)), \operatorname{entry}(b, p, k)) \Rightarrow \\ \quad \quad \quad (d(c) = \operatorname{add}(r)(\operatorname{mul}(r)(s, d(a)), d(b))) \wedge \\ \quad (a \in \operatorname{mat}(n, n, \operatorname{carr}(r))) \wedge \\ \quad (b \in \operatorname{mat}(n, n, \operatorname{carr}(r))) \wedge \\ \quad (c \in \operatorname{mat}(n, n, \operatorname{carr}(r))) \wedge \\ \quad (p \in \operatorname{interval}(1, n)) \wedge \\ \quad (s \in \operatorname{carr}(r)) \wedge \\ \quad \forall i, k \in \operatorname{interval}(1, n). \\ \quad \quad \neg (i = p) \Rightarrow  \\ \quad \quad \quad (\operatorname{entry}(b, i, k) = \operatorname{entry}(a, i, k)) \wedge \\ \quad \forall i, k \in \operatorname{interval}(1, n). \\ \quad \quad \neg (i = p) \Rightarrow  \\ \quad \quad \quad (\operatorname{entry}(c, i, k) = \operatorname{entry}(a, i, k)) \wedge \\ \quad \forall k \in \operatorname{interval}(1, n). \\ \quad \quad \operatorname{entry}(c, p, k) \\ \quad \quad =\; \operatorname{add}(r)(\operatorname{mul}(r)(s, \operatorname{entry}(a, p, k)), \operatorname{entry}(b, p, k)) \Rightarrow \\ \quad \quad (\lambda l \in \operatorname{mat}(n, n, \operatorname{carr}(r)).\; d(\operatorname{border}(r, \operatorname{one}(r), l, n, n)))(c) \\ \quad \quad =\; \operatorname{add}(r)(\operatorname{mul}(r)(s, \\ \quad \quad \quad \quad (\lambda l \in \operatorname{mat}(n, n, \operatorname{carr}(r)).\; d(\operatorname{border}(r, \operatorname{one}(r), l, n, n)))(a)), \\ \quad \quad \quad (\lambda l \in \operatorname{mat}(n, n, \operatorname{carr}(r)).\; d(\operatorname{border}(r, \operatorname{one}(r), l, n, n)))(b)) \end{array}}}
\vnbentry{\texttt{daf-\allowbreak{}bd-\allowbreak{}alt}}{\vnbstmt{\begin{array}{@{}l@{}} \forall r, n, d, a, p, q. \\ \quad \operatorname{is\text{-}commutative\text{-}ring}(r) \wedge \\ \quad (n \in \mathbb{N}) \wedge \\ \quad (d \in (\operatorname{mat}((n + 1), (n + 1), \operatorname{carr}(r)) \to \operatorname{carr}(r))) \wedge \\ \quad \forall a, p, q. \\ \quad \quad (a \in \operatorname{mat}((n + 1), (n + 1), \operatorname{carr}(r))) \wedge \\ \quad \quad (p \in \operatorname{interval}(1, (n + 1))) \wedge \\ \quad \quad (q \in \operatorname{interval}(1, (n + 1))) \wedge \\ \quad \quad \neg (p = q) \wedge \\ \quad \quad \forall k \in \operatorname{interval}(1, (n + 1)). \\ \quad \quad \quad (\operatorname{entry}(a, p, k) = \operatorname{entry}(a, q, k)) \Rightarrow \\ \quad \quad \quad (d(a) = \operatorname{zero}(r)) \wedge \\ \quad (a \in \operatorname{mat}(n, n, \operatorname{carr}(r))) \wedge \\ \quad (p \in \operatorname{interval}(1, n)) \wedge \\ \quad (q \in \operatorname{interval}(1, n)) \wedge \\ \quad \neg (p = q) \wedge \\ \quad \forall k \in \operatorname{interval}(1, n). \\ \quad \quad (\operatorname{entry}(a, p, k) = \operatorname{entry}(a, q, k)) \Rightarrow \\ \quad \quad (\lambda l \in \operatorname{mat}(n, n, \operatorname{carr}(r)).\; d(\operatorname{border}(r, \operatorname{one}(r), l, n, n)))(a) \\ \quad \quad =\; \operatorname{zero}(r) \end{array}}}
\vnbentry{\texttt{daf-\allowbreak{}swap}}{\vnbstmt{\begin{array}{@{}l@{}} \forall r, n, d, a, b, p, q. \\ \quad \operatorname{is\text{-}commutative\text{-}ring}(r) \wedge \\ \quad (n \in \mathbb{N}) \wedge \\ \quad (d \in (\operatorname{mat}(n, n, \operatorname{carr}(r)) \to \operatorname{carr}(r))) \wedge \\ \quad \forall \mathit{dfa}, \mathit{dfb}, c, \mathit{dfp}, s. \\ \quad \quad (\mathit{dfa} \in \operatorname{mat}(n, n, \operatorname{carr}(r))) \wedge \\ \quad \quad (\mathit{dfb} \in \operatorname{mat}(n, n, \operatorname{carr}(r))) \wedge \\ \quad \quad (c \in \operatorname{mat}(n, n, \operatorname{carr}(r))) \wedge \\ \quad \quad (\mathit{dfp} \in \operatorname{interval}(1, n)) \wedge \\ \quad \quad (s \in \operatorname{carr}(r)) \wedge \\ \quad \quad \forall i, k \in \operatorname{interval}(1, n). \\ \quad \quad \quad \neg (i = \mathit{dfp}) \Rightarrow  \\ \quad \quad \quad \quad (\operatorname{entry}(\mathit{dfb}, i, k) = \operatorname{entry}(\mathit{dfa}, i, k)) \wedge \\ \quad \quad \forall i, k \in \operatorname{interval}(1, n). \\ \quad \quad \quad \neg (i = \mathit{dfp}) \Rightarrow  \\ \quad \quad \quad \quad (\operatorname{entry}(c, i, k) = \operatorname{entry}(\mathit{dfa}, i, k)) \wedge \\ \quad \quad \forall k \in \operatorname{interval}(1, n). \\ \quad \quad \quad \operatorname{entry}(c, \mathit{dfp}, k) \\ \quad \quad \quad =\; \operatorname{add}(r)(\operatorname{mul}(r)(s, \operatorname{entry}(\mathit{dfa}, \mathit{dfp}, k)), \operatorname{entry}(\mathit{dfb}, \mathit{dfp}, k)) \Rightarrow \\ \quad \quad \quad (d(c) = \operatorname{add}(r)(\operatorname{mul}(r)(s, d(\mathit{dfa})), d(\mathit{dfb}))) \wedge \\ \quad \forall \mathit{dfa}, \mathit{dfp}, \mathit{dfq}. \\ \quad \quad (\mathit{dfa} \in \operatorname{mat}(n, n, \operatorname{carr}(r))) \wedge \\ \quad \quad (\mathit{dfp} \in \operatorname{interval}(1, n)) \wedge \\ \quad \quad (\mathit{dfq} \in \operatorname{interval}(1, n)) \wedge \\ \quad \quad \neg (\mathit{dfp} = \mathit{dfq}) \wedge \\ \quad \quad \forall k \in \operatorname{interval}(1, n). \\ \quad \quad \quad \operatorname{entry}(\mathit{dfa}, \mathit{dfp}, k) \\ \quad \quad \quad =\; \operatorname{entry}(\mathit{dfa}, \mathit{dfq}, k) \Rightarrow \\ \quad \quad \quad (d(\mathit{dfa}) = \operatorname{zero}(r)) \wedge \\ \quad (a \in \operatorname{mat}(n, n, \operatorname{carr}(r))) \wedge \\ \quad (b \in \operatorname{mat}(n, n, \operatorname{carr}(r))) \wedge \\ \quad (p \in \operatorname{interval}(1, n)) \wedge \\ \quad (q \in \operatorname{interval}(1, n)) \wedge \\ \quad \neg (p = q) \wedge \\ \quad \forall i, k \in \operatorname{interval}(1, n). \\ \quad \quad (\neg (i = p) \wedge \neg (i = q)) \Rightarrow  \\ \quad \quad \quad (\operatorname{entry}(b, i, k) = \operatorname{entry}(a, i, k)) \wedge \\ \quad \forall k \in \operatorname{interval}(1, n). \\ \quad \quad (\operatorname{entry}(b, p, k) = \operatorname{entry}(a, q, k)) \wedge \\ \quad \forall k \in \operatorname{interval}(1, n). \\ \quad \quad (\operatorname{entry}(b, q, k) = \operatorname{entry}(a, p, k)) \Rightarrow \\ \quad \quad (d(b) = \operatorname{neg}(r)(d(a))) \end{array}}}
\vnbentry{\texttt{daf-\allowbreak{}elem-\allowbreak{}f}}{\vnbstmt{\begin{array}{@{}l@{}} \forall r, n, d, k, l. \\ \quad \operatorname{is\text{-}commutative\text{-}ring}(r) \wedge \\ \quad (n \in \mathbb{N}) \wedge \\ \quad (d \in (\operatorname{mat}(n, n, \operatorname{carr}(r)) \to \operatorname{carr}(r))) \wedge \\ \quad \forall a, b, c, p, s. \\ \quad \quad (a \in \operatorname{mat}(n, n, \operatorname{carr}(r))) \wedge \\ \quad \quad (b \in \operatorname{mat}(n, n, \operatorname{carr}(r))) \wedge \\ \quad \quad (c \in \operatorname{mat}(n, n, \operatorname{carr}(r))) \wedge \\ \quad \quad (p \in \operatorname{interval}(1, n)) \wedge \\ \quad \quad (s \in \operatorname{carr}(r)) \wedge \\ \quad \quad \forall i, \mathit{drk} \in \operatorname{interval}(1, n). \\ \quad \quad \quad \neg (i = p) \Rightarrow  \\ \quad \quad \quad \quad (\operatorname{entry}(b, i, \mathit{drk}) = \operatorname{entry}(a, i, \mathit{drk})) \wedge \\ \quad \quad \forall i, \mathit{drk} \in \operatorname{interval}(1, n). \\ \quad \quad \quad \neg (i = p) \Rightarrow  \\ \quad \quad \quad \quad (\operatorname{entry}(c, i, \mathit{drk}) = \operatorname{entry}(a, i, \mathit{drk})) \wedge \\ \quad \quad \forall \mathit{drk} \in \operatorname{interval}(1, n). \\ \quad \quad \quad \operatorname{entry}(c, p, \mathit{drk}) \\ \quad \quad \quad =\; \operatorname{add}(r)(\operatorname{mul}(r)(s, \operatorname{entry}(a, p, \mathit{drk})), \operatorname{entry}(b, p, \mathit{drk})) \Rightarrow \\ \quad \quad \quad (d(c) = \operatorname{add}(r)(\operatorname{mul}(r)(s, d(a)), d(b))) \wedge \\ \quad \forall a, p, q. \\ \quad \quad (a \in \operatorname{mat}(n, n, \operatorname{carr}(r))) \wedge \\ \quad \quad (p \in \operatorname{interval}(1, n)) \wedge \\ \quad \quad (q \in \operatorname{interval}(1, n)) \wedge \\ \quad \quad \neg (p = q) \wedge \\ \quad \quad \forall \mathit{drk} \in \operatorname{interval}(1, n). \\ \quad \quad \quad (\operatorname{entry}(a, p, \mathit{drk}) = \operatorname{entry}(a, q, \mathit{drk})) \Rightarrow \\ \quad \quad \quad (d(a) = \operatorname{zero}(r)) \wedge \\ \quad (k \in \operatorname{interval}(1, n)) \wedge \\ \quad (l \in \operatorname{interval}(1, n)) \wedge \\ \quad \neg (k = l) \Rightarrow \\ \quad \quad d(\operatorname{elem\text{-}f}(r, n, k, l)) \\ \quad \quad =\; \operatorname{neg}(r)(d(\operatorname{identmat}(r, n))) \end{array}}}
\vnbentry{\texttt{daf-\allowbreak{}clear-\allowbreak{}ind}}{\vnbstmt{\begin{array}{@{}l@{}} \forall m \in \mathbb{N}, r, n, d, y. \\ \quad \operatorname{is\text{-}commutative\text{-}ring}(r) \wedge \\ \quad (n \in \mathbb{N}) \wedge \\ \quad (d \in (\operatorname{mat}((n + 1), (n + 1), \operatorname{carr}(r)) \to \operatorname{carr}(r))) \wedge \\ \quad \forall a, b, c, p, s. \\ \quad \quad (a \in \operatorname{mat}((n + 1), (n + 1), \operatorname{carr}(r))) \wedge \\ \quad \quad (b \in \operatorname{mat}((n + 1), (n + 1), \operatorname{carr}(r))) \wedge \\ \quad \quad (c \in \operatorname{mat}((n + 1), (n + 1), \operatorname{carr}(r))) \wedge \\ \quad \quad (p \in \operatorname{interval}(1, (n + 1))) \wedge \\ \quad \quad (s \in \operatorname{carr}(r)) \wedge \\ \quad \quad \forall i, k \in \operatorname{interval}(1, (n + 1)). \\ \quad \quad \quad \neg (i = p) \Rightarrow  \\ \quad \quad \quad \quad (\operatorname{entry}(b, i, k) = \operatorname{entry}(a, i, k)) \wedge \\ \quad \quad \forall i, k \in \operatorname{interval}(1, (n + 1)). \\ \quad \quad \quad \neg (i = p) \Rightarrow  \\ \quad \quad \quad \quad (\operatorname{entry}(c, i, k) = \operatorname{entry}(a, i, k)) \wedge \\ \quad \quad \forall k \in \operatorname{interval}(1, (n + 1)). \\ \quad \quad \quad \operatorname{entry}(c, p, k) \\ \quad \quad \quad =\; \operatorname{add}(r)(\operatorname{mul}(r)(s, \operatorname{entry}(a, p, k)), \operatorname{entry}(b, p, k)) \Rightarrow \\ \quad \quad \quad (d(c) = \operatorname{add}(r)(\operatorname{mul}(r)(s, d(a)), d(b))) \wedge \\ \quad \forall a, p, q. \\ \quad \quad (a \in \operatorname{mat}((n + 1), (n + 1), \operatorname{carr}(r))) \wedge \\ \quad \quad (p \in \operatorname{interval}(1, (n + 1))) \wedge \\ \quad \quad (q \in \operatorname{interval}(1, (n + 1))) \wedge \\ \quad \quad \neg (p = q) \wedge \\ \quad \quad \forall k \in \operatorname{interval}(1, (n + 1)). \\ \quad \quad \quad (\operatorname{entry}(a, p, k) = \operatorname{entry}(a, q, k)) \Rightarrow \\ \quad \quad \quad (d(a) = \operatorname{zero}(r)) \wedge \\ \quad (y \in \operatorname{mat}((n + 1), (n + 1), \operatorname{carr}(r))) \wedge \\ \quad (\operatorname{entry}(y, 1, 1) = \operatorname{one}(r)) \wedge \\ \quad \forall k \in \operatorname{interval}(1, (n + 1)). \\ \quad \quad \neg (k = 1) \Rightarrow  \\ \quad \quad \quad (\operatorname{entry}(y, 1, k) = \operatorname{zero}(r)) \wedge \\ \quad (m \leq n) \Rightarrow \\ \quad \quad d(\operatorname{matof}((n + 1), (n + 1), \\ \quad \quad \quad \quad (\lambda \langle \operatorname{dav\_}, \operatorname{daw\_} \rangle \in (\operatorname{interval}(1, (n + 1)) \times \operatorname{interval}(1, (n + 1))).\; \\ \quad \quad \quad \quad \quad \left(\text{if }(\operatorname{daw\_} = 1)\text{ then }\left(\text{if }(\operatorname{dav\_} = 1)\text{ then }\operatorname{entry}(y, \operatorname{dav\_}, \operatorname{daw\_})\text{ else }\left(\text{if }(\operatorname{dav\_} \leq (m + 1))\text{ then }\operatorname{zero}(r)\text{ else }\operatorname{entry}(y, \operatorname{dav\_}, \operatorname{daw\_})\right)\right)\text{ else }\operatorname{entry}(y, \operatorname{dav\_}, \operatorname{daw\_})\right)))) \\ \quad \quad =\; d(y) \end{array}}}
\vnbentry{\texttt{daf-\allowbreak{}clear}}{\vnbstmt{\begin{array}{@{}l@{}} \forall r, n, d, y. \\ \quad \operatorname{is\text{-}commutative\text{-}ring}(r) \wedge \\ \quad (n \in \mathbb{N}) \wedge \\ \quad (d \in (\operatorname{mat}((n + 1), (n + 1), \operatorname{carr}(r)) \to \operatorname{carr}(r))) \wedge \\ \quad \forall a, b, c, p, s. \\ \quad \quad (a \in \operatorname{mat}((n + 1), (n + 1), \operatorname{carr}(r))) \wedge \\ \quad \quad (b \in \operatorname{mat}((n + 1), (n + 1), \operatorname{carr}(r))) \wedge \\ \quad \quad (c \in \operatorname{mat}((n + 1), (n + 1), \operatorname{carr}(r))) \wedge \\ \quad \quad (p \in \operatorname{interval}(1, (n + 1))) \wedge \\ \quad \quad (s \in \operatorname{carr}(r)) \wedge \\ \quad \quad \forall i, k \in \operatorname{interval}(1, (n + 1)). \\ \quad \quad \quad \neg (i = p) \Rightarrow  \\ \quad \quad \quad \quad (\operatorname{entry}(b, i, k) = \operatorname{entry}(a, i, k)) \wedge \\ \quad \quad \forall i, k \in \operatorname{interval}(1, (n + 1)). \\ \quad \quad \quad \neg (i = p) \Rightarrow  \\ \quad \quad \quad \quad (\operatorname{entry}(c, i, k) = \operatorname{entry}(a, i, k)) \wedge \\ \quad \quad \forall k \in \operatorname{interval}(1, (n + 1)). \\ \quad \quad \quad \operatorname{entry}(c, p, k) \\ \quad \quad \quad =\; \operatorname{add}(r)(\operatorname{mul}(r)(s, \operatorname{entry}(a, p, k)), \operatorname{entry}(b, p, k)) \Rightarrow \\ \quad \quad \quad (d(c) = \operatorname{add}(r)(\operatorname{mul}(r)(s, d(a)), d(b))) \wedge \\ \quad \forall a, p, q. \\ \quad \quad (a \in \operatorname{mat}((n + 1), (n + 1), \operatorname{carr}(r))) \wedge \\ \quad \quad (p \in \operatorname{interval}(1, (n + 1))) \wedge \\ \quad \quad (q \in \operatorname{interval}(1, (n + 1))) \wedge \\ \quad \quad \neg (p = q) \wedge \\ \quad \quad \forall k \in \operatorname{interval}(1, (n + 1)). \\ \quad \quad \quad (\operatorname{entry}(a, p, k) = \operatorname{entry}(a, q, k)) \Rightarrow \\ \quad \quad \quad (d(a) = \operatorname{zero}(r)) \wedge \\ \quad (y \in \operatorname{mat}((n + 1), (n + 1), \operatorname{carr}(r))) \wedge \\ \quad (\operatorname{entry}(y, 1, 1) = \operatorname{one}(r)) \wedge \\ \quad \forall k \in \operatorname{interval}(1, (n + 1)). \\ \quad \quad \neg (k = 1) \Rightarrow  \\ \quad \quad \quad (\operatorname{entry}(y, 1, k) = \operatorname{zero}(r)) \Rightarrow \\ \quad \quad d(\operatorname{matof}((n + 1), (n + 1), \\ \quad \quad \quad \quad (\lambda \langle \operatorname{dav\_}, \operatorname{daw\_} \rangle \in (\operatorname{interval}(1, (n + 1)) \times \operatorname{interval}(1, (n + 1))).\; \\ \quad \quad \quad \quad \quad \left(\text{if }(\operatorname{daw\_} = 1)\text{ then }\left(\text{if }(\operatorname{dav\_} = 1)\text{ then }\operatorname{entry}(y, \operatorname{dav\_}, \operatorname{daw\_})\text{ else }\left(\text{if }(\operatorname{dav\_} \leq (n + 1))\text{ then }\operatorname{zero}(r)\text{ else }\operatorname{entry}(y, \operatorname{dav\_}, \operatorname{daw\_})\right)\right)\text{ else }\operatorname{entry}(y, \operatorname{dav\_}, \operatorname{daw\_})\right)))) \\ \quad \quad =\; d(y) \end{array}}}
\vnbentry{\texttt{daf-\allowbreak{}bordered}}{\vnbstmt{\begin{array}{@{}l@{}} \forall r, n, y. \\ \quad \operatorname{is\text{-}commutative\text{-}ring}(r) \wedge \\ \quad (n \in \mathbb{N}) \wedge \\ \quad (y \in \operatorname{mat}((n + 1), (n + 1), \operatorname{carr}(r))) \wedge \\ \quad (\operatorname{entry}(y, 1, 1) = \operatorname{one}(r)) \wedge \\ \quad \forall k \in \operatorname{interval}(1, (n + 1)). \\ \quad \quad \neg (k = 1) \Rightarrow  \\ \quad \quad \quad (\operatorname{entry}(y, 1, k) = \operatorname{zero}(r)) \Rightarrow \\ \quad \quad \operatorname{matof}((n + 1), (n + 1), \\ \quad \quad \quad (\lambda \langle \operatorname{dav\_}, \operatorname{daw\_} \rangle \in (\operatorname{interval}(1, (n + 1)) \times \operatorname{interval}(1, (n + 1))).\; \\ \quad \quad \quad \quad \left(\text{if }(\operatorname{daw\_} = 1)\text{ then }\left(\text{if }(\operatorname{dav\_} = 1)\text{ then }\operatorname{entry}(y, \operatorname{dav\_}, \operatorname{daw\_})\text{ else }\left(\text{if }(\operatorname{dav\_} \leq (n + 1))\text{ then }\operatorname{zero}(r)\text{ else }\operatorname{entry}(y, \operatorname{dav\_}, \operatorname{daw\_})\right)\right)\text{ else }\operatorname{entry}(y, \operatorname{dav\_}, \operatorname{daw\_})\right))) \\ \quad \quad =\; \operatorname{border}(r, \operatorname{one}(r), \\ \quad \quad \quad \operatorname{submat}(\operatorname{matof}((n + 1), (n + 1), \\ \quad \quad \quad \quad \quad (\lambda \langle \operatorname{dav\_}, \operatorname{daw\_} \rangle \in (\operatorname{interval}(1, (n + 1)) \times \operatorname{interval}(1, (n + 1))).\; \\ \quad \quad \quad \quad \quad \quad \left(\text{if }(\operatorname{daw\_} = 1)\text{ then }\left(\text{if }(\operatorname{dav\_} = 1)\text{ then }\operatorname{entry}(y, \operatorname{dav\_}, \operatorname{daw\_})\text{ else }\left(\text{if }(\operatorname{dav\_} \leq (n + 1))\text{ then }\operatorname{zero}(r)\text{ else }\operatorname{entry}(y, \operatorname{dav\_}, \operatorname{daw\_})\right)\right)\text{ else }\operatorname{entry}(y, \operatorname{dav\_}, \operatorname{daw\_})\right))), \\ \quad \quad \quad \quad n, n), \\ \quad \quad \quad n, n) \end{array}}}
\vnbentry{\texttt{daf-\allowbreak{}cl}}{\vnbstmt{\begin{array}{@{}l@{}} \forall n \in \mathbb{N}, r, d, y. \\ \quad \forall r, d.; \forall x \in \operatorname{mat}(n, n, \operatorname{carr}(r)). \\ \quad \quad \operatorname{is\text{-}commutative\text{-}ring}(r) \wedge \\ \quad \quad (d \in (\operatorname{mat}(n, n, \operatorname{carr}(r)) \to \operatorname{carr}(r))) \wedge \\ \quad \quad \forall a, b, c, p, s. \\ \quad \quad \quad (a \in \operatorname{mat}(n, n, \operatorname{carr}(r))) \wedge \\ \quad \quad \quad (b \in \operatorname{mat}(n, n, \operatorname{carr}(r))) \wedge \\ \quad \quad \quad (c \in \operatorname{mat}(n, n, \operatorname{carr}(r))) \wedge \\ \quad \quad \quad (p \in \operatorname{interval}(1, n)) \wedge \\ \quad \quad \quad (s \in \operatorname{carr}(r)) \wedge \\ \quad \quad \quad \forall i, k \in \operatorname{interval}(1, n). \\ \quad \quad \quad \quad \neg (i = p) \Rightarrow  \\ \quad \quad \quad \quad \quad (\operatorname{entry}(b, i, k) = \operatorname{entry}(a, i, k)) \wedge \\ \quad \quad \quad \forall i, k \in \operatorname{interval}(1, n). \\ \quad \quad \quad \quad \neg (i = p) \Rightarrow  \\ \quad \quad \quad \quad \quad (\operatorname{entry}(c, i, k) = \operatorname{entry}(a, i, k)) \wedge \\ \quad \quad \quad \forall k \in \operatorname{interval}(1, n). \\ \quad \quad \quad \quad \operatorname{entry}(c, p, k) \\ \quad \quad \quad \quad =\; \operatorname{add}(r)(\operatorname{mul}(r)(s, \operatorname{entry}(a, p, k)), \operatorname{entry}(b, p, k)) \Rightarrow \\ \quad \quad \quad \quad (d(c) = \operatorname{add}(r)(\operatorname{mul}(r)(s, d(a)), d(b))) \wedge \\ \quad \quad \forall a, p, q. \\ \quad \quad \quad (a \in \operatorname{mat}(n, n, \operatorname{carr}(r))) \wedge \\ \quad \quad \quad (p \in \operatorname{interval}(1, n)) \wedge \\ \quad \quad \quad (q \in \operatorname{interval}(1, n)) \wedge \\ \quad \quad \quad \neg (p = q) \wedge \\ \quad \quad \quad \forall k \in \operatorname{interval}(1, n). \\ \quad \quad \quad \quad (\operatorname{entry}(a, p, k) = \operatorname{entry}(a, q, k)) \Rightarrow \\ \quad \quad \quad \quad (d(a) = \operatorname{zero}(r)) \Rightarrow \\ \quad \quad \quad (d(x) = \operatorname{mul}(r)(\operatorname{det}(r, n, x), d(\operatorname{identmat}(r, n)))) \wedge \\ \quad \operatorname{is\text{-}commutative\text{-}ring}(r) \wedge \\ \quad (d \in (\operatorname{mat}((n + 1), (n + 1), \operatorname{carr}(r)) \to \operatorname{carr}(r))) \wedge \\ \quad \forall a, b, c, p, s. \\ \quad \quad (a \in \operatorname{mat}((n + 1), (n + 1), \operatorname{carr}(r))) \wedge \\ \quad \quad (b \in \operatorname{mat}((n + 1), (n + 1), \operatorname{carr}(r))) \wedge \\ \quad \quad (c \in \operatorname{mat}((n + 1), (n + 1), \operatorname{carr}(r))) \wedge \\ \quad \quad (p \in \operatorname{interval}(1, (n + 1))) \wedge \\ \quad \quad (s \in \operatorname{carr}(r)) \wedge \\ \quad \quad \forall i, k \in \operatorname{interval}(1, (n + 1)). \\ \quad \quad \quad \neg (i = p) \Rightarrow  \\ \quad \quad \quad \quad (\operatorname{entry}(b, i, k) = \operatorname{entry}(a, i, k)) \wedge \\ \quad \quad \forall i, k \in \operatorname{interval}(1, (n + 1)). \\ \quad \quad \quad \neg (i = p) \Rightarrow  \\ \quad \quad \quad \quad (\operatorname{entry}(c, i, k) = \operatorname{entry}(a, i, k)) \wedge \\ \quad \quad \forall k \in \operatorname{interval}(1, (n + 1)). \\ \quad \quad \quad \operatorname{entry}(c, p, k) \\ \quad \quad \quad =\; \operatorname{add}(r)(\operatorname{mul}(r)(s, \operatorname{entry}(a, p, k)), \operatorname{entry}(b, p, k)) \Rightarrow \\ \quad \quad \quad (d(c) = \operatorname{add}(r)(\operatorname{mul}(r)(s, d(a)), d(b))) \wedge \\ \quad \forall a, p, q. \\ \quad \quad (a \in \operatorname{mat}((n + 1), (n + 1), \operatorname{carr}(r))) \wedge \\ \quad \quad (p \in \operatorname{interval}(1, (n + 1))) \wedge \\ \quad \quad (q \in \operatorname{interval}(1, (n + 1))) \wedge \\ \quad \quad \neg (p = q) \wedge \\ \quad \quad \forall k \in \operatorname{interval}(1, (n + 1)). \\ \quad \quad \quad (\operatorname{entry}(a, p, k) = \operatorname{entry}(a, q, k)) \Rightarrow \\ \quad \quad \quad (d(a) = \operatorname{zero}(r)) \wedge \\ \quad (y \in \operatorname{mat}((n + 1), (n + 1), \operatorname{carr}(r))) \wedge \\ \quad (\operatorname{entry}(y, 1, 1) = \operatorname{one}(r)) \wedge \\ \quad \forall k \in \operatorname{interval}(1, (n + 1)). \\ \quad \quad \neg (k = 1) \Rightarrow  \\ \quad \quad \quad (\operatorname{entry}(y, 1, k) = \operatorname{zero}(r)) \Rightarrow \\ \quad \quad d(y) \\ \quad \quad =\; \operatorname{mul}(r)(\operatorname{det}(r, n, \\ \quad \quad \quad \quad \operatorname{submat}(\operatorname{matof}((n + 1), (n + 1), \\ \quad \quad \quad \quad \quad \quad (\lambda \langle \operatorname{dav\_}, \operatorname{daw\_} \rangle \in (\operatorname{interval}(1, (n + 1)) \times \operatorname{interval}(1, (n + 1))).\; \\ \quad \quad \quad \quad \quad \quad \quad \left(\text{if }(\operatorname{daw\_} = 1)\text{ then }\left(\text{if }(\operatorname{dav\_} = 1)\text{ then }\operatorname{entry}(y, \operatorname{dav\_}, \operatorname{daw\_})\text{ else }\left(\text{if }(\operatorname{dav\_} \leq (n + 1))\text{ then }\operatorname{zero}(r)\text{ else }\operatorname{entry}(y, \operatorname{dav\_}, \operatorname{daw\_})\right)\right)\text{ else }\operatorname{entry}(y, \operatorname{dav\_}, \operatorname{daw\_})\right))), \\ \quad \quad \quad \quad \quad n, n)), \\ \quad \quad \quad d(\operatorname{identmat}(r, (n + 1)))) \end{array}}}
\vnbentry{\texttt{daf-\allowbreak{}basis}}{\vnbstmt{\begin{array}{@{}l@{}} \forall n \in \mathbb{N}, r, d, x, j. \\ \quad \forall r, d.; \forall \mathit{dfx} \in \operatorname{mat}(n, n, \operatorname{carr}(r)). \\ \quad \quad \operatorname{is\text{-}commutative\text{-}ring}(r) \wedge \\ \quad \quad (d \in (\operatorname{mat}(n, n, \operatorname{carr}(r)) \to \operatorname{carr}(r))) \wedge \\ \quad \quad \forall a, b, c, p, s. \\ \quad \quad \quad (a \in \operatorname{mat}(n, n, \operatorname{carr}(r))) \wedge \\ \quad \quad \quad (b \in \operatorname{mat}(n, n, \operatorname{carr}(r))) \wedge \\ \quad \quad \quad (c \in \operatorname{mat}(n, n, \operatorname{carr}(r))) \wedge \\ \quad \quad \quad (p \in \operatorname{interval}(1, n)) \wedge \\ \quad \quad \quad (s \in \operatorname{carr}(r)) \wedge \\ \quad \quad \quad \forall i, k \in \operatorname{interval}(1, n). \\ \quad \quad \quad \quad \neg (i = p) \Rightarrow  \\ \quad \quad \quad \quad \quad (\operatorname{entry}(b, i, k) = \operatorname{entry}(a, i, k)) \wedge \\ \quad \quad \quad \forall i, k \in \operatorname{interval}(1, n). \\ \quad \quad \quad \quad \neg (i = p) \Rightarrow  \\ \quad \quad \quad \quad \quad (\operatorname{entry}(c, i, k) = \operatorname{entry}(a, i, k)) \wedge \\ \quad \quad \quad \forall k \in \operatorname{interval}(1, n). \\ \quad \quad \quad \quad \operatorname{entry}(c, p, k) \\ \quad \quad \quad \quad =\; \operatorname{add}(r)(\operatorname{mul}(r)(s, \operatorname{entry}(a, p, k)), \operatorname{entry}(b, p, k)) \Rightarrow \\ \quad \quad \quad \quad (d(c) = \operatorname{add}(r)(\operatorname{mul}(r)(s, d(a)), d(b))) \wedge \\ \quad \quad \forall a, p, q. \\ \quad \quad \quad (a \in \operatorname{mat}(n, n, \operatorname{carr}(r))) \wedge \\ \quad \quad \quad (p \in \operatorname{interval}(1, n)) \wedge \\ \quad \quad \quad (q \in \operatorname{interval}(1, n)) \wedge \\ \quad \quad \quad \neg (p = q) \wedge \\ \quad \quad \quad \forall k \in \operatorname{interval}(1, n). \\ \quad \quad \quad \quad (\operatorname{entry}(a, p, k) = \operatorname{entry}(a, q, k)) \Rightarrow \\ \quad \quad \quad \quad (d(a) = \operatorname{zero}(r)) \Rightarrow \\ \quad \quad \quad d(\mathit{dfx}) \\ \quad \quad \quad =\; \operatorname{mul}(r)(\operatorname{det}(r, n, \mathit{dfx}), d(\operatorname{identmat}(r, n))) \wedge \\ \quad \operatorname{is\text{-}commutative\text{-}ring}(r) \wedge \\ \quad (d \in (\operatorname{mat}((n + 1), (n + 1), \operatorname{carr}(r)) \to \operatorname{carr}(r))) \wedge \\ \quad \forall a, b, c, p, s. \\ \quad \quad (a \in \operatorname{mat}((n + 1), (n + 1), \operatorname{carr}(r))) \wedge \\ \quad \quad (b \in \operatorname{mat}((n + 1), (n + 1), \operatorname{carr}(r))) \wedge \\ \quad \quad (c \in \operatorname{mat}((n + 1), (n + 1), \operatorname{carr}(r))) \wedge \\ \quad \quad (p \in \operatorname{interval}(1, (n + 1))) \wedge \\ \quad \quad (s \in \operatorname{carr}(r)) \wedge \\ \quad \quad \forall i, k \in \operatorname{interval}(1, (n + 1)). \\ \quad \quad \quad \neg (i = p) \Rightarrow  \\ \quad \quad \quad \quad (\operatorname{entry}(b, i, k) = \operatorname{entry}(a, i, k)) \wedge \\ \quad \quad \forall i, k \in \operatorname{interval}(1, (n + 1)). \\ \quad \quad \quad \neg (i = p) \Rightarrow  \\ \quad \quad \quad \quad (\operatorname{entry}(c, i, k) = \operatorname{entry}(a, i, k)) \wedge \\ \quad \quad \forall k \in \operatorname{interval}(1, (n + 1)). \\ \quad \quad \quad \operatorname{entry}(c, p, k) \\ \quad \quad \quad =\; \operatorname{add}(r)(\operatorname{mul}(r)(s, \operatorname{entry}(a, p, k)), \operatorname{entry}(b, p, k)) \Rightarrow \\ \quad \quad \quad (d(c) = \operatorname{add}(r)(\operatorname{mul}(r)(s, d(a)), d(b))) \wedge \\ \quad \forall a, p, q. \\ \quad \quad (a \in \operatorname{mat}((n + 1), (n + 1), \operatorname{carr}(r))) \wedge \\ \quad \quad (p \in \operatorname{interval}(1, (n + 1))) \wedge \\ \quad \quad (q \in \operatorname{interval}(1, (n + 1))) \wedge \\ \quad \quad \neg (p = q) \wedge \\ \quad \quad \forall k \in \operatorname{interval}(1, (n + 1)). \\ \quad \quad \quad (\operatorname{entry}(a, p, k) = \operatorname{entry}(a, q, k)) \Rightarrow \\ \quad \quad \quad (d(a) = \operatorname{zero}(r)) \wedge \\ \quad (x \in \operatorname{mat}((n + 1), (n + 1), \operatorname{carr}(r))) \wedge \\ \quad (j \in \operatorname{interval}(1, (n + 1))) \wedge \\ \quad \forall k \in \operatorname{interval}(1, (n + 1)). \\ \quad \quad \operatorname{entry}(x, 1, k) \\ \quad \quad =\; \left(\text{if }(k = j)\text{ then }\operatorname{one}(r)\text{ else }\operatorname{zero}(r)\right) \Rightarrow \\ \quad \quad d(x) \\ \quad \quad =\; \operatorname{mul}(r)(\operatorname{det}(r, (n + 1), x), d(\operatorname{identmat}(r, (n + 1)))) \end{array}}}
\vnbentry{\texttt{daf-\allowbreak{}psum-\allowbreak{}ind}}{\vnbstmt{\begin{array}{@{}l@{}} \forall m \in \mathbb{N}, r, n, d, a. \\ \quad \operatorname{is\text{-}commutative\text{-}ring}(r) \wedge \\ \quad (n \in \mathbb{N}) \wedge \\ \quad (d \in (\operatorname{mat}((n + 1), (n + 1), \operatorname{carr}(r)) \to \operatorname{carr}(r))) \wedge \\ \quad \forall \mathit{dfa}, b, c, p, s. \\ \quad \quad (\mathit{dfa} \in \operatorname{mat}((n + 1), (n + 1), \operatorname{carr}(r))) \wedge \\ \quad \quad (b \in \operatorname{mat}((n + 1), (n + 1), \operatorname{carr}(r))) \wedge \\ \quad \quad (c \in \operatorname{mat}((n + 1), (n + 1), \operatorname{carr}(r))) \wedge \\ \quad \quad (p \in \operatorname{interval}(1, (n + 1))) \wedge \\ \quad \quad (s \in \operatorname{carr}(r)) \wedge \\ \quad \quad \forall i, k \in \operatorname{interval}(1, (n + 1)). \\ \quad \quad \quad \neg (i = p) \Rightarrow  \\ \quad \quad \quad \quad (\operatorname{entry}(b, i, k) = \operatorname{entry}(\mathit{dfa}, i, k)) \wedge \\ \quad \quad \forall i, k \in \operatorname{interval}(1, (n + 1)). \\ \quad \quad \quad \neg (i = p) \Rightarrow  \\ \quad \quad \quad \quad (\operatorname{entry}(c, i, k) = \operatorname{entry}(\mathit{dfa}, i, k)) \wedge \\ \quad \quad \forall k \in \operatorname{interval}(1, (n + 1)). \\ \quad \quad \quad \operatorname{entry}(c, p, k) \\ \quad \quad \quad =\; \operatorname{add}(r)(\operatorname{mul}(r)(s, \operatorname{entry}(\mathit{dfa}, p, k)), \operatorname{entry}(b, p, k)) \Rightarrow \\ \quad \quad \quad (d(c) = \operatorname{add}(r)(\operatorname{mul}(r)(s, d(\mathit{dfa})), d(b))) \wedge \\ \quad \forall \mathit{dfa}, p, q. \\ \quad \quad (\mathit{dfa} \in \operatorname{mat}((n + 1), (n + 1), \operatorname{carr}(r))) \wedge \\ \quad \quad (p \in \operatorname{interval}(1, (n + 1))) \wedge \\ \quad \quad (q \in \operatorname{interval}(1, (n + 1))) \wedge \\ \quad \quad \neg (p = q) \wedge \\ \quad \quad \forall k \in \operatorname{interval}(1, (n + 1)). \\ \quad \quad \quad (\operatorname{entry}(\mathit{dfa}, p, k) = \operatorname{entry}(\mathit{dfa}, q, k)) \Rightarrow \\ \quad \quad \quad (d(\mathit{dfa}) = \operatorname{zero}(r)) \wedge \\ \quad (a \in \operatorname{mat}((n + 1), (n + 1), \operatorname{carr}(r))) \wedge \\ \quad \forall j \in \operatorname{interval}(1, (n + 1)). \\ \quad \quad d(\operatorname{matof}((n + 1), (n + 1), \\ \quad \quad \quad \quad (\lambda \langle \operatorname{dav\_}, \operatorname{daw\_} \rangle \in (\operatorname{interval}(1, (n + 1)) \times \operatorname{interval}(1, (n + 1))).\; \\ \quad \quad \quad \quad \quad \left(\text{if }(\operatorname{dav\_} = 1)\text{ then }\left(\text{if }(\operatorname{daw\_} = j)\text{ then }\operatorname{one}(r)\text{ else }\operatorname{zero}(r)\right)\text{ else }\operatorname{entry}(a, \operatorname{dav\_}, \operatorname{daw\_})\right)))) \\ \quad \quad =\; \operatorname{mul}(r)(\operatorname{det}(r, (n + 1), \\ \quad \quad \quad \quad \operatorname{matof}((n + 1), (n + 1), \\ \quad \quad \quad \quad \quad (\lambda \langle \operatorname{dav\_}, \operatorname{daw\_} \rangle \in (\operatorname{interval}(1, (n + 1)) \times \operatorname{interval}(1, (n + 1))).\; \\ \quad \quad \quad \quad \quad \quad \left(\text{if }(\operatorname{dav\_} = 1)\text{ then }\left(\text{if }(\operatorname{daw\_} = j)\text{ then }\operatorname{one}(r)\text{ else }\operatorname{zero}(r)\right)\text{ else }\operatorname{entry}(a, \operatorname{dav\_}, \operatorname{daw\_})\right)))), \\ \quad \quad \quad d(\operatorname{identmat}(r, (n + 1)))) \wedge \\ \quad (m \leq (n + 1)) \Rightarrow \\ \quad \quad d(\operatorname{matof}((n + 1), (n + 1), \\ \quad \quad \quad \quad (\lambda \langle \operatorname{dav\_}, \operatorname{daw\_} \rangle \in (\operatorname{interval}(1, (n + 1)) \times \operatorname{interval}(1, (n + 1))).\; \\ \quad \quad \quad \quad \quad \left(\text{if }(\operatorname{dav\_} = 1)\text{ then }\left(\text{if }(\operatorname{daw\_} \leq m)\text{ then }\operatorname{entry}(a, 1, \operatorname{daw\_})\text{ else }\operatorname{zero}(r)\right)\text{ else }\operatorname{entry}(a, \operatorname{dav\_}, \operatorname{daw\_})\right)))) \\ \quad \quad =\; \operatorname{mul}(r)(\operatorname{det}(r, (n + 1), \\ \quad \quad \quad \quad \operatorname{matof}((n + 1), (n + 1), \\ \quad \quad \quad \quad \quad (\lambda \langle \operatorname{dav\_}, \operatorname{daw\_} \rangle \in (\operatorname{interval}(1, (n + 1)) \times \operatorname{interval}(1, (n + 1))).\; \\ \quad \quad \quad \quad \quad \quad \left(\text{if }(\operatorname{dav\_} = 1)\text{ then }\left(\text{if }(\operatorname{daw\_} \leq m)\text{ then }\operatorname{entry}(a, 1, \operatorname{daw\_})\text{ else }\operatorname{zero}(r)\right)\text{ else }\operatorname{entry}(a, \operatorname{dav\_}, \operatorname{daw\_})\right)))), \\ \quad \quad \quad d(\operatorname{identmat}(r, (n + 1)))) \end{array}}}
\vnbentry{\texttt{daf-\allowbreak{}psum}}{\vnbstmt{\begin{array}{@{}l@{}} \forall r, n, d, a. \\ \quad \operatorname{is\text{-}commutative\text{-}ring}(r) \wedge \\ \quad (n \in \mathbb{N}) \wedge \\ \quad (d \in (\operatorname{mat}((n + 1), (n + 1), \operatorname{carr}(r)) \to \operatorname{carr}(r))) \wedge \\ \quad \forall \mathit{dfa}, b, c, p, s. \\ \quad \quad (\mathit{dfa} \in \operatorname{mat}((n + 1), (n + 1), \operatorname{carr}(r))) \wedge \\ \quad \quad (b \in \operatorname{mat}((n + 1), (n + 1), \operatorname{carr}(r))) \wedge \\ \quad \quad (c \in \operatorname{mat}((n + 1), (n + 1), \operatorname{carr}(r))) \wedge \\ \quad \quad (p \in \operatorname{interval}(1, (n + 1))) \wedge \\ \quad \quad (s \in \operatorname{carr}(r)) \wedge \\ \quad \quad \forall i, k \in \operatorname{interval}(1, (n + 1)). \\ \quad \quad \quad \neg (i = p) \Rightarrow  \\ \quad \quad \quad \quad (\operatorname{entry}(b, i, k) = \operatorname{entry}(\mathit{dfa}, i, k)) \wedge \\ \quad \quad \forall i, k \in \operatorname{interval}(1, (n + 1)). \\ \quad \quad \quad \neg (i = p) \Rightarrow  \\ \quad \quad \quad \quad (\operatorname{entry}(c, i, k) = \operatorname{entry}(\mathit{dfa}, i, k)) \wedge \\ \quad \quad \forall k \in \operatorname{interval}(1, (n + 1)). \\ \quad \quad \quad \operatorname{entry}(c, p, k) \\ \quad \quad \quad =\; \operatorname{add}(r)(\operatorname{mul}(r)(s, \operatorname{entry}(\mathit{dfa}, p, k)), \operatorname{entry}(b, p, k)) \Rightarrow \\ \quad \quad \quad (d(c) = \operatorname{add}(r)(\operatorname{mul}(r)(s, d(\mathit{dfa})), d(b))) \wedge \\ \quad \forall \mathit{dfa}, p, q. \\ \quad \quad (\mathit{dfa} \in \operatorname{mat}((n + 1), (n + 1), \operatorname{carr}(r))) \wedge \\ \quad \quad (p \in \operatorname{interval}(1, (n + 1))) \wedge \\ \quad \quad (q \in \operatorname{interval}(1, (n + 1))) \wedge \\ \quad \quad \neg (p = q) \wedge \\ \quad \quad \forall k \in \operatorname{interval}(1, (n + 1)). \\ \quad \quad \quad (\operatorname{entry}(\mathit{dfa}, p, k) = \operatorname{entry}(\mathit{dfa}, q, k)) \Rightarrow \\ \quad \quad \quad (d(\mathit{dfa}) = \operatorname{zero}(r)) \wedge \\ \quad (a \in \operatorname{mat}((n + 1), (n + 1), \operatorname{carr}(r))) \wedge \\ \quad \forall j \in \operatorname{interval}(1, (n + 1)). \\ \quad \quad d(\operatorname{matof}((n + 1), (n + 1), \\ \quad \quad \quad \quad (\lambda \langle \operatorname{dav\_}, \operatorname{daw\_} \rangle \in (\operatorname{interval}(1, (n + 1)) \times \operatorname{interval}(1, (n + 1))).\; \\ \quad \quad \quad \quad \quad \left(\text{if }(\operatorname{dav\_} = 1)\text{ then }\left(\text{if }(\operatorname{daw\_} = j)\text{ then }\operatorname{one}(r)\text{ else }\operatorname{zero}(r)\right)\text{ else }\operatorname{entry}(a, \operatorname{dav\_}, \operatorname{daw\_})\right)))) \\ \quad \quad =\; \operatorname{mul}(r)(\operatorname{det}(r, (n + 1), \\ \quad \quad \quad \quad \operatorname{matof}((n + 1), (n + 1), \\ \quad \quad \quad \quad \quad (\lambda \langle \operatorname{dav\_}, \operatorname{daw\_} \rangle \in (\operatorname{interval}(1, (n + 1)) \times \operatorname{interval}(1, (n + 1))).\; \\ \quad \quad \quad \quad \quad \quad \left(\text{if }(\operatorname{dav\_} = 1)\text{ then }\left(\text{if }(\operatorname{daw\_} = j)\text{ then }\operatorname{one}(r)\text{ else }\operatorname{zero}(r)\right)\text{ else }\operatorname{entry}(a, \operatorname{dav\_}, \operatorname{daw\_})\right)))), \\ \quad \quad \quad d(\operatorname{identmat}(r, (n + 1)))) \Rightarrow \\ \quad \quad d(a) \\ \quad \quad =\; \operatorname{mul}(r)(\operatorname{det}(r, (n + 1), a), d(\operatorname{identmat}(r, (n + 1)))) \end{array}}}
\vnbentry{\texttt{det-\allowbreak{}alternating-\allowbreak{}form-\allowbreak{}ind}}{\vnbstmt{\begin{array}{@{}l@{}} \forall n \in \mathbb{N}, r, d.; \forall x \in \operatorname{mat}(n, n, \operatorname{carr}(r)). \\ \quad \operatorname{is\text{-}commutative\text{-}ring}(r) \wedge \\ \quad (d \in (\operatorname{mat}(n, n, \operatorname{carr}(r)) \to \operatorname{carr}(r))) \wedge \\ \quad \forall a, b, c, p, s. \\ \quad \quad (a \in \operatorname{mat}(n, n, \operatorname{carr}(r))) \wedge \\ \quad \quad (b \in \operatorname{mat}(n, n, \operatorname{carr}(r))) \wedge \\ \quad \quad (c \in \operatorname{mat}(n, n, \operatorname{carr}(r))) \wedge \\ \quad \quad (p \in \operatorname{interval}(1, n)) \wedge \\ \quad \quad (s \in \operatorname{carr}(r)) \wedge \\ \quad \quad \forall i, k \in \operatorname{interval}(1, n). \\ \quad \quad \quad \neg (i = p) \Rightarrow  \\ \quad \quad \quad \quad (\operatorname{entry}(b, i, k) = \operatorname{entry}(a, i, k)) \wedge \\ \quad \quad \forall i, k \in \operatorname{interval}(1, n). \\ \quad \quad \quad \neg (i = p) \Rightarrow  \\ \quad \quad \quad \quad (\operatorname{entry}(c, i, k) = \operatorname{entry}(a, i, k)) \wedge \\ \quad \quad \forall k \in \operatorname{interval}(1, n). \\ \quad \quad \quad \operatorname{entry}(c, p, k) \\ \quad \quad \quad =\; \operatorname{add}(r)(\operatorname{mul}(r)(s, \operatorname{entry}(a, p, k)), \operatorname{entry}(b, p, k)) \Rightarrow \\ \quad \quad \quad (d(c) = \operatorname{add}(r)(\operatorname{mul}(r)(s, d(a)), d(b))) \wedge \\ \quad \forall a, p, q. \\ \quad \quad (a \in \operatorname{mat}(n, n, \operatorname{carr}(r))) \wedge \\ \quad \quad (p \in \operatorname{interval}(1, n)) \wedge \\ \quad \quad (q \in \operatorname{interval}(1, n)) \wedge \\ \quad \quad \neg (p = q) \wedge \\ \quad \quad \forall k \in \operatorname{interval}(1, n). \\ \quad \quad \quad (\operatorname{entry}(a, p, k) = \operatorname{entry}(a, q, k)) \Rightarrow \\ \quad \quad \quad (d(a) = \operatorname{zero}(r)) \Rightarrow \\ \quad \quad (d(x) = \operatorname{mul}(r)(\operatorname{det}(r, n, x), d(\operatorname{identmat}(r, n)))) \end{array}}}
\vnbentry{\texttt{det-\allowbreak{}alternating-\allowbreak{}form}}{\vnbstmt{\begin{array}{@{}l@{}} \forall r, n, d, a. \\ \quad \operatorname{is\text{-}commutative\text{-}ring}(r) \wedge \\ \quad (n \in \mathbb{N}) \wedge \\ \quad (d \in (\operatorname{mat}(n, n, \operatorname{carr}(r)) \to \operatorname{carr}(r))) \wedge \\ \quad \forall \mathit{dfa}, b, c, p, s. \\ \quad \quad (\mathit{dfa} \in \operatorname{mat}(n, n, \operatorname{carr}(r))) \wedge \\ \quad \quad (b \in \operatorname{mat}(n, n, \operatorname{carr}(r))) \wedge \\ \quad \quad (c \in \operatorname{mat}(n, n, \operatorname{carr}(r))) \wedge \\ \quad \quad (p \in \operatorname{interval}(1, n)) \wedge \\ \quad \quad (s \in \operatorname{carr}(r)) \wedge \\ \quad \quad \forall i, k \in \operatorname{interval}(1, n). \\ \quad \quad \quad \neg (i = p) \Rightarrow  \\ \quad \quad \quad \quad (\operatorname{entry}(b, i, k) = \operatorname{entry}(\mathit{dfa}, i, k)) \wedge \\ \quad \quad \forall i, k \in \operatorname{interval}(1, n). \\ \quad \quad \quad \neg (i = p) \Rightarrow  \\ \quad \quad \quad \quad (\operatorname{entry}(c, i, k) = \operatorname{entry}(\mathit{dfa}, i, k)) \wedge \\ \quad \quad \forall k \in \operatorname{interval}(1, n). \\ \quad \quad \quad \operatorname{entry}(c, p, k) \\ \quad \quad \quad =\; \operatorname{add}(r)(\operatorname{mul}(r)(s, \operatorname{entry}(\mathit{dfa}, p, k)), \operatorname{entry}(b, p, k)) \Rightarrow \\ \quad \quad \quad (d(c) = \operatorname{add}(r)(\operatorname{mul}(r)(s, d(\mathit{dfa})), d(b))) \wedge \\ \quad \forall \mathit{dfa}, p, q. \\ \quad \quad (\mathit{dfa} \in \operatorname{mat}(n, n, \operatorname{carr}(r))) \wedge \\ \quad \quad (p \in \operatorname{interval}(1, n)) \wedge \\ \quad \quad (q \in \operatorname{interval}(1, n)) \wedge \\ \quad \quad \neg (p = q) \wedge \\ \quad \quad \forall k \in \operatorname{interval}(1, n). \\ \quad \quad \quad (\operatorname{entry}(\mathit{dfa}, p, k) = \operatorname{entry}(\mathit{dfa}, q, k)) \Rightarrow \\ \quad \quad \quad (d(\mathit{dfa}) = \operatorname{zero}(r)) \wedge \\ \quad (a \in \operatorname{mat}(n, n, \operatorname{carr}(r))) \Rightarrow \\ \quad \quad (d(a) = \operatorname{mul}(r)(\operatorname{det}(r, n, a), d(\operatorname{identmat}(r, n)))) \end{array}}}

\vnbfile{theorem-library/det-multiplicative.scm}{theorem-library/det-multiplicative.scm}
\noindent{\small det-multiplicative.scm -- det(P Q) = det(P) det(Q) over a commutative ring (Hoffman-Kunze 5.3 Theorem 3; the notes' proposition "det is multiplicative"), and the notes' corollary 3.20 (R a field, V invertible: det(A V) = det A det V). Batch 26-B, 2026-09-24.}\par
\vnbentry{\texttt{det-\allowbreak{}multiplicative}}{\vnbstmt{\begin{array}{@{}l@{}} \forall r, n, p, q. \\ \quad \operatorname{is\text{-}commutative\text{-}ring}(r) \wedge \\ \quad (n \in \mathbb{N}) \wedge \\ \quad (p \in \operatorname{mat}(n, n, \operatorname{carr}(r))) \wedge \\ \quad (q \in \operatorname{mat}(n, n, \operatorname{carr}(r))) \Rightarrow \\ \quad \quad \operatorname{det}(r, n, \operatorname{matmul}(r, p, q)) \\ \quad \quad =\; \operatorname{mul}(r)(\operatorname{det}(r, n, p), \operatorname{det}(r, n, q)) \end{array}}}
\vnbentry{\texttt{det-\allowbreak{}mul-\allowbreak{}invertible}}{\vnbstmt{\begin{array}{@{}l@{}} \forall r, n, a, v. \\ \quad \operatorname{is\text{-}field\text{-}ring}(r) \wedge \\ \quad (n \in \mathbb{N}) \wedge \\ \quad (a \in \operatorname{mat}(n, n, \operatorname{carr}(r))) \wedge \\ \quad \operatorname{is\text{-}invertible\text{-}mat}(r, n, v) \Rightarrow \\ \quad \quad \operatorname{det}(r, n, \operatorname{matmul}(r, a, v)) \\ \quad \quad =\; \operatorname{mul}(r)(\operatorname{det}(r, n, a), \operatorname{det}(r, n, v)) \end{array}}}

\vnbfile{theorem-library/det-transpose.scm}{theorem-library/det-transpose.scm}
\noindent{\small det-transpose.scm -- the transpose of a matrix and det(A\textasciicircum{}t) = det A; the COLUMN clauses of the linear-algebra notes' 3.18 (linearity in each column, two equal columns give 0, interchanging two columns negates det), 3.19 (det under an elementary COLUMN operation) and 3.22 (expansion along any column).  Batch 32}\par
\vnbentry{\texttt{transpose-\allowbreak{}in-\allowbreak{}mat}}{\vnbstmt{\begin{array}{@{}l@{}} \forall a, m, n, x. \\ \quad (a \in \operatorname{mat}(m, n, x)) \Rightarrow  \\ \quad \quad (\operatorname{transpose}(a, m, n) \in \operatorname{mat}(n, m, x)) \end{array}}}
\vnbentry{\texttt{transpose-\allowbreak{}entry}}{\vnbstmt{\begin{array}{@{}l@{}} \forall a, m, n, x, i, j. \\ \quad (a \in \operatorname{mat}(m, n, x)) \wedge \\ \quad (i \in \operatorname{interval}(1, n)) \wedge \\ \quad (j \in \operatorname{interval}(1, m)) \Rightarrow \\ \quad \quad (\operatorname{entry}(\operatorname{transpose}(a, m, n), i, j) = \operatorname{entry}(a, j, i)) \end{array}}}
\vnbentry{\texttt{transpose-\allowbreak{}transpose}}{\vnbstmt{\begin{array}{@{}l@{}} \forall a, m, n, x. \\ \quad (a \in \operatorname{mat}(m, n, x)) \Rightarrow  \\ \quad \quad (\operatorname{transpose}(\operatorname{transpose}(a, m, n), n, m) = a) \end{array}}}
\vnbentry{\texttt{transpose-\allowbreak{}identity}}{\vnbstmt{\begin{array}{@{}l@{}} \forall r, n. \\ \quad (\operatorname{is\text{-}ring}(r) \wedge (n \in \mathbb{N})) \Rightarrow  \\ \quad \quad (\operatorname{transpose}(\operatorname{identmat}(r, n), n, n) = \operatorname{identmat}(r, n)) \end{array}}}
\vnbentry{\texttt{transpose-\allowbreak{}matmul}}{\vnbstmt{\begin{array}{@{}l@{}} \forall r, m, n, c, p, q. \\ \quad \operatorname{is\text{-}commutative\text{-}ring}(r) \wedge \\ \quad (p \in \operatorname{mat}(m, n, \operatorname{carr}(r))) \wedge \\ \quad (q \in \operatorname{mat}(n, c, \operatorname{carr}(r))) \wedge \\ \quad ((n = 0) \Rightarrow ((m = 0) \vee (c = 0))) \Rightarrow \\ \quad \quad \operatorname{transpose}(\operatorname{matmul}(r, p, q), m, c) \\ \quad \quad =\; \operatorname{matmul}(r, \operatorname{transpose}(q, n, c), \operatorname{transpose}(p, m, n)) \end{array}}}
\vnbentry{\texttt{dtt-\allowbreak{}ring-\allowbreak{}mul-\allowbreak{}neg1}}{\vnbstmt{\begin{array}{@{}l@{}} \forall r.; \forall x \in \operatorname{carr}(r). \\ \quad \operatorname{is\text{-}commutative\text{-}ring}(r) \Rightarrow  \\ \quad \quad (\operatorname{mul}(r)(x, \operatorname{neg}(r)(\operatorname{one}(r))) = \operatorname{neg}(r)(x)) \end{array}}}
\vnbentry{\texttt{dtt-\allowbreak{}ring-\allowbreak{}mul-\allowbreak{}one}}{\vnbstmt{\begin{array}{@{}l@{}} \forall r.; \forall x \in \operatorname{carr}(r). \\ \quad \operatorname{is\text{-}commutative\text{-}ring}(r) \Rightarrow  \\ \quad \quad (\operatorname{mul}(r)(x, \operatorname{one}(r)) = x) \end{array}}}
\vnbentry{\texttt{dtt-\allowbreak{}ring-\allowbreak{}neg-\allowbreak{}lin}}{\vnbstmt{\begin{array}{@{}l@{}} \forall r, s, a, b. \\ \quad \operatorname{is\text{-}commutative\text{-}ring}(r) \wedge \\ \quad (s \in \operatorname{carr}(r)) \wedge \\ \quad (a \in \operatorname{carr}(r)) \wedge \\ \quad (b \in \operatorname{carr}(r)) \Rightarrow \\ \quad \quad \operatorname{neg}(r)(\operatorname{add}(r)(\operatorname{mul}(r)(s, \operatorname{neg}(r)(a)), \operatorname{neg}(r)(b))) \\ \quad \quad =\; \operatorname{add}(r)(\operatorname{mul}(r)(s, a), b) \end{array}}}
\vnbentry{\texttt{dtt-\allowbreak{}ring-\allowbreak{}neg-\allowbreak{}neg}}{\vnbstmt{\begin{array}{@{}l@{}} \forall r.; \forall x \in \operatorname{carr}(r). \\ \quad \operatorname{is\text{-}commutative\text{-}ring}(r) \Rightarrow  \\ \quad \quad (\operatorname{neg}(r)(\operatorname{neg}(r)(x)) = x) \end{array}}}
\vnbentry{\texttt{dtt-\allowbreak{}ring-\allowbreak{}negx-\allowbreak{}x}}{\vnbstmt{\begin{array}{@{}l@{}} \forall r.; \forall x \in \operatorname{carr}(r). \\ \quad \operatorname{is\text{-}commutative\text{-}ring}(r) \Rightarrow  \\ \quad \quad \operatorname{add}(r)(\operatorname{mul}(r)(\operatorname{neg}(r)(\operatorname{one}(r)), x), x) \\ \quad \quad =\; \operatorname{zero}(r) \end{array}}}
\vnbentry{\texttt{dtt-\allowbreak{}ring-\allowbreak{}zero-\allowbreak{}lin}}{\vnbstmt{\begin{array}{@{}l@{}} \forall r. \\ \quad \operatorname{is\text{-}commutative\text{-}ring}(r) \Rightarrow  \\ \quad \quad \operatorname{zero}(r) \\ \quad \quad =\; \operatorname{add}(r)(\operatorname{mul}(r)(\operatorname{neg}(r)(\operatorname{one}(r)), \operatorname{zero}(r)), \operatorname{zero}(r)) \end{array}}}
\vnbentry{\texttt{dtt-\allowbreak{}ring-\allowbreak{}x-\allowbreak{}xneg}}{\vnbstmt{\begin{array}{@{}l@{}} \forall r.; \forall x \in \operatorname{carr}(r). \\ \quad \operatorname{is\text{-}commutative\text{-}ring}(r) \Rightarrow  \\ \quad \quad \operatorname{add}(r)(x, \operatorname{mul}(r)(x, \operatorname{neg}(r)(\operatorname{one}(r)))) \\ \quad \quad =\; \operatorname{zero}(r) \end{array}}}
\vnbentry{\texttt{det-\allowbreak{}mul-\allowbreak{}elem-\allowbreak{}f-\allowbreak{}col}}{\vnbstmt{\begin{array}{@{}l@{}} \forall r, n, a, k, l. \\ \quad \operatorname{is\text{-}commutative\text{-}ring}(r) \wedge \\ \quad (n \in \mathbb{N}) \wedge \\ \quad (a \in \operatorname{mat}(n, n, \operatorname{carr}(r))) \wedge \\ \quad (k \in \operatorname{interval}(1, n)) \wedge \\ \quad (l \in \operatorname{interval}(1, n)) \wedge \\ \quad \neg (k = l) \Rightarrow \\ \quad \quad \operatorname{det}(r, n, \operatorname{matmul}(r, a, \operatorname{elem\text{-}f}(r, n, k, l))) \\ \quad \quad =\; \operatorname{mul}(r)(\operatorname{det}(r, n, a), \operatorname{det}(r, n, \operatorname{elem\text{-}f}(r, n, k, l))) \end{array}}}
\vnbentry{\texttt{det-\allowbreak{}elem-\allowbreak{}f-\allowbreak{}col}}{\vnbstmt{\begin{array}{@{}l@{}} \forall r, n, a, k, l. \\ \quad \operatorname{is\text{-}commutative\text{-}ring}(r) \wedge \\ \quad (n \in \mathbb{N}) \wedge \\ \quad (a \in \operatorname{mat}(n, n, \operatorname{carr}(r))) \wedge \\ \quad (k \in \operatorname{interval}(1, n)) \wedge \\ \quad (l \in \operatorname{interval}(1, n)) \wedge \\ \quad \neg (k = l) \Rightarrow \\ \quad \quad \operatorname{det}(r, n, \operatorname{matmul}(r, a, \operatorname{elem\text{-}f}(r, n, k, l))) \\ \quad \quad =\; \operatorname{neg}(r)(\operatorname{det}(r, n, a)) \end{array}}}
\vnbentry{\texttt{det-\allowbreak{}mul-\allowbreak{}elem-\allowbreak{}g-\allowbreak{}col}}{\vnbstmt{\begin{array}{@{}l@{}} \forall r, n, a, c, k, l. \\ \quad \operatorname{is\text{-}commutative\text{-}ring}(r) \wedge \\ \quad (n \in \mathbb{N}) \wedge \\ \quad (a \in \operatorname{mat}(n, n, \operatorname{carr}(r))) \wedge \\ \quad (c \in \operatorname{carr}(r)) \wedge \\ \quad (k \in \operatorname{interval}(1, n)) \wedge \\ \quad (l \in \operatorname{interval}(1, n)) \wedge \\ \quad \neg (k = l) \Rightarrow \\ \quad \quad \operatorname{det}(r, n, \operatorname{matmul}(r, a, \operatorname{elem\text{-}g}(r, n, c, k, l))) \\ \quad \quad =\; \operatorname{mul}(r)(\operatorname{det}(r, n, a), \operatorname{det}(r, n, \operatorname{elem\text{-}g}(r, n, c, k, l))) \end{array}}}
\vnbentry{\texttt{det-\allowbreak{}elem-\allowbreak{}g-\allowbreak{}col}}{\vnbstmt{\begin{array}{@{}l@{}} \forall r, n, a, c, k, l. \\ \quad \operatorname{is\text{-}commutative\text{-}ring}(r) \wedge \\ \quad (n \in \mathbb{N}) \wedge \\ \quad (a \in \operatorname{mat}(n, n, \operatorname{carr}(r))) \wedge \\ \quad (c \in \operatorname{carr}(r)) \wedge \\ \quad (k \in \operatorname{interval}(1, n)) \wedge \\ \quad (l \in \operatorname{interval}(1, n)) \wedge \\ \quad \neg (k = l) \Rightarrow \\ \quad \quad \operatorname{det}(r, n, \operatorname{matmul}(r, a, \operatorname{elem\text{-}g}(r, n, c, k, l))) \\ \quad \quad =\; \operatorname{det}(r, n, a) \end{array}}}
\vnbentry{\texttt{det-\allowbreak{}mul-\allowbreak{}elem-\allowbreak{}h-\allowbreak{}col}}{\vnbstmt{\begin{array}{@{}l@{}} \forall r, n, a, c, k. \\ \quad \operatorname{is\text{-}commutative\text{-}ring}(r) \wedge \\ \quad (n \in \mathbb{N}) \wedge \\ \quad (a \in \operatorname{mat}(n, n, \operatorname{carr}(r))) \wedge \\ \quad (c \in \operatorname{carr}(r)) \wedge \\ \quad (k \in \operatorname{interval}(1, n)) \Rightarrow \\ \quad \quad \operatorname{det}(r, n, \operatorname{matmul}(r, a, \operatorname{elem\text{-}h}(r, n, c, k))) \\ \quad \quad =\; \operatorname{mul}(r)(\operatorname{det}(r, n, a), \operatorname{det}(r, n, \operatorname{elem\text{-}h}(r, n, c, k))) \end{array}}}
\vnbentry{\texttt{det-\allowbreak{}elem-\allowbreak{}h-\allowbreak{}col}}{\vnbstmt{\begin{array}{@{}l@{}} \forall r, n, a, c, k. \\ \quad \operatorname{is\text{-}commutative\text{-}ring}(r) \wedge \\ \quad (n \in \mathbb{N}) \wedge \\ \quad (a \in \operatorname{mat}(n, n, \operatorname{carr}(r))) \wedge \\ \quad (c \in \operatorname{carr}(r)) \wedge \\ \quad (k \in \operatorname{interval}(1, n)) \Rightarrow \\ \quad \quad \operatorname{det}(r, n, \operatorname{matmul}(r, a, \operatorname{elem\text{-}h}(r, n, c, k))) \\ \quad \quad =\; \operatorname{mul}(r)(c, \operatorname{det}(r, n, a)) \end{array}}}
\vnbentry{\texttt{dtt-\allowbreak{}minor1-\allowbreak{}colcongr}}{\vnbstmt{\begin{array}{@{}l@{}} \forall r, n, x, y, p. \\ \quad \operatorname{is\text{-}ring}(r) \wedge \\ \quad (n \in \mathbb{N}) \wedge \\ \quad (x \in \operatorname{mat}((n + 1), (n + 1), \operatorname{carr}(r))) \wedge \\ \quad (y \in \operatorname{mat}((n + 1), (n + 1), \operatorname{carr}(r))) \wedge \\ \quad (p = 1) \wedge \\ \quad \forall j, i \in \operatorname{interval}(1, (n + 1)). \\ \quad \quad \neg (j = p) \Rightarrow  \\ \quad \quad \quad (\operatorname{entry}(x, i, j) = \operatorname{entry}(y, i, j)) \Rightarrow \\ \quad \quad (\operatorname{minor}(x, 1, p, n) = \operatorname{minor}(y, 1, p, n)) \end{array}}}
\vnbentry{\texttt{dtt-\allowbreak{}col1-\allowbreak{}ind}}{\vnbstmt{\begin{array}{@{}l@{}} \forall n \in \mathbb{N}, r, a, b, c, p, \mathit{dsc}. \\ \quad \operatorname{is\text{-}commutative\text{-}ring}(r) \wedge \\ \quad (a \in \operatorname{mat}(n, n, \operatorname{carr}(r))) \wedge \\ \quad (b \in \operatorname{mat}(n, n, \operatorname{carr}(r))) \wedge \\ \quad (c \in \operatorname{mat}(n, n, \operatorname{carr}(r))) \wedge \\ \quad (p \in \operatorname{interval}(1, n)) \wedge \\ \quad (p = 1) \wedge \\ \quad (\mathit{dsc} \in \operatorname{carr}(r)) \wedge \\ \quad \forall j, i \in \operatorname{interval}(1, n). \\ \quad \quad \neg (j = p) \Rightarrow  \\ \quad \quad \quad (\operatorname{entry}(b, i, j) = \operatorname{entry}(a, i, j)) \wedge \\ \quad \forall j, i \in \operatorname{interval}(1, n). \\ \quad \quad \neg (j = p) \Rightarrow  \\ \quad \quad \quad (\operatorname{entry}(c, i, j) = \operatorname{entry}(a, i, j)) \wedge \\ \quad \forall i \in \operatorname{interval}(1, n). \\ \quad \quad \operatorname{entry}(c, i, p) \\ \quad \quad =\; \operatorname{add}(r)(\operatorname{mul}(r)(\mathit{dsc}, \operatorname{entry}(a, i, p)), \operatorname{entry}(b, i, p)) \Rightarrow \\ \quad \quad \operatorname{det}(r, n, c) \\ \quad \quad =\; \operatorname{add}(r)(\operatorname{mul}(r)(\mathit{dsc}, \operatorname{det}(r, n, a)), \operatorname{det}(r, n, b)) \end{array}}}
\vnbentry{\texttt{det-\allowbreak{}col-\allowbreak{}linear}}{\vnbstmt{\begin{array}{@{}l@{}} \forall r, n, a, b, c, p, d. \\ \quad \operatorname{is\text{-}commutative\text{-}ring}(r) \wedge \\ \quad (n \in \mathbb{N}) \wedge \\ \quad (a \in \operatorname{mat}(n, n, \operatorname{carr}(r))) \wedge \\ \quad (b \in \operatorname{mat}(n, n, \operatorname{carr}(r))) \wedge \\ \quad (c \in \operatorname{mat}(n, n, \operatorname{carr}(r))) \wedge \\ \quad (p \in \operatorname{interval}(1, n)) \wedge \\ \quad (d \in \operatorname{carr}(r)) \wedge \\ \quad \forall j, i \in \operatorname{interval}(1, n). \\ \quad \quad \neg (j = p) \Rightarrow  \\ \quad \quad \quad (\operatorname{entry}(b, i, j) = \operatorname{entry}(a, i, j)) \wedge \\ \quad \forall j, i \in \operatorname{interval}(1, n). \\ \quad \quad \neg (j = p) \Rightarrow  \\ \quad \quad \quad (\operatorname{entry}(c, i, j) = \operatorname{entry}(a, i, j)) \wedge \\ \quad \forall i \in \operatorname{interval}(1, n). \\ \quad \quad \operatorname{entry}(c, i, p) \\ \quad \quad =\; \operatorname{add}(r)(\operatorname{mul}(r)(d, \operatorname{entry}(a, i, p)), \operatorname{entry}(b, i, p)) \Rightarrow \\ \quad \quad \operatorname{det}(r, n, c) \\ \quad \quad =\; \operatorname{add}(r)(\operatorname{mul}(r)(d, \operatorname{det}(r, n, a)), \operatorname{det}(r, n, b)) \end{array}}}
\vnbentry{\texttt{dtt-\allowbreak{}zero-\allowbreak{}col}}{\vnbstmt{\begin{array}{@{}l@{}} \forall r, n, c, q. \\ \quad \operatorname{is\text{-}commutative\text{-}ring}(r) \wedge \\ \quad (n \in \mathbb{N}) \wedge \\ \quad (c \in \operatorname{mat}(n, n, \operatorname{carr}(r))) \wedge \\ \quad (q \in \operatorname{interval}(1, n)) \wedge \\ \quad \forall i \in \operatorname{interval}(1, n). \\ \quad \quad (\operatorname{entry}(c, i, q) = \operatorname{zero}(r)) \Rightarrow \\ \quad \quad (\operatorname{det}(r, n, c) = \operatorname{zero}(r)) \end{array}}}
\vnbentry{\texttt{det-\allowbreak{}equal-\allowbreak{}cols-\allowbreak{}zero}}{\vnbstmt{\begin{array}{@{}l@{}} \forall r, n, a, p, q. \\ \quad \operatorname{is\text{-}commutative\text{-}ring}(r) \wedge \\ \quad (n \in \mathbb{N}) \wedge \\ \quad (a \in \operatorname{mat}(n, n, \operatorname{carr}(r))) \wedge \\ \quad (p \in \operatorname{interval}(1, n)) \wedge \\ \quad (q \in \operatorname{interval}(1, n)) \wedge \\ \quad \neg (p = q) \wedge \\ \quad \forall i \in \operatorname{interval}(1, n). \\ \quad \quad (\operatorname{entry}(a, i, p) = \operatorname{entry}(a, i, q)) \Rightarrow \\ \quad \quad (\operatorname{det}(r, n, a) = \operatorname{zero}(r)) \end{array}}}
\vnbentry{\texttt{det-\allowbreak{}swap-\allowbreak{}cols}}{\vnbstmt{\begin{array}{@{}l@{}} \forall r, n, a, b, p, q. \\ \quad \operatorname{is\text{-}commutative\text{-}ring}(r) \wedge \\ \quad (n \in \mathbb{N}) \wedge \\ \quad (a \in \operatorname{mat}(n, n, \operatorname{carr}(r))) \wedge \\ \quad (b \in \operatorname{mat}(n, n, \operatorname{carr}(r))) \wedge \\ \quad (p \in \operatorname{interval}(1, n)) \wedge \\ \quad (q \in \operatorname{interval}(1, n)) \wedge \\ \quad \neg (p = q) \wedge \\ \quad \forall j, i \in \operatorname{interval}(1, n). \\ \quad \quad (\neg (j = p) \wedge \neg (j = q)) \Rightarrow  \\ \quad \quad \quad (\operatorname{entry}(b, i, j) = \operatorname{entry}(a, i, j)) \wedge \\ \quad \forall i \in \operatorname{interval}(1, n). \\ \quad \quad (\operatorname{entry}(b, i, p) = \operatorname{entry}(a, i, q)) \wedge \\ \quad \forall i \in \operatorname{interval}(1, n). \\ \quad \quad (\operatorname{entry}(b, i, q) = \operatorname{entry}(a, i, p)) \Rightarrow \\ \quad \quad (\operatorname{det}(r, n, b) = \operatorname{neg}(r)(\operatorname{det}(r, n, a))) \end{array}}}
\vnbentry{\texttt{dtt-\allowbreak{}tr-\allowbreak{}fun}}{\vnbstmt{\begin{array}{@{}l@{}} \forall r, n. \\ \quad (\operatorname{is\text{-}commutative\text{-}ring}(r) \wedge (n \in \mathbb{N})) \Rightarrow  \\ \quad \quad (\lambda x \in \operatorname{mat}(n, n, \operatorname{carr}(r)).\; \\ \quad \quad \quad \operatorname{det}(r, n, \operatorname{transpose}(x, n, n))) \\ \quad \quad \in\; (\operatorname{mat}(n, n, \operatorname{carr}(r)) \to \operatorname{carr}(r)) \end{array}}}
\vnbentry{\texttt{dtt-\allowbreak{}tr-\allowbreak{}lin}}{\vnbstmt{\begin{array}{@{}l@{}} \forall r, n, a, b, c, p, s. \\ \quad \operatorname{is\text{-}commutative\text{-}ring}(r) \wedge \\ \quad (n \in \mathbb{N}) \wedge \\ \quad (a \in \operatorname{mat}(n, n, \operatorname{carr}(r))) \wedge \\ \quad (b \in \operatorname{mat}(n, n, \operatorname{carr}(r))) \wedge \\ \quad (c \in \operatorname{mat}(n, n, \operatorname{carr}(r))) \wedge \\ \quad (p \in \operatorname{interval}(1, n)) \wedge \\ \quad (s \in \operatorname{carr}(r)) \wedge \\ \quad \forall i, k \in \operatorname{interval}(1, n). \\ \quad \quad \neg (i = p) \Rightarrow  \\ \quad \quad \quad (\operatorname{entry}(b, i, k) = \operatorname{entry}(a, i, k)) \wedge \\ \quad \forall i, k \in \operatorname{interval}(1, n). \\ \quad \quad \neg (i = p) \Rightarrow  \\ \quad \quad \quad (\operatorname{entry}(c, i, k) = \operatorname{entry}(a, i, k)) \wedge \\ \quad \forall k \in \operatorname{interval}(1, n). \\ \quad \quad \operatorname{entry}(c, p, k) \\ \quad \quad =\; \operatorname{add}(r)(\operatorname{mul}(r)(s, \operatorname{entry}(a, p, k)), \operatorname{entry}(b, p, k)) \Rightarrow \\ \quad \quad (\lambda x \in \operatorname{mat}(n, n, \operatorname{carr}(r)).\; \operatorname{det}(r, n, \operatorname{transpose}(x, n, n)))(c) \\ \quad \quad =\; \operatorname{add}(r)(\operatorname{mul}(r)(s, \\ \quad \quad \quad \quad (\lambda x \in \operatorname{mat}(n, n, \operatorname{carr}(r)).\; \operatorname{det}(r, n, \operatorname{transpose}(x, n, n)))(a)), \\ \quad \quad \quad (\lambda x \in \operatorname{mat}(n, n, \operatorname{carr}(r)).\; \operatorname{det}(r, n, \operatorname{transpose}(x, n, n)))(b)) \end{array}}}
\vnbentry{\texttt{dtt-\allowbreak{}tr-\allowbreak{}alt}}{\vnbstmt{\begin{array}{@{}l@{}} \forall r, n, a, p, q. \\ \quad \operatorname{is\text{-}commutative\text{-}ring}(r) \wedge \\ \quad (n \in \mathbb{N}) \wedge \\ \quad (a \in \operatorname{mat}(n, n, \operatorname{carr}(r))) \wedge \\ \quad (p \in \operatorname{interval}(1, n)) \wedge \\ \quad (q \in \operatorname{interval}(1, n)) \wedge \\ \quad \neg (p = q) \wedge \\ \quad \forall k \in \operatorname{interval}(1, n). \\ \quad \quad (\operatorname{entry}(a, p, k) = \operatorname{entry}(a, q, k)) \Rightarrow \\ \quad \quad (\lambda x \in \operatorname{mat}(n, n, \operatorname{carr}(r)).\; \operatorname{det}(r, n, \operatorname{transpose}(x, n, n)))(a) \\ \quad \quad =\; \operatorname{zero}(r) \end{array}}}
\vnbentry{\texttt{det-\allowbreak{}transpose}}{\vnbstmt{\begin{array}{@{}l@{}} \forall r, n, a. \\ \quad \operatorname{is\text{-}commutative\text{-}ring}(r) \wedge \\ \quad (n \in \mathbb{N}) \wedge \\ \quad (a \in \operatorname{mat}(n, n, \operatorname{carr}(r))) \Rightarrow \\ \quad \quad (\operatorname{det}(r, n, \operatorname{transpose}(a, n, n)) = \operatorname{det}(r, n, a)) \end{array}}}
\vnbentry{\texttt{dtt-\allowbreak{}minor-\allowbreak{}transpose}}{\vnbstmt{\begin{array}{@{}l@{}} \forall r, n, s, p, q. \\ \quad \operatorname{is\text{-}ring}(r) \wedge \\ \quad (n \in \mathbb{N}) \wedge \\ \quad (s \in \operatorname{mat}((n + 1), (n + 1), \operatorname{carr}(r))) \wedge \\ \quad (p \in \mathbb{N}) \wedge \\ \quad (q \in \mathbb{N}) \Rightarrow \\ \quad \quad \operatorname{minor}(\operatorname{transpose}(s, (n + 1), (n + 1)), q, p, n) \\ \quad \quad =\; \operatorname{transpose}(\operatorname{minor}(s, p, q, n), n, n) \end{array}}}
\vnbentry{\texttt{det-\allowbreak{}expand-\allowbreak{}col}}{\vnbstmt{\begin{array}{@{}l@{}} \forall r, n, a, j. \\ \quad \operatorname{is\text{-}commutative\text{-}ring}(r) \wedge \\ \quad (n \in \mathbb{N}) \wedge \\ \quad (a \in \operatorname{mat}((n + 1), (n + 1), \operatorname{carr}(r))) \wedge \\ \quad (j \in \operatorname{interval}(1, (n + 1))) \Rightarrow \\ \quad \quad \operatorname{det}(r, (n + 1), a) \\ \quad \quad =\; \operatorname{finsum}(\operatorname{ring\text{-}additive\text{-}ag}(r), \\ \quad \quad \quad (\lambda e \in \operatorname{interval}(1, (n + 1)).\; \\ \quad \quad \quad \quad \operatorname{mul}(r)(\operatorname{mpow}(\operatorname{ring\text{-}multiplicative\text{-}monoid}(r), \operatorname{neg}(r)(\operatorname{one}(r)), (e + j)), \\ \quad \quad \quad \quad \quad \operatorname{mul}(r)(\operatorname{entry}(a, e, j), \operatorname{det}(r, n, \operatorname{minor}(a, e, j, n))))), \\ \quad \quad \quad \operatorname{interval}(1, (n + 1))) \end{array}}}

\vnbfile{theorem-library/clear-pivot-cross-proof.scm}{theorem-library/clear-pivot-cross-proof.scm}
\noindent{\small clear-pivot-cross -- over a euclidean ring, a matrix P (in MAT(succ p, succ q)) with a nonzero entry is \textasciitilde{}-equivalent to a matrix C2 in BORDER block form [[b,0],[0,C']] with b = C2\_\{1,1\} nonzero.  This is the per-level Smith step: place a class-minimal pivot at (1,1), clear its row and column, and read the}\par
\vnbentry{\texttt{clear-\allowbreak{}pivot-\allowbreak{}cross}}{\vnbstmt{\begin{array}{@{}l@{}} \forall a, p, q, m. \\ \quad \operatorname{is\text{-}euclidean\text{-}ring}(a) \wedge \\ \quad (p \in \mathbb{N}) \wedge \\ \quad (q \in \mathbb{N}) \wedge \\ \quad (m \in \operatorname{mat}((p + 1), (q + 1), \operatorname{carr}(a))) \wedge \\ \quad \exists b \in \operatorname{interval}(1, (p + 1)), c \in \operatorname{interval}(1, (q + 1)). \\ \quad \quad \neg (\operatorname{entry}(m, b, c) = \operatorname{zero}(a)) \Rightarrow \\ \quad \quad \exists d. \\ \quad \quad \quad \operatorname{mat\text{-}equiv}(a, (p + 1), (q + 1), m, d) \wedge \\ \quad \quad \quad (d = \operatorname{border}(a, \operatorname{entry}(d, 1, 1), \operatorname{submat}(d, p, q), p, q)) \wedge \\ \quad \quad \quad \neg (\operatorname{entry}(d, 1, 1) = \operatorname{zero}(a)) \end{array}}}

\vnbfile{theorem-library/smith-diagonalization-proof.scm}{theorem-library/smith-diagonalization-proof.scm}
\noindent{\small smith-diagonalization -- every matrix P in MAT(m,n) over a euclidean ring is \textasciitilde{}-equivalent to a DIAGONAL matrix (Smith normal form, DIAGONALIZATION ONLY -- no divisibility chain).  The capstone of the LA arc.  Induction (ni) on the row dim k, inner statement INNER(k) = forall n P over MAT(k,n) exists diagonal}\par
\vnbentry{\texttt{smith-\allowbreak{}diagonalization}}{\vnbstmt{\begin{array}{@{}l@{}} \forall a.; \forall k \in \mathbb{N}, n, p. \\ \quad \operatorname{is\text{-}euclidean\text{-}ring}(a) \wedge \\ \quad (n \in \mathbb{N}) \wedge \\ \quad (p \in \operatorname{mat}(k, n, \operatorname{carr}(a))) \Rightarrow \\ \quad \quad \exists d. \\ \quad \quad \quad \operatorname{mat\text{-}equiv}(a, k, n, p, d) \wedge \\ \quad \quad \quad \operatorname{is\text{-}diagonal}(a, k, n, d) \end{array}}}

\vnbfile{theorem-library/smith-staircase-proof.scm}{theorem-library/smith-staircase-proof.scm}
\noindent{\small smith-staircase-proof.scm -- the Smith normal form with its nonzero diagonal entries recorded as an INITIAL SEGMENT.}\par
\vnbentry{\texttt{border-\allowbreak{}staircase}}{\vnbstmt{\begin{array}{@{}l@{}} \forall a, b, w, p, q, c. \\ \quad \operatorname{is\text{-}ring}(a) \wedge \\ \quad (p \in \mathbb{N}) \wedge \\ \quad (q \in \mathbb{N}) \wedge \\ \quad (c \in \mathbb{N}) \wedge \\ \quad (b \in \operatorname{carr}(a)) \wedge \\ \quad \neg (b = \operatorname{zero}(a)) \wedge \\ \quad \operatorname{smith\text{-}staircase}(a, p, q, w, c) \wedge \\ \quad (w \in \operatorname{mat}(p, q, \operatorname{carr}(a))) \Rightarrow \\ \quad \quad \operatorname{smith\text{-}staircase}(a, (p + 1), (q + 1), \operatorname{border}(a, b, w, p, q), (c + 1)) \end{array}}}
\vnbentry{\texttt{smith-\allowbreak{}normal-\allowbreak{}form}}{\vnbstmt{\begin{array}{@{}l@{}} \forall a.; \forall k \in \mathbb{N}, n, p. \\ \quad \operatorname{is\text{-}euclidean\text{-}ring}(a) \wedge \\ \quad (n \in \mathbb{N}) \wedge \\ \quad (p \in \operatorname{mat}(k, n, \operatorname{carr}(a))) \Rightarrow \\ \quad \quad \exists d.; \exists b \in \mathbb{N}. \\ \quad \quad \quad \operatorname{mat\text{-}equiv}(a, k, n, p, d) \wedge \\ \quad \quad \quad \operatorname{smith\text{-}staircase}(a, k, n, d, b) \end{array}}}

\vnbfile{theorem-library/matact-row-linear-proof.scm}{theorem-library/matact-row-linear-proof.scm}
\noindent{\small matact-row-linear-proof.scm -- the coefficient row acts LINEARLY.}\par
\vnbentry{\texttt{lincomb-\allowbreak{}unfold}}{\vnbstmt{\begin{array}{@{}l@{}} \forall d, n, c, u. \\ \quad (\operatorname{lincomb}(d, n, c, u) \simeq \operatorname{finsum}(\operatorname{module\text{-}vector\text{-}ag}(d), (\lambda j \in \operatorname{interval}(1, n).\; \operatorname{act}(d)(\operatorname{entry}(c, 1, j), \operatorname{entry}(u, j, 1))), \operatorname{interval}(1, n))) \end{array}}}
\vnbentry{\texttt{lincomb-\allowbreak{}type}}{\vnbstmt{\begin{array}{@{}l@{}} \forall d, n, c, u. \\ \quad \operatorname{is\text{-}module}(d) \wedge \\ \quad (c \in \operatorname{mat}(1, n, \operatorname{carr}(\operatorname{scal}(d)))) \wedge \\ \quad (u \in \operatorname{mat}(n, 1, \operatorname{vec}(d))) \Rightarrow \\ \quad \quad (\operatorname{lincomb}(d, n, c, u) \in \operatorname{vec}(d)) \end{array}}}
\vnbentry{\texttt{matact-\allowbreak{}lincomb}}{\vnbstmt{\begin{array}{@{}l@{}} \forall d, n, c, u. \\ \quad \operatorname{is\text{-}module}(d) \wedge \\ \quad (c \in \operatorname{mat}(1, n, \operatorname{carr}(\operatorname{scal}(d)))) \wedge \\ \quad (u \in \operatorname{mat}(n, 1, \operatorname{vec}(d))) \wedge \\ \quad (1 \leq n) \Rightarrow \\ \quad \quad (\operatorname{entry}(\operatorname{matact}(d, c, u), 1, 1) = \operatorname{lincomb}(d, n, c, u)) \end{array}}}
\vnbentry{\texttt{lincomb-\allowbreak{}row-\allowbreak{}add}}{\vnbstmt{\begin{array}{@{}l@{}} \forall d, n, a, b, u. \\ \quad \operatorname{is\text{-}module}(d) \wedge \\ \quad (a \in \operatorname{mat}(1, n, \operatorname{carr}(\operatorname{scal}(d)))) \wedge \\ \quad (b \in \operatorname{mat}(1, n, \operatorname{carr}(\operatorname{scal}(d)))) \wedge \\ \quad (u \in \operatorname{mat}(n, 1, \operatorname{vec}(d))) \Rightarrow \\ \quad \quad \operatorname{lincomb}(d, n, \operatorname{matadd}(\operatorname{scal}(d), a, b), u) \\ \quad \quad =\; \operatorname{vadd}(d)(\operatorname{lincomb}(d, n, a, u), \operatorname{lincomb}(d, n, b, u)) \end{array}}}
\vnbentry{\texttt{lincomb-\allowbreak{}row-\allowbreak{}scale}}{\vnbstmt{\begin{array}{@{}l@{}} \forall d, n, r, c, u. \\ \quad \operatorname{is\text{-}module}(d) \wedge \\ \quad (r \in \operatorname{carr}(\operatorname{scal}(d))) \wedge \\ \quad (c \in \operatorname{mat}(1, n, \operatorname{carr}(\operatorname{scal}(d)))) \wedge \\ \quad (u \in \operatorname{mat}(n, 1, \operatorname{vec}(d))) \Rightarrow \\ \quad \quad \operatorname{lincomb}(d, n, \operatorname{matscale}(\operatorname{scal}(d), r, c), u) \\ \quad \quad =\; \operatorname{act}(d)(r, \operatorname{lincomb}(d, n, c, u)) \end{array}}}

\vnbfile{theorem-library/span-bricks-proof.scm}{theorem-library/span-bricks-proof.scm}
\noindent{\small span-bricks-proof.scm -- the module-action bricks under spans-submodule-fg.}\par
\vnbentry{\texttt{module-\allowbreak{}act-\allowbreak{}neg-\allowbreak{}one}}{\vnbstmt{\begin{array}{@{}l@{}} \forall d.; \forall x \in \operatorname{vec}(d). \\ \quad \operatorname{is\text{-}module}(d) \Rightarrow  \\ \quad \quad \operatorname{act}(d)(\operatorname{neg}(\operatorname{scal}(d))(\operatorname{one}(\operatorname{scal}(d))), x) \\ \quad \quad =\; \operatorname{vneg}(d)(x) \end{array}}}
\vnbentry{\texttt{lincomb-\allowbreak{}zerorow}}{\vnbstmt{\begin{array}{@{}l@{}} \forall d, n.; \forall u \in \operatorname{mat}(n, 1, \operatorname{vec}(d)). \\ \quad \operatorname{is\text{-}module}(d) \Rightarrow  \\ \quad \quad \operatorname{lincomb}(d, n, \operatorname{zeromat}(\operatorname{scal}(d), 1, n), u) \\ \quad \quad =\; \operatorname{vzero}(d) \end{array}}}
\vnbentry{\texttt{lincomb-\allowbreak{}unitrow}}{\vnbstmt{\begin{array}{@{}l@{}} \forall d, n, i.; \forall u \in \operatorname{mat}(n, 1, \operatorname{vec}(d)). \\ \quad (\operatorname{is\text{-}module}(d) \wedge (i \in \operatorname{interval}(1, n))) \Rightarrow  \\ \quad \quad \operatorname{lincomb}(d, n, \operatorname{unitrow}(\operatorname{scal}(d), n, i), u) \\ \quad \quad =\; \operatorname{entry}(u, i, 1) \end{array}}}
\vnbentry{\texttt{lincomb-\allowbreak{}empty}}{\vnbstmt{\begin{array}{@{}l@{}} \forall d, c, u. \\ \quad \operatorname{is\text{-}module}(d) \Rightarrow  \\ \quad \quad (\operatorname{lincomb}(d, 0, c, u) = \operatorname{vzero}(d)) \end{array}}}
\vnbentry{\texttt{span-\allowbreak{}is-\allowbreak{}submodule}}{\vnbstmt{\begin{array}{@{}l@{}} \forall d, n.; \forall u \in \operatorname{mat}(n, 1, \operatorname{vec}(d)). \\ \quad \operatorname{is\text{-}module}(d) \Rightarrow  \\ \quad \quad \operatorname{is\text{-}submodule}(d, \operatorname{span}(d, n, u)) \end{array}}}
\vnbentry{\texttt{spans-\allowbreak{}span}}{\vnbstmt{\begin{array}{@{}l@{}} \forall d, n.; \forall u \in \operatorname{mat}(n, 1, \operatorname{vec}(d)). \\ \quad \operatorname{is\text{-}module}(d) \Rightarrow  \\ \quad \quad \operatorname{spans}(d, n, u, \operatorname{span}(d, n, u)) \end{array}}}

\vnbfile{theorem-library/span-bricks2-proof.scm}{theorem-library/span-bricks2-proof.scm}
\noindent{\small span-bricks2-proof.scm -- bricks 4, 5, 6 under spans-submodule-fg.}\par
\vnbentry{\texttt{lincomb-\allowbreak{}row-\allowbreak{}peel}}{\vnbstmt{\begin{array}{@{}l@{}} \forall d.; \forall n \in \mathbb{N}, c \in \operatorname{mat}(1, (n + 1), \operatorname{carr}(\operatorname{scal}(d))), u \in \operatorname{mat}((n + 1), 1, \operatorname{vec}(d)). \\ \quad \operatorname{is\text{-}module}(d) \Rightarrow  \\ \quad \quad \operatorname{lincomb}(d, (n + 1), c, u) \\ \quad \quad =\; \operatorname{vadd}(d)(\operatorname{lincomb}(d, n, \operatorname{block}(c, 1, n), \operatorname{block}(u, n, 1)), \\ \quad \quad \quad \operatorname{act}(d)(\operatorname{entry}(c, 1, (n + 1)), \operatorname{entry}(u, (n + 1), 1))) \end{array}}}
\vnbentry{\texttt{lincomb-\allowbreak{}snoc}}{\vnbstmt{\begin{array}{@{}l@{}} \forall d.; \forall n \in \mathbb{N}, c \in \operatorname{mat}(1, n, \operatorname{carr}(\operatorname{scal}(d))), u \in \operatorname{mat}(n, 1, \operatorname{vec}(d)), r \in \operatorname{carr}(\operatorname{scal}(d)), x \in \operatorname{vec}(d). \\ \quad \operatorname{is\text{-}module}(d) \Rightarrow  \\ \quad \quad \operatorname{lincomb}(d, (n + 1), \operatorname{snoc\text{-}row}(c, n, r), \operatorname{snoc\text{-}col}(u, n, x)) \\ \quad \quad =\; \operatorname{vadd}(d)(\operatorname{lincomb}(d, n, c, u), \operatorname{act}(d)(r, x)) \end{array}}}
\vnbentry{\texttt{submodule-\allowbreak{}intersection}}{\vnbstmt{\begin{array}{@{}l@{}} \forall d, a, b. \\ \quad \operatorname{is\text{-}module}(d) \wedge \\ \quad \operatorname{is\text{-}submodule}(d, a) \wedge \\ \quad \operatorname{is\text{-}submodule}(d, b) \Rightarrow \\ \quad \quad \operatorname{is\text{-}submodule}(d, (a \cap b)) \end{array}}}
\vnbentry{\texttt{mat-\allowbreak{}1-\allowbreak{}0-\allowbreak{}nonempty}}{\vnbstmt{\forall x.\; \exists p.\; (p \in \operatorname{mat}(1, 0, x))}}
\vnbentry{\texttt{mat-\allowbreak{}0-\allowbreak{}1-\allowbreak{}nonempty}}{\vnbstmt{\forall x.\; \exists p.\; (p \in \operatorname{mat}(0, 1, x))}}

\vnbfile{theorem-library/lastcoeff-ideal-proof.scm}{theorem-library/lastcoeff-ideal-proof.scm}
\noindent{\small lastcoeff-ideal-proof.scm -- the two self-contained lemmas the spans-submodule-fg descent leans on (no induction here).}\par
\vnbentry{\texttt{lastcoeff-\allowbreak{}zero-\allowbreak{}in-\allowbreak{}span}}{\vnbstmt{\begin{array}{@{}l@{}} \forall d.; \forall p \in \mathbb{N}, u \in \operatorname{mat}((p + 1), 1, \operatorname{vec}(d)), c \in \operatorname{mat}(1, (p + 1), \operatorname{carr}(\operatorname{scal}(d))). \\ \quad \operatorname{is\text{-}module}(d) \wedge \\ \quad (\operatorname{entry}(c, 1, (p + 1)) = \operatorname{zero}(\operatorname{scal}(d))) \Rightarrow \\ \quad \quad \operatorname{lincomb}(d, (p + 1), c, u) \\ \quad \quad \in\; \operatorname{span}(d, p, \operatorname{block}(u, p, 1)) \end{array}}}
\vnbentry{\texttt{lastcoeff-\allowbreak{}set-\allowbreak{}is-\allowbreak{}ideal}}{\vnbstmt{\begin{array}{@{}l@{}} \forall d.; \forall p \in \mathbb{N}, u \in \operatorname{mat}((p + 1), 1, \operatorname{vec}(d)), m. \\ \quad \operatorname{is\text{-}module}(d) \wedge \\ \quad \operatorname{is\text{-}commutative\text{-}ring}(\operatorname{scal}(d)) \wedge \\ \quad \operatorname{is\text{-}submodule}(d, m) \Rightarrow \\ \quad \quad \operatorname{is\text{-}ideal}(\operatorname{scal}(d), \operatorname{lastcoeff\text{-}set}(d, p, u, m)) \end{array}}}
\vnbentry{\texttt{descent-\allowbreak{}remainder}}{\vnbstmt{\begin{array}{@{}l@{}} \forall d.; \forall n \in \mathbb{N}, u \in \operatorname{mat}((n + 1), 1, \operatorname{vec}(d)), c, a \in \operatorname{mat}(1, (n + 1), \operatorname{carr}(\operatorname{scal}(d))), q, b \in \operatorname{carr}(\operatorname{scal}(d)). \\ \quad \operatorname{is\text{-}module}(d) \wedge \\ \quad \operatorname{is\text{-}commutative\text{-}ring}(\operatorname{scal}(d)) \wedge \\ \quad (\operatorname{entry}(a, 1, (n + 1)) = b) \wedge \\ \quad (\operatorname{entry}(c, 1, (n + 1)) = \operatorname{mul}(\operatorname{scal}(d))(q, b)) \Rightarrow \\ \quad \quad \exists y \in \operatorname{span}(d, n, \operatorname{block}(u, n, 1)). \\ \quad \quad \quad \operatorname{lincomb}(d, (n + 1), c, u) \\ \quad \quad \quad =\; \operatorname{vadd}(d)(y, \operatorname{act}(d)(q, \operatorname{lincomb}(d, (n + 1), a, u))) \wedge \\ \quad \quad \quad y \\ \quad \quad \quad =\; \operatorname{vadd}(d)(\operatorname{lincomb}(d, (n + 1), c, u), \\ \quad \quad \quad \quad \operatorname{act}(d)(\operatorname{neg}(\operatorname{scal}(d))(q), \operatorname{lincomb}(d, (n + 1), a, u))) \end{array}}}

\vnbfile{theorem-library/spans-submodule-fg-proof.scm}{theorem-library/spans-submodule-fg-proof.scm}
\noindent{\small spans-submodule-fg-proof.scm -- THE DESCENT.}\par
\vnbentry{\texttt{spans-\allowbreak{}fg-\allowbreak{}base}}{\vnbstmt{\begin{array}{@{}l@{}} \forall d.; \forall u \in \operatorname{mat}(0, 1, \operatorname{vec}(d)), m, a. \\ \quad \operatorname{is\text{-}module}(d) \wedge \\ \quad \operatorname{is\text{-}euclidean\text{-}ring}(\operatorname{scal}(d)) \wedge \\ \quad \operatorname{is\text{-}submodule}(d, m) \wedge \\ \quad \operatorname{spans}(d, 0, u, m) \wedge \\ \quad \operatorname{is\text{-}submodule}(d, a) \wedge \\ \quad (a \subseteq m) \Rightarrow \\ \quad \quad \exists k. \\ \quad \quad \quad (k \in \mathbb{N}) \wedge \\ \quad \quad \quad (k \leq 0) \wedge \\ \quad \quad \quad \exists w \in \operatorname{mat}(k, 1, \operatorname{vec}(d)). \\ \quad \quad \quad \quad \operatorname{spans}(d, k, w, a) \end{array}}}
\vnbentry{\texttt{spans-\allowbreak{}fg-\allowbreak{}step}}{\vnbstmt{\begin{array}{@{}l@{}} \forall d.; \forall n \in \mathbb{N}, u \in \operatorname{mat}((n + 1), 1, \operatorname{vec}(d)), m, a. \\ \quad \operatorname{is\text{-}module}(d) \wedge \\ \quad \operatorname{is\text{-}euclidean\text{-}ring}(\operatorname{scal}(d)) \wedge \\ \quad \forall u \in \operatorname{mat}(n, 1, \operatorname{vec}(d)), m, a. \\ \quad \quad \operatorname{is\text{-}submodule}(d, m) \wedge \\ \quad \quad \operatorname{spans}(d, n, u, m) \wedge \\ \quad \quad \operatorname{is\text{-}submodule}(d, a) \wedge \\ \quad \quad (a \subseteq m) \Rightarrow \\ \quad \quad \quad \exists k. \\ \quad \quad \quad \quad (k \in \mathbb{N}) \wedge \\ \quad \quad \quad \quad (k \leq n) \wedge \\ \quad \quad \quad \quad \exists w \in \operatorname{mat}(k, 1, \operatorname{vec}(d)). \\ \quad \quad \quad \quad \quad \operatorname{spans}(d, k, w, a) \wedge \\ \quad \operatorname{is\text{-}submodule}(d, m) \wedge \\ \quad \operatorname{spans}(d, (n + 1), u, m) \wedge \\ \quad \operatorname{is\text{-}submodule}(d, a) \wedge \\ \quad (a \subseteq m) \Rightarrow \\ \quad \quad \exists k. \\ \quad \quad \quad (k \in \mathbb{N}) \wedge \\ \quad \quad \quad (k \leq (n + 1)) \wedge \\ \quad \quad \quad \exists w \in \operatorname{mat}(k, 1, \operatorname{vec}(d)). \\ \quad \quad \quad \quad \operatorname{spans}(d, k, w, a) \end{array}}}
\vnbentry{\texttt{spans-\allowbreak{}submodule-\allowbreak{}fg}}{\vnbstmt{\begin{array}{@{}l@{}} \forall d.; \forall n \in \mathbb{N}, u \in \operatorname{mat}(n, 1, \operatorname{vec}(d)), m, a. \\ \quad \operatorname{is\text{-}module}(d) \wedge \\ \quad \operatorname{is\text{-}euclidean\text{-}ring}(\operatorname{scal}(d)) \wedge \\ \quad \operatorname{is\text{-}submodule}(d, m) \wedge \\ \quad \operatorname{spans}(d, n, u, m) \wedge \\ \quad \operatorname{is\text{-}submodule}(d, a) \wedge \\ \quad (a \subseteq m) \Rightarrow \\ \quad \quad \exists k. \\ \quad \quad \quad (k \in \mathbb{N}) \wedge \\ \quad \quad \quad (k \leq n) \wedge \\ \quad \quad \quad \exists w \in \operatorname{mat}(k, 1, \operatorname{vec}(d)). \\ \quad \quad \quad \quad \operatorname{spans}(d, k, w, a) \end{array}}}

\vnbfile{theorem-library/submodule-fg-proof.scm}{theorem-library/submodule-fg-proof.scm}
\noindent{\small submodule-fg-proof.scm -- submodule-fg, PROVEN as a corollary of the SPANS-relativized statement it should always have had.}\par
\vnbentry{\texttt{whole-\allowbreak{}module-\allowbreak{}is-\allowbreak{}submodule}}{\vnbstmt{\begin{array}{@{}l@{}} \forall d. \\ \quad \operatorname{is\text{-}module}(d) \Rightarrow  \\ \quad \quad \operatorname{is\text{-}submodule}(d, \operatorname{vec}(d)) \end{array}}}
\vnbentry{\texttt{generates-\allowbreak{}implies-\allowbreak{}spans-\allowbreak{}vec}}{\vnbstmt{\begin{array}{@{}l@{}} \forall d.; \forall n \in \mathbb{N}, u \in \operatorname{mat}(n, 1, \operatorname{vec}(d)). \\ \quad (\operatorname{is\text{-}module}(d) \wedge \operatorname{generates}(d, n, u)) \Rightarrow  \\ \quad \quad \operatorname{spans}(d, n, u, \operatorname{vec}(d)) \end{array}}}
\vnbentry{\texttt{submodule-\allowbreak{}fg}}{\vnbstmt{\begin{array}{@{}l@{}} \forall d.; \forall n \in \mathbb{N}, u \in \operatorname{mat}(n, 1, \operatorname{vec}(d)), m. \\ \quad \operatorname{is\text{-}module}(d) \wedge \\ \quad \operatorname{is\text{-}euclidean\text{-}ring}(\operatorname{scal}(d)) \wedge \\ \quad \operatorname{generates}(d, n, u) \wedge \\ \quad \operatorname{is\text{-}submodule}(d, m) \Rightarrow \\ \quad \quad \exists k. \\ \quad \quad \quad (k \in \mathbb{N}) \wedge \\ \quad \quad \quad (k \leq n) \wedge \\ \quad \quad \quad \exists w \in \operatorname{mat}(k, 1, \operatorname{vec}(d)). \\ \quad \quad \quad \quad \operatorname{spans}(d, k, w, m) \end{array}}}

\vnbfile{theorem-library/matact-assoc-proof.scm}{theorem-library/matact-assoc-proof.scm}
\noindent{\small matact-assoc-proof.scm -- Remark 3.39: the matrix action on a sequence of module elements is associative.}\par
\vnbentry{\texttt{matact-\allowbreak{}triple-\allowbreak{}left}}{\vnbstmt{\begin{array}{@{}l@{}} \forall d, m, n, k, l, p, q, u, w, a. \\ \quad \operatorname{is\text{-}module}(d) \wedge \\ \quad (p \in \operatorname{mat}(m, n, \operatorname{carr}(\operatorname{scal}(d)))) \wedge \\ \quad (q \in \operatorname{mat}(n, k, \operatorname{carr}(\operatorname{scal}(d)))) \wedge \\ \quad (u \in \operatorname{mat}(k, l, \operatorname{vec}(d))) \wedge \\ \quad (1 \leq n) \wedge \\ \quad (1 \leq k) \wedge \\ \quad (w \in \operatorname{interval}(1, m)) \wedge \\ \quad (a \in \operatorname{interval}(1, l)) \Rightarrow \\ \quad \quad \operatorname{entry}(\operatorname{matact}(d, \operatorname{matmul}(\operatorname{scal}(d), p, q), u), w, a) \\ \quad \quad =\; \operatorname{finsum}(\operatorname{module\text{-}vector\text{-}ag}(d), \\ \quad \quad \quad (\lambda c \in \operatorname{interval}(1, k).\; \\ \quad \quad \quad \quad \operatorname{finsum}(\operatorname{module\text{-}vector\text{-}ag}(d), \\ \quad \quad \quad \quad \quad (\lambda j \in \operatorname{interval}(1, n).\; \\ \quad \quad \quad \quad \quad \quad (\lambda z \in (\operatorname{interval}(1, k) \times \operatorname{interval}(1, n)).\; \operatorname{act}(d)(\operatorname{mul}(\operatorname{scal}(d))(\operatorname{entry}(p, w, \operatorname{nth}(2, z)), \operatorname{entry}(q, \operatorname{nth}(2, z), \operatorname{nth}(1, z))), \operatorname{entry}(u, \operatorname{nth}(1, z), a)))(\langle c, j \rangle)), \\ \quad \quad \quad \quad \quad \operatorname{interval}(1, n))), \\ \quad \quad \quad \operatorname{interval}(1, k)) \end{array}}}
\vnbentry{\texttt{matact-\allowbreak{}triple-\allowbreak{}right}}{\vnbstmt{\begin{array}{@{}l@{}} \forall d, m, n, k, l, p, q, u, w, a. \\ \quad \operatorname{is\text{-}module}(d) \wedge \\ \quad (p \in \operatorname{mat}(m, n, \operatorname{carr}(\operatorname{scal}(d)))) \wedge \\ \quad (q \in \operatorname{mat}(n, k, \operatorname{carr}(\operatorname{scal}(d)))) \wedge \\ \quad (u \in \operatorname{mat}(k, l, \operatorname{vec}(d))) \wedge \\ \quad (1 \leq n) \wedge \\ \quad (1 \leq k) \wedge \\ \quad (w \in \operatorname{interval}(1, m)) \wedge \\ \quad (a \in \operatorname{interval}(1, l)) \Rightarrow \\ \quad \quad \operatorname{entry}(\operatorname{matact}(d, p, \operatorname{matact}(d, q, u)), w, a) \\ \quad \quad =\; \operatorname{finsum}(\operatorname{module\text{-}vector\text{-}ag}(d), \\ \quad \quad \quad (\lambda j \in \operatorname{interval}(1, n).\; \\ \quad \quad \quad \quad \operatorname{finsum}(\operatorname{module\text{-}vector\text{-}ag}(d), \\ \quad \quad \quad \quad \quad (\lambda c \in \operatorname{interval}(1, k).\; \\ \quad \quad \quad \quad \quad \quad (\lambda z \in (\operatorname{interval}(1, k) \times \operatorname{interval}(1, n)).\; \operatorname{act}(d)(\operatorname{mul}(\operatorname{scal}(d))(\operatorname{entry}(p, w, \operatorname{nth}(2, z)), \operatorname{entry}(q, \operatorname{nth}(2, z), \operatorname{nth}(1, z))), \operatorname{entry}(u, \operatorname{nth}(1, z), a)))(\langle c, j \rangle)), \\ \quad \quad \quad \quad \quad \operatorname{interval}(1, k))), \\ \quad \quad \quad \operatorname{interval}(1, n)) \end{array}}}
\vnbentry{\texttt{matact-\allowbreak{}assoc}}{\vnbstmt{\begin{array}{@{}l@{}} \forall d, m, n, k, l, p, q, u. \\ \quad \operatorname{is\text{-}module}(d) \wedge \\ \quad (p \in \operatorname{mat}(m, n, \operatorname{carr}(\operatorname{scal}(d)))) \wedge \\ \quad (q \in \operatorname{mat}(n, k, \operatorname{carr}(\operatorname{scal}(d)))) \wedge \\ \quad (u \in \operatorname{mat}(k, l, \operatorname{vec}(d))) \wedge \\ \quad ((n = 0) \Rightarrow ((m = 0) \vee ((k = 0) \wedge (l = 0)))) \wedge \\ \quad ((k = 0) \Rightarrow ((l = 0) \vee ((n = 0) \wedge (m = 0)))) \Rightarrow \\ \quad \quad \operatorname{matact}(d, \operatorname{matmul}(\operatorname{scal}(d), p, q), u) \\ \quad \quad =\; \operatorname{matact}(d, p, \operatorname{matact}(d, q, u)) \end{array}}}

\vnbfile{theorem-library/mod-basis-proof.scm}{theorem-library/mod-basis-proof.scm}
\noindent{\small mod-basis-proof.scm -- Lemma 3.40 (algebraic-numbers.pdf ch.3 sec 8.2): an invertible matrix carries generating sequences to generating sequences and relation-free sequences to relation-free sequences.  Plus the identity action I . u = u that says MATACT is a genuine action.}\par
\vnbentry{\texttt{matact-\allowbreak{}identmat}}{\vnbstmt{\begin{array}{@{}l@{}} \forall d, n, q.; \forall u \in \operatorname{mat}(n, q, \operatorname{vec}(d)). \\ \quad \operatorname{is\text{-}module}(d) \Rightarrow  \\ \quad \quad (\operatorname{matact}(d, \operatorname{identmat}(\operatorname{scal}(d), n), u) = u) \end{array}}}
\vnbentry{\texttt{generates-\allowbreak{}transport}}{\vnbstmt{\begin{array}{@{}l@{}} \forall d, n, u, m. \\ \quad \operatorname{is\text{-}module}(d) \wedge \\ \quad (u \in \operatorname{mat}(n, 1, \operatorname{vec}(d))) \wedge \\ \quad \operatorname{is\text{-}invertible\text{-}mat}(\operatorname{scal}(d), n, m) \wedge \\ \quad \operatorname{generates}(d, n, u) \Rightarrow \\ \quad \quad \operatorname{generates}(d, n, \operatorname{matact}(d, m, u)) \end{array}}}
\vnbentry{\texttt{free-\allowbreak{}transport}}{\vnbstmt{\begin{array}{@{}l@{}} \forall d, m, v, a. \\ \quad \operatorname{is\text{-}module}(d) \wedge \\ \quad (v \in \operatorname{mat}(m, 1, \operatorname{vec}(d))) \wedge \\ \quad \operatorname{is\text{-}invertible\text{-}mat}(\operatorname{scal}(d), m, a) \wedge \\ \quad \operatorname{rel\text{-}free}(d, m, v) \Rightarrow \\ \quad \quad \operatorname{rel\text{-}free}(d, m, \operatorname{matact}(d, a, v)) \end{array}}}

\vnbfile{theorem-library/inverse-invertible-proof.scm}{theorem-library/inverse-invertible-proof.scm}
\noindent{\small inverse-invertible-proof.scm -- the inverse of an invertible matrix is invertible.}\par
\vnbentry{\texttt{inverse-\allowbreak{}is-\allowbreak{}invertible}}{\vnbstmt{\begin{array}{@{}l@{}} \forall a, n, m, b. \\ \quad \operatorname{is\text{-}ring}(a) \wedge \\ \quad (m \in \operatorname{mat}(n, n, \operatorname{carr}(a))) \wedge \\ \quad (b \in \operatorname{mat}(n, n, \operatorname{carr}(a))) \wedge \\ \quad (\operatorname{matmul}(a, m, b) = \operatorname{identmat}(a, n)) \wedge \\ \quad (\operatorname{matmul}(a, b, m) = \operatorname{identmat}(a, n)) \Rightarrow \\ \quad \quad \operatorname{is\text{-}invertible\text{-}mat}(a, n, b) \end{array}}}

\vnbfile{theorem-library/spans-transport-proof.scm}{theorem-library/spans-transport-proof.scm}
\noindent{\small spans-transport-proof.scm -- the SPANS-relativized generates-transport.}\par
\vnbentry{\texttt{spans-\allowbreak{}transport}}{\vnbstmt{\begin{array}{@{}l@{}} \forall d, n, u, m, a. \\ \quad \operatorname{is\text{-}module}(d) \wedge \\ \quad (u \in \operatorname{mat}(n, 1, \operatorname{vec}(d))) \wedge \\ \quad \operatorname{is\text{-}invertible\text{-}mat}(\operatorname{scal}(d), n, m) \wedge \\ \quad \operatorname{is\text{-}submodule}(d, a) \wedge \\ \quad \operatorname{spans}(d, n, u, a) \Rightarrow \\ \quad \quad \operatorname{spans}(d, n, \operatorname{matact}(d, m, u), a) \end{array}}}

\vnbfile{theorem-library/rank-bound-proof.scm}{theorem-library/rank-bound-proof.scm}
\noindent{\small rank-bound-proof.scm -- Proposition 3.41 (algebraic-numbers.pdf ch.3 sec 8.2), the payoff of the whole Smith arc:}\par
\vnbentry{\texttt{free-\allowbreak{}length-\allowbreak{}le-\allowbreak{}generators}}{\vnbstmt{\begin{array}{@{}l@{}} \forall d, n, m, u, v. \\ \quad \operatorname{is\text{-}module}(d) \wedge \\ \quad \operatorname{is\text{-}euclidean\text{-}ring}(\operatorname{scal}(d)) \wedge \\ \quad (n \in \mathbb{N}) \wedge \\ \quad (m \in \mathbb{N}) \wedge \\ \quad (u \in \operatorname{mat}(n, 1, \operatorname{vec}(d))) \wedge \\ \quad (v \in \operatorname{mat}(m, 1, \operatorname{vec}(d))) \wedge \\ \quad \operatorname{generates}(d, n, u) \wedge \\ \quad \operatorname{rel\text{-}free}(d, m, v) \Rightarrow \\ \quad \quad (m \leq n) \end{array}}}

\vnbfile{theorem-library/binomial-proof.scm}{theorem-library/binomial-proof.scm}
\noindent{\small binomial-proof.scm -- machine proof of the Binomial Theorem (SUM form), PROVEN by ordinary induction, IMPS-style (docs/algebra.t).}\par
\vnbentry{\texttt{sum-\allowbreak{}expansion}}{\vnbstmt{\begin{array}{@{}l@{}} \forall n \in \mathbb{N}, r.; \forall x, y \in \operatorname{carr}(r), g \in (\mathbb{Z} \to \operatorname{carr}(r)). \\ \quad \operatorname{is\text{-}commutative\text{-}ring}(r) \Rightarrow  \\ \quad \quad \operatorname{sum}(r, \\ \quad \quad \quad (\lambda k \in \mathbb{Z}.\; \\ \quad \quad \quad \quad \operatorname{add}(r)(\operatorname{mul}(r)(x, g((k - 1))), \operatorname{mul}(r)(y, g(k)))), \\ \quad \quad \quad (n + 1)) \\ \quad \quad =\; \operatorname{add}(r)(\operatorname{add}(r)(\operatorname{mul}(r)(\operatorname{add}(r)(x, y), \operatorname{sum}(r, g, n)), \operatorname{mul}(r)(x, g((0 - 1)))), \\ \quad \quad \quad \operatorname{mul}(r)(y, g(n))) \end{array}}}
\vnbentry{\texttt{binomial-\allowbreak{}theorem}}{\vnbstmt{\begin{array}{@{}l@{}} \forall n \in \mathbb{N}, r.; \forall x, y \in \operatorname{carr}(r). \\ \quad \operatorname{is\text{-}commutative\text{-}ring}(r) \Rightarrow  \\ \quad \quad \operatorname{ring\text{-}power}(r, \operatorname{add}(r)(x, y), n) \\ \quad \quad =\; \operatorname{sum}(r, \operatorname{comb\text{-}kk}(r, x, y, n), (n + 1)) \end{array}}}

\vnbfile{theorem-library/poly-is-ring-proof.scm}{theorem-library/poly-is-ring-proof.scm}
\noindent{\small poly-is-ring-proof.scm -- POLY(A) is a ring when A is a ring.}\par
\vnbentry{\texttt{poly-\allowbreak{}is-\allowbreak{}ring}}{\vnbstmt{\begin{array}{@{}l@{}} \forall a. \\ \quad \operatorname{is\text{-}ring}(a) \Rightarrow  \\ \quad \quad \operatorname{is\text{-}ring}(\operatorname{poly}(a)) \end{array}}}

\vnbfile{theorem-library/ccint-basics.scm}{theorem-library/ccint-basics.scm}
\noindent{\small ccint-basics.scm -- CCINT's membership law, PROVEN from its definition.}\par
\vnbentry{\texttt{ccint-\allowbreak{}membership}}{\vnbstmt{\begin{array}{@{}l@{}} \forall a, b, x. \\ \quad ((x \in \operatorname{ccint}(a, b)) \Leftrightarrow ((x \in \mathbb{R}) \wedge ((a \leq x) \wedge (x \leq b)))) \end{array}}}
\vnbentry{\texttt{ccint-\allowbreak{}elt-\allowbreak{}in-\allowbreak{}rr}}{\vnbstmt{\forall a, b.; \forall x \in \operatorname{ccint}(a, b).\; (x \in \mathbb{R})}}

\vnbfile{theorem-library/ivt-proof.scm}{theorem-library/ivt-proof.scm}
\noindent{\small ivt-proof.scm -- the INTERMEDIATE VALUE THEOREM, PROVEN.}\par
\vnbentry{\texttt{ivt}}{\vnbstmt{\begin{array}{@{}l@{}} \forall f, a, b, c. \\ \quad (f \in (\mathbb{R} \to \mathbb{R})) \wedge \\ \quad (a \in \mathbb{R}) \wedge \\ \quad (b \in \mathbb{R}) \wedge \\ \quad (c \in \mathbb{R}) \wedge \\ \quad (a \leq b) \wedge \\ \quad \forall x \in \operatorname{ccint}(a, b). \\ \quad \quad \operatorname{is\text{-}continuous\text{-}at}(\operatorname{rr\text{-}ms}, \operatorname{rr\text{-}ms}, f, x) \wedge \\ \quad (f(a) \leq c) \wedge \\ \quad (c \leq f(b)) \Rightarrow \\ \quad \quad \exists p \in \mathbb{R}. \\ \quad \quad \quad ((p \in \operatorname{ccint}(a, b)) \wedge (f(p) = c)) \end{array}}}

\vnbfile{theorem-library/monotone-inverse.scm}{theorem-library/monotone-inverse.scm}
\noindent{\small monotone-inverse.scm -- MONOTONICITY VOCABULARY FOR REAL FUNCTIONS, and the EXISTENCE HALF of the inverse function theorem on a closed interval. Everything here is `modulo 0'.}\par
\vnbentry{\texttt{ccint-\allowbreak{}subset-\allowbreak{}rr}}{\vnbstmt{\forall a, b.\; (\operatorname{ccint}(a, b) \subseteq \mathbb{R})}}
\vnbentry{\texttt{strict-\allowbreak{}increase-\allowbreak{}injective}}{\vnbstmt{\begin{array}{@{}l@{}} \forall f, s.; \forall u, v \in s. \\ \quad (\operatorname{is\text{-}strictly\text{-}increasing\text{-}on}(f, s) \wedge (f(u) = f(v))) \Rightarrow  \\ \quad \quad (u = v) \end{array}}}
\vnbentry{\texttt{monotone-\allowbreak{}ccint-\allowbreak{}inverse-\allowbreak{}exists}}{\vnbstmt{\begin{array}{@{}l@{}} \forall f, a, b.; \forall y \in \operatorname{ccint}(f(a), f(b)). \\ \quad (f \in (\mathbb{R} \to \mathbb{R})) \wedge \\ \quad (a \in \mathbb{R}) \wedge \\ \quad (b \in \mathbb{R}) \wedge \\ \quad (a \leq b) \wedge \\ \quad \forall x \in \operatorname{ccint}(a, b). \\ \quad \quad \operatorname{is\text{-}continuous\text{-}at}(\operatorname{rr\text{-}ms}, \operatorname{rr\text{-}ms}, f, x) \wedge \\ \quad \operatorname{is\text{-}strictly\text{-}increasing\text{-}on}(f, \operatorname{ccint}(a, b)) \Rightarrow \\ \quad \quad \exists p \in \operatorname{ccint}(a, b). \\ \quad \quad \quad ((f(p) = y) \wedge \forall q \in \operatorname{ccint}(a, b).\; ((f(q) = y) \Rightarrow (p = q))) \end{array}}}
\vnbentry{\texttt{monotone-\allowbreak{}ccint-\allowbreak{}inverse-\allowbreak{}order}}{\vnbstmt{\begin{array}{@{}l@{}} \forall f, a, b.; \forall u, v \in \operatorname{ccint}(a, b). \\ \quad \operatorname{is\text{-}strictly\text{-}increasing\text{-}on}(f, \operatorname{ccint}(a, b)) \wedge \\ \quad (f(u) < f(v)) \Rightarrow \\ \quad \quad (u < v) \end{array}}}

\vnbfile{theorem-library/sq-continuous.scm}{theorem-library/sq-continuous.scm}
\noindent{\small sq-continuous.scm -- z |-$>$ z*z is a function RR -$>$ RR and is CONTINUOUS at every real point, PROVEN.}\par
\vnbentry{\texttt{sq-\allowbreak{}fun-\allowbreak{}in-\allowbreak{}fun}}{\vnbstmt{((\lambda z \in \mathbb{R}.\; (z \cdot z)) \in (\mathbb{R} \to \mathbb{R}))}}
\vnbentry{\texttt{sq-\allowbreak{}continuous-\allowbreak{}at}}{\vnbstmt{\begin{array}{@{}l@{}} \forall t \in \mathbb{R}. \\ \quad \operatorname{is\text{-}continuous\text{-}at}(\operatorname{rr\text{-}ms}, \operatorname{rr\text{-}ms}, (\lambda z \in \mathbb{R}.\; (z \cdot z)), t) \end{array}}}

\vnbfile{theorem-library/sqrt-defined.scm}{theorem-library/sqrt-defined.scm}
\noindent{\small sqrt-defined.scm -- SQRT DEFINED by the intermediate value theorem, and the five supports that used to characterise it RETIRED as theorems.}\par
\vnbentry{\texttt{sqrt-\allowbreak{}exists}}{\vnbstmt{\begin{array}{@{}l@{}} \forall a \in \mathbb{R}. \\ \quad (0 \leq a) \Rightarrow  \\ \quad \quad \exists x \in \mathbb{R}.\; ((0 \leq x) \wedge ((x \cdot x) = a)) \end{array}}}
\vnbentry{\texttt{sqrt-\allowbreak{}unique}}{\vnbstmt{\begin{array}{@{}l@{}} \forall x, y. \\ \quad (x \in \mathbb{R}) \wedge \\ \quad (y \in \mathbb{R}) \wedge \\ \quad (0 \leq x) \wedge \\ \quad (0 \leq y) \wedge \\ \quad ((x \cdot x) = (y \cdot y)) \Rightarrow \\ \quad \quad (x = y) \end{array}}}
\vnbentry{\texttt{sqrt-\allowbreak{}prop}}{\vnbstmt{\begin{array}{@{}l@{}} \forall a \in \mathbb{R}. \\ \quad (0 \leq a) \Rightarrow  \\ \quad \quad ((\sqrt{a} \in \mathbb{R}) \wedge ((0 \leq \sqrt{a}) \wedge ((\sqrt{a} \cdot \sqrt{a}) = a))) \end{array}}}
\vnbentry{\texttt{sqrt-\allowbreak{}char}}{\vnbstmt{\begin{array}{@{}l@{}} \forall a, x. \\ \quad (a \in \mathbb{R}) \wedge \\ \quad (0 \leq a) \wedge \\ \quad (x \in \mathbb{R}) \wedge \\ \quad (0 \leq x) \wedge \\ \quad ((x \cdot x) = a) \Rightarrow \\ \quad \quad (\sqrt{a} = x) \end{array}}}
\vnbentry{\texttt{sqrt-\allowbreak{}nonneg}}{\vnbstmt{\begin{array}{@{}l@{}} \forall a \in \mathbb{R}. \\ \quad ((0 \leq a) \Rightarrow ((\sqrt{a} \in \mathbb{R}) \wedge (0 \leq \sqrt{a}))) \end{array}}}
\vnbentry{\texttt{sqrt-\allowbreak{}sq}}{\vnbstmt{\begin{array}{@{}l@{}} \forall a \in \mathbb{R}. \\ \quad ((0 \leq a) \Rightarrow ((\sqrt{a} \cdot \sqrt{a}) = a)) \end{array}}}
\vnbentry{\texttt{sqrt-\allowbreak{}of-\allowbreak{}sq}}{\vnbstmt{\forall a \in \mathbb{R}.\; (\sqrt{(a \cdot a)} = \lvert a \rvert)}}
\vnbentry{\texttt{sqrt-\allowbreak{}mono}}{\vnbstmt{\begin{array}{@{}l@{}} \forall a, b \in \mathbb{R}. \\ \quad (((0 \leq a) \wedge (a \leq b)) \Rightarrow (\sqrt{a} \leq \sqrt{b})) \end{array}}}
\vnbentry{\texttt{sqrt-\allowbreak{}mul}}{\vnbstmt{\begin{array}{@{}l@{}} \forall a, b \in \mathbb{R}. \\ \quad ((0 \leq a) \wedge (0 \leq b)) \Rightarrow  \\ \quad \quad (\sqrt{(a \cdot b)} = (\sqrt{a} \cdot \sqrt{b})) \end{array}}}

\vnbfile{theorem-library/cc-magnitude.scm}{theorem-library/cc-magnitude.scm}
\noindent{\small cc-magnitude.scm -- the complex modulus, PROVEN from its definition.}\par
\vnbentry{\texttt{cc-\allowbreak{}magnitude-\allowbreak{}closed}}{\vnbstmt{\forall z \in \mathbb{C}.\; (\operatorname{magnitude}(z) \in \mathbb{R})}}
\vnbentry{\texttt{cc-\allowbreak{}magnitude-\allowbreak{}nonneg}}{\vnbstmt{\forall z \in \mathbb{C}.\; (0 \leq \operatorname{magnitude}(z))}}
\vnbentry{\texttt{sqrt-\allowbreak{}zero}}{\vnbstmt{(\sqrt{0} = 0)}}
\vnbentry{\texttt{rr-\allowbreak{}le-\allowbreak{}of-\allowbreak{}sq-\allowbreak{}le}}{\vnbstmt{\begin{array}{@{}l@{}} \forall t, s. \\ \quad (t \in \mathbb{R}) \wedge \\ \quad (s \in \mathbb{R}) \wedge \\ \quad (0 \leq s) \wedge \\ \quad ((t \cdot t) \leq (s \cdot s)) \Rightarrow \\ \quad \quad (t \leq s) \end{array}}}
\vnbentry{\texttt{cc-\allowbreak{}magnitude-\allowbreak{}zero-\allowbreak{}iff}}{\vnbstmt{\begin{array}{@{}l@{}} \forall z \in \mathbb{C}. \\ \quad ((\operatorname{magnitude}(z) = 0) \Leftrightarrow (z = 0)) \end{array}}}
\vnbentry{\texttt{cc-\allowbreak{}magnitude-\allowbreak{}neg}}{\vnbstmt{\begin{array}{@{}l@{}} \forall z \in \mathbb{C}. \\ \quad (\operatorname{magnitude}(-z) = \operatorname{magnitude}(z)) \end{array}}}
\vnbentry{\texttt{rr-\allowbreak{}magnitude-\allowbreak{}is-\allowbreak{}abs}}{\vnbstmt{\forall a \in \mathbb{R}.\; (\operatorname{magnitude}(a) = \lvert a \rvert)}}
\vnbentry{\texttt{cc-\allowbreak{}magnitude-\allowbreak{}mul}}{\vnbstmt{\begin{array}{@{}l@{}} \forall a, b. \\ \quad ((a \in \mathbb{C}) \wedge (b \in \mathbb{C})) \Rightarrow  \\ \quad \quad \operatorname{magnitude}((a \cdot b)) \\ \quad \quad =\; (\operatorname{magnitude}(a) \cdot \operatorname{magnitude}(b)) \end{array}}}
\vnbentry{\texttt{cc-\allowbreak{}magnitude-\allowbreak{}triangle}}{\vnbstmt{\begin{array}{@{}l@{}} \forall a, b. \\ \quad ((a \in \mathbb{C}) \wedge (b \in \mathbb{C})) \Rightarrow  \\ \quad \quad (\operatorname{magnitude}((a + b)) \leq (\operatorname{magnitude}(a) + \operatorname{magnitude}(b))) \end{array}}}

\vnbfile{theorem-library/cc-normed-field.scm}{theorem-library/cc-normed-field.scm}
\noindent{\small cc-normed-field.scm -- IS-NORMED-FIELD(CC-NORMED-FIELD), PROVED.}\par
\vnbentry{\texttt{cc-\allowbreak{}is-\allowbreak{}normed-\allowbreak{}field}}{\vnbstmt{\operatorname{is\text{-}normed\text{-}field}(\operatorname{cc\text{-}normed\text{-}field})}}

\vnbfile{theorem-library/cc-metric-space-proof.scm}{theorem-library/cc-metric-space-proof.scm}
\noindent{\small cc-metric-space-proof.scm -- IS-METRIC-SPACE(CC-MS), PROVEN.}\par
\vnbentry{\texttt{cc-\allowbreak{}sub-\allowbreak{}in-\allowbreak{}cc}}{\vnbstmt{\begin{array}{@{}l@{}} \forall a, b. \\ \quad ((a \in \mathbb{C}) \wedge (b \in \mathbb{C})) \Rightarrow  \\ \quad \quad ((a - b) \in \mathbb{C}) \end{array}}}
\vnbentry{\texttt{cc-\allowbreak{}sub-\allowbreak{}zero-\allowbreak{}eq}}{\vnbstmt{\begin{array}{@{}l@{}} \forall a, b. \\ \quad ((a \in \mathbb{C}) \wedge ((b \in \mathbb{C}) \wedge ((a - b) = 0))) \Rightarrow  \\ \quad \quad (a = b) \end{array}}}
\vnbentry{\texttt{cc-\allowbreak{}is-\allowbreak{}metric-\allowbreak{}space}}{\vnbstmt{\operatorname{is\text{-}metric\text{-}space}(\operatorname{cc\text{-}ms})}}

\vnbfile{theorem-library/complex-inner-product-laws.scm}{theorem-library/complex-inner-product-laws.scm}
\noindent{\small complex-inner-product-laws.scm -- what a sesquilinear form does in its SECOND argument, and the two expansion identities that follow.  All PROVEN modulo 0.}\par
\vnbentry{\texttt{cc-\allowbreak{}conjugate-\allowbreak{}neg}}{\vnbstmt{\begin{array}{@{}l@{}} \forall a \in \mathbb{C}. \\ \quad (\operatorname{conjugate}(-a) = -\operatorname{conjugate}(a)) \end{array}}}
\vnbentry{\texttt{cips-\allowbreak{}ip-\allowbreak{}add-\allowbreak{}right}}{\vnbstmt{\begin{array}{@{}l@{}} \forall v.; \forall x, y, z \in \operatorname{vec}(v). \\ \quad \operatorname{is\text{-}complex\text{-}inner\text{-}product\text{-}space}(v) \Rightarrow  \\ \quad \quad \langle\cdot,\cdot\rangle_{v}(x, \operatorname{vadd}(v)(y, z)) \\ \quad \quad =\; (\langle\cdot,\cdot\rangle_{v}(x, y) + \langle\cdot,\cdot\rangle_{v}(x, z)) \end{array}}}
\vnbentry{\texttt{cips-\allowbreak{}ip-\allowbreak{}conj-\allowbreak{}homog-\allowbreak{}right}}{\vnbstmt{\begin{array}{@{}l@{}} \forall v.; \forall a \in \mathbb{C}, x, y \in \operatorname{vec}(v). \\ \quad \operatorname{is\text{-}complex\text{-}inner\text{-}product\text{-}space}(v) \Rightarrow  \\ \quad \quad \langle\cdot,\cdot\rangle_{v}(x, \operatorname{act}(v)(a, y)) \\ \quad \quad =\; (\operatorname{conjugate}(a) \cdot \langle\cdot,\cdot\rangle_{v}(x, y)) \end{array}}}
\vnbentry{\texttt{cips-\allowbreak{}ip-\allowbreak{}zero-\allowbreak{}right}}{\vnbstmt{\begin{array}{@{}l@{}} \forall v.; \forall x \in \operatorname{vec}(v). \\ \quad \operatorname{is\text{-}complex\text{-}inner\text{-}product\text{-}space}(v) \Rightarrow  \\ \quad \quad (\langle\cdot,\cdot\rangle_{v}(x, \operatorname{vzero}(v)) = 0) \end{array}}}
\vnbentry{\texttt{cips-\allowbreak{}expand-\allowbreak{}sum}}{\vnbstmt{\begin{array}{@{}l@{}} \forall v.; \forall x, y \in \operatorname{vec}(v). \\ \quad \operatorname{is\text{-}complex\text{-}inner\text{-}product\text{-}space}(v) \Rightarrow  \\ \quad \quad \langle\cdot,\cdot\rangle_{v}(\operatorname{vadd}(v)(x, y), \operatorname{vadd}(v)(x, y)) \\ \quad \quad =\; ((\langle\cdot,\cdot\rangle_{v}(x, x) + \langle\cdot,\cdot\rangle_{v}(x, y)) + (\operatorname{conjugate}(\langle\cdot,\cdot\rangle_{v}(x, y)) + \langle\cdot,\cdot\rangle_{v}(y, y))) \end{array}}}
\vnbentry{\texttt{cips-\allowbreak{}expand-\allowbreak{}two}}{\vnbstmt{\begin{array}{@{}l@{}} \forall v, a, b, x, y. \\ \quad \operatorname{is\text{-}complex\text{-}inner\text{-}product\text{-}space}(v) \wedge \\ \quad (a \in \mathbb{C}) \wedge \\ \quad (b \in \mathbb{C}) \wedge \\ \quad (x \in \operatorname{vec}(v)) \wedge \\ \quad (y \in \operatorname{vec}(v)) \Rightarrow \\ \quad \quad \langle\cdot,\cdot\rangle_{v}(\operatorname{vadd}(v)(\operatorname{act}(v)(a, x), \operatorname{act}(v)(b, y)), \\ \quad \quad \quad \operatorname{vadd}(v)(\operatorname{act}(v)(a, x), \operatorname{act}(v)(b, y))) \\ \quad \quad =\; ((((a \cdot \operatorname{conjugate}(a)) \cdot \langle\cdot,\cdot\rangle_{v}(x, x)) + ((a \cdot \operatorname{conjugate}(b)) \cdot \langle\cdot,\cdot\rangle_{v}(x, y))) + (((b \cdot \operatorname{conjugate}(a)) \cdot \operatorname{conjugate}(\langle\cdot,\cdot\rangle_{v}(x, y))) + ((b \cdot \operatorname{conjugate}(b)) \cdot \langle\cdot,\cdot\rangle_{v}(y, y)))) \end{array}}}

\vnbfile{theorem-library/inner-product-inequalities.scm}{theorem-library/inner-product-inequalities.scm}
\noindent{\small inner-product-inequalities.scm -- CAUCHY-SCHWARZ and MINKOWSKI at the INNER-PRODUCT level, over a COMPLEX-INNER-PRODUCT-SPACE.  Both PROVEN modulo 0.}\par
\vnbentry{\texttt{cips-\allowbreak{}schwarz}}{\vnbstmt{\begin{array}{@{}l@{}} \forall v, x, y. \\ \quad \operatorname{is\text{-}complex\text{-}inner\text{-}product\text{-}space}(v) \wedge \\ \quad (x \in \operatorname{vec}(v)) \wedge \\ \quad (y \in \operatorname{vec}(v)) \Rightarrow \\ \quad \quad (\langle\cdot,\cdot\rangle_{v}(x, y) \cdot \operatorname{conjugate}(\langle\cdot,\cdot\rangle_{v}(x, y))) \\ \quad \quad \leq\; (\langle\cdot,\cdot\rangle_{v}(x, x) \cdot \langle\cdot,\cdot\rangle_{v}(y, y)) \end{array}}}
\vnbentry{\texttt{cips-\allowbreak{}minkowski}}{\vnbstmt{\begin{array}{@{}l@{}} \forall v, x, y. \\ \quad \operatorname{is\text{-}complex\text{-}inner\text{-}product\text{-}space}(v) \wedge \\ \quad (x \in \operatorname{vec}(v)) \wedge \\ \quad (y \in \operatorname{vec}(v)) \Rightarrow \\ \quad \quad (\lVert \operatorname{vadd}(v)(x, y) \rVert \leq (\lVert x \rVert + \lVert y \rVert)) \end{array}}}

\vnbfile{theorem-library/rake-ip-normed-ag.scm}{theorem-library/rake-ip-normed-ag.scm}
\noindent{\small rake-ip-normed-ag.scm -- rake batch 5b E: the induced norm of a complex inner product makes the additive group of its vectors a NORMED ABELIAN GROUP.  The support proved here is}\par
\vnbentry{\texttt{ip-\allowbreak{}norm-\allowbreak{}unfold}}{\vnbstmt{\begin{array}{@{}l@{}} \forall v, x. \\ \quad (\lVert x \rVert \simeq \sqrt{\langle\cdot,\cdot\rangle_{v}(x, x)}) \end{array}}}
\vnbentry{\texttt{cips-\allowbreak{}vneg-\allowbreak{}type}}{\vnbstmt{\begin{array}{@{}l@{}} \forall v.; \forall x \in \operatorname{vec}(v). \\ \quad \operatorname{is\text{-}complex\text{-}inner\text{-}product\text{-}space}(v) \Rightarrow  \\ \quad \quad (\operatorname{vneg}(v)(x) \in \operatorname{vec}(v)) \end{array}}}
\vnbentry{\texttt{cips-\allowbreak{}vadd-\allowbreak{}vneg}}{\vnbstmt{\begin{array}{@{}l@{}} \forall v.; \forall x \in \operatorname{vec}(v). \\ \quad \operatorname{is\text{-}complex\text{-}inner\text{-}product\text{-}space}(v) \Rightarrow  \\ \quad \quad (\operatorname{vadd}(v)(x, \operatorname{vneg}(v)(x)) = \operatorname{vzero}(v)) \end{array}}}
\vnbentry{\texttt{cips-\allowbreak{}ip-\allowbreak{}zero-\allowbreak{}left}}{\vnbstmt{\begin{array}{@{}l@{}} \forall v.; \forall z \in \operatorname{vec}(v). \\ \quad \operatorname{is\text{-}complex\text{-}inner\text{-}product\text{-}space}(v) \Rightarrow  \\ \quad \quad (\langle\cdot,\cdot\rangle_{v}(\operatorname{vzero}(v), z) = 0) \end{array}}}
\vnbentry{\texttt{cips-\allowbreak{}ip-\allowbreak{}neg-\allowbreak{}left}}{\vnbstmt{\begin{array}{@{}l@{}} \forall v.; \forall x, z \in \operatorname{vec}(v). \\ \quad \operatorname{is\text{-}complex\text{-}inner\text{-}product\text{-}space}(v) \Rightarrow  \\ \quad \quad (\langle\cdot,\cdot\rangle_{v}(\operatorname{vneg}(v)(x), z) = -\langle\cdot,\cdot\rangle_{v}(x, z)) \end{array}}}
\vnbentry{\texttt{cips-\allowbreak{}ip-\allowbreak{}neg-\allowbreak{}self}}{\vnbstmt{\begin{array}{@{}l@{}} \forall v.; \forall x \in \operatorname{vec}(v). \\ \quad \operatorname{is\text{-}complex\text{-}inner\text{-}product\text{-}space}(v) \Rightarrow  \\ \quad \quad \langle\cdot,\cdot\rangle_{v}(\operatorname{vneg}(v)(x), \operatorname{vneg}(v)(x)) \\ \quad \quad =\; \langle\cdot,\cdot\rangle_{v}(x, x) \end{array}}}
\vnbentry{\texttt{ip-\allowbreak{}normed-\allowbreak{}ag-\allowbreak{}is-\allowbreak{}normed-\allowbreak{}ag}}{\vnbstmt{\begin{array}{@{}l@{}} \forall v. \\ \quad \operatorname{is\text{-}complex\text{-}inner\text{-}product\text{-}space}(v) \Rightarrow  \\ \quad \quad \operatorname{is\text{-}normed\text{-}ag}(\operatorname{ip\text{-}normed\text{-}ag}(v)) \end{array}}}

\vnbfile{theorem-library/monotone-convergence-proof.scm}{theorem-library/monotone-convergence-proof.scm}
\noindent{\small monotone-convergence-proof.scm -- MONOTONE CONVERGENCE on RR, PROVEN.}\par
\vnbentry{\texttt{nn-\allowbreak{}monotone-\allowbreak{}step-\allowbreak{}implies-\allowbreak{}le}}{\vnbstmt{\begin{array}{@{}l@{}} \forall f \in (\mathbb{N} \to \mathbb{R}), n, m \in \mathbb{N}. \\ \quad (\forall k \in \mathbb{N}.\; (f(k) \leq f((k + 1))) \wedge (m \leq n)) \Rightarrow  \\ \quad \quad (f(m) \leq f(n)) \end{array}}}
\vnbentry{\texttt{monotone-\allowbreak{}convergence-\allowbreak{}rr}}{\vnbstmt{\begin{array}{@{}l@{}} \forall f \in (\mathbb{N} \to \mathbb{R}). \\ \quad \forall k \in \mathbb{N}.\; (f(k) \leq f((k + 1))) \wedge \\ \quad \exists d \in \mathbb{R}.\; \forall k \in \mathbb{N}.\; (f(k) \leq d) \Rightarrow \\ \quad \quad \operatorname{converges}(\operatorname{rr\text{-}ms}, f) \end{array}}}

\vnbfile{theorem-library/comparison-test-proof.scm}{theorem-library/comparison-test-proof.scm}
\noindent{\small comparison-test-proof.scm -- THE COMPARISON TEST for real series, PROVEN, together with the two order lemmas it rests on and the readouts that make a real partial sum computable.}\par
\vnbentry{\texttt{rr-\allowbreak{}additive-\allowbreak{}ag-\allowbreak{}opr}}{\vnbstmt{(\operatorname{opr}(\operatorname{normed\text{-}field\text{-}additive\text{-}ag}(\operatorname{rr\text{-}normed\text{-}field})) \simeq (\lambda \langle \operatorname{x\_}, \operatorname{y\_} \rangle \in (\mathbb{R} \times \mathbb{R}).\; (\operatorname{x\_} + \operatorname{y\_})))}}
\vnbentry{\texttt{rr-\allowbreak{}additive-\allowbreak{}ag-\allowbreak{}iden}}{\vnbstmt{(\operatorname{iden}(\operatorname{normed\text{-}field\text{-}additive\text{-}ag}(\operatorname{rr\text{-}normed\text{-}field})) \simeq 0)}}
\vnbentry{\texttt{series-\allowbreak{}partial-\allowbreak{}sum-\allowbreak{}zero}}{\vnbstmt{\forall f.\; (\operatorname{series\text{-}partial\text{-}sum}(f, 0) \simeq 0)}}
\vnbentry{\texttt{sum-\allowbreak{}ag-\allowbreak{}rr-\allowbreak{}in-\allowbreak{}rr}}{\vnbstmt{\begin{array}{@{}l@{}} \forall k \in \mathbb{N}, f \in (\mathbb{N} \to \mathbb{R}). \\ \quad \operatorname{sum\text{-}ag}(\operatorname{normed\text{-}field\text{-}additive\text{-}ag}(\operatorname{rr\text{-}normed\text{-}field}), \\ \quad \quad f, k) \\ \quad \in\; \mathbb{R} \end{array}}}
\vnbentry{\texttt{series-\allowbreak{}partial-\allowbreak{}sum-\allowbreak{}in-\allowbreak{}rr}}{\vnbstmt{\begin{array}{@{}l@{}} \forall k \in \mathbb{N}, f \in (\mathbb{N} \to \mathbb{R}). \\ \quad (\operatorname{series\text{-}partial\text{-}sum}(f, k) \in \mathbb{R}) \end{array}}}
\vnbentry{\texttt{series-\allowbreak{}partial-\allowbreak{}sum-\allowbreak{}unfold}}{\vnbstmt{\begin{array}{@{}l@{}} \forall f, k. \\ \quad (\operatorname{series\text{-}partial\text{-}sum}(f, k) \simeq \operatorname{sum\text{-}ag}(\operatorname{normed\text{-}field\text{-}additive\text{-}ag}(\operatorname{rr\text{-}normed\text{-}field}), f, k)) \end{array}}}
\vnbentry{\texttt{series-\allowbreak{}partial-\allowbreak{}sum-\allowbreak{}succ}}{\vnbstmt{\begin{array}{@{}l@{}} \forall f.; \forall k \in \mathbb{N}. \\ \quad (\operatorname{series\text{-}partial\text{-}sum}(f, k) \in \mathbb{R}) \wedge \\ \quad (f(k) \in \mathbb{R}) \Rightarrow \\ \quad \quad (\operatorname{series\text{-}partial\text{-}sum}(f, (k + 1)) \simeq (\operatorname{series\text{-}partial\text{-}sum}(f, k) + f(k))) \end{array}}}
\vnbentry{\texttt{series-\allowbreak{}partial-\allowbreak{}sum-\allowbreak{}monotone-\allowbreak{}nonneg}}{\vnbstmt{\begin{array}{@{}l@{}} \forall f \in (\mathbb{N} \to \mathbb{R}), k \in \mathbb{N}. \\ \quad \forall n \in \mathbb{N}.\; (0 \leq f(n)) \Rightarrow  \\ \quad \quad \operatorname{series\text{-}partial\text{-}sum}(f, k) \\ \quad \quad \leq\; \operatorname{series\text{-}partial\text{-}sum}(f, (k + 1)) \end{array}}}
\vnbentry{\texttt{series-\allowbreak{}partial-\allowbreak{}sum-\allowbreak{}le-\allowbreak{}termwise}}{\vnbstmt{\begin{array}{@{}l@{}} \forall f, g \in (\mathbb{N} \to \mathbb{R}), k \in \mathbb{N}. \\ \quad \forall n \in \mathbb{N}.\; (f(n) \leq g(n)) \Rightarrow  \\ \quad \quad \operatorname{series\text{-}partial\text{-}sum}(f, k) \\ \quad \quad \leq\; \operatorname{series\text{-}partial\text{-}sum}(g, k) \end{array}}}
\vnbentry{\texttt{monotone-\allowbreak{}convergent-\allowbreak{}bounded-\allowbreak{}above}}{\vnbstmt{\begin{array}{@{}l@{}} \forall f \in (\mathbb{N} \to \mathbb{R}). \\ \quad \forall k \in \mathbb{N}.\; (f(k) \leq f((k + 1))) \wedge \\ \quad \operatorname{converges}(\operatorname{rr\text{-}ms}, f) \Rightarrow \\ \quad \quad \exists d \in \mathbb{R}.\; \forall k \in \mathbb{N}.\; (f(k) \leq d) \end{array}}}
\vnbentry{\texttt{series-\allowbreak{}partial-\allowbreak{}sum-\allowbreak{}seq-\allowbreak{}apply}}{\vnbstmt{\begin{array}{@{}l@{}} \forall f.; \forall j \in \mathbb{N}. \\ \quad ((\lambda k \in \mathbb{N}.\; \operatorname{series\text{-}partial\text{-}sum}(f, k))(j) \simeq \operatorname{series\text{-}partial\text{-}sum}(f, j)) \end{array}}}
\vnbentry{\texttt{series-\allowbreak{}partial-\allowbreak{}sum-\allowbreak{}seq-\allowbreak{}in-\allowbreak{}fun}}{\vnbstmt{\begin{array}{@{}l@{}} \forall f \in (\mathbb{N} \to \mathbb{R}). \\ \quad (\lambda k \in \mathbb{N}.\; \operatorname{series\text{-}partial\text{-}sum}(f, k)) \\ \quad \in\; (\mathbb{N} \to \mathbb{R}) \end{array}}}
\vnbentry{\texttt{abs-\allowbreak{}seq-\allowbreak{}in-\allowbreak{}fun}}{\vnbstmt{\begin{array}{@{}l@{}} \forall f \in (\mathbb{N} \to \mathbb{R}). \\ \quad ((\lambda n \in \mathbb{N}.\; \lvert f(n) \rvert) \in (\mathbb{N} \to \mathbb{R})) \end{array}}}
\vnbentry{\texttt{comparison-\allowbreak{}test}}{\vnbstmt{\begin{array}{@{}l@{}} \forall f, g \in (\mathbb{N} \to \mathbb{R}). \\ \quad \forall n \in \mathbb{N}.\; ((0 \leq f(n)) \wedge (f(n) \leq g(n))) \wedge \\ \quad \operatorname{series\text{-}converges}(g) \Rightarrow \\ \quad \quad \operatorname{series\text{-}converges}(f) \end{array}}}
\vnbentry{\texttt{series-\allowbreak{}partial-\allowbreak{}sum-\allowbreak{}abs-\allowbreak{}le}}{\vnbstmt{\begin{array}{@{}l@{}} \forall k \in \mathbb{N}, f \in (\mathbb{N} \to \mathbb{R}). \\ \quad \lvert \operatorname{series\text{-}partial\text{-}sum}(f, k) \rvert \\ \quad \leq\; \operatorname{series\text{-}partial\text{-}sum}((\lambda n \in \mathbb{N}.\; \lvert f(n) \rvert), k) \end{array}}}

\vnbfile{theorem-library/series-abs-triangle.scm}{theorem-library/series-abs-triangle.scm}
\noindent{\small series-abs-triangle.scm -- |Sum\_\{n$<$k\} f(n)| $<$= Sum\_\{n$<$k\} |f(n)|, PROVEN, and `abs(0) = 0', which the tree did not state.}\par
\vnbentry{\texttt{rr-\allowbreak{}abs-\allowbreak{}zero-\allowbreak{}value}}{\vnbstmt{(\lvert 0 \rvert = 0)}}
\vnbentry{\texttt{series-\allowbreak{}abs-\allowbreak{}triangle}}{\vnbstmt{\begin{array}{@{}l@{}} \forall k \in \mathbb{N}, f, g \in (\mathbb{N} \to \mathbb{R}). \\ \quad \forall j \in \mathbb{N}.\; (g(j) = \lvert f(j) \rvert) \Rightarrow  \\ \quad \quad \lvert \operatorname{series\text{-}partial\text{-}sum}(f, k) \rvert \\ \quad \quad \leq\; \operatorname{series\text{-}partial\text{-}sum}(g, k) \end{array}}}

\vnbfile{theorem-library/series-block-abs.scm}{theorem-library/series-block-abs.scm}
\noindent{\small series-block-abs.scm -- THE BLOCK TRIANGLE INEQUALITY, in the form a Cauchy argument can use, and the NN gap lemma under it.}\par
\vnbentry{\texttt{series-\allowbreak{}abs-\allowbreak{}triangle-\allowbreak{}block}}{\vnbstmt{\begin{array}{@{}l@{}} \forall d, m \in \mathbb{N}, f, g \in (\mathbb{N} \to \mathbb{R}). \\ \quad \forall j \in \mathbb{N}.\; (g(j) = \lvert f(j) \rvert) \Rightarrow  \\ \quad \quad \lvert (\operatorname{series\text{-}partial\text{-}sum}(f, (m + d)) - \operatorname{series\text{-}partial\text{-}sum}(f, m)) \rvert \\ \quad \quad \leq\; (\operatorname{series\text{-}partial\text{-}sum}(g, (m + d)) - \operatorname{series\text{-}partial\text{-}sum}(g, m)) \end{array}}}
\vnbentry{\texttt{nn-\allowbreak{}le-\allowbreak{}gap}}{\vnbstmt{\begin{array}{@{}l@{}} \forall k, m \in \mathbb{N}. \\ \quad ((m \leq k) \Rightarrow \exists d \in \mathbb{N}.\; (k = (m + d))) \end{array}}}
\vnbentry{\texttt{series-\allowbreak{}abs-\allowbreak{}triangle-\allowbreak{}le}}{\vnbstmt{\begin{array}{@{}l@{}} \forall k, m \in \mathbb{N}, f, g \in (\mathbb{N} \to \mathbb{R}). \\ \quad (\forall j \in \mathbb{N}.\; (g(j) = \lvert f(j) \rvert) \wedge (m \leq k)) \Rightarrow  \\ \quad \quad \lvert (\operatorname{series\text{-}partial\text{-}sum}(f, k) - \operatorname{series\text{-}partial\text{-}sum}(f, m)) \rvert \\ \quad \quad \leq\; (\operatorname{series\text{-}partial\text{-}sum}(g, k) - \operatorname{series\text{-}partial\text{-}sum}(g, m)) \end{array}}}

\vnbfile{theorem-library/rake-finsum-typing-nn.scm}{theorem-library/rake-finsum-typing-nn.scm}
\noindent{\small rake-finsum-typing-nn.scm -- two NN typing leaves of the rake, PROVEN.}\par
\vnbentry{\texttt{nn-\allowbreak{}minus-\allowbreak{}in-\allowbreak{}nn}}{\vnbstmt{\begin{array}{@{}l@{}} \forall k, l \in \mathbb{N}. \\ \quad (\operatorname{nn\text{-}minus}(k, l) \in \mathbb{N}) \end{array}}}
\vnbentry{\texttt{falling-\allowbreak{}in-\allowbreak{}nn-\allowbreak{}ind}}{\vnbstmt{\begin{array}{@{}l@{}} \forall m, n \in \mathbb{N}. \\ \quad (\operatorname{falling}(n, m) \in \mathbb{N}) \wedge \\ \quad (((n - m) \in \mathbb{N}) \vee (\operatorname{falling}(n, m) = 0)) \end{array}}}
\vnbentry{\texttt{falling-\allowbreak{}in-\allowbreak{}nn}}{\vnbstmt{\forall n, m \in \mathbb{N}.\; (\operatorname{falling}(n, m) \in \mathbb{N})}}

\vnbfile{theorem-library/series-linearity.scm}{theorem-library/series-linearity.scm}
\noindent{\small series-linearity.scm -- LINEARITY OF THE REAL PARTIAL SUM, in POINTWISE form.}\par
\vnbentry{\texttt{series-\allowbreak{}partial-\allowbreak{}sum-\allowbreak{}in-\allowbreak{}rr-\allowbreak{}ptwise}}{\vnbstmt{\begin{array}{@{}l@{}} \forall n \in \mathbb{N}, f. \\ \quad \forall k \in \mathbb{N}.\; (f(k) \in \mathbb{R}) \Rightarrow  \\ \quad \quad (\operatorname{series\text{-}partial\text{-}sum}(f, n) \in \mathbb{R}) \end{array}}}
\vnbentry{\texttt{series-\allowbreak{}partial-\allowbreak{}sum-\allowbreak{}add-\allowbreak{}ptwise}}{\vnbstmt{\begin{array}{@{}l@{}} \forall n \in \mathbb{N}, f, g, h. \\ \quad \forall k \in \mathbb{N}.\; (f(k) \in \mathbb{R}) \wedge \\ \quad \forall k \in \mathbb{N}.\; (g(k) \in \mathbb{R}) \wedge \\ \quad \forall k \in \mathbb{N}.\; (h(k) = (f(k) + g(k))) \Rightarrow \\ \quad \quad \operatorname{series\text{-}partial\text{-}sum}(h, n) \\ \quad \quad =\; (\operatorname{series\text{-}partial\text{-}sum}(f, n) + \operatorname{series\text{-}partial\text{-}sum}(g, n)) \end{array}}}
\vnbentry{\texttt{series-\allowbreak{}partial-\allowbreak{}sum-\allowbreak{}scale-\allowbreak{}ptwise}}{\vnbstmt{\begin{array}{@{}l@{}} \forall n \in \mathbb{N}, c, f, h. \\ \quad (c \in \mathbb{R}) \wedge \\ \quad \forall k \in \mathbb{N}.\; (f(k) \in \mathbb{R}) \wedge \\ \quad \forall k \in \mathbb{N}.\; (h(k) = (c \cdot f(k))) \Rightarrow \\ \quad \quad \operatorname{series\text{-}partial\text{-}sum}(h, n) \\ \quad \quad =\; (c \cdot \operatorname{series\text{-}partial\text{-}sum}(f, n)) \end{array}}}
\vnbentry{\texttt{series-\allowbreak{}partial-\allowbreak{}sum-\allowbreak{}weighted-\allowbreak{}expansion}}{\vnbstmt{\begin{array}{@{}l@{}} \forall n \in \mathbb{N}, x, y, g, w, p, q, s. \\ \quad (x \in \mathbb{R}) \wedge \\ \quad (y \in \mathbb{R}) \wedge \\ \quad \forall k \in \mathbb{Z}.\; (g(k) \in \mathbb{R}) \wedge \\ \quad \forall k \in \mathbb{Z}.\; (w(k) \in \mathbb{R}) \wedge \\ \quad \forall k \in \mathbb{N}.\; (q(k) \in \mathbb{R}) \wedge \\ \quad \forall k \in \mathbb{N}.\; (s(k) \in \mathbb{R}) \wedge \\ \quad \forall k \in \mathbb{N}. \\ \quad \quad (p(k) = (w(k) \cdot ((x \cdot g((k - 1))) + (y \cdot g(k))))) \wedge \\ \quad \forall k \in \mathbb{N}.\; (q(k) = (w((k + 1)) \cdot g(k))) \wedge \\ \quad \forall k \in \mathbb{N}.\; (s(k) = (w(k) \cdot g(k))) \Rightarrow \\ \quad \quad \operatorname{series\text{-}partial\text{-}sum}(p, (n + 1)) \\ \quad \quad =\; ((x \cdot ((w(0) \cdot g((0 - 1))) + \operatorname{series\text{-}partial\text{-}sum}(q, n))) + (y \cdot (\operatorname{series\text{-}partial\text{-}sum}(s, n) + (w(n) \cdot g(n))))) \end{array}}}
\vnbentry{\texttt{series-\allowbreak{}partial-\allowbreak{}sum-\allowbreak{}le-\allowbreak{}termwise-\allowbreak{}ptwise}}{\vnbstmt{\begin{array}{@{}l@{}} \forall n \in \mathbb{N}, f, g. \\ \quad \forall k \in \mathbb{N}.\; (f(k) \in \mathbb{R}) \wedge \\ \quad \forall k \in \mathbb{N}.\; (g(k) \in \mathbb{R}) \wedge \\ \quad \forall k \in \mathbb{N}.\; (f(k) \leq g(k)) \Rightarrow \\ \quad \quad \operatorname{series\text{-}partial\text{-}sum}(f, n) \\ \quad \quad \leq\; \operatorname{series\text{-}partial\text{-}sum}(g, n) \end{array}}}

\vnbfile{theorem-library/bernstein-basis.scm}{theorem-library/bernstein-basis.scm}
\noindent{\small bernstein-basis.scm -- the Bernstein basis polynomials, and the first of the two identities Proposition 5.4 of docs/calculus.pdf asks for.}\par
\vnbentry{\texttt{ring-\allowbreak{}power-\allowbreak{}of-\allowbreak{}one}}{\vnbstmt{\begin{array}{@{}l@{}} \forall n \in \mathbb{N}, r. \\ \quad \operatorname{is\text{-}commutative\text{-}ring}(r) \Rightarrow  \\ \quad \quad (\operatorname{ring\text{-}power}(r, \operatorname{one}(r), n) = \operatorname{one}(r)) \end{array}}}
\vnbentry{\texttt{series-\allowbreak{}partial-\allowbreak{}sum-\allowbreak{}is-\allowbreak{}ring-\allowbreak{}sum}}{\vnbstmt{\begin{array}{@{}l@{}} \forall n \in \mathbb{N}, f. \\ \quad (\operatorname{series\text{-}partial\text{-}sum}(f, n) \simeq \operatorname{sum}(\operatorname{normed\text{-}field\text{-}as\text{-}commutative\text{-}ring}(\operatorname{rr\text{-}normed\text{-}field}), f, n)) \end{array}}}
\vnbentry{\texttt{binomial-\allowbreak{}sum-\allowbreak{}value}}{\vnbstmt{\begin{array}{@{}l@{}} \forall n, r, x, y. \\ \quad (n \in \mathbb{N}) \wedge \\ \quad \operatorname{is\text{-}commutative\text{-}ring}(r) \wedge \\ \quad (x \in \operatorname{carr}(r)) \wedge \\ \quad (y \in \operatorname{carr}(r)) \Rightarrow \\ \quad \quad \operatorname{sum}(r, \operatorname{comb\text{-}kk}(r, x, y, n), (n + 1)) \\ \quad \quad =\; \operatorname{ring\text{-}power}(r, \operatorname{add}(r)(x, y), n) \end{array}}}
\vnbentry{\texttt{bernstein-\allowbreak{}basis-\allowbreak{}unfold}}{\vnbstmt{\begin{array}{@{}l@{}} \forall n, x. \\ \quad (\operatorname{bernstein\text{-}basis}(n, x) \simeq \operatorname{comb\text{-}kk}(\operatorname{normed\text{-}field\text{-}as\text{-}commutative\text{-}ring}(\operatorname{rr\text{-}normed\text{-}field}), x, (1 - x), n)) \end{array}}}
\vnbentry{\texttt{bernstein-\allowbreak{}basis-\allowbreak{}in-\allowbreak{}fun}}{\vnbstmt{\begin{array}{@{}l@{}} \forall n \in \mathbb{N}, x \in \mathbb{R}. \\ \quad (\operatorname{bernstein\text{-}basis}(n, x) \in (\mathbb{Z} \to \mathbb{R})) \end{array}}}
\vnbentry{\texttt{bernstein-\allowbreak{}partition}}{\vnbstmt{\begin{array}{@{}l@{}} \forall n \in \mathbb{N}, x \in \mathbb{R}. \\ \quad \operatorname{series\text{-}partial\text{-}sum}(\operatorname{bernstein\text{-}basis}(n, x), (n + 1)) \\ \quad =\; 1 \end{array}}}

\vnbfile{theorem-library/bernstein-moments.scm}{theorem-library/bernstein-moments.scm}
\noindent{\small bernstein-moments.scm -- the moments of the Bernstein basis, and identity (76) of docs/calculus.pdf Proposition 5.4.}\par
\vnbentry{\texttt{bernstein-\allowbreak{}basis-\allowbreak{}ptwise-\allowbreak{}in-\allowbreak{}rr}}{\vnbstmt{\begin{array}{@{}l@{}} \forall n, x, k. \\ \quad ((n \in \mathbb{N}) \wedge ((x \in \mathbb{R}) \wedge (k \in \mathbb{Z}))) \Rightarrow  \\ \quad \quad (\operatorname{bernstein\text{-}basis}(n, x)(k) \in \mathbb{R}) \end{array}}}
\vnbentry{\texttt{bernstein-\allowbreak{}basis-\allowbreak{}succ}}{\vnbstmt{\begin{array}{@{}l@{}} \forall n, x, k. \\ \quad ((n \in \mathbb{N}) \wedge ((x \in \mathbb{R}) \wedge (k \in \mathbb{N}))) \Rightarrow  \\ \quad \quad \operatorname{bernstein\text{-}basis}((n + 1), x)(k) \\ \quad \quad =\; ((x \cdot \operatorname{bernstein\text{-}basis}(n, x)((k - 1))) + ((1 - x) \cdot \operatorname{bernstein\text{-}basis}(n, x)(k))) \end{array}}}
\vnbentry{\texttt{bernstein-\allowbreak{}basis-\allowbreak{}null}}{\vnbstmt{\begin{array}{@{}l@{}} \forall n, x, k. \\ \quad ((n \in \mathbb{N}) \wedge ((x \in \mathbb{R}) \wedge ((k \in \mathbb{Z}) \wedge (k < 0)))) \Rightarrow  \\ \quad \quad (\operatorname{bernstein\text{-}basis}(n, x)(k) = 0) \end{array}}}
\vnbentry{\texttt{bernstein-\allowbreak{}basis-\allowbreak{}above}}{\vnbstmt{\begin{array}{@{}l@{}} \forall n, x, k. \\ \quad ((n \in \mathbb{N}) \wedge ((x \in \mathbb{R}) \wedge ((k \in \mathbb{Z}) \wedge (n < k)))) \Rightarrow  \\ \quad \quad (\operatorname{bernstein\text{-}basis}(n, x)(k) = 0) \end{array}}}
\vnbentry{\texttt{bernstein-\allowbreak{}weighted-\allowbreak{}step}}{\vnbstmt{\begin{array}{@{}l@{}} \forall n, x, w, p, q, s. \\ \quad (n \in \mathbb{N}) \wedge \\ \quad (x \in \mathbb{R}) \wedge \\ \quad \forall k \in \mathbb{Z}.\; (w(k) \in \mathbb{R}) \wedge \\ \quad \forall k \in \mathbb{N}.\; (q(k) \in \mathbb{R}) \wedge \\ \quad \forall k \in \mathbb{N}.\; (s(k) \in \mathbb{R}) \wedge \\ \quad \forall k \in \mathbb{N}. \\ \quad \quad (p(k) = (w(k) \cdot \operatorname{bernstein\text{-}basis}((n + 1), x)(k))) \wedge \\ \quad \forall k \in \mathbb{N}. \\ \quad \quad (q(k) = (w((k + 1)) \cdot \operatorname{bernstein\text{-}basis}(n, x)(k))) \wedge \\ \quad \forall k \in \mathbb{N}. \\ \quad \quad (s(k) = (w(k) \cdot \operatorname{bernstein\text{-}basis}(n, x)(k))) \Rightarrow \\ \quad \quad \operatorname{series\text{-}partial\text{-}sum}(p, ((n + 1) + 1)) \\ \quad \quad =\; ((x \cdot \operatorname{series\text{-}partial\text{-}sum}(q, (n + 1))) + ((1 - x) \cdot \operatorname{series\text{-}partial\text{-}sum}(s, (n + 1)))) \end{array}}}
\vnbentry{\texttt{bernstein-\allowbreak{}moment-\allowbreak{}1}}{\vnbstmt{\begin{array}{@{}l@{}} \forall n \in \mathbb{N}, x, h. \\ \quad (x \in \mathbb{R}) \wedge \\ \quad \forall k \in \mathbb{N}. \\ \quad \quad (h(k) = (k \cdot \operatorname{bernstein\text{-}basis}(n, x)(k))) \Rightarrow \\ \quad \quad (\operatorname{series\text{-}partial\text{-}sum}(h, (n + 1)) = (n \cdot x)) \end{array}}}
\vnbentry{\texttt{bernstein-\allowbreak{}moment-\allowbreak{}2}}{\vnbstmt{\begin{array}{@{}l@{}} \forall n \in \mathbb{N}, x, h. \\ \quad (x \in \mathbb{R}) \wedge \\ \quad \forall k \in \mathbb{N}. \\ \quad \quad (h(k) = ((k \cdot (k - 1)) \cdot \operatorname{bernstein\text{-}basis}(n, x)(k))) \Rightarrow \\ \quad \quad \operatorname{series\text{-}partial\text{-}sum}(h, (n + 1)) \\ \quad \quad =\; ((n \cdot (n - 1)) \cdot (x \cdot x)) \end{array}}}
\vnbentry{\texttt{bernstein-\allowbreak{}variance}}{\vnbstmt{\begin{array}{@{}l@{}} \forall n, x, h. \\ \quad (n \in \mathbb{N}) \wedge \\ \quad (x \in \mathbb{R}) \wedge \\ \quad \forall k \in \mathbb{N}. \\ \quad \quad h(k) \\ \quad \quad =\; (((k - (n \cdot x)) \cdot (k - (n \cdot x))) \cdot \operatorname{bernstein\text{-}basis}(n, x)(k)) \Rightarrow \\ \quad \quad (\operatorname{series\text{-}partial\text{-}sum}(h, (n + 1)) = ((n \cdot x) \cdot (1 - x))) \end{array}}}
\vnbentry{\texttt{bernstein-\allowbreak{}variance-\allowbreak{}bound}}{\vnbstmt{\begin{array}{@{}l@{}} \forall n, x, h. \\ \quad (n \in \mathbb{N}) \wedge \\ \quad \neg (n = 0) \wedge \\ \quad (x \in \mathbb{R}) \wedge \\ \quad \forall k \in \mathbb{N}. \\ \quad \quad h(k) \\ \quad \quad =\; (((k - (n \cdot x)) \cdot (k - (n \cdot x))) \cdot \operatorname{bernstein\text{-}basis}(n, x)(k)) \Rightarrow \\ \quad \quad ((/(1, n) \cdot \operatorname{series\text{-}partial\text{-}sum}(h, (n + 1))) \leq 1) \end{array}}}

\vnbfile{theorem-library/series-cauchy-proof.scm}{theorem-library/series-cauchy-proof.scm}
\noindent{\small series-cauchy-proof.scm -- THE CAUCHY CRITERION FOR REAL SERIES, PROVEN.}\par
\vnbentry{\texttt{series-\allowbreak{}cauchy-\allowbreak{}criterion}}{\vnbstmt{\begin{array}{@{}l@{}} \forall f \in (\mathbb{N} \to \mathbb{R}), s. \\ \quad (\operatorname{series\text{-}converges}(f) \wedge \operatorname{pos\text{-}rr}(s)) \Rightarrow  \\ \quad \quad \exists d \in \mathbb{N}. \forall m, n \in \mathbb{N}. \\ \quad \quad \quad ((d \leq m) \wedge (m \leq n)) \Rightarrow  \\ \quad \quad \quad \quad \lvert (\operatorname{series\text{-}partial\text{-}sum}(f, n) - \operatorname{series\text{-}partial\text{-}sum}(f, m)) \rvert \\ \quad \quad \quad \quad \leq\; s \end{array}}}

\vnbfile{theorem-library/dominated-convergence.scm}{theorem-library/dominated-convergence.scm}
\noindent{\small dominated-convergence.scm -- DOMINATED CONVERGENCE, PROVEN, in two forms: for real series (the content) and in ELL-ONE (the statement the user asked for), together with the twenty limit / partial-sum / complex-magnitude lemmas the two are assembled from.}\par
\vnbentry{\texttt{rr-\allowbreak{}null-\allowbreak{}sum}}{\vnbstmt{\begin{array}{@{}l@{}} \forall f, g, h \in (\mathbb{N} \to \mathbb{R}). \\ \quad \operatorname{converges\text{-}to}(\operatorname{rr\text{-}ms}, f, 0) \wedge \\ \quad \operatorname{converges\text{-}to}(\operatorname{rr\text{-}ms}, g, 0) \wedge \\ \quad \forall j \in \mathbb{N}.\; (h(j) = (f(j) + g(j))) \Rightarrow \\ \quad \quad \operatorname{converges\text{-}to}(\operatorname{rr\text{-}ms}, h, 0) \end{array}}}
\vnbentry{\texttt{rr-\allowbreak{}limit-\allowbreak{}abs-\allowbreak{}le}}{\vnbstmt{\begin{array}{@{}l@{}} \forall f \in (\mathbb{N} \to \mathbb{R}), v, d \in \mathbb{R}. \\ \quad \operatorname{converges\text{-}to}(\operatorname{rr\text{-}ms}, f, v) \wedge \\ \quad \forall n \in \mathbb{N}.\; (\lvert f(n) \rvert \leq d) \Rightarrow \\ \quad \quad (\lvert v \rvert \leq d) \end{array}}}
\vnbentry{\texttt{series-\allowbreak{}block-\allowbreak{}le}}{\vnbstmt{\begin{array}{@{}l@{}} \forall k, n \in \mathbb{N}, f, g \in (\mathbb{N} \to \mathbb{R}). \\ \quad (\forall i \in \mathbb{N}.\; ((0 \leq f(i)) \wedge (f(i) \leq g(i))) \wedge (n \leq k)) \Rightarrow  \\ \quad \quad (\operatorname{series\text{-}partial\text{-}sum}(f, k) - \operatorname{series\text{-}partial\text{-}sum}(f, n)) \\ \quad \quad \leq\; (\operatorname{series\text{-}partial\text{-}sum}(g, k) - \operatorname{series\text{-}partial\text{-}sum}(g, n)) \end{array}}}
\vnbentry{\texttt{series-\allowbreak{}head-\allowbreak{}null}}{\vnbstmt{\begin{array}{@{}l@{}} \forall n \in \mathbb{N}, u.; \forall d \in (\mathbb{N} \to \mathbb{R}). \\ \quad \forall j \in \mathbb{N}.\; (u(j) \in (\mathbb{N} \to \mathbb{R})) \wedge \\ \quad \forall i \in \mathbb{N}. \\ \quad \quad \operatorname{converges\text{-}to}(\operatorname{rr\text{-}ms}, (\lambda j \in \mathbb{N}.\; u(j)(i)), 0) \wedge \\ \quad \forall j \in \mathbb{N}. \\ \quad \quad (d(j) = \operatorname{series\text{-}partial\text{-}sum}(u(j), n)) \Rightarrow \\ \quad \quad \operatorname{converges\text{-}to}(\operatorname{rr\text{-}ms}, d, 0) \end{array}}}
\vnbentry{\texttt{series-\allowbreak{}partial-\allowbreak{}sum-\allowbreak{}nonneg}}{\vnbstmt{\begin{array}{@{}l@{}} \forall k \in \mathbb{N}, f \in (\mathbb{N} \to \mathbb{R}). \\ \quad \forall i \in \mathbb{N}.\; (0 \leq f(i)) \Rightarrow  \\ \quad \quad (0 \leq \operatorname{series\text{-}partial\text{-}sum}(f, k)) \end{array}}}
\vnbentry{\texttt{series-\allowbreak{}partial-\allowbreak{}sum-\allowbreak{}mono}}{\vnbstmt{\begin{array}{@{}l@{}} \forall k, m \in \mathbb{N}, f \in (\mathbb{N} \to \mathbb{R}). \\ \quad (\forall i \in \mathbb{N}.\; (0 \leq f(i)) \wedge (m \leq k)) \Rightarrow  \\ \quad \quad \operatorname{series\text{-}partial\text{-}sum}(f, m) \\ \quad \quad \leq\; \operatorname{series\text{-}partial\text{-}sum}(f, k) \end{array}}}
\vnbentry{\texttt{series-\allowbreak{}tail-\allowbreak{}small}}{\vnbstmt{\begin{array}{@{}l@{}} \forall g \in (\mathbb{N} \to \mathbb{R}), s. \\ \quad (\operatorname{series\text{-}converges}(g) \wedge \operatorname{pos\text{-}rr}(s)) \Rightarrow  \\ \quad \quad \exists d \in \mathbb{N}. \forall k \in \mathbb{N}. \\ \quad \quad \quad (d \leq k) \Rightarrow  \\ \quad \quad \quad \quad (\operatorname{series\text{-}partial\text{-}sum}(g, k) - \operatorname{series\text{-}partial\text{-}sum}(g, d)) \\ \quad \quad \quad \quad \leq\; s \end{array}}}
\vnbentry{\texttt{dominated-\allowbreak{}null-\allowbreak{}series}}{\vnbstmt{\begin{array}{@{}l@{}} \forall g, e \in (\mathbb{N} \to \mathbb{R}), u. \\ \quad \forall j \in \mathbb{N}.\; (u(j) \in (\mathbb{N} \to \mathbb{R})) \wedge \\ \quad \forall j, i \in \mathbb{N}.\; ((0 \leq u(j)(i)) \wedge (u(j)(i) \leq g(i))) \wedge \\ \quad \operatorname{series\text{-}converges}(g) \wedge \\ \quad \forall i \in \mathbb{N}. \\ \quad \quad \operatorname{converges\text{-}to}(\operatorname{rr\text{-}ms}, (\lambda j \in \mathbb{N}.\; u(j)(i)), 0) \wedge \\ \quad \forall j \in \mathbb{N}. \\ \quad \quad \operatorname{series\text{-}converges\text{-}to}(u(j), e(j)) \Rightarrow \\ \quad \quad \operatorname{converges\text{-}to}(\operatorname{rr\text{-}ms}, e, 0) \end{array}}}
\vnbentry{\texttt{series-\allowbreak{}limit-\allowbreak{}converges-\allowbreak{}to}}{\vnbstmt{\begin{array}{@{}l@{}} \forall f \in (\mathbb{N} \to \mathbb{R}). \\ \quad \operatorname{series\text{-}converges}(f) \Rightarrow  \\ \quad \quad \operatorname{series\text{-}converges\text{-}to}(f, \operatorname{series\text{-}limit}(f)) \end{array}}}
\vnbentry{\texttt{series-\allowbreak{}limit-\allowbreak{}in-\allowbreak{}rr}}{\vnbstmt{\begin{array}{@{}l@{}} \forall f \in (\mathbb{N} \to \mathbb{R}). \\ \quad \operatorname{series\text{-}converges}(f) \Rightarrow  \\ \quad \quad (\operatorname{series\text{-}limit}(f) \in \mathbb{R}) \end{array}}}
\vnbentry{\texttt{series-\allowbreak{}partial-\allowbreak{}sum-\allowbreak{}add}}{\vnbstmt{\begin{array}{@{}l@{}} \forall k \in \mathbb{N}, f, g, h \in (\mathbb{N} \to \mathbb{R}). \\ \quad \forall i \in \mathbb{N}.\; (h(i) = (f(i) + g(i))) \Rightarrow  \\ \quad \quad \operatorname{series\text{-}partial\text{-}sum}(h, k) \\ \quad \quad =\; (\operatorname{series\text{-}partial\text{-}sum}(f, k) + \operatorname{series\text{-}partial\text{-}sum}(g, k)) \end{array}}}
\vnbentry{\texttt{series-\allowbreak{}partial-\allowbreak{}sum-\allowbreak{}bounded}}{\vnbstmt{\begin{array}{@{}l@{}} \forall f \in (\mathbb{N} \to \mathbb{R}). \\ \quad (\forall i \in \mathbb{N}.\; (0 \leq f(i)) \wedge \operatorname{series\text{-}converges}(f)) \Rightarrow  \\ \quad \quad \exists d \in \mathbb{R}. \forall x \in \mathbb{N}. \\ \quad \quad \quad (\operatorname{series\text{-}partial\text{-}sum}(f, x) \leq d) \end{array}}}
\vnbentry{\texttt{series-\allowbreak{}converges-\allowbreak{}sum}}{\vnbstmt{\begin{array}{@{}l@{}} \forall f, g, h \in (\mathbb{N} \to \mathbb{R}). \\ \quad \forall i \in \mathbb{N}.\; (0 \leq f(i)) \wedge \\ \quad \forall i \in \mathbb{N}.\; (0 \leq g(i)) \wedge \\ \quad \forall i \in \mathbb{N}.\; (h(i) = (f(i) + g(i))) \wedge \\ \quad \operatorname{series\text{-}converges}(f) \wedge \\ \quad \operatorname{series\text{-}converges}(g) \Rightarrow \\ \quad \quad \operatorname{series\text{-}converges}(h) \end{array}}}
\vnbentry{\texttt{cc-\allowbreak{}ms-\allowbreak{}dist}}{\vnbstmt{\begin{array}{@{}l@{}} \forall u, v. \\ \quad ((u \in \mathbb{C}) \wedge (v \in \mathbb{C})) \Rightarrow  \\ \quad \quad (\operatorname{dist}(\operatorname{cc\text{-}ms})(u, v) \simeq \operatorname{magnitude}((u - v))) \end{array}}}
\vnbentry{\texttt{cc-\allowbreak{}magnitude-\allowbreak{}sub-\allowbreak{}triangle}}{\vnbstmt{\begin{array}{@{}l@{}} \forall p, q \in \mathbb{C}. \\ \quad (\operatorname{magnitude}((p - q)) \leq (\operatorname{magnitude}(p) + \operatorname{magnitude}(q))) \end{array}}}
\vnbentry{\texttt{cc-\allowbreak{}magnitude-\allowbreak{}sub-\allowbreak{}sym}}{\vnbstmt{\begin{array}{@{}l@{}} \forall p, q \in \mathbb{C}. \\ \quad (\operatorname{magnitude}((p - q)) = \operatorname{magnitude}((q - p))) \end{array}}}
\vnbentry{\texttt{cc-\allowbreak{}magnitude-\allowbreak{}le-\allowbreak{}add}}{\vnbstmt{\begin{array}{@{}l@{}} \forall p, q \in \mathbb{C}. \\ \quad (\operatorname{magnitude}(p) \leq (\operatorname{magnitude}((p - q)) + \operatorname{magnitude}(q))) \end{array}}}
\vnbentry{\texttt{cc-\allowbreak{}limit-\allowbreak{}magnitude-\allowbreak{}le}}{\vnbstmt{\begin{array}{@{}l@{}} \forall s \in (\mathbb{N} \to \mathbb{C}), v \in \mathbb{C}, d \in \mathbb{R}. \\ \quad \operatorname{converges\text{-}to}(\operatorname{cc\text{-}ms}, s, v) \wedge \\ \quad \forall n \in \mathbb{N}.\; (\operatorname{magnitude}(s(n)) \leq d) \Rightarrow \\ \quad \quad (\operatorname{magnitude}(v) \leq d) \end{array}}}
\vnbentry{\texttt{cc-\allowbreak{}converges-\allowbreak{}magnitude-\allowbreak{}null}}{\vnbstmt{\begin{array}{@{}l@{}} \forall s \in (\mathbb{N} \to \mathbb{C}), v \in \mathbb{C}, m \in (\mathbb{N} \to \mathbb{R}). \\ \quad \operatorname{converges\text{-}to}(\operatorname{cc\text{-}ms}, s, v) \wedge \\ \quad \forall j \in \mathbb{N}.\; (m(j) = \operatorname{magnitude}((s(j) - v))) \Rightarrow \\ \quad \quad \operatorname{converges\text{-}to}(\operatorname{rr\text{-}ms}, m, 0) \end{array}}}
\vnbentry{\texttt{ell-\allowbreak{}one-\allowbreak{}membership}}{\vnbstmt{\begin{array}{@{}l@{}} \forall y. \\ \quad ((y \in \operatorname{ell\text{-}one}) \Leftrightarrow ((y \in \mathbb{C}^{\mathbb{N}}) \wedge \operatorname{series\text{-}converges}((\lambda k \in \mathbb{N}.\; \operatorname{magnitude}(y(k)))))) \end{array}}}
\vnbentry{\texttt{ell-\allowbreak{}one-\allowbreak{}dominated-\allowbreak{}convergence}}{\vnbstmt{\begin{array}{@{}l@{}} \forall a \in (\mathbb{N} \to \mathbb{R}), x \in \mathbb{C}^{\mathbb{N}}, e \in (\mathbb{N} \to \mathbb{R}), s. \\ \quad \forall j \in \mathbb{N}.\; (s(j) \in \mathbb{C}^{\mathbb{N}}) \wedge \\ \quad \operatorname{series\text{-}converges}(a) \wedge \\ \quad \forall j, k \in \mathbb{N}.\; (\operatorname{magnitude}(s(j)(k)) \leq a(k)) \wedge \\ \quad \forall k \in \mathbb{N}. \\ \quad \quad \operatorname{converges\text{-}to}(\operatorname{cc\text{-}ms}, (\lambda j \in \mathbb{N}.\; s(j)(k)), x(k)) \wedge \\ \quad \forall j \in \mathbb{N}. \\ \quad \quad e(j) \\ \quad \quad =\; \operatorname{series\text{-}limit}((\lambda k \in \mathbb{N}.\; \operatorname{magnitude}((s(j)(k) - x(k))))) \Rightarrow \\ \quad \quad (x \in \operatorname{ell\text{-}one}) \wedge \\ \quad \quad \operatorname{converges\text{-}to}(\operatorname{rr\text{-}ms}, e, 0) \end{array}}}

\vnbfile{theorem-library/cc-complete-proof.scm}{theorem-library/cc-complete-proof.scm}
\noindent{\small cc-complete-proof.scm -- CC IS A COMPLETE METRIC SPACE.}\par
\vnbentry{\texttt{cc-\allowbreak{}re-\allowbreak{}sub}}{\vnbstmt{\begin{array}{@{}l@{}} \forall z, w. \\ \quad ((z \in \mathbb{C}) \wedge (w \in \mathbb{C})) \Rightarrow  \\ \quad \quad \operatorname{real\text{-}part}((z - w)) \\ \quad \quad =\; (\operatorname{real\text{-}part}(z) - \operatorname{real\text{-}part}(w)) \end{array}}}
\vnbentry{\texttt{cc-\allowbreak{}im-\allowbreak{}sub}}{\vnbstmt{\begin{array}{@{}l@{}} \forall z, w. \\ \quad ((z \in \mathbb{C}) \wedge (w \in \mathbb{C})) \Rightarrow  \\ \quad \quad \operatorname{imag\text{-}part}((z - w)) \\ \quad \quad =\; (\operatorname{imag\text{-}part}(z) - \operatorname{imag\text{-}part}(w)) \end{array}}}
\vnbentry{\texttt{cc-\allowbreak{}abs-\allowbreak{}re-\allowbreak{}le-\allowbreak{}magnitude}}{\vnbstmt{\begin{array}{@{}l@{}} \forall z \in \mathbb{C}. \\ \quad (\lvert \operatorname{real\text{-}part}(z) \rvert \leq \operatorname{magnitude}(z)) \end{array}}}
\vnbentry{\texttt{cc-\allowbreak{}abs-\allowbreak{}im-\allowbreak{}le-\allowbreak{}magnitude}}{\vnbstmt{\begin{array}{@{}l@{}} \forall z \in \mathbb{C}. \\ \quad (\lvert \operatorname{imag\text{-}part}(z) \rvert \leq \operatorname{magnitude}(z)) \end{array}}}
\vnbentry{\texttt{cc-\allowbreak{}magnitude-\allowbreak{}le-\allowbreak{}re-\allowbreak{}im}}{\vnbstmt{\begin{array}{@{}l@{}} \forall z \in \mathbb{C}. \\ \quad \operatorname{magnitude}(z) \\ \quad \leq\; (\lvert \operatorname{real\text{-}part}(z) \rvert + \lvert \operatorname{imag\text{-}part}(z) \rvert) \end{array}}}
\vnbentry{\texttt{cc-\allowbreak{}cauchy-\allowbreak{}of-\allowbreak{}dominated}}{\vnbstmt{\begin{array}{@{}l@{}} \forall f \in (\mathbb{N} \to \mathbb{C}), p \in (\mathbb{N} \to \mathbb{R}). \\ \quad \operatorname{is\text{-}cauchy\text{-}seq}(\operatorname{cc\text{-}ms}, f) \wedge \\ \quad \forall m, n \in \mathbb{N}. \\ \quad \quad (\lvert (p(m) - p(n)) \rvert \leq \operatorname{magnitude}((f(m) - f(n)))) \Rightarrow \\ \quad \quad \operatorname{is\text{-}cauchy\text{-}seq}(\operatorname{rr\text{-}ms}, p) \end{array}}}
\vnbentry{\texttt{cc-\allowbreak{}converges-\allowbreak{}of-\allowbreak{}coords}}{\vnbstmt{\begin{array}{@{}l@{}} \forall f \in (\mathbb{N} \to \mathbb{C}), g, q \in (\mathbb{N} \to \mathbb{R}), v \in \mathbb{C}. \\ \quad \forall n \in \mathbb{N}.\; (\operatorname{real\text{-}part}(f(n)) = g(n)) \wedge \\ \quad \forall n \in \mathbb{N}.\; (\operatorname{imag\text{-}part}(f(n)) = q(n)) \wedge \\ \quad \operatorname{converges\text{-}to}(\operatorname{rr\text{-}ms}, g, \operatorname{real\text{-}part}(v)) \wedge \\ \quad \operatorname{converges\text{-}to}(\operatorname{rr\text{-}ms}, q, \operatorname{imag\text{-}part}(v)) \Rightarrow \\ \quad \quad \operatorname{converges\text{-}to}(\operatorname{cc\text{-}ms}, f, v) \end{array}}}
\vnbentry{\texttt{cc-\allowbreak{}complete}}{\vnbstmt{\operatorname{is\text{-}complete}(\operatorname{cc\text{-}ms})}}

\vnbfile{theorem-library/series-abs-converges.scm}{theorem-library/series-abs-converges.scm}
\noindent{\small series-abs-converges.scm -- ABSOLUTE CONVERGENCE IMPLIES CONVERGENCE, for real series.}\par
\vnbentry{\texttt{series-\allowbreak{}abs-\allowbreak{}converges}}{\vnbstmt{\begin{array}{@{}l@{}} \forall f, g \in (\mathbb{N} \to \mathbb{R}). \\ \quad \forall j \in \mathbb{N}.\; (g(j) = \lvert f(j) \rvert) \wedge \\ \quad \operatorname{series\text{-}converges}(g) \Rightarrow \\ \quad \quad \operatorname{series\text{-}converges}(f) \end{array}}}

\vnbfile{theorem-library/cc-series.scm}{theorem-library/cc-series.scm}
\noindent{\small cc-series.scm -- THE COMPLEX SERIES LAYER: ordered partial sums in CC, and their convergence in CC-MS.}\par
\vnbentry{\texttt{cc-\allowbreak{}additive-\allowbreak{}ag-\allowbreak{}opr}}{\vnbstmt{(\operatorname{opr}(\operatorname{normed\text{-}field\text{-}additive\text{-}ag}(\operatorname{cc\text{-}normed\text{-}field})) \simeq (\lambda \langle \operatorname{x\_}, \operatorname{y\_} \rangle \in (\mathbb{C} \times \mathbb{C}).\; (\operatorname{x\_} + \operatorname{y\_})))}}
\vnbentry{\texttt{cc-\allowbreak{}additive-\allowbreak{}ag-\allowbreak{}iden}}{\vnbstmt{(\operatorname{iden}(\operatorname{normed\text{-}field\text{-}additive\text{-}ag}(\operatorname{cc\text{-}normed\text{-}field})) \simeq 0)}}
\vnbentry{\texttt{cc-\allowbreak{}series-\allowbreak{}partial-\allowbreak{}sum-\allowbreak{}zero}}{\vnbstmt{\begin{array}{@{}l@{}} \forall f. \\ \quad (\operatorname{cc\text{-}series\text{-}partial\text{-}sum}(f, 0) \simeq 0) \end{array}}}
\vnbentry{\texttt{sum-\allowbreak{}ag-\allowbreak{}cc-\allowbreak{}in-\allowbreak{}cc}}{\vnbstmt{\begin{array}{@{}l@{}} \forall k \in \mathbb{N}, f \in (\mathbb{N} \to \mathbb{C}). \\ \quad \operatorname{sum\text{-}ag}(\operatorname{normed\text{-}field\text{-}additive\text{-}ag}(\operatorname{cc\text{-}normed\text{-}field}), \\ \quad \quad f, k) \\ \quad \in\; \mathbb{C} \end{array}}}
\vnbentry{\texttt{cc-\allowbreak{}series-\allowbreak{}partial-\allowbreak{}sum-\allowbreak{}in-\allowbreak{}cc}}{\vnbstmt{\begin{array}{@{}l@{}} \forall k \in \mathbb{N}, f \in (\mathbb{N} \to \mathbb{C}). \\ \quad (\operatorname{cc\text{-}series\text{-}partial\text{-}sum}(f, k) \in \mathbb{C}) \end{array}}}
\vnbentry{\texttt{cc-\allowbreak{}series-\allowbreak{}partial-\allowbreak{}sum-\allowbreak{}unfold}}{\vnbstmt{\begin{array}{@{}l@{}} \forall f, k. \\ \quad (\operatorname{cc\text{-}series\text{-}partial\text{-}sum}(f, k) \simeq \operatorname{sum\text{-}ag}(\operatorname{normed\text{-}field\text{-}additive\text{-}ag}(\operatorname{cc\text{-}normed\text{-}field}), f, k)) \end{array}}}
\vnbentry{\texttt{cc-\allowbreak{}series-\allowbreak{}partial-\allowbreak{}sum-\allowbreak{}succ}}{\vnbstmt{\begin{array}{@{}l@{}} \forall f.; \forall k \in \mathbb{N}. \\ \quad (\operatorname{cc\text{-}series\text{-}partial\text{-}sum}(f, k) \in \mathbb{C}) \wedge \\ \quad (f(k) \in \mathbb{C}) \Rightarrow \\ \quad \quad (\operatorname{cc\text{-}series\text{-}partial\text{-}sum}(f, (k + 1)) \simeq (\operatorname{cc\text{-}series\text{-}partial\text{-}sum}(f, k) + f(k))) \end{array}}}
\vnbentry{\texttt{cc-\allowbreak{}magnitude-\allowbreak{}zero}}{\vnbstmt{(\operatorname{magnitude}(0) = 0)}}
\vnbentry{\texttt{cc-\allowbreak{}series-\allowbreak{}partial-\allowbreak{}sum-\allowbreak{}magnitude-\allowbreak{}le}}{\vnbstmt{\begin{array}{@{}l@{}} \forall k \in \mathbb{N}, f \in (\mathbb{N} \to \mathbb{C}), g \in (\mathbb{N} \to \mathbb{R}). \\ \quad \forall j \in \mathbb{N}.\; (g(j) = \operatorname{magnitude}(f(j))) \Rightarrow  \\ \quad \quad \operatorname{magnitude}(\operatorname{cc\text{-}series\text{-}partial\text{-}sum}(f, k)) \\ \quad \quad \leq\; \operatorname{series\text{-}partial\text{-}sum}(g, k) \end{array}}}

\vnbfile{theorem-library/mono-le-limit.scm}{theorem-library/mono-le-limit.scm}
\noindent{\small mono-le-limit.scm -- A TERM OF A NONDECREASING CONVERGENT SEQUENCE IS AT MOST ITS LIMIT, and the corollary the product metric's zero law needs.}\par
\vnbentry{\texttt{rr-\allowbreak{}mono-\allowbreak{}le-\allowbreak{}limit}}{\vnbstmt{\begin{array}{@{}l@{}} \forall f \in (\mathbb{N} \to \mathbb{R}), \ell.; \forall m \in \mathbb{N}. \\ \quad \forall k \in \mathbb{N}.\; (f(k) \leq f((k + 1))) \wedge \\ \quad \operatorname{converges\text{-}to}(\operatorname{rr\text{-}ms}, f, \ell) \Rightarrow \\ \quad \quad (f(m) \leq \ell) \end{array}}}
\vnbentry{\texttt{series-\allowbreak{}term-\allowbreak{}le-\allowbreak{}sum}}{\vnbstmt{\begin{array}{@{}l@{}} \forall f \in (\mathbb{N} \to \mathbb{R}), \ell.; \forall n \in \mathbb{N}. \\ \quad \forall n \in \mathbb{N}.\; (0 \leq f(n)) \wedge \\ \quad \operatorname{series\text{-}converges\text{-}to}(f, \ell) \Rightarrow \\ \quad \quad (f(n) \leq \ell) \end{array}}}
\vnbentry{\texttt{series-\allowbreak{}nonneg-\allowbreak{}zero-\allowbreak{}sum}}{\vnbstmt{\begin{array}{@{}l@{}} \forall f \in (\mathbb{N} \to \mathbb{R}), n \in \mathbb{N}. \\ \quad \forall n \in \mathbb{N}.\; (0 \leq f(n)) \wedge \\ \quad \operatorname{series\text{-}converges\text{-}to}(f, 0) \Rightarrow \\ \quad \quad (f(n) = 0) \end{array}}}

\vnbfile{theorem-library/limit-arithmetic.scm}{theorem-library/limit-arithmetic.scm}
\noindent{\small limit-arithmetic.scm -- THE TWO LIMIT LAWS THE SERIES AND PRODUCT LANES WERE MISSING: `$<$=' carried across TWO limits, and the limit of a pointwise sum.  Plus the series readings of the second.}\par
\vnbentry{\texttt{rr-\allowbreak{}limit-\allowbreak{}le}}{\vnbstmt{\begin{array}{@{}l@{}} \forall f, g \in (\mathbb{N} \to \mathbb{R}), v, a \in \mathbb{R}. \\ \quad \forall n \in \mathbb{N}.\; (f(n) \leq g(n)) \wedge \\ \quad \operatorname{converges\text{-}to}(\operatorname{rr\text{-}ms}, f, v) \wedge \\ \quad \operatorname{converges\text{-}to}(\operatorname{rr\text{-}ms}, g, a) \Rightarrow \\ \quad \quad (v \leq a) \end{array}}}
\vnbentry{\texttt{rr-\allowbreak{}limit-\allowbreak{}add}}{\vnbstmt{\begin{array}{@{}l@{}} \forall f, g, h \in (\mathbb{N} \to \mathbb{R}), v, a \in \mathbb{R}. \\ \quad \forall j \in \mathbb{N}.\; (h(j) = (f(j) + g(j))) \wedge \\ \quad \operatorname{converges\text{-}to}(\operatorname{rr\text{-}ms}, f, v) \wedge \\ \quad \operatorname{converges\text{-}to}(\operatorname{rr\text{-}ms}, g, a) \Rightarrow \\ \quad \quad \operatorname{converges\text{-}to}(\operatorname{rr\text{-}ms}, h, (v + a)) \end{array}}}
\vnbentry{\texttt{series-\allowbreak{}converges-\allowbreak{}to-\allowbreak{}converges}}{\vnbstmt{\begin{array}{@{}l@{}} \forall f \in (\mathbb{N} \to \mathbb{R}), v \in \mathbb{R}. \\ \quad \operatorname{series\text{-}converges\text{-}to}(f, v) \Rightarrow  \\ \quad \quad \operatorname{series\text{-}converges}(f) \end{array}}}
\vnbentry{\texttt{series-\allowbreak{}converges-\allowbreak{}to-\allowbreak{}add}}{\vnbstmt{\begin{array}{@{}l@{}} \forall f, g, h \in (\mathbb{N} \to \mathbb{R}), v, a \in \mathbb{R}. \\ \quad \forall i \in \mathbb{N}.\; (h(i) = (f(i) + g(i))) \wedge \\ \quad \operatorname{series\text{-}converges\text{-}to}(f, v) \wedge \\ \quad \operatorname{series\text{-}converges\text{-}to}(g, a) \Rightarrow \\ \quad \quad \operatorname{series\text{-}converges\text{-}to}(h, (v + a)) \end{array}}}
\vnbentry{\texttt{series-\allowbreak{}limit-\allowbreak{}add}}{\vnbstmt{\begin{array}{@{}l@{}} \forall f, g, h \in (\mathbb{N} \to \mathbb{R}). \\ \quad \forall i \in \mathbb{N}.\; (h(i) = (f(i) + g(i))) \wedge \\ \quad \operatorname{series\text{-}converges}(f) \wedge \\ \quad \operatorname{series\text{-}converges}(g) \Rightarrow \\ \quad \quad \operatorname{series\text{-}limit}(h) \\ \quad \quad =\; (\operatorname{series\text{-}limit}(f) + \operatorname{series\text{-}limit}(g)) \end{array}}}

\vnbfile{theorem-library/ell-two.scm}{theorem-library/ell-two.scm}
\noindent{\small ell-two.scm -- ELL-TWO, the square-summable complex sequences, and the closure of that carrier under the pointwise vector operations.}\par
\vnbentry{\texttt{rr-\allowbreak{}sq-\allowbreak{}add-\allowbreak{}le}}{\vnbstmt{\begin{array}{@{}l@{}} \forall u, v \in \mathbb{R}. \\ \quad (((u + v) \cdot (u + v)) \leq (((u \cdot u) + (u \cdot u)) + ((v \cdot v) + (v \cdot v)))) \end{array}}}
\vnbentry{\texttt{cc-\allowbreak{}magnitude-\allowbreak{}sq-\allowbreak{}add-\allowbreak{}le}}{\vnbstmt{\begin{array}{@{}l@{}} \forall p, q \in \mathbb{C}. \\ \quad (\operatorname{magnitude}((p + q)) \cdot \operatorname{magnitude}((p + q))) \\ \quad \leq\; (((\operatorname{magnitude}(p) \cdot \operatorname{magnitude}(p)) + (\operatorname{magnitude}(p) \cdot \operatorname{magnitude}(p))) + ((\operatorname{magnitude}(q) \cdot \operatorname{magnitude}(q)) + (\operatorname{magnitude}(q) \cdot \operatorname{magnitude}(q)))) \end{array}}}
\vnbentry{\texttt{ell-\allowbreak{}two-\allowbreak{}membership}}{\vnbstmt{\begin{array}{@{}l@{}} \forall y. \\ \quad ((y \in \operatorname{ell\text{-}two}) \Leftrightarrow ((y \in \mathbb{C}^{\mathbb{N}}) \wedge \operatorname{series\text{-}converges}((\lambda k \in \mathbb{N}.\; (\operatorname{magnitude}(y(k)) \cdot \operatorname{magnitude}(y(k))))))) \end{array}}}
\vnbentry{\texttt{ell-\allowbreak{}two-\allowbreak{}add-\allowbreak{}closed}}{\vnbstmt{\begin{array}{@{}l@{}} \forall x, y \in \operatorname{ell\text{-}two}, z \in \mathbb{C}^{\mathbb{N}}. \\ \quad \forall k \in \mathbb{N}.\; (z(k) = (x(k) + y(k))) \Rightarrow  \\ \quad \quad (z \in \operatorname{ell\text{-}two}) \end{array}}}
\vnbentry{\texttt{zero-\allowbreak{}seq-\allowbreak{}in-\allowbreak{}fun}}{\vnbstmt{((\lambda n \in \mathbb{N}.\; 0) \in (\mathbb{N} \to \mathbb{R}))}}
\vnbentry{\texttt{series-\allowbreak{}partial-\allowbreak{}sum-\allowbreak{}zero-\allowbreak{}seq}}{\vnbstmt{\begin{array}{@{}l@{}} \forall k \in \mathbb{N}. \\ \quad (\operatorname{series\text{-}partial\text{-}sum}((\lambda n \in \mathbb{N}.\; 0), k) = 0) \end{array}}}
\vnbentry{\texttt{zero-\allowbreak{}series-\allowbreak{}converges}}{\vnbstmt{\operatorname{series\text{-}converges}((\lambda n \in \mathbb{N}.\; 0))}}
\vnbentry{\texttt{ell-\allowbreak{}two-\allowbreak{}zero-\allowbreak{}in}}{\vnbstmt{((\lambda n \in \mathbb{N}.\; 0) \in \operatorname{ell\text{-}two})}}
\vnbentry{\texttt{ell-\allowbreak{}two-\allowbreak{}neg-\allowbreak{}closed}}{\vnbstmt{\begin{array}{@{}l@{}} \forall x \in \operatorname{ell\text{-}two}, z \in \mathbb{C}^{\mathbb{N}}. \\ \quad \forall k \in \mathbb{N}.\; (z(k) = -x(k)) \Rightarrow  \\ \quad \quad (z \in \operatorname{ell\text{-}two}) \end{array}}}

\vnbfile{theorem-library/dyadic-weights.scm}{theorem-library/dyadic-weights.scm}
\noindent{\small dyadic-weights.scm -- THE CANONICAL PRODUCT-METRIC WEIGHTS ARE SUMMABLE, proven -- retiring `product-metric-default-summable' (structure-library/product-metric.scm), the last asserted support of that file that is not blocked on a bridge somebody else is proving:}\par
\vnbentry{\texttt{power-\allowbreak{}closed-\allowbreak{}at}}{\vnbstmt{\begin{array}{@{}l@{}} \forall n \in \mathbb{N}, x \in \mathbb{R}. \\ \quad (\operatorname{power}(x, n) \in \mathbb{R}) \end{array}}}
\vnbentry{\texttt{power-\allowbreak{}real-\allowbreak{}closed}}{\vnbstmt{\begin{array}{@{}l@{}} \forall x \in \mathbb{R}, n \in \mathbb{N}. \\ \quad (\operatorname{power}(x, n) \in \mathbb{R}) \end{array}}}
\vnbentry{\texttt{rr-\allowbreak{}zero-\allowbreak{}lt-\allowbreak{}two}}{\vnbstmt{(0 < 2)}}
\vnbentry{\texttt{power-\allowbreak{}two-\allowbreak{}pos}}{\vnbstmt{\forall k \in \mathbb{N}.\; (0 < \operatorname{power}(2, k))}}
\vnbentry{\texttt{rr-\allowbreak{}scaled-\allowbreak{}inverse-\allowbreak{}unique}}{\vnbstmt{\begin{array}{@{}l@{}} \forall a, u, v \in \mathbb{R}. \\ \quad (((a \cdot u) = 1) \wedge (((2 \cdot a) \cdot v) = 1)) \Rightarrow  \\ \quad \quad (u = (2 \cdot v)) \end{array}}}
\vnbentry{\texttt{recip-\allowbreak{}power-\allowbreak{}two-\allowbreak{}halves}}{\vnbstmt{\begin{array}{@{}l@{}} \forall k \in \mathbb{N}. \\ \quad \operatorname{recip}(\operatorname{power}(2, k)) \\ \quad =\; (2 \cdot \operatorname{recip}((2 \cdot \operatorname{power}(2, k)))) \end{array}}}
\vnbentry{\texttt{dyadic-\allowbreak{}weight-\allowbreak{}apply}}{\vnbstmt{\begin{array}{@{}l@{}} \forall k \in \mathbb{N}. \\ \quad (\lambda n \in \mathbb{N}.\; /(1, \operatorname{power}(2, (n + 1))))(k) \\ \quad =\; \operatorname{recip}(\operatorname{power}(2, (k + 1))) \end{array}}}
\vnbentry{\texttt{rr-\allowbreak{}halving-\allowbreak{}identity}}{\vnbstmt{\forall b \in \mathbb{R}.\; (((1 - (2 \cdot b)) + b) = (1 - b))}}
\vnbentry{\texttt{dyadic-\allowbreak{}weight-\allowbreak{}in-\allowbreak{}fun}}{\vnbstmt{\begin{array}{@{}l@{}} (\lambda n \in \mathbb{N}.\; /(1, \operatorname{power}(2, (n + 1)))) \\ \in\; (\mathbb{N} \to \mathbb{R}) \end{array}}}
\vnbentry{\texttt{dyadic-\allowbreak{}weight-\allowbreak{}pos}}{\vnbstmt{\begin{array}{@{}l@{}} \forall a \in \mathbb{N}. \\ \quad (0 < (\lambda n \in \mathbb{N}.\; /(1, \operatorname{power}(2, (n + 1))))(a)) \end{array}}}
\vnbentry{\texttt{dyadic-\allowbreak{}partial-\allowbreak{}sum}}{\vnbstmt{\begin{array}{@{}l@{}} \forall k \in \mathbb{N}. \\ \quad \operatorname{series\text{-}partial\text{-}sum}((\lambda n \in \mathbb{N}.\; /(1, \operatorname{power}(2, (n + 1)))), k) \\ \quad =\; (1 - \operatorname{recip}(\operatorname{power}(2, k))) \end{array}}}
\vnbentry{\texttt{product-\allowbreak{}metric-\allowbreak{}default-\allowbreak{}summable}}{\vnbstmt{\operatorname{summable\text{-}weight}((\lambda n \in \mathbb{N}.\; /(1, \operatorname{power}(2, (n + 1)))))}}

\vnbfile{theorem-library/rake-inequalities.scm}{theorem-library/rake-inequalities.scm}
\noindent{\small rake-inequalities.scm -- the real-inequality leaves of the rake (batch 5, batch Q), PROVEN.  Nine asserted supports, and the bridge that four of them turned out to need.}\par
\vnbentry{\texttt{rr-\allowbreak{}young-\allowbreak{}2}}{\vnbstmt{\begin{array}{@{}l@{}} \forall a, b \in \mathbb{R}. \\ \quad ((a \cdot b) \leq /(((a \cdot a) + (b \cdot b)), 2)) \end{array}}}
\vnbentry{\texttt{rr-\allowbreak{}amgm-\allowbreak{}2}}{\vnbstmt{\begin{array}{@{}l@{}} \forall a, b \in \mathbb{R}. \\ \quad ((4 \cdot (a \cdot b)) \leq ((a + b) \cdot (a + b))) \end{array}}}
\vnbentry{\texttt{rr-\allowbreak{}qm-\allowbreak{}am-\allowbreak{}2}}{\vnbstmt{\begin{array}{@{}l@{}} \forall a, b \in \mathbb{R}. \\ \quad (((a + b) \cdot (a + b)) \leq (2 \cdot ((a \cdot a) + (b \cdot b)))) \end{array}}}
\vnbentry{\texttt{rr-\allowbreak{}cauchy-\allowbreak{}schwarz-\allowbreak{}2}}{\vnbstmt{\begin{array}{@{}l@{}} \forall a, b, c, d \in \mathbb{R}. \\ \quad (((a \cdot c) + (b \cdot d)) \cdot ((a \cdot c) + (b \cdot d))) \\ \quad \leq\; (((a \cdot a) + (b \cdot b)) \cdot ((c \cdot c) + (d \cdot d))) \end{array}}}
\vnbentry{\texttt{rr-\allowbreak{}pos-\allowbreak{}shrink}}{\vnbstmt{\begin{array}{@{}l@{}} \forall s. \\ \quad \operatorname{pos\text{-}rr}(s) \Rightarrow  \\ \quad \quad \exists d.\; (\operatorname{pos\text{-}rr}(d) \wedge (d < s)) \end{array}}}
\vnbentry{\texttt{rr-\allowbreak{}bernoulli-\allowbreak{}ind}}{\vnbstmt{\begin{array}{@{}l@{}} \forall n \in \mathbb{N}, x \in \mathbb{R}. \\ \quad (-1 \leq x) \Rightarrow  \\ \quad \quad ((1 + (n \cdot x)) \leq \operatorname{power}((1 + x), n)) \end{array}}}
\vnbentry{\texttt{rr-\allowbreak{}bernoulli}}{\vnbstmt{\begin{array}{@{}l@{}} \forall x \in \mathbb{R}, n \in \mathbb{N}. \\ \quad (-1 \leq x) \Rightarrow  \\ \quad \quad ((1 + (n \cdot x)) \leq \operatorname{power}((1 + x), n)) \end{array}}}
\vnbentry{\texttt{cra-\allowbreak{}ag-\allowbreak{}rr-\allowbreak{}carr}}{\vnbstmt{(\operatorname{carr}(\operatorname{commutative\text{-}ring\text{-}additive\text{-}ag}(\operatorname{rr\text{-}normed\text{-}field})) \simeq \mathbb{R})}}
\vnbentry{\texttt{cra-\allowbreak{}ag-\allowbreak{}rr-\allowbreak{}iden}}{\vnbstmt{(\operatorname{iden}(\operatorname{commutative\text{-}ring\text{-}additive\text{-}ag}(\operatorname{rr\text{-}normed\text{-}field})) \simeq 0)}}
\vnbentry{\texttt{cra-\allowbreak{}ag-\allowbreak{}rr-\allowbreak{}opr}}{\vnbstmt{(\operatorname{opr}(\operatorname{commutative\text{-}ring\text{-}additive\text{-}ag}(\operatorname{rr\text{-}normed\text{-}field})) \simeq (\lambda \langle \operatorname{x\_}, \operatorname{y\_} \rangle \in (\mathbb{R} \times \mathbb{R}).\; (\operatorname{x\_} + \operatorname{y\_})))}}
\vnbentry{\texttt{cra-\allowbreak{}ag-\allowbreak{}is-\allowbreak{}nf-\allowbreak{}ag}}{\vnbstmt{\begin{array}{@{}l@{}} \forall r. \\ \quad (\operatorname{commutative\text{-}ring\text{-}additive\text{-}ag}(r) \simeq \operatorname{normed\text{-}field\text{-}additive\text{-}ag}(r)) \end{array}}}
\vnbentry{\texttt{cra-\allowbreak{}ag-\allowbreak{}rr-\allowbreak{}is-\allowbreak{}abelian-\allowbreak{}group}}{\vnbstmt{\operatorname{is\text{-}abelian\text{-}group}(\operatorname{commutative\text{-}ring\text{-}additive\text{-}ag}(\operatorname{rr\text{-}normed\text{-}field}))}}
\vnbentry{\texttt{cra-\allowbreak{}ag-\allowbreak{}rr-\allowbreak{}opr-\allowbreak{}apply}}{\vnbstmt{\begin{array}{@{}l@{}} \forall u, v \in \mathbb{R}. \\ \quad (\operatorname{opr}(\operatorname{commutative\text{-}ring\text{-}additive\text{-}ag}(\operatorname{rr\text{-}normed\text{-}field}))(u, v) \simeq (u + v)) \end{array}}}
\vnbentry{\texttt{rr-\allowbreak{}nonneg-\allowbreak{}has-\allowbreak{}zero}}{\vnbstmt{\begin{array}{@{}l@{}} \operatorname{iden}(\operatorname{commutative\text{-}ring\text{-}additive\text{-}ag}(\operatorname{rr\text{-}normed\text{-}field})) \\ \in\; \{t \in \mathbb{R} : (0 \leq t)\} \end{array}}}
\vnbentry{\texttt{rr-\allowbreak{}nonneg-\allowbreak{}opr-\allowbreak{}closed}}{\vnbstmt{\begin{array}{@{}l@{}} \forall x, y \in \{t \in \mathbb{R} : (0 \leq t)\}. \\ \quad \operatorname{opr}(\operatorname{commutative\text{-}ring\text{-}additive\text{-}ag}(\operatorname{rr\text{-}normed\text{-}field}))(x, y) \\ \quad \in\; \{t \in \mathbb{R} : (0 \leq t)\} \end{array}}}
\vnbentry{\texttt{finsum-\allowbreak{}rr-\allowbreak{}in-\allowbreak{}rr}}{\vnbstmt{\begin{array}{@{}l@{}} \forall s \in \mathbf{Set}, f \in (s \to \mathbb{R}). \\ \quad (|s| \in \mathbb{N}) \Rightarrow  \\ \quad \quad \operatorname{finsum}(\operatorname{commutative\text{-}ring\text{-}additive\text{-}ag}(\operatorname{rr\text{-}normed\text{-}field}), \\ \quad \quad \quad f, s) \\ \quad \quad \in\; \mathbb{R} \end{array}}}
\vnbentry{\texttt{finsum-\allowbreak{}rr-\allowbreak{}nonneg}}{\vnbstmt{\begin{array}{@{}l@{}} \forall s \in \mathbf{Set}, f \in (s \to \mathbb{R}). \\ \quad ((|s| \in \mathbb{N}) \wedge \forall z \in s.\; (0 \leq f(z))) \Rightarrow  \\ \quad \quad 0 \\ \quad \quad \leq\; \operatorname{finsum}(\operatorname{commutative\text{-}ring\text{-}additive\text{-}ag}(\operatorname{rr\text{-}normed\text{-}field}), \\ \quad \quad \quad f, s) \end{array}}}
\vnbentry{\texttt{finsum-\allowbreak{}sq-\allowbreak{}nonneg}}{\vnbstmt{\begin{array}{@{}l@{}} \forall s \in \mathbf{Set}, a \in (s \to \mathbb{R}). \\ \quad (|s| \in \mathbb{N}) \Rightarrow  \\ \quad \quad 0 \\ \quad \quad \leq\; \operatorname{finsum}(\operatorname{commutative\text{-}ring\text{-}additive\text{-}ag}(\operatorname{rr\text{-}normed\text{-}field}), \\ \quad \quad \quad (\lambda i \in s.\; (a(i) \cdot a(i))), s) \end{array}}}
\vnbentry{\texttt{finsum-\allowbreak{}le-\allowbreak{}termwise}}{\vnbstmt{\begin{array}{@{}l@{}} \forall s \in \mathbf{Set}, a, b \in (s \to \mathbb{R}). \\ \quad ((|s| \in \mathbb{N}) \wedge \forall i \in s.\; (a(i) \leq b(i))) \Rightarrow  \\ \quad \quad \operatorname{finsum}(\operatorname{commutative\text{-}ring\text{-}additive\text{-}ag}(\operatorname{rr\text{-}normed\text{-}field}), \\ \quad \quad \quad a, s) \\ \quad \quad \leq\; \operatorname{finsum}(\operatorname{commutative\text{-}ring\text{-}additive\text{-}ag}(\operatorname{rr\text{-}normed\text{-}field}), \\ \quad \quad \quad b, s) \end{array}}}
\vnbentry{\texttt{finsum-\allowbreak{}abs-\allowbreak{}triangle}}{\vnbstmt{\begin{array}{@{}l@{}} \forall s \in \mathbf{Set}, a \in (s \to \mathbb{R}). \\ \quad (|s| \in \mathbb{N}) \Rightarrow  \\ \quad \quad \lvert \operatorname{finsum}(\operatorname{commutative\text{-}ring\text{-}additive\text{-}ag}(\operatorname{rr\text{-}normed\text{-}field}), a, s) \rvert \\ \quad \quad \leq\; \operatorname{finsum}(\operatorname{commutative\text{-}ring\text{-}additive\text{-}ag}(\operatorname{rr\text{-}normed\text{-}field}), \\ \quad \quad \quad (\lambda i \in s.\; \lvert a(i) \rvert), s) \end{array}}}

\vnbfile{theorem-library/deriv-power.scm}{theorem-library/deriv-power.scm}
\noindent{\small deriv-power.scm -- the POWER RULE for the real monomial, PROVEN `modulo 0'.}\par
\vnbentry{\texttt{pw-\allowbreak{}step-\allowbreak{}identity}}{\vnbstmt{\begin{array}{@{}l@{}} \forall c, u, v. \\ \quad ((c \in \mathbb{R}) \wedge ((u \in \mathbb{R}) \wedge (v \in \mathbb{R}))) \Rightarrow  \\ \quad \quad (((c + 1) \cdot (u \cdot v)) = ((1 \cdot (u \cdot v)) + (u \cdot (c \cdot v)))) \end{array}}}
\vnbentry{\texttt{pow-\allowbreak{}lam-\allowbreak{}in-\allowbreak{}fun}}{\vnbstmt{\begin{array}{@{}l@{}} \forall n \in \mathbb{N}. \\ \quad ((\lambda x \in \mathbb{R}.\; \operatorname{power}(x, n)) \in (\mathbb{R} \to \mathbb{R})) \end{array}}}
\vnbentry{\texttt{deriv-\allowbreak{}power}}{\vnbstmt{\begin{array}{@{}l@{}} \forall n \in \mathbb{N}, a \in \mathbb{R}. \\ \quad \operatorname{is\text{-}diff\text{-}at}((\lambda x \in \mathbb{R}.\; \operatorname{power}(x, (n + 1))), a, ((n + 1) \cdot \operatorname{power}(a, n))) \end{array}}}

\vnbfile{theorem-library/deriv-polynomial.scm}{theorem-library/deriv-polynomial.scm}
\noindent{\small deriv-polynomial.scm -- rung 1 of the integration arc: the DERIVATIVE OF A POLYNOMIAL FUNCTION, in the coefficient-lambda representation.}\par
\vnbentry{\texttt{poly-\allowbreak{}term-\allowbreak{}lam-\allowbreak{}in-\allowbreak{}fun}}{\vnbstmt{\begin{array}{@{}l@{}} \forall f \in (\mathbb{N} \to \mathbb{R}), x \in \mathbb{R}. \\ \quad (\lambda k \in \mathbb{N}.\; (f(k) \cdot \operatorname{power}(x, k))) \\ \quad \in\; (\mathbb{N} \to \mathbb{R}) \end{array}}}
\vnbentry{\texttt{poly-\allowbreak{}dseq-\allowbreak{}lam-\allowbreak{}in-\allowbreak{}fun}}{\vnbstmt{\begin{array}{@{}l@{}} \forall f \in (\mathbb{N} \to \mathbb{R}), t \in \mathbb{R}. \\ \quad (\lambda j \in \mathbb{N}.\; (((j + 1) \cdot f((j + 1))) \cdot \operatorname{power}(t, j))) \\ \quad \in\; (\mathbb{N} \to \mathbb{R}) \end{array}}}
\vnbentry{\texttt{poly-\allowbreak{}lam-\allowbreak{}in-\allowbreak{}fun}}{\vnbstmt{\begin{array}{@{}l@{}} \forall f \in (\mathbb{N} \to \mathbb{R}), n \in \mathbb{N}. \\ \quad (\lambda x \in \mathbb{R}.\; \\ \quad \quad \operatorname{series\text{-}partial\text{-}sum}((\lambda k \in \mathbb{N}.\; (f(k) \cdot \operatorname{power}(x, k))), n)) \\ \quad \in\; (\mathbb{R} \to \mathbb{R}) \end{array}}}
\vnbentry{\texttt{poly-\allowbreak{}step-\allowbreak{}identity}}{\vnbstmt{\begin{array}{@{}l@{}} \forall c, u, v. \\ \quad ((c \in \mathbb{R}) \wedge ((u \in \mathbb{R}) \wedge (v \in \mathbb{R}))) \Rightarrow  \\ \quad \quad (((c \cdot u) \cdot v) = (u \cdot (c \cdot v))) \end{array}}}
\vnbentry{\texttt{deriv-\allowbreak{}scalar-\allowbreak{}mult}}{\vnbstmt{\begin{array}{@{}l@{}} \forall c \in \mathbb{R}, f, t, l. \\ \quad \operatorname{is\text{-}diff\text{-}at}(f, t, l) \Rightarrow  \\ \quad \quad \operatorname{is\text{-}diff\text{-}at}((\lambda x \in \mathbb{R}.\; (c \cdot f(x))), t, (c \cdot l)) \end{array}}}
\vnbentry{\texttt{deriv-\allowbreak{}coef-\allowbreak{}monomial}}{\vnbstmt{\begin{array}{@{}l@{}} \forall c \in \mathbb{R}, g \in \mathbb{N}, t \in \mathbb{R}. \\ \quad \operatorname{is\text{-}diff\text{-}at}((\lambda x \in \mathbb{R}.\; (c \cdot \operatorname{power}(x, (g + 1)))), t, \\ \quad \quad (((g + 1) \cdot c) \cdot \operatorname{power}(t, g))) \end{array}}}
\vnbentry{\texttt{deriv-\allowbreak{}polynomial}}{\vnbstmt{\begin{array}{@{}l@{}} \forall n \in \mathbb{N}, f \in (\mathbb{N} \to \mathbb{R}), t \in \mathbb{R}. \\ \quad \operatorname{is\text{-}diff\text{-}at}((\lambda x \in \mathbb{R}.\; \\ \quad \quad \quad \operatorname{series\text{-}partial\text{-}sum}((\lambda k \in \mathbb{N}.\; (f(k) \cdot \operatorname{power}(x, k))), (n + 1))), \\ \quad \quad t, \\ \quad \quad \operatorname{series\text{-}partial\text{-}sum}((\lambda j \in \mathbb{N}.\; (((j + 1) \cdot f((j + 1))) \cdot \operatorname{power}(t, j))), n)) \end{array}}}

\vnbfile{theorem-library/poly-antiderivative.scm}{theorem-library/poly-antiderivative.scm}
\noindent{\small poly-antiderivative.scm -- calculus.pdf EXAMPLE 4.7, "every polynomial function is antiderivable", with the antiderivative written out:}\par
\vnbentry{\texttt{anti-\allowbreak{}term-\allowbreak{}in-\allowbreak{}rr}}{\vnbstmt{\begin{array}{@{}l@{}} \forall f \in (\mathbb{N} \to \mathbb{R}), k \in \mathbb{N}, x \in \mathbb{R}. \\ \quad (((\operatorname{recip}((k + 1)) \cdot f(k)) \cdot \operatorname{power}(x, (k + 1))) \in \mathbb{R}) \end{array}}}
\vnbentry{\texttt{anti-\allowbreak{}term-\allowbreak{}lam-\allowbreak{}in-\allowbreak{}fun}}{\vnbstmt{\begin{array}{@{}l@{}} \forall f \in (\mathbb{N} \to \mathbb{R}), x \in \mathbb{R}. \\ \quad (\lambda k \in \mathbb{N}.\; \\ \quad \quad ((\operatorname{recip}((k + 1)) \cdot f(k)) \cdot \operatorname{power}(x, (k + 1)))) \\ \quad \in\; (\mathbb{N} \to \mathbb{R}) \end{array}}}
\vnbentry{\texttt{anti-\allowbreak{}lam-\allowbreak{}in-\allowbreak{}fun}}{\vnbstmt{\begin{array}{@{}l@{}} \forall f \in (\mathbb{N} \to \mathbb{R}), m \in \mathbb{N}. \\ \quad (\lambda x \in \mathbb{R}.\; \\ \quad \quad \operatorname{series\text{-}partial\text{-}sum}((\lambda k \in \mathbb{N}.\; \\ \quad \quad \quad \quad ((\operatorname{recip}((k + 1)) \cdot f(k)) \cdot \operatorname{power}(x, (k + 1)))), \\ \quad \quad \quad m)) \\ \quad \in\; (\mathbb{R} \to \mathbb{R}) \end{array}}}
\vnbentry{\texttt{rr-\allowbreak{}zero-\allowbreak{}plus}}{\vnbstmt{\forall t \in \mathbb{R}.\; ((0 + t) = t)}}
\vnbentry{\texttt{recip-\allowbreak{}succ-\allowbreak{}cancel}}{\vnbstmt{\begin{array}{@{}l@{}} \forall m \in \mathbb{N}, c, v \in \mathbb{R}. \\ \quad ((c \cdot v) = (((m + 1) \cdot (\operatorname{recip}((m + 1)) \cdot c)) \cdot v)) \end{array}}}
\vnbentry{\texttt{deriv-\allowbreak{}anti-\allowbreak{}monomial}}{\vnbstmt{\begin{array}{@{}l@{}} \forall c \in \mathbb{R}, m \in \mathbb{N}, t \in \mathbb{R}. \\ \quad \operatorname{is\text{-}diff\text{-}at}((\lambda x \in \mathbb{R}.\; \\ \quad \quad \quad ((\operatorname{recip}((m + 1)) \cdot c) \cdot \operatorname{power}(x, (m + 1)))), \\ \quad \quad t, (c \cdot \operatorname{power}(t, m))) \end{array}}}
\vnbentry{\texttt{poly-\allowbreak{}antiderivative}}{\vnbstmt{\begin{array}{@{}l@{}} \forall n \in \mathbb{N}, f \in (\mathbb{N} \to \mathbb{R}), t \in \mathbb{R}. \\ \quad \operatorname{is\text{-}diff\text{-}at}((\lambda x \in \mathbb{R}.\; \\ \quad \quad \quad \operatorname{series\text{-}partial\text{-}sum}((\lambda k \in \mathbb{N}.\; \\ \quad \quad \quad \quad \quad ((\operatorname{recip}((k + 1)) \cdot f(k)) \cdot \operatorname{power}(x, (k + 1)))), \\ \quad \quad \quad \quad (n + 1))), \\ \quad \quad t, \\ \quad \quad \operatorname{series\text{-}partial\text{-}sum}((\lambda k \in \mathbb{N}.\; (f(k) \cdot \operatorname{power}(t, k))), (n + 1))) \end{array}}}

\vnbfile{theorem-library/product-summable.scm}{theorem-library/product-summable.scm}
\noindent{\small product-summable.scm -- THE DEFINING SERIES OF THE PRODUCT METRIC CONVERGES, proven -- retiring the asserted `product-weighted-summable' (structure-library/product-metric.scm), which is what makes the IOTA in PRODUCT-METRIC-W a well-defined distance, and with it the second asserted}\par
\vnbentry{\texttt{product-\allowbreak{}carrier-\allowbreak{}unfold}}{\vnbstmt{\begin{array}{@{}l@{}} \forall s. \\ \quad (\operatorname{product\text{-}carrier}(s) \simeq \{x \in (\mathbb{N} \to \operatorname{big\text{-}union}(n, \mathbb{N}, \operatorname{pts}(s(n)))) : \forall n \in \mathbb{N}.\; (x(n) \in \operatorname{pts}(s(n)))\}) \end{array}}}
\vnbentry{\texttt{product-\allowbreak{}carrier-\allowbreak{}coord}}{\vnbstmt{\begin{array}{@{}l@{}} \forall s.; \forall x \in \operatorname{product\text{-}carrier}(s), n \in \mathbb{N}. \\ \quad (x(n) \in \operatorname{pts}(s(n))) \end{array}}}
\vnbentry{\texttt{bdd-\allowbreak{}metric-\allowbreak{}weight-\allowbreak{}bound}}{\vnbstmt{\begin{array}{@{}l@{}} \forall s.; \forall u, v \in \operatorname{pts}(s), n \in \mathbb{R}. \\ \quad (\operatorname{is\text{-}metric\text{-}space}(s) \wedge (0 \leq n)) \Rightarrow  \\ \quad \quad ((n \cdot \operatorname{dist}(\operatorname{bdd\text{-}metric}(s))(u, v)) \in \mathbb{R}) \wedge \\ \quad \quad (0 \leq (n \cdot \operatorname{dist}(\operatorname{bdd\text{-}metric}(s))(u, v))) \wedge \\ \quad \quad ((n \cdot \operatorname{dist}(\operatorname{bdd\text{-}metric}(s))(u, v)) \leq n) \end{array}}}
\vnbentry{\texttt{product-\allowbreak{}metric-\allowbreak{}carrier}}{\vnbstmt{\begin{array}{@{}l@{}} \forall s, w. \\ \quad (\operatorname{pts}(\operatorname{product\text{-}metric\text{-}w}(s, w)) \simeq \operatorname{product\text{-}carrier}(s)) \end{array}}}
\vnbentry{\texttt{product-\allowbreak{}term-\allowbreak{}seq-\allowbreak{}in-\allowbreak{}fun}}{\vnbstmt{\begin{array}{@{}l@{}} \forall s, w.; \forall x, y \in \operatorname{product\text{-}carrier}(s). \\ \quad (\operatorname{is\text{-}ms\text{-}sequence}(s) \wedge \operatorname{summable\text{-}weight}(w)) \Rightarrow  \\ \quad \quad (\lambda n \in \mathbb{N}.\; \\ \quad \quad \quad (w(n) \cdot \operatorname{dist}(\operatorname{bdd\text{-}metric}(s(n)))(x(n), y(n)))) \\ \quad \quad \in\; (\mathbb{N} \to \mathbb{R}) \end{array}}}
\vnbentry{\texttt{product-\allowbreak{}weighted-\allowbreak{}summable}}{\vnbstmt{\begin{array}{@{}l@{}} \forall s, w.; \forall x, y \in \operatorname{product\text{-}carrier}(s). \\ \quad (\operatorname{is\text{-}ms\text{-}sequence}(s) \wedge \operatorname{summable\text{-}weight}(w)) \Rightarrow  \\ \quad \quad \operatorname{series\text{-}converges}((\lambda n \in \mathbb{N}.\; \\ \quad \quad \quad \quad (w(n) \cdot \operatorname{dist}(\operatorname{bdd\text{-}metric}(s(n)))(x(n), y(n))))) \end{array}}}
\vnbentry{\texttt{product-\allowbreak{}metric-\allowbreak{}distance}}{\vnbstmt{\begin{array}{@{}l@{}} \forall s, w.; \forall x, y \in \operatorname{product\text{-}carrier}(s). \\ \quad (\operatorname{dist}(\operatorname{product\text{-}metric\text{-}w}(s, w))(x, y) \simeq \operatorname{series\text{-}limit}((\lambda n \in \mathbb{N}.\; (w(n) \cdot \operatorname{dist}(\operatorname{bdd\text{-}metric}(s(n)))(x(n), y(n)))))) \end{array}}}
\vnbentry{\texttt{product-\allowbreak{}metric-\allowbreak{}dist-\allowbreak{}converges}}{\vnbstmt{\begin{array}{@{}l@{}} \forall s, w.; \forall x, y \in \operatorname{product\text{-}carrier}(s). \\ \quad (\operatorname{is\text{-}ms\text{-}sequence}(s) \wedge \operatorname{summable\text{-}weight}(w)) \Rightarrow  \\ \quad \quad \operatorname{series\text{-}converges\text{-}to}((\lambda n \in \mathbb{N}.\; \\ \quad \quad \quad \quad (w(n) \cdot \operatorname{dist}(\operatorname{bdd\text{-}metric}(s(n)))(x(n), y(n)))), \\ \quad \quad \quad \operatorname{dist}(\operatorname{product\text{-}metric\text{-}w}(s, w))(x, y)) \end{array}}}

\vnbfile{theorem-library/product-is-metric-space.scm}{theorem-library/product-is-metric-space.scm}
\noindent{\small product-is-metric-space.scm -- THE COUNTABLE PRODUCT OF METRIC SPACES IS A METRIC SPACE, proven -- retiring the asserted `product-is-metric-space' of structure-library/product-metric.scm (line 79).}\par
\vnbentry{\texttt{series-\allowbreak{}limit-\allowbreak{}le}}{\vnbstmt{\begin{array}{@{}l@{}} \forall f, g \in (\mathbb{N} \to \mathbb{R}), v, a \in \mathbb{R}. \\ \quad \forall n \in \mathbb{N}.\; (f(n) \leq g(n)) \wedge \\ \quad \operatorname{series\text{-}converges\text{-}to}(f, v) \wedge \\ \quad \operatorname{series\text{-}converges\text{-}to}(g, a) \Rightarrow \\ \quad \quad (v \leq a) \end{array}}}
\vnbentry{\texttt{series-\allowbreak{}limit-\allowbreak{}iota}}{\vnbstmt{\forall f.\; (\iota(l) \simeq \operatorname{series\text{-}limit}(f))}}
\vnbentry{\texttt{product-\allowbreak{}dist-\allowbreak{}in-\allowbreak{}rr}}{\vnbstmt{\begin{array}{@{}l@{}} \forall s, w.; \forall x, y \in \operatorname{product\text{-}carrier}(s). \\ \quad (\operatorname{is\text{-}ms\text{-}sequence}(s) \wedge \operatorname{summable\text{-}weight}(w)) \Rightarrow  \\ \quad \quad (\operatorname{dist}(\operatorname{product\text{-}metric\text{-}w}(s, w))(x, y) \in \mathbb{R}) \end{array}}}
\vnbentry{\texttt{product-\allowbreak{}term-\allowbreak{}nonneg}}{\vnbstmt{\begin{array}{@{}l@{}} \forall s, w.; \forall x, y \in \operatorname{product\text{-}carrier}(s), a \in \mathbb{N}. \\ \quad (\operatorname{is\text{-}ms\text{-}sequence}(s) \wedge \operatorname{summable\text{-}weight}(w)) \Rightarrow  \\ \quad \quad 0 \\ \quad \quad \leq\; (\lambda n \in \mathbb{N}.\; (w(n) \cdot \operatorname{dist}(\operatorname{bdd\text{-}metric}(s(n)))(x(n), y(n))))(a) \end{array}}}
\vnbentry{\texttt{product-\allowbreak{}term-\allowbreak{}self-\allowbreak{}zero}}{\vnbstmt{\begin{array}{@{}l@{}} \forall s, w.; \forall x \in \operatorname{product\text{-}carrier}(s), a \in \mathbb{N}. \\ \quad (\operatorname{is\text{-}ms\text{-}sequence}(s) \wedge \operatorname{summable\text{-}weight}(w)) \Rightarrow  \\ \quad \quad (\lambda n \in \mathbb{N}.\; (w(n) \cdot \operatorname{dist}(\operatorname{bdd\text{-}metric}(s(n)))(x(n), x(n))))(a) \\ \quad \quad =\; 0 \end{array}}}
\vnbentry{\texttt{product-\allowbreak{}term-\allowbreak{}sym}}{\vnbstmt{\begin{array}{@{}l@{}} \forall s, w.; \forall x, y \in \operatorname{product\text{-}carrier}(s), a \in \mathbb{N}. \\ \quad (\operatorname{is\text{-}ms\text{-}sequence}(s) \wedge \operatorname{summable\text{-}weight}(w)) \Rightarrow  \\ \quad \quad (\lambda n \in \mathbb{N}.\; (w(n) \cdot \operatorname{dist}(\operatorname{bdd\text{-}metric}(s(n)))(x(n), y(n))))(a) \\ \quad \quad \leq\; (\lambda n \in \mathbb{N}.\; (w(n) \cdot \operatorname{dist}(\operatorname{bdd\text{-}metric}(s(n)))(y(n), x(n))))(a) \end{array}}}
\vnbentry{\texttt{product-\allowbreak{}sum-\allowbreak{}seq-\allowbreak{}in-\allowbreak{}fun}}{\vnbstmt{\begin{array}{@{}l@{}} \forall s, w.; \forall x, y, z \in \operatorname{product\text{-}carrier}(s). \\ \quad (\operatorname{is\text{-}ms\text{-}sequence}(s) \wedge \operatorname{summable\text{-}weight}(w)) \Rightarrow  \\ \quad \quad (\lambda n \in \mathbb{N}.\; \\ \quad \quad \quad ((w(n) \cdot \operatorname{dist}(\operatorname{bdd\text{-}metric}(s(n)))(x(n), y(n))) + (w(n) \cdot \operatorname{dist}(\operatorname{bdd\text{-}metric}(s(n)))(y(n), z(n))))) \\ \quad \quad \in\; (\mathbb{N} \to \mathbb{R}) \end{array}}}
\vnbentry{\texttt{product-\allowbreak{}sum-\allowbreak{}seq-\allowbreak{}split}}{\vnbstmt{\begin{array}{@{}l@{}} \forall s, w.; \forall x, y, z \in \operatorname{product\text{-}carrier}(s), i \in \mathbb{N}. \\ \quad (\operatorname{is\text{-}ms\text{-}sequence}(s) \wedge \operatorname{summable\text{-}weight}(w)) \Rightarrow  \\ \quad \quad (\lambda n \in \mathbb{N}.\; ((w(n) \cdot \operatorname{dist}(\operatorname{bdd\text{-}metric}(s(n)))(x(n), y(n))) + (w(n) \cdot \operatorname{dist}(\operatorname{bdd\text{-}metric}(s(n)))(y(n), z(n)))))(i) \\ \quad \quad =\; ((\lambda n \in \mathbb{N}.\; (w(n) \cdot \operatorname{dist}(\operatorname{bdd\text{-}metric}(s(n)))(x(n), y(n))))(i) + (\lambda n \in \mathbb{N}.\; (w(n) \cdot \operatorname{dist}(\operatorname{bdd\text{-}metric}(s(n)))(y(n), z(n))))(i)) \end{array}}}
\vnbentry{\texttt{product-\allowbreak{}term-\allowbreak{}triangle}}{\vnbstmt{\begin{array}{@{}l@{}} \forall s, w.; \forall x, y, z \in \operatorname{product\text{-}carrier}(s), a \in \mathbb{N}. \\ \quad (\operatorname{is\text{-}ms\text{-}sequence}(s) \wedge \operatorname{summable\text{-}weight}(w)) \Rightarrow  \\ \quad \quad (\lambda n \in \mathbb{N}.\; (w(n) \cdot \operatorname{dist}(\operatorname{bdd\text{-}metric}(s(n)))(x(n), z(n))))(a) \\ \quad \quad \leq\; (\lambda n \in \mathbb{N}.\; ((w(n) \cdot \operatorname{dist}(\operatorname{bdd\text{-}metric}(s(n)))(x(n), y(n))) + (w(n) \cdot \operatorname{dist}(\operatorname{bdd\text{-}metric}(s(n)))(y(n), z(n)))))(a) \end{array}}}
\vnbentry{\texttt{product-\allowbreak{}is-\allowbreak{}metric-\allowbreak{}space}}{\vnbstmt{\begin{array}{@{}l@{}} \forall s, w. \\ \quad (\operatorname{is\text{-}ms\text{-}sequence}(s) \wedge \operatorname{summable\text{-}weight}(w)) \Rightarrow  \\ \quad \quad \operatorname{is\text{-}metric\text{-}space}(\operatorname{product\text{-}metric\text{-}w}(s, w)) \end{array}}}

\vnbfile{theorem-library/rake-block-tower.scm}{theorem-library/rake-block-tower.scm}
\noindent{\small theorem-library/rake-block-tower.scm ==================================================================== convergence-block-tower (theorem-library/seq-compact-product.scm:91) -- the refinement-tower core of the cross-coordinate diagonal argument, and}\par
\vnbentry{\texttt{block-\allowbreak{}step-\allowbreak{}converges}}{\vnbstmt{\begin{array}{@{}l@{}} \forall s.; \forall h \in (\mathbb{N} \to \operatorname{pts}(s)), j \in \mathrm{Inf}(\mathbb{N}). \\ \quad \operatorname{seq\text{-}compact}(s) \Rightarrow  \\ \quad \quad \exists b \in \mathrm{Inf}(\mathbb{N}). \\ \quad \quad \quad (b \subseteq j) \wedge \\ \quad \quad \quad \exists p \in \operatorname{pts}(s). \\ \quad \quad \quad \quad \operatorname{converges\text{-}along}(s, h, b, p) \end{array}}}
\vnbentry{\texttt{convergence-\allowbreak{}block-\allowbreak{}tower}}{\vnbstmt{\begin{array}{@{}l@{}} \forall a.; \forall q \in (\mathbb{N} \to \operatorname{product\text{-}carrier}(a)). \\ \quad \operatorname{is\text{-}ms\text{-}sequence}(a) \wedge \\ \quad \forall n \in \mathbb{N}.\; \operatorname{seq\text{-}compact}(a(n)) \Rightarrow \\ \quad \quad \exists s \in (\mathbb{N} \to \mathrm{Inf}(\mathbb{N})). \\ \quad \quad \quad \forall k \in \mathbb{N}.\; (s((k + 1)) \subseteq s(k)) \wedge \\ \quad \quad \quad \exists l \in \operatorname{product\text{-}carrier}(a). \forall n \in \mathbb{N}. \\ \quad \quad \quad \quad \operatorname{converges\text{-}along}(a(n), (\lambda i \in \mathbb{N}.\; q(i)(n)), s(n), l(n)) \end{array}}}

\vnbfile{theorem-library/rr-null-scale.scm}{theorem-library/rr-null-scale.scm}
\noindent{\small rr-null-scale.scm -- a NONNEGATIVE CONSTANT MULTIPLE OF A NULL SEQUENCE IS NULL, proven, together with the eps/delta lemma that carries it.}\par
\vnbentry{\texttt{rr-\allowbreak{}scale-\allowbreak{}eps}}{\vnbstmt{\begin{array}{@{}l@{}} \forall c, s \in \mathbb{R}. \\ \quad ((0 \leq c) \wedge \operatorname{pos\text{-}rr}(s)) \Rightarrow  \\ \quad \quad \exists d. \\ \quad \quad \quad \operatorname{pos\text{-}rr}(d) \wedge \\ \quad \quad \quad \forall t \in \mathbb{R}. \\ \quad \quad \quad \quad (((0 \leq t) \wedge (t \leq d)) \Rightarrow ((c \cdot t) \leq s)) \end{array}}}
\vnbentry{\texttt{rr-\allowbreak{}null-\allowbreak{}scale}}{\vnbstmt{\begin{array}{@{}l@{}} \forall c \in \mathbb{R}, f, h \in (\mathbb{N} \to \mathbb{R}). \\ \quad (0 \leq c) \wedge \\ \quad \operatorname{converges\text{-}to}(\operatorname{rr\text{-}ms}, f, 0) \wedge \\ \quad \forall j \in \mathbb{N}.\; (h(j) = (c \cdot f(j))) \Rightarrow \\ \quad \quad \operatorname{converges\text{-}to}(\operatorname{rr\text{-}ms}, h, 0) \end{array}}}

\vnbfile{theorem-library/bdd-metric-convergence.scm}{theorem-library/bdd-metric-convergence.scm}
\noindent{\small bdd-metric-convergence.scm -- CONVERGENCE IS THE SAME IN d AND IN d/(1+d), proven: a sequence converges to L in BDD-METRIC(s) if and only if it converges to L in s.}\par
\vnbentry{\texttt{bdd-\allowbreak{}fn-\allowbreak{}reflect}}{\vnbstmt{\begin{array}{@{}l@{}} \forall a, b \in \mathbb{R}. \\ \quad ((0 \leq a) \wedge ((0 \leq b) \wedge (/(a, (1 + a)) \leq /(b, (1 + b))))) \Rightarrow  \\ \quad \quad (a \leq b) \end{array}}}
\vnbentry{\texttt{bdd-\allowbreak{}metric-\allowbreak{}dist-\allowbreak{}le}}{\vnbstmt{\begin{array}{@{}l@{}} \forall s.; \forall x, y \in \operatorname{pts}(s). \\ \quad \operatorname{is\text{-}metric\text{-}space}(s) \Rightarrow  \\ \quad \quad (\operatorname{dist}(\operatorname{bdd\text{-}metric}(s))(x, y) \leq \operatorname{dist}(s)(x, y)) \end{array}}}
\vnbentry{\texttt{bdd-\allowbreak{}metric-\allowbreak{}dist-\allowbreak{}reflect}}{\vnbstmt{\begin{array}{@{}l@{}} \forall s.; \forall x, y \in \operatorname{pts}(s), e \in \mathbb{R}. \\ \quad \operatorname{is\text{-}metric\text{-}space}(s) \wedge \\ \quad (0 \leq e) \wedge \\ \quad (\operatorname{dist}(\operatorname{bdd\text{-}metric}(s))(x, y) \leq /(e, (1 + e))) \Rightarrow \\ \quad \quad (\operatorname{dist}(s)(x, y) \leq e) \end{array}}}
\vnbentry{\texttt{bdd-\allowbreak{}metric-\allowbreak{}converges-\allowbreak{}fwd}}{\vnbstmt{\begin{array}{@{}l@{}} \forall s.; \forall f \in (\mathbb{N} \to \operatorname{pts}(s)), \ell \in \operatorname{pts}(s). \\ \quad \operatorname{is\text{-}metric\text{-}space}(s) \wedge \\ \quad \operatorname{converges\text{-}to}(s, f, \ell) \Rightarrow \\ \quad \quad \operatorname{converges\text{-}to}(\operatorname{bdd\text{-}metric}(s), f, \ell) \end{array}}}
\vnbentry{\texttt{bdd-\allowbreak{}metric-\allowbreak{}converges-\allowbreak{}bwd}}{\vnbstmt{\begin{array}{@{}l@{}} \forall s.; \forall f \in (\mathbb{N} \to \operatorname{pts}(s)), \ell \in \operatorname{pts}(s). \\ \quad \operatorname{is\text{-}metric\text{-}space}(s) \wedge \\ \quad \operatorname{converges\text{-}to}(\operatorname{bdd\text{-}metric}(s), f, \ell) \Rightarrow \\ \quad \quad \operatorname{converges\text{-}to}(s, f, \ell) \end{array}}}
\vnbentry{\texttt{bdd-\allowbreak{}metric-\allowbreak{}converges-\allowbreak{}iff}}{\vnbstmt{\begin{array}{@{}l@{}} \forall s.; \forall f \in (\mathbb{N} \to \operatorname{pts}(s)), \ell \in \operatorname{pts}(s). \\ \quad \operatorname{is\text{-}metric\text{-}space}(s) \Rightarrow  \\ \quad \quad (\operatorname{converges\text{-}to}(\operatorname{bdd\text{-}metric}(s), f, \ell) \Leftrightarrow \operatorname{converges\text{-}to}(s, f, \ell)) \end{array}}}

\vnbfile{theorem-library/rake-metric-constructions.scm}{theorem-library/rake-metric-constructions.scm}
\noindent{\small rake-metric-constructions.scm -- BATCH S of the 2026-09-18 rake: the metric CONSTRUCTIONS (bounded metric, ball cover, countable product).  Every statement below is its support's statement UNCHANGED, copied from the support site, except where the header says otherwise.}\par
\vnbentry{\texttt{bdd-\allowbreak{}metric-\allowbreak{}id-\allowbreak{}bicontinuous}}{\vnbstmt{\begin{array}{@{}l@{}} \forall s. \\ \quad \operatorname{is\text{-}metric\text{-}space}(s) \Rightarrow  \\ \quad \quad \operatorname{is\text{-}continuous}(s, \operatorname{bdd\text{-}metric}(s), (\lambda x \in \operatorname{pts}(s).\; x)) \wedge \\ \quad \quad \operatorname{is\text{-}continuous}(\operatorname{bdd\text{-}metric}(s), s, (\lambda x \in \operatorname{pts}(\operatorname{bdd\text{-}metric}(s)).\; x)) \end{array}}}
\vnbentry{\texttt{bdd-\allowbreak{}rr-\allowbreak{}ms-\allowbreak{}is-\allowbreak{}metric-\allowbreak{}space}}{\vnbstmt{\operatorname{is\text{-}metric\text{-}space}(\operatorname{bdd\text{-}metric}(\operatorname{rr\text{-}ms}))}}
\vnbentry{\texttt{bdd-\allowbreak{}rr-\allowbreak{}ms-\allowbreak{}bounded}}{\vnbstmt{\begin{array}{@{}l@{}} \forall x, y \in \mathbb{R}. \\ \quad (\operatorname{dist}(\operatorname{bdd\text{-}metric}(\operatorname{rr\text{-}ms}))(x, y) < 1) \end{array}}}
\vnbentry{\texttt{bdd-\allowbreak{}rr-\allowbreak{}equivalent}}{\vnbstmt{\begin{array}{@{}l@{}} \operatorname{is\text{-}continuous}(\operatorname{rr\text{-}ms}, \operatorname{bdd\text{-}metric}(\operatorname{rr\text{-}ms}), (\lambda x \in \mathbb{R}.\; x)) \wedge \\ \operatorname{is\text{-}continuous}(\operatorname{bdd\text{-}metric}(\operatorname{rr\text{-}ms}), \operatorname{rr\text{-}ms}, (\lambda x \in \mathbb{R}.\; x)) \end{array}}}
\vnbentry{\texttt{ball-\allowbreak{}cover-\allowbreak{}is-\allowbreak{}open-\allowbreak{}cover}}{\vnbstmt{\begin{array}{@{}l@{}} \forall s.; \forall r \in \mathbb{R}. \\ \quad (\operatorname{is\text{-}metric\text{-}space}(s) \wedge ((0 \leq r) \wedge \neg (0 = r))) \Rightarrow  \\ \quad \quad \operatorname{is\text{-}open\text{-}cover}(s, \operatorname{ball\text{-}cover}(s, r)) \end{array}}}
\vnbentry{\texttt{product-\allowbreak{}projection-\allowbreak{}continuous}}{\vnbstmt{\begin{array}{@{}l@{}} \forall s, w.; \forall n \in \mathbb{N}. \\ \quad (\operatorname{is\text{-}ms\text{-}sequence}(s) \wedge \operatorname{summable\text{-}weight}(w)) \Rightarrow  \\ \quad \quad \operatorname{is\text{-}continuous}(\operatorname{product\text{-}metric\text{-}w}(s, w), s(n), \operatorname{product\text{-}proj}(s, n)) \end{array}}}

\vnbfile{theorem-library/rake-rr-bounded-ms.scm}{theorem-library/rake-rr-bounded-ms.scm}
\noindent{\small rake-rr-bounded-ms.scm -- the three facts about the constant RR-BOUNDED-MS (rake batch 5, integrator, 2026-09-18).}\par
\vnbentry{\texttt{rr-\allowbreak{}bounded-\allowbreak{}ms-\allowbreak{}is-\allowbreak{}metric-\allowbreak{}space}}{\vnbstmt{\operatorname{is\text{-}metric\text{-}space}(\operatorname{rr\text{-}bounded\text{-}ms})}}
\vnbentry{\texttt{rr-\allowbreak{}bounded-\allowbreak{}ms-\allowbreak{}bounded}}{\vnbstmt{\begin{array}{@{}l@{}} \forall x, y \in \mathbb{R}. \\ \quad (\operatorname{dist}(\operatorname{rr\text{-}bounded\text{-}ms})(x, y) < 1) \end{array}}}
\vnbentry{\texttt{rr-\allowbreak{}bounded-\allowbreak{}equivalent}}{\vnbstmt{\begin{array}{@{}l@{}} \operatorname{is\text{-}continuous}(\operatorname{rr\text{-}ms}, \operatorname{rr\text{-}bounded\text{-}ms}, (\lambda x \in \mathbb{R}.\; x)) \wedge \\ \operatorname{is\text{-}continuous}(\operatorname{rr\text{-}bounded\text{-}ms}, \operatorname{rr\text{-}ms}, (\lambda x \in \mathbb{R}.\; x)) \end{array}}}

\vnbfile{calculus/compact-tb-proof.scm}{calculus/compact-tb-proof.scm}
\noindent{\small compact-tb-proof.scm -- machine proof of  compact =$>$ totally bounded (calculus.pdf Prop 3.12, the (1)=$>$(4) half), assembled from the two warranted ball-cover lemmas (compactness.scm).  Run: ./prover calculus/compact-tb-proof.scm}\par
\vnbentry{\texttt{compact-\allowbreak{}implies-\allowbreak{}totally-\allowbreak{}bounded}}{\vnbstmt{\begin{array}{@{}l@{}} \forall s. \\ \quad \operatorname{is\text{-}compact}(s) \Rightarrow  \\ \quad \quad \operatorname{totally\text{-}bounded}(s) \end{array}}}

\vnbfile{theorem-library/rake-compact-equivalences.scm}{theorem-library/rake-compact-equivalences.scm}
\noindent{\small rake-compact-equivalences.scm -- Prop 3.12, the two remaining equivalences (batch 8, assignment 8-A; helper prefix `r9a-').}\par
\vnbentry{\texttt{compact-\allowbreak{}iff-\allowbreak{}cluster-\allowbreak{}point}}{\vnbstmt{\begin{array}{@{}l@{}} \forall s. \\ \quad \operatorname{is\text{-}metric\text{-}space}(s) \Rightarrow  \\ \quad \quad (\operatorname{is\text{-}compact}(s) \Leftrightarrow \forall f \in (\mathbb{N} \to \operatorname{pts}(s)).\; \exists x.\; \operatorname{cluster\text{-}point}(s, f, x)) \end{array}}}
\vnbentry{\texttt{compact-\allowbreak{}iff-\allowbreak{}tb-\allowbreak{}complete}}{\vnbstmt{\begin{array}{@{}l@{}} \forall s. \\ \quad \operatorname{is\text{-}metric\text{-}space}(s) \Rightarrow  \\ \quad \quad (\operatorname{is\text{-}compact}(s) \Leftrightarrow (\operatorname{totally\text{-}bounded}(s) \wedge \operatorname{is\text{-}complete}(s))) \end{array}}}

\vnbfile{theorem-library/rake-compact-fip.scm}{theorem-library/rake-compact-fip.scm}
\noindent{\small rake-compact-fip.scm -- Prop 3.12 (1) $<$=$>$ (2), the finite-intersection characterisation of compactness, PROVEN (2026-09-20, batch 11-E; helper prefix `rcf-').}\par
\vnbentry{\texttt{compact-\allowbreak{}implies-\allowbreak{}fip}}{\vnbstmt{\begin{array}{@{}l@{}} \forall s, c. \\ \quad \operatorname{is\text{-}metric\text{-}space}(s) \wedge \\ \quad \operatorname{is\text{-}compact}(s) \wedge \\ \quad \operatorname{has\text{-}fip}(s, c) \Rightarrow \\ \quad \quad \exists p.\; (p \in \operatorname{intersection\text{-}of}(c)) \end{array}}}
\vnbentry{\texttt{fip-\allowbreak{}implies-\allowbreak{}compact}}{\vnbstmt{\begin{array}{@{}l@{}} \forall s. \\ \quad \operatorname{is\text{-}metric\text{-}space}(s) \wedge \\ \quad \forall c. \\ \quad \quad \operatorname{has\text{-}fip}(s, c) \Rightarrow  \\ \quad \quad \quad \exists p.\; (p \in \operatorname{intersection\text{-}of}(c)) \Rightarrow \\ \quad \quad \operatorname{is\text{-}compact}(s) \end{array}}}
\vnbentry{\texttt{compact-\allowbreak{}iff-\allowbreak{}fip}}{\vnbstmt{\begin{array}{@{}l@{}} \forall s. \\ \quad \operatorname{is\text{-}metric\text{-}space}(s) \Rightarrow  \\ \quad \quad (\operatorname{is\text{-}compact}(s) \Leftrightarrow \forall c.\; (\operatorname{has\text{-}fip}(s, c) \Rightarrow \exists p.\; (p \in \operatorname{intersection\text{-}of}(c)))) \end{array}}}

\vnbfile{theorem-library/compact-separable-proof.scm}{theorem-library/compact-separable-proof.scm}
\noindent{\small theorem-library/compact-separable-proof.scm}\par
\vnbentry{\texttt{compact-\allowbreak{}metric-\allowbreak{}is-\allowbreak{}separable}}{\vnbstmt{\begin{array}{@{}l@{}} \forall s. \\ \quad (\operatorname{is\text{-}compact}(s) \wedge \exists x.\; (x \in \operatorname{pts}(s))) \Rightarrow  \\ \quad \quad \operatorname{is\text{-}separable}(s) \end{array}}}

\vnbfile{theorem-library/rake-ascoli.scm}{theorem-library/rake-ascoli.scm}
\noindent{\small theorem-library/rake-ascoli.scm ==================================================================== Ascoli-Arzela (sequential, RR-valued) -- batch 8-B.}\par
\vnbentry{\texttt{equicont-\allowbreak{}subseq}}{\vnbstmt{\begin{array}{@{}l@{}} \forall s, t, m, i. \\ \quad \operatorname{is\text{-}equicontinuous}(s, t, m) \wedge \\ \quad \operatorname{strictly\text{-}mono\text{-}nn}(i) \Rightarrow \\ \quad \quad \operatorname{is\text{-}equicontinuous}(s, t, \operatorname{subseq}(m, i)) \end{array}}}
\vnbentry{\texttt{ascoli-\allowbreak{}sequential-\allowbreak{}from-\allowbreak{}diagonal}}{\vnbstmt{\begin{array}{@{}l@{}} \forall s, m. \\ \quad \operatorname{is\text{-}compact}(s) \wedge \\ \quad \exists x.\; (x \in \operatorname{pts}(s)) \wedge \\ \quad (m \in (\mathbb{N} \to (\operatorname{pts}(s) \to \mathbb{R}))) \wedge \\ \quad \forall k \in \mathbb{N}. \\ \quad \quad \operatorname{is\text{-}continuous}(s, \operatorname{rr\text{-}ms}, m(k)) \wedge \\ \quad \operatorname{is\text{-}equicontinuous}(s, \operatorname{rr\text{-}ms}, m) \wedge \\ \quad \operatorname{pointwise\text{-}bounded}(s, m) \wedge \\ \quad \forall q. \\ \quad \quad \operatorname{is\text{-}dense\text{-}seq}(s, q) \Rightarrow  \\ \quad \quad \quad \exists l. \\ \quad \quad \quad \quad \operatorname{strictly\text{-}mono\text{-}nn}(l) \wedge \\ \quad \quad \quad \quad \operatorname{converges\text{-}on}(s, \operatorname{subseq}(m, l), q) \Rightarrow \\ \quad \quad \exists l. \\ \quad \quad \quad \operatorname{strictly\text{-}mono\text{-}nn}(l) \wedge \\ \quad \quad \quad \exists g \in (\operatorname{pts}(s) \to \mathbb{R}). \\ \quad \quad \quad \quad \operatorname{is\text{-}continuous}(s, \operatorname{rr\text{-}ms}, g) \wedge \\ \quad \quad \quad \quad \operatorname{converges\text{-}uniformly}(s, \operatorname{subseq}(m, l), g) \end{array}}}

\vnbfile{theorem-library/converges-dist-null.scm}{theorem-library/converges-dist-null.scm}
\noindent{\small converges-dist-null.scm -- CONVERGENCE IN A METRIC SPACE IS A NULL SEQUENCE OF REAL DISTANCES, proven both ways:}\par
\vnbentry{\texttt{dist-\allowbreak{}seq-\allowbreak{}apply}}{\vnbstmt{\begin{array}{@{}l@{}} \forall t, f, v.; \forall n \in \mathbb{N}. \\ \quad ((\lambda j \in \mathbb{N}.\; \operatorname{dist}(t)(f(j), v))(n) \simeq \operatorname{dist}(t)(f(n), v)) \end{array}}}
\vnbentry{\texttt{dist-\allowbreak{}seq-\allowbreak{}in-\allowbreak{}fun}}{\vnbstmt{\begin{array}{@{}l@{}} \forall t.; \forall f \in (\mathbb{N} \to \operatorname{pts}(t)), v \in \operatorname{pts}(t). \\ \quad \operatorname{is\text{-}metric\text{-}space}(t) \Rightarrow  \\ \quad \quad ((\lambda j \in \mathbb{N}.\; \operatorname{dist}(t)(f(j), v)) \in (\mathbb{N} \to \mathbb{R})) \end{array}}}
\vnbentry{\texttt{converges-\allowbreak{}dist-\allowbreak{}null-\allowbreak{}fwd}}{\vnbstmt{\begin{array}{@{}l@{}} \forall t.; \forall f \in (\mathbb{N} \to \operatorname{pts}(t)), v \in \operatorname{pts}(t). \\ \quad (\operatorname{is\text{-}metric\text{-}space}(t) \wedge \operatorname{converges\text{-}to}(t, f, v)) \Rightarrow  \\ \quad \quad \operatorname{converges\text{-}to}(\operatorname{rr\text{-}ms}, (\lambda j \in \mathbb{N}.\; \operatorname{dist}(t)(f(j), v)), 0) \end{array}}}
\vnbentry{\texttt{converges-\allowbreak{}dist-\allowbreak{}null-\allowbreak{}bwd}}{\vnbstmt{\begin{array}{@{}l@{}} \forall t.; \forall f \in (\mathbb{N} \to \operatorname{pts}(t)), v \in \operatorname{pts}(t). \\ \quad \operatorname{is\text{-}metric\text{-}space}(t) \wedge \\ \quad \operatorname{converges\text{-}to}(\operatorname{rr\text{-}ms}, (\lambda j \in \mathbb{N}.\; \operatorname{dist}(t)(f(j), v)), 0) \Rightarrow \\ \quad \quad \operatorname{converges\text{-}to}(t, f, v) \end{array}}}
\vnbentry{\texttt{converges-\allowbreak{}iff-\allowbreak{}dist-\allowbreak{}null}}{\vnbstmt{\begin{array}{@{}l@{}} \forall t.; \forall f \in (\mathbb{N} \to \operatorname{pts}(t)), v \in \operatorname{pts}(t). \\ \quad \operatorname{is\text{-}metric\text{-}space}(t) \Rightarrow  \\ \quad \quad (\operatorname{converges\text{-}to}(t, f, v) \Leftrightarrow \operatorname{converges\text{-}to}(\operatorname{rr\text{-}ms}, (\lambda j \in \mathbb{N}.\; \operatorname{dist}(t)(f(j), v)), 0)) \end{array}}}

\vnbfile{theorem-library/sequential-continuity.scm}{theorem-library/sequential-continuity.scm}
\noindent{\small sequential-continuity.scm -- THE SEQUENTIAL CHARACTERIZATION OF CONTINUITY, both directions PROVEN:}\par
\vnbentry{\texttt{continuous-\allowbreak{}at-\allowbreak{}implies-\allowbreak{}sequential}}{\vnbstmt{\begin{array}{@{}l@{}} \forall s, t, f, a.; \forall q \in \operatorname{pts}(s)^{\mathbb{N}}. \\ \quad \operatorname{is\text{-}metric\text{-}space}(s) \wedge \\ \quad \operatorname{is\text{-}metric\text{-}space}(t) \wedge \\ \quad (f \in (\operatorname{pts}(s) \to \operatorname{pts}(t))) \wedge \\ \quad (a \in \operatorname{pts}(s)) \wedge \\ \quad \operatorname{is\text{-}continuous\text{-}at}(s, t, f, a) \wedge \\ \quad \operatorname{converges\text{-}to}(s, q, a) \Rightarrow \\ \quad \quad \operatorname{converges\text{-}to}(t, \operatorname{compose}(f, q), f(a)) \end{array}}}
\vnbentry{\texttt{not-\allowbreak{}continuous-\allowbreak{}witness}}{\vnbstmt{\begin{array}{@{}l@{}} \forall s, t, f, a, b, l. \\ \quad \neg \forall \beta \in \operatorname{pts}(s).\; ((\operatorname{dist}(s)(a, \beta) \leq l) \Rightarrow (\operatorname{dist}(t)(f(a), f(\beta)) \leq b)) \Rightarrow \\ \quad \quad \exists \beta \in \operatorname{pts}(s). \\ \quad \quad \quad (\operatorname{dist}(s)(a, \beta) \leq l) \wedge \\ \quad \quad \quad \neg (\operatorname{dist}(t)(f(a), f(\beta)) \leq b) \end{array}}}
\vnbentry{\texttt{rr-\allowbreak{}recip-\allowbreak{}antitone}}{\vnbstmt{\begin{array}{@{}l@{}} \forall u, v \in \mathbb{R}. \\ \quad ((0 < u) \wedge (u \leq v)) \Rightarrow  \\ \quad \quad (\operatorname{recip}(v) \leq \operatorname{recip}(u)) \end{array}}}
\vnbentry{\texttt{nn-\allowbreak{}recip-\allowbreak{}succ-\allowbreak{}antitone}}{\vnbstmt{\begin{array}{@{}l@{}} \forall m, n \in \mathbb{N}. \\ \quad (m \leq n) \Rightarrow  \\ \quad \quad (\operatorname{recip}((n + 1)) \leq \operatorname{recip}((m + 1))) \end{array}}}
\vnbentry{\texttt{sequential-\allowbreak{}implies-\allowbreak{}continuous-\allowbreak{}at}}{\vnbstmt{\begin{array}{@{}l@{}} \forall s, t, f, a. \\ \quad \operatorname{is\text{-}metric\text{-}space}(s) \wedge \\ \quad \operatorname{is\text{-}metric\text{-}space}(t) \wedge \\ \quad (f \in (\operatorname{pts}(s) \to \operatorname{pts}(t))) \wedge \\ \quad (a \in \operatorname{pts}(s)) \wedge \\ \quad \forall q \in \operatorname{pts}(s)^{\mathbb{N}}. \\ \quad \quad \operatorname{converges\text{-}to}(s, q, a) \Rightarrow  \\ \quad \quad \quad \operatorname{converges\text{-}to}(t, \operatorname{compose}(f, q), f(a)) \Rightarrow \\ \quad \quad \operatorname{is\text{-}continuous\text{-}at}(s, t, f, a) \end{array}}}
\vnbentry{\texttt{continuous-\allowbreak{}at-\allowbreak{}iff-\allowbreak{}sequential}}{\vnbstmt{\begin{array}{@{}l@{}} \forall s, t, f, a. \\ \quad \operatorname{is\text{-}metric\text{-}space}(s) \wedge \\ \quad \operatorname{is\text{-}metric\text{-}space}(t) \wedge \\ \quad (f \in (\operatorname{pts}(s) \to \operatorname{pts}(t))) \wedge \\ \quad (a \in \operatorname{pts}(s)) \Rightarrow \\ \quad \quad (\operatorname{is\text{-}continuous\text{-}at}(s, t, f, a) \Leftrightarrow \forall q \in \operatorname{pts}(s)^{\mathbb{N}}.\; (\operatorname{converges\text{-}to}(s, q, a) \Rightarrow \operatorname{converges\text{-}to}(t, \operatorname{compose}(f, q), f(a)))) \end{array}}}

\vnbfile{theorem-library/recip-succ-null.scm}{theorem-library/recip-succ-null.scm}
\noindent{\small recip-succ-null.scm -- 1/(n+1) -$>$ 0, PROVEN.}\par
\vnbentry{\texttt{rr-\allowbreak{}recip-\allowbreak{}succ-\allowbreak{}converges-\allowbreak{}to-\allowbreak{}zero}}{\vnbstmt{\operatorname{converges\text{-}to}(\operatorname{rr\text{-}ms}, (\lambda k \in \mathbb{N}.\; \operatorname{recip}((k + 1))), 0)}}

\vnbfile{theorem-library/product-convergence.scm}{theorem-library/product-convergence.scm}
\noindent{\small product-convergence.scm -- CONVERGENCE IN THE COUNTABLE PRODUCT IS EXACTLY COORDINATEWISE CONVERGENCE, proven both ways, retiring the asserted `product-convergence-coordinatewise' of structure-library/product-metric.scm -- the defining property of the PRODUCT TOPOLOGY, and the rung every later}\par
\vnbentry{\texttt{rr-\allowbreak{}null-\allowbreak{}squeeze}}{\vnbstmt{\begin{array}{@{}l@{}} \forall f, h \in (\mathbb{N} \to \mathbb{R}). \\ \quad \operatorname{converges\text{-}to}(\operatorname{rr\text{-}ms}, f, 0) \wedge \\ \quad \forall j \in \mathbb{N}.\; ((0 \leq h(j)) \wedge (h(j) \leq f(j))) \Rightarrow \\ \quad \quad \operatorname{converges\text{-}to}(\operatorname{rr\text{-}ms}, h, 0) \end{array}}}
\vnbentry{\texttt{converges-\allowbreak{}to-\allowbreak{}transfer}}{\vnbstmt{\begin{array}{@{}l@{}} \forall t.; \forall f, h \in (\mathbb{N} \to \operatorname{pts}(t)), v \in \operatorname{pts}(t). \\ \quad \operatorname{is\text{-}metric\text{-}space}(t) \wedge \\ \quad \operatorname{converges\text{-}to}(t, f, v) \wedge \\ \quad \forall j \in \mathbb{N}.\; (h(j) = f(j)) \Rightarrow \\ \quad \quad \operatorname{converges\text{-}to}(t, h, v) \end{array}}}
\vnbentry{\texttt{product-\allowbreak{}coord-\allowbreak{}weighted-\allowbreak{}null}}{\vnbstmt{\begin{array}{@{}l@{}} \forall s, w.; \forall q \in (\mathbb{N} \to \operatorname{product\text{-}carrier}(s)), v \in \operatorname{product\text{-}carrier}(s), n \in \mathbb{N}. \\ \quad \operatorname{is\text{-}ms\text{-}sequence}(s) \wedge \\ \quad \operatorname{summable\text{-}weight}(w) \wedge \\ \quad \operatorname{converges\text{-}to}(s(n), (\lambda k \in \mathbb{N}.\; q(k)(n)), v(n)) \Rightarrow \\ \quad \quad \operatorname{converges\text{-}to}(\operatorname{rr\text{-}ms}, \\ \quad \quad \quad (\lambda j \in \mathbb{N}.\; \\ \quad \quad \quad \quad (w(n) \cdot \operatorname{dist}(\operatorname{bdd\text{-}metric}(s(n)))(q(j)(n), v(n)))), \\ \quad \quad \quad 0) \end{array}}}
\vnbentry{\texttt{product-\allowbreak{}coord-\allowbreak{}from-\allowbreak{}weighted-\allowbreak{}null}}{\vnbstmt{\begin{array}{@{}l@{}} \forall s, w.; \forall q \in (\mathbb{N} \to \operatorname{product\text{-}carrier}(s)), v \in \operatorname{product\text{-}carrier}(s), n \in \mathbb{N}. \\ \quad \operatorname{is\text{-}ms\text{-}sequence}(s) \wedge \\ \quad \operatorname{summable\text{-}weight}(w) \wedge \\ \quad \operatorname{converges\text{-}to}(\operatorname{rr\text{-}ms}, \\ \quad \quad (\lambda j \in \mathbb{N}.\; \\ \quad \quad \quad (w(n) \cdot \operatorname{dist}(\operatorname{bdd\text{-}metric}(s(n)))(q(j)(n), v(n)))), \\ \quad \quad 0) \Rightarrow \\ \quad \quad \operatorname{converges\text{-}to}(s(n), (\lambda k \in \mathbb{N}.\; q(k)(n)), v(n)) \end{array}}}
\vnbentry{\texttt{product-\allowbreak{}convergence-\allowbreak{}coordinatewise-\allowbreak{}bwd}}{\vnbstmt{\begin{array}{@{}l@{}} \forall s, w.; \forall q \in (\mathbb{N} \to \operatorname{product\text{-}carrier}(s)), v \in \operatorname{product\text{-}carrier}(s). \\ \quad \operatorname{is\text{-}ms\text{-}sequence}(s) \wedge \\ \quad \operatorname{summable\text{-}weight}(w) \wedge \\ \quad \forall n \in \mathbb{N}. \\ \quad \quad \operatorname{converges\text{-}to}(s(n), (\lambda k \in \mathbb{N}.\; q(k)(n)), v(n)) \Rightarrow \\ \quad \quad \operatorname{converges\text{-}to}(\operatorname{product\text{-}metric\text{-}w}(s, w), q, v) \end{array}}}
\vnbentry{\texttt{product-\allowbreak{}convergence-\allowbreak{}coordinatewise-\allowbreak{}fwd}}{\vnbstmt{\begin{array}{@{}l@{}} \forall s, w.; \forall q \in (\mathbb{N} \to \operatorname{product\text{-}carrier}(s)), v \in \operatorname{product\text{-}carrier}(s), n \in \mathbb{N}. \\ \quad \operatorname{is\text{-}ms\text{-}sequence}(s) \wedge \\ \quad \operatorname{summable\text{-}weight}(w) \wedge \\ \quad \operatorname{converges\text{-}to}(\operatorname{product\text{-}metric\text{-}w}(s, w), q, v) \Rightarrow \\ \quad \quad \operatorname{converges\text{-}to}(s(n), (\lambda k \in \mathbb{N}.\; q(k)(n)), v(n)) \end{array}}}
\vnbentry{\texttt{product-\allowbreak{}convergence-\allowbreak{}coordinatewise}}{\vnbstmt{\begin{array}{@{}l@{}} \forall s, w.; \forall q \in (\mathbb{N} \to \operatorname{product\text{-}carrier}(s)), v \in \operatorname{product\text{-}carrier}(s). \\ \quad (\operatorname{is\text{-}ms\text{-}sequence}(s) \wedge \operatorname{summable\text{-}weight}(w)) \Rightarrow  \\ \quad \quad (\operatorname{converges\text{-}to}(\operatorname{product\text{-}metric\text{-}w}(s, w), q, v) \Leftrightarrow \forall n \in \mathbb{N}.\; \operatorname{converges\text{-}to}(s(n), (\lambda k \in \mathbb{N}.\; q(k)(n)), v(n))) \end{array}}}

\vnbfile{theorem-library/rake-diagonal-subseq.scm}{theorem-library/rake-diagonal-subseq.scm}
\noindent{\small theorem-library/rake-diagonal-subseq.scm ==================================================================== The cross-coordinate DIAGONAL argument, and the sequential Tychonoff assembly on top of it.  Two supports of theorem-library/seq-compact-product}\par
\vnbentry{\texttt{coordinatewise-\allowbreak{}diagonal-\allowbreak{}subseq}}{\vnbstmt{\begin{array}{@{}l@{}} \forall s.; \forall q \in (\mathbb{N} \to \operatorname{product\text{-}carrier}(s)). \\ \quad \operatorname{is\text{-}ms\text{-}sequence}(s) \wedge \\ \quad \forall n \in \mathbb{N}.\; \operatorname{seq\text{-}compact}(s(n)) \Rightarrow \\ \quad \quad \exists a. \\ \quad \quad \quad \operatorname{strictly\text{-}mono\text{-}nn}(a) \wedge \\ \quad \quad \quad \exists l \in \operatorname{product\text{-}carrier}(s). \forall n \in \mathbb{N}. \\ \quad \quad \quad \quad \operatorname{converges\text{-}to}(s(n), (\lambda k \in \mathbb{N}.\; \operatorname{subseq}(q, a)(k)(n)), l(n)) \end{array}}}
\vnbentry{\texttt{seq-\allowbreak{}compact-\allowbreak{}countable-\allowbreak{}product}}{\vnbstmt{\begin{array}{@{}l@{}} \forall s. \\ \quad \operatorname{is\text{-}ms\text{-}sequence}(s) \wedge \\ \quad \forall n \in \mathbb{N}.\; \operatorname{seq\text{-}compact}(s(n)) \Rightarrow \\ \quad \quad \operatorname{seq\text{-}compact}(\operatorname{product\text{-}metric}(s)) \end{array}}}

\vnbfile{theorem-library/tychonoff-proof.scm}{theorem-library/tychonoff-proof.scm}
\noindent{\small theorem-library/tychonoff-proof.scm ==================================================================== compact-countable-product (countable Tychonoff headline) -- PROVEN.}\par
\vnbentry{\texttt{compact-\allowbreak{}countable-\allowbreak{}product}}{\vnbstmt{\begin{array}{@{}l@{}} \forall s. \\ \quad \operatorname{is\text{-}ms\text{-}sequence}(s) \wedge \\ \quad \forall n \in \mathbb{N}.\; \operatorname{is\text{-}compact}(s(n)) \Rightarrow \\ \quad \quad \operatorname{is\text{-}compact}(\operatorname{product\text{-}metric}(s)) \end{array}}}

\vnbfile{theorem-library/rake-setoid2.scm}{theorem-library/rake-setoid2.scm}
\noindent{\small rake-setoid2.scm -- BATCH O of the 2026-09-17 rake: the SETOID quotient and the METRIC COMPLETION.  Twelve of the thirteen asserted leaves of the batch are PROVEN here; the thirteenth is a STATEMENT DEFECT and is recorded below. Every statement is its support's statement UNCHANGED (copied from the}\par
\vnbentry{\texttt{rko2-\allowbreak{}related-\allowbreak{}unfold}}{\vnbstmt{\begin{array}{@{}l@{}} \forall s, a, b. \\ \quad (\operatorname{related}(s, a, b) \Leftrightarrow (\langle a, b \rangle \in \operatorname{rel}(s))) \end{array}}}
\vnbentry{\texttt{rko2-\allowbreak{}refl}}{\vnbstmt{\begin{array}{@{}l@{}} \forall s.; \forall a \in \operatorname{pts}(s). \\ \quad (\operatorname{is\text{-}setoid}(s) \Rightarrow \operatorname{related}(s, a, a)) \end{array}}}
\vnbentry{\texttt{rko2-\allowbreak{}sym}}{\vnbstmt{\begin{array}{@{}l@{}} \forall s, a, b. \\ \quad \operatorname{is\text{-}setoid}(s) \wedge \\ \quad (a \in \operatorname{pts}(s)) \wedge \\ \quad (b \in \operatorname{pts}(s)) \wedge \\ \quad \operatorname{related}(s, a, b) \Rightarrow \\ \quad \quad \operatorname{related}(s, b, a) \end{array}}}
\vnbentry{\texttt{rko2-\allowbreak{}trans}}{\vnbstmt{\begin{array}{@{}l@{}} \forall s, a, b, c. \\ \quad \operatorname{is\text{-}setoid}(s) \wedge \\ \quad (a \in \operatorname{pts}(s)) \wedge \\ \quad (b \in \operatorname{pts}(s)) \wedge \\ \quad (c \in \operatorname{pts}(s)) \wedge \\ \quad \operatorname{related}(s, a, b) \wedge \\ \quad \operatorname{related}(s, b, c) \Rightarrow \\ \quad \quad \operatorname{related}(s, a, c) \end{array}}}
\vnbentry{\texttt{class-\allowbreak{}self}}{\vnbstmt{\begin{array}{@{}l@{}} \forall s.; \forall a \in \operatorname{pts}(s). \\ \quad \operatorname{is\text{-}setoid}(s) \Rightarrow  \\ \quad \quad (a \in \operatorname{class}(s, a)) \end{array}}}
\vnbentry{\texttt{class-\allowbreak{}in-\allowbreak{}quotient}}{\vnbstmt{\begin{array}{@{}l@{}} \forall s.; \forall a \in \operatorname{pts}(s). \\ \quad \operatorname{is\text{-}setoid}(s) \Rightarrow  \\ \quad \quad (\operatorname{class}(s, a) \in \operatorname{quotient}(s)) \end{array}}}
\vnbentry{\texttt{quotient-\allowbreak{}rep}}{\vnbstmt{\begin{array}{@{}l@{}} \forall s.; \forall x \in \operatorname{quotient}(s). \\ \quad \operatorname{is\text{-}setoid}(s) \Rightarrow  \\ \quad \quad \exists a \in \operatorname{pts}(s).\; (x = \operatorname{class}(s, a)) \end{array}}}
\vnbentry{\texttt{class-\allowbreak{}subset-\allowbreak{}carrier}}{\vnbstmt{\begin{array}{@{}l@{}} \forall s.; \forall a \in \operatorname{pts}(s). \\ \quad \operatorname{is\text{-}setoid}(s) \Rightarrow  \\ \quad \quad (\operatorname{class}(s, a) \subseteq \operatorname{pts}(s)) \end{array}}}
\vnbentry{\texttt{class-\allowbreak{}eq-\allowbreak{}iff}}{\vnbstmt{\begin{array}{@{}l@{}} \forall s.; \forall a, b \in \operatorname{pts}(s). \\ \quad \operatorname{is\text{-}setoid}(s) \Rightarrow  \\ \quad \quad ((\operatorname{class}(s, a) = \operatorname{class}(s, b)) \Leftrightarrow \operatorname{related}(s, a, b)) \end{array}}}
\vnbentry{\texttt{class-\allowbreak{}disjoint}}{\vnbstmt{\begin{array}{@{}l@{}} \forall s.; \forall a, b \in \operatorname{pts}(s). \\ \quad \operatorname{is\text{-}setoid}(s) \Rightarrow  \\ \quad \quad (\operatorname{class}(s, a) = \operatorname{class}(s, b)) \vee \\ \quad \quad ((\operatorname{class}(s, a) \cap \operatorname{class}(s, b)) = \emptyset) \end{array}}}
\vnbentry{\texttt{descend-\allowbreak{}computes}}{\vnbstmt{\begin{array}{@{}l@{}} \forall s, z.; \forall f \in (\operatorname{pts}(s) \to z), a \in \operatorname{pts}(s). \\ \quad (\operatorname{is\text{-}setoid}(s) \wedge \operatorname{respects}(s, f)) \Rightarrow  \\ \quad \quad (\operatorname{descend}(s, f)(\operatorname{class}(s, a)) = f(a)) \end{array}}}
\vnbentry{\texttt{descend2-\allowbreak{}computes}}{\vnbstmt{\begin{array}{@{}l@{}} \forall s, z.; \forall f \in ((\operatorname{pts}(s) \times \operatorname{pts}(s)) \to z), a, b. \\ \quad \operatorname{is\text{-}setoid}(s) \wedge \\ \quad \operatorname{respects2}(s, f) \wedge \\ \quad (a \in \operatorname{pts}(s)) \wedge \\ \quad (b \in \operatorname{pts}(s)) \Rightarrow \\ \quad \quad (\operatorname{descend2}(s, f)(\operatorname{class}(s, a), \operatorname{class}(s, b)) = f(a, b)) \end{array}}}
\vnbentry{\texttt{quotient-\allowbreak{}universal}}{\vnbstmt{\begin{array}{@{}l@{}} \forall s, z.; \forall f \in (\operatorname{pts}(s) \to z). \\ \quad (\operatorname{is\text{-}setoid}(s) \wedge \operatorname{respects}(s, f)) \Rightarrow  \\ \quad \quad \exists g \in (\operatorname{quotient}(s) \to z). \\ \quad \quad \quad \forall a \in \operatorname{pts}(s).\; (g(\operatorname{class}(s, a)) = f(a)) \wedge \\ \quad \quad \quad \forall b \in (\operatorname{quotient}(s) \to z). \\ \quad \quad \quad \quad \forall a \in \operatorname{pts}(s).\; (b(\operatorname{class}(s, a)) = f(a)) \Rightarrow  \\ \quad \quad \quad \quad \quad (b = g) \end{array}}}
\vnbentry{\texttt{rko2-\allowbreak{}cseq-\allowbreak{}unfold}}{\vnbstmt{\begin{array}{@{}l@{}} \forall m. \\ \quad (\operatorname{cseq}(m) \simeq \{f \in (\mathbb{N} \to \operatorname{pts}(m)) : \operatorname{is\text{-}cauchy\text{-}seq}(m, f)\}) \end{array}}}
\vnbentry{\texttt{rko2-\allowbreak{}cseq-\allowbreak{}in-\allowbreak{}fun}}{\vnbstmt{\begin{array}{@{}l@{}} \forall m.; \forall f \in \operatorname{cseq}(m). \\ \quad (f \in (\mathbb{N} \to \operatorname{pts}(m))) \end{array}}}
\vnbentry{\texttt{rko2-\allowbreak{}completion-\allowbreak{}dist}}{\vnbstmt{\begin{array}{@{}l@{}} \forall m. \\ \quad (\operatorname{dist}(\operatorname{completion}(m)) \simeq \operatorname{completion\text{-}dist}(m)) \end{array}}}
\vnbentry{\texttt{rko2-\allowbreak{}cauchy-\allowbreak{}setoid-\allowbreak{}rel}}{\vnbstmt{\begin{array}{@{}l@{}} \forall m. \\ \quad (\operatorname{rel}(\operatorname{cauchy\text{-}setoid}(m)) \simeq \operatorname{crel}(m)) \end{array}}}
\vnbentry{\texttt{rko2-\allowbreak{}crel-\allowbreak{}unfold}}{\vnbstmt{\begin{array}{@{}l@{}} \forall m. \\ \quad (\operatorname{crel}(m) \simeq \{p \in (\operatorname{cseq}(m) \times \operatorname{cseq}(m)) : \operatorname{cseq\text{-}equiv}(m, \operatorname{nth}(1, p), \operatorname{nth}(2, p))\}) \end{array}}}
\vnbentry{\texttt{rko2-\allowbreak{}crel-\allowbreak{}intro}}{\vnbstmt{\begin{array}{@{}l@{}} \forall m, p, q. \\ \quad (p \in \operatorname{cseq}(m)) \wedge \\ \quad (q \in \operatorname{cseq}(m)) \wedge \\ \quad \operatorname{cseq\text{-}equiv}(m, p, q) \Rightarrow \\ \quad \quad (\langle p, q \rangle \in \operatorname{crel}(m)) \end{array}}}
\vnbentry{\texttt{rko2-\allowbreak{}crel-\allowbreak{}elim}}{\vnbstmt{\begin{array}{@{}l@{}} \forall m, p, q. \\ \quad (\langle p, q \rangle \in \operatorname{crel}(m)) \Rightarrow  \\ \quad \quad \operatorname{cseq\text{-}equiv}(m, p, q) \end{array}}}
\vnbentry{\texttt{rko2-\allowbreak{}cauchy-\allowbreak{}setoid-\allowbreak{}length}}{\vnbstmt{\forall m.\; (\operatorname{length}(\operatorname{cauchy\text{-}setoid}(m)) = 2)}}
\vnbentry{\texttt{rko2-\allowbreak{}crel-\allowbreak{}is-\allowbreak{}set}}{\vnbstmt{\begin{array}{@{}l@{}} \forall m. \\ \quad \operatorname{is\text{-}metric\text{-}space}(m) \Rightarrow  \\ \quad \quad (\operatorname{crel}(m) \in \mathbf{Set}) \end{array}}}
\vnbentry{\texttt{rko2-\allowbreak{}dist-\allowbreak{}seq-\allowbreak{}unfold}}{\vnbstmt{\begin{array}{@{}l@{}} \forall m, f, g. \\ \quad (\operatorname{dist\text{-}seq}(m, f, g) \simeq (\lambda n \in \mathbb{N}.\; \operatorname{dist}(m)(f(n), g(n)))) \end{array}}}
\vnbentry{\texttt{rko2-\allowbreak{}dist-\allowbreak{}seq-\allowbreak{}type}}{\vnbstmt{\begin{array}{@{}l@{}} \forall m, f, g. \\ \quad \operatorname{is\text{-}metric\text{-}space}(m) \wedge \\ \quad (f \in (\mathbb{N} \to \operatorname{pts}(m))) \wedge \\ \quad (g \in (\mathbb{N} \to \operatorname{pts}(m))) \Rightarrow \\ \quad \quad (\operatorname{dist\text{-}seq}(m, f, g) \in (\mathbb{N} \to \mathbb{R})) \end{array}}}
\vnbentry{\texttt{rko2-\allowbreak{}dist-\allowbreak{}seq-\allowbreak{}at}}{\vnbstmt{\begin{array}{@{}l@{}} \forall m, f, g.; \forall n \in \mathbb{N}. \\ \quad (\operatorname{dist\text{-}seq}(m, f, g)(n) \simeq \operatorname{dist}(m)(f(n), g(n))) \end{array}}}
\vnbentry{\texttt{rko2-\allowbreak{}null-\allowbreak{}refl}}{\vnbstmt{\begin{array}{@{}l@{}} \forall m, f. \\ \quad \operatorname{is\text{-}metric\text{-}space}(m) \wedge \\ \quad (f \in (\mathbb{N} \to \operatorname{pts}(m))) \Rightarrow \\ \quad \quad \operatorname{converges\text{-}to}(\operatorname{rr\text{-}ms}, \operatorname{dist\text{-}seq}(m, f, f), 0) \end{array}}}
\vnbentry{\texttt{rko2-\allowbreak{}null-\allowbreak{}sym}}{\vnbstmt{\begin{array}{@{}l@{}} \forall m, f, g. \\ \quad \operatorname{is\text{-}metric\text{-}space}(m) \wedge \\ \quad (f \in (\mathbb{N} \to \operatorname{pts}(m))) \wedge \\ \quad (g \in (\mathbb{N} \to \operatorname{pts}(m))) \wedge \\ \quad \operatorname{converges\text{-}to}(\operatorname{rr\text{-}ms}, \operatorname{dist\text{-}seq}(m, f, g), 0) \Rightarrow \\ \quad \quad \operatorname{converges\text{-}to}(\operatorname{rr\text{-}ms}, \operatorname{dist\text{-}seq}(m, g, f), 0) \end{array}}}
\vnbentry{\texttt{rko2-\allowbreak{}null-\allowbreak{}trans}}{\vnbstmt{\begin{array}{@{}l@{}} \forall m, f, g, h. \\ \quad \operatorname{is\text{-}metric\text{-}space}(m) \wedge \\ \quad (f \in (\mathbb{N} \to \operatorname{pts}(m))) \wedge \\ \quad (g \in (\mathbb{N} \to \operatorname{pts}(m))) \wedge \\ \quad (h \in (\mathbb{N} \to \operatorname{pts}(m))) \wedge \\ \quad \operatorname{converges\text{-}to}(\operatorname{rr\text{-}ms}, \operatorname{dist\text{-}seq}(m, f, g), 0) \wedge \\ \quad \operatorname{converges\text{-}to}(\operatorname{rr\text{-}ms}, \operatorname{dist\text{-}seq}(m, g, h), 0) \Rightarrow \\ \quad \quad \operatorname{converges\text{-}to}(\operatorname{rr\text{-}ms}, \operatorname{dist\text{-}seq}(m, f, h), 0) \end{array}}}
\vnbentry{\texttt{cauchy-\allowbreak{}setoid-\allowbreak{}is-\allowbreak{}setoid}}{\vnbstmt{\begin{array}{@{}l@{}} \forall m. \\ \quad \operatorname{is\text{-}metric\text{-}space}(m) \Rightarrow  \\ \quad \quad \operatorname{is\text{-}setoid}(\operatorname{cauchy\text{-}setoid}(m)) \end{array}}}
\vnbentry{\texttt{rko2-\allowbreak{}quad}}{\vnbstmt{\begin{array}{@{}l@{}} \forall m, a, b, c, e. \\ \quad \operatorname{is\text{-}metric\text{-}space}(m) \wedge \\ \quad (a \in \operatorname{pts}(m)) \wedge \\ \quad (b \in \operatorname{pts}(m)) \wedge \\ \quad (c \in \operatorname{pts}(m)) \wedge \\ \quad (e \in \operatorname{pts}(m)) \Rightarrow \\ \quad \quad \lvert (\operatorname{dist}(m)(a, b) - \operatorname{dist}(m)(c, e)) \rvert \\ \quad \quad \leq\; (\operatorname{dist}(m)(a, c) + \operatorname{dist}(m)(b, e)) \end{array}}}
\vnbentry{\texttt{rko2-\allowbreak{}dist-\allowbreak{}converges}}{\vnbstmt{\begin{array}{@{}l@{}} \forall m, f, g, a, b. \\ \quad \operatorname{is\text{-}metric\text{-}space}(m) \wedge \\ \quad (f \in (\mathbb{N} \to \operatorname{pts}(m))) \wedge \\ \quad (g \in (\mathbb{N} \to \operatorname{pts}(m))) \wedge \\ \quad (a \in \operatorname{pts}(m)) \wedge \\ \quad (b \in \operatorname{pts}(m)) \wedge \\ \quad \operatorname{converges\text{-}to}(m, f, a) \wedge \\ \quad \operatorname{converges\text{-}to}(m, g, b) \Rightarrow \\ \quad \quad \operatorname{converges\text{-}to}(\operatorname{rr\text{-}ms}, \operatorname{dist\text{-}seq}(m, f, g), \operatorname{dist}(m)(a, b)) \end{array}}}
\vnbentry{\texttt{rko2-\allowbreak{}const-\allowbreak{}seq-\allowbreak{}converges}}{\vnbstmt{\begin{array}{@{}l@{}} \forall m, u. \\ \quad (\operatorname{is\text{-}metric\text{-}space}(m) \wedge (u \in \operatorname{pts}(m))) \Rightarrow  \\ \quad \quad \operatorname{converges\text{-}to}(m, \operatorname{embed\text{-}seq}(m, u), u) \end{array}}}
\vnbentry{\texttt{rko2-\allowbreak{}class-\allowbreak{}eq-\allowbreak{}null}}{\vnbstmt{\begin{array}{@{}l@{}} \forall m, p, q. \\ \quad \operatorname{is\text{-}metric\text{-}space}(m) \wedge \\ \quad (p \in \operatorname{cseq}(m)) \wedge \\ \quad (q \in \operatorname{cseq}(m)) \wedge \\ \quad \operatorname{class}(\operatorname{cauchy\text{-}setoid}(m), p) \\ \quad =\; \operatorname{class}(\operatorname{cauchy\text{-}setoid}(m), q) \Rightarrow \\ \quad \quad \operatorname{converges\text{-}to}(\operatorname{rr\text{-}ms}, \operatorname{dist\text{-}seq}(m, p, q), 0) \end{array}}}
\vnbentry{\texttt{rko2-\allowbreak{}const-\allowbreak{}equiv-\allowbreak{}converges}}{\vnbstmt{\begin{array}{@{}l@{}} \forall m, u, f. \\ \quad \operatorname{is\text{-}metric\text{-}space}(m) \wedge \\ \quad (u \in \operatorname{pts}(m)) \wedge \\ \quad (f \in (\mathbb{N} \to \operatorname{pts}(m))) \wedge \\ \quad \operatorname{converges\text{-}to}(\operatorname{rr\text{-}ms}, \operatorname{dist\text{-}seq}(m, \operatorname{embed\text{-}seq}(m, u), f), 0) \Rightarrow \\ \quad \quad \operatorname{converges\text{-}to}(m, f, u) \end{array}}}
\vnbentry{\texttt{embed-\allowbreak{}isometry}}{\vnbstmt{\begin{array}{@{}l@{}} \forall m.; \forall u, v \in \operatorname{pts}(m). \\ \quad \operatorname{is\text{-}metric\text{-}space}(m) \Rightarrow  \\ \quad \quad \operatorname{dist}(\operatorname{completion}(m))(\operatorname{embed}(m)(u), \operatorname{embed}(m)(v)) \\ \quad \quad =\; \operatorname{dist}(m)(u, v) \end{array}}}

\vnbfile{theorem-library/rake-completion-ms.scm}{theorem-library/rake-completion-ms.scm}
\noindent{\small rake-completion-ms.scm -- BATCH 5c-Q of the 2026-09-18 rake: completion-is-metric-space (structure-library/metric-completion.scm:107, asserted `well-known').  PROVEN `modulo 0'.  The statement is the support's, copied from the support site and unchanged:}\par
\vnbentry{\texttt{r7q-\allowbreak{}cseq-\allowbreak{}cauchy}}{\vnbstmt{\begin{array}{@{}l@{}} \forall m.; \forall f \in \operatorname{cseq}(m). \\ \quad \operatorname{is\text{-}cauchy\text{-}seq}(m, f) \end{array}}}
\vnbentry{\texttt{r7q-\allowbreak{}dist-\allowbreak{}seq-\allowbreak{}cauchy}}{\vnbstmt{\begin{array}{@{}l@{}} \forall m.; \forall f, g \in \operatorname{cseq}(m). \\ \quad \operatorname{is\text{-}metric\text{-}space}(m) \Rightarrow  \\ \quad \quad \operatorname{is\text{-}cauchy\text{-}seq}(\operatorname{rr\text{-}ms}, \operatorname{dist\text{-}seq}(m, f, g)) \end{array}}}
\vnbentry{\texttt{r7q-\allowbreak{}dist-\allowbreak{}seq-\allowbreak{}converges}}{\vnbstmt{\begin{array}{@{}l@{}} \forall m.; \forall f, g \in \operatorname{cseq}(m). \\ \quad \operatorname{is\text{-}metric\text{-}space}(m) \Rightarrow  \\ \quad \quad \operatorname{converges}(\operatorname{rr\text{-}ms}, \operatorname{dist\text{-}seq}(m, f, g)) \end{array}}}
\vnbentry{\texttt{r7q-\allowbreak{}dist-\allowbreak{}seq-\allowbreak{}sym}}{\vnbstmt{\begin{array}{@{}l@{}} \forall m.; \forall f, g \in (\mathbb{N} \to \operatorname{pts}(m)), v. \\ \quad \operatorname{is\text{-}metric\text{-}space}(m) \wedge \\ \quad \operatorname{converges\text{-}to}(\operatorname{rr\text{-}ms}, \operatorname{dist\text{-}seq}(m, f, g), v) \Rightarrow \\ \quad \quad \operatorname{converges\text{-}to}(\operatorname{rr\text{-}ms}, \operatorname{dist\text{-}seq}(m, g, f), v) \end{array}}}
\vnbentry{\texttt{r7q-\allowbreak{}null-\allowbreak{}transfer}}{\vnbstmt{\begin{array}{@{}l@{}} \forall t.; \forall f, h \in (\mathbb{N} \to \operatorname{pts}(t)), v \in \operatorname{pts}(t). \\ \quad \operatorname{is\text{-}metric\text{-}space}(t) \wedge \\ \quad \operatorname{converges\text{-}to}(t, f, v) \wedge \\ \quad \operatorname{converges\text{-}to}(\operatorname{rr\text{-}ms}, \operatorname{dist\text{-}seq}(t, f, h), 0) \Rightarrow \\ \quad \quad \operatorname{converges\text{-}to}(t, h, v) \end{array}}}
\vnbentry{\texttt{r7q-\allowbreak{}quad-\allowbreak{}null}}{\vnbstmt{\begin{array}{@{}l@{}} \forall m.; \forall f, g, p, q \in (\mathbb{N} \to \operatorname{pts}(m)). \\ \quad \operatorname{is\text{-}metric\text{-}space}(m) \wedge \\ \quad \operatorname{converges\text{-}to}(\operatorname{rr\text{-}ms}, \operatorname{dist\text{-}seq}(m, f, p), 0) \wedge \\ \quad \operatorname{converges\text{-}to}(\operatorname{rr\text{-}ms}, \operatorname{dist\text{-}seq}(m, g, q), 0) \Rightarrow \\ \quad \quad \operatorname{converges\text{-}to}(\operatorname{rr\text{-}ms}, \\ \quad \quad \quad \operatorname{dist\text{-}seq}(\operatorname{rr\text{-}ms}, \operatorname{dist\text{-}seq}(m, f, g), \operatorname{dist\text{-}seq}(m, p, q)), \\ \quad \quad \quad 0) \end{array}}}
\vnbentry{\texttt{r7q-\allowbreak{}limit-\allowbreak{}indep}}{\vnbstmt{\begin{array}{@{}l@{}} \forall m.; \forall f, g, p, q \in (\mathbb{N} \to \operatorname{pts}(m)), v. \\ \quad \operatorname{is\text{-}metric\text{-}space}(m) \wedge \\ \quad \operatorname{converges\text{-}to}(\operatorname{rr\text{-}ms}, \operatorname{dist\text{-}seq}(m, f, p), 0) \wedge \\ \quad \operatorname{converges\text{-}to}(\operatorname{rr\text{-}ms}, \operatorname{dist\text{-}seq}(m, g, q), 0) \wedge \\ \quad \operatorname{converges\text{-}to}(\operatorname{rr\text{-}ms}, \operatorname{dist\text{-}seq}(m, f, g), v) \Rightarrow \\ \quad \quad \operatorname{converges\text{-}to}(\operatorname{rr\text{-}ms}, \operatorname{dist\text{-}seq}(m, p, q), v) \end{array}}}
\vnbentry{\texttt{r7q-\allowbreak{}iota-\allowbreak{}dhat}}{\vnbstmt{\begin{array}{@{}l@{}} \forall m, x, y. \\ \quad (x \in \operatorname{quotient}(\operatorname{cauchy\text{-}setoid}(m))) \wedge \\ \quad (y \in \operatorname{quotient}(\operatorname{cauchy\text{-}setoid}(m))) \Rightarrow \\ \quad \quad (\iota(l) \simeq \operatorname{dist}(\operatorname{completion}(m))(x, y)) \end{array}}}
\vnbentry{\texttt{r7q-\allowbreak{}dhat-\allowbreak{}iota}}{\vnbstmt{\begin{array}{@{}l@{}} \forall m, x, y. \\ \quad (x \in \operatorname{quotient}(\operatorname{cauchy\text{-}setoid}(m))) \wedge \\ \quad (y \in \operatorname{quotient}(\operatorname{cauchy\text{-}setoid}(m))) \Rightarrow \\ \quad \quad (\operatorname{dist}(\operatorname{completion}(m))(x, y) \simeq \iota(l)) \end{array}}}
\vnbentry{\texttt{r7q-\allowbreak{}dist-\allowbreak{}value}}{\vnbstmt{\begin{array}{@{}l@{}} \forall m.; \forall f, g \in \operatorname{cseq}(m), v. \\ \quad \operatorname{is\text{-}metric\text{-}space}(m) \wedge \\ \quad \operatorname{converges\text{-}to}(\operatorname{rr\text{-}ms}, \operatorname{dist\text{-}seq}(m, f, g), v) \Rightarrow \\ \quad \quad \operatorname{dist}(\operatorname{completion}(m))(\operatorname{class}(\operatorname{cauchy\text{-}setoid}(m), f), \operatorname{class}(\operatorname{cauchy\text{-}setoid}(m), g)) \\ \quad \quad =\; v \end{array}}}
\vnbentry{\texttt{r7q-\allowbreak{}rr-\allowbreak{}pts}}{\vnbstmt{(\mathbb{R} \simeq \operatorname{pts}(\operatorname{rr\text{-}ms}))}}
\vnbentry{\texttt{r7q-\allowbreak{}dhat-\allowbreak{}in-\allowbreak{}rr}}{\vnbstmt{\begin{array}{@{}l@{}} \forall m.; \forall x, y \in \operatorname{quotient}(\operatorname{cauchy\text{-}setoid}(m)). \\ \quad \operatorname{is\text{-}metric\text{-}space}(m) \Rightarrow  \\ \quad \quad (\operatorname{dist}(\operatorname{completion}(m))(x, y) \in \mathbb{R}) \end{array}}}
\vnbentry{\texttt{r7q-\allowbreak{}dhat-\allowbreak{}self-\allowbreak{}zero}}{\vnbstmt{\begin{array}{@{}l@{}} \forall m.; \forall x \in \operatorname{quotient}(\operatorname{cauchy\text{-}setoid}(m)). \\ \quad \operatorname{is\text{-}metric\text{-}space}(m) \Rightarrow  \\ \quad \quad (\operatorname{dist}(\operatorname{completion}(m))(x, x) = 0) \end{array}}}
\vnbentry{\texttt{r7q-\allowbreak{}dhat-\allowbreak{}nonneg}}{\vnbstmt{\begin{array}{@{}l@{}} \forall m.; \forall x, y \in \operatorname{quotient}(\operatorname{cauchy\text{-}setoid}(m)). \\ \quad \operatorname{is\text{-}metric\text{-}space}(m) \Rightarrow  \\ \quad \quad (0 \leq \operatorname{dist}(\operatorname{completion}(m))(x, y)) \end{array}}}
\vnbentry{\texttt{r7q-\allowbreak{}dhat-\allowbreak{}zero-\allowbreak{}eq}}{\vnbstmt{\begin{array}{@{}l@{}} \forall m.; \forall x, y \in \operatorname{quotient}(\operatorname{cauchy\text{-}setoid}(m)). \\ \quad \operatorname{is\text{-}metric\text{-}space}(m) \wedge \\ \quad (\operatorname{dist}(\operatorname{completion}(m))(x, y) = 0) \Rightarrow \\ \quad \quad (x = y) \end{array}}}
\vnbentry{\texttt{r7q-\allowbreak{}dhat-\allowbreak{}sym}}{\vnbstmt{\begin{array}{@{}l@{}} \forall m.; \forall x, y \in \operatorname{quotient}(\operatorname{cauchy\text{-}setoid}(m)). \\ \quad \operatorname{is\text{-}metric\text{-}space}(m) \Rightarrow  \\ \quad \quad \operatorname{dist}(\operatorname{completion}(m))(x, y) \\ \quad \quad =\; \operatorname{dist}(\operatorname{completion}(m))(y, x) \end{array}}}
\vnbentry{\texttt{r7q-\allowbreak{}dhat-\allowbreak{}triangle}}{\vnbstmt{\begin{array}{@{}l@{}} \forall m.; \forall x, y, z \in \operatorname{quotient}(\operatorname{cauchy\text{-}setoid}(m)). \\ \quad \operatorname{is\text{-}metric\text{-}space}(m) \Rightarrow  \\ \quad \quad \operatorname{dist}(\operatorname{completion}(m))(x, z) \\ \quad \quad \leq\; (\operatorname{dist}(\operatorname{completion}(m))(x, y) + \operatorname{dist}(\operatorname{completion}(m))(y, z)) \end{array}}}
\vnbentry{\texttt{completion-\allowbreak{}is-\allowbreak{}metric-\allowbreak{}space}}{\vnbstmt{\begin{array}{@{}l@{}} \forall m. \\ \quad \operatorname{is\text{-}metric\text{-}space}(m) \Rightarrow  \\ \quad \quad \operatorname{is\text{-}metric\text{-}space}(\operatorname{completion}(m)) \end{array}}}

\vnbfile{theorem-library/rake-completion-complete.scm}{theorem-library/rake-completion-complete.scm}
\noindent{\small rake-completion-complete.scm -- BATCH 12-D (2026-09-20): `completion-is-complete' (structure-library/metric-completion.scm:110, asserted `well-known') and the general bricks its diagonal argument needs.}\par
\vnbentry{\texttt{rr-\allowbreak{}limit-\allowbreak{}tail-\allowbreak{}le}}{\vnbstmt{\begin{array}{@{}l@{}} \forall f, v, c. \\ \quad (f \in (\mathbb{N} \to \mathbb{R})) \wedge \\ \quad (v \in \mathbb{R}) \wedge \\ \quad (c \in \mathbb{R}) \wedge \\ \quad \operatorname{converges\text{-}to}(\operatorname{rr\text{-}ms}, f, v) \wedge \\ \quad \operatorname{pos\text{-}rr}(c) \Rightarrow \\ \quad \quad \exists t \in \mathbb{N}. \forall k \in \mathbb{N}. \\ \quad \quad \quad ((t \leq k) \Rightarrow (f(k) \leq (v + c))) \end{array}}}
\vnbentry{\texttt{rr-\allowbreak{}limit-\allowbreak{}le-\allowbreak{}of-\allowbreak{}tail}}{\vnbstmt{\begin{array}{@{}l@{}} \forall f, v, c, t. \\ \quad (f \in (\mathbb{N} \to \mathbb{R})) \wedge \\ \quad (v \in \mathbb{R}) \wedge \\ \quad (c \in \mathbb{R}) \wedge \\ \quad (t \in \mathbb{N}) \wedge \\ \quad \operatorname{converges\text{-}to}(\operatorname{rr\text{-}ms}, f, v) \wedge \\ \quad \forall k \in \mathbb{N}.\; ((t \leq k) \Rightarrow (f(k) \leq c)) \Rightarrow \\ \quad \quad (v \leq c) \end{array}}}
\vnbentry{\texttt{r12d-\allowbreak{}repfam}}{\vnbstmt{\begin{array}{@{}l@{}} \forall m.; \forall f \in (\mathbb{N} \to \operatorname{quotient}(\operatorname{cauchy\text{-}setoid}(m))). \\ \quad \operatorname{is\text{-}metric\text{-}space}(m) \Rightarrow  \\ \quad \quad \exists a \in (\mathbb{N} \to \operatorname{cseq}(m)). \forall p \in \mathbb{N}. \\ \quad \quad \quad (f(p) = \operatorname{class}(\operatorname{cauchy\text{-}setoid}(m), a(p))) \end{array}}}
\vnbentry{\texttt{r12d-\allowbreak{}rapid-\allowbreak{}indices}}{\vnbstmt{\begin{array}{@{}l@{}} \forall m, f, d. \\ \quad \operatorname{is\text{-}metric\text{-}space}(m) \wedge \\ \quad (f \in (\mathbb{N} \to \operatorname{cseq}(m))) \wedge \\ \quad (d \in (\mathbb{N} \to \mathbb{R})) \wedge \\ \quad \forall p \in \mathbb{N}.\; \operatorname{pos\text{-}rr}(d(p)) \Rightarrow \\ \quad \quad \exists a \in (\mathbb{N} \to \mathbb{N}). \forall p \in \mathbb{N}, i, j. \\ \quad \quad \quad ((i \in \mathbb{N}) \wedge ((j \in \mathbb{N}) \wedge ((a(p) \leq i) \wedge (a(p) \leq j)))) \Rightarrow  \\ \quad \quad \quad \quad (\operatorname{dist}(m)(f(p)(i), f(p)(j)) \leq d(p)) \end{array}}}
\vnbentry{\texttt{r12d-\allowbreak{}key}}{\vnbstmt{\begin{array}{@{}l@{}} \forall m, f, a, p, q, b, c. \\ \quad \operatorname{is\text{-}metric\text{-}space}(m) \wedge \\ \quad (f \in (\mathbb{N} \to \operatorname{quotient}(\operatorname{cauchy\text{-}setoid}(m)))) \wedge \\ \quad (a \in (\mathbb{N} \to \operatorname{cseq}(m))) \wedge \\ \quad \forall p \in \mathbb{N}. \\ \quad \quad (f(p) = \operatorname{class}(\operatorname{cauchy\text{-}setoid}(m), a(p))) \wedge \\ \quad (p \in \mathbb{N}) \wedge \\ \quad (q \in \mathbb{N}) \wedge \\ \quad (b \in \mathbb{N}) \wedge \\ \quad (c \in \mathbb{R}) \wedge \\ \quad \operatorname{pos\text{-}rr}(c) \Rightarrow \\ \quad \quad \exists n \in \mathbb{N}. \\ \quad \quad \quad (b \leq n) \wedge \\ \quad \quad \quad \operatorname{dist}(m)(a(p)(n), a(q)(n)) \\ \quad \quad \quad \leq\; (\operatorname{dist}(\operatorname{completion}(m))(f(p), f(q)) + c) \end{array}}}
\vnbentry{\texttt{r12d-\allowbreak{}diag-\allowbreak{}at}}{\vnbstmt{\begin{array}{@{}l@{}} \forall f, m.; \forall p \in \mathbb{N}. \\ \quad ((\lambda v \in \mathbb{N}.\; f(v)(m(v)))(p) \simeq f(p)(m(p))) \end{array}}}
\vnbentry{\texttt{r12d-\allowbreak{}diag-\allowbreak{}in-\allowbreak{}fun}}{\vnbstmt{\begin{array}{@{}l@{}} \forall m, f, a. \\ \quad \operatorname{is\text{-}metric\text{-}space}(m) \wedge \\ \quad (f \in (\mathbb{N} \to \operatorname{cseq}(m))) \wedge \\ \quad (a \in (\mathbb{N} \to \mathbb{N})) \Rightarrow \\ \quad \quad ((\lambda v \in \mathbb{N}.\; f(v)(a(v))) \in (\mathbb{N} \to \operatorname{pts}(m))) \end{array}}}
\vnbentry{\texttt{r12d-\allowbreak{}rad-\allowbreak{}tail}}{\vnbstmt{\begin{array}{@{}l@{}} \forall d, v. \\ \quad (\operatorname{null\text{-}rr\text{-}seq}(d) \wedge \operatorname{pos\text{-}rr}(v)) \Rightarrow  \\ \quad \quad \exists r \in \mathbb{N}. \forall k \in \mathbb{N}. \\ \quad \quad \quad ((r \leq k) \Rightarrow (d(k) \leq v)) \end{array}}}
\vnbentry{\texttt{r12d-\allowbreak{}diag-\allowbreak{}est}}{\vnbstmt{\begin{array}{@{}l@{}} \forall m, f, a, d, b, p, q, c. \\ \quad \operatorname{is\text{-}metric\text{-}space}(m) \wedge \\ \quad (f \in (\mathbb{N} \to \operatorname{quotient}(\operatorname{cauchy\text{-}setoid}(m)))) \wedge \\ \quad (a \in (\mathbb{N} \to \operatorname{cseq}(m))) \wedge \\ \quad (d \in (\mathbb{N} \to \mathbb{R})) \wedge \\ \quad (b \in (\mathbb{N} \to \mathbb{N})) \wedge \\ \quad \forall p \in \mathbb{N}. \\ \quad \quad (f(p) = \operatorname{class}(\operatorname{cauchy\text{-}setoid}(m), a(p))) \wedge \\ \quad \forall p \in \mathbb{N}, i, j. \\ \quad \quad ((i \in \mathbb{N}) \wedge ((j \in \mathbb{N}) \wedge ((b(p) \leq i) \wedge (b(p) \leq j)))) \Rightarrow  \\ \quad \quad \quad (\operatorname{dist}(m)(a(p)(i), a(p)(j)) \leq d(p)) \wedge \\ \quad (p \in \mathbb{N}) \wedge \\ \quad (q \in \mathbb{N}) \wedge \\ \quad (c \in \mathbb{R}) \wedge \\ \quad \operatorname{pos\text{-}rr}(c) \Rightarrow \\ \quad \quad \operatorname{dist}(m)(a(p)(b(p)), a(q)(b(q))) \\ \quad \quad \leq\; ((d(p) + d(q)) + (\operatorname{dist}(\operatorname{completion}(m))(f(p), f(q)) + c)) \end{array}}}
\vnbentry{\texttt{r12d-\allowbreak{}cross-\allowbreak{}est}}{\vnbstmt{\begin{array}{@{}l@{}} \forall m, f, a, d, b, p, q, c. \\ \quad \operatorname{is\text{-}metric\text{-}space}(m) \wedge \\ \quad (f \in (\mathbb{N} \to \operatorname{quotient}(\operatorname{cauchy\text{-}setoid}(m)))) \wedge \\ \quad (a \in (\mathbb{N} \to \operatorname{cseq}(m))) \wedge \\ \quad (d \in (\mathbb{N} \to \mathbb{R})) \wedge \\ \quad (b \in (\mathbb{N} \to \mathbb{N})) \wedge \\ \quad \forall p \in \mathbb{N}. \\ \quad \quad (f(p) = \operatorname{class}(\operatorname{cauchy\text{-}setoid}(m), a(p))) \wedge \\ \quad \forall p \in \mathbb{N}, i, j. \\ \quad \quad ((i \in \mathbb{N}) \wedge ((j \in \mathbb{N}) \wedge ((b(p) \leq i) \wedge (b(p) \leq j)))) \Rightarrow  \\ \quad \quad \quad (\operatorname{dist}(m)(a(p)(i), a(p)(j)) \leq d(p)) \wedge \\ \quad (p \in \mathbb{N}) \wedge \\ \quad (q \in \mathbb{N}) \wedge \\ \quad (c \in \mathbb{R}) \wedge \\ \quad \operatorname{pos\text{-}rr}(c) \wedge \\ \quad (b(p) \leq q) \Rightarrow \\ \quad \quad \operatorname{dist}(m)(a(q)(b(q)), a(p)(q)) \\ \quad \quad \leq\; ((d(q) + d(p)) + (\operatorname{dist}(\operatorname{completion}(m))(f(q), f(p)) + c)) \end{array}}}
\vnbentry{\texttt{r12d-\allowbreak{}diag-\allowbreak{}cauchy}}{\vnbstmt{\begin{array}{@{}l@{}} \forall m, f, a, d, b. \\ \quad \operatorname{is\text{-}metric\text{-}space}(m) \wedge \\ \quad (f \in (\mathbb{N} \to \operatorname{quotient}(\operatorname{cauchy\text{-}setoid}(m)))) \wedge \\ \quad (a \in (\mathbb{N} \to \operatorname{cseq}(m))) \wedge \\ \quad \operatorname{null\text{-}rr\text{-}seq}(d) \wedge \\ \quad (b \in (\mathbb{N} \to \mathbb{N})) \wedge \\ \quad \forall p \in \mathbb{N}. \\ \quad \quad (f(p) = \operatorname{class}(\operatorname{cauchy\text{-}setoid}(m), a(p))) \wedge \\ \quad \forall p \in \mathbb{N}, i, j. \\ \quad \quad ((i \in \mathbb{N}) \wedge ((j \in \mathbb{N}) \wedge ((b(p) \leq i) \wedge (b(p) \leq j)))) \Rightarrow  \\ \quad \quad \quad (\operatorname{dist}(m)(a(p)(i), a(p)(j)) \leq d(p)) \wedge \\ \quad \operatorname{is\text{-}cauchy\text{-}seq}(\operatorname{completion}(m), f) \Rightarrow \\ \quad \quad \operatorname{is\text{-}cauchy\text{-}seq}(m, (\lambda v \in \mathbb{N}.\; a(v)(b(v)))) \end{array}}}
\vnbentry{\texttt{r12d-\allowbreak{}diag-\allowbreak{}in-\allowbreak{}cseq}}{\vnbstmt{\begin{array}{@{}l@{}} \forall m, f, a, d, b. \\ \quad \operatorname{is\text{-}metric\text{-}space}(m) \wedge \\ \quad (f \in (\mathbb{N} \to \operatorname{quotient}(\operatorname{cauchy\text{-}setoid}(m)))) \wedge \\ \quad (a \in (\mathbb{N} \to \operatorname{cseq}(m))) \wedge \\ \quad \operatorname{null\text{-}rr\text{-}seq}(d) \wedge \\ \quad (b \in (\mathbb{N} \to \mathbb{N})) \wedge \\ \quad \forall p \in \mathbb{N}. \\ \quad \quad (f(p) = \operatorname{class}(\operatorname{cauchy\text{-}setoid}(m), a(p))) \wedge \\ \quad \forall p \in \mathbb{N}, i, j. \\ \quad \quad ((i \in \mathbb{N}) \wedge ((j \in \mathbb{N}) \wedge ((b(p) \leq i) \wedge (b(p) \leq j)))) \Rightarrow  \\ \quad \quad \quad (\operatorname{dist}(m)(a(p)(i), a(p)(j)) \leq d(p)) \wedge \\ \quad \operatorname{is\text{-}cauchy\text{-}seq}(\operatorname{completion}(m), f) \Rightarrow \\ \quad \quad ((\lambda v \in \mathbb{N}.\; a(v)(b(v))) \in \operatorname{cseq}(m)) \end{array}}}
\vnbentry{\texttt{r12d-\allowbreak{}diag-\allowbreak{}limit}}{\vnbstmt{\begin{array}{@{}l@{}} \forall m, f, a, d, b. \\ \quad \operatorname{is\text{-}metric\text{-}space}(m) \wedge \\ \quad (f \in (\mathbb{N} \to \operatorname{quotient}(\operatorname{cauchy\text{-}setoid}(m)))) \wedge \\ \quad (a \in (\mathbb{N} \to \operatorname{cseq}(m))) \wedge \\ \quad \operatorname{null\text{-}rr\text{-}seq}(d) \wedge \\ \quad (b \in (\mathbb{N} \to \mathbb{N})) \wedge \\ \quad \forall p \in \mathbb{N}. \\ \quad \quad (f(p) = \operatorname{class}(\operatorname{cauchy\text{-}setoid}(m), a(p))) \wedge \\ \quad \forall p \in \mathbb{N}, i, j. \\ \quad \quad ((i \in \mathbb{N}) \wedge ((j \in \mathbb{N}) \wedge ((b(p) \leq i) \wedge (b(p) \leq j)))) \Rightarrow  \\ \quad \quad \quad (\operatorname{dist}(m)(a(p)(i), a(p)(j)) \leq d(p)) \wedge \\ \quad \operatorname{is\text{-}cauchy\text{-}seq}(\operatorname{completion}(m), f) \Rightarrow \\ \quad \quad \operatorname{converges\text{-}to}(\operatorname{completion}(m), f, \\ \quad \quad \quad \operatorname{class}(\operatorname{cauchy\text{-}setoid}(m), (\lambda v \in \mathbb{N}.\; a(v)(b(v))))) \end{array}}}
\vnbentry{\texttt{completion-\allowbreak{}is-\allowbreak{}complete}}{\vnbstmt{\begin{array}{@{}l@{}} \forall m. \\ \quad \operatorname{is\text{-}metric\text{-}space}(m) \Rightarrow  \\ \quad \quad \operatorname{is\text{-}complete}(\operatorname{completion}(m)) \end{array}}}

\vnbfile{theorem-library/product-weights.scm}{theorem-library/product-weights.scm}
\noindent{\small product-weights.scm -- ANY TWO SUMMABLE WEIGHT SEQUENCES GIVE THE SAME TOPOLOGY ON THE COUNTABLE PRODUCT, proven -- retiring the asserted `product-weights-equivalent' of structure-library/product-metric.scm, the last statement of that file about the product's topology.}\par
\vnbentry{\texttt{product-\allowbreak{}identity-\allowbreak{}continuous}}{\vnbstmt{\begin{array}{@{}l@{}} \forall s, w, a. \\ \quad \operatorname{is\text{-}ms\text{-}sequence}(s) \wedge \\ \quad \operatorname{summable\text{-}weight}(w) \wedge \\ \quad \operatorname{summable\text{-}weight}(a) \Rightarrow \\ \quad \quad \operatorname{is\text{-}continuous}(\operatorname{product\text{-}metric\text{-}w}(s, w), \operatorname{product\text{-}metric\text{-}w}(s, a), \\ \quad \quad \quad (\lambda x \in \operatorname{pts}(\operatorname{product\text{-}metric\text{-}w}(s, w)).\; x)) \end{array}}}
\vnbentry{\texttt{product-\allowbreak{}weights-\allowbreak{}equivalent}}{\vnbstmt{\begin{array}{@{}l@{}} \forall s, w, a. \\ \quad \operatorname{is\text{-}ms\text{-}sequence}(s) \wedge \\ \quad \operatorname{summable\text{-}weight}(w) \wedge \\ \quad \operatorname{summable\text{-}weight}(a) \Rightarrow \\ \quad \quad \operatorname{is\text{-}continuous}(\operatorname{product\text{-}metric\text{-}w}(s, w), \operatorname{product\text{-}metric\text{-}w}(s, a), \\ \quad \quad \quad (\lambda x \in \operatorname{pts}(\operatorname{product\text{-}metric\text{-}w}(s, w)).\; x)) \wedge \\ \quad \quad \operatorname{is\text{-}continuous}(\operatorname{product\text{-}metric\text{-}w}(s, a), \operatorname{product\text{-}metric\text{-}w}(s, w), \\ \quad \quad \quad (\lambda x \in \operatorname{pts}(\operatorname{product\text{-}metric\text{-}w}(s, a)).\; x)) \end{array}}}

\vnbfile{theorem-library/c-metric-summable.scm}{theorem-library/c-metric-summable.scm}
\noindent{\small c-metric-summable.scm -- THE SERIES DEFINING THE CANONICAL METRIC OF A COUNTABLY-METRISED SPACE CONVERGES, proven, together with the plumbing every later argument about C-METRIC needs.}\par
\vnbentry{\texttt{c-\allowbreak{}metric-\allowbreak{}pseudometric-\allowbreak{}at}}{\vnbstmt{\begin{array}{@{}l@{}} \forall s.; \forall k \in \mathbb{N}. \\ \quad \operatorname{is\text{-}c\text{-}metric\text{-}space}(s) \Rightarrow  \\ \quad \quad \operatorname{is\text{-}pseudometric}(\operatorname{dists}(s)(k), \operatorname{pts}(s)) \end{array}}}
\vnbentry{\texttt{c-\allowbreak{}metric-\allowbreak{}dist-\allowbreak{}real}}{\vnbstmt{\begin{array}{@{}l@{}} \forall s.; \forall k \in \mathbb{N}, u, v \in \operatorname{pts}(s). \\ \quad \operatorname{is\text{-}c\text{-}metric\text{-}space}(s) \Rightarrow  \\ \quad \quad (\operatorname{dists}(s)(k)(u, v) \in \mathbb{R}) \end{array}}}
\vnbentry{\texttt{c-\allowbreak{}metric-\allowbreak{}term-\allowbreak{}bound}}{\vnbstmt{\begin{array}{@{}l@{}} \forall s, k, u, v, a. \\ \quad \operatorname{is\text{-}c\text{-}metric\text{-}space}(s) \wedge \\ \quad (k \in \mathbb{N}) \wedge \\ \quad (u \in \operatorname{pts}(s)) \wedge \\ \quad (v \in \operatorname{pts}(s)) \wedge \\ \quad (a \in \mathbb{R}) \wedge \\ \quad (0 \leq a) \Rightarrow \\ \quad \quad ((a \cdot \operatorname{min}(1, \operatorname{dists}(s)(k)(u, v))) \in \mathbb{R}) \wedge \\ \quad \quad (0 \leq (a \cdot \operatorname{min}(1, \operatorname{dists}(s)(k)(u, v)))) \wedge \\ \quad \quad ((a \cdot \operatorname{min}(1, \operatorname{dists}(s)(k)(u, v))) \leq a) \end{array}}}
\vnbentry{\texttt{c-\allowbreak{}metric-\allowbreak{}term-\allowbreak{}seq-\allowbreak{}in-\allowbreak{}fun}}{\vnbstmt{\begin{array}{@{}l@{}} \forall s, w, u, v. \\ \quad \operatorname{is\text{-}c\text{-}metric\text{-}space}(s) \wedge \\ \quad \operatorname{summable\text{-}weight}(w) \wedge \\ \quad (u \in \operatorname{pts}(s)) \wedge \\ \quad (v \in \operatorname{pts}(s)) \Rightarrow \\ \quad \quad (\lambda k \in \mathbb{N}.\; (w(k) \cdot \operatorname{min}(1, \operatorname{dists}(s)(k)(u, v)))) \\ \quad \quad \in\; (\mathbb{N} \to \mathbb{R}) \end{array}}}
\vnbentry{\texttt{c-\allowbreak{}metric-\allowbreak{}summable}}{\vnbstmt{\begin{array}{@{}l@{}} \forall s, w, u, v. \\ \quad \operatorname{is\text{-}c\text{-}metric\text{-}space}(s) \wedge \\ \quad \operatorname{summable\text{-}weight}(w) \wedge \\ \quad (u \in \operatorname{pts}(s)) \wedge \\ \quad (v \in \operatorname{pts}(s)) \Rightarrow \\ \quad \quad \operatorname{series\text{-}converges}((\lambda k \in \mathbb{N}.\; (w(k) \cdot \operatorname{min}(1, \operatorname{dists}(s)(k)(u, v))))) \end{array}}}
\vnbentry{\texttt{c-\allowbreak{}metric-\allowbreak{}carrier}}{\vnbstmt{\begin{array}{@{}l@{}} \forall s, w. \\ \quad (\operatorname{pts}(\operatorname{c\text{-}metric\text{-}w}(s, w)) \simeq \operatorname{pts}(s)) \end{array}}}
\vnbentry{\texttt{c-\allowbreak{}metric-\allowbreak{}distance}}{\vnbstmt{\begin{array}{@{}l@{}} \forall s, w, u, v. \\ \quad ((u \in \operatorname{pts}(s)) \wedge (v \in \operatorname{pts}(s))) \Rightarrow  \\ \quad \quad (\operatorname{dist}(\operatorname{c\text{-}metric\text{-}w}(s, w))(u, v) \simeq \operatorname{series\text{-}limit}((\lambda k \in \mathbb{N}.\; (w(k) \cdot \operatorname{min}(1, \operatorname{dists}(s)(k)(u, v)))))) \end{array}}}
\vnbentry{\texttt{c-\allowbreak{}metric-\allowbreak{}dist-\allowbreak{}converges}}{\vnbstmt{\begin{array}{@{}l@{}} \forall s, w, u, v. \\ \quad \operatorname{is\text{-}c\text{-}metric\text{-}space}(s) \wedge \\ \quad \operatorname{summable\text{-}weight}(w) \wedge \\ \quad (u \in \operatorname{pts}(s)) \wedge \\ \quad (v \in \operatorname{pts}(s)) \Rightarrow \\ \quad \quad \operatorname{series\text{-}converges\text{-}to}((\lambda k \in \mathbb{N}.\; (w(k) \cdot \operatorname{min}(1, \operatorname{dists}(s)(k)(u, v)))), \\ \quad \quad \quad \operatorname{dist}(\operatorname{c\text{-}metric\text{-}w}(s, w))(u, v)) \end{array}}}
\vnbentry{\texttt{c-\allowbreak{}metric-\allowbreak{}canonical-\allowbreak{}summable}}{\vnbstmt{\begin{array}{@{}l@{}} \forall s, u, v. \\ \quad \operatorname{is\text{-}c\text{-}metric\text{-}space}(s) \wedge \\ \quad (u \in \operatorname{pts}(s)) \wedge \\ \quad (v \in \operatorname{pts}(s)) \Rightarrow \\ \quad \quad \operatorname{series\text{-}converges}((\lambda k \in \mathbb{N}.\; \\ \quad \quad \quad \quad ((\lambda n \in \mathbb{N}.\; /(1, \operatorname{power}(2, (n + 1))))(k) \cdot \operatorname{min}(1, \operatorname{dists}(s)(k)(u, v))))) \end{array}}}

\vnbfile{theorem-library/c-metric-is-metric.scm}{theorem-library/c-metric-is-metric.scm}
\noindent{\small c-metric-is-metric.scm -- THE CANONICAL METRIC OF A COUNTABLY-METRISED SPACE IS A METRIC SPACE, proven.}\par
\vnbentry{\texttt{cmim-\allowbreak{}dist-\allowbreak{}in-\allowbreak{}rr}}{\vnbstmt{\begin{array}{@{}l@{}} \forall s, w.; \forall x, y \in \operatorname{pts}(s). \\ \quad \operatorname{is\text{-}c\text{-}metric\text{-}space}(s) \wedge \\ \quad \operatorname{summable\text{-}weight}(w) \Rightarrow \\ \quad \quad (\operatorname{dist}(\operatorname{c\text{-}metric\text{-}w}(s, w))(x, y) \in \mathbb{R}) \end{array}}}
\vnbentry{\texttt{cmim-\allowbreak{}term-\allowbreak{}nonneg}}{\vnbstmt{\begin{array}{@{}l@{}} \forall s, w.; \forall x, y \in \operatorname{pts}(s), n \in \mathbb{N}. \\ \quad \operatorname{is\text{-}c\text{-}metric\text{-}space}(s) \wedge \\ \quad \operatorname{summable\text{-}weight}(w) \Rightarrow \\ \quad \quad 0 \\ \quad \quad \leq\; (\lambda k \in \mathbb{N}.\; (w(k) \cdot \operatorname{min}(1, \operatorname{dists}(s)(k)(x, y))))(n) \end{array}}}
\vnbentry{\texttt{cmim-\allowbreak{}dist-\allowbreak{}nonneg}}{\vnbstmt{\begin{array}{@{}l@{}} \forall s, w.; \forall x, y \in \operatorname{pts}(s). \\ \quad \operatorname{is\text{-}c\text{-}metric\text{-}space}(s) \wedge \\ \quad \operatorname{summable\text{-}weight}(w) \Rightarrow \\ \quad \quad (0 \leq \operatorname{dist}(\operatorname{c\text{-}metric\text{-}w}(s, w))(x, y)) \end{array}}}
\vnbentry{\texttt{cmim-\allowbreak{}term-\allowbreak{}self-\allowbreak{}zero}}{\vnbstmt{\begin{array}{@{}l@{}} \forall s, w.; \forall x \in \operatorname{pts}(s), n \in \mathbb{N}. \\ \quad \operatorname{is\text{-}c\text{-}metric\text{-}space}(s) \wedge \\ \quad \operatorname{summable\text{-}weight}(w) \Rightarrow \\ \quad \quad (\lambda k \in \mathbb{N}.\; (w(k) \cdot \operatorname{min}(1, \operatorname{dists}(s)(k)(x, x))))(n) \\ \quad \quad =\; 0 \end{array}}}
\vnbentry{\texttt{cmim-\allowbreak{}dist-\allowbreak{}self-\allowbreak{}zero}}{\vnbstmt{\begin{array}{@{}l@{}} \forall s, w.; \forall x \in \operatorname{pts}(s). \\ \quad \operatorname{is\text{-}c\text{-}metric\text{-}space}(s) \wedge \\ \quad \operatorname{summable\text{-}weight}(w) \Rightarrow \\ \quad \quad (\operatorname{dist}(\operatorname{c\text{-}metric\text{-}w}(s, w))(x, x) = 0) \end{array}}}
\vnbentry{\texttt{cmim-\allowbreak{}term-\allowbreak{}zero}}{\vnbstmt{\begin{array}{@{}l@{}} \forall s, w.; \forall x, y \in \operatorname{pts}(s), n \in \mathbb{N}. \\ \quad \operatorname{is\text{-}c\text{-}metric\text{-}space}(s) \wedge \\ \quad \operatorname{summable\text{-}weight}(w) \wedge \\ \quad (\lambda k \in \mathbb{N}.\; (w(k) \cdot \operatorname{min}(1, \operatorname{dists}(s)(k)(x, y))))(n) \\ \quad =\; 0 \Rightarrow \\ \quad \quad (\operatorname{dists}(s)(n)(x, y) = 0) \end{array}}}
\vnbentry{\texttt{cmim-\allowbreak{}dist-\allowbreak{}zero-\allowbreak{}eq}}{\vnbstmt{\begin{array}{@{}l@{}} \forall s, w.; \forall x, y \in \operatorname{pts}(s). \\ \quad \operatorname{is\text{-}c\text{-}metric\text{-}space}(s) \wedge \\ \quad \operatorname{summable\text{-}weight}(w) \wedge \\ \quad (\operatorname{dist}(\operatorname{c\text{-}metric\text{-}w}(s, w))(x, y) = 0) \Rightarrow \\ \quad \quad (x = y) \end{array}}}
\vnbentry{\texttt{cmim-\allowbreak{}term-\allowbreak{}sym}}{\vnbstmt{\begin{array}{@{}l@{}} \forall s, w.; \forall x, y \in \operatorname{pts}(s), n \in \mathbb{N}. \\ \quad \operatorname{is\text{-}c\text{-}metric\text{-}space}(s) \wedge \\ \quad \operatorname{summable\text{-}weight}(w) \Rightarrow \\ \quad \quad (\lambda k \in \mathbb{N}.\; (w(k) \cdot \operatorname{min}(1, \operatorname{dists}(s)(k)(x, y))))(n) \\ \quad \quad \leq\; (\lambda k \in \mathbb{N}.\; (w(k) \cdot \operatorname{min}(1, \operatorname{dists}(s)(k)(y, x))))(n) \end{array}}}
\vnbentry{\texttt{cmim-\allowbreak{}dist-\allowbreak{}sym}}{\vnbstmt{\begin{array}{@{}l@{}} \forall s, w.; \forall x, y \in \operatorname{pts}(s). \\ \quad \operatorname{is\text{-}c\text{-}metric\text{-}space}(s) \wedge \\ \quad \operatorname{summable\text{-}weight}(w) \Rightarrow \\ \quad \quad \operatorname{dist}(\operatorname{c\text{-}metric\text{-}w}(s, w))(x, y) \\ \quad \quad =\; \operatorname{dist}(\operatorname{c\text{-}metric\text{-}w}(s, w))(y, x) \end{array}}}
\vnbentry{\texttt{cmim-\allowbreak{}sum-\allowbreak{}seq-\allowbreak{}in-\allowbreak{}fun}}{\vnbstmt{\begin{array}{@{}l@{}} \forall s, w.; \forall x, y, z \in \operatorname{pts}(s). \\ \quad \operatorname{is\text{-}c\text{-}metric\text{-}space}(s) \wedge \\ \quad \operatorname{summable\text{-}weight}(w) \Rightarrow \\ \quad \quad (\lambda k \in \mathbb{N}.\; \\ \quad \quad \quad ((w(k) \cdot \operatorname{min}(1, \operatorname{dists}(s)(k)(x, y))) + (w(k) \cdot \operatorname{min}(1, \operatorname{dists}(s)(k)(y, z))))) \\ \quad \quad \in\; (\mathbb{N} \to \mathbb{R}) \end{array}}}
\vnbentry{\texttt{cmim-\allowbreak{}sum-\allowbreak{}seq-\allowbreak{}split}}{\vnbstmt{\begin{array}{@{}l@{}} \forall s, w.; \forall x, y, z \in \operatorname{pts}(s), i \in \mathbb{N}. \\ \quad \operatorname{is\text{-}c\text{-}metric\text{-}space}(s) \wedge \\ \quad \operatorname{summable\text{-}weight}(w) \Rightarrow \\ \quad \quad (\lambda k \in \mathbb{N}.\; ((w(k) \cdot \operatorname{min}(1, \operatorname{dists}(s)(k)(x, y))) + (w(k) \cdot \operatorname{min}(1, \operatorname{dists}(s)(k)(y, z)))))(i) \\ \quad \quad =\; ((\lambda k \in \mathbb{N}.\; (w(k) \cdot \operatorname{min}(1, \operatorname{dists}(s)(k)(x, y))))(i) + (\lambda k \in \mathbb{N}.\; (w(k) \cdot \operatorname{min}(1, \operatorname{dists}(s)(k)(y, z))))(i)) \end{array}}}
\vnbentry{\texttt{cmim-\allowbreak{}term-\allowbreak{}triangle}}{\vnbstmt{\begin{array}{@{}l@{}} \forall s, w.; \forall x, y, z \in \operatorname{pts}(s), n \in \mathbb{N}. \\ \quad \operatorname{is\text{-}c\text{-}metric\text{-}space}(s) \wedge \\ \quad \operatorname{summable\text{-}weight}(w) \Rightarrow \\ \quad \quad (\lambda k \in \mathbb{N}.\; (w(k) \cdot \operatorname{min}(1, \operatorname{dists}(s)(k)(x, z))))(n) \\ \quad \quad \leq\; (\lambda k \in \mathbb{N}.\; ((w(k) \cdot \operatorname{min}(1, \operatorname{dists}(s)(k)(x, y))) + (w(k) \cdot \operatorname{min}(1, \operatorname{dists}(s)(k)(y, z)))))(n) \end{array}}}
\vnbentry{\texttt{cmim-\allowbreak{}dist-\allowbreak{}triangle}}{\vnbstmt{\begin{array}{@{}l@{}} \forall s, w.; \forall x, y, z \in \operatorname{pts}(s). \\ \quad \operatorname{is\text{-}c\text{-}metric\text{-}space}(s) \wedge \\ \quad \operatorname{summable\text{-}weight}(w) \Rightarrow \\ \quad \quad \operatorname{dist}(\operatorname{c\text{-}metric\text{-}w}(s, w))(x, z) \\ \quad \quad \leq\; (\operatorname{dist}(\operatorname{c\text{-}metric\text{-}w}(s, w))(x, y) + \operatorname{dist}(\operatorname{c\text{-}metric\text{-}w}(s, w))(y, z)) \end{array}}}
\vnbentry{\texttt{c-\allowbreak{}metric-\allowbreak{}w-\allowbreak{}is-\allowbreak{}metric-\allowbreak{}space}}{\vnbstmt{\begin{array}{@{}l@{}} \forall s, w. \\ \quad \operatorname{is\text{-}c\text{-}metric\text{-}space}(s) \wedge \\ \quad \operatorname{summable\text{-}weight}(w) \Rightarrow \\ \quad \quad \operatorname{is\text{-}metric\text{-}space}(\operatorname{c\text{-}metric\text{-}w}(s, w)) \end{array}}}
\vnbentry{\texttt{c-\allowbreak{}metric-\allowbreak{}is-\allowbreak{}metric-\allowbreak{}space}}{\vnbstmt{\begin{array}{@{}l@{}} \forall s. \\ \quad \operatorname{is\text{-}c\text{-}metric\text{-}space}(s) \Rightarrow  \\ \quad \quad \operatorname{is\text{-}metric\text{-}space}(\operatorname{c\text{-}metric}(s)) \end{array}}}

\vnbfile{theorem-library/ccint-creep.scm}{theorem-library/ccint-creep.scm}
\noindent{\small ccint-creep.scm -- the CREEPING PRINCIPLE on a closed interval, PROVEN.}\par
\vnbentry{\texttt{ccint-\allowbreak{}creep}}{\vnbstmt{\begin{array}{@{}l@{}} \forall a, b, g. \\ \quad (a \in \mathbb{R}) \wedge \\ \quad (b \in \mathbb{R}) \wedge \\ \quad (a \leq b) \wedge \\ \quad (g \subseteq \mathbb{R}) \wedge \\ \quad (a \in g) \wedge \\ \quad \operatorname{rr\text{-}upper\text{-}bound}(g, b) \wedge \\ \quad \forall t \in \operatorname{ccint}(a, b). \exists d \in \mathbb{R}. \\ \quad \quad (0 < d) \wedge \\ \quad \quad \forall y \in \operatorname{ccint}(a, b). \\ \quad \quad \quad (\exists w \in g.\; ((t - d) < w) \wedge (y \leq (t + d))) \Rightarrow  \\ \quad \quad \quad \quad (y \in g) \Rightarrow \\ \quad \quad (b \in g) \end{array}}}

\vnbfile{theorem-library/ccint-bounded.scm}{theorem-library/ccint-bounded.scm}
\noindent{\small ccint-bounded.scm -- a continuous function on [a,b] is BOUNDED ABOVE, PROVEN.}\par
\vnbentry{\texttt{continuous-\allowbreak{}bounded-\allowbreak{}above-\allowbreak{}on-\allowbreak{}ccint}}{\vnbstmt{\begin{array}{@{}l@{}} \forall f, a, b. \\ \quad (f \in (\mathbb{R} \to \mathbb{R})) \wedge \\ \quad (a \in \mathbb{R}) \wedge \\ \quad (b \in \mathbb{R}) \wedge \\ \quad (a \leq b) \wedge \\ \quad \forall x \in \operatorname{ccint}(a, b). \\ \quad \quad \operatorname{is\text{-}continuous\text{-}at}(\operatorname{rr\text{-}ms}, \operatorname{rr\text{-}ms}, f, x) \Rightarrow \\ \quad \quad \exists k \in \mathbb{R}. \forall z \in \operatorname{ccint}(a, b). \\ \quad \quad \quad (f(z) \leq k) \end{array}}}

\vnbfile{theorem-library/uniform-continuity-ccint.scm}{theorem-library/uniform-continuity-ccint.scm}
\noindent{\small uniform-continuity-ccint.scm -- a continuous function on a CLOSED BOUNDED INTERVAL is UNIFORMLY continuous there, PROVEN.}\par
\vnbentry{\texttt{continuous-\allowbreak{}uniformly-\allowbreak{}continuous-\allowbreak{}on-\allowbreak{}ccint}}{\vnbstmt{\begin{array}{@{}l@{}} \forall f, a, b, s. \\ \quad (f \in (\mathbb{R} \to \mathbb{R})) \wedge \\ \quad (a \in \mathbb{R}) \wedge \\ \quad (b \in \mathbb{R}) \wedge \\ \quad (a \leq b) \wedge \\ \quad \operatorname{pos\text{-}rr}(s) \wedge \\ \quad \forall x \in \operatorname{ccint}(a, b). \\ \quad \quad \operatorname{is\text{-}continuous\text{-}at}(\operatorname{rr\text{-}ms}, \operatorname{rr\text{-}ms}, f, x) \Rightarrow \\ \quad \quad \exists d \in \mathbb{R}. \\ \quad \quad \quad (0 < d) \wedge \\ \quad \quad \quad \forall p, q \in \operatorname{ccint}(a, b). \\ \quad \quad \quad \quad (\lvert (p - q) \rvert \leq d) \Rightarrow  \\ \quad \quad \quad \quad \quad (\lvert (f(p) - f(q)) \rvert \leq s) \end{array}}}

\vnbfile{theorem-library/evt-proof.scm}{theorem-library/evt-proof.scm}
\noindent{\small evt-proof.scm -- the EXTREME VALUE THEOREM (max form), PROVEN.}\par
\vnbentry{\texttt{extreme-\allowbreak{}value-\allowbreak{}max}}{\vnbstmt{\begin{array}{@{}l@{}} \forall f, a, b. \\ \quad (f \in (\mathbb{R} \to \mathbb{R})) \wedge \\ \quad (a \in \mathbb{R}) \wedge \\ \quad (b \in \mathbb{R}) \wedge \\ \quad (a \leq b) \wedge \\ \quad \forall x \in \operatorname{ccint}(a, b). \\ \quad \quad \operatorname{is\text{-}continuous\text{-}at}(\operatorname{rr\text{-}ms}, \operatorname{rr\text{-}ms}, f, x) \Rightarrow \\ \quad \quad \exists c \in \operatorname{ccint}(a, b). \forall x \in \operatorname{ccint}(a, b). \\ \quad \quad \quad (f(x) \leq f(c)) \end{array}}}

\vnbfile{theorem-library/heine-borel-baby.scm}{theorem-library/heine-borel-baby.scm}
\noindent{\small NOTE 2026-09-20 (batch 12-B).  Where this file says the tree has no metric SUBSPACE structure, that is history: SUBSPACE-MS(s, A) and RESTRICT(f, A) are defined in structure-library/metric-subspace.scm, their laws are in theorem-library/metric-subspace-laws.scm and compact-subspace.scm, and}\par
\vnbentry{\texttt{empty-\allowbreak{}subset-\allowbreak{}any}}{\vnbstmt{\forall b.\; (\emptyset \subseteq b)}}
\vnbentry{\texttt{union-\allowbreak{}singleton-\allowbreak{}mem}}{\vnbstmt{\forall y \in \mathbf{Set}, s.\; (y \in (\{y\} \cup s))}}
\vnbentry{\texttt{union-\allowbreak{}right-\allowbreak{}subset}}{\vnbstmt{\forall a, s.\; (s \subseteq (a \cup s))}}
\vnbentry{\texttt{union-\allowbreak{}singleton-\allowbreak{}subset}}{\vnbstmt{\begin{array}{@{}l@{}} \forall b, y, s. \\ \quad (((y \in b) \wedge (s \subseteq b)) \Rightarrow ((\{y\} \cup s) \subseteq b)) \end{array}}}
\vnbentry{\texttt{big-\allowbreak{}union-\allowbreak{}mono}}{\vnbstmt{\begin{array}{@{}l@{}} \forall a, b. \\ \quad (a \subseteq b) \Rightarrow  \\ \quad \quad (\operatorname{big\text{-}union}(u, a, u) \subseteq \operatorname{big\text{-}union}(u, b, u)) \end{array}}}
\vnbentry{\texttt{card-\allowbreak{}union-\allowbreak{}singleton-\allowbreak{}nn}}{\vnbstmt{\begin{array}{@{}l@{}} \forall s, y \in \mathbf{Set}. \\ \quad ((|s| \in \mathbb{N}) \Rightarrow (|(\{y\} \cup s)| \in \mathbb{N})) \end{array}}}
\vnbentry{\texttt{heine-\allowbreak{}borel-\allowbreak{}baby}}{\vnbstmt{\begin{array}{@{}l@{}} \forall a, b, f. \\ \quad (a \in \mathbb{R}) \wedge \\ \quad (b \in \mathbb{R}) \wedge \\ \quad (a \leq b) \wedge \\ \quad \forall v \in f.\; \operatorname{is\text{-}open}(\operatorname{rr\text{-}ms}, v) \wedge \\ \quad (\operatorname{ccint}(a, b) \subseteq \operatorname{big\text{-}union}(u, f, u)) \Rightarrow \\ \quad \quad \exists g. \\ \quad \quad \quad (g \subseteq f) \wedge \\ \quad \quad \quad (g \in \mathbf{Set}) \wedge \\ \quad \quad \quad (|g| \in \mathbb{N}) \wedge \\ \quad \quad \quad (\operatorname{ccint}(a, b) \subseteq \operatorname{big\text{-}union}(u, g, u)) \end{array}}}

\vnbfile{theorem-library/evt-min-proof.scm}{theorem-library/evt-min-proof.scm}
\noindent{\small evt-min-proof.scm -- the EXTREME VALUE THEOREM (min form), PROVEN.}\par
\vnbentry{\texttt{extreme-\allowbreak{}value-\allowbreak{}min}}{\vnbstmt{\begin{array}{@{}l@{}} \forall f, a, b. \\ \quad (f \in (\mathbb{R} \to \mathbb{R})) \wedge \\ \quad (a \in \mathbb{R}) \wedge \\ \quad (b \in \mathbb{R}) \wedge \\ \quad (a \leq b) \wedge \\ \quad \forall x \in \operatorname{ccint}(a, b). \\ \quad \quad \operatorname{is\text{-}continuous\text{-}at}(\operatorname{rr\text{-}ms}, \operatorname{rr\text{-}ms}, f, x) \Rightarrow \\ \quad \quad \exists c \in \operatorname{ccint}(a, b). \forall x \in \operatorname{ccint}(a, b). \\ \quad \quad \quad (f(c) \leq f(x)) \end{array}}}

\vnbfile{theorem-library/ccint-abs-bounded.scm}{theorem-library/ccint-abs-bounded.scm}
\noindent{\small ccint-abs-bounded.scm -- a continuous function on a CLOSED BOUNDED INTERVAL is bounded in ABSOLUTE VALUE there, PROVEN modulo 0.}\par
\vnbentry{\texttt{continuous-\allowbreak{}abs-\allowbreak{}bounded-\allowbreak{}on-\allowbreak{}ccint}}{\vnbstmt{\begin{array}{@{}l@{}} \forall f, a, b. \\ \quad (f \in (\mathbb{R} \to \mathbb{R})) \wedge \\ \quad (a \in \mathbb{R}) \wedge \\ \quad (b \in \mathbb{R}) \wedge \\ \quad (a \leq b) \wedge \\ \quad \forall x \in \operatorname{ccint}(a, b). \\ \quad \quad \operatorname{is\text{-}continuous\text{-}at}(\operatorname{rr\text{-}ms}, \operatorname{rr\text{-}ms}, f, x) \Rightarrow \\ \quad \quad \exists m \in \mathbb{R}. \\ \quad \quad \quad ((0 \leq m) \wedge \forall z \in \operatorname{ccint}(a, b).\; (\lvert f(z) \rvert \leq m)) \end{array}}}

\vnbfile{theorem-library/bernstein-density.scm}{theorem-library/bernstein-density.scm}
\noindent{\small bernstein-density.scm -- THEOREM 5.2 of docs/calculus.pdf: a continuous function on [0,1] is the uniform limit of its Bernstein polynomials.}\par
\vnbentry{\texttt{bernstein-\allowbreak{}basis-\allowbreak{}nonneg}}{\vnbstmt{\begin{array}{@{}l@{}} \forall n \in \mathbb{N}, x, k. \\ \quad ((x \in \mathbb{R}) \wedge ((0 \leq x) \wedge ((x \leq 1) \wedge (k \in \mathbb{N})))) \Rightarrow  \\ \quad \quad (0 \leq \operatorname{bernstein\text{-}basis}(n, x)(k)) \end{array}}}
\vnbentry{\texttt{bernstein-\allowbreak{}term-\allowbreak{}bound}}{\vnbstmt{\begin{array}{@{}l@{}} \forall f, n, x, h, a, m, c, l. \\ \quad (f \in (\mathbb{R} \to \mathbb{R})) \wedge \\ \quad (n \in \mathbb{N}) \wedge \\ \quad \neg (n = 0) \wedge \\ \quad (x \in \operatorname{ccint}(0, 1)) \wedge \\ \quad (h \in \mathbb{R}) \wedge \\ \quad (0 < h) \wedge \\ \quad (a \in \mathbb{R}) \wedge \\ \quad (0 < a) \wedge \\ \quad (m \in \mathbb{R}) \wedge \\ \quad (0 \leq m) \wedge \\ \quad (c \in \mathbb{R}) \wedge \\ \quad (0 \leq c) \wedge \\ \quad (((n \cdot n) \cdot ((a \cdot a) \cdot c)) = (m + m)) \wedge \\ \quad \forall p, q \in \operatorname{ccint}(0, 1). \\ \quad \quad (\lvert (p - q) \rvert \leq a) \Rightarrow  \\ \quad \quad \quad (\lvert (f(p) - f(q)) \rvert \leq h) \wedge \\ \quad \forall z \in \operatorname{ccint}(0, 1).\; (\lvert f(z) \rvert \leq m) \wedge \\ \quad (l \in \mathbb{N}) \Rightarrow \\ \quad \quad (\operatorname{bernstein\text{-}basis}(n, x)(l) \cdot (f(x) - f(/(l, n)))) \\ \quad \quad \leq\; (\operatorname{bernstein\text{-}basis}(n, x)(l) \cdot (h + (c \cdot ((l - (n \cdot x)) \cdot (l - (n \cdot x)))))) \wedge \\ \quad \quad (\operatorname{bernstein\text{-}basis}(n, x)(l) \cdot ((0 - 1) \cdot (h + (c \cdot ((l - (n \cdot x)) \cdot (l - (n \cdot x))))))) \\ \quad \quad \leq\; (\operatorname{bernstein\text{-}basis}(n, x)(l) \cdot (f(x) - f(/(l, n)))) \end{array}}}
\vnbentry{\texttt{bernstein-\allowbreak{}approx-\allowbreak{}bound}}{\vnbstmt{\begin{array}{@{}l@{}} \forall f, n, x, h, a, m, c. \\ \quad (f \in (\mathbb{R} \to \mathbb{R})) \wedge \\ \quad (n \in \mathbb{N}) \wedge \\ \quad \neg (n = 0) \wedge \\ \quad (x \in \operatorname{ccint}(0, 1)) \wedge \\ \quad (h \in \mathbb{R}) \wedge \\ \quad (0 < h) \wedge \\ \quad (a \in \mathbb{R}) \wedge \\ \quad (0 < a) \wedge \\ \quad (m \in \mathbb{R}) \wedge \\ \quad (0 \leq m) \wedge \\ \quad (c \in \mathbb{R}) \wedge \\ \quad (0 \leq c) \wedge \\ \quad (((n \cdot n) \cdot ((a \cdot a) \cdot c)) = (m + m)) \wedge \\ \quad ((c \cdot n) \leq h) \wedge \\ \quad \forall p, q \in \operatorname{ccint}(0, 1). \\ \quad \quad (\lvert (p - q) \rvert \leq a) \Rightarrow  \\ \quad \quad \quad (\lvert (f(p) - f(q)) \rvert \leq h) \wedge \\ \quad \forall z \in \operatorname{ccint}(0, 1).\; (\lvert f(z) \rvert \leq m) \Rightarrow \\ \quad \quad (\lvert (f(x) - \operatorname{bernstein\text{-}poly}(f, n, x)) \rvert \leq (h + h)) \end{array}}}
\vnbentry{\texttt{bernstein-\allowbreak{}uniform-\allowbreak{}approximation}}{\vnbstmt{\begin{array}{@{}l@{}} \forall f, s. \\ \quad (f \in (\mathbb{R} \to \mathbb{R})) \wedge \\ \quad \forall x \in \operatorname{ccint}(0, 1). \\ \quad \quad \operatorname{is\text{-}continuous\text{-}at}(\operatorname{rr\text{-}ms}, \operatorname{rr\text{-}ms}, f, x) \wedge \\ \quad \operatorname{pos\text{-}rr}(s) \Rightarrow \\ \quad \quad \exists p \in \mathbb{N}. \forall n, a. \\ \quad \quad \quad ((n \in \mathbb{N}) \wedge ((p \leq n) \wedge (a \in \operatorname{ccint}(0, 1)))) \Rightarrow  \\ \quad \quad \quad \quad (\lvert (f(a) - \operatorname{bernstein\text{-}poly}(f, n, a)) \rvert \leq s) \end{array}}}

\vnbfile{theorem-library/mvt-cluster-readoffs.scm}{theorem-library/mvt-cluster-readoffs.scm}
\noindent{\small theorem-library/mvt-cluster-readoffs.scm}\par
\vnbentry{\texttt{diff-\allowbreak{}value-\allowbreak{}real}}{\vnbstmt{\begin{array}{@{}l@{}} \forall f, a, l. \\ \quad \operatorname{is\text{-}diff\text{-}at}(f, a, l) \Rightarrow  \\ \quad \quad (l \in \mathbb{R}) \end{array}}}
\vnbentry{\texttt{deriv-\allowbreak{}neg}}{\vnbstmt{\begin{array}{@{}l@{}} \forall f, a, l. \\ \quad \operatorname{is\text{-}diff\text{-}at}(f, a, l) \Rightarrow  \\ \quad \quad \operatorname{is\text{-}diff\text{-}at}((\lambda z \in \mathbb{R}.\; -f(z)), a, -l) \end{array}}}
\vnbentry{\texttt{max-\allowbreak{}val-\allowbreak{}strict-\allowbreak{}interior}}{\vnbstmt{\begin{array}{@{}l@{}} \forall h, a, b.; \forall c \in \operatorname{ccint}(a, b). \\ \quad (((h(a) = h(b)) \wedge (h(a) < h(c))) \Rightarrow ((a < c) \wedge (c < b))) \end{array}}}
\vnbentry{\texttt{min-\allowbreak{}val-\allowbreak{}strict-\allowbreak{}interior}}{\vnbstmt{\begin{array}{@{}l@{}} \forall h, a, b.; \forall c \in \operatorname{ccint}(a, b). \\ \quad (((h(a) = h(b)) \wedge (h(c) < h(a))) \Rightarrow ((a < c) \wedge (c < b))) \end{array}}}

\vnbfile{theorem-library/deriv-difference.scm}{theorem-library/deriv-difference.scm}
\noindent{\small deriv-difference.scm -- the DIFFERENCE rule for the Caratheodory derivative IS-DIFF-AT: if f has derivative L at a and g has derivative M at a, then x |-$>$ f(x) - g(x) has derivative L - M at a.  PROVEN, `modulo 0'.}\par
\vnbentry{\texttt{deriv-\allowbreak{}difference}}{\vnbstmt{\begin{array}{@{}l@{}} \forall f, g, a, l, m. \\ \quad \operatorname{is\text{-}diff\text{-}at}(f, a, l) \wedge \\ \quad \operatorname{is\text{-}diff\text{-}at}(g, a, m) \Rightarrow \\ \quad \quad \operatorname{is\text{-}diff\text{-}at}((\lambda x \in \mathbb{R}.\; (f(x) - g(x))), a, (l - m)) \end{array}}}

\vnbfile{theorem-library/mvt-aux-guarded.scm}{theorem-library/mvt-aux-guarded.scm}
\noindent{\small theorem-library/mvt-aux-guarded.scm -- the four MVT auxiliary-function supports, GUARDED and PROVEN.}\par
\vnbentry{\texttt{mvt-\allowbreak{}aux-\allowbreak{}diff}}{\vnbstmt{\begin{array}{@{}l@{}} \forall f, a, b, x, l. \\ \quad (a \in \mathbb{R}) \wedge \\ \quad (b \in \mathbb{R}) \wedge \\ \quad \operatorname{is\text{-}diff\text{-}at}(f, x, l) \Rightarrow \\ \quad \quad \operatorname{is\text{-}diff\text{-}at}((\lambda z \in \mathbb{R}.\; ((f(z) \cdot (b - a)) - (z \cdot (f(b) - f(a))))), x, \\ \quad \quad \quad ((l \cdot (b - a)) - (f(b) - f(a)))) \end{array}}}
\vnbentry{\texttt{mvt-\allowbreak{}aux-\allowbreak{}cont}}{\vnbstmt{\begin{array}{@{}l@{}} \forall f, a, b, x. \\ \quad (a \in \mathbb{R}) \wedge \\ \quad (b \in \mathbb{R}) \wedge \\ \quad \operatorname{is\text{-}continuous\text{-}at}(\operatorname{rr\text{-}ms}, \operatorname{rr\text{-}ms}, f, x) \Rightarrow \\ \quad \quad \operatorname{is\text{-}continuous\text{-}at}(\operatorname{rr\text{-}ms}, \operatorname{rr\text{-}ms}, \\ \quad \quad \quad (\lambda z \in \mathbb{R}.\; ((f(z) \cdot (b - a)) - (z \cdot (f(b) - f(a))))), x) \end{array}}}
\vnbentry{\texttt{gmvt-\allowbreak{}aux-\allowbreak{}diff}}{\vnbstmt{\begin{array}{@{}l@{}} \forall f, g, a, b, x, l, m. \\ \quad (a \in \mathbb{R}) \wedge \\ \quad (b \in \mathbb{R}) \wedge \\ \quad \operatorname{is\text{-}diff\text{-}at}(f, x, l) \wedge \\ \quad \operatorname{is\text{-}diff\text{-}at}(g, x, m) \Rightarrow \\ \quad \quad \operatorname{is\text{-}diff\text{-}at}((\lambda z \in \mathbb{R}.\; ((f(z) \cdot (g(b) - g(a))) - (g(z) \cdot (f(b) - f(a))))), \\ \quad \quad \quad x, ((l \cdot (g(b) - g(a))) - (m \cdot (f(b) - f(a))))) \end{array}}}
\vnbentry{\texttt{gmvt-\allowbreak{}aux-\allowbreak{}cont}}{\vnbstmt{\begin{array}{@{}l@{}} \forall f, g, a, b, x. \\ \quad (a \in \mathbb{R}) \wedge \\ \quad (b \in \mathbb{R}) \wedge \\ \quad \operatorname{is\text{-}continuous\text{-}at}(\operatorname{rr\text{-}ms}, \operatorname{rr\text{-}ms}, f, x) \wedge \\ \quad \operatorname{is\text{-}continuous\text{-}at}(\operatorname{rr\text{-}ms}, \operatorname{rr\text{-}ms}, g, x) \Rightarrow \\ \quad \quad \operatorname{is\text{-}continuous\text{-}at}(\operatorname{rr\text{-}ms}, \operatorname{rr\text{-}ms}, \\ \quad \quad \quad (\lambda z \in \mathbb{R}.\; ((f(z) \cdot (g(b) - g(a))) - (g(z) \cdot (f(b) - f(a))))), \\ \quad \quad \quad x) \end{array}}}

\vnbfile{theorem-library/interior-extremum-proof.scm}{theorem-library/interior-extremum-proof.scm}
\noindent{\small theorem-library/interior-extremum-proof.scm ==================================================================== Prop 2.10 (interior max): interior-max-deriv-zero, MACHINE-PROVEN.}\par
\vnbentry{\texttt{interior-\allowbreak{}max-\allowbreak{}deriv-\allowbreak{}zero}}{\vnbstmt{\begin{array}{@{}l@{}} \forall f, a, b, c, l. \\ \quad (f \in (\mathbb{R} \to \mathbb{R})) \wedge \\ \quad (a \in \mathbb{R}) \wedge \\ \quad (b \in \mathbb{R}) \wedge \\ \quad (c \in \mathbb{R}) \wedge \\ \quad (a < c) \wedge \\ \quad (c < b) \wedge \\ \quad \forall x \in \operatorname{ccint}(a, b).\; (f(x) \leq f(c)) \wedge \\ \quad \operatorname{is\text{-}diff\text{-}at}(f, c, l) \Rightarrow \\ \quad \quad (l = 0) \end{array}}}
\vnbentry{\texttt{interior-\allowbreak{}min-\allowbreak{}deriv-\allowbreak{}zero}}{\vnbstmt{\begin{array}{@{}l@{}} \forall f, a, b, c, l. \\ \quad (f \in (\mathbb{R} \to \mathbb{R})) \wedge \\ \quad (a \in \mathbb{R}) \wedge \\ \quad (b \in \mathbb{R}) \wedge \\ \quad (c \in \mathbb{R}) \wedge \\ \quad (a < c) \wedge \\ \quad (c < b) \wedge \\ \quad \forall x \in \operatorname{ccint}(a, b).\; (f(c) \leq f(x)) \wedge \\ \quad \operatorname{is\text{-}diff\text{-}at}(f, c, l) \Rightarrow \\ \quad \quad (l = 0) \end{array}}}

\vnbfile{theorem-library/rolle-proof.scm}{theorem-library/rolle-proof.scm}
\noindent{\small theorem-library/rolle-proof.scm -- Rolle's theorem, MACHINE-PROVEN. EVT gives argmax c / argmin d on [a,b]; if either is interior, the proven Fermat theorem (interior-max/min-deriv-zero) gives h'=0; else h is constant and a midpoint (rr-midpoint-between) is an interior max.}\par
\vnbentry{\texttt{rolle}}{\vnbstmt{\begin{array}{@{}l@{}} \forall h, a, b. \\ \quad (h \in (\mathbb{R} \to \mathbb{R})) \wedge \\ \quad (a \in \mathbb{R}) \wedge \\ \quad (b \in \mathbb{R}) \wedge \\ \quad (a < b) \wedge \\ \quad \forall x \in \operatorname{ccint}(a, b). \\ \quad \quad \operatorname{is\text{-}continuous\text{-}at}(\operatorname{rr\text{-}ms}, \operatorname{rr\text{-}ms}, h, x) \wedge \\ \quad \forall x \in \mathbb{R}. \\ \quad \quad ((a < x) \wedge (x < b)) \Rightarrow  \\ \quad \quad \quad \exists l.\; \operatorname{is\text{-}diff\text{-}at}(h, x, l) \wedge \\ \quad (h(a) = h(b)) \Rightarrow \\ \quad \quad \exists c \in \mathbb{R}. \\ \quad \quad \quad ((a < c) \wedge ((c < b) \wedge \operatorname{is\text{-}diff\text{-}at}(h, c, 0))) \end{array}}}

\vnbfile{theorem-library/mvt-proof.scm}{theorem-library/mvt-proof.scm}
\noindent{\small theorem-library/mvt-proof.scm -- the Mean Value Theorem, MACHINE-PROVEN. Strategy: apply Rolle's theorem to the auxiliary h(z) = f(z)*(b-a) - z*(f(b)-f(a)) which satisfies h(a)=h(b); Rolle gives an interior theta with h'(theta)=0,}\par
\vnbentry{\texttt{mvt}}{\vnbstmt{\begin{array}{@{}l@{}} \forall f, a, b. \\ \quad (f \in (\mathbb{R} \to \mathbb{R})) \wedge \\ \quad (a \in \mathbb{R}) \wedge \\ \quad (b \in \mathbb{R}) \wedge \\ \quad (a < b) \wedge \\ \quad \forall x \in \operatorname{ccint}(a, b). \\ \quad \quad \operatorname{is\text{-}continuous\text{-}at}(\operatorname{rr\text{-}ms}, \operatorname{rr\text{-}ms}, f, x) \wedge \\ \quad \forall x \in \mathbb{R}. \\ \quad \quad ((a < x) \wedge (x < b)) \Rightarrow  \\ \quad \quad \quad \exists l.\; \operatorname{is\text{-}diff\text{-}at}(f, x, l) \Rightarrow \\ \quad \quad \exists c \in \mathbb{R}. \\ \quad \quad \quad (a < c) \wedge \\ \quad \quad \quad (c < b) \wedge \\ \quad \quad \quad \exists l. \\ \quad \quad \quad \quad (\operatorname{is\text{-}diff\text{-}at}(f, c, l) \wedge ((l \cdot (b - a)) = (f(b) - f(a)))) \end{array}}}

\vnbfile{theorem-library/deriv-constant-proof.scm}{theorem-library/deriv-constant-proof.scm}
\noindent{\small theorem-library/deriv-constant-proof.scm -- Cor 2.15, MACHINE-PROVEN. f continuous on [a,b], f'=0 on (a,b)  =$>$  f constant on [a,b]. Strategy: for u$<$v in [a,b], the proven MVT on [u,v] gives an interior theta with f'(theta)(v-u)=f(v)-f(u); the zero-derivative hypothesis + derivative-}\par
\vnbentry{\texttt{deriv-\allowbreak{}zero-\allowbreak{}const-\allowbreak{}up}}{\vnbstmt{\begin{array}{@{}l@{}} \forall f, a, b, u, v. \\ \quad (f \in (\mathbb{R} \to \mathbb{R})) \wedge \\ \quad (a \in \mathbb{R}) \wedge \\ \quad (b \in \mathbb{R}) \wedge \\ \quad (a < b) \wedge \\ \quad \forall x \in \operatorname{ccint}(a, b). \\ \quad \quad \operatorname{is\text{-}continuous\text{-}at}(\operatorname{rr\text{-}ms}, \operatorname{rr\text{-}ms}, f, x) \wedge \\ \quad \forall x \in \mathbb{R}. \\ \quad \quad ((a < x) \wedge (x < b)) \Rightarrow  \\ \quad \quad \quad \operatorname{is\text{-}diff\text{-}at}(f, x, 0) \wedge \\ \quad (u \in \operatorname{ccint}(a, b)) \wedge \\ \quad (v \in \operatorname{ccint}(a, b)) \wedge \\ \quad (u < v) \Rightarrow \\ \quad \quad (f(u) = f(v)) \end{array}}}
\vnbentry{\texttt{deriv-\allowbreak{}zero-\allowbreak{}implies-\allowbreak{}constant}}{\vnbstmt{\begin{array}{@{}l@{}} \forall f, a, b, u, v. \\ \quad (f \in (\mathbb{R} \to \mathbb{R})) \wedge \\ \quad (a \in \mathbb{R}) \wedge \\ \quad (b \in \mathbb{R}) \wedge \\ \quad (a < b) \wedge \\ \quad \forall x \in \operatorname{ccint}(a, b). \\ \quad \quad \operatorname{is\text{-}continuous\text{-}at}(\operatorname{rr\text{-}ms}, \operatorname{rr\text{-}ms}, f, x) \wedge \\ \quad \forall x \in \mathbb{R}. \\ \quad \quad ((a < x) \wedge (x < b)) \Rightarrow  \\ \quad \quad \quad \operatorname{is\text{-}diff\text{-}at}(f, x, 0) \wedge \\ \quad (u \in \operatorname{ccint}(a, b)) \wedge \\ \quad (v \in \operatorname{ccint}(a, b)) \Rightarrow \\ \quad \quad (f(u) = f(v)) \end{array}}}

\vnbfile{theorem-library/mvt-bounds-proof.scm}{theorem-library/mvt-bounds-proof.scm}
\noindent{\small theorem-library/mvt-bounds-proof.scm -- Cor 2.14, MACHINE-PROVEN. f' $<$= M on (a,b)  =$>$  f(b)-f(a) $<$= M(b-a)     (mvt-upper-bound) m $<$= f' on (a,b)  =$>$  m(b-a) $<$= f(b)-f(a)     (mvt-lower-bound) Strategy: MVT gives theta with f(b)-f(a) = f'(theta)(b-a) = L(b-a); the}\par
\vnbentry{\texttt{mvt-\allowbreak{}upper-\allowbreak{}bound}}{\vnbstmt{\begin{array}{@{}l@{}} \forall f, a, b, m. \\ \quad (f \in (\mathbb{R} \to \mathbb{R})) \wedge \\ \quad (a \in \mathbb{R}) \wedge \\ \quad (b \in \mathbb{R}) \wedge \\ \quad (m \in \mathbb{R}) \wedge \\ \quad (a < b) \wedge \\ \quad \forall x \in \operatorname{ccint}(a, b). \\ \quad \quad \operatorname{is\text{-}continuous\text{-}at}(\operatorname{rr\text{-}ms}, \operatorname{rr\text{-}ms}, f, x) \wedge \\ \quad \forall x \in \mathbb{R}. \\ \quad \quad ((a < x) \wedge (x < b)) \Rightarrow  \\ \quad \quad \quad \exists l. \\ \quad \quad \quad \quad (\operatorname{is\text{-}diff\text{-}at}(f, x, l) \wedge (l \leq m)) \Rightarrow \\ \quad \quad ((f(b) - f(a)) \leq (m \cdot (b - a))) \end{array}}}
\vnbentry{\texttt{mvt-\allowbreak{}lower-\allowbreak{}bound}}{\vnbstmt{\begin{array}{@{}l@{}} \forall f, a, b, m. \\ \quad (f \in (\mathbb{R} \to \mathbb{R})) \wedge \\ \quad (a \in \mathbb{R}) \wedge \\ \quad (b \in \mathbb{R}) \wedge \\ \quad (m \in \mathbb{R}) \wedge \\ \quad (a < b) \wedge \\ \quad \forall x \in \operatorname{ccint}(a, b). \\ \quad \quad \operatorname{is\text{-}continuous\text{-}at}(\operatorname{rr\text{-}ms}, \operatorname{rr\text{-}ms}, f, x) \wedge \\ \quad \forall x \in \mathbb{R}. \\ \quad \quad ((a < x) \wedge (x < b)) \Rightarrow  \\ \quad \quad \quad \exists l. \\ \quad \quad \quad \quad (\operatorname{is\text{-}diff\text{-}at}(f, x, l) \wedge (m \leq l)) \Rightarrow \\ \quad \quad ((m \cdot (b - a)) \leq (f(b) - f(a))) \end{array}}}

\vnbfile{theorem-library/mvt-abs-bound.scm}{theorem-library/mvt-abs-bound.scm}
\noindent{\small mvt-abs-bound.scm -- the TWO-SIDED corollary of the Mean Value Theorem:}\par
\vnbentry{\texttt{mvt-\allowbreak{}abs-\allowbreak{}bound}}{\vnbstmt{\begin{array}{@{}l@{}} \forall f, a, b, m. \\ \quad (f \in (\mathbb{R} \to \mathbb{R})) \wedge \\ \quad (a \in \mathbb{R}) \wedge \\ \quad (b \in \mathbb{R}) \wedge \\ \quad (m \in \mathbb{R}) \wedge \\ \quad (a < b) \wedge \\ \quad \forall x \in \operatorname{ccint}(a, b). \\ \quad \quad \operatorname{is\text{-}continuous\text{-}at}(\operatorname{rr\text{-}ms}, \operatorname{rr\text{-}ms}, f, x) \wedge \\ \quad \forall x \in \mathbb{R}. \\ \quad \quad ((a < x) \wedge (x < b)) \Rightarrow  \\ \quad \quad \quad \exists l \in \mathbb{R}. \\ \quad \quad \quad \quad (\operatorname{is\text{-}diff\text{-}at}(f, x, l) \wedge (\lvert l \rvert \leq m)) \Rightarrow \\ \quad \quad (\lvert (f(b) - f(a)) \rvert \leq (m \cdot (b - a))) \end{array}}}

\vnbfile{theorem-library/continuity-local.scm}{theorem-library/continuity-local.scm}
\noindent{\small continuity-local.scm -- CONTINUITY AT A POINT READS ONLY A NEIGHBOURHOOD OF IT: a map agreeing with a continuous map on a BALL about the point, and saying nothing whatever elsewhere, is continuous at that point.  PROVEN.}\par
\vnbentry{\texttt{continuous-\allowbreak{}at-\allowbreak{}local}}{\vnbstmt{\begin{array}{@{}l@{}} \forall f, g \in (\mathbb{R} \to \mathbb{R}), t, o \in \mathbb{R}. \\ \quad (0 < o) \wedge \\ \quad \operatorname{is\text{-}continuous\text{-}at}(\operatorname{rr\text{-}ms}, \operatorname{rr\text{-}ms}, g, t) \wedge \\ \quad \forall x \in \mathbb{R}. \\ \quad \quad ((\lvert (x - t) \rvert \leq o) \Rightarrow (f(x) = g(x))) \Rightarrow \\ \quad \quad \operatorname{is\text{-}continuous\text{-}at}(\operatorname{rr\text{-}ms}, \operatorname{rr\text{-}ms}, f, t) \end{array}}}

\vnbfile{theorem-library/recip-star.scm}{theorem-library/recip-star.scm}
\noindent{\small recip-star.scm -- THE TOTALISED RECIPROCAL, and the reason it exists.}\par
\vnbentry{\texttt{recip-\allowbreak{}star-\allowbreak{}in-\allowbreak{}rr}}{\vnbstmt{\forall u \in \mathbb{R}.\; ({u}^{\ast\!-1} \in \mathbb{R})}}
\vnbentry{\texttt{recip-\allowbreak{}star-\allowbreak{}lam-\allowbreak{}in-\allowbreak{}fun}}{\vnbstmt{((\lambda z \in \mathbb{R}.\; {z}^{\ast\!-1}) \in (\mathbb{R} \to \mathbb{R}))}}
\vnbentry{\texttt{recip-\allowbreak{}star-\allowbreak{}value}}{\vnbstmt{\begin{array}{@{}l@{}} \forall u \in \mathbb{R}. \\ \quad (\neg (u = 0) \Rightarrow ({u}^{\ast\!-1} = \operatorname{recip}(u))) \end{array}}}
\vnbentry{\texttt{recip-\allowbreak{}star-\allowbreak{}zero}}{\vnbstmt{({0}^{\ast\!-1} = 0)}}
\vnbentry{\texttt{recip-\allowbreak{}star-\allowbreak{}inverse}}{\vnbstmt{\begin{array}{@{}l@{}} \forall u \in \mathbb{R}. \\ \quad (\neg (u = 0) \Rightarrow ((u \cdot {u}^{\ast\!-1}) = 1)) \end{array}}}
\vnbentry{\texttt{rr-\allowbreak{}half-\allowbreak{}ball-\allowbreak{}nonzero}}{\vnbstmt{\begin{array}{@{}l@{}} \forall c, m, y \in \mathbb{R}. \\ \quad ((0 < m) \wedge (((m + m) = \lvert c \rvert) \wedge (\lvert (y - c) \rvert \leq m))) \Rightarrow  \\ \quad \quad \neg (y = 0) \end{array}}}
\vnbentry{\texttt{nonzero-\allowbreak{}ident-\allowbreak{}lam-\allowbreak{}in-\allowbreak{}fun}}{\vnbstmt{\begin{array}{@{}l@{}} (\lambda w \in \mathbb{R}.\; \left(\text{if }(w = 0)\text{ then }1\text{ else }w\right)) \\ \in\; (\mathbb{R} \to \mathbb{R}) \end{array}}}
\vnbentry{\texttt{nonzero-\allowbreak{}ident-\allowbreak{}nowhere-\allowbreak{}zero}}{\vnbstmt{\begin{array}{@{}l@{}} \forall x \in \mathbb{R}. \\ \quad \neg ((\lambda w \in \mathbb{R}.\; \left(\text{if }(w = 0)\text{ then }1\text{ else }w\right))(x) = 0) \end{array}}}
\vnbentry{\texttt{nonzero-\allowbreak{}ident-\allowbreak{}continuous-\allowbreak{}at}}{\vnbstmt{\begin{array}{@{}l@{}} \forall c \in \mathbb{R}. \\ \quad \neg (c = 0) \Rightarrow  \\ \quad \quad \operatorname{is\text{-}continuous\text{-}at}(\operatorname{rr\text{-}ms}, \operatorname{rr\text{-}ms}, (\lambda w \in \mathbb{R}.\; \left(\text{if }(w = 0)\text{ then }1\text{ else }w\right)), c) \end{array}}}
\vnbentry{\texttt{recip-\allowbreak{}star-\allowbreak{}continuous-\allowbreak{}at}}{\vnbstmt{\begin{array}{@{}l@{}} \forall c \in \mathbb{R}. \\ \quad \neg (c = 0) \Rightarrow  \\ \quad \quad \operatorname{is\text{-}continuous\text{-}at}(\operatorname{rr\text{-}ms}, \operatorname{rr\text{-}ms}, (\lambda z \in \mathbb{R}.\; {z}^{\ast\!-1}), c) \end{array}}}
\vnbentry{\texttt{recip-\allowbreak{}star-\allowbreak{}continuous-\allowbreak{}on-\allowbreak{}ccint}}{\vnbstmt{\begin{array}{@{}l@{}} \forall a, b, x \in \mathbb{R}. \\ \quad ((0 < a) \wedge (x \in \operatorname{ccint}(a, b))) \Rightarrow  \\ \quad \quad \operatorname{is\text{-}continuous\text{-}at}(\operatorname{rr\text{-}ms}, \operatorname{rr\text{-}ms}, (\lambda z \in \mathbb{R}.\; {z}^{\ast\!-1}), x) \end{array}}}

\vnbfile{theorem-library/clamp.scm}{theorem-library/clamp.scm}
\noindent{\small clamp.scm -- CLAMP(a,b,x) = min(max(x,a),b), the retraction of the line onto the interval [a,b], and the ONE mechanism that dissolves the endpoint problem of docs/calculus.pdf Prop 4.16.}\par
\vnbentry{\texttt{clamp-\allowbreak{}in-\allowbreak{}rr}}{\vnbstmt{\forall a, b, x \in \mathbb{R}.\; (\operatorname{clamp}(a, b, x) \in \mathbb{R})}}
\vnbentry{\texttt{clamp-\allowbreak{}in-\allowbreak{}ccint}}{\vnbstmt{\begin{array}{@{}l@{}} \forall a, b, x \in \mathbb{R}. \\ \quad (a \leq b) \Rightarrow  \\ \quad \quad (\operatorname{clamp}(a, b, x) \in \operatorname{ccint}(a, b)) \end{array}}}
\vnbentry{\texttt{clamp-\allowbreak{}fixes}}{\vnbstmt{\begin{array}{@{}l@{}} \forall a, b, x \in \mathbb{R}. \\ \quad (((a \leq x) \wedge (x \leq b)) \Rightarrow (\operatorname{clamp}(a, b, x) = x)) \end{array}}}
\vnbentry{\texttt{rr-\allowbreak{}max-\allowbreak{}lipschitz}}{\vnbstmt{\begin{array}{@{}l@{}} \forall u, v, c \in \mathbb{R}. \\ \quad (\lvert (\operatorname{max}(u, c) - \operatorname{max}(v, c)) \rvert \leq \lvert (u - v) \rvert) \end{array}}}
\vnbentry{\texttt{rr-\allowbreak{}min-\allowbreak{}lipschitz}}{\vnbstmt{\begin{array}{@{}l@{}} \forall u, v, c \in \mathbb{R}. \\ \quad (\lvert (\operatorname{min}(u, c) - \operatorname{min}(v, c)) \rvert \leq \lvert (u - v) \rvert) \end{array}}}
\vnbentry{\texttt{clamp-\allowbreak{}lipschitz}}{\vnbstmt{\begin{array}{@{}l@{}} \forall a, b, y, x \in \mathbb{R}. \\ \quad \lvert (\operatorname{clamp}(a, b, y) - \operatorname{clamp}(a, b, x)) \rvert \\ \quad \leq\; \lvert (y - x) \rvert \end{array}}}
\vnbentry{\texttt{clamp-\allowbreak{}compose-\allowbreak{}lam-\allowbreak{}in-\allowbreak{}fun}}{\vnbstmt{\begin{array}{@{}l@{}} \forall f \in (\mathbb{R} \to \mathbb{R}), a, b \in \mathbb{R}. \\ \quad ((\lambda z \in \mathbb{R}.\; f(\operatorname{clamp}(a, b, z))) \in (\mathbb{R} \to \mathbb{R})) \end{array}}}
\vnbentry{\texttt{clamp-\allowbreak{}compose-\allowbreak{}fixes}}{\vnbstmt{\begin{array}{@{}l@{}} \forall f \in (\mathbb{R} \to \mathbb{R}), a, b, x \in \mathbb{R}. \\ \quad ((a \leq x) \wedge (x \leq b)) \Rightarrow  \\ \quad \quad ((\lambda z \in \mathbb{R}.\; f(\operatorname{clamp}(a, b, z)))(x) = f(x)) \end{array}}}
\vnbentry{\texttt{clamp-\allowbreak{}compose-\allowbreak{}continuous-\allowbreak{}at}}{\vnbstmt{\begin{array}{@{}l@{}} \forall f \in (\mathbb{R} \to \mathbb{R}), a, b, t \in \mathbb{R}. \\ \quad (a \leq b) \wedge \\ \quad \forall x \in \mathbb{R}. \\ \quad \quad (x \in \operatorname{ccint}(a, b)) \Rightarrow  \\ \quad \quad \quad \operatorname{is\text{-}continuous\text{-}at}(\operatorname{rr\text{-}ms}, \operatorname{rr\text{-}ms}, f, x) \Rightarrow \\ \quad \quad \operatorname{is\text{-}continuous\text{-}at}(\operatorname{rr\text{-}ms}, \operatorname{rr\text{-}ms}, (\lambda z \in \mathbb{R}.\; f(\operatorname{clamp}(a, b, z))), t) \end{array}}}

\vnbfile{theorem-library/seq-limit.scm}{theorem-library/seq-limit.scm}
\noindent{\small theorem-library/seq-limit.scm -- THE LIMIT OF A REAL SEQUENCE AS A TERM, and the two pieces of limit arithmetic the tree was missing.}\par
\vnbentry{\texttt{rr-\allowbreak{}limit-\allowbreak{}scale}}{\vnbstmt{\begin{array}{@{}l@{}} \forall c \in \mathbb{R}, f, h \in (\mathbb{N} \to \mathbb{R}), v \in \mathbb{R}. \\ \quad \forall j \in \mathbb{N}.\; (h(j) = (f(j) \cdot c)) \wedge \\ \quad \operatorname{converges\text{-}to}(\operatorname{rr\text{-}ms}, f, v) \Rightarrow \\ \quad \quad \operatorname{converges\text{-}to}(\operatorname{rr\text{-}ms}, h, (v \cdot c)) \end{array}}}
\vnbentry{\texttt{rr-\allowbreak{}limit-\allowbreak{}sub}}{\vnbstmt{\begin{array}{@{}l@{}} \forall f, g, h \in (\mathbb{N} \to \mathbb{R}), v, a \in \mathbb{R}. \\ \quad \forall j \in \mathbb{N}.\; (h(j) = (f(j) - g(j))) \wedge \\ \quad \operatorname{converges\text{-}to}(\operatorname{rr\text{-}ms}, f, v) \wedge \\ \quad \operatorname{converges\text{-}to}(\operatorname{rr\text{-}ms}, g, a) \Rightarrow \\ \quad \quad \operatorname{converges\text{-}to}(\operatorname{rr\text{-}ms}, h, (v - a)) \end{array}}}

\vnbfile{theorem-library/onesided-limits.scm}{theorem-library/onesided-limits.scm}
\noindent{\small onesided-limits.scm -- docs/calculus.pdf Ch 4 Sec 1. Continuity at x gives BOTH one-sided limits at x, with value f(x).  Both `modulo 0'.  The two drivers are the same text with IS-RIGHT-LIMIT swapped for IS-LEFT-LIMIT: the interval side is the only difference, and `ineq!'}\par
\vnbentry{\texttt{continuous-\allowbreak{}at-\allowbreak{}right-\allowbreak{}limit}}{\vnbstmt{\begin{array}{@{}l@{}} \forall f, x. \\ \quad (f \in (\mathbb{R} \to \mathbb{R})) \wedge \\ \quad (x \in \mathbb{R}) \wedge \\ \quad \operatorname{is\text{-}continuous\text{-}at}(\operatorname{rr\text{-}ms}, \operatorname{rr\text{-}ms}, f, x) \Rightarrow \\ \quad \quad \operatorname{is\text{-}right\text{-}limit}(f, x, f(x)) \end{array}}}
\vnbentry{\texttt{continuous-\allowbreak{}at-\allowbreak{}left-\allowbreak{}limit}}{\vnbstmt{\begin{array}{@{}l@{}} \forall f, x. \\ \quad (f \in (\mathbb{R} \to \mathbb{R})) \wedge \\ \quad (x \in \mathbb{R}) \wedge \\ \quad \operatorname{is\text{-}continuous\text{-}at}(\operatorname{rr\text{-}ms}, \operatorname{rr\text{-}ms}, f, x) \Rightarrow \\ \quad \quad \operatorname{is\text{-}left\text{-}limit}(f, x, f(x)) \end{array}}}

\vnbfile{theorem-library/recip-succ-small.scm}{theorem-library/recip-succ-small.scm}
\noindent{\small recip-succ-small.scm -- the eventual form of the Archimedean bound.}\par
\vnbentry{\texttt{recip-\allowbreak{}succ-\allowbreak{}small}}{\vnbstmt{\begin{array}{@{}l@{}} \forall d. \\ \quad \operatorname{pos\text{-}rr}(d) \Rightarrow  \\ \quad \quad \exists n \in \mathbb{N}. \forall m \in \mathbb{N}. \\ \quad \quad \quad ((n \leq m) \Rightarrow (\operatorname{recip}((m + 1)) \leq d)) \end{array}}}

\vnbfile{theorem-library/continuous-is-regulated.scm}{theorem-library/continuous-is-regulated.scm}
\noindent{\small continuous-is-regulated.scm -- docs/calculus.pdf Ch 4, the bridge from Section 1 to the finished Section 2/3 arc.}\par
\vnbentry{\texttt{continuous-\allowbreak{}is-\allowbreak{}regulated}}{\vnbstmt{\begin{array}{@{}l@{}} \forall f, a, b. \\ \quad (f \in (\mathbb{R} \to \mathbb{R})) \wedge \\ \quad (a \in \mathbb{R}) \wedge \\ \quad (b \in \mathbb{R}) \wedge \\ \quad (a < b) \wedge \\ \quad \forall x \in \operatorname{ccint}(a, b). \\ \quad \quad \operatorname{is\text{-}continuous\text{-}at}(\operatorname{rr\text{-}ms}, \operatorname{rr\text{-}ms}, f, x) \Rightarrow \\ \quad \quad \operatorname{is\text{-}regulated}(f, a, b) \end{array}}}

\vnbfile{theorem-library/cauchy-criterion-left.scm}{theorem-library/cauchy-criterion-left.scm}
\noindent{\small cauchy-criterion-left.scm -- docs/calculus.pdf Ch 4 Sec 1.}\par
\vnbentry{\texttt{cauchy-\allowbreak{}criterion-\allowbreak{}left}}{\vnbstmt{\begin{array}{@{}l@{}} \forall f, x. \\ \quad (f \in (\mathbb{R} \to \mathbb{R})) \wedge \\ \quad (x \in \mathbb{R}) \wedge \\ \quad \forall s. \\ \quad \quad \operatorname{pos\text{-}rr}(s) \Rightarrow  \\ \quad \quad \quad \exists a. \\ \quad \quad \quad \quad \operatorname{pos\text{-}rr}(a) \wedge \\ \quad \quad \quad \quad \forall b, t. \\ \quad \quad \quad \quad \quad (b \in \mathbb{R}) \wedge \\ \quad \quad \quad \quad \quad ((x - a) \leq b) \wedge \\ \quad \quad \quad \quad \quad (b < x) \wedge \\ \quad \quad \quad \quad \quad (t \in \mathbb{R}) \wedge \\ \quad \quad \quad \quad \quad ((x - a) \leq t) \wedge \\ \quad \quad \quad \quad \quad (t < x) \Rightarrow \\ \quad \quad \quad \quad \quad \quad (\lvert (f(b) - f(t)) \rvert \leq s) \Rightarrow \\ \quad \quad \exists l.\; \operatorname{is\text{-}left\text{-}limit}(f, x, l) \end{array}}}

\vnbfile{theorem-library/cauchy-criterion-right.scm}{theorem-library/cauchy-criterion-right.scm}
\noindent{\small cauchy-criterion-right.scm -- docs/calculus.pdf Ch 4 Sec 1.}\par
\vnbentry{\texttt{cauchy-\allowbreak{}criterion-\allowbreak{}right}}{\vnbstmt{\begin{array}{@{}l@{}} \forall f, x. \\ \quad (f \in (\mathbb{R} \to \mathbb{R})) \wedge \\ \quad (x \in \mathbb{R}) \wedge \\ \quad \forall s. \\ \quad \quad \operatorname{pos\text{-}rr}(s) \Rightarrow  \\ \quad \quad \quad \exists a. \\ \quad \quad \quad \quad \operatorname{pos\text{-}rr}(a) \wedge \\ \quad \quad \quad \quad \forall b, t. \\ \quad \quad \quad \quad \quad (b \in \mathbb{R}) \wedge \\ \quad \quad \quad \quad \quad (x < b) \wedge \\ \quad \quad \quad \quad \quad (b \leq (x + a)) \wedge \\ \quad \quad \quad \quad \quad (t \in \mathbb{R}) \wedge \\ \quad \quad \quad \quad \quad (x < t) \wedge \\ \quad \quad \quad \quad \quad (t \leq (x + a)) \Rightarrow \\ \quad \quad \quad \quad \quad \quad (\lvert (f(b) - f(t)) \rvert \leq s) \Rightarrow \\ \quad \quad \exists l.\; \operatorname{is\text{-}right\text{-}limit}(f, x, l) \end{array}}}

\vnbfile{theorem-library/partition-locates-point.scm}{theorem-library/partition-locates-point.scm}
\noindent{\small partition-locates-point.scm -- docs/calculus.pdf Ch 4 Sec 1.}\par
\vnbentry{\texttt{partition-\allowbreak{}locates-\allowbreak{}point}}{\vnbstmt{\begin{array}{@{}l@{}} \forall n \in \mathbb{N}, p, a, b, x. \\ \quad (\operatorname{is\text{-}partition}(p, n, a, b) \wedge (x \in \operatorname{ccint}(a, b))) \Rightarrow  \\ \quad \quad \exists i \in \operatorname{interval}(0, n). \\ \quad \quad \quad ((i < n) \wedge ((p(i) \leq x) \wedge (x \leq p((i + 1))))) \end{array}}}

\vnbfile{theorem-library/uniform-limit-regulated.scm}{theorem-library/uniform-limit-regulated.scm}
\noindent{\small uniform-limit-regulated.scm -- a uniform limit of regulated functions is regulated.  docs/calculus.pdf Prop 4.5, the sufficiency half.}\par
\vnbentry{\texttt{uniform-\allowbreak{}limit-\allowbreak{}of-\allowbreak{}regulated}}{\vnbstmt{\begin{array}{@{}l@{}} \forall m, g, a, b. \\ \quad \operatorname{is\text{-}unif\text{-}limit\text{-}on}(m, g, a, b) \wedge \\ \quad \forall k \in \mathbb{N}.\; \operatorname{is\text{-}regulated}(m(k), a, b) \Rightarrow \\ \quad \quad \operatorname{is\text{-}regulated}(g, a, b) \end{array}}}

\vnbfile{theorem-library/uniform-limit-local.scm}{theorem-library/uniform-limit-local.scm}
\noindent{\small theorem-library/uniform-limit-local.scm -- A UNIFORM LIMIT IS CONTINUOUS AT A POINT, from uniform convergence on a BALL about that point only.}\par
\vnbentry{\texttt{uniform-\allowbreak{}limit-\allowbreak{}continuous-\allowbreak{}at}}{\vnbstmt{\begin{array}{@{}l@{}} \forall m \in (\mathbb{N} \to (\mathbb{R} \to \mathbb{R})), g \in (\mathbb{R} \to \mathbb{R}), t, o \in \mathbb{R}. \\ \quad (0 < o) \wedge \\ \quad \forall k \in \mathbb{N}. \\ \quad \quad \operatorname{is\text{-}continuous\text{-}at}(\operatorname{rr\text{-}ms}, \operatorname{rr\text{-}ms}, m(k), t) \wedge \\ \quad \forall s. \\ \quad \quad \operatorname{pos\text{-}rr}(s) \Rightarrow  \\ \quad \quad \quad \exists p \in \mathbb{N}. \forall k \in \mathbb{N}, x \in \mathbb{R}. \\ \quad \quad \quad \quad ((p \leq k) \wedge (\lvert (x - t) \rvert \leq o)) \Rightarrow  \\ \quad \quad \quad \quad \quad (\lvert (g(x) - m(k)(x)) \rvert \leq s) \Rightarrow \\ \quad \quad \operatorname{is\text{-}continuous\text{-}at}(\operatorname{rr\text{-}ms}, \operatorname{rr\text{-}ms}, g, t) \end{array}}}

\vnbfile{theorem-library/diff-at-local.scm}{theorem-library/diff-at-local.scm}
\noindent{\small theorem-library/diff-at-local.scm -- DIFFERENTIABILITY IS LOCAL, and the Caratheodory transfer of a POINTWISE LIMIT.}\par
\vnbentry{\texttt{diff-\allowbreak{}at-\allowbreak{}local}}{\vnbstmt{\begin{array}{@{}l@{}} \forall f, i \in (\mathbb{R} \to \mathbb{R}), t, v, o \in \mathbb{R}. \\ \quad (0 < o) \wedge \\ \quad \operatorname{is\text{-}continuous\text{-}at}(\operatorname{rr\text{-}ms}, \operatorname{rr\text{-}ms}, i, t) \wedge \\ \quad (i(t) = v) \wedge \\ \quad \forall x \in \mathbb{R}. \\ \quad \quad (\lvert (x - t) \rvert \leq o) \Rightarrow  \\ \quad \quad \quad ((f(x) - f(t)) = (i(x) \cdot (x - t))) \Rightarrow \\ \quad \quad \operatorname{is\text{-}diff\text{-}at}(f, t, v) \end{array}}}
\vnbentry{\texttt{diff-\allowbreak{}at-\allowbreak{}ptwise-\allowbreak{}limit}}{\vnbstmt{\begin{array}{@{}l@{}} \forall m, a \in (\mathbb{N} \to (\mathbb{R} \to \mathbb{R})), f, i \in (\mathbb{R} \to \mathbb{R}), t, o \in \mathbb{R}. \\ \quad (0 < o) \wedge \\ \quad \operatorname{is\text{-}continuous\text{-}at}(\operatorname{rr\text{-}ms}, \operatorname{rr\text{-}ms}, i, t) \wedge \\ \quad \forall k \in \mathbb{N}, y \in \mathbb{R}. \\ \quad \quad (\lvert (y - t) \rvert \leq o) \Rightarrow  \\ \quad \quad \quad ((m(k)(y) - m(k)(t)) = (a(k)(y) \cdot (y - t))) \wedge \\ \quad \forall y \in \mathbb{R}. \\ \quad \quad (\lvert (y - t) \rvert \leq o) \Rightarrow  \\ \quad \quad \quad \operatorname{converges\text{-}to}(\operatorname{rr\text{-}ms}, (\lambda b \in \mathbb{N}.\; m(b)(y)), f(y)) \wedge \\ \quad \forall y \in \mathbb{R}. \\ \quad \quad (\lvert (y - t) \rvert \leq o) \Rightarrow  \\ \quad \quad \quad \operatorname{converges\text{-}to}(\operatorname{rr\text{-}ms}, (\lambda b \in \mathbb{N}.\; a(b)(y)), i(y)) \Rightarrow \\ \quad \quad \operatorname{is\text{-}diff\text{-}at}(f, t, i(t)) \end{array}}}

\vnbfile{theorem-library/generalized-mvt-proof.scm}{theorem-library/generalized-mvt-proof.scm}
\noindent{\small theorem-library/generalized-mvt-proof.scm -- Cauchy MVT (Thm 2.11), PROVEN. L*(g(b)-g(a)) = M*(f(b)-f(a))  for some interior theta with f'=L, g'=M. Strategy: apply Rolle to the auxiliary h(z) = f(z)*(g(b)-g(a)) - g(z)*(f(b)-f(a))}\par
\vnbentry{\texttt{generalized-\allowbreak{}mvt}}{\vnbstmt{\begin{array}{@{}l@{}} \forall f, g, a, b. \\ \quad (f \in (\mathbb{R} \to \mathbb{R})) \wedge \\ \quad (g \in (\mathbb{R} \to \mathbb{R})) \wedge \\ \quad (a \in \mathbb{R}) \wedge \\ \quad (b \in \mathbb{R}) \wedge \\ \quad (a < b) \wedge \\ \quad \forall x \in \operatorname{ccint}(a, b). \\ \quad \quad \operatorname{is\text{-}continuous\text{-}at}(\operatorname{rr\text{-}ms}, \operatorname{rr\text{-}ms}, f, x) \wedge \\ \quad \quad \operatorname{is\text{-}continuous\text{-}at}(\operatorname{rr\text{-}ms}, \operatorname{rr\text{-}ms}, g, x) \wedge \\ \quad \forall x \in \mathbb{R}. \\ \quad \quad ((a < x) \wedge (x < b)) \Rightarrow  \\ \quad \quad \quad \exists l.\; \operatorname{is\text{-}diff\text{-}at}(f, x, l) \wedge \\ \quad \quad \quad \exists m.\; \operatorname{is\text{-}diff\text{-}at}(g, x, m) \Rightarrow \\ \quad \quad \exists c \in \mathbb{R}. \\ \quad \quad \quad (a < c) \wedge \\ \quad \quad \quad (c < b) \wedge \\ \quad \quad \quad \exists l, m. \\ \quad \quad \quad \quad \operatorname{is\text{-}diff\text{-}at}(f, c, l) \wedge \\ \quad \quad \quad \quad \operatorname{is\text{-}diff\text{-}at}(g, c, m) \wedge \\ \quad \quad \quad \quad ((l \cdot (g(b) - g(a))) = (m \cdot (f(b) - f(a)))) \end{array}}}

\vnbfile{theorem-library/deriv-monotone-proof.scm}{theorem-library/deriv-monotone-proof.scm}
\noindent{\small theorem-library/deriv-monotone-proof.scm -- increasing function theorem. f continuous on [a,b], f' $>$ 0 on (a,b)  =$>$  f strictly increasing on [a,b]. Strategy: for u$<$v in [a,b], MVT on [u,v] gives theta with f(v)-f(u) = f'(theta)(v-u);  f'(theta) $>$ 0 and v-u $>$ 0, so f(v)-f(u) $>$ 0.}\par
\vnbentry{\texttt{deriv-\allowbreak{}pos-\allowbreak{}strictly-\allowbreak{}increasing}}{\vnbstmt{\begin{array}{@{}l@{}} \forall f, a, b, u, v. \\ \quad (f \in (\mathbb{R} \to \mathbb{R})) \wedge \\ \quad (a \in \mathbb{R}) \wedge \\ \quad (b \in \mathbb{R}) \wedge \\ \quad (a < b) \wedge \\ \quad \forall x \in \operatorname{ccint}(a, b). \\ \quad \quad \operatorname{is\text{-}continuous\text{-}at}(\operatorname{rr\text{-}ms}, \operatorname{rr\text{-}ms}, f, x) \wedge \\ \quad \forall x \in \mathbb{R}. \\ \quad \quad ((a < x) \wedge (x < b)) \Rightarrow  \\ \quad \quad \quad \exists l.\; (\operatorname{is\text{-}diff\text{-}at}(f, x, l) \wedge (0 < l)) \wedge \\ \quad (u \in \operatorname{ccint}(a, b)) \wedge \\ \quad (v \in \operatorname{ccint}(a, b)) \wedge \\ \quad (u < v) \Rightarrow \\ \quad \quad (f(u) < f(v)) \end{array}}}

\vnbfile{theorem-library/taylor-proof.scm}{theorem-library/taylor-proof.scm}
\noindent{\small RETIRED 2026-09-17 (proven): power-zero-base (was an add-to-pss here) -- theorem-library/rake-combinatorics.scm theorem-library/taylor-proof.scm -- Taylor's theorem, Lagrange remainder. f\textasciicircum{}(n) C\textasciicircum{}0 on [a,x], f\textasciicircum{}(n+1) exists on (a,x), a$<$x  =$>$  there is theta in (a,x) with  (n+1)! * R\_n(x) = f\textasciicircum{}(n+1)(theta) * (x-a)\textasciicircum{}(n+1),}\par
\vnbentry{\texttt{rr-\allowbreak{}recip-\allowbreak{}one}}{\vnbstmt{(\operatorname{recip}(1) = 1)}}
\vnbentry{\texttt{rr-\allowbreak{}zero-\allowbreak{}mul}}{\vnbstmt{\forall v \in \mathbb{R}.\; ((0 \cdot v) = 0)}}
\vnbentry{\texttt{rr-\allowbreak{}mul-\allowbreak{}zero-\allowbreak{}mid}}{\vnbstmt{\forall u, v \in \mathbb{R}.\; (((u \cdot 0) \cdot v) = 0)}}
\vnbentry{\texttt{rr-\allowbreak{}mul-\allowbreak{}one-\allowbreak{}mid}}{\vnbstmt{\forall u \in \mathbb{R}.\; (((u \cdot 1) \cdot 1) = u)}}
\vnbentry{\texttt{factorial-\allowbreak{}real-\allowbreak{}pos}}{\vnbstmt{\begin{array}{@{}l@{}} \forall n \in \mathbb{N}. \\ \quad ((\operatorname{factorial}(n) \in \mathbb{R}) \wedge (0 < \operatorname{factorial}(n))) \end{array}}}
\vnbentry{\texttt{recip-\allowbreak{}factorial-\allowbreak{}in-\allowbreak{}rr}}{\vnbstmt{\begin{array}{@{}l@{}} \forall n \in \mathbb{N}. \\ \quad (\operatorname{recip}(\operatorname{factorial}(n)) \in \mathbb{R}) \end{array}}}
\vnbentry{\texttt{taylor-\allowbreak{}center-\allowbreak{}term-\allowbreak{}at-\allowbreak{}zero}}{\vnbstmt{\begin{array}{@{}l@{}} \forall f.; \forall x \in \mathbb{R}. \\ \quad (f(x) \in \mathbb{R}) \Rightarrow  \\ \quad \quad (\lambda k \in \mathbb{N}.\; ((\operatorname{nth\text{-}deriv}(f, k)(x) \cdot \operatorname{power}((x - x), k)) \cdot \operatorname{recip}(\operatorname{factorial}(k))))(0) \\ \quad \quad =\; f(x) \end{array}}}
\vnbentry{\texttt{taylor-\allowbreak{}center-\allowbreak{}term-\allowbreak{}at-\allowbreak{}succ}}{\vnbstmt{\begin{array}{@{}l@{}} \forall f.; \forall x \in \mathbb{R}, m \in \mathbb{N}. \\ \quad (\operatorname{nth\text{-}deriv}(f, (m + 1))(x) \in \mathbb{R}) \Rightarrow  \\ \quad \quad (\lambda k \in \mathbb{N}.\; ((\operatorname{nth\text{-}deriv}(f, k)(x) \cdot \operatorname{power}((x - x), k)) \cdot \operatorname{recip}(\operatorname{factorial}(k))))((m + 1)) \\ \quad \quad =\; 0 \end{array}}}
\vnbentry{\texttt{taylor-\allowbreak{}center-\allowbreak{}partial-\allowbreak{}sum}}{\vnbstmt{\begin{array}{@{}l@{}} \forall n \in \mathbb{N}, f.; \forall x \in \mathbb{R}. \\ \quad \forall k \in \mathbb{N}. \\ \quad \quad (k \leq n) \Rightarrow  \\ \quad \quad \quad (\operatorname{nth\text{-}deriv}(f, k)(x) \in \mathbb{R}) \Rightarrow \\ \quad \quad \operatorname{series\text{-}partial\text{-}sum}((\lambda k \in \mathbb{N}.\; \\ \quad \quad \quad \quad ((\operatorname{nth\text{-}deriv}(f, k)(x) \cdot \operatorname{power}((x - x), k)) \cdot \operatorname{recip}(\operatorname{factorial}(k)))), \\ \quad \quad \quad (n + 1)) \\ \quad \quad =\; f(x) \end{array}}}
\vnbentry{\texttt{taylor-\allowbreak{}poly-\allowbreak{}at-\allowbreak{}center}}{\vnbstmt{\begin{array}{@{}l@{}} \forall f, x, n. \\ \quad (x \in \mathbb{R}) \wedge \\ \quad (n \in \mathbb{N}) \wedge \\ \quad \forall k \in \mathbb{N}. \\ \quad \quad (k \leq n) \Rightarrow  \\ \quad \quad \quad (\operatorname{nth\text{-}deriv}(f, k)(x) \in \mathbb{R}) \Rightarrow \\ \quad \quad (\operatorname{taylor\text{-}poly}(f, x, n, x) = f(x)) \end{array}}}
\vnbentry{\texttt{taylor-\allowbreak{}derivs-\allowbreak{}in-\allowbreak{}fun}}{\vnbstmt{\begin{array}{@{}l@{}} \forall f, a, x, n.; \forall k \in \mathbb{N}. \\ \quad \operatorname{taylor\text{-}differentiable}(f, a, x, n) \wedge \\ \quad (a \in \mathbb{R}) \wedge \\ \quad (x \in \mathbb{R}) \wedge \\ \quad (a < x) \wedge \\ \quad (k \leq n) \Rightarrow \\ \quad \quad (\operatorname{nth\text{-}deriv}(f, k) \in (\mathbb{R} \to \mathbb{R})) \end{array}}}
\vnbentry{\texttt{taylor-\allowbreak{}derivs-\allowbreak{}values-\allowbreak{}real}}{\vnbstmt{\begin{array}{@{}l@{}} \forall f, n.; \forall t \in \mathbb{R}, k \in \mathbb{N}. \\ \quad \forall k \in \mathbb{N}. \\ \quad \quad (k \leq n) \Rightarrow  \\ \quad \quad \quad (\operatorname{nth\text{-}deriv}(f, k) \in (\mathbb{R} \to \mathbb{R})) \wedge \\ \quad (k \leq n) \Rightarrow \\ \quad \quad (\operatorname{nth\text{-}deriv}(f, k)(t) \in \mathbb{R}) \end{array}}}
\vnbentry{\texttt{taylor-\allowbreak{}poly-\allowbreak{}in-\allowbreak{}rr-\allowbreak{}at}}{\vnbstmt{\begin{array}{@{}l@{}} \forall n \in \mathbb{N}, f.; \forall a, x \in \mathbb{R}. \\ \quad \forall k \in \mathbb{N}. \\ \quad \quad (k \leq n) \Rightarrow  \\ \quad \quad \quad (\operatorname{nth\text{-}deriv}(f, k)(a) \in \mathbb{R}) \Rightarrow \\ \quad \quad \operatorname{series\text{-}partial\text{-}sum}((\lambda k \in \mathbb{N}.\; \\ \quad \quad \quad \quad ((\operatorname{nth\text{-}deriv}(f, k)(a) \cdot \operatorname{power}((x - a), k)) \cdot \operatorname{recip}(\operatorname{factorial}(k)))), \\ \quad \quad \quad (n + 1)) \\ \quad \quad \in\; \mathbb{R} \end{array}}}
\vnbentry{\texttt{taylor-\allowbreak{}poly-\allowbreak{}in-\allowbreak{}rr}}{\vnbstmt{\begin{array}{@{}l@{}} \forall f \in (\mathbb{R} \to \mathbb{R}), a \in \mathbb{R}, n \in \mathbb{N}, x \in \mathbb{R}. \\ \quad \forall k \in \mathbb{N}. \\ \quad \quad (k \leq n) \Rightarrow  \\ \quad \quad \quad (\operatorname{nth\text{-}deriv}(f, k) \in (\mathbb{R} \to \mathbb{R})) \Rightarrow \\ \quad \quad (\operatorname{taylor\text{-}poly}(f, a, n, x) \in \mathbb{R}) \end{array}}}
\vnbentry{\texttt{taylor-\allowbreak{}g-\allowbreak{}in-\allowbreak{}fun}}{\vnbstmt{\begin{array}{@{}l@{}} \forall n \in \mathbb{N}, f \in (\mathbb{R} \to \mathbb{R}), x \in \mathbb{R}. \\ \quad \forall k \in \mathbb{N}. \\ \quad \quad (k \leq n) \Rightarrow  \\ \quad \quad \quad (\operatorname{nth\text{-}deriv}(f, k) \in (\mathbb{R} \to \mathbb{R})) \Rightarrow \\ \quad \quad (\lambda z \in \mathbb{R}.\; (f(x) - \operatorname{taylor\text{-}poly}(f, z, n, x))) \\ \quad \quad \in\; (\mathbb{R} \to \mathbb{R}) \end{array}}}
\vnbentry{\texttt{taylor-\allowbreak{}h-\allowbreak{}in-\allowbreak{}fun}}{\vnbstmt{\begin{array}{@{}l@{}} \forall n \in \mathbb{N}, x \in \mathbb{R}. \\ \quad (\lambda z \in \mathbb{R}.\; \operatorname{power}((x - z), (n + 1))) \\ \quad \in\; (\mathbb{R} \to \mathbb{R}) \end{array}}}
\vnbentry{\texttt{taylor-\allowbreak{}h-\allowbreak{}at-\allowbreak{}a}}{\vnbstmt{\begin{array}{@{}l@{}} \forall a, x, n. \\ \quad ((a \in \mathbb{R}) \wedge ((x \in \mathbb{R}) \wedge (n \in \mathbb{N}))) \Rightarrow  \\ \quad \quad (\lambda z \in \mathbb{R}.\; \operatorname{power}((x - z), (n + 1)))(a) \\ \quad \quad =\; \operatorname{power}((x - a), (n + 1)) \end{array}}}
\vnbentry{\texttt{taylor-\allowbreak{}h-\allowbreak{}at-\allowbreak{}x}}{\vnbstmt{\begin{array}{@{}l@{}} \forall x \in \mathbb{R}, n \in \mathbb{N}. \\ \quad ((\lambda z \in \mathbb{R}.\; \operatorname{power}((x - z), (n + 1)))(x) = 0) \end{array}}}
\vnbentry{\texttt{taylor-\allowbreak{}g-\allowbreak{}at-\allowbreak{}a}}{\vnbstmt{\begin{array}{@{}l@{}} \forall f, a, x, n. \\ \quad (f \in (\mathbb{R} \to \mathbb{R})) \wedge \\ \quad (a \in \mathbb{R}) \wedge \\ \quad (x \in \mathbb{R}) \wedge \\ \quad (n \in \mathbb{N}) \wedge \\ \quad \forall k \in \mathbb{N}. \\ \quad \quad (k \leq n) \Rightarrow  \\ \quad \quad \quad (\operatorname{nth\text{-}deriv}(f, k) \in (\mathbb{R} \to \mathbb{R})) \Rightarrow \\ \quad \quad (\lambda z \in \mathbb{R}.\; (f(x) - \operatorname{taylor\text{-}poly}(f, z, n, x)))(a) \\ \quad \quad =\; (f(x) - \operatorname{taylor\text{-}poly}(f, a, n, x)) \end{array}}}
\vnbentry{\texttt{taylor-\allowbreak{}g-\allowbreak{}at-\allowbreak{}x}}{\vnbstmt{\begin{array}{@{}l@{}} \forall f, x.; \forall n \in \mathbb{N}. \\ \quad (f \in (\mathbb{R} \to \mathbb{R})) \wedge \\ \quad (x \in \mathbb{R}) \wedge \\ \quad \forall k \in \mathbb{N}. \\ \quad \quad (k \leq n) \Rightarrow  \\ \quad \quad \quad (\operatorname{nth\text{-}deriv}(f, k) \in (\mathbb{R} \to \mathbb{R})) \Rightarrow \\ \quad \quad ((\lambda z \in \mathbb{R}.\; (f(x) - \operatorname{taylor\text{-}poly}(f, z, n, x)))(x) = 0) \end{array}}}
\vnbentry{\texttt{taylor-\allowbreak{}deriv-\allowbreak{}real}}{\vnbstmt{\begin{array}{@{}l@{}} \forall f, a, x, n, t. \\ \quad \operatorname{taylor\text{-}differentiable}(f, a, x, n) \wedge \\ \quad (a < t) \wedge \\ \quad (t < x) \wedge \\ \quad (n \in \mathbb{N}) \Rightarrow \\ \quad \quad (\operatorname{nth\text{-}deriv}(f, (n + 1))(t) \in \mathbb{R}) \end{array}}}
\vnbentry{\texttt{tg-\allowbreak{}neg-\allowbreak{}scale}}{\vnbstmt{\begin{array}{@{}l@{}} \forall s, p \in \mathbb{R}. \\ \quad ((0 - (s \cdot p)) = ((s \cdot p) \cdot (0 + -1))) \end{array}}}
\vnbentry{\texttt{taylor-\allowbreak{}h-\allowbreak{}diff}}{\vnbstmt{\begin{array}{@{}l@{}} \forall x, n, t. \\ \quad ((x \in \mathbb{R}) \wedge ((t \in \mathbb{R}) \wedge (n \in \mathbb{N}))) \Rightarrow  \\ \quad \quad \operatorname{is\text{-}diff\text{-}at}((\lambda z \in \mathbb{R}.\; \operatorname{power}((x - z), (n + 1))), t, \\ \quad \quad \quad (0 - ((n + 1) \cdot \operatorname{power}((x - t), n)))) \end{array}}}
\vnbentry{\texttt{tgd-\allowbreak{}recip-\allowbreak{}factorial-\allowbreak{}succ}}{\vnbstmt{\begin{array}{@{}l@{}} \forall m \in \mathbb{N}. \\ \quad \operatorname{recip}(\operatorname{factorial}(m)) \\ \quad =\; ((m + 1) \cdot \operatorname{recip}(\operatorname{factorial}((m + 1)))) \end{array}}}
\vnbentry{\texttt{tgd-\allowbreak{}sub-\allowbreak{}diff}}{\vnbstmt{\begin{array}{@{}l@{}} \forall x, t \in \mathbb{R}. \\ \quad \operatorname{is\text{-}diff\text{-}at}((\lambda z \in \mathbb{R}.\; (x - z)), t, -1) \end{array}}}
\vnbentry{\texttt{tgd-\allowbreak{}power-\allowbreak{}sub-\allowbreak{}diff}}{\vnbstmt{\begin{array}{@{}l@{}} \forall m \in \mathbb{N}, x, t \in \mathbb{R}. \\ \quad \operatorname{is\text{-}diff\text{-}at}((\lambda z \in \mathbb{R}.\; \operatorname{power}((x - z), (m + 1))), t, \\ \quad \quad (((m + 1) \cdot \operatorname{power}((x - t), m)) \cdot -1)) \end{array}}}
\vnbentry{\texttt{tgd-\allowbreak{}step-\allowbreak{}identity}}{\vnbstmt{\begin{array}{@{}l@{}} \forall s, r, a, b, p, q \in \mathbb{R}. \\ \quad ((((b \cdot q) + (a \cdot ((s \cdot p) \cdot -1))) \cdot r) + ((a \cdot q) \cdot 0)) \\ \quad =\; (((r \cdot b) \cdot q) - (((s \cdot r) \cdot a) \cdot p)) \end{array}}}
\vnbentry{\texttt{tgd-\allowbreak{}telescope-\allowbreak{}identity}}{\vnbstmt{\forall u, v \in \mathbb{R}.\; (v = (u + (v - u)))}}
\vnbentry{\texttt{tgd-\allowbreak{}term-\allowbreak{}diff}}{\vnbstmt{\begin{array}{@{}l@{}} \forall m \in \mathbb{N}, f.; \forall x, t \in \mathbb{R}. \\ \quad \operatorname{is\text{-}diff\text{-}at}(\operatorname{nth\text{-}deriv}(f, (m + 1)), t, \operatorname{nth\text{-}deriv}(f, ((m + 1) + 1))(t)) \Rightarrow \\ \quad \quad \operatorname{is\text{-}diff\text{-}at}((\lambda z \in \mathbb{R}.\; \\ \quad \quad \quad \quad ((\operatorname{nth\text{-}deriv}(f, (m + 1))(z) \cdot \operatorname{power}((x - z), (m + 1))) \cdot \operatorname{recip}(\operatorname{factorial}((m + 1))))), \\ \quad \quad \quad t, \\ \quad \quad \quad (((\operatorname{recip}(\operatorname{factorial}((m + 1))) \cdot \operatorname{nth\text{-}deriv}(f, ((m + 1) + 1))(t)) \cdot \operatorname{power}((x - t), (m + 1))) - ((\operatorname{recip}(\operatorname{factorial}(m)) \cdot \operatorname{nth\text{-}deriv}(f, (m + 1))(t)) \cdot \operatorname{power}((x - t), m)))) \end{array}}}
\vnbentry{\texttt{tgd-\allowbreak{}base-\allowbreak{}sum-\allowbreak{}value}}{\vnbstmt{\begin{array}{@{}l@{}} \forall f \in (\mathbb{R} \to \mathbb{R}), x, z \in \mathbb{R}. \\ \quad (\operatorname{series\text{-}partial\text{-}sum}((\lambda k \in \mathbb{N}.\; ((\operatorname{nth\text{-}deriv}(f, k)(z) \cdot \operatorname{power}((x - z), k)) \cdot \operatorname{recip}(\operatorname{factorial}(k)))), (0 + 1)) \simeq f(z)) \end{array}}}
\vnbentry{\texttt{tgd-\allowbreak{}sum-\allowbreak{}diff}}{\vnbstmt{\begin{array}{@{}l@{}} \forall m \in \mathbb{N}, f.; \forall x, t \in \mathbb{R}. \\ \quad \forall k \in \mathbb{N}. \\ \quad \quad (k \leq m) \Rightarrow  \\ \quad \quad \quad \operatorname{is\text{-}diff\text{-}at}(\operatorname{nth\text{-}deriv}(f, k), t, \operatorname{nth\text{-}deriv}(f, (k + 1))(t)) \Rightarrow \\ \quad \quad \operatorname{is\text{-}diff\text{-}at}((\lambda z \in \mathbb{R}.\; \\ \quad \quad \quad \quad \operatorname{series\text{-}partial\text{-}sum}((\lambda k \in \mathbb{N}.\; \\ \quad \quad \quad \quad \quad \quad ((\operatorname{nth\text{-}deriv}(f, k)(z) \cdot \operatorname{power}((x - z), k)) \cdot \operatorname{recip}(\operatorname{factorial}(k)))), \\ \quad \quad \quad \quad \quad (m + 1))), \\ \quad \quad \quad t, \\ \quad \quad \quad ((\operatorname{recip}(\operatorname{factorial}(m)) \cdot \operatorname{nth\text{-}deriv}(f, (m + 1))(t)) \cdot \operatorname{power}((x - t), m))) \end{array}}}
\vnbentry{\texttt{taylor-\allowbreak{}g-\allowbreak{}diff}}{\vnbstmt{\begin{array}{@{}l@{}} \forall f, a, x, n, t. \\ \quad (n \in \mathbb{N}) \wedge \\ \quad (x \in \mathbb{R}) \wedge \\ \quad \operatorname{taylor\text{-}differentiable}(f, a, x, n) \wedge \\ \quad (a < t) \wedge \\ \quad (t < x) \Rightarrow \\ \quad \quad \operatorname{is\text{-}diff\text{-}at}((\lambda z \in \mathbb{R}.\; (f(x) - \operatorname{taylor\text{-}poly}(f, z, n, x))), t, \\ \quad \quad \quad (0 - ((\operatorname{recip}(\operatorname{factorial}(n)) \cdot \operatorname{nth\text{-}deriv}(f, (n + 1))(t)) \cdot \operatorname{power}((x - t), n)))) \end{array}}}
\vnbentry{\texttt{taylor-\allowbreak{}h-\allowbreak{}cont}}{\vnbstmt{\begin{array}{@{}l@{}} \forall x, n, t. \\ \quad ((x \in \mathbb{R}) \wedge ((t \in \mathbb{R}) \wedge (n \in \mathbb{N}))) \Rightarrow  \\ \quad \quad \operatorname{is\text{-}continuous\text{-}at}(\operatorname{rr\text{-}ms}, \operatorname{rr\text{-}ms}, (\lambda z \in \mathbb{R}.\; \operatorname{power}((x - z), (n + 1))), t) \end{array}}}
\vnbentry{\texttt{tgc-\allowbreak{}term-\allowbreak{}cont}}{\vnbstmt{\begin{array}{@{}l@{}} \forall m \in \mathbb{N}, f.; \forall x, t \in \mathbb{R}. \\ \quad \operatorname{is\text{-}continuous\text{-}at}(\operatorname{rr\text{-}ms}, \operatorname{rr\text{-}ms}, \operatorname{nth\text{-}deriv}(f, (m + 1)), t) \Rightarrow \\ \quad \quad \operatorname{is\text{-}continuous\text{-}at}(\operatorname{rr\text{-}ms}, \operatorname{rr\text{-}ms}, \\ \quad \quad \quad (\lambda z \in \mathbb{R}.\; \\ \quad \quad \quad \quad ((\operatorname{nth\text{-}deriv}(f, (m + 1))(z) \cdot \operatorname{power}((x - z), (m + 1))) \cdot \operatorname{recip}(\operatorname{factorial}((m + 1))))), \\ \quad \quad \quad t) \end{array}}}
\vnbentry{\texttt{tgc-\allowbreak{}sum-\allowbreak{}cont}}{\vnbstmt{\begin{array}{@{}l@{}} \forall m \in \mathbb{N}, f.; \forall x, t \in \mathbb{R}. \\ \quad \forall k \in \mathbb{N}. \\ \quad \quad (k \leq m) \Rightarrow  \\ \quad \quad \quad \operatorname{is\text{-}continuous\text{-}at}(\operatorname{rr\text{-}ms}, \operatorname{rr\text{-}ms}, \operatorname{nth\text{-}deriv}(f, k), t) \Rightarrow \\ \quad \quad \operatorname{is\text{-}continuous\text{-}at}(\operatorname{rr\text{-}ms}, \operatorname{rr\text{-}ms}, \\ \quad \quad \quad (\lambda z \in \mathbb{R}.\; \\ \quad \quad \quad \quad \operatorname{series\text{-}partial\text{-}sum}((\lambda k \in \mathbb{N}.\; \\ \quad \quad \quad \quad \quad \quad ((\operatorname{nth\text{-}deriv}(f, k)(z) \cdot \operatorname{power}((x - z), k)) \cdot \operatorname{recip}(\operatorname{factorial}(k)))), \\ \quad \quad \quad \quad \quad (m + 1))), \\ \quad \quad \quad t) \end{array}}}
\vnbentry{\texttt{taylor-\allowbreak{}g-\allowbreak{}cont}}{\vnbstmt{\begin{array}{@{}l@{}} \forall f, a, x, n, t. \\ \quad (n \in \mathbb{N}) \wedge \\ \quad (x \in \mathbb{R}) \wedge \\ \quad \operatorname{taylor\text{-}differentiable}(f, a, x, n) \wedge \\ \quad (t \in \operatorname{ccint}(a, x)) \Rightarrow \\ \quad \quad \operatorname{is\text{-}continuous\text{-}at}(\operatorname{rr\text{-}ms}, \operatorname{rr\text{-}ms}, (\lambda z \in \mathbb{R}.\; (f(x) - \operatorname{taylor\text{-}poly}(f, z, n, x))), \\ \quad \quad \quad t) \end{array}}}
\vnbentry{\texttt{taylor-\allowbreak{}gmvt-\allowbreak{}cont}}{\vnbstmt{\begin{array}{@{}l@{}} \forall f, a, x, n.; \forall t \in \operatorname{ccint}(a, x). \\ \quad \operatorname{taylor\text{-}differentiable}(f, a, x, n) \wedge \\ \quad (x \in \mathbb{R}) \wedge \\ \quad (n \in \mathbb{N}) \Rightarrow \\ \quad \quad \operatorname{is\text{-}continuous\text{-}at}(\operatorname{rr\text{-}ms}, \operatorname{rr\text{-}ms}, (\lambda z \in \mathbb{R}.\; (f(x) - \operatorname{taylor\text{-}poly}(f, z, n, x))), \\ \quad \quad \quad t) \wedge \\ \quad \quad \operatorname{is\text{-}continuous\text{-}at}(\operatorname{rr\text{-}ms}, \operatorname{rr\text{-}ms}, (\lambda z \in \mathbb{R}.\; \operatorname{power}((x - z), (n + 1))), t) \end{array}}}
\vnbentry{\texttt{taylor-\allowbreak{}gmvt-\allowbreak{}diff}}{\vnbstmt{\begin{array}{@{}l@{}} \forall f, a, x, n.; \forall t \in \mathbb{R}. \\ \quad \operatorname{taylor\text{-}differentiable}(f, a, x, n) \wedge \\ \quad (n \in \mathbb{N}) \wedge \\ \quad (x \in \mathbb{R}) \wedge \\ \quad (a < t) \wedge \\ \quad (t < x) \Rightarrow \\ \quad \quad \exists l. \\ \quad \quad \quad \operatorname{is\text{-}diff\text{-}at}((\lambda z \in \mathbb{R}.\; (f(x) - \operatorname{taylor\text{-}poly}(f, z, n, x))), t, l) \wedge \\ \quad \quad \exists m. \\ \quad \quad \quad \operatorname{is\text{-}diff\text{-}at}((\lambda z \in \mathbb{R}.\; \operatorname{power}((x - z), (n + 1))), t, m) \end{array}}}
\vnbentry{\texttt{rr-\allowbreak{}power-\allowbreak{}pos}}{\vnbstmt{\begin{array}{@{}l@{}} \forall d \in \mathbb{R}, n \in \mathbb{N}. \\ \quad ((0 < d) \Rightarrow (0 < \operatorname{power}(d, n))) \end{array}}}
\vnbentry{\texttt{rr-\allowbreak{}cancel-\allowbreak{}mul-\allowbreak{}left}}{\vnbstmt{\begin{array}{@{}l@{}} \forall c, u, v \in \mathbb{R}. \\ \quad ((\neg (c = 0) \wedge ((c \cdot u) = (c \cdot v))) \Rightarrow (u = v)) \end{array}}}
\vnbentry{\texttt{rr-\allowbreak{}recip-\allowbreak{}factorial}}{\vnbstmt{\begin{array}{@{}l@{}} \forall n \in \mathbb{N}. \\ \quad ((\operatorname{factorial}(n) \cdot \operatorname{recip}(\operatorname{factorial}(n))) = 1) \end{array}}}
\vnbentry{\texttt{taylor-\allowbreak{}clear}}{\vnbstmt{\begin{array}{@{}l@{}} \forall n_{1}, w, d, r \in \mathbb{R}, x, a, b, c.; \forall n \in \mathbb{N}. \\ \quad (0 < w) \wedge \\ \quad (x = 0) \wedge \\ \quad (b = 0) \wedge \\ \quad (c = d) \wedge \\ \quad (a = r) \wedge \\ \quad ((0 - ((\operatorname{recip}(\operatorname{factorial}(n)) \cdot n_{1}) \cdot w)) \cdot (b - c)) \\ \quad =\; ((0 - ((n + 1) \cdot w)) \cdot (x - a)) \Rightarrow \\ \quad \quad ((\operatorname{factorial}((n + 1)) \cdot r) = (n_{1} \cdot d)) \end{array}}}
\vnbentry{\texttt{taylor-\allowbreak{}lagrange}}{\vnbstmt{\begin{array}{@{}l@{}} \forall f, a, x, n. \\ \quad (f \in (\mathbb{R} \to \mathbb{R})) \wedge \\ \quad (a \in \mathbb{R}) \wedge \\ \quad (x \in \mathbb{R}) \wedge \\ \quad (n \in \mathbb{N}) \wedge \\ \quad (a < x) \wedge \\ \quad \operatorname{taylor\text{-}differentiable}(f, a, x, n) \Rightarrow \\ \quad \quad \exists b \in \mathbb{R}. \\ \quad \quad \quad (a < b) \wedge \\ \quad \quad \quad (b < x) \wedge \\ \quad \quad \quad (\operatorname{factorial}((n + 1)) \cdot (f(x) - \operatorname{taylor\text{-}poly}(f, a, n, x))) \\ \quad \quad \quad =\; (\operatorname{nth\text{-}deriv}(f, (n + 1))(b) \cdot \operatorname{power}((x - a), (n + 1))) \end{array}}}

\vnbfile{theorem-library/rake-hb-leaves.scm}{theorem-library/rake-hb-leaves.scm}
\noindent{\small theorem-library/rake-hb-leaves.scm -- nine asserted leaves of the Hahn-Banach cluster, PROVEN.}\par
\vnbentry{\texttt{extends-\allowbreak{}on-\allowbreak{}trans}}{\vnbstmt{\begin{array}{@{}l@{}} \forall s, t, a, b, f. \\ \quad (s \subseteq t) \wedge \\ \quad \operatorname{extends\text{-}on}(t, a, b) \wedge \\ \quad \operatorname{extends\text{-}on}(s, b, f) \Rightarrow \\ \quad \quad \operatorname{extends\text{-}on}(s, a, f) \end{array}}}
\vnbentry{\texttt{vnrm-\allowbreak{}nonneg}}{\vnbstmt{\begin{array}{@{}l@{}} \forall m, w. \\ \quad (\operatorname{is\text{-}normed\text{-}vector\text{-}space}(m) \wedge (w \in \operatorname{vec}(m))) \Rightarrow  \\ \quad \quad (0 \leq \operatorname{vnrm}(m)(w)) \end{array}}}
\vnbentry{\texttt{linfun-\allowbreak{}on-\allowbreak{}vec-\allowbreak{}is-\allowbreak{}linfun}}{\vnbstmt{\begin{array}{@{}l@{}} \forall m, f. \\ \quad \operatorname{is\text{-}linear\text{-}functional\text{-}on}(m, \operatorname{vec}(m), f) \Rightarrow  \\ \quad \quad \operatorname{is\text{-}linear\text{-}functional}(m, f) \end{array}}}
\vnbentry{\texttt{blf-\allowbreak{}app-\allowbreak{}real}}{\vnbstmt{\begin{array}{@{}l@{}} \forall m, g, v. \\ \quad \operatorname{is\text{-}bounded\text{-}linear\text{-}functional}(m, g) \wedge \\ \quad (v \in \operatorname{vec}(m)) \Rightarrow \\ \quad \quad (g(v) \in \mathbb{R}) \end{array}}}
\vnbentry{\texttt{bdd-\allowbreak{}linfun-\allowbreak{}abs-\allowbreak{}real}}{\vnbstmt{\begin{array}{@{}l@{}} \forall m, f, x. \\ \quad \operatorname{is\text{-}bounded\text{-}linear\text{-}functional}(m, f) \wedge \\ \quad (x \in \operatorname{vec}(m)) \Rightarrow \\ \quad \quad (\lvert f(x) \rvert \in \mathbb{R}) \end{array}}}
\vnbentry{\texttt{linfun-\allowbreak{}app-\allowbreak{}abs-\allowbreak{}real}}{\vnbstmt{\begin{array}{@{}l@{}} \forall m, t, g, w. \\ \quad (\operatorname{is\text{-}linear\text{-}functional\text{-}on}(m, t, g) \wedge (w \in t)) \Rightarrow  \\ \quad \quad (\lvert g(w) \rvert \in \mathbb{R}) \end{array}}}
\vnbentry{\texttt{abs-\allowbreak{}nonneg-\allowbreak{}le}}{\vnbstmt{\begin{array}{@{}l@{}} \forall a, c. \\ \quad ((a \in \mathbb{R}) \wedge ((0 \leq a) \wedge (\lvert a \rvert \leq c))) \Rightarrow  \\ \quad \quad (a \leq c) \end{array}}}
\vnbentry{\texttt{rr-\allowbreak{}mul-\allowbreak{}le-\allowbreak{}right}}{\vnbstmt{\begin{array}{@{}l@{}} \forall a, b, c \in \mathbb{R}. \\ \quad (((a \leq b) \wedge (0 \leq c)) \Rightarrow ((a \cdot c) \leq (b \cdot c))) \end{array}}}
\vnbentry{\texttt{proper-\allowbreak{}subset-\allowbreak{}witness}}{\vnbstmt{\begin{array}{@{}l@{}} \forall a, b. \\ \quad ((a \subseteq b) \wedge \neg (a = b)) \Rightarrow  \\ \quad \quad \exists x \in b.\; \neg (x \in a) \end{array}}}

\vnbfile{theorem-library/rake-dual-norm-spec.scm}{theorem-library/rake-dual-norm-spec.scm}
\noindent{\small rake-dual-norm-spec.scm -- the DUAL-NORM / DUAL-NORM-ON descriptions DENOTE, and the corollaries that follow from them.  Rake batch 7, assignment 7-A. Helper prefix `r8a-'.}\par
\vnbentry{\texttt{rr-\allowbreak{}le-\allowbreak{}of-\allowbreak{}le-\allowbreak{}add-\allowbreak{}all-\allowbreak{}pos}}{\vnbstmt{\begin{array}{@{}l@{}} \forall a, b \in \mathbb{R}. \\ \quad \forall d.\; (\operatorname{pos\text{-}rr}(d) \Rightarrow (a \leq (b + d))) \Rightarrow  \\ \quad \quad (a \leq b) \end{array}}}
\vnbentry{\texttt{rr-\allowbreak{}recip-\allowbreak{}cancel-\allowbreak{}right}}{\vnbstmt{\begin{array}{@{}l@{}} \forall a, n \in \mathbb{R}. \\ \quad (\neg (n = 0) \Rightarrow (a = ((a \cdot \operatorname{recip}(n)) \cdot n))) \end{array}}}
\vnbentry{\texttt{rr-\allowbreak{}mul-\allowbreak{}recip-\allowbreak{}cancel}}{\vnbstmt{\begin{array}{@{}l@{}} \forall a, n \in \mathbb{R}. \\ \quad (\neg (n = 0) \Rightarrow (a = ((a \cdot n) \cdot \operatorname{recip}(n)))) \end{array}}}
\vnbentry{\texttt{linfun-\allowbreak{}is-\allowbreak{}linfun-\allowbreak{}on-\allowbreak{}vec}}{\vnbstmt{\begin{array}{@{}l@{}} \forall m, f. \\ \quad \operatorname{is\text{-}linear\text{-}functional}(m, f) \Rightarrow  \\ \quad \quad \operatorname{is\text{-}linear\text{-}functional\text{-}on}(m, \operatorname{vec}(m), f) \end{array}}}
\vnbentry{\texttt{blf-\allowbreak{}is-\allowbreak{}blf-\allowbreak{}on-\allowbreak{}vec}}{\vnbstmt{\begin{array}{@{}l@{}} \forall m, f. \\ \quad \operatorname{is\text{-}bounded\text{-}linear\text{-}functional}(m, f) \Rightarrow  \\ \quad \quad \operatorname{is\text{-}bounded\text{-}linear\text{-}functional\text{-}on}(m, \operatorname{vec}(m), f) \end{array}}}
\vnbentry{\texttt{dual-\allowbreak{}norm-\allowbreak{}on-\allowbreak{}lub-\allowbreak{}exists}}{\vnbstmt{\begin{array}{@{}l@{}} \forall m, s, f. \\ \quad \operatorname{is\text{-}normed\text{-}vector\text{-}space}(m) \wedge \\ \quad (s \subseteq \operatorname{vec}(m)) \wedge \\ \quad \operatorname{is\text{-}bounded\text{-}linear\text{-}functional\text{-}on}(m, s, f) \Rightarrow \\ \quad \quad \exists c \in \mathbb{R}. \\ \quad \quad \quad (0 \leq c) \wedge \\ \quad \quad \quad \forall x \in s.\; (\lvert f(x) \rvert \leq (c \cdot \operatorname{vnrm}(m)(x))) \wedge \\ \quad \quad \quad \forall d \in \mathbb{R}. \\ \quad \quad \quad \quad ((0 \leq d) \wedge \forall x \in s.\; (\lvert f(x) \rvert \leq (d \cdot \operatorname{vnrm}(m)(x)))) \Rightarrow  \\ \quad \quad \quad \quad \quad (c \leq d) \end{array}}}
\vnbentry{\texttt{dual-\allowbreak{}norm-\allowbreak{}on-\allowbreak{}lub-\allowbreak{}unique}}{\vnbstmt{\begin{array}{@{}l@{}} \forall m, s, f, c_{1}, c_{2}. \\ \quad (c_{1} \in \mathbb{R}) \wedge \\ \quad (0 \leq c_{1}) \wedge \\ \quad \forall x \in s. \\ \quad \quad (\lvert f(x) \rvert \leq (c_{1} \cdot \operatorname{vnrm}(m)(x))) \wedge \\ \quad \forall d \in \mathbb{R}. \\ \quad \quad ((0 \leq d) \wedge \forall x \in s.\; (\lvert f(x) \rvert \leq (d \cdot \operatorname{vnrm}(m)(x)))) \Rightarrow  \\ \quad \quad \quad (c_{1} \leq d) \wedge \\ \quad (c_{2} \in \mathbb{R}) \wedge \\ \quad (0 \leq c_{2}) \wedge \\ \quad \forall x \in s. \\ \quad \quad (\lvert f(x) \rvert \leq (c_{2} \cdot \operatorname{vnrm}(m)(x))) \wedge \\ \quad \forall d \in \mathbb{R}. \\ \quad \quad ((0 \leq d) \wedge \forall x \in s.\; (\lvert f(x) \rvert \leq (d \cdot \operatorname{vnrm}(m)(x)))) \Rightarrow  \\ \quad \quad \quad (c_{2} \leq d) \Rightarrow \\ \quad \quad (c_{1} = c_{2}) \end{array}}}
\vnbentry{\texttt{dual-\allowbreak{}norm-\allowbreak{}on-\allowbreak{}spec}}{\vnbstmt{\begin{array}{@{}l@{}} \forall m, s, f. \\ \quad \operatorname{is\text{-}normed\text{-}vector\text{-}space}(m) \wedge \\ \quad (s \subseteq \operatorname{vec}(m)) \wedge \\ \quad \operatorname{is\text{-}bounded\text{-}linear\text{-}functional\text{-}on}(m, s, f) \Rightarrow \\ \quad \quad (\operatorname{dual\text{-}norm\text{-}on}(m, s, f) \in \mathbb{R}) \wedge \\ \quad \quad (0 \leq \operatorname{dual\text{-}norm\text{-}on}(m, s, f)) \wedge \\ \quad \quad \forall x \in s. \\ \quad \quad \quad \lvert f(x) \rvert \\ \quad \quad \quad \leq\; (\operatorname{dual\text{-}norm\text{-}on}(m, s, f) \cdot \operatorname{vnrm}(m)(x)) \wedge \\ \quad \quad \forall d \in \mathbb{R}. \\ \quad \quad \quad ((0 \leq d) \wedge \forall x \in s.\; (\lvert f(x) \rvert \leq (d \cdot \operatorname{vnrm}(m)(x)))) \Rightarrow  \\ \quad \quad \quad \quad (\operatorname{dual\text{-}norm\text{-}on}(m, s, f) \leq d) \end{array}}}
\vnbentry{\texttt{dual-\allowbreak{}norm-\allowbreak{}spec}}{\vnbstmt{\begin{array}{@{}l@{}} \forall m, f. \\ \quad \operatorname{is\text{-}normed\text{-}vector\text{-}space}(m) \wedge \\ \quad \operatorname{is\text{-}bounded\text{-}linear\text{-}functional}(m, f) \Rightarrow \\ \quad \quad (\operatorname{dual\text{-}norm}(m, f) \in \mathbb{R}) \wedge \\ \quad \quad (0 \leq \operatorname{dual\text{-}norm}(m, f)) \wedge \\ \quad \quad \forall x \in \operatorname{vec}(m). \\ \quad \quad \quad (\lvert f(x) \rvert \leq (\operatorname{dual\text{-}norm}(m, f) \cdot \operatorname{vnrm}(m)(x))) \wedge \\ \quad \quad \forall d \in \mathbb{R}. \\ \quad \quad \quad (0 \leq d) \wedge \\ \quad \quad \quad \forall x \in \operatorname{vec}(m). \\ \quad \quad \quad \quad (\lvert f(x) \rvert \leq (d \cdot \operatorname{vnrm}(m)(x))) \Rightarrow \\ \quad \quad \quad \quad (\operatorname{dual\text{-}norm}(m, f) \leq d) \end{array}}}
\vnbentry{\texttt{dual-\allowbreak{}norm-\allowbreak{}on-\allowbreak{}nonneg-\allowbreak{}nvs}}{\vnbstmt{\begin{array}{@{}l@{}} \forall m, s, f. \\ \quad \operatorname{is\text{-}normed\text{-}vector\text{-}space}(m) \wedge \\ \quad (s \subseteq \operatorname{vec}(m)) \wedge \\ \quad \operatorname{is\text{-}bounded\text{-}linear\text{-}functional\text{-}on}(m, s, f) \Rightarrow \\ \quad \quad (\operatorname{dual\text{-}norm\text{-}on}(m, s, f) \in \mathbb{R}) \wedge \\ \quad \quad (0 \leq \operatorname{dual\text{-}norm\text{-}on}(m, s, f)) \end{array}}}
\vnbentry{\texttt{dual-\allowbreak{}norm-\allowbreak{}on-\allowbreak{}le-\allowbreak{}bound-\allowbreak{}nvs}}{\vnbstmt{\begin{array}{@{}l@{}} \forall m, s, g, c. \\ \quad \operatorname{is\text{-}normed\text{-}vector\text{-}space}(m) \wedge \\ \quad (s \subseteq \operatorname{vec}(m)) \wedge \\ \quad \operatorname{is\text{-}linear\text{-}functional\text{-}on}(m, s, g) \wedge \\ \quad (c \in \mathbb{R}) \wedge \\ \quad (0 \leq c) \wedge \\ \quad \forall w \in s.\; (\lvert g(w) \rvert \leq (c \cdot \operatorname{vnrm}(m)(w))) \Rightarrow \\ \quad \quad (\operatorname{dual\text{-}norm\text{-}on}(m, s, g) \leq c) \end{array}}}
\vnbentry{\texttt{dual-\allowbreak{}norm-\allowbreak{}on-\allowbreak{}is-\allowbreak{}bound-\allowbreak{}nvs}}{\vnbstmt{\begin{array}{@{}l@{}} \forall m, s, f, x. \\ \quad \operatorname{is\text{-}normed\text{-}vector\text{-}space}(m) \wedge \\ \quad (s \subseteq \operatorname{vec}(m)) \wedge \\ \quad \operatorname{is\text{-}bounded\text{-}linear\text{-}functional\text{-}on}(m, s, f) \wedge \\ \quad (x \in s) \Rightarrow \\ \quad \quad \lvert f(x) \rvert \\ \quad \quad \leq\; (\operatorname{dual\text{-}norm\text{-}on}(m, s, f) \cdot \operatorname{vnrm}(m)(x)) \end{array}}}
\vnbentry{\texttt{dual-\allowbreak{}norm-\allowbreak{}nonneg-\allowbreak{}nvs}}{\vnbstmt{\begin{array}{@{}l@{}} \forall m, f. \\ \quad \operatorname{is\text{-}normed\text{-}vector\text{-}space}(m) \wedge \\ \quad \operatorname{is\text{-}bounded\text{-}linear\text{-}functional}(m, f) \Rightarrow \\ \quad \quad (\operatorname{dual\text{-}norm}(m, f) \in \mathbb{R}) \wedge \\ \quad \quad (0 \leq \operatorname{dual\text{-}norm}(m, f)) \end{array}}}
\vnbentry{\texttt{dual-\allowbreak{}norm-\allowbreak{}is-\allowbreak{}bound-\allowbreak{}nvs}}{\vnbstmt{\begin{array}{@{}l@{}} \forall m, f, x. \\ \quad \operatorname{is\text{-}normed\text{-}vector\text{-}space}(m) \wedge \\ \quad \operatorname{is\text{-}bounded\text{-}linear\text{-}functional}(m, f) \wedge \\ \quad (x \in \operatorname{vec}(m)) \Rightarrow \\ \quad \quad (\lvert f(x) \rvert \leq (\operatorname{dual\text{-}norm}(m, f) \cdot \operatorname{vnrm}(m)(x))) \end{array}}}
\vnbentry{\texttt{dual-\allowbreak{}norm-\allowbreak{}le-\allowbreak{}bound-\allowbreak{}nvs}}{\vnbstmt{\begin{array}{@{}l@{}} \forall m, f, c. \\ \quad \operatorname{is\text{-}normed\text{-}vector\text{-}space}(m) \wedge \\ \quad \operatorname{is\text{-}linear\text{-}functional}(m, f) \wedge \\ \quad (c \in \mathbb{R}) \wedge \\ \quad (0 \leq c) \wedge \\ \quad \forall w \in \operatorname{vec}(m). \\ \quad \quad (\lvert f(w) \rvert \leq (c \cdot \operatorname{vnrm}(m)(w))) \Rightarrow \\ \quad \quad (\operatorname{dual\text{-}norm}(m, f) \leq c) \end{array}}}
\vnbentry{\texttt{good-\allowbreak{}sub-\allowbreak{}self-\allowbreak{}nvs}}{\vnbstmt{\begin{array}{@{}l@{}} \forall m, s, f. \\ \quad \operatorname{is\text{-}normed\text{-}vector\text{-}space}(m) \wedge \\ \quad \operatorname{is\text{-}submodule}(m, s) \wedge \\ \quad \operatorname{is\text{-}bounded\text{-}linear\text{-}functional\text{-}on}(m, s, f) \Rightarrow \\ \quad \quad \operatorname{good\text{-}sub}(m, s, f, s) \end{array}}}

\vnbfile{theorem-library/nvs-module-view.scm}{theorem-library/nvs-module-view.scm}
\noindent{\small nvs-module-view.scm -- THE MODULE BRIDGE OF A NORMED VECTOR SPACE, PROVEN.}\par
\vnbentry{\texttt{nvs-\allowbreak{}module-\allowbreak{}view-\allowbreak{}scal}}{\vnbstmt{\begin{array}{@{}l@{}} \forall m. \\ \quad (\operatorname{scal}(\operatorname{normed\text{-}vector\text{-}space\text{-}as\text{-}module}(m)) \simeq \operatorname{scal}(m)) \end{array}}}
\vnbentry{\texttt{nvs-\allowbreak{}module-\allowbreak{}view-\allowbreak{}vec}}{\vnbstmt{\begin{array}{@{}l@{}} \forall m. \\ \quad (\operatorname{vec}(\operatorname{normed\text{-}vector\text{-}space\text{-}as\text{-}module}(m)) \simeq \operatorname{vec}(m)) \end{array}}}
\vnbentry{\texttt{nvs-\allowbreak{}module-\allowbreak{}view-\allowbreak{}vadd}}{\vnbstmt{\begin{array}{@{}l@{}} \forall m. \\ \quad (\operatorname{vadd}(\operatorname{normed\text{-}vector\text{-}space\text{-}as\text{-}module}(m)) \simeq \operatorname{vadd}(m)) \end{array}}}
\vnbentry{\texttt{nvs-\allowbreak{}module-\allowbreak{}view-\allowbreak{}vzero}}{\vnbstmt{\begin{array}{@{}l@{}} \forall m. \\ \quad (\operatorname{vzero}(\operatorname{normed\text{-}vector\text{-}space\text{-}as\text{-}module}(m)) \simeq \operatorname{vzero}(m)) \end{array}}}
\vnbentry{\texttt{nvs-\allowbreak{}module-\allowbreak{}view-\allowbreak{}vneg}}{\vnbstmt{\begin{array}{@{}l@{}} \forall m. \\ \quad (\operatorname{vneg}(\operatorname{normed\text{-}vector\text{-}space\text{-}as\text{-}module}(m)) \simeq \operatorname{vneg}(m)) \end{array}}}
\vnbentry{\texttt{nvs-\allowbreak{}module-\allowbreak{}view-\allowbreak{}act}}{\vnbstmt{\begin{array}{@{}l@{}} \forall m. \\ \quad (\operatorname{act}(\operatorname{normed\text{-}vector\text{-}space\text{-}as\text{-}module}(m)) \simeq \operatorname{act}(m)) \end{array}}}
\vnbentry{\texttt{submodule-\allowbreak{}nvs-\allowbreak{}module-\allowbreak{}view}}{\vnbstmt{\begin{array}{@{}l@{}} \forall m, t. \\ \quad \operatorname{is\text{-}submodule}(m, t) \Rightarrow  \\ \quad \quad \operatorname{is\text{-}submodule}(\operatorname{normed\text{-}vector\text{-}space\text{-}as\text{-}module}(m), t) \end{array}}}

\vnbfile{theorem-library/nvs-act-laws.scm}{theorem-library/nvs-act-laws.scm}
\noindent{\small nvs-act-laws.scm -- THE GROUP AND SCALAR-ACTION LAWS OF A REAL NORMED VECTOR SPACE.  Since 2026-09-19 (rake batch 8, assignment 8-K1) this file is their ONE home: section (4) below holds the laws that three to five agents had each proved separately on 2026-09-19 under their own prefixes.  The NORM}\par
\vnbentry{\texttt{nvs-\allowbreak{}scal-\allowbreak{}carr}}{\vnbstmt{\begin{array}{@{}l@{}} \forall m. \\ \quad \operatorname{is\text{-}normed\text{-}vector\text{-}space}(m) \Rightarrow  \\ \quad \quad (\operatorname{carr}(\operatorname{scal}(m)) \simeq \mathbb{R}) \end{array}}}
\vnbentry{\texttt{nvs-\allowbreak{}scal-\allowbreak{}add}}{\vnbstmt{\begin{array}{@{}l@{}} \forall m. \\ \quad \operatorname{is\text{-}normed\text{-}vector\text{-}space}(m) \Rightarrow  \\ \quad \quad (\operatorname{add}(\operatorname{scal}(m)) \simeq (\lambda \langle \operatorname{x\_}, \operatorname{y\_} \rangle \in (\mathbb{R} \times \mathbb{R}).\; (\operatorname{x\_} + \operatorname{y\_}))) \end{array}}}
\vnbentry{\texttt{nvs-\allowbreak{}scal-\allowbreak{}mul}}{\vnbstmt{\begin{array}{@{}l@{}} \forall m. \\ \quad \operatorname{is\text{-}normed\text{-}vector\text{-}space}(m) \Rightarrow  \\ \quad \quad (\operatorname{mul}(\operatorname{scal}(m)) \simeq (\lambda \langle \operatorname{x\_}, \operatorname{y\_} \rangle \in (\mathbb{R} \times \mathbb{R}).\; (\operatorname{x\_} \cdot \operatorname{y\_}))) \end{array}}}
\vnbentry{\texttt{nvs-\allowbreak{}scal-\allowbreak{}zero}}{\vnbstmt{\begin{array}{@{}l@{}} \forall m. \\ \quad \operatorname{is\text{-}normed\text{-}vector\text{-}space}(m) \Rightarrow  \\ \quad \quad (\operatorname{zero}(\operatorname{scal}(m)) \simeq 0) \end{array}}}
\vnbentry{\texttt{nvs-\allowbreak{}scal-\allowbreak{}one}}{\vnbstmt{\begin{array}{@{}l@{}} \forall m. \\ \quad \operatorname{is\text{-}normed\text{-}vector\text{-}space}(m) \Rightarrow  \\ \quad \quad (\operatorname{one}(\operatorname{scal}(m)) \simeq 1) \end{array}}}
\vnbentry{\texttt{nvs-\allowbreak{}act-\allowbreak{}distrib-\allowbreak{}scalar}}{\vnbstmt{\begin{array}{@{}l@{}} \forall m.; \forall r, s \in \operatorname{carr}(\operatorname{scal}(m)), x \in \operatorname{vec}(m). \\ \quad \operatorname{is\text{-}normed\text{-}vector\text{-}space}(m) \Rightarrow  \\ \quad \quad \operatorname{act}(m)(\operatorname{add}(\operatorname{scal}(m))(r, s), x) \\ \quad \quad =\; \operatorname{vadd}(m)(\operatorname{act}(m)(r, x), \operatorname{act}(m)(s, x)) \end{array}}}
\vnbentry{\texttt{nvs-\allowbreak{}act-\allowbreak{}mul-\allowbreak{}compat}}{\vnbstmt{\begin{array}{@{}l@{}} \forall m.; \forall r, s \in \operatorname{carr}(\operatorname{scal}(m)), x \in \operatorname{vec}(m). \\ \quad \operatorname{is\text{-}normed\text{-}vector\text{-}space}(m) \Rightarrow  \\ \quad \quad \operatorname{act}(m)(\operatorname{mul}(\operatorname{scal}(m))(r, s), x) \\ \quad \quad =\; \operatorname{act}(m)(r, \operatorname{act}(m)(s, x)) \end{array}}}
\vnbentry{\texttt{nvs-\allowbreak{}act-\allowbreak{}unital}}{\vnbstmt{\begin{array}{@{}l@{}} \forall m.; \forall x \in \operatorname{vec}(m). \\ \quad \operatorname{is\text{-}normed\text{-}vector\text{-}space}(m) \Rightarrow  \\ \quad \quad (\operatorname{act}(m)(\operatorname{one}(\operatorname{scal}(m)), x) = x) \end{array}}}
\vnbentry{\texttt{nvs-\allowbreak{}vadd-\allowbreak{}assoc}}{\vnbstmt{\begin{array}{@{}l@{}} \forall m.; \forall u, v, w \in \operatorname{vec}(m). \\ \quad \operatorname{is\text{-}normed\text{-}vector\text{-}space}(m) \Rightarrow  \\ \quad \quad \operatorname{vadd}(m)(\operatorname{vadd}(m)(u, v), w) \\ \quad \quad =\; \operatorname{vadd}(m)(u, \operatorname{vadd}(m)(v, w)) \end{array}}}
\vnbentry{\texttt{nvs-\allowbreak{}vzero-\allowbreak{}right}}{\vnbstmt{\begin{array}{@{}l@{}} \forall m.; \forall u \in \operatorname{vec}(m). \\ \quad \operatorname{is\text{-}normed\text{-}vector\text{-}space}(m) \Rightarrow  \\ \quad \quad (\operatorname{vadd}(m)(u, \operatorname{vzero}(m)) = u) \end{array}}}
\vnbentry{\texttt{nvs-\allowbreak{}act-\allowbreak{}type}}{\vnbstmt{\begin{array}{@{}l@{}} \forall m.; \forall r \in \operatorname{carr}(\operatorname{scal}(m)), x \in \operatorname{vec}(m). \\ \quad \operatorname{is\text{-}normed\text{-}vector\text{-}space}(m) \Rightarrow  \\ \quad \quad (\operatorname{act}(m)(r, x) \in \operatorname{vec}(m)) \end{array}}}
\vnbentry{\texttt{nvs-\allowbreak{}act-\allowbreak{}in-\allowbreak{}vec}}{\vnbstmt{\begin{array}{@{}l@{}} \forall m.; \forall r \in \mathbb{R}, x \in \operatorname{vec}(m). \\ \quad \operatorname{is\text{-}normed\text{-}vector\text{-}space}(m) \Rightarrow  \\ \quad \quad (\operatorname{act}(m)(r, x) \in \operatorname{vec}(m)) \end{array}}}
\vnbentry{\texttt{nvs-\allowbreak{}act-\allowbreak{}scale-\allowbreak{}assoc}}{\vnbstmt{\begin{array}{@{}l@{}} \forall m.; \forall r, s \in \mathbb{R}, x \in \operatorname{vec}(m). \\ \quad \operatorname{is\text{-}normed\text{-}vector\text{-}space}(m) \Rightarrow  \\ \quad \quad (\operatorname{act}(m)(r, \operatorname{act}(m)(s, x)) = \operatorname{act}(m)((r \cdot s), x)) \end{array}}}
\vnbentry{\texttt{nvs-\allowbreak{}act-\allowbreak{}collect}}{\vnbstmt{\begin{array}{@{}l@{}} \forall m.; \forall r, s \in \mathbb{R}, x, y \in \operatorname{vec}(m). \\ \quad \operatorname{is\text{-}normed\text{-}vector\text{-}space}(m) \Rightarrow  \\ \quad \quad \operatorname{vadd}(m)(\operatorname{vadd}(m)(y, \operatorname{act}(m)(r, x)), \operatorname{act}(m)(s, x)) \\ \quad \quad =\; \operatorname{vadd}(m)(y, \operatorname{act}(m)((r + s), x)) \end{array}}}
\vnbentry{\texttt{nvs-\allowbreak{}act-\allowbreak{}zero}}{\vnbstmt{\begin{array}{@{}l@{}} \forall m.; \forall x, y \in \operatorname{vec}(m). \\ \quad \operatorname{is\text{-}normed\text{-}vector\text{-}space}(m) \Rightarrow  \\ \quad \quad (\operatorname{vadd}(m)(y, \operatorname{act}(m)(0, x)) = y) \end{array}}}
\vnbentry{\texttt{nvs-\allowbreak{}act-\allowbreak{}one}}{\vnbstmt{\begin{array}{@{}l@{}} \forall m.; \forall x, y \in \operatorname{vec}(m). \\ \quad \operatorname{is\text{-}normed\text{-}vector\text{-}space}(m) \Rightarrow  \\ \quad \quad (\operatorname{vadd}(m)(y, \operatorname{act}(m)(1, x)) = \operatorname{vadd}(m)(y, x)) \end{array}}}
\vnbentry{\texttt{nvs-\allowbreak{}scal-\allowbreak{}neg}}{\vnbstmt{\begin{array}{@{}l@{}} \forall m. \\ \quad \operatorname{is\text{-}normed\text{-}vector\text{-}space}(m) \Rightarrow  \\ \quad \quad (\operatorname{neg}(\operatorname{scal}(m)) \simeq (\lambda x \in \mathbb{R}.\; -x)) \end{array}}}
\vnbentry{\texttt{nvs-\allowbreak{}vzero-\allowbreak{}left}}{\vnbstmt{\begin{array}{@{}l@{}} \forall m.; \forall u \in \operatorname{vec}(m). \\ \quad \operatorname{is\text{-}normed\text{-}vector\text{-}space}(m) \Rightarrow  \\ \quad \quad (\operatorname{vadd}(m)(\operatorname{vzero}(m), u) = u) \end{array}}}
\vnbentry{\texttt{nvs-\allowbreak{}vadd-\allowbreak{}comm}}{\vnbstmt{\begin{array}{@{}l@{}} \forall m.; \forall u, w \in \operatorname{vec}(m). \\ \quad \operatorname{is\text{-}normed\text{-}vector\text{-}space}(m) \Rightarrow  \\ \quad \quad (\operatorname{vadd}(m)(u, w) = \operatorname{vadd}(m)(w, u)) \end{array}}}
\vnbentry{\texttt{nvs-\allowbreak{}act-\allowbreak{}distrib-\allowbreak{}vec}}{\vnbstmt{\begin{array}{@{}l@{}} \forall m.; \forall r \in \operatorname{carr}(\operatorname{scal}(m)), x, y \in \operatorname{vec}(m). \\ \quad \operatorname{is\text{-}normed\text{-}vector\text{-}space}(m) \Rightarrow  \\ \quad \quad \operatorname{act}(m)(r, \operatorname{vadd}(m)(x, y)) \\ \quad \quad =\; \operatorname{vadd}(m)(\operatorname{act}(m)(r, x), \operatorname{act}(m)(r, y)) \end{array}}}
\vnbentry{\texttt{nvs-\allowbreak{}zero-\allowbreak{}act}}{\vnbstmt{\begin{array}{@{}l@{}} \forall m.; \forall x \in \operatorname{vec}(m). \\ \quad \operatorname{is\text{-}normed\text{-}vector\text{-}space}(m) \Rightarrow  \\ \quad \quad (\operatorname{act}(m)(0, x) = \operatorname{vzero}(m)) \end{array}}}
\vnbentry{\texttt{nvs-\allowbreak{}neg-\allowbreak{}one-\allowbreak{}act}}{\vnbstmt{\begin{array}{@{}l@{}} \forall m.; \forall x \in \operatorname{vec}(m). \\ \quad \operatorname{is\text{-}normed\text{-}vector\text{-}space}(m) \Rightarrow  \\ \quad \quad (\operatorname{act}(m)(-1, x) = \operatorname{vneg}(m)(x)) \end{array}}}
\vnbentry{\texttt{nvs-\allowbreak{}act-\allowbreak{}add-\allowbreak{}scalars}}{\vnbstmt{\begin{array}{@{}l@{}} \forall m.; \forall r, s \in \mathbb{R}, x \in \operatorname{vec}(m). \\ \quad \operatorname{is\text{-}normed\text{-}vector\text{-}space}(m) \Rightarrow  \\ \quad \quad \operatorname{vadd}(m)(\operatorname{act}(m)(r, x), \operatorname{act}(m)(s, x)) \\ \quad \quad =\; \operatorname{act}(m)((r + s), x) \end{array}}}
\vnbentry{\texttt{nvs-\allowbreak{}act-\allowbreak{}vneg}}{\vnbstmt{\begin{array}{@{}l@{}} \forall m.; \forall r \in \mathbb{R}, x \in \operatorname{vec}(m). \\ \quad \operatorname{is\text{-}normed\text{-}vector\text{-}space}(m) \Rightarrow  \\ \quad \quad (\operatorname{vneg}(m)(\operatorname{act}(m)(r, x)) = \operatorname{act}(m)((-1 \cdot r), x)) \end{array}}}
\vnbentry{\texttt{nvs-\allowbreak{}vneg-\allowbreak{}vadd}}{\vnbstmt{\begin{array}{@{}l@{}} \forall m.; \forall x, y \in \operatorname{vec}(m). \\ \quad \operatorname{is\text{-}normed\text{-}vector\text{-}space}(m) \Rightarrow  \\ \quad \quad \operatorname{vneg}(m)(\operatorname{vadd}(m)(x, y)) \\ \quad \quad =\; \operatorname{vadd}(m)(\operatorname{vneg}(m)(x), \operatorname{vneg}(m)(y)) \end{array}}}
\vnbentry{\texttt{nvs-\allowbreak{}vadd-\allowbreak{}shuffle}}{\vnbstmt{\begin{array}{@{}l@{}} \forall m.; \forall a, b, c, d \in \operatorname{vec}(m). \\ \quad \operatorname{is\text{-}normed\text{-}vector\text{-}space}(m) \Rightarrow  \\ \quad \quad \operatorname{vadd}(m)(\operatorname{vadd}(m)(a, c), \operatorname{vadd}(m)(b, d)) \\ \quad \quad =\; \operatorname{vadd}(m)(\operatorname{vadd}(m)(a, b), \operatorname{vadd}(m)(c, d)) \end{array}}}
\vnbentry{\texttt{nvs-\allowbreak{}vzero-\allowbreak{}in-\allowbreak{}vec}}{\vnbstmt{\begin{array}{@{}l@{}} \forall m. \\ \quad \operatorname{is\text{-}normed\text{-}vector\text{-}space}(m) \Rightarrow  \\ \quad \quad (\operatorname{vzero}(m) \in \operatorname{vec}(m)) \end{array}}}
\vnbentry{\texttt{nvs-\allowbreak{}vneg-\allowbreak{}in-\allowbreak{}vec}}{\vnbstmt{\begin{array}{@{}l@{}} \forall m.; \forall x \in \operatorname{vec}(m). \\ \quad \operatorname{is\text{-}normed\text{-}vector\text{-}space}(m) \Rightarrow  \\ \quad \quad (\operatorname{vneg}(m)(x) \in \operatorname{vec}(m)) \end{array}}}
\vnbentry{\texttt{nvs-\allowbreak{}vneg-\allowbreak{}left}}{\vnbstmt{\begin{array}{@{}l@{}} \forall m.; \forall u \in \operatorname{vec}(m). \\ \quad \operatorname{is\text{-}normed\text{-}vector\text{-}space}(m) \Rightarrow  \\ \quad \quad (\operatorname{vadd}(m)(\operatorname{vneg}(m)(u), u) = \operatorname{vzero}(m)) \end{array}}}
\vnbentry{\texttt{nvs-\allowbreak{}act-\allowbreak{}distrib-\allowbreak{}vec-\allowbreak{}rr}}{\vnbstmt{\begin{array}{@{}l@{}} \forall m.; \forall r \in \mathbb{R}, x, y \in \operatorname{vec}(m). \\ \quad \operatorname{is\text{-}normed\text{-}vector\text{-}space}(m) \Rightarrow  \\ \quad \quad \operatorname{act}(m)(r, \operatorname{vadd}(m)(x, y)) \\ \quad \quad =\; \operatorname{vadd}(m)(\operatorname{act}(m)(r, x), \operatorname{act}(m)(r, y)) \end{array}}}
\vnbentry{\texttt{nvs-\allowbreak{}one-\allowbreak{}act}}{\vnbstmt{\begin{array}{@{}l@{}} \forall m.; \forall x \in \operatorname{vec}(m). \\ \quad \operatorname{is\text{-}normed\text{-}vector\text{-}space}(m) \Rightarrow  \\ \quad \quad (\operatorname{act}(m)(1, x) = x) \end{array}}}
\vnbentry{\texttt{nvs-\allowbreak{}vadd-\allowbreak{}cancel}}{\vnbstmt{\begin{array}{@{}l@{}} \forall m.; \forall x, z, w \in \operatorname{vec}(m). \\ \quad \operatorname{is\text{-}normed\text{-}vector\text{-}space}(m) \wedge \\ \quad (\operatorname{vadd}(m)(x, z) = \operatorname{vadd}(m)(x, w)) \Rightarrow \\ \quad \quad (z = w) \end{array}}}
\vnbentry{\texttt{nvs-\allowbreak{}vadd-\allowbreak{}cancel-\allowbreak{}right}}{\vnbstmt{\begin{array}{@{}l@{}} \forall m.; \forall x, z, w \in \operatorname{vec}(m). \\ \quad \operatorname{is\text{-}normed\text{-}vector\text{-}space}(m) \wedge \\ \quad (\operatorname{vadd}(m)(z, x) = \operatorname{vadd}(m)(w, x)) \Rightarrow \\ \quad \quad (z = w) \end{array}}}
\vnbentry{\texttt{nvs-\allowbreak{}vneg-\allowbreak{}right}}{\vnbstmt{\begin{array}{@{}l@{}} \forall m.; \forall u \in \operatorname{vec}(m). \\ \quad \operatorname{is\text{-}normed\text{-}vector\text{-}space}(m) \Rightarrow  \\ \quad \quad (\operatorname{vadd}(m)(u, \operatorname{vneg}(m)(u)) = \operatorname{vzero}(m)) \end{array}}}
\vnbentry{\texttt{nvs-\allowbreak{}vadd-\allowbreak{}cancel-\allowbreak{}zero}}{\vnbstmt{\begin{array}{@{}l@{}} \forall m.; \forall x, z \in \operatorname{vec}(m). \\ \quad \operatorname{is\text{-}normed\text{-}vector\text{-}space}(m) \wedge \\ \quad (\operatorname{vadd}(m)(x, z) = x) \Rightarrow \\ \quad \quad (z = \operatorname{vzero}(m)) \end{array}}}
\vnbentry{\texttt{nvs-\allowbreak{}neg-\allowbreak{}one-\allowbreak{}act-\allowbreak{}sub}}{\vnbstmt{\begin{array}{@{}l@{}} \forall m.; \forall u \in \operatorname{vec}(m). \\ \quad \operatorname{is\text{-}normed\text{-}vector\text{-}space}(m) \Rightarrow  \\ \quad \quad (\operatorname{act}(m)((0 - 1), u) = \operatorname{vneg}(m)(u)) \end{array}}}

\vnbfile{theorem-library/rake-norm-metrics.scm}{theorem-library/rake-norm-metrics.scm}
\noindent{\small rake-norm-metrics.scm -- rake batch 5 R: continuity projections and the norm metrics, PROVEN.}\par
\vnbentry{\texttt{ag-\allowbreak{}norm-\allowbreak{}metric-\allowbreak{}is-\allowbreak{}metric-\allowbreak{}space}}{\vnbstmt{\begin{array}{@{}l@{}} \forall g, m. \\ \quad \operatorname{is\text{-}abelian\text{-}group}(g) \wedge \\ \quad \operatorname{is\text{-}group\text{-}norm}(m, \operatorname{opr}(g), \operatorname{inv}(g), \operatorname{iden}(g), \operatorname{carr}(g)) \Rightarrow \\ \quad \quad \operatorname{is\text{-}metric\text{-}space}(\langle \operatorname{carr}(g), (\lambda \langle u, v \rangle \in (\operatorname{carr}(g) \times \operatorname{carr}(g)).\; m(\operatorname{opr}(g)(u, \operatorname{inv}(g)(v)))) \rangle) \end{array}}}
\vnbentry{\texttt{continuous-\allowbreak{}is-\allowbreak{}continuous-\allowbreak{}at}}{\vnbstmt{\begin{array}{@{}l@{}} \forall s, t, f.; \forall a \in \operatorname{pts}(s). \\ \quad \operatorname{is\text{-}continuous}(s, t, f) \Rightarrow  \\ \quad \quad \operatorname{is\text{-}continuous\text{-}at}(s, t, f, a) \end{array}}}
\vnbentry{\texttt{uniformly-\allowbreak{}continuous-\allowbreak{}is-\allowbreak{}continuous}}{\vnbstmt{\begin{array}{@{}l@{}} \forall s, t, f. \\ \quad \operatorname{is\text{-}uniformly\text{-}continuous}(s, t, f) \Rightarrow  \\ \quad \quad \operatorname{is\text{-}continuous}(s, t, f) \end{array}}}
\vnbentry{\texttt{nag-\allowbreak{}metric-\allowbreak{}carrier}}{\vnbstmt{\begin{array}{@{}l@{}} \forall g. \\ \quad (\operatorname{pts}(\operatorname{nag\text{-}metric\text{-}space}(g)) \simeq \operatorname{carr}(g)) \end{array}}}
\vnbentry{\texttt{nag-\allowbreak{}metric-\allowbreak{}distance}}{\vnbstmt{\begin{array}{@{}l@{}} \forall g.; \forall u, v \in \operatorname{carr}(g). \\ \quad \operatorname{is\text{-}normed\text{-}ag}(g) \Rightarrow  \\ \quad \quad \operatorname{dist}(\operatorname{nag\text{-}metric\text{-}space}(g))(u, v) \\ \quad \quad =\; \operatorname{nrm}(g)(\operatorname{opr}(g)(u, \operatorname{inv}(g)(v))) \end{array}}}
\vnbentry{\texttt{nag-\allowbreak{}metric-\allowbreak{}space-\allowbreak{}is-\allowbreak{}metric-\allowbreak{}space}}{\vnbstmt{\begin{array}{@{}l@{}} \forall g. \\ \quad \operatorname{is\text{-}normed\text{-}ag}(g) \Rightarrow  \\ \quad \quad \operatorname{is\text{-}metric\text{-}space}(\operatorname{nag\text{-}metric\text{-}space}(g)) \end{array}}}
\vnbentry{\texttt{nvs-\allowbreak{}metric-\allowbreak{}is-\allowbreak{}ms}}{\vnbstmt{\begin{array}{@{}l@{}} \forall m. \\ \quad \operatorname{is\text{-}normed\text{-}vector\text{-}space}(m) \Rightarrow  \\ \quad \quad \operatorname{is\text{-}metric\text{-}space}(\operatorname{nvs\text{-}metric\text{-}space}(m)) \end{array}}}

\vnbfile{theorem-library/nvs-norm-laws.scm}{theorem-library/nvs-norm-laws.scm}
\noindent{\small nvs-norm-laws.scm -- THE NORM LAWS OF A REAL NORMED VECTOR SPACE, in ONE place.  Rake batch 8, assignment 8-K1.}\par
\vnbentry{\texttt{nvs-\allowbreak{}vnrm-\allowbreak{}triangle}}{\vnbstmt{\begin{array}{@{}l@{}} \forall m.; \forall x, y \in \operatorname{vec}(m). \\ \quad \operatorname{is\text{-}normed\text{-}vector\text{-}space}(m) \Rightarrow  \\ \quad \quad \operatorname{vnrm}(m)(\operatorname{vadd}(m)(x, y)) \\ \quad \quad \leq\; (\operatorname{vnrm}(m)(x) + \operatorname{vnrm}(m)(y)) \end{array}}}
\vnbentry{\texttt{nvs-\allowbreak{}vnrm-\allowbreak{}homog-\allowbreak{}c}}{\vnbstmt{\begin{array}{@{}l@{}} \forall m.; \forall r \in \operatorname{carr}(\operatorname{scal}(m)), x \in \operatorname{vec}(m). \\ \quad \operatorname{is\text{-}normed\text{-}vector\text{-}space}(m) \Rightarrow  \\ \quad \quad \operatorname{vnrm}(m)(\operatorname{act}(m)(r, x)) \\ \quad \quad =\; (\lvert r \rvert \cdot \operatorname{vnrm}(m)(x)) \end{array}}}
\vnbentry{\texttt{nvs-\allowbreak{}vnrm-\allowbreak{}homog}}{\vnbstmt{\begin{array}{@{}l@{}} \forall m.; \forall r \in \mathbb{R}, x \in \operatorname{vec}(m). \\ \quad \operatorname{is\text{-}normed\text{-}vector\text{-}space}(m) \Rightarrow  \\ \quad \quad \operatorname{vnrm}(m)(\operatorname{act}(m)(r, x)) \\ \quad \quad =\; (\lvert r \rvert \cdot \operatorname{vnrm}(m)(x)) \end{array}}}
\vnbentry{\texttt{nvs-\allowbreak{}vnrm-\allowbreak{}vzero}}{\vnbstmt{\begin{array}{@{}l@{}} \forall m. \\ \quad \operatorname{is\text{-}normed\text{-}vector\text{-}space}(m) \Rightarrow  \\ \quad \quad (\operatorname{vnrm}(m)(\operatorname{vzero}(m)) = 0) \end{array}}}
\vnbentry{\texttt{nvs-\allowbreak{}ms-\allowbreak{}pts}}{\vnbstmt{\begin{array}{@{}l@{}} \forall m. \\ \quad (\operatorname{pts}(\operatorname{nvs\text{-}metric\text{-}space}(m)) \simeq \operatorname{vec}(m)) \end{array}}}
\vnbentry{\texttt{nvs-\allowbreak{}metric-\allowbreak{}distance}}{\vnbstmt{\begin{array}{@{}l@{}} \forall m.; \forall u, v \in \operatorname{vec}(m). \\ \quad \operatorname{is\text{-}normed\text{-}vector\text{-}space}(m) \Rightarrow  \\ \quad \quad \operatorname{dist}(\operatorname{nvs\text{-}metric\text{-}space}(m))(u, v) \\ \quad \quad =\; \operatorname{vnrm}(m)(\operatorname{vadd}(m)(u, \operatorname{vneg}(m)(v))) \end{array}}}
\vnbentry{\texttt{nvs-\allowbreak{}metric-\allowbreak{}distance-\allowbreak{}q}}{\vnbstmt{\begin{array}{@{}l@{}} \forall m.; \forall u, v \in \operatorname{vec}(m). \\ \quad \operatorname{is\text{-}normed\text{-}vector\text{-}space}(m) \Rightarrow  \\ \quad \quad (\operatorname{dist}(\operatorname{nvs\text{-}metric\text{-}space}(m))(u, v) \simeq \operatorname{vnrm}(m)(\operatorname{vadd}(m)(u, \operatorname{vneg}(m)(v)))) \end{array}}}
\vnbentry{\texttt{nvs-\allowbreak{}dist-\allowbreak{}shift}}{\vnbstmt{\begin{array}{@{}l@{}} \forall m.; \forall u, w \in \operatorname{vec}(m). \\ \quad \operatorname{is\text{-}normed\text{-}vector\text{-}space}(m) \Rightarrow  \\ \quad \quad \operatorname{dist}(\operatorname{nvs\text{-}metric\text{-}space}(m))(u, \operatorname{vadd}(m)(u, w)) \\ \quad \quad =\; \operatorname{vnrm}(m)(w) \end{array}}}

\vnbfile{theorem-library/rake-hb-submodules.scm}{theorem-library/rake-hb-submodules.scm}
\noindent{\small rake-hb-submodules.scm -- the two SUBMODULE leaves of the Hahn-Banach cluster, PROVEN, and the normed-vector-space action bricks they need.}\par
\vnbentry{\texttt{line-\allowbreak{}unfold}}{\vnbstmt{\begin{array}{@{}l@{}} \forall m, v. \\ \quad (\operatorname{line}(m, v) \simeq \{y \in \operatorname{vec}(m) : \exists r \in \mathbb{R}.\; (y = \operatorname{act}(m)(r, v))\}) \end{array}}}
\vnbentry{\texttt{line-\allowbreak{}is-\allowbreak{}submodule}}{\vnbstmt{\begin{array}{@{}l@{}} \forall m, v. \\ \quad (\operatorname{is\text{-}normed\text{-}vector\text{-}space}(m) \wedge (v \in \operatorname{vec}(m))) \Rightarrow  \\ \quad \quad \operatorname{is\text{-}submodule}(m, \operatorname{line}(m, v)) \end{array}}}
\vnbentry{\texttt{span-\allowbreak{}add-\allowbreak{}one-\allowbreak{}submodule}}{\vnbstmt{\begin{array}{@{}l@{}} \forall m, t, v. \\ \quad \operatorname{is\text{-}normed\text{-}vector\text{-}space}(m) \wedge \\ \quad \operatorname{is\text{-}submodule}(m, t) \wedge \\ \quad (v \in \operatorname{vec}(m)) \Rightarrow \\ \quad \quad \operatorname{is\text{-}submodule}(m, \operatorname{span\text{-}add\text{-}one}(m, t, v)) \end{array}}}

\vnbfile{theorem-library/rake-span-add-one-guarded.scm}{theorem-library/rake-span-add-one-guarded.scm}
\noindent{\small rake-span-add-one-guarded.scm -- the GUARDED forms of span-add-one-superset and span-add-one-has-v, PROVEN.}\par
\vnbentry{\texttt{span-\allowbreak{}add-\allowbreak{}one-\allowbreak{}superset-\allowbreak{}nvs}}{\vnbstmt{\begin{array}{@{}l@{}} \forall m, t, v. \\ \quad \operatorname{is\text{-}normed\text{-}vector\text{-}space}(m) \wedge \\ \quad \operatorname{is\text{-}submodule}(m, t) \wedge \\ \quad (v \in \operatorname{vec}(m)) \Rightarrow \\ \quad \quad (t \subseteq \operatorname{span\text{-}add\text{-}one}(m, t, v)) \end{array}}}
\vnbentry{\texttt{span-\allowbreak{}add-\allowbreak{}one-\allowbreak{}has-\allowbreak{}v-\allowbreak{}nvs}}{\vnbstmt{\begin{array}{@{}l@{}} \forall m, t, v. \\ \quad \operatorname{is\text{-}normed\text{-}vector\text{-}space}(m) \wedge \\ \quad \operatorname{is\text{-}submodule}(m, t) \wedge \\ \quad (v \in \operatorname{vec}(m)) \Rightarrow \\ \quad \quad (v \in \operatorname{span\text{-}add\text{-}one}(m, t, v)) \end{array}}}

\vnbfile{theorem-library/rake-hb-leaves-2.scm}{theorem-library/rake-hb-leaves-2.scm}
\noindent{\small theorem-library/rake-hb-leaves-2.scm -- three Hahn-Banach leaves, PROVEN modulo 0: good-sub-submodule, good-sub-in-power, line-has-v.}\par
\vnbentry{\texttt{good-\allowbreak{}sub-\allowbreak{}submodule}}{\vnbstmt{\begin{array}{@{}l@{}} \forall m, s, f, t. \\ \quad \operatorname{good\text{-}sub}(m, s, f, t) \Rightarrow  \\ \quad \quad \operatorname{is\text{-}submodule}(m, t) \end{array}}}
\vnbentry{\texttt{good-\allowbreak{}sub-\allowbreak{}in-\allowbreak{}power}}{\vnbstmt{\begin{array}{@{}l@{}} \forall m, s, f, t. \\ \quad \operatorname{good\text{-}sub}(m, s, f, t) \Rightarrow  \\ \quad \quad (t \in \operatorname{power}(\operatorname{vec}(m))) \end{array}}}
\vnbentry{\texttt{line-\allowbreak{}has-\allowbreak{}v}}{\vnbstmt{\begin{array}{@{}l@{}} \forall m, v. \\ \quad (\operatorname{is\text{-}normed\text{-}vector\text{-}space}(m) \wedge (v \in \operatorname{vec}(m))) \Rightarrow  \\ \quad \quad (v \in \operatorname{line}(m, v)) \end{array}}}

\vnbfile{theorem-library/rake-hb-gap.scm}{theorem-library/rake-hb-gap.scm}
\noindent{\small rake-hb-gap.scm -- hb-gap (theorem-library/hahn-banach-proof.scm:24), the analytic core of the one-dimension Hahn-Banach step, PROVEN.}\par
\vnbentry{\texttt{hbg-\allowbreak{}scale-\allowbreak{}pos}}{\vnbstmt{\begin{array}{@{}l@{}} \forall d, t, u, a, o \in \mathbb{R}. \\ \quad ((0 \leq o) \wedge (((-u - (d \cdot t)) \leq a) \wedge (a \leq ((d \cdot t) - u)))) \Rightarrow  \\ \quad \quad (\lvert ((o \cdot u) + (o \cdot a)) \rvert \leq (d \cdot (o \cdot t))) \end{array}}}
\vnbentry{\texttt{hbg-\allowbreak{}scale-\allowbreak{}neg}}{\vnbstmt{\begin{array}{@{}l@{}} \forall d, t, u, a, o \in \mathbb{R}. \\ \quad ((o \leq 0) \wedge (((-u - (d \cdot t)) \leq a) \wedge (a \leq ((d \cdot t) - u)))) \Rightarrow  \\ \quad \quad (\lvert ((o \cdot u) + (o \cdot a)) \rvert \leq (d \cdot (-o \cdot t))) \end{array}}}
\vnbentry{\texttt{hbg-\allowbreak{}neg-\allowbreak{}one-\allowbreak{}cancel}}{\vnbstmt{\begin{array}{@{}l@{}} \forall m.; \forall y, v \in \operatorname{vec}(m). \\ \quad \operatorname{is\text{-}normed\text{-}vector\text{-}space}(m) \Rightarrow  \\ \quad \quad (\operatorname{vadd}(m)(\operatorname{vadd}(m)(y, v), \operatorname{act}(m)(-1, v)) = y) \end{array}}}
\vnbentry{\texttt{hbg-\allowbreak{}shift}}{\vnbstmt{\begin{array}{@{}l@{}} \forall m.; \forall w, z, v \in \operatorname{vec}(m). \\ \quad \operatorname{is\text{-}normed\text{-}vector\text{-}space}(m) \Rightarrow  \\ \quad \quad \operatorname{vadd}(m)(w, \operatorname{act}(m)(-1, z)) \\ \quad \quad =\; \operatorname{vadd}(m)(\operatorname{vadd}(m)(w, v), \operatorname{act}(m)(-1, \operatorname{vadd}(m)(z, v))) \end{array}}}
\vnbentry{\texttt{hbg-\allowbreak{}diff-\allowbreak{}norm}}{\vnbstmt{\begin{array}{@{}l@{}} \forall m.; \forall w, z, v \in \operatorname{vec}(m). \\ \quad \operatorname{is\text{-}normed\text{-}vector\text{-}space}(m) \Rightarrow  \\ \quad \quad \operatorname{vnrm}(m)(\operatorname{vadd}(m)(w, \operatorname{act}(m)(-1, z))) \\ \quad \quad \leq\; (\operatorname{vnrm}(m)(\operatorname{vadd}(m)(w, v)) + \operatorname{vnrm}(m)(\operatorname{vadd}(m)(z, v))) \end{array}}}
\vnbentry{\texttt{hbg-\allowbreak{}scale-\allowbreak{}back}}{\vnbstmt{\begin{array}{@{}l@{}} \forall m.; \forall y \in \operatorname{vec}(m), r \in \mathbb{R}. \\ \quad (\operatorname{is\text{-}normed\text{-}vector\text{-}space}(m) \wedge \neg (r = 0)) \Rightarrow  \\ \quad \quad (\operatorname{act}(m)(r, \operatorname{act}(m)(\operatorname{recip}(r), y)) = y) \end{array}}}
\vnbentry{\texttt{hbg-\allowbreak{}scale-\allowbreak{}decomp}}{\vnbstmt{\begin{array}{@{}l@{}} \forall m.; \forall y, v \in \operatorname{vec}(m), r \in \mathbb{R}. \\ \quad (\operatorname{is\text{-}normed\text{-}vector\text{-}space}(m) \wedge \neg (r = 0)) \Rightarrow  \\ \quad \quad \operatorname{vadd}(m)(y, \operatorname{act}(m)(r, v)) \\ \quad \quad =\; \operatorname{act}(m)(r, \operatorname{vadd}(m)(\operatorname{act}(m)(\operatorname{recip}(r), y), v)) \end{array}}}
\vnbentry{\texttt{hbg-\allowbreak{}linfun-\allowbreak{}app-\allowbreak{}real}}{\vnbstmt{\begin{array}{@{}l@{}} \forall m, s, f, x. \\ \quad (\operatorname{is\text{-}linear\text{-}functional\text{-}on}(m, s, f) \wedge (x \in s)) \Rightarrow  \\ \quad \quad (f(x) \in \mathbb{R}) \end{array}}}
\vnbentry{\texttt{hbg-\allowbreak{}sep}}{\vnbstmt{\begin{array}{@{}l@{}} \forall m, s, f, v, c.; \forall z, w \in s. \\ \quad \operatorname{is\text{-}normed\text{-}vector\text{-}space}(m) \wedge \\ \quad \operatorname{is\text{-}submodule}(m, s) \wedge \\ \quad \operatorname{is\text{-}linear\text{-}functional\text{-}on}(m, s, f) \wedge \\ \quad (v \in \operatorname{vec}(m)) \wedge \\ \quad (c \in \mathbb{R}) \wedge \\ \quad \forall x \in s.\; (\lvert f(x) \rvert \leq (c \cdot \operatorname{vnrm}(m)(x))) \wedge \\ \quad (0 \leq c) \Rightarrow \\ \quad \quad (-f(z) - (c \cdot \operatorname{vnrm}(m)(\operatorname{vadd}(m)(z, v)))) \\ \quad \quad \leq\; ((c \cdot \operatorname{vnrm}(m)(\operatorname{vadd}(m)(w, v))) - f(w)) \end{array}}}
\vnbentry{\texttt{hbg-\allowbreak{}core}}{\vnbstmt{\begin{array}{@{}l@{}} \forall m, s, f, v, c, a.; \forall y \in s, r \in \mathbb{R}. \\ \quad \operatorname{is\text{-}normed\text{-}vector\text{-}space}(m) \wedge \\ \quad \operatorname{is\text{-}submodule}(m, s) \wedge \\ \quad \operatorname{is\text{-}linear\text{-}functional\text{-}on}(m, s, f) \wedge \\ \quad (v \in \operatorname{vec}(m)) \wedge \\ \quad (c \in \mathbb{R}) \wedge \\ \quad (a \in \mathbb{R}) \wedge \\ \quad \forall x \in s.\; (\lvert f(x) \rvert \leq (c \cdot \operatorname{vnrm}(m)(x))) \wedge \\ \quad \forall u \in s. \\ \quad \quad ((-f(u) - (c \cdot \operatorname{vnrm}(m)(\operatorname{vadd}(m)(u, v)))) \leq a) \wedge \\ \quad \forall u \in s. \\ \quad \quad (a \leq ((c \cdot \operatorname{vnrm}(m)(\operatorname{vadd}(m)(u, v))) - f(u))) \Rightarrow \\ \quad \quad \lvert (f(y) + (r \cdot a)) \rvert \\ \quad \quad \leq\; (c \cdot \operatorname{vnrm}(m)(\operatorname{vadd}(m)(y, \operatorname{act}(m)(r, v)))) \end{array}}}
\vnbentry{\texttt{hb-\allowbreak{}gap-\allowbreak{}bound}}{\vnbstmt{\begin{array}{@{}l@{}} \forall m, s, f, v, c. \\ \quad \operatorname{is\text{-}normed\text{-}vector\text{-}space}(m) \wedge \\ \quad \operatorname{is\text{-}submodule}(m, s) \wedge \\ \quad \operatorname{is\text{-}linear\text{-}functional\text{-}on}(m, s, f) \wedge \\ \quad (v \in \operatorname{vec}(m)) \wedge \\ \quad (c \in \mathbb{R}) \wedge \\ \quad \forall x \in s.\; (\lvert f(x) \rvert \leq (c \cdot \operatorname{vnrm}(m)(x))) \wedge \\ \quad (0 \leq c) \Rightarrow \\ \quad \quad \exists a \in \mathbb{R}. \forall y \in s, r \in \mathbb{R}. \\ \quad \quad \quad \lvert (f(y) + (r \cdot a)) \rvert \\ \quad \quad \quad \leq\; (c \cdot \operatorname{vnrm}(m)(\operatorname{vadd}(m)(y, \operatorname{act}(m)(r, v)))) \end{array}}}
\vnbentry{\texttt{hb-\allowbreak{}gap}}{\vnbstmt{\begin{array}{@{}l@{}} \forall m, s, f, v. \\ \quad \operatorname{is\text{-}normed\text{-}vector\text{-}space}(m) \wedge \\ \quad \operatorname{is\text{-}submodule}(m, s) \wedge \\ \quad \operatorname{is\text{-}bounded\text{-}linear\text{-}functional\text{-}on}(m, s, f) \wedge \\ \quad (v \in \operatorname{vec}(m)) \wedge \\ \quad \neg (v \in s) \Rightarrow \\ \quad \quad \exists a \in \mathbb{R}. \forall y \in s, r \in \mathbb{R}. \\ \quad \quad \quad \lvert (f(y) + (r \cdot a)) \rvert \\ \quad \quad \quad \leq\; (\operatorname{dual\text{-}norm\text{-}on}(m, s, f) \cdot \operatorname{vnrm}(m)(\operatorname{vadd}(m)(y, \operatorname{act}(m)(r, v)))) \end{array}}}

\vnbfile{theorem-library/rake-hb-extend-construct.scm}{theorem-library/rake-hb-extend-construct.scm}
\noindent{\small rake-hb-extend-construct.scm -- hb-extend-construct (theorem-library/ hahn-banach-proof.scm:44), rake batch 7, assignment 7-C, second leaf.}\par
\vnbentry{\texttt{hbx-\allowbreak{}s-\allowbreak{}in-\allowbreak{}vec}}{\vnbstmt{\begin{array}{@{}l@{}} \forall m, s, p. \\ \quad \operatorname{is\text{-}normed\text{-}vector\text{-}space}(m) \wedge \\ \quad \operatorname{is\text{-}submodule}(m, s) \wedge \\ \quad (p \in s) \Rightarrow \\ \quad \quad (p \in \operatorname{vec}(m)) \end{array}}}
\vnbentry{\texttt{hbx-\allowbreak{}span-\allowbreak{}in-\allowbreak{}vec}}{\vnbstmt{\begin{array}{@{}l@{}} \forall m, s, v.; \forall w \in \operatorname{span\text{-}add\text{-}one}(m, s, v). \\ \quad (w \in \operatorname{vec}(m)) \end{array}}}
\vnbentry{\texttt{hbx-\allowbreak{}span-\allowbreak{}decomp}}{\vnbstmt{\begin{array}{@{}l@{}} \forall m, s, v.; \forall w \in \operatorname{span\text{-}add\text{-}one}(m, s, v). \exists y \in s, r \in \mathbb{R}. \\ \quad (w = \operatorname{vadd}(m)(y, \operatorname{act}(m)(r, v))) \end{array}}}
\vnbentry{\texttt{hbx-\allowbreak{}span-\allowbreak{}mem-\allowbreak{}eq}}{\vnbstmt{\begin{array}{@{}l@{}} \forall m, s, v, w, p, c. \\ \quad \operatorname{is\text{-}normed\text{-}vector\text{-}space}(m) \wedge \\ \quad \operatorname{is\text{-}submodule}(m, s) \wedge \\ \quad (v \in \operatorname{vec}(m)) \wedge \\ \quad (p \in s) \wedge \\ \quad (c \in \mathbb{R}) \wedge \\ \quad (w = \operatorname{vadd}(m)(p, \operatorname{act}(m)(c, v))) \Rightarrow \\ \quad \quad (w \in \operatorname{span\text{-}add\text{-}one}(m, s, v)) \end{array}}}
\vnbentry{\texttt{hbx-\allowbreak{}decomp-\allowbreak{}unique}}{\vnbstmt{\begin{array}{@{}l@{}} \forall m, s, v, p_{1}, c_{1}, p_{2}, c_{2}. \\ \quad \operatorname{is\text{-}normed\text{-}vector\text{-}space}(m) \wedge \\ \quad \operatorname{is\text{-}submodule}(m, s) \wedge \\ \quad (v \in \operatorname{vec}(m)) \wedge \\ \quad \neg (v \in s) \wedge \\ \quad (p_{1} \in s) \wedge \\ \quad (c_{1} \in \mathbb{R}) \wedge \\ \quad (p_{2} \in s) \wedge \\ \quad (c_{2} \in \mathbb{R}) \wedge \\ \quad \operatorname{vadd}(m)(p_{1}, \operatorname{act}(m)(c_{1}, v)) \\ \quad =\; \operatorname{vadd}(m)(p_{2}, \operatorname{act}(m)(c_{2}, v)) \Rightarrow \\ \quad \quad ((c_{1} = c_{2}) \wedge (p_{1} = p_{2})) \end{array}}}
\vnbentry{\texttt{hbx-\allowbreak{}vadd-\allowbreak{}exchange}}{\vnbstmt{\begin{array}{@{}l@{}} \forall m, a, b, c, d. \\ \quad \operatorname{is\text{-}normed\text{-}vector\text{-}space}(m) \wedge \\ \quad (a \in \operatorname{vec}(m)) \wedge \\ \quad (b \in \operatorname{vec}(m)) \wedge \\ \quad (c \in \operatorname{vec}(m)) \wedge \\ \quad (d \in \operatorname{vec}(m)) \Rightarrow \\ \quad \quad \operatorname{vadd}(m)(\operatorname{vadd}(m)(a, b), \operatorname{vadd}(m)(c, d)) \\ \quad \quad =\; \operatorname{vadd}(m)(\operatorname{vadd}(m)(a, c), \operatorname{vadd}(m)(b, d)) \end{array}}}
\vnbentry{\texttt{hbx-\allowbreak{}add-\allowbreak{}decomp}}{\vnbstmt{\begin{array}{@{}l@{}} \forall m, v, y_{1}, c_{1}, y_{2}, c_{2}. \\ \quad \operatorname{is\text{-}normed\text{-}vector\text{-}space}(m) \wedge \\ \quad (v \in \operatorname{vec}(m)) \wedge \\ \quad (y_{1} \in \operatorname{vec}(m)) \wedge \\ \quad (c_{1} \in \mathbb{R}) \wedge \\ \quad (y_{2} \in \operatorname{vec}(m)) \wedge \\ \quad (c_{2} \in \mathbb{R}) \Rightarrow \\ \quad \quad \operatorname{vadd}(m)(\operatorname{vadd}(m)(y_{1}, \operatorname{act}(m)(c_{1}, v)), \operatorname{vadd}(m)(y_{2}, \operatorname{act}(m)(c_{2}, v))) \\ \quad \quad =\; \operatorname{vadd}(m)(\operatorname{vadd}(m)(y_{1}, y_{2}), \operatorname{act}(m)((c_{1} + c_{2}), v)) \end{array}}}
\vnbentry{\texttt{hbx-\allowbreak{}f-\allowbreak{}add}}{\vnbstmt{\begin{array}{@{}l@{}} \forall m, s, f, x, y. \\ \quad \operatorname{is\text{-}linear\text{-}functional\text{-}on}(m, s, f) \wedge \\ \quad (x \in s) \wedge \\ \quad (y \in s) \Rightarrow \\ \quad \quad (f(\operatorname{vadd}(m)(x, y)) = (f(x) + f(y))) \end{array}}}
\vnbentry{\texttt{hbx-\allowbreak{}f-\allowbreak{}act}}{\vnbstmt{\begin{array}{@{}l@{}} \forall m, s, f, r, x. \\ \quad \operatorname{is\text{-}linear\text{-}functional\text{-}on}(m, s, f) \wedge \\ \quad (r \in \mathbb{R}) \wedge \\ \quad (x \in s) \Rightarrow \\ \quad \quad (f(\operatorname{act}(m)(r, x)) = (r \cdot f(x))) \end{array}}}
\vnbentry{\texttt{hbx-\allowbreak{}value}}{\vnbstmt{\begin{array}{@{}l@{}} \forall m, s, f, v, a, u, p, c. \\ \quad \operatorname{is\text{-}normed\text{-}vector\text{-}space}(m) \wedge \\ \quad \operatorname{is\text{-}submodule}(m, s) \wedge \\ \quad \operatorname{is\text{-}linear\text{-}functional\text{-}on}(m, s, f) \wedge \\ \quad (v \in \operatorname{vec}(m)) \wedge \\ \quad \neg (v \in s) \wedge \\ \quad (a \in \mathbb{R}) \wedge \\ \quad (p \in s) \wedge \\ \quad (c \in \mathbb{R}) \wedge \\ \quad (u = \operatorname{vadd}(m)(p, \operatorname{act}(m)(c, v))) \Rightarrow \\ \quad \quad (\iota(\mathit{tv}) = (f(p) + (c \cdot a))) \end{array}}}
\vnbentry{\texttt{hb-\allowbreak{}extend-\allowbreak{}construct}}{\vnbstmt{\begin{array}{@{}l@{}} \forall m, s, f, v, a. \\ \quad \operatorname{is\text{-}normed\text{-}vector\text{-}space}(m) \wedge \\ \quad \operatorname{is\text{-}submodule}(m, s) \wedge \\ \quad \operatorname{is\text{-}linear\text{-}functional\text{-}on}(m, s, f) \wedge \\ \quad (v \in \operatorname{vec}(m)) \wedge \\ \quad \neg (v \in s) \wedge \\ \quad (a \in \mathbb{R}) \Rightarrow \\ \quad \quad \exists g. \\ \quad \quad \quad \operatorname{is\text{-}linear\text{-}functional\text{-}on}(m, \operatorname{span\text{-}add\text{-}one}(m, s, v), g) \wedge \\ \quad \quad \quad \operatorname{extends\text{-}on}(s, g, f) \wedge \\ \quad \quad \quad \forall y \in s, r \in \mathbb{R}. \\ \quad \quad \quad \quad (g(\operatorname{vadd}(m)(y, \operatorname{act}(m)(r, v))) = (f(y) + (r \cdot a))) \end{array}}}

\vnbfile{theorem-library/rake-line-functional.scm}{theorem-library/rake-line-functional.scm}
\noindent{\small rake-line-functional.scm -- line-functional-exists (theorem-library/ norm-as-sup-proof.scm:43), rake batch 7, assignment 7-C.}\par
\vnbentry{\texttt{lfe-\allowbreak{}line-\allowbreak{}membership}}{\vnbstmt{\begin{array}{@{}l@{}} \forall m, v, w. \\ \quad ((w \in \operatorname{line}(m, v)) \Leftrightarrow ((w \in \operatorname{vec}(m)) \wedge \exists r \in \mathbb{R}.\; (w = \operatorname{act}(m)(r, v)))) \end{array}}}
\vnbentry{\texttt{lfe-\allowbreak{}line-\allowbreak{}in-\allowbreak{}vec}}{\vnbstmt{\begin{array}{@{}l@{}} \forall m, v.; \forall w \in \operatorname{line}(m, v). \\ \quad (w \in \operatorname{vec}(m)) \end{array}}}
\vnbentry{\texttt{lfe-\allowbreak{}line-\allowbreak{}subset}}{\vnbstmt{\forall m, v.\; (\operatorname{line}(m, v) \subseteq \operatorname{vec}(m))}}
\vnbentry{\texttt{lfe-\allowbreak{}line-\allowbreak{}coef}}{\vnbstmt{\begin{array}{@{}l@{}} \forall m, v.; \forall w \in \operatorname{line}(m, v). \exists r \in \mathbb{R}. \\ \quad (w = \operatorname{act}(m)(r, v)) \end{array}}}
\vnbentry{\texttt{lfe-\allowbreak{}line-\allowbreak{}mem-\allowbreak{}eq}}{\vnbstmt{\begin{array}{@{}l@{}} \forall m, v, w, r. \\ \quad \operatorname{is\text{-}normed\text{-}vector\text{-}space}(m) \wedge \\ \quad (v \in \operatorname{vec}(m)) \wedge \\ \quad (r \in \mathbb{R}) \wedge \\ \quad (w = \operatorname{act}(m)(r, v)) \Rightarrow \\ \quad \quad (w \in \operatorname{line}(m, v)) \end{array}}}
\vnbentry{\texttt{lfe-\allowbreak{}coef-\allowbreak{}unique}}{\vnbstmt{\begin{array}{@{}l@{}} \forall m, v.; \forall a, b \in \mathbb{R}. \\ \quad \operatorname{is\text{-}normed\text{-}vector\text{-}space}(m) \wedge \\ \quad (v \in \operatorname{vec}(m)) \wedge \\ \quad (\operatorname{act}(m)(a, v) = \operatorname{act}(m)(b, v)) \Rightarrow \\ \quad \quad ((a \cdot \operatorname{vnrm}(m)(v)) = (b \cdot \operatorname{vnrm}(m)(v))) \end{array}}}
\vnbentry{\texttt{lfe-\allowbreak{}coef-\allowbreak{}value}}{\vnbstmt{\begin{array}{@{}l@{}} \forall m, v, u, c. \\ \quad \operatorname{is\text{-}normed\text{-}vector\text{-}space}(m) \wedge \\ \quad (v \in \operatorname{vec}(m)) \wedge \\ \quad (c \in \mathbb{R}) \wedge \\ \quad (u = \operatorname{act}(m)(c, v)) \Rightarrow \\ \quad \quad (\iota(t) = (c \cdot \operatorname{vnrm}(m)(v))) \end{array}}}
\vnbentry{\texttt{line-\allowbreak{}functional-\allowbreak{}exists}}{\vnbstmt{\begin{array}{@{}l@{}} \forall m, v. \\ \quad (\operatorname{is\text{-}normed\text{-}vector\text{-}space}(m) \wedge (v \in \operatorname{vec}(m))) \Rightarrow  \\ \quad \quad \exists f. \\ \quad \quad \quad \operatorname{is\text{-}bounded\text{-}linear\text{-}functional\text{-}on}(m, \operatorname{line}(m, v), f) \wedge \\ \quad \quad \quad (\operatorname{dual\text{-}norm\text{-}on}(m, \operatorname{line}(m, v), f) \leq 1) \wedge \\ \quad \quad \quad (f(v) = \operatorname{vnrm}(m)(v)) \end{array}}}

\vnbfile{theorem-library/hahn-banach-proof.scm}{theorem-library/hahn-banach-proof.scm}
\noindent{\small hahn-banach-proof.scm -- Hahn-Banach, the one-dimension-higher extension step (real normed vector space), MACHINE-PROVEN modulo the warranted core.}\par
\vnbentry{\texttt{hahn-\allowbreak{}banach-\allowbreak{}extend-\allowbreak{}one}}{\vnbstmt{\begin{array}{@{}l@{}} \forall m, s, f, v. \\ \quad \operatorname{is\text{-}normed\text{-}vector\text{-}space}(m) \wedge \\ \quad \operatorname{is\text{-}submodule}(m, s) \wedge \\ \quad \operatorname{is\text{-}bounded\text{-}linear\text{-}functional\text{-}on}(m, s, f) \wedge \\ \quad (v \in \operatorname{vec}(m)) \wedge \\ \quad \neg (v \in s) \Rightarrow \\ \quad \quad \exists g. \\ \quad \quad \quad \operatorname{is\text{-}linear\text{-}functional\text{-}on}(m, \operatorname{span\text{-}add\text{-}one}(m, s, v), g) \wedge \\ \quad \quad \quad \operatorname{extends\text{-}on}(s, g, f) \wedge \\ \quad \quad \quad \forall w \in \operatorname{span\text{-}add\text{-}one}(m, s, v). \\ \quad \quad \quad \quad \lvert g(w) \rvert \\ \quad \quad \quad \quad \leq\; (\operatorname{dual\text{-}norm\text{-}on}(m, s, f) \cdot \operatorname{vnrm}(m)(w)) \end{array}}}

\vnbfile{theorem-library/noetherian-maximal-proof.scm}{theorem-library/noetherian-maximal-proof.scm}
\noindent{\small noetherian-maximal-proof.scm -- the noetherian maximal-element principle, MACHINE-PROVEN, and its specialisation hb-good-has-maximal (the input that makes finite-dimensional Hahn-Banach terminate).}\par
\vnbentry{\texttt{noetherian-\allowbreak{}set-\allowbreak{}has-\allowbreak{}maximal}}{\vnbstmt{\begin{array}{@{}l@{}} \forall m.; \forall g \in \mathbf{Set}. \\ \quad \operatorname{is\text{-}noetherian}(m) \wedge \\ \quad (g \subseteq \operatorname{power}(\operatorname{vec}(m))) \wedge \\ \quad \forall t \in g.\; \operatorname{is\text{-}submodule}(m, t) \wedge \\ \quad \exists a.\; (a \in g) \Rightarrow \\ \quad \quad \exists t \in g.\; \forall u \in g.\; ((t \subseteq u) \Rightarrow (t = u)) \end{array}}}
\vnbentry{\texttt{hb-\allowbreak{}good-\allowbreak{}has-\allowbreak{}maximal}}{\vnbstmt{\begin{array}{@{}l@{}} \forall m, s, f. \\ \quad \operatorname{is\text{-}normed\text{-}vector\text{-}space}(m) \wedge \\ \quad \operatorname{is\text{-}finite\text{-}dimensional}(\operatorname{normed\text{-}vector\text{-}space\text{-}as\text{-}module}(m)) \wedge \\ \quad \operatorname{is\text{-}submodule}(m, s) \wedge \\ \quad \operatorname{is\text{-}bounded\text{-}linear\text{-}functional\text{-}on}(m, s, f) \Rightarrow \\ \quad \quad \exists t. \\ \quad \quad \quad \operatorname{good\text{-}sub}(m, s, f, t) \wedge \\ \quad \quad \quad \forall u. \\ \quad \quad \quad \quad (\operatorname{good\text{-}sub}(m, s, f, u) \wedge (t \subseteq u)) \Rightarrow  \\ \quad \quad \quad \quad \quad (t = u) \end{array}}}

\vnbfile{theorem-library/hahn-banach-full-proof.scm}{theorem-library/hahn-banach-full-proof.scm}
\noindent{\small RETIRED 2026-09-17 (proven): le-bound-mono (was an add-to-pss here) -- theorem-library/rake-analysis2.scm hahn-banach-full-proof.scm -- FULL Hahn-Banach for a finite-dimensional real normed vector space, MACHINE-PROVEN modulo a warranted plumbing core.}\par
\vnbentry{\texttt{good-\allowbreak{}step}}{\vnbstmt{\begin{array}{@{}l@{}} \forall m, s, f, t, x. \\ \quad \operatorname{is\text{-}normed\text{-}vector\text{-}space}(m) \wedge \\ \quad \operatorname{is\text{-}bounded\text{-}linear\text{-}functional\text{-}on}(m, s, f) \wedge \\ \quad \operatorname{good\text{-}sub}(m, s, f, t) \wedge \\ \quad (x \in \operatorname{vec}(m)) \wedge \\ \quad \neg (x \in t) \Rightarrow \\ \quad \quad \operatorname{good\text{-}sub}(m, s, f, \operatorname{span\text{-}add\text{-}one}(m, t, x)) \end{array}}}
\vnbentry{\texttt{hahn-\allowbreak{}banach}}{\vnbstmt{\begin{array}{@{}l@{}} \forall m, s, f. \\ \quad \operatorname{is\text{-}normed\text{-}vector\text{-}space}(m) \wedge \\ \quad \operatorname{is\text{-}finite\text{-}dimensional}(\operatorname{normed\text{-}vector\text{-}space\text{-}as\text{-}module}(m)) \wedge \\ \quad \operatorname{is\text{-}submodule}(m, s) \wedge \\ \quad \operatorname{is\text{-}bounded\text{-}linear\text{-}functional\text{-}on}(m, s, f) \Rightarrow \\ \quad \quad \exists g. \\ \quad \quad \quad \operatorname{is\text{-}linear\text{-}functional\text{-}on}(m, \operatorname{vec}(m), g) \wedge \\ \quad \quad \quad \operatorname{extends\text{-}on}(s, g, f) \wedge \\ \quad \quad \quad \forall w \in \operatorname{vec}(m). \\ \quad \quad \quad \quad \lvert g(w) \rvert \\ \quad \quad \quad \quad \leq\; (\operatorname{dual\text{-}norm\text{-}on}(m, s, f) \cdot \operatorname{vnrm}(m)(w)) \end{array}}}

\vnbfile{theorem-library/norm-as-sup-proof.scm}{theorem-library/norm-as-sup-proof.scm}
\noindent{\small norm-as-sup-proof.scm -- the norm as a supremum of bounded functionals: for a finite-dimensional real NVS m and x in VEC(m), ||x|| = sup \{ |f(x)| : f a bounded linear functional, ||f|| $<$= 1 \}.}\par
\vnbentry{\texttt{norm-\allowbreak{}bounded-\allowbreak{}by-\allowbreak{}functionals}}{\vnbstmt{\begin{array}{@{}l@{}} \forall m, f, x. \\ \quad \operatorname{is\text{-}normed\text{-}vector\text{-}space}(m) \wedge \\ \quad \operatorname{is\text{-}bounded\text{-}linear\text{-}functional}(m, f) \wedge \\ \quad (x \in \operatorname{vec}(m)) \wedge \\ \quad (\operatorname{dual\text{-}norm}(m, f) \leq 1) \Rightarrow \\ \quad \quad (\lvert f(x) \rvert \leq \operatorname{vnrm}(m)(x)) \end{array}}}
\vnbentry{\texttt{norm-\allowbreak{}attained-\allowbreak{}by-\allowbreak{}functional}}{\vnbstmt{\begin{array}{@{}l@{}} \forall m, x. \\ \quad \operatorname{is\text{-}normed\text{-}vector\text{-}space}(m) \wedge \\ \quad \operatorname{is\text{-}finite\text{-}dimensional}(\operatorname{normed\text{-}vector\text{-}space\text{-}as\text{-}module}(m)) \wedge \\ \quad (x \in \operatorname{vec}(m)) \Rightarrow \\ \quad \quad \exists g. \\ \quad \quad \quad \operatorname{is\text{-}bounded\text{-}linear\text{-}functional}(m, g) \wedge \\ \quad \quad \quad (\operatorname{dual\text{-}norm}(m, g) \leq 1) \wedge \\ \quad \quad \quad (g(x) = \operatorname{vnrm}(m)(x)) \end{array}}}
\vnbentry{\texttt{norm-\allowbreak{}as-\allowbreak{}sup}}{\vnbstmt{\begin{array}{@{}l@{}} \forall m, x. \\ \quad \operatorname{is\text{-}normed\text{-}vector\text{-}space}(m) \wedge \\ \quad \operatorname{is\text{-}finite\text{-}dimensional}(\operatorname{normed\text{-}vector\text{-}space\text{-}as\text{-}module}(m)) \wedge \\ \quad (x \in \operatorname{vec}(m)) \Rightarrow \\ \quad \quad \forall f. \\ \quad \quad \quad \operatorname{is\text{-}bounded\text{-}linear\text{-}functional}(m, f) \wedge \\ \quad \quad \quad (\operatorname{dual\text{-}norm}(m, f) \leq 1) \Rightarrow \\ \quad \quad \quad \quad (\lvert f(x) \rvert \leq \operatorname{vnrm}(m)(x)) \wedge \\ \quad \quad \forall d \in \mathbb{R}. \\ \quad \quad \quad \forall f. \\ \quad \quad \quad \quad \operatorname{is\text{-}bounded\text{-}linear\text{-}functional}(m, f) \wedge \\ \quad \quad \quad \quad (\operatorname{dual\text{-}norm}(m, f) \leq 1) \Rightarrow \\ \quad \quad \quad \quad \quad (\lvert f(x) \rvert \leq d) \Rightarrow \\ \quad \quad \quad \quad (\operatorname{vnrm}(m)(x) \leq d) \end{array}}}

\vnbfile{theorem-library/vector-taylor-proof.scm}{theorem-library/vector-taylor-proof.scm}
\noindent{\small vector-taylor-proof.scm -- vector-valued Taylor with a norm remainder bound, REDUCED TO THE SCALAR CASE via a norm-attaining functional (Hahn-Banach payoff).}\par
\vnbentry{\texttt{diff-\allowbreak{}v-\allowbreak{}value-\allowbreak{}in-\allowbreak{}vec}}{\vnbstmt{\begin{array}{@{}l@{}} \forall m, f, a, l. \\ \quad \operatorname{is\text{-}diff\text{-}at\text{-}v}(m, f, a, l) \Rightarrow  \\ \quad \quad (l \in \operatorname{vec}(m)) \end{array}}}
\vnbentry{\texttt{taylor-\allowbreak{}v-\allowbreak{}derivs-\allowbreak{}in-\allowbreak{}fun}}{\vnbstmt{\begin{array}{@{}l@{}} \forall m, f, a, x, n.; \forall j \in \mathbb{N}. \\ \quad \operatorname{taylor\text{-}differentiable\text{-}v}(m, f, a, x, n) \wedge \\ \quad (a \in \mathbb{R}) \wedge \\ \quad (x \in \mathbb{R}) \wedge \\ \quad (a < x) \wedge \\ \quad (j \leq n) \Rightarrow \\ \quad \quad (\operatorname{nth\text{-}deriv\text{-}v}(m, f, j) \in (\mathbb{R} \to \operatorname{vec}(m))) \end{array}}}
\vnbentry{\texttt{dfun-\allowbreak{}v-\allowbreak{}mono}}{\vnbstmt{\begin{array}{@{}l@{}} \forall m, f, n, k.; \forall j \in \mathbb{N}. \\ \quad \forall j \in \mathbb{N}. \\ \quad \quad (j \leq n) \Rightarrow  \\ \quad \quad \quad (\operatorname{nth\text{-}deriv\text{-}v}(m, f, j) \in (\mathbb{R} \to \operatorname{vec}(m))) \wedge \\ \quad (n \in \mathbb{N}) \wedge \\ \quad (k \in \mathbb{N}) \wedge \\ \quad (k \leq n) \wedge \\ \quad (j \leq k) \Rightarrow \\ \quad \quad (\operatorname{nth\text{-}deriv\text{-}v}(m, f, j) \in (\mathbb{R} \to \operatorname{vec}(m))) \end{array}}}
\vnbentry{\texttt{nth-\allowbreak{}deriv-\allowbreak{}v-\allowbreak{}in-\allowbreak{}vec}}{\vnbstmt{\begin{array}{@{}l@{}} \forall m, f, k, t. \\ \quad \operatorname{is\text{-}normed\text{-}vector\text{-}space}(m) \wedge \\ \quad (f \in (\mathbb{R} \to \operatorname{vec}(m))) \wedge \\ \quad (k \in \mathbb{N}) \wedge \\ \quad (t \in \mathbb{R}) \wedge \\ \quad \forall j \in \mathbb{N}. \\ \quad \quad (j \leq k) \Rightarrow  \\ \quad \quad \quad (\operatorname{nth\text{-}deriv\text{-}v}(m, f, j) \in (\mathbb{R} \to \operatorname{vec}(m))) \Rightarrow \\ \quad \quad (\operatorname{nth\text{-}deriv\text{-}v}(m, f, k)(t) \in \operatorname{vec}(m)) \end{array}}}
\vnbentry{\texttt{taylor-\allowbreak{}v-\allowbreak{}deriv-\allowbreak{}in-\allowbreak{}vec}}{\vnbstmt{\begin{array}{@{}l@{}} \forall m, f, a, x, n, t. \\ \quad \operatorname{taylor\text{-}differentiable\text{-}v}(m, f, a, x, n) \wedge \\ \quad (a < t) \wedge \\ \quad (t < x) \wedge \\ \quad (n \in \mathbb{N}) \Rightarrow \\ \quad \quad (\operatorname{nth\text{-}deriv\text{-}v}(m, f, (n + 1))(t) \in \operatorname{vec}(m)) \end{array}}}
\vnbentry{\texttt{gof-\allowbreak{}in-\allowbreak{}fun}}{\vnbstmt{\begin{array}{@{}l@{}} \forall m, f, g. \\ \quad \operatorname{is\text{-}bounded\text{-}linear\text{-}functional}(m, g) \wedge \\ \quad (f \in (\mathbb{R} \to \operatorname{vec}(m))) \Rightarrow \\ \quad \quad (\operatorname{compose}(g, f) \in (\mathbb{R} \to \mathbb{R})) \end{array}}}
\vnbentry{\texttt{rkt-\allowbreak{}vtaylor-\allowbreak{}poly-\allowbreak{}in-\allowbreak{}vec-\allowbreak{}ind}}{\vnbstmt{\begin{array}{@{}l@{}} \forall n \in \mathbb{N}, m, f, a, x. \\ \quad \operatorname{is\text{-}normed\text{-}vector\text{-}space}(m) \wedge \\ \quad (f \in (\mathbb{R} \to \operatorname{vec}(m))) \wedge \\ \quad (a \in \mathbb{R}) \wedge \\ \quad (x \in \mathbb{R}) \wedge \\ \quad \forall j \in \mathbb{N}. \\ \quad \quad (j \leq n) \Rightarrow  \\ \quad \quad \quad (\operatorname{nth\text{-}deriv\text{-}v}(m, f, j) \in (\mathbb{R} \to \operatorname{vec}(m))) \Rightarrow \\ \quad \quad (\operatorname{taylor\text{-}poly\text{-}v}(m, f, a, x, n) \in \operatorname{vec}(m)) \end{array}}}
\vnbentry{\texttt{vtaylor-\allowbreak{}poly-\allowbreak{}in-\allowbreak{}vec}}{\vnbstmt{\begin{array}{@{}l@{}} \forall m, f, a, x, n. \\ \quad \operatorname{is\text{-}normed\text{-}vector\text{-}space}(m) \wedge \\ \quad (f \in (\mathbb{R} \to \operatorname{vec}(m))) \wedge \\ \quad (a \in \mathbb{R}) \wedge \\ \quad (x \in \mathbb{R}) \wedge \\ \quad (n \in \mathbb{N}) \wedge \\ \quad \forall j \in \mathbb{N}. \\ \quad \quad (j \leq n) \Rightarrow  \\ \quad \quad \quad (\operatorname{nth\text{-}deriv\text{-}v}(m, f, j) \in (\mathbb{R} \to \operatorname{vec}(m))) \Rightarrow \\ \quad \quad (\operatorname{taylor\text{-}poly\text{-}v}(m, f, a, x, n) \in \operatorname{vec}(m)) \end{array}}}
\vnbentry{\texttt{vtaylor-\allowbreak{}remainder-\allowbreak{}in-\allowbreak{}vec}}{\vnbstmt{\begin{array}{@{}l@{}} \forall m, f, a, x, n. \\ \quad \operatorname{is\text{-}normed\text{-}vector\text{-}space}(m) \wedge \\ \quad (f \in (\mathbb{R} \to \operatorname{vec}(m))) \wedge \\ \quad (a \in \mathbb{R}) \wedge \\ \quad (x \in \mathbb{R}) \wedge \\ \quad (n \in \mathbb{N}) \wedge \\ \quad \forall j \in \mathbb{N}. \\ \quad \quad (j \leq n) \Rightarrow  \\ \quad \quad \quad (\operatorname{nth\text{-}deriv\text{-}v}(m, f, j) \in (\mathbb{R} \to \operatorname{vec}(m))) \Rightarrow \\ \quad \quad \operatorname{vadd}(m)(f(x), \operatorname{vneg}(m)(\operatorname{taylor\text{-}poly\text{-}v}(m, f, a, x, n))) \\ \quad \quad \in\; \operatorname{vec}(m) \end{array}}}
\vnbentry{\texttt{r8e-\allowbreak{}blf-\allowbreak{}vneg}}{\vnbstmt{\begin{array}{@{}l@{}} \forall m, g.; \forall u \in \operatorname{vec}(m). \\ \quad \operatorname{is\text{-}normed\text{-}vector\text{-}space}(m) \wedge \\ \quad \operatorname{is\text{-}bounded\text{-}linear\text{-}functional}(m, g) \Rightarrow \\ \quad \quad (g(\operatorname{vneg}(m)(u)) = (0 - g(u))) \end{array}}}
\vnbentry{\texttt{r8e-\allowbreak{}blf-\allowbreak{}sub}}{\vnbstmt{\begin{array}{@{}l@{}} \forall m, g.; \forall u, v \in \operatorname{vec}(m). \\ \quad \operatorname{is\text{-}normed\text{-}vector\text{-}space}(m) \wedge \\ \quad \operatorname{is\text{-}bounded\text{-}linear\text{-}functional}(m, g) \Rightarrow \\ \quad \quad (g(\operatorname{vadd}(m)(u, \operatorname{vneg}(m)(v))) = (g(u) - g(v))) \end{array}}}
\vnbentry{\texttt{r8e-\allowbreak{}eps-\allowbreak{}over}}{\vnbstmt{\begin{array}{@{}l@{}} \forall c \in \mathbb{R}, e. \\ \quad ((0 \leq c) \wedge \operatorname{pos\text{-}rr}(e)) \Rightarrow  \\ \quad \quad \exists d. \\ \quad \quad \quad \operatorname{pos\text{-}rr}(d) \wedge \\ \quad \quad \quad \forall z \in \mathbb{R}.\; ((z \leq d) \Rightarrow ((c \cdot z) \leq e)) \end{array}}}
\vnbentry{\texttt{r8e-\allowbreak{}blf-\allowbreak{}compose-\allowbreak{}continuous}}{\vnbstmt{\begin{array}{@{}l@{}} \forall m, g, h, t. \\ \quad \operatorname{is\text{-}normed\text{-}vector\text{-}space}(m) \wedge \\ \quad \operatorname{is\text{-}bounded\text{-}linear\text{-}functional}(m, g) \wedge \\ \quad (h \in (\mathbb{R} \to \operatorname{vec}(m))) \wedge \\ \quad \operatorname{is\text{-}continuous\text{-}at}(\operatorname{rr\text{-}ms}, \operatorname{nvs\text{-}metric\text{-}space}(m), h, t) \Rightarrow \\ \quad \quad \operatorname{is\text{-}continuous\text{-}at}(\operatorname{rr\text{-}ms}, \operatorname{rr\text{-}ms}, \operatorname{compose}(g, h), t) \end{array}}}
\vnbentry{\texttt{r8e-\allowbreak{}blf-\allowbreak{}diff-\allowbreak{}commute}}{\vnbstmt{\begin{array}{@{}l@{}} \forall m, g, h, t, l. \\ \quad \operatorname{is\text{-}bounded\text{-}linear\text{-}functional}(m, g) \wedge \\ \quad \operatorname{is\text{-}diff\text{-}at\text{-}v}(m, h, t, l) \Rightarrow \\ \quad \quad \operatorname{is\text{-}diff\text{-}at}(\operatorname{compose}(g, h), t, g(l)) \end{array}}}
\vnbentry{\texttt{r8e-\allowbreak{}nth-\allowbreak{}deriv-\allowbreak{}v-\allowbreak{}diff}}{\vnbstmt{\begin{array}{@{}l@{}} \forall m, f, j, x. \\ \quad (j \in \mathbb{N}) \wedge \\ \quad (x \in \mathbb{R}) \wedge \\ \quad (\operatorname{nth\text{-}deriv\text{-}v}(m, f, j) \in (\mathbb{R} \to \operatorname{vec}(m))) \wedge \\ \quad (\operatorname{nth\text{-}deriv\text{-}v}(m, f, (j + 1)) \in (\mathbb{R} \to \operatorname{vec}(m))) \Rightarrow \\ \quad \quad \operatorname{is\text{-}diff\text{-}at\text{-}v}(m, \operatorname{nth\text{-}deriv\text{-}v}(m, f, j), x, \operatorname{nth\text{-}deriv\text{-}v}(m, f, (j + 1))(x)) \end{array}}}
\vnbentry{\texttt{r8e-\allowbreak{}gof-\allowbreak{}nth-\allowbreak{}deriv-\allowbreak{}fn}}{\vnbstmt{\begin{array}{@{}l@{}} \forall k \in \mathbb{N}, m, f, g. \\ \quad \operatorname{is\text{-}bounded\text{-}linear\text{-}functional}(m, g) \wedge \\ \quad (f \in (\mathbb{R} \to \operatorname{vec}(m))) \wedge \\ \quad \forall j \in \mathbb{N}. \\ \quad \quad (j \leq k) \Rightarrow  \\ \quad \quad \quad (\operatorname{nth\text{-}deriv\text{-}v}(m, f, j) \in (\mathbb{R} \to \operatorname{vec}(m))) \Rightarrow \\ \quad \quad \operatorname{nth\text{-}deriv}(\operatorname{compose}(g, f), k) \\ \quad \quad =\; \operatorname{compose}(g, \operatorname{nth\text{-}deriv\text{-}v}(m, f, k)) \end{array}}}
\vnbentry{\texttt{gof-\allowbreak{}nth-\allowbreak{}deriv-\allowbreak{}dfun}}{\vnbstmt{\begin{array}{@{}l@{}} \forall m, f, g, n, t. \\ \quad \operatorname{is\text{-}bounded\text{-}linear\text{-}functional}(m, g) \wedge \\ \quad (f \in (\mathbb{R} \to \operatorname{vec}(m))) \wedge \\ \quad (n \in \mathbb{N}) \wedge \\ \quad (t \in \mathbb{R}) \wedge \\ \quad \forall j \in \mathbb{N}. \\ \quad \quad (j \leq n) \Rightarrow  \\ \quad \quad \quad (\operatorname{nth\text{-}deriv\text{-}v}(m, f, j) \in (\mathbb{R} \to \operatorname{vec}(m))) \wedge \\ \quad \operatorname{is\text{-}diff\text{-}at\text{-}v}(m, \operatorname{nth\text{-}deriv\text{-}v}(m, f, n), t, \operatorname{nth\text{-}deriv\text{-}v}(m, f, (n + 1))(t)) \Rightarrow \\ \quad \quad \operatorname{nth\text{-}deriv}(\operatorname{compose}(g, f), (n + 1))(t) \\ \quad \quad =\; g(\operatorname{nth\text{-}deriv\text{-}v}(m, f, (n + 1))(t)) \end{array}}}
\vnbentry{\texttt{gof-\allowbreak{}taylor-\allowbreak{}diff-\allowbreak{}dfun}}{\vnbstmt{\begin{array}{@{}l@{}} \forall m, f, g, a, x, n. \\ \quad \operatorname{is\text{-}normed\text{-}vector\text{-}space}(m) \wedge \\ \quad \operatorname{is\text{-}bounded\text{-}linear\text{-}functional}(m, g) \wedge \\ \quad (f \in (\mathbb{R} \to \operatorname{vec}(m))) \wedge \\ \quad (n \in \mathbb{N}) \wedge \\ \quad \operatorname{taylor\text{-}differentiable\text{-}v}(m, f, a, x, n) \wedge \\ \quad \forall j \in \mathbb{N}. \\ \quad \quad (j \leq n) \Rightarrow  \\ \quad \quad \quad (\operatorname{nth\text{-}deriv\text{-}v}(m, f, j) \in (\mathbb{R} \to \operatorname{vec}(m))) \Rightarrow \\ \quad \quad \operatorname{taylor\text{-}differentiable}(\operatorname{compose}(g, f), a, x, n) \end{array}}}
\vnbentry{\texttt{r8e-\allowbreak{}td-\allowbreak{}v-\allowbreak{}diff-\allowbreak{}at}}{\vnbstmt{\begin{array}{@{}l@{}} \forall m, f, a, x, n, t. \\ \quad \operatorname{taylor\text{-}differentiable\text{-}v}(m, f, a, x, n) \wedge \\ \quad (n \in \mathbb{N}) \wedge \\ \quad (a < t) \wedge \\ \quad (t < x) \Rightarrow \\ \quad \quad \operatorname{is\text{-}diff\text{-}at\text{-}v}(m, \operatorname{nth\text{-}deriv\text{-}v}(m, f, n), t, \operatorname{nth\text{-}deriv\text{-}v}(m, f, (n + 1))(t)) \end{array}}}
\vnbentry{\texttt{r8e-\allowbreak{}tp-\allowbreak{}unfold}}{\vnbstmt{\begin{array}{@{}l@{}} \forall f, a, d, x. \\ \quad (\operatorname{taylor\text{-}poly}(f, a, d, x) \simeq \operatorname{series\text{-}partial\text{-}sum}((\lambda k \in \mathbb{N}.\; ((\operatorname{nth\text{-}deriv}(f, k)(a) \cdot \operatorname{power}((x - a), k)) \cdot \operatorname{recip}(\operatorname{factorial}(k)))), (d + 1))) \end{array}}}
\vnbentry{\texttt{r8e-\allowbreak{}coef-\allowbreak{}comm}}{\vnbstmt{\begin{array}{@{}l@{}} \forall p, q, v \in \mathbb{R}. \\ \quad (((p \cdot q) \cdot v) = ((v \cdot p) \cdot q)) \end{array}}}
\vnbentry{\texttt{r8e-\allowbreak{}g-\allowbreak{}of-\allowbreak{}poly-\allowbreak{}ind}}{\vnbstmt{\begin{array}{@{}l@{}} \forall n \in \mathbb{N}, m, f, g, a, x. \\ \quad \operatorname{is\text{-}normed\text{-}vector\text{-}space}(m) \wedge \\ \quad \operatorname{is\text{-}bounded\text{-}linear\text{-}functional}(m, g) \wedge \\ \quad (f \in (\mathbb{R} \to \operatorname{vec}(m))) \wedge \\ \quad (a \in \mathbb{R}) \wedge \\ \quad (x \in \mathbb{R}) \wedge \\ \quad \forall j \in \mathbb{N}. \\ \quad \quad (j \leq n) \Rightarrow  \\ \quad \quad \quad (\operatorname{nth\text{-}deriv\text{-}v}(m, f, j) \in (\mathbb{R} \to \operatorname{vec}(m))) \Rightarrow \\ \quad \quad g(\operatorname{taylor\text{-}poly\text{-}v}(m, f, a, x, n)) \\ \quad \quad =\; \operatorname{taylor\text{-}poly}(\operatorname{compose}(g, f), a, n, x) \end{array}}}
\vnbentry{\texttt{g-\allowbreak{}of-\allowbreak{}remainder-\allowbreak{}dfun}}{\vnbstmt{\begin{array}{@{}l@{}} \forall m, f, g, a, x, n. \\ \quad \operatorname{is\text{-}normed\text{-}vector\text{-}space}(m) \wedge \\ \quad \operatorname{is\text{-}bounded\text{-}linear\text{-}functional}(m, g) \wedge \\ \quad (f \in (\mathbb{R} \to \operatorname{vec}(m))) \wedge \\ \quad (a \in \mathbb{R}) \wedge \\ \quad (x \in \mathbb{R}) \wedge \\ \quad (n \in \mathbb{N}) \wedge \\ \quad \forall j \in \mathbb{N}. \\ \quad \quad (j \leq n) \Rightarrow  \\ \quad \quad \quad (\operatorname{nth\text{-}deriv\text{-}v}(m, f, j) \in (\mathbb{R} \to \operatorname{vec}(m))) \Rightarrow \\ \quad \quad g(\operatorname{vadd}(m)(f(x), \operatorname{vneg}(m)(\operatorname{taylor\text{-}poly\text{-}v}(m, f, a, x, n)))) \\ \quad \quad =\; (\operatorname{compose}(g, f)(x) - \operatorname{taylor\text{-}poly}(\operatorname{compose}(g, f), a, n, x)) \end{array}}}
\vnbentry{\texttt{vector-\allowbreak{}taylor-\allowbreak{}remainder-\allowbreak{}bound}}{\vnbstmt{\begin{array}{@{}l@{}} \forall m, f, a, x, n. \\ \quad \operatorname{is\text{-}normed\text{-}vector\text{-}space}(m) \wedge \\ \quad \operatorname{is\text{-}finite\text{-}dimensional}(\operatorname{normed\text{-}vector\text{-}space\text{-}as\text{-}module}(m)) \wedge \\ \quad (f \in (\mathbb{R} \to \operatorname{vec}(m))) \wedge \\ \quad (a \in \mathbb{R}) \wedge \\ \quad (x \in \mathbb{R}) \wedge \\ \quad (n \in \mathbb{N}) \wedge \\ \quad (a < x) \wedge \\ \quad \operatorname{taylor\text{-}differentiable\text{-}v}(m, f, a, x, n) \Rightarrow \\ \quad \quad \exists b \in \mathbb{R}. \\ \quad \quad \quad (a < b) \wedge \\ \quad \quad \quad (b < x) \wedge \\ \quad \quad \quad (\operatorname{factorial}((n + 1)) \cdot \operatorname{vnrm}(m)(\operatorname{vadd}(m)(f(x), \operatorname{vneg}(m)(\operatorname{taylor\text{-}poly\text{-}v}(m, f, a, x, n))))) \\ \quad \quad \quad \leq\; (\operatorname{vnrm}(m)(\operatorname{nth\text{-}deriv\text{-}v}(m, f, (n + 1))(b)) \cdot \operatorname{power}((x - a), (n + 1))) \end{array}}}

\vnbfile{theorem-library/rake-nf-norm.scm}{theorem-library/rake-nf-norm.scm}
\noindent{\small rake-nf-norm.scm -- rake batch 5b-D: the norm of a negative, and the metric space underlying a normed field.}\par
\vnbentry{\texttt{nf-\allowbreak{}norm-\allowbreak{}neg}}{\vnbstmt{\begin{array}{@{}l@{}} \forall f.; \forall a \in \operatorname{carr}(f). \\ \quad \operatorname{is\text{-}normed\text{-}field}(f) \Rightarrow  \\ \quad \quad (\operatorname{fnrm}(f)(\operatorname{neg}(f)(a)) = \operatorname{fnrm}(f)(a)) \end{array}}}
\vnbentry{\texttt{nf-\allowbreak{}metric-\allowbreak{}space-\allowbreak{}is-\allowbreak{}metric-\allowbreak{}space}}{\vnbstmt{\begin{array}{@{}l@{}} \forall f. \\ \quad \operatorname{is\text{-}normed\text{-}field}(f) \Rightarrow  \\ \quad \quad \operatorname{is\text{-}metric\text{-}space}(\operatorname{nf\text{-}metric\text{-}space}(f)) \end{array}}}

\vnbfile{theorem-library/diff-on-laws.scm}{theorem-library/diff-on-laws.scm}
\noindent{\small diff-on-laws.scm -- the laws of IS-DIFF-ON / DERIV-ON / HOLOMORPHIC-ON (structure-library/diff-on.scm).  Everything here is PROVEN; the definition file asserts nothing.}\par
\vnbentry{\texttt{diff-\allowbreak{}on-\allowbreak{}normed-\allowbreak{}field}}{\vnbstmt{\begin{array}{@{}l@{}} \forall k, u, f, a, l. \\ \quad \operatorname{is\text{-}diff\text{-}on}(k, u, f, a, l) \Rightarrow  \\ \quad \quad \operatorname{is\text{-}normed\text{-}field}(k) \end{array}}}
\vnbentry{\texttt{diff-\allowbreak{}on-\allowbreak{}open}}{\vnbstmt{\begin{array}{@{}l@{}} \forall k, u, f, a, l. \\ \quad \operatorname{is\text{-}diff\text{-}on}(k, u, f, a, l) \Rightarrow  \\ \quad \quad \operatorname{is\text{-}open}(\operatorname{nf\text{-}metric\text{-}space}(k), u) \end{array}}}
\vnbentry{\texttt{diff-\allowbreak{}on-\allowbreak{}in-\allowbreak{}fun}}{\vnbstmt{\begin{array}{@{}l@{}} \forall k, u, f, a, l. \\ \quad \operatorname{is\text{-}diff\text{-}on}(k, u, f, a, l) \Rightarrow  \\ \quad \quad (f \in (u \to \operatorname{carr}(k))) \end{array}}}
\vnbentry{\texttt{diff-\allowbreak{}on-\allowbreak{}pt-\allowbreak{}in}}{\vnbstmt{\begin{array}{@{}l@{}} \forall k, u, f, a, l. \\ \quad (\operatorname{is\text{-}diff\text{-}on}(k, u, f, a, l) \Rightarrow (a \in u)) \end{array}}}
\vnbentry{\texttt{diff-\allowbreak{}on-\allowbreak{}deriv-\allowbreak{}in-\allowbreak{}carr}}{\vnbstmt{\begin{array}{@{}l@{}} \forall k, u, f, a, l. \\ \quad \operatorname{is\text{-}diff\text{-}on}(k, u, f, a, l) \Rightarrow  \\ \quad \quad (l \in \operatorname{carr}(k)) \end{array}}}
\vnbentry{\texttt{nf-\allowbreak{}open-\allowbreak{}subset-\allowbreak{}carr}}{\vnbstmt{\begin{array}{@{}l@{}} \forall k, u. \\ \quad \operatorname{is\text{-}open}(\operatorname{nf\text{-}metric\text{-}space}(k), u) \Rightarrow  \\ \quad \quad (u \subseteq \operatorname{carr}(k)) \end{array}}}
\vnbentry{\texttt{diff-\allowbreak{}on-\allowbreak{}value-\allowbreak{}in-\allowbreak{}carr}}{\vnbstmt{\begin{array}{@{}l@{}} \forall k, u, f, a, l.; \forall x \in u. \\ \quad \operatorname{is\text{-}diff\text{-}on}(k, u, f, a, l) \Rightarrow  \\ \quad \quad (f(x) \in \operatorname{carr}(k)) \end{array}}}
\vnbentry{\texttt{rr-\allowbreak{}nf-\allowbreak{}carr}}{\vnbstmt{(\operatorname{carr}(\operatorname{rr\text{-}normed\text{-}field}) \simeq \mathbb{R})}}
\vnbentry{\texttt{cc-\allowbreak{}nf-\allowbreak{}carr}}{\vnbstmt{(\operatorname{carr}(\operatorname{cc\text{-}normed\text{-}field}) \simeq \mathbb{C})}}
\vnbentry{\texttt{rr-\allowbreak{}nf-\allowbreak{}add}}{\vnbstmt{(\operatorname{add}(\operatorname{rr\text{-}normed\text{-}field}) \simeq (\lambda \langle \operatorname{x\_}, \operatorname{y\_} \rangle \in (\mathbb{R} \times \mathbb{R}).\; (\operatorname{x\_} + \operatorname{y\_})))}}
\vnbentry{\texttt{rr-\allowbreak{}nf-\allowbreak{}mul}}{\vnbstmt{(\operatorname{mul}(\operatorname{rr\text{-}normed\text{-}field}) \simeq (\lambda \langle \operatorname{x\_}, \operatorname{y\_} \rangle \in (\mathbb{R} \times \mathbb{R}).\; (\operatorname{x\_} \cdot \operatorname{y\_})))}}
\vnbentry{\texttt{rr-\allowbreak{}nf-\allowbreak{}neg}}{\vnbstmt{(\operatorname{neg}(\operatorname{rr\text{-}normed\text{-}field}) \simeq (\lambda x \in \mathbb{R}.\; -x))}}
\vnbentry{\texttt{rr-\allowbreak{}nf-\allowbreak{}fnrm}}{\vnbstmt{(\operatorname{fnrm}(\operatorname{rr\text{-}normed\text{-}field}) \simeq (\lambda x \in \mathbb{R}.\; \lvert x \rvert))}}
\vnbentry{\texttt{cc-\allowbreak{}nf-\allowbreak{}fnrm}}{\vnbstmt{(\operatorname{fnrm}(\operatorname{cc\text{-}normed\text{-}field}) \simeq (\lambda x \in \mathbb{C}.\; \operatorname{magnitude}(x)))}}
\vnbentry{\texttt{cc-\allowbreak{}nf-\allowbreak{}zero}}{\vnbstmt{(\operatorname{zero}(\operatorname{cc\text{-}normed\text{-}field}) \simeq 0)}}
\vnbentry{\texttt{rr-\allowbreak{}nf-\allowbreak{}zero}}{\vnbstmt{(\operatorname{zero}(\operatorname{rr\text{-}normed\text{-}field}) \simeq 0)}}
\vnbentry{\texttt{rr-\allowbreak{}nf-\allowbreak{}add-\allowbreak{}apply}}{\vnbstmt{\begin{array}{@{}l@{}} \forall u, v. \\ \quad ((u \in \mathbb{R}) \wedge (v \in \mathbb{R})) \Rightarrow  \\ \quad \quad (\operatorname{add}(\operatorname{rr\text{-}normed\text{-}field})(u, v) \simeq (u + v)) \end{array}}}
\vnbentry{\texttt{rr-\allowbreak{}nf-\allowbreak{}mul-\allowbreak{}apply}}{\vnbstmt{\begin{array}{@{}l@{}} \forall u, v. \\ \quad ((u \in \mathbb{R}) \wedge (v \in \mathbb{R})) \Rightarrow  \\ \quad \quad (\operatorname{mul}(\operatorname{rr\text{-}normed\text{-}field})(u, v) \simeq (u \cdot v)) \end{array}}}
\vnbentry{\texttt{rr-\allowbreak{}nf-\allowbreak{}neg-\allowbreak{}apply}}{\vnbstmt{\begin{array}{@{}l@{}} \forall u \in \mathbb{R}. \\ \quad (\operatorname{neg}(\operatorname{rr\text{-}normed\text{-}field})(u) \simeq -u) \end{array}}}
\vnbentry{\texttt{rr-\allowbreak{}nf-\allowbreak{}fnrm-\allowbreak{}apply}}{\vnbstmt{\begin{array}{@{}l@{}} \forall u \in \mathbb{R}. \\ \quad (\operatorname{fnrm}(\operatorname{rr\text{-}normed\text{-}field})(u) \simeq \lvert u \rvert) \end{array}}}
\vnbentry{\texttt{cc-\allowbreak{}nf-\allowbreak{}fnrm-\allowbreak{}apply}}{\vnbstmt{\begin{array}{@{}l@{}} \forall u \in \mathbb{C}. \\ \quad (\operatorname{fnrm}(\operatorname{cc\text{-}normed\text{-}field})(u) \simeq \operatorname{magnitude}(u)) \end{array}}}
\vnbentry{\texttt{ms-\allowbreak{}agree-\allowbreak{}ball}}{\vnbstmt{\begin{array}{@{}l@{}} \forall a, b.; \forall c \in \operatorname{pts}(a), r. \\ \quad \operatorname{is\text{-}metric\text{-}space}(a) \wedge \\ \quad \operatorname{is\text{-}metric\text{-}space}(b) \wedge \\ \quad (\operatorname{pts}(b) \simeq \operatorname{pts}(a)) \wedge \\ \quad \forall u, v \in \operatorname{pts}(a). \\ \quad \quad (\operatorname{dist}(b)(u, v) \simeq \operatorname{dist}(a)(u, v)) \Rightarrow \\ \quad \quad (\operatorname{ball}(b, c, r) = \operatorname{ball}(a, c, r)) \end{array}}}
\vnbentry{\texttt{ms-\allowbreak{}agree-\allowbreak{}open}}{\vnbstmt{\begin{array}{@{}l@{}} \forall a, b, w. \\ \quad \operatorname{is\text{-}metric\text{-}space}(b) \wedge \\ \quad (\operatorname{pts}(b) \simeq \operatorname{pts}(a)) \wedge \\ \quad \forall u, v \in \operatorname{pts}(a). \\ \quad \quad (\operatorname{dist}(b)(u, v) \simeq \operatorname{dist}(a)(u, v)) \wedge \\ \quad \operatorname{is\text{-}open}(a, w) \Rightarrow \\ \quad \quad \operatorname{is\text{-}open}(b, w) \end{array}}}
\vnbentry{\texttt{ms-\allowbreak{}agree-\allowbreak{}continuous-\allowbreak{}at}}{\vnbstmt{\begin{array}{@{}l@{}} \forall a, b, \mathit{mta}, \mathit{mtb}, g, p. \\ \quad \operatorname{is\text{-}metric\text{-}space}(b) \wedge \\ \quad \operatorname{is\text{-}metric\text{-}space}(\mathit{mtb}) \wedge \\ \quad (\operatorname{pts}(b) \simeq \operatorname{pts}(a)) \wedge \\ \quad (\operatorname{pts}(\mathit{mtb}) \simeq \operatorname{pts}(\mathit{mta})) \wedge \\ \quad \forall u, v \in \operatorname{pts}(a). \\ \quad \quad (\operatorname{dist}(b)(u, v) \simeq \operatorname{dist}(a)(u, v)) \wedge \\ \quad \forall u, v \in \operatorname{pts}(\mathit{mta}). \\ \quad \quad (\operatorname{dist}(\mathit{mtb})(u, v) \simeq \operatorname{dist}(\mathit{mta})(u, v)) \wedge \\ \quad \operatorname{is\text{-}continuous\text{-}at}(a, \mathit{mta}, g, p) \Rightarrow \\ \quad \quad \operatorname{is\text{-}continuous\text{-}at}(b, \mathit{mtb}, g, p) \end{array}}}
\vnbentry{\texttt{ms-\allowbreak{}agree-\allowbreak{}pts-\allowbreak{}sym}}{\vnbstmt{\begin{array}{@{}l@{}} \forall a, b. \\ \quad (\operatorname{pts}(b) \simeq \operatorname{pts}(a)) \Rightarrow  \\ \quad \quad (\operatorname{pts}(a) \simeq \operatorname{pts}(b)) \end{array}}}
\vnbentry{\texttt{ms-\allowbreak{}agree-\allowbreak{}dist-\allowbreak{}sym}}{\vnbstmt{\begin{array}{@{}l@{}} \forall a, b.; \forall u, v \in \operatorname{pts}(b). \\ \quad (\operatorname{pts}(b) \simeq \operatorname{pts}(a)) \wedge \\ \quad \forall u, v \in \operatorname{pts}(a). \\ \quad \quad (\operatorname{dist}(b)(u, v) \simeq \operatorname{dist}(a)(u, v)) \Rightarrow \\ \quad \quad (\operatorname{dist}(a)(u, v) \simeq \operatorname{dist}(b)(u, v)) \end{array}}}
\vnbentry{\texttt{nf-\allowbreak{}rr-\allowbreak{}is-\allowbreak{}metric-\allowbreak{}space}}{\vnbstmt{\operatorname{is\text{-}metric\text{-}space}(\operatorname{nf\text{-}metric\text{-}space}(\operatorname{rr\text{-}normed\text{-}field}))}}
\vnbentry{\texttt{nf-\allowbreak{}rr-\allowbreak{}agree-\allowbreak{}pts}}{\vnbstmt{(\operatorname{pts}(\operatorname{nf\text{-}metric\text{-}space}(\operatorname{rr\text{-}normed\text{-}field})) \simeq \operatorname{pts}(\operatorname{rr\text{-}ms}))}}
\vnbentry{\texttt{nf-\allowbreak{}rr-\allowbreak{}agree-\allowbreak{}dist}}{\vnbstmt{\begin{array}{@{}l@{}} \forall u, v \in \operatorname{pts}(\operatorname{rr\text{-}ms}). \\ \quad (\operatorname{dist}(\operatorname{nf\text{-}metric\text{-}space}(\operatorname{rr\text{-}normed\text{-}field}))(u, v) \simeq \operatorname{dist}(\operatorname{rr\text{-}ms})(u, v)) \end{array}}}
\vnbentry{\texttt{rr-\allowbreak{}open-\allowbreak{}in-\allowbreak{}nf}}{\vnbstmt{\operatorname{is\text{-}open}(\operatorname{nf\text{-}metric\text{-}space}(\operatorname{rr\text{-}normed\text{-}field}), \mathbb{R})}}
\vnbentry{\texttt{rr-\allowbreak{}subset-\allowbreak{}nf-\allowbreak{}pts}}{\vnbstmt{\begin{array}{@{}l@{}} \mathbb{R} \\ \subseteq\; \operatorname{pts}(\operatorname{nf\text{-}metric\text{-}space}(\operatorname{rr\text{-}normed\text{-}field})) \end{array}}}
\vnbentry{\texttt{sub-\allowbreak{}rr-\allowbreak{}is-\allowbreak{}metric-\allowbreak{}space}}{\vnbstmt{\operatorname{is\text{-}metric\text{-}space}(\operatorname{subspace\text{-}ms}(\operatorname{nf\text{-}metric\text{-}space}(\operatorname{rr\text{-}normed\text{-}field}), \mathbb{R}))}}
\vnbentry{\texttt{sub-\allowbreak{}rr-\allowbreak{}agree-\allowbreak{}pts}}{\vnbstmt{(\operatorname{pts}(\operatorname{subspace\text{-}ms}(\operatorname{nf\text{-}metric\text{-}space}(\operatorname{rr\text{-}normed\text{-}field}), \mathbb{R})) \simeq \operatorname{pts}(\operatorname{rr\text{-}ms}))}}
\vnbentry{\texttt{sub-\allowbreak{}rr-\allowbreak{}agree-\allowbreak{}dist}}{\vnbstmt{\begin{array}{@{}l@{}} \forall u, v \in \operatorname{pts}(\operatorname{rr\text{-}ms}). \\ \quad (\operatorname{dist}(\operatorname{subspace\text{-}ms}(\operatorname{nf\text{-}metric\text{-}space}(\operatorname{rr\text{-}normed\text{-}field}), \mathbb{R}))(u, v) \simeq \operatorname{dist}(\operatorname{rr\text{-}ms})(u, v)) \end{array}}}
\vnbentry{\texttt{rr-\allowbreak{}nf-\allowbreak{}sub}}{\vnbstmt{\begin{array}{@{}l@{}} \forall u, v. \\ \quad ((u \in \mathbb{R}) \wedge (v \in \mathbb{R})) \Rightarrow  \\ \quad \quad (\operatorname{add}(\operatorname{rr\text{-}normed\text{-}field})(u, \operatorname{neg}(\operatorname{rr\text{-}normed\text{-}field})(v)) \simeq (u - v)) \end{array}}}
\vnbentry{\texttt{rr-\allowbreak{}nf-\allowbreak{}fun-\allowbreak{}class}}{\vnbstmt{((\mathbb{R} \to \operatorname{carr}(\operatorname{rr\text{-}normed\text{-}field})) \simeq (\mathbb{R} \to \mathbb{R}))}}
\vnbentry{\texttt{diff-\allowbreak{}at-\allowbreak{}iff-\allowbreak{}diff-\allowbreak{}on}}{\vnbstmt{\begin{array}{@{}l@{}} \forall f, a, l. \\ \quad (\operatorname{is\text{-}diff\text{-}at}(f, a, l) \Leftrightarrow \operatorname{is\text{-}diff\text{-}on}(\operatorname{rr\text{-}normed\text{-}field}, \mathbb{R}, f, a, l)) \end{array}}}
\vnbentry{\texttt{nf-\allowbreak{}norm-\allowbreak{}in-\allowbreak{}rr}}{\vnbstmt{\begin{array}{@{}l@{}} \forall k.; \forall u \in \operatorname{carr}(k). \\ \quad \operatorname{is\text{-}normed\text{-}field}(k) \Rightarrow  \\ \quad \quad (\operatorname{fnrm}(k)(u) \in \mathbb{R}) \end{array}}}
\vnbentry{\texttt{nf-\allowbreak{}norm-\allowbreak{}nonneg}}{\vnbstmt{\begin{array}{@{}l@{}} \forall k.; \forall u \in \operatorname{carr}(k). \\ \quad \operatorname{is\text{-}normed\text{-}field}(k) \Rightarrow  \\ \quad \quad (0 \leq \operatorname{fnrm}(k)(u)) \end{array}}}
\vnbentry{\texttt{nf-\allowbreak{}norm-\allowbreak{}mult}}{\vnbstmt{\begin{array}{@{}l@{}} \forall k, u, v. \\ \quad \operatorname{is\text{-}normed\text{-}field}(k) \wedge \\ \quad (u \in \operatorname{carr}(k)) \wedge \\ \quad (v \in \operatorname{carr}(k)) \Rightarrow \\ \quad \quad \operatorname{fnrm}(k)(\operatorname{mul}(k)(u, v)) \\ \quad \quad =\; (\operatorname{fnrm}(k)(u) \cdot \operatorname{fnrm}(k)(v)) \end{array}}}
\vnbentry{\texttt{nf-\allowbreak{}norm-\allowbreak{}subadd}}{\vnbstmt{\begin{array}{@{}l@{}} \forall k, u, v. \\ \quad \operatorname{is\text{-}normed\text{-}field}(k) \wedge \\ \quad (u \in \operatorname{carr}(k)) \wedge \\ \quad (v \in \operatorname{carr}(k)) \Rightarrow \\ \quad \quad \operatorname{fnrm}(k)(\operatorname{add}(k)(u, v)) \\ \quad \quad \leq\; (\operatorname{fnrm}(k)(u) + \operatorname{fnrm}(k)(v)) \end{array}}}
\vnbentry{\texttt{nf-\allowbreak{}add-\allowbreak{}in-\allowbreak{}carr}}{\vnbstmt{\begin{array}{@{}l@{}} \forall k, u, v. \\ \quad \operatorname{is\text{-}normed\text{-}field}(k) \wedge \\ \quad (u \in \operatorname{carr}(k)) \wedge \\ \quad (v \in \operatorname{carr}(k)) \Rightarrow \\ \quad \quad (\operatorname{add}(k)(u, v) \in \operatorname{carr}(k)) \end{array}}}
\vnbentry{\texttt{nf-\allowbreak{}mul-\allowbreak{}in-\allowbreak{}carr}}{\vnbstmt{\begin{array}{@{}l@{}} \forall k, u, v. \\ \quad \operatorname{is\text{-}normed\text{-}field}(k) \wedge \\ \quad (u \in \operatorname{carr}(k)) \wedge \\ \quad (v \in \operatorname{carr}(k)) \Rightarrow \\ \quad \quad (\operatorname{mul}(k)(u, v) \in \operatorname{carr}(k)) \end{array}}}
\vnbentry{\texttt{nf-\allowbreak{}neg-\allowbreak{}in-\allowbreak{}carr}}{\vnbstmt{\begin{array}{@{}l@{}} \forall k.; \forall u \in \operatorname{carr}(k). \\ \quad \operatorname{is\text{-}normed\text{-}field}(k) \Rightarrow  \\ \quad \quad (\operatorname{neg}(k)(u) \in \operatorname{carr}(k)) \end{array}}}
\vnbentry{\texttt{nf-\allowbreak{}sub-\allowbreak{}add-\allowbreak{}back}}{\vnbstmt{\begin{array}{@{}l@{}} \forall k, u, v. \\ \quad \operatorname{is\text{-}normed\text{-}field}(k) \wedge \\ \quad (u \in \operatorname{carr}(k)) \wedge \\ \quad (v \in \operatorname{carr}(k)) \Rightarrow \\ \quad \quad (\operatorname{add}(k)(\operatorname{add}(k)(u, \operatorname{neg}(k)(v)), v) = u) \end{array}}}
\vnbentry{\texttt{nf-\allowbreak{}norm-\allowbreak{}le-\allowbreak{}add}}{\vnbstmt{\begin{array}{@{}l@{}} \forall k, u, v. \\ \quad \operatorname{is\text{-}normed\text{-}field}(k) \wedge \\ \quad (u \in \operatorname{carr}(k)) \wedge \\ \quad (v \in \operatorname{carr}(k)) \Rightarrow \\ \quad \quad \operatorname{fnrm}(k)(u) \\ \quad \quad \leq\; (\operatorname{fnrm}(k)(\operatorname{add}(k)(u, \operatorname{neg}(k)(v))) + \operatorname{fnrm}(k)(v)) \end{array}}}
\vnbentry{\texttt{diff-\allowbreak{}on-\allowbreak{}implies-\allowbreak{}continuous}}{\vnbstmt{\begin{array}{@{}l@{}} \forall k, u, f, a, l. \\ \quad \operatorname{is\text{-}diff\text{-}on}(k, u, f, a, l) \Rightarrow  \\ \quad \quad \operatorname{is\text{-}continuous\text{-}at}(\operatorname{subspace\text{-}ms}(\operatorname{nf\text{-}metric\text{-}space}(k), u), \operatorname{nf\text{-}metric\text{-}space}(k), f, a) \end{array}}}
\vnbentry{\texttt{deriv-\allowbreak{}on-\allowbreak{}unfold}}{\vnbstmt{\begin{array}{@{}l@{}} \forall k, u, f, a. \\ \quad (\operatorname{deriv\text{-}on}(k, u, f, a) \simeq \iota(l)) \end{array}}}
\vnbentry{\texttt{holomorphic-\allowbreak{}on-\allowbreak{}open}}{\vnbstmt{\begin{array}{@{}l@{}} \forall u, f. \\ \quad \operatorname{holomorphic\text{-}on}(u, f) \Rightarrow  \\ \quad \quad \operatorname{is\text{-}open}(\operatorname{nf\text{-}metric\text{-}space}(\operatorname{cc\text{-}normed\text{-}field}), u) \end{array}}}
\vnbentry{\texttt{holomorphic-\allowbreak{}on-\allowbreak{}in-\allowbreak{}fun}}{\vnbstmt{\begin{array}{@{}l@{}} \forall u, f. \\ \quad \operatorname{holomorphic\text{-}on}(u, f) \Rightarrow  \\ \quad \quad (f \in (u \to \mathbb{C})) \end{array}}}
\vnbentry{\texttt{holomorphic-\allowbreak{}on-\allowbreak{}diff}}{\vnbstmt{\begin{array}{@{}l@{}} \forall u, f.; \forall a \in u. \\ \quad \operatorname{holomorphic\text{-}on}(u, f) \Rightarrow  \\ \quad \quad \exists l. \\ \quad \quad \quad \operatorname{is\text{-}diff\text{-}on}(\operatorname{cc\text{-}normed\text{-}field}, u, f, a, l) \end{array}}}

\vnbfile{theorem-library/ms-continuity-algebra.scm}{theorem-library/ms-continuity-algebra.scm}
\noindent{\small ms-continuity-algebra.scm -- THE CONTINUITY ALGEBRA OFF THE REAL LINE.}\par
\vnbentry{\texttt{ms-\allowbreak{}pts-\allowbreak{}is-\allowbreak{}set}}{\vnbstmt{\begin{array}{@{}l@{}} \forall s. \\ \quad \operatorname{is\text{-}metric\text{-}space}(s) \Rightarrow  \\ \quad \quad (\operatorname{pts}(s) \in \mathbf{Set}) \end{array}}}
\vnbentry{\texttt{ms-\allowbreak{}const-\allowbreak{}lam-\allowbreak{}in-\allowbreak{}fun}}{\vnbstmt{\begin{array}{@{}l@{}} \forall s, t.; \forall c \in \operatorname{pts}(t). \\ \quad \operatorname{is\text{-}metric\text{-}space}(s) \Rightarrow  \\ \quad \quad ((\lambda z \in \operatorname{pts}(s).\; c) \in (\operatorname{pts}(s) \to \operatorname{pts}(t))) \end{array}}}
\vnbentry{\texttt{ms-\allowbreak{}const-\allowbreak{}continuous-\allowbreak{}at}}{\vnbstmt{\begin{array}{@{}l@{}} \forall s, t.; \forall c \in \operatorname{pts}(t), a \in \operatorname{pts}(s). \\ \quad \operatorname{is\text{-}metric\text{-}space}(s) \wedge \\ \quad \operatorname{is\text{-}metric\text{-}space}(t) \Rightarrow \\ \quad \quad \operatorname{is\text{-}continuous\text{-}at}(s, t, (\lambda z \in \operatorname{pts}(s).\; c), a) \end{array}}}
\vnbentry{\texttt{ms-\allowbreak{}ident-\allowbreak{}lam-\allowbreak{}in-\allowbreak{}fun}}{\vnbstmt{\begin{array}{@{}l@{}} \forall s. \\ \quad \operatorname{is\text{-}metric\text{-}space}(s) \Rightarrow  \\ \quad \quad ((\lambda z \in \operatorname{pts}(s).\; z) \in (\operatorname{pts}(s) \to \operatorname{pts}(s))) \end{array}}}
\vnbentry{\texttt{ms-\allowbreak{}identity-\allowbreak{}continuous-\allowbreak{}at}}{\vnbstmt{\begin{array}{@{}l@{}} \forall s.; \forall a \in \operatorname{pts}(s). \\ \quad \operatorname{is\text{-}metric\text{-}space}(s) \Rightarrow  \\ \quad \quad \operatorname{is\text{-}continuous\text{-}at}(s, s, (\lambda z \in \operatorname{pts}(s).\; z), a) \end{array}}}
\vnbentry{\texttt{ms-\allowbreak{}cont-\allowbreak{}transfer-\allowbreak{}ptwise-\allowbreak{}eq}}{\vnbstmt{\begin{array}{@{}l@{}} \forall s, t, g, f, a. \\ \quad \operatorname{is\text{-}continuous\text{-}at}(s, t, g, a) \wedge \\ \quad (f \in (\operatorname{pts}(s) \to \operatorname{pts}(t))) \wedge \\ \quad \forall z \in \operatorname{pts}(s).\; (f(z) = g(z)) \Rightarrow \\ \quad \quad \operatorname{is\text{-}continuous\text{-}at}(s, t, f, a) \end{array}}}
\vnbentry{\texttt{ms-\allowbreak{}compose-\allowbreak{}continuous-\allowbreak{}at}}{\vnbstmt{\begin{array}{@{}l@{}} \forall s, t, u, f, g, a. \\ \quad \operatorname{is\text{-}continuous\text{-}at}(s, t, f, a) \wedge \\ \quad \operatorname{is\text{-}continuous\text{-}at}(t, u, g, f(a)) \Rightarrow \\ \quad \quad \operatorname{is\text{-}continuous\text{-}at}(s, u, (\lambda z \in \operatorname{pts}(s).\; g(f(z))), a) \end{array}}}
\vnbentry{\texttt{ms-\allowbreak{}cont-\allowbreak{}agree-\allowbreak{}off-\allowbreak{}pt}}{\vnbstmt{\begin{array}{@{}l@{}} \forall s, t, f, g, p. \\ \quad \operatorname{is\text{-}continuous\text{-}at}(s, t, f, p) \wedge \\ \quad \operatorname{is\text{-}continuous\text{-}at}(s, t, g, p) \wedge \\ \quad \forall r. \\ \quad \quad \operatorname{pos\text{-}rr}(r) \Rightarrow  \\ \quad \quad \quad \exists b \in \operatorname{pts}(s). \\ \quad \quad \quad \quad (\neg (b = p) \wedge (\operatorname{dist}(s)(p, b) \leq r)) \wedge \\ \quad \forall y \in \operatorname{pts}(s).\; (\neg (y = p) \Rightarrow (f(y) = g(y))) \Rightarrow \\ \quad \quad (f(p) = g(p)) \end{array}}}
\vnbentry{\texttt{nf-\allowbreak{}sub-\allowbreak{}split}}{\vnbstmt{\begin{array}{@{}l@{}} \forall k, u, v, p, q. \\ \quad \operatorname{is\text{-}normed\text{-}field}(k) \wedge \\ \quad (u \in \operatorname{carr}(k)) \wedge \\ \quad (v \in \operatorname{carr}(k)) \wedge \\ \quad (p \in \operatorname{carr}(k)) \wedge \\ \quad (q \in \operatorname{carr}(k)) \Rightarrow \\ \quad \quad \operatorname{add}(k)(\operatorname{add}(k)(u, v), \operatorname{neg}(k)(\operatorname{add}(k)(p, q))) \\ \quad \quad =\; \operatorname{add}(k)(\operatorname{add}(k)(u, \operatorname{neg}(k)(p)), \operatorname{add}(k)(v, \operatorname{neg}(k)(q))) \end{array}}}
\vnbentry{\texttt{nf-\allowbreak{}mul-\allowbreak{}sub-\allowbreak{}split}}{\vnbstmt{\begin{array}{@{}l@{}} \forall k, u, v, p, q. \\ \quad \operatorname{is\text{-}normed\text{-}field}(k) \wedge \\ \quad (u \in \operatorname{carr}(k)) \wedge \\ \quad (v \in \operatorname{carr}(k)) \wedge \\ \quad (p \in \operatorname{carr}(k)) \wedge \\ \quad (q \in \operatorname{carr}(k)) \Rightarrow \\ \quad \quad \operatorname{add}(k)(\operatorname{mul}(k)(u, v), \operatorname{neg}(k)(\operatorname{mul}(k)(p, q))) \\ \quad \quad =\; \operatorname{add}(k)(\operatorname{mul}(k)(u, \operatorname{add}(k)(v, \operatorname{neg}(k)(q))), \operatorname{mul}(k)(q, \operatorname{add}(k)(u, \operatorname{neg}(k)(p)))) \end{array}}}
\vnbentry{\texttt{nf-\allowbreak{}neg-\allowbreak{}sub-\allowbreak{}split}}{\vnbstmt{\begin{array}{@{}l@{}} \forall k, u, v. \\ \quad \operatorname{is\text{-}normed\text{-}field}(k) \wedge \\ \quad (u \in \operatorname{carr}(k)) \wedge \\ \quad (v \in \operatorname{carr}(k)) \Rightarrow \\ \quad \quad \operatorname{add}(k)(\operatorname{neg}(k)(u), \operatorname{neg}(k)(\operatorname{neg}(k)(v))) \\ \quad \quad =\; \operatorname{neg}(k)(\operatorname{add}(k)(u, \operatorname{neg}(k)(v))) \end{array}}}
\vnbentry{\texttt{nf-\allowbreak{}sub-\allowbreak{}anti}}{\vnbstmt{\begin{array}{@{}l@{}} \forall k, u, v. \\ \quad \operatorname{is\text{-}normed\text{-}field}(k) \wedge \\ \quad (u \in \operatorname{carr}(k)) \wedge \\ \quad (v \in \operatorname{carr}(k)) \Rightarrow \\ \quad \quad \operatorname{neg}(k)(\operatorname{add}(k)(u, \operatorname{neg}(k)(v))) \\ \quad \quad =\; \operatorname{add}(k)(v, \operatorname{neg}(k)(u)) \end{array}}}
\vnbentry{\texttt{nf-\allowbreak{}norm-\allowbreak{}sub-\allowbreak{}sym}}{\vnbstmt{\begin{array}{@{}l@{}} \forall k, u, v. \\ \quad \operatorname{is\text{-}normed\text{-}field}(k) \wedge \\ \quad (u \in \operatorname{carr}(k)) \wedge \\ \quad (v \in \operatorname{carr}(k)) \Rightarrow \\ \quad \quad \operatorname{fnrm}(k)(\operatorname{add}(k)(u, \operatorname{neg}(k)(v))) \\ \quad \quad =\; \operatorname{fnrm}(k)(\operatorname{add}(k)(v, \operatorname{neg}(k)(u))) \end{array}}}
\vnbentry{\texttt{nf-\allowbreak{}norm-\allowbreak{}neg-\allowbreak{}diff}}{\vnbstmt{\begin{array}{@{}l@{}} \forall k, u, v. \\ \quad \operatorname{is\text{-}normed\text{-}field}(k) \wedge \\ \quad (u \in \operatorname{carr}(k)) \wedge \\ \quad (v \in \operatorname{carr}(k)) \Rightarrow \\ \quad \quad \operatorname{fnrm}(k)(\operatorname{add}(k)(\operatorname{neg}(k)(u), \operatorname{neg}(k)(\operatorname{neg}(k)(v)))) \\ \quad \quad =\; \operatorname{fnrm}(k)(\operatorname{add}(k)(u, \operatorname{neg}(k)(v))) \end{array}}}
\vnbentry{\texttt{nf-\allowbreak{}norm-\allowbreak{}sum-\allowbreak{}diff}}{\vnbstmt{\begin{array}{@{}l@{}} \forall k, u, v, p, q. \\ \quad \operatorname{is\text{-}normed\text{-}field}(k) \wedge \\ \quad (u \in \operatorname{carr}(k)) \wedge \\ \quad (v \in \operatorname{carr}(k)) \wedge \\ \quad (p \in \operatorname{carr}(k)) \wedge \\ \quad (q \in \operatorname{carr}(k)) \Rightarrow \\ \quad \quad \operatorname{fnrm}(k)(\operatorname{add}(k)(\operatorname{add}(k)(u, v), \operatorname{neg}(k)(\operatorname{add}(k)(p, q)))) \\ \quad \quad \leq\; (\operatorname{fnrm}(k)(\operatorname{add}(k)(u, \operatorname{neg}(k)(p))) + \operatorname{fnrm}(k)(\operatorname{add}(k)(v, \operatorname{neg}(k)(q)))) \end{array}}}
\vnbentry{\texttt{nf-\allowbreak{}norm-\allowbreak{}prod-\allowbreak{}bound}}{\vnbstmt{\begin{array}{@{}l@{}} \forall k, u, v, p, q, m, \mathit{nfk}, t, e. \\ \quad \operatorname{is\text{-}normed\text{-}field}(k) \wedge \\ \quad (u \in \operatorname{carr}(k)) \wedge \\ \quad (v \in \operatorname{carr}(k)) \wedge \\ \quad (p \in \operatorname{carr}(k)) \wedge \\ \quad (q \in \operatorname{carr}(k)) \wedge \\ \quad (m \in \mathbb{R}) \wedge \\ \quad (\mathit{nfk} \in \mathbb{R}) \wedge \\ \quad (t \in \mathbb{R}) \wedge \\ \quad (e \in \mathbb{R}) \wedge \\ \quad (\operatorname{fnrm}(k)(u) \leq m) \wedge \\ \quad (\operatorname{fnrm}(k)(q) \leq \mathit{nfk}) \wedge \\ \quad (\operatorname{fnrm}(k)(\operatorname{add}(k)(u, \operatorname{neg}(k)(p))) \leq t) \wedge \\ \quad (\operatorname{fnrm}(k)(\operatorname{add}(k)(v, \operatorname{neg}(k)(q))) \leq t) \wedge \\ \quad (((m + \mathit{nfk}) \cdot t) \leq e) \Rightarrow \\ \quad \quad \operatorname{fnrm}(k)(\operatorname{add}(k)(\operatorname{mul}(k)(u, v), \operatorname{neg}(k)(\operatorname{mul}(k)(p, q)))) \\ \quad \quad \leq\; e \end{array}}}
\vnbentry{\texttt{nf-\allowbreak{}sum-\allowbreak{}lam-\allowbreak{}in-\allowbreak{}fun}}{\vnbstmt{\begin{array}{@{}l@{}} \forall k, s.; \forall g, h \in (\operatorname{pts}(s) \to \operatorname{pts}(\operatorname{nf\text{-}metric\text{-}space}(k))). \\ \quad \operatorname{is\text{-}normed\text{-}field}(k) \wedge \\ \quad \operatorname{is\text{-}metric\text{-}space}(s) \Rightarrow \\ \quad \quad (\lambda z \in \operatorname{pts}(s).\; \operatorname{add}(k)(g(z), h(z))) \\ \quad \quad \in\; (\operatorname{pts}(s) \to \operatorname{pts}(\operatorname{nf\text{-}metric\text{-}space}(k))) \end{array}}}
\vnbentry{\texttt{nf-\allowbreak{}prod-\allowbreak{}lam-\allowbreak{}in-\allowbreak{}fun}}{\vnbstmt{\begin{array}{@{}l@{}} \forall k, s.; \forall g, h \in (\operatorname{pts}(s) \to \operatorname{pts}(\operatorname{nf\text{-}metric\text{-}space}(k))). \\ \quad \operatorname{is\text{-}normed\text{-}field}(k) \wedge \\ \quad \operatorname{is\text{-}metric\text{-}space}(s) \Rightarrow \\ \quad \quad (\lambda z \in \operatorname{pts}(s).\; \operatorname{mul}(k)(g(z), h(z))) \\ \quad \quad \in\; (\operatorname{pts}(s) \to \operatorname{pts}(\operatorname{nf\text{-}metric\text{-}space}(k))) \end{array}}}
\vnbentry{\texttt{nf-\allowbreak{}neg-\allowbreak{}lam-\allowbreak{}in-\allowbreak{}fun}}{\vnbstmt{\begin{array}{@{}l@{}} \forall k, s.; \forall g \in (\operatorname{pts}(s) \to \operatorname{pts}(\operatorname{nf\text{-}metric\text{-}space}(k))). \\ \quad \operatorname{is\text{-}normed\text{-}field}(k) \wedge \\ \quad \operatorname{is\text{-}metric\text{-}space}(s) \Rightarrow \\ \quad \quad (\lambda z \in \operatorname{pts}(s).\; \operatorname{neg}(k)(g(z))) \\ \quad \quad \in\; (\operatorname{pts}(s) \to \operatorname{pts}(\operatorname{nf\text{-}metric\text{-}space}(k))) \end{array}}}
\vnbentry{\texttt{nf-\allowbreak{}sum-\allowbreak{}continuous-\allowbreak{}at}}{\vnbstmt{\begin{array}{@{}l@{}} \forall k, s, g, h, a. \\ \quad \operatorname{is\text{-}normed\text{-}field}(k) \wedge \\ \quad \operatorname{is\text{-}continuous\text{-}at}(s, \operatorname{nf\text{-}metric\text{-}space}(k), g, a) \wedge \\ \quad \operatorname{is\text{-}continuous\text{-}at}(s, \operatorname{nf\text{-}metric\text{-}space}(k), h, a) \Rightarrow \\ \quad \quad \operatorname{is\text{-}continuous\text{-}at}(s, \operatorname{nf\text{-}metric\text{-}space}(k), (\lambda z \in \operatorname{pts}(s).\; \operatorname{add}(k)(g(z), h(z))), \\ \quad \quad \quad a) \end{array}}}
\vnbentry{\texttt{nf-\allowbreak{}neg-\allowbreak{}continuous-\allowbreak{}at}}{\vnbstmt{\begin{array}{@{}l@{}} \forall k, s, g, a. \\ \quad \operatorname{is\text{-}normed\text{-}field}(k) \wedge \\ \quad \operatorname{is\text{-}continuous\text{-}at}(s, \operatorname{nf\text{-}metric\text{-}space}(k), g, a) \Rightarrow \\ \quad \quad \operatorname{is\text{-}continuous\text{-}at}(s, \operatorname{nf\text{-}metric\text{-}space}(k), (\lambda z \in \operatorname{pts}(s).\; \operatorname{neg}(k)(g(z))), a) \end{array}}}
\vnbentry{\texttt{nf-\allowbreak{}product-\allowbreak{}continuous-\allowbreak{}at}}{\vnbstmt{\begin{array}{@{}l@{}} \forall k, s, g, h, a. \\ \quad \operatorname{is\text{-}normed\text{-}field}(k) \wedge \\ \quad \operatorname{is\text{-}continuous\text{-}at}(s, \operatorname{nf\text{-}metric\text{-}space}(k), g, a) \wedge \\ \quad \operatorname{is\text{-}continuous\text{-}at}(s, \operatorname{nf\text{-}metric\text{-}space}(k), h, a) \Rightarrow \\ \quad \quad \operatorname{is\text{-}continuous\text{-}at}(s, \operatorname{nf\text{-}metric\text{-}space}(k), (\lambda z \in \operatorname{pts}(s).\; \operatorname{mul}(k)(g(z), h(z))), \\ \quad \quad \quad a) \end{array}}}

\vnbfile{theorem-library/categories.scm}{theorem-library/categories.scm}
\noindent{\small categories.scm -- FOUR CATEGORIES BESIDE THE DEFAULTS: Lipschitz, uniformly continuous and continuous maps of metric spaces, bounded linear maps of normed vector spaces; their laws, their Hom functors, and the inclusions.}\par
\vnbentry{\texttt{hom-\allowbreak{}lipschitz-\allowbreak{}member-\allowbreak{}iff}}{\vnbstmt{\begin{array}{@{}l@{}} \forall a, b, f. \\ \quad ((f \in \operatorname{hom\text{-}lipschitz}(a, b)) \Leftrightarrow ((f \in (\operatorname{pts}(a) \to \operatorname{pts}(b))) \wedge \operatorname{is\text{-}lipschitz\text{-}arrow}(a, b, f))) \end{array}}}
\vnbentry{\texttt{hom-\allowbreak{}lipschitz-\allowbreak{}in-\allowbreak{}set}}{\vnbstmt{\begin{array}{@{}l@{}} \forall a, b. \\ \quad \operatorname{is\text{-}metric\text{-}space}(a) \wedge \\ \quad \operatorname{is\text{-}metric\text{-}space}(b) \Rightarrow \\ \quad \quad (\operatorname{hom\text{-}lipschitz}(a, b) \in \mathbf{Set}) \end{array}}}
\vnbentry{\texttt{hom-\allowbreak{}uniformly-\allowbreak{}continuous-\allowbreak{}member-\allowbreak{}iff}}{\vnbstmt{\begin{array}{@{}l@{}} \forall a, b, f. \\ \quad ((f \in \operatorname{hom\text{-}uniformly\text{-}continuous}(a, b)) \Leftrightarrow ((f \in (\operatorname{pts}(a) \to \operatorname{pts}(b))) \wedge \operatorname{is\text{-}uniformly\text{-}continuous\text{-}arrow}(a, b, f))) \end{array}}}
\vnbentry{\texttt{hom-\allowbreak{}uniformly-\allowbreak{}continuous-\allowbreak{}in-\allowbreak{}set}}{\vnbstmt{\begin{array}{@{}l@{}} \forall a, b. \\ \quad \operatorname{is\text{-}metric\text{-}space}(a) \wedge \\ \quad \operatorname{is\text{-}metric\text{-}space}(b) \Rightarrow \\ \quad \quad (\operatorname{hom\text{-}uniformly\text{-}continuous}(a, b) \in \mathbf{Set}) \end{array}}}
\vnbentry{\texttt{hom-\allowbreak{}continuous-\allowbreak{}member-\allowbreak{}iff}}{\vnbstmt{\begin{array}{@{}l@{}} \forall a, b, f. \\ \quad ((f \in \operatorname{hom\text{-}continuous}(a, b)) \Leftrightarrow ((f \in (\operatorname{pts}(a) \to \operatorname{pts}(b))) \wedge \operatorname{is\text{-}continuous\text{-}arrow}(a, b, f))) \end{array}}}
\vnbentry{\texttt{hom-\allowbreak{}continuous-\allowbreak{}in-\allowbreak{}set}}{\vnbstmt{\begin{array}{@{}l@{}} \forall a, b. \\ \quad \operatorname{is\text{-}metric\text{-}space}(a) \wedge \\ \quad \operatorname{is\text{-}metric\text{-}space}(b) \Rightarrow \\ \quad \quad (\operatorname{hom\text{-}continuous}(a, b) \in \mathbf{Set}) \end{array}}}
\vnbentry{\texttt{hom-\allowbreak{}bounded-\allowbreak{}linear-\allowbreak{}member-\allowbreak{}iff}}{\vnbstmt{\begin{array}{@{}l@{}} \forall a, b, f. \\ \quad ((f \in \operatorname{hom\text{-}bounded\text{-}linear}(a, b)) \Leftrightarrow ((f \in (\operatorname{vec}(a) \to \operatorname{vec}(b))) \wedge \operatorname{is\text{-}bounded\text{-}linear\text{-}arrow}(a, b, f))) \end{array}}}
\vnbentry{\texttt{hom-\allowbreak{}bounded-\allowbreak{}linear-\allowbreak{}in-\allowbreak{}set}}{\vnbstmt{\begin{array}{@{}l@{}} \forall a, b. \\ \quad \operatorname{is\text{-}normed\text{-}vector\text{-}space}(a) \wedge \\ \quad \operatorname{is\text{-}normed\text{-}vector\text{-}space}(b) \Rightarrow \\ \quad \quad (\operatorname{hom\text{-}bounded\text{-}linear}(a, b) \in \mathbf{Set}) \end{array}}}
\vnbentry{\texttt{hom-\allowbreak{}metric-\allowbreak{}space-\allowbreak{}is-\allowbreak{}lipschitz-\allowbreak{}arrow}}{\vnbstmt{\begin{array}{@{}l@{}} \forall a, b, f. \\ \quad \operatorname{is\text{-}hom\text{-}metric\text{-}space}(a, b, f) \Rightarrow  \\ \quad \quad \operatorname{is\text{-}lipschitz\text{-}arrow}(a, b, f) \end{array}}}
\vnbentry{\texttt{lipschitz-\allowbreak{}arrow-\allowbreak{}is-\allowbreak{}uniformly-\allowbreak{}continuous-\allowbreak{}arrow}}{\vnbstmt{\begin{array}{@{}l@{}} \forall a, b, f. \\ \quad \operatorname{is\text{-}lipschitz\text{-}arrow}(a, b, f) \Rightarrow  \\ \quad \quad \operatorname{is\text{-}uniformly\text{-}continuous\text{-}arrow}(a, b, f) \end{array}}}
\vnbentry{\texttt{uniformly-\allowbreak{}continuous-\allowbreak{}arrow-\allowbreak{}is-\allowbreak{}continuous-\allowbreak{}arrow}}{\vnbstmt{\begin{array}{@{}l@{}} \forall a, b, f. \\ \quad \operatorname{is\text{-}uniformly\text{-}continuous\text{-}arrow}(a, b, f) \Rightarrow  \\ \quad \quad \operatorname{is\text{-}continuous\text{-}arrow}(a, b, f) \end{array}}}
\vnbentry{\texttt{hom-\allowbreak{}normed-\allowbreak{}vector-\allowbreak{}space-\allowbreak{}is-\allowbreak{}bounded-\allowbreak{}linear-\allowbreak{}arrow}}{\vnbstmt{\begin{array}{@{}l@{}} \forall a, b, f. \\ \quad \operatorname{is\text{-}hom\text{-}normed\text{-}vector\text{-}space}(a, b, f) \Rightarrow  \\ \quad \quad \operatorname{is\text{-}bounded\text{-}linear\text{-}arrow}(a, b, f) \end{array}}}
\vnbentry{\texttt{hom-\allowbreak{}lipschitz-\allowbreak{}id}}{\vnbstmt{\begin{array}{@{}l@{}} \forall a. \\ \quad \operatorname{is\text{-}metric\text{-}space}(a) \Rightarrow  \\ \quad \quad \operatorname{is\text{-}lipschitz\text{-}arrow}(a, a, \operatorname{id\text{-}fun}(\operatorname{pts}(a))) \end{array}}}
\vnbentry{\texttt{hom-\allowbreak{}uniformly-\allowbreak{}continuous-\allowbreak{}id}}{\vnbstmt{\begin{array}{@{}l@{}} \forall a. \\ \quad \operatorname{is\text{-}metric\text{-}space}(a) \Rightarrow  \\ \quad \quad \operatorname{is\text{-}uniformly\text{-}continuous\text{-}arrow}(a, a, \operatorname{id\text{-}fun}(\operatorname{pts}(a))) \end{array}}}
\vnbentry{\texttt{hom-\allowbreak{}continuous-\allowbreak{}id}}{\vnbstmt{\begin{array}{@{}l@{}} \forall a. \\ \quad \operatorname{is\text{-}metric\text{-}space}(a) \Rightarrow  \\ \quad \quad \operatorname{is\text{-}continuous\text{-}arrow}(a, a, \operatorname{id\text{-}fun}(\operatorname{pts}(a))) \end{array}}}
\vnbentry{\texttt{hom-\allowbreak{}bounded-\allowbreak{}linear-\allowbreak{}id}}{\vnbstmt{\begin{array}{@{}l@{}} \forall a. \\ \quad \operatorname{is\text{-}normed\text{-}vector\text{-}space}(a) \Rightarrow  \\ \quad \quad \operatorname{is\text{-}bounded\text{-}linear\text{-}arrow}(a, a, \operatorname{id\text{-}fun}(\operatorname{vec}(a))) \end{array}}}
\vnbentry{\texttt{hom-\allowbreak{}lipschitz-\allowbreak{}compose}}{\vnbstmt{\begin{array}{@{}l@{}} \forall a, b, c, f, g. \\ \quad \operatorname{is\text{-}lipschitz\text{-}arrow}(a, b, f) \wedge \\ \quad \operatorname{is\text{-}lipschitz\text{-}arrow}(b, c, g) \Rightarrow \\ \quad \quad \operatorname{is\text{-}lipschitz\text{-}arrow}(a, c, \operatorname{compose}(g, f)) \end{array}}}
\vnbentry{\texttt{hom-\allowbreak{}uniformly-\allowbreak{}continuous-\allowbreak{}compose}}{\vnbstmt{\begin{array}{@{}l@{}} \forall a, b, c, f, g. \\ \quad \operatorname{is\text{-}uniformly\text{-}continuous\text{-}arrow}(a, b, f) \wedge \\ \quad \operatorname{is\text{-}uniformly\text{-}continuous\text{-}arrow}(b, c, g) \Rightarrow \\ \quad \quad \operatorname{is\text{-}uniformly\text{-}continuous\text{-}arrow}(a, c, \operatorname{compose}(g, f)) \end{array}}}
\vnbentry{\texttt{hom-\allowbreak{}continuous-\allowbreak{}compose}}{\vnbstmt{\begin{array}{@{}l@{}} \forall a, b, c, f, g. \\ \quad \operatorname{is\text{-}continuous\text{-}arrow}(a, b, f) \wedge \\ \quad \operatorname{is\text{-}continuous\text{-}arrow}(b, c, g) \Rightarrow \\ \quad \quad \operatorname{is\text{-}continuous\text{-}arrow}(a, c, \operatorname{compose}(g, f)) \end{array}}}
\vnbentry{\texttt{hom-\allowbreak{}bounded-\allowbreak{}linear-\allowbreak{}compose}}{\vnbstmt{\begin{array}{@{}l@{}} \forall a, b, c, f, g. \\ \quad \operatorname{is\text{-}bounded\text{-}linear\text{-}arrow}(a, b, f) \wedge \\ \quad \operatorname{is\text{-}bounded\text{-}linear\text{-}arrow}(b, c, g) \Rightarrow \\ \quad \quad \operatorname{is\text{-}bounded\text{-}linear\text{-}arrow}(a, c, \operatorname{compose}(g, f)) \end{array}}}
\vnbentry{\texttt{hom-\allowbreak{}lipschitz-\allowbreak{}post-\allowbreak{}type}}{\vnbstmt{\begin{array}{@{}l@{}} \forall a, b, c, h. \\ \quad \operatorname{is\text{-}metric\text{-}space}(a) \wedge \\ \quad (h \in \operatorname{hom\text{-}lipschitz}(b, c)) \Rightarrow \\ \quad \quad (\lambda f \in \operatorname{hom\text{-}lipschitz}(a, b).\; \operatorname{compose}(h, f)) \\ \quad \quad \in\; (\operatorname{hom\text{-}lipschitz}(a, b) \to \operatorname{hom\text{-}lipschitz}(a, c)) \end{array}}}
\vnbentry{\texttt{hom-\allowbreak{}lipschitz-\allowbreak{}pre-\allowbreak{}type}}{\vnbstmt{\begin{array}{@{}l@{}} \forall a, b, c, k. \\ \quad \operatorname{is\text{-}metric\text{-}space}(b) \wedge \\ \quad (k \in \operatorname{hom\text{-}lipschitz}(c, a)) \Rightarrow \\ \quad \quad (\lambda f \in \operatorname{hom\text{-}lipschitz}(a, b).\; \operatorname{compose}(f, k)) \\ \quad \quad \in\; (\operatorname{hom\text{-}lipschitz}(a, b) \to \operatorname{hom\text{-}lipschitz}(c, b)) \end{array}}}
\vnbentry{\texttt{hom-\allowbreak{}uniformly-\allowbreak{}continuous-\allowbreak{}post-\allowbreak{}type}}{\vnbstmt{\begin{array}{@{}l@{}} \forall a, b, c, h. \\ \quad \operatorname{is\text{-}metric\text{-}space}(a) \wedge \\ \quad (h \in \operatorname{hom\text{-}uniformly\text{-}continuous}(b, c)) \Rightarrow \\ \quad \quad (\lambda f \in \operatorname{hom\text{-}uniformly\text{-}continuous}(a, b).\; \\ \quad \quad \quad \operatorname{compose}(h, f)) \\ \quad \quad \in\; (\operatorname{hom\text{-}uniformly\text{-}continuous}(a, b) \to \operatorname{hom\text{-}uniformly\text{-}continuous}(a, c)) \end{array}}}
\vnbentry{\texttt{hom-\allowbreak{}uniformly-\allowbreak{}continuous-\allowbreak{}pre-\allowbreak{}type}}{\vnbstmt{\begin{array}{@{}l@{}} \forall a, b, c, k. \\ \quad \operatorname{is\text{-}metric\text{-}space}(b) \wedge \\ \quad (k \in \operatorname{hom\text{-}uniformly\text{-}continuous}(c, a)) \Rightarrow \\ \quad \quad (\lambda f \in \operatorname{hom\text{-}uniformly\text{-}continuous}(a, b).\; \\ \quad \quad \quad \operatorname{compose}(f, k)) \\ \quad \quad \in\; (\operatorname{hom\text{-}uniformly\text{-}continuous}(a, b) \to \operatorname{hom\text{-}uniformly\text{-}continuous}(c, b)) \end{array}}}
\vnbentry{\texttt{hom-\allowbreak{}continuous-\allowbreak{}post-\allowbreak{}type}}{\vnbstmt{\begin{array}{@{}l@{}} \forall a, b, c, h. \\ \quad \operatorname{is\text{-}metric\text{-}space}(a) \wedge \\ \quad (h \in \operatorname{hom\text{-}continuous}(b, c)) \Rightarrow \\ \quad \quad (\lambda f \in \operatorname{hom\text{-}continuous}(a, b).\; \operatorname{compose}(h, f)) \\ \quad \quad \in\; (\operatorname{hom\text{-}continuous}(a, b) \to \operatorname{hom\text{-}continuous}(a, c)) \end{array}}}
\vnbentry{\texttt{hom-\allowbreak{}continuous-\allowbreak{}pre-\allowbreak{}type}}{\vnbstmt{\begin{array}{@{}l@{}} \forall a, b, c, k. \\ \quad \operatorname{is\text{-}metric\text{-}space}(b) \wedge \\ \quad (k \in \operatorname{hom\text{-}continuous}(c, a)) \Rightarrow \\ \quad \quad (\lambda f \in \operatorname{hom\text{-}continuous}(a, b).\; \operatorname{compose}(f, k)) \\ \quad \quad \in\; (\operatorname{hom\text{-}continuous}(a, b) \to \operatorname{hom\text{-}continuous}(c, b)) \end{array}}}
\vnbentry{\texttt{hom-\allowbreak{}bounded-\allowbreak{}linear-\allowbreak{}post-\allowbreak{}type}}{\vnbstmt{\begin{array}{@{}l@{}} \forall a, b, c, h. \\ \quad \operatorname{is\text{-}normed\text{-}vector\text{-}space}(a) \wedge \\ \quad (h \in \operatorname{hom\text{-}bounded\text{-}linear}(b, c)) \Rightarrow \\ \quad \quad (\lambda f \in \operatorname{hom\text{-}bounded\text{-}linear}(a, b).\; \\ \quad \quad \quad \operatorname{compose}(h, f)) \\ \quad \quad \in\; (\operatorname{hom\text{-}bounded\text{-}linear}(a, b) \to \operatorname{hom\text{-}bounded\text{-}linear}(a, c)) \end{array}}}
\vnbentry{\texttt{hom-\allowbreak{}bounded-\allowbreak{}linear-\allowbreak{}pre-\allowbreak{}type}}{\vnbstmt{\begin{array}{@{}l@{}} \forall a, b, c, k. \\ \quad \operatorname{is\text{-}normed\text{-}vector\text{-}space}(b) \wedge \\ \quad (k \in \operatorname{hom\text{-}bounded\text{-}linear}(c, a)) \Rightarrow \\ \quad \quad (\lambda f \in \operatorname{hom\text{-}bounded\text{-}linear}(a, b).\; \\ \quad \quad \quad \operatorname{compose}(f, k)) \\ \quad \quad \in\; (\operatorname{hom\text{-}bounded\text{-}linear}(a, b) \to \operatorname{hom\text{-}bounded\text{-}linear}(c, b)) \end{array}}}
\vnbentry{\texttt{hom-\allowbreak{}metric-\allowbreak{}space-\allowbreak{}subset-\allowbreak{}hom-\allowbreak{}lipschitz}}{\vnbstmt{\begin{array}{@{}l@{}} \forall a, b. \\ \quad \operatorname{hom\text{-}metric\text{-}space}(a, b) \\ \quad \subseteq\; \operatorname{hom\text{-}lipschitz}(a, b) \end{array}}}
\vnbentry{\texttt{hom-\allowbreak{}lipschitz-\allowbreak{}subset-\allowbreak{}hom-\allowbreak{}uniformly-\allowbreak{}continuous}}{\vnbstmt{\begin{array}{@{}l@{}} \forall a, b. \\ \quad \operatorname{hom\text{-}lipschitz}(a, b) \\ \quad \subseteq\; \operatorname{hom\text{-}uniformly\text{-}continuous}(a, b) \end{array}}}
\vnbentry{\texttt{hom-\allowbreak{}uniformly-\allowbreak{}continuous-\allowbreak{}subset-\allowbreak{}hom-\allowbreak{}continuous}}{\vnbstmt{\begin{array}{@{}l@{}} \forall a, b. \\ \quad \operatorname{hom\text{-}uniformly\text{-}continuous}(a, b) \\ \quad \subseteq\; \operatorname{hom\text{-}continuous}(a, b) \end{array}}}
\vnbentry{\texttt{hom-\allowbreak{}normed-\allowbreak{}vector-\allowbreak{}space-\allowbreak{}subset-\allowbreak{}hom-\allowbreak{}bounded-\allowbreak{}linear}}{\vnbstmt{\begin{array}{@{}l@{}} \forall a, b. \\ \quad \operatorname{hom\text{-}normed\text{-}vector\text{-}space}(a, b) \\ \quad \subseteq\; \operatorname{hom\text{-}bounded\text{-}linear}(a, b) \end{array}}}

\vnbfile{theorem-library/c-continuous-category.scm}{theorem-library/c-continuous-category.scm}
\noindent{\small c-continuous-category.scm -- THE CATEGORY OF COUNTABLY-METRISED SPACES AND CONTINUOUS MAPS: objects the C-METRIC-SPACEs, arrows the maps f : PTS(a) -$>$ PTS(b) continuous for the canonical metrics C-METRIC(a), C-METRIC(b).}\par
\vnbentry{\texttt{c-\allowbreak{}metric-\allowbreak{}canonical-\allowbreak{}carrier}}{\vnbstmt{\begin{array}{@{}l@{}} \forall s. \\ \quad (\operatorname{pts}(\operatorname{c\text{-}metric}(s)) \simeq \operatorname{pts}(s)) \end{array}}}
\vnbentry{\texttt{hom-\allowbreak{}c-\allowbreak{}continuous-\allowbreak{}member-\allowbreak{}iff}}{\vnbstmt{\begin{array}{@{}l@{}} \forall a, b, f. \\ \quad ((f \in \operatorname{hom\text{-}c\text{-}continuous}(a, b)) \Leftrightarrow ((f \in (\operatorname{pts}(a) \to \operatorname{pts}(b))) \wedge \operatorname{is\text{-}c\text{-}continuous\text{-}arrow}(a, b, f))) \end{array}}}
\vnbentry{\texttt{hom-\allowbreak{}c-\allowbreak{}continuous-\allowbreak{}in-\allowbreak{}set}}{\vnbstmt{\begin{array}{@{}l@{}} \forall a, b. \\ \quad \operatorname{is\text{-}c\text{-}metric\text{-}space}(a) \wedge \\ \quad \operatorname{is\text{-}c\text{-}metric\text{-}space}(b) \Rightarrow \\ \quad \quad (\operatorname{hom\text{-}c\text{-}continuous}(a, b) \in \mathbf{Set}) \end{array}}}
\vnbentry{\texttt{hom-\allowbreak{}c-\allowbreak{}continuous-\allowbreak{}compose}}{\vnbstmt{\begin{array}{@{}l@{}} \forall a, b, c, f, g. \\ \quad \operatorname{is\text{-}c\text{-}continuous\text{-}arrow}(a, b, f) \wedge \\ \quad \operatorname{is\text{-}c\text{-}continuous\text{-}arrow}(b, c, g) \Rightarrow \\ \quad \quad \operatorname{is\text{-}c\text{-}continuous\text{-}arrow}(a, c, \operatorname{compose}(g, f)) \end{array}}}
\vnbentry{\texttt{hom-\allowbreak{}c-\allowbreak{}continuous-\allowbreak{}id-\allowbreak{}of-\allowbreak{}metric}}{\vnbstmt{\begin{array}{@{}l@{}} \forall a. \\ \quad \operatorname{is\text{-}c\text{-}metric\text{-}space}(a) \wedge \\ \quad \operatorname{is\text{-}metric\text{-}space}(\operatorname{c\text{-}metric}(a)) \Rightarrow \\ \quad \quad \operatorname{is\text{-}c\text{-}continuous\text{-}arrow}(a, a, \operatorname{id\text{-}fun}(\operatorname{pts}(a))) \end{array}}}
\vnbentry{\texttt{hom-\allowbreak{}c-\allowbreak{}continuous-\allowbreak{}post-\allowbreak{}type}}{\vnbstmt{\begin{array}{@{}l@{}} \forall a, b, c, h. \\ \quad \operatorname{is\text{-}c\text{-}metric\text{-}space}(a) \wedge \\ \quad (h \in \operatorname{hom\text{-}c\text{-}continuous}(b, c)) \Rightarrow \\ \quad \quad (\lambda f \in \operatorname{hom\text{-}c\text{-}continuous}(a, b).\; \operatorname{compose}(h, f)) \\ \quad \quad \in\; (\operatorname{hom\text{-}c\text{-}continuous}(a, b) \to \operatorname{hom\text{-}c\text{-}continuous}(a, c)) \end{array}}}
\vnbentry{\texttt{hom-\allowbreak{}c-\allowbreak{}continuous-\allowbreak{}pre-\allowbreak{}type}}{\vnbstmt{\begin{array}{@{}l@{}} \forall a, b, c, k. \\ \quad \operatorname{is\text{-}c\text{-}metric\text{-}space}(b) \wedge \\ \quad (k \in \operatorname{hom\text{-}c\text{-}continuous}(c, a)) \Rightarrow \\ \quad \quad (\lambda f \in \operatorname{hom\text{-}c\text{-}continuous}(a, b).\; \operatorname{compose}(f, k)) \\ \quad \quad \in\; (\operatorname{hom\text{-}c\text{-}continuous}(a, b) \to \operatorname{hom\text{-}c\text{-}continuous}(c, b)) \end{array}}}
\vnbentry{\texttt{hom-\allowbreak{}c-\allowbreak{}continuous-\allowbreak{}id}}{\vnbstmt{\begin{array}{@{}l@{}} \forall a. \\ \quad \operatorname{is\text{-}c\text{-}metric\text{-}space}(a) \Rightarrow  \\ \quad \quad \operatorname{is\text{-}c\text{-}continuous\text{-}arrow}(a, a, \operatorname{id\text{-}fun}(\operatorname{pts}(a))) \end{array}}}

\vnbfile{theorem-library/diff-on-laws-2.scm}{theorem-library/diff-on-laws-2.scm}
\noindent{\small diff-on-laws-2.scm -- THE DERIVATIVE LAWS ON AN OPEN SET, over a normed field: uniqueness, the constant and the identity, the sum, the product and the chain rule.  Continues theorem-library/diff-on-laws.scm (batch 12-G), which proved the read-offs, the bridge `diff-at-iff-diff-on' and}\par
\vnbentry{\texttt{ms-\allowbreak{}const-\allowbreak{}continuous-\allowbreak{}on}}{\vnbstmt{\begin{array}{@{}l@{}} \forall s, d, t.; \forall c \in \operatorname{pts}(t), a \in \operatorname{pts}(s). \\ \quad (\operatorname{pts}(s) \simeq d) \wedge \\ \quad \operatorname{is\text{-}metric\text{-}space}(s) \wedge \\ \quad \operatorname{is\text{-}metric\text{-}space}(t) \Rightarrow \\ \quad \quad \operatorname{is\text{-}continuous\text{-}at}(s, t, (\lambda z \in d.\; c), a) \end{array}}}
\vnbentry{\texttt{ms-\allowbreak{}identity-\allowbreak{}continuous-\allowbreak{}on}}{\vnbstmt{\begin{array}{@{}l@{}} \forall s, d.; \forall a \in \operatorname{pts}(s). \\ \quad ((\operatorname{pts}(s) \simeq d) \wedge \operatorname{is\text{-}metric\text{-}space}(s)) \Rightarrow  \\ \quad \quad \operatorname{is\text{-}continuous\text{-}at}(s, s, (\lambda z \in d.\; z), a) \end{array}}}
\vnbentry{\texttt{nf-\allowbreak{}sum-\allowbreak{}continuous-\allowbreak{}on}}{\vnbstmt{\begin{array}{@{}l@{}} \forall k, s, d, g, h, a. \\ \quad \operatorname{is\text{-}normed\text{-}field}(k) \wedge \\ \quad (\operatorname{pts}(s) \simeq d) \wedge \\ \quad \operatorname{is\text{-}continuous\text{-}at}(s, \operatorname{nf\text{-}metric\text{-}space}(k), g, a) \wedge \\ \quad \operatorname{is\text{-}continuous\text{-}at}(s, \operatorname{nf\text{-}metric\text{-}space}(k), h, a) \Rightarrow \\ \quad \quad \operatorname{is\text{-}continuous\text{-}at}(s, \operatorname{nf\text{-}metric\text{-}space}(k), (\lambda z \in d.\; \operatorname{add}(k)(g(z), h(z))), a) \end{array}}}
\vnbentry{\texttt{nf-\allowbreak{}product-\allowbreak{}continuous-\allowbreak{}on}}{\vnbstmt{\begin{array}{@{}l@{}} \forall k, s, d, g, h, a. \\ \quad \operatorname{is\text{-}normed\text{-}field}(k) \wedge \\ \quad (\operatorname{pts}(s) \simeq d) \wedge \\ \quad \operatorname{is\text{-}continuous\text{-}at}(s, \operatorname{nf\text{-}metric\text{-}space}(k), g, a) \wedge \\ \quad \operatorname{is\text{-}continuous\text{-}at}(s, \operatorname{nf\text{-}metric\text{-}space}(k), h, a) \Rightarrow \\ \quad \quad \operatorname{is\text{-}continuous\text{-}at}(s, \operatorname{nf\text{-}metric\text{-}space}(k), (\lambda z \in d.\; \operatorname{mul}(k)(g(z), h(z))), a) \end{array}}}
\vnbentry{\texttt{ms-\allowbreak{}compose-\allowbreak{}continuous-\allowbreak{}on}}{\vnbstmt{\begin{array}{@{}l@{}} \forall s, d, t, u, f, g, a. \\ \quad (\operatorname{pts}(s) \simeq d) \wedge \\ \quad \operatorname{is\text{-}continuous\text{-}at}(s, t, f, a) \wedge \\ \quad \operatorname{is\text{-}continuous\text{-}at}(t, u, g, f(a)) \Rightarrow \\ \quad \quad \operatorname{is\text{-}continuous\text{-}at}(s, u, (\lambda z \in d.\; g(f(z))), a) \end{array}}}
\vnbentry{\texttt{ms-\allowbreak{}cont-\allowbreak{}transfer-\allowbreak{}on}}{\vnbstmt{\begin{array}{@{}l@{}} \forall s, t, d, c, g, f, a. \\ \quad (\operatorname{pts}(s) \simeq d) \wedge \\ \quad (\operatorname{pts}(t) \simeq c) \wedge \\ \quad \operatorname{is\text{-}continuous\text{-}at}(s, t, g, a) \wedge \\ \quad (f \in (d \to c)) \wedge \\ \quad \forall y \in d.\; (f(y) = g(y)) \Rightarrow \\ \quad \quad \operatorname{is\text{-}continuous\text{-}at}(s, t, f, a) \end{array}}}
\vnbentry{\texttt{nf-\allowbreak{}lam-\allowbreak{}const-\allowbreak{}in-\allowbreak{}fun}}{\vnbstmt{\begin{array}{@{}l@{}} \forall k.; \forall u \in \mathbf{Set}, c \in \operatorname{carr}(k). \\ \quad \operatorname{is\text{-}normed\text{-}field}(k) \Rightarrow  \\ \quad \quad ((\lambda x \in u.\; c) \in (u \to \operatorname{carr}(k))) \end{array}}}
\vnbentry{\texttt{nf-\allowbreak{}lam-\allowbreak{}id-\allowbreak{}in-\allowbreak{}fun}}{\vnbstmt{\begin{array}{@{}l@{}} \forall k.; \forall u \in \mathbf{Set}. \\ \quad (\operatorname{is\text{-}normed\text{-}field}(k) \wedge (u \subseteq \operatorname{carr}(k))) \Rightarrow  \\ \quad \quad ((\lambda x \in u.\; x) \in (u \to \operatorname{carr}(k))) \end{array}}}
\vnbentry{\texttt{nf-\allowbreak{}lam-\allowbreak{}sum-\allowbreak{}in-\allowbreak{}fun}}{\vnbstmt{\begin{array}{@{}l@{}} \forall k.; \forall u \in \mathbf{Set}, f, g \in (u \to \operatorname{carr}(k)). \\ \quad \operatorname{is\text{-}normed\text{-}field}(k) \Rightarrow  \\ \quad \quad ((\lambda x \in u.\; \operatorname{add}(k)(f(x), g(x))) \in (u \to \operatorname{carr}(k))) \end{array}}}
\vnbentry{\texttt{nf-\allowbreak{}lam-\allowbreak{}prod-\allowbreak{}in-\allowbreak{}fun}}{\vnbstmt{\begin{array}{@{}l@{}} \forall k.; \forall u \in \mathbf{Set}, f, g \in (u \to \operatorname{carr}(k)). \\ \quad \operatorname{is\text{-}normed\text{-}field}(k) \Rightarrow  \\ \quad \quad ((\lambda x \in u.\; \operatorname{mul}(k)(f(x), g(x))) \in (u \to \operatorname{carr}(k))) \end{array}}}
\vnbentry{\texttt{nf-\allowbreak{}norm-\allowbreak{}zero-\allowbreak{}iff}}{\vnbstmt{\begin{array}{@{}l@{}} \forall k.; \forall u \in \operatorname{carr}(k). \\ \quad \operatorname{is\text{-}normed\text{-}field}(k) \Rightarrow  \\ \quad \quad ((\operatorname{fnrm}(k)(u) = 0) \Leftrightarrow (u = \operatorname{zero}(k))) \end{array}}}
\vnbentry{\texttt{nf-\allowbreak{}norm-\allowbreak{}zero-\allowbreak{}implies}}{\vnbstmt{\begin{array}{@{}l@{}} \forall k.; \forall u \in \operatorname{carr}(k). \\ \quad (\operatorname{is\text{-}normed\text{-}field}(k) \wedge (\operatorname{fnrm}(k)(u) = 0)) \Rightarrow  \\ \quad \quad (u = \operatorname{zero}(k)) \end{array}}}
\vnbentry{\texttt{rr-\allowbreak{}nf-\allowbreak{}small-\allowbreak{}elements}}{\vnbstmt{\begin{array}{@{}l@{}} \forall e. \\ \quad \operatorname{pos\text{-}rr}(e) \Rightarrow  \\ \quad \quad \exists h \in \operatorname{carr}(\operatorname{rr\text{-}normed\text{-}field}). \\ \quad \quad \quad \neg (h = \operatorname{zero}(\operatorname{rr\text{-}normed\text{-}field})) \wedge \\ \quad \quad \quad (\operatorname{fnrm}(\operatorname{rr\text{-}normed\text{-}field})(h) < e) \end{array}}}
\vnbentry{\texttt{nf-\allowbreak{}open-\allowbreak{}limit-\allowbreak{}point}}{\vnbstmt{\begin{array}{@{}l@{}} \forall k, u.; \forall a \in u, r. \\ \quad \operatorname{is\text{-}normed\text{-}field}(k) \wedge \\ \quad \forall e. \\ \quad \quad \operatorname{pos\text{-}rr}(e) \Rightarrow  \\ \quad \quad \quad \exists h \in \operatorname{carr}(k). \\ \quad \quad \quad \quad (\neg (h = \operatorname{zero}(k)) \wedge (\operatorname{fnrm}(k)(h) < e)) \wedge \\ \quad \operatorname{is\text{-}open}(\operatorname{nf\text{-}metric\text{-}space}(k), u) \wedge \\ \quad \operatorname{pos\text{-}rr}(r) \Rightarrow \\ \quad \quad \exists b \in u. \\ \quad \quad \quad \neg (b = a) \wedge \\ \quad \quad \quad \operatorname{dist}(\operatorname{subspace\text{-}ms}(\operatorname{nf\text{-}metric\text{-}space}(k), u))(a, b) \\ \quad \quad \quad \leq\; r \end{array}}}
\vnbentry{\texttt{nf-\allowbreak{}zero-\allowbreak{}in-\allowbreak{}carr}}{\vnbstmt{\begin{array}{@{}l@{}} \forall k. \\ \quad \operatorname{is\text{-}normed\text{-}field}(k) \Rightarrow  \\ \quad \quad (\operatorname{zero}(k) \in \operatorname{carr}(k)) \end{array}}}
\vnbentry{\texttt{nf-\allowbreak{}norm-\allowbreak{}of-\allowbreak{}zero}}{\vnbstmt{\begin{array}{@{}l@{}} \forall k. \\ \quad \operatorname{is\text{-}normed\text{-}field}(k) \Rightarrow  \\ \quad \quad (\operatorname{fnrm}(k)(\operatorname{zero}(k)) = 0) \end{array}}}
\vnbentry{\texttt{nf-\allowbreak{}sub-\allowbreak{}zero-\allowbreak{}eq}}{\vnbstmt{\begin{array}{@{}l@{}} \forall k, u, v. \\ \quad \operatorname{is\text{-}normed\text{-}field}(k) \wedge \\ \quad (u \in \operatorname{carr}(k)) \wedge \\ \quad (v \in \operatorname{carr}(k)) \wedge \\ \quad (\operatorname{add}(k)(u, \operatorname{neg}(k)(v)) = \operatorname{zero}(k)) \Rightarrow \\ \quad \quad (u = v) \end{array}}}
\vnbentry{\texttt{nf-\allowbreak{}mul-\allowbreak{}cancel}}{\vnbstmt{\begin{array}{@{}l@{}} \forall k, u, v, w. \\ \quad \operatorname{is\text{-}normed\text{-}field}(k) \wedge \\ \quad (u \in \operatorname{carr}(k)) \wedge \\ \quad (v \in \operatorname{carr}(k)) \wedge \\ \quad (w \in \operatorname{carr}(k)) \wedge \\ \quad \neg (w = \operatorname{zero}(k)) \wedge \\ \quad (\operatorname{mul}(k)(u, w) = \operatorname{mul}(k)(v, w)) \Rightarrow \\ \quad \quad (u = v) \end{array}}}
\vnbentry{\texttt{diff-\allowbreak{}on-\allowbreak{}unique}}{\vnbstmt{\begin{array}{@{}l@{}} \forall k, u, f, a, l, m. \\ \quad \operatorname{is\text{-}normed\text{-}field}(k) \wedge \\ \quad \forall e. \\ \quad \quad \operatorname{pos\text{-}rr}(e) \Rightarrow  \\ \quad \quad \quad \exists h \in \operatorname{carr}(k). \\ \quad \quad \quad \quad (\neg (h = \operatorname{zero}(k)) \wedge (\operatorname{fnrm}(k)(h) < e)) \wedge \\ \quad \operatorname{is\text{-}diff\text{-}on}(k, u, f, a, l) \wedge \\ \quad \operatorname{is\text{-}diff\text{-}on}(k, u, f, a, m) \Rightarrow \\ \quad \quad (l = m) \end{array}}}
\vnbentry{\texttt{nf-\allowbreak{}one-\allowbreak{}in-\allowbreak{}carr}}{\vnbstmt{\begin{array}{@{}l@{}} \forall k. \\ \quad \operatorname{is\text{-}normed\text{-}field}(k) \Rightarrow  \\ \quad \quad (\operatorname{one}(k) \in \operatorname{carr}(k)) \end{array}}}
\vnbentry{\texttt{nf-\allowbreak{}open-\allowbreak{}is-\allowbreak{}set}}{\vnbstmt{\begin{array}{@{}l@{}} \forall k, u. \\ \quad \operatorname{is\text{-}normed\text{-}field}(k) \wedge \\ \quad \operatorname{is\text{-}open}(\operatorname{nf\text{-}metric\text{-}space}(k), u) \Rightarrow \\ \quad \quad (u \in \mathbf{Set}) \end{array}}}
\vnbentry{\texttt{diff-\allowbreak{}on-\allowbreak{}const}}{\vnbstmt{\begin{array}{@{}l@{}} \forall k, u.; \forall c \in \operatorname{carr}(k), a \in u. \\ \quad \operatorname{is\text{-}normed\text{-}field}(k) \wedge \\ \quad \operatorname{is\text{-}open}(\operatorname{nf\text{-}metric\text{-}space}(k), u) \Rightarrow \\ \quad \quad \operatorname{is\text{-}diff\text{-}on}(k, u, (\lambda x \in u.\; c), a, \operatorname{zero}(k)) \end{array}}}
\vnbentry{\texttt{diff-\allowbreak{}on-\allowbreak{}identity}}{\vnbstmt{\begin{array}{@{}l@{}} \forall k, u.; \forall a \in u. \\ \quad \operatorname{is\text{-}normed\text{-}field}(k) \wedge \\ \quad \operatorname{is\text{-}open}(\operatorname{nf\text{-}metric\text{-}space}(k), u) \Rightarrow \\ \quad \quad \operatorname{is\text{-}diff\text{-}on}(k, u, (\lambda x \in u.\; x), a, \operatorname{one}(k)) \end{array}}}
\vnbentry{\texttt{diff-\allowbreak{}on-\allowbreak{}sum}}{\vnbstmt{\begin{array}{@{}l@{}} \forall k, u, f, g, a, l, m. \\ \quad \operatorname{is\text{-}normed\text{-}field}(k) \wedge \\ \quad \operatorname{is\text{-}diff\text{-}on}(k, u, f, a, l) \wedge \\ \quad \operatorname{is\text{-}diff\text{-}on}(k, u, g, a, m) \Rightarrow \\ \quad \quad \operatorname{is\text{-}diff\text{-}on}(k, u, (\lambda x \in u.\; \operatorname{add}(k)(f(x), g(x))), a, \operatorname{add}(k)(l, m)) \end{array}}}
\vnbentry{\texttt{nf-\allowbreak{}lam-\allowbreak{}const-\allowbreak{}z-\allowbreak{}in-\allowbreak{}fun}}{\vnbstmt{\begin{array}{@{}l@{}} \forall k.; \forall u \in \mathbf{Set}, c \in \operatorname{carr}(k). \\ \quad \operatorname{is\text{-}normed\text{-}field}(k) \Rightarrow  \\ \quad \quad ((\lambda z \in u.\; c) \in (u \to \operatorname{carr}(k))) \end{array}}}
\vnbentry{\texttt{nf-\allowbreak{}lam-\allowbreak{}prod-\allowbreak{}v-\allowbreak{}in-\allowbreak{}fun}}{\vnbstmt{\begin{array}{@{}l@{}} \forall k.; \forall u \in \mathbf{Set}, f, g \in (u \to \operatorname{carr}(k)). \\ \quad \operatorname{is\text{-}normed\text{-}field}(k) \Rightarrow  \\ \quad \quad ((\lambda v \in u.\; \operatorname{mul}(k)(f(v), g(v))) \in (u \to \operatorname{carr}(k))) \end{array}}}
\vnbentry{\texttt{nf-\allowbreak{}product-\allowbreak{}continuous-\allowbreak{}v}}{\vnbstmt{\begin{array}{@{}l@{}} \forall k, s, d, g, h, a. \\ \quad \operatorname{is\text{-}normed\text{-}field}(k) \wedge \\ \quad (\operatorname{pts}(s) \simeq d) \wedge \\ \quad \operatorname{is\text{-}continuous\text{-}at}(s, \operatorname{nf\text{-}metric\text{-}space}(k), g, a) \wedge \\ \quad \operatorname{is\text{-}continuous\text{-}at}(s, \operatorname{nf\text{-}metric\text{-}space}(k), h, a) \Rightarrow \\ \quad \quad \operatorname{is\text{-}continuous\text{-}at}(s, \operatorname{nf\text{-}metric\text{-}space}(k), (\lambda v \in d.\; \operatorname{mul}(k)(g(v), h(v))), a) \end{array}}}
\vnbentry{\texttt{diff-\allowbreak{}on-\allowbreak{}product}}{\vnbstmt{\begin{array}{@{}l@{}} \forall k, u, f, g, a, l, m. \\ \quad \operatorname{is\text{-}normed\text{-}field}(k) \wedge \\ \quad \operatorname{is\text{-}diff\text{-}on}(k, u, f, a, l) \wedge \\ \quad \operatorname{is\text{-}diff\text{-}on}(k, u, g, a, m) \Rightarrow \\ \quad \quad \operatorname{is\text{-}diff\text{-}on}(k, u, (\lambda x \in u.\; \operatorname{mul}(k)(f(x), g(x))), a, \\ \quad \quad \quad \operatorname{add}(k)(\operatorname{mul}(k)(l, g(a)), \operatorname{mul}(k)(m, f(a)))) \end{array}}}
\vnbentry{\texttt{ms-\allowbreak{}corestrict-\allowbreak{}continuous}}{\vnbstmt{\begin{array}{@{}l@{}} \forall s, t, \mathit{d2s}, f, a. \\ \quad \operatorname{is\text{-}metric\text{-}space}(t) \wedge \\ \quad (\mathit{d2s} \subseteq \operatorname{pts}(t)) \wedge \\ \quad \operatorname{is\text{-}continuous\text{-}at}(s, t, f, a) \wedge \\ \quad \forall y \in \operatorname{pts}(s).\; (f(y) \in \mathit{d2s}) \Rightarrow \\ \quad \quad \operatorname{is\text{-}continuous\text{-}at}(s, \operatorname{subspace\text{-}ms}(t, \mathit{d2s}), f, a) \end{array}}}
\vnbentry{\texttt{nf-\allowbreak{}lam-\allowbreak{}comp-\allowbreak{}z-\allowbreak{}in-\allowbreak{}fun}}{\vnbstmt{\begin{array}{@{}l@{}} \forall k.; \forall u \in \mathbf{Set}, s, f.; \forall g \in (s \to \operatorname{carr}(k)). \\ \quad (\operatorname{is\text{-}normed\text{-}field}(k) \wedge \forall y \in u.\; (f(y) \in s)) \Rightarrow  \\ \quad \quad ((\lambda z \in u.\; g(f(z))) \in (u \to \operatorname{carr}(k))) \end{array}}}
\vnbentry{\texttt{diff-\allowbreak{}on-\allowbreak{}chain}}{\vnbstmt{\begin{array}{@{}l@{}} \forall k, u, s, f, g, a, l, m. \\ \quad \operatorname{is\text{-}normed\text{-}field}(k) \wedge \\ \quad \operatorname{is\text{-}diff\text{-}on}(k, u, f, a, l) \wedge \\ \quad \operatorname{is\text{-}diff\text{-}on}(k, s, g, f(a), m) \wedge \\ \quad \forall y \in u.\; (f(y) \in s) \Rightarrow \\ \quad \quad \operatorname{is\text{-}diff\text{-}on}(k, u, (\lambda x \in u.\; g(f(x))), a, \operatorname{mul}(k)(m, l)) \end{array}}}

\vnbfile{theorem-library/holomorphic-basics.scm}{theorem-library/holomorphic-basics.scm}
\noindent{\small holomorphic-basics.scm -- the COMPLEX side of IS-DIFF-ON / HOLOMORPHIC-ON.}\par
\vnbentry{\texttt{cc-\allowbreak{}nf-\allowbreak{}add-\allowbreak{}apply}}{\vnbstmt{\begin{array}{@{}l@{}} \forall u, v. \\ \quad ((u \in \mathbb{C}) \wedge (v \in \mathbb{C})) \Rightarrow  \\ \quad \quad (\operatorname{add}(\operatorname{cc\text{-}normed\text{-}field})(u, v) \simeq (u + v)) \end{array}}}
\vnbentry{\texttt{cc-\allowbreak{}nf-\allowbreak{}mul-\allowbreak{}apply}}{\vnbstmt{\begin{array}{@{}l@{}} \forall u, v. \\ \quad ((u \in \mathbb{C}) \wedge (v \in \mathbb{C})) \Rightarrow  \\ \quad \quad (\operatorname{mul}(\operatorname{cc\text{-}normed\text{-}field})(u, v) \simeq (u \cdot v)) \end{array}}}
\vnbentry{\texttt{cc-\allowbreak{}nf-\allowbreak{}neg-\allowbreak{}apply}}{\vnbstmt{\begin{array}{@{}l@{}} \forall u \in \mathbb{C}. \\ \quad (\operatorname{neg}(\operatorname{cc\text{-}normed\text{-}field})(u) \simeq -u) \end{array}}}
\vnbentry{\texttt{cc-\allowbreak{}nf-\allowbreak{}sub}}{\vnbstmt{\begin{array}{@{}l@{}} \forall u, v. \\ \quad ((u \in \mathbb{C}) \wedge (v \in \mathbb{C})) \Rightarrow  \\ \quad \quad (\operatorname{add}(\operatorname{cc\text{-}normed\text{-}field})(u, \operatorname{neg}(\operatorname{cc\text{-}normed\text{-}field})(v)) \simeq (u - v)) \end{array}}}
\vnbentry{\texttt{nf-\allowbreak{}cc-\allowbreak{}is-\allowbreak{}metric-\allowbreak{}space}}{\vnbstmt{\operatorname{is\text{-}metric\text{-}space}(\operatorname{nf\text{-}metric\text{-}space}(\operatorname{cc\text{-}normed\text{-}field}))}}
\vnbentry{\texttt{nf-\allowbreak{}cc-\allowbreak{}agree-\allowbreak{}pts}}{\vnbstmt{(\operatorname{pts}(\operatorname{nf\text{-}metric\text{-}space}(\operatorname{cc\text{-}normed\text{-}field})) \simeq \operatorname{pts}(\operatorname{cc\text{-}ms}))}}
\vnbentry{\texttt{nf-\allowbreak{}cc-\allowbreak{}agree-\allowbreak{}dist}}{\vnbstmt{\begin{array}{@{}l@{}} \forall u, v \in \operatorname{pts}(\operatorname{cc\text{-}ms}). \\ \quad (\operatorname{dist}(\operatorname{nf\text{-}metric\text{-}space}(\operatorname{cc\text{-}normed\text{-}field}))(u, v) \simeq \operatorname{dist}(\operatorname{cc\text{-}ms})(u, v)) \end{array}}}
\vnbentry{\texttt{cc-\allowbreak{}ms-\allowbreak{}open-\allowbreak{}iff}}{\vnbstmt{\begin{array}{@{}l@{}} \forall w. \\ \quad (\operatorname{is\text{-}open}(\operatorname{nf\text{-}metric\text{-}space}(\operatorname{cc\text{-}normed\text{-}field}), w) \Leftrightarrow \operatorname{is\text{-}open}(\operatorname{cc\text{-}ms}, w)) \end{array}}}
\vnbentry{\texttt{cc-\allowbreak{}nf-\allowbreak{}small-\allowbreak{}elements}}{\vnbstmt{\begin{array}{@{}l@{}} \forall s. \\ \quad \operatorname{pos\text{-}rr}(s) \Rightarrow  \\ \quad \quad \exists h \in \operatorname{carr}(\operatorname{cc\text{-}normed\text{-}field}). \\ \quad \quad \quad \neg (h = \operatorname{zero}(\operatorname{cc\text{-}normed\text{-}field})) \wedge \\ \quad \quad \quad (\operatorname{fnrm}(\operatorname{cc\text{-}normed\text{-}field})(h) < s) \end{array}}}
\vnbentry{\texttt{holomorphic-\allowbreak{}on-\allowbreak{}open-\allowbreak{}cc-\allowbreak{}ms}}{\vnbstmt{\begin{array}{@{}l@{}} \forall u, f. \\ \quad \operatorname{holomorphic\text{-}on}(u, f) \Rightarrow  \\ \quad \quad \operatorname{is\text{-}open}(\operatorname{cc\text{-}ms}, u) \end{array}}}
\vnbentry{\texttt{holomorphic-\allowbreak{}on-\allowbreak{}iff-\allowbreak{}cc-\allowbreak{}ms}}{\vnbstmt{\begin{array}{@{}l@{}} \forall u, f. \\ \quad (\operatorname{holomorphic\text{-}on}(u, f) \Leftrightarrow (\operatorname{is\text{-}open}(\operatorname{cc\text{-}ms}, u) \wedge ((f \in (u \to \mathbb{C})) \wedge \forall a \in u.\; \exists l.\; \operatorname{is\text{-}diff\text{-}on}(\operatorname{cc\text{-}normed\text{-}field}, u, f, a, l)))) \end{array}}}

\vnbfile{theorem-library/deriv-on-read-off.scm}{theorem-library/deriv-on-read-off.scm}
\noindent{\small deriv-on-read-off.scm -- DERIV-ON reads off IS-DIFF-ON. Batch 27-A follow-on, 2026-09-24.}\par
\vnbentry{\texttt{deriv-\allowbreak{}on-\allowbreak{}of-\allowbreak{}is-\allowbreak{}diff-\allowbreak{}on}}{\vnbstmt{\begin{array}{@{}l@{}} \forall k, u, f, a, l. \\ \quad \operatorname{is\text{-}normed\text{-}field}(k) \wedge \\ \quad \forall e. \\ \quad \quad \operatorname{pos\text{-}rr}(e) \Rightarrow  \\ \quad \quad \quad \exists h \in \operatorname{carr}(k). \\ \quad \quad \quad \quad (\neg (h = \operatorname{zero}(k)) \wedge (\operatorname{fnrm}(k)(h) < e)) \wedge \\ \quad \operatorname{is\text{-}diff\text{-}on}(k, u, f, a, l) \Rightarrow \\ \quad \quad (\operatorname{deriv\text{-}on}(k, u, f, a) = l) \end{array}}}
\vnbentry{\texttt{cc-\allowbreak{}deriv-\allowbreak{}on-\allowbreak{}of-\allowbreak{}is-\allowbreak{}diff-\allowbreak{}on}}{\vnbstmt{\begin{array}{@{}l@{}} \forall u, f, a, l. \\ \quad \operatorname{is\text{-}diff\text{-}on}(\operatorname{cc\text{-}normed\text{-}field}, u, f, a, l) \Rightarrow  \\ \quad \quad (\operatorname{deriv\text{-}on}(\operatorname{cc\text{-}normed\text{-}field}, u, f, a) = l) \end{array}}}

\vnbfile{theorem-library/holomorphic-basics-2.scm}{theorem-library/holomorphic-basics-2.scm}
\noindent{\small holomorphic-basics-2.scm -- section 2.1 of the user's complex-analysis notes, finished: the constant and the identity are holomorphic on any open set, sums and products of holomorphic functions are holomorphic, and so is z |-$>$ z\textasciicircum{}n for every natural n.}\par
\vnbentry{\texttt{diff-\allowbreak{}on-\allowbreak{}transfer-\allowbreak{}ptwise-\allowbreak{}eq}}{\vnbstmt{\begin{array}{@{}l@{}} \forall k, u, f, g, a, l. \\ \quad (f \in (u \to \operatorname{carr}(k))) \wedge \\ \quad \operatorname{is\text{-}diff\text{-}on}(k, u, g, a, l) \wedge \\ \quad \forall x \in u.\; (f(x) \simeq g(x)) \Rightarrow \\ \quad \quad \operatorname{is\text{-}diff\text{-}on}(k, u, f, a, l) \end{array}}}
\vnbentry{\texttt{holomorphic-\allowbreak{}transfer-\allowbreak{}ptwise-\allowbreak{}eq}}{\vnbstmt{\begin{array}{@{}l@{}} \forall u, f, g. \\ \quad \operatorname{holomorphic\text{-}on}(u, g) \wedge \\ \quad (f \in (u \to \mathbb{C})) \wedge \\ \quad \forall x \in u.\; (f(x) \simeq g(x)) \Rightarrow \\ \quad \quad \operatorname{holomorphic\text{-}on}(u, f) \end{array}}}
\vnbentry{\texttt{holomorphic-\allowbreak{}const}}{\vnbstmt{\begin{array}{@{}l@{}} \forall u.; \forall c \in \mathbb{C}. \\ \quad \operatorname{is\text{-}open}(\operatorname{cc\text{-}ms}, u) \Rightarrow  \\ \quad \quad \operatorname{holomorphic\text{-}on}(u, (\lambda x \in u.\; c)) \end{array}}}
\vnbentry{\texttt{holomorphic-\allowbreak{}identity}}{\vnbstmt{\begin{array}{@{}l@{}} \forall u. \\ \quad \operatorname{is\text{-}open}(\operatorname{cc\text{-}ms}, u) \Rightarrow  \\ \quad \quad \operatorname{holomorphic\text{-}on}(u, (\lambda x \in u.\; x)) \end{array}}}
\vnbentry{\texttt{holomorphic-\allowbreak{}sum}}{\vnbstmt{\begin{array}{@{}l@{}} \forall u, f, g, m. \\ \quad \operatorname{holomorphic\text{-}on}(u, f) \wedge \\ \quad \operatorname{holomorphic\text{-}on}(u, g) \wedge \\ \quad (m \in (u \to \mathbb{C})) \wedge \\ \quad \forall x \in u.\; (m(x) \simeq (f(x) + g(x))) \Rightarrow \\ \quad \quad \operatorname{holomorphic\text{-}on}(u, m) \end{array}}}
\vnbentry{\texttt{holomorphic-\allowbreak{}product}}{\vnbstmt{\begin{array}{@{}l@{}} \forall u, f, g, m. \\ \quad \operatorname{holomorphic\text{-}on}(u, f) \wedge \\ \quad \operatorname{holomorphic\text{-}on}(u, g) \wedge \\ \quad (m \in (u \to \mathbb{C})) \wedge \\ \quad \forall x \in u.\; (m(x) \simeq (f(x) \cdot g(x))) \Rightarrow \\ \quad \quad \operatorname{holomorphic\text{-}on}(u, m) \end{array}}}
\vnbentry{\texttt{holomorphic-\allowbreak{}power}}{\vnbstmt{\begin{array}{@{}l@{}} \forall n \in \mathbb{N}, u.; \forall w \in (u \to \mathbb{C}). \\ \quad \operatorname{is\text{-}open}(\operatorname{cc\text{-}ms}, u) \wedge \\ \quad \forall z \in u.\; (w(z) \simeq \operatorname{power}(z, n)) \Rightarrow \\ \quad \quad \operatorname{holomorphic\text{-}on}(u, w) \end{array}}}

\vnbfile{theorem-library/cc-reciprocal.scm}{theorem-library/cc-reciprocal.scm}
\noindent{\small cc-reciprocal.scm -- z |-$>$ 1/z: the punctured plane, the continuity of the reciprocal, its derivative, and the reciprocal and quotient rules for HOLOMORPHIC-ON.}\par
\vnbentry{\texttt{cc-\allowbreak{}nonzero-\allowbreak{}membership}}{\vnbstmt{\begin{array}{@{}l@{}} \forall w. \\ \quad ((w \in \operatorname{difference}(\mathbb{C}, \operatorname{singleton}(0))) \Leftrightarrow ((w \in \mathbb{C}) \wedge \neg (w = 0))) \end{array}}}
\vnbentry{\texttt{cc-\allowbreak{}nonzero-\allowbreak{}subset}}{\vnbstmt{(\operatorname{difference}(\mathbb{C}, \operatorname{singleton}(0)) \subseteq \mathbb{C})}}
\vnbentry{\texttt{cc-\allowbreak{}nonzero-\allowbreak{}is-\allowbreak{}set}}{\vnbstmt{(\operatorname{difference}(\mathbb{C}, \operatorname{singleton}(0)) \in \mathbf{Set})}}
\vnbentry{\texttt{cc-\allowbreak{}nonzero-\allowbreak{}open}}{\vnbstmt{\operatorname{is\text{-}open}(\operatorname{cc\text{-}ms}, \operatorname{difference}(\mathbb{C}, \operatorname{singleton}(0)))}}
\vnbentry{\texttt{cc-\allowbreak{}nonzero-\allowbreak{}open-\allowbreak{}nf}}{\vnbstmt{\operatorname{is\text{-}open}(\operatorname{nf\text{-}metric\text{-}space}(\operatorname{cc\text{-}normed\text{-}field}), \operatorname{difference}(\mathbb{C}, \operatorname{singleton}(0)))}}
\vnbentry{\texttt{cc-\allowbreak{}recip-\allowbreak{}lam-\allowbreak{}in-\allowbreak{}fun}}{\vnbstmt{\begin{array}{@{}l@{}} (\lambda z \in \operatorname{difference}(\mathbb{C}, \operatorname{singleton}(0)).\; \\ \quad \operatorname{recip}(z)) \\ \in\; (\operatorname{difference}(\mathbb{C}, \operatorname{singleton}(0)) \to \mathbb{C}) \end{array}}}
\vnbentry{\texttt{cc-\allowbreak{}recip-\allowbreak{}continuous-\allowbreak{}at}}{\vnbstmt{\begin{array}{@{}l@{}} \forall a \in \operatorname{difference}(\mathbb{C}, \operatorname{singleton}(0)). \\ \quad \operatorname{is\text{-}continuous\text{-}at}(\operatorname{subspace\text{-}ms}(\operatorname{nf\text{-}metric\text{-}space}(\operatorname{cc\text{-}normed\text{-}field}), \operatorname{difference}(\mathbb{C}, \operatorname{singleton}(0))), \\ \quad \quad \operatorname{nf\text{-}metric\text{-}space}(\operatorname{cc\text{-}normed\text{-}field}), \\ \quad \quad (\lambda z \in \operatorname{difference}(\mathbb{C}, \operatorname{singleton}(0)).\; \\ \quad \quad \quad \operatorname{recip}(z)), \\ \quad \quad a) \end{array}}}
\vnbentry{\texttt{cc-\allowbreak{}mul-\allowbreak{}nonzero}}{\vnbstmt{\begin{array}{@{}l@{}} \forall u, v \in \mathbb{C}. \\ \quad ((\neg (u = 0) \wedge \neg (v = 0)) \Rightarrow \neg ((u \cdot v) = 0)) \end{array}}}
\vnbentry{\texttt{cc-\allowbreak{}recip-\allowbreak{}mul}}{\vnbstmt{\begin{array}{@{}l@{}} \forall u, v \in \mathbb{C}. \\ \quad (\neg (u = 0) \wedge \neg (v = 0)) \Rightarrow  \\ \quad \quad (\operatorname{recip}((u \cdot v)) = (\operatorname{recip}(u) \cdot \operatorname{recip}(v))) \end{array}}}
\vnbentry{\texttt{cc-\allowbreak{}recip-\allowbreak{}diff-\allowbreak{}on}}{\vnbstmt{\begin{array}{@{}l@{}} \forall a \in \operatorname{difference}(\mathbb{C}, \operatorname{singleton}(0)). \\ \quad \operatorname{is\text{-}diff\text{-}on}(\operatorname{cc\text{-}normed\text{-}field}, \operatorname{difference}(\mathbb{C}, \operatorname{singleton}(0)), \\ \quad \quad (\lambda z \in \operatorname{difference}(\mathbb{C}, \operatorname{singleton}(0)).\; \\ \quad \quad \quad \operatorname{recip}(z)), \\ \quad \quad a, -\operatorname{recip}((a \cdot a))) \end{array}}}
\vnbentry{\texttt{holomorphic-\allowbreak{}recip}}{\vnbstmt{\begin{array}{@{}l@{}} \operatorname{holomorphic\text{-}on}(\operatorname{difference}(\mathbb{C}, \operatorname{singleton}(0)), \\ \quad (\lambda z \in \operatorname{difference}(\mathbb{C}, \operatorname{singleton}(0)).\; \\ \quad \quad \operatorname{recip}(z))) \end{array}}}
\vnbentry{\texttt{diff-\allowbreak{}on-\allowbreak{}recip-\allowbreak{}cc}}{\vnbstmt{\begin{array}{@{}l@{}} \forall u, f, a, l, q. \\ \quad \operatorname{is\text{-}diff\text{-}on}(\operatorname{cc\text{-}normed\text{-}field}, u, f, a, l) \wedge \\ \quad \forall y \in u.\; \neg (f(y) = 0) \wedge \\ \quad (q \in (u \to \mathbb{C})) \wedge \\ \quad \forall y \in u.\; (q(y) \simeq \operatorname{recip}(f(y))) \Rightarrow \\ \quad \quad \operatorname{is\text{-}diff\text{-}on}(\operatorname{cc\text{-}normed\text{-}field}, u, q, a, -(l \cdot \operatorname{recip}((f(a) \cdot f(a))))) \end{array}}}
\vnbentry{\texttt{holomorphic-\allowbreak{}recip-\allowbreak{}of}}{\vnbstmt{\begin{array}{@{}l@{}} \forall u, f, q. \\ \quad \operatorname{holomorphic\text{-}on}(u, f) \wedge \\ \quad \forall y \in u.\; \neg (f(y) = 0) \wedge \\ \quad (q \in (u \to \mathbb{C})) \wedge \\ \quad \forall y \in u.\; (q(y) \simeq \operatorname{recip}(f(y))) \Rightarrow \\ \quad \quad \operatorname{holomorphic\text{-}on}(u, q) \end{array}}}
\vnbentry{\texttt{holomorphic-\allowbreak{}quotient}}{\vnbstmt{\begin{array}{@{}l@{}} \forall u, f, g, q. \\ \quad \operatorname{holomorphic\text{-}on}(u, f) \wedge \\ \quad \operatorname{holomorphic\text{-}on}(u, g) \wedge \\ \quad \forall y \in u.\; \neg (g(y) = 0) \wedge \\ \quad (q \in (u \to \mathbb{C})) \wedge \\ \quad \forall y \in u.\; (q(y) \simeq /(f(y), g(y))) \Rightarrow \\ \quad \quad \operatorname{holomorphic\text{-}on}(u, q) \end{array}}}

\vnbfile{theorem-library/interval-calculus-laws.scm}{theorem-library/interval-calculus-laws.scm}
\noindent{\small interval-calculus-laws.scm -- the laws of the vocabulary of structure-library/interval-calculus.scm: OOINT, IS-CONTINUOUS-ON, HAS-DERIV-AT and EXTEND-CONST.  Everything here is PROVEN; the definition file asserts nothing.}\par
\vnbentry{\texttt{ooint-\allowbreak{}membership}}{\vnbstmt{\begin{array}{@{}l@{}} \forall a, b, y. \\ \quad ((y \in \operatorname{ooint}(a, b)) \Leftrightarrow ((y \in \mathbb{R}) \wedge ((a < y) \wedge (y < b)))) \end{array}}}
\vnbentry{\texttt{ooint-\allowbreak{}elt-\allowbreak{}in-\allowbreak{}rr}}{\vnbstmt{\forall a, b.; \forall y \in \operatorname{ooint}(a, b).\; (y \in \mathbb{R})}}
\vnbentry{\texttt{ooint-\allowbreak{}subset-\allowbreak{}rr}}{\vnbstmt{\forall a, b.\; (\operatorname{ooint}(a, b) \subseteq \mathbb{R})}}
\vnbentry{\texttt{ooint-\allowbreak{}subset-\allowbreak{}ccint}}{\vnbstmt{\begin{array}{@{}l@{}} \forall a, b. \\ \quad ((a \in \mathbb{R}) \wedge (b \in \mathbb{R})) \Rightarrow  \\ \quad \quad (\operatorname{ooint}(a, b) \subseteq \operatorname{ccint}(a, b)) \end{array}}}
\vnbentry{\texttt{ooint-\allowbreak{}center}}{\vnbstmt{\begin{array}{@{}l@{}} \forall y, e. \\ \quad ((y \in \mathbb{R}) \wedge \operatorname{pos\text{-}rr}(e)) \Rightarrow  \\ \quad \quad (y \in \operatorname{ooint}((y - e), (y + e))) \end{array}}}
\vnbentry{\texttt{ooint-\allowbreak{}open}}{\vnbstmt{\begin{array}{@{}l@{}} \forall a, b. \\ \quad ((a \in \mathbb{R}) \wedge (b \in \mathbb{R})) \Rightarrow  \\ \quad \quad \operatorname{is\text{-}open}(\operatorname{rr\text{-}ms}, \operatorname{ooint}(a, b)) \end{array}}}
\vnbentry{\texttt{ooint-\allowbreak{}open-\allowbreak{}nf}}{\vnbstmt{\begin{array}{@{}l@{}} \forall a, b. \\ \quad ((a \in \mathbb{R}) \wedge (b \in \mathbb{R})) \Rightarrow  \\ \quad \quad \operatorname{is\text{-}open}(\operatorname{nf\text{-}metric\text{-}space}(\operatorname{rr\text{-}normed\text{-}field}), \operatorname{ooint}(a, b)) \end{array}}}
\vnbentry{\texttt{subspace-\allowbreak{}restrict-\allowbreak{}continuous-\allowbreak{}at}}{\vnbstmt{\begin{array}{@{}l@{}} \forall s, u, v, t, f.; \forall z \in v. \\ \quad \operatorname{is\text{-}metric\text{-}space}(s) \wedge \\ \quad (u \subseteq \operatorname{pts}(s)) \wedge \\ \quad (v \subseteq u) \wedge \\ \quad \operatorname{is\text{-}metric\text{-}space}(t) \wedge \\ \quad \operatorname{is\text{-}continuous\text{-}at}(\operatorname{subspace\text{-}ms}(s, u), t, f, z) \Rightarrow \\ \quad \quad \operatorname{is\text{-}continuous\text{-}at}(\operatorname{subspace\text{-}ms}(s, v), t, \operatorname{restrict}(f, v), z) \end{array}}}
\vnbentry{\texttt{diff-\allowbreak{}on-\allowbreak{}shrink}}{\vnbstmt{\begin{array}{@{}l@{}} \forall k, u, f, a, l, v. \\ \quad \operatorname{is\text{-}diff\text{-}on}(k, u, f, a, l) \wedge \\ \quad \operatorname{is\text{-}open}(\operatorname{nf\text{-}metric\text{-}space}(k), v) \wedge \\ \quad (v \subseteq u) \wedge \\ \quad (a \in v) \Rightarrow \\ \quad \quad \operatorname{is\text{-}diff\text{-}on}(k, v, \operatorname{restrict}(f, v), a, l) \end{array}}}
\vnbentry{\texttt{restrict-\allowbreak{}in-\allowbreak{}fun-\allowbreak{}of-\allowbreak{}restrict}}{\vnbstmt{\begin{array}{@{}l@{}} \forall f, u, c, v. \\ \quad ((\operatorname{restrict}(f, u) \in (u \to c)) \wedge (v \subseteq u)) \Rightarrow  \\ \quad \quad (\operatorname{restrict}(f, v) \in (v \to c)) \end{array}}}
\vnbentry{\texttt{diff-\allowbreak{}on-\allowbreak{}shrink-\allowbreak{}restrict}}{\vnbstmt{\begin{array}{@{}l@{}} \forall k, u, f, a, l, v. \\ \quad \operatorname{is\text{-}diff\text{-}on}(k, u, \operatorname{restrict}(f, u), a, l) \wedge \\ \quad \operatorname{is\text{-}open}(\operatorname{nf\text{-}metric\text{-}space}(k), v) \wedge \\ \quad (v \subseteq u) \wedge \\ \quad (a \in v) \Rightarrow \\ \quad \quad \operatorname{is\text{-}diff\text{-}on}(k, v, \operatorname{restrict}(f, v), a, l) \end{array}}}
\vnbentry{\texttt{diff-\allowbreak{}on-\allowbreak{}ooint-\allowbreak{}shrink}}{\vnbstmt{\begin{array}{@{}l@{}} \forall a \in \mathbb{R}, e, w, f, l. \\ \quad \operatorname{pos\text{-}rr}(e) \wedge \\ \quad \operatorname{pos\text{-}rr}(w) \wedge \\ \quad (w \leq e) \wedge \\ \quad \operatorname{is\text{-}diff\text{-}on}(\operatorname{rr\text{-}normed\text{-}field}, \operatorname{ooint}((a - e), (a + e)), \\ \quad \quad \operatorname{restrict}(f, \operatorname{ooint}((a - e), (a + e))), a, l) \Rightarrow \\ \quad \quad \operatorname{is\text{-}diff\text{-}on}(\operatorname{rr\text{-}normed\text{-}field}, \operatorname{ooint}((a - w), (a + w)), \\ \quad \quad \quad \operatorname{restrict}(f, \operatorname{ooint}((a - w), (a + w))), a, l) \end{array}}}
\vnbentry{\texttt{has-\allowbreak{}deriv-\allowbreak{}at-\allowbreak{}intro}}{\vnbstmt{\begin{array}{@{}l@{}} \forall f, a, l, e. \\ \quad \operatorname{pos\text{-}rr}(e) \wedge \\ \quad \operatorname{is\text{-}diff\text{-}on}(\operatorname{rr\text{-}normed\text{-}field}, \operatorname{ooint}((a - e), (a + e)), \\ \quad \quad \operatorname{restrict}(f, \operatorname{ooint}((a - e), (a + e))), a, l) \Rightarrow \\ \quad \quad \operatorname{has\text{-}deriv\text{-}at}(f, a, l) \end{array}}}
\vnbentry{\texttt{has-\allowbreak{}deriv-\allowbreak{}at-\allowbreak{}shrink}}{\vnbstmt{\begin{array}{@{}l@{}} \forall f, a, l, r. \\ \quad (\operatorname{has\text{-}deriv\text{-}at}(f, a, l) \wedge \operatorname{pos\text{-}rr}(r)) \Rightarrow  \\ \quad \quad \exists w. \\ \quad \quad \quad \operatorname{pos\text{-}rr}(w) \wedge \\ \quad \quad \quad (w \leq r) \wedge \\ \quad \quad \quad \operatorname{is\text{-}diff\text{-}on}(\operatorname{rr\text{-}normed\text{-}field}, \operatorname{ooint}((a - w), (a + w)), \\ \quad \quad \quad \quad \operatorname{restrict}(f, \operatorname{ooint}((a - w), (a + w))), a, l) \end{array}}}
\vnbentry{\texttt{has-\allowbreak{}deriv-\allowbreak{}at-\allowbreak{}pt-\allowbreak{}in-\allowbreak{}rr}}{\vnbstmt{\begin{array}{@{}l@{}} \forall f, a, l. \\ \quad \operatorname{has\text{-}deriv\text{-}at}(f, a, l) \Rightarrow  \\ \quad \quad (a \in \mathbb{R}) \end{array}}}
\vnbentry{\texttt{has-\allowbreak{}deriv-\allowbreak{}at-\allowbreak{}in-\allowbreak{}rr}}{\vnbstmt{\begin{array}{@{}l@{}} \forall f, a, l. \\ \quad \operatorname{has\text{-}deriv\text{-}at}(f, a, l) \Rightarrow  \\ \quad \quad (l \in \mathbb{R}) \end{array}}}
\vnbentry{\texttt{has-\allowbreak{}deriv-\allowbreak{}at-\allowbreak{}unique}}{\vnbstmt{\begin{array}{@{}l@{}} \forall f, a, l, m. \\ \quad \operatorname{has\text{-}deriv\text{-}at}(f, a, l) \wedge \\ \quad \operatorname{has\text{-}deriv\text{-}at}(f, a, m) \Rightarrow \\ \quad \quad (l = m) \end{array}}}
\vnbentry{\texttt{diff-\allowbreak{}on-\allowbreak{}nbhd-\allowbreak{}implies-\allowbreak{}diff-\allowbreak{}at}}{\vnbstmt{\begin{array}{@{}l@{}} \forall f \in (\mathbb{R} \to \mathbb{R}), a \in \mathbb{R}, e, l. \\ \quad \operatorname{pos\text{-}rr}(e) \wedge \\ \quad \operatorname{is\text{-}diff\text{-}on}(\operatorname{rr\text{-}normed\text{-}field}, \operatorname{ooint}((a - e), (a + e)), \\ \quad \quad \operatorname{restrict}(f, \operatorname{ooint}((a - e), (a + e))), a, l) \Rightarrow \\ \quad \quad \operatorname{is\text{-}diff\text{-}at}(f, a, l) \end{array}}}
\vnbentry{\texttt{has-\allowbreak{}deriv-\allowbreak{}at-\allowbreak{}iff-\allowbreak{}diff-\allowbreak{}at}}{\vnbstmt{\begin{array}{@{}l@{}} \forall f \in (\mathbb{R} \to \mathbb{R}), a, l. \\ \quad (\operatorname{has\text{-}deriv\text{-}at}(f, a, l) \Leftrightarrow \operatorname{is\text{-}diff\text{-}at}(f, a, l)) \end{array}}}
\vnbentry{\texttt{extend-\allowbreak{}const-\allowbreak{}value}}{\vnbstmt{\begin{array}{@{}l@{}} \forall f, a, b.; \forall z \in \mathbb{R}. \\ \quad (\operatorname{extend\text{-}const}(f, a, b)(z) \simeq f(\operatorname{clamp}(a, b, z))) \end{array}}}
\vnbentry{\texttt{extend-\allowbreak{}const-\allowbreak{}in-\allowbreak{}fun}}{\vnbstmt{\begin{array}{@{}l@{}} \forall a, b \in \mathbb{R}, f \in (\operatorname{ccint}(a, b) \to \mathbb{R}). \\ \quad (a \leq b) \Rightarrow  \\ \quad \quad (\operatorname{extend\text{-}const}(f, a, b) \in (\mathbb{R} \to \mathbb{R})) \end{array}}}
\vnbentry{\texttt{extend-\allowbreak{}const-\allowbreak{}fixes}}{\vnbstmt{\begin{array}{@{}l@{}} \forall a, b, z \in \mathbb{R}, f. \\ \quad ((a \leq z) \wedge (z \leq b)) \Rightarrow  \\ \quad \quad (\operatorname{extend\text{-}const}(f, a, b)(z) \simeq f(z)) \end{array}}}
\vnbentry{\texttt{extend-\allowbreak{}const-\allowbreak{}ccint-\allowbreak{}fixes}}{\vnbstmt{\begin{array}{@{}l@{}} \forall a, b \in \mathbb{R}, z \in \operatorname{ccint}(a, b), f. \\ \quad (\operatorname{extend\text{-}const}(f, a, b)(z) \simeq f(z)) \end{array}}}
\vnbentry{\texttt{extend-\allowbreak{}const-\allowbreak{}continuous-\allowbreak{}at}}{\vnbstmt{\begin{array}{@{}l@{}} \forall a, b \in \mathbb{R}, f.; \forall t \in \mathbb{R}. \\ \quad ((a \leq b) \wedge \operatorname{is\text{-}continuous\text{-}on}(f, \operatorname{ccint}(a, b))) \Rightarrow  \\ \quad \quad \operatorname{is\text{-}continuous\text{-}at}(\operatorname{rr\text{-}ms}, \operatorname{rr\text{-}ms}, \operatorname{extend\text{-}const}(f, a, b), t) \end{array}}}
\vnbentry{\texttt{extend-\allowbreak{}const-\allowbreak{}deriv-\allowbreak{}fwd}}{\vnbstmt{\begin{array}{@{}l@{}} \forall a, b \in \mathbb{R}, f \in (\operatorname{ccint}(a, b) \to \mathbb{R}), x \in \operatorname{ooint}(a, b), l. \\ \quad ((a \leq b) \wedge \operatorname{has\text{-}deriv\text{-}at}(f, x, l)) \Rightarrow  \\ \quad \quad \operatorname{is\text{-}diff\text{-}at}(\operatorname{extend\text{-}const}(f, a, b), x, l) \end{array}}}
\vnbentry{\texttt{extend-\allowbreak{}const-\allowbreak{}deriv-\allowbreak{}bwd}}{\vnbstmt{\begin{array}{@{}l@{}} \forall a, b \in \mathbb{R}, f \in (\operatorname{ccint}(a, b) \to \mathbb{R}), x \in \operatorname{ooint}(a, b), l. \\ \quad (a \leq b) \wedge \\ \quad \operatorname{is\text{-}diff\text{-}at}(\operatorname{extend\text{-}const}(f, a, b), x, l) \Rightarrow \\ \quad \quad \operatorname{has\text{-}deriv\text{-}at}(f, x, l) \end{array}}}
\vnbentry{\texttt{extend-\allowbreak{}const-\allowbreak{}deriv-\allowbreak{}iff}}{\vnbstmt{\begin{array}{@{}l@{}} \forall a, b \in \mathbb{R}, f \in (\operatorname{ccint}(a, b) \to \mathbb{R}), x \in \operatorname{ooint}(a, b), l. \\ \quad (a \leq b) \Rightarrow  \\ \quad \quad (\operatorname{has\text{-}deriv\text{-}at}(f, x, l) \Leftrightarrow \operatorname{is\text{-}diff\text{-}at}(\operatorname{extend\text{-}const}(f, a, b), x, l)) \end{array}}}
\vnbentry{\texttt{has-\allowbreak{}deriv-\allowbreak{}at-\allowbreak{}local}}{\vnbstmt{\begin{array}{@{}l@{}} \forall r, f, g, x, l. \\ \quad \operatorname{pos\text{-}rr}(r) \wedge \\ \quad \operatorname{restrict}(f, \operatorname{ooint}((x - r), (x + r))) \\ \quad \in\; (\operatorname{ooint}((x - r), (x + r)) \to \mathbb{R}) \wedge \\ \quad \forall \mathit{hbx} \in \operatorname{ooint}((x - r), (x + r)). \\ \quad \quad (f(\mathit{hbx}) \simeq g(\mathit{hbx})) \wedge \\ \quad \operatorname{has\text{-}deriv\text{-}at}(g, x, l) \Rightarrow \\ \quad \quad \operatorname{has\text{-}deriv\text{-}at}(f, x, l) \end{array}}}
\vnbentry{\texttt{has-\allowbreak{}deriv-\allowbreak{}at-\allowbreak{}sum}}{\vnbstmt{\begin{array}{@{}l@{}} \forall r, f, g, h, x, l, m. \\ \quad \operatorname{pos\text{-}rr}(r) \wedge \\ \quad \operatorname{restrict}(h, \operatorname{ooint}((x - r), (x + r))) \\ \quad \in\; (\operatorname{ooint}((x - r), (x + r)) \to \mathbb{R}) \wedge \\ \quad \forall \mathit{hbx} \in \operatorname{ooint}((x - r), (x + r)). \\ \quad \quad (h(\mathit{hbx}) \simeq (f(\mathit{hbx}) + g(\mathit{hbx}))) \wedge \\ \quad \operatorname{has\text{-}deriv\text{-}at}(f, x, l) \wedge \\ \quad \operatorname{has\text{-}deriv\text{-}at}(g, x, m) \Rightarrow \\ \quad \quad \operatorname{has\text{-}deriv\text{-}at}(h, x, (l + m)) \end{array}}}
\vnbentry{\texttt{has-\allowbreak{}deriv-\allowbreak{}at-\allowbreak{}product}}{\vnbstmt{\begin{array}{@{}l@{}} \forall r, f, g, h, x, l, m. \\ \quad \operatorname{pos\text{-}rr}(r) \wedge \\ \quad \operatorname{restrict}(h, \operatorname{ooint}((x - r), (x + r))) \\ \quad \in\; (\operatorname{ooint}((x - r), (x + r)) \to \mathbb{R}) \wedge \\ \quad \forall \mathit{hbx} \in \operatorname{ooint}((x - r), (x + r)). \\ \quad \quad (h(\mathit{hbx}) \simeq (f(\mathit{hbx}) \cdot g(\mathit{hbx}))) \wedge \\ \quad \operatorname{has\text{-}deriv\text{-}at}(f, x, l) \wedge \\ \quad \operatorname{has\text{-}deriv\text{-}at}(g, x, m) \Rightarrow \\ \quad \quad \operatorname{has\text{-}deriv\text{-}at}(h, x, ((l \cdot g(x)) + (m \cdot f(x)))) \end{array}}}

\vnbfile{theorem-library/interval-mvt.scm}{theorem-library/interval-mvt.scm}
\noindent{\small interval-mvt.scm -- ROLLE, THE TWO MEAN VALUE THEOREMS AND THE THREE BOUNDS, STATED FOR A FUNCTION ON ITS INTERVAL.}\par
\vnbentry{\texttt{rolle-\allowbreak{}on-\allowbreak{}interval}}{\vnbstmt{\begin{array}{@{}l@{}} \forall a, b \in \mathbb{R}, h \in (\operatorname{ccint}(a, b) \to \mathbb{R}). \\ \quad (a < b) \wedge \\ \quad \operatorname{is\text{-}continuous\text{-}on}(h, \operatorname{ccint}(a, b)) \wedge \\ \quad \forall x \in \operatorname{ooint}(a, b). \exists l. \\ \quad \quad \operatorname{has\text{-}deriv\text{-}at}(h, x, l) \wedge \\ \quad (h(a) = h(b)) \Rightarrow \\ \quad \quad \exists c \in \operatorname{ooint}(a, b). \\ \quad \quad \quad \operatorname{has\text{-}deriv\text{-}at}(h, c, 0) \end{array}}}
\vnbentry{\texttt{generalized-\allowbreak{}mvt-\allowbreak{}on-\allowbreak{}interval}}{\vnbstmt{\begin{array}{@{}l@{}} \forall a, b \in \mathbb{R}, f, g. \\ \quad (a < b) \wedge \\ \quad (f \in (\operatorname{ccint}(a, b) \to \mathbb{R})) \wedge \\ \quad (g \in (\operatorname{ccint}(a, b) \to \mathbb{R})) \wedge \\ \quad \operatorname{is\text{-}continuous\text{-}on}(f, \operatorname{ccint}(a, b)) \wedge \\ \quad \operatorname{is\text{-}continuous\text{-}on}(g, \operatorname{ccint}(a, b)) \wedge \\ \quad \forall x \in \operatorname{ooint}(a, b). \exists l. \\ \quad \quad \operatorname{has\text{-}deriv\text{-}at}(f, x, l) \wedge \\ \quad \forall x \in \operatorname{ooint}(a, b). \exists m. \\ \quad \quad \operatorname{has\text{-}deriv\text{-}at}(g, x, m) \Rightarrow \\ \quad \quad \exists c \in \operatorname{ooint}(a, b), l, m. \\ \quad \quad \quad \operatorname{has\text{-}deriv\text{-}at}(f, c, l) \wedge \\ \quad \quad \quad \operatorname{has\text{-}deriv\text{-}at}(g, c, m) \wedge \\ \quad \quad \quad ((l \cdot (g(b) - g(a))) = (m \cdot (f(b) - f(a)))) \end{array}}}
\vnbentry{\texttt{mvt-\allowbreak{}on-\allowbreak{}interval}}{\vnbstmt{\begin{array}{@{}l@{}} \forall a, b \in \mathbb{R}, f \in (\operatorname{ccint}(a, b) \to \mathbb{R}). \\ \quad (a < b) \wedge \\ \quad \operatorname{is\text{-}continuous\text{-}on}(f, \operatorname{ccint}(a, b)) \wedge \\ \quad \forall x \in \operatorname{ooint}(a, b). \exists l. \\ \quad \quad \operatorname{has\text{-}deriv\text{-}at}(f, x, l) \Rightarrow \\ \quad \quad \exists c \in \operatorname{ooint}(a, b), l. \\ \quad \quad \quad \operatorname{has\text{-}deriv\text{-}at}(f, c, l) \wedge \\ \quad \quad \quad (l = ((f(b) - f(a)) \cdot \operatorname{recip}((b - a)))) \end{array}}}
\vnbentry{\texttt{mvt-\allowbreak{}upper-\allowbreak{}bound-\allowbreak{}on-\allowbreak{}interval}}{\vnbstmt{\begin{array}{@{}l@{}} \forall a, b, m \in \mathbb{R}, f \in (\operatorname{ccint}(a, b) \to \mathbb{R}). \\ \quad (a < b) \wedge \\ \quad \operatorname{is\text{-}continuous\text{-}on}(f, \operatorname{ccint}(a, b)) \wedge \\ \quad \forall x \in \operatorname{ooint}(a, b). \exists l. \\ \quad \quad (\operatorname{has\text{-}deriv\text{-}at}(f, x, l) \wedge (l \leq m)) \Rightarrow \\ \quad \quad ((f(b) - f(a)) \leq (m \cdot (b - a))) \end{array}}}
\vnbentry{\texttt{mvt-\allowbreak{}lower-\allowbreak{}bound-\allowbreak{}on-\allowbreak{}interval}}{\vnbstmt{\begin{array}{@{}l@{}} \forall a, b, m \in \mathbb{R}, f \in (\operatorname{ccint}(a, b) \to \mathbb{R}). \\ \quad (a < b) \wedge \\ \quad \operatorname{is\text{-}continuous\text{-}on}(f, \operatorname{ccint}(a, b)) \wedge \\ \quad \forall x \in \operatorname{ooint}(a, b). \exists l. \\ \quad \quad (\operatorname{has\text{-}deriv\text{-}at}(f, x, l) \wedge (m \leq l)) \Rightarrow \\ \quad \quad ((m \cdot (b - a)) \leq (f(b) - f(a))) \end{array}}}
\vnbentry{\texttt{mvt-\allowbreak{}abs-\allowbreak{}bound-\allowbreak{}on-\allowbreak{}interval}}{\vnbstmt{\begin{array}{@{}l@{}} \forall a, b, m \in \mathbb{R}, f \in (\operatorname{ccint}(a, b) \to \mathbb{R}). \\ \quad (a < b) \wedge \\ \quad \operatorname{is\text{-}continuous\text{-}on}(f, \operatorname{ccint}(a, b)) \wedge \\ \quad \forall x \in \operatorname{ooint}(a, b). \exists l. \\ \quad \quad (\operatorname{has\text{-}deriv\text{-}at}(f, x, l) \wedge (\lvert l \rvert \leq m)) \Rightarrow \\ \quad \quad (\lvert (f(b) - f(a)) \rvert \leq (m \cdot (b - a))) \end{array}}}

\vnbfile{theorem-library/interval-extremum.scm}{theorem-library/interval-extremum.scm}
\noindent{\small interval-extremum.scm -- PROPOSITION 2.10 FOR A FUNCTION ON ITS INTERVAL.}\par
\vnbentry{\texttt{interval-\allowbreak{}local-\allowbreak{}max-\allowbreak{}deriv-\allowbreak{}zero}}{\vnbstmt{\begin{array}{@{}l@{}} \forall a, b \in \mathbb{R}, f.; \forall c \in \operatorname{ooint}(a, b), l. \\ \quad \operatorname{is\text{-}local\text{-}max\text{-}at}(f, \operatorname{ccint}(a, b), c) \wedge \\ \quad \operatorname{has\text{-}deriv\text{-}at}(f, c, l) \Rightarrow \\ \quad \quad (l = 0) \end{array}}}
\vnbentry{\texttt{interval-\allowbreak{}local-\allowbreak{}min-\allowbreak{}deriv-\allowbreak{}zero}}{\vnbstmt{\begin{array}{@{}l@{}} \forall a, b \in \mathbb{R}, f.; \forall c \in \operatorname{ooint}(a, b), l. \\ \quad \operatorname{is\text{-}local\text{-}min\text{-}at}(f, \operatorname{ccint}(a, b), c) \wedge \\ \quad \operatorname{has\text{-}deriv\text{-}at}(f, c, l) \Rightarrow \\ \quad \quad (l = 0) \end{array}}}

\vnbfile{theorem-library/one-sided-derivative-laws.scm}{theorem-library/one-sided-derivative-laws.scm}
\noindent{\small one-sided-derivative-laws.scm -- the laws of structure-library/one-sided-derivative.scm: COINT, OCINT, HAS-DERIV-AT-WITHIN, HAS-RIGHT-DERIV-AT and HAS-LEFT-DERIV-AT.  Everything here is PROVEN; the definition file asserts nothing.}\par
\vnbentry{\texttt{coint-\allowbreak{}membership}}{\vnbstmt{\begin{array}{@{}l@{}} \forall a, b, y. \\ \quad ((y \in \operatorname{coint}(a, b)) \Leftrightarrow ((y \in \mathbb{R}) \wedge ((a \leq y) \wedge (y < b)))) \end{array}}}
\vnbentry{\texttt{ocint-\allowbreak{}membership}}{\vnbstmt{\begin{array}{@{}l@{}} \forall a, b, y. \\ \quad ((y \in \operatorname{ocint}(a, b)) \Leftrightarrow ((y \in \mathbb{R}) \wedge ((a < y) \wedge (y \leq b)))) \end{array}}}
\vnbentry{\texttt{coint-\allowbreak{}elt-\allowbreak{}in-\allowbreak{}rr}}{\vnbstmt{\forall a, b.; \forall y \in \operatorname{coint}(a, b).\; (y \in \mathbb{R})}}
\vnbentry{\texttt{ocint-\allowbreak{}elt-\allowbreak{}in-\allowbreak{}rr}}{\vnbstmt{\forall a, b.; \forall y \in \operatorname{ocint}(a, b).\; (y \in \mathbb{R})}}
\vnbentry{\texttt{coint-\allowbreak{}subset-\allowbreak{}rr}}{\vnbstmt{\forall a, b.\; (\operatorname{coint}(a, b) \subseteq \mathbb{R})}}
\vnbentry{\texttt{ocint-\allowbreak{}subset-\allowbreak{}rr}}{\vnbstmt{\forall a, b.\; (\operatorname{ocint}(a, b) \subseteq \mathbb{R})}}
\vnbentry{\texttt{coint-\allowbreak{}left-\allowbreak{}endpoint}}{\vnbstmt{\begin{array}{@{}l@{}} \forall x, r. \\ \quad ((x \in \mathbb{R}) \wedge \operatorname{pos\text{-}rr}(r)) \Rightarrow  \\ \quad \quad (x \in \operatorname{coint}(x, (x + r))) \end{array}}}
\vnbentry{\texttt{ocint-\allowbreak{}right-\allowbreak{}endpoint}}{\vnbstmt{\begin{array}{@{}l@{}} \forall x, r. \\ \quad ((x \in \mathbb{R}) \wedge \operatorname{pos\text{-}rr}(r)) \Rightarrow  \\ \quad \quad (x \in \operatorname{ocint}((x - r), x)) \end{array}}}
\vnbentry{\texttt{coint-\allowbreak{}subset-\allowbreak{}ooint}}{\vnbstmt{\begin{array}{@{}l@{}} \forall x, r. \\ \quad ((x \in \mathbb{R}) \wedge \operatorname{pos\text{-}rr}(r)) \Rightarrow  \\ \quad \quad (\operatorname{coint}(x, (x + r)) \subseteq \operatorname{ooint}((x - r), (x + r))) \end{array}}}
\vnbentry{\texttt{ocint-\allowbreak{}subset-\allowbreak{}ooint}}{\vnbstmt{\begin{array}{@{}l@{}} \forall x, r. \\ \quad ((x \in \mathbb{R}) \wedge \operatorname{pos\text{-}rr}(r)) \Rightarrow  \\ \quad \quad (\operatorname{ocint}((x - r), x) \subseteq \operatorname{ooint}((x - r), (x + r))) \end{array}}}
\vnbentry{\texttt{coint-\allowbreak{}shrink}}{\vnbstmt{\begin{array}{@{}l@{}} \forall x, r, v. \\ \quad (x \in \mathbb{R}) \wedge \\ \quad \operatorname{pos\text{-}rr}(r) \wedge \\ \quad \operatorname{pos\text{-}rr}(v) \wedge \\ \quad (v \leq r) \Rightarrow \\ \quad \quad (\operatorname{coint}(x, (x + v)) \subseteq \operatorname{coint}(x, (x + r))) \end{array}}}
\vnbentry{\texttt{ocint-\allowbreak{}shrink}}{\vnbstmt{\begin{array}{@{}l@{}} \forall x, r, v. \\ \quad (x \in \mathbb{R}) \wedge \\ \quad \operatorname{pos\text{-}rr}(r) \wedge \\ \quad \operatorname{pos\text{-}rr}(v) \wedge \\ \quad (v \leq r) \Rightarrow \\ \quad \quad (\operatorname{ocint}((x - v), x) \subseteq \operatorname{ocint}((x - r), x)) \end{array}}}
\vnbentry{\texttt{coint-\allowbreak{}left-\allowbreak{}limit-\allowbreak{}point}}{\vnbstmt{\begin{array}{@{}l@{}} \forall x, d, r. \\ \quad ((x \in \mathbb{R}) \wedge (\operatorname{pos\text{-}rr}(d) \wedge \operatorname{pos\text{-}rr}(r))) \Rightarrow  \\ \quad \quad \exists b \in \operatorname{coint}(x, (x + d)). \\ \quad \quad \quad (\neg (b = x) \wedge (\lvert (x - b) \rvert \leq r)) \end{array}}}
\vnbentry{\texttt{ocint-\allowbreak{}right-\allowbreak{}limit-\allowbreak{}point}}{\vnbstmt{\begin{array}{@{}l@{}} \forall x, d, r. \\ \quad ((x \in \mathbb{R}) \wedge (\operatorname{pos\text{-}rr}(d) \wedge \operatorname{pos\text{-}rr}(r))) \Rightarrow  \\ \quad \quad \exists b \in \operatorname{ocint}((x - d), x). \\ \quad \quad \quad (\neg (b = x) \wedge (\lvert (x - b) \rvert \leq r)) \end{array}}}
\vnbentry{\texttt{restrict-\allowbreak{}fun-\allowbreak{}domain-\allowbreak{}subset}}{\vnbstmt{\begin{array}{@{}l@{}} \forall f, d, c, w, b. \\ \quad ((f \in (d \to c)) \wedge (\operatorname{restrict}(f, w) \in (w \to b))) \Rightarrow  \\ \quad \quad (w \subseteq d) \end{array}}}
\vnbentry{\texttt{deriv-\allowbreak{}within-\allowbreak{}in-\allowbreak{}rr}}{\vnbstmt{\begin{array}{@{}l@{}} \forall f, w, x, l. \\ \quad \operatorname{has\text{-}deriv\text{-}at\text{-}within}(f, w, x, l) \Rightarrow  \\ \quad \quad (l \in \mathbb{R}) \end{array}}}
\vnbentry{\texttt{deriv-\allowbreak{}within-\allowbreak{}pt-\allowbreak{}in}}{\vnbstmt{\begin{array}{@{}l@{}} \forall f, w, x, l. \\ \quad \operatorname{has\text{-}deriv\text{-}at\text{-}within}(f, w, x, l) \Rightarrow  \\ \quad \quad (x \in w) \end{array}}}
\vnbentry{\texttt{deriv-\allowbreak{}within-\allowbreak{}restrict-\allowbreak{}in-\allowbreak{}fun}}{\vnbstmt{\begin{array}{@{}l@{}} \forall f, w, x, l. \\ \quad \operatorname{has\text{-}deriv\text{-}at\text{-}within}(f, w, x, l) \Rightarrow  \\ \quad \quad (\operatorname{restrict}(f, w) \in (w \to \mathbb{R})) \end{array}}}
\vnbentry{\texttt{deriv-\allowbreak{}within-\allowbreak{}shrink}}{\vnbstmt{\begin{array}{@{}l@{}} \forall f, w, v, x, l. \\ \quad (w \subseteq \mathbb{R}) \wedge \\ \quad \operatorname{has\text{-}deriv\text{-}at\text{-}within}(f, w, x, l) \wedge \\ \quad (v \subseteq w) \wedge \\ \quad (x \in v) \Rightarrow \\ \quad \quad \operatorname{has\text{-}deriv\text{-}at\text{-}within}(f, v, x, l) \end{array}}}
\vnbentry{\texttt{deriv-\allowbreak{}within-\allowbreak{}local}}{\vnbstmt{\begin{array}{@{}l@{}} \forall f, g, w, x, l. \\ \quad (\operatorname{restrict}(f, w) \in (w \to \mathbb{R})) \wedge \\ \quad \forall y \in w.\; (f(y) \simeq g(y)) \wedge \\ \quad \operatorname{has\text{-}deriv\text{-}at\text{-}within}(g, w, x, l) \Rightarrow \\ \quad \quad \operatorname{has\text{-}deriv\text{-}at\text{-}within}(f, w, x, l) \end{array}}}
\vnbentry{\texttt{deriv-\allowbreak{}within-\allowbreak{}unique}}{\vnbstmt{\begin{array}{@{}l@{}} \forall f, w, x, l, m. \\ \quad (w \subseteq \mathbb{R}) \wedge \\ \quad \forall r. \\ \quad \quad \operatorname{pos\text{-}rr}(r) \Rightarrow  \\ \quad \quad \quad \exists b \in w.\; (\neg (b = x) \wedge (\lvert (x - b) \rvert \leq r)) \wedge \\ \quad \operatorname{has\text{-}deriv\text{-}at\text{-}within}(f, w, x, l) \wedge \\ \quad \operatorname{has\text{-}deriv\text{-}at\text{-}within}(f, w, x, m) \Rightarrow \\ \quad \quad (l = m) \end{array}}}
\vnbentry{\texttt{rr-\allowbreak{}mul-\allowbreak{}neg-\allowbreak{}neg-\allowbreak{}pos}}{\vnbstmt{\begin{array}{@{}l@{}} \forall y, b. \\ \quad ((y \in \mathbb{R}) \wedge ((b \in \mathbb{R}) \wedge ((y < 0) \wedge (b < 0)))) \Rightarrow  \\ \quad \quad (0 < (y \cdot b)) \end{array}}}
\vnbentry{\texttt{has-\allowbreak{}deriv-\allowbreak{}at-\allowbreak{}of-\allowbreak{}within-\allowbreak{}ooint}}{\vnbstmt{\begin{array}{@{}l@{}} \forall x, r, f, l. \\ \quad \operatorname{pos\text{-}rr}(r) \wedge \\ \quad \operatorname{has\text{-}deriv\text{-}at\text{-}within}(f, \operatorname{ooint}((x - r), (x + r)), x, l) \Rightarrow \\ \quad \quad \operatorname{has\text{-}deriv\text{-}at}(f, x, l) \end{array}}}
\vnbentry{\texttt{has-\allowbreak{}deriv-\allowbreak{}at-\allowbreak{}to-\allowbreak{}within-\allowbreak{}ooint}}{\vnbstmt{\begin{array}{@{}l@{}} \forall f, x, l, r. \\ \quad (\operatorname{has\text{-}deriv\text{-}at}(f, x, l) \wedge \operatorname{pos\text{-}rr}(r)) \Rightarrow  \\ \quad \quad \exists v. \\ \quad \quad \quad \operatorname{pos\text{-}rr}(v) \wedge \\ \quad \quad \quad (v \leq r) \wedge \\ \quad \quad \quad \operatorname{has\text{-}deriv\text{-}at\text{-}within}(f, \operatorname{ooint}((x - v), (x + v)), x, l) \end{array}}}
\vnbentry{\texttt{has-\allowbreak{}deriv-\allowbreak{}at-\allowbreak{}iff-\allowbreak{}within-\allowbreak{}ooint}}{\vnbstmt{\begin{array}{@{}l@{}} \forall f, x, l. \\ \quad (\operatorname{has\text{-}deriv\text{-}at}(f, x, l) \Leftrightarrow \exists v.\; (\operatorname{pos\text{-}rr}(v) \wedge \operatorname{has\text{-}deriv\text{-}at\text{-}within}(f, \operatorname{ooint}((x - v), (x + v)), x, l))) \end{array}}}
\vnbentry{\texttt{has-\allowbreak{}right-\allowbreak{}deriv-\allowbreak{}at-\allowbreak{}intro}}{\vnbstmt{\begin{array}{@{}l@{}} \forall f, x, l, r. \\ \quad \operatorname{pos\text{-}rr}(r) \wedge \\ \quad \operatorname{has\text{-}deriv\text{-}at\text{-}within}(f, \operatorname{coint}(x, (x + r)), x, l) \Rightarrow \\ \quad \quad \operatorname{has\text{-}right\text{-}deriv\text{-}at}(f, x, l) \end{array}}}
\vnbentry{\texttt{has-\allowbreak{}left-\allowbreak{}deriv-\allowbreak{}at-\allowbreak{}intro}}{\vnbstmt{\begin{array}{@{}l@{}} \forall f, x, l, r. \\ \quad \operatorname{pos\text{-}rr}(r) \wedge \\ \quad \operatorname{has\text{-}deriv\text{-}at\text{-}within}(f, \operatorname{ocint}((x - r), x), x, l) \Rightarrow \\ \quad \quad \operatorname{has\text{-}left\text{-}deriv\text{-}at}(f, x, l) \end{array}}}
\vnbentry{\texttt{has-\allowbreak{}right-\allowbreak{}deriv-\allowbreak{}at-\allowbreak{}pt-\allowbreak{}in-\allowbreak{}rr}}{\vnbstmt{\begin{array}{@{}l@{}} \forall f, x, l. \\ \quad \operatorname{has\text{-}right\text{-}deriv\text{-}at}(f, x, l) \Rightarrow  \\ \quad \quad (x \in \mathbb{R}) \end{array}}}
\vnbentry{\texttt{has-\allowbreak{}right-\allowbreak{}deriv-\allowbreak{}at-\allowbreak{}in-\allowbreak{}rr}}{\vnbstmt{\begin{array}{@{}l@{}} \forall f, x, l. \\ \quad \operatorname{has\text{-}right\text{-}deriv\text{-}at}(f, x, l) \Rightarrow  \\ \quad \quad (l \in \mathbb{R}) \end{array}}}
\vnbentry{\texttt{has-\allowbreak{}left-\allowbreak{}deriv-\allowbreak{}at-\allowbreak{}pt-\allowbreak{}in-\allowbreak{}rr}}{\vnbstmt{\begin{array}{@{}l@{}} \forall f, x, l. \\ \quad \operatorname{has\text{-}left\text{-}deriv\text{-}at}(f, x, l) \Rightarrow  \\ \quad \quad (x \in \mathbb{R}) \end{array}}}
\vnbentry{\texttt{has-\allowbreak{}left-\allowbreak{}deriv-\allowbreak{}at-\allowbreak{}in-\allowbreak{}rr}}{\vnbstmt{\begin{array}{@{}l@{}} \forall f, x, l. \\ \quad \operatorname{has\text{-}left\text{-}deriv\text{-}at}(f, x, l) \Rightarrow  \\ \quad \quad (l \in \mathbb{R}) \end{array}}}
\vnbentry{\texttt{has-\allowbreak{}right-\allowbreak{}deriv-\allowbreak{}at-\allowbreak{}shrink}}{\vnbstmt{\begin{array}{@{}l@{}} \forall f, x, l, r. \\ \quad \operatorname{has\text{-}right\text{-}deriv\text{-}at}(f, x, l) \wedge \\ \quad \operatorname{pos\text{-}rr}(r) \Rightarrow \\ \quad \quad \exists v. \\ \quad \quad \quad \operatorname{pos\text{-}rr}(v) \wedge \\ \quad \quad \quad (v \leq r) \wedge \\ \quad \quad \quad \operatorname{has\text{-}deriv\text{-}at\text{-}within}(f, \operatorname{coint}(x, (x + v)), x, l) \end{array}}}
\vnbentry{\texttt{has-\allowbreak{}left-\allowbreak{}deriv-\allowbreak{}at-\allowbreak{}shrink}}{\vnbstmt{\begin{array}{@{}l@{}} \forall f, x, l, r. \\ \quad \operatorname{has\text{-}left\text{-}deriv\text{-}at}(f, x, l) \wedge \\ \quad \operatorname{pos\text{-}rr}(r) \Rightarrow \\ \quad \quad \exists v. \\ \quad \quad \quad \operatorname{pos\text{-}rr}(v) \wedge \\ \quad \quad \quad (v \leq r) \wedge \\ \quad \quad \quad \operatorname{has\text{-}deriv\text{-}at\text{-}within}(f, \operatorname{ocint}((x - v), x), x, l) \end{array}}}
\vnbentry{\texttt{has-\allowbreak{}right-\allowbreak{}deriv-\allowbreak{}at-\allowbreak{}unique}}{\vnbstmt{\begin{array}{@{}l@{}} \forall f, x, l, m. \\ \quad \operatorname{has\text{-}right\text{-}deriv\text{-}at}(f, x, l) \wedge \\ \quad \operatorname{has\text{-}right\text{-}deriv\text{-}at}(f, x, m) \Rightarrow \\ \quad \quad (l = m) \end{array}}}
\vnbentry{\texttt{has-\allowbreak{}left-\allowbreak{}deriv-\allowbreak{}at-\allowbreak{}unique}}{\vnbstmt{\begin{array}{@{}l@{}} \forall f, x, l, m. \\ \quad \operatorname{has\text{-}left\text{-}deriv\text{-}at}(f, x, l) \wedge \\ \quad \operatorname{has\text{-}left\text{-}deriv\text{-}at}(f, x, m) \Rightarrow \\ \quad \quad (l = m) \end{array}}}
\vnbentry{\texttt{has-\allowbreak{}right-\allowbreak{}deriv-\allowbreak{}at-\allowbreak{}local}}{\vnbstmt{\begin{array}{@{}l@{}} \forall r, f, g, x, l. \\ \quad \operatorname{pos\text{-}rr}(r) \wedge \\ \quad \operatorname{restrict}(f, \operatorname{coint}(x, (x + r))) \\ \quad \in\; (\operatorname{coint}(x, (x + r)) \to \mathbb{R}) \wedge \\ \quad \forall y \in \operatorname{coint}(x, (x + r)).\; (f(y) \simeq g(y)) \wedge \\ \quad \operatorname{has\text{-}right\text{-}deriv\text{-}at}(g, x, l) \Rightarrow \\ \quad \quad \operatorname{has\text{-}right\text{-}deriv\text{-}at}(f, x, l) \end{array}}}
\vnbentry{\texttt{has-\allowbreak{}left-\allowbreak{}deriv-\allowbreak{}at-\allowbreak{}local}}{\vnbstmt{\begin{array}{@{}l@{}} \forall r, f, g, x, l. \\ \quad \operatorname{pos\text{-}rr}(r) \wedge \\ \quad \operatorname{restrict}(f, \operatorname{ocint}((x - r), x)) \\ \quad \in\; (\operatorname{ocint}((x - r), x) \to \mathbb{R}) \wedge \\ \quad \forall y \in \operatorname{ocint}((x - r), x).\; (f(y) \simeq g(y)) \wedge \\ \quad \operatorname{has\text{-}left\text{-}deriv\text{-}at}(g, x, l) \Rightarrow \\ \quad \quad \operatorname{has\text{-}left\text{-}deriv\text{-}at}(f, x, l) \end{array}}}
\vnbentry{\texttt{has-\allowbreak{}deriv-\allowbreak{}at-\allowbreak{}implies-\allowbreak{}right}}{\vnbstmt{\begin{array}{@{}l@{}} \forall f, x, l. \\ \quad \operatorname{has\text{-}deriv\text{-}at}(f, x, l) \Rightarrow  \\ \quad \quad \operatorname{has\text{-}right\text{-}deriv\text{-}at}(f, x, l) \end{array}}}
\vnbentry{\texttt{has-\allowbreak{}deriv-\allowbreak{}at-\allowbreak{}implies-\allowbreak{}left}}{\vnbstmt{\begin{array}{@{}l@{}} \forall f, x, l. \\ \quad \operatorname{has\text{-}deriv\text{-}at}(f, x, l) \Rightarrow  \\ \quad \quad \operatorname{has\text{-}left\text{-}deriv\text{-}at}(f, x, l) \end{array}}}
\vnbentry{\texttt{local-\allowbreak{}max-\allowbreak{}right-\allowbreak{}deriv-\allowbreak{}nonpos}}{\vnbstmt{\begin{array}{@{}l@{}} \forall f, d, x, l. \\ \quad \operatorname{is\text{-}local\text{-}max\text{-}at}(f, d, x) \wedge \\ \quad \operatorname{has\text{-}right\text{-}deriv\text{-}at}(f, x, l) \Rightarrow \\ \quad \quad (l \leq 0) \end{array}}}
\vnbentry{\texttt{local-\allowbreak{}max-\allowbreak{}left-\allowbreak{}deriv-\allowbreak{}nonneg}}{\vnbstmt{\begin{array}{@{}l@{}} \forall f, d, x, l. \\ \quad \operatorname{is\text{-}local\text{-}max\text{-}at}(f, d, x) \wedge \\ \quad \operatorname{has\text{-}left\text{-}deriv\text{-}at}(f, x, l) \Rightarrow \\ \quad \quad (0 \leq l) \end{array}}}
\vnbentry{\texttt{interval-\allowbreak{}local-\allowbreak{}max-\allowbreak{}right-\allowbreak{}deriv-\allowbreak{}nonpos}}{\vnbstmt{\begin{array}{@{}l@{}} \forall a, b \in \mathbb{R}, f, x, l. \\ \quad \operatorname{is\text{-}local\text{-}max\text{-}at}(f, \operatorname{ccint}(a, b), x) \wedge \\ \quad \operatorname{has\text{-}right\text{-}deriv\text{-}at}(f, x, l) \Rightarrow \\ \quad \quad (l \leq 0) \end{array}}}
\vnbentry{\texttt{interval-\allowbreak{}local-\allowbreak{}max-\allowbreak{}left-\allowbreak{}deriv-\allowbreak{}nonneg}}{\vnbstmt{\begin{array}{@{}l@{}} \forall a, b \in \mathbb{R}, f, x, l. \\ \quad \operatorname{is\text{-}local\text{-}max\text{-}at}(f, \operatorname{ccint}(a, b), x) \wedge \\ \quad \operatorname{has\text{-}left\text{-}deriv\text{-}at}(f, x, l) \Rightarrow \\ \quad \quad (0 \leq l) \end{array}}}
\vnbentry{\texttt{rr-\allowbreak{}max-\allowbreak{}left-\allowbreak{}le}}{\vnbstmt{\begin{array}{@{}l@{}} \forall y, b. \\ \quad ((y \in \mathbb{R}) \wedge ((b \in \mathbb{R}) \wedge (b \leq y))) \Rightarrow  \\ \quad \quad (\operatorname{max}(y, b) = y) \end{array}}}
\vnbentry{\texttt{rr-\allowbreak{}min-\allowbreak{}left-\allowbreak{}le}}{\vnbstmt{\begin{array}{@{}l@{}} \forall y, b. \\ \quad ((y \in \mathbb{R}) \wedge ((b \in \mathbb{R}) \wedge (y \leq b))) \Rightarrow  \\ \quad \quad (\operatorname{min}(y, b) = y) \end{array}}}
\vnbentry{\texttt{rr-\allowbreak{}max-\allowbreak{}right-\allowbreak{}le}}{\vnbstmt{\begin{array}{@{}l@{}} \forall y, b. \\ \quad ((y \in \mathbb{R}) \wedge ((b \in \mathbb{R}) \wedge (y \leq b))) \Rightarrow  \\ \quad \quad (\operatorname{max}(y, b) = b) \end{array}}}
\vnbentry{\texttt{rr-\allowbreak{}min-\allowbreak{}right-\allowbreak{}le}}{\vnbstmt{\begin{array}{@{}l@{}} \forall y, b. \\ \quad ((y \in \mathbb{R}) \wedge ((b \in \mathbb{R}) \wedge (b \leq y))) \Rightarrow  \\ \quad \quad (\operatorname{min}(y, b) = b) \end{array}}}
\vnbentry{\texttt{rr-\allowbreak{}max-\allowbreak{}lt}}{\vnbstmt{\begin{array}{@{}l@{}} \forall y, b, c. \\ \quad (y \in \mathbb{R}) \wedge \\ \quad (b \in \mathbb{R}) \wedge \\ \quad (c \in \mathbb{R}) \wedge \\ \quad (y < c) \wedge \\ \quad (b < c) \Rightarrow \\ \quad \quad (\operatorname{max}(y, b) < c) \end{array}}}
\vnbentry{\texttt{rr-\allowbreak{}lt-\allowbreak{}min}}{\vnbstmt{\begin{array}{@{}l@{}} \forall c, y, b. \\ \quad (c \in \mathbb{R}) \wedge \\ \quad (y \in \mathbb{R}) \wedge \\ \quad (b \in \mathbb{R}) \wedge \\ \quad (c < y) \wedge \\ \quad (c < b) \Rightarrow \\ \quad \quad (c < \operatorname{min}(y, b)) \end{array}}}
\vnbentry{\texttt{rr-\allowbreak{}max-\allowbreak{}retract-\allowbreak{}close}}{\vnbstmt{\begin{array}{@{}l@{}} \forall y, c. \\ \quad ((y \in \mathbb{R}) \wedge (c \in \mathbb{R})) \Rightarrow  \\ \quad \quad (\lvert (c - \operatorname{max}(y, c)) \rvert \leq \lvert (c - y) \rvert) \end{array}}}
\vnbentry{\texttt{rr-\allowbreak{}min-\allowbreak{}retract-\allowbreak{}close}}{\vnbstmt{\begin{array}{@{}l@{}} \forall y, c. \\ \quad ((y \in \mathbb{R}) \wedge (c \in \mathbb{R})) \Rightarrow  \\ \quad \quad (\lvert (c - \operatorname{min}(y, c)) \rvert \leq \lvert (c - y) \rvert) \end{array}}}
\vnbentry{\texttt{coint-\allowbreak{}of-\allowbreak{}ooint}}{\vnbstmt{\begin{array}{@{}l@{}} \forall x, r.; \forall y \in \operatorname{ooint}((x - r), (x + r)). \\ \quad ((x \in \mathbb{R}) \wedge (\operatorname{pos\text{-}rr}(r) \wedge (x \leq y))) \Rightarrow  \\ \quad \quad (y \in \operatorname{coint}(x, (x + r))) \end{array}}}
\vnbentry{\texttt{ocint-\allowbreak{}of-\allowbreak{}ooint}}{\vnbstmt{\begin{array}{@{}l@{}} \forall x, r.; \forall y \in \operatorname{ooint}((x - r), (x + r)). \\ \quad ((x \in \mathbb{R}) \wedge (\operatorname{pos\text{-}rr}(r) \wedge (y \leq x))) \Rightarrow  \\ \quad \quad (y \in \operatorname{ocint}((x - r), x)) \end{array}}}
\vnbentry{\texttt{has-\allowbreak{}deriv-\allowbreak{}at-\allowbreak{}of-\allowbreak{}one-\allowbreak{}sided}}{\vnbstmt{\begin{array}{@{}l@{}} \forall f, x, l. \\ \quad \operatorname{has\text{-}right\text{-}deriv\text{-}at}(f, x, l) \wedge \\ \quad \operatorname{has\text{-}left\text{-}deriv\text{-}at}(f, x, l) \Rightarrow \\ \quad \quad \operatorname{has\text{-}deriv\text{-}at}(f, x, l) \end{array}}}
\vnbentry{\texttt{has-\allowbreak{}deriv-\allowbreak{}at-\allowbreak{}iff-\allowbreak{}one-\allowbreak{}sided}}{\vnbstmt{\begin{array}{@{}l@{}} \forall f, x, l. \\ \quad (\operatorname{has\text{-}deriv\text{-}at}(f, x, l) \Leftrightarrow (\operatorname{has\text{-}right\text{-}deriv\text{-}at}(f, x, l) \wedge \operatorname{has\text{-}left\text{-}deriv\text{-}at}(f, x, l))) \end{array}}}
\vnbentry{\texttt{deriv-\allowbreak{}within-\allowbreak{}sum}}{\vnbstmt{\begin{array}{@{}l@{}} \forall f, g, h, w, x, l, m. \\ \quad (w \subseteq \mathbb{R}) \wedge \\ \quad (\operatorname{restrict}(h, w) \in (w \to \mathbb{R})) \wedge \\ \quad \forall y \in w.\; (h(y) \simeq (f(y) + g(y))) \wedge \\ \quad \operatorname{has\text{-}deriv\text{-}at\text{-}within}(f, w, x, l) \wedge \\ \quad \operatorname{has\text{-}deriv\text{-}at\text{-}within}(g, w, x, m) \Rightarrow \\ \quad \quad \operatorname{has\text{-}deriv\text{-}at\text{-}within}(h, w, x, (l + m)) \end{array}}}
\vnbentry{\texttt{has-\allowbreak{}right-\allowbreak{}deriv-\allowbreak{}at-\allowbreak{}sum}}{\vnbstmt{\begin{array}{@{}l@{}} \forall r, f, g, h, x, l, m. \\ \quad \operatorname{pos\text{-}rr}(r) \wedge \\ \quad \operatorname{restrict}(h, \operatorname{coint}(x, (x + r))) \\ \quad \in\; (\operatorname{coint}(x, (x + r)) \to \mathbb{R}) \wedge \\ \quad \forall y \in \operatorname{coint}(x, (x + r)).\; (h(y) \simeq (f(y) + g(y))) \wedge \\ \quad \operatorname{has\text{-}right\text{-}deriv\text{-}at}(f, x, l) \wedge \\ \quad \operatorname{has\text{-}right\text{-}deriv\text{-}at}(g, x, m) \Rightarrow \\ \quad \quad \operatorname{has\text{-}right\text{-}deriv\text{-}at}(h, x, (l + m)) \end{array}}}
\vnbentry{\texttt{has-\allowbreak{}left-\allowbreak{}deriv-\allowbreak{}at-\allowbreak{}sum}}{\vnbstmt{\begin{array}{@{}l@{}} \forall r, f, g, h, x, l, m. \\ \quad \operatorname{pos\text{-}rr}(r) \wedge \\ \quad \operatorname{restrict}(h, \operatorname{ocint}((x - r), x)) \\ \quad \in\; (\operatorname{ocint}((x - r), x) \to \mathbb{R}) \wedge \\ \quad \forall y \in \operatorname{ocint}((x - r), x).\; (h(y) \simeq (f(y) + g(y))) \wedge \\ \quad \operatorname{has\text{-}left\text{-}deriv\text{-}at}(f, x, l) \wedge \\ \quad \operatorname{has\text{-}left\text{-}deriv\text{-}at}(g, x, m) \Rightarrow \\ \quad \quad \operatorname{has\text{-}left\text{-}deriv\text{-}at}(h, x, (l + m)) \end{array}}}

\vnbfile{theorem-library/directional-derivative.scm}{theorem-library/directional-derivative.scm}
\noindent{\small directional-derivative.scm -- calculus.pdf Chapter 2 section 8 "Multivariable Calculus", Def 2.23 through Prop 2.27, over an arbitrary real normed vector space.}\par
\vnbentry{\texttt{seg-\allowbreak{}curve-\allowbreak{}apply}}{\vnbstmt{\begin{array}{@{}l@{}} \forall m, f, a, b.; \forall t \in \mathbb{R}. \\ \quad (\operatorname{seg\text{-}curve}(m, f, a, b)(t) \simeq f(\operatorname{vadd}(m)(a, \operatorname{act}(m)(t, b)))) \end{array}}}
\vnbentry{\texttt{diff-\allowbreak{}at-\allowbreak{}in-\allowbreak{}fun}}{\vnbstmt{\begin{array}{@{}l@{}} \forall f, a, l. \\ \quad \operatorname{is\text{-}diff\text{-}at}(f, a, l) \Rightarrow  \\ \quad \quad (f \in (\mathbb{R} \to \mathbb{R})) \end{array}}}
\vnbentry{\texttt{affine-\allowbreak{}lam-\allowbreak{}in-\allowbreak{}fun}}{\vnbstmt{\begin{array}{@{}l@{}} \forall c, m. \\ \quad ((c \in \mathbb{R}) \wedge (m \in \mathbb{R})) \Rightarrow  \\ \quad \quad ((\lambda x \in \mathbb{R}.\; (c + (m \cdot x))) \in (\mathbb{R} \to \mathbb{R})) \end{array}}}
\vnbentry{\texttt{deriv-\allowbreak{}affine}}{\vnbstmt{\begin{array}{@{}l@{}} \forall c, m, a. \\ \quad ((c \in \mathbb{R}) \wedge ((m \in \mathbb{R}) \wedge (a \in \mathbb{R}))) \Rightarrow  \\ \quad \quad \operatorname{is\text{-}diff\text{-}at}((\lambda x \in \mathbb{R}.\; (c + (m \cdot x))), a, m) \end{array}}}
\vnbentry{\texttt{dir-\allowbreak{}diff-\allowbreak{}reparam}}{\vnbstmt{\begin{array}{@{}l@{}} \forall f, g, c, m, l. \\ \quad (f \in (\mathbb{R} \to \mathbb{R})) \wedge \\ \quad (c \in \mathbb{R}) \wedge \\ \quad (m \in \mathbb{R}) \wedge \\ \quad \operatorname{is\text{-}diff\text{-}at}(g, c, l) \wedge \\ \quad \forall x \in \mathbb{R}.\; (f(x) \simeq g((c + (m \cdot x)))) \Rightarrow \\ \quad \quad \operatorname{is\text{-}diff\text{-}at}(f, 0, (l \cdot m)) \end{array}}}
\vnbentry{\texttt{dir-\allowbreak{}diff-\allowbreak{}scale}}{\vnbstmt{\begin{array}{@{}l@{}} \forall f, g, m, l. \\ \quad (f \in (\mathbb{R} \to \mathbb{R})) \wedge \\ \quad (m \in \mathbb{R}) \wedge \\ \quad \operatorname{is\text{-}diff\text{-}at}(g, 0, l) \wedge \\ \quad \forall x \in \mathbb{R}.\; (f(x) \simeq g((m \cdot x))) \Rightarrow \\ \quad \quad \operatorname{is\text{-}diff\text{-}at}(f, 0, (m \cdot l)) \end{array}}}
\vnbentry{\texttt{dir-\allowbreak{}diff-\allowbreak{}shift}}{\vnbstmt{\begin{array}{@{}l@{}} \forall f, g, c, l. \\ \quad (f \in (\mathbb{R} \to \mathbb{R})) \wedge \\ \quad (c \in \mathbb{R}) \wedge \\ \quad \operatorname{is\text{-}diff\text{-}at}(g, c, l) \wedge \\ \quad \forall x \in \mathbb{R}.\; (f(x) \simeq g((c + x))) \Rightarrow \\ \quad \quad \operatorname{is\text{-}diff\text{-}at}(f, 0, l) \end{array}}}
\vnbentry{\texttt{seg-\allowbreak{}curve-\allowbreak{}in-\allowbreak{}fun}}{\vnbstmt{\begin{array}{@{}l@{}} \forall m, f, a, b. \\ \quad \operatorname{is\text{-}normed\text{-}vector\text{-}space}(m) \wedge \\ \quad (f \in (\operatorname{vec}(m) \to \mathbb{R})) \wedge \\ \quad (a \in \operatorname{vec}(m)) \wedge \\ \quad (b \in \operatorname{vec}(m)) \Rightarrow \\ \quad \quad (\operatorname{seg\text{-}curve}(m, f, a, b) \in (\mathbb{R} \to \mathbb{R})) \end{array}}}
\vnbentry{\texttt{dir-\allowbreak{}deriv-\allowbreak{}homogeneous}}{\vnbstmt{\begin{array}{@{}l@{}} \forall m, f, a, b, c, l. \\ \quad \operatorname{is\text{-}normed\text{-}vector\text{-}space}(m) \wedge \\ \quad (f \in (\operatorname{vec}(m) \to \mathbb{R})) \wedge \\ \quad (a \in \operatorname{vec}(m)) \wedge \\ \quad (b \in \operatorname{vec}(m)) \wedge \\ \quad (c \in \mathbb{R}) \wedge \\ \quad \operatorname{is\text{-}dir\text{-}diff\text{-}at}(m, f, a, b, l) \Rightarrow \\ \quad \quad \operatorname{is\text{-}dir\text{-}diff\text{-}at}(m, f, a, \operatorname{act}(m)(c, b), (c \cdot l)) \end{array}}}
\vnbentry{\texttt{dir-\allowbreak{}deriv-\allowbreak{}mvt}}{\vnbstmt{\begin{array}{@{}l@{}} \forall m, f, a, b. \\ \quad \operatorname{is\text{-}normed\text{-}vector\text{-}space}(m) \wedge \\ \quad (f \in (\operatorname{vec}(m) \to \mathbb{R})) \wedge \\ \quad (a \in \operatorname{vec}(m)) \wedge \\ \quad (b \in \operatorname{vec}(m)) \wedge \\ \quad \forall x \in \operatorname{ccint}(0, 1). \\ \quad \quad \operatorname{is\text{-}continuous\text{-}at}(\operatorname{rr\text{-}ms}, \operatorname{rr\text{-}ms}, \operatorname{seg\text{-}curve}(m, f, a, b), x) \wedge \\ \quad \forall x \in \mathbb{R}. \\ \quad \quad ((0 < x) \wedge (x < 1)) \Rightarrow  \\ \quad \quad \quad \exists l. \\ \quad \quad \quad \quad \operatorname{is\text{-}diff\text{-}at}(\operatorname{seg\text{-}curve}(m, f, a, b), x, l) \Rightarrow \\ \quad \quad \exists h \in \mathbb{R}. \\ \quad \quad \quad (0 < h) \wedge \\ \quad \quad \quad (h < 1) \wedge \\ \quad \quad \quad \exists v. \\ \quad \quad \quad \quad \operatorname{is\text{-}dir\text{-}diff\text{-}at}(m, f, \operatorname{vadd}(m)(a, \operatorname{act}(m)(h, b)), b, v) \wedge \\ \quad \quad \quad \quad (v = (f(\operatorname{vadd}(m)(a, b)) - f(a))) \end{array}}}
\vnbentry{\texttt{dir-\allowbreak{}deriv-\allowbreak{}value}}{\vnbstmt{\begin{array}{@{}l@{}} \forall m, f, a, b, l. \\ \quad \operatorname{is\text{-}dir\text{-}diff\text{-}at}(m, f, a, b, l) \Rightarrow  \\ \quad \quad (\operatorname{dir\text{-}deriv}(m, f, a, b) = l) \end{array}}}
\vnbentry{\texttt{diff-\allowbreak{}at-\allowbreak{}shift-\allowbreak{}back}}{\vnbstmt{\begin{array}{@{}l@{}} \forall f, g, c, l. \\ \quad (f \in (\mathbb{R} \to \mathbb{R})) \wedge \\ \quad (c \in \mathbb{R}) \wedge \\ \quad \operatorname{is\text{-}diff\text{-}at}(g, 0, l) \wedge \\ \quad \forall x \in \mathbb{R}.\; (g(x) \simeq f((c + x))) \Rightarrow \\ \quad \quad \operatorname{is\text{-}diff\text{-}at}(f, c, l) \end{array}}}
\vnbentry{\texttt{seg-\allowbreak{}curve-\allowbreak{}diff-\allowbreak{}at}}{\vnbstmt{\begin{array}{@{}l@{}} \forall m.; \forall f \in (\operatorname{vec}(m) \to \mathbb{R}), b, a \in \operatorname{vec}(m), s \in \mathbb{R}, k. \\ \quad \operatorname{is\text{-}normed\text{-}vector\text{-}space}(m) \wedge \\ \quad \operatorname{is\text{-}dir\text{-}diff\text{-}at}(m, f, \operatorname{vadd}(m)(b, \operatorname{act}(m)(s, a)), a, k) \Rightarrow \\ \quad \quad \operatorname{is\text{-}diff\text{-}at}(\operatorname{seg\text{-}curve}(m, f, b, a), s, k) \end{array}}}
\vnbentry{\texttt{dir-\allowbreak{}deriv-\allowbreak{}increment}}{\vnbstmt{\begin{array}{@{}l@{}} \forall m.; \forall f \in (\operatorname{vec}(m) \to \mathbb{R}), b, a \in \operatorname{vec}(m). \\ \quad \operatorname{is\text{-}normed\text{-}vector\text{-}space}(m) \wedge \\ \quad \forall x \in \operatorname{vec}(m). \exists k. \\ \quad \quad \operatorname{is\text{-}dir\text{-}diff\text{-}at}(m, f, x, a, k) \Rightarrow \\ \quad \quad \exists h \in \mathbb{R}. \\ \quad \quad \quad (0 < h) \wedge \\ \quad \quad \quad (h < 1) \wedge \\ \quad \quad \quad (f(\operatorname{vadd}(m)(b, a)) - f(b)) \\ \quad \quad \quad =\; \operatorname{dir\text{-}deriv}(m, f, \operatorname{vadd}(m)(b, \operatorname{act}(m)(h, a)), a) \end{array}}}
\vnbentry{\texttt{dir-\allowbreak{}diff-\allowbreak{}value-\allowbreak{}in-\allowbreak{}rr}}{\vnbstmt{\begin{array}{@{}l@{}} \forall m, f, a, b, l. \\ \quad \operatorname{is\text{-}dir\text{-}diff\text{-}at}(m, f, a, b, l) \Rightarrow  \\ \quad \quad (l \in \mathbb{R}) \end{array}}}
\vnbentry{\texttt{dir-\allowbreak{}deriv-\allowbreak{}increment-\allowbreak{}scaled}}{\vnbstmt{\begin{array}{@{}l@{}} \forall m.; \forall f \in (\operatorname{vec}(m) \to \mathbb{R}), b, a \in \operatorname{vec}(m), t \in \mathbb{R}. \\ \quad \operatorname{is\text{-}normed\text{-}vector\text{-}space}(m) \wedge \\ \quad \forall x \in \operatorname{vec}(m). \exists k. \\ \quad \quad \operatorname{is\text{-}dir\text{-}diff\text{-}at}(m, f, x, a, k) \Rightarrow \\ \quad \quad \exists h \in \mathbb{R}. \\ \quad \quad \quad (0 < h) \wedge \\ \quad \quad \quad (h < 1) \wedge \\ \quad \quad \quad (f(\operatorname{vadd}(m)(b, \operatorname{act}(m)(t, a))) - f(b)) \\ \quad \quad \quad =\; (t \cdot \operatorname{dir\text{-}deriv}(m, f, \operatorname{vadd}(m)(b, \operatorname{act}(m)((h \cdot t), a)), a)) \end{array}}}
\vnbentry{\texttt{rr-\allowbreak{}abs-\allowbreak{}unit-\allowbreak{}scale}}{\vnbstmt{\begin{array}{@{}l@{}} \forall s, t \in \mathbb{R}. \\ \quad ((0 < s) \wedge (s < 1)) \Rightarrow  \\ \quad \quad (\lvert (s \cdot t) \rvert \leq \lvert t \rvert) \end{array}}}
\vnbentry{\texttt{nvs-\allowbreak{}norm-\allowbreak{}two-\allowbreak{}term-\allowbreak{}bound}}{\vnbstmt{\begin{array}{@{}l@{}} \forall m.; \forall a, b \in \operatorname{vec}(m), c, d, f \in \mathbb{R}. \\ \quad \operatorname{is\text{-}normed\text{-}vector\text{-}space}(m) \wedge \\ \quad (0 \leq f) \wedge \\ \quad (\lvert c \rvert \leq f) \wedge \\ \quad (\lvert d \rvert \leq f) \Rightarrow \\ \quad \quad \operatorname{vnrm}(m)(\operatorname{vadd}(m)(\operatorname{act}(m)(c, a), \operatorname{act}(m)(d, b))) \\ \quad \quad \leq\; (f \cdot (1 + (\operatorname{vnrm}(m)(b) + \operatorname{vnrm}(m)(a)))) \end{array}}}
\vnbentry{\texttt{nvs-\allowbreak{}norm-\allowbreak{}one-\allowbreak{}term-\allowbreak{}bound}}{\vnbstmt{\begin{array}{@{}l@{}} \forall m.; \forall a, b \in \operatorname{vec}(m), c, d \in \mathbb{R}. \\ \quad \operatorname{is\text{-}normed\text{-}vector\text{-}space}(m) \wedge \\ \quad (0 \leq d) \wedge \\ \quad (\lvert c \rvert \leq d) \Rightarrow \\ \quad \quad \operatorname{vnrm}(m)(\operatorname{act}(m)(c, a)) \\ \quad \quad \leq\; (d \cdot (1 + (\operatorname{vnrm}(m)(b) + \operatorname{vnrm}(m)(a)))) \end{array}}}
\vnbentry{\texttt{rr-\allowbreak{}diff-\allowbreak{}split-\allowbreak{}mid}}{\vnbstmt{\forall x, y, z \in \mathbb{R}.\; ((x - z) = ((x - y) + (y - z)))}}
\vnbentry{\texttt{dir-\allowbreak{}deriv-\allowbreak{}corner-\allowbreak{}split}}{\vnbstmt{\begin{array}{@{}l@{}} \forall m.; \forall f \in (\operatorname{vec}(m) \to \mathbb{R}), a, b, i \in \operatorname{vec}(m), t \in \mathbb{R}. \\ \quad \operatorname{is\text{-}normed\text{-}vector\text{-}space}(m) \wedge \\ \quad \forall x \in \operatorname{vec}(m). \exists k. \\ \quad \quad \operatorname{is\text{-}dir\text{-}diff\text{-}at}(m, f, x, b, k) \wedge \\ \quad \forall x \in \operatorname{vec}(m). \exists k. \\ \quad \quad \operatorname{is\text{-}dir\text{-}diff\text{-}at}(m, f, x, i, k) \Rightarrow \\ \quad \quad \exists c, d \in \operatorname{vec}(m). \\ \quad \quad \quad \operatorname{dist}(\operatorname{nvs\text{-}metric\text{-}space}(m))(a, c) \\ \quad \quad \quad \leq\; (\lvert t \rvert \cdot (1 + (\operatorname{vnrm}(m)(b) + \operatorname{vnrm}(m)(i)))) \wedge \\ \quad \quad \quad \operatorname{dist}(\operatorname{nvs\text{-}metric\text{-}space}(m))(a, d) \\ \quad \quad \quad \leq\; (\lvert t \rvert \cdot (1 + (\operatorname{vnrm}(m)(b) + \operatorname{vnrm}(m)(i)))) \wedge \\ \quad \quad \quad (\operatorname{seg\text{-}curve}(m, f, a, \operatorname{vadd}(m)(b, i))(t) - \operatorname{seg\text{-}curve}(m, f, a, \operatorname{vadd}(m)(b, i))(0)) \\ \quad \quad \quad =\; (t \cdot (\operatorname{dir\text{-}deriv}(m, f, c, b) + \operatorname{dir\text{-}deriv}(m, f, d, i))) \end{array}}}
\vnbentry{\texttt{rr-\allowbreak{}factor-\allowbreak{}shift}}{\vnbstmt{\begin{array}{@{}l@{}} \forall d, u, r \in \mathbb{R}. \\ \quad (((d \cdot r) \cdot (u - 0)) = (d \cdot (u \cdot r))) \end{array}}}
\vnbentry{\texttt{dir-\allowbreak{}quotient-\allowbreak{}in-\allowbreak{}fun}}{\vnbstmt{\begin{array}{@{}l@{}} \forall m.; \forall f \in (\operatorname{vec}(m) \to \mathbb{R}), a, b, i \in \operatorname{vec}(m), l, c \in \mathbb{R}. \\ \quad \operatorname{is\text{-}normed\text{-}vector\text{-}space}(m) \Rightarrow  \\ \quad \quad (\lambda z \in \mathbb{R}.\; \\ \quad \quad \quad \left(\text{if }(z = 0)\text{ then }(l + c)\text{ else }((\operatorname{seg\text{-}curve}(m, f, a, \operatorname{vadd}(m)(b, i))(z) - \operatorname{seg\text{-}curve}(m, f, a, \operatorname{vadd}(m)(b, i))(0)) \cdot \operatorname{recip}(z))\right)) \\ \quad \quad \in\; (\mathbb{R} \to \mathbb{R}) \end{array}}}
\vnbentry{\texttt{dir-\allowbreak{}quotient-\allowbreak{}at-\allowbreak{}zero}}{\vnbstmt{\begin{array}{@{}l@{}} \forall m.; \forall f \in (\operatorname{vec}(m) \to \mathbb{R}), a, b, i \in \operatorname{vec}(m), l, c \in \mathbb{R}. \\ \quad \operatorname{is\text{-}normed\text{-}vector\text{-}space}(m) \Rightarrow  \\ \quad \quad (\lambda z \in \mathbb{R}.\; \left(\text{if }(z = 0)\text{ then }(l + c)\text{ else }((\operatorname{seg\text{-}curve}(m, f, a, \operatorname{vadd}(m)(b, i))(z) - \operatorname{seg\text{-}curve}(m, f, a, \operatorname{vadd}(m)(b, i))(0)) \cdot \operatorname{recip}(z))\right))(0) \\ \quad \quad =\; (l + c) \end{array}}}
\vnbentry{\texttt{dir-\allowbreak{}quotient-\allowbreak{}identity}}{\vnbstmt{\begin{array}{@{}l@{}} \forall m.; \forall f \in (\operatorname{vec}(m) \to \mathbb{R}), a, b, i \in \operatorname{vec}(m), l, c, x \in \mathbb{R}. \\ \quad \operatorname{is\text{-}normed\text{-}vector\text{-}space}(m) \Rightarrow  \\ \quad \quad (\operatorname{seg\text{-}curve}(m, f, a, \operatorname{vadd}(m)(b, i))(x) - \operatorname{seg\text{-}curve}(m, f, a, \operatorname{vadd}(m)(b, i))(0)) \\ \quad \quad =\; ((\lambda z \in \mathbb{R}.\; \left(\text{if }(z = 0)\text{ then }(l + c)\text{ else }((\operatorname{seg\text{-}curve}(m, f, a, \operatorname{vadd}(m)(b, i))(z) - \operatorname{seg\text{-}curve}(m, f, a, \operatorname{vadd}(m)(b, i))(0)) \cdot \operatorname{recip}(z))\right))(x) \cdot (x - 0)) \end{array}}}
\vnbentry{\texttt{rr-\allowbreak{}mul-\allowbreak{}rotate}}{\vnbstmt{\begin{array}{@{}l@{}} \forall s, u, r \in \mathbb{R}. \\ \quad (((u \cdot s) \cdot r) = (s \cdot (u \cdot r))) \end{array}}}
\vnbentry{\texttt{dir-\allowbreak{}quotient-\allowbreak{}continuous}}{\vnbstmt{\begin{array}{@{}l@{}} \forall m.; \forall f \in (\operatorname{vec}(m) \to \mathbb{R}), a, b, i \in \operatorname{vec}(m), l, c. \\ \quad \operatorname{is\text{-}normed\text{-}vector\text{-}space}(m) \wedge \\ \quad \operatorname{is\text{-}dir\text{-}diff\text{-}at}(m, f, a, b, l) \wedge \\ \quad \operatorname{is\text{-}dir\text{-}diff\text{-}at}(m, f, a, i, c) \wedge \\ \quad \forall x \in \operatorname{vec}(m). \exists k. \\ \quad \quad \operatorname{is\text{-}dir\text{-}diff\text{-}at}(m, f, x, b, k) \wedge \\ \quad \forall x \in \operatorname{vec}(m). \exists k. \\ \quad \quad \operatorname{is\text{-}dir\text{-}diff\text{-}at}(m, f, x, i, k) \wedge \\ \quad \operatorname{is\text{-}continuous\text{-}at}(\operatorname{nvs\text{-}metric\text{-}space}(m), \operatorname{rr\text{-}ms}, \\ \quad \quad (\lambda x \in \operatorname{vec}(m).\; \operatorname{dir\text{-}deriv}(m, f, x, b)), a) \wedge \\ \quad \operatorname{is\text{-}continuous\text{-}at}(\operatorname{nvs\text{-}metric\text{-}space}(m), \operatorname{rr\text{-}ms}, \\ \quad \quad (\lambda x \in \operatorname{vec}(m).\; \operatorname{dir\text{-}deriv}(m, f, x, i)), a) \Rightarrow \\ \quad \quad \operatorname{is\text{-}continuous\text{-}at}(\operatorname{rr\text{-}ms}, \operatorname{rr\text{-}ms}, \\ \quad \quad \quad (\lambda z \in \mathbb{R}.\; \\ \quad \quad \quad \quad \left(\text{if }(z = 0)\text{ then }(l + c)\text{ else }((\operatorname{seg\text{-}curve}(m, f, a, \operatorname{vadd}(m)(b, i))(z) - \operatorname{seg\text{-}curve}(m, f, a, \operatorname{vadd}(m)(b, i))(0)) \cdot \operatorname{recip}(z))\right)), \\ \quad \quad \quad 0) \end{array}}}
\vnbentry{\texttt{dir-\allowbreak{}deriv-\allowbreak{}additive}}{\vnbstmt{\begin{array}{@{}l@{}} \forall m, f, a, b, i, l, c. \\ \quad \operatorname{is\text{-}normed\text{-}vector\text{-}space}(m) \wedge \\ \quad (f \in (\operatorname{vec}(m) \to \mathbb{R})) \wedge \\ \quad (a \in \operatorname{vec}(m)) \wedge \\ \quad (b \in \operatorname{vec}(m)) \wedge \\ \quad (i \in \operatorname{vec}(m)) \wedge \\ \quad \operatorname{is\text{-}dir\text{-}diff\text{-}at}(m, f, a, b, l) \wedge \\ \quad \operatorname{is\text{-}dir\text{-}diff\text{-}at}(m, f, a, i, c) \wedge \\ \quad \forall x \in \operatorname{vec}(m). \exists k. \\ \quad \quad \operatorname{is\text{-}dir\text{-}diff\text{-}at}(m, f, x, b, k) \wedge \\ \quad \forall x \in \operatorname{vec}(m). \exists k. \\ \quad \quad \operatorname{is\text{-}dir\text{-}diff\text{-}at}(m, f, x, i, k) \wedge \\ \quad \operatorname{is\text{-}continuous\text{-}at}(\operatorname{nvs\text{-}metric\text{-}space}(m), \operatorname{rr\text{-}ms}, \\ \quad \quad (\lambda x \in \operatorname{vec}(m).\; \operatorname{dir\text{-}deriv}(m, f, x, b)), a) \wedge \\ \quad \operatorname{is\text{-}continuous\text{-}at}(\operatorname{nvs\text{-}metric\text{-}space}(m), \operatorname{rr\text{-}ms}, \\ \quad \quad (\lambda x \in \operatorname{vec}(m).\; \operatorname{dir\text{-}deriv}(m, f, x, i)), a) \Rightarrow \\ \quad \quad \operatorname{is\text{-}dir\text{-}diff\text{-}at}(m, f, a, \operatorname{vadd}(m)(b, i), (l + c)) \end{array}}}
\vnbentry{\texttt{dir-\allowbreak{}deriv-\allowbreak{}linear}}{\vnbstmt{\begin{array}{@{}l@{}} \forall m, f, a, b, i, r, s, l, c. \\ \quad \operatorname{is\text{-}normed\text{-}vector\text{-}space}(m) \wedge \\ \quad (f \in (\operatorname{vec}(m) \to \mathbb{R})) \wedge \\ \quad (a \in \operatorname{vec}(m)) \wedge \\ \quad (b \in \operatorname{vec}(m)) \wedge \\ \quad (i \in \operatorname{vec}(m)) \wedge \\ \quad (r \in \mathbb{R}) \wedge \\ \quad (s \in \mathbb{R}) \wedge \\ \quad \operatorname{is\text{-}dir\text{-}diff\text{-}at}(m, f, a, b, l) \wedge \\ \quad \operatorname{is\text{-}dir\text{-}diff\text{-}at}(m, f, a, i, c) \wedge \\ \quad \forall x \in \operatorname{vec}(m). \exists k. \\ \quad \quad \operatorname{is\text{-}dir\text{-}diff\text{-}at}(m, f, x, \operatorname{act}(m)(r, b), k) \wedge \\ \quad \forall x \in \operatorname{vec}(m). \exists k. \\ \quad \quad \operatorname{is\text{-}dir\text{-}diff\text{-}at}(m, f, x, \operatorname{act}(m)(s, i), k) \wedge \\ \quad \operatorname{is\text{-}continuous\text{-}at}(\operatorname{nvs\text{-}metric\text{-}space}(m), \operatorname{rr\text{-}ms}, \\ \quad \quad (\lambda x \in \operatorname{vec}(m).\; \\ \quad \quad \quad \operatorname{dir\text{-}deriv}(m, f, x, \operatorname{act}(m)(r, b))), \\ \quad \quad a) \wedge \\ \quad \operatorname{is\text{-}continuous\text{-}at}(\operatorname{nvs\text{-}metric\text{-}space}(m), \operatorname{rr\text{-}ms}, \\ \quad \quad (\lambda x \in \operatorname{vec}(m).\; \\ \quad \quad \quad \operatorname{dir\text{-}deriv}(m, f, x, \operatorname{act}(m)(s, i))), \\ \quad \quad a) \Rightarrow \\ \quad \quad \operatorname{is\text{-}dir\text{-}diff\text{-}at}(m, f, a, \operatorname{vadd}(m)(\operatorname{act}(m)(r, b), \operatorname{act}(m)(s, i)), ((r \cdot l) + (s \cdot c))) \end{array}}}

\vnbfile{theorem-library/nvs-taylor-proof.scm}{theorem-library/nvs-taylor-proof.scm}
\noindent{\small nvs-taylor-proof.scm -- Taylor's theorem with remainder bound for a map between two finite-dimensional real normed vector spaces.  PROVED.}\par
\vnbentry{\texttt{rr-\allowbreak{}power-\allowbreak{}one}}{\vnbstmt{\forall k \in \mathbb{N}.\; (\operatorname{power}(1, k) = 1)}}
\vnbentry{\texttt{nt-\allowbreak{}unit-\allowbreak{}factor}}{\vnbstmt{\begin{array}{@{}l@{}} \forall y, k. \\ \quad ((y \in \mathbb{R}) \wedge (k \in \mathbb{N})) \Rightarrow  \\ \quad \quad ((y \cdot \operatorname{power}((1 - 0), (k + 1))) = y) \end{array}}}
\vnbentry{\texttt{nvs-\allowbreak{}taylor-\allowbreak{}remainder-\allowbreak{}bound}}{\vnbstmt{\begin{array}{@{}l@{}} \forall m, b, f, a, h, n. \\ \quad \operatorname{is\text{-}normed\text{-}vector\text{-}space}(m) \wedge \\ \quad \operatorname{is\text{-}finite\text{-}dimensional}(\operatorname{normed\text{-}vector\text{-}space\text{-}as\text{-}module}(m)) \wedge \\ \quad \operatorname{is\text{-}normed\text{-}vector\text{-}space}(b) \wedge \\ \quad \operatorname{is\text{-}finite\text{-}dimensional}(\operatorname{normed\text{-}vector\text{-}space\text{-}as\text{-}module}(b)) \wedge \\ \quad (f \in (\operatorname{vec}(m) \to \operatorname{vec}(b))) \wedge \\ \quad (a \in \operatorname{vec}(m)) \wedge \\ \quad (h \in \operatorname{vec}(m)) \wedge \\ \quad (n \in \mathbb{N}) \wedge \\ \quad \operatorname{taylor\text{-}differentiable\text{-}v}(b, (\lambda t \in \mathbb{R}.\; f(\operatorname{vadd}(m)(a, \operatorname{act}(m)(t, h)))), 0, 1, n) \Rightarrow \\ \quad \quad \exists c. \\ \quad \quad \quad (0 < c) \wedge \\ \quad \quad \quad (c < 1) \wedge \\ \quad \quad \quad (\operatorname{factorial}((n + 1)) \cdot \operatorname{vnrm}(b)(\operatorname{vadd}(b)(f(\operatorname{vadd}(m)(a, h)), \operatorname{vneg}(b)(\operatorname{taylor\text{-}poly\text{-}v}(b, (\lambda t \in \mathbb{R}.\; f(\operatorname{vadd}(m)(a, \operatorname{act}(m)(t, h)))), 0, 1, n))))) \\ \quad \quad \quad \leq\; \operatorname{vnrm}(b)(\operatorname{nth\text{-}deriv\text{-}v}(b, (\lambda t \in \mathbb{R}.\; f(\operatorname{vadd}(m)(a, \operatorname{act}(m)(t, h)))), (n + 1))(c)) \end{array}}}

\vnbfile{theorem-library/antiderivative.scm}{theorem-library/antiderivative.scm}
\noindent{\small antiderivative.scm -- docs/calculus.pdf CHAPTER 4 SECTION 2: Definition 4.6 (antiderivative / antiderivable), Proposition 4.8 (closed under sums; the scalar half is handed off), Proposition 4.10 (two antiderivatives of the same function differ by a constant) and Corollary 4.11 (so f(b) - f(a) does}\par
\vnbentry{\texttt{antiderivative-\allowbreak{}map-\allowbreak{}in-\allowbreak{}fun}}{\vnbstmt{\begin{array}{@{}l@{}} \forall f, i, a, b. \\ \quad \operatorname{is\text{-}antiderivative}(f, i, a, b) \Rightarrow  \\ \quad \quad (f \in (\mathbb{R} \to \mathbb{R})) \end{array}}}
\vnbentry{\texttt{antiderivative-\allowbreak{}fn-\allowbreak{}in-\allowbreak{}fun}}{\vnbstmt{\begin{array}{@{}l@{}} \forall f, i, a, b. \\ \quad \operatorname{is\text{-}antiderivative}(f, i, a, b) \Rightarrow  \\ \quad \quad (i \in (\mathbb{R} \to \mathbb{R})) \end{array}}}
\vnbentry{\texttt{antiderivative-\allowbreak{}endpoints}}{\vnbstmt{\begin{array}{@{}l@{}} \forall f, i, a, b. \\ \quad \operatorname{is\text{-}antiderivative}(f, i, a, b) \Rightarrow  \\ \quad \quad ((a \in \mathbb{R}) \wedge ((b \in \mathbb{R}) \wedge (a < b))) \end{array}}}
\vnbentry{\texttt{deriv-\allowbreak{}sub}}{\vnbstmt{\begin{array}{@{}l@{}} \forall f, g, a, l, m. \\ \quad \operatorname{is\text{-}diff\text{-}at}(f, a, l) \wedge \\ \quad \operatorname{is\text{-}diff\text{-}at}(g, a, m) \Rightarrow \\ \quad \quad \operatorname{is\text{-}diff\text{-}at}((\lambda x \in \mathbb{R}.\; (f(x) - g(x))), a, (l - m)) \end{array}}}
\vnbentry{\texttt{antiderivative-\allowbreak{}differ-\allowbreak{}by-\allowbreak{}constant}}{\vnbstmt{\begin{array}{@{}l@{}} \forall f, g, i, a, b. \\ \quad \operatorname{is\text{-}antiderivative}(f, i, a, b) \wedge \\ \quad \operatorname{is\text{-}antiderivative}(g, i, a, b) \Rightarrow \\ \quad \quad \exists t \in \mathbb{R}. \forall s \in \operatorname{ccint}(a, b). \\ \quad \quad \quad (f(s) = (g(s) + t)) \end{array}}}
\vnbentry{\texttt{antiderivative-\allowbreak{}endpoint-\allowbreak{}difference}}{\vnbstmt{\begin{array}{@{}l@{}} \forall f, g, i, a, b. \\ \quad \operatorname{is\text{-}antiderivative}(f, i, a, b) \wedge \\ \quad \operatorname{is\text{-}antiderivative}(g, i, a, b) \Rightarrow \\ \quad \quad ((f(b) - f(a)) = (g(b) - g(a))) \end{array}}}
\vnbentry{\texttt{poly-\allowbreak{}is-\allowbreak{}antiderivable}}{\vnbstmt{\begin{array}{@{}l@{}} \forall n \in \mathbb{N}, f \in (\mathbb{N} \to \mathbb{R}), a, b. \\ \quad ((a \in \mathbb{R}) \wedge ((b \in \mathbb{R}) \wedge (a < b))) \Rightarrow  \\ \quad \quad \operatorname{is\text{-}antiderivable}((\lambda x \in \mathbb{R}.\; \\ \quad \quad \quad \quad \operatorname{series\text{-}partial\text{-}sum}((\lambda k \in \mathbb{N}.\; (f(k) \cdot \operatorname{power}(x, k))), (n + 1))), \\ \quad \quad \quad a, b) \end{array}}}
\vnbentry{\texttt{antiderivative-\allowbreak{}add}}{\vnbstmt{\begin{array}{@{}l@{}} \forall f, g, i, c, a, b. \\ \quad \operatorname{is\text{-}antiderivative}(f, i, a, b) \wedge \\ \quad \operatorname{is\text{-}antiderivative}(g, c, a, b) \Rightarrow \\ \quad \quad \operatorname{is\text{-}antiderivative}((\lambda x \in \mathbb{R}.\; (f(x) + g(x))), (\lambda x \in \mathbb{R}.\; (i(x) + c(x))), \\ \quad \quad \quad a, b) \end{array}}}
\vnbentry{\texttt{antiderivable-\allowbreak{}add}}{\vnbstmt{\begin{array}{@{}l@{}} \forall i, c, a, b. \\ \quad \operatorname{is\text{-}antiderivable}(i, a, b) \wedge \\ \quad \operatorname{is\text{-}antiderivable}(c, a, b) \Rightarrow \\ \quad \quad \operatorname{is\text{-}antiderivable}((\lambda x \in \mathbb{R}.\; (i(x) + c(x))), a, b) \end{array}}}
\vnbentry{\texttt{antiderivative-\allowbreak{}scale}}{\vnbstmt{\begin{array}{@{}l@{}} \forall c, f, i, a, b. \\ \quad ((c \in \mathbb{R}) \wedge \operatorname{is\text{-}antiderivative}(f, i, a, b)) \Rightarrow  \\ \quad \quad \operatorname{is\text{-}antiderivative}((\lambda x \in \mathbb{R}.\; (c \cdot f(x))), (\lambda x \in \mathbb{R}.\; (c \cdot i(x))), a, b) \end{array}}}
\vnbentry{\texttt{antiderivable-\allowbreak{}scale}}{\vnbstmt{\begin{array}{@{}l@{}} \forall c, i, a, b. \\ \quad ((c \in \mathbb{R}) \wedge \operatorname{is\text{-}antiderivable}(i, a, b)) \Rightarrow  \\ \quad \quad \operatorname{is\text{-}antiderivable}((\lambda x \in \mathbb{R}.\; (c \cdot i(x))), a, b) \end{array}}}

\vnbfile{theorem-library/antiderivative-transfer.scm}{theorem-library/antiderivative-transfer.scm}
\noindent{\small antiderivative-transfer.scm -- the POINTWISE TRANSFER for Def 4.6, and the DIFFERENCE of two antiderivatives.  All three `modulo 0'.}\par
\vnbentry{\texttt{antiderivative-\allowbreak{}transfer-\allowbreak{}ptwise-\allowbreak{}eq}}{\vnbstmt{\begin{array}{@{}l@{}} \forall f, i \in (\mathbb{R} \to \mathbb{R}), g, c, a, b. \\ \quad \operatorname{is\text{-}antiderivative}(g, c, a, b) \wedge \\ \quad \forall x \in \mathbb{R}.\; (f(x) = g(x)) \wedge \\ \quad \forall x \in \mathbb{R}.\; (i(x) = c(x)) \Rightarrow \\ \quad \quad \operatorname{is\text{-}antiderivative}(f, i, a, b) \end{array}}}
\vnbentry{\texttt{antiderivative-\allowbreak{}sub}}{\vnbstmt{\begin{array}{@{}l@{}} \forall f, g, i, c, a, b. \\ \quad \operatorname{is\text{-}antiderivative}(f, i, a, b) \wedge \\ \quad \operatorname{is\text{-}antiderivative}(g, c, a, b) \Rightarrow \\ \quad \quad \operatorname{is\text{-}antiderivative}((\lambda x \in \mathbb{R}.\; (f(x) - g(x))), (\lambda x \in \mathbb{R}.\; (i(x) - c(x))), \\ \quad \quad \quad a, b) \end{array}}}
\vnbentry{\texttt{antiderivable-\allowbreak{}sub}}{\vnbstmt{\begin{array}{@{}l@{}} \forall i, c, a, b. \\ \quad \operatorname{is\text{-}antiderivable}(i, a, b) \wedge \\ \quad \operatorname{is\text{-}antiderivable}(c, a, b) \Rightarrow \\ \quad \quad \operatorname{is\text{-}antiderivable}((\lambda x \in \mathbb{R}.\; (i(x) - c(x))), a, b) \end{array}}}

\vnbfile{theorem-library/antiderivative-lipschitz.scm}{theorem-library/antiderivative-lipschitz.scm}
\noindent{\small antiderivative-lipschitz.scm -- THE WITNESS ESTIMATE of docs/calculus.pdf Prop 4.16: a bound on the INTEGRAND bounds the CARATHEODORY WITNESS, at every point of the interval and by the same constant.}\par
\vnbentry{\texttt{antiderivative-\allowbreak{}increment-\allowbreak{}bound}}{\vnbstmt{\begin{array}{@{}l@{}} \forall h, i \in (\mathbb{R} \to \mathbb{R}), a, b, m, u, v \in \mathbb{R}. \\ \quad \operatorname{is\text{-}antiderivative}(h, i, a, b) \wedge \\ \quad \forall x \in \mathbb{R}. \\ \quad \quad ((x \in \operatorname{ccint}(a, b)) \Rightarrow (\lvert i(x) \rvert \leq m)) \wedge \\ \quad (u \in \operatorname{ccint}(a, b)) \wedge \\ \quad (v \in \operatorname{ccint}(a, b)) \wedge \\ \quad (u < v) \Rightarrow \\ \quad \quad (\lvert (h(v) - h(u)) \rvert \leq (m \cdot (v - u))) \end{array}}}
\vnbentry{\texttt{antiderivative-\allowbreak{}lipschitz}}{\vnbstmt{\begin{array}{@{}l@{}} \forall h, i \in (\mathbb{R} \to \mathbb{R}), a, b, m, u, v \in \mathbb{R}. \\ \quad \operatorname{is\text{-}antiderivative}(h, i, a, b) \wedge \\ \quad \forall x \in \mathbb{R}. \\ \quad \quad ((x \in \operatorname{ccint}(a, b)) \Rightarrow (\lvert i(x) \rvert \leq m)) \wedge \\ \quad (u \in \operatorname{ccint}(a, b)) \wedge \\ \quad (v \in \operatorname{ccint}(a, b)) \Rightarrow \\ \quad \quad (\lvert (h(v) - h(u)) \rvert \leq (m \cdot \lvert (v - u) \rvert)) \end{array}}}
\vnbentry{\texttt{caratheodory-\allowbreak{}witness-\allowbreak{}bound}}{\vnbstmt{\begin{array}{@{}l@{}} \forall h, i, c \in (\mathbb{R} \to \mathbb{R}), a, b, t, m, x \in \mathbb{R}. \\ \quad \operatorname{is\text{-}antiderivative}(h, i, a, b) \wedge \\ \quad \forall x \in \mathbb{R}. \\ \quad \quad ((x \in \operatorname{ccint}(a, b)) \Rightarrow (\lvert i(x) \rvert \leq m)) \wedge \\ \quad (t \in \operatorname{ccint}(a, b)) \wedge \\ \quad (c(t) = i(t)) \wedge \\ \quad \forall x \in \mathbb{R}.\; ((h(x) - h(t)) = (c(x) \cdot (x - t))) \wedge \\ \quad (x \in \operatorname{ccint}(a, b)) \Rightarrow \\ \quad \quad (\lvert c(x) \rvert \leq m) \end{array}}}
\vnbentry{\texttt{antiderivative-\allowbreak{}pair-\allowbreak{}close}}{\vnbstmt{\begin{array}{@{}l@{}} \forall f, g, i, c \in (\mathbb{R} \to \mathbb{R}), a, b, m, x \in \mathbb{R}. \\ \quad \operatorname{is\text{-}antiderivative}(f, i, a, b) \wedge \\ \quad \operatorname{is\text{-}antiderivative}(g, c, a, b) \wedge \\ \quad (f(a) = g(a)) \wedge \\ \quad \forall x \in \mathbb{R}. \\ \quad \quad (x \in \operatorname{ccint}(a, b)) \Rightarrow  \\ \quad \quad \quad (\lvert (i(x) - c(x)) \rvert \leq m) \wedge \\ \quad (x \in \operatorname{ccint}(a, b)) \Rightarrow \\ \quad \quad (\lvert (f(x) - g(x)) \rvert \leq (m \cdot (b - a))) \end{array}}}

\vnbfile{theorem-library/witness-family-choice.scm}{theorem-library/witness-family-choice.scm}
\noindent{\small witness-family-choice.scm -- TWO FAMILIES BUILT BY CHOICE, and nothing else.}\par
\vnbentry{\texttt{antiderivable-\allowbreak{}family-\allowbreak{}choice}}{\vnbstmt{\begin{array}{@{}l@{}} \forall m \in (\mathbb{N} \to (\mathbb{R} \to \mathbb{R})), a, b \in \mathbb{R}. \\ \quad \forall k \in \mathbb{N}.\; \operatorname{is\text{-}antiderivable}(m(k), a, b) \Rightarrow  \\ \quad \quad \exists c \in (\mathbb{N} \to (\mathbb{R} \to \mathbb{R})). \forall k \in \mathbb{N}. \\ \quad \quad \quad \operatorname{is\text{-}antiderivative}(c(k), m(k), a, b) \end{array}}}
\vnbentry{\texttt{diff-\allowbreak{}at-\allowbreak{}witness-\allowbreak{}family}}{\vnbstmt{\begin{array}{@{}l@{}} \forall m, a, t. \\ \quad (m \in (\mathbb{N} \to (\mathbb{R} \to \mathbb{R}))) \wedge \\ \quad (a \in (\mathbb{N} \to \mathbb{R})) \wedge \\ \quad (t \in \mathbb{R}) \wedge \\ \quad \forall k \in \mathbb{N}. \\ \quad \quad \operatorname{is\text{-}diff\text{-}at}(m(k), t, a(k)) \Rightarrow \\ \quad \quad \exists b \in (\mathbb{N} \to (\mathbb{R} \to \mathbb{R})). \forall k \in \mathbb{N}. \\ \quad \quad \quad \operatorname{is\text{-}continuous\text{-}at}(\operatorname{rr\text{-}ms}, \operatorname{rr\text{-}ms}, b(k), t) \wedge \\ \quad \quad \quad (b(k)(t) = a(k)) \wedge \\ \quad \quad \quad \forall x \in \mathbb{R}.\; ((m(k)(x) - m(k)(t)) = (b(k)(x) \cdot (x - t))) \end{array}}}

\vnbfile{theorem-library/antiderivative-normalize.scm}{theorem-library/antiderivative-normalize.scm}
\noindent{\small antiderivative-normalize.scm -- SHIFTING AN ANTIDERIVATIVE BY A CONSTANT, and the NORMALISED antiderivative family that docs/calculus.pdf Cor 4.17 needs and that nothing in the tree supplied.}\par
\vnbentry{\texttt{antiderivative-\allowbreak{}shift-\allowbreak{}const}}{\vnbstmt{\begin{array}{@{}l@{}} \forall f, i, a, b, c. \\ \quad (f \in (\mathbb{R} \to \mathbb{R})) \wedge \\ \quad (i \in (\mathbb{R} \to \mathbb{R})) \wedge \\ \quad (a \in \mathbb{R}) \wedge \\ \quad (b \in \mathbb{R}) \wedge \\ \quad (c \in \mathbb{R}) \wedge \\ \quad \operatorname{is\text{-}antiderivative}(f, i, a, b) \Rightarrow \\ \quad \quad \operatorname{is\text{-}antiderivative}((\lambda x \in \mathbb{R}.\; (f(x) - c)), i, a, b) \end{array}}}
\vnbentry{\texttt{antiderivable-\allowbreak{}family-\allowbreak{}normalized}}{\vnbstmt{\begin{array}{@{}l@{}} \forall m \in (\mathbb{N} \to (\mathbb{R} \to \mathbb{R})), a, b \in \mathbb{R}. \\ \quad \forall k \in \mathbb{N}.\; \operatorname{is\text{-}antiderivable}(m(k), a, b) \Rightarrow  \\ \quad \quad \exists c \in (\mathbb{N} \to (\mathbb{R} \to \mathbb{R})). \forall k \in \mathbb{N}. \\ \quad \quad \quad (\operatorname{is\text{-}antiderivative}(c(k), m(k), a, b) \wedge (c(k)(a) = 0)) \end{array}}}

\vnbfile{theorem-library/antiderivable-uniform-limit.scm}{theorem-library/antiderivable-uniform-limit.scm}
\noindent{\small theorem-library/antiderivable-uniform-limit.scm -- docs/calculus.pdf PROP 4.16 and COR 4.17: THE ANTIDERIVABLE FUNCTIONS ARE CLOSED UNDER UNIFORM LIMITS.  Rung 3 of the integration arc.}\par
\vnbentry{\texttt{rr-\allowbreak{}converges-\allowbreak{}to-\allowbreak{}abs}}{\vnbstmt{\begin{array}{@{}l@{}} \forall f \in (\mathbb{N} \to \mathbb{R}), v \in \mathbb{R}. \\ \quad \forall s. \\ \quad \quad \operatorname{pos\text{-}rr}(s) \Rightarrow  \\ \quad \quad \quad \exists n \in \mathbb{N}. \forall k \in \mathbb{N}. \\ \quad \quad \quad \quad ((n \leq k) \Rightarrow (\lvert (f(k) - v) \rvert \leq s)) \Rightarrow \\ \quad \quad \operatorname{converges\text{-}to}(\operatorname{rr\text{-}ms}, f, v) \end{array}}}
\vnbentry{\texttt{rr-\allowbreak{}limit-\allowbreak{}ptwise-\allowbreak{}eq}}{\vnbstmt{\begin{array}{@{}l@{}} \forall f, h \in (\mathbb{N} \to \mathbb{R}), v \in \mathbb{R}. \\ \quad \forall j \in \mathbb{N}.\; (h(j) = f(j)) \wedge \\ \quad \operatorname{converges\text{-}to}(\operatorname{rr\text{-}ms}, f, v) \Rightarrow \\ \quad \quad \operatorname{converges\text{-}to}(\operatorname{rr\text{-}ms}, h, v) \end{array}}}
\vnbentry{\texttt{uniform-\allowbreak{}cauchy-\allowbreak{}limit}}{\vnbstmt{\begin{array}{@{}l@{}} \forall g \in (\mathbb{N} \to (\mathbb{R} \to \mathbb{R})), s. \\ \quad \forall a. \\ \quad \quad \operatorname{pos\text{-}rr}(a) \Rightarrow  \\ \quad \quad \quad \exists n \in \mathbb{N}. \forall k, l \in \mathbb{N}, x \in \mathbb{R}. \\ \quad \quad \quad \quad ((n \leq k) \wedge ((n \leq l) \wedge (x \in s))) \Rightarrow  \\ \quad \quad \quad \quad \quad (\lvert (g(k)(x) - g(l)(x)) \rvert \leq a) \Rightarrow \\ \quad \quad \exists h \in (\mathbb{R} \to \mathbb{R}). \\ \quad \quad \quad \forall x \in \mathbb{R}. \\ \quad \quad \quad \quad (x \in s) \Rightarrow  \\ \quad \quad \quad \quad \quad \operatorname{converges\text{-}to}(\operatorname{rr\text{-}ms}, (\lambda k \in \mathbb{N}.\; g(k)(x)), h(x)) \wedge \\ \quad \quad \quad \forall a. \\ \quad \quad \quad \quad \operatorname{pos\text{-}rr}(a) \Rightarrow  \\ \quad \quad \quad \quad \quad \exists n \in \mathbb{N}. \forall k \in \mathbb{N}, x \in \mathbb{R}. \\ \quad \quad \quad \quad \quad \quad ((n \leq k) \wedge (x \in s)) \Rightarrow  \\ \quad \quad \quad \quad \quad \quad \quad (\lvert (h(x) - g(k)(x)) \rvert \leq a) \end{array}}}
\vnbentry{\texttt{caratheodory-\allowbreak{}witness-\allowbreak{}pair-\allowbreak{}bound}}{\vnbstmt{\begin{array}{@{}l@{}} \forall c, d, f, g, h, k \in (\mathbb{R} \to \mathbb{R}), a, b, t, m, x \in \mathbb{R}. \\ \quad \operatorname{is\text{-}antiderivative}(c, f, a, b) \wedge \\ \quad \operatorname{is\text{-}antiderivative}(d, g, a, b) \wedge \\ \quad \forall x \in \mathbb{R}. \\ \quad \quad (x \in \operatorname{ccint}(a, b)) \Rightarrow  \\ \quad \quad \quad (\lvert (f(x) - g(x)) \rvert \leq m) \wedge \\ \quad (t \in \operatorname{ccint}(a, b)) \wedge \\ \quad (h(t) = f(t)) \wedge \\ \quad (k(t) = g(t)) \wedge \\ \quad \forall x \in \mathbb{R}.\; ((c(x) - c(t)) = (h(x) \cdot (x - t))) \wedge \\ \quad \forall x \in \mathbb{R}.\; ((d(x) - d(t)) = (k(x) \cdot (x - t))) \wedge \\ \quad (x \in \operatorname{ccint}(a, b)) \Rightarrow \\ \quad \quad (\lvert (h(x) - k(x)) \rvert \leq m) \end{array}}}
\vnbentry{\texttt{uniform-\allowbreak{}witness-\allowbreak{}diff-\allowbreak{}at}}{\vnbstmt{\begin{array}{@{}l@{}} \forall m, c \in (\mathbb{N} \to (\mathbb{R} \to \mathbb{R})), d \in (\mathbb{R} \to \mathbb{R}), v, a, b, t, o \in \mathbb{R}. \\ \quad (0 < o) \wedge \\ \quad \forall y \in \mathbb{R}. \\ \quad \quad ((\lvert (y - t) \rvert \leq o) \Rightarrow (y \in \operatorname{ccint}(a, b))) \wedge \\ \quad (t \in \operatorname{ccint}(a, b)) \wedge \\ \quad \forall k \in \mathbb{N}. \\ \quad \quad \operatorname{is\text{-}continuous\text{-}at}(\operatorname{rr\text{-}ms}, \operatorname{rr\text{-}ms}, c(k), t) \wedge \\ \quad \forall k \in \mathbb{N}, y \in \mathbb{R}. \\ \quad \quad (y \in \operatorname{ccint}(a, b)) \Rightarrow  \\ \quad \quad \quad ((m(k)(y) - m(k)(t)) = (c(k)(y) \cdot (y - t))) \wedge \\ \quad \forall y \in \mathbb{R}. \\ \quad \quad (y \in \operatorname{ccint}(a, b)) \Rightarrow  \\ \quad \quad \quad \operatorname{converges\text{-}to}(\operatorname{rr\text{-}ms}, (\lambda k \in \mathbb{N}.\; m(k)(y)), d(y)) \wedge \\ \quad \forall s. \\ \quad \quad \operatorname{pos\text{-}rr}(s) \Rightarrow  \\ \quad \quad \quad \exists n \in \mathbb{N}. \forall k, l \in \mathbb{N}, x \in \mathbb{R}. \\ \quad \quad \quad \quad ((n \leq k) \wedge ((n \leq l) \wedge (x \in \operatorname{ccint}(a, b)))) \Rightarrow  \\ \quad \quad \quad \quad \quad (\lvert (c(k)(x) - c(l)(x)) \rvert \leq s) \wedge \\ \quad \operatorname{converges\text{-}to}(\operatorname{rr\text{-}ms}, (\lambda k \in \mathbb{N}.\; c(k)(t)), v) \Rightarrow \\ \quad \quad \operatorname{is\text{-}diff\text{-}at}(d, t, v) \end{array}}}
\vnbentry{\texttt{antiderivative-\allowbreak{}family-\allowbreak{}uniform-\allowbreak{}limit}}{\vnbstmt{\begin{array}{@{}l@{}} \forall m, c \in (\mathbb{N} \to (\mathbb{R} \to \mathbb{R})), i \in (\mathbb{R} \to \mathbb{R}), a, b \in \mathbb{R}. \\ \quad (a < b) \wedge \\ \quad \forall k \in \mathbb{N}. \\ \quad \quad (\operatorname{is\text{-}antiderivative}(m(k), c(k), a, b) \wedge (m(k)(a) = 0)) \wedge \\ \quad \forall s. \\ \quad \quad \operatorname{pos\text{-}rr}(s) \Rightarrow  \\ \quad \quad \quad \exists n \in \mathbb{N}. \forall k \in \mathbb{N}, x \in \mathbb{R}. \\ \quad \quad \quad \quad ((n \leq k) \wedge (x \in \operatorname{ccint}(a, b))) \Rightarrow  \\ \quad \quad \quad \quad \quad (\lvert (i(x) - c(k)(x)) \rvert \leq s) \Rightarrow \\ \quad \quad \exists f \in (\mathbb{R} \to \mathbb{R}). \\ \quad \quad \quad \operatorname{is\text{-}antiderivative}(f, i, a, b) \wedge \\ \quad \quad \quad \forall s. \\ \quad \quad \quad \quad \operatorname{pos\text{-}rr}(s) \Rightarrow  \\ \quad \quad \quad \quad \quad \exists n \in \mathbb{N}. \forall k \in \mathbb{N}, x \in \mathbb{R}. \\ \quad \quad \quad \quad \quad \quad ((n \leq k) \wedge (x \in \operatorname{ccint}(a, b))) \Rightarrow  \\ \quad \quad \quad \quad \quad \quad \quad (\lvert (f(x) - m(k)(x)) \rvert \leq s) \end{array}}}
\vnbentry{\texttt{antiderivable-\allowbreak{}uniform-\allowbreak{}limit}}{\vnbstmt{\begin{array}{@{}l@{}} \forall m \in (\mathbb{N} \to (\mathbb{R} \to \mathbb{R})), i \in (\mathbb{R} \to \mathbb{R}), a, b \in \mathbb{R}. \\ \quad \forall k \in \mathbb{N}.\; \operatorname{is\text{-}antiderivable}(m(k), a, b) \wedge \\ \quad \forall s. \\ \quad \quad \operatorname{pos\text{-}rr}(s) \Rightarrow  \\ \quad \quad \quad \exists n \in \mathbb{N}. \forall k \in \mathbb{N}, x \in \mathbb{R}. \\ \quad \quad \quad \quad ((n \leq k) \wedge (x \in \operatorname{ccint}(a, b))) \Rightarrow  \\ \quad \quad \quad \quad \quad (\lvert (i(x) - m(k)(x)) \rvert \leq s) \Rightarrow \\ \quad \quad \operatorname{is\text{-}antiderivable}(i, a, b) \end{array}}}

\vnbfile{theorem-library/series-antiderivable.scm}{theorem-library/series-antiderivable.scm}
\noindent{\small series-antiderivable.scm -- A FINITE SUM OF ANTIDERIVABLE TERMS IS ANTIDERIVABLE, and the vocabulary that makes such a statement usable.}\par
\vnbentry{\texttt{antiderivable-\allowbreak{}fn-\allowbreak{}in-\allowbreak{}fun}}{\vnbstmt{\begin{array}{@{}l@{}} \forall i, a, b. \\ \quad \operatorname{is\text{-}antiderivable}(i, a, b) \Rightarrow  \\ \quad \quad (i \in (\mathbb{R} \to \mathbb{R})) \end{array}}}
\vnbentry{\texttt{antiderivative-\allowbreak{}integrand-\allowbreak{}transfer}}{\vnbstmt{\begin{array}{@{}l@{}} \forall f, i, c, a, b. \\ \quad (c \in (\mathbb{R} \to \mathbb{R})) \wedge \\ \quad \operatorname{is\text{-}antiderivative}(f, i, a, b) \wedge \\ \quad \forall x \in \mathbb{R}.\; (c(x) \simeq i(x)) \Rightarrow \\ \quad \quad \operatorname{is\text{-}antiderivative}(f, c, a, b) \end{array}}}
\vnbentry{\texttt{antiderivable-\allowbreak{}integrand-\allowbreak{}transfer}}{\vnbstmt{\begin{array}{@{}l@{}} \forall i, c, a, b. \\ \quad (c \in (\mathbb{R} \to \mathbb{R})) \wedge \\ \quad \operatorname{is\text{-}antiderivable}(i, a, b) \wedge \\ \quad \forall x \in \mathbb{R}.\; (c(x) \simeq i(x)) \Rightarrow \\ \quad \quad \operatorname{is\text{-}antiderivable}(c, a, b) \end{array}}}
\vnbentry{\texttt{antiderivable-\allowbreak{}from-\allowbreak{}deriv}}{\vnbstmt{\begin{array}{@{}l@{}} \forall f, i, a, b. \\ \quad (f \in (\mathbb{R} \to \mathbb{R})) \wedge \\ \quad (i \in (\mathbb{R} \to \mathbb{R})) \wedge \\ \quad (a \in \mathbb{R}) \wedge \\ \quad (b \in \mathbb{R}) \wedge \\ \quad (a < b) \wedge \\ \quad \forall x \in \mathbb{R}.\; \operatorname{is\text{-}diff\text{-}at}(f, x, i(x)) \Rightarrow \\ \quad \quad \operatorname{is\text{-}antiderivable}(i, a, b) \end{array}}}
\vnbentry{\texttt{zero-\allowbreak{}is-\allowbreak{}antiderivable}}{\vnbstmt{\begin{array}{@{}l@{}} \forall a, b. \\ \quad ((a \in \mathbb{R}) \wedge ((b \in \mathbb{R}) \wedge (a < b))) \Rightarrow  \\ \quad \quad \operatorname{is\text{-}antiderivable}((\lambda x \in \mathbb{R}.\; 0), a, b) \end{array}}}
\vnbentry{\texttt{series-\allowbreak{}partial-\allowbreak{}sum-\allowbreak{}antiderivable}}{\vnbstmt{\begin{array}{@{}l@{}} \forall n \in \mathbb{N}, s, h, a, b. \\ \quad (a \in \mathbb{R}) \wedge \\ \quad (b \in \mathbb{R}) \wedge \\ \quad (a < b) \wedge \\ \quad (h \in (\mathbb{R} \to \mathbb{R})) \wedge \\ \quad \forall j \in \mathbb{N}.\; \operatorname{is\text{-}antiderivable}(s(j), a, b) \wedge \\ \quad \forall x \in \mathbb{R}. \\ \quad \quad (h(x) \simeq \operatorname{series\text{-}partial\text{-}sum}((\lambda j \in \mathbb{N}.\; s(j)(x)), n)) \Rightarrow \\ \quad \quad \operatorname{is\text{-}antiderivable}(h, a, b) \end{array}}}
\vnbentry{\texttt{power-\allowbreak{}antiderivable}}{\vnbstmt{\begin{array}{@{}l@{}} \forall m, a, b. \\ \quad ((m \in \mathbb{N}) \wedge ((a \in \mathbb{R}) \wedge ((b \in \mathbb{R}) \wedge (a < b)))) \Rightarrow  \\ \quad \quad \operatorname{is\text{-}antiderivable}((\lambda x \in \mathbb{R}.\; \operatorname{power}(x, m)), a, b) \end{array}}}
\vnbentry{\texttt{poly-\allowbreak{}fn-\allowbreak{}antiderivable}}{\vnbstmt{\begin{array}{@{}l@{}} \forall n \in \mathbb{N}, f, h, a, b. \\ \quad (f \in (\mathbb{N} \to \mathbb{R})) \wedge \\ \quad (a \in \mathbb{R}) \wedge \\ \quad (b \in \mathbb{R}) \wedge \\ \quad (a < b) \wedge \\ \quad (h \in (\mathbb{R} \to \mathbb{R})) \wedge \\ \quad \forall x \in \mathbb{R}. \\ \quad \quad (h(x) \simeq \operatorname{series\text{-}partial\text{-}sum}((\lambda k \in \mathbb{N}.\; (f(k) \cdot \operatorname{power}(x, k))), n)) \Rightarrow \\ \quad \quad \operatorname{is\text{-}antiderivable}(h, a, b) \end{array}}}

\vnbfile{theorem-library/c-int.scm}{theorem-library/c-int.scm}
\noindent{\small c-int.scm -- docs/calculus.pdf CHAPTER 4, equation (64): the DEFINITE INTEGRAL of a continuous function, `C-INT'.  Rung 4.}\par
\vnbentry{\texttt{c-\allowbreak{}int-\allowbreak{}value}}{\vnbstmt{\begin{array}{@{}l@{}} \forall f, i, a, b. \\ \quad \operatorname{is\text{-}antiderivative}(f, i, a, b) \Rightarrow  \\ \quad \quad (\int_{a}^{b} i = (f(b) - f(a))) \end{array}}}
\vnbentry{\texttt{c-\allowbreak{}int-\allowbreak{}in-\allowbreak{}rr}}{\vnbstmt{\begin{array}{@{}l@{}} \forall i, a, b. \\ \quad \operatorname{is\text{-}antiderivable}(i, a, b) \Rightarrow  \\ \quad \quad (\int_{a}^{b} i \in \mathbb{R}) \end{array}}}
\vnbentry{\texttt{c-\allowbreak{}int-\allowbreak{}add}}{\vnbstmt{\begin{array}{@{}l@{}} \forall i, c, a, b. \\ \quad \operatorname{is\text{-}antiderivable}(i, a, b) \wedge \\ \quad \operatorname{is\text{-}antiderivable}(c, a, b) \Rightarrow \\ \quad \quad (\int_{a}^{b} (\lambda x \in \mathbb{R}.\; (i(x) + c(x))) = (\int_{a}^{b} i + \int_{a}^{b} c)) \end{array}}}
\vnbentry{\texttt{c-\allowbreak{}int-\allowbreak{}scale}}{\vnbstmt{\begin{array}{@{}l@{}} \forall c, i, a, b. \\ \quad ((c \in \mathbb{R}) \wedge \operatorname{is\text{-}antiderivable}(i, a, b)) \Rightarrow  \\ \quad \quad (\int_{a}^{b} (\lambda x \in \mathbb{R}.\; (c \cdot i(x))) = (c \cdot \int_{a}^{b} i)) \end{array}}}

\vnbfile{theorem-library/c-int-oriented.scm}{theorem-library/c-int-oriented.scm}
\noindent{\small c-int-oriented.scm -- THE ORIENTED DEFINITE INTEGRAL, the sub-interval restriction of Def 4.6 (docs/calculus.pdf REMARK 4.9), and ADDITIVITY IN THE BOUNDS.}\par
\vnbentry{\texttt{rr-\allowbreak{}lt-\allowbreak{}asymm}}{\vnbstmt{\forall a, b \in \mathbb{R}.\; ((a < b) \Rightarrow \neg (b < a))}}
\vnbentry{\texttt{antiderivable-\allowbreak{}endpoints}}{\vnbstmt{\begin{array}{@{}l@{}} \forall i, a, b. \\ \quad \operatorname{is\text{-}antiderivable}(i, a, b) \Rightarrow  \\ \quad \quad ((a \in \mathbb{R}) \wedge ((b \in \mathbb{R}) \wedge (a < b))) \end{array}}}
\vnbentry{\texttt{antiderivable-\allowbreak{}not-\allowbreak{}reverse}}{\vnbstmt{\begin{array}{@{}l@{}} \forall i, a, b. \\ \quad \operatorname{is\text{-}antiderivable}(i, a, b) \Rightarrow  \\ \quad \quad \neg \operatorname{is\text{-}antiderivable}(i, b, a) \end{array}}}
\vnbentry{\texttt{antiderivable-\allowbreak{}not-\allowbreak{}degenerate}}{\vnbstmt{\forall i, a.\; \neg \operatorname{is\text{-}antiderivable}(i, a, a)}}
\vnbentry{\texttt{c-\allowbreak{}int-\allowbreak{}or-\allowbreak{}anti}}{\vnbstmt{\begin{array}{@{}l@{}} \forall i, a, b. \\ \quad \operatorname{is\text{-}antiderivable}(i, a, b) \Rightarrow  \\ \quad \quad (\int_{a}^{b} i = \int_{a}^{b} i) \end{array}}}
\vnbentry{\texttt{c-\allowbreak{}int-\allowbreak{}or-\allowbreak{}flip}}{\vnbstmt{\begin{array}{@{}l@{}} \forall i, a, b. \\ \quad \operatorname{is\text{-}antiderivable}(i, a, b) \Rightarrow  \\ \quad \quad (\int_{b}^{a} i = -\int_{a}^{b} i) \end{array}}}
\vnbentry{\texttt{c-\allowbreak{}int-\allowbreak{}or-\allowbreak{}none}}{\vnbstmt{\begin{array}{@{}l@{}} \forall i, a, b. \\ \quad \neg \operatorname{is\text{-}antiderivable}(i, a, b) \wedge \\ \quad \neg \operatorname{is\text{-}antiderivable}(i, b, a) \Rightarrow \\ \quad \quad (\int_{a}^{b} i = 0) \end{array}}}
\vnbentry{\texttt{c-\allowbreak{}int-\allowbreak{}or-\allowbreak{}degenerate}}{\vnbstmt{\forall i, a.\; (\int_{a}^{a} i = 0)}}
\vnbentry{\texttt{c-\allowbreak{}int-\allowbreak{}or-\allowbreak{}in-\allowbreak{}rr}}{\vnbstmt{\forall i, a, b.\; (\int_{a}^{b} i \in \mathbb{R})}}
\vnbentry{\texttt{c-\allowbreak{}int-\allowbreak{}or-\allowbreak{}reverse}}{\vnbstmt{\forall i, a, b.\; (\int_{b}^{a} i = -\int_{a}^{b} i)}}
\vnbentry{\texttt{antiderivative-\allowbreak{}subinterval}}{\vnbstmt{\begin{array}{@{}l@{}} \forall f, i, a, b, o, c. \\ \quad \operatorname{is\text{-}antiderivative}(f, i, a, b) \wedge \\ \quad (o \in \mathbb{R}) \wedge \\ \quad (c \in \mathbb{R}) \wedge \\ \quad (a \leq o) \wedge \\ \quad (o < c) \wedge \\ \quad (c \leq b) \Rightarrow \\ \quad \quad \operatorname{is\text{-}antiderivative}(f, i, o, c) \end{array}}}
\vnbentry{\texttt{antiderivable-\allowbreak{}subinterval}}{\vnbstmt{\begin{array}{@{}l@{}} \forall i, a, b, o, c. \\ \quad \operatorname{is\text{-}antiderivable}(i, a, b) \wedge \\ \quad (o \in \mathbb{R}) \wedge \\ \quad (c \in \mathbb{R}) \wedge \\ \quad (a \leq o) \wedge \\ \quad (o < c) \wedge \\ \quad (c \leq b) \Rightarrow \\ \quad \quad \operatorname{is\text{-}antiderivable}(i, o, c) \end{array}}}
\vnbentry{\texttt{c-\allowbreak{}int-\allowbreak{}or-\allowbreak{}value}}{\vnbstmt{\begin{array}{@{}l@{}} \forall f, i, p, q, a, b. \\ \quad \operatorname{is\text{-}antiderivative}(f, i, p, q) \wedge \\ \quad (a \in \operatorname{ccint}(p, q)) \wedge \\ \quad (b \in \operatorname{ccint}(p, q)) \Rightarrow \\ \quad \quad (\int_{a}^{b} i = (f(b) - f(a))) \end{array}}}
\vnbentry{\texttt{c-\allowbreak{}int-\allowbreak{}or-\allowbreak{}additive}}{\vnbstmt{\begin{array}{@{}l@{}} \forall i, p, q, a, b, c. \\ \quad \operatorname{is\text{-}antiderivable}(i, p, q) \wedge \\ \quad (a \in \operatorname{ccint}(p, q)) \wedge \\ \quad (b \in \operatorname{ccint}(p, q)) \wedge \\ \quad (c \in \operatorname{ccint}(p, q)) \Rightarrow \\ \quad \quad (\int_{a}^{c} i = (\int_{a}^{b} i + \int_{b}^{c} i)) \end{array}}}
\vnbentry{\texttt{c-\allowbreak{}int-\allowbreak{}add-\allowbreak{}bounds}}{\vnbstmt{\begin{array}{@{}l@{}} \forall i, a, b, c. \\ \quad \operatorname{is\text{-}antiderivable}(i, a, c) \wedge \\ \quad (b \in \mathbb{R}) \wedge \\ \quad (a < b) \wedge \\ \quad (b < c) \Rightarrow \\ \quad \quad (\int_{a}^{c} i = (\int_{a}^{b} i + \int_{b}^{c} i)) \end{array}}}

\vnbfile{theorem-library/bernstein-ccint.scm}{theorem-library/bernstein-ccint.scm}
\noindent{\small bernstein-ccint.scm -- THEOREM 5.2 of docs/calculus.pdf moved off [0,1] and onto an arbitrary nondegenerate [a,b], by the affine change of variable}\par
\vnbentry{\texttt{ccint-\allowbreak{}parts}}{\vnbstmt{\begin{array}{@{}l@{}} \forall o, i.; \forall t \in \operatorname{ccint}(o, i). \\ \quad ((t \in \mathbb{R}) \wedge ((o \leq t) \wedge (t \leq i))) \end{array}}}
\vnbentry{\texttt{ccint-\allowbreak{}affine-\allowbreak{}in}}{\vnbstmt{\begin{array}{@{}l@{}} \forall a, b, t. \\ \quad (a \in \mathbb{R}) \wedge \\ \quad (b \in \mathbb{R}) \wedge \\ \quad (a < b) \wedge \\ \quad (t \in \operatorname{ccint}(0, 1)) \Rightarrow \\ \quad \quad ((a + ((b - a) \cdot t)) \in \operatorname{ccint}(a, b)) \end{array}}}
\vnbentry{\texttt{ccint-\allowbreak{}coaffine-\allowbreak{}in}}{\vnbstmt{\begin{array}{@{}l@{}} \forall a, b, x. \\ \quad (a \in \mathbb{R}) \wedge \\ \quad (b \in \mathbb{R}) \wedge \\ \quad (a < b) \wedge \\ \quad (x \in \operatorname{ccint}(a, b)) \Rightarrow \\ \quad \quad (((x - a) \cdot \operatorname{recip}((b - a))) \in \operatorname{ccint}(0, 1)) \end{array}}}
\vnbentry{\texttt{ccint-\allowbreak{}affine-\allowbreak{}inverse}}{\vnbstmt{\begin{array}{@{}l@{}} \forall a, b, x. \\ \quad ((a \in \mathbb{R}) \wedge ((b \in \mathbb{R}) \wedge ((a < b) \wedge (x \in \mathbb{R})))) \Rightarrow  \\ \quad \quad ((a + ((b - a) \cdot ((x - a) \cdot \operatorname{recip}((b - a))))) = x) \end{array}}}
\vnbentry{\texttt{affine-\allowbreak{}continuous-\allowbreak{}at}}{\vnbstmt{\begin{array}{@{}l@{}} \forall c, m, t. \\ \quad ((c \in \mathbb{R}) \wedge ((m \in \mathbb{R}) \wedge (t \in \mathbb{R}))) \Rightarrow  \\ \quad \quad \operatorname{is\text{-}continuous\text{-}at}(\operatorname{rr\text{-}ms}, \operatorname{rr\text{-}ms}, (\lambda x \in \mathbb{R}.\; (c + (m \cdot x))), t) \end{array}}}
\vnbentry{\texttt{bernstein-\allowbreak{}uniform-\allowbreak{}approximation-\allowbreak{}ccint}}{\vnbstmt{\begin{array}{@{}l@{}} \forall n, a, b, p. \\ \quad (n \in (\mathbb{R} \to \mathbb{R})) \wedge \\ \quad (a \in \mathbb{R}) \wedge \\ \quad (b \in \mathbb{R}) \wedge \\ \quad (a < b) \wedge \\ \quad \forall c \in \operatorname{ccint}(a, b). \\ \quad \quad \operatorname{is\text{-}continuous\text{-}at}(\operatorname{rr\text{-}ms}, \operatorname{rr\text{-}ms}, n, c) \wedge \\ \quad \operatorname{pos\text{-}rr}(p) \Rightarrow \\ \quad \quad \exists \mathit{cap} \in \mathbb{N}. \forall d \in \mathbb{N}, c \in \operatorname{ccint}(a, b). \\ \quad \quad \quad (\mathit{cap} \leq d) \Rightarrow  \\ \quad \quad \quad \quad \lvert (n(c) - \operatorname{bernstein\text{-}poly}(\operatorname{compose}(n, (\lambda x \in \mathbb{R}.\; (a + ((b - a) \cdot x)))), d, ((c - a) \cdot \operatorname{recip}((b - a))))) \rvert \\ \quad \quad \quad \quad \leq\; p \end{array}}}

\vnbfile{theorem-library/rake-bolzano-weierstrass.scm}{theorem-library/rake-bolzano-weierstrass.scm}
\noindent{\small theorem-library/rake-bolzano-weierstrass.scm ==================================================================== The eps-GRID FINITE COVER of a closed interval -- batch 8-E, item (1) -- and the scaled archimedean property it needs.}\par
\vnbentry{\texttt{rr-\allowbreak{}archimedean-\allowbreak{}multiple}}{\vnbstmt{\begin{array}{@{}l@{}} \forall v.; \forall \mathit{xv} \in \mathbb{R}. \\ \quad \operatorname{pos\text{-}rr}(v) \Rightarrow  \\ \quad \quad \exists \mathit{nv} \in \mathbb{N}.\; (\mathit{xv} \leq (\mathit{nv} \cdot v)) \end{array}}}
\vnbentry{\texttt{ccint-\allowbreak{}eps-\allowbreak{}grid-\allowbreak{}cover-\allowbreak{}at}}{\vnbstmt{\begin{array}{@{}l@{}} \forall v \in \mathbb{N}, \mathit{av}, \mathit{bv} \in \mathbb{R}, \mathit{ev}. \\ \quad \operatorname{pos\text{-}rr}(\mathit{ev}) \wedge \\ \quad (\mathit{bv} \leq (\mathit{av} + (v \cdot \mathit{ev}))) \Rightarrow \\ \quad \quad \exists \mathit{cv}. \\ \quad \quad \quad \operatorname{is\text{-}finite\text{-}cover}(\mathit{cv}, \operatorname{ccint}(\mathit{av}, \mathit{bv})) \wedge \\ \quad \quad \quad \forall \mathit{uv} \in \mathit{cv}, \mathit{pv}, \mathit{qv} \in \mathit{uv}. \\ \quad \quad \quad \quad (\mathit{pv} \in \mathbb{R}) \wedge \\ \quad \quad \quad \quad (\mathit{qv} \in \mathbb{R}) \wedge \\ \quad \quad \quad \quad (\lvert (\mathit{pv} - \mathit{qv}) \rvert \leq \mathit{ev}) \end{array}}}
\vnbentry{\texttt{ccint-\allowbreak{}eps-\allowbreak{}grid-\allowbreak{}cover}}{\vnbstmt{\begin{array}{@{}l@{}} \forall v, \mathit{bv} \in \mathbb{R}, \mathit{ev}. \\ \quad \operatorname{pos\text{-}rr}(\mathit{ev}) \Rightarrow  \\ \quad \quad \exists \mathit{cv}. \\ \quad \quad \quad \operatorname{is\text{-}finite\text{-}cover}(\mathit{cv}, \operatorname{ccint}(v, \mathit{bv})) \wedge \\ \quad \quad \quad \forall \mathit{uv} \in \mathit{cv}, \mathit{pv}, \mathit{qv} \in \mathit{uv}. \\ \quad \quad \quad \quad (\mathit{pv} \in \mathbb{R}) \wedge \\ \quad \quad \quad \quad (\mathit{qv} \in \mathbb{R}) \wedge \\ \quad \quad \quad \quad (\lvert (\mathit{pv} - \mathit{qv}) \rvert \leq \mathit{ev}) \end{array}}}
\vnbentry{\texttt{ccint-\allowbreak{}grid-\allowbreak{}cover-\allowbreak{}seq}}{\vnbstmt{\begin{array}{@{}l@{}} \forall v, \mathit{bv} \in \mathbb{R}, d. \\ \quad \forall \mathit{kv} \in \mathbb{N}.\; \operatorname{pos\text{-}rr}(d(\mathit{kv})) \Rightarrow  \\ \quad \quad \exists w. \forall \mathit{kv} \in \mathbb{N}. \\ \quad \quad \quad \operatorname{is\text{-}finite\text{-}cover}(w(\mathit{kv}), \operatorname{ccint}(v, \mathit{bv})) \wedge \\ \quad \quad \quad \forall \mathit{uv} \in w(\mathit{kv}), \mathit{pv}, \mathit{qv} \in \mathit{uv}. \\ \quad \quad \quad \quad (\mathit{pv} \in \mathbb{R}) \wedge \\ \quad \quad \quad \quad (\mathit{qv} \in \mathbb{R}) \wedge \\ \quad \quad \quad \quad (\lvert (\mathit{pv} - \mathit{qv}) \rvert \leq d(\mathit{kv})) \end{array}}}

\vnbfile{theorem-library/rake-bolzano-weierstrass-2.scm}{theorem-library/rake-bolzano-weierstrass-2.scm}
\noindent{\small NOTE 2026-09-20 (batch 12-B).  Where this file says the tree has no metric SUBSPACE structure, that is history: SUBSPACE-MS(s, A) and RESTRICT(f, A) are defined in structure-library/metric-subspace.scm, their laws are in theorem-library/metric-subspace-laws.scm and compact-subspace.scm, and}\par
\vnbentry{\texttt{rr-\allowbreak{}interval-\allowbreak{}seq-\allowbreak{}has-\allowbreak{}convergent-\allowbreak{}subseq}}{\vnbstmt{\begin{array}{@{}l@{}} \forall v, \mathit{hiv} \in \mathbb{R}, \mathit{hqv} \in (\mathbb{N} \to \operatorname{ccint}(v, \mathit{hiv})). \exists \mathit{phv}. \\ \quad \operatorname{strictly\text{-}mono\text{-}nn}(\mathit{phv}) \wedge \\ \quad \exists \mathit{ptv} \in \mathbb{R}. \\ \quad \quad \operatorname{converges\text{-}to}(\operatorname{rr\text{-}ms}, \operatorname{subseq}(\mathit{hqv}, \mathit{phv}), \mathit{ptv}) \end{array}}}
\vnbentry{\texttt{rr-\allowbreak{}bolzano-\allowbreak{}weierstrass}}{\vnbstmt{\begin{array}{@{}l@{}} \forall v \in (\mathbb{N} \to \mathbb{R}), \mathit{bdv} \in \mathbb{R}. \\ \quad \forall \mathit{iqv} \in \mathbb{N}.\; (\lvert v(\mathit{iqv}) \rvert \leq \mathit{bdv}) \Rightarrow  \\ \quad \quad \exists \mathit{phv}. \\ \quad \quad \quad \operatorname{strictly\text{-}mono\text{-}nn}(\mathit{phv}) \wedge \\ \quad \quad \quad \exists \mathit{ptv} \in \mathbb{R}. \\ \quad \quad \quad \quad \operatorname{converges\text{-}to}(\operatorname{rr\text{-}ms}, \operatorname{subseq}(v, \mathit{phv}), \mathit{ptv}) \end{array}}}
\vnbentry{\texttt{bounded-\allowbreak{}block-\allowbreak{}converges}}{\vnbstmt{\begin{array}{@{}l@{}} \forall v \in (\mathbb{N} \to \mathbb{R}), \mathit{bdv} \in \mathbb{R}, \mathit{jqv} \in \mathrm{Inf}(\mathbb{N}). \\ \quad \forall \mathit{iqv} \in \mathbb{N}.\; (\lvert v(\mathit{iqv}) \rvert \leq \mathit{bdv}) \Rightarrow  \\ \quad \quad \exists \mathit{bqv} \in \mathrm{Inf}(\mathbb{N}). \\ \quad \quad \quad (\mathit{bqv} \subseteq \mathit{jqv}) \wedge \\ \quad \quad \quad \exists \mathit{ptv} \in \mathbb{R}. \\ \quad \quad \quad \quad \operatorname{converges\text{-}along}(\operatorname{rr\text{-}ms}, v, \mathit{bqv}, \mathit{ptv}) \end{array}}}

\vnbfile{theorem-library/rake-ascoli-diagonal.scm}{theorem-library/rake-ascoli-diagonal.scm}
\noindent{\small theorem-library/rake-ascoli-diagonal.scm ==================================================================== ASCOLI-ARZELA, the last stretch -- batch 11-F.  Two theorems:}\par
\vnbentry{\texttt{ascoli-\allowbreak{}pointwise-\allowbreak{}diagonal}}{\vnbstmt{\begin{array}{@{}l@{}} \forall s, m, q. \\ \quad \operatorname{pointwise\text{-}bounded}(s, m) \wedge \\ \quad \operatorname{is\text{-}dense\text{-}seq}(s, q) \Rightarrow \\ \quad \quad \exists l. \\ \quad \quad \quad \operatorname{strictly\text{-}mono\text{-}nn}(l) \wedge \\ \quad \quad \quad \operatorname{converges\text{-}on}(s, \operatorname{subseq}(m, l), q) \end{array}}}
\vnbentry{\texttt{ascoli-\allowbreak{}arzela-\allowbreak{}sequential}}{\vnbstmt{\begin{array}{@{}l@{}} \forall s, m. \\ \quad \operatorname{is\text{-}compact}(s) \wedge \\ \quad \exists x.\; (x \in \operatorname{pts}(s)) \wedge \\ \quad (m \in (\mathbb{N} \to (\operatorname{pts}(s) \to \mathbb{R}))) \wedge \\ \quad \forall k \in \mathbb{N}. \\ \quad \quad \operatorname{is\text{-}continuous}(s, \operatorname{rr\text{-}ms}, m(k)) \wedge \\ \quad \operatorname{is\text{-}equicontinuous}(s, \operatorname{rr\text{-}ms}, m) \wedge \\ \quad \operatorname{pointwise\text{-}bounded}(s, m) \Rightarrow \\ \quad \quad \exists l. \\ \quad \quad \quad \operatorname{strictly\text{-}mono\text{-}nn}(l) \wedge \\ \quad \quad \quad \exists g \in (\operatorname{pts}(s) \to \mathbb{R}). \\ \quad \quad \quad \quad \operatorname{is\text{-}continuous}(s, \operatorname{rr\text{-}ms}, g) \wedge \\ \quad \quad \quad \quad \operatorname{converges\text{-}uniformly}(s, \operatorname{subseq}(m, l), g) \end{array}}}

\vnbfile{theorem-library/bernstein-antiderivable.scm}{theorem-library/bernstein-antiderivable.scm}
\noindent{\small bernstein-antiderivable.scm -- EVERY BERNSTEIN APPROXIMANT IS ANTIDERIVABLE.}\par
\vnbentry{\texttt{bernstein-\allowbreak{}basis-\allowbreak{}degree-\allowbreak{}zero-\allowbreak{}diag}}{\vnbstmt{\begin{array}{@{}l@{}} \forall x, k. \\ \quad ((x \in \mathbb{R}) \wedge ((k \in \mathbb{N}) \wedge (k = 0))) \Rightarrow  \\ \quad \quad (\operatorname{bernstein\text{-}basis}(0, x)(k) = 1) \end{array}}}
\vnbentry{\texttt{bernstein-\allowbreak{}basis-\allowbreak{}degree-\allowbreak{}zero-\allowbreak{}off}}{\vnbstmt{\begin{array}{@{}l@{}} \forall x, k. \\ \quad ((x \in \mathbb{R}) \wedge ((k \in \mathbb{N}) \wedge \neg (k = 0))) \Rightarrow  \\ \quad \quad (\operatorname{bernstein\text{-}basis}(0, x)(k) = 0) \end{array}}}
\vnbentry{\texttt{bernstein-\allowbreak{}term-\allowbreak{}antiderivable}}{\vnbstmt{\begin{array}{@{}l@{}} \forall n \in \mathbb{N}, a, b, k, j. \\ \quad (a \in \mathbb{R}) \wedge \\ \quad (b \in \mathbb{R}) \wedge \\ \quad (a < b) \wedge \\ \quad (k \in \mathbb{N}) \wedge \\ \quad (j \in \mathbb{N}) \Rightarrow \\ \quad \quad \operatorname{is\text{-}antiderivable}((\lambda x \in \mathbb{R}.\; \\ \quad \quad \quad \quad (\operatorname{power}(x, j) \cdot \operatorname{bernstein\text{-}basis}(n, x)(k))), \\ \quad \quad \quad a, b) \end{array}}}
\vnbentry{\texttt{bernstein-\allowbreak{}basis-\allowbreak{}antiderivable}}{\vnbstmt{\begin{array}{@{}l@{}} \forall n, k, a, b. \\ \quad (n \in \mathbb{N}) \wedge \\ \quad (k \in \mathbb{N}) \wedge \\ \quad (a \in \mathbb{R}) \wedge \\ \quad (b \in \mathbb{R}) \wedge \\ \quad (a < b) \Rightarrow \\ \quad \quad \operatorname{is\text{-}antiderivable}((\lambda x \in \mathbb{R}.\; \operatorname{bernstein\text{-}basis}(n, x)(k)), a, b) \end{array}}}
\vnbentry{\texttt{bernstein-\allowbreak{}basis-\allowbreak{}scaled-\allowbreak{}antiderivable}}{\vnbstmt{\begin{array}{@{}l@{}} \forall c, n, k, a, b. \\ \quad (c \in \mathbb{R}) \wedge \\ \quad (n \in \mathbb{N}) \wedge \\ \quad (k \in \mathbb{N}) \wedge \\ \quad (a \in \mathbb{R}) \wedge \\ \quad (b \in \mathbb{R}) \wedge \\ \quad (a < b) \Rightarrow \\ \quad \quad \operatorname{is\text{-}antiderivable}((\lambda x \in \mathbb{R}.\; (c \cdot \operatorname{bernstein\text{-}basis}(n, x)(k))), a, b) \end{array}}}
\vnbentry{\texttt{bernstein-\allowbreak{}poly-\allowbreak{}antiderivable}}{\vnbstmt{\begin{array}{@{}l@{}} \forall f, n, a, b. \\ \quad (f \in (\mathbb{R} \to \mathbb{R})) \wedge \\ \quad (n \in \mathbb{N}) \wedge \\ \quad \neg (n = 0) \wedge \\ \quad (a \in \mathbb{R}) \wedge \\ \quad (b \in \mathbb{R}) \wedge \\ \quad (a < b) \Rightarrow \\ \quad \quad \operatorname{is\text{-}antiderivable}((\lambda x \in \mathbb{R}.\; \operatorname{bernstein\text{-}poly}(f, n, x)), a, b) \end{array}}}

\vnbfile{theorem-library/antiderivable-affine-subst.scm}{theorem-library/antiderivable-affine-subst.scm}
\noindent{\small antiderivable-affine-subst.scm -- ANTIDERIVABILITY SURVIVES AN AFFINE CHANGE OF VARIABLE, and with it the rung-3 form of the Bernstein approximants.}\par
\vnbentry{\texttt{antiderivable-\allowbreak{}affine-\allowbreak{}subst}}{\vnbstmt{\begin{array}{@{}l@{}} \forall i, k, m, p, q, a, b. \\ \quad (i \in (\mathbb{R} \to \mathbb{R})) \wedge \\ \quad (k \in \mathbb{R}) \wedge \\ \quad (m \in \mathbb{R}) \wedge \\ \quad \neg (m = 0) \wedge \\ \quad (p \in \mathbb{R}) \wedge \\ \quad (q \in \mathbb{R}) \wedge \\ \quad (p < q) \wedge \\ \quad \operatorname{is\text{-}antiderivable}(i, a, b) \wedge \\ \quad \forall t \in \operatorname{ccint}(p, q). \\ \quad \quad ((k + (m \cdot t)) \in \operatorname{ccint}(a, b)) \wedge \\ \quad \forall t \in \mathbb{R}. \\ \quad \quad ((p < t) \wedge (t < q)) \Rightarrow  \\ \quad \quad \quad ((a < (k + (m \cdot t))) \wedge ((k + (m \cdot t)) < b)) \Rightarrow \\ \quad \quad \operatorname{is\text{-}antiderivable}((\lambda x \in \mathbb{R}.\; i((k + (m \cdot x)))), p, q) \end{array}}}
\vnbentry{\texttt{antiderivable-\allowbreak{}affine-\allowbreak{}subst-\allowbreak{}pos}}{\vnbstmt{\begin{array}{@{}l@{}} \forall i, k, m, p, q. \\ \quad (i \in (\mathbb{R} \to \mathbb{R})) \wedge \\ \quad (k \in \mathbb{R}) \wedge \\ \quad (m \in \mathbb{R}) \wedge \\ \quad (0 < m) \wedge \\ \quad (p \in \mathbb{R}) \wedge \\ \quad (q \in \mathbb{R}) \wedge \\ \quad (p < q) \wedge \\ \quad \operatorname{is\text{-}antiderivable}(i, (k + (m \cdot p)), (k + (m \cdot q))) \Rightarrow \\ \quad \quad \operatorname{is\text{-}antiderivable}((\lambda x \in \mathbb{R}.\; i((k + (m \cdot x)))), p, q) \end{array}}}
\vnbentry{\texttt{antiderivable-\allowbreak{}affine-\allowbreak{}subst-\allowbreak{}pos-\allowbreak{}transfer}}{\vnbstmt{\begin{array}{@{}l@{}} \forall i, a, k, m, p, q. \\ \quad (i \in (\mathbb{R} \to \mathbb{R})) \wedge \\ \quad (a \in (\mathbb{R} \to \mathbb{R})) \wedge \\ \quad (k \in \mathbb{R}) \wedge \\ \quad (m \in \mathbb{R}) \wedge \\ \quad (0 < m) \wedge \\ \quad (p \in \mathbb{R}) \wedge \\ \quad (q \in \mathbb{R}) \wedge \\ \quad (p < q) \wedge \\ \quad \operatorname{is\text{-}antiderivable}(i, (k + (m \cdot p)), (k + (m \cdot q))) \wedge \\ \quad \forall x \in \mathbb{R}.\; (a(x) \simeq i((k + (m \cdot x)))) \Rightarrow \\ \quad \quad \operatorname{is\text{-}antiderivable}(a, p, q) \end{array}}}
\vnbentry{\texttt{bernstein-\allowbreak{}poly-\allowbreak{}unit-\allowbreak{}coordinate-\allowbreak{}antiderivable}}{\vnbstmt{\begin{array}{@{}l@{}} \forall g, n, a, b. \\ \quad (g \in (\mathbb{R} \to \mathbb{R})) \wedge \\ \quad (n \in \mathbb{N}) \wedge \\ \quad \neg (n = 0) \wedge \\ \quad (a \in \mathbb{R}) \wedge \\ \quad (b \in \mathbb{R}) \wedge \\ \quad (a < b) \Rightarrow \\ \quad \quad \operatorname{is\text{-}antiderivable}((\lambda x \in \mathbb{R}.\; \\ \quad \quad \quad \quad \operatorname{bernstein\text{-}poly}(g, n, ((x - a) \cdot \operatorname{recip}((b - a))))), \\ \quad \quad \quad a, b) \end{array}}}
\vnbentry{\texttt{bernstein-\allowbreak{}ccint-\allowbreak{}approximant-\allowbreak{}antiderivable}}{\vnbstmt{\begin{array}{@{}l@{}} \forall f, n, a, b. \\ \quad (f \in (\mathbb{R} \to \mathbb{R})) \wedge \\ \quad (n \in \mathbb{N}) \wedge \\ \quad \neg (n = 0) \wedge \\ \quad (a \in \mathbb{R}) \wedge \\ \quad (b \in \mathbb{R}) \wedge \\ \quad (a < b) \Rightarrow \\ \quad \quad \operatorname{is\text{-}antiderivable}((\lambda x \in \mathbb{R}.\; \\ \quad \quad \quad \quad \operatorname{bernstein\text{-}poly}(\operatorname{compose}(f, (\lambda c \in \mathbb{R}.\; (a + ((b - a) \cdot c)))), n, \\ \quad \quad \quad \quad \quad ((x - a) \cdot \operatorname{recip}((b - a))))), \\ \quad \quad \quad a, b) \end{array}}}

\vnbfile{theorem-library/continuous-antiderivable.scm}{theorem-library/continuous-antiderivable.scm}
\noindent{\small continuous-antiderivable.scm -- RUNG 3, AND THE INTEGRAL THAT FOLLOWS.}\par
\vnbentry{\texttt{continuous-\allowbreak{}is-\allowbreak{}antiderivable}}{\vnbstmt{\begin{array}{@{}l@{}} \forall f, a, b. \\ \quad (f \in (\mathbb{R} \to \mathbb{R})) \wedge \\ \quad (a \in \mathbb{R}) \wedge \\ \quad (b \in \mathbb{R}) \wedge \\ \quad (a < b) \wedge \\ \quad \forall x \in \operatorname{ccint}(a, b). \\ \quad \quad \operatorname{is\text{-}continuous\text{-}at}(\operatorname{rr\text{-}ms}, \operatorname{rr\text{-}ms}, f, x) \Rightarrow \\ \quad \quad \operatorname{is\text{-}antiderivable}(f, a, b) \end{array}}}
\vnbentry{\texttt{c-\allowbreak{}int-\allowbreak{}defined-\allowbreak{}for-\allowbreak{}continuous}}{\vnbstmt{\begin{array}{@{}l@{}} \forall f, a, b. \\ \quad (f \in (\mathbb{R} \to \mathbb{R})) \wedge \\ \quad (a \in \mathbb{R}) \wedge \\ \quad (b \in \mathbb{R}) \wedge \\ \quad (a < b) \wedge \\ \quad \forall x \in \operatorname{ccint}(a, b). \\ \quad \quad \operatorname{is\text{-}continuous\text{-}at}(\operatorname{rr\text{-}ms}, \operatorname{rr\text{-}ms}, f, x) \Rightarrow \\ \quad \quad (\int_{a}^{b} f \in \mathbb{R}) \end{array}}}
\vnbentry{\texttt{c-\allowbreak{}int-\allowbreak{}value-\allowbreak{}continuous}}{\vnbstmt{\begin{array}{@{}l@{}} \forall f, a, b. \\ \quad (f \in (\mathbb{R} \to \mathbb{R})) \wedge \\ \quad (a \in \mathbb{R}) \wedge \\ \quad (b \in \mathbb{R}) \wedge \\ \quad (a < b) \wedge \\ \quad \forall x \in \operatorname{ccint}(a, b). \\ \quad \quad \operatorname{is\text{-}continuous\text{-}at}(\operatorname{rr\text{-}ms}, \operatorname{rr\text{-}ms}, f, x) \Rightarrow \\ \quad \quad \exists g. \\ \quad \quad \quad (\operatorname{is\text{-}antiderivative}(g, f, a, b) \wedge (\int_{a}^{b} f = (g(b) - g(a)))) \end{array}}}
\vnbentry{\texttt{c-\allowbreak{}int-\allowbreak{}add-\allowbreak{}continuous}}{\vnbstmt{\begin{array}{@{}l@{}} \forall i, c, a, b. \\ \quad (i \in (\mathbb{R} \to \mathbb{R})) \wedge \\ \quad (c \in (\mathbb{R} \to \mathbb{R})) \wedge \\ \quad (a \in \mathbb{R}) \wedge \\ \quad (b \in \mathbb{R}) \wedge \\ \quad (a < b) \wedge \\ \quad \forall d \in \operatorname{ccint}(a, b). \\ \quad \quad \operatorname{is\text{-}continuous\text{-}at}(\operatorname{rr\text{-}ms}, \operatorname{rr\text{-}ms}, i, d) \wedge \\ \quad \forall d \in \operatorname{ccint}(a, b). \\ \quad \quad \operatorname{is\text{-}continuous\text{-}at}(\operatorname{rr\text{-}ms}, \operatorname{rr\text{-}ms}, c, d) \Rightarrow \\ \quad \quad (\int_{a}^{b} (\lambda x \in \mathbb{R}.\; (i(x) + c(x))) = (\int_{a}^{b} i + \int_{a}^{b} c)) \end{array}}}
\vnbentry{\texttt{c-\allowbreak{}int-\allowbreak{}scale-\allowbreak{}continuous}}{\vnbstmt{\begin{array}{@{}l@{}} \forall c, i, a, b. \\ \quad (c \in \mathbb{R}) \wedge \\ \quad (i \in (\mathbb{R} \to \mathbb{R})) \wedge \\ \quad (a \in \mathbb{R}) \wedge \\ \quad (b \in \mathbb{R}) \wedge \\ \quad (a < b) \wedge \\ \quad \forall d \in \operatorname{ccint}(a, b). \\ \quad \quad \operatorname{is\text{-}continuous\text{-}at}(\operatorname{rr\text{-}ms}, \operatorname{rr\text{-}ms}, i, d) \Rightarrow \\ \quad \quad (\int_{a}^{b} (\lambda x \in \mathbb{R}.\; (c \cdot i(x))) = (c \cdot \int_{a}^{b} i)) \end{array}}}

\vnbfile{theorem-library/c-int-change-of-variable.scm}{theorem-library/c-int-change-of-variable.scm}
\noindent{\small c-int-change-of-variable.scm -- CHANGE OF VARIABLES FOR `C-INT', in the INCREASING case, docs/calculus.pdf's substitution rule:}\par
\vnbentry{\texttt{change-\allowbreak{}of-\allowbreak{}variable-\allowbreak{}antiderivative}}{\vnbstmt{\begin{array}{@{}l@{}} \forall w, i, p, r, a, q. \\ \quad (p \in (\mathbb{R} \to \mathbb{R})) \wedge \\ \quad (r \in (\mathbb{R} \to \mathbb{R})) \wedge \\ \quad (a \in \mathbb{R}) \wedge \\ \quad (q \in \mathbb{R}) \wedge \\ \quad (a < q) \wedge \\ \quad \operatorname{is\text{-}antiderivative}(w, i, p(a), p(q)) \wedge \\ \quad \forall u \in \operatorname{ccint}(a, q). \\ \quad \quad \operatorname{is\text{-}continuous\text{-}at}(\operatorname{rr\text{-}ms}, \operatorname{rr\text{-}ms}, p, u) \wedge \\ \quad \forall b \in \mathbb{R}. \\ \quad \quad ((a < b) \wedge (b < q)) \Rightarrow  \\ \quad \quad \quad \operatorname{is\text{-}diff\text{-}at}(p, b, r(b)) \wedge \\ \quad \forall u \in \operatorname{ccint}(a, q). \\ \quad \quad (p(u) \in \operatorname{ccint}(p(a), p(q))) \wedge \\ \quad \forall b \in \mathbb{R}. \\ \quad \quad (((a < b) \wedge (b < q)) \Rightarrow ((p(a) < p(b)) \wedge (p(b) < p(q)))) \Rightarrow \\ \quad \quad \operatorname{is\text{-}antiderivative}(\operatorname{compose}(w, p), (\lambda x \in \mathbb{R}.\; (i(p(x)) \cdot r(x))), a, q) \end{array}}}
\vnbentry{\texttt{change-\allowbreak{}of-\allowbreak{}variable-\allowbreak{}antiderivable}}{\vnbstmt{\begin{array}{@{}l@{}} \forall i, p, r, a, q. \\ \quad (p \in (\mathbb{R} \to \mathbb{R})) \wedge \\ \quad (r \in (\mathbb{R} \to \mathbb{R})) \wedge \\ \quad (a \in \mathbb{R}) \wedge \\ \quad (q \in \mathbb{R}) \wedge \\ \quad (a < q) \wedge \\ \quad \operatorname{is\text{-}antiderivable}(i, p(a), p(q)) \wedge \\ \quad \forall u \in \operatorname{ccint}(a, q). \\ \quad \quad \operatorname{is\text{-}continuous\text{-}at}(\operatorname{rr\text{-}ms}, \operatorname{rr\text{-}ms}, p, u) \wedge \\ \quad \forall w \in \mathbb{R}. \\ \quad \quad ((a < w) \wedge (w < q)) \Rightarrow  \\ \quad \quad \quad \operatorname{is\text{-}diff\text{-}at}(p, w, r(w)) \wedge \\ \quad \forall u \in \operatorname{ccint}(a, q). \\ \quad \quad (p(u) \in \operatorname{ccint}(p(a), p(q))) \wedge \\ \quad \forall w \in \mathbb{R}. \\ \quad \quad (((a < w) \wedge (w < q)) \Rightarrow ((p(a) < p(w)) \wedge (p(w) < p(q)))) \Rightarrow \\ \quad \quad \operatorname{is\text{-}antiderivable}((\lambda x \in \mathbb{R}.\; (i(p(x)) \cdot r(x))), a, q) \end{array}}}
\vnbentry{\texttt{c-\allowbreak{}int-\allowbreak{}change-\allowbreak{}of-\allowbreak{}variable}}{\vnbstmt{\begin{array}{@{}l@{}} \forall i, p, r, a, q. \\ \quad (p \in (\mathbb{R} \to \mathbb{R})) \wedge \\ \quad (r \in (\mathbb{R} \to \mathbb{R})) \wedge \\ \quad (a \in \mathbb{R}) \wedge \\ \quad (q \in \mathbb{R}) \wedge \\ \quad (a < q) \wedge \\ \quad \operatorname{is\text{-}antiderivable}(i, p(a), p(q)) \wedge \\ \quad \forall u \in \operatorname{ccint}(a, q). \\ \quad \quad \operatorname{is\text{-}continuous\text{-}at}(\operatorname{rr\text{-}ms}, \operatorname{rr\text{-}ms}, p, u) \wedge \\ \quad \forall w \in \mathbb{R}. \\ \quad \quad ((a < w) \wedge (w < q)) \Rightarrow  \\ \quad \quad \quad \operatorname{is\text{-}diff\text{-}at}(p, w, r(w)) \wedge \\ \quad \forall u \in \operatorname{ccint}(a, q). \\ \quad \quad (p(u) \in \operatorname{ccint}(p(a), p(q))) \wedge \\ \quad \forall w \in \mathbb{R}. \\ \quad \quad (((a < w) \wedge (w < q)) \Rightarrow ((p(a) < p(w)) \wedge (p(w) < p(q)))) \Rightarrow \\ \quad \quad (\int_{a}^{q} (\lambda x \in \mathbb{R}.\; (i(p(x)) \cdot r(x))) = \int_{p(a)}^{p(q)} i) \end{array}}}
\vnbentry{\texttt{c-\allowbreak{}int-\allowbreak{}change-\allowbreak{}of-\allowbreak{}variable-\allowbreak{}transfer}}{\vnbstmt{\begin{array}{@{}l@{}} \forall i, a, p, r, b, q. \\ \quad (p \in (\mathbb{R} \to \mathbb{R})) \wedge \\ \quad (r \in (\mathbb{R} \to \mathbb{R})) \wedge \\ \quad (b \in \mathbb{R}) \wedge \\ \quad (q \in \mathbb{R}) \wedge \\ \quad (b < q) \wedge \\ \quad (a \in (\mathbb{R} \to \mathbb{R})) \wedge \\ \quad \operatorname{is\text{-}antiderivable}(i, p(b), p(q)) \wedge \\ \quad \forall u \in \operatorname{ccint}(b, q). \\ \quad \quad \operatorname{is\text{-}continuous\text{-}at}(\operatorname{rr\text{-}ms}, \operatorname{rr\text{-}ms}, p, u) \wedge \\ \quad \forall w \in \mathbb{R}. \\ \quad \quad ((b < w) \wedge (w < q)) \Rightarrow  \\ \quad \quad \quad \operatorname{is\text{-}diff\text{-}at}(p, w, r(w)) \wedge \\ \quad \forall u \in \operatorname{ccint}(b, q). \\ \quad \quad (p(u) \in \operatorname{ccint}(p(b), p(q))) \wedge \\ \quad \forall w \in \mathbb{R}. \\ \quad \quad (((b < w) \wedge (w < q)) \Rightarrow ((p(b) < p(w)) \wedge (p(w) < p(q)))) \wedge \\ \quad \forall x \in \mathbb{R}.\; (a(x) \simeq (i(p(x)) \cdot r(x))) \Rightarrow \\ \quad \quad (\int_{b}^{q} a = \int_{p(b)}^{p(q)} i) \end{array}}}

\vnbfile{theorem-library/log.scm}{theorem-library/log.scm}
\noindent{\small theorem-library/log.scm -- THE LOGARITHM, DEFINED AS AN INTEGRAL, and its basic laws.  The one that costs anything is the functional equation, and what it costs is a CHANGE OF VARIABLES.}\par
\vnbentry{\texttt{log-\allowbreak{}unfold}}{\vnbstmt{\begin{array}{@{}l@{}} \forall x. \\ \quad (\log\left(x\right) \simeq \int_{1}^{x} (\lambda z \in \mathbb{R}.\; {z}^{\ast\!-1})) \end{array}}}
\vnbentry{\texttt{log-\allowbreak{}in-\allowbreak{}rr}}{\vnbstmt{\forall x.\; (\log\left(x\right) \in \mathbb{R})}}
\vnbentry{\texttt{log-\allowbreak{}lam-\allowbreak{}in-\allowbreak{}fun}}{\vnbstmt{((\lambda x \in \mathbb{R}.\; \log\left(x\right)) \in (\mathbb{R} \to \mathbb{R}))}}
\vnbentry{\texttt{log-\allowbreak{}one}}{\vnbstmt{(\log\left(1\right) = 0)}}
\vnbentry{\texttt{recip-\allowbreak{}star-\allowbreak{}antiderivable-\allowbreak{}pos}}{\vnbstmt{\begin{array}{@{}l@{}} \forall a, b. \\ \quad ((a \in \mathbb{R}) \wedge ((b \in \mathbb{R}) \wedge ((0 < a) \wedge (a < b)))) \Rightarrow  \\ \quad \quad \operatorname{is\text{-}antiderivable}((\lambda z \in \mathbb{R}.\; {z}^{\ast\!-1}), a, b) \end{array}}}
\vnbentry{\texttt{recip-\allowbreak{}star-\allowbreak{}scale}}{\vnbstmt{\begin{array}{@{}l@{}} \forall m, u. \\ \quad ((m \in \mathbb{R}) \wedge ((u \in \mathbb{R}) \wedge \neg (m = 0))) \Rightarrow  \\ \quad \quad (({(m \cdot u)}^{\ast\!-1} \cdot m) = {u}^{\ast\!-1}) \end{array}}}
\vnbentry{\texttt{c-\allowbreak{}int-\allowbreak{}or-\allowbreak{}scale-\allowbreak{}ordered}}{\vnbstmt{\begin{array}{@{}l@{}} \forall m, a, b. \\ \quad (m \in \mathbb{R}) \wedge \\ \quad (a \in \mathbb{R}) \wedge \\ \quad (b \in \mathbb{R}) \wedge \\ \quad (0 < m) \wedge \\ \quad (0 < a) \wedge \\ \quad (a < b) \Rightarrow \\ \quad \quad \int_{(m \cdot a)}^{(m \cdot b)} (\lambda z \in \mathbb{R}.\; {z}^{\ast\!-1}) \\ \quad \quad =\; \int_{a}^{b} (\lambda z \in \mathbb{R}.\; {z}^{\ast\!-1}) \end{array}}}
\vnbentry{\texttt{c-\allowbreak{}int-\allowbreak{}or-\allowbreak{}scale}}{\vnbstmt{\begin{array}{@{}l@{}} \forall m, a, b. \\ \quad (m \in \mathbb{R}) \wedge \\ \quad (a \in \mathbb{R}) \wedge \\ \quad (b \in \mathbb{R}) \wedge \\ \quad (0 < m) \wedge \\ \quad (0 < a) \wedge \\ \quad (0 < b) \Rightarrow \\ \quad \quad \int_{(m \cdot a)}^{(m \cdot b)} (\lambda z \in \mathbb{R}.\; {z}^{\ast\!-1}) \\ \quad \quad =\; \int_{a}^{b} (\lambda z \in \mathbb{R}.\; {z}^{\ast\!-1}) \end{array}}}
\vnbentry{\texttt{recip-\allowbreak{}star-\allowbreak{}int-\allowbreak{}additive}}{\vnbstmt{\begin{array}{@{}l@{}} \forall u, v, w. \\ \quad (u \in \mathbb{R}) \wedge \\ \quad (v \in \mathbb{R}) \wedge \\ \quad (w \in \mathbb{R}) \wedge \\ \quad (0 < u) \wedge \\ \quad (0 < v) \wedge \\ \quad (0 < w) \Rightarrow \\ \quad \quad \int_{u}^{w} (\lambda z \in \mathbb{R}.\; {z}^{\ast\!-1}) \\ \quad \quad =\; (\int_{u}^{v} (\lambda z \in \mathbb{R}.\; {z}^{\ast\!-1}) + \int_{v}^{w} (\lambda z \in \mathbb{R}.\; {z}^{\ast\!-1})) \end{array}}}
\vnbentry{\texttt{log-\allowbreak{}diff}}{\vnbstmt{\begin{array}{@{}l@{}} \forall u, v. \\ \quad ((u \in \mathbb{R}) \wedge ((v \in \mathbb{R}) \wedge ((0 < u) \wedge (0 < v)))) \Rightarrow  \\ \quad \quad \int_{u}^{v} (\lambda z \in \mathbb{R}.\; {z}^{\ast\!-1}) \\ \quad \quad =\; (\log\left(v\right) - \log\left(u\right)) \end{array}}}
\vnbentry{\texttt{log-\allowbreak{}mul}}{\vnbstmt{\begin{array}{@{}l@{}} \forall x, y. \\ \quad ((x \in \mathbb{R}) \wedge ((y \in \mathbb{R}) \wedge ((0 < x) \wedge (0 < y)))) \Rightarrow  \\ \quad \quad (\log\left((x \cdot y)\right) = (\log\left(x\right) + \log\left(y\right))) \end{array}}}
\vnbentry{\texttt{log-\allowbreak{}strictly-\allowbreak{}increasing}}{\vnbstmt{\begin{array}{@{}l@{}} \forall u, v. \\ \quad ((u \in \mathbb{R}) \wedge ((v \in \mathbb{R}) \wedge ((0 < u) \wedge (u < v)))) \Rightarrow  \\ \quad \quad (\log\left(u\right) < \log\left(v\right)) \end{array}}}
\vnbentry{\texttt{log-\allowbreak{}neg-\allowbreak{}below-\allowbreak{}one}}{\vnbstmt{\begin{array}{@{}l@{}} \forall x \in \mathbb{R}. \\ \quad (((0 < x) \wedge (x < 1)) \Rightarrow (\log\left(x\right) < 0)) \end{array}}}
\vnbentry{\texttt{log-\allowbreak{}pos-\allowbreak{}above-\allowbreak{}one}}{\vnbstmt{\forall x \in \mathbb{R}.\; ((1 < x) \Rightarrow (0 < \log\left(x\right)))}}
\vnbentry{\texttt{log-\allowbreak{}recip}}{\vnbstmt{\begin{array}{@{}l@{}} \forall x \in \mathbb{R}. \\ \quad ((0 < x) \Rightarrow (\log\left({x}^{\ast\!-1}\right) = -\log\left(x\right))) \end{array}}}
\vnbentry{\texttt{log-\allowbreak{}deriv}}{\vnbstmt{\begin{array}{@{}l@{}} \forall c \in \mathbb{R}. \\ \quad (0 < c) \Rightarrow  \\ \quad \quad \operatorname{is\text{-}diff\text{-}at}((\lambda x \in \mathbb{R}.\; \log\left(x\right)), c, {c}^{\ast\!-1}) \end{array}}}

\vnbfile{theorem-library/r-exp.scm}{theorem-library/r-exp.scm}
\noindent{\small theorem-library/r-exp.scm -- THE REAL EXPONENTIAL, DEFINED AS THE INVERSE OF THE LOGARITHM, and its laws.}\par
\vnbentry{\texttt{log-\allowbreak{}two-\allowbreak{}pos}}{\vnbstmt{(0 < \log\left(2\right))}}
\vnbentry{\texttt{log-\allowbreak{}unbounded-\allowbreak{}above}}{\vnbstmt{\begin{array}{@{}l@{}} \forall n \in \mathbb{N}. \exists b \in \mathbb{R}. \\ \quad ((0 < b) \wedge ((n \cdot \log\left(2\right)) \leq \log\left(b\right))) \end{array}}}
\vnbentry{\texttt{log-\allowbreak{}above-\allowbreak{}any}}{\vnbstmt{\begin{array}{@{}l@{}} \forall y \in \mathbb{R}. \exists b \in \mathbb{R}. \\ \quad ((0 < b) \wedge (y \leq \log\left(b\right))) \end{array}}}
\vnbentry{\texttt{log-\allowbreak{}below-\allowbreak{}any}}{\vnbstmt{\begin{array}{@{}l@{}} \forall y \in \mathbb{R}. \exists a \in \mathbb{R}. \\ \quad ((0 < a) \wedge (\log\left(a\right) \leq y)) \end{array}}}
\vnbentry{\texttt{log-\allowbreak{}surjective}}{\vnbstmt{\begin{array}{@{}l@{}} \forall y \in \mathbb{R}. \exists x \in \mathbb{R}. \\ \quad ((0 < x) \wedge (\log\left(x\right) = y)) \end{array}}}
\vnbentry{\texttt{log-\allowbreak{}injective-\allowbreak{}pos}}{\vnbstmt{\begin{array}{@{}l@{}} \forall u, v. \\ \quad (u \in \mathbb{R}) \wedge \\ \quad (v \in \mathbb{R}) \wedge \\ \quad (0 < u) \wedge \\ \quad (0 < v) \wedge \\ \quad (\log\left(u\right) = \log\left(v\right)) \Rightarrow \\ \quad \quad (u = v) \end{array}}}
\vnbentry{\texttt{r-\allowbreak{}exp-\allowbreak{}prop}}{\vnbstmt{\begin{array}{@{}l@{}} \forall y \in \mathbb{R}. \\ \quad (\exp\left(y\right) \in \mathbb{R}) \wedge \\ \quad (0 < \exp\left(y\right)) \wedge \\ \quad (\log\left(\exp\left(y\right)\right) = y) \end{array}}}
\vnbentry{\texttt{r-\allowbreak{}exp-\allowbreak{}in-\allowbreak{}rr}}{\vnbstmt{\forall y \in \mathbb{R}.\; (\exp\left(y\right) \in \mathbb{R})}}
\vnbentry{\texttt{r-\allowbreak{}exp-\allowbreak{}pos}}{\vnbstmt{\forall y \in \mathbb{R}.\; (0 < \exp\left(y\right))}}
\vnbentry{\texttt{log-\allowbreak{}r-\allowbreak{}exp}}{\vnbstmt{\forall y \in \mathbb{R}.\; (\log\left(\exp\left(y\right)\right) = y)}}
\vnbentry{\texttt{r-\allowbreak{}exp-\allowbreak{}char}}{\vnbstmt{\begin{array}{@{}l@{}} \forall y, x. \\ \quad ((y \in \mathbb{R}) \wedge ((x \in \mathbb{R}) \wedge ((0 < x) \wedge (\log\left(x\right) = y)))) \Rightarrow  \\ \quad \quad (\exp\left(y\right) = x) \end{array}}}
\vnbentry{\texttt{r-\allowbreak{}exp-\allowbreak{}lam-\allowbreak{}in-\allowbreak{}fun}}{\vnbstmt{((\lambda y \in \mathbb{R}.\; \exp\left(y\right)) \in (\mathbb{R} \to \mathbb{R}))}}
\vnbentry{\texttt{r-\allowbreak{}exp-\allowbreak{}zero}}{\vnbstmt{(\exp\left(0\right) = 1)}}
\vnbentry{\texttt{r-\allowbreak{}exp-\allowbreak{}log}}{\vnbstmt{\begin{array}{@{}l@{}} \forall x \in \mathbb{R}. \\ \quad ((0 < x) \Rightarrow (\exp\left(\log\left(x\right)\right) = x)) \end{array}}}
\vnbentry{\texttt{r-\allowbreak{}exp-\allowbreak{}add}}{\vnbstmt{\begin{array}{@{}l@{}} \forall u, v. \\ \quad ((u \in \mathbb{R}) \wedge (v \in \mathbb{R})) \Rightarrow  \\ \quad \quad (\exp\left((u + v)\right) = (\exp\left(u\right) \cdot \exp\left(v\right))) \end{array}}}
\vnbentry{\texttt{r-\allowbreak{}exp-\allowbreak{}strictly-\allowbreak{}increasing}}{\vnbstmt{\begin{array}{@{}l@{}} \forall u, v. \\ \quad ((u \in \mathbb{R}) \wedge ((v \in \mathbb{R}) \wedge (u < v))) \Rightarrow  \\ \quad \quad (\exp\left(u\right) < \exp\left(v\right)) \end{array}}}
\vnbentry{\texttt{r-\allowbreak{}exp-\allowbreak{}mono}}{\vnbstmt{\begin{array}{@{}l@{}} \forall u, v. \\ \quad ((u \in \mathbb{R}) \wedge ((v \in \mathbb{R}) \wedge (u \leq v))) \Rightarrow  \\ \quad \quad (\exp\left(u\right) \leq \exp\left(v\right)) \end{array}}}
\vnbentry{\texttt{r-\allowbreak{}exp-\allowbreak{}neg}}{\vnbstmt{\begin{array}{@{}l@{}} \forall y \in \mathbb{R}. \\ \quad (\exp\left(-y\right) = {\exp\left(y\right)}^{\ast\!-1}) \end{array}}}
\vnbentry{\texttt{r-\allowbreak{}exp-\allowbreak{}log-\allowbreak{}not-\allowbreak{}global}}{\vnbstmt{\neg \forall x \in \mathbb{R}.\; (\exp\left(\log\left(x\right)\right) = x)}}
\vnbentry{\texttt{r-\allowbreak{}exp-\allowbreak{}continuous-\allowbreak{}at}}{\vnbstmt{\begin{array}{@{}l@{}} \forall t \in \mathbb{R}. \\ \quad \operatorname{is\text{-}continuous\text{-}at}(\operatorname{rr\text{-}ms}, \operatorname{rr\text{-}ms}, (\lambda y \in \mathbb{R}.\; \exp\left(y\right)), t) \end{array}}}
\vnbentry{\texttt{rr-\allowbreak{}recip-\allowbreak{}recip-\allowbreak{}pos}}{\vnbstmt{\begin{array}{@{}l@{}} \forall c \in \mathbb{R}. \\ \quad ((0 < c) \Rightarrow (\operatorname{recip}(\operatorname{recip}(c)) = c)) \end{array}}}
\vnbentry{\texttt{deriv-\allowbreak{}right-\allowbreak{}inverse}}{\vnbstmt{\begin{array}{@{}l@{}} \forall m, \mathit{gm}, \mathit{am}, v. \\ \quad \operatorname{is\text{-}diff\text{-}at}(m, \mathit{am}, v) \wedge \\ \quad \neg (v = 0) \wedge \\ \quad (\mathit{gm} \in (\mathbb{R} \to \mathbb{R})) \wedge \\ \quad \forall z \in \mathbb{R}.\; (m(\operatorname{gm}(z)) = z) \wedge \\ \quad (\operatorname{gm}(m(\mathit{am})) = \mathit{am}) \wedge \\ \quad \operatorname{is\text{-}continuous\text{-}at}(\operatorname{rr\text{-}ms}, \operatorname{rr\text{-}ms}, \mathit{gm}, m(\mathit{am})) \Rightarrow \\ \quad \quad \operatorname{is\text{-}diff\text{-}at}(\mathit{gm}, m(\mathit{am}), \operatorname{recip}(v)) \end{array}}}
\vnbentry{\texttt{r-\allowbreak{}exp-\allowbreak{}diff-\allowbreak{}at-\allowbreak{}log}}{\vnbstmt{\begin{array}{@{}l@{}} \forall c, w. \\ \quad ((c \in \mathbb{R}) \wedge ((w \in \mathbb{R}) \wedge ((0 < c) \wedge (w = \log\left(c\right))))) \Rightarrow  \\ \quad \quad \operatorname{is\text{-}diff\text{-}at}((\lambda y \in \mathbb{R}.\; \exp\left(y\right)), w, c) \end{array}}}
\vnbentry{\texttt{r-\allowbreak{}exp-\allowbreak{}deriv}}{\vnbstmt{\begin{array}{@{}l@{}} \forall t \in \mathbb{R}. \\ \quad \operatorname{is\text{-}diff\text{-}at}((\lambda y \in \mathbb{R}.\; \exp\left(y\right)), t, \exp\left(t\right)) \end{array}}}

\vnbfile{theorem-library/rpow-star.scm}{theorem-library/rpow-star.scm}
\noindent{\small theorem-library/rpow-star.scm -- THE REAL POWER, DEFINED, with the base-zero conventions carried:}\par
\vnbentry{\texttt{rpow-\allowbreak{}star-\allowbreak{}unfold}}{\vnbstmt{\begin{array}{@{}l@{}} \forall x, s. \\ \quad ({x}^{s} \simeq \left(\text{if }(0 < x)\text{ then }\exp\left((s \cdot \log\left(x\right))\right)\text{ else }\left(\text{if }(0 < s)\text{ then }0\text{ else }1\right)\right)) \end{array}}}
\vnbentry{\texttt{rpow-\allowbreak{}star-\allowbreak{}value}}{\vnbstmt{\begin{array}{@{}l@{}} \forall x.; \forall s \in \mathbb{R}. \\ \quad ((0 < x) \Rightarrow ({x}^{s} = \exp\left((s \cdot \log\left(x\right))\right))) \end{array}}}
\vnbentry{\texttt{rpow-\allowbreak{}star-\allowbreak{}zero-\allowbreak{}base}}{\vnbstmt{\forall s \in \mathbb{R}.\; ((0 < s) \Rightarrow ({0}^{s} = 0))}}
\vnbentry{\texttt{rpow-\allowbreak{}star-\allowbreak{}zero-\allowbreak{}zero}}{\vnbstmt{({0}^{0} = 1)}}
\vnbentry{\texttt{rpow-\allowbreak{}star-\allowbreak{}in-\allowbreak{}rr}}{\vnbstmt{\begin{array}{@{}l@{}} \forall x.; \forall s \in \mathbb{R}. \\ \quad ((0 < x) \Rightarrow ({x}^{s} \in \mathbb{R})) \end{array}}}
\vnbentry{\texttt{rpow-\allowbreak{}star-\allowbreak{}pos}}{\vnbstmt{\forall x.; \forall s \in \mathbb{R}.\; ((0 < x) \Rightarrow (0 < {x}^{s}))}}
\vnbentry{\texttt{log-\allowbreak{}rpow-\allowbreak{}star}}{\vnbstmt{\begin{array}{@{}l@{}} \forall x.; \forall s \in \mathbb{R}. \\ \quad ((0 < x) \Rightarrow (\log\left({x}^{s}\right) = (s \cdot \log\left(x\right)))) \end{array}}}
\vnbentry{\texttt{rpow-\allowbreak{}star-\allowbreak{}zero}}{\vnbstmt{\forall x \in \mathbb{R}.\; ((0 < x) \Rightarrow ({x}^{0} = 1))}}
\vnbentry{\texttt{rpow-\allowbreak{}star-\allowbreak{}one}}{\vnbstmt{\forall x \in \mathbb{R}.\; ((0 < x) \Rightarrow ({x}^{1} = x))}}
\vnbentry{\texttt{rpow-\allowbreak{}star-\allowbreak{}add}}{\vnbstmt{\begin{array}{@{}l@{}} \forall x, s, t. \\ \quad ((x \in \mathbb{R}) \wedge ((s \in \mathbb{R}) \wedge ((t \in \mathbb{R}) \wedge (0 < x)))) \Rightarrow  \\ \quad \quad ({x}^{(s + t)} = ({x}^{s} \cdot {x}^{t})) \end{array}}}
\vnbentry{\texttt{rpow-\allowbreak{}star-\allowbreak{}mul-\allowbreak{}base}}{\vnbstmt{\begin{array}{@{}l@{}} \forall x, y, s. \\ \quad (x \in \mathbb{R}) \wedge \\ \quad (y \in \mathbb{R}) \wedge \\ \quad (s \in \mathbb{R}) \wedge \\ \quad (0 < x) \wedge \\ \quad (0 < y) \Rightarrow \\ \quad \quad ({(x \cdot y)}^{s} = ({x}^{s} \cdot {y}^{s})) \end{array}}}
\vnbentry{\texttt{rpow-\allowbreak{}star-\allowbreak{}pow}}{\vnbstmt{\begin{array}{@{}l@{}} \forall x, s, t. \\ \quad ((x \in \mathbb{R}) \wedge ((s \in \mathbb{R}) \wedge ((t \in \mathbb{R}) \wedge (0 < x)))) \Rightarrow  \\ \quad \quad ({{x}^{s}}^{t} = {x}^{(s \cdot t)}) \end{array}}}
\vnbentry{\texttt{rpow-\allowbreak{}star-\allowbreak{}neg}}{\vnbstmt{\begin{array}{@{}l@{}} \forall x, s. \\ \quad ((x \in \mathbb{R}) \wedge ((s \in \mathbb{R}) \wedge (0 < x))) \Rightarrow  \\ \quad \quad ({x}^{-s} = {{x}^{s}}^{\ast\!-1}) \end{array}}}
\vnbentry{\texttt{rpow-\allowbreak{}star-\allowbreak{}mono-\allowbreak{}exp}}{\vnbstmt{\begin{array}{@{}l@{}} \forall x, s, t. \\ \quad (x \in \mathbb{R}) \wedge \\ \quad (s \in \mathbb{R}) \wedge \\ \quad (t \in \mathbb{R}) \wedge \\ \quad (1 < x) \wedge \\ \quad (s < t) \Rightarrow \\ \quad \quad ({x}^{s} < {x}^{t}) \end{array}}}
\vnbentry{\texttt{rpow-\allowbreak{}star-\allowbreak{}mono-\allowbreak{}base}}{\vnbstmt{\begin{array}{@{}l@{}} \forall x, y, s. \\ \quad (x \in \mathbb{R}) \wedge \\ \quad (y \in \mathbb{R}) \wedge \\ \quad (s \in \mathbb{R}) \wedge \\ \quad (0 < x) \wedge \\ \quad (x < y) \wedge \\ \quad (0 < s) \Rightarrow \\ \quad \quad ({x}^{s} < {y}^{s}) \end{array}}}

\vnbfile{theorem-library/rpow-star-convexity.scm}{theorem-library/rpow-star-convexity.scm}
\noindent{\small theorem-library/rpow-star-convexity.scm -- the two REAL inequalities of the axiomatic RPOW, proved for the DEFINED power RPOW-STAR (rpow-star.scm):}\par
\vnbentry{\texttt{r-\allowbreak{}exp-\allowbreak{}ge-\allowbreak{}one-\allowbreak{}plus}}{\vnbstmt{\forall z \in \mathbb{R}.\; ((1 + z) \leq \exp\left(z\right))}}
\vnbentry{\texttt{log-\allowbreak{}le-\allowbreak{}minus-\allowbreak{}one}}{\vnbstmt{\begin{array}{@{}l@{}} \forall w \in \mathbb{R}. \\ \quad ((0 < w) \Rightarrow (\log\left(w\right) \leq (w - 1))) \end{array}}}
\vnbentry{\texttt{log-\allowbreak{}jensen-\allowbreak{}two}}{\vnbstmt{\begin{array}{@{}l@{}} \forall u, v, s, t \in \mathbb{R}. \\ \quad ((0 < u) \wedge ((0 < v) \wedge ((0 < s) \wedge ((0 < t) \wedge ((s + t) = 1))))) \Rightarrow  \\ \quad \quad ((s \cdot \log\left(u\right)) + (t \cdot \log\left(v\right))) \\ \quad \quad \leq\; \log\left(((s \cdot u) + (t \cdot v))\right) \end{array}}}
\vnbentry{\texttt{young-\allowbreak{}inequality-\allowbreak{}star}}{\vnbstmt{\begin{array}{@{}l@{}} \forall a, b, p, q \in \mathbb{R}. \\ \quad ((0 < a) \wedge ((0 < b) \wedge ((1 < p) \wedge ((1 < q) \wedge ((/(1, p) + /(1, q)) = 1))))) \Rightarrow  \\ \quad \quad ((a \cdot b) \leq (/({a}^{p}, p) + /({b}^{q}, q))) \end{array}}}
\vnbentry{\texttt{bernoulli-\allowbreak{}rpow-\allowbreak{}star}}{\vnbstmt{\begin{array}{@{}l@{}} \forall x, b \in \mathbb{R}. \\ \quad ((0 < (1 + x)) \wedge (1 \leq b)) \Rightarrow  \\ \quad \quad ((1 + (b \cdot x)) \leq {(1 + x)}^{b}) \end{array}}}

\vnbfile{structure-library/functor-invariance.scm}{structure-library/functor-invariance.scm}
\noindent{\small functor-invariance.scm -- a construction that cannot see the functor.}\par
\vnbentry{\texttt{metric-\allowbreak{}top@preimage}}{\vnbstmt{\begin{array}{@{}l@{}} \forall s, f, v. \\ \quad (\operatorname{preimage}(\operatorname{metric\text{-}top}(s), f, v) \simeq \operatorname{preimage}(s, f, v)) \end{array}}}
\vnbentry{\texttt{trunc-\allowbreak{}metric@preimage}}{\vnbstmt{\begin{array}{@{}l@{}} \forall s, f, v. \\ \quad (\operatorname{preimage}(\operatorname{trunc\text{-}metric}(s), f, v) \simeq \operatorname{preimage}(s, f, v)) \end{array}}}

\vnbfile{theorem-library/metric-top-proof.scm}{theorem-library/metric-top-proof.scm}
\noindent{\small metric-top-proof.scm -- the metric open sets form a topology.}\par
\vnbentry{\texttt{metric-\allowbreak{}top-\allowbreak{}is-\allowbreak{}top-\allowbreak{}space}}{\vnbstmt{\begin{array}{@{}l@{}} \forall d. \\ \quad \operatorname{is\text{-}metric\text{-}space}(d) \Rightarrow  \\ \quad \quad \operatorname{is\text{-}top\text{-}space}(\operatorname{metric\text{-}top}(d)) \end{array}}}
\vnbentry{\texttt{metric-\allowbreak{}top-\allowbreak{}is-\allowbreak{}metrizable-\allowbreak{}top-\allowbreak{}space}}{\vnbstmt{\begin{array}{@{}l@{}} \forall d. \\ \quad \operatorname{is\text{-}metric\text{-}space}(d) \Rightarrow  \\ \quad \quad \operatorname{is\text{-}metrizable\text{-}top\text{-}space}(\operatorname{metric\text{-}top}(d)) \end{array}}}

\vnbfile{theorem-library/metric-top-functorial-proof.scm}{theorem-library/metric-top-functorial-proof.scm}
\noindent{\small metric-top-functorial-proof.scm -- the Met -$>$ Top functor acts on morphisms.}\par
\vnbentry{\texttt{metric-\allowbreak{}top-\allowbreak{}functorial}}{\vnbstmt{\begin{array}{@{}l@{}} \forall a, b, f. \\ \quad \operatorname{is\text{-}hom\text{-}metric\text{-}space}(a, b, f) \Rightarrow  \\ \quad \quad \operatorname{is\text{-}hom\text{-}metrizable\text{-}top\text{-}space}(\operatorname{metric\text{-}top}(a), \operatorname{metric\text{-}top}(b), f) \end{array}}}

\vnbfile{theorem-library/rake-bdd-metric.scm}{theorem-library/rake-bdd-metric.scm}
\noindent{\small rake-bdd-metric.scm -- the three asserted leaves still on the bill of `metrizable-iff-bounded-metrizable' (theorem-library/metrizable-bounded-proof.scm), PROVEN.  Statements copied VERBATIM from their support sites:}\par
\vnbentry{\texttt{metrizable-\allowbreak{}has-\allowbreak{}metric-\allowbreak{}top}}{\vnbstmt{\begin{array}{@{}l@{}} \forall s. \\ \quad \operatorname{is\text{-}metrizable\text{-}top\text{-}space}(s) \Rightarrow  \\ \quad \quad \exists d. \\ \quad \quad \quad \operatorname{is\text{-}metric\text{-}space}(d) \wedge \\ \quad \quad \quad (\operatorname{metric\text{-}top}(d) \simeq s) \end{array}}}
\vnbentry{\texttt{bdd-\allowbreak{}metric-\allowbreak{}is-\allowbreak{}bounded-\allowbreak{}metric-\allowbreak{}space}}{\vnbstmt{\begin{array}{@{}l@{}} \forall s. \\ \quad \operatorname{is\text{-}metric\text{-}space}(s) \Rightarrow  \\ \quad \quad \operatorname{is\text{-}bounded\text{-}metric\text{-}space}(\operatorname{bdd\text{-}metric}(s)) \end{array}}}
\vnbentry{\texttt{bdd-\allowbreak{}metric-\allowbreak{}open-\allowbreak{}bwd}}{\vnbstmt{\begin{array}{@{}l@{}} \forall s, u. \\ \quad \operatorname{is\text{-}metric\text{-}space}(s) \wedge \\ \quad \operatorname{is\text{-}open}(\operatorname{bdd\text{-}metric}(s), u) \Rightarrow \\ \quad \quad \operatorname{is\text{-}open}(s, u) \end{array}}}
\vnbentry{\texttt{bdd-\allowbreak{}metric-\allowbreak{}open-\allowbreak{}fwd}}{\vnbstmt{\begin{array}{@{}l@{}} \forall s, u. \\ \quad (\operatorname{is\text{-}metric\text{-}space}(s) \wedge \operatorname{is\text{-}open}(s, u)) \Rightarrow  \\ \quad \quad \operatorname{is\text{-}open}(\operatorname{bdd\text{-}metric}(s), u) \end{array}}}
\vnbentry{\texttt{bdd-\allowbreak{}metric-\allowbreak{}preserves-\allowbreak{}metric-\allowbreak{}top}}{\vnbstmt{\begin{array}{@{}l@{}} \forall s. \\ \quad \operatorname{is\text{-}metric\text{-}space}(s) \Rightarrow  \\ \quad \quad (\operatorname{metric\text{-}top}(\operatorname{bdd\text{-}metric}(s)) \simeq \operatorname{metric\text{-}top}(s)) \end{array}}}

\vnbfile{theorem-library/metrizable-compact-image.scm}{theorem-library/metrizable-compact-image.scm}
\noindent{\small metrizable-compact-image.scm -- THE CONTINUOUS IMAGE OF A COMPACT SPACE IS COMPACT, IN THE CATEGORY OF METRIZABLE TOPOLOGICAL SPACES, PROVED THROUGH SEQUENCES (the user's request, 2026-10-01: "in that category ... it can be done using sequences.  I like using sequences to prove things whenever}\par
\vnbentry{\texttt{metric-\allowbreak{}top-\allowbreak{}opens-\allowbreak{}iff}}{\vnbstmt{\begin{array}{@{}l@{}} \forall d, u. \\ \quad \operatorname{is\text{-}metric\text{-}space}(d) \Rightarrow  \\ \quad \quad ((u \in \operatorname{opens}(\operatorname{metric\text{-}top}(d))) \Leftrightarrow \operatorname{is\text{-}open}(d, u)) \end{array}}}
\vnbentry{\texttt{metric-\allowbreak{}top-\allowbreak{}compact}}{\vnbstmt{\begin{array}{@{}l@{}} \forall d. \\ \quad \operatorname{is\text{-}metric\text{-}space}(d) \Rightarrow  \\ \quad \quad (\operatorname{is\text{-}compact\text{-}t}(\operatorname{metric\text{-}top}(d)) \Leftrightarrow \operatorname{is\text{-}compact}(d)) \end{array}}}
\vnbentry{\texttt{metric-\allowbreak{}top-\allowbreak{}hom}}{\vnbstmt{\begin{array}{@{}l@{}} \forall s, t, f. \\ \quad \operatorname{is\text{-}metric\text{-}space}(s) \wedge \\ \quad \operatorname{is\text{-}metric\text{-}space}(t) \Rightarrow \\ \quad \quad (\operatorname{is\text{-}hom\text{-}top\text{-}space}(\operatorname{metric\text{-}top}(s), \operatorname{metric\text{-}top}(t), f) \Leftrightarrow \operatorname{is\text{-}continuous}(s, t, f)) \end{array}}}
\vnbentry{\texttt{seq-\allowbreak{}compact-\allowbreak{}image}}{\vnbstmt{\begin{array}{@{}l@{}} \forall d, t, f. \\ \quad \operatorname{is\text{-}metric\text{-}space}(d) \wedge \\ \quad \operatorname{is\text{-}metric\text{-}space}(t) \wedge \\ \quad \operatorname{is\text{-}continuous}(d, t, f) \wedge \\ \quad \operatorname{seq\text{-}compact}(d) \Rightarrow \\ \quad \quad \operatorname{seq\text{-}compact}(\operatorname{subspace\text{-}ms}(t, \operatorname{image}(f, \operatorname{pts}(d)))) \end{array}}}
\vnbentry{\texttt{metrizable-\allowbreak{}compact-\allowbreak{}t-\allowbreak{}iff-\allowbreak{}seq-\allowbreak{}compact}}{\vnbstmt{\begin{array}{@{}l@{}} \forall s, d. \\ \quad \operatorname{is\text{-}metric\text{-}space}(d) \wedge \\ \quad (\operatorname{metric\text{-}top}(d) \simeq s) \Rightarrow \\ \quad \quad (\operatorname{is\text{-}compact\text{-}t}(s) \Leftrightarrow \operatorname{seq\text{-}compact}(d)) \end{array}}}
\vnbentry{\texttt{metrizable-\allowbreak{}hom-\allowbreak{}iff-\allowbreak{}sequential}}{\vnbstmt{\begin{array}{@{}l@{}} \forall s, t, d, b, f. \\ \quad \operatorname{is\text{-}metric\text{-}space}(d) \wedge \\ \quad \operatorname{is\text{-}metric\text{-}space}(b) \wedge \\ \quad (\operatorname{metric\text{-}top}(d) \simeq s) \wedge \\ \quad (\operatorname{metric\text{-}top}(b) \simeq t) \Rightarrow \\ \quad \quad (\operatorname{is\text{-}hom\text{-}top\text{-}space}(s, t, f) \Leftrightarrow ((f \in (\operatorname{pts}(d) \to \operatorname{pts}(b))) \wedge \forall a \in \operatorname{pts}(d), q \in \operatorname{pts}(d)^{\mathbb{N}}.\; (\operatorname{converges\text{-}to}(d, q, a) \Rightarrow \operatorname{converges\text{-}to}(b, \operatorname{compose}(f, q), f(a))))) \end{array}}}
\vnbentry{\texttt{metrizable-\allowbreak{}compact-\allowbreak{}image-\allowbreak{}ms}}{\vnbstmt{\begin{array}{@{}l@{}} \forall s, t, f, a. \\ \quad \operatorname{is\text{-}metrizable\text{-}top\text{-}space}(s) \wedge \\ \quad \operatorname{is\text{-}metrizable\text{-}top\text{-}space}(t) \wedge \\ \quad \operatorname{is\text{-}hom\text{-}top\text{-}space}(s, t, f) \wedge \\ \quad \operatorname{is\text{-}compact\text{-}t}(s) \wedge \\ \quad \operatorname{is\text{-}metric\text{-}space}(a) \wedge \\ \quad (\operatorname{metric\text{-}top}(a) \simeq t) \Rightarrow \\ \quad \quad \operatorname{is\text{-}compact}(\operatorname{subspace\text{-}ms}(a, \operatorname{image}(f, \operatorname{pts}(s)))) \end{array}}}
\vnbentry{\texttt{subspace-\allowbreak{}top-\allowbreak{}pts}}{\vnbstmt{\begin{array}{@{}l@{}} \forall t, a. \\ \quad (\operatorname{pts}(\operatorname{subspace\text{-}top}(t, a)) \simeq a) \end{array}}}
\vnbentry{\texttt{subspace-\allowbreak{}top-\allowbreak{}opens}}{\vnbstmt{\begin{array}{@{}l@{}} \forall t, a. \\ \quad (\operatorname{opens}(\operatorname{subspace\text{-}top}(t, a)) \simeq \{w \in \operatorname{power}(a) : \exists u \in \operatorname{opens}(t).\; (w = (u \cap a))\}) \end{array}}}
\vnbentry{\texttt{subspace-\allowbreak{}top-\allowbreak{}is-\allowbreak{}top-\allowbreak{}space}}{\vnbstmt{\begin{array}{@{}l@{}} \forall t, a. \\ \quad (\operatorname{is\text{-}top\text{-}space}(t) \wedge (a \subseteq \operatorname{pts}(t))) \Rightarrow  \\ \quad \quad \operatorname{is\text{-}top\text{-}space}(\operatorname{subspace\text{-}top}(t, a)) \end{array}}}
\vnbentry{\texttt{subspace-\allowbreak{}top-\allowbreak{}of-\allowbreak{}metric-\allowbreak{}top}}{\vnbstmt{\begin{array}{@{}l@{}} \forall t, a. \\ \quad (\operatorname{is\text{-}metric\text{-}space}(t) \wedge (a \subseteq \operatorname{pts}(t))) \Rightarrow  \\ \quad \quad (\operatorname{subspace\text{-}top}(\operatorname{metric\text{-}top}(t), a) \simeq \operatorname{metric\text{-}top}(\operatorname{subspace\text{-}ms}(t, a))) \end{array}}}
\vnbentry{\texttt{metrizable-\allowbreak{}compact-\allowbreak{}image}}{\vnbstmt{\begin{array}{@{}l@{}} \forall s, t, f. \\ \quad \operatorname{is\text{-}metrizable\text{-}top\text{-}space}(s) \wedge \\ \quad \operatorname{is\text{-}metrizable\text{-}top\text{-}space}(t) \wedge \\ \quad \operatorname{is\text{-}hom\text{-}top\text{-}space}(s, t, f) \wedge \\ \quad \operatorname{is\text{-}compact\text{-}t}(s) \Rightarrow \\ \quad \quad \operatorname{is\text{-}compact\text{-}t}(\operatorname{subspace\text{-}top}(t, \operatorname{image}(f, \operatorname{pts}(s)))) \end{array}}}

\vnbfile{theorem-library/metrizable-bounded-proof.scm}{theorem-library/metrizable-bounded-proof.scm}
\noindent{\small metrizable-bounded-proof.scm -- T1: a topological space is metrizable iff it is metrizable by a BOUNDED metric.}\par
\vnbentry{\texttt{metrizable-\allowbreak{}iff-\allowbreak{}bounded-\allowbreak{}metrizable}}{\vnbstmt{\begin{array}{@{}l@{}} \forall s. \\ \quad \operatorname{is\text{-}top\text{-}space}(s) \Rightarrow  \\ \quad \quad (\operatorname{is\text{-}metrizable\text{-}top\text{-}space}(s) \Leftrightarrow \exists d.\; (\operatorname{is\text{-}bounded\text{-}metric\text{-}space}(d) \wedge (\operatorname{metric\text{-}top}(d) \simeq s))) \end{array}}}

\vnbfile{theorem-library/ringoid-setoid-proof.scm}{theorem-library/ringoid-setoid-proof.scm}
\noindent{\small ringoid-setoid-proof.scm}\par
\vnbentry{\texttt{ringoid-\allowbreak{}rel-\allowbreak{}is-\allowbreak{}equivalence}}{\vnbstmt{\begin{array}{@{}l@{}} \forall r. \\ \quad \operatorname{is\text{-}ringoid}(r) \Rightarrow  \\ \quad \quad \operatorname{is\text{-}equivalence}(\operatorname{ringoid\text{-}rel}(r), \operatorname{carr}(r)) \end{array}}}
\vnbentry{\texttt{ringoid-\allowbreak{}setoid-\allowbreak{}is-\allowbreak{}setoid}}{\vnbstmt{\begin{array}{@{}l@{}} \forall r. \\ \quad \operatorname{is\text{-}ringoid}(r) \Rightarrow  \\ \quad \quad \operatorname{is\text{-}setoid}(\operatorname{ringoid\text{-}setoid}(r)) \end{array}}}

\vnbfile{theorem-library/rake-ringoid.scm}{theorem-library/rake-ringoid.scm}
\noindent{\small rake-ringoid.scm -- rq-add-computes: addition on R/I computes on classes.}\par
\vnbentry{\texttt{r6b-\allowbreak{}setoid-\allowbreak{}rel}}{\vnbstmt{\begin{array}{@{}l@{}} \forall r. \\ \quad (\operatorname{rel}(\operatorname{ringoid\text{-}setoid}(r)) \simeq \operatorname{ringoid\text{-}rel}(r)) \end{array}}}
\vnbentry{\texttt{r6b-\allowbreak{}related-\allowbreak{}out}}{\vnbstmt{\begin{array}{@{}l@{}} \forall r, a, b. \\ \quad \operatorname{is\text{-}ringoid}(r) \wedge \\ \quad (a \in \operatorname{carr}(r)) \wedge \\ \quad (b \in \operatorname{carr}(r)) \wedge \\ \quad \operatorname{related}(\operatorname{ringoid\text{-}setoid}(r), a, b) \Rightarrow \\ \quad \quad (\operatorname{add}(r)(a, \operatorname{neg}(r)(b)) \in \operatorname{idl}(r)) \end{array}}}
\vnbentry{\texttt{r6b-\allowbreak{}related-\allowbreak{}in}}{\vnbstmt{\begin{array}{@{}l@{}} \forall r, a, b. \\ \quad \operatorname{is\text{-}ringoid}(r) \wedge \\ \quad (a \in \operatorname{carr}(r)) \wedge \\ \quad (b \in \operatorname{carr}(r)) \wedge \\ \quad (\operatorname{add}(r)(a, \operatorname{neg}(r)(b)) \in \operatorname{idl}(r)) \Rightarrow \\ \quad \quad \operatorname{related}(\operatorname{ringoid\text{-}setoid}(r), a, b) \end{array}}}
\vnbentry{\texttt{r6b-\allowbreak{}rq-\allowbreak{}add-\allowbreak{}lam-\allowbreak{}type}}{\vnbstmt{\begin{array}{@{}l@{}} \forall r. \\ \quad \operatorname{is\text{-}ringoid}(r) \Rightarrow  \\ \quad \quad (\lambda \langle a, b \rangle \in (\operatorname{carr}(r) \times \operatorname{carr}(r)).\; \\ \quad \quad \quad \operatorname{class}(\operatorname{ringoid\text{-}setoid}(r), \operatorname{add}(r)(a, b))) \\ \quad \quad \in\; ((\operatorname{pts}(\operatorname{ringoid\text{-}setoid}(r)) \times \operatorname{pts}(\operatorname{ringoid\text{-}setoid}(r))) \to \operatorname{quotient}(\operatorname{ringoid\text{-}setoid}(r))) \end{array}}}
\vnbentry{\texttt{r6b-\allowbreak{}add-\allowbreak{}respects2}}{\vnbstmt{\begin{array}{@{}l@{}} \forall r. \\ \quad \operatorname{is\text{-}ringoid}(r) \Rightarrow  \\ \quad \quad \operatorname{respects2}(\operatorname{ringoid\text{-}setoid}(r), \\ \quad \quad \quad (\lambda \langle a, b \rangle \in (\operatorname{carr}(r) \times \operatorname{carr}(r)).\; \\ \quad \quad \quad \quad \operatorname{class}(\operatorname{ringoid\text{-}setoid}(r), \operatorname{add}(r)(a, b)))) \end{array}}}
\vnbentry{\texttt{rq-\allowbreak{}add-\allowbreak{}computes}}{\vnbstmt{\begin{array}{@{}l@{}} \forall r, a, b. \\ \quad \operatorname{is\text{-}ringoid}(r) \wedge \\ \quad (a \in \operatorname{carr}(r)) \wedge \\ \quad (b \in \operatorname{carr}(r)) \Rightarrow \\ \quad \quad \operatorname{add}(\operatorname{ringoid\text{-}quotient\text{-}ring}(r))(\operatorname{class}(\operatorname{ringoid\text{-}setoid}(r), a), \operatorname{class}(\operatorname{ringoid\text{-}setoid}(r), b)) \\ \quad \quad =\; \operatorname{class}(\operatorname{ringoid\text{-}setoid}(r), \operatorname{add}(r)(a, b)) \end{array}}}

\vnbfile{theorem-library/rake-ringoid-quotient.scm}{theorem-library/rake-ringoid-quotient.scm}
\noindent{\small rake-ringoid-quotient.scm -- MUL and NEG compute on R/I.}\par
\vnbentry{\texttt{r7n-\allowbreak{}rmul-\allowbreak{}type}}{\vnbstmt{\begin{array}{@{}l@{}} \forall r, a, b. \\ \quad \operatorname{is\text{-}ringoid}(r) \wedge \\ \quad (a \in \operatorname{carr}(r)) \wedge \\ \quad (b \in \operatorname{carr}(r)) \Rightarrow \\ \quad \quad (\operatorname{mul}(r)(a, b) \in \operatorname{carr}(r)) \end{array}}}
\vnbentry{\texttt{r7n-\allowbreak{}rone-\allowbreak{}type}}{\vnbstmt{\begin{array}{@{}l@{}} \forall r. \\ \quad \operatorname{is\text{-}ringoid}(r) \Rightarrow  \\ \quad \quad (\operatorname{one}(r) \in \operatorname{carr}(r)) \end{array}}}
\vnbentry{\texttt{r7n-\allowbreak{}radd-\allowbreak{}idem-\allowbreak{}zero}}{\vnbstmt{\begin{array}{@{}l@{}} \forall r.; \forall a \in \operatorname{carr}(r). \\ \quad (\operatorname{is\text{-}ringoid}(r) \wedge (\operatorname{add}(r)(a, a) = a)) \Rightarrow  \\ \quad \quad (a = \operatorname{zero}(r)) \end{array}}}
\vnbentry{\texttt{r7n-\allowbreak{}rzero-\allowbreak{}mul-\allowbreak{}left}}{\vnbstmt{\begin{array}{@{}l@{}} \forall r.; \forall b \in \operatorname{carr}(r). \\ \quad \operatorname{is\text{-}ringoid}(r) \Rightarrow  \\ \quad \quad (\operatorname{mul}(r)(\operatorname{zero}(r), b) = \operatorname{zero}(r)) \end{array}}}
\vnbentry{\texttt{r7n-\allowbreak{}rmul-\allowbreak{}zero-\allowbreak{}right}}{\vnbstmt{\begin{array}{@{}l@{}} \forall r.; \forall a \in \operatorname{carr}(r). \\ \quad \operatorname{is\text{-}ringoid}(r) \Rightarrow  \\ \quad \quad (\operatorname{mul}(r)(a, \operatorname{zero}(r)) = \operatorname{zero}(r)) \end{array}}}
\vnbentry{\texttt{r7n-\allowbreak{}rneg-\allowbreak{}mul-\allowbreak{}left}}{\vnbstmt{\begin{array}{@{}l@{}} \forall r, a, b. \\ \quad \operatorname{is\text{-}ringoid}(r) \wedge \\ \quad (a \in \operatorname{carr}(r)) \wedge \\ \quad (b \in \operatorname{carr}(r)) \Rightarrow \\ \quad \quad \operatorname{mul}(r)(\operatorname{neg}(r)(a), b) \\ \quad \quad =\; \operatorname{neg}(r)(\operatorname{mul}(r)(a, b)) \end{array}}}
\vnbentry{\texttt{r7n-\allowbreak{}rmul-\allowbreak{}neg-\allowbreak{}right}}{\vnbstmt{\begin{array}{@{}l@{}} \forall r, a, b. \\ \quad \operatorname{is\text{-}ringoid}(r) \wedge \\ \quad (a \in \operatorname{carr}(r)) \wedge \\ \quad (b \in \operatorname{carr}(r)) \Rightarrow \\ \quad \quad \operatorname{mul}(r)(a, \operatorname{neg}(r)(b)) \\ \quad \quad =\; \operatorname{neg}(r)(\operatorname{mul}(r)(a, b)) \end{array}}}
\vnbentry{\texttt{r7n-\allowbreak{}rdiff-\allowbreak{}mul}}{\vnbstmt{\begin{array}{@{}l@{}} \forall r, a, b, u, v. \\ \quad \operatorname{is\text{-}ringoid}(r) \wedge \\ \quad (a \in \operatorname{carr}(r)) \wedge \\ \quad (b \in \operatorname{carr}(r)) \wedge \\ \quad (u \in \operatorname{carr}(r)) \wedge \\ \quad (v \in \operatorname{carr}(r)) \Rightarrow \\ \quad \quad \operatorname{add}(r)(\operatorname{mul}(r)(a, b), \operatorname{neg}(r)(\operatorname{mul}(r)(u, v))) \\ \quad \quad =\; \operatorname{add}(r)(\operatorname{mul}(r)(\operatorname{add}(r)(a, \operatorname{neg}(r)(u)), b), \operatorname{mul}(r)(u, \operatorname{add}(r)(b, \operatorname{neg}(r)(v)))) \end{array}}}
\vnbentry{\texttt{r7n-\allowbreak{}rq-\allowbreak{}mul-\allowbreak{}lam-\allowbreak{}type}}{\vnbstmt{\begin{array}{@{}l@{}} \forall r. \\ \quad \operatorname{is\text{-}ringoid}(r) \Rightarrow  \\ \quad \quad (\lambda \langle a, b \rangle \in (\operatorname{carr}(r) \times \operatorname{carr}(r)).\; \\ \quad \quad \quad \operatorname{class}(\operatorname{ringoid\text{-}setoid}(r), \operatorname{mul}(r)(a, b))) \\ \quad \quad \in\; ((\operatorname{pts}(\operatorname{ringoid\text{-}setoid}(r)) \times \operatorname{pts}(\operatorname{ringoid\text{-}setoid}(r))) \to \operatorname{quotient}(\operatorname{ringoid\text{-}setoid}(r))) \end{array}}}
\vnbentry{\texttt{r7n-\allowbreak{}mul-\allowbreak{}respects2}}{\vnbstmt{\begin{array}{@{}l@{}} \forall r. \\ \quad \operatorname{is\text{-}ringoid}(r) \Rightarrow  \\ \quad \quad \operatorname{respects2}(\operatorname{ringoid\text{-}setoid}(r), \\ \quad \quad \quad (\lambda \langle a, b \rangle \in (\operatorname{carr}(r) \times \operatorname{carr}(r)).\; \\ \quad \quad \quad \quad \operatorname{class}(\operatorname{ringoid\text{-}setoid}(r), \operatorname{mul}(r)(a, b)))) \end{array}}}
\vnbentry{\texttt{rq-\allowbreak{}mul-\allowbreak{}computes}}{\vnbstmt{\begin{array}{@{}l@{}} \forall r, a, b. \\ \quad \operatorname{is\text{-}ringoid}(r) \wedge \\ \quad (a \in \operatorname{carr}(r)) \wedge \\ \quad (b \in \operatorname{carr}(r)) \Rightarrow \\ \quad \quad \operatorname{mul}(\operatorname{ringoid\text{-}quotient\text{-}ring}(r))(\operatorname{class}(\operatorname{ringoid\text{-}setoid}(r), a), \operatorname{class}(\operatorname{ringoid\text{-}setoid}(r), b)) \\ \quad \quad =\; \operatorname{class}(\operatorname{ringoid\text{-}setoid}(r), \operatorname{mul}(r)(a, b)) \end{array}}}
\vnbentry{\texttt{r7n-\allowbreak{}rq-\allowbreak{}neg-\allowbreak{}lam-\allowbreak{}type}}{\vnbstmt{\begin{array}{@{}l@{}} \forall r. \\ \quad \operatorname{is\text{-}ringoid}(r) \Rightarrow  \\ \quad \quad (\lambda a \in \operatorname{carr}(r).\; \\ \quad \quad \quad \operatorname{class}(\operatorname{ringoid\text{-}setoid}(r), \operatorname{neg}(r)(a))) \\ \quad \quad \in\; (\operatorname{pts}(\operatorname{ringoid\text{-}setoid}(r)) \to \operatorname{quotient}(\operatorname{ringoid\text{-}setoid}(r))) \end{array}}}
\vnbentry{\texttt{r7n-\allowbreak{}neg-\allowbreak{}respects}}{\vnbstmt{\begin{array}{@{}l@{}} \forall r. \\ \quad \operatorname{is\text{-}ringoid}(r) \Rightarrow  \\ \quad \quad \operatorname{respects}(\operatorname{ringoid\text{-}setoid}(r), \\ \quad \quad \quad (\lambda a \in \operatorname{carr}(r).\; \\ \quad \quad \quad \quad \operatorname{class}(\operatorname{ringoid\text{-}setoid}(r), \operatorname{neg}(r)(a)))) \end{array}}}
\vnbentry{\texttt{rq-\allowbreak{}neg-\allowbreak{}computes}}{\vnbstmt{\begin{array}{@{}l@{}} \forall r.; \forall a \in \operatorname{carr}(r). \\ \quad \operatorname{is\text{-}ringoid}(r) \Rightarrow  \\ \quad \quad \operatorname{neg}(\operatorname{ringoid\text{-}quotient\text{-}ring}(r))(\operatorname{class}(\operatorname{ringoid\text{-}setoid}(r), a)) \\ \quad \quad =\; \operatorname{class}(\operatorname{ringoid\text{-}setoid}(r), \operatorname{neg}(r)(a)) \end{array}}}
\vnbentry{\texttt{ringoid-\allowbreak{}quotient-\allowbreak{}is-\allowbreak{}ring}}{\vnbstmt{\begin{array}{@{}l@{}} \forall r. \\ \quad \operatorname{is\text{-}ringoid}(r) \Rightarrow  \\ \quad \quad \operatorname{is\text{-}ring}(\operatorname{ringoid\text{-}quotient\text{-}ring}(r)) \end{array}}}

\vnbfile{theorem-library/geometric-series.scm}{theorem-library/geometric-series.scm}
\noindent{\small geometric-series.scm -- THE GEOMETRIC MAJORANT, PROVEN, and the four small facts about powers it rests on.}\par
\vnbentry{\texttt{rr-\allowbreak{}abs-\allowbreak{}power}}{\vnbstmt{\begin{array}{@{}l@{}} \forall k \in \mathbb{N}, r \in \mathbb{R}. \\ \quad (\lvert \operatorname{power}(r, k) \rvert = \operatorname{power}(\lvert r \rvert, k)) \end{array}}}
\vnbentry{\texttt{geometric-\allowbreak{}partial-\allowbreak{}sum-\allowbreak{}mul}}{\vnbstmt{\begin{array}{@{}l@{}} \forall k \in \mathbb{N}, r \in \mathbb{R}. \\ \quad ((1 - r) \cdot \operatorname{series\text{-}partial\text{-}sum}((\lambda n \in \mathbb{N}.\; \operatorname{power}(r, n)), k)) \\ \quad =\; (1 - \operatorname{power}(r, k)) \end{array}}}
\vnbentry{\texttt{rr-\allowbreak{}power-\allowbreak{}nonneg}}{\vnbstmt{\begin{array}{@{}l@{}} \forall k \in \mathbb{N}, r \in \mathbb{R}. \\ \quad ((0 \leq r) \Rightarrow (0 \leq \operatorname{power}(r, k))) \end{array}}}
\vnbentry{\texttt{geometric-\allowbreak{}partial-\allowbreak{}sum}}{\vnbstmt{\begin{array}{@{}l@{}} \forall r \in \mathbb{R}, k \in \mathbb{N}. \\ \quad \neg (r = 1) \Rightarrow  \\ \quad \quad \operatorname{series\text{-}partial\text{-}sum}((\lambda n \in \mathbb{N}.\; \operatorname{power}(r, n)), k) \\ \quad \quad =\; ((1 - \operatorname{power}(r, k)) \cdot \operatorname{recip}((1 - r))) \end{array}}}
\vnbentry{\texttt{geometric-\allowbreak{}series-\allowbreak{}converges}}{\vnbstmt{\begin{array}{@{}l@{}} \forall r \in \mathbb{R}. \\ \quad ((0 \leq r) \wedge (r < 1)) \Rightarrow  \\ \quad \quad \operatorname{series\text{-}converges}((\lambda n \in \mathbb{N}.\; \operatorname{power}(r, n))) \end{array}}}

\vnbfile{theorem-library/limsup-tests.scm}{theorem-library/limsup-tests.scm}
\noindent{\small limsup-tests.scm -- THE UPPER LIMIT OF A SEQUENCE IN [0, +inf], THE ROOT TEST AND THE RATIO TEST, as the notes state them. Batch 25-A, 2026-09-24.}\par
\vnbentry{\texttt{einf-\allowbreak{}lower}}{\vnbstmt{\begin{array}{@{}l@{}} \forall s.; \forall x \in s. \\ \quad (s \subseteq \operatorname{rr\text{-}pos\text{-}star}) \Rightarrow  \\ \quad \quad (\operatorname{einf}(s) \leq x) \end{array}}}
\vnbentry{\texttt{einf-\allowbreak{}greatest}}{\vnbstmt{\begin{array}{@{}l@{}} \forall s.; \forall b \in \operatorname{rr\text{-}pos\text{-}star}. \\ \quad ((s \subseteq \operatorname{rr\text{-}pos\text{-}star}) \wedge \forall x \in s.\; (b \leq x)) \Rightarrow  \\ \quad \quad (b \leq \operatorname{einf}(s)) \end{array}}}
\vnbentry{\texttt{einf-\allowbreak{}in}}{\vnbstmt{\begin{array}{@{}l@{}} \forall s. \\ \quad (s \subseteq \operatorname{rr\text{-}pos\text{-}star}) \Rightarrow  \\ \quad \quad (\operatorname{einf}(s) \in \operatorname{rr\text{-}pos\text{-}star}) \end{array}}}
\vnbentry{\texttt{etail-\allowbreak{}subset-\allowbreak{}rr-\allowbreak{}pos-\allowbreak{}star}}{\vnbstmt{\begin{array}{@{}l@{}} \forall f, n. \\ \quad (\forall k \in \mathbb{N}.\; (f(k) \in \mathbb{R}) \wedge \forall k \in \mathbb{N}.\; (0 \leq f(k))) \Rightarrow  \\ \quad \quad (\operatorname{etail}(f, n) \subseteq \operatorname{rr\text{-}pos\text{-}star}) \end{array}}}
\vnbentry{\texttt{elimsup-\allowbreak{}eventually-\allowbreak{}below}}{\vnbstmt{\begin{array}{@{}l@{}} \forall f.; \forall r \in \mathbb{R}. \\ \quad \forall k \in \mathbb{N}.\; (f(k) \in \mathbb{R}) \wedge \\ \quad \forall k \in \mathbb{N}.\; (0 \leq f(k)) \wedge \\ \quad (\operatorname{elimsup}(f) < r) \Rightarrow \\ \quad \quad \exists n \in \mathbb{N}. \forall k \in \mathbb{N}. \\ \quad \quad \quad ((n \leq k) \Rightarrow (f(k) < r)) \end{array}}}
\vnbentry{\texttt{elimsup-\allowbreak{}frequently-\allowbreak{}above}}{\vnbstmt{\begin{array}{@{}l@{}} \forall f.; \forall r \in \mathbb{R}, n \in \mathbb{N}. \\ \quad \forall k \in \mathbb{N}.\; (f(k) \in \mathbb{R}) \wedge \\ \quad \forall k \in \mathbb{N}.\; (0 \leq f(k)) \wedge \\ \quad (0 \leq r) \wedge \\ \quad (r < \operatorname{elimsup}(f)) \Rightarrow \\ \quad \quad \exists k \in \mathbb{N}.\; ((n \leq k) \wedge (r < f(k))) \end{array}}}
\vnbentry{\texttt{eliminf-\allowbreak{}eventually-\allowbreak{}above}}{\vnbstmt{\begin{array}{@{}l@{}} \forall f.; \forall r \in \mathbb{R}. \\ \quad \forall k \in \mathbb{N}.\; (f(k) \in \mathbb{R}) \wedge \\ \quad \forall k \in \mathbb{N}.\; (0 \leq f(k)) \wedge \\ \quad (0 \leq r) \wedge \\ \quad (r < \operatorname{eliminf}(f)) \Rightarrow \\ \quad \quad \exists n \in \mathbb{N}. \forall k \in \mathbb{N}. \\ \quad \quad \quad ((n \leq k) \Rightarrow (r < f(k))) \end{array}}}
\vnbentry{\texttt{elimsup-\allowbreak{}in-\allowbreak{}rr-\allowbreak{}pos-\allowbreak{}star}}{\vnbstmt{\begin{array}{@{}l@{}} \forall f. \\ \quad (\forall k \in \mathbb{N}.\; (f(k) \in \mathbb{R}) \wedge \forall k \in \mathbb{N}.\; (0 \leq f(k))) \Rightarrow  \\ \quad \quad (\operatorname{elimsup}(f) \in \operatorname{rr\text{-}pos\text{-}star}) \end{array}}}
\vnbentry{\texttt{eliminf-\allowbreak{}in-\allowbreak{}rr-\allowbreak{}pos-\allowbreak{}star}}{\vnbstmt{\begin{array}{@{}l@{}} \forall f. \\ \quad (\forall k \in \mathbb{N}.\; (f(k) \in \mathbb{R}) \wedge \forall k \in \mathbb{N}.\; (0 \leq f(k))) \Rightarrow  \\ \quad \quad (\operatorname{eliminf}(f) \in \operatorname{rr\text{-}pos\text{-}star}) \end{array}}}
\vnbentry{\texttt{rr-\allowbreak{}power-\allowbreak{}mono-\allowbreak{}base}}{\vnbstmt{\begin{array}{@{}l@{}} \forall k \in \mathbb{N}, x, y \in \mathbb{R}. \\ \quad ((0 \leq x) \wedge (x \leq y)) \Rightarrow  \\ \quad \quad (\operatorname{power}(x, k) \leq \operatorname{power}(y, k)) \end{array}}}
\vnbentry{\texttt{rpow-\allowbreak{}star-\allowbreak{}nat}}{\vnbstmt{\begin{array}{@{}l@{}} \forall k \in \mathbb{N}, x \in \mathbb{R}. \\ \quad ((0 < x) \Rightarrow ({x}^{k} = \operatorname{power}(x, k))) \end{array}}}
\vnbentry{\texttt{recip-\allowbreak{}star-\allowbreak{}nn-\allowbreak{}pos}}{\vnbstmt{\forall k \in \mathbb{N}.\; ((1 \leq k) \Rightarrow (0 < {k}^{\ast\!-1}))}}
\vnbentry{\texttt{recip-\allowbreak{}star-\allowbreak{}nn-\allowbreak{}nonneg}}{\vnbstmt{\forall k \in \mathbb{N}.\; (0 \leq {k}^{\ast\!-1})}}
\vnbentry{\texttt{rpow-\allowbreak{}star-\allowbreak{}root-\allowbreak{}power}}{\vnbstmt{\begin{array}{@{}l@{}} \forall k \in \mathbb{N}, x \in \mathbb{R}. \\ \quad ((1 \leq k) \wedge (0 < x)) \Rightarrow  \\ \quad \quad (\operatorname{power}({x}^{{k}^{\ast\!-1}}, k) = x) \end{array}}}
\vnbentry{\texttt{rpow-\allowbreak{}star-\allowbreak{}nonneg-\allowbreak{}in-\allowbreak{}rr}}{\vnbstmt{\begin{array}{@{}l@{}} \forall x, s \in \mathbb{R}. \\ \quad (((0 \leq x) \wedge (0 \leq s)) \Rightarrow ({x}^{s} \in \mathbb{R})) \end{array}}}
\vnbentry{\texttt{rpow-\allowbreak{}star-\allowbreak{}nonneg-\allowbreak{}nonneg}}{\vnbstmt{\begin{array}{@{}l@{}} \forall x, s \in \mathbb{R}. \\ \quad (((0 \leq x) \wedge (0 \leq s)) \Rightarrow (0 \leq {x}^{s})) \end{array}}}
\vnbentry{\texttt{series-\allowbreak{}bounded-\allowbreak{}converges}}{\vnbstmt{\begin{array}{@{}l@{}} \forall f \in (\mathbb{N} \to \mathbb{R}), b \in \mathbb{R}. \\ \quad \forall k \in \mathbb{N}.\; (0 \leq f(k)) \wedge \\ \quad \forall k \in \mathbb{N}. \\ \quad \quad (\operatorname{series\text{-}partial\text{-}sum}(f, k) \leq b) \Rightarrow \\ \quad \quad \operatorname{series\text{-}converges}(f) \end{array}}}
\vnbentry{\texttt{series-\allowbreak{}eventually-\allowbreak{}geometric-\allowbreak{}converges}}{\vnbstmt{\begin{array}{@{}l@{}} \forall f \in (\mathbb{N} \to \mathbb{R}), q, c \in \mathbb{R}, n \in \mathbb{N}. \\ \quad \forall k \in \mathbb{N}.\; (0 \leq f(k)) \wedge \\ \quad (0 \leq q) \wedge \\ \quad (q < 1) \wedge \\ \quad (0 \leq c) \wedge \\ \quad \forall k \in \mathbb{N}. \\ \quad \quad ((n \leq k) \Rightarrow (f(k) \leq (c \cdot \operatorname{power}(q, k)))) \Rightarrow \\ \quad \quad \operatorname{series\text{-}converges}(f) \end{array}}}
\vnbentry{\texttt{series-\allowbreak{}frequently-\allowbreak{}large-\allowbreak{}diverges}}{\vnbstmt{\begin{array}{@{}l@{}} \forall f \in (\mathbb{N} \to \mathbb{R}), c \in \mathbb{R}. \\ \quad (0 < c) \wedge \\ \quad \forall n \in \mathbb{N}. \exists k \in \mathbb{N}. \\ \quad \quad ((n \leq k) \wedge (c \leq f(k))) \Rightarrow \\ \quad \quad \neg \operatorname{series\text{-}converges}(f) \end{array}}}
\vnbentry{\texttt{root-\allowbreak{}test-\allowbreak{}converges}}{\vnbstmt{\begin{array}{@{}l@{}} \forall h \in (\mathbb{N} \to \mathbb{R}). \\ \quad \forall k \in \mathbb{N}.\; (0 \leq h(k)) \wedge \\ \quad (\operatorname{elimsup}((\lambda j \in \mathbb{N}.\; {h(j)}^{{j}^{\ast\!-1}})) < 1) \Rightarrow \\ \quad \quad \operatorname{series\text{-}converges}(h) \end{array}}}
\vnbentry{\texttt{root-\allowbreak{}test-\allowbreak{}diverges}}{\vnbstmt{\begin{array}{@{}l@{}} \forall h \in (\mathbb{N} \to \mathbb{R}). \\ \quad \forall k \in \mathbb{N}.\; (0 \leq h(k)) \wedge \\ \quad (1 < \operatorname{elimsup}((\lambda j \in \mathbb{N}.\; {h(j)}^{{j}^{\ast\!-1}}))) \Rightarrow \\ \quad \quad \neg \operatorname{series\text{-}converges}(h) \end{array}}}
\vnbentry{\texttt{root-\allowbreak{}test}}{\vnbstmt{\begin{array}{@{}l@{}} \forall h \in (\mathbb{N} \to \mathbb{R}). \\ \quad \forall k \in \mathbb{N}.\; (0 \leq h(k)) \Rightarrow  \\ \quad \quad (\operatorname{elimsup}((\lambda j \in \mathbb{N}.\; {h(j)}^{{j}^{\ast\!-1}})) < 1) \Rightarrow  \\ \quad \quad \quad \operatorname{series\text{-}converges}(h) \wedge \\ \quad \quad (1 < \operatorname{elimsup}((\lambda j \in \mathbb{N}.\; {h(j)}^{{j}^{\ast\!-1}}))) \Rightarrow  \\ \quad \quad \quad \neg \operatorname{series\text{-}converges}(h) \end{array}}}
\vnbentry{\texttt{ratio-\allowbreak{}test-\allowbreak{}converges-\allowbreak{}limsup}}{\vnbstmt{\begin{array}{@{}l@{}} \forall h \in (\mathbb{N} \to \mathbb{R}). \\ \quad \forall k \in \mathbb{N}.\; (0 < h(k)) \wedge \\ \quad \operatorname{elimsup}((\lambda j \in \mathbb{N}.\; (h((j + 1)) \cdot \operatorname{recip}(h(j))))) \\ \quad <\; 1 \Rightarrow \\ \quad \quad \operatorname{series\text{-}converges}(h) \end{array}}}
\vnbentry{\texttt{ratio-\allowbreak{}test-\allowbreak{}diverges-\allowbreak{}liminf}}{\vnbstmt{\begin{array}{@{}l@{}} \forall h \in (\mathbb{N} \to \mathbb{R}). \\ \quad \forall k \in \mathbb{N}.\; (0 < h(k)) \wedge \\ \quad 1 \\ \quad <\; \operatorname{eliminf}((\lambda j \in \mathbb{N}.\; (h((j + 1)) \cdot \operatorname{recip}(h(j))))) \Rightarrow \\ \quad \quad \neg \operatorname{series\text{-}converges}(h) \end{array}}}
\vnbentry{\texttt{ratio-\allowbreak{}test}}{\vnbstmt{\begin{array}{@{}l@{}} \forall h \in (\mathbb{N} \to \mathbb{R}). \\ \quad \forall k \in \mathbb{N}.\; (0 < h(k)) \Rightarrow  \\ \quad \quad \operatorname{elimsup}((\lambda j \in \mathbb{N}.\; (h((j + 1)) \cdot \operatorname{recip}(h(j))))) \\ \quad \quad <\; 1 \Rightarrow \\ \quad \quad \quad \operatorname{series\text{-}converges}(h) \wedge \\ \quad \quad 1 \\ \quad \quad <\; \operatorname{eliminf}((\lambda j \in \mathbb{N}.\; (h((j + 1)) \cdot \operatorname{recip}(h(j))))) \Rightarrow \\ \quad \quad \quad \neg \operatorname{series\text{-}converges}(h) \end{array}}}

\vnbfile{theorem-library/rpow-defined.scm}{theorem-library/rpow-defined.scm}
\noindent{\small theorem-library/rpow-defined.scm -- THE AXIOMATIC REAL POWER RETIRED: RPOW is DEFINED as RPOW-STAR, and its laws are THEOREMS under their old names and with their old statements (the user's decision, 2026-10-03).}\par
\vnbentry{\texttt{rpow-\allowbreak{}pos}}{\vnbstmt{\begin{array}{@{}l@{}} \forall a \in \mathbb{R}, b \in \mathbb{Q}. \\ \quad (0 < a) \Rightarrow  \\ \quad \quad ((\operatorname{rpow}(a, b) \in \mathbb{R}) \wedge (0 < \operatorname{rpow}(a, b))) \end{array}}}
\vnbentry{\texttt{rpow-\allowbreak{}zero}}{\vnbstmt{\begin{array}{@{}l@{}} \forall a \in \mathbb{R}. \\ \quad ((0 < a) \Rightarrow (\operatorname{rpow}(a, 0) = 1)) \end{array}}}
\vnbentry{\texttt{rpow-\allowbreak{}one}}{\vnbstmt{\begin{array}{@{}l@{}} \forall a \in \mathbb{R}. \\ \quad ((0 < a) \Rightarrow (\operatorname{rpow}(a, 1) = a)) \end{array}}}
\vnbentry{\texttt{rpow-\allowbreak{}add}}{\vnbstmt{\begin{array}{@{}l@{}} \forall a \in \mathbb{R}, b, d \in \mathbb{Q}. \\ \quad (0 < a) \Rightarrow  \\ \quad \quad (\operatorname{rpow}(a, (b + d)) = (\operatorname{rpow}(a, b) \cdot \operatorname{rpow}(a, d))) \end{array}}}
\vnbentry{\texttt{rpow-\allowbreak{}mul-\allowbreak{}base}}{\vnbstmt{\begin{array}{@{}l@{}} \forall a, c \in \mathbb{R}, b \in \mathbb{Q}. \\ \quad ((0 < a) \wedge (0 < c)) \Rightarrow  \\ \quad \quad (\operatorname{rpow}((a \cdot c), b) = (\operatorname{rpow}(a, b) \cdot \operatorname{rpow}(c, b))) \end{array}}}
\vnbentry{\texttt{rpow-\allowbreak{}pow}}{\vnbstmt{\begin{array}{@{}l@{}} \forall a \in \mathbb{R}, b, d \in \mathbb{Q}. \\ \quad (0 < a) \Rightarrow  \\ \quad \quad (\operatorname{rpow}(\operatorname{rpow}(a, b), d) = \operatorname{rpow}(a, (b \cdot d))) \end{array}}}
\vnbentry{\texttt{rpow-\allowbreak{}neg}}{\vnbstmt{\begin{array}{@{}l@{}} \forall a \in \mathbb{R}, b \in \mathbb{Q}. \\ \quad (0 < a) \Rightarrow  \\ \quad \quad (\operatorname{rpow}(a, (0 - b)) = \operatorname{recip}(\operatorname{rpow}(a, b))) \end{array}}}
\vnbentry{\texttt{rpow-\allowbreak{}nat}}{\vnbstmt{\begin{array}{@{}l@{}} \forall a \in \mathbb{R}, n \in \mathbb{N}. \\ \quad ((0 < a) \Rightarrow (\operatorname{rpow}(a, n) = \operatorname{power}(a, n))) \end{array}}}
\vnbentry{\texttt{rpow-\allowbreak{}zero-\allowbreak{}base}}{\vnbstmt{\begin{array}{@{}l@{}} \forall b \in \mathbb{Q}. \\ \quad ((0 < b) \Rightarrow (\operatorname{rpow}(0, b) = 0)) \end{array}}}
\vnbentry{\texttt{rpow-\allowbreak{}zero-\allowbreak{}zero}}{\vnbstmt{(\operatorname{rpow}(0, 0) = 1)}}
\vnbentry{\texttt{rpow-\allowbreak{}mono-\allowbreak{}base}}{\vnbstmt{\begin{array}{@{}l@{}} \forall a, c \in \mathbb{R}, b \in \mathbb{Q}. \\ \quad ((0 < a) \wedge ((a \leq c) \wedge (0 \leq b))) \Rightarrow  \\ \quad \quad (\operatorname{rpow}(a, b) \leq \operatorname{rpow}(c, b)) \end{array}}}
\vnbentry{\texttt{rpow-\allowbreak{}mono-\allowbreak{}exp-\allowbreak{}ge1}}{\vnbstmt{\begin{array}{@{}l@{}} \forall a \in \mathbb{R}, b, d \in \mathbb{Q}. \\ \quad ((1 \leq a) \wedge (b \leq d)) \Rightarrow  \\ \quad \quad (\operatorname{rpow}(a, b) \leq \operatorname{rpow}(a, d)) \end{array}}}
\vnbentry{\texttt{rpow-\allowbreak{}mono-\allowbreak{}exp-\allowbreak{}le1}}{\vnbstmt{\begin{array}{@{}l@{}} \forall a \in \mathbb{R}, b, d \in \mathbb{Q}. \\ \quad (((0 < a) \wedge (a \leq 1)) \wedge (b \leq d)) \Rightarrow  \\ \quad \quad (\operatorname{rpow}(a, d) \leq \operatorname{rpow}(a, b)) \end{array}}}
\vnbentry{\texttt{sqrt-\allowbreak{}rpow}}{\vnbstmt{\begin{array}{@{}l@{}} \forall a \in \mathbb{R}. \\ \quad ((0 < a) \Rightarrow (\sqrt{a} = \operatorname{rpow}(a, /(1, 2)))) \end{array}}}
\vnbentry{\texttt{amgm-\allowbreak{}2-\allowbreak{}sqrt}}{\vnbstmt{\begin{array}{@{}l@{}} \forall a, b \in \mathbb{R}. \\ \quad ((0 \leq a) \wedge (0 \leq b)) \Rightarrow  \\ \quad \quad (\sqrt{(a \cdot b)} \leq /((a + b), 2)) \end{array}}}
\vnbentry{\texttt{young-\allowbreak{}inequality}}{\vnbstmt{\begin{array}{@{}l@{}} \forall a, b \in \mathbb{R}, p, q \in \mathbb{Q}. \\ \quad ((0 < a) \wedge ((0 < b) \wedge ((1 < p) \wedge ((1 < q) \wedge ((/(1, p) + /(1, q)) = 1))))) \Rightarrow  \\ \quad \quad ((a \cdot b) \leq (/(\operatorname{rpow}(a, p), p) + /(\operatorname{rpow}(b, q), q))) \end{array}}}
\vnbentry{\texttt{bernoulli-\allowbreak{}rpow}}{\vnbstmt{\begin{array}{@{}l@{}} \forall x \in \mathbb{R}, b \in \mathbb{Q}. \\ \quad ((0 < (1 + x)) \wedge (1 \leq b)) \Rightarrow  \\ \quad \quad ((1 + (b \cdot x)) \leq \operatorname{rpow}((1 + x), b)) \end{array}}}

\vnbfile{theorem-library/measure-laws.scm}{theorem-library/measure-laws.scm}
\noindent{\small measure-laws.scm -- Chapter 1 of the measure notes (Thayer, Construction of Measures, Sec. 3-4): the elementary laws of a countably additive measure. Agent M-1 of the October roadmap (docs/roadmap-2026-10-03.md, line (b), decision 4).  Helper prefix mlw-.  25 theorems, every one `modulo 0'.}\par
\vnbentry{\texttt{rps-\allowbreak{}prefix-\allowbreak{}zero}}{\vnbstmt{\begin{array}{@{}l@{}} \forall f. \\ \quad \operatorname{finsum}(\operatorname{rr\text{-}pos\text{-}star\text{-}add\text{-}monoid}, f, \mathrm{OS}(0)) \\ \quad =\; 0 \end{array}}}
\vnbentry{\texttt{rps-\allowbreak{}prefix-\allowbreak{}succ}}{\vnbstmt{\begin{array}{@{}l@{}} \forall f \in (\mathbb{N} \to \operatorname{rr\text{-}pos\text{-}star}), n \in \mathbb{N}. \\ \quad \operatorname{finsum}(\operatorname{rr\text{-}pos\text{-}star\text{-}add\text{-}monoid}, f, \mathrm{OS}((n + 1))) \\ \quad =\; \operatorname{eplus}(\operatorname{finsum}(\operatorname{rr\text{-}pos\text{-}star\text{-}add\text{-}monoid}, f, \mathrm{OS}(n)), \\ \quad \quad f(n)) \end{array}}}
\vnbentry{\texttt{rps-\allowbreak{}finsum-\allowbreak{}mono-\allowbreak{}set}}{\vnbstmt{\begin{array}{@{}l@{}} \forall f, s.; \forall t \in \mathbf{Set}. \\ \quad (|t| \in \mathbb{N}) \wedge \\ \quad (s \subseteq t) \wedge \\ \quad \forall z \in t.\; (f(z) \in \operatorname{rr\text{-}pos\text{-}star}) \Rightarrow \\ \quad \quad \operatorname{finsum}(\operatorname{rr\text{-}pos\text{-}star\text{-}add\text{-}monoid}, f, s) \\ \quad \quad \leq\; \operatorname{finsum}(\operatorname{rr\text{-}pos\text{-}star\text{-}add\text{-}monoid}, f, t) \end{array}}}
\vnbentry{\texttt{esum-\allowbreak{}nn-\allowbreak{}prefix-\allowbreak{}le}}{\vnbstmt{\begin{array}{@{}l@{}} \forall f \in (\mathbb{N} \to \operatorname{rr\text{-}pos\text{-}star}), n \in \mathbb{N}. \\ \quad \operatorname{finsum}(\operatorname{rr\text{-}pos\text{-}star\text{-}add\text{-}monoid}, f, \mathrm{OS}(n)) \\ \quad \leq\; \operatorname{esum}(f) \end{array}}}
\vnbentry{\texttt{esum-\allowbreak{}nn-\allowbreak{}le-\allowbreak{}of-\allowbreak{}prefixes}}{\vnbstmt{\begin{array}{@{}l@{}} \forall f \in (\mathbb{N} \to \operatorname{rr\text{-}pos\text{-}star}), b \in \operatorname{rr\text{-}pos\text{-}star}. \\ \quad \forall n \in \mathbb{N}. \\ \quad \quad \operatorname{finsum}(\operatorname{rr\text{-}pos\text{-}star\text{-}add\text{-}monoid}, f, \mathrm{OS}(n)) \\ \quad \quad \leq\; b \Rightarrow \\ \quad \quad (\operatorname{esum}(f) \leq b) \end{array}}}
\vnbentry{\texttt{rps-\allowbreak{}prefix-\allowbreak{}type}}{\vnbstmt{\begin{array}{@{}l@{}} \forall f \in (\mathbb{N} \to \operatorname{rr\text{-}pos\text{-}star}), n \in \mathbb{N}. \\ \quad \operatorname{finsum}(\operatorname{rr\text{-}pos\text{-}star\text{-}add\text{-}monoid}, f, \mathrm{OS}(n)) \\ \quad \in\; \operatorname{rr\text{-}pos\text{-}star} \end{array}}}
\vnbentry{\texttt{rps-\allowbreak{}prefix-\allowbreak{}le-\allowbreak{}termwise}}{\vnbstmt{\begin{array}{@{}l@{}} \forall n \in \mathbb{N}, f, g. \\ \quad (f \in (\mathbb{N} \to \operatorname{rr\text{-}pos\text{-}star})) \wedge \\ \quad (g \in (\mathbb{N} \to \operatorname{rr\text{-}pos\text{-}star})) \wedge \\ \quad \forall k \in \mathbb{N}.\; (f(k) \leq g(k)) \Rightarrow \\ \quad \quad \operatorname{finsum}(\operatorname{rr\text{-}pos\text{-}star\text{-}add\text{-}monoid}, f, \mathrm{OS}(n)) \\ \quad \quad \leq\; \operatorname{finsum}(\operatorname{rr\text{-}pos\text{-}star\text{-}add\text{-}monoid}, g, \mathrm{OS}(n)) \end{array}}}
\vnbentry{\texttt{esum-\allowbreak{}nn-\allowbreak{}le-\allowbreak{}termwise}}{\vnbstmt{\begin{array}{@{}l@{}} \forall f, g. \\ \quad (f \in (\mathbb{N} \to \operatorname{rr\text{-}pos\text{-}star})) \wedge \\ \quad (g \in (\mathbb{N} \to \operatorname{rr\text{-}pos\text{-}star})) \wedge \\ \quad \forall k \in \mathbb{N}.\; (f(k) \leq g(k)) \Rightarrow \\ \quad \quad (\operatorname{esum}(f) \leq \operatorname{esum}(g)) \end{array}}}
\vnbentry{\texttt{esum-\allowbreak{}nn-\allowbreak{}two-\allowbreak{}support}}{\vnbstmt{\begin{array}{@{}l@{}} \forall h \in (\mathbb{N} \to \operatorname{rr\text{-}pos\text{-}star}). \\ \quad \forall q \in \mathbb{N}.\; (h(((q + 1) + 1)) = 0) \Rightarrow  \\ \quad \quad (\operatorname{esum}(h) = \operatorname{eplus}(h(0), h((0 + 1)))) \end{array}}}
\vnbentry{\texttt{measure-\allowbreak{}finitely-\allowbreak{}additive}}{\vnbstmt{\begin{array}{@{}l@{}} \forall a, b, u, \alpha, \beta. \\ \quad \operatorname{is\text{-}measure}(a, b, u) \wedge \\ \quad (\alpha \in b) \wedge \\ \quad (\beta \in b) \wedge \\ \quad ((\alpha \cap \beta) = \emptyset) \Rightarrow \\ \quad \quad (u((\alpha \cup \beta)) = \operatorname{eplus}(u(\alpha), u(\beta))) \end{array}}}
\vnbentry{\texttt{measure-\allowbreak{}monotone}}{\vnbstmt{\begin{array}{@{}l@{}} \forall a, b, u, \alpha, \beta. \\ \quad \operatorname{is\text{-}measure}(a, b, u) \wedge \\ \quad (\alpha \in b) \wedge \\ \quad (\beta \in b) \wedge \\ \quad (\alpha \subseteq \beta) \Rightarrow \\ \quad \quad (u(\alpha) \leq u(\beta)) \end{array}}}
\vnbentry{\texttt{bu-\allowbreak{}seg-\allowbreak{}zero}}{\vnbstmt{\begin{array}{@{}l@{}} \forall f. \\ \quad (\operatorname{big\text{-}union}(j, \mathrm{OS}(0), f(j)) = \emptyset) \end{array}}}
\vnbentry{\texttt{bu-\allowbreak{}seg-\allowbreak{}succ}}{\vnbstmt{\begin{array}{@{}l@{}} \forall k \in \mathbb{N}, c.; \forall f \in (\mathbb{N} \to c). \\ \quad \operatorname{big\text{-}union}(j, \mathrm{OS}((k + 1)), f(j)) \\ \quad =\; (\operatorname{big\text{-}union}(j, \mathrm{OS}(k), f(j)) \cup f(k)) \end{array}}}
\vnbentry{\texttt{sigma-\allowbreak{}algebra-\allowbreak{}seg-\allowbreak{}union-\allowbreak{}in}}{\vnbstmt{\begin{array}{@{}l@{}} \forall k \in \mathbb{N}, a, b, f. \\ \quad (\operatorname{is\text{-}sigma\text{-}algebra}(a, b) \wedge (f \in (\mathbb{N} \to b))) \Rightarrow  \\ \quad \quad (\operatorname{big\text{-}union}(j, \mathrm{OS}(k), f(j)) \in b) \end{array}}}
\vnbentry{\texttt{disjointify-\allowbreak{}type}}{\vnbstmt{\begin{array}{@{}l@{}} \forall a, b, f. \\ \quad (\operatorname{is\text{-}sigma\text{-}algebra}(a, b) \wedge (f \in (\mathbb{N} \to b))) \Rightarrow  \\ \quad \quad (\lambda k \in \mathbb{N}.\; \\ \quad \quad \quad \operatorname{complement\text{-}in}(f(k), \operatorname{big\text{-}union}(j, \mathrm{OS}(k), f(j)))) \\ \quad \quad \in\; (\mathbb{N} \to b) \end{array}}}
\vnbentry{\texttt{disjointify-\allowbreak{}disjoint}}{\vnbstmt{\begin{array}{@{}l@{}} \forall f, c, i, j. \\ \quad (f \in (\mathbb{N} \to c)) \wedge \\ \quad (i \in \mathbb{N}) \wedge \\ \quad (j \in \mathbb{N}) \wedge \\ \quad \neg (i = j) \Rightarrow \\ \quad \quad ((\lambda k \in \mathbb{N}.\; \operatorname{complement\text{-}in}(f(k), \operatorname{big\text{-}union}(\mathit{uj}, \mathrm{OS}(k), f(\mathit{uj}))))(i) \cap (\lambda k \in \mathbb{N}.\; \operatorname{complement\text{-}in}(f(k), \operatorname{big\text{-}union}(\mathit{uj}, \mathrm{OS}(k), f(\mathit{uj}))))(j)) \\ \quad \quad =\; \emptyset \end{array}}}
\vnbentry{\texttt{disjointify-\allowbreak{}union}}{\vnbstmt{\begin{array}{@{}l@{}} \forall f, c. \\ \quad (f \in (\mathbb{N} \to c)) \Rightarrow  \\ \quad \quad \operatorname{big\text{-}union}(n, \mathbb{N}, \\ \quad \quad \quad (\lambda k \in \mathbb{N}.\; \operatorname{complement\text{-}in}(f(k), \operatorname{big\text{-}union}(j, \mathrm{OS}(k), f(j))))(n)) \\ \quad \quad =\; \operatorname{big\text{-}union}(n, \mathbb{N}, f(n)) \end{array}}}
\vnbentry{\texttt{measure-\allowbreak{}countably-\allowbreak{}subadditive}}{\vnbstmt{\begin{array}{@{}l@{}} \forall a, b, u, f. \\ \quad (\operatorname{is\text{-}measure}(a, b, u) \wedge (f \in (\mathbb{N} \to b))) \Rightarrow  \\ \quad \quad u(\operatorname{big\text{-}union}(n, \mathbb{N}, f(n))) \\ \quad \quad \leq\; \operatorname{esum}((\lambda n \in \mathbb{N}.\; u(f(n)))) \end{array}}}
\vnbentry{\texttt{disjointify-\allowbreak{}prefix}}{\vnbstmt{\begin{array}{@{}l@{}} \forall n \in \mathbb{N}, a, b, u, f. \\ \quad (\operatorname{is\text{-}measure}(a, b, u) \wedge (f \in (\mathbb{N} \to b))) \Rightarrow  \\ \quad \quad \operatorname{finsum}(\operatorname{rr\text{-}pos\text{-}star\text{-}add\text{-}monoid}, \\ \quad \quad \quad (\lambda c \in \mathbb{N}.\; \\ \quad \quad \quad \quad u((\lambda k \in \mathbb{N}.\; \operatorname{complement\text{-}in}(f(k), \operatorname{big\text{-}union}(j, \mathrm{OS}(k), f(j))))(c))), \\ \quad \quad \quad \mathrm{OS}(n)) \\ \quad \quad =\; u(\operatorname{big\text{-}union}(j, \mathrm{OS}(n), f(j))) \end{array}}}
\vnbentry{\texttt{nondecreasing-\allowbreak{}seg-\allowbreak{}subset}}{\vnbstmt{\begin{array}{@{}l@{}} \forall n \in \mathbb{N}, c.; \forall f \in (\mathbb{N} \to c), j \in \mathrm{OS}(n). \\ \quad \forall a \in \mathbb{N}.\; (f(a) \subseteq f((a + 1))) \Rightarrow  \\ \quad \quad (f(j) \subseteq f(n)) \end{array}}}
\vnbentry{\texttt{measure-\allowbreak{}continuity-\allowbreak{}from-\allowbreak{}below}}{\vnbstmt{\begin{array}{@{}l@{}} \forall a, b, u, f. \\ \quad \operatorname{is\text{-}measure}(a, b, u) \wedge \\ \quad (f \in (\mathbb{N} \to b)) \wedge \\ \quad \forall n \in \mathbb{N}.\; (f(n) \subseteq f((n + 1))) \Rightarrow \\ \quad \quad u(\operatorname{big\text{-}union}(n, \mathbb{N}, f(n))) \\ \quad \quad =\; \operatorname{esup}(\operatorname{image}((\lambda n \in \mathbb{N}.\; u(f(n))), \mathbb{N})) \end{array}}}
\vnbentry{\texttt{measure-\allowbreak{}finite-\allowbreak{}subset}}{\vnbstmt{\begin{array}{@{}l@{}} \forall a, b, u, \alpha, \beta. \\ \quad \operatorname{is\text{-}measure}(a, b, u) \wedge \\ \quad (\alpha \in b) \wedge \\ \quad (\beta \in b) \wedge \\ \quad (\alpha \subseteq \beta) \wedge \\ \quad (u(\beta) \in \mathbb{R}) \Rightarrow \\ \quad \quad (u(\alpha) \in \mathbb{R}) \end{array}}}
\vnbentry{\texttt{measure-\allowbreak{}subtractive}}{\vnbstmt{\begin{array}{@{}l@{}} \forall a, b, u, \alpha, \beta. \\ \quad \operatorname{is\text{-}measure}(a, b, u) \wedge \\ \quad (\alpha \in b) \wedge \\ \quad (\beta \in b) \wedge \\ \quad (\alpha \subseteq \beta) \wedge \\ \quad (u(\beta) \in \mathbb{R}) \Rightarrow \\ \quad \quad (u(\operatorname{complement\text{-}in}(\beta, \alpha)) = (u(\beta) - u(\alpha))) \end{array}}}
\vnbentry{\texttt{measure-\allowbreak{}continuity-\allowbreak{}from-\allowbreak{}above}}{\vnbstmt{\begin{array}{@{}l@{}} \forall a, b, u, f. \\ \quad \operatorname{is\text{-}measure}(a, b, u) \wedge \\ \quad (f \in (\mathbb{N} \to b)) \wedge \\ \quad \forall n \in \mathbb{N}.\; (f((n + 1)) \subseteq f(n)) \wedge \\ \quad (u(f(0)) \in \mathbb{R}) \Rightarrow \\ \quad \quad u(\operatorname{complement\text{-}in}(a, \operatorname{big\text{-}union}(n, \mathbb{N}, \operatorname{complement\text{-}in}(a, f(n))))) \\ \quad \quad =\; \operatorname{einf}(\operatorname{image}((\lambda n \in \mathbb{N}.\; u(f(n))), \mathbb{N})) \end{array}}}
\vnbentry{\texttt{measure-\allowbreak{}pushforward}}{\vnbstmt{\begin{array}{@{}l@{}} \forall a, b, u, a_{2}, c, f. \\ \quad \operatorname{is\text{-}measure}(a, b, u) \wedge \\ \quad \operatorname{is\text{-}measurable\text{-}map}(a, b, a_{2}, c, f) \Rightarrow \\ \quad \quad \operatorname{is\text{-}measure}(a_{2}, c, (\lambda \mathit{pb} \in c.\; u(\{x \in a : (f(x) \in \mathit{pb})\}))) \end{array}}}

\vnbfile{theorem-library/ps-series-bridges.scm}{theorem-library/ps-series-bridges.scm}
\noindent{\small ps-series-bridges.scm -- THE THREE PS/SERIES BRIDGES, PROVEN, and the CONVERGES twin of rr-limit-ptwise-eq that the third one needed.}\par
\vnbentry{\texttt{ps-\allowbreak{}abs-\allowbreak{}term}}{\vnbstmt{\begin{array}{@{}l@{}} \forall f \in (\mathbb{N} \to \mathbb{R}), x \in \mathbb{R}, n \in \mathbb{N}. \\ \quad \lvert (f(n) \cdot \operatorname{power}(x, n)) \rvert \\ \quad =\; (\lvert f(n) \rvert \cdot \operatorname{power}(\lvert x \rvert, n)) \end{array}}}
\vnbentry{\texttt{ps-\allowbreak{}partial-\allowbreak{}sum-\allowbreak{}as-\allowbreak{}series}}{\vnbstmt{\begin{array}{@{}l@{}} \forall f \in (\mathbb{N} \to \mathbb{R}), x \in \mathbb{R}, k \in \mathbb{N}. \\ \quad \operatorname{ps\text{-}partial\text{-}sum}(f, x, k) \\ \quad =\; \operatorname{series\text{-}partial\text{-}sum}((\lambda n \in \mathbb{N}.\; (f(n) \cdot \operatorname{power}(x, n))), k) \end{array}}}
\vnbentry{\texttt{rr-\allowbreak{}converges-\allowbreak{}ptwise-\allowbreak{}eq}}{\vnbstmt{\begin{array}{@{}l@{}} \forall f, h \in (\mathbb{N} \to \mathbb{R}). \\ \quad \forall j \in \mathbb{N}.\; (h(j) = f(j)) \wedge \\ \quad \operatorname{converges}(\operatorname{rr\text{-}ms}, f) \Rightarrow \\ \quad \quad \operatorname{converges}(\operatorname{rr\text{-}ms}, h) \end{array}}}
\vnbentry{\texttt{ps-\allowbreak{}converges-\allowbreak{}as-\allowbreak{}series}}{\vnbstmt{\begin{array}{@{}l@{}} \forall f \in (\mathbb{N} \to \mathbb{R}), x \in \mathbb{R}. \\ \quad (\operatorname{ps\text{-}converges\text{-}at}(f, x) \Leftrightarrow \operatorname{series\text{-}converges}((\lambda n \in \mathbb{N}.\; (f(n) \cdot \operatorname{power}(x, n))))) \end{array}}}

\vnbfile{theorem-library/ratio-test-proof.scm}{theorem-library/ratio-test-proof.scm}
\noindent{\small ratio-test-proof.scm -- THE RATIO TEST FOR A REAL POWER SERIES (ratio-test-converges, structure-library/power-series.scm), PROVEN. Batch 27-A, 2026-09-24.}\par
\vnbentry{\texttt{elimsup-\allowbreak{}eventually-\allowbreak{}le}}{\vnbstmt{\begin{array}{@{}l@{}} \forall f.; \forall r \in \mathbb{R}, n \in \mathbb{N}. \\ \quad \forall k \in \mathbb{N}.\; (f(k) \in \mathbb{R}) \wedge \\ \quad \forall k \in \mathbb{N}.\; (0 \leq f(k)) \wedge \\ \quad \forall a \in \mathbb{N}.\; ((n \leq a) \Rightarrow (f(a) \leq r)) \Rightarrow \\ \quad \quad (\operatorname{elimsup}(f) \leq r) \end{array}}}
\vnbentry{\texttt{ratio-\allowbreak{}test-\allowbreak{}converges}}{\vnbstmt{\begin{array}{@{}l@{}} \forall f \in (\mathbb{N} \to \mathbb{R}), l, x \in \mathbb{R}. \\ \quad (0 \leq l) \wedge \\ \quad \forall n \in \mathbb{N}.\; \neg (f(n) = 0) \wedge \\ \quad \operatorname{ps\text{-}ratio\text{-}limit}(f, l) \wedge \\ \quad ((\lvert x \rvert \cdot l) < 1) \Rightarrow \\ \quad \quad \operatorname{ps\text{-}converges\text{-}at}(f, x) \end{array}}}

\vnbfile{theorem-library/rake-series.scm}{theorem-library/rake-series.scm}
\noindent{\small rake-series.scm -- rake batch 5 T (2026-09-18): the little-o algebra, the power-series absolute-convergence bridge, and the Taylor clearing step.}\par
\vnbentry{\texttt{vtaylor-\allowbreak{}clear}}{\vnbstmt{\begin{array}{@{}l@{}} \forall v, a, b, w \in \mathbb{R}, n \in \mathbb{N}. \\ \quad (((\operatorname{factorial}((n + 1)) \cdot v) = (a \cdot w)) \wedge ((a \leq b) \wedge (0 \leq w))) \Rightarrow  \\ \quad \quad ((\operatorname{factorial}((n + 1)) \cdot v) \leq (b \cdot w)) \end{array}}}
\vnbentry{\texttt{little-\allowbreak{}o-\allowbreak{}scalar}}{\vnbstmt{\begin{array}{@{}l@{}} \forall c, g, a. \\ \quad ((c \in \mathbb{R}) \wedge \operatorname{little\text{-}o\text{-}at}(g, a)) \Rightarrow  \\ \quad \quad \operatorname{little\text{-}o\text{-}at}((\lambda x \in \mathbb{R}.\; (c \cdot g(x))), a) \end{array}}}
\vnbentry{\texttt{little-\allowbreak{}o-\allowbreak{}sum}}{\vnbstmt{\begin{array}{@{}l@{}} \forall g, h, a. \\ \quad \operatorname{little\text{-}o\text{-}at}(g, a) \wedge \\ \quad \operatorname{little\text{-}o\text{-}at}(h, a) \Rightarrow \\ \quad \quad \operatorname{little\text{-}o\text{-}at}((\lambda x \in \mathbb{R}.\; (g(x) + h(x))), a) \end{array}}}
\vnbentry{\texttt{shift-\allowbreak{}lam-\allowbreak{}in-\allowbreak{}fun}}{\vnbstmt{\begin{array}{@{}l@{}} \forall k \in \mathbb{R}, f \in (\mathbb{R} \to \mathbb{R}). \\ \quad ((\lambda x \in \mathbb{R}.\; (f(x) + k)) \in (\mathbb{R} \to \mathbb{R})) \end{array}}}
\vnbentry{\texttt{shift-\allowbreak{}continuous-\allowbreak{}at}}{\vnbstmt{\begin{array}{@{}l@{}} \forall k \in \mathbb{R}, f, a. \\ \quad \operatorname{is\text{-}continuous\text{-}at}(\operatorname{rr\text{-}ms}, \operatorname{rr\text{-}ms}, f, a) \Rightarrow \\ \quad \quad \operatorname{is\text{-}continuous\text{-}at}(\operatorname{rr\text{-}ms}, \operatorname{rr\text{-}ms}, (\lambda x \in \mathbb{R}.\; (f(x) + k)), a) \end{array}}}
\vnbentry{\texttt{diff-\allowbreak{}iff-\allowbreak{}little-\allowbreak{}o}}{\vnbstmt{\begin{array}{@{}l@{}} \forall f, a, l. \\ \quad (\operatorname{is\text{-}diff\text{-}at}(f, a, l) \Leftrightarrow ((f \in (\mathbb{R} \to \mathbb{R})) \wedge ((a \in \mathbb{R}) \wedge ((l \in \mathbb{R}) \wedge \operatorname{little\text{-}o\text{-}at}((\lambda x \in \mathbb{R}.\; ((f(x) - f(a)) - (l \cdot (x - a)))), a))))) \end{array}}}
\vnbentry{\texttt{ps-\allowbreak{}absolute-\allowbreak{}implies-\allowbreak{}convergent}}{\vnbstmt{\begin{array}{@{}l@{}} \forall f \in (\mathbb{N} \to \mathbb{R}), x \in \mathbb{R}. \\ \quad \operatorname{ps\text{-}absolutely\text{-}converges\text{-}at}(f, x) \Rightarrow  \\ \quad \quad \operatorname{ps\text{-}converges\text{-}at}(f, x) \end{array}}}
\vnbentry{\texttt{esum-\allowbreak{}finite-\allowbreak{}implies-\allowbreak{}bounded}}{\vnbstmt{\begin{array}{@{}l@{}} \forall f \in (\operatorname{dom}(f) \to \operatorname{rr\text{-}pos\text{-}star}). \\ \quad (\operatorname{esum}(f) \in \mathbb{R}) \Rightarrow  \\ \quad \quad \exists m \in \mathbb{R}. \forall s \in \mathbf{Set}. \\ \quad \quad \quad ((|s| \in \mathbb{N}) \wedge (s \subseteq \operatorname{dom}(f))) \Rightarrow  \\ \quad \quad \quad \quad \operatorname{finsum}(\operatorname{rr\text{-}pos\text{-}star\text{-}add\text{-}monoid}, f, s) \\ \quad \quad \quad \quad \leq\; m \end{array}}}

\vnbfile{theorem-library/rake-series2.scm}{theorem-library/rake-series2.scm}
\noindent{\small rake-series2.scm -- rake batch 5c (2026-09-18), assignment 5c-R: the two series leaves batch T left BLOCKED, and the four bricks they were missing.}\par
\vnbentry{\texttt{metric-\allowbreak{}chain-\allowbreak{}bound}}{\vnbstmt{\begin{array}{@{}l@{}} \forall n \in \mathbb{N}, s, f, d, m. \\ \quad \operatorname{is\text{-}metric\text{-}space}(s) \wedge \\ \quad (f \in (\mathbb{N} \to \operatorname{pts}(s))) \wedge \\ \quad (d \in (\mathbb{N} \to \mathbb{R})) \wedge \\ \quad \forall k \in \mathbb{N}.\; (\operatorname{dist}(s)(f(k), f((k + 1))) \leq d(k)) \wedge \\ \quad (m \in \mathbb{N}) \wedge \\ \quad (m \leq n) \Rightarrow \\ \quad \quad \operatorname{dist}(s)(f(m), f(n)) \\ \quad \quad \leq\; (\operatorname{series\text{-}partial\text{-}sum}(d, n) - \operatorname{series\text{-}partial\text{-}sum}(d, m)) \end{array}}}
\vnbentry{\texttt{metric-\allowbreak{}chain-\allowbreak{}triangle}}{\vnbstmt{\begin{array}{@{}l@{}} \forall s, f.; \forall m, n \in \mathbb{N}. \\ \quad \operatorname{is\text{-}metric\text{-}space}(s) \wedge \\ \quad (f \in (\mathbb{N} \to \operatorname{pts}(s))) \wedge \\ \quad (m \leq n) \Rightarrow \\ \quad \quad \operatorname{dist}(s)(f(m), f(n)) \\ \quad \quad \leq\; (\operatorname{series\text{-}partial\text{-}sum}((\lambda k \in \mathbb{N}.\; \operatorname{dist}(s)(f(k), f((k + 1)))), n) - \operatorname{series\text{-}partial\text{-}sum}((\lambda k \in \mathbb{N}.\; \operatorname{dist}(s)(f(k), f((k + 1)))), m)) \end{array}}}
\vnbentry{\texttt{summable-\allowbreak{}bound-\allowbreak{}implies-\allowbreak{}cauchy}}{\vnbstmt{\begin{array}{@{}l@{}} \forall s, f, d. \\ \quad \operatorname{is\text{-}metric\text{-}space}(s) \wedge \\ \quad (f \in (\mathbb{N} \to \operatorname{pts}(s))) \wedge \\ \quad (d \in (\mathbb{N} \to \mathbb{R})) \wedge \\ \quad \operatorname{series\text{-}converges}(d) \wedge \\ \quad \forall k \in \mathbb{N}.\; (\operatorname{dist}(s)(f(k), f((k + 1))) \leq d(k)) \Rightarrow \\ \quad \quad \operatorname{is\text{-}cauchy\text{-}seq}(s, f) \end{array}}}
\vnbentry{\texttt{rr-\allowbreak{}power-\allowbreak{}le-\allowbreak{}one}}{\vnbstmt{\begin{array}{@{}l@{}} \forall k \in \mathbb{N}, r \in \mathbb{R}. \\ \quad (((0 \leq r) \wedge (r \leq 1)) \Rightarrow (\operatorname{power}(r, k) \leq 1)) \end{array}}}
\vnbentry{\texttt{rr-\allowbreak{}power-\allowbreak{}shrink-\allowbreak{}bound}}{\vnbstmt{\begin{array}{@{}l@{}} \forall k \in \mathbb{N}, r \in \mathbb{R}. \\ \quad ((0 \leq r) \wedge (r \leq 1)) \Rightarrow  \\ \quad \quad (((k \cdot (1 - r)) \cdot \operatorname{power}(r, k)) \leq 1) \end{array}}}
\vnbentry{\texttt{power-\allowbreak{}small-\allowbreak{}threshold}}{\vnbstmt{\begin{array}{@{}l@{}} \forall r \in \mathbb{R}, s. \\ \quad (((0 \leq r) \wedge (r < 1)) \wedge \operatorname{pos\text{-}rr}(s)) \Rightarrow  \\ \quad \quad \exists p \in \mathbb{N}. \forall k \in \mathbb{N}. \\ \quad \quad \quad ((p \leq k) \Rightarrow (\operatorname{power}(r, k) \leq s)) \end{array}}}
\vnbentry{\texttt{power-\allowbreak{}tends-\allowbreak{}to-\allowbreak{}zero}}{\vnbstmt{\begin{array}{@{}l@{}} \forall r \in \mathbb{R}. \\ \quad (\lvert r \rvert < 1) \Rightarrow  \\ \quad \quad \operatorname{converges\text{-}to}(\operatorname{rr\text{-}ms}, (\lambda n \in \mathbb{N}.\; \operatorname{power}(r, n)), 0) \end{array}}}
\vnbentry{\texttt{rr-\allowbreak{}const-\allowbreak{}converges-\allowbreak{}to}}{\vnbstmt{\begin{array}{@{}l@{}} \forall c \in \mathbb{R}. \\ \quad \operatorname{converges\text{-}to}(\operatorname{rr\text{-}ms}, (\lambda n \in \mathbb{N}.\; c), c) \end{array}}}
\vnbentry{\texttt{geometric-\allowbreak{}series-\allowbreak{}converges-\allowbreak{}to}}{\vnbstmt{\begin{array}{@{}l@{}} \forall r \in \mathbb{R}. \\ \quad (\lvert r \rvert < 1) \Rightarrow  \\ \quad \quad \operatorname{series\text{-}converges\text{-}to}((\lambda n \in \mathbb{N}.\; \operatorname{power}(r, n)), \operatorname{recip}((1 - r))) \end{array}}}

\vnbfile{theorem-library/heine-borel-interval.scm}{theorem-library/heine-borel-interval.scm}
\noindent{\small heine-borel-interval.scm -- HEINE-BOREL for a closed interval AS A SPACE:}\par
\vnbentry{\texttt{rr-\allowbreak{}limit-\allowbreak{}in-\allowbreak{}ccint}}{\vnbstmt{\begin{array}{@{}l@{}} \forall o, i \in \mathbb{R}, q \in (\mathbb{N} \to \operatorname{ccint}(o, i)), t \in \mathbb{R}. \\ \quad \operatorname{converges\text{-}to}(\operatorname{rr\text{-}ms}, q, t) \Rightarrow  \\ \quad \quad (t \in \operatorname{ccint}(o, i)) \end{array}}}
\vnbentry{\texttt{heine-\allowbreak{}borel-\allowbreak{}ccint}}{\vnbstmt{\begin{array}{@{}l@{}} \forall v, \mathit{bv} \in \mathbb{R}. \\ \quad \operatorname{is\text{-}compact}(\operatorname{subspace\text{-}ms}(\operatorname{rr\text{-}ms}, \operatorname{ccint}(v, \mathit{bv}))) \end{array}}}

\vnbfile{theorem-library/cc-coords-laws.scm}{theorem-library/cc-coords-laws.scm}
\noindent{\small cc-coords-laws.scm -- CC and cartesian(RR, RR), explicitly.}\par
\vnbentry{\texttt{apply-\allowbreak{}congruence-\allowbreak{}1}}{\vnbstmt{\forall f, g, u.\; ((f \simeq g) \Rightarrow (f(u) \simeq g(u)))}}
\vnbentry{\texttt{cc-\allowbreak{}coords-\allowbreak{}apply}}{\vnbstmt{\begin{array}{@{}l@{}} \forall z \in \mathbb{C}. \\ \quad (\operatorname{cc\text{-}coords}(z) \simeq \langle \operatorname{real\text{-}part}(z), \operatorname{imag\text{-}part}(z) \rangle) \end{array}}}
\vnbentry{\texttt{cc-\allowbreak{}coords-\allowbreak{}in-\allowbreak{}fun}}{\vnbstmt{(\operatorname{cc\text{-}coords} \in (\mathbb{C} \to (\mathbb{R} \times \mathbb{R})))}}
\vnbentry{\texttt{cc-\allowbreak{}of-\allowbreak{}pair-\allowbreak{}apply}}{\vnbstmt{\begin{array}{@{}l@{}} \forall a, b \in \mathbb{R}. \\ \quad (\operatorname{cc\text{-}of\text{-}pair}(\langle a, b \rangle) \simeq (a + (b \cdot +i))) \end{array}}}
\vnbentry{\texttt{cc-\allowbreak{}of-\allowbreak{}pair-\allowbreak{}in-\allowbreak{}fun}}{\vnbstmt{(\operatorname{cc\text{-}of\text{-}pair} \in ((\mathbb{R} \times \mathbb{R}) \to \mathbb{C}))}}
\vnbentry{\texttt{cc-\allowbreak{}of-\allowbreak{}pair-\allowbreak{}of-\allowbreak{}coords}}{\vnbstmt{\begin{array}{@{}l@{}} \forall z \in \mathbb{C}. \\ \quad (\operatorname{cc\text{-}of\text{-}pair}(\operatorname{cc\text{-}coords}(z)) = z) \end{array}}}
\vnbentry{\texttt{cc-\allowbreak{}coords-\allowbreak{}of-\allowbreak{}pair}}{\vnbstmt{\begin{array}{@{}l@{}} \forall r \in (\mathbb{R} \times \mathbb{R}). \\ \quad (\operatorname{cc\text{-}coords}(\operatorname{cc\text{-}of\text{-}pair}(r)) = r) \end{array}}}
\vnbentry{\texttt{cc-\allowbreak{}coords-\allowbreak{}is-\allowbreak{}bijection}}{\vnbstmt{(\operatorname{cc\text{-}coords} \in \mathrm{Bij}(\mathbb{C}, (\mathbb{R} \times \mathbb{R})))}}
\vnbentry{\texttt{cc-\allowbreak{}of-\allowbreak{}pair-\allowbreak{}is-\allowbreak{}bijection}}{\vnbstmt{\begin{array}{@{}l@{}} \operatorname{cc\text{-}of\text{-}pair} \\ \in\; \mathrm{Bij}((\mathbb{R} \times \mathbb{R}), \mathbb{C}) \end{array}}}
\vnbentry{\texttt{cc-\allowbreak{}re-\allowbreak{}add}}{\vnbstmt{\begin{array}{@{}l@{}} \forall z, w \in \mathbb{C}. \\ \quad \operatorname{real\text{-}part}((z + w)) \\ \quad =\; (\operatorname{real\text{-}part}(z) + \operatorname{real\text{-}part}(w)) \end{array}}}
\vnbentry{\texttt{cc-\allowbreak{}im-\allowbreak{}add}}{\vnbstmt{\begin{array}{@{}l@{}} \forall z, w \in \mathbb{C}. \\ \quad \operatorname{imag\text{-}part}((z + w)) \\ \quad =\; (\operatorname{imag\text{-}part}(z) + \operatorname{imag\text{-}part}(w)) \end{array}}}
\vnbentry{\texttt{cc-\allowbreak{}re-\allowbreak{}real-\allowbreak{}mul}}{\vnbstmt{\begin{array}{@{}l@{}} \forall c \in \mathbb{R}, z \in \mathbb{C}. \\ \quad (\operatorname{real\text{-}part}((c \cdot z)) = (c \cdot \operatorname{real\text{-}part}(z))) \end{array}}}
\vnbentry{\texttt{cc-\allowbreak{}im-\allowbreak{}real-\allowbreak{}mul}}{\vnbstmt{\begin{array}{@{}l@{}} \forall c \in \mathbb{R}, z \in \mathbb{C}. \\ \quad (\operatorname{imag\text{-}part}((c \cdot z)) = (c \cdot \operatorname{imag\text{-}part}(z))) \end{array}}}
\vnbentry{\texttt{cc-\allowbreak{}re-\allowbreak{}mul}}{\vnbstmt{\begin{array}{@{}l@{}} \forall z, w \in \mathbb{C}. \\ \quad \operatorname{real\text{-}part}((z \cdot w)) \\ \quad =\; ((\operatorname{real\text{-}part}(z) \cdot \operatorname{real\text{-}part}(w)) - (\operatorname{imag\text{-}part}(z) \cdot \operatorname{imag\text{-}part}(w))) \end{array}}}
\vnbentry{\texttt{cc-\allowbreak{}im-\allowbreak{}mul}}{\vnbstmt{\begin{array}{@{}l@{}} \forall z, w \in \mathbb{C}. \\ \quad \operatorname{imag\text{-}part}((z \cdot w)) \\ \quad =\; ((\operatorname{real\text{-}part}(z) \cdot \operatorname{imag\text{-}part}(w)) + (\operatorname{imag\text{-}part}(z) \cdot \operatorname{real\text{-}part}(w))) \end{array}}}
\vnbentry{\texttt{cc-\allowbreak{}coords-\allowbreak{}add}}{\vnbstmt{\begin{array}{@{}l@{}} \forall z, w \in \mathbb{C}. \\ \quad \operatorname{cc\text{-}coords}((z + w)) \\ \quad =\; \langle (\operatorname{real\text{-}part}(z) + \operatorname{real\text{-}part}(w)), (\operatorname{imag\text{-}part}(z) + \operatorname{imag\text{-}part}(w)) \rangle \end{array}}}
\vnbentry{\texttt{cc-\allowbreak{}coords-\allowbreak{}real-\allowbreak{}mul}}{\vnbstmt{\begin{array}{@{}l@{}} \forall c \in \mathbb{R}, z \in \mathbb{C}. \\ \quad \operatorname{cc\text{-}coords}((c \cdot z)) \\ \quad =\; \langle (c \cdot \operatorname{real\text{-}part}(z)), (c \cdot \operatorname{imag\text{-}part}(z)) \rangle \end{array}}}
\vnbentry{\texttt{cc-\allowbreak{}of-\allowbreak{}pair-\allowbreak{}add}}{\vnbstmt{\begin{array}{@{}l@{}} \forall a_{1}, b_{1}, a_{2}, b_{2} \in \mathbb{R}. \\ \quad \operatorname{cc\text{-}of\text{-}pair}(\langle (a_{1} + a_{2}), (b_{1} + b_{2}) \rangle) \\ \quad =\; (\operatorname{cc\text{-}of\text{-}pair}(\langle a_{1}, b_{1} \rangle) + \operatorname{cc\text{-}of\text{-}pair}(\langle a_{2}, b_{2} \rangle)) \end{array}}}
\vnbentry{\texttt{cc-\allowbreak{}of-\allowbreak{}pair-\allowbreak{}real-\allowbreak{}mul}}{\vnbstmt{\begin{array}{@{}l@{}} \forall c, a, b \in \mathbb{R}. \\ \quad \operatorname{cc\text{-}of\text{-}pair}(\langle (c \cdot a), (c \cdot b) \rangle) \\ \quad =\; (c \cdot \operatorname{cc\text{-}of\text{-}pair}(\langle a, b \rangle)) \end{array}}}
\vnbentry{\texttt{cc-\allowbreak{}re-\allowbreak{}dist-\allowbreak{}le}}{\vnbstmt{\begin{array}{@{}l@{}} \forall z, w \in \mathbb{C}. \\ \quad \lvert (\operatorname{real\text{-}part}(z) - \operatorname{real\text{-}part}(w)) \rvert \\ \quad \leq\; \operatorname{magnitude}((z - w)) \end{array}}}
\vnbentry{\texttt{cc-\allowbreak{}im-\allowbreak{}dist-\allowbreak{}le}}{\vnbstmt{\begin{array}{@{}l@{}} \forall z, w \in \mathbb{C}. \\ \quad \lvert (\operatorname{imag\text{-}part}(z) - \operatorname{imag\text{-}part}(w)) \rvert \\ \quad \leq\; \operatorname{magnitude}((z - w)) \end{array}}}
\vnbentry{\texttt{cc-\allowbreak{}dist-\allowbreak{}le-\allowbreak{}coords}}{\vnbstmt{\begin{array}{@{}l@{}} \forall z, w \in \mathbb{C}. \\ \quad \operatorname{magnitude}((z - w)) \\ \quad \leq\; (\lvert (\operatorname{real\text{-}part}(z) - \operatorname{real\text{-}part}(w)) \rvert + \lvert (\operatorname{imag\text{-}part}(z) - \operatorname{imag\text{-}part}(w)) \rvert) \end{array}}}
\vnbentry{\texttt{cc-\allowbreak{}coords-\allowbreak{}of-\allowbreak{}converges}}{\vnbstmt{\begin{array}{@{}l@{}} \forall f \in (\mathbb{N} \to \mathbb{C}), g, q \in (\mathbb{N} \to \mathbb{R}), v \in \mathbb{C}. \\ \quad \forall n \in \mathbb{N}.\; (\operatorname{real\text{-}part}(f(n)) = g(n)) \wedge \\ \quad \forall n \in \mathbb{N}.\; (\operatorname{imag\text{-}part}(f(n)) = q(n)) \wedge \\ \quad \operatorname{converges\text{-}to}(\operatorname{cc\text{-}ms}, f, v) \Rightarrow \\ \quad \quad \operatorname{converges\text{-}to}(\operatorname{rr\text{-}ms}, g, \operatorname{real\text{-}part}(v)) \wedge \\ \quad \quad \operatorname{converges\text{-}to}(\operatorname{rr\text{-}ms}, q, \operatorname{imag\text{-}part}(v)) \end{array}}}
\vnbentry{\texttt{cc-\allowbreak{}converges-\allowbreak{}iff-\allowbreak{}coords}}{\vnbstmt{\begin{array}{@{}l@{}} \forall f \in (\mathbb{N} \to \mathbb{C}), g, q \in (\mathbb{N} \to \mathbb{R}), v \in \mathbb{C}. \\ \quad \forall n \in \mathbb{N}.\; (\operatorname{real\text{-}part}(f(n)) = g(n)) \wedge \\ \quad \forall n \in \mathbb{N}.\; (\operatorname{imag\text{-}part}(f(n)) = q(n)) \Rightarrow \\ \quad \quad (\operatorname{converges\text{-}to}(\operatorname{cc\text{-}ms}, f, v) \Leftrightarrow (\operatorname{converges\text{-}to}(\operatorname{rr\text{-}ms}, g, \operatorname{real\text{-}part}(v)) \wedge \operatorname{converges\text{-}to}(\operatorname{rr\text{-}ms}, q, \operatorname{imag\text{-}part}(v)))) \end{array}}}
\vnbentry{\texttt{cc-\allowbreak{}continuous-\allowbreak{}at-\allowbreak{}iff-\allowbreak{}coords}}{\vnbstmt{\begin{array}{@{}l@{}} \forall s, f, r, i, a. \\ \quad \operatorname{is\text{-}metric\text{-}space}(s) \wedge \\ \quad (f \in (\operatorname{pts}(s) \to \mathbb{C})) \wedge \\ \quad (r \in (\operatorname{pts}(s) \to \mathbb{R})) \wedge \\ \quad (i \in (\operatorname{pts}(s) \to \mathbb{R})) \wedge \\ \quad (a \in \operatorname{pts}(s)) \wedge \\ \quad \forall u \in \operatorname{pts}(s). \\ \quad \quad (\operatorname{real\text{-}part}(f(u)) = r(u)) \wedge \\ \quad \forall u \in \operatorname{pts}(s). \\ \quad \quad (\operatorname{imag\text{-}part}(f(u)) = i(u)) \Rightarrow \\ \quad \quad (\operatorname{is\text{-}continuous\text{-}at}(s, \operatorname{cc\text{-}ms}, f, a) \Leftrightarrow (\operatorname{is\text{-}continuous\text{-}at}(s, \operatorname{rr\text{-}ms}, r, a) \wedge \operatorname{is\text{-}continuous\text{-}at}(s, \operatorname{rr\text{-}ms}, i, a))) \end{array}}}
\vnbentry{\texttt{cc-\allowbreak{}diff-\allowbreak{}at-\allowbreak{}iff-\allowbreak{}coords}}{\vnbstmt{\begin{array}{@{}l@{}} \forall g \in (\mathbb{R} \to \mathbb{C}), r, i \in (\mathbb{R} \to \mathbb{R}), v \in \mathbb{R}, a \in \mathbb{C}. \\ \quad \forall x \in \mathbb{R}.\; (\operatorname{real\text{-}part}(g(x)) = r(x)) \wedge \\ \quad \forall x \in \mathbb{R}.\; (\operatorname{imag\text{-}part}(g(x)) = i(x)) \Rightarrow \\ \quad \quad (\operatorname{is\text{-}cc\text{-}diff\text{-}at}(g, v, a) \Leftrightarrow (\operatorname{is\text{-}diff\text{-}at}(r, v, \operatorname{real\text{-}part}(a)) \wedge \operatorname{is\text{-}diff\text{-}at}(i, v, \operatorname{imag\text{-}part}(a)))) \end{array}}}
\vnbentry{\texttt{cc-\allowbreak{}diff-\allowbreak{}at-\allowbreak{}sum}}{\vnbstmt{\begin{array}{@{}l@{}} \forall g, h, m \in (\mathbb{R} \to \mathbb{C}), v \in \mathbb{R}, a, b \in \mathbb{C}. \\ \quad \forall x \in \mathbb{R}.\; (m(x) = (g(x) + h(x))) \wedge \\ \quad \operatorname{is\text{-}cc\text{-}diff\text{-}at}(g, v, a) \wedge \\ \quad \operatorname{is\text{-}cc\text{-}diff\text{-}at}(h, v, b) \Rightarrow \\ \quad \quad \operatorname{is\text{-}cc\text{-}diff\text{-}at}(m, v, (a + b)) \end{array}}}
\vnbentry{\texttt{cc-\allowbreak{}diff-\allowbreak{}at-\allowbreak{}real-\allowbreak{}mul}}{\vnbstmt{\begin{array}{@{}l@{}} \forall c \in \mathbb{R}, g, d \in (\mathbb{R} \to \mathbb{C}), v \in \mathbb{R}, a \in \mathbb{C}. \\ \quad \forall x \in \mathbb{R}.\; (d(x) = (c \cdot g(x))) \wedge \\ \quad \operatorname{is\text{-}cc\text{-}diff\text{-}at}(g, v, a) \Rightarrow \\ \quad \quad \operatorname{is\text{-}cc\text{-}diff\text{-}at}(d, v, (c \cdot a)) \end{array}}}
\vnbentry{\texttt{cc-\allowbreak{}diff-\allowbreak{}at-\allowbreak{}product}}{\vnbstmt{\begin{array}{@{}l@{}} \forall g, h, r \in (\mathbb{R} \to \mathbb{C}), v \in \mathbb{R}, a, b \in \mathbb{C}. \\ \quad \forall x \in \mathbb{R}.\; (r(x) = (g(x) \cdot h(x))) \wedge \\ \quad \operatorname{is\text{-}cc\text{-}diff\text{-}at}(g, v, a) \wedge \\ \quad \operatorname{is\text{-}cc\text{-}diff\text{-}at}(h, v, b) \Rightarrow \\ \quad \quad \operatorname{is\text{-}cc\text{-}diff\text{-}at}(r, v, ((a \cdot h(v)) + (g(v) \cdot b))) \end{array}}}

\vnbfile{theorem-library/cc-diff-laws-2.scm}{theorem-library/cc-diff-laws-2.scm}
\noindent{\small cc-diff-laws-2.scm -- the rest of the elementary theory of IS-CC-DIFF-AT, the derivative of a map RR -$>$ CC (structure-library/cc-coords.scm). Continues theorem-library/cc-coords-laws.scm (batch 14-A), which defined the predicate coordinatewise and proved the transfer equivalence}\par
\vnbentry{\texttt{cc-\allowbreak{}diff-\allowbreak{}at-\allowbreak{}implies-\allowbreak{}continuous}}{\vnbstmt{\begin{array}{@{}l@{}} \forall g \in (\mathbb{R} \to \mathbb{C}), v \in \mathbb{R}, a \in \mathbb{C}. \\ \quad \operatorname{is\text{-}cc\text{-}diff\text{-}at}(g, v, a) \Rightarrow  \\ \quad \quad \operatorname{is\text{-}continuous\text{-}at}(\operatorname{rr\text{-}ms}, \operatorname{cc\text{-}ms}, g, v) \end{array}}}
\vnbentry{\texttt{cc-\allowbreak{}diff-\allowbreak{}at-\allowbreak{}const}}{\vnbstmt{\begin{array}{@{}l@{}} \forall v \in \mathbb{C}, g \in (\mathbb{R} \to \mathbb{C}), \mathit{tv} \in \mathbb{R}. \\ \quad \forall x \in \mathbb{R}.\; (g(x) = v) \Rightarrow  \\ \quad \quad \operatorname{is\text{-}cc\text{-}diff\text{-}at}(g, \mathit{tv}, 0) \end{array}}}
\vnbentry{\texttt{cc-\allowbreak{}diff-\allowbreak{}at-\allowbreak{}cc-\allowbreak{}mul}}{\vnbstmt{\begin{array}{@{}l@{}} \forall v \in \mathbb{C}, g, d \in (\mathbb{R} \to \mathbb{C}), \mathit{tv} \in \mathbb{R}, a \in \mathbb{C}. \\ \quad \forall x \in \mathbb{R}.\; (d(x) = (v \cdot g(x))) \wedge \\ \quad \operatorname{is\text{-}cc\text{-}diff\text{-}at}(g, \mathit{tv}, a) \Rightarrow \\ \quad \quad \operatorname{is\text{-}cc\text{-}diff\text{-}at}(d, \mathit{tv}, (v \cdot a)) \end{array}}}
\vnbentry{\texttt{cc-\allowbreak{}diff-\allowbreak{}at-\allowbreak{}difference}}{\vnbstmt{\begin{array}{@{}l@{}} \forall g, h, b \in (\mathbb{R} \to \mathbb{C}), v \in \mathbb{R}, a, c \in \mathbb{C}. \\ \quad \forall x \in \mathbb{R}.\; (b(x) = (g(x) - h(x))) \wedge \\ \quad \operatorname{is\text{-}cc\text{-}diff\text{-}at}(g, v, a) \wedge \\ \quad \operatorname{is\text{-}cc\text{-}diff\text{-}at}(h, v, c) \Rightarrow \\ \quad \quad \operatorname{is\text{-}cc\text{-}diff\text{-}at}(b, v, (a - c)) \end{array}}}
\vnbentry{\texttt{cc-\allowbreak{}diff-\allowbreak{}at-\allowbreak{}unique}}{\vnbstmt{\begin{array}{@{}l@{}} \forall g \in (\mathbb{R} \to \mathbb{C}), v \in \mathbb{R}, a, b \in \mathbb{C}. \\ \quad \operatorname{is\text{-}cc\text{-}diff\text{-}at}(g, v, a) \wedge \\ \quad \operatorname{is\text{-}cc\text{-}diff\text{-}at}(g, v, b) \Rightarrow \\ \quad \quad (a = b) \end{array}}}
\vnbentry{\texttt{cc-\allowbreak{}diff-\allowbreak{}at-\allowbreak{}chain-\allowbreak{}real}}{\vnbstmt{\begin{array}{@{}l@{}} \forall g \in (\mathbb{R} \to \mathbb{C}), h \in (\mathbb{R} \to \mathbb{R}), m \in (\mathbb{R} \to \mathbb{C}), v, a \in \mathbb{R}, b \in \mathbb{C}. \\ \quad \forall x \in \mathbb{R}.\; (m(x) = g(h(x))) \wedge \\ \quad \operatorname{is\text{-}diff\text{-}at}(h, v, a) \wedge \\ \quad \operatorname{is\text{-}cc\text{-}diff\text{-}at}(g, h(v), b) \Rightarrow \\ \quad \quad \operatorname{is\text{-}cc\text{-}diff\text{-}at}(m, v, (a \cdot b)) \end{array}}}

\vnbfile{theorem-library/regulated-primitive-laws.scm}{theorem-library/regulated-primitive-laws.scm}
\noindent{\small regulated-primitive-laws.scm -- the laws of IS-COUNTABLE, IS-RIGHT-LIMIT-WITHIN / IS-LEFT-LIMIT-WITHIN, IS-REGULATED-ON and IS-PRIMITIVE (structure-library/regulated-primitive.scm).  Batch 20-D, 2026-09-23.  THE SPECIFICATION is docs/roads-regulated-primitives-2026-09-23.md.}\par
\vnbentry{\texttt{countable-\allowbreak{}of-\allowbreak{}enum}}{\vnbstmt{\begin{array}{@{}l@{}} \forall d, e. \\ \quad (d \in \mathbf{Set}) \wedge \\ \quad (e \in (\mathbb{N} \to d)) \wedge \\ \quad \forall y \in d.\; \exists n \in \mathbb{N}.\; (y = e(n)) \Rightarrow \\ \quad \quad \operatorname{is\text{-}countable}(d) \end{array}}}
\vnbentry{\texttt{countable-\allowbreak{}empty}}{\vnbstmt{\operatorname{is\text{-}countable}(\emptyset)}}
\vnbentry{\texttt{rr-\allowbreak{}abs-\allowbreak{}lt-\allowbreak{}parts}}{\vnbstmt{\begin{array}{@{}l@{}} \forall x, c \in \mathbb{R}. \\ \quad ((\lvert x \rvert < c) \Rightarrow ((-c < x) \wedge (x < c))) \end{array}}}
\vnbentry{\texttt{countable-\allowbreak{}singleton}}{\vnbstmt{\forall x \in \mathbf{Set}.\; \operatorname{is\text{-}countable}(\{x\})}}
\vnbentry{\texttt{countable-\allowbreak{}subset}}{\vnbstmt{\begin{array}{@{}l@{}} \forall a, d. \\ \quad (\operatorname{is\text{-}countable}(d) \wedge (a \subseteq d)) \Rightarrow  \\ \quad \quad \operatorname{is\text{-}countable}(a) \end{array}}}
\vnbentry{\texttt{countable-\allowbreak{}cases}}{\vnbstmt{\begin{array}{@{}l@{}} \forall d. \\ \quad \operatorname{is\text{-}countable}(d) \Rightarrow  \\ \quad \quad (d \in \mathbf{Set}) \wedge \\ \quad \quad (d = \emptyset) \vee \\ \quad \quad \exists e \in (\mathbb{N} \to d). \forall y \in d. \exists n \in \mathbb{N}. \\ \quad \quad \quad (y = e(n)) \end{array}}}
\vnbentry{\texttt{countable-\allowbreak{}of-\allowbreak{}two-\allowbreak{}enum}}{\vnbstmt{\begin{array}{@{}l@{}} \forall u, f, g. \\ \quad (u \in \mathbf{Set}) \wedge \\ \quad (f \in (\mathbb{N} \to u)) \wedge \\ \quad (g \in (\mathbb{N} \to u)) \wedge \\ \quad \forall y \in u. \\ \quad \quad (\exists j \in \mathbb{N}.\; (y = f(j)) \vee \exists j \in \mathbb{N}.\; (y = g(j))) \Rightarrow \\ \quad \quad \operatorname{is\text{-}countable}(u) \end{array}}}
\vnbentry{\texttt{countable-\allowbreak{}union-\allowbreak{}2}}{\vnbstmt{\begin{array}{@{}l@{}} \forall a, b. \\ \quad (\operatorname{is\text{-}countable}(a) \wedge \operatorname{is\text{-}countable}(b)) \Rightarrow  \\ \quad \quad \operatorname{is\text{-}countable}((a \cup b)) \end{array}}}
\vnbentry{\texttt{finite-\allowbreak{}implies-\allowbreak{}countable}}{\vnbstmt{\begin{array}{@{}l@{}} \forall d \in \mathbf{Set}. \\ \quad ((|d| \in \mathbb{N}) \Rightarrow \operatorname{is\text{-}countable}(d)) \end{array}}}
\vnbentry{\texttt{countable-\allowbreak{}image}}{\vnbstmt{\begin{array}{@{}l@{}} \forall d, h, a, b. \\ \quad (\operatorname{is\text{-}countable}(d) \wedge ((h \in (a \to b)) \wedge (d \subseteq a))) \Rightarrow  \\ \quad \quad \operatorname{is\text{-}countable}(\operatorname{image}(h, d)) \end{array}}}
\vnbentry{\texttt{rr-\allowbreak{}abs-\allowbreak{}le-\allowbreak{}parts}}{\vnbstmt{\begin{array}{@{}l@{}} \forall x, c \in \mathbb{R}. \\ \quad ((\lvert x \rvert \leq c) \Rightarrow ((-c \leq x) \wedge (x \leq c))) \end{array}}}
\vnbentry{\texttt{rr-\allowbreak{}abs-\allowbreak{}le-\allowbreak{}of-\allowbreak{}parts}}{\vnbstmt{\begin{array}{@{}l@{}} \forall x, c \in \mathbb{R}. \\ \quad (((-c \leq x) \wedge (x \leq c)) \Rightarrow (\lvert x \rvert \leq c)) \end{array}}}
\vnbentry{\texttt{rr-\allowbreak{}abs-\allowbreak{}lt-\allowbreak{}of-\allowbreak{}parts}}{\vnbstmt{\begin{array}{@{}l@{}} \forall x, c \in \mathbb{R}. \\ \quad (((-c < x) \wedge (x < c)) \Rightarrow (\lvert x \rvert < c)) \end{array}}}
\vnbentry{\texttt{continuous-\allowbreak{}on-\allowbreak{}right-\allowbreak{}limit-\allowbreak{}within}}{\vnbstmt{\begin{array}{@{}l@{}} \forall f, w.; \forall x \in w. \\ \quad ((w \subseteq \mathbb{R}) \wedge \operatorname{is\text{-}continuous\text{-}on}(f, w)) \Rightarrow  \\ \quad \quad \operatorname{is\text{-}right\text{-}limit\text{-}within}(f, w, x, f(x)) \end{array}}}
\vnbentry{\texttt{continuous-\allowbreak{}on-\allowbreak{}left-\allowbreak{}limit-\allowbreak{}within}}{\vnbstmt{\begin{array}{@{}l@{}} \forall f, w.; \forall x \in w. \\ \quad ((w \subseteq \mathbb{R}) \wedge \operatorname{is\text{-}continuous\text{-}on}(f, w)) \Rightarrow  \\ \quad \quad \operatorname{is\text{-}left\text{-}limit\text{-}within}(f, w, x, f(x)) \end{array}}}
\vnbentry{\texttt{continuous-\allowbreak{}on-\allowbreak{}implies-\allowbreak{}regulated-\allowbreak{}on}}{\vnbstmt{\begin{array}{@{}l@{}} \forall f, a, b. \\ \quad (a \in \mathbb{R}) \wedge \\ \quad (b \in \mathbb{R}) \wedge \\ \quad (a < b) \wedge \\ \quad \operatorname{is\text{-}continuous\text{-}on}(f, \operatorname{ccint}(a, b)) \Rightarrow \\ \quad \quad \operatorname{is\text{-}regulated\text{-}on}(f, a, b) \end{array}}}
\vnbentry{\texttt{right-\allowbreak{}limit-\allowbreak{}within-\allowbreak{}le}}{\vnbstmt{\begin{array}{@{}l@{}} \forall f, w, x, l, m. \\ \quad \operatorname{is\text{-}right\text{-}limit\text{-}within}(f, w, x, l) \wedge \\ \quad \operatorname{is\text{-}right\text{-}limit\text{-}within}(f, w, x, m) \wedge \\ \quad \forall r. \\ \quad \quad \operatorname{pos\text{-}rr}(r) \Rightarrow  \\ \quad \quad \quad \exists t \in w.\; ((x < t) \wedge (t < (x + r))) \Rightarrow \\ \quad \quad ((l - m) \leq 0) \end{array}}}
\vnbentry{\texttt{left-\allowbreak{}limit-\allowbreak{}within-\allowbreak{}le}}{\vnbstmt{\begin{array}{@{}l@{}} \forall f, w, x, l, m. \\ \quad \operatorname{is\text{-}left\text{-}limit\text{-}within}(f, w, x, l) \wedge \\ \quad \operatorname{is\text{-}left\text{-}limit\text{-}within}(f, w, x, m) \wedge \\ \quad \forall r. \\ \quad \quad \operatorname{pos\text{-}rr}(r) \Rightarrow  \\ \quad \quad \quad \exists t \in w.\; (((x - r) < t) \wedge (t < x)) \Rightarrow \\ \quad \quad ((l - m) \leq 0) \end{array}}}
\vnbentry{\texttt{right-\allowbreak{}limit-\allowbreak{}within-\allowbreak{}unique}}{\vnbstmt{\begin{array}{@{}l@{}} \forall f, w, x, l, m. \\ \quad \operatorname{is\text{-}right\text{-}limit\text{-}within}(f, w, x, l) \wedge \\ \quad \operatorname{is\text{-}right\text{-}limit\text{-}within}(f, w, x, m) \wedge \\ \quad \forall r. \\ \quad \quad \operatorname{pos\text{-}rr}(r) \Rightarrow  \\ \quad \quad \quad \exists t \in w.\; ((x < t) \wedge (t < (x + r))) \Rightarrow \\ \quad \quad (l = m) \end{array}}}
\vnbentry{\texttt{left-\allowbreak{}limit-\allowbreak{}within-\allowbreak{}unique}}{\vnbstmt{\begin{array}{@{}l@{}} \forall f, w, x, l, m. \\ \quad \operatorname{is\text{-}left\text{-}limit\text{-}within}(f, w, x, l) \wedge \\ \quad \operatorname{is\text{-}left\text{-}limit\text{-}within}(f, w, x, m) \wedge \\ \quad \forall r. \\ \quad \quad \operatorname{pos\text{-}rr}(r) \Rightarrow  \\ \quad \quad \quad \exists t \in w.\; (((x - r) < t) \wedge (t < x)) \Rightarrow \\ \quad \quad (l = m) \end{array}}}
\vnbentry{\texttt{ccint-\allowbreak{}right-\allowbreak{}approach}}{\vnbstmt{\begin{array}{@{}l@{}} \forall a, b, x \in \mathbb{R}, r. \\ \quad ((a \leq x) \wedge ((x < b) \wedge \operatorname{pos\text{-}rr}(r))) \Rightarrow  \\ \quad \quad \exists t \in \operatorname{ccint}(a, b).\; ((x < t) \wedge (t < (x + r))) \end{array}}}
\vnbentry{\texttt{ccint-\allowbreak{}left-\allowbreak{}approach}}{\vnbstmt{\begin{array}{@{}l@{}} \forall a, b, x \in \mathbb{R}, r. \\ \quad ((a < x) \wedge ((x \leq b) \wedge \operatorname{pos\text{-}rr}(r))) \Rightarrow  \\ \quad \quad \exists t \in \operatorname{ccint}(a, b).\; (((x - r) < t) \wedge (t < x)) \end{array}}}
\vnbentry{\texttt{right-\allowbreak{}limit-\allowbreak{}within-\allowbreak{}ccint-\allowbreak{}unique}}{\vnbstmt{\begin{array}{@{}l@{}} \forall f.; \forall a, b, x \in \mathbb{R}, l, m. \\ \quad (a \leq x) \wedge \\ \quad (x < b) \wedge \\ \quad \operatorname{is\text{-}right\text{-}limit\text{-}within}(f, \operatorname{ccint}(a, b), x, l) \wedge \\ \quad \operatorname{is\text{-}right\text{-}limit\text{-}within}(f, \operatorname{ccint}(a, b), x, m) \Rightarrow \\ \quad \quad (l = m) \end{array}}}
\vnbentry{\texttt{left-\allowbreak{}limit-\allowbreak{}within-\allowbreak{}ccint-\allowbreak{}unique}}{\vnbstmt{\begin{array}{@{}l@{}} \forall f.; \forall a, b, x \in \mathbb{R}, l, m. \\ \quad (a < x) \wedge \\ \quad (x \leq b) \wedge \\ \quad \operatorname{is\text{-}left\text{-}limit\text{-}within}(f, \operatorname{ccint}(a, b), x, l) \wedge \\ \quad \operatorname{is\text{-}left\text{-}limit\text{-}within}(f, \operatorname{ccint}(a, b), x, m) \Rightarrow \\ \quad \quad (l = m) \end{array}}}
\vnbentry{\texttt{right-\allowbreak{}limit-\allowbreak{}within-\allowbreak{}subset}}{\vnbstmt{\begin{array}{@{}l@{}} \forall f, w, x, l, g, v. \\ \quad \operatorname{is\text{-}right\text{-}limit\text{-}within}(f, w, x, l) \wedge \\ \quad (v \subseteq w) \wedge \\ \quad (g \in (v \to \mathbb{R})) \wedge \\ \quad \forall z \in v.\; (g(z) = f(z)) \Rightarrow \\ \quad \quad \operatorname{is\text{-}right\text{-}limit\text{-}within}(g, v, x, l) \end{array}}}
\vnbentry{\texttt{left-\allowbreak{}limit-\allowbreak{}within-\allowbreak{}subset}}{\vnbstmt{\begin{array}{@{}l@{}} \forall f, w, x, l, g, v. \\ \quad \operatorname{is\text{-}left\text{-}limit\text{-}within}(f, w, x, l) \wedge \\ \quad (v \subseteq w) \wedge \\ \quad (g \in (v \to \mathbb{R})) \wedge \\ \quad \forall z \in v.\; (g(z) = f(z)) \Rightarrow \\ \quad \quad \operatorname{is\text{-}left\text{-}limit\text{-}within}(g, v, x, l) \end{array}}}
\vnbentry{\texttt{regulated-\allowbreak{}on-\allowbreak{}restrict}}{\vnbstmt{\begin{array}{@{}l@{}} \forall f, a, b, c, d, g. \\ \quad \operatorname{is\text{-}regulated\text{-}on}(f, a, b) \wedge \\ \quad (c \in \mathbb{R}) \wedge \\ \quad (d \in \mathbb{R}) \wedge \\ \quad (a \leq c) \wedge \\ \quad (c < d) \wedge \\ \quad (d \leq b) \wedge \\ \quad (g \in (\operatorname{ccint}(c, d) \to \mathbb{R})) \wedge \\ \quad \forall z \in \operatorname{ccint}(c, d).\; (g(z) = f(z)) \Rightarrow \\ \quad \quad \operatorname{is\text{-}regulated\text{-}on}(g, c, d) \end{array}}}
\vnbentry{\texttt{primitive-\allowbreak{}endpoints}}{\vnbstmt{\begin{array}{@{}l@{}} \forall f, i, a, b. \\ \quad \operatorname{is\text{-}primitive}(f, i, a, b) \Rightarrow  \\ \quad \quad ((a \in \mathbb{R}) \wedge ((b \in \mathbb{R}) \wedge (a < b))) \end{array}}}
\vnbentry{\texttt{primitive-\allowbreak{}in-\allowbreak{}fun}}{\vnbstmt{\begin{array}{@{}l@{}} \forall f, i, a, b. \\ \quad \operatorname{is\text{-}primitive}(f, i, a, b) \Rightarrow  \\ \quad \quad (f \in (\operatorname{ccint}(a, b) \to \mathbb{R})) \end{array}}}
\vnbentry{\texttt{primitive-\allowbreak{}integrand-\allowbreak{}in-\allowbreak{}fun}}{\vnbstmt{\begin{array}{@{}l@{}} \forall f, i, a, b. \\ \quad \operatorname{is\text{-}primitive}(f, i, a, b) \Rightarrow  \\ \quad \quad (i \in (\operatorname{ccint}(a, b) \to \mathbb{R})) \end{array}}}
\vnbentry{\texttt{primitive-\allowbreak{}continuous}}{\vnbstmt{\begin{array}{@{}l@{}} \forall f, i, a, b. \\ \quad \operatorname{is\text{-}primitive}(f, i, a, b) \Rightarrow  \\ \quad \quad \operatorname{is\text{-}continuous\text{-}on}(f, \operatorname{ccint}(a, b)) \end{array}}}
\vnbentry{\texttt{primitive-\allowbreak{}exceptional-\allowbreak{}set}}{\vnbstmt{\begin{array}{@{}l@{}} \forall f, i, a, b. \\ \quad \operatorname{is\text{-}primitive}(f, i, a, b) \Rightarrow  \\ \quad \quad \exists s. \\ \quad \quad \quad \operatorname{is\text{-}countable}(s) \wedge \\ \quad \quad \quad (s \subseteq \operatorname{ccint}(a, b)) \wedge \\ \quad \quad \quad \forall t \in \operatorname{ooint}(a, b). \\ \quad \quad \quad \quad \neg (t \in s) \Rightarrow  \\ \quad \quad \quad \quad \quad \operatorname{has\text{-}deriv\text{-}at}(f, t, i(t)) \end{array}}}

\vnbfile{theorem-library/regulated-step-approx.scm}{theorem-library/regulated-step-approx.scm}
\noindent{\small regulated-step-approx.scm -- Dieudonne (7.6.1), NECESSITY: a function regulated on [a, b] is the uniform limit on [a, b] of step functions. (Batch 22-A, 2026-09-23.)}\par
\vnbentry{\texttt{strict-\allowbreak{}seq-\allowbreak{}append-\allowbreak{}point}}{\vnbstmt{\begin{array}{@{}l@{}} \forall p, n, d. \\ \quad (p \in (\mathbb{N} \to \mathbb{R})) \wedge \\ \quad (n \in \mathbb{N}) \wedge \\ \quad (d \in \mathbb{R}) \wedge \\ \quad \forall i, j. \\ \quad \quad (i \in \operatorname{interval}(0, n)) \wedge \\ \quad \quad (j \in \operatorname{interval}(0, n)) \wedge \\ \quad \quad (i < j) \Rightarrow \\ \quad \quad \quad (p(i) < p(j)) \wedge \\ \quad (p(n) < d) \Rightarrow \\ \quad \quad \exists q \in (\mathbb{N} \to \mathbb{R}). \\ \quad \quad \quad (q((n + 1)) = d) \wedge \\ \quad \quad \quad \forall k \in \mathbb{N}.\; ((k \leq n) \Rightarrow (q(k) = p(k))) \wedge \\ \quad \quad \quad \forall i, j. \\ \quad \quad \quad \quad (i \in \operatorname{interval}(0, (n + 1))) \wedge \\ \quad \quad \quad \quad (j \in \operatorname{interval}(0, (n + 1))) \wedge \\ \quad \quad \quad \quad (i < j) \Rightarrow \\ \quad \quad \quad \quad \quad (q(i) < q(j)) \end{array}}}
\vnbentry{\texttt{fun-\allowbreak{}append-\allowbreak{}constant-\allowbreak{}piece}}{\vnbstmt{\begin{array}{@{}l@{}} \forall h \in (\mathbb{R} \to \mathbb{R}), c, d, k, v \in \mathbb{R}. \\ \quad (c < d) \Rightarrow  \\ \quad \quad \exists g \in (\mathbb{R} \to \mathbb{R}). \\ \quad \quad \quad \forall t \in \mathbb{R}.\; ((t \leq c) \Rightarrow (g(t) = h(t))) \wedge \\ \quad \quad \quad \forall t \in \mathbb{R}.\; (((c < t) \wedge (t < d)) \Rightarrow (g(t) = k)) \wedge \\ \quad \quad \quad (g(d) = v) \end{array}}}
\vnbentry{\texttt{step-\allowbreak{}pieces-\allowbreak{}append}}{\vnbstmt{\begin{array}{@{}l@{}} \forall p, n, h, q, g, c, d, k.; \forall i \in \operatorname{interval}(0, (n + 1)). \\ \quad (p \in (\mathbb{N} \to \mathbb{R})) \wedge \\ \quad (n \in \mathbb{N}) \wedge \\ \quad (h \in (\mathbb{R} \to \mathbb{R})) \wedge \\ \quad (q \in (\mathbb{N} \to \mathbb{R})) \wedge \\ \quad (g \in (\mathbb{R} \to \mathbb{R})) \wedge \\ \quad (c \in \mathbb{R}) \wedge \\ \quad (d \in \mathbb{R}) \wedge \\ \quad (k \in \mathbb{R}) \wedge \\ \quad \forall \mathit{rsi}, j. \\ \quad \quad (\mathit{rsi} \in \operatorname{interval}(0, n)) \wedge \\ \quad \quad (j \in \operatorname{interval}(0, n)) \wedge \\ \quad \quad (\mathit{rsi} < j) \Rightarrow \\ \quad \quad \quad (p(\mathit{rsi}) < p(j)) \wedge \\ \quad \forall i \in \operatorname{interval}(0, n). \\ \quad \quad (i < n) \Rightarrow  \\ \quad \quad \quad \exists \mathit{rspc} \in \mathbb{R}. \forall t \in \mathbb{R}. \\ \quad \quad \quad \quad (((p(i) < t) \wedge (t < p((i + 1)))) \Rightarrow (h(t) = \mathit{rspc})) \wedge \\ \quad (p(n) = c) \wedge \\ \quad \forall j \in \mathbb{N}.\; ((j \leq n) \Rightarrow (q(j) = p(j))) \wedge \\ \quad (q((n + 1)) = d) \wedge \\ \quad \forall \mathit{rst} \in \mathbb{R}. \\ \quad \quad ((\mathit{rst} \leq c) \Rightarrow (g(\mathit{rst}) = h(\mathit{rst}))) \wedge \\ \quad \forall \mathit{rst} \in \mathbb{R}. \\ \quad \quad ((c < \mathit{rst}) \wedge (\mathit{rst} < d)) \Rightarrow  \\ \quad \quad \quad (g(\mathit{rst}) = k) \wedge \\ \quad (i < (n + 1)) \Rightarrow \\ \quad \quad \exists \mathit{rspc} \in \mathbb{R}. \forall t \in \mathbb{R}. \\ \quad \quad \quad (((q(i) < t) \wedge (t < q((i + 1)))) \Rightarrow (g(t) = \mathit{rspc})) \end{array}}}
\vnbentry{\texttt{step-\allowbreak{}approx-\allowbreak{}start}}{\vnbstmt{\begin{array}{@{}l@{}} \forall a, f, e. \\ \quad (a \in \mathbb{R}) \wedge \\ \quad (f \in (\mathbb{R} \to \mathbb{R})) \wedge \\ \quad (e \in \mathbb{R}) \wedge \\ \quad (0 \leq e) \Rightarrow \\ \quad \quad \exists n \in \mathbb{N}, p \in (\mathbb{N} \to \mathbb{R}), h \in (\mathbb{R} \to \mathbb{R}). \\ \quad \quad \quad (p(0) = a) \wedge \\ \quad \quad \quad (p(n) = a) \wedge \\ \quad \quad \quad \forall i, j. \\ \quad \quad \quad \quad (i \in \operatorname{interval}(0, n)) \wedge \\ \quad \quad \quad \quad (j \in \operatorname{interval}(0, n)) \wedge \\ \quad \quad \quad \quad (i < j) \Rightarrow \\ \quad \quad \quad \quad \quad (p(i) < p(j)) \wedge \\ \quad \quad \quad \forall \mathit{rspi} \in \operatorname{interval}(0, n). \\ \quad \quad \quad \quad (\mathit{rspi} < n) \Rightarrow  \\ \quad \quad \quad \quad \quad \exists c \in \mathbb{R}. \forall t \in \mathbb{R}. \\ \quad \quad \quad \quad \quad \quad ((p(\mathit{rspi}) < t) \wedge (t < p((\mathit{rspi} + 1)))) \Rightarrow  \\ \quad \quad \quad \quad \quad \quad \quad (h(t) = c) \wedge \\ \quad \quad \quad \forall z \in \operatorname{ccint}(a, a).\; (\lvert (f(z) - h(z)) \rvert \leq e) \end{array}}}
\vnbentry{\texttt{step-\allowbreak{}approx-\allowbreak{}append}}{\vnbstmt{\begin{array}{@{}l@{}} \forall a, x, y, f, e, k. \\ \quad (a \in \mathbb{R}) \wedge \\ \quad (x \in \mathbb{R}) \wedge \\ \quad (y \in \mathbb{R}) \wedge \\ \quad (f \in (\mathbb{R} \to \mathbb{R})) \wedge \\ \quad (e \in \mathbb{R}) \wedge \\ \quad (k \in \mathbb{R}) \wedge \\ \quad (0 \leq e) \wedge \\ \quad \exists n \in \mathbb{N}, p \in (\mathbb{N} \to \mathbb{R}), h \in (\mathbb{R} \to \mathbb{R}). \\ \quad \quad (p(0) = a) \wedge \\ \quad \quad (p(n) = x) \wedge \\ \quad \quad \forall i, j. \\ \quad \quad \quad (i \in \operatorname{interval}(0, n)) \wedge \\ \quad \quad \quad (j \in \operatorname{interval}(0, n)) \wedge \\ \quad \quad \quad (i < j) \Rightarrow \\ \quad \quad \quad \quad (p(i) < p(j)) \wedge \\ \quad \quad \forall \mathit{rspi} \in \operatorname{interval}(0, n). \\ \quad \quad \quad (\mathit{rspi} < n) \Rightarrow  \\ \quad \quad \quad \quad \exists c \in \mathbb{R}. \forall t \in \mathbb{R}. \\ \quad \quad \quad \quad \quad ((p(\mathit{rspi}) < t) \wedge (t < p((\mathit{rspi} + 1)))) \Rightarrow  \\ \quad \quad \quad \quad \quad \quad (h(t) = c) \wedge \\ \quad \quad \forall z \in \operatorname{ccint}(a, x).\; (\lvert (f(z) - h(z)) \rvert \leq e) \wedge \\ \quad (x < y) \wedge \\ \quad \forall s \in \mathbb{R}. \\ \quad \quad (((x < s) \wedge (s < y)) \Rightarrow (\lvert (f(s) - k) \rvert \leq e)) \Rightarrow \\ \quad \quad \exists n \in \mathbb{N}, p \in (\mathbb{N} \to \mathbb{R}), h \in (\mathbb{R} \to \mathbb{R}). \\ \quad \quad \quad (p(0) = a) \wedge \\ \quad \quad \quad (p(n) = y) \wedge \\ \quad \quad \quad \forall i, j. \\ \quad \quad \quad \quad (i \in \operatorname{interval}(0, n)) \wedge \\ \quad \quad \quad \quad (j \in \operatorname{interval}(0, n)) \wedge \\ \quad \quad \quad \quad (i < j) \Rightarrow \\ \quad \quad \quad \quad \quad (p(i) < p(j)) \wedge \\ \quad \quad \quad \forall \mathit{rspi} \in \operatorname{interval}(0, n). \\ \quad \quad \quad \quad (\mathit{rspi} < n) \Rightarrow  \\ \quad \quad \quad \quad \quad \exists c \in \mathbb{R}. \forall t \in \mathbb{R}. \\ \quad \quad \quad \quad \quad \quad ((p(\mathit{rspi}) < t) \wedge (t < p((\mathit{rspi} + 1)))) \Rightarrow  \\ \quad \quad \quad \quad \quad \quad \quad (h(t) = c) \wedge \\ \quad \quad \quad \forall z \in \operatorname{ccint}(a, y).\; (\lvert (f(z) - h(z)) \rvert \leq e) \end{array}}}
\vnbentry{\texttt{regulated-\allowbreak{}on-\allowbreak{}right-\allowbreak{}window}}{\vnbstmt{\begin{array}{@{}l@{}} \forall a, b, f, g, e, t. \\ \quad \operatorname{is\text{-}regulated\text{-}on}(f, a, b) \wedge \\ \quad (g \in (\mathbb{R} \to \mathbb{R})) \wedge \\ \quad \forall z \in \operatorname{ccint}(a, b).\; (g(z) = f(z)) \wedge \\ \quad \operatorname{pos\text{-}rr}(e) \wedge \\ \quad (t \in \operatorname{ccint}(a, b)) \Rightarrow \\ \quad \quad \exists d \in \mathbb{R}. \\ \quad \quad \quad (0 < d) \wedge \\ \quad \quad \quad \exists l \in \mathbb{R}. \forall s \in \mathbb{R}. \\ \quad \quad \quad \quad ((t < s) \wedge ((s < (t + d)) \wedge (s \leq b))) \Rightarrow  \\ \quad \quad \quad \quad \quad (\lvert (g(s) - l) \rvert \leq e) \end{array}}}
\vnbentry{\texttt{regulated-\allowbreak{}on-\allowbreak{}left-\allowbreak{}window}}{\vnbstmt{\begin{array}{@{}l@{}} \forall a, b, f, g, e, t. \\ \quad \operatorname{is\text{-}regulated\text{-}on}(f, a, b) \wedge \\ \quad (g \in (\mathbb{R} \to \mathbb{R})) \wedge \\ \quad \forall z \in \operatorname{ccint}(a, b).\; (g(z) = f(z)) \wedge \\ \quad \operatorname{pos\text{-}rr}(e) \wedge \\ \quad (t \in \operatorname{ccint}(a, b)) \Rightarrow \\ \quad \quad \exists d \in \mathbb{R}. \\ \quad \quad \quad (0 < d) \wedge \\ \quad \quad \quad \exists l \in \mathbb{R}. \forall s \in \mathbb{R}. \\ \quad \quad \quad \quad ((a \leq s) \wedge (((t - d) < s) \wedge (s < t))) \Rightarrow  \\ \quad \quad \quad \quad \quad (\lvert (g(s) - l) \rvert \leq e) \end{array}}}
\vnbentry{\texttt{regulated-\allowbreak{}on-\allowbreak{}step-\allowbreak{}approx-\allowbreak{}creep}}{\vnbstmt{\begin{array}{@{}l@{}} \forall a, b, f, g, e. \\ \quad \operatorname{is\text{-}regulated\text{-}on}(f, a, b) \wedge \\ \quad (g \in (\mathbb{R} \to \mathbb{R})) \wedge \\ \quad \forall z \in \operatorname{ccint}(a, b).\; (g(z) = f(z)) \wedge \\ \quad \operatorname{pos\text{-}rr}(e) \Rightarrow \\ \quad \quad \exists n \in \mathbb{N}, p \in (\mathbb{N} \to \mathbb{R}), h \in (\mathbb{R} \to \mathbb{R}). \\ \quad \quad \quad (p(0) = a) \wedge \\ \quad \quad \quad (p(n) = b) \wedge \\ \quad \quad \quad \forall i, j. \\ \quad \quad \quad \quad (i \in \operatorname{interval}(0, n)) \wedge \\ \quad \quad \quad \quad (j \in \operatorname{interval}(0, n)) \wedge \\ \quad \quad \quad \quad (i < j) \Rightarrow \\ \quad \quad \quad \quad \quad (p(i) < p(j)) \wedge \\ \quad \quad \quad \forall \mathit{rspi} \in \operatorname{interval}(0, n). \\ \quad \quad \quad \quad (\mathit{rspi} < n) \Rightarrow  \\ \quad \quad \quad \quad \quad \exists c \in \mathbb{R}. \forall t \in \mathbb{R}. \\ \quad \quad \quad \quad \quad \quad ((p(\mathit{rspi}) < t) \wedge (t < p((\mathit{rspi} + 1)))) \Rightarrow  \\ \quad \quad \quad \quad \quad \quad \quad (h(t) = c) \wedge \\ \quad \quad \quad \forall \mathit{rsqz} \in \operatorname{ccint}(a, b). \\ \quad \quad \quad \quad (\lvert (g(\mathit{rsqz}) - h(\mathit{rsqz})) \rvert \leq e) \end{array}}}
\vnbentry{\texttt{step-\allowbreak{}approx-\allowbreak{}q-\allowbreak{}step-\allowbreak{}fn}}{\vnbstmt{\begin{array}{@{}l@{}} \forall a, b, f, e. \\ \quad (a \in \mathbb{R}) \wedge \\ \quad (b \in \mathbb{R}) \wedge \\ \quad (a < b) \wedge \\ \quad \exists n \in \mathbb{N}, p \in (\mathbb{N} \to \mathbb{R}), h \in (\mathbb{R} \to \mathbb{R}). \\ \quad \quad (p(0) = a) \wedge \\ \quad \quad (p(n) = b) \wedge \\ \quad \quad \forall i, j. \\ \quad \quad \quad (i \in \operatorname{interval}(0, n)) \wedge \\ \quad \quad \quad (j \in \operatorname{interval}(0, n)) \wedge \\ \quad \quad \quad (i < j) \Rightarrow \\ \quad \quad \quad \quad (p(i) < p(j)) \wedge \\ \quad \quad \forall \mathit{rspi} \in \operatorname{interval}(0, n). \\ \quad \quad \quad (\mathit{rspi} < n) \Rightarrow  \\ \quad \quad \quad \quad \exists c \in \mathbb{R}. \forall t \in \mathbb{R}. \\ \quad \quad \quad \quad \quad ((p(\mathit{rspi}) < t) \wedge (t < p((\mathit{rspi} + 1)))) \Rightarrow  \\ \quad \quad \quad \quad \quad \quad (h(t) = c) \wedge \\ \quad \quad \forall z \in \operatorname{ccint}(a, b).\; (\lvert (f(z) - h(z)) \rvert \leq e) \Rightarrow \\ \quad \quad \exists \mathit{rsth} \in (\mathbb{R} \to \mathbb{R}). \\ \quad \quad \quad \operatorname{is\text{-}step\text{-}fn}(\mathit{rsth}, a, b) \wedge \\ \quad \quad \quad \forall z \in \operatorname{ccint}(a, b). \\ \quad \quad \quad \quad (\lvert (f(z) - \operatorname{rsth}(z)) \rvert \leq e) \end{array}}}
\vnbentry{\texttt{regulated-\allowbreak{}on-\allowbreak{}step-\allowbreak{}approx}}{\vnbstmt{\begin{array}{@{}l@{}} \forall a, b, f, e. \\ \quad (\operatorname{is\text{-}regulated\text{-}on}(f, a, b) \wedge \operatorname{pos\text{-}rr}(e)) \Rightarrow  \\ \quad \quad \exists h \in (\mathbb{R} \to \mathbb{R}). \\ \quad \quad \quad \operatorname{is\text{-}step\text{-}fn}(h, a, b) \wedge \\ \quad \quad \quad \forall z \in \operatorname{ccint}(a, b). \\ \quad \quad \quad \quad (\lvert (\operatorname{extend\text{-}const}(f, a, b)(z) - h(z)) \rvert \leq e) \end{array}}}
\vnbentry{\texttt{regulated-\allowbreak{}on-\allowbreak{}is-\allowbreak{}step-\allowbreak{}limit}}{\vnbstmt{\begin{array}{@{}l@{}} \forall a, b \in \mathbb{R}, f. \\ \quad \operatorname{is\text{-}regulated\text{-}on}(f, a, b) \Rightarrow  \\ \quad \quad \exists m \in (\mathbb{N} \to (\mathbb{R} \to \mathbb{R})). \\ \quad \quad \quad \forall k \in \mathbb{N}.\; \operatorname{is\text{-}step\text{-}fn}(m(k), a, b) \wedge \\ \quad \quad \quad \operatorname{is\text{-}unif\text{-}limit\text{-}on}(m, \operatorname{extend\text{-}const}(f, a, b), a, b) \end{array}}}

\vnbfile{theorem-library/has-deriv-at-more.scm}{theorem-library/has-deriv-at-more.scm}
\noindent{\small has-deriv-at-more.scm -- the four HAS-DERIV-AT laws that theorem-library/interval-calculus-laws.scm did not prove: the VALUE typing, the AFFINE function (and with it the constant and the identity), the REAL MULTIPLE and the DIFFERENCE.}\par
\vnbentry{\texttt{has-\allowbreak{}deriv-\allowbreak{}at-\allowbreak{}value-\allowbreak{}in-\allowbreak{}rr}}{\vnbstmt{\begin{array}{@{}l@{}} \forall f, x, l. \\ \quad \operatorname{has\text{-}deriv\text{-}at}(f, x, l) \Rightarrow  \\ \quad \quad (f(x) \in \mathbb{R}) \end{array}}}
\vnbentry{\texttt{has-\allowbreak{}deriv-\allowbreak{}at-\allowbreak{}affine}}{\vnbstmt{\begin{array}{@{}l@{}} \forall r, f.; \forall c, m, x \in \mathbb{R}. \\ \quad \operatorname{pos\text{-}rr}(r) \wedge \\ \quad \operatorname{restrict}(f, \operatorname{ooint}((x - r), (x + r))) \\ \quad \in\; (\operatorname{ooint}((x - r), (x + r)) \to \mathbb{R}) \wedge \\ \quad \forall \mathit{hbx} \in \operatorname{ooint}((x - r), (x + r)). \\ \quad \quad (f(\mathit{hbx}) \simeq (c + (m \cdot \mathit{hbx}))) \Rightarrow \\ \quad \quad \operatorname{has\text{-}deriv\text{-}at}(f, x, m) \end{array}}}
\vnbentry{\texttt{has-\allowbreak{}deriv-\allowbreak{}at-\allowbreak{}const}}{\vnbstmt{\begin{array}{@{}l@{}} \forall r, f.; \forall c, x \in \mathbb{R}. \\ \quad \operatorname{pos\text{-}rr}(r) \wedge \\ \quad \operatorname{restrict}(f, \operatorname{ooint}((x - r), (x + r))) \\ \quad \in\; (\operatorname{ooint}((x - r), (x + r)) \to \mathbb{R}) \wedge \\ \quad \forall \mathit{hbx} \in \operatorname{ooint}((x - r), (x + r)). \\ \quad \quad (f(\mathit{hbx}) \simeq c) \Rightarrow \\ \quad \quad \operatorname{has\text{-}deriv\text{-}at}(f, x, 0) \end{array}}}
\vnbentry{\texttt{has-\allowbreak{}deriv-\allowbreak{}at-\allowbreak{}identity}}{\vnbstmt{\begin{array}{@{}l@{}} \forall r, f.; \forall x \in \mathbb{R}. \\ \quad \operatorname{pos\text{-}rr}(r) \wedge \\ \quad \operatorname{restrict}(f, \operatorname{ooint}((x - r), (x + r))) \\ \quad \in\; (\operatorname{ooint}((x - r), (x + r)) \to \mathbb{R}) \wedge \\ \quad \forall \mathit{hbx} \in \operatorname{ooint}((x - r), (x + r)). \\ \quad \quad (f(\mathit{hbx}) \simeq \mathit{hbx}) \Rightarrow \\ \quad \quad \operatorname{has\text{-}deriv\text{-}at}(f, x, 1) \end{array}}}
\vnbentry{\texttt{has-\allowbreak{}deriv-\allowbreak{}at-\allowbreak{}real-\allowbreak{}mul}}{\vnbstmt{\begin{array}{@{}l@{}} \forall r, g, h.; \forall c \in \mathbb{R}, x, m. \\ \quad \operatorname{pos\text{-}rr}(r) \wedge \\ \quad \operatorname{restrict}(h, \operatorname{ooint}((x - r), (x + r))) \\ \quad \in\; (\operatorname{ooint}((x - r), (x + r)) \to \mathbb{R}) \wedge \\ \quad \forall \mathit{hbx} \in \operatorname{ooint}((x - r), (x + r)). \\ \quad \quad (h(\mathit{hbx}) \simeq (c \cdot g(\mathit{hbx}))) \wedge \\ \quad \operatorname{has\text{-}deriv\text{-}at}(g, x, m) \Rightarrow \\ \quad \quad \operatorname{has\text{-}deriv\text{-}at}(h, x, (c \cdot m)) \end{array}}}
\vnbentry{\texttt{has-\allowbreak{}deriv-\allowbreak{}at-\allowbreak{}difference}}{\vnbstmt{\begin{array}{@{}l@{}} \forall r, f, g, h, x, l, m. \\ \quad \operatorname{pos\text{-}rr}(r) \wedge \\ \quad \operatorname{restrict}(h, \operatorname{ooint}((x - r), (x + r))) \\ \quad \in\; (\operatorname{ooint}((x - r), (x + r)) \to \mathbb{R}) \wedge \\ \quad \operatorname{restrict}(g, \operatorname{ooint}((x - r), (x + r))) \\ \quad \in\; (\operatorname{ooint}((x - r), (x + r)) \to \mathbb{R}) \wedge \\ \quad \forall \mathit{hbx} \in \operatorname{ooint}((x - r), (x + r)). \\ \quad \quad (h(\mathit{hbx}) \simeq (f(\mathit{hbx}) - g(\mathit{hbx}))) \wedge \\ \quad \operatorname{has\text{-}deriv\text{-}at}(f, x, l) \wedge \\ \quad \operatorname{has\text{-}deriv\text{-}at}(g, x, m) \Rightarrow \\ \quad \quad \operatorname{has\text{-}deriv\text{-}at}(h, x, (l - m)) \end{array}}}

\vnbfile{theorem-library/has-deriv-at-chain.scm}{theorem-library/has-deriv-at-chain.scm}
\noindent{\small has-deriv-at-chain.scm -- THE CHAIN RULE FOR THE LOCAL DERIVATIVE, and the corollary at an affine inner map.}\par
\vnbentry{\texttt{has-\allowbreak{}deriv-\allowbreak{}at-\allowbreak{}maps-\allowbreak{}into}}{\vnbstmt{\begin{array}{@{}l@{}} \forall f, x, l, d, b. \\ \quad \operatorname{has\text{-}deriv\text{-}at}(f, x, l) \wedge \\ \quad \operatorname{pos\text{-}rr}(d) \wedge \\ \quad \operatorname{pos\text{-}rr}(b) \Rightarrow \\ \quad \quad \exists w. \\ \quad \quad \quad \operatorname{pos\text{-}rr}(w) \wedge \\ \quad \quad \quad (w \leq b) \wedge \\ \quad \quad \quad \forall y \in \operatorname{ooint}((x - w), (x + w)). \\ \quad \quad \quad \quad (f(y) \in \operatorname{ooint}((f(x) - d), (f(x) + d))) \end{array}}}
\vnbentry{\texttt{has-\allowbreak{}deriv-\allowbreak{}at-\allowbreak{}chain}}{\vnbstmt{\begin{array}{@{}l@{}} \forall r, f, g, h, x, l, m. \\ \quad \operatorname{pos\text{-}rr}(r) \wedge \\ \quad \operatorname{restrict}(h, \operatorname{ooint}((x - r), (x + r))) \\ \quad \in\; (\operatorname{ooint}((x - r), (x + r)) \to \mathbb{R}) \wedge \\ \quad \forall y \in \operatorname{ooint}((x - r), (x + r)).\; (h(y) \simeq g(f(y))) \wedge \\ \quad \operatorname{has\text{-}deriv\text{-}at}(f, x, l) \wedge \\ \quad \operatorname{has\text{-}deriv\text{-}at}(g, f(x), m) \Rightarrow \\ \quad \quad \operatorname{has\text{-}deriv\text{-}at}(h, x, (m \cdot l)) \end{array}}}
\vnbentry{\texttt{has-\allowbreak{}deriv-\allowbreak{}at-\allowbreak{}affine-\allowbreak{}chain}}{\vnbstmt{\begin{array}{@{}l@{}} \forall r, g, h.; \forall p, q, x \in \mathbb{R}, m. \\ \quad \operatorname{pos\text{-}rr}(r) \wedge \\ \quad \operatorname{restrict}(h, \operatorname{ooint}((x - r), (x + r))) \\ \quad \in\; (\operatorname{ooint}((x - r), (x + r)) \to \mathbb{R}) \wedge \\ \quad \forall y \in \operatorname{ooint}((x - r), (x + r)). \\ \quad \quad (h(y) \simeq g((p + (q \cdot y)))) \wedge \\ \quad \operatorname{has\text{-}deriv\text{-}at}(g, (p + (q \cdot x)), m) \Rightarrow \\ \quad \quad \operatorname{has\text{-}deriv\text{-}at}(h, x, (m \cdot q)) \end{array}}}

\vnbfile{theorem-library/interval-taylor.scm}{theorem-library/interval-taylor.scm}
\noindent{\small interval-taylor.scm -- TAYLOR'S FORMULA (2.18) FOR A FUNCTION ON ITS INTERVAL, AND ITS TWO COROLLARIES.}\par
\vnbentry{\texttt{gmvt-\allowbreak{}with-\allowbreak{}known-\allowbreak{}derivative}}{\vnbstmt{\begin{array}{@{}l@{}} \forall a, b \in \mathbb{R}, f, g.; \forall v \in (\operatorname{ooint}(a, b) \to \mathbb{R}). \\ \quad (a < b) \wedge \\ \quad (f \in (\operatorname{ccint}(a, b) \to \mathbb{R})) \wedge \\ \quad \operatorname{is\text{-}continuous\text{-}on}(f, \operatorname{ccint}(a, b)) \wedge \\ \quad \forall x \in \operatorname{ooint}(a, b). \\ \quad \quad \operatorname{has\text{-}deriv\text{-}at}(f, x, v(x)) \wedge \\ \quad (g \in (\operatorname{ccint}(a, b) \to \mathbb{R})) \wedge \\ \quad \operatorname{is\text{-}continuous\text{-}on}(g, \operatorname{ccint}(a, b)) \wedge \\ \quad \forall x \in \operatorname{ooint}(a, b). \exists m. \\ \quad \quad \operatorname{has\text{-}deriv\text{-}at}(g, x, m) \Rightarrow \\ \quad \quad \exists i \in \operatorname{ooint}(a, b), m. \\ \quad \quad \quad \operatorname{has\text{-}deriv\text{-}at}(g, i, m) \wedge \\ \quad \quad \quad ((v(i) \cdot (g(b) - g(a))) = (m \cdot (f(b) - f(a)))) \end{array}}}
\vnbentry{\texttt{taylor-\allowbreak{}family-\allowbreak{}order-\allowbreak{}in-\allowbreak{}nn}}{\vnbstmt{\begin{array}{@{}l@{}} \forall d, n, a, b. \\ \quad \operatorname{is\text{-}taylor\text{-}family}(d, n, a, b) \Rightarrow  \\ \quad \quad (n \in \mathbb{N}) \end{array}}}
\vnbentry{\texttt{taylor-\allowbreak{}family-\allowbreak{}order-\allowbreak{}pos}}{\vnbstmt{\begin{array}{@{}l@{}} \forall d, n, a, b. \\ \quad \operatorname{is\text{-}taylor\text{-}family}(d, n, a, b) \Rightarrow  \\ \quad \quad (1 \leq n) \end{array}}}
\vnbentry{\texttt{taylor-\allowbreak{}family-\allowbreak{}left-\allowbreak{}in-\allowbreak{}rr}}{\vnbstmt{\begin{array}{@{}l@{}} \forall d, n, a, b. \\ \quad \operatorname{is\text{-}taylor\text{-}family}(d, n, a, b) \Rightarrow  \\ \quad \quad (a \in \mathbb{R}) \end{array}}}
\vnbentry{\texttt{taylor-\allowbreak{}family-\allowbreak{}right-\allowbreak{}in-\allowbreak{}rr}}{\vnbstmt{\begin{array}{@{}l@{}} \forall d, n, a, b. \\ \quad \operatorname{is\text{-}taylor\text{-}family}(d, n, a, b) \Rightarrow  \\ \quad \quad (b \in \mathbb{R}) \end{array}}}
\vnbentry{\texttt{taylor-\allowbreak{}family-\allowbreak{}endpoints-\allowbreak{}lt}}{\vnbstmt{\begin{array}{@{}l@{}} \forall d, n, a, b. \\ \quad (\operatorname{is\text{-}taylor\text{-}family}(d, n, a, b) \Rightarrow (a < b)) \end{array}}}
\vnbentry{\texttt{taylor-\allowbreak{}family-\allowbreak{}top-\allowbreak{}in-\allowbreak{}fun}}{\vnbstmt{\begin{array}{@{}l@{}} \forall d, n, a, b. \\ \quad \operatorname{is\text{-}taylor\text{-}family}(d, n, a, b) \Rightarrow  \\ \quad \quad (d(n) \in (\operatorname{ooint}(a, b) \to \mathbb{R})) \end{array}}}
\vnbentry{\texttt{taylor-\allowbreak{}family-\allowbreak{}in-\allowbreak{}fun}}{\vnbstmt{\begin{array}{@{}l@{}} \forall d, n, a, b.; \forall k \in \mathbb{N}. \\ \quad (\operatorname{is\text{-}taylor\text{-}family}(d, n, a, b) \wedge (k < n)) \Rightarrow  \\ \quad \quad (d(k) \in (\operatorname{ccint}(a, b) \to \mathbb{R})) \end{array}}}
\vnbentry{\texttt{taylor-\allowbreak{}family-\allowbreak{}continuous}}{\vnbstmt{\begin{array}{@{}l@{}} \forall d, n, a, b.; \forall k \in \mathbb{N}. \\ \quad (\operatorname{is\text{-}taylor\text{-}family}(d, n, a, b) \wedge (k < n)) \Rightarrow  \\ \quad \quad \operatorname{is\text{-}continuous\text{-}on}(d(k), \operatorname{ccint}(a, b)) \end{array}}}
\vnbentry{\texttt{taylor-\allowbreak{}family-\allowbreak{}deriv}}{\vnbstmt{\begin{array}{@{}l@{}} \forall d, n, a, b.; \forall k \in \mathbb{N}, x \in \operatorname{ooint}(a, b). \\ \quad (\operatorname{is\text{-}taylor\text{-}family}(d, n, a, b) \wedge (k < n)) \Rightarrow  \\ \quad \quad \operatorname{has\text{-}deriv\text{-}at}(d(k), x, d((k + 1))(x)) \end{array}}}
\vnbentry{\texttt{taylor-\allowbreak{}aux-\allowbreak{}value}}{\vnbstmt{\begin{array}{@{}l@{}} \forall d, n, a, b.; \forall z \in \operatorname{ccint}(a, b). \\ \quad (\operatorname{taylor\text{-}aux}(d, n, a, b)(z) \simeq \operatorname{taylor\text{-}sum}(d, n, b, z)) \end{array}}}
\vnbentry{\texttt{taylor-\allowbreak{}clear-\allowbreak{}recip}}{\vnbstmt{\begin{array}{@{}l@{}} \forall p, q, r, s, u, w \in \mathbb{R}. \\ \quad (((s \cdot r) = 1) \wedge ((((p \cdot q) \cdot r) \cdot w) = u)) \Rightarrow  \\ \quad \quad (((u \cdot s) = ((p \cdot q) \cdot w)) \wedge (u = (((p \cdot r) \cdot q) \cdot w))) \end{array}}}
\vnbentry{\texttt{series-\allowbreak{}partial-\allowbreak{}sum-\allowbreak{}in-\allowbreak{}rr-\allowbreak{}below}}{\vnbstmt{\begin{array}{@{}l@{}} \forall n \in \mathbb{N}, f. \\ \quad \forall k \in \mathbb{N}.\; ((k < n) \Rightarrow (f(k) \in \mathbb{R})) \Rightarrow  \\ \quad \quad (\operatorname{series\text{-}partial\text{-}sum}(f, n) \in \mathbb{R}) \end{array}}}
\vnbentry{\texttt{taylor-\allowbreak{}term-\allowbreak{}value}}{\vnbstmt{\begin{array}{@{}l@{}} \forall d, b, v.; \forall k \in \mathbb{N}. \\ \quad (\operatorname{taylor\text{-}term}(d, b, v)(k) \simeq ((d(k)(v) \cdot \operatorname{power}((b - v), k)) \cdot \operatorname{recip}(\operatorname{factorial}(k)))) \end{array}}}
\vnbentry{\texttt{taylor-\allowbreak{}sum-\allowbreak{}unfold}}{\vnbstmt{\begin{array}{@{}l@{}} \forall d, n, b, v. \\ \quad (\operatorname{taylor\text{-}sum}(d, n, b, v) \simeq \operatorname{series\text{-}partial\text{-}sum}(\operatorname{taylor\text{-}term}(d, b, v), n)) \end{array}}}
\vnbentry{\texttt{taylor-\allowbreak{}sum-\allowbreak{}in-\allowbreak{}rr}}{\vnbstmt{\begin{array}{@{}l@{}} \forall d, n, a, b.; \forall j \in \mathbb{N}, z \in \operatorname{ccint}(a, b). \\ \quad (\operatorname{is\text{-}taylor\text{-}family}(d, n, a, b) \wedge (j \leq n)) \Rightarrow  \\ \quad \quad (\operatorname{taylor\text{-}sum}(d, j, b, z) \in \mathbb{R}) \end{array}}}
\vnbentry{\texttt{taylor-\allowbreak{}aux-\allowbreak{}in-\allowbreak{}fun}}{\vnbstmt{\begin{array}{@{}l@{}} \forall d, n, a, b.; \forall j \in \mathbb{N}. \\ \quad (\operatorname{is\text{-}taylor\text{-}family}(d, n, a, b) \wedge (j \leq n)) \Rightarrow  \\ \quad \quad (\operatorname{taylor\text{-}aux}(d, j, a, b) \in (\operatorname{ccint}(a, b) \to \mathbb{R})) \end{array}}}
\vnbentry{\texttt{taylor-\allowbreak{}sum-\allowbreak{}one}}{\vnbstmt{\begin{array}{@{}l@{}} \forall d.; \forall b, v \in \mathbb{R}. \\ \quad (d(0)(v) \in \mathbb{R}) \Rightarrow  \\ \quad \quad (\operatorname{taylor\text{-}sum}(d, (0 + 1), b, v) = d(0)(v)) \end{array}}}
\vnbentry{\texttt{taylor-\allowbreak{}sum-\allowbreak{}at-\allowbreak{}endpoint}}{\vnbstmt{\begin{array}{@{}l@{}} \forall m \in \mathbb{N}, d.; \forall b \in \mathbb{R}. \\ \quad \forall k \in \mathbb{N}.\; ((k \leq m) \Rightarrow (d(k)(b) \in \mathbb{R})) \Rightarrow  \\ \quad \quad (\operatorname{taylor\text{-}sum}(d, (m + 1), b, b) = d(0)(b)) \end{array}}}
\vnbentry{\texttt{taylor-\allowbreak{}family-\allowbreak{}sum-\allowbreak{}at-\allowbreak{}endpoint}}{\vnbstmt{\begin{array}{@{}l@{}} \forall d, n, a, b.; \forall m \in \mathbb{N}. \\ \quad (\operatorname{is\text{-}taylor\text{-}family}(d, n, a, b) \wedge ((m + 1) \leq n)) \Rightarrow  \\ \quad \quad (\operatorname{taylor\text{-}sum}(d, (m + 1), b, b) = d(0)(b)) \end{array}}}
\vnbentry{\texttt{cont-\allowbreak{}on-\allowbreak{}in-\allowbreak{}fun}}{\vnbstmt{\begin{array}{@{}l@{}} \forall f, m. \\ \quad \operatorname{is\text{-}continuous\text{-}on}(f, m) \Rightarrow  \\ \quad \quad (f \in (m \to \mathbb{R})) \end{array}}}
\vnbentry{\texttt{cont-\allowbreak{}on-\allowbreak{}transfer-\allowbreak{}ptwise-\allowbreak{}eq}}{\vnbstmt{\begin{array}{@{}l@{}} \forall m, f, g. \\ \quad (m \subseteq \mathbb{R}) \wedge \\ \quad (f \in (m \to \mathbb{R})) \wedge \\ \quad \operatorname{is\text{-}continuous\text{-}on}(g, m) \wedge \\ \quad \forall z \in m.\; (f(z) = g(z)) \Rightarrow \\ \quad \quad \operatorname{is\text{-}continuous\text{-}on}(f, m) \end{array}}}
\vnbentry{\texttt{cont-\allowbreak{}on-\allowbreak{}of-\allowbreak{}total}}{\vnbstmt{\begin{array}{@{}l@{}} \forall m.; \forall g \in (\mathbb{R} \to \mathbb{R}), f \in (m \to \mathbb{R}). \\ \quad (m \subseteq \mathbb{R}) \wedge \\ \quad \forall y \in \mathbb{R}. \\ \quad \quad \operatorname{is\text{-}continuous\text{-}at}(\operatorname{rr\text{-}ms}, \operatorname{rr\text{-}ms}, g, y) \wedge \\ \quad \forall z \in m.\; (f(z) = g(z)) \Rightarrow \\ \quad \quad \operatorname{is\text{-}continuous\text{-}on}(f, m) \end{array}}}
\vnbentry{\texttt{power-\allowbreak{}sub-\allowbreak{}succ-\allowbreak{}continuous}}{\vnbstmt{\begin{array}{@{}l@{}} \forall b \in \mathbb{R}, k \in \mathbb{N}, y \in \mathbb{R}. \\ \quad \operatorname{is\text{-}continuous\text{-}at}(\operatorname{rr\text{-}ms}, \operatorname{rr\text{-}ms}, (\lambda z \in \mathbb{R}.\; \operatorname{power}((b - z), (k + 1))), y) \end{array}}}
\vnbentry{\texttt{taylor-\allowbreak{}step-\allowbreak{}continuous}}{\vnbstmt{\begin{array}{@{}l@{}} \forall b, c \in \mathbb{R}, k \in \mathbb{N}, g, f \in (\mathbb{R} \to \mathbb{R}), y \in \mathbb{R}. \\ \quad \operatorname{is\text{-}continuous\text{-}at}(\operatorname{rr\text{-}ms}, \operatorname{rr\text{-}ms}, g, y) \wedge \\ \quad \operatorname{is\text{-}continuous\text{-}at}(\operatorname{rr\text{-}ms}, \operatorname{rr\text{-}ms}, f, y) \Rightarrow \\ \quad \quad \operatorname{is\text{-}continuous\text{-}at}(\operatorname{rr\text{-}ms}, \operatorname{rr\text{-}ms}, \\ \quad \quad \quad (\lambda w \in \mathbb{R}.\; (g(w) + ((f(w) \cdot \operatorname{power}((b - w), (k + 1))) \cdot c))), \\ \quad \quad \quad y) \end{array}}}
\vnbentry{\texttt{taylor-\allowbreak{}aux-\allowbreak{}step}}{\vnbstmt{\begin{array}{@{}l@{}} \forall d, n, a, b.; \forall j \in \mathbb{N}, z \in \operatorname{ccint}(a, b). \\ \quad (\operatorname{is\text{-}taylor\text{-}family}(d, n, a, b) \wedge ((j + 1) \leq n)) \Rightarrow  \\ \quad \quad \operatorname{taylor\text{-}aux}(d, (j + 1), a, b)(z) \\ \quad \quad =\; (\operatorname{taylor\text{-}aux}(d, j, a, b)(z) + ((d(j)(z) \cdot \operatorname{power}((b - z), j)) \cdot \operatorname{recip}(\operatorname{factorial}(j)))) \end{array}}}
\vnbentry{\texttt{taylor-\allowbreak{}aux-\allowbreak{}continuous-\allowbreak{}on}}{\vnbstmt{\begin{array}{@{}l@{}} \forall m \in \mathbb{N}, d, n, a, b. \\ \quad (\operatorname{is\text{-}taylor\text{-}family}(d, n, a, b) \wedge ((m + 1) \leq n)) \Rightarrow  \\ \quad \quad \operatorname{is\text{-}continuous\text{-}on}(\operatorname{taylor\text{-}aux}(d, (m + 1), a, b), \operatorname{ccint}(a, b)) \end{array}}}
\vnbentry{\texttt{ooint-\allowbreak{}interior-\allowbreak{}radius}}{\vnbstmt{\begin{array}{@{}l@{}} \forall a, b \in \mathbb{R}, x \in \operatorname{ooint}(a, b). \exists r. \\ \quad \operatorname{pos\text{-}rr}(r) \wedge \\ \quad (\operatorname{ooint}((x - r), (x + r)) \subseteq \operatorname{ooint}(a, b)) \end{array}}}
\vnbentry{\texttt{power-\allowbreak{}sub-\allowbreak{}succ-\allowbreak{}has-\allowbreak{}deriv}}{\vnbstmt{\begin{array}{@{}l@{}} \forall b \in \mathbb{R}, k \in \mathbb{N}, x \in \mathbb{R}. \\ \quad \operatorname{has\text{-}deriv\text{-}at}((\lambda z \in \mathbb{R}.\; \operatorname{power}((b - z), (k + 1))), x, \\ \quad \quad (((k + 1) \cdot \operatorname{power}((b - x), k)) \cdot -1)) \end{array}}}
\vnbentry{\texttt{taylor-\allowbreak{}step-\allowbreak{}identity}}{\vnbstmt{\begin{array}{@{}l@{}} \forall a, b, p, q, r, s \in \mathbb{R}. \\ \quad (((a \cdot p) \cdot (s \cdot r)) + (r \cdot ((b \cdot q) + (((s \cdot p) \cdot -1) \cdot a)))) \\ \quad =\; ((b \cdot q) \cdot r) \end{array}}}
\vnbentry{\texttt{taylor-\allowbreak{}aux-\allowbreak{}has-\allowbreak{}deriv}}{\vnbstmt{\begin{array}{@{}l@{}} \forall m \in \mathbb{N}, d, n, a, b.; \forall x \in \operatorname{ooint}(a, b). \\ \quad (\operatorname{is\text{-}taylor\text{-}family}(d, n, a, b) \wedge ((m + 1) \leq n)) \Rightarrow  \\ \quad \quad \operatorname{has\text{-}deriv\text{-}at}(\operatorname{taylor\text{-}aux}(d, (m + 1), a, b), x, \\ \quad \quad \quad ((d((m + 1))(x) \cdot \operatorname{power}((b - x), m)) \cdot \operatorname{recip}(\operatorname{factorial}(m)))) \end{array}}}
\vnbentry{\texttt{taylor-\allowbreak{}formula-\allowbreak{}on-\allowbreak{}interval-\allowbreak{}from-\allowbreak{}aux}}{\vnbstmt{\begin{array}{@{}l@{}} \forall a, b \in \mathbb{R}, d.; \forall m \in \mathbb{N}, g \in (\operatorname{ccint}(a, b) \to \mathbb{R}). \\ \quad \operatorname{is\text{-}taylor\text{-}family}(d, (m + 1), a, b) \wedge \\ \quad \operatorname{is\text{-}continuous\text{-}on}(\operatorname{taylor\text{-}aux}(d, (m + 1), a, b), \operatorname{ccint}(a, b)) \wedge \\ \quad \forall x \in \operatorname{ooint}(a, b). \\ \quad \quad \operatorname{has\text{-}deriv\text{-}at}(\operatorname{taylor\text{-}aux}(d, (m + 1), a, b), x, \\ \quad \quad \quad ((d((m + 1))(x) \cdot \operatorname{power}((b - x), m)) \cdot \operatorname{recip}(\operatorname{factorial}(m)))) \wedge \\ \quad (\operatorname{taylor\text{-}sum}(d, (m + 1), b, b) = d(0)(b)) \wedge \\ \quad \operatorname{is\text{-}continuous\text{-}on}(g, \operatorname{ccint}(a, b)) \wedge \\ \quad \forall x \in \operatorname{ooint}(a, b). \exists l. \\ \quad \quad \operatorname{has\text{-}deriv\text{-}at}(g, x, l) \Rightarrow \\ \quad \quad \exists i \in \operatorname{ooint}(a, b), u. \\ \quad \quad \quad \operatorname{has\text{-}deriv\text{-}at}(g, i, u) \wedge \\ \quad \quad \quad ((u \cdot \operatorname{taylor\text{-}rem}(d, (m + 1), a, b)) \cdot \operatorname{factorial}(m)) \\ \quad \quad \quad =\; ((d((m + 1))(i) \cdot \operatorname{power}((b - i), m)) \cdot (g(b) - g(a))) \wedge \\ \quad \quad \quad (u \cdot \operatorname{taylor\text{-}rem}(d, (m + 1), a, b)) \\ \quad \quad \quad =\; (((d((m + 1))(i) \cdot \operatorname{recip}(\operatorname{factorial}(m))) \cdot \operatorname{power}((b - i), m)) \cdot (g(b) - g(a))) \end{array}}}
\vnbentry{\texttt{taylor-\allowbreak{}formula-\allowbreak{}on-\allowbreak{}interval}}{\vnbstmt{\begin{array}{@{}l@{}} \forall a, b \in \mathbb{R}, d.; \forall m \in \mathbb{N}, g \in (\operatorname{ccint}(a, b) \to \mathbb{R}). \\ \quad \operatorname{is\text{-}taylor\text{-}family}(d, (m + 1), a, b) \wedge \\ \quad \operatorname{is\text{-}continuous\text{-}on}(g, \operatorname{ccint}(a, b)) \wedge \\ \quad \forall x \in \operatorname{ooint}(a, b). \exists l. \\ \quad \quad \operatorname{has\text{-}deriv\text{-}at}(g, x, l) \Rightarrow \\ \quad \quad \exists i \in \operatorname{ooint}(a, b), u. \\ \quad \quad \quad \operatorname{has\text{-}deriv\text{-}at}(g, i, u) \wedge \\ \quad \quad \quad ((u \cdot \operatorname{taylor\text{-}rem}(d, (m + 1), a, b)) \cdot \operatorname{factorial}(m)) \\ \quad \quad \quad =\; ((d((m + 1))(i) \cdot \operatorname{power}((b - i), m)) \cdot (g(b) - g(a))) \wedge \\ \quad \quad \quad (u \cdot \operatorname{taylor\text{-}rem}(d, (m + 1), a, b)) \\ \quad \quad \quad =\; (((d((m + 1))(i) \cdot \operatorname{recip}(\operatorname{factorial}(m))) \cdot \operatorname{power}((b - i), m)) \cdot (g(b) - g(a))) \end{array}}}
\vnbentry{\texttt{taylor-\allowbreak{}lagrange-\allowbreak{}on-\allowbreak{}interval}}{\vnbstmt{\begin{array}{@{}l@{}} \forall a, b \in \mathbb{R}, d.; \forall m \in \mathbb{N}. \\ \quad \operatorname{is\text{-}taylor\text{-}family}(d, (m + 1), a, b) \Rightarrow  \\ \quad \quad \exists i \in \operatorname{ooint}(a, b). \\ \quad \quad \quad \operatorname{taylor\text{-}rem}(d, (m + 1), a, b) \\ \quad \quad \quad =\; ((d((m + 1))(i) \cdot \operatorname{recip}(\operatorname{factorial}((m + 1)))) \cdot \operatorname{power}((b - a), (m + 1))) \end{array}}}

\vnbfile{theorem-library/zero-deriv-off-countable.scm}{theorem-library/zero-deriv-off-countable.scm}
\noindent{\small zero-deriv-off-countable.scm -- CONSTANCY AND THE MEAN VALUE INEQUALITY WITH A COUNTABLE EXCEPTIONAL SET (batch 20-A, 2026-09-23).}\par
\vnbentry{\texttt{pw-\allowbreak{}lt-\allowbreak{}ne}}{\vnbstmt{\begin{array}{@{}l@{}} \forall u, v \in \mathbb{R}. \\ \quad ((u < v) \Rightarrow (\neg (u = v) \wedge \neg (v = u))) \end{array}}}
\vnbentry{\texttt{nn-\allowbreak{}antitone-\allowbreak{}step-\allowbreak{}implies-\allowbreak{}le}}{\vnbstmt{\begin{array}{@{}l@{}} \forall f \in (\mathbb{N} \to \mathbb{R}), n, m \in \mathbb{N}. \\ \quad (\forall k \in \mathbb{N}.\; (f((k + 1)) \leq f(k)) \wedge (m \leq n)) \Rightarrow  \\ \quad \quad (f(n) \leq f(m)) \end{array}}}
\vnbentry{\texttt{rr-\allowbreak{}nested-\allowbreak{}intervals}}{\vnbstmt{\begin{array}{@{}l@{}} \forall o, i \in (\mathbb{N} \to \mathbb{R}). \\ \quad \forall k \in \mathbb{N}.\; (o(k) \leq o((k + 1))) \wedge \\ \quad \forall k \in \mathbb{N}.\; (i((k + 1)) \leq i(k)) \wedge \\ \quad \forall k \in \mathbb{N}.\; (o(k) \leq i(k)) \Rightarrow \\ \quad \quad \exists y \in \mathbb{R}. \forall n \in \mathbb{N}. \\ \quad \quad \quad ((o(n) \leq y) \wedge (y \leq i(n))) \end{array}}}
\vnbentry{\texttt{rr-\allowbreak{}interval-\allowbreak{}avoids-\allowbreak{}sequence}}{\vnbstmt{\begin{array}{@{}l@{}} \forall e \in (\mathbb{N} \to \mathbb{R}), p, q \in \mathbb{R}. \\ \quad (p < q) \Rightarrow  \\ \quad \quad \exists y \in \mathbb{R}. \\ \quad \quad \quad ((p < y) \wedge ((y < q) \wedge \forall n \in \mathbb{N}.\; \neg (e(n) = y))) \end{array}}}
\vnbentry{\texttt{countable-\allowbreak{}image-\allowbreak{}avoids-\allowbreak{}interval}}{\vnbstmt{\begin{array}{@{}l@{}} \forall c, m, h.; \forall p, q \in \mathbb{R}. \\ \quad \operatorname{is\text{-}countable}(c) \wedge \\ \quad (c \subseteq m) \wedge \\ \quad (h \in (m \to \mathbb{R})) \wedge \\ \quad (p < q) \Rightarrow \\ \quad \quad \exists y \in \mathbb{R}. \\ \quad \quad \quad ((p < y) \wedge ((y < q) \wedge \forall x \in c.\; \neg (h(x) = y))) \end{array}}}
\vnbentry{\texttt{neg-\allowbreak{}deriv-\allowbreak{}off-\allowbreak{}countable-\allowbreak{}nonincreasing}}{\vnbstmt{\begin{array}{@{}l@{}} \forall h \in (\mathbb{R} \to \mathbb{R}), u, v \in \mathbb{R}, c. \\ \quad (u < v) \wedge \\ \quad \forall x \in \operatorname{ccint}(u, v). \\ \quad \quad \operatorname{is\text{-}continuous\text{-}at}(\operatorname{rr\text{-}ms}, \operatorname{rr\text{-}ms}, h, x) \wedge \\ \quad \operatorname{is\text{-}countable}(c) \wedge \\ \quad (c \subseteq \mathbb{R}) \wedge \\ \quad \forall t \in \operatorname{ooint}(u, v). \\ \quad \quad \neg (t \in c) \Rightarrow  \\ \quad \quad \quad \exists l.\; (\operatorname{is\text{-}diff\text{-}at}(h, t, l) \wedge (l < 0)) \Rightarrow \\ \quad \quad (h(v) \leq h(u)) \end{array}}}
\vnbentry{\texttt{mvi-\allowbreak{}off-\allowbreak{}countable}}{\vnbstmt{\begin{array}{@{}l@{}} \forall h \in (\mathbb{R} \to \mathbb{R}), m, u, v \in \mathbb{R}, c. \\ \quad (u < v) \wedge \\ \quad \forall x \in \operatorname{ccint}(u, v). \\ \quad \quad \operatorname{is\text{-}continuous\text{-}at}(\operatorname{rr\text{-}ms}, \operatorname{rr\text{-}ms}, h, x) \wedge \\ \quad \operatorname{is\text{-}countable}(c) \wedge \\ \quad (c \subseteq \mathbb{R}) \wedge \\ \quad \forall t \in \operatorname{ooint}(u, v). \\ \quad \quad \neg (t \in c) \Rightarrow  \\ \quad \quad \quad \exists l. \\ \quad \quad \quad \quad (\operatorname{is\text{-}diff\text{-}at}(h, t, l) \wedge (l \leq m)) \Rightarrow \\ \quad \quad ((h(v) - h(u)) \leq (m \cdot (v - u))) \end{array}}}
\vnbentry{\texttt{mvi-\allowbreak{}abs-\allowbreak{}off-\allowbreak{}countable}}{\vnbstmt{\begin{array}{@{}l@{}} \forall h \in (\mathbb{R} \to \mathbb{R}), m, u, v \in \mathbb{R}, c. \\ \quad (u < v) \wedge \\ \quad \forall x \in \operatorname{ccint}(u, v). \\ \quad \quad \operatorname{is\text{-}continuous\text{-}at}(\operatorname{rr\text{-}ms}, \operatorname{rr\text{-}ms}, h, x) \wedge \\ \quad \operatorname{is\text{-}countable}(c) \wedge \\ \quad (c \subseteq \mathbb{R}) \wedge \\ \quad \forall t \in \operatorname{ooint}(u, v). \\ \quad \quad \neg (t \in c) \Rightarrow  \\ \quad \quad \quad \exists l. \\ \quad \quad \quad \quad (\operatorname{is\text{-}diff\text{-}at}(h, t, l) \wedge (\lvert l \rvert \leq m)) \Rightarrow \\ \quad \quad (\lvert (h(v) - h(u)) \rvert \leq (m \cdot (v - u))) \end{array}}}
\vnbentry{\texttt{zero-\allowbreak{}deriv-\allowbreak{}off-\allowbreak{}countable}}{\vnbstmt{\begin{array}{@{}l@{}} \forall s.; \forall h \in (\mathbb{R} \to \mathbb{R}), u, v \in \mathbb{R}. \\ \quad \operatorname{is\text{-}countable}(s) \wedge \\ \quad (s \subseteq \mathbb{R}) \wedge \\ \quad \forall x \in \mathbb{R}. \\ \quad \quad \operatorname{is\text{-}continuous\text{-}at}(\operatorname{rr\text{-}ms}, \operatorname{rr\text{-}ms}, h, x) \wedge \\ \quad (u < v) \wedge \\ \quad \forall t \in \operatorname{ooint}(u, v). \\ \quad \quad (\neg (t \in s) \Rightarrow \operatorname{is\text{-}diff\text{-}at}(h, t, 0)) \Rightarrow \\ \quad \quad (h(u) = h(v)) \end{array}}}
\vnbentry{\texttt{mvi-\allowbreak{}abs-\allowbreak{}off-\allowbreak{}countable-\allowbreak{}on-\allowbreak{}interval}}{\vnbstmt{\begin{array}{@{}l@{}} \forall a, b, m \in \mathbb{R}, f \in (\operatorname{ccint}(a, b) \to \mathbb{R}), d. \\ \quad (a < b) \wedge \\ \quad \operatorname{is\text{-}continuous\text{-}on}(f, \operatorname{ccint}(a, b)) \wedge \\ \quad \operatorname{is\text{-}countable}(d) \wedge \\ \quad (d \subseteq \operatorname{ccint}(a, b)) \wedge \\ \quad \forall t \in \operatorname{ooint}(a, b). \\ \quad \quad \neg (t \in d) \Rightarrow  \\ \quad \quad \quad \exists l. \\ \quad \quad \quad \quad (\operatorname{has\text{-}deriv\text{-}at}(f, t, l) \wedge (\lvert l \rvert \leq m)) \Rightarrow \\ \quad \quad (\lvert (f(b) - f(a)) \rvert \leq (m \cdot (b - a))) \end{array}}}
\vnbentry{\texttt{zero-\allowbreak{}deriv-\allowbreak{}off-\allowbreak{}countable-\allowbreak{}constant}}{\vnbstmt{\begin{array}{@{}l@{}} \forall a, b \in \mathbb{R}, g \in (\operatorname{ccint}(a, b) \to \mathbb{R}), d.; \forall s, w \in \operatorname{ccint}(a, b). \\ \quad (a < b) \wedge \\ \quad \operatorname{is\text{-}continuous\text{-}on}(g, \operatorname{ccint}(a, b)) \wedge \\ \quad \operatorname{is\text{-}countable}(d) \wedge \\ \quad (d \subseteq \operatorname{ccint}(a, b)) \wedge \\ \quad \forall t \in \operatorname{ooint}(a, b). \\ \quad \quad (\neg (t \in d) \Rightarrow \operatorname{has\text{-}deriv\text{-}at}(g, t, 0)) \Rightarrow \\ \quad \quad (g(s) = g(w)) \end{array}}}

\vnbfile{theorem-library/pw-antiderivative-laws.scm}{theorem-library/pw-antiderivative-laws.scm}
\noindent{\small pw-antiderivative-laws.scm -- the laws of IS-PW-ANTIDERIVATIVE and PW-INT (structure-library/path-integral.scm), with the exceptional set a FINITE SET.}\par
\vnbentry{\texttt{pw-\allowbreak{}antiderivative-\allowbreak{}value-\allowbreak{}in-\allowbreak{}rr}}{\vnbstmt{\begin{array}{@{}l@{}} \forall f, i, a, b.; \forall y \in \operatorname{ccint}(a, b). \\ \quad \operatorname{is\text{-}primitive}(f, i, a, b) \Rightarrow  \\ \quad \quad (f(y) \in \mathbb{R}) \end{array}}}
\vnbentry{\texttt{pw-\allowbreak{}not-\allowbreak{}in-\allowbreak{}insert}}{\vnbstmt{\begin{array}{@{}l@{}} \forall y \in \mathbf{Set}, w, t. \\ \quad ((\neg (t \in w) \wedge \neg (t = y)) \Rightarrow \neg (t \in (w \cup \{y\}))) \end{array}}}
\vnbentry{\texttt{ooint-\allowbreak{}inner-\allowbreak{}radius}}{\vnbstmt{\begin{array}{@{}l@{}} \forall u, v \in \mathbb{R}, t \in \operatorname{ooint}(u, v). \exists r. \\ \quad \operatorname{pos\text{-}rr}(r) \wedge \\ \quad (\operatorname{ooint}((t - r), (t + r)) \subseteq \operatorname{ccint}(u, v)) \end{array}}}
\vnbentry{\texttt{pw-\allowbreak{}zero-\allowbreak{}deriv-\allowbreak{}total}}{\vnbstmt{\begin{array}{@{}l@{}} \forall h \in (\mathbb{R} \to \mathbb{R}), u, v \in \mathbb{R}. \\ \quad \forall x \in \mathbb{R}. \\ \quad \quad \operatorname{is\text{-}continuous\text{-}at}(\operatorname{rr\text{-}ms}, \operatorname{rr\text{-}ms}, h, x) \wedge \\ \quad (u < v) \wedge \\ \quad \forall t \in \operatorname{ooint}(u, v). \\ \quad \quad \operatorname{is\text{-}diff\text{-}at}(h, t, 0) \Rightarrow \\ \quad \quad (h(u) = h(v)) \end{array}}}
\vnbentry{\texttt{deriv-\allowbreak{}zero-\allowbreak{}constant-\allowbreak{}on-\allowbreak{}interval}}{\vnbstmt{\begin{array}{@{}l@{}} \forall a, b \in \mathbb{R}, h \in (\operatorname{ccint}(a, b) \to \mathbb{R}), u, v \in \operatorname{ccint}(a, b). \\ \quad (a < b) \wedge \\ \quad \operatorname{is\text{-}continuous\text{-}on}(h, \operatorname{ccint}(a, b)) \wedge \\ \quad \forall t \in \operatorname{ooint}(a, b). \\ \quad \quad \operatorname{has\text{-}deriv\text{-}at}(h, t, 0) \Rightarrow \\ \quad \quad (h(u) = h(v)) \end{array}}}
\vnbentry{\texttt{pw-\allowbreak{}zero-\allowbreak{}deriv-\allowbreak{}off-\allowbreak{}finite-\allowbreak{}set}}{\vnbstmt{\begin{array}{@{}l@{}} \forall s \in \mathbf{Set}, h \in (\mathbb{R} \to \mathbb{R}), u, v \in \mathbb{R}. \\ \quad (|s| \in \mathbb{N}) \wedge \\ \quad \forall x \in \mathbb{R}. \\ \quad \quad \operatorname{is\text{-}continuous\text{-}at}(\operatorname{rr\text{-}ms}, \operatorname{rr\text{-}ms}, h, x) \wedge \\ \quad (u < v) \wedge \\ \quad \forall t \in \operatorname{ooint}(u, v). \\ \quad \quad (\neg (t \in s) \Rightarrow \operatorname{is\text{-}diff\text{-}at}(h, t, 0)) \Rightarrow \\ \quad \quad (h(u) = h(v)) \end{array}}}
\vnbentry{\texttt{pw-\allowbreak{}antiderivative-\allowbreak{}endpoint-\allowbreak{}difference}}{\vnbstmt{\begin{array}{@{}l@{}} \forall f, w, i, a, b. \\ \quad \operatorname{is\text{-}primitive}(f, i, a, b) \wedge \\ \quad \operatorname{is\text{-}primitive}(w, i, a, b) \Rightarrow \\ \quad \quad ((f(b) - f(a)) = (w(b) - w(a))) \end{array}}}
\vnbentry{\texttt{pw-\allowbreak{}int-\allowbreak{}value}}{\vnbstmt{\begin{array}{@{}l@{}} \forall f, i, a, b. \\ \quad \operatorname{is\text{-}primitive}(f, i, a, b) \Rightarrow  \\ \quad \quad (\int_{a}^{b} i = (f(b) - f(a))) \end{array}}}
\vnbentry{\texttt{pw-\allowbreak{}int-\allowbreak{}in-\allowbreak{}rr}}{\vnbstmt{\begin{array}{@{}l@{}} \forall f, i, a, b. \\ \quad \operatorname{is\text{-}primitive}(f, i, a, b) \Rightarrow  \\ \quad \quad (\int_{a}^{b} i \in \mathbb{R}) \end{array}}}
\vnbentry{\texttt{pw-\allowbreak{}antiderivative-\allowbreak{}integrand-\allowbreak{}congruence}}{\vnbstmt{\begin{array}{@{}l@{}} \forall f, i, \mathit{ppsi}, a, b. \\ \quad \operatorname{is\text{-}primitive}(f, i, a, b) \wedge \\ \quad (\mathit{ppsi} \in (\operatorname{ccint}(a, b) \to \mathbb{R})) \wedge \\ \quad \forall y \in \operatorname{ccint}(a, b).\; (i(y) \simeq \operatorname{ppsi}(y)) \Rightarrow \\ \quad \quad \operatorname{is\text{-}primitive}(f, \mathit{ppsi}, a, b) \end{array}}}
\vnbentry{\texttt{pw-\allowbreak{}int-\allowbreak{}congruence}}{\vnbstmt{\begin{array}{@{}l@{}} \forall f, i, \mathit{ppsi}, a, b. \\ \quad \operatorname{is\text{-}primitive}(f, i, a, b) \wedge \\ \quad (\mathit{ppsi} \in (\operatorname{ccint}(a, b) \to \mathbb{R})) \wedge \\ \quad \forall y \in \operatorname{ccint}(a, b).\; (i(y) \simeq \operatorname{ppsi}(y)) \Rightarrow \\ \quad \quad (\int_{a}^{b} i = \int_{a}^{b} \mathit{ppsi}) \end{array}}}
\vnbentry{\texttt{pw-\allowbreak{}antiderivative-\allowbreak{}one-\allowbreak{}piece}}{\vnbstmt{\begin{array}{@{}l@{}} \forall a, b \in \mathbb{R}, f, i. \\ \quad (a < b) \wedge \\ \quad (f \in (\operatorname{ccint}(a, b) \to \mathbb{R})) \wedge \\ \quad (i \in (\operatorname{ccint}(a, b) \to \mathbb{R})) \wedge \\ \quad \operatorname{is\text{-}continuous\text{-}on}(f, \operatorname{ccint}(a, b)) \wedge \\ \quad \forall t \in \operatorname{ooint}(a, b). \\ \quad \quad \operatorname{has\text{-}deriv\text{-}at}(f, t, i(t)) \Rightarrow \\ \quad \quad \operatorname{is\text{-}primitive}(f, i, a, b) \end{array}}}
\vnbentry{\texttt{pw-\allowbreak{}restrict-\allowbreak{}continuous-\allowbreak{}on}}{\vnbstmt{\begin{array}{@{}l@{}} \forall a, b \in \mathbb{R}, f, \mathit{pwf}. \\ \quad (a < b) \wedge \\ \quad (f \in (\mathbb{R} \to \mathbb{R})) \wedge \\ \quad \forall y \in \mathbb{R}. \\ \quad \quad \operatorname{is\text{-}continuous\text{-}at}(\operatorname{rr\text{-}ms}, \operatorname{rr\text{-}ms}, f, y) \wedge \\ \quad (\mathit{pwf} \in (\operatorname{ccint}(a, b) \to \mathbb{R})) \wedge \\ \quad \forall y \in \operatorname{ccint}(a, b).\; (\operatorname{pwf}(y) \simeq f(y)) \Rightarrow \\ \quad \quad \operatorname{is\text{-}continuous\text{-}on}(\mathit{pwf}, \operatorname{ccint}(a, b)) \end{array}}}
\vnbentry{\texttt{pw-\allowbreak{}affine-\allowbreak{}antiderivative}}{\vnbstmt{\begin{array}{@{}l@{}} \forall a, b, c, m \in \mathbb{R}, f, i. \\ \quad (a < b) \wedge \\ \quad (f \in (\operatorname{ccint}(a, b) \to \mathbb{R})) \wedge \\ \quad (i \in (\operatorname{ccint}(a, b) \to \mathbb{R})) \wedge \\ \quad \forall y \in \operatorname{ccint}(a, b).\; (f(y) \simeq (c + (m \cdot y))) \wedge \\ \quad \forall y \in \operatorname{ccint}(a, b).\; (i(y) \simeq m) \Rightarrow \\ \quad \quad \operatorname{is\text{-}primitive}(f, i, a, b) \end{array}}}
\vnbentry{\texttt{pw-\allowbreak{}int-\allowbreak{}const}}{\vnbstmt{\begin{array}{@{}l@{}} \forall a, b, c \in \mathbb{R}, i \in (\operatorname{ccint}(a, b) \to \mathbb{R}). \\ \quad ((a < b) \wedge \forall y \in \operatorname{ccint}(a, b).\; (i(y) \simeq c)) \Rightarrow  \\ \quad \quad (\int_{a}^{b} i = (c \cdot (b - a))) \end{array}}}
\vnbentry{\texttt{pw-\allowbreak{}antiderivative-\allowbreak{}sum}}{\vnbstmt{\begin{array}{@{}l@{}} \forall a, b \in \mathbb{R}, f, w, i, \mathit{ppsi}.; \forall h, \mathit{pchi} \in (\operatorname{ccint}(a, b) \to \mathbb{R}). \\ \quad \operatorname{is\text{-}primitive}(f, i, a, b) \wedge \\ \quad \operatorname{is\text{-}primitive}(w, \mathit{ppsi}, a, b) \wedge \\ \quad \forall y \in \operatorname{ccint}(a, b).\; (h(y) \simeq (f(y) + w(y))) \wedge \\ \quad \forall y \in \operatorname{ccint}(a, b). \\ \quad \quad (\operatorname{pchi}(y) \simeq (i(y) + \operatorname{ppsi}(y))) \Rightarrow \\ \quad \quad \operatorname{is\text{-}primitive}(h, \mathit{pchi}, a, b) \wedge \\ \quad \quad (\int_{a}^{b} \mathit{pchi} = (\int_{a}^{b} i + \int_{a}^{b} \mathit{ppsi})) \end{array}}}

\vnbfile{theorem-library/cc-int-laws.scm}{theorem-library/cc-int-laws.scm}
\noindent{\small cc-int-laws.scm -- the read-offs of CC-INT, TRACE and LINE-INT (structure-library/path-integral.scm), sections 3.3 and 3.4 of docs/paths-and-line-integrals-2026-09-21.md.}\par
\vnbentry{\texttt{cc-\allowbreak{}int-\allowbreak{}unfold}}{\vnbstmt{\begin{array}{@{}l@{}} \forall i, a, b. \\ \quad (\int_{a}^{b} i \simeq (\int_{a}^{b} (\lambda t \in \operatorname{ccint}(a, b).\; \operatorname{real\text{-}part}(i(t))) + (\int_{a}^{b} (\lambda t \in \operatorname{ccint}(a, b).\; \operatorname{imag\text{-}part}(i(t))) \cdot +i))) \end{array}}}
\vnbentry{\texttt{line-\allowbreak{}int-\allowbreak{}unfold}}{\vnbstmt{\begin{array}{@{}l@{}} \forall f, m, c, a, b. \\ \quad (\operatorname{line\text{-}int}(f, m, c, a, b) \simeq \int_{a}^{b} (\lambda t \in \operatorname{ccint}(a, b).\; (f(m(t)) \cdot c(t)))) \end{array}}}
\vnbentry{\texttt{trace-\allowbreak{}unfold}}{\vnbstmt{\begin{array}{@{}l@{}} \forall m, a, b. \\ \quad (\operatorname{trace}(m, a, b) \simeq \operatorname{image}(m, \operatorname{ccint}(a, b))) \end{array}}}
\vnbentry{\texttt{trace-\allowbreak{}value-\allowbreak{}in}}{\vnbstmt{\begin{array}{@{}l@{}} \forall m, a, b.; \forall t \in \operatorname{ccint}(a, b). \\ \quad (m \in (\operatorname{ccint}(a, b) \to \mathbb{C})) \Rightarrow  \\ \quad \quad (m(t) \in \operatorname{trace}(m, a, b)) \end{array}}}
\vnbentry{\texttt{trace-\allowbreak{}subset-\allowbreak{}cc}}{\vnbstmt{\begin{array}{@{}l@{}} \forall m, a, b. \\ \quad (m \in (\operatorname{ccint}(a, b) \to \mathbb{C})) \Rightarrow  \\ \quad \quad (\operatorname{trace}(m, a, b) \subseteq \mathbb{C}) \end{array}}}
\vnbentry{\texttt{trace-\allowbreak{}is-\allowbreak{}set}}{\vnbstmt{\begin{array}{@{}l@{}} \forall m, a, b. \\ \quad ((a \in \mathbb{R}) \wedge (b \in \mathbb{R})) \Rightarrow  \\ \quad \quad (\operatorname{trace}(m, a, b) \in \mathbf{Set}) \end{array}}}
\vnbentry{\texttt{line-\allowbreak{}int-\allowbreak{}integrand-\allowbreak{}in-\allowbreak{}fun}}{\vnbstmt{\begin{array}{@{}l@{}} \forall a, b, m, c, f, d. \\ \quad (a \in \mathbb{R}) \wedge \\ \quad (b \in \mathbb{R}) \wedge \\ \quad (m \in (\operatorname{ccint}(a, b) \to \mathbb{C})) \wedge \\ \quad (c \in (\operatorname{ccint}(a, b) \to \mathbb{C})) \wedge \\ \quad (f \in (d \to \mathbb{C})) \wedge \\ \quad (\operatorname{trace}(m, a, b) \subseteq d) \Rightarrow \\ \quad \quad (\lambda t \in \operatorname{ccint}(a, b).\; (f(m(t)) \cdot c(t))) \\ \quad \quad \in\; (\operatorname{ccint}(a, b) \to \mathbb{C}) \end{array}}}

\vnbfile{theorem-library/road-laws.scm}{theorem-library/road-laws.scm}
\noindent{\small road-laws.scm -- the laws of IS-PW-CONTINUOUS-ON, IS-PATH and IS-ROAD (structure-library/path-integral.scm), section 3.4 of docs/paths-and-line-integrals-2026-09-21.md and Dieudonne 9.6.}\par
\vnbentry{\texttt{pw-\allowbreak{}continuous-\allowbreak{}endpoints}}{\vnbstmt{\begin{array}{@{}l@{}} \forall i, a, b. \\ \quad \operatorname{is\text{-}pw\text{-}continuous\text{-}on}(i, a, b) \Rightarrow  \\ \quad \quad ((a \in \mathbb{R}) \wedge ((b \in \mathbb{R}) \wedge (a < b))) \end{array}}}
\vnbentry{\texttt{pw-\allowbreak{}continuous-\allowbreak{}in-\allowbreak{}fun}}{\vnbstmt{\begin{array}{@{}l@{}} \forall i, a, b. \\ \quad \operatorname{is\text{-}pw\text{-}continuous\text{-}on}(i, a, b) \Rightarrow  \\ \quad \quad (i \in (\operatorname{ccint}(a, b) \to \mathbb{R})) \end{array}}}
\vnbentry{\texttt{is-\allowbreak{}path-\allowbreak{}endpoints}}{\vnbstmt{\begin{array}{@{}l@{}} \forall m, a, b. \\ \quad \operatorname{is\text{-}path}(m, a, b) \Rightarrow  \\ \quad \quad ((a \in \mathbb{R}) \wedge ((b \in \mathbb{R}) \wedge (a < b))) \end{array}}}
\vnbentry{\texttt{is-\allowbreak{}path-\allowbreak{}in-\allowbreak{}fun}}{\vnbstmt{\begin{array}{@{}l@{}} \forall m, a, b. \\ \quad \operatorname{is\text{-}path}(m, a, b) \Rightarrow  \\ \quad \quad (m \in (\operatorname{ccint}(a, b) \to \mathbb{C})) \end{array}}}
\vnbentry{\texttt{is-\allowbreak{}path-\allowbreak{}continuous}}{\vnbstmt{\begin{array}{@{}l@{}} \forall m, a, b.; \forall x \in \operatorname{ccint}(a, b). \\ \quad \operatorname{is\text{-}path}(m, a, b) \Rightarrow  \\ \quad \quad \operatorname{is\text{-}continuous\text{-}at}(\operatorname{subspace\text{-}ms}(\operatorname{rr\text{-}ms}, \operatorname{ccint}(a, b)), \operatorname{cc\text{-}ms}, m, x) \end{array}}}
\vnbentry{\texttt{is-\allowbreak{}road-\allowbreak{}is-\allowbreak{}path}}{\vnbstmt{\begin{array}{@{}l@{}} \forall m, c, a, b. \\ \quad \operatorname{is\text{-}road}(m, c, a, b) \Rightarrow  \\ \quad \quad \operatorname{is\text{-}path}(m, a, b) \end{array}}}
\vnbentry{\texttt{is-\allowbreak{}road-\allowbreak{}dgam-\allowbreak{}in-\allowbreak{}fun}}{\vnbstmt{\begin{array}{@{}l@{}} \forall m, c, a, b. \\ \quad \operatorname{is\text{-}road}(m, c, a, b) \Rightarrow  \\ \quad \quad (c \in (\operatorname{ccint}(a, b) \to \mathbb{C})) \end{array}}}
\vnbentry{\texttt{is-\allowbreak{}road-\allowbreak{}re-\allowbreak{}antiderivative}}{\vnbstmt{\begin{array}{@{}l@{}} \forall m, c, a, b. \\ \quad \operatorname{is\text{-}road}(m, c, a, b) \Rightarrow  \\ \quad \quad \operatorname{is\text{-}primitive}((\lambda t \in \operatorname{ccint}(a, b).\; \operatorname{real\text{-}part}(m(t))), \\ \quad \quad \quad (\lambda t \in \operatorname{ccint}(a, b).\; \operatorname{real\text{-}part}(c(t))), a, b) \end{array}}}
\vnbentry{\texttt{is-\allowbreak{}road-\allowbreak{}im-\allowbreak{}antiderivative}}{\vnbstmt{\begin{array}{@{}l@{}} \forall m, c, a, b. \\ \quad \operatorname{is\text{-}road}(m, c, a, b) \Rightarrow  \\ \quad \quad \operatorname{is\text{-}primitive}((\lambda t \in \operatorname{ccint}(a, b).\; \operatorname{imag\text{-}part}(m(t))), \\ \quad \quad \quad (\lambda t \in \operatorname{ccint}(a, b).\; \operatorname{imag\text{-}part}(c(t))), a, b) \end{array}}}
\vnbentry{\texttt{is-\allowbreak{}road-\allowbreak{}re-\allowbreak{}regulated}}{\vnbstmt{\begin{array}{@{}l@{}} \forall m, c, a, b. \\ \quad \operatorname{is\text{-}road}(m, c, a, b) \Rightarrow  \\ \quad \quad \operatorname{is\text{-}regulated\text{-}on}((\lambda t \in \operatorname{ccint}(a, b).\; \operatorname{real\text{-}part}(c(t))), a, b) \end{array}}}
\vnbentry{\texttt{is-\allowbreak{}road-\allowbreak{}im-\allowbreak{}regulated}}{\vnbstmt{\begin{array}{@{}l@{}} \forall m, c, a, b. \\ \quad \operatorname{is\text{-}road}(m, c, a, b) \Rightarrow  \\ \quad \quad \operatorname{is\text{-}regulated\text{-}on}((\lambda t \in \operatorname{ccint}(a, b).\; \operatorname{imag\text{-}part}(c(t))), a, b) \end{array}}}
\vnbentry{\texttt{pw-\allowbreak{}continuous-\allowbreak{}one-\allowbreak{}piece}}{\vnbstmt{\begin{array}{@{}l@{}} \forall a, b \in \mathbb{R}, i \in (\operatorname{ccint}(a, b) \to \mathbb{R}). \\ \quad ((a < b) \wedge \operatorname{is\text{-}continuous\text{-}on}(i, \operatorname{ccint}(a, b))) \Rightarrow  \\ \quad \quad \operatorname{is\text{-}pw\text{-}continuous\text{-}on}(i, a, b) \end{array}}}
\vnbentry{\texttt{is-\allowbreak{}path-\allowbreak{}of-\allowbreak{}coords}}{\vnbstmt{\begin{array}{@{}l@{}} \forall a, b \in \mathbb{R}, m, r, i. \\ \quad (a < b) \wedge \\ \quad (m \in (\operatorname{ccint}(a, b) \to \mathbb{C})) \wedge \\ \quad (r \in (\operatorname{ccint}(a, b) \to \mathbb{R})) \wedge \\ \quad (i \in (\operatorname{ccint}(a, b) \to \mathbb{R})) \wedge \\ \quad \forall y \in \operatorname{ccint}(a, b). \\ \quad \quad (\operatorname{real\text{-}part}(m(y)) = r(y)) \wedge \\ \quad \forall y \in \operatorname{ccint}(a, b). \\ \quad \quad (\operatorname{imag\text{-}part}(m(y)) = i(y)) \wedge \\ \quad \operatorname{is\text{-}continuous\text{-}on}(r, \operatorname{ccint}(a, b)) \wedge \\ \quad \operatorname{is\text{-}continuous\text{-}on}(i, \operatorname{ccint}(a, b)) \Rightarrow \\ \quad \quad \operatorname{is\text{-}path}(m, a, b) \end{array}}}
\vnbentry{\texttt{segment-\allowbreak{}is-\allowbreak{}road}}{\vnbstmt{\begin{array}{@{}l@{}} \forall z_{0}, w. \\ \quad ((z_{0} \in \mathbb{C}) \wedge (w \in \mathbb{C})) \Rightarrow  \\ \quad \quad \operatorname{is\text{-}road}((\lambda x \in \operatorname{ccint}(0, 1).\; (z_{0} + (x \cdot w))), \\ \quad \quad \quad (\lambda x \in \operatorname{ccint}(0, 1).\; w), 0, 1) \end{array}}}
\vnbentry{\texttt{line-\allowbreak{}int-\allowbreak{}of-\allowbreak{}one}}{\vnbstmt{\begin{array}{@{}l@{}} \forall m, c, a, b, f, d. \\ \quad \operatorname{is\text{-}road}(m, c, a, b) \wedge \\ \quad (f \in (d \to \mathbb{C})) \wedge \\ \quad (\operatorname{trace}(m, a, b) \subseteq d) \wedge \\ \quad \forall y \in d.\; (f(y) = 1) \Rightarrow \\ \quad \quad (\operatorname{line\text{-}int}(f, m, c, a, b) = (m(b) - m(a))) \end{array}}}

\vnbfile{theorem-library/pw-int-laws-2.scm}{theorem-library/pw-int-laws-2.scm}
\noindent{\small pw-int-laws-2.scm -- the LINEARITY and the ADDITIVITY of the integral on [a, b], continuing theorem-library/pw-antiderivative-laws.scm.}\par
\vnbentry{\texttt{pw-\allowbreak{}antiderivative-\allowbreak{}real-\allowbreak{}mul}}{\vnbstmt{\begin{array}{@{}l@{}} \forall a, b \in \mathbb{R}, f, i.; \forall c \in \mathbb{R}, h, \mathit{pchi} \in (\operatorname{ccint}(a, b) \to \mathbb{R}). \\ \quad \operatorname{is\text{-}primitive}(f, i, a, b) \wedge \\ \quad \forall y \in \operatorname{ccint}(a, b).\; (h(y) \simeq (c \cdot f(y))) \wedge \\ \quad \forall y \in \operatorname{ccint}(a, b). \\ \quad \quad (\operatorname{pchi}(y) \simeq (c \cdot i(y))) \Rightarrow \\ \quad \quad \operatorname{is\text{-}primitive}(h, \mathit{pchi}, a, b) \wedge \\ \quad \quad (\int_{a}^{b} \mathit{pchi} = (c \cdot \int_{a}^{b} i)) \end{array}}}
\vnbentry{\texttt{pw-\allowbreak{}antiderivative-\allowbreak{}differ-\allowbreak{}by-\allowbreak{}constant}}{\vnbstmt{\begin{array}{@{}l@{}} \forall f, w, i, a, b.; \forall s \in \operatorname{ccint}(a, b). \\ \quad \operatorname{is\text{-}primitive}(f, i, a, b) \wedge \\ \quad \operatorname{is\text{-}primitive}(w, i, a, b) \Rightarrow \\ \quad \quad (f(s) = (w(s) + (f(a) - w(a)))) \end{array}}}
\vnbentry{\texttt{pw-\allowbreak{}antiderivative-\allowbreak{}differ-\allowbreak{}by-\allowbreak{}some-\allowbreak{}constant}}{\vnbstmt{\begin{array}{@{}l@{}} \forall f, w, i, a, b. \\ \quad \operatorname{is\text{-}primitive}(f, i, a, b) \wedge \\ \quad \operatorname{is\text{-}primitive}(w, i, a, b) \Rightarrow \\ \quad \quad \exists c \in \mathbb{R}. \forall s \in \operatorname{ccint}(a, b). \\ \quad \quad \quad (f(s) = (w(s) + c)) \end{array}}}
\vnbentry{\texttt{ccint-\allowbreak{}subset-\allowbreak{}ccint}}{\vnbstmt{\begin{array}{@{}l@{}} \forall a, b, v, \mathit{pdv} \in \mathbb{R}. \\ \quad ((a \leq v) \wedge (\mathit{pdv} \leq b)) \Rightarrow  \\ \quad \quad (\operatorname{ccint}(v, \mathit{pdv}) \subseteq \operatorname{ccint}(a, b)) \end{array}}}
\vnbentry{\texttt{pw-\allowbreak{}antiderivative-\allowbreak{}restrict}}{\vnbstmt{\begin{array}{@{}l@{}} \forall a, b, v, \mathit{pdv} \in \mathbb{R}, f, i. \\ \quad \operatorname{is\text{-}primitive}(f, i, a, b) \wedge \\ \quad (a \leq v) \wedge \\ \quad (v < \mathit{pdv}) \wedge \\ \quad (\mathit{pdv} \leq b) \Rightarrow \\ \quad \quad \operatorname{is\text{-}primitive}(\operatorname{restrict}(f, \operatorname{ccint}(v, \mathit{pdv})), \operatorname{restrict}(i, \operatorname{ccint}(v, \mathit{pdv})), v, \mathit{pdv}) \end{array}}}
\vnbentry{\texttt{pw-\allowbreak{}int-\allowbreak{}adjacent}}{\vnbstmt{\begin{array}{@{}l@{}} \forall a, b \in \mathbb{R}, v \in \operatorname{ooint}(a, b), f, i. \\ \quad \operatorname{is\text{-}primitive}(f, i, a, b) \Rightarrow  \\ \quad \quad \int_{a}^{b} i \\ \quad \quad =\; (\int_{a}^{v} \operatorname{restrict}(i, \operatorname{ccint}(a, v)) + \int_{v}^{b} \operatorname{restrict}(i, \operatorname{ccint}(v, b))) \end{array}}}
\vnbentry{\texttt{cc-\allowbreak{}int-\allowbreak{}value}}{\vnbstmt{\begin{array}{@{}l@{}} \forall a, b, i, f, w. \\ \quad (a \in \mathbb{R}) \wedge \\ \quad (b \in \mathbb{R}) \wedge \\ \quad \operatorname{is\text{-}primitive}(f, (\lambda t \in \operatorname{ccint}(a, b).\; \operatorname{real\text{-}part}(i(t))), a, b) \wedge \\ \quad \operatorname{is\text{-}primitive}(w, (\lambda t \in \operatorname{ccint}(a, b).\; \operatorname{imag\text{-}part}(i(t))), a, b) \Rightarrow \\ \quad \quad (\int_{a}^{b} i = ((f(b) - f(a)) + ((w(b) - w(a)) \cdot +i))) \end{array}}}
\vnbentry{\texttt{cc-\allowbreak{}int-\allowbreak{}in-\allowbreak{}cc}}{\vnbstmt{\begin{array}{@{}l@{}} \forall a, b, i, f, w. \\ \quad (a \in \mathbb{R}) \wedge \\ \quad (b \in \mathbb{R}) \wedge \\ \quad \operatorname{is\text{-}primitive}(f, (\lambda t \in \operatorname{ccint}(a, b).\; \operatorname{real\text{-}part}(i(t))), a, b) \wedge \\ \quad \operatorname{is\text{-}primitive}(w, (\lambda t \in \operatorname{ccint}(a, b).\; \operatorname{imag\text{-}part}(i(t))), a, b) \Rightarrow \\ \quad \quad (\int_{a}^{b} i \in \mathbb{C}) \end{array}}}
\vnbentry{\texttt{cc-\allowbreak{}int-\allowbreak{}sum}}{\vnbstmt{\begin{array}{@{}l@{}} \forall a, b, i, \mathit{ppsi}, \mathit{pchi}, f, w, h, \mathit{pbw}. \\ \quad (a \in \mathbb{R}) \wedge \\ \quad (b \in \mathbb{R}) \wedge \\ \quad (i \in (\operatorname{ccint}(a, b) \to \mathbb{C})) \wedge \\ \quad (\mathit{ppsi} \in (\operatorname{ccint}(a, b) \to \mathbb{C})) \wedge \\ \quad (\mathit{pchi} \in (\operatorname{ccint}(a, b) \to \mathbb{C})) \wedge \\ \quad \forall y \in \operatorname{ccint}(a, b). \\ \quad \quad (\operatorname{pchi}(y) \simeq (i(y) + \operatorname{ppsi}(y))) \wedge \\ \quad \operatorname{is\text{-}primitive}(f, (\lambda t \in \operatorname{ccint}(a, b).\; \operatorname{real\text{-}part}(i(t))), a, b) \wedge \\ \quad \operatorname{is\text{-}primitive}(w, (\lambda t \in \operatorname{ccint}(a, b).\; \operatorname{imag\text{-}part}(i(t))), a, b) \wedge \\ \quad \operatorname{is\text{-}primitive}(h, \\ \quad \quad (\lambda t \in \operatorname{ccint}(a, b).\; \\ \quad \quad \quad \operatorname{real\text{-}part}(\operatorname{ppsi}(t))), \\ \quad \quad a, b) \wedge \\ \quad \operatorname{is\text{-}primitive}(\mathit{pbw}, \\ \quad \quad (\lambda t \in \operatorname{ccint}(a, b).\; \\ \quad \quad \quad \operatorname{imag\text{-}part}(\operatorname{ppsi}(t))), \\ \quad \quad a, b) \Rightarrow \\ \quad \quad (\int_{a}^{b} \mathit{pchi} = (\int_{a}^{b} i + \int_{a}^{b} \mathit{ppsi})) \end{array}}}
\vnbentry{\texttt{cc-\allowbreak{}int-\allowbreak{}real-\allowbreak{}mul}}{\vnbstmt{\begin{array}{@{}l@{}} \forall a, b, i, \mathit{pchi}.; \forall c \in \mathbb{R}, f, w. \\ \quad (a \in \mathbb{R}) \wedge \\ \quad (b \in \mathbb{R}) \wedge \\ \quad (i \in (\operatorname{ccint}(a, b) \to \mathbb{C})) \wedge \\ \quad (\mathit{pchi} \in (\operatorname{ccint}(a, b) \to \mathbb{C})) \wedge \\ \quad \forall y \in \operatorname{ccint}(a, b). \\ \quad \quad (\operatorname{pchi}(y) \simeq (c \cdot i(y))) \wedge \\ \quad \operatorname{is\text{-}primitive}(f, (\lambda t \in \operatorname{ccint}(a, b).\; \operatorname{real\text{-}part}(i(t))), a, b) \wedge \\ \quad \operatorname{is\text{-}primitive}(w, (\lambda t \in \operatorname{ccint}(a, b).\; \operatorname{imag\text{-}part}(i(t))), a, b) \Rightarrow \\ \quad \quad (\int_{a}^{b} \mathit{pchi} = (c \cdot \int_{a}^{b} i)) \end{array}}}

\vnbfile{theorem-library/taylor-cluster.scm}{theorem-library/taylor-cluster.scm}
\noindent{\small taylor-cluster.scm -- THE TAYLOR CLUSTER OF THE CALCULUS NOTES (calculus.pdf, ch. 2, sections 1, 6 and 7): 2.3, 2.16, 2.20, 2.21, 2.22.  Batch 25-B, 2026-09-24.}\par
\vnbentry{\texttt{nth-\allowbreak{}deriv-\allowbreak{}succ-\allowbreak{}inner}}{\vnbstmt{\begin{array}{@{}l@{}} \forall k \in \mathbb{N}, f. \\ \quad (\operatorname{nth\text{-}deriv}(f, (k + 1)) \simeq \operatorname{nth\text{-}deriv}((\lambda x \in \mathbb{R}.\; \operatorname{deriv}(f, x)), k)) \end{array}}}
\vnbentry{\texttt{nth-\allowbreak{}deriv-\allowbreak{}monomial-\allowbreak{}core}}{\vnbstmt{\begin{array}{@{}l@{}} \forall k, m \in \mathbb{N}. \\ \quad (\operatorname{nth\text{-}deriv}((\lambda x \in \mathbb{R}.\; \operatorname{power}(x, (m + k))), k) \simeq (\lambda x \in \mathbb{R}.\; ((\operatorname{factorial}((m + k)) \cdot \operatorname{recip}(\operatorname{factorial}(m))) \cdot \operatorname{power}(x, m)))) \end{array}}}
\vnbentry{\texttt{nth-\allowbreak{}deriv-\allowbreak{}monomial-\allowbreak{}gap}}{\vnbstmt{\begin{array}{@{}l@{}} \forall k, m \in \mathbb{N}, p. \\ \quad (p = (m + k)) \Rightarrow  \\ \quad \quad (\operatorname{nth\text{-}deriv}((\lambda x \in \mathbb{R}.\; \operatorname{power}(x, p)), k) \simeq (\lambda x \in \mathbb{R}.\; ((\operatorname{factorial}(p) \cdot \operatorname{recip}(\operatorname{factorial}(m))) \cdot \operatorname{power}(x, m)))) \end{array}}}
\vnbentry{\texttt{nth-\allowbreak{}deriv-\allowbreak{}monomial-\allowbreak{}beyond}}{\vnbstmt{\begin{array}{@{}l@{}} \forall j, n \in \mathbb{N}. \\ \quad (\operatorname{nth\text{-}deriv}((\lambda x \in \mathbb{R}.\; \operatorname{power}(x, n)), ((n + j) + 1)) \simeq (\lambda x \in \mathbb{R}.\; 0)) \end{array}}}
\vnbentry{\texttt{nth-\allowbreak{}deriv-\allowbreak{}monomial-\allowbreak{}vanish}}{\vnbstmt{\begin{array}{@{}l@{}} \forall n, k \in \mathbb{N}. \\ \quad (n < k) \Rightarrow  \\ \quad \quad (\operatorname{nth\text{-}deriv}((\lambda x \in \mathbb{R}.\; \operatorname{power}(x, n)), k) \simeq (\lambda x \in \mathbb{R}.\; 0)) \end{array}}}
\vnbentry{\texttt{nth-\allowbreak{}deriv-\allowbreak{}monomial}}{\vnbstmt{\begin{array}{@{}l@{}} \forall n, k \in \mathbb{N}. \\ \quad (k \leq n) \Rightarrow  \\ \quad \quad (\operatorname{nth\text{-}deriv}((\lambda x \in \mathbb{R}.\; \operatorname{power}(x, n)), k) \simeq (\lambda x \in \mathbb{R}.\; ((\operatorname{factorial}(n) \cdot \operatorname{recip}(\operatorname{factorial}(\operatorname{nn\text{-}minus}(n, k)))) \cdot \operatorname{power}(x, \operatorname{nn\text{-}minus}(n, k))))) \end{array}}}
\vnbentry{\texttt{poly-\allowbreak{}deriv-\allowbreak{}fn}}{\vnbstmt{\begin{array}{@{}l@{}} \forall n \in \mathbb{N}, f, d \in (\mathbb{N} \to \mathbb{R}). \\ \quad \forall j \in \mathbb{N}.\; (d(j) = ((j + 1) \cdot f((j + 1)))) \Rightarrow  \\ \quad \quad ((\lambda \mathit{tcx} \in \mathbb{R}.\; \operatorname{deriv}((\lambda x \in \mathbb{R}.\; \operatorname{series\text{-}partial\text{-}sum}((\lambda k \in \mathbb{N}.\; (f(k) \cdot \operatorname{power}(x, k))), (n + 1))), \mathit{tcx})) \simeq (\lambda x \in \mathbb{R}.\; \operatorname{series\text{-}partial\text{-}sum}((\lambda k \in \mathbb{N}.\; (d(k) \cdot \operatorname{power}(x, k))), n))) \end{array}}}
\vnbentry{\texttt{poly-\allowbreak{}deriv-\allowbreak{}coef-\allowbreak{}exists}}{\vnbstmt{\begin{array}{@{}l@{}} \forall f \in (\mathbb{N} \to \mathbb{R}). \exists d \in (\mathbb{N} \to \mathbb{R}). \forall j \in \mathbb{N}. \\ \quad (d(j) = ((j + 1) \cdot f((j + 1)))) \end{array}}}
\vnbentry{\texttt{poly-\allowbreak{}deriv-\allowbreak{}fn-\allowbreak{}zero}}{\vnbstmt{\begin{array}{@{}l@{}} \forall f \in (\mathbb{N} \to \mathbb{R}). \\ \quad ((\lambda \mathit{tcx} \in \mathbb{R}.\; \operatorname{deriv}((\lambda x \in \mathbb{R}.\; \operatorname{series\text{-}partial\text{-}sum}((\lambda k \in \mathbb{N}.\; (f(k) \cdot \operatorname{power}(x, k))), 0)), \mathit{tcx})) \simeq (\lambda x \in \mathbb{R}.\; \operatorname{series\text{-}partial\text{-}sum}((\lambda k \in \mathbb{N}.\; (f(k) \cdot \operatorname{power}(x, k))), 0))) \end{array}}}
\vnbentry{\texttt{poly-\allowbreak{}nth-\allowbreak{}deriv-\allowbreak{}in-\allowbreak{}fun}}{\vnbstmt{\begin{array}{@{}l@{}} \forall \mathit{tck}, n \in \mathbb{N}, f \in (\mathbb{N} \to \mathbb{R}). \\ \quad \operatorname{nth\text{-}deriv}((\lambda x \in \mathbb{R}.\; \\ \quad \quad \quad \operatorname{series\text{-}partial\text{-}sum}((\lambda k \in \mathbb{N}.\; (f(k) \cdot \operatorname{power}(x, k))), n)), \\ \quad \quad \mathit{tck}) \\ \quad \in\; (\mathbb{R} \to \mathbb{R}) \end{array}}}
\vnbentry{\texttt{poly-\allowbreak{}value-\allowbreak{}at-\allowbreak{}zero}}{\vnbstmt{\begin{array}{@{}l@{}} \forall n \in \mathbb{N}, f \in (\mathbb{N} \to \mathbb{R}). \\ \quad \operatorname{series\text{-}partial\text{-}sum}((\lambda k \in \mathbb{N}.\; (f(k) \cdot \operatorname{power}(0, k))), (n + 1)) \\ \quad =\; f(0) \end{array}}}
\vnbentry{\texttt{poly-\allowbreak{}nth-\allowbreak{}deriv-\allowbreak{}at-\allowbreak{}zero}}{\vnbstmt{\begin{array}{@{}l@{}} \forall \mathit{tck}, n \in \mathbb{N}, f \in (\mathbb{N} \to \mathbb{R}). \\ \quad (\mathit{tck} \leq n) \Rightarrow  \\ \quad \quad \operatorname{nth\text{-}deriv}((\lambda x \in \mathbb{R}.\; \operatorname{series\text{-}partial\text{-}sum}((\lambda k \in \mathbb{N}.\; (f(k) \cdot \operatorname{power}(x, k))), (n + 1))), \mathit{tck})(0) \\ \quad \quad =\; (\operatorname{factorial}(\mathit{tck}) \cdot f(\mathit{tck})) \end{array}}}
\vnbentry{\texttt{taylor-\allowbreak{}coef-\allowbreak{}from-\allowbreak{}derivs}}{\vnbstmt{\begin{array}{@{}l@{}} \forall n \in \mathbb{N}, f.; \forall a \in (\mathbb{N} \to \mathbb{R}), \mathit{tck} \in \mathbb{N}. \\ \quad \forall \mathit{tck} \in \mathbb{N}. \\ \quad \quad (\mathit{tck} \leq n) \Rightarrow  \\ \quad \quad \quad \operatorname{nth\text{-}deriv}((\lambda x \in \mathbb{R}.\; \operatorname{series\text{-}partial\text{-}sum}((\lambda k \in \mathbb{N}.\; (a(k) \cdot \operatorname{power}(x, k))), (n + 1))), \mathit{tck})(0) \\ \quad \quad \quad =\; \operatorname{nth\text{-}deriv}(f, \mathit{tck})(0) \wedge \\ \quad (\mathit{tck} \leq n) \Rightarrow \\ \quad \quad a(\mathit{tck}) \\ \quad \quad =\; (\operatorname{nth\text{-}deriv}(f, \mathit{tck})(0) \cdot \operatorname{recip}(\operatorname{factorial}(\mathit{tck}))) \end{array}}}
\vnbentry{\texttt{taylor-\allowbreak{}derivs-\allowbreak{}from-\allowbreak{}coef}}{\vnbstmt{\begin{array}{@{}l@{}} \forall n \in \mathbb{N}, f.; \forall a \in (\mathbb{N} \to \mathbb{R}), \mathit{tck} \in \mathbb{N}. \\ \quad \forall \mathit{tck} \in \mathbb{N}. \\ \quad \quad (\mathit{tck} \leq n) \Rightarrow  \\ \quad \quad \quad (\operatorname{nth\text{-}deriv}(f, \mathit{tck})(0) \in \mathbb{R}) \wedge \\ \quad \forall \mathit{tck} \in \mathbb{N}. \\ \quad \quad (\mathit{tck} \leq n) \Rightarrow  \\ \quad \quad \quad a(\mathit{tck}) \\ \quad \quad \quad =\; (\operatorname{nth\text{-}deriv}(f, \mathit{tck})(0) \cdot \operatorname{recip}(\operatorname{factorial}(\mathit{tck}))) \wedge \\ \quad (\mathit{tck} \leq n) \Rightarrow \\ \quad \quad \operatorname{nth\text{-}deriv}((\lambda x \in \mathbb{R}.\; \operatorname{series\text{-}partial\text{-}sum}((\lambda k \in \mathbb{N}.\; (a(k) \cdot \operatorname{power}(x, k))), (n + 1))), \mathit{tck})(0) \\ \quad \quad =\; \operatorname{nth\text{-}deriv}(f, \mathit{tck})(0) \end{array}}}
\vnbentry{\texttt{taylor-\allowbreak{}interpolant-\allowbreak{}at-\allowbreak{}zero}}{\vnbstmt{\begin{array}{@{}l@{}} \forall n \in \mathbb{N}, f. \\ \quad \forall \mathit{tck} \in \mathbb{N}. \\ \quad \quad (\mathit{tck} \leq n) \Rightarrow  \\ \quad \quad \quad (\operatorname{nth\text{-}deriv}(f, \mathit{tck})(0) \in \mathbb{R}) \Rightarrow \\ \quad \quad \exists a \in (\mathbb{N} \to \mathbb{R}). \\ \quad \quad \quad \forall \mathit{tck} \in \mathbb{N}. \\ \quad \quad \quad \quad (\mathit{tck} \leq n) \Rightarrow  \\ \quad \quad \quad \quad \quad a(\mathit{tck}) \\ \quad \quad \quad \quad \quad =\; (\operatorname{nth\text{-}deriv}(f, \mathit{tck})(0) \cdot \operatorname{recip}(\operatorname{factorial}(\mathit{tck}))) \wedge \\ \quad \quad \quad \forall \mathit{tck} \in \mathbb{N}. \\ \quad \quad \quad \quad (\mathit{tck} \leq n) \Rightarrow  \\ \quad \quad \quad \quad \quad \operatorname{nth\text{-}deriv}((\lambda x \in \mathbb{R}.\; \operatorname{series\text{-}partial\text{-}sum}((\lambda k \in \mathbb{N}.\; (a(k) \cdot \operatorname{power}(x, k))), (n + 1))), \mathit{tck})(0) \\ \quad \quad \quad \quad \quad =\; \operatorname{nth\text{-}deriv}(f, \mathit{tck})(0) \wedge \\ \quad \quad \quad \forall g \in (\mathbb{N} \to \mathbb{R}), \mathit{tck} \in \mathbb{N}. \\ \quad \quad \quad \quad \forall \mathit{tck} \in \mathbb{N}. \\ \quad \quad \quad \quad \quad (\mathit{tck} \leq n) \Rightarrow  \\ \quad \quad \quad \quad \quad \quad \operatorname{nth\text{-}deriv}((\lambda x \in \mathbb{R}.\; \operatorname{series\text{-}partial\text{-}sum}((\lambda k \in \mathbb{N}.\; (g(k) \cdot \operatorname{power}(x, k))), (n + 1))), \mathit{tck})(0) \\ \quad \quad \quad \quad \quad \quad =\; \operatorname{nth\text{-}deriv}(f, \mathit{tck})(0) \wedge \\ \quad \quad \quad \quad (\mathit{tck} \leq n) \Rightarrow \\ \quad \quad \quad \quad \quad (g(\mathit{tck}) = a(\mathit{tck})) \end{array}}}
\vnbentry{\texttt{series-\allowbreak{}partial-\allowbreak{}sum-\allowbreak{}shift}}{\vnbstmt{\begin{array}{@{}l@{}} \forall n \in \mathbb{N}, u \in (\mathbb{N} \to \mathbb{R}). \\ \quad \operatorname{series\text{-}partial\text{-}sum}(u, (n + 1)) \\ \quad =\; (u(0) + \operatorname{series\text{-}partial\text{-}sum}((\lambda i \in \mathbb{N}.\; u((i + 1))), n)) \end{array}}}
\vnbentry{\texttt{poly-\allowbreak{}degree-\allowbreak{}zero-\allowbreak{}value}}{\vnbstmt{\begin{array}{@{}l@{}} \forall f \in (\mathbb{N} \to \mathbb{R}), v \in \mathbb{R}. \\ \quad \operatorname{series\text{-}partial\text{-}sum}((\lambda k \in \mathbb{N}.\; (f(k) \cdot \operatorname{power}(v, k))), (0 + 1)) \\ \quad =\; f(0) \end{array}}}
\vnbentry{\texttt{taylor-\allowbreak{}exact-\allowbreak{}polynomial}}{\vnbstmt{\begin{array}{@{}l@{}} \forall m \in \mathbb{N}, f \in (\mathbb{N} \to \mathbb{R}), a, h \in \mathbb{R}. \\ \quad (\lambda x \in \mathbb{R}.\; \operatorname{series\text{-}partial\text{-}sum}((\lambda k \in \mathbb{N}.\; (f(k) \cdot \operatorname{power}(x, k))), (m + 1)))((a + h)) \\ \quad =\; \operatorname{series\text{-}partial\text{-}sum}((\lambda i \in \mathbb{N}.\; \\ \quad \quad \quad ((\operatorname{nth\text{-}deriv}((\lambda x \in \mathbb{R}.\; \operatorname{series\text{-}partial\text{-}sum}((\lambda k \in \mathbb{N}.\; (f(k) \cdot \operatorname{power}(x, k))), (m + 1))), i)(a) \cdot \operatorname{power}(h, i)) \cdot \operatorname{recip}(\operatorname{factorial}(i)))), \\ \quad \quad (m + 1)) \end{array}}}
\vnbentry{\texttt{taylor-\allowbreak{}remainder-\allowbreak{}sup-\allowbreak{}bound}}{\vnbstmt{\begin{array}{@{}l@{}} \forall a, b \in \mathbb{R}, f.; \forall m \in \mathbb{N}. \\ \quad \operatorname{is\text{-}taylor\text{-}family}(f, (m + 1), a, b) \wedge \\ \quad \operatorname{rr\text{-}bounded\text{-}above}(\operatorname{image}((\lambda v \in \operatorname{ooint}(a, b).\; \lvert f((m + 1))(v) \rvert), \operatorname{ooint}(a, b))) \Rightarrow \\ \quad \quad \lvert \operatorname{taylor\text{-}rem}(f, (m + 1), a, b) \rvert \\ \quad \quad \leq\; ((\operatorname{power}((b - a), (m + 1)) \cdot \operatorname{recip}(\operatorname{factorial}((m + 1)))) \cdot \operatorname{sup}(\operatorname{image}((\lambda v \in \operatorname{ooint}(a, b).\; \lvert f((m + 1))(v) \rvert), \operatorname{ooint}(a, b)))) \end{array}}}
\vnbentry{\texttt{series-\allowbreak{}partial-\allowbreak{}sum-\allowbreak{}agree-\allowbreak{}below}}{\vnbstmt{\begin{array}{@{}l@{}} \forall n \in \mathbb{N}, u, g. \\ \quad \forall k \in \mathbb{N}.\; ((k < n) \Rightarrow (g(k) \in \mathbb{R})) \wedge \\ \quad \forall k \in \mathbb{N}.\; ((k < n) \Rightarrow (u(k) = g(k))) \Rightarrow \\ \quad \quad \operatorname{series\text{-}partial\text{-}sum}(u, n) \\ \quad \quad =\; \operatorname{series\text{-}partial\text{-}sum}(g, n) \end{array}}}
\vnbentry{\texttt{taylor-\allowbreak{}lagrange-\allowbreak{}subinterval}}{\vnbstmt{\begin{array}{@{}l@{}} \forall a, b \in \mathbb{R}, f.; \forall m \in \mathbb{N}, t \in \mathbb{R}. \\ \quad \forall k \in \mathbb{N}. \\ \quad \quad (k \leq (m + 1)) \Rightarrow  \\ \quad \quad \quad (f(k) \in (\operatorname{ccint}(a, b) \to \mathbb{R})) \wedge \\ \quad \forall k \in \mathbb{N}. \\ \quad \quad (k < (m + 1)) \Rightarrow  \\ \quad \quad \quad \operatorname{is\text{-}continuous\text{-}on}(f(k), \operatorname{ccint}(a, b)) \wedge \\ \quad \forall k \in \mathbb{N}, x \in \operatorname{ooint}(a, b). \\ \quad \quad (k < (m + 1)) \Rightarrow  \\ \quad \quad \quad \operatorname{has\text{-}deriv\text{-}at}(f(k), x, f((k + 1))(x)) \wedge \\ \quad (a < t) \wedge \\ \quad (t \leq b) \Rightarrow \\ \quad \quad \exists i \in \operatorname{ooint}(a, t). \\ \quad \quad \quad \operatorname{taylor\text{-}rem}(f, (m + 1), a, t) \\ \quad \quad \quad =\; ((f((m + 1))(i) \cdot \operatorname{recip}(\operatorname{factorial}((m + 1)))) \cdot \operatorname{power}((t - a), (m + 1))) \end{array}}}
\vnbentry{\texttt{taylor-\allowbreak{}sum-\allowbreak{}at-\allowbreak{}base}}{\vnbstmt{\begin{array}{@{}l@{}} \forall j \in \mathbb{N}, f.; \forall a \in \mathbb{R}. \\ \quad \forall k \in \mathbb{N}.\; ((k \leq j) \Rightarrow (f(k)(a) \in \mathbb{R})) \Rightarrow  \\ \quad \quad \operatorname{series\text{-}partial\text{-}sum}((\lambda \mathit{tk} \in \mathbb{N}.\; \\ \quad \quad \quad \quad ((f(\mathit{tk})(a) \cdot \operatorname{power}((a - a), \mathit{tk})) \cdot \operatorname{recip}(\operatorname{factorial}(\mathit{tk})))), \\ \quad \quad \quad (j + 1)) \\ \quad \quad =\; f(0)(a) \end{array}}}
\vnbentry{\texttt{taylor-\allowbreak{}factor-\allowbreak{}value}}{\vnbstmt{\begin{array}{@{}l@{}} \forall a, b \in \mathbb{R}, f.; \forall m \in \mathbb{N}, t \in \mathbb{R}. \\ \quad \forall k \in \mathbb{N}. \\ \quad \quad (k \leq (m + 1)) \Rightarrow  \\ \quad \quad \quad (f(k) \in (\operatorname{ccint}(a, b) \to \mathbb{R})) \wedge \\ \quad \forall k \in \mathbb{N}. \\ \quad \quad (k < (m + 1)) \Rightarrow  \\ \quad \quad \quad \operatorname{is\text{-}continuous\text{-}on}(f(k), \operatorname{ccint}(a, b)) \wedge \\ \quad \forall k \in \mathbb{N}, x \in \operatorname{ooint}(a, b). \\ \quad \quad (k < (m + 1)) \Rightarrow  \\ \quad \quad \quad \operatorname{has\text{-}deriv\text{-}at}(f(k), x, f((k + 1))(x)) \wedge \\ \quad (a < t) \wedge \\ \quad (t \leq b) \Rightarrow \\ \quad \quad \exists i \in \operatorname{ooint}(a, t). \\ \quad \quad \quad ((\operatorname{taylor\text{-}rem}(f, (m + 1), a, t) \cdot \operatorname{factorial}((m + 1))) \cdot \operatorname{recip}(\operatorname{power}((t - a), (m + 1)))) \\ \quad \quad \quad =\; f((m + 1))(i) \end{array}}}
\vnbentry{\texttt{taylor-\allowbreak{}sum-\allowbreak{}continuous}}{\vnbstmt{\begin{array}{@{}l@{}} \forall a \in \mathbb{R}, f.; \forall m \in \mathbb{N}, p \in \mathbb{R}. \\ \quad \forall k \in \mathbb{N}.\; ((k < (m + 1)) \Rightarrow (f(k)(a) \in \mathbb{R})) \Rightarrow  \\ \quad \quad \operatorname{is\text{-}continuous\text{-}at}(\operatorname{rr\text{-}ms}, \operatorname{rr\text{-}ms}, (\lambda s \in \mathbb{R}.\; \operatorname{taylor\text{-}sum}(f, (m + 1), s, a)), p) \end{array}}}
\vnbentry{\texttt{taylor-\allowbreak{}factor-\allowbreak{}fn-\allowbreak{}continuous}}{\vnbstmt{\begin{array}{@{}l@{}} \forall a, b \in \mathbb{R}, f.; \forall m \in \mathbb{N}, p \in \mathbb{R}. \\ \quad (a \leq b) \wedge \\ \quad \operatorname{is\text{-}continuous\text{-}on}(f(0), \operatorname{ccint}(a, b)) \wedge \\ \quad \forall k \in \mathbb{N}.\; ((k < (m + 1)) \Rightarrow (f(k)(a) \in \mathbb{R})) \wedge \\ \quad (a < p) \Rightarrow \\ \quad \quad \operatorname{is\text{-}continuous\text{-}at}(\operatorname{rr\text{-}ms}, \operatorname{rr\text{-}ms}, \\ \quad \quad \quad (\lambda s \in \mathbb{R}.\; \\ \quad \quad \quad \quad (((\operatorname{extend\text{-}const}(f(0), a, b)(s) - \operatorname{taylor\text{-}sum}(f, (m + 1), s, a)) \cdot \operatorname{factorial}((m + 1))) \cdot {\operatorname{power}((s - a), (m + 1))}^{\ast\!-1})), \\ \quad \quad \quad p) \end{array}}}
\vnbentry{\texttt{taylor-\allowbreak{}factor-\allowbreak{}fn-\allowbreak{}value}}{\vnbstmt{\begin{array}{@{}l@{}} \forall a, b \in \mathbb{R}, f.; \forall m \in \mathbb{N}, t \in \mathbb{R}. \\ \quad \forall k \in \mathbb{N}. \\ \quad \quad (k \leq (m + 1)) \Rightarrow  \\ \quad \quad \quad (f(k) \in (\operatorname{ccint}(a, b) \to \mathbb{R})) \wedge \\ \quad \forall k \in \mathbb{N}. \\ \quad \quad (k < (m + 1)) \Rightarrow  \\ \quad \quad \quad \operatorname{is\text{-}continuous\text{-}on}(f(k), \operatorname{ccint}(a, b)) \wedge \\ \quad \forall k \in \mathbb{N}, x \in \operatorname{ooint}(a, b). \\ \quad \quad (k < (m + 1)) \Rightarrow  \\ \quad \quad \quad \operatorname{has\text{-}deriv\text{-}at}(f(k), x, f((k + 1))(x)) \wedge \\ \quad (a < t) \wedge \\ \quad (t \leq b) \Rightarrow \\ \quad \quad \exists i \in \operatorname{ooint}(a, t). \\ \quad \quad \quad (\lambda s \in \mathbb{R}.\; (((\operatorname{extend\text{-}const}(f(0), a, b)(s) - \operatorname{taylor\text{-}sum}(f, (m + 1), s, a)) \cdot \operatorname{factorial}((m + 1))) \cdot {\operatorname{power}((s - a), (m + 1))}^{\ast\!-1}))(t) \\ \quad \quad \quad =\; f((m + 1))(i) \end{array}}}
\vnbentry{\texttt{taylor-\allowbreak{}factor-\allowbreak{}glue-\allowbreak{}continuous}}{\vnbstmt{\begin{array}{@{}l@{}} \forall a, b \in \mathbb{R}, f.; \forall m \in \mathbb{N}, p \in \mathbb{R}. \\ \quad (a < b) \wedge \\ \quad \forall k \in \mathbb{N}. \\ \quad \quad (k \leq (m + 1)) \Rightarrow  \\ \quad \quad \quad (f(k) \in (\operatorname{ccint}(a, b) \to \mathbb{R})) \wedge \\ \quad \forall k \in \mathbb{N}. \\ \quad \quad (k < (m + 1)) \Rightarrow  \\ \quad \quad \quad \operatorname{is\text{-}continuous\text{-}on}(f(k), \operatorname{ccint}(a, b)) \wedge \\ \quad \forall k \in \mathbb{N}, x \in \operatorname{ooint}(a, b). \\ \quad \quad (k < (m + 1)) \Rightarrow  \\ \quad \quad \quad \operatorname{has\text{-}deriv\text{-}at}(f(k), x, f((k + 1))(x)) \wedge \\ \quad \operatorname{is\text{-}continuous\text{-}on}(f((m + 1)), \operatorname{ccint}(a, b)) \Rightarrow \\ \quad \quad \operatorname{is\text{-}continuous\text{-}at}(\operatorname{rr\text{-}ms}, \operatorname{rr\text{-}ms}, \\ \quad \quad \quad (\lambda s \in \mathbb{R}.\; \\ \quad \quad \quad \quad \left(\text{if }(s \leq a)\text{ then }f((m + 1))(a)\text{ else }(((\operatorname{extend\text{-}const}(f(0), a, b)(s) - \operatorname{taylor\text{-}sum}(f, (m + 1), s, a)) \cdot \operatorname{factorial}((m + 1))) \cdot {\operatorname{power}((s - a), (m + 1))}^{\ast\!-1})\right)), \\ \quad \quad \quad p) \end{array}}}
\vnbentry{\texttt{taylor-\allowbreak{}remainder-\allowbreak{}factor}}{\vnbstmt{\begin{array}{@{}l@{}} \forall a, b \in \mathbb{R}, f.; \forall m \in \mathbb{N}. \\ \quad (a < b) \wedge \\ \quad \forall k \in \mathbb{N}. \\ \quad \quad (k \leq (m + 1)) \Rightarrow  \\ \quad \quad \quad (f(k) \in (\operatorname{ccint}(a, b) \to \mathbb{R})) \wedge \\ \quad \forall k \in \mathbb{N}. \\ \quad \quad (k < (m + 1)) \Rightarrow  \\ \quad \quad \quad \operatorname{is\text{-}continuous\text{-}on}(f(k), \operatorname{ccint}(a, b)) \wedge \\ \quad \forall k \in \mathbb{N}, x \in \operatorname{ooint}(a, b). \\ \quad \quad (k < (m + 1)) \Rightarrow  \\ \quad \quad \quad \operatorname{has\text{-}deriv\text{-}at}(f(k), x, f((k + 1))(x)) \wedge \\ \quad \operatorname{is\text{-}continuous\text{-}on}(f((m + 1)), \operatorname{ccint}(a, b)) \Rightarrow \\ \quad \quad \exists s \in (\operatorname{ccint}(a, b) \to \mathbb{R}). \\ \quad \quad \quad \operatorname{is\text{-}continuous\text{-}on}(s, \operatorname{ccint}(a, b)) \wedge \\ \quad \quad \quad (s(a) = f((m + 1))(a)) \wedge \\ \quad \quad \quad \forall t \in \operatorname{ccint}(a, b). \\ \quad \quad \quad \quad \operatorname{taylor\text{-}rem}(f, (m + 1), a, t) \\ \quad \quad \quad \quad =\; ((s(t) \cdot \operatorname{power}((t - a), (m + 1))) \cdot \operatorname{recip}(\operatorname{factorial}((m + 1)))) \wedge \\ \quad \quad \quad \forall t \in \operatorname{ccint}(a, b). \\ \quad \quad \quad \quad \lvert s(t) \rvert \\ \quad \quad \quad \quad \leq\; \operatorname{sup}(\operatorname{image}((\lambda v \in \operatorname{ccint}(a, b).\; \lvert f((m + 1))(v) \rvert), \operatorname{ccint}(a, b))) \end{array}}}

\vnbfile{theorem-library/primitive-glue.scm}{theorem-library/primitive-glue.scm}
\noindent{\small primitive-glue.scm -- PRIMITIVES ON ADJACENT INTERVALS GLUE (batch 22-C).}\par
\vnbentry{\texttt{clamp-\allowbreak{}left-\allowbreak{}const}}{\vnbstmt{\begin{array}{@{}l@{}} \forall a, b, x \in \mathbb{R}. \\ \quad (((x \leq a) \wedge (a \leq b)) \Rightarrow (\operatorname{clamp}(a, b, x) = a)) \end{array}}}
\vnbentry{\texttt{clamp-\allowbreak{}right-\allowbreak{}const}}{\vnbstmt{\begin{array}{@{}l@{}} \forall a, b, x \in \mathbb{R}. \\ \quad (((a \leq b) \wedge (b \leq x)) \Rightarrow (\operatorname{clamp}(a, b, x) = b)) \end{array}}}
\vnbentry{\texttt{primitive-\allowbreak{}integrand-\allowbreak{}interior-\allowbreak{}congruence}}{\vnbstmt{\begin{array}{@{}l@{}} \forall f, i, \mathit{ppsi}, a, b. \\ \quad \operatorname{is\text{-}primitive}(f, i, a, b) \wedge \\ \quad (\mathit{ppsi} \in (\operatorname{ccint}(a, b) \to \mathbb{R})) \wedge \\ \quad \forall y \in \operatorname{ooint}(a, b).\; (i(y) \simeq \operatorname{ppsi}(y)) \Rightarrow \\ \quad \quad \operatorname{is\text{-}primitive}(f, \mathit{ppsi}, a, b) \end{array}}}
\vnbentry{\texttt{primitive-\allowbreak{}glue}}{\vnbstmt{\begin{array}{@{}l@{}} \forall a, c, b \in \mathbb{R}, d, h, k, m.; \forall g, f \in (\operatorname{ccint}(a, b) \to \mathbb{R}). \\ \quad (a < c) \wedge \\ \quad (c < b) \wedge \\ \quad \operatorname{is\text{-}primitive}(d, h, a, c) \wedge \\ \quad \operatorname{is\text{-}primitive}(k, m, c, b) \wedge \\ \quad (d(c) = k(c)) \wedge \\ \quad \forall t \in \operatorname{ccint}(a, c).\; (g(t) \simeq d(t)) \wedge \\ \quad \forall t \in \operatorname{ccint}(c, b).\; (g(t) \simeq k(t)) \wedge \\ \quad \forall t \in \operatorname{ooint}(a, c).\; (f(t) \simeq h(t)) \wedge \\ \quad \forall t \in \operatorname{ooint}(c, b).\; (f(t) \simeq m(t)) \Rightarrow \\ \quad \quad \operatorname{is\text{-}primitive}(g, f, a, b) \end{array}}}
\vnbentry{\texttt{primitive-\allowbreak{}of-\allowbreak{}constant-\allowbreak{}piece}}{\vnbstmt{\begin{array}{@{}l@{}} \forall u, v \in \mathbb{R}, f \in (\mathbb{R} \to \mathbb{R}), a, k \in \mathbb{R}. \\ \quad ((u < v) \wedge \forall t \in \mathbb{R}.\; (((u < t) \wedge (t < v)) \Rightarrow (f(t) = a))) \Rightarrow  \\ \quad \quad \operatorname{is\text{-}primitive}((\lambda s \in \operatorname{ccint}(u, v).\; (k + (a \cdot (s - u)))), \\ \quad \quad \quad \operatorname{restrict}(f, \operatorname{ccint}(u, v)), u, v) \end{array}}}
\vnbentry{\texttt{step-\allowbreak{}fn-\allowbreak{}has-\allowbreak{}primitive-\allowbreak{}ind}}{\vnbstmt{\begin{array}{@{}l@{}} \forall n \in \mathbb{N}, p, a, b.; \forall f \in (\mathbb{R} \to \mathbb{R}). \\ \quad \operatorname{is\text{-}partition}(p, n, a, b) \wedge \\ \quad \forall i \in \operatorname{interval}(0, n). \\ \quad \quad (i < n) \Rightarrow  \\ \quad \quad \quad \exists c \in \mathbb{R}. \forall t \in \mathbb{R}. \\ \quad \quad \quad \quad (((p(i) < t) \wedge (t < p((i + 1)))) \Rightarrow (f(t) = c)) \Rightarrow \\ \quad \quad \exists g. \\ \quad \quad \quad \operatorname{is\text{-}primitive}(g, \operatorname{restrict}(f, \operatorname{ccint}(a, b)), a, b) \end{array}}}
\vnbentry{\texttt{step-\allowbreak{}fn-\allowbreak{}has-\allowbreak{}primitive}}{\vnbstmt{\begin{array}{@{}l@{}} \forall a, b \in \mathbb{R}, f. \\ \quad \operatorname{is\text{-}step\text{-}fn}(f, a, b) \Rightarrow  \\ \quad \quad \exists g. \\ \quad \quad \quad \operatorname{is\text{-}primitive}(g, \operatorname{restrict}(f, \operatorname{ccint}(a, b)), a, b) \end{array}}}

\vnbfile{theorem-library/pw-int-laws-3.scm}{theorem-library/pw-int-laws-3.scm}
\noindent{\small pw-int-laws-3.scm -- LINEARITY over the reals as ONE law, the COMPLEX scalar, and the additivity of the complex integral over adjacent intervals.  Continues theorem-library/pw-int-laws-2.scm.}\par
\vnbentry{\texttt{pw-\allowbreak{}antiderivative-\allowbreak{}linear}}{\vnbstmt{\begin{array}{@{}l@{}} \forall a, b, c, d \in \mathbb{R}, f, w, i, \mathit{ppsi}.; \forall h, \mathit{pchi} \in (\operatorname{ccint}(a, b) \to \mathbb{R}). \\ \quad \operatorname{is\text{-}primitive}(f, i, a, b) \wedge \\ \quad \operatorname{is\text{-}primitive}(w, \mathit{ppsi}, a, b) \wedge \\ \quad \forall y \in \operatorname{ccint}(a, b). \\ \quad \quad (h(y) \simeq ((c \cdot f(y)) + (d \cdot w(y)))) \wedge \\ \quad \forall y \in \operatorname{ccint}(a, b). \\ \quad \quad (\operatorname{pchi}(y) \simeq ((c \cdot i(y)) + (d \cdot \operatorname{ppsi}(y)))) \Rightarrow \\ \quad \quad \operatorname{is\text{-}primitive}(h, \mathit{pchi}, a, b) \wedge \\ \quad \quad (\int_{a}^{b} \mathit{pchi} = ((c \cdot \int_{a}^{b} i) + (d \cdot \int_{a}^{b} \mathit{ppsi}))) \end{array}}}
\vnbentry{\texttt{cc-\allowbreak{}int-\allowbreak{}complex-\allowbreak{}mul}}{\vnbstmt{\begin{array}{@{}l@{}} \forall a, b, i, \mathit{pchi}.; \forall w \in \mathbb{C}, f, \mathit{paw}. \\ \quad (a \in \mathbb{R}) \wedge \\ \quad (b \in \mathbb{R}) \wedge \\ \quad (i \in (\operatorname{ccint}(a, b) \to \mathbb{C})) \wedge \\ \quad (\mathit{pchi} \in (\operatorname{ccint}(a, b) \to \mathbb{C})) \wedge \\ \quad \forall y \in \operatorname{ccint}(a, b). \\ \quad \quad (\operatorname{pchi}(y) \simeq (w \cdot i(y))) \wedge \\ \quad \operatorname{is\text{-}primitive}(f, (\lambda t \in \operatorname{ccint}(a, b).\; \operatorname{real\text{-}part}(i(t))), a, b) \wedge \\ \quad \operatorname{is\text{-}primitive}(\mathit{paw}, (\lambda t \in \operatorname{ccint}(a, b).\; \operatorname{imag\text{-}part}(i(t))), a, b) \Rightarrow \\ \quad \quad (\int_{a}^{b} \mathit{pchi} = (w \cdot \int_{a}^{b} i)) \end{array}}}
\vnbentry{\texttt{cc-\allowbreak{}int-\allowbreak{}adjacent}}{\vnbstmt{\begin{array}{@{}l@{}} \forall a, b.; \forall v \in \operatorname{ooint}(a, b), i, f, w. \\ \quad (a \in \mathbb{R}) \wedge \\ \quad (b \in \mathbb{R}) \wedge \\ \quad (i \in (\operatorname{ccint}(a, b) \to \mathbb{C})) \wedge \\ \quad \operatorname{is\text{-}primitive}(f, (\lambda t \in \operatorname{ccint}(a, b).\; \operatorname{real\text{-}part}(i(t))), a, b) \wedge \\ \quad \operatorname{is\text{-}primitive}(w, (\lambda t \in \operatorname{ccint}(a, b).\; \operatorname{imag\text{-}part}(i(t))), a, b) \Rightarrow \\ \quad \quad \int_{a}^{b} i \\ \quad \quad =\; (\int_{a}^{v} \operatorname{restrict}(i, \operatorname{ccint}(a, v)) + \int_{v}^{b} \operatorname{restrict}(i, \operatorname{ccint}(v, b))) \end{array}}}

\vnbfile{theorem-library/primitive-uniform-limit.scm}{theorem-library/primitive-uniform-limit.scm}
\noindent{\small theorem-library/primitive-uniform-limit.scm -- Dieudonne (8.6.4) on a compact interval [a,b], with COUNTABLE exceptional sets: A UNIFORM LIMIT OF PRIMITIVES IS A PRIMITIVE.  Batch 22-B, 2026-09-23.}\par
\vnbentry{\texttt{countable-\allowbreak{}union-\allowbreak{}nn}}{\vnbstmt{\begin{array}{@{}l@{}} \forall d.; \forall u \in \mathbf{Set}. \\ \quad \forall k \in \mathbb{N}.\; \operatorname{is\text{-}countable}(d(k)) \wedge \\ \quad \forall y \in u.\; \exists k \in \mathbb{N}.\; (y \in d(k)) \Rightarrow \\ \quad \quad \operatorname{is\text{-}countable}(u) \end{array}}}
\vnbentry{\texttt{primitive-\allowbreak{}pair-\allowbreak{}lipschitz}}{\vnbstmt{\begin{array}{@{}l@{}} \forall a, b, g_{1}, f_{1}, g_{2}, f_{2}.; \forall m \in \mathbb{R}, x, y \in \operatorname{ccint}(a, b). \\ \quad \operatorname{is\text{-}primitive}(g_{1}, f_{1}, a, b) \wedge \\ \quad \operatorname{is\text{-}primitive}(g_{2}, f_{2}, a, b) \wedge \\ \quad \forall x \in \operatorname{ccint}(a, b). \\ \quad \quad (\lvert (\operatorname{f1}(x) - \operatorname{f2}(x)) \rvert \leq m) \wedge \\ \quad (x < y) \Rightarrow \\ \quad \quad \lvert ((\operatorname{g1}(y) - \operatorname{g2}(y)) - (\operatorname{g1}(x) - \operatorname{g2}(x))) \rvert \\ \quad \quad \leq\; (m \cdot (y - x)) \end{array}}}
\vnbentry{\texttt{rr-\allowbreak{}mul-\allowbreak{}le-\allowbreak{}cancel-\allowbreak{}pos}}{\vnbstmt{\begin{array}{@{}l@{}} \forall a, c, m \in \mathbb{R}. \\ \quad (((0 < c) \wedge ((a \cdot c) \leq (m \cdot c))) \Rightarrow (a \leq m)) \end{array}}}
\vnbentry{\texttt{primitive-\allowbreak{}pair-\allowbreak{}lipschitz-\allowbreak{}abs}}{\vnbstmt{\begin{array}{@{}l@{}} \forall a, b, g_{1}, f_{1}, g_{2}, f_{2}.; \forall m \in \mathbb{R}, x, y \in \operatorname{ccint}(a, b). \\ \quad \operatorname{is\text{-}primitive}(g_{1}, f_{1}, a, b) \wedge \\ \quad \operatorname{is\text{-}primitive}(g_{2}, f_{2}, a, b) \wedge \\ \quad \forall x \in \operatorname{ccint}(a, b). \\ \quad \quad (\lvert (\operatorname{f1}(x) - \operatorname{f2}(x)) \rvert \leq m) \Rightarrow \\ \quad \quad \lvert ((\operatorname{g1}(y) - \operatorname{g2}(y)) - (\operatorname{g1}(x) - \operatorname{g2}(x))) \rvert \\ \quad \quad \leq\; (m \cdot \lvert (y - x) \rvert) \end{array}}}
\vnbentry{\texttt{primitive-\allowbreak{}family-\allowbreak{}limit}}{\vnbstmt{\begin{array}{@{}l@{}} \forall a, b \in \mathbb{R}, m, c \in (\mathbb{N} \to (\operatorname{ccint}(a, b) \to \mathbb{R})). \\ \quad (a < b) \wedge \\ \quad \forall k \in \mathbb{N}.\; \operatorname{is\text{-}primitive}(m(k), c(k), a, b) \wedge \\ \quad \forall s. \\ \quad \quad \operatorname{pos\text{-}rr}(s) \Rightarrow  \\ \quad \quad \quad \exists n \in \mathbb{N}. \forall k, l \in \mathbb{N}, x \in \operatorname{ccint}(a, b). \\ \quad \quad \quad \quad ((n \leq k) \wedge (n \leq l)) \Rightarrow  \\ \quad \quad \quad \quad \quad (\lvert (c(k)(x) - c(l)(x)) \rvert \leq s) \wedge \\ \quad \forall k \in \mathbb{N}.\; (m(k)(a) = 0) \Rightarrow \\ \quad \quad \exists \mathit{pufc} \in (\mathbb{R} \to \mathbb{R}). \\ \quad \quad \quad \forall x \in \operatorname{ccint}(a, b). \\ \quad \quad \quad \quad \operatorname{converges\text{-}to}(\operatorname{rr\text{-}ms}, (\lambda k \in \mathbb{N}.\; m(k)(x)), \operatorname{pufc}(x)) \wedge \\ \quad \quad \quad \forall t \in \mathbb{R}. \\ \quad \quad \quad \quad \operatorname{is\text{-}continuous\text{-}at}(\operatorname{rr\text{-}ms}, \operatorname{rr\text{-}ms}, \mathit{pufc}, t) \end{array}}}
\vnbentry{\texttt{primitive-\allowbreak{}family-\allowbreak{}deriv-\allowbreak{}at}}{\vnbstmt{\begin{array}{@{}l@{}} \forall a, b \in \mathbb{R}, m, c \in (\mathbb{N} \to (\operatorname{ccint}(a, b) \to \mathbb{R})), \mathit{pufc} \in (\mathbb{R} \to \mathbb{R}), t \in \operatorname{ooint}(a, b), v \in \mathbb{R}. \\ \quad (a < b) \wedge \\ \quad \forall k \in \mathbb{N}.\; \operatorname{is\text{-}primitive}(m(k), c(k), a, b) \wedge \\ \quad \forall s. \\ \quad \quad \operatorname{pos\text{-}rr}(s) \Rightarrow  \\ \quad \quad \quad \exists n \in \mathbb{N}. \forall k, l \in \mathbb{N}, x \in \operatorname{ccint}(a, b). \\ \quad \quad \quad \quad ((n \leq k) \wedge (n \leq l)) \Rightarrow  \\ \quad \quad \quad \quad \quad (\lvert (c(k)(x) - c(l)(x)) \rvert \leq s) \wedge \\ \quad \forall x \in \operatorname{ccint}(a, b). \\ \quad \quad \operatorname{converges\text{-}to}(\operatorname{rr\text{-}ms}, (\lambda k \in \mathbb{N}.\; m(k)(x)), \operatorname{pufc}(x)) \wedge \\ \quad \forall k \in \mathbb{N}. \\ \quad \quad \operatorname{has\text{-}deriv\text{-}at}(m(k), t, c(k)(t)) \wedge \\ \quad \operatorname{converges\text{-}to}(\operatorname{rr\text{-}ms}, (\lambda k \in \mathbb{N}.\; c(k)(t)), v) \Rightarrow \\ \quad \quad \operatorname{is\text{-}diff\text{-}at}(\mathit{pufc}, t, v) \end{array}}}
\vnbentry{\texttt{primitive-\allowbreak{}uniform-\allowbreak{}limit}}{\vnbstmt{\begin{array}{@{}l@{}} \forall a, b \in \mathbb{R}, m, c \in (\mathbb{N} \to (\operatorname{ccint}(a, b) \to \mathbb{R})), f \in (\operatorname{ccint}(a, b) \to \mathbb{R}). \\ \quad (a < b) \wedge \\ \quad \forall k \in \mathbb{N}.\; \operatorname{is\text{-}primitive}(m(k), c(k), a, b) \wedge \\ \quad \forall s. \\ \quad \quad \operatorname{pos\text{-}rr}(s) \Rightarrow  \\ \quad \quad \quad \exists n \in \mathbb{N}. \forall k \in \mathbb{N}, x \in \operatorname{ccint}(a, b). \\ \quad \quad \quad \quad ((n \leq k) \Rightarrow (\lvert (c(k)(x) - f(x)) \rvert \leq s)) \wedge \\ \quad \forall k \in \mathbb{N}.\; (m(k)(a) = 0) \Rightarrow \\ \quad \quad \exists g. \\ \quad \quad \quad \operatorname{is\text{-}primitive}(g, f, a, b) \wedge \\ \quad \quad \quad \forall x \in \operatorname{ccint}(a, b). \\ \quad \quad \quad \quad \operatorname{converges\text{-}to}(\operatorname{rr\text{-}ms}, (\lambda k \in \mathbb{N}.\; m(k)(x)), g(x)) \end{array}}}

\vnbfile{theorem-library/cc-power-series-laws.scm}{theorem-library/cc-power-series-laws.scm}
\noindent{\small cc-power-series-laws.scm -- COMPLEX POWER SERIES: THE RADIUS OF CONVERGENCE (Prop 2.7) AND WHAT FOLLOWS FROM IT, as the notes state them. Batch 25-A, 2026-09-24.}\par
\vnbentry{\texttt{rr-\allowbreak{}power-\allowbreak{}mul-\allowbreak{}base}}{\vnbstmt{\begin{array}{@{}l@{}} \forall k \in \mathbb{N}, x, y \in \mathbb{R}. \\ \quad \operatorname{power}((x \cdot y), k) \\ \quad =\; (\operatorname{power}(x, k) \cdot \operatorname{power}(y, k)) \end{array}}}
\vnbentry{\texttt{cps-\allowbreak{}root-\allowbreak{}conv}}{\vnbstmt{\begin{array}{@{}l@{}} \forall f \in (\mathbb{N} \to \mathbb{C}), s, r \in \mathbb{R}. \\ \quad \operatorname{elimsup}((\lambda j \in \mathbb{N}.\; {\operatorname{magnitude}(f(j))}^{{j}^{\ast\!-1}})) \\ \quad <\; s \wedge \\ \quad (0 \leq r) \wedge \\ \quad ((r \cdot s) < 1) \Rightarrow \\ \quad \quad \operatorname{series\text{-}converges}((\lambda k \in \mathbb{N}.\; (\operatorname{magnitude}(f(k)) \cdot \operatorname{power}(r, k)))) \end{array}}}
\vnbentry{\texttt{cps-\allowbreak{}root-\allowbreak{}div}}{\vnbstmt{\begin{array}{@{}l@{}} \forall f \in (\mathbb{N} \to \mathbb{C}), s, r \in \mathbb{R}. \\ \quad (0 \leq s) \wedge \\ \quad s \\ \quad <\; \operatorname{elimsup}((\lambda j \in \mathbb{N}.\; {\operatorname{magnitude}(f(j))}^{{j}^{\ast\!-1}})) \wedge \\ \quad (0 \leq r) \wedge \\ \quad (1 \leq (r \cdot s)) \Rightarrow \\ \quad \quad \neg \operatorname{series\text{-}converges}((\lambda k \in \mathbb{N}.\; (\operatorname{magnitude}(f(k)) \cdot \operatorname{power}(r, k)))) \end{array}}}
\vnbentry{\texttt{cps-\allowbreak{}small-\allowbreak{}factor}}{\vnbstmt{\begin{array}{@{}l@{}} \forall r \in \mathbb{R}. \\ \quad (0 \leq r) \Rightarrow  \\ \quad \quad \exists s \in \mathbb{R}.\; ((0 < s) \wedge ((r \cdot s) < 1)) \end{array}}}
\vnbentry{\texttt{cps-\allowbreak{}radius-\allowbreak{}infinite}}{\vnbstmt{\begin{array}{@{}l@{}} \forall f \in (\mathbb{N} \to \mathbb{C}). \\ \quad \operatorname{elimsup}((\lambda j \in \mathbb{N}.\; {\operatorname{magnitude}(f(j))}^{{j}^{\ast\!-1}})) \\ \quad =\; 0 \Rightarrow \\ \quad \quad (\operatorname{cps\text{-}radius}(f) = \operatorname{pos\text{-}inf}) \end{array}}}
\vnbentry{\texttt{cps-\allowbreak{}radius-\allowbreak{}zero}}{\vnbstmt{\begin{array}{@{}l@{}} \forall f \in (\mathbb{N} \to \mathbb{C}). \\ \quad \operatorname{elimsup}((\lambda j \in \mathbb{N}.\; {\operatorname{magnitude}(f(j))}^{{j}^{\ast\!-1}})) \\ \quad =\; \operatorname{pos\text{-}inf} \Rightarrow \\ \quad \quad (\operatorname{cps\text{-}radius}(f) = 0) \end{array}}}
\vnbentry{\texttt{cps-\allowbreak{}radius-\allowbreak{}recip}}{\vnbstmt{\begin{array}{@{}l@{}} \forall f \in (\mathbb{N} \to \mathbb{C}). \\ \quad \operatorname{elimsup}((\lambda j \in \mathbb{N}.\; {\operatorname{magnitude}(f(j))}^{{j}^{\ast\!-1}})) \\ \quad \in\; \mathbb{R} \wedge \\ \quad 0 \\ \quad <\; \operatorname{elimsup}((\lambda j \in \mathbb{N}.\; {\operatorname{magnitude}(f(j))}^{{j}^{\ast\!-1}})) \Rightarrow \\ \quad \quad \operatorname{cps\text{-}radius}(f) \\ \quad \quad =\; \operatorname{recip}(\operatorname{elimsup}((\lambda j \in \mathbb{N}.\; {\operatorname{magnitude}(f(j))}^{{j}^{\ast\!-1}}))) \end{array}}}
\vnbentry{\texttt{cps-\allowbreak{}radius-\allowbreak{}root}}{\vnbstmt{\begin{array}{@{}l@{}} \forall f \in (\mathbb{N} \to \mathbb{C}). \\ \quad \operatorname{elimsup}((\lambda j \in \mathbb{N}.\; {\operatorname{magnitude}(f(j))}^{{j}^{\ast\!-1}})) \\ \quad =\; 0 \Rightarrow \\ \quad \quad (\operatorname{cps\text{-}radius}(f) = \operatorname{pos\text{-}inf}) \wedge \\ \quad \operatorname{elimsup}((\lambda j \in \mathbb{N}.\; {\operatorname{magnitude}(f(j))}^{{j}^{\ast\!-1}})) \\ \quad =\; \operatorname{pos\text{-}inf} \Rightarrow \\ \quad \quad (\operatorname{cps\text{-}radius}(f) = 0) \wedge \\ \quad \operatorname{elimsup}((\lambda j \in \mathbb{N}.\; {\operatorname{magnitude}(f(j))}^{{j}^{\ast\!-1}})) \\ \quad \in\; \mathbb{R} \wedge \\ \quad 0 \\ \quad <\; \operatorname{elimsup}((\lambda j \in \mathbb{N}.\; {\operatorname{magnitude}(f(j))}^{{j}^{\ast\!-1}})) \Rightarrow \\ \quad \quad \operatorname{cps\text{-}radius}(f) \\ \quad \quad =\; \operatorname{recip}(\operatorname{elimsup}((\lambda j \in \mathbb{N}.\; {\operatorname{magnitude}(f(j))}^{{j}^{\ast\!-1}}))) \end{array}}}
\vnbentry{\texttt{cps-\allowbreak{}radius-\allowbreak{}unfold}}{\vnbstmt{\begin{array}{@{}l@{}} \forall f. \\ \quad (\operatorname{cps\text{-}radius}(f) \simeq \operatorname{esup}(\{r \in \mathbb{R} : ((0 \leq r) \wedge \operatorname{series\text{-}converges}((\lambda k \in \mathbb{N}.\; (\operatorname{magnitude}(f(k)) \cdot \operatorname{power}(r, k)))))\})) \end{array}}}
\vnbentry{\texttt{cc-\allowbreak{}closed-\allowbreak{}ball-\allowbreak{}origin}}{\vnbstmt{\begin{array}{@{}l@{}} \forall d.; \forall z \in \operatorname{closed\text{-}ball}(\operatorname{cc\text{-}ms}, 0, d). \\ \quad ((z \in \mathbb{C}) \wedge (\operatorname{magnitude}(z) \leq d)) \end{array}}}
\vnbentry{\texttt{cc-\allowbreak{}magnitude-\allowbreak{}power}}{\vnbstmt{\begin{array}{@{}l@{}} \forall k \in \mathbb{N}, z \in \mathbb{C}. \\ \quad \operatorname{magnitude}(\operatorname{power}(z, k)) \\ \quad =\; \operatorname{power}(\operatorname{magnitude}(z), k) \end{array}}}
\vnbentry{\texttt{cc-\allowbreak{}sup-\allowbreak{}norm-\allowbreak{}bound}}{\vnbstmt{\begin{array}{@{}l@{}} \forall u, s.; \forall m \in \mathbb{R}. \\ \quad (0 \leq m) \wedge \\ \quad \forall x \in s.\; (u(x) \in \mathbb{C}) \wedge \\ \quad \forall x \in s.\; (\operatorname{magnitude}(u(x)) \leq m) \Rightarrow \\ \quad \quad (\operatorname{cc\text{-}sup\text{-}norm}(u, s) \in \mathbb{R}) \wedge \\ \quad \quad (0 \leq \operatorname{cc\text{-}sup\text{-}norm}(u, s)) \wedge \\ \quad \quad (\operatorname{cc\text{-}sup\text{-}norm}(u, s) \leq m) \end{array}}}
\vnbentry{\texttt{cps-\allowbreak{}radius-\allowbreak{}approx}}{\vnbstmt{\begin{array}{@{}l@{}} \forall f \in (\mathbb{N} \to \mathbb{C}), d \in \mathbb{R}. \\ \quad ((0 \leq d) \wedge (d < \operatorname{cps\text{-}radius}(f))) \Rightarrow  \\ \quad \quad \exists t \in \mathbb{R}. \\ \quad \quad \quad (d < t) \wedge \\ \quad \quad \quad \operatorname{series\text{-}converges}((\lambda k \in \mathbb{N}.\; (\operatorname{magnitude}(f(k)) \cdot \operatorname{power}(t, k)))) \end{array}}}
\vnbentry{\texttt{cps-\allowbreak{}normally-\allowbreak{}convergent}}{\vnbstmt{\begin{array}{@{}l@{}} \forall f \in (\mathbb{N} \to \mathbb{C}), d \in \mathbb{R}. \\ \quad ((0 \leq d) \wedge (d < \operatorname{cps\text{-}radius}(f))) \Rightarrow  \\ \quad \quad \operatorname{is\text{-}normally\text{-}convergent}((\lambda k \in \mathbb{N}.\; \\ \quad \quad \quad \quad (\lambda z \in \operatorname{closed\text{-}ball}(\operatorname{cc\text{-}ms}, 0, d).\; \\ \quad \quad \quad \quad \quad (f(k) \cdot \operatorname{power}(z, k)))), \\ \quad \quad \quad \operatorname{closed\text{-}ball}(\operatorname{cc\text{-}ms}, 0, d)) \end{array}}}
\vnbentry{\texttt{cps-\allowbreak{}abs-\allowbreak{}conv-\allowbreak{}closed-\allowbreak{}ball}}{\vnbstmt{\begin{array}{@{}l@{}} \forall f \in (\mathbb{N} \to \mathbb{C}), d \in \mathbb{R}, z \in \operatorname{closed\text{-}ball}(\operatorname{cc\text{-}ms}, 0, d). \\ \quad ((0 \leq d) \wedge (d < \operatorname{cps\text{-}radius}(f))) \Rightarrow  \\ \quad \quad \operatorname{series\text{-}converges}((\lambda k \in \mathbb{N}.\; \operatorname{magnitude}((f(k) \cdot \operatorname{power}(z, k))))) \end{array}}}
\vnbentry{\texttt{cc-\allowbreak{}series-\allowbreak{}block-\allowbreak{}le}}{\vnbstmt{\begin{array}{@{}l@{}} \forall k, m \in \mathbb{N}, f \in (\mathbb{N} \to \mathbb{C}), g \in (\mathbb{N} \to \mathbb{R}). \\ \quad (\forall j \in \mathbb{N}.\; (g(j) = \operatorname{magnitude}(f(j))) \wedge (m \leq k)) \Rightarrow  \\ \quad \quad \operatorname{magnitude}((\operatorname{cc\text{-}series\text{-}partial\text{-}sum}(f, k) - \operatorname{cc\text{-}series\text{-}partial\text{-}sum}(f, m))) \\ \quad \quad \leq\; (\operatorname{series\text{-}partial\text{-}sum}(g, k) - \operatorname{series\text{-}partial\text{-}sum}(g, m)) \end{array}}}
\vnbentry{\texttt{cc-\allowbreak{}series-\allowbreak{}abs-\allowbreak{}converges}}{\vnbstmt{\begin{array}{@{}l@{}} \forall f \in (\mathbb{N} \to \mathbb{C}), g \in (\mathbb{N} \to \mathbb{R}). \\ \quad \forall j \in \mathbb{N}.\; (g(j) = \operatorname{magnitude}(f(j))) \wedge \\ \quad \operatorname{series\text{-}converges}(g) \Rightarrow \\ \quad \quad \operatorname{cc\text{-}series\text{-}converges}(f) \end{array}}}
\vnbentry{\texttt{cc-\allowbreak{}series-\allowbreak{}limit-\allowbreak{}converges-\allowbreak{}to}}{\vnbstmt{\begin{array}{@{}l@{}} \forall f \in (\mathbb{N} \to \mathbb{C}). \\ \quad \operatorname{cc\text{-}series\text{-}converges}(f) \Rightarrow  \\ \quad \quad \operatorname{cc\text{-}series\text{-}converges\text{-}to}(f, \operatorname{cc\text{-}series\text{-}limit}(f)) \end{array}}}
\vnbentry{\texttt{cc-\allowbreak{}series-\allowbreak{}limit-\allowbreak{}in-\allowbreak{}cc}}{\vnbstmt{\begin{array}{@{}l@{}} \forall f \in (\mathbb{N} \to \mathbb{C}). \\ \quad \operatorname{cc\text{-}series\text{-}converges}(f) \Rightarrow  \\ \quad \quad (\operatorname{cc\text{-}series\text{-}limit}(f) \in \mathbb{C}) \end{array}}}
\vnbentry{\texttt{cps-\allowbreak{}converges-\allowbreak{}closed-\allowbreak{}ball}}{\vnbstmt{\begin{array}{@{}l@{}} \forall f \in (\mathbb{N} \to \mathbb{C}), d \in \mathbb{R}, z \in \operatorname{closed\text{-}ball}(\operatorname{cc\text{-}ms}, 0, d). \\ \quad ((0 \leq d) \wedge (d < \operatorname{cps\text{-}radius}(f))) \Rightarrow  \\ \quad \quad \operatorname{cc\text{-}series\text{-}converges}((\lambda k \in \mathbb{N}.\; (f(k) \cdot \operatorname{power}(z, k)))) \end{array}}}
\vnbentry{\texttt{cps-\allowbreak{}uniform-\allowbreak{}closed-\allowbreak{}ball}}{\vnbstmt{\begin{array}{@{}l@{}} \forall f \in (\mathbb{N} \to \mathbb{C}), d \in \mathbb{R}, e. \\ \quad (0 \leq d) \wedge \\ \quad (d < \operatorname{cps\text{-}radius}(f)) \wedge \\ \quad \operatorname{pos\text{-}rr}(e) \Rightarrow \\ \quad \quad \exists n \in \mathbb{N}. \forall m \in \mathbb{N}, z \in \operatorname{closed\text{-}ball}(\operatorname{cc\text{-}ms}, 0, d). \\ \quad \quad \quad (n \leq m) \Rightarrow  \\ \quad \quad \quad \quad \operatorname{magnitude}((\operatorname{cc\text{-}series\text{-}limit}((\lambda k \in \mathbb{N}.\; (f(k) \cdot \operatorname{power}(z, k)))) - \operatorname{cc\text{-}series\text{-}partial\text{-}sum}((\lambda k \in \mathbb{N}.\; (f(k) \cdot \operatorname{power}(z, k))), m))) \\ \quad \quad \quad \quad \leq\; e \end{array}}}
\vnbentry{\texttt{rr-\allowbreak{}power-\allowbreak{}power}}{\vnbstmt{\begin{array}{@{}l@{}} \forall k, m \in \mathbb{N}, x \in \mathbb{R}. \\ \quad \operatorname{power}(\operatorname{power}(x, k), m) \\ \quad =\; \operatorname{power}(\operatorname{power}(x, m), k) \end{array}}}
\vnbentry{\texttt{cps-\allowbreak{}poly-\allowbreak{}geom-\allowbreak{}bound}}{\vnbstmt{\begin{array}{@{}l@{}} \forall m \in \mathbb{N}, q \in \mathbb{R}. \\ \quad ((0 \leq q) \wedge (q < 1)) \Rightarrow  \\ \quad \quad \exists b \in \mathbb{R}. \\ \quad \quad \quad (0 \leq b) \wedge \\ \quad \quad \quad \forall k \in \mathbb{N}. \\ \quad \quad \quad \quad ((\operatorname{power}(k, m) \cdot \operatorname{power}(q, k)) \leq b) \end{array}}}
\vnbentry{\texttt{series-\allowbreak{}scale-\allowbreak{}converges}}{\vnbstmt{\begin{array}{@{}l@{}} \forall f \in (\mathbb{N} \to \mathbb{R}), c \in \mathbb{R}, h \in (\mathbb{N} \to \mathbb{R}). \\ \quad \forall k \in \mathbb{N}.\; (0 \leq f(k)) \wedge \\ \quad (0 \leq c) \wedge \\ \quad \forall k \in \mathbb{N}.\; (h(k) = (c \cdot f(k))) \wedge \\ \quad \operatorname{series\text{-}converges}(f) \Rightarrow \\ \quad \quad \operatorname{series\text{-}converges}(h) \end{array}}}
\vnbentry{\texttt{cps-\allowbreak{}factor-\allowbreak{}conv}}{\vnbstmt{\begin{array}{@{}l@{}} \forall f, g, a \in (\mathbb{N} \to \mathbb{C}), c \in \mathbb{R}, m \in \mathbb{N}, t, d \in \mathbb{R}. \\ \quad \forall k \in \mathbb{N}.\; (a(k) = (f(k) \cdot g(k))) \wedge \\ \quad \forall k \in \mathbb{N}. \\ \quad \quad (\operatorname{magnitude}(g(k)) \leq (c \cdot \operatorname{power}(k, m))) \wedge \\ \quad (0 \leq d) \wedge \\ \quad (d < t) \wedge \\ \quad \operatorname{series\text{-}converges}((\lambda k \in \mathbb{N}.\; (\operatorname{magnitude}(f(k)) \cdot \operatorname{power}(t, k)))) \Rightarrow \\ \quad \quad \operatorname{series\text{-}converges}((\lambda k \in \mathbb{N}.\; (\operatorname{magnitude}(a(k)) \cdot \operatorname{power}(d, k)))) \end{array}}}
\vnbentry{\texttt{cps-\allowbreak{}radius-\allowbreak{}factor}}{\vnbstmt{\begin{array}{@{}l@{}} \forall f, g, a \in (\mathbb{N} \to \mathbb{C}), c \in \mathbb{R}, m \in \mathbb{N}. \\ \quad \forall k \in \mathbb{N}.\; (a(k) = (f(k) \cdot g(k))) \wedge \\ \quad \forall k \in \mathbb{N}. \\ \quad \quad (\operatorname{magnitude}(g(k)) \leq (c \cdot \operatorname{power}(k, m))) \Rightarrow \\ \quad \quad (\operatorname{cps\text{-}radius}(f) \leq \operatorname{cps\text{-}radius}(a)) \end{array}}}

\vnbfile{theorem-library/series-domination.scm}{theorem-library/series-domination.scm}
\noindent{\small series-domination.scm -- SHIFTING AND EVENTUAL DOMINATION FOR REAL SERIES. Batch 26-A, 2026-09-24.  The two lemmas the batch-25 report named as standing in the way of Proposition 2.10 (theorem-library/cc-power-series-deriv.scm):}\par
\vnbentry{\texttt{series-\allowbreak{}shift-\allowbreak{}partial-\allowbreak{}sum}}{\vnbstmt{\begin{array}{@{}l@{}} \forall k \in \mathbb{N}, f, g \in (\mathbb{N} \to \mathbb{R}). \\ \quad \forall j \in \mathbb{N}.\; (g(j) = f((j + 1))) \Rightarrow  \\ \quad \quad \operatorname{series\text{-}partial\text{-}sum}(f, (k + 1)) \\ \quad \quad =\; (f(0) + \operatorname{series\text{-}partial\text{-}sum}(g, k)) \end{array}}}
\vnbentry{\texttt{series-\allowbreak{}shift-\allowbreak{}converges-\allowbreak{}iff}}{\vnbstmt{\begin{array}{@{}l@{}} \forall f, g \in (\mathbb{N} \to \mathbb{R}). \\ \quad (\forall j \in \mathbb{N}.\; (0 \leq f(j)) \wedge \forall j \in \mathbb{N}.\; (g(j) = f((j + 1)))) \Rightarrow  \\ \quad \quad (\operatorname{series\text{-}converges}(f) \Leftrightarrow \operatorname{series\text{-}converges}(g)) \end{array}}}
\vnbentry{\texttt{series-\allowbreak{}dominated-\allowbreak{}partial-\allowbreak{}bound}}{\vnbstmt{\begin{array}{@{}l@{}} \forall k \in \mathbb{N}, f, g \in (\mathbb{N} \to \mathbb{R}), n \in \mathbb{N}. \\ \quad \forall j \in \mathbb{N}.\; (0 \leq f(j)) \wedge \\ \quad \forall j \in \mathbb{N}.\; (0 \leq g(j)) \wedge \\ \quad \forall j \in \mathbb{N}.\; ((n \leq j) \Rightarrow (f(j) \leq g(j))) \Rightarrow \\ \quad \quad \operatorname{series\text{-}partial\text{-}sum}(f, k) \\ \quad \quad \leq\; (\operatorname{series\text{-}partial\text{-}sum}(f, n) + \operatorname{series\text{-}partial\text{-}sum}(g, k)) \end{array}}}
\vnbentry{\texttt{series-\allowbreak{}eventually-\allowbreak{}dominated-\allowbreak{}converges}}{\vnbstmt{\begin{array}{@{}l@{}} \forall f, g \in (\mathbb{N} \to \mathbb{R}), n \in \mathbb{N}. \\ \quad \forall j \in \mathbb{N}.\; (0 \leq f(j)) \wedge \\ \quad \forall j \in \mathbb{N}.\; (0 \leq g(j)) \wedge \\ \quad \forall j \in \mathbb{N}.\; ((n \leq j) \Rightarrow (f(j) \leq g(j))) \wedge \\ \quad \operatorname{series\text{-}converges}(g) \Rightarrow \\ \quad \quad \operatorname{series\text{-}converges}(f) \end{array}}}

\vnbfile{theorem-library/cc-power-series-deriv.scm}{theorem-library/cc-power-series-deriv.scm}
\noindent{\small cc-power-series-deriv.scm -- PROPOSITION 2.10: THE SUM OF A COMPLEX POWER SERIES IS HOLOMORPHIC INSIDE THE RADIUS, AND ITS DERIVATIVE IS THE TERMWISE DERIVED SERIES, WHICH HAS THE SAME RADIUS. Batch 26-A, 2026-09-24.}\par
\vnbentry{\texttt{cpd-\allowbreak{}cc-\allowbreak{}squeeze}}{\vnbstmt{\begin{array}{@{}l@{}} \forall s \in (\mathbb{N} \to \mathbb{C}), v \in \mathbb{C}, m \in (\mathbb{N} \to \mathbb{R}). \\ \quad \operatorname{converges\text{-}to}(\operatorname{rr\text{-}ms}, m, 0) \wedge \\ \quad \forall j \in \mathbb{N}.\; (\operatorname{magnitude}((s(j) - v)) \leq m(j)) \Rightarrow \\ \quad \quad \operatorname{converges\text{-}to}(\operatorname{cc\text{-}ms}, s, v) \end{array}}}
\vnbentry{\texttt{cc-\allowbreak{}limit-\allowbreak{}lincomb}}{\vnbstmt{\begin{array}{@{}l@{}} \forall s, t, h \in (\mathbb{N} \to \mathbb{C}), v, a, c, d \in \mathbb{C}. \\ \quad \forall j \in \mathbb{N}.\; (h(j) = ((c \cdot s(j)) + (d \cdot t(j)))) \wedge \\ \quad \operatorname{converges\text{-}to}(\operatorname{cc\text{-}ms}, s, v) \wedge \\ \quad \operatorname{converges\text{-}to}(\operatorname{cc\text{-}ms}, t, a) \Rightarrow \\ \quad \quad \operatorname{converges\text{-}to}(\operatorname{cc\text{-}ms}, h, ((c \cdot v) + (d \cdot a))) \end{array}}}
\vnbentry{\texttt{cc-\allowbreak{}limit-\allowbreak{}shift}}{\vnbstmt{\begin{array}{@{}l@{}} \forall s, h \in (\mathbb{N} \to \mathbb{C}), v \in \mathbb{C}. \\ \quad \forall j \in \mathbb{N}.\; (h(j) = s((j + 1))) \wedge \\ \quad \operatorname{converges\text{-}to}(\operatorname{cc\text{-}ms}, s, v) \Rightarrow \\ \quad \quad \operatorname{converges\text{-}to}(\operatorname{cc\text{-}ms}, h, v) \end{array}}}
\vnbentry{\texttt{cc-\allowbreak{}series-\allowbreak{}partial-\allowbreak{}sum-\allowbreak{}lincomb}}{\vnbstmt{\begin{array}{@{}l@{}} \forall k \in \mathbb{N}, f, g, h \in (\mathbb{N} \to \mathbb{C}), c, d \in \mathbb{C}. \\ \quad \forall j \in \mathbb{N}.\; (h(j) = ((c \cdot f(j)) + (d \cdot g(j)))) \Rightarrow  \\ \quad \quad \operatorname{cc\text{-}series\text{-}partial\text{-}sum}(h, k) \\ \quad \quad =\; ((c \cdot \operatorname{cc\text{-}series\text{-}partial\text{-}sum}(f, k)) + (d \cdot \operatorname{cc\text{-}series\text{-}partial\text{-}sum}(g, k))) \end{array}}}
\vnbentry{\texttt{cc-\allowbreak{}series-\allowbreak{}limit-\allowbreak{}linear}}{\vnbstmt{\begin{array}{@{}l@{}} \forall f, g, h \in (\mathbb{N} \to \mathbb{C}), c, d \in \mathbb{C}. \\ \quad \forall j \in \mathbb{N}.\; (h(j) = ((c \cdot f(j)) + (d \cdot g(j)))) \wedge \\ \quad \operatorname{cc\text{-}series\text{-}converges}(f) \wedge \\ \quad \operatorname{cc\text{-}series\text{-}converges}(g) \Rightarrow \\ \quad \quad \operatorname{cc\text{-}series\text{-}converges}(h) \wedge \\ \quad \quad \operatorname{cc\text{-}series\text{-}limit}(h) \\ \quad \quad =\; ((c \cdot \operatorname{cc\text{-}series\text{-}limit}(f)) + (d \cdot \operatorname{cc\text{-}series\text{-}limit}(g))) \end{array}}}
\vnbentry{\texttt{cpd-\allowbreak{}power-\allowbreak{}second-\allowbreak{}order}}{\vnbstmt{\begin{array}{@{}l@{}} \forall k \in \mathbb{N}, z, u \in \mathbb{C}, r \in \mathbb{R}. \\ \quad ((\operatorname{magnitude}(z) \leq r) \wedge (\operatorname{magnitude}(u) \leq r)) \Rightarrow  \\ \quad \quad (r \cdot \operatorname{magnitude}(((\operatorname{power}(z, (k + 1)) - \operatorname{power}(u, (k + 1))) - (((k + 1) \cdot \operatorname{power}(u, k)) \cdot (z - u))))) \\ \quad \quad \leq\; ((((k + 1) \cdot (k + 1)) \cdot \operatorname{power}(r, k)) \cdot (\operatorname{magnitude}((z - u)) \cdot \operatorname{magnitude}((z - u)))) \end{array}}}
\vnbentry{\texttt{cc-\allowbreak{}power-\allowbreak{}diff-\allowbreak{}bound}}{\vnbstmt{\begin{array}{@{}l@{}} \forall k \in \mathbb{N}, z, u \in \mathbb{C}, r \in \mathbb{R}. \\ \quad ((\operatorname{magnitude}(z) \leq r) \wedge (\operatorname{magnitude}(u) \leq r)) \Rightarrow  \\ \quad \quad \operatorname{magnitude}((\operatorname{power}(z, (k + 1)) - \operatorname{power}(u, (k + 1)))) \\ \quad \quad \leq\; (((k + 1) \cdot \operatorname{power}(r, k)) \cdot \operatorname{magnitude}((z - u))) \end{array}}}
\vnbentry{\texttt{cps-\allowbreak{}shift-\allowbreak{}conv}}{\vnbstmt{\begin{array}{@{}l@{}} \forall f, g \in (\mathbb{N} \to \mathbb{C}), t \in \mathbb{R}. \\ \quad \forall j \in \mathbb{N}. \\ \quad \quad (\operatorname{magnitude}(g(j)) = \operatorname{magnitude}(f((j + 1)))) \wedge \\ \quad (0 < t) \Rightarrow \\ \quad \quad (\operatorname{series\text{-}converges}((\lambda k \in \mathbb{N}.\; (\operatorname{magnitude}(f(k)) \cdot \operatorname{power}(t, k)))) \Leftrightarrow \operatorname{series\text{-}converges}((\lambda k \in \mathbb{N}.\; (\operatorname{magnitude}(g(k)) \cdot \operatorname{power}(t, k))))) \end{array}}}
\vnbentry{\texttt{cps-\allowbreak{}radius-\allowbreak{}in}}{\vnbstmt{\begin{array}{@{}l@{}} \forall f. \\ \quad (\operatorname{cps\text{-}radius}(f) \in \operatorname{rr\text{-}pos\text{-}star}) \end{array}}}
\vnbentry{\texttt{cps-\allowbreak{}radius-\allowbreak{}mono}}{\vnbstmt{\begin{array}{@{}l@{}} \forall f, g \in (\mathbb{N} \to \mathbb{C}). \\ \quad \forall t \in \mathbb{R}. \\ \quad \quad (0 < t) \wedge \\ \quad \quad \operatorname{series\text{-}converges}((\lambda k \in \mathbb{N}.\; (\operatorname{magnitude}(f(k)) \cdot \operatorname{power}(t, k)))) \Rightarrow \\ \quad \quad \quad \operatorname{series\text{-}converges}((\lambda k \in \mathbb{N}.\; (\operatorname{magnitude}(g(k)) \cdot \operatorname{power}(t, k)))) \Rightarrow \\ \quad \quad (\operatorname{cps\text{-}radius}(f) \leq \operatorname{cps\text{-}radius}(g)) \end{array}}}
\vnbentry{\texttt{cps-\allowbreak{}radius-\allowbreak{}derived}}{\vnbstmt{\begin{array}{@{}l@{}} \forall f, a \in (\mathbb{N} \to \mathbb{C}). \\ \quad \forall j \in \mathbb{N}.\; (a(j) = ((j + 1) \cdot f((j + 1)))) \Rightarrow  \\ \quad \quad (\operatorname{cps\text{-}radius}(a) = \operatorname{cps\text{-}radius}(f)) \end{array}}}
\vnbentry{\texttt{cc-\allowbreak{}diff-\allowbreak{}on-\allowbreak{}quadratic}}{\vnbstmt{\begin{array}{@{}l@{}} \forall s.; \forall f \in (s \to \mathbb{C}), a \in s, l \in \mathbb{C}, m \in \mathbb{R}. \\ \quad \operatorname{is\text{-}open}(\operatorname{nf\text{-}metric\text{-}space}(\operatorname{cc\text{-}normed\text{-}field}), s) \wedge \\ \quad (0 \leq m) \wedge \\ \quad \forall x \in s. \\ \quad \quad \operatorname{magnitude}(((f(x) - f(a)) - ((x - a) \cdot l))) \\ \quad \quad \leq\; (m \cdot (\operatorname{magnitude}((x - a)) \cdot \operatorname{magnitude}((x - a)))) \Rightarrow \\ \quad \quad \operatorname{is\text{-}diff\text{-}on}(\operatorname{cc\text{-}normed\text{-}field}, s, f, a, l) \end{array}}}
\vnbentry{\texttt{cps-\allowbreak{}partial-\allowbreak{}error-\allowbreak{}bound}}{\vnbstmt{\begin{array}{@{}l@{}} \forall k \in \mathbb{N}, f, a \in (\mathbb{N} \to \mathbb{C}), z, u \in \mathbb{C}, r \in \mathbb{R}, \mathit{cpdtz}, \mathit{cpdtu}, d \in (\mathbb{N} \to \mathbb{C}), v \in (\mathbb{N} \to \mathbb{R}). \\ \quad \forall j \in \mathbb{N}.\; (a(j) = ((j + 1) \cdot f((j + 1)))) \wedge \\ \quad (\operatorname{magnitude}(z) \leq r) \wedge \\ \quad (\operatorname{magnitude}(u) \leq r) \wedge \\ \quad \forall j \in \mathbb{N}. \\ \quad \quad (\operatorname{cpdtz}(j) = (f(j) \cdot \operatorname{power}(z, j))) \wedge \\ \quad \forall j \in \mathbb{N}. \\ \quad \quad (\operatorname{cpdtu}(j) = (f(j) \cdot \operatorname{power}(u, j))) \wedge \\ \quad \forall j \in \mathbb{N}.\; (d(j) = (a(j) \cdot \operatorname{power}(u, j))) \wedge \\ \quad \forall j \in \mathbb{N}. \\ \quad \quad v(j) \\ \quad \quad =\; (\operatorname{magnitude}((f((j + 1)) \cdot ((j + 1) \cdot (j + 1)))) \cdot \operatorname{power}(r, (j + 1))) \Rightarrow \\ \quad \quad ((r \cdot r) \cdot \operatorname{magnitude}(((\operatorname{cc\text{-}series\text{-}partial\text{-}sum}(\mathit{cpdtz}, (k + 1)) - \operatorname{cc\text{-}series\text{-}partial\text{-}sum}(\mathit{cpdtu}, (k + 1))) - ((z - u) \cdot \operatorname{cc\text{-}series\text{-}partial\text{-}sum}(d, k))))) \\ \quad \quad \leq\; ((\operatorname{magnitude}((z - u)) \cdot \operatorname{magnitude}((z - u))) \cdot \operatorname{series\text{-}partial\text{-}sum}(v, k)) \end{array}}}
\vnbentry{\texttt{cc-\allowbreak{}partial-\allowbreak{}error-\allowbreak{}limit}}{\vnbstmt{\begin{array}{@{}l@{}} \forall f, g, h \in (\mathbb{N} \to \mathbb{C}), w \in \mathbb{C}. \\ \quad \operatorname{cc\text{-}series\text{-}converges}(f) \wedge \\ \quad \operatorname{cc\text{-}series\text{-}converges}(g) \wedge \\ \quad \operatorname{cc\text{-}series\text{-}converges}(h) \Rightarrow \\ \quad \quad \operatorname{converges\text{-}to}(\operatorname{cc\text{-}ms}, \\ \quad \quad \quad (\lambda j \in \mathbb{N}.\; \\ \quad \quad \quad \quad ((\operatorname{cc\text{-}series\text{-}partial\text{-}sum}(f, (j + 1)) - \operatorname{cc\text{-}series\text{-}partial\text{-}sum}(g, (j + 1))) - (w \cdot \operatorname{cc\text{-}series\text{-}partial\text{-}sum}(h, j)))), \\ \quad \quad \quad ((\operatorname{cc\text{-}series\text{-}limit}(f) - \operatorname{cc\text{-}series\text{-}limit}(g)) - (w \cdot \operatorname{cc\text{-}series\text{-}limit}(h)))) \end{array}}}
\vnbentry{\texttt{cps-\allowbreak{}diff-\allowbreak{}on-\allowbreak{}ball}}{\vnbstmt{\begin{array}{@{}l@{}} \forall f, a \in (\mathbb{N} \to \mathbb{C}), r \in \mathbb{R}, \mathit{cpda} \in \operatorname{ball}(\operatorname{cc\text{-}ms}, 0, r). \\ \quad \forall j \in \mathbb{N}.\; (a(j) = ((j + 1) \cdot f((j + 1)))) \wedge \\ \quad (0 < r) \wedge \\ \quad (r < \operatorname{cps\text{-}radius}(f)) \Rightarrow \\ \quad \quad \operatorname{is\text{-}diff\text{-}on}(\operatorname{cc\text{-}normed\text{-}field}, \operatorname{ball}(\operatorname{cc\text{-}ms}, 0, r), \\ \quad \quad \quad (\lambda z \in \operatorname{ball}(\operatorname{cc\text{-}ms}, 0, r).\; \\ \quad \quad \quad \quad \operatorname{cc\text{-}series\text{-}limit}((\lambda k \in \mathbb{N}.\; (f(k) \cdot \operatorname{power}(z, k))))), \\ \quad \quad \quad \mathit{cpda}, \\ \quad \quad \quad \operatorname{cc\text{-}series\text{-}limit}((\lambda k \in \mathbb{N}.\; (a(k) \cdot \operatorname{power}(\mathit{cpda}, k))))) \end{array}}}
\vnbentry{\texttt{cps-\allowbreak{}holomorphic-\allowbreak{}ball}}{\vnbstmt{\begin{array}{@{}l@{}} \forall f \in (\mathbb{N} \to \mathbb{C}), r \in \mathbb{R}. \\ \quad ((0 < r) \wedge (r < \operatorname{cps\text{-}radius}(f))) \Rightarrow  \\ \quad \quad \operatorname{holomorphic\text{-}on}(\operatorname{ball}(\operatorname{cc\text{-}ms}, 0, r), \\ \quad \quad \quad (\lambda z \in \operatorname{ball}(\operatorname{cc\text{-}ms}, 0, r).\; \\ \quad \quad \quad \quad \operatorname{cc\text{-}series\text{-}limit}((\lambda k \in \mathbb{N}.\; (f(k) \cdot \operatorname{power}(z, k)))))) \wedge \\ \quad \quad \operatorname{cps\text{-}radius}((\lambda n \in \mathbb{N}.\; ((n + 1) \cdot f((n + 1))))) \\ \quad \quad =\; \operatorname{cps\text{-}radius}(f) \wedge \\ \quad \quad \forall a \in \operatorname{ball}(\operatorname{cc\text{-}ms}, 0, r). \\ \quad \quad \quad \operatorname{is\text{-}diff\text{-}on}(\operatorname{cc\text{-}normed\text{-}field}, \operatorname{ball}(\operatorname{cc\text{-}ms}, 0, r), \\ \quad \quad \quad \quad (\lambda z \in \operatorname{ball}(\operatorname{cc\text{-}ms}, 0, r).\; \\ \quad \quad \quad \quad \quad \operatorname{cc\text{-}series\text{-}limit}((\lambda k \in \mathbb{N}.\; (f(k) \cdot \operatorname{power}(z, k))))), \\ \quad \quad \quad \quad a, \\ \quad \quad \quad \quad \operatorname{cc\text{-}series\text{-}limit}((\lambda k \in \mathbb{N}.\; (((k + 1) \cdot f((k + 1))) \cdot \operatorname{power}(a, k))))) \end{array}}}

\vnbfile{theorem-library/cps-taylor-coefficients.scm}{theorem-library/cps-taylor-coefficients.scm}
\noindent{\small cps-taylor-coefficients.scm -- COROLLARY 2.14: THE TAYLOR COEFFICIENTS OF A COMPLEX POWER SERIES.  f\textasciicircum{}(k)(0) = k! a\_k. Batch 27-A, 2026-09-24.}\par
\vnbentry{\texttt{cps-\allowbreak{}kth-\allowbreak{}derived-\allowbreak{}step}}{\vnbstmt{\begin{array}{@{}l@{}} \forall f \in (\mathbb{N} \to \mathbb{C}), k \in \mathbb{N}, g, h \in (\mathbb{N} \to \mathbb{C}), j \in \mathbb{N}. \\ \quad \forall n \in \mathbb{N}. \\ \quad \quad g(n) \\ \quad \quad =\; ((\operatorname{factorial}((n + k)) \cdot \operatorname{recip}(\operatorname{factorial}(n))) \cdot f((n + k))) \wedge \\ \quad \forall n \in \mathbb{N}. \\ \quad \quad h(n) \\ \quad \quad =\; ((\operatorname{factorial}((n + (k + 1))) \cdot \operatorname{recip}(\operatorname{factorial}(n))) \cdot f((n + (k + 1)))) \Rightarrow \\ \quad \quad (h(j) = ((j + 1) \cdot g((j + 1)))) \end{array}}}
\vnbentry{\texttt{cps-\allowbreak{}kth-\allowbreak{}derived-\allowbreak{}radius}}{\vnbstmt{\begin{array}{@{}l@{}} \forall f \in (\mathbb{N} \to \mathbb{C}), k \in \mathbb{N}, g \in (\mathbb{N} \to \mathbb{C}). \\ \quad \forall n \in \mathbb{N}. \\ \quad \quad g(n) \\ \quad \quad =\; ((\operatorname{factorial}((n + k)) \cdot \operatorname{recip}(\operatorname{factorial}(n))) \cdot f((n + k))) \Rightarrow \\ \quad \quad (\operatorname{cps\text{-}radius}(g) = \operatorname{cps\text{-}radius}(f)) \end{array}}}
\vnbentry{\texttt{cps-\allowbreak{}sum-\allowbreak{}at-\allowbreak{}zero}}{\vnbstmt{\begin{array}{@{}l@{}} \forall g \in (\mathbb{N} \to \mathbb{C}). \\ \quad \operatorname{cc\text{-}series\text{-}converges}((\lambda k \in \mathbb{N}.\; (g(k) \cdot \operatorname{power}(0, k)))) \wedge \\ \quad \operatorname{cc\text{-}series\text{-}limit}((\lambda k \in \mathbb{N}.\; (g(k) \cdot \operatorname{power}(0, k)))) \\ \quad =\; g(0) \end{array}}}
\vnbentry{\texttt{cps-\allowbreak{}taylor-\allowbreak{}coefficients}}{\vnbstmt{\begin{array}{@{}l@{}} \forall f \in (\mathbb{N} \to \mathbb{C}), k \in \mathbb{N}, g, h \in (\mathbb{N} \to \mathbb{C}), r \in \mathbb{R}. \\ \quad \forall n \in \mathbb{N}. \\ \quad \quad g(n) \\ \quad \quad =\; ((\operatorname{factorial}((n + k)) \cdot \operatorname{recip}(\operatorname{factorial}(n))) \cdot f((n + k))) \wedge \\ \quad \forall n \in \mathbb{N}. \\ \quad \quad h(n) \\ \quad \quad =\; ((\operatorname{factorial}((n + (k + 1))) \cdot \operatorname{recip}(\operatorname{factorial}(n))) \cdot f((n + (k + 1)))) \wedge \\ \quad (0 < r) \wedge \\ \quad (r < \operatorname{cps\text{-}radius}(f)) \Rightarrow \\ \quad \quad \operatorname{holomorphic\text{-}on}(\operatorname{ball}(\operatorname{cc\text{-}ms}, 0, r), \\ \quad \quad \quad (\lambda z \in \operatorname{ball}(\operatorname{cc\text{-}ms}, 0, r).\; \\ \quad \quad \quad \quad \operatorname{cc\text{-}series\text{-}limit}((\lambda \mathit{cpsk} \in \mathbb{N}.\; \\ \quad \quad \quad \quad \quad \quad (g(\mathit{cpsk}) \cdot \operatorname{power}(z, \mathit{cpsk})))))) \wedge \\ \quad \quad (\lambda z \in \operatorname{ball}(\operatorname{cc\text{-}ms}, 0, r).\; \operatorname{cc\text{-}series\text{-}limit}((\lambda \mathit{cpsk} \in \mathbb{N}.\; (g(\mathit{cpsk}) \cdot \operatorname{power}(z, \mathit{cpsk})))))(0) \\ \quad \quad =\; (\operatorname{factorial}(k) \cdot f(k)) \wedge \\ \quad \quad \forall a \in \operatorname{ball}(\operatorname{cc\text{-}ms}, 0, r). \\ \quad \quad \quad \operatorname{deriv\text{-}on}(\operatorname{cc\text{-}normed\text{-}field}, \operatorname{ball}(\operatorname{cc\text{-}ms}, 0, r), \\ \quad \quad \quad \quad (\lambda z \in \operatorname{ball}(\operatorname{cc\text{-}ms}, 0, r).\; \\ \quad \quad \quad \quad \quad \operatorname{cc\text{-}series\text{-}limit}((\lambda \mathit{cpsk} \in \mathbb{N}.\; \\ \quad \quad \quad \quad \quad \quad \quad (g(\mathit{cpsk}) \cdot \operatorname{power}(z, \mathit{cpsk}))))), \\ \quad \quad \quad \quad a) \\ \quad \quad \quad =\; (\lambda z \in \operatorname{ball}(\operatorname{cc\text{-}ms}, 0, r).\; \operatorname{cc\text{-}series\text{-}limit}((\lambda \mathit{cpsk} \in \mathbb{N}.\; (h(\mathit{cpsk}) \cdot \operatorname{power}(z, \mathit{cpsk})))))(a) \end{array}}}
\vnbentry{\texttt{cps-\allowbreak{}coefficient-\allowbreak{}formula}}{\vnbstmt{\begin{array}{@{}l@{}} \forall f \in (\mathbb{N} \to \mathbb{C}), k \in \mathbb{N}, g \in (\mathbb{N} \to \mathbb{C}), r \in \mathbb{R}. \\ \quad \forall n \in \mathbb{N}. \\ \quad \quad g(n) \\ \quad \quad =\; ((\operatorname{factorial}((n + k)) \cdot \operatorname{recip}(\operatorname{factorial}(n))) \cdot f((n + k))) \wedge \\ \quad (0 < r) \Rightarrow \\ \quad \quad f(k) \\ \quad \quad =\; ((\lambda z \in \operatorname{ball}(\operatorname{cc\text{-}ms}, 0, r).\; \operatorname{cc\text{-}series\text{-}limit}((\lambda \mathit{cpsk} \in \mathbb{N}.\; (g(\mathit{cpsk}) \cdot \operatorname{power}(z, \mathit{cpsk})))))(0) \cdot \operatorname{recip}(\operatorname{factorial}(k))) \end{array}}}

\vnbfile{theorem-library/nth-deriv-laws.scm}{theorem-library/nth-deriv-laws.scm}
\noindent{\small nth-deriv-laws.scm -- THE LAWS OF NTH-DERIV-ON, AND COROLLARY 2.14 RESTATED LITERALLY: f\textasciicircum{}(k)(0) = k! a\_k. Batch 36, 2026-09-26.}\par
\vnbentry{\texttt{nth-\allowbreak{}deriv-\allowbreak{}on-\allowbreak{}succ-\allowbreak{}apply}}{\vnbstmt{\begin{array}{@{}l@{}} \forall f, u, \mathit{ndlf}.; \forall j \in \mathbb{N}, z \in u. \\ \quad (\mathit{ndlf}^{((j + 1))}(z) \simeq \operatorname{deriv\text{-}on}(f, u, \mathit{ndlf}^{(j)}, z)) \end{array}}}
\vnbentry{\texttt{nth-\allowbreak{}deriv-\allowbreak{}on-\allowbreak{}of-\allowbreak{}is-\allowbreak{}diff-\allowbreak{}on}}{\vnbstmt{\begin{array}{@{}l@{}} \forall f, u, \mathit{ndlf}.; \forall j \in \mathbb{N}, g.; \forall w \in u. \\ \quad \operatorname{is\text{-}normed\text{-}field}(f) \wedge \\ \quad \forall e. \\ \quad \quad \operatorname{pos\text{-}rr}(e) \Rightarrow  \\ \quad \quad \quad \exists h \in \operatorname{carr}(f). \\ \quad \quad \quad \quad (\neg (h = \operatorname{zero}(f)) \wedge (\operatorname{fnrm}(f)(h) < e)) \wedge \\ \quad \forall z \in u. \\ \quad \quad \operatorname{is\text{-}diff\text{-}on}(f, u, \mathit{ndlf}^{(j)}, z, g(z)) \Rightarrow \\ \quad \quad (\mathit{ndlf}^{((j + 1))}(w) = g(w)) \end{array}}}
\vnbentry{\texttt{nth-\allowbreak{}deriv-\allowbreak{}on-\allowbreak{}succ-\allowbreak{}eq}}{\vnbstmt{\begin{array}{@{}l@{}} \forall f, u, \mathit{ndlf}.; \forall j \in \mathbb{N}, g \in (u \to \operatorname{carr}(f)). \\ \quad \operatorname{is\text{-}normed\text{-}field}(f) \wedge \\ \quad \forall e. \\ \quad \quad \operatorname{pos\text{-}rr}(e) \Rightarrow  \\ \quad \quad \quad \exists h \in \operatorname{carr}(f). \\ \quad \quad \quad \quad (\neg (h = \operatorname{zero}(f)) \wedge (\operatorname{fnrm}(f)(h) < e)) \wedge \\ \quad \forall z \in u. \\ \quad \quad \operatorname{is\text{-}diff\text{-}on}(f, u, \mathit{ndlf}^{(j)}, z, g(z)) \Rightarrow \\ \quad \quad (\mathit{ndlf}^{((j + 1))} = g) \end{array}}}
\vnbentry{\texttt{nth-\allowbreak{}deriv-\allowbreak{}on-\allowbreak{}in-\allowbreak{}fun}}{\vnbstmt{\begin{array}{@{}l@{}} \forall f, u.; \forall \mathit{ndlf} \in (u \to \operatorname{carr}(f)), j \in \mathbb{N}. \\ \quad \operatorname{is\text{-}normed\text{-}field}(f) \wedge \\ \quad \forall e. \\ \quad \quad \operatorname{pos\text{-}rr}(e) \Rightarrow  \\ \quad \quad \quad \exists h \in \operatorname{carr}(f). \\ \quad \quad \quad \quad (\neg (h = \operatorname{zero}(f)) \wedge (\operatorname{fnrm}(f)(h) < e)) \wedge \\ \quad \forall i \in \mathbb{N}, z \in u. \\ \quad \quad (i < j) \Rightarrow  \\ \quad \quad \quad \exists l. \\ \quad \quad \quad \quad \operatorname{is\text{-}diff\text{-}on}(f, u, \mathit{ndlf}^{(i)}, z, l) \Rightarrow \\ \quad \quad (\mathit{ndlf}^{(j)} \in (u \to \operatorname{carr}(f))) \end{array}}}
\vnbentry{\texttt{cps-\allowbreak{}nth-\allowbreak{}deriv}}{\vnbstmt{\begin{array}{@{}l@{}} \forall f \in (\mathbb{N} \to \mathbb{C}), r \in \mathbb{R}, k \in \mathbb{N}, g \in (\mathbb{N} \to \mathbb{C}). \\ \quad (0 < r) \wedge \\ \quad (r < \operatorname{cps\text{-}radius}(f)) \wedge \\ \quad \forall n \in \mathbb{N}. \\ \quad \quad g(n) \\ \quad \quad =\; ((\operatorname{factorial}((n + k)) \cdot \operatorname{recip}(\operatorname{factorial}(n))) \cdot f((n + k))) \Rightarrow \\ \quad \quad (\lambda z \in \operatorname{ball}(\operatorname{cc\text{-}ms}, 0, r).\; \operatorname{cc\text{-}series\text{-}limit}((\lambda \mathit{cpsk} \in \mathbb{N}.\; (f(\mathit{cpsk}) \cdot \operatorname{power}(z, \mathit{cpsk})))))^{(k)} \\ \quad \quad =\; (\lambda z \in \operatorname{ball}(\operatorname{cc\text{-}ms}, 0, r).\; \\ \quad \quad \quad \operatorname{cc\text{-}series\text{-}limit}((\lambda \mathit{cpsk} \in \mathbb{N}.\; \\ \quad \quad \quad \quad \quad (g(\mathit{cpsk}) \cdot \operatorname{power}(z, \mathit{cpsk}))))) \end{array}}}
\vnbentry{\texttt{taylor-\allowbreak{}coefficients}}{\vnbstmt{\begin{array}{@{}l@{}} \forall f \in (\mathbb{N} \to \mathbb{C}), r \in \mathbb{R}, k \in \mathbb{N}. \\ \quad ((0 < r) \wedge (r < \operatorname{cps\text{-}radius}(f))) \Rightarrow  \\ \quad \quad (\lambda z \in \operatorname{ball}(\operatorname{cc\text{-}ms}, 0, r).\; \operatorname{cc\text{-}series\text{-}limit}((\lambda \mathit{cpsk} \in \mathbb{N}.\; (f(\mathit{cpsk}) \cdot \operatorname{power}(z, \mathit{cpsk})))))^{(k)}(0) \\ \quad \quad =\; (\operatorname{factorial}(k) \cdot f(k)) \end{array}}}
\vnbentry{\texttt{taylor-\allowbreak{}coefficient-\allowbreak{}formula}}{\vnbstmt{\begin{array}{@{}l@{}} \forall f \in (\mathbb{N} \to \mathbb{C}), r \in \mathbb{R}, k \in \mathbb{N}. \\ \quad ((0 < r) \wedge (r < \operatorname{cps\text{-}radius}(f))) \Rightarrow  \\ \quad \quad f(k) \\ \quad \quad =\; ((\lambda z \in \operatorname{ball}(\operatorname{cc\text{-}ms}, 0, r).\; \operatorname{cc\text{-}series\text{-}limit}((\lambda \mathit{cpsk} \in \mathbb{N}.\; (f(\mathit{cpsk}) \cdot \operatorname{power}(z, \mathit{cpsk})))))^{(k)}(0) \cdot \operatorname{recip}(\operatorname{factorial}(k))) \end{array}}}
\vnbentry{\texttt{nth-\allowbreak{}deriv-\allowbreak{}on-\allowbreak{}holomorphic-\allowbreak{}ball}}{\vnbstmt{\begin{array}{@{}l@{}} \forall f \in (\mathbb{N} \to \mathbb{C}), r \in \mathbb{R}, k \in \mathbb{N}. \\ \quad ((0 < r) \wedge (r < \operatorname{cps\text{-}radius}(f))) \Rightarrow  \\ \quad \quad \operatorname{holomorphic\text{-}on}(\operatorname{ball}(\operatorname{cc\text{-}ms}, 0, r), \\ \quad \quad \quad (\lambda z \in \operatorname{ball}(\operatorname{cc\text{-}ms}, 0, r).\; \operatorname{cc\text{-}series\text{-}limit}((\lambda \mathit{cpsk} \in \mathbb{N}.\; (f(\mathit{cpsk}) \cdot \operatorname{power}(z, \mathit{cpsk})))))^{(k)}) \end{array}}}

\vnbfile{theorem-library/regulated-has-primitive.scm}{theorem-library/regulated-has-primitive.scm}
\noindent{\small regulated-has-primitive.scm -- EVERY REGULATED FUNCTION HAS A PRIMITIVE on [a, b] (batch 22, stage 2: the assembly).}\par
\vnbentry{\texttt{primitive-\allowbreak{}shift-\allowbreak{}const}}{\vnbstmt{\begin{array}{@{}l@{}} \forall g, f, a, b.; \forall c \in \mathbb{R}. \\ \quad \operatorname{is\text{-}primitive}(g, f, a, b) \Rightarrow  \\ \quad \quad \operatorname{is\text{-}primitive}((\lambda y \in \operatorname{ccint}(a, b).\; (g(y) - c)), f, a, b) \end{array}}}
\vnbentry{\texttt{primitive-\allowbreak{}family-\allowbreak{}normalised-\allowbreak{}choice}}{\vnbstmt{\begin{array}{@{}l@{}} \forall a, b \in \mathbb{R}, f. \\ \quad \forall k \in \mathbb{N}.\; \exists g.\; \operatorname{is\text{-}primitive}(g, f(k), a, b) \Rightarrow  \\ \quad \quad \exists \mathit{rhpgf} \in (\mathbb{N} \to (\operatorname{ccint}(a, b) \to \mathbb{R})). \forall k \in \mathbb{N}. \\ \quad \quad \quad \operatorname{is\text{-}primitive}(\operatorname{rhpgf}(k), f(k), a, b) \wedge \\ \quad \quad \quad (\operatorname{rhpgf}(k)(a) = 0) \end{array}}}
\vnbentry{\texttt{regulated-\allowbreak{}on-\allowbreak{}has-\allowbreak{}primitive}}{\vnbstmt{\begin{array}{@{}l@{}} \forall a, b \in \mathbb{R}, f. \\ \quad \operatorname{is\text{-}regulated\text{-}on}(f, a, b) \Rightarrow  \\ \quad \quad \exists g.\; \operatorname{is\text{-}primitive}(g, f, a, b) \end{array}}}

\vnbfile{theorem-library/road-laws-2.scm}{theorem-library/road-laws-2.scm}
\noindent{\small road-laws-2.scm -- THE OPPOSITE ROAD AND DIEUDONNE (9.6.1).}\par
\vnbentry{\texttt{ccint-\allowbreak{}reflect-\allowbreak{}in}}{\vnbstmt{\begin{array}{@{}l@{}} \forall a, b \in \mathbb{R}, y \in \operatorname{ccint}(a, b). \\ \quad (((a + b) - y) \in \operatorname{ccint}(a, b)) \end{array}}}
\vnbentry{\texttt{ooint-\allowbreak{}reflect-\allowbreak{}in}}{\vnbstmt{\begin{array}{@{}l@{}} \forall a, b \in \mathbb{R}, y \in \operatorname{ooint}(a, b). \\ \quad (((a + b) - y) \in \operatorname{ooint}(a, b)) \end{array}}}
\vnbentry{\texttt{reflect-\allowbreak{}continuous-\allowbreak{}at-\allowbreak{}sub}}{\vnbstmt{\begin{array}{@{}l@{}} \forall a, b \in \mathbb{R}, f.; \forall c \in \mathbb{R}, q \in (\operatorname{ccint}(a, b) \to \mathbb{R}), t \in \operatorname{ccint}(a, b). \\ \quad \forall y \in \operatorname{ccint}(a, b). \\ \quad \quad (q(y) \simeq (c \cdot f(((a + b) - y)))) \wedge \\ \quad \operatorname{is\text{-}continuous\text{-}at}(\operatorname{subspace\text{-}ms}(\operatorname{rr\text{-}ms}, \operatorname{ccint}(a, b)), \operatorname{rr\text{-}ms}, f, ((a + b) - t)) \Rightarrow \\ \quad \quad \operatorname{is\text{-}continuous\text{-}at}(\operatorname{subspace\text{-}ms}(\operatorname{rr\text{-}ms}, \operatorname{ccint}(a, b)), \operatorname{rr\text{-}ms}, q, t) \end{array}}}
\vnbentry{\texttt{reflect-\allowbreak{}continuous-\allowbreak{}on}}{\vnbstmt{\begin{array}{@{}l@{}} \forall a, b \in \mathbb{R}, f.; \forall c \in \mathbb{R}, q \in (\operatorname{ccint}(a, b) \to \mathbb{R}). \\ \quad \operatorname{is\text{-}continuous\text{-}on}(f, \operatorname{ccint}(a, b)) \wedge \\ \quad \forall y \in \operatorname{ccint}(a, b). \\ \quad \quad (q(y) \simeq (c \cdot f(((a + b) - y)))) \Rightarrow \\ \quad \quad \operatorname{is\text{-}continuous\text{-}on}(q, \operatorname{ccint}(a, b)) \end{array}}}
\vnbentry{\texttt{regulated-\allowbreak{}on-\allowbreak{}reflect}}{\vnbstmt{\begin{array}{@{}l@{}} \forall a, b \in \mathbb{R}, f.; \forall c \in \mathbb{R}, q \in (\operatorname{ccint}(a, b) \to \mathbb{R}). \\ \quad \operatorname{is\text{-}regulated\text{-}on}(f, a, b) \wedge \\ \quad \forall y \in \operatorname{ccint}(a, b). \\ \quad \quad (q(y) \simeq (c \cdot f(((a + b) - y)))) \Rightarrow \\ \quad \quad \operatorname{is\text{-}regulated\text{-}on}(q, a, b) \end{array}}}
\vnbentry{\texttt{reflect-\allowbreak{}image-\allowbreak{}finite}}{\vnbstmt{\begin{array}{@{}l@{}} \forall a, b.; \forall d \in \mathbf{Set}. \\ \quad (a \in \mathbb{R}) \wedge \\ \quad (b \in \mathbb{R}) \wedge \\ \quad (|d| \in \mathbb{N}) \wedge \\ \quad (d \subseteq \operatorname{ccint}(a, b)) \Rightarrow \\ \quad \quad \operatorname{image}((\lambda x \in \operatorname{ccint}(a, b).\; ((a + b) - x)), d) \\ \quad \quad \in\; \mathbf{Set} \wedge \\ \quad \quad |\operatorname{image}((\lambda x \in \operatorname{ccint}(a, b).\; ((a + b) - x)), d)| \\ \quad \quad \in\; \mathbb{N} \wedge \\ \quad \quad \operatorname{image}((\lambda x \in \operatorname{ccint}(a, b).\; ((a + b) - x)), d) \\ \quad \quad \subseteq\; \operatorname{ccint}(a, b) \end{array}}}
\vnbentry{\texttt{reflect-\allowbreak{}image-\allowbreak{}countable}}{\vnbstmt{\begin{array}{@{}l@{}} \forall a, b, d. \\ \quad (a \in \mathbb{R}) \wedge \\ \quad (b \in \mathbb{R}) \wedge \\ \quad \operatorname{is\text{-}countable}(d) \wedge \\ \quad (d \subseteq \operatorname{ccint}(a, b)) \Rightarrow \\ \quad \quad \operatorname{is\text{-}countable}(\operatorname{image}((\lambda x \in \operatorname{ccint}(a, b).\; ((a + b) - x)), d)) \wedge \\ \quad \quad \operatorname{image}((\lambda x \in \operatorname{ccint}(a, b).\; ((a + b) - x)), d) \\ \quad \quad \subseteq\; \operatorname{ccint}(a, b) \end{array}}}
\vnbentry{\texttt{reflect-\allowbreak{}not-\allowbreak{}in-\allowbreak{}image}}{\vnbstmt{\begin{array}{@{}l@{}} \forall a, b, d.; \forall t \in \operatorname{ccint}(a, b). \\ \quad (a \in \mathbb{R}) \wedge \\ \quad (b \in \mathbb{R}) \wedge \\ \quad (d \subseteq \operatorname{ccint}(a, b)) \wedge \\ \quad \neg (t \in \operatorname{image}((\lambda x \in \operatorname{ccint}(a, b).\; ((a + b) - x)), d)) \Rightarrow \\ \quad \quad \neg (((a + b) - t) \in d) \end{array}}}
\vnbentry{\texttt{pw-\allowbreak{}antiderivative-\allowbreak{}reflect}}{\vnbstmt{\begin{array}{@{}l@{}} \forall a, b \in \mathbb{R}, f, i.; \forall g, k \in (\operatorname{ccint}(a, b) \to \mathbb{R}). \\ \quad \operatorname{is\text{-}primitive}(f, i, a, b) \wedge \\ \quad \forall y \in \operatorname{ccint}(a, b).\; (g(y) \simeq f(((a + b) - y))) \wedge \\ \quad \forall y \in \operatorname{ccint}(a, b). \\ \quad \quad (k(y) \simeq (-1 \cdot i(((a + b) - y)))) \Rightarrow \\ \quad \quad \operatorname{is\text{-}primitive}(g, k, a, b) \end{array}}}
\vnbentry{\texttt{pw-\allowbreak{}int-\allowbreak{}reflect}}{\vnbstmt{\begin{array}{@{}l@{}} \forall a, b \in \mathbb{R}, f, i.; \forall h, s \in (\operatorname{ccint}(a, b) \to \mathbb{R}). \\ \quad \operatorname{is\text{-}primitive}(f, i, a, b) \wedge \\ \quad \forall y \in \operatorname{ccint}(a, b). \\ \quad \quad (h(y) \simeq (-1 \cdot f(((a + b) - y)))) \wedge \\ \quad \forall y \in \operatorname{ccint}(a, b).\; (s(y) \simeq i(((a + b) - y))) \Rightarrow \\ \quad \quad (\operatorname{is\text{-}primitive}(h, s, a, b) \wedge (\int_{a}^{b} s = \int_{a}^{b} i)) \end{array}}}
\vnbentry{\texttt{pw-\allowbreak{}continuous-\allowbreak{}reflect}}{\vnbstmt{\begin{array}{@{}l@{}} \forall a, b \in \mathbb{R}, i.; \forall c \in \mathbb{R}, q \in (\operatorname{ccint}(a, b) \to \mathbb{R}). \\ \quad \operatorname{is\text{-}pw\text{-}continuous\text{-}on}(i, a, b) \wedge \\ \quad \forall y \in \operatorname{ccint}(a, b). \\ \quad \quad (q(y) \simeq (c \cdot i(((a + b) - y)))) \Rightarrow \\ \quad \quad \operatorname{is\text{-}pw\text{-}continuous\text{-}on}(q, a, b) \end{array}}}
\vnbentry{\texttt{trace-\allowbreak{}opposite-\allowbreak{}subset}}{\vnbstmt{\begin{array}{@{}l@{}} \forall a, b \in \mathbb{R}, m, \mathit{prgb}. \\ \quad (m \in (\operatorname{ccint}(a, b) \to \mathbb{C})) \wedge \\ \quad (\mathit{prgb} \in (\operatorname{ccint}(a, b) \to \mathbb{C})) \wedge \\ \quad \forall y \in \operatorname{ccint}(a, b). \\ \quad \quad (\operatorname{prgb}(y) \simeq m(((a + b) - y))) \Rightarrow \\ \quad \quad (\operatorname{trace}(\mathit{prgb}, a, b) \subseteq \operatorname{trace}(m, a, b)) \end{array}}}
\vnbentry{\texttt{trace-\allowbreak{}opposite}}{\vnbstmt{\begin{array}{@{}l@{}} \forall a, b \in \mathbb{R}, m, \mathit{prgb}. \\ \quad (m \in (\operatorname{ccint}(a, b) \to \mathbb{C})) \wedge \\ \quad (\mathit{prgb} \in (\operatorname{ccint}(a, b) \to \mathbb{C})) \wedge \\ \quad \forall y \in \operatorname{ccint}(a, b). \\ \quad \quad (\operatorname{prgb}(y) \simeq m(((a + b) - y))) \Rightarrow \\ \quad \quad (\operatorname{trace}(\mathit{prgb}, a, b) \subseteq \operatorname{trace}(m, a, b)) \wedge \\ \quad \quad (\operatorname{trace}(m, a, b) \subseteq \operatorname{trace}(\mathit{prgb}, a, b)) \end{array}}}
\vnbentry{\texttt{opposite-\allowbreak{}is-\allowbreak{}road}}{\vnbstmt{\begin{array}{@{}l@{}} \forall a, b \in \mathbb{R}, m, c, \mathit{prgb}, \mathit{prdb}. \\ \quad \operatorname{is\text{-}road}(m, c, a, b) \wedge \\ \quad (\mathit{prgb} \in (\operatorname{ccint}(a, b) \to \mathbb{C})) \wedge \\ \quad (\mathit{prdb} \in (\operatorname{ccint}(a, b) \to \mathbb{C})) \wedge \\ \quad \forall y \in \operatorname{ccint}(a, b). \\ \quad \quad (\operatorname{prgb}(y) \simeq m(((a + b) - y))) \wedge \\ \quad \forall y \in \operatorname{ccint}(a, b). \\ \quad \quad (\operatorname{prdb}(y) \simeq (-1 \cdot c(((a + b) - y)))) \Rightarrow \\ \quad \quad \operatorname{is\text{-}road}(\mathit{prgb}, \mathit{prdb}, a, b) \end{array}}}
\vnbentry{\texttt{cc-\allowbreak{}int-\allowbreak{}reflect}}{\vnbstmt{\begin{array}{@{}l@{}} \forall a, b, i, \mathit{prb}, f, w. \\ \quad (a \in \mathbb{R}) \wedge \\ \quad (b \in \mathbb{R}) \wedge \\ \quad (i \in (\operatorname{ccint}(a, b) \to \mathbb{C})) \wedge \\ \quad (\mathit{prb} \in (\operatorname{ccint}(a, b) \to \mathbb{C})) \wedge \\ \quad \forall y \in \operatorname{ccint}(a, b). \\ \quad \quad (\operatorname{prb}(y) \simeq i(((a + b) - y))) \wedge \\ \quad \operatorname{is\text{-}primitive}(f, (\lambda t \in \operatorname{ccint}(a, b).\; \operatorname{real\text{-}part}(i(t))), a, b) \wedge \\ \quad \operatorname{is\text{-}primitive}(w, (\lambda t \in \operatorname{ccint}(a, b).\; \operatorname{imag\text{-}part}(i(t))), a, b) \Rightarrow \\ \quad \quad (\int_{a}^{b} \mathit{prb} = \int_{a}^{b} i) \end{array}}}
\vnbentry{\texttt{line-\allowbreak{}int-\allowbreak{}opposite}}{\vnbstmt{\begin{array}{@{}l@{}} \forall a, b, m, c, f, d, \mathit{prgb}, \mathit{prdb}, \mathit{pwf}, w. \\ \quad (a \in \mathbb{R}) \wedge \\ \quad (b \in \mathbb{R}) \wedge \\ \quad \operatorname{is\text{-}road}(m, c, a, b) \wedge \\ \quad (f \in (d \to \mathbb{C})) \wedge \\ \quad (\operatorname{trace}(m, a, b) \subseteq d) \wedge \\ \quad (\mathit{prgb} \in (\operatorname{ccint}(a, b) \to \mathbb{C})) \wedge \\ \quad (\mathit{prdb} \in (\operatorname{ccint}(a, b) \to \mathbb{C})) \wedge \\ \quad \forall y \in \operatorname{ccint}(a, b). \\ \quad \quad (\operatorname{prgb}(y) \simeq m(((a + b) - y))) \wedge \\ \quad \forall y \in \operatorname{ccint}(a, b). \\ \quad \quad (\operatorname{prdb}(y) \simeq (-1 \cdot c(((a + b) - y)))) \wedge \\ \quad \operatorname{is\text{-}primitive}(\mathit{pwf}, \\ \quad \quad (\lambda t \in \operatorname{ccint}(a, b).\; \\ \quad \quad \quad \operatorname{real\text{-}part}((\lambda t \in \operatorname{ccint}(a, b).\; (f(m(t)) \cdot c(t)))(t))), \\ \quad \quad a, b) \wedge \\ \quad \operatorname{is\text{-}primitive}(w, \\ \quad \quad (\lambda t \in \operatorname{ccint}(a, b).\; \\ \quad \quad \quad \operatorname{imag\text{-}part}((\lambda t \in \operatorname{ccint}(a, b).\; (f(m(t)) \cdot c(t)))(t))), \\ \quad \quad a, b) \Rightarrow \\ \quad \quad \operatorname{line\text{-}int}(f, \mathit{prgb}, \mathit{prdb}, a, b) \\ \quad \quad =\; (-1 \cdot \operatorname{line\text{-}int}(f, m, c, a, b)) \end{array}}}

\vnbfile{theorem-library/line-int-fundamental.scm}{theorem-library/line-int-fundamental.scm}
\noindent{\small line-int-fundamental.scm -- THE FUNDAMENTAL THEOREM FOR INTEGRALS ALONG A ROAD (the user's notes, \textasciitilde{}/docs/complex-analysis.pdf, section 3.1, Proposition 3.2 and equation (50)), and the chain rule it rests on.}\par
\vnbentry{\texttt{cc-\allowbreak{}compose-\allowbreak{}continuous-\allowbreak{}at}}{\vnbstmt{\begin{array}{@{}l@{}} \forall s, u, g, i, x. \\ \quad \operatorname{is\text{-}metric\text{-}space}(s) \wedge \\ \quad (u \subseteq \mathbb{C}) \wedge \\ \quad \operatorname{is\text{-}continuous\text{-}at}(s, \operatorname{cc\text{-}ms}, g, x) \wedge \\ \quad \forall y \in \operatorname{pts}(s).\; (g(y) \in u) \wedge \\ \quad \operatorname{is\text{-}continuous\text{-}at}(\operatorname{subspace\text{-}ms}(\operatorname{nf\text{-}metric\text{-}space}(\operatorname{cc\text{-}normed\text{-}field}), u), \\ \quad \quad \operatorname{nf\text{-}metric\text{-}space}(\operatorname{cc\text{-}normed\text{-}field}), i, g(x)) \Rightarrow \\ \quad \quad \operatorname{is\text{-}continuous\text{-}at}(s, \operatorname{cc\text{-}ms}, (\lambda z \in \operatorname{pts}(s).\; i(g(z))), x) \end{array}}}
\vnbentry{\texttt{cc-\allowbreak{}product-\allowbreak{}continuous-\allowbreak{}at}}{\vnbstmt{\begin{array}{@{}l@{}} \forall s, p, q, x. \\ \quad \operatorname{is\text{-}continuous\text{-}at}(s, \operatorname{cc\text{-}ms}, p, x) \wedge \\ \quad \operatorname{is\text{-}continuous\text{-}at}(s, \operatorname{cc\text{-}ms}, q, x) \Rightarrow \\ \quad \quad \operatorname{is\text{-}continuous\text{-}at}(s, \operatorname{cc\text{-}ms}, (\lambda z \in \operatorname{pts}(s).\; (p(z) \cdot q(z))), x) \end{array}}}
\vnbentry{\texttt{lfn-\allowbreak{}cc-\allowbreak{}factor-\allowbreak{}identity}}{\vnbstmt{\begin{array}{@{}l@{}} \forall a, b, c, e, h \in \mathbb{C}. \\ \quad ((a \cdot ((b \cdot h) + (e \cdot (c \cdot h)))) = (h \cdot (a \cdot (b + (c \cdot e))))) \end{array}}}
\vnbentry{\texttt{holomorphic-\allowbreak{}chain-\allowbreak{}within}}{\vnbstmt{\begin{array}{@{}l@{}} \forall m, w, g, x, u, f, l, d. \\ \quad (m \subseteq \mathbb{R}) \wedge \\ \quad (w \subseteq m) \wedge \\ \quad (g \in (m \to \mathbb{C})) \wedge \\ \quad \operatorname{is\text{-}continuous\text{-}at}(\operatorname{subspace\text{-}ms}(\operatorname{rr\text{-}ms}, m), \operatorname{cc\text{-}ms}, g, x) \wedge \\ \quad \forall y \in m.\; (g(y) \in u) \wedge \\ \quad \operatorname{is\text{-}diff\text{-}on}(\operatorname{cc\text{-}normed\text{-}field}, u, f, g(x), l) \wedge \\ \quad (d \in \mathbb{C}) \wedge \\ \quad \operatorname{has\text{-}deriv\text{-}at\text{-}within}((\lambda t \in m.\; \operatorname{real\text{-}part}(g(t))), w, x, \operatorname{real\text{-}part}(d)) \wedge \\ \quad \operatorname{has\text{-}deriv\text{-}at\text{-}within}((\lambda t \in m.\; \operatorname{imag\text{-}part}(g(t))), w, x, \operatorname{imag\text{-}part}(d)) \Rightarrow \\ \quad \quad \operatorname{has\text{-}deriv\text{-}at\text{-}within}((\lambda t \in m.\; \operatorname{real\text{-}part}(f(g(t)))), w, x, \operatorname{real\text{-}part}((l \cdot d))) \wedge \\ \quad \quad \operatorname{has\text{-}deriv\text{-}at\text{-}within}((\lambda t \in m.\; \operatorname{imag\text{-}part}(f(g(t)))), w, x, \operatorname{imag\text{-}part}((l \cdot d))) \end{array}}}
\vnbentry{\texttt{ooint-\allowbreak{}center-\allowbreak{}nested}}{\vnbstmt{\begin{array}{@{}l@{}} \forall x \in \mathbb{R}, v.; \forall r \in \mathbb{R}. \\ \quad (\operatorname{pos\text{-}rr}(v) \wedge (v \leq r)) \Rightarrow  \\ \quad \quad (\operatorname{ooint}((x - v), (x + v)) \subseteq \operatorname{ooint}((x - r), (x + r))) \end{array}}}
\vnbentry{\texttt{holomorphic-\allowbreak{}chain}}{\vnbstmt{\begin{array}{@{}l@{}} \forall m, g, x, u, f, l, d. \\ \quad (m \subseteq \mathbb{R}) \wedge \\ \quad (g \in (m \to \mathbb{C})) \wedge \\ \quad \operatorname{is\text{-}continuous\text{-}at}(\operatorname{subspace\text{-}ms}(\operatorname{rr\text{-}ms}, m), \operatorname{cc\text{-}ms}, g, x) \wedge \\ \quad \forall y \in m.\; (g(y) \in u) \wedge \\ \quad \operatorname{is\text{-}diff\text{-}on}(\operatorname{cc\text{-}normed\text{-}field}, u, f, g(x), l) \wedge \\ \quad (d \in \mathbb{C}) \wedge \\ \quad \operatorname{has\text{-}deriv\text{-}at}((\lambda t \in m.\; \operatorname{real\text{-}part}(g(t))), x, \operatorname{real\text{-}part}(d)) \wedge \\ \quad \operatorname{has\text{-}deriv\text{-}at}((\lambda t \in m.\; \operatorname{imag\text{-}part}(g(t))), x, \operatorname{imag\text{-}part}(d)) \Rightarrow \\ \quad \quad \operatorname{has\text{-}deriv\text{-}at}((\lambda t \in m.\; \operatorname{real\text{-}part}(f(g(t)))), x, \operatorname{real\text{-}part}((l \cdot d))) \wedge \\ \quad \quad \operatorname{has\text{-}deriv\text{-}at}((\lambda t \in m.\; \operatorname{imag\text{-}part}(f(g(t)))), x, \operatorname{imag\text{-}part}((l \cdot d))) \end{array}}}
\vnbentry{\texttt{line-\allowbreak{}int-\allowbreak{}of-\allowbreak{}derivative}}{\vnbstmt{\begin{array}{@{}l@{}} \forall u, f, c, m, d, a, b. \\ \quad \operatorname{holomorphic\text{-}on}(u, f) \wedge \\ \quad (c \in (u \to \mathbb{C})) \wedge \\ \quad \forall z \in u. \\ \quad \quad \operatorname{is\text{-}diff\text{-}on}(\operatorname{cc\text{-}normed\text{-}field}, u, f, z, c(z)) \wedge \\ \quad \operatorname{is\text{-}road}(m, d, a, b) \wedge \\ \quad (\operatorname{trace}(m, a, b) \subseteq u) \Rightarrow \\ \quad \quad (\operatorname{line\text{-}int}(c, m, d, a, b) = (f(m(b)) - f(m(a)))) \end{array}}}

\vnbfile{theorem-library/goursat-bricks.scm}{theorem-library/goursat-bricks.scm}
\noindent{\small goursat-bricks.scm -- the definition-free bricks of Cauchy-Goursat (complex-analysis.pdf ch. 3, s. 3.2).  Batch 29-B, 2026-09-25.}\par
\vnbentry{\texttt{quarter-\allowbreak{}step}}{\vnbstmt{\begin{array}{@{}l@{}} \forall j, a, b, c, d \in \mathbb{R}. \\ \quad (0 \leq a) \wedge \\ \quad (0 \leq b) \wedge \\ \quad (0 \leq c) \wedge \\ \quad (0 \leq d) \wedge \\ \quad (j \leq (((a + b) + c) + d)) \Rightarrow \\ \quad \quad (j \leq (4 \cdot a)) \vee \\ \quad \quad (j \leq (4 \cdot b)) \vee \\ \quad \quad (j \leq (4 \cdot c)) \vee \\ \quad \quad (j \leq (4 \cdot d)) \end{array}}}
\vnbentry{\texttt{four-\allowbreak{}power-\allowbreak{}halving}}{\vnbstmt{\begin{array}{@{}l@{}} \forall n \in \mathbb{N}. \\ \quad (\operatorname{power}(4, n) \cdot (\operatorname{recip}(\operatorname{power}(2, n)) \cdot \operatorname{recip}(\operatorname{power}(2, n)))) \\ \quad =\; 1 \end{array}}}
\vnbentry{\texttt{line-\allowbreak{}int-\allowbreak{}closed-\allowbreak{}road-\allowbreak{}of-\allowbreak{}derivative}}{\vnbstmt{\begin{array}{@{}l@{}} \forall u, f, c, m, d, a, b. \\ \quad \operatorname{holomorphic\text{-}on}(u, f) \wedge \\ \quad (c \in (u \to \mathbb{C})) \wedge \\ \quad \forall z \in u. \\ \quad \quad \operatorname{is\text{-}diff\text{-}on}(\operatorname{cc\text{-}normed\text{-}field}, u, f, z, c(z)) \wedge \\ \quad \operatorname{is\text{-}road}(m, d, a, b) \wedge \\ \quad (\operatorname{trace}(m, a, b) \subseteq u) \wedge \\ \quad (m(b) = m(a)) \Rightarrow \\ \quad \quad (\operatorname{line\text{-}int}(c, m, d, a, b) = 0) \end{array}}}
\vnbentry{\texttt{affine-\allowbreak{}has-\allowbreak{}primitive-\allowbreak{}cc}}{\vnbstmt{\begin{array}{@{}l@{}} \forall p, q, w \in \mathbb{C}, f \in (\mathbb{C} \to \mathbb{C}). \\ \quad \forall z \in \mathbb{C}. \\ \quad \quad (f(z) \simeq ((p \cdot z) + ((q \cdot ((z - w) \cdot (z - w))) \cdot \operatorname{recip}(2)))) \Rightarrow \\ \quad \quad \operatorname{holomorphic\text{-}on}(\mathbb{C}, f) \wedge \\ \quad \quad \forall z \in \mathbb{C}. \\ \quad \quad \quad \operatorname{is\text{-}diff\text{-}on}(\operatorname{cc\text{-}normed\text{-}field}, \mathbb{C}, f, z, (p + (q \cdot (z - w)))) \end{array}}}
\vnbentry{\texttt{affine-\allowbreak{}closed-\allowbreak{}road-\allowbreak{}integral}}{\vnbstmt{\begin{array}{@{}l@{}} \forall p, q, w \in \mathbb{C}, g \in (\mathbb{C} \to \mathbb{C}), m, c, a, b. \\ \quad \forall z \in \mathbb{C}.\; (g(z) \simeq (p + (q \cdot (z - w)))) \wedge \\ \quad \operatorname{is\text{-}road}(m, c, a, b) \wedge \\ \quad (m(b) = m(a)) \Rightarrow \\ \quad \quad (\operatorname{line\text{-}int}(g, m, c, a, b) = 0) \end{array}}}
\vnbentry{\texttt{caratheodory-\allowbreak{}remainder-\allowbreak{}small}}{\vnbstmt{\begin{array}{@{}l@{}} \forall s, f, a, l, e. \\ \quad \operatorname{is\text{-}diff\text{-}on}(\operatorname{cc\text{-}normed\text{-}field}, s, f, a, l) \wedge \\ \quad \operatorname{pos\text{-}rr}(e) \Rightarrow \\ \quad \quad \exists r. \\ \quad \quad \quad \operatorname{pos\text{-}rr}(r) \wedge \\ \quad \quad \quad \forall z \in s. \\ \quad \quad \quad \quad (\operatorname{magnitude}((z - a)) < r) \Rightarrow  \\ \quad \quad \quad \quad \quad \operatorname{magnitude}((f(z) - (f(a) + (l \cdot (z - a))))) \\ \quad \quad \quad \quad \quad \leq\; (e \cdot \operatorname{magnitude}((z - a))) \end{array}}}
\vnbentry{\texttt{nn-\allowbreak{}succ-\allowbreak{}le-\allowbreak{}power-\allowbreak{}two}}{\vnbstmt{\forall k \in \mathbb{N}.\; ((k + 1) \leq \operatorname{power}(2, k))}}
\vnbentry{\texttt{cc-\allowbreak{}telescoping-\allowbreak{}limit-\allowbreak{}bound}}{\vnbstmt{\begin{array}{@{}l@{}} \forall s \in (\mathbb{N} \to \mathbb{C}), c \in \mathbb{R}. \\ \quad (0 \leq c) \wedge \\ \quad \forall k \in \mathbb{N}. \\ \quad \quad \operatorname{magnitude}((s((k + 1)) - s(k))) \\ \quad \quad \leq\; (c \cdot \operatorname{recip}(\operatorname{power}(2, k))) \Rightarrow \\ \quad \quad \exists v \in \mathbb{C}. \\ \quad \quad \quad \operatorname{converges\text{-}to}(\operatorname{cc\text{-}ms}, s, v) \wedge \\ \quad \quad \quad \forall n \in \mathbb{N}. \\ \quad \quad \quad \quad \operatorname{magnitude}((s(n) - v)) \\ \quad \quad \quad \quad \leq\; ((2 \cdot c) \cdot \operatorname{recip}(\operatorname{power}(2, n))) \end{array}}}

\vnbfile{theorem-library/regulated-algebra.scm}{theorem-library/regulated-algebra.scm}
\noindent{\small regulated-algebra.scm -- THE ALGEBRA OF REGULATED FUNCTIONS ON [a, b], THE INTEGRAND OF A ROAD IS REGULATED, AND THE INTEGRAL ALONG A ROAD EXISTS. Batch 23-A, 2026-09-23.}\par
\vnbentry{\texttt{right-\allowbreak{}limit-\allowbreak{}within-\allowbreak{}sum}}{\vnbstmt{\begin{array}{@{}l@{}} \forall f, g, h, w, x, l, m. \\ \quad \operatorname{is\text{-}right\text{-}limit\text{-}within}(f, w, x, l) \wedge \\ \quad \operatorname{is\text{-}right\text{-}limit\text{-}within}(g, w, x, m) \wedge \\ \quad (h \in (w \to \mathbb{R})) \wedge \\ \quad \forall z \in w.\; (h(z) \simeq (f(z) + g(z))) \Rightarrow \\ \quad \quad \operatorname{is\text{-}right\text{-}limit\text{-}within}(h, w, x, (l + m)) \end{array}}}
\vnbentry{\texttt{left-\allowbreak{}limit-\allowbreak{}within-\allowbreak{}sum}}{\vnbstmt{\begin{array}{@{}l@{}} \forall f, g, h, w, x, l, m. \\ \quad \operatorname{is\text{-}left\text{-}limit\text{-}within}(f, w, x, l) \wedge \\ \quad \operatorname{is\text{-}left\text{-}limit\text{-}within}(g, w, x, m) \wedge \\ \quad (h \in (w \to \mathbb{R})) \wedge \\ \quad \forall z \in w.\; (h(z) \simeq (f(z) + g(z))) \Rightarrow \\ \quad \quad \operatorname{is\text{-}left\text{-}limit\text{-}within}(h, w, x, (l + m)) \end{array}}}
\vnbentry{\texttt{right-\allowbreak{}limit-\allowbreak{}within-\allowbreak{}sub}}{\vnbstmt{\begin{array}{@{}l@{}} \forall f, g, h, w, x, l, m. \\ \quad \operatorname{is\text{-}right\text{-}limit\text{-}within}(f, w, x, l) \wedge \\ \quad \operatorname{is\text{-}right\text{-}limit\text{-}within}(g, w, x, m) \wedge \\ \quad (h \in (w \to \mathbb{R})) \wedge \\ \quad \forall z \in w.\; (h(z) \simeq (f(z) - g(z))) \Rightarrow \\ \quad \quad \operatorname{is\text{-}right\text{-}limit\text{-}within}(h, w, x, (l - m)) \end{array}}}
\vnbentry{\texttt{left-\allowbreak{}limit-\allowbreak{}within-\allowbreak{}sub}}{\vnbstmt{\begin{array}{@{}l@{}} \forall f, g, h, w, x, l, m. \\ \quad \operatorname{is\text{-}left\text{-}limit\text{-}within}(f, w, x, l) \wedge \\ \quad \operatorname{is\text{-}left\text{-}limit\text{-}within}(g, w, x, m) \wedge \\ \quad (h \in (w \to \mathbb{R})) \wedge \\ \quad \forall z \in w.\; (h(z) \simeq (f(z) - g(z))) \Rightarrow \\ \quad \quad \operatorname{is\text{-}left\text{-}limit\text{-}within}(h, w, x, (l - m)) \end{array}}}
\vnbentry{\texttt{right-\allowbreak{}limit-\allowbreak{}within-\allowbreak{}mul}}{\vnbstmt{\begin{array}{@{}l@{}} \forall f, g, h, w, x, l, m. \\ \quad \operatorname{is\text{-}right\text{-}limit\text{-}within}(f, w, x, l) \wedge \\ \quad \operatorname{is\text{-}right\text{-}limit\text{-}within}(g, w, x, m) \wedge \\ \quad (h \in (w \to \mathbb{R})) \wedge \\ \quad \forall z \in w.\; (h(z) \simeq (f(z) \cdot g(z))) \Rightarrow \\ \quad \quad \operatorname{is\text{-}right\text{-}limit\text{-}within}(h, w, x, (l \cdot m)) \end{array}}}
\vnbentry{\texttt{left-\allowbreak{}limit-\allowbreak{}within-\allowbreak{}mul}}{\vnbstmt{\begin{array}{@{}l@{}} \forall f, g, h, w, x, l, m. \\ \quad \operatorname{is\text{-}left\text{-}limit\text{-}within}(f, w, x, l) \wedge \\ \quad \operatorname{is\text{-}left\text{-}limit\text{-}within}(g, w, x, m) \wedge \\ \quad (h \in (w \to \mathbb{R})) \wedge \\ \quad \forall z \in w.\; (h(z) \simeq (f(z) \cdot g(z))) \Rightarrow \\ \quad \quad \operatorname{is\text{-}left\text{-}limit\text{-}within}(h, w, x, (l \cdot m)) \end{array}}}
\vnbentry{\texttt{right-\allowbreak{}limit-\allowbreak{}within-\allowbreak{}scalar}}{\vnbstmt{\begin{array}{@{}l@{}} \forall f, h, w, x, l.; \forall c \in \mathbb{R}. \\ \quad \operatorname{is\text{-}right\text{-}limit\text{-}within}(f, w, x, l) \wedge \\ \quad (h \in (w \to \mathbb{R})) \wedge \\ \quad \forall z \in w.\; (h(z) \simeq (c \cdot f(z))) \Rightarrow \\ \quad \quad \operatorname{is\text{-}right\text{-}limit\text{-}within}(h, w, x, (c \cdot l)) \end{array}}}
\vnbentry{\texttt{left-\allowbreak{}limit-\allowbreak{}within-\allowbreak{}scalar}}{\vnbstmt{\begin{array}{@{}l@{}} \forall f, h, w, x, l.; \forall c \in \mathbb{R}. \\ \quad \operatorname{is\text{-}left\text{-}limit\text{-}within}(f, w, x, l) \wedge \\ \quad (h \in (w \to \mathbb{R})) \wedge \\ \quad \forall z \in w.\; (h(z) \simeq (c \cdot f(z))) \Rightarrow \\ \quad \quad \operatorname{is\text{-}left\text{-}limit\text{-}within}(h, w, x, (c \cdot l)) \end{array}}}
\vnbentry{\texttt{regulated-\allowbreak{}on-\allowbreak{}sum}}{\vnbstmt{\begin{array}{@{}l@{}} \forall f, g, h, a, b. \\ \quad \operatorname{is\text{-}regulated\text{-}on}(f, a, b) \wedge \\ \quad \operatorname{is\text{-}regulated\text{-}on}(g, a, b) \wedge \\ \quad (h \in (\operatorname{ccint}(a, b) \to \mathbb{R})) \wedge \\ \quad \forall z \in \operatorname{ccint}(a, b).\; (h(z) \simeq (f(z) + g(z))) \Rightarrow \\ \quad \quad \operatorname{is\text{-}regulated\text{-}on}(h, a, b) \end{array}}}
\vnbentry{\texttt{regulated-\allowbreak{}on-\allowbreak{}sub}}{\vnbstmt{\begin{array}{@{}l@{}} \forall f, g, h, a, b. \\ \quad \operatorname{is\text{-}regulated\text{-}on}(f, a, b) \wedge \\ \quad \operatorname{is\text{-}regulated\text{-}on}(g, a, b) \wedge \\ \quad (h \in (\operatorname{ccint}(a, b) \to \mathbb{R})) \wedge \\ \quad \forall z \in \operatorname{ccint}(a, b).\; (h(z) \simeq (f(z) - g(z))) \Rightarrow \\ \quad \quad \operatorname{is\text{-}regulated\text{-}on}(h, a, b) \end{array}}}
\vnbentry{\texttt{regulated-\allowbreak{}on-\allowbreak{}mul}}{\vnbstmt{\begin{array}{@{}l@{}} \forall f, g, h, a, b. \\ \quad \operatorname{is\text{-}regulated\text{-}on}(f, a, b) \wedge \\ \quad \operatorname{is\text{-}regulated\text{-}on}(g, a, b) \wedge \\ \quad (h \in (\operatorname{ccint}(a, b) \to \mathbb{R})) \wedge \\ \quad \forall z \in \operatorname{ccint}(a, b).\; (h(z) \simeq (f(z) \cdot g(z))) \Rightarrow \\ \quad \quad \operatorname{is\text{-}regulated\text{-}on}(h, a, b) \end{array}}}
\vnbentry{\texttt{regulated-\allowbreak{}on-\allowbreak{}scalar}}{\vnbstmt{\begin{array}{@{}l@{}} \forall f, h, a, b.; \forall c \in \mathbb{R}. \\ \quad \operatorname{is\text{-}regulated\text{-}on}(f, a, b) \wedge \\ \quad (h \in (\operatorname{ccint}(a, b) \to \mathbb{R})) \wedge \\ \quad \forall z \in \operatorname{ccint}(a, b).\; (h(z) \simeq (c \cdot f(z))) \Rightarrow \\ \quad \quad \operatorname{is\text{-}regulated\text{-}on}(h, a, b) \end{array}}}
\vnbentry{\texttt{regulated-\allowbreak{}on-\allowbreak{}congruence}}{\vnbstmt{\begin{array}{@{}l@{}} \forall f, g, a, b. \\ \quad \operatorname{is\text{-}regulated\text{-}on}(g, a, b) \wedge \\ \quad (f \in (\operatorname{ccint}(a, b) \to \mathbb{R})) \wedge \\ \quad \forall z \in \operatorname{ccint}(a, b).\; (f(z) \simeq g(z)) \Rightarrow \\ \quad \quad \operatorname{is\text{-}regulated\text{-}on}(f, a, b) \end{array}}}
\vnbentry{\texttt{continuous-\allowbreak{}on-\allowbreak{}compose-\allowbreak{}path}}{\vnbstmt{\begin{array}{@{}l@{}} \forall u, f, m, a, b. \\ \quad (u \subseteq \mathbb{C}) \wedge \\ \quad (f \in (u \to \mathbb{C})) \wedge \\ \quad \forall y \in u. \\ \quad \quad \operatorname{is\text{-}continuous\text{-}at}(\operatorname{subspace\text{-}ms}(\operatorname{nf\text{-}metric\text{-}space}(\operatorname{cc\text{-}normed\text{-}field}), u), \\ \quad \quad \quad \operatorname{nf\text{-}metric\text{-}space}(\operatorname{cc\text{-}normed\text{-}field}), f, y) \wedge \\ \quad \operatorname{is\text{-}path}(m, a, b) \wedge \\ \quad (\operatorname{trace}(m, a, b) \subseteq u) \Rightarrow \\ \quad \quad \operatorname{is\text{-}continuous\text{-}on}((\lambda t \in \operatorname{ccint}(a, b).\; \operatorname{real\text{-}part}(f(m(t)))), \operatorname{ccint}(a, b)) \wedge \\ \quad \quad \operatorname{is\text{-}continuous\text{-}on}((\lambda t \in \operatorname{ccint}(a, b).\; \operatorname{imag\text{-}part}(f(m(t)))), \operatorname{ccint}(a, b)) \end{array}}}
\vnbentry{\texttt{continuous-\allowbreak{}on-\allowbreak{}compose-\allowbreak{}regulated}}{\vnbstmt{\begin{array}{@{}l@{}} \forall u, f, m, a, b. \\ \quad (u \subseteq \mathbb{C}) \wedge \\ \quad (f \in (u \to \mathbb{C})) \wedge \\ \quad \forall y \in u. \\ \quad \quad \operatorname{is\text{-}continuous\text{-}at}(\operatorname{subspace\text{-}ms}(\operatorname{nf\text{-}metric\text{-}space}(\operatorname{cc\text{-}normed\text{-}field}), u), \\ \quad \quad \quad \operatorname{nf\text{-}metric\text{-}space}(\operatorname{cc\text{-}normed\text{-}field}), f, y) \wedge \\ \quad \operatorname{is\text{-}path}(m, a, b) \wedge \\ \quad (\operatorname{trace}(m, a, b) \subseteq u) \Rightarrow \\ \quad \quad \operatorname{is\text{-}regulated\text{-}on}((\lambda t \in \operatorname{ccint}(a, b).\; \operatorname{real\text{-}part}(f(m(t)))), a, b) \wedge \\ \quad \quad \operatorname{is\text{-}regulated\text{-}on}((\lambda t \in \operatorname{ccint}(a, b).\; \operatorname{imag\text{-}part}(f(m(t)))), a, b) \end{array}}}
\vnbentry{\texttt{road-\allowbreak{}integrand-\allowbreak{}regulated}}{\vnbstmt{\begin{array}{@{}l@{}} \forall u, f, m, c, a, b. \\ \quad (u \subseteq \mathbb{C}) \wedge \\ \quad (f \in (u \to \mathbb{C})) \wedge \\ \quad \forall y \in u. \\ \quad \quad \operatorname{is\text{-}continuous\text{-}at}(\operatorname{subspace\text{-}ms}(\operatorname{nf\text{-}metric\text{-}space}(\operatorname{cc\text{-}normed\text{-}field}), u), \\ \quad \quad \quad \operatorname{nf\text{-}metric\text{-}space}(\operatorname{cc\text{-}normed\text{-}field}), f, y) \wedge \\ \quad \operatorname{is\text{-}road}(m, c, a, b) \wedge \\ \quad (\operatorname{trace}(m, a, b) \subseteq u) \Rightarrow \\ \quad \quad \operatorname{is\text{-}regulated\text{-}on}((\lambda t \in \operatorname{ccint}(a, b).\; \operatorname{real\text{-}part}((f(m(t)) \cdot c(t)))), a, b) \wedge \\ \quad \quad \operatorname{is\text{-}regulated\text{-}on}((\lambda t \in \operatorname{ccint}(a, b).\; \operatorname{imag\text{-}part}((f(m(t)) \cdot c(t)))), a, b) \end{array}}}
\vnbentry{\texttt{line-\allowbreak{}int-\allowbreak{}exists}}{\vnbstmt{\begin{array}{@{}l@{}} \forall u, f, m, c, a, b. \\ \quad (u \subseteq \mathbb{C}) \wedge \\ \quad (f \in (u \to \mathbb{C})) \wedge \\ \quad \forall y \in u. \\ \quad \quad \operatorname{is\text{-}continuous\text{-}at}(\operatorname{subspace\text{-}ms}(\operatorname{nf\text{-}metric\text{-}space}(\operatorname{cc\text{-}normed\text{-}field}), u), \\ \quad \quad \quad \operatorname{nf\text{-}metric\text{-}space}(\operatorname{cc\text{-}normed\text{-}field}), f, y) \wedge \\ \quad \operatorname{is\text{-}road}(m, c, a, b) \wedge \\ \quad (\operatorname{trace}(m, a, b) \subseteq u) \Rightarrow \\ \quad \quad \exists \mathit{pwf}, w. \\ \quad \quad \quad \operatorname{is\text{-}primitive}(\mathit{pwf}, \\ \quad \quad \quad \quad (\lambda t \in \operatorname{ccint}(a, b).\; \\ \quad \quad \quad \quad \quad \operatorname{real\text{-}part}((\lambda t \in \operatorname{ccint}(a, b).\; (f(m(t)) \cdot c(t)))(t))), \\ \quad \quad \quad \quad a, b) \wedge \\ \quad \quad \quad \operatorname{is\text{-}primitive}(w, \\ \quad \quad \quad \quad (\lambda t \in \operatorname{ccint}(a, b).\; \\ \quad \quad \quad \quad \quad \operatorname{imag\text{-}part}((\lambda t \in \operatorname{ccint}(a, b).\; (f(m(t)) \cdot c(t)))(t))), \\ \quad \quad \quad \quad a, b) \end{array}}}
\vnbentry{\texttt{line-\allowbreak{}int-\allowbreak{}in-\allowbreak{}cc}}{\vnbstmt{\begin{array}{@{}l@{}} \forall u, f, m, c, a, b. \\ \quad (u \subseteq \mathbb{C}) \wedge \\ \quad (f \in (u \to \mathbb{C})) \wedge \\ \quad \forall y \in u. \\ \quad \quad \operatorname{is\text{-}continuous\text{-}at}(\operatorname{subspace\text{-}ms}(\operatorname{nf\text{-}metric\text{-}space}(\operatorname{cc\text{-}normed\text{-}field}), u), \\ \quad \quad \quad \operatorname{nf\text{-}metric\text{-}space}(\operatorname{cc\text{-}normed\text{-}field}), f, y) \wedge \\ \quad \operatorname{is\text{-}road}(m, c, a, b) \wedge \\ \quad (\operatorname{trace}(m, a, b) \subseteq u) \Rightarrow \\ \quad \quad (\operatorname{line\text{-}int}(f, m, c, a, b) \in \mathbb{C}) \end{array}}}

\vnbfile{theorem-library/compact-image.scm}{theorem-library/compact-image.scm}
\noindent{\small compact-image.scm -- THE CONTINUOUS IMAGE OF A COMPACT SPACE IS COMPACT (calculus.pdf, Proposition 3.19), and the COMPACTNESS OF THE TRACE of a path, which follows from it (complex-analysis.pdf 3.1.1: "gamma\_* is sometimes called the trace of gamma.  It is always a compact set").}\par
\vnbentry{\texttt{preimage-\allowbreak{}in-\allowbreak{}power}}{\vnbstmt{\begin{array}{@{}l@{}} \forall s, f, u. \\ \quad (\operatorname{pts}(s) \in \mathbf{Set}) \Rightarrow  \\ \quad \quad (\operatorname{preimage}(s, f, u) \in \operatorname{power}(\operatorname{pts}(s))) \end{array}}}
\vnbentry{\texttt{image-\allowbreak{}of-\allowbreak{}preimage-\allowbreak{}onto}}{\vnbstmt{\begin{array}{@{}l@{}} \forall s, f, u. \\ \quad (\operatorname{pts}(s) \in \mathbf{Set}) \wedge \\ \quad (u \subseteq \operatorname{image}(f, \operatorname{pts}(s))) \Rightarrow \\ \quad \quad (\operatorname{image}(f, \operatorname{preimage}(s, f, u)) = u) \end{array}}}
\vnbentry{\texttt{preimage-\allowbreak{}subspace-\allowbreak{}open}}{\vnbstmt{\begin{array}{@{}l@{}} \forall s, t, f, b, u. \\ \quad \operatorname{is\text{-}continuous}(s, t, f) \wedge \\ \quad (b \subseteq \operatorname{pts}(t)) \wedge \\ \quad \forall p \in \operatorname{pts}(s).\; (f(p) \in b) \wedge \\ \quad \operatorname{is\text{-}open}(\operatorname{subspace\text{-}ms}(t, b), u) \Rightarrow \\ \quad \quad \operatorname{is\text{-}open}(s, \operatorname{preimage}(s, f, u)) \end{array}}}
\vnbentry{\texttt{compact-\allowbreak{}image}}{\vnbstmt{\begin{array}{@{}l@{}} \forall s, t, f. \\ \quad (\operatorname{is\text{-}continuous}(s, t, f) \wedge \operatorname{is\text{-}compact}(s)) \Rightarrow  \\ \quad \quad \operatorname{is\text{-}compact}(\operatorname{subspace\text{-}ms}(t, \operatorname{image}(f, \operatorname{pts}(s)))) \end{array}}}
\vnbentry{\texttt{image-\allowbreak{}restrict}}{\vnbstmt{\begin{array}{@{}l@{}} \forall a \in \mathbf{Set}, f. \\ \quad (\operatorname{image}(\operatorname{restrict}(f, a), a) = \operatorname{image}(f, a)) \end{array}}}
\vnbentry{\texttt{continuous-\allowbreak{}image-\allowbreak{}of-\allowbreak{}compact}}{\vnbstmt{\begin{array}{@{}l@{}} \forall s, t, f, a. \\ \quad \operatorname{is\text{-}continuous}(s, t, f) \wedge \\ \quad (a \subseteq \operatorname{pts}(s)) \wedge \\ \quad \operatorname{is\text{-}compact}(\operatorname{subspace\text{-}ms}(s, a)) \Rightarrow \\ \quad \quad \operatorname{is\text{-}compact}(\operatorname{subspace\text{-}ms}(t, \operatorname{image}(f, a))) \end{array}}}
\vnbentry{\texttt{path-\allowbreak{}trace-\allowbreak{}compact}}{\vnbstmt{\begin{array}{@{}l@{}} \forall m, v, \mathit{bv}. \\ \quad \operatorname{is\text{-}path}(m, v, \mathit{bv}) \Rightarrow  \\ \quad \quad \operatorname{is\text{-}compact}(\operatorname{subspace\text{-}ms}(\operatorname{cc\text{-}ms}, \operatorname{trace}(m, v, \mathit{bv}))) \end{array}}}

\vnbfile{theorem-library/pw-int-order.scm}{theorem-library/pw-int-order.scm}
\noindent{\small pw-int-order.scm -- ORDER PROPERTIES OF THE PIECEWISE INTEGRAL.}\par
\vnbentry{\texttt{nondecreasing-\allowbreak{}total}}{\vnbstmt{\begin{array}{@{}l@{}} \forall h \in (\mathbb{R} \to \mathbb{R}), u, v \in \mathbb{R}. \\ \quad \forall x \in \mathbb{R}. \\ \quad \quad \operatorname{is\text{-}continuous\text{-}at}(\operatorname{rr\text{-}ms}, \operatorname{rr\text{-}ms}, h, x) \wedge \\ \quad (u < v) \wedge \\ \quad \forall t \in \operatorname{ooint}(u, v). \exists l. \\ \quad \quad (\operatorname{is\text{-}diff\text{-}at}(h, t, l) \wedge (0 \leq l)) \Rightarrow \\ \quad \quad (h(u) \leq h(v)) \end{array}}}
\vnbentry{\texttt{nondecreasing-\allowbreak{}off-\allowbreak{}finite-\allowbreak{}set}}{\vnbstmt{\begin{array}{@{}l@{}} \forall s \in \mathbf{Set}, h \in (\mathbb{R} \to \mathbb{R}), u, v \in \mathbb{R}. \\ \quad (|s| \in \mathbb{N}) \wedge \\ \quad \forall x \in \mathbb{R}. \\ \quad \quad \operatorname{is\text{-}continuous\text{-}at}(\operatorname{rr\text{-}ms}, \operatorname{rr\text{-}ms}, h, x) \wedge \\ \quad (u < v) \wedge \\ \quad \forall t \in \operatorname{ooint}(u, v). \\ \quad \quad \neg (t \in s) \Rightarrow  \\ \quad \quad \quad \exists l. \\ \quad \quad \quad \quad (\operatorname{is\text{-}diff\text{-}at}(h, t, l) \wedge (0 \leq l)) \Rightarrow \\ \quad \quad (h(u) \leq h(v)) \end{array}}}
\vnbentry{\texttt{pw-\allowbreak{}nondecreasing-\allowbreak{}on-\allowbreak{}interval}}{\vnbstmt{\begin{array}{@{}l@{}} \forall a, b \in \mathbb{R}, h \in (\operatorname{ccint}(a, b) \to \mathbb{R}), s \in \mathbf{Set}. \\ \quad (a < b) \wedge \\ \quad \operatorname{is\text{-}continuous\text{-}on}(h, \operatorname{ccint}(a, b)) \wedge \\ \quad (|s| \in \mathbb{N}) \wedge \\ \quad \forall t \in \operatorname{ooint}(a, b). \\ \quad \quad \neg (t \in s) \Rightarrow  \\ \quad \quad \quad \exists l. \\ \quad \quad \quad \quad (\operatorname{has\text{-}deriv\text{-}at}(h, t, l) \wedge (0 \leq l)) \Rightarrow \\ \quad \quad (h(a) \leq h(b)) \end{array}}}
\vnbentry{\texttt{nondecreasing-\allowbreak{}off-\allowbreak{}countable}}{\vnbstmt{\begin{array}{@{}l@{}} \forall h \in (\mathbb{R} \to \mathbb{R}), u, v \in \mathbb{R}, s. \\ \quad (u < v) \wedge \\ \quad \forall x \in \operatorname{ccint}(u, v). \\ \quad \quad \operatorname{is\text{-}continuous\text{-}at}(\operatorname{rr\text{-}ms}, \operatorname{rr\text{-}ms}, h, x) \wedge \\ \quad \operatorname{is\text{-}countable}(s) \wedge \\ \quad (s \subseteq \mathbb{R}) \wedge \\ \quad \forall t \in \operatorname{ooint}(u, v). \\ \quad \quad \neg (t \in s) \Rightarrow  \\ \quad \quad \quad \exists l. \\ \quad \quad \quad \quad (\operatorname{is\text{-}diff\text{-}at}(h, t, l) \wedge (0 \leq l)) \Rightarrow \\ \quad \quad (h(u) \leq h(v)) \end{array}}}
\vnbentry{\texttt{nondecreasing-\allowbreak{}off-\allowbreak{}countable-\allowbreak{}on-\allowbreak{}interval}}{\vnbstmt{\begin{array}{@{}l@{}} \forall a, b \in \mathbb{R}, h \in (\operatorname{ccint}(a, b) \to \mathbb{R}), s. \\ \quad (a < b) \wedge \\ \quad \operatorname{is\text{-}continuous\text{-}on}(h, \operatorname{ccint}(a, b)) \wedge \\ \quad \operatorname{is\text{-}countable}(s) \wedge \\ \quad (s \subseteq \operatorname{ccint}(a, b)) \wedge \\ \quad \forall t \in \operatorname{ooint}(a, b). \\ \quad \quad \neg (t \in s) \Rightarrow  \\ \quad \quad \quad \exists l. \\ \quad \quad \quad \quad (\operatorname{has\text{-}deriv\text{-}at}(h, t, l) \wedge (0 \leq l)) \Rightarrow \\ \quad \quad (h(a) \leq h(b)) \end{array}}}
\vnbentry{\texttt{pw-\allowbreak{}int-\allowbreak{}nonneg}}{\vnbstmt{\begin{array}{@{}l@{}} \forall f, i, a, b. \\ \quad \operatorname{is\text{-}primitive}(f, i, a, b) \wedge \\ \quad \forall y \in \operatorname{ccint}(a, b).\; (0 \leq i(y)) \Rightarrow \\ \quad \quad (0 \leq \int_{a}^{b} i) \end{array}}}
\vnbentry{\texttt{pw-\allowbreak{}int-\allowbreak{}monotone}}{\vnbstmt{\begin{array}{@{}l@{}} \forall f, i, w, \mathit{ppsi}, a, b. \\ \quad \operatorname{is\text{-}primitive}(f, i, a, b) \wedge \\ \quad \operatorname{is\text{-}primitive}(w, \mathit{ppsi}, a, b) \wedge \\ \quad \forall y \in \operatorname{ccint}(a, b).\; (i(y) \leq \operatorname{ppsi}(y)) \Rightarrow \\ \quad \quad (\int_{a}^{b} i \leq \int_{a}^{b} \mathit{ppsi}) \end{array}}}
\vnbentry{\texttt{pw-\allowbreak{}int-\allowbreak{}abs-\allowbreak{}bound}}{\vnbstmt{\begin{array}{@{}l@{}} \forall f, i, w, s, a, b. \\ \quad \operatorname{is\text{-}primitive}(f, i, a, b) \wedge \\ \quad \operatorname{is\text{-}primitive}(w, s, a, b) \wedge \\ \quad \forall y \in \operatorname{ccint}(a, b).\; (s(y) \simeq \lvert i(y) \rvert) \Rightarrow \\ \quad \quad (\lvert \int_{a}^{b} i \rvert \leq \int_{a}^{b} s) \end{array}}}

\vnbfile{theorem-library/cc-int-estimate.scm}{theorem-library/cc-int-estimate.scm}
\noindent{\small cc-int-estimate.scm -- THE NOTES' ESTIMATE (45).}\par
\vnbentry{\texttt{cc-\allowbreak{}magnitude-\allowbreak{}sq}}{\vnbstmt{\begin{array}{@{}l@{}} \forall z \in \mathbb{C}. \\ \quad (\operatorname{magnitude}(z) \cdot \operatorname{magnitude}(z)) \\ \quad =\; ((\operatorname{real\text{-}part}(z) \cdot \operatorname{real\text{-}part}(z)) + (\operatorname{imag\text{-}part}(z) \cdot \operatorname{imag\text{-}part}(z))) \end{array}}}
\vnbentry{\texttt{cc-\allowbreak{}re-\allowbreak{}inner-\allowbreak{}bound}}{\vnbstmt{\begin{array}{@{}l@{}} \forall w, z \in \mathbb{C}. \\ \quad ((\operatorname{real\text{-}part}(w) \cdot \operatorname{real\text{-}part}(z)) + (\operatorname{imag\text{-}part}(w) \cdot \operatorname{imag\text{-}part}(z))) \\ \quad \leq\; (\operatorname{magnitude}(w) \cdot \operatorname{magnitude}(z)) \end{array}}}
\vnbentry{\texttt{cc-\allowbreak{}int-\allowbreak{}abs-\allowbreak{}bound}}{\vnbstmt{\begin{array}{@{}l@{}} \forall i, f, w, s, h, a, b. \\ \quad (i \in (\operatorname{ccint}(a, b) \to \mathbb{C})) \wedge \\ \quad \operatorname{is\text{-}primitive}(f, (\lambda t \in \operatorname{ccint}(a, b).\; \operatorname{real\text{-}part}(i(t))), a, b) \wedge \\ \quad \operatorname{is\text{-}primitive}(w, (\lambda t \in \operatorname{ccint}(a, b).\; \operatorname{imag\text{-}part}(i(t))), a, b) \wedge \\ \quad \operatorname{is\text{-}primitive}(h, s, a, b) \wedge \\ \quad \forall y \in \operatorname{ccint}(a, b). \\ \quad \quad (s(y) \simeq \operatorname{magnitude}(i(y))) \Rightarrow \\ \quad \quad (\operatorname{magnitude}(\int_{a}^{b} i) \leq \int_{a}^{b} s) \end{array}}}

\vnbfile{theorem-library/line-int-laws.scm}{theorem-library/line-int-laws.scm}
\noindent{\small line-int-laws.scm -- THE LAWS OF THE INTEGRAL ALONG A ROAD, for a road and a CONTINUOUS f, with no primitive in the statement. Batch 24-A, 2026-09-23.}\par
\vnbentry{\texttt{line-\allowbreak{}int-\allowbreak{}value}}{\vnbstmt{\begin{array}{@{}l@{}} \forall f, m, c, a, b, g_{1}, g_{2}. \\ \quad \operatorname{is\text{-}primitive}(g_{1}, (\lambda t \in \operatorname{ccint}(a, b).\; \operatorname{real\text{-}part}((f(m(t)) \cdot c(t)))), a, b) \wedge \\ \quad \operatorname{is\text{-}primitive}(g_{2}, (\lambda t \in \operatorname{ccint}(a, b).\; \operatorname{imag\text{-}part}((f(m(t)) \cdot c(t)))), a, b) \Rightarrow \\ \quad \quad \operatorname{line\text{-}int}(f, m, c, a, b) \\ \quad \quad =\; ((\operatorname{g1}(b) - \operatorname{g1}(a)) + ((\operatorname{g2}(b) - \operatorname{g2}(a)) \cdot +i)) \end{array}}}
\vnbentry{\texttt{line-\allowbreak{}int-\allowbreak{}sum}}{\vnbstmt{\begin{array}{@{}l@{}} \forall u, f, g, m, c, a, b.; \forall h \in (u \to \mathbb{C}). \\ \quad (u \subseteq \mathbb{C}) \wedge \\ \quad (f \in (u \to \mathbb{C})) \wedge \\ \quad \forall y \in u. \\ \quad \quad \operatorname{is\text{-}continuous\text{-}at}(\operatorname{subspace\text{-}ms}(\operatorname{nf\text{-}metric\text{-}space}(\operatorname{cc\text{-}normed\text{-}field}), u), \\ \quad \quad \quad \operatorname{nf\text{-}metric\text{-}space}(\operatorname{cc\text{-}normed\text{-}field}), f, y) \wedge \\ \quad (g \in (u \to \mathbb{C})) \wedge \\ \quad \forall y \in u. \\ \quad \quad \operatorname{is\text{-}continuous\text{-}at}(\operatorname{subspace\text{-}ms}(\operatorname{nf\text{-}metric\text{-}space}(\operatorname{cc\text{-}normed\text{-}field}), u), \\ \quad \quad \quad \operatorname{nf\text{-}metric\text{-}space}(\operatorname{cc\text{-}normed\text{-}field}), g, y) \wedge \\ \quad \operatorname{is\text{-}road}(m, c, a, b) \wedge \\ \quad (\operatorname{trace}(m, a, b) \subseteq u) \wedge \\ \quad \forall z \in u.\; (h(z) \simeq (f(z) + g(z))) \Rightarrow \\ \quad \quad \operatorname{line\text{-}int}(h, m, c, a, b) \\ \quad \quad =\; (\operatorname{line\text{-}int}(f, m, c, a, b) + \operatorname{line\text{-}int}(g, m, c, a, b)) \end{array}}}
\vnbentry{\texttt{line-\allowbreak{}int-\allowbreak{}scalar}}{\vnbstmt{\begin{array}{@{}l@{}} \forall u, f, m, c, a, b.; \forall \mathit{lic} \in \mathbb{C}, h \in (u \to \mathbb{C}). \\ \quad (u \subseteq \mathbb{C}) \wedge \\ \quad (f \in (u \to \mathbb{C})) \wedge \\ \quad \forall y \in u. \\ \quad \quad \operatorname{is\text{-}continuous\text{-}at}(\operatorname{subspace\text{-}ms}(\operatorname{nf\text{-}metric\text{-}space}(\operatorname{cc\text{-}normed\text{-}field}), u), \\ \quad \quad \quad \operatorname{nf\text{-}metric\text{-}space}(\operatorname{cc\text{-}normed\text{-}field}), f, y) \wedge \\ \quad \operatorname{is\text{-}road}(m, c, a, b) \wedge \\ \quad (\operatorname{trace}(m, a, b) \subseteq u) \wedge \\ \quad \forall z \in u.\; (h(z) \simeq (\mathit{lic} \cdot f(z))) \Rightarrow \\ \quad \quad \operatorname{line\text{-}int}(h, m, c, a, b) \\ \quad \quad =\; (\mathit{lic} \cdot \operatorname{line\text{-}int}(f, m, c, a, b)) \end{array}}}
\vnbentry{\texttt{line-\allowbreak{}int-\allowbreak{}opposite-\allowbreak{}road}}{\vnbstmt{\begin{array}{@{}l@{}} \forall u, f, m, c, a, b, g, d. \\ \quad (u \subseteq \mathbb{C}) \wedge \\ \quad (f \in (u \to \mathbb{C})) \wedge \\ \quad \forall y \in u. \\ \quad \quad \operatorname{is\text{-}continuous\text{-}at}(\operatorname{subspace\text{-}ms}(\operatorname{nf\text{-}metric\text{-}space}(\operatorname{cc\text{-}normed\text{-}field}), u), \\ \quad \quad \quad \operatorname{nf\text{-}metric\text{-}space}(\operatorname{cc\text{-}normed\text{-}field}), f, y) \wedge \\ \quad \operatorname{is\text{-}road}(m, c, a, b) \wedge \\ \quad (\operatorname{trace}(m, a, b) \subseteq u) \wedge \\ \quad (g \in (\operatorname{ccint}(a, b) \to \mathbb{C})) \wedge \\ \quad (d \in (\operatorname{ccint}(a, b) \to \mathbb{C})) \wedge \\ \quad \forall \mathit{pay} \in \operatorname{ccint}(a, b). \\ \quad \quad (g(\mathit{pay}) \simeq m(((a + b) - \mathit{pay}))) \wedge \\ \quad \forall \mathit{pay} \in \operatorname{ccint}(a, b). \\ \quad \quad (d(\mathit{pay}) \simeq (-1 \cdot c(((a + b) - \mathit{pay})))) \Rightarrow \\ \quad \quad \operatorname{line\text{-}int}(f, g, d, a, b) \\ \quad \quad =\; (-1 \cdot \operatorname{line\text{-}int}(f, m, c, a, b)) \end{array}}}
\vnbentry{\texttt{road-\allowbreak{}restrict}}{\vnbstmt{\begin{array}{@{}l@{}} \forall m, c, a, b.; \forall v, \mathit{lidv} \in \mathbb{R}. \\ \quad \operatorname{is\text{-}road}(m, c, a, b) \wedge \\ \quad (a \leq v) \wedge \\ \quad (v < \mathit{lidv}) \wedge \\ \quad (\mathit{lidv} \leq b) \Rightarrow \\ \quad \quad \operatorname{is\text{-}road}(\operatorname{restrict}(m, \operatorname{ccint}(v, \mathit{lidv})), \operatorname{restrict}(c, \operatorname{ccint}(v, \mathit{lidv})), v, \mathit{lidv}) \end{array}}}
\vnbentry{\texttt{cc-\allowbreak{}int-\allowbreak{}congruence}}{\vnbstmt{\begin{array}{@{}l@{}} \forall i, \mathit{ppsi}, a, b, g_{1}, g_{2}. \\ \quad \operatorname{is\text{-}primitive}(g_{1}, (\lambda t \in \operatorname{ccint}(a, b).\; \operatorname{real\text{-}part}(i(t))), a, b) \wedge \\ \quad \operatorname{is\text{-}primitive}(g_{2}, (\lambda t \in \operatorname{ccint}(a, b).\; \operatorname{imag\text{-}part}(i(t))), a, b) \wedge \\ \quad (\mathit{ppsi} \in (\operatorname{ccint}(a, b) \to \mathbb{C})) \wedge \\ \quad \forall y \in \operatorname{ccint}(a, b).\; (i(y) \simeq \operatorname{ppsi}(y)) \Rightarrow \\ \quad \quad (\int_{a}^{b} i = \int_{a}^{b} \mathit{ppsi}) \end{array}}}
\vnbentry{\texttt{line-\allowbreak{}int-\allowbreak{}adjacent}}{\vnbstmt{\begin{array}{@{}l@{}} \forall u, f, m, c, a, b.; \forall v \in \operatorname{ooint}(a, b). \\ \quad (u \subseteq \mathbb{C}) \wedge \\ \quad (f \in (u \to \mathbb{C})) \wedge \\ \quad \forall y \in u. \\ \quad \quad \operatorname{is\text{-}continuous\text{-}at}(\operatorname{subspace\text{-}ms}(\operatorname{nf\text{-}metric\text{-}space}(\operatorname{cc\text{-}normed\text{-}field}), u), \\ \quad \quad \quad \operatorname{nf\text{-}metric\text{-}space}(\operatorname{cc\text{-}normed\text{-}field}), f, y) \wedge \\ \quad \operatorname{is\text{-}road}(m, c, a, b) \wedge \\ \quad (\operatorname{trace}(m, a, b) \subseteq u) \Rightarrow \\ \quad \quad \operatorname{line\text{-}int}(f, m, c, a, b) \\ \quad \quad =\; (\operatorname{line\text{-}int}(f, \operatorname{restrict}(m, \operatorname{ccint}(a, v)), \operatorname{restrict}(c, \operatorname{ccint}(a, v)), a, v) + \operatorname{line\text{-}int}(f, \operatorname{restrict}(m, \operatorname{ccint}(v, b)), \operatorname{restrict}(c, \operatorname{ccint}(v, b)), v, b)) \end{array}}}
\vnbentry{\texttt{right-\allowbreak{}limit-\allowbreak{}within-\allowbreak{}magnitude}}{\vnbstmt{\begin{array}{@{}l@{}} \forall f, g, h, w, x, l, k. \\ \quad \operatorname{is\text{-}right\text{-}limit\text{-}within}(f, w, x, l) \wedge \\ \quad \operatorname{is\text{-}right\text{-}limit\text{-}within}(g, w, x, k) \wedge \\ \quad (h \in (w \to \mathbb{R})) \wedge \\ \quad \forall z \in w. \\ \quad \quad (h(z) \simeq \operatorname{magnitude}((f(z) + (g(z) \cdot +i)))) \Rightarrow \\ \quad \quad \operatorname{is\text{-}right\text{-}limit\text{-}within}(h, w, x, \operatorname{magnitude}((l + (k \cdot +i)))) \end{array}}}
\vnbentry{\texttt{left-\allowbreak{}limit-\allowbreak{}within-\allowbreak{}magnitude}}{\vnbstmt{\begin{array}{@{}l@{}} \forall f, g, h, w, x, l, k. \\ \quad \operatorname{is\text{-}left\text{-}limit\text{-}within}(f, w, x, l) \wedge \\ \quad \operatorname{is\text{-}left\text{-}limit\text{-}within}(g, w, x, k) \wedge \\ \quad (h \in (w \to \mathbb{R})) \wedge \\ \quad \forall z \in w. \\ \quad \quad (h(z) \simeq \operatorname{magnitude}((f(z) + (g(z) \cdot +i)))) \Rightarrow \\ \quad \quad \operatorname{is\text{-}left\text{-}limit\text{-}within}(h, w, x, \operatorname{magnitude}((l + (k \cdot +i)))) \end{array}}}
\vnbentry{\texttt{regulated-\allowbreak{}on-\allowbreak{}magnitude}}{\vnbstmt{\begin{array}{@{}l@{}} \forall f, g, h, a, b. \\ \quad \operatorname{is\text{-}regulated\text{-}on}(f, a, b) \wedge \\ \quad \operatorname{is\text{-}regulated\text{-}on}(g, a, b) \wedge \\ \quad (h \in (\operatorname{ccint}(a, b) \to \mathbb{R})) \wedge \\ \quad \forall z \in \operatorname{ccint}(a, b). \\ \quad \quad (h(z) \simeq \operatorname{magnitude}((f(z) + (g(z) \cdot +i)))) \Rightarrow \\ \quad \quad \operatorname{is\text{-}regulated\text{-}on}(h, a, b) \end{array}}}
\vnbentry{\texttt{line-\allowbreak{}int-\allowbreak{}abs-\allowbreak{}bound}}{\vnbstmt{\begin{array}{@{}l@{}} \forall u, f, m, c, a, b.; \forall \mathit{liab} \in (\operatorname{ccint}(a, b) \to \mathbb{R}). \\ \quad (u \subseteq \mathbb{C}) \wedge \\ \quad (f \in (u \to \mathbb{C})) \wedge \\ \quad \forall y \in u. \\ \quad \quad \operatorname{is\text{-}continuous\text{-}at}(\operatorname{subspace\text{-}ms}(\operatorname{nf\text{-}metric\text{-}space}(\operatorname{cc\text{-}normed\text{-}field}), u), \\ \quad \quad \quad \operatorname{nf\text{-}metric\text{-}space}(\operatorname{cc\text{-}normed\text{-}field}), f, y) \wedge \\ \quad \operatorname{is\text{-}road}(m, c, a, b) \wedge \\ \quad (\operatorname{trace}(m, a, b) \subseteq u) \wedge \\ \quad \forall \mathit{liy} \in \operatorname{ccint}(a, b). \\ \quad \quad (\operatorname{liab}(\mathit{liy}) \simeq (\operatorname{magnitude}(f(m(\mathit{liy}))) \cdot \operatorname{magnitude}(c(\mathit{liy})))) \Rightarrow \\ \quad \quad \operatorname{magnitude}(\operatorname{line\text{-}int}(f, m, c, a, b)) \\ \quad \quad \leq\; \int_{a}^{b} \mathit{liab} \end{array}}}

\vnbfile{theorem-library/line-int-reparam.scm}{theorem-library/line-int-reparam.scm}
\noindent{\small line-int-reparam.scm -- THE AFFINE CHANGE OF PARAMETER (the notes' (48), affine case) AND THE SEGMENT LEMMAS OF GOURSAT'S DEFINITION-FREE HALF. Batch 29-A, 2026-09-25.}\par
\vnbentry{\texttt{affine-\allowbreak{}reparam-\allowbreak{}ccint-\allowbreak{}in}}{\vnbstmt{\begin{array}{@{}l@{}} \forall a, b, c, d, l, m \in \mathbb{R}, z \in \operatorname{ccint}(c, d). \\ \quad ((0 < l) \wedge (((m + (l \cdot c)) = a) \wedge ((m + (l \cdot d)) = b))) \Rightarrow  \\ \quad \quad ((m + (l \cdot z)) \in \operatorname{ccint}(a, b)) \end{array}}}
\vnbentry{\texttt{affine-\allowbreak{}reparam-\allowbreak{}ooint-\allowbreak{}in}}{\vnbstmt{\begin{array}{@{}l@{}} \forall a, b, c, d, l, m \in \mathbb{R}, z \in \operatorname{ooint}(c, d). \\ \quad ((0 < l) \wedge (((m + (l \cdot c)) = a) \wedge ((m + (l \cdot d)) = b))) \Rightarrow  \\ \quad \quad ((m + (l \cdot z)) \in \operatorname{ooint}(a, b)) \end{array}}}
\vnbentry{\texttt{affine-\allowbreak{}reparam-\allowbreak{}solve}}{\vnbstmt{\begin{array}{@{}l@{}} \forall a, b, c, d, l, m \in \mathbb{R}, z \in \operatorname{ccint}(a, b), q \in \mathbb{R}. \\ \quad (0 < l) \wedge \\ \quad ((m + (l \cdot c)) = a) \wedge \\ \quad ((m + (l \cdot d)) = b) \wedge \\ \quad ((z - m) = (q \cdot l)) \Rightarrow \\ \quad \quad ((q \in \operatorname{ccint}(c, d)) \wedge ((m + (l \cdot q)) = z)) \end{array}}}
\vnbentry{\texttt{affine-\allowbreak{}reparam-\allowbreak{}inverse}}{\vnbstmt{\begin{array}{@{}l@{}} \forall a, b, c, d, l, m \in \mathbb{R}, z \in \operatorname{ccint}(a, b). \\ \quad ((0 < l) \wedge (((m + (l \cdot c)) = a) \wedge ((m + (l \cdot d)) = b))) \Rightarrow  \\ \quad \quad (((z - m) \cdot \operatorname{recip}(l)) \in \operatorname{ccint}(c, d)) \wedge \\ \quad \quad ((m + (l \cdot ((z - m) \cdot \operatorname{recip}(l)))) = z) \end{array}}}
\vnbentry{\texttt{affine-\allowbreak{}reparam-\allowbreak{}inverse-\allowbreak{}value}}{\vnbstmt{\begin{array}{@{}l@{}} \forall l, m, z \in \mathbb{R}. \\ \quad (0 < l) \Rightarrow  \\ \quad \quad ((((m + (l \cdot z)) - m) \cdot \operatorname{recip}(l)) = z) \end{array}}}
\vnbentry{\texttt{affine-\allowbreak{}reparam-\allowbreak{}continuous-\allowbreak{}at-\allowbreak{}sub}}{\vnbstmt{\begin{array}{@{}l@{}} \forall a, b, c, d, l, m \in \mathbb{R}, f.; \forall g \in (\operatorname{ccint}(c, d) \to \mathbb{R}), t \in \operatorname{ccint}(c, d). \\ \quad (0 < l) \wedge \\ \quad ((m + (l \cdot c)) = a) \wedge \\ \quad ((m + (l \cdot d)) = b) \wedge \\ \quad \forall z \in \operatorname{ccint}(c, d).\; (g(z) \simeq f((m + (l \cdot z)))) \wedge \\ \quad \operatorname{is\text{-}continuous\text{-}at}(\operatorname{subspace\text{-}ms}(\operatorname{rr\text{-}ms}, \operatorname{ccint}(a, b)), \operatorname{rr\text{-}ms}, f, (m + (l \cdot t))) \Rightarrow \\ \quad \quad \operatorname{is\text{-}continuous\text{-}at}(\operatorname{subspace\text{-}ms}(\operatorname{rr\text{-}ms}, \operatorname{ccint}(c, d)), \operatorname{rr\text{-}ms}, g, t) \end{array}}}
\vnbentry{\texttt{affine-\allowbreak{}reparam-\allowbreak{}continuous-\allowbreak{}on}}{\vnbstmt{\begin{array}{@{}l@{}} \forall a, b, c, d, l, m \in \mathbb{R}, f.; \forall g \in (\operatorname{ccint}(c, d) \to \mathbb{R}). \\ \quad (0 < l) \wedge \\ \quad ((m + (l \cdot c)) = a) \wedge \\ \quad ((m + (l \cdot d)) = b) \wedge \\ \quad \operatorname{is\text{-}continuous\text{-}on}(f, \operatorname{ccint}(a, b)) \wedge \\ \quad \forall z \in \operatorname{ccint}(c, d).\; (g(z) \simeq f((m + (l \cdot z)))) \Rightarrow \\ \quad \quad \operatorname{is\text{-}continuous\text{-}on}(g, \operatorname{ccint}(c, d)) \end{array}}}
\vnbentry{\texttt{affine-\allowbreak{}reparam-\allowbreak{}image-\allowbreak{}countable}}{\vnbstmt{\begin{array}{@{}l@{}} \forall a, b, c, d, l, m \in \mathbb{R}, e. \\ \quad (0 < l) \wedge \\ \quad ((m + (l \cdot c)) = a) \wedge \\ \quad ((m + (l \cdot d)) = b) \wedge \\ \quad \operatorname{is\text{-}countable}(e) \wedge \\ \quad (e \subseteq \operatorname{ccint}(a, b)) \Rightarrow \\ \quad \quad \operatorname{is\text{-}countable}(\operatorname{image}((\lambda v \in \operatorname{ccint}(a, b).\; ((v - m) \cdot \operatorname{recip}(l))), e)) \wedge \\ \quad \quad \operatorname{image}((\lambda v \in \operatorname{ccint}(a, b).\; ((v - m) \cdot \operatorname{recip}(l))), e) \\ \quad \quad \subseteq\; \operatorname{ccint}(c, d) \end{array}}}
\vnbentry{\texttt{affine-\allowbreak{}reparam-\allowbreak{}not-\allowbreak{}in-\allowbreak{}image}}{\vnbstmt{\begin{array}{@{}l@{}} \forall a, b, c, d, l, m \in \mathbb{R}, e.; \forall t \in \operatorname{ccint}(c, d). \\ \quad (0 < l) \wedge \\ \quad ((m + (l \cdot c)) = a) \wedge \\ \quad ((m + (l \cdot d)) = b) \wedge \\ \quad (e \subseteq \operatorname{ccint}(a, b)) \wedge \\ \quad \neg (t \in \operatorname{image}((\lambda v \in \operatorname{ccint}(a, b).\; ((v - m) \cdot \operatorname{recip}(l))), e)) \Rightarrow \\ \quad \quad \neg ((m + (l \cdot t)) \in e) \end{array}}}
\vnbentry{\texttt{affine-\allowbreak{}reparam-\allowbreak{}regulated}}{\vnbstmt{\begin{array}{@{}l@{}} \forall a, b, c, d, l, m \in \mathbb{R}, f.; \forall g \in (\operatorname{ccint}(c, d) \to \mathbb{R}). \\ \quad (0 < l) \wedge \\ \quad ((m + (l \cdot c)) = a) \wedge \\ \quad ((m + (l \cdot d)) = b) \wedge \\ \quad (c < d) \wedge \\ \quad \operatorname{is\text{-}regulated\text{-}on}(f, a, b) \wedge \\ \quad \forall z \in \operatorname{ccint}(c, d).\; (g(z) \simeq f((m + (l \cdot z)))) \Rightarrow \\ \quad \quad \operatorname{is\text{-}regulated\text{-}on}(g, c, d) \end{array}}}
\vnbentry{\texttt{pw-\allowbreak{}int-\allowbreak{}affine-\allowbreak{}subst}}{\vnbstmt{\begin{array}{@{}l@{}} \forall a, b, c, d, l, m \in \mathbb{R}, f, i.; \forall h, s \in (\operatorname{ccint}(c, d) \to \mathbb{R}). \\ \quad (0 < l) \wedge \\ \quad ((m + (l \cdot c)) = a) \wedge \\ \quad ((m + (l \cdot d)) = b) \wedge \\ \quad (c < d) \wedge \\ \quad \operatorname{is\text{-}primitive}(f, i, a, b) \wedge \\ \quad \forall z \in \operatorname{ccint}(c, d).\; (h(z) \simeq f((m + (l \cdot z)))) \wedge \\ \quad \forall z \in \operatorname{ccint}(c, d). \\ \quad \quad (s(z) \simeq (l \cdot i((m + (l \cdot z))))) \Rightarrow \\ \quad \quad (\operatorname{is\text{-}primitive}(h, s, c, d) \wedge (\int_{c}^{d} s = \int_{a}^{b} i)) \end{array}}}
\vnbentry{\texttt{cc-\allowbreak{}int-\allowbreak{}affine-\allowbreak{}subst}}{\vnbstmt{\begin{array}{@{}l@{}} \forall a, b, c, d, l, m \in \mathbb{R}, i, \mathit{lrpb}, f, w. \\ \quad (0 < l) \wedge \\ \quad ((m + (l \cdot c)) = a) \wedge \\ \quad ((m + (l \cdot d)) = b) \wedge \\ \quad (c < d) \wedge \\ \quad (i \in (\operatorname{ccint}(a, b) \to \mathbb{C})) \wedge \\ \quad (\mathit{lrpb} \in (\operatorname{ccint}(c, d) \to \mathbb{C})) \wedge \\ \quad \forall z \in \operatorname{ccint}(c, d). \\ \quad \quad (\operatorname{lrpb}(z) \simeq (l \cdot i((m + (l \cdot z))))) \wedge \\ \quad \operatorname{is\text{-}primitive}(f, (\lambda t \in \operatorname{ccint}(a, b).\; \operatorname{real\text{-}part}(i(t))), a, b) \wedge \\ \quad \operatorname{is\text{-}primitive}(w, (\lambda t \in \operatorname{ccint}(a, b).\; \operatorname{imag\text{-}part}(i(t))), a, b) \Rightarrow \\ \quad \quad (\int_{c}^{d} \mathit{lrpb} = \int_{a}^{b} i) \end{array}}}
\vnbentry{\texttt{trace-\allowbreak{}affine-\allowbreak{}reparam}}{\vnbstmt{\begin{array}{@{}l@{}} \forall a, b, c, d, l, m \in \mathbb{R}, f, r. \\ \quad (0 < l) \wedge \\ \quad ((m + (l \cdot c)) = a) \wedge \\ \quad ((m + (l \cdot d)) = b) \wedge \\ \quad (f \in (\operatorname{ccint}(a, b) \to \mathbb{C})) \wedge \\ \quad (r \in (\operatorname{ccint}(c, d) \to \mathbb{C})) \wedge \\ \quad \forall z \in \operatorname{ccint}(c, d).\; (r(z) \simeq f((m + (l \cdot z)))) \Rightarrow \\ \quad \quad (\operatorname{trace}(r, c, d) \subseteq \operatorname{trace}(f, a, b)) \wedge \\ \quad \quad (\operatorname{trace}(f, a, b) \subseteq \operatorname{trace}(r, c, d)) \end{array}}}
\vnbentry{\texttt{road-\allowbreak{}affine-\allowbreak{}reparam}}{\vnbstmt{\begin{array}{@{}l@{}} \forall a, b, c, d, l, m \in \mathbb{R}, f, g, r, k. \\ \quad (0 < l) \wedge \\ \quad ((m + (l \cdot c)) = a) \wedge \\ \quad ((m + (l \cdot d)) = b) \wedge \\ \quad (c < d) \wedge \\ \quad \operatorname{is\text{-}road}(f, g, a, b) \wedge \\ \quad (r \in (\operatorname{ccint}(c, d) \to \mathbb{C})) \wedge \\ \quad (k \in (\operatorname{ccint}(c, d) \to \mathbb{C})) \wedge \\ \quad \forall z \in \operatorname{ccint}(c, d).\; (r(z) \simeq f((m + (l \cdot z)))) \wedge \\ \quad \forall z \in \operatorname{ccint}(c, d). \\ \quad \quad (k(z) \simeq (l \cdot g((m + (l \cdot z))))) \Rightarrow \\ \quad \quad \operatorname{is\text{-}road}(r, k, c, d) \end{array}}}
\vnbentry{\texttt{line-\allowbreak{}int-\allowbreak{}affine-\allowbreak{}reparam}}{\vnbstmt{\begin{array}{@{}l@{}} \forall u, f, m, c, a, b.; \forall \mathit{lrpc}, d, l, \mathit{lrpm} \in \mathbb{R}, r, k. \\ \quad (u \subseteq \mathbb{C}) \wedge \\ \quad (f \in (u \to \mathbb{C})) \wedge \\ \quad \forall y \in u. \\ \quad \quad \operatorname{is\text{-}continuous\text{-}at}(\operatorname{subspace\text{-}ms}(\operatorname{nf\text{-}metric\text{-}space}(\operatorname{cc\text{-}normed\text{-}field}), u), \\ \quad \quad \quad \operatorname{nf\text{-}metric\text{-}space}(\operatorname{cc\text{-}normed\text{-}field}), f, y) \wedge \\ \quad \operatorname{is\text{-}road}(m, c, a, b) \wedge \\ \quad (\operatorname{trace}(m, a, b) \subseteq u) \wedge \\ \quad (0 < l) \wedge \\ \quad ((\mathit{lrpm} + (l \cdot \mathit{lrpc})) = a) \wedge \\ \quad ((\mathit{lrpm} + (l \cdot d)) = b) \wedge \\ \quad (\mathit{lrpc} < d) \wedge \\ \quad (r \in (\operatorname{ccint}(\mathit{lrpc}, d) \to \mathbb{C})) \wedge \\ \quad (k \in (\operatorname{ccint}(\mathit{lrpc}, d) \to \mathbb{C})) \wedge \\ \quad \forall z \in \operatorname{ccint}(\mathit{lrpc}, d). \\ \quad \quad (r(z) \simeq m((\mathit{lrpm} + (l \cdot z)))) \wedge \\ \quad \forall z \in \operatorname{ccint}(\mathit{lrpc}, d). \\ \quad \quad (k(z) \simeq (l \cdot c((\mathit{lrpm} + (l \cdot z))))) \Rightarrow \\ \quad \quad \operatorname{line\text{-}int}(f, r, k, \mathit{lrpc}, d) \\ \quad \quad =\; \operatorname{line\text{-}int}(f, m, c, a, b) \end{array}}}
\vnbentry{\texttt{trace-\allowbreak{}restrict-\allowbreak{}subset}}{\vnbstmt{\begin{array}{@{}l@{}} \forall m, a, b.; \forall c, d \in \mathbb{R}. \\ \quad (m \in (\operatorname{ccint}(a, b) \to \mathbb{C})) \wedge \\ \quad (a \in \mathbb{R}) \wedge \\ \quad (b \in \mathbb{R}) \wedge \\ \quad (a \leq c) \wedge \\ \quad (d \leq b) \Rightarrow \\ \quad \quad \operatorname{trace}(\operatorname{restrict}(m, \operatorname{ccint}(c, d)), c, d) \\ \quad \quad \subseteq\; \operatorname{trace}(m, a, b) \end{array}}}
\vnbentry{\texttt{segment-\allowbreak{}int-\allowbreak{}split}}{\vnbstmt{\begin{array}{@{}l@{}} \forall u, f.; \forall z_{0}, w \in \mathbb{C}, t \in \operatorname{ooint}(0, 1). \\ \quad (u \subseteq \mathbb{C}) \wedge \\ \quad (f \in (u \to \mathbb{C})) \wedge \\ \quad \forall y \in u. \\ \quad \quad \operatorname{is\text{-}continuous\text{-}at}(\operatorname{subspace\text{-}ms}(\operatorname{nf\text{-}metric\text{-}space}(\operatorname{cc\text{-}normed\text{-}field}), u), \\ \quad \quad \quad \operatorname{nf\text{-}metric\text{-}space}(\operatorname{cc\text{-}normed\text{-}field}), f, y) \wedge \\ \quad \operatorname{trace}((\lambda x \in \operatorname{ccint}(0, 1).\; (z_{0} + (x \cdot w))), 0, 1) \\ \quad \subseteq\; u \Rightarrow \\ \quad \quad \operatorname{line\text{-}int}(f, (\lambda x \in \operatorname{ccint}(0, 1).\; (z_{0} + (x \cdot w))), \\ \quad \quad \quad (\lambda x \in \operatorname{ccint}(0, 1).\; w), 0, 1) \\ \quad \quad =\; (\operatorname{line\text{-}int}(f, (\lambda x \in \operatorname{ccint}(0, 1).\; (z_{0} + (x \cdot (t \cdot w)))), (\lambda x \in \operatorname{ccint}(0, 1).\; (t \cdot w)), 0, 1) + \operatorname{line\text{-}int}(f, (\lambda x \in \operatorname{ccint}(0, 1).\; ((z_{0} + (t \cdot w)) + (x \cdot ((1 - t) \cdot w)))), (\lambda x \in \operatorname{ccint}(0, 1).\; ((1 - t) \cdot w)), 0, 1)) \end{array}}}
\vnbentry{\texttt{segment-\allowbreak{}int-\allowbreak{}abs-\allowbreak{}bound}}{\vnbstmt{\begin{array}{@{}l@{}} \forall u, f.; \forall z_{0}, w \in \mathbb{C}, n \in \mathbb{R}. \\ \quad (u \subseteq \mathbb{C}) \wedge \\ \quad (f \in (u \to \mathbb{C})) \wedge \\ \quad \forall y \in u. \\ \quad \quad \operatorname{is\text{-}continuous\text{-}at}(\operatorname{subspace\text{-}ms}(\operatorname{nf\text{-}metric\text{-}space}(\operatorname{cc\text{-}normed\text{-}field}), u), \\ \quad \quad \quad \operatorname{nf\text{-}metric\text{-}space}(\operatorname{cc\text{-}normed\text{-}field}), f, y) \wedge \\ \quad \operatorname{trace}((\lambda x \in \operatorname{ccint}(0, 1).\; (z_{0} + (x \cdot w))), 0, 1) \\ \quad \subseteq\; u \wedge \\ \quad \forall \mathit{lrpy} \in \operatorname{trace}((\lambda x \in \operatorname{ccint}(0, 1).\; (z_{0} + (x \cdot w))), 0, 1). \\ \quad \quad (\operatorname{magnitude}(f(\mathit{lrpy})) \leq n) \Rightarrow \\ \quad \quad \operatorname{magnitude}(\operatorname{line\text{-}int}(f, (\lambda x \in \operatorname{ccint}(0, 1).\; (z_{0} + (x \cdot w))), \\ \quad \quad \quad \quad (\lambda x \in \operatorname{ccint}(0, 1).\; w), 0, 1)) \\ \quad \quad \leq\; (\operatorname{magnitude}(w) \cdot n) \end{array}}}
\vnbentry{\texttt{segment-\allowbreak{}int-\allowbreak{}reverse}}{\vnbstmt{\begin{array}{@{}l@{}} \forall u, f.; \forall z_{0}, w \in \mathbb{C}. \\ \quad (u \subseteq \mathbb{C}) \wedge \\ \quad (f \in (u \to \mathbb{C})) \wedge \\ \quad \forall y \in u. \\ \quad \quad \operatorname{is\text{-}continuous\text{-}at}(\operatorname{subspace\text{-}ms}(\operatorname{nf\text{-}metric\text{-}space}(\operatorname{cc\text{-}normed\text{-}field}), u), \\ \quad \quad \quad \operatorname{nf\text{-}metric\text{-}space}(\operatorname{cc\text{-}normed\text{-}field}), f, y) \wedge \\ \quad \operatorname{trace}((\lambda x \in \operatorname{ccint}(0, 1).\; (z_{0} + (x \cdot w))), 0, 1) \\ \quad \subseteq\; u \Rightarrow \\ \quad \quad \operatorname{line\text{-}int}(f, (\lambda x \in \operatorname{ccint}(0, 1).\; ((z_{0} + w) + (x \cdot (-1 \cdot w)))), \\ \quad \quad \quad (\lambda x \in \operatorname{ccint}(0, 1).\; (-1 \cdot w)), 0, 1) \\ \quad \quad =\; (-1 \cdot \operatorname{line\text{-}int}(f, (\lambda x \in \operatorname{ccint}(0, 1).\; (z_{0} + (x \cdot w))), (\lambda x \in \operatorname{ccint}(0, 1).\; w), 0, 1)) \end{array}}}

\vnbfile{theorem-library/goursat.scm}{theorem-library/goursat.scm}
\noindent{\small goursat.scm -- CAUCHY-GOURSAT FOR A TRIANGLE AND CAUCHY ON A CONVEX SET (complex-analysis.pdf ch. 3, s. 3.2: Lemma 3.4, Lemma 3.5, Prop 3.6, Cor 3.7, Prop 3.8, Cor 3.9, all for f holomorphic on the whole open set).  Batch 30, 2026-09-25.}\par
\vnbentry{\texttt{seg-\allowbreak{}int-\allowbreak{}unfold}}{\vnbstmt{\begin{array}{@{}l@{}} \forall f, a, b. \\ \quad (\operatorname{seg\text{-}int}(f, a, b) \simeq \operatorname{line\text{-}int}(f, \operatorname{seg\text{-}path}(a, b), \operatorname{seg\text{-}deriv}(a, b), 0, 1)) \end{array}}}
\vnbentry{\texttt{tri-\allowbreak{}int-\allowbreak{}unfold}}{\vnbstmt{\begin{array}{@{}l@{}} \forall f, a, b, c. \\ \quad (\operatorname{tri\text{-}int}(f, a, b, c) \simeq ((\operatorname{seg\text{-}int}(f, a, b) + \operatorname{seg\text{-}int}(f, b, c)) + \operatorname{seg\text{-}int}(f, c, a))) \end{array}}}
\vnbentry{\texttt{cc-\allowbreak{}mid-\allowbreak{}unfold}}{\vnbstmt{\begin{array}{@{}l@{}} \forall x, y. \\ \quad (\operatorname{cc\text{-}mid}(x, y) \simeq ((x + y) \cdot \operatorname{recip}(2))) \end{array}}}
\vnbentry{\texttt{seg-\allowbreak{}path-\allowbreak{}bridge}}{\vnbstmt{\begin{array}{@{}l@{}} \forall a, b, p, w. \\ \quad ((a = p) \wedge ((b - a) = w)) \Rightarrow  \\ \quad \quad (\operatorname{seg\text{-}path}(a, b) \simeq (\lambda x \in \operatorname{ccint}(0, 1).\; (p + (x \cdot w)))) \end{array}}}
\vnbentry{\texttt{seg-\allowbreak{}deriv-\allowbreak{}bridge}}{\vnbstmt{\begin{array}{@{}l@{}} \forall a, b, w. \\ \quad ((b - a) = w) \Rightarrow  \\ \quad \quad (\operatorname{seg\text{-}deriv}(a, b) \simeq (\lambda x \in \operatorname{ccint}(0, 1).\; w)) \end{array}}}
\vnbentry{\texttt{seg-\allowbreak{}int-\allowbreak{}bridge}}{\vnbstmt{\begin{array}{@{}l@{}} \forall f, a, b, p, w. \\ \quad ((a = p) \wedge ((b - a) = w)) \Rightarrow  \\ \quad \quad (\operatorname{seg\text{-}int}(f, a, b) \simeq \operatorname{line\text{-}int}(f, (\lambda x \in \operatorname{ccint}(0, 1).\; (p + (x \cdot w))), (\lambda x \in \operatorname{ccint}(0, 1).\; w), 0, 1)) \end{array}}}
\vnbentry{\texttt{seg-\allowbreak{}trace-\allowbreak{}bridge}}{\vnbstmt{\begin{array}{@{}l@{}} \forall a, b, p, w. \\ \quad ((a = p) \wedge ((b - a) = w)) \Rightarrow  \\ \quad \quad (\operatorname{trace}(\operatorname{seg\text{-}path}(a, b), 0, 1) \simeq \operatorname{trace}((\lambda x \in \operatorname{ccint}(0, 1).\; (p + (x \cdot w))), 0, 1)) \end{array}}}
\vnbentry{\texttt{seg-\allowbreak{}path-\allowbreak{}is-\allowbreak{}road}}{\vnbstmt{\begin{array}{@{}l@{}} \forall a, b \in \mathbb{C}. \\ \quad \operatorname{is\text{-}road}(\operatorname{seg\text{-}path}(a, b), \operatorname{seg\text{-}deriv}(a, b), 0, 1) \end{array}}}
\vnbentry{\texttt{in-\allowbreak{}transport-\allowbreak{}eq}}{\vnbstmt{\forall x, y, s.\; (((x \in s) \wedge (x = y)) \Rightarrow (y \in s))}}
\vnbentry{\texttt{seg-\allowbreak{}path-\allowbreak{}apply}}{\vnbstmt{\begin{array}{@{}l@{}} \forall a, b \in \mathbb{C}, t \in \operatorname{ccint}(0, 1). \\ \quad (\operatorname{seg\text{-}path}(a, b)(t) = (a + (t \cdot (b - a)))) \end{array}}}
\vnbentry{\texttt{seg-\allowbreak{}trace-\allowbreak{}intro}}{\vnbstmt{\begin{array}{@{}l@{}} \forall a, b \in \mathbb{C}, t \in \operatorname{ccint}(0, 1). \\ \quad ((a + (t \cdot (b - a))) \in \operatorname{trace}(\operatorname{seg\text{-}path}(a, b), 0, 1)) \end{array}}}
\vnbentry{\texttt{seg-\allowbreak{}trace-\allowbreak{}elim}}{\vnbstmt{\begin{array}{@{}l@{}} \forall a, b \in \mathbb{C}, z \in \operatorname{trace}(\operatorname{seg\text{-}path}(a, b), 0, 1). \exists t \in \operatorname{ccint}(0, 1). \\ \quad (z = (a + (t \cdot (b - a)))) \end{array}}}
\vnbentry{\texttt{conv3-\allowbreak{}member-\allowbreak{}iff}}{\vnbstmt{\begin{array}{@{}l@{}} \forall a, b, c, z. \\ \quad ((z \in \operatorname{conv3}(a, b, c)) \Leftrightarrow ((z \in \mathbb{C}) \wedge \exists l, m, n \in \operatorname{ccint}(0, 1).\; ((((l + m) + n) = 1) \wedge (z = (((l \cdot a) + (m \cdot b)) + (n \cdot c)))))) \end{array}}}
\vnbentry{\texttt{conv3-\allowbreak{}intro}}{\vnbstmt{\begin{array}{@{}l@{}} \forall a, b, c \in \mathbb{C}, l, m, n \in \operatorname{ccint}(0, 1), z. \\ \quad ((((l + m) + n) = 1) \wedge (z = (((l \cdot a) + (m \cdot b)) + (n \cdot c)))) \Rightarrow  \\ \quad \quad (z \in \operatorname{conv3}(a, b, c)) \end{array}}}
\vnbentry{\texttt{is-\allowbreak{}convex-\allowbreak{}unfold}}{\vnbstmt{\begin{array}{@{}l@{}} \forall u.; \forall x, y \in u, t \in \operatorname{ccint}(0, 1). \\ \quad (\operatorname{is\text{-}convex}(u) \Rightarrow ((x + (t \cdot (y - x))) \in u)) \end{array}}}
\vnbentry{\texttt{conv3-\allowbreak{}is-\allowbreak{}convex}}{\vnbstmt{\begin{array}{@{}l@{}} \forall a, b, c \in \mathbb{C}. \\ \quad \operatorname{is\text{-}convex}(\operatorname{conv3}(a, b, c)) \end{array}}}
\vnbentry{\texttt{seg-\allowbreak{}trace-\allowbreak{}in-\allowbreak{}convex}}{\vnbstmt{\begin{array}{@{}l@{}} \forall u.; \forall a, b \in \mathbb{C}. \\ \quad (\operatorname{is\text{-}convex}(u) \wedge ((a \in u) \wedge (b \in u))) \Rightarrow  \\ \quad \quad (\operatorname{trace}(\operatorname{seg\text{-}path}(a, b), 0, 1) \subseteq u) \end{array}}}
\vnbentry{\texttt{conv3-\allowbreak{}in-\allowbreak{}convex}}{\vnbstmt{\begin{array}{@{}l@{}} \forall u.; \forall a, b, c \in \mathbb{C}. \\ \quad (\operatorname{is\text{-}convex}(u) \wedge ((a \in u) \wedge ((b \in u) \wedge (c \in u)))) \Rightarrow  \\ \quad \quad (\operatorname{conv3}(a, b, c) \subseteq u) \end{array}}}
\vnbentry{\texttt{conv3-\allowbreak{}vertices-\allowbreak{}in}}{\vnbstmt{\begin{array}{@{}l@{}} \forall a, b, c \in \mathbb{C}. \\ \quad (a \in \operatorname{conv3}(a, b, c)) \wedge \\ \quad (b \in \operatorname{conv3}(a, b, c)) \wedge \\ \quad (c \in \operatorname{conv3}(a, b, c)) \end{array}}}
\vnbentry{\texttt{conv3-\allowbreak{}seg-\allowbreak{}trace}}{\vnbstmt{\begin{array}{@{}l@{}} \forall a, b, c \in \mathbb{C}, x, y \in \operatorname{conv3}(a, b, c). \\ \quad \operatorname{trace}(\operatorname{seg\text{-}path}(x, y), 0, 1) \\ \quad \subseteq\; \operatorname{conv3}(a, b, c) \end{array}}}
\vnbentry{\texttt{conv3-\allowbreak{}sub-\allowbreak{}hull}}{\vnbstmt{\begin{array}{@{}l@{}} \forall a, b, c \in \mathbb{C}, x, y, z \in \operatorname{conv3}(a, b, c). \\ \quad (\operatorname{conv3}(x, y, z) \subseteq \operatorname{conv3}(a, b, c)) \end{array}}}
\vnbentry{\texttt{cc-\allowbreak{}mid-\allowbreak{}affine}}{\vnbstmt{\begin{array}{@{}l@{}} \forall x, y \in \mathbb{C}. \\ \quad (\operatorname{cc\text{-}mid}(x, y) = (x + (1/2 \cdot (y - x)))) \wedge \\ \quad (\operatorname{cc\text{-}mid}(x, y) = (y + (1/2 \cdot (x - y)))) \end{array}}}
\vnbentry{\texttt{cc-\allowbreak{}mid-\allowbreak{}in-\allowbreak{}seg-\allowbreak{}trace}}{\vnbstmt{\begin{array}{@{}l@{}} \forall x, y \in \mathbb{C}. \\ \quad \operatorname{cc\text{-}mid}(x, y) \\ \quad \in\; \operatorname{trace}(\operatorname{seg\text{-}path}(x, y), 0, 1) \wedge \\ \quad \operatorname{cc\text{-}mid}(x, y) \\ \quad \in\; \operatorname{trace}(\operatorname{seg\text{-}path}(y, x), 0, 1) \end{array}}}
\vnbentry{\texttt{conv3-\allowbreak{}mid-\allowbreak{}in}}{\vnbstmt{\begin{array}{@{}l@{}} \forall a, b, c \in \mathbb{C}, x, y \in \operatorname{conv3}(a, b, c). \\ \quad (\operatorname{cc\text{-}mid}(x, y) \in \operatorname{conv3}(a, b, c)) \end{array}}}
\vnbentry{\texttt{seg-\allowbreak{}trace-\allowbreak{}in-\allowbreak{}conv3}}{\vnbstmt{\begin{array}{@{}l@{}} \forall a, b, c \in \mathbb{C}. \\ \quad \operatorname{trace}(\operatorname{seg\text{-}path}(a, b), 0, 1) \\ \quad \subseteq\; \operatorname{conv3}(a, b, c) \wedge \\ \quad \operatorname{trace}(\operatorname{seg\text{-}path}(b, c), 0, 1) \\ \quad \subseteq\; \operatorname{conv3}(a, b, c) \wedge \\ \quad \operatorname{trace}(\operatorname{seg\text{-}path}(c, a), 0, 1) \\ \quad \subseteq\; \operatorname{conv3}(a, b, c) \end{array}}}
\vnbentry{\texttt{conv3-\allowbreak{}sub-\allowbreak{}triangle}}{\vnbstmt{\begin{array}{@{}l@{}} \forall a, b, c \in \mathbb{C}. \\ \quad \operatorname{conv3}(a, \operatorname{cc\text{-}mid}(a, b), \operatorname{cc\text{-}mid}(a, c)) \\ \quad \subseteq\; \operatorname{conv3}(a, b, c) \wedge \\ \quad \operatorname{conv3}(\operatorname{cc\text{-}mid}(a, b), b, \operatorname{cc\text{-}mid}(b, c)) \\ \quad \subseteq\; \operatorname{conv3}(a, b, c) \wedge \\ \quad \operatorname{conv3}(\operatorname{cc\text{-}mid}(a, c), \operatorname{cc\text{-}mid}(b, c), c) \\ \quad \subseteq\; \operatorname{conv3}(a, b, c) \wedge \\ \quad \operatorname{conv3}(\operatorname{cc\text{-}mid}(a, b), \operatorname{cc\text{-}mid}(b, c), \operatorname{cc\text{-}mid}(a, c)) \\ \quad \subseteq\; \operatorname{conv3}(a, b, c) \end{array}}}
\vnbentry{\texttt{seg-\allowbreak{}int-\allowbreak{}in-\allowbreak{}cc}}{\vnbstmt{\begin{array}{@{}l@{}} \forall u, f.; \forall a, b \in \mathbb{C}. \\ \quad (u \subseteq \mathbb{C}) \wedge \\ \quad (f \in (u \to \mathbb{C})) \wedge \\ \quad \forall y \in u. \\ \quad \quad \operatorname{is\text{-}continuous\text{-}at}(\operatorname{subspace\text{-}ms}(\operatorname{nf\text{-}metric\text{-}space}(\operatorname{cc\text{-}normed\text{-}field}), u), \\ \quad \quad \quad \operatorname{nf\text{-}metric\text{-}space}(\operatorname{cc\text{-}normed\text{-}field}), f, y) \wedge \\ \quad (\operatorname{trace}(\operatorname{seg\text{-}path}(a, b), 0, 1) \subseteq u) \Rightarrow \\ \quad \quad (\operatorname{seg\text{-}int}(f, a, b) \in \mathbb{C}) \end{array}}}
\vnbentry{\texttt{seg-\allowbreak{}int-\allowbreak{}degenerate}}{\vnbstmt{\begin{array}{@{}l@{}} \forall u, f.; \forall z \in u. \\ \quad (u \subseteq \mathbb{C}) \wedge \\ \quad (f \in (u \to \mathbb{C})) \wedge \\ \quad \forall y \in u. \\ \quad \quad \operatorname{is\text{-}continuous\text{-}at}(\operatorname{subspace\text{-}ms}(\operatorname{nf\text{-}metric\text{-}space}(\operatorname{cc\text{-}normed\text{-}field}), u), \\ \quad \quad \quad \operatorname{nf\text{-}metric\text{-}space}(\operatorname{cc\text{-}normed\text{-}field}), f, y) \Rightarrow \\ \quad \quad (\operatorname{seg\text{-}int}(f, z, z) = 0) \end{array}}}
\vnbentry{\texttt{seg-\allowbreak{}int-\allowbreak{}reverse-\allowbreak{}at}}{\vnbstmt{\begin{array}{@{}l@{}} \forall u, f.; \forall a, b \in \mathbb{C}. \\ \quad (u \subseteq \mathbb{C}) \wedge \\ \quad (f \in (u \to \mathbb{C})) \wedge \\ \quad \forall y \in u. \\ \quad \quad \operatorname{is\text{-}continuous\text{-}at}(\operatorname{subspace\text{-}ms}(\operatorname{nf\text{-}metric\text{-}space}(\operatorname{cc\text{-}normed\text{-}field}), u), \\ \quad \quad \quad \operatorname{nf\text{-}metric\text{-}space}(\operatorname{cc\text{-}normed\text{-}field}), f, y) \wedge \\ \quad (\operatorname{trace}(\operatorname{seg\text{-}path}(a, b), 0, 1) \subseteq u) \Rightarrow \\ \quad \quad (\operatorname{seg\text{-}int}(f, b, a) = (-1 \cdot \operatorname{seg\text{-}int}(f, a, b))) \end{array}}}
\vnbentry{\texttt{seg-\allowbreak{}int-\allowbreak{}split-\allowbreak{}at}}{\vnbstmt{\begin{array}{@{}l@{}} \forall u, f.; \forall a, b \in \mathbb{C}, c \in \operatorname{trace}(\operatorname{seg\text{-}path}(a, b), 0, 1). \\ \quad (u \subseteq \mathbb{C}) \wedge \\ \quad (f \in (u \to \mathbb{C})) \wedge \\ \quad \forall y \in u. \\ \quad \quad \operatorname{is\text{-}continuous\text{-}at}(\operatorname{subspace\text{-}ms}(\operatorname{nf\text{-}metric\text{-}space}(\operatorname{cc\text{-}normed\text{-}field}), u), \\ \quad \quad \quad \operatorname{nf\text{-}metric\text{-}space}(\operatorname{cc\text{-}normed\text{-}field}), f, y) \wedge \\ \quad (\operatorname{trace}(\operatorname{seg\text{-}path}(a, b), 0, 1) \subseteq u) \Rightarrow \\ \quad \quad \operatorname{seg\text{-}int}(f, a, b) \\ \quad \quad =\; (\operatorname{seg\text{-}int}(f, a, c) + \operatorname{seg\text{-}int}(f, c, b)) \end{array}}}
\vnbentry{\texttt{tri-\allowbreak{}int-\allowbreak{}subdivision}}{\vnbstmt{\begin{array}{@{}l@{}} \forall u, f.; \forall a, b, c \in \mathbb{C}. \\ \quad (u \subseteq \mathbb{C}) \wedge \\ \quad (f \in (u \to \mathbb{C})) \wedge \\ \quad \forall y \in u. \\ \quad \quad \operatorname{is\text{-}continuous\text{-}at}(\operatorname{subspace\text{-}ms}(\operatorname{nf\text{-}metric\text{-}space}(\operatorname{cc\text{-}normed\text{-}field}), u), \\ \quad \quad \quad \operatorname{nf\text{-}metric\text{-}space}(\operatorname{cc\text{-}normed\text{-}field}), f, y) \wedge \\ \quad (\operatorname{conv3}(a, b, c) \subseteq u) \Rightarrow \\ \quad \quad \operatorname{tri\text{-}int}(f, a, b, c) \\ \quad \quad =\; (((\operatorname{tri\text{-}int}(f, a, \operatorname{cc\text{-}mid}(a, b), \operatorname{cc\text{-}mid}(a, c)) + \operatorname{tri\text{-}int}(f, \operatorname{cc\text{-}mid}(a, b), b, \operatorname{cc\text{-}mid}(b, c))) + \operatorname{tri\text{-}int}(f, \operatorname{cc\text{-}mid}(a, c), \operatorname{cc\text{-}mid}(b, c), c)) + \operatorname{tri\text{-}int}(f, \operatorname{cc\text{-}mid}(a, b), \operatorname{cc\text{-}mid}(b, c), \operatorname{cc\text{-}mid}(a, c))) \end{array}}}
\vnbentry{\texttt{tri-\allowbreak{}int-\allowbreak{}abs-\allowbreak{}bound}}{\vnbstmt{\begin{array}{@{}l@{}} \forall u, f.; \forall a, b, c \in \mathbb{C}, m \in \mathbb{R}. \\ \quad (u \subseteq \mathbb{C}) \wedge \\ \quad (f \in (u \to \mathbb{C})) \wedge \\ \quad \forall y \in u. \\ \quad \quad \operatorname{is\text{-}continuous\text{-}at}(\operatorname{subspace\text{-}ms}(\operatorname{nf\text{-}metric\text{-}space}(\operatorname{cc\text{-}normed\text{-}field}), u), \\ \quad \quad \quad \operatorname{nf\text{-}metric\text{-}space}(\operatorname{cc\text{-}normed\text{-}field}), f, y) \wedge \\ \quad (\operatorname{conv3}(a, b, c) \subseteq u) \wedge \\ \quad \forall z \in \operatorname{conv3}(a, b, c). \\ \quad \quad (\operatorname{magnitude}(f(z)) \leq m) \Rightarrow \\ \quad \quad \operatorname{magnitude}(\operatorname{tri\text{-}int}(f, a, b, c)) \\ \quad \quad \leq\; (((\operatorname{magnitude}((b - a)) + \operatorname{magnitude}((c - b))) + \operatorname{magnitude}((a - c))) \cdot m) \end{array}}}
\vnbentry{\texttt{tri-\allowbreak{}int-\allowbreak{}sum}}{\vnbstmt{\begin{array}{@{}l@{}} \forall u, f, g.; \forall h \in (u \to \mathbb{C}), a, b, c \in \mathbb{C}. \\ \quad (u \subseteq \mathbb{C}) \wedge \\ \quad (f \in (u \to \mathbb{C})) \wedge \\ \quad \forall y \in u. \\ \quad \quad \operatorname{is\text{-}continuous\text{-}at}(\operatorname{subspace\text{-}ms}(\operatorname{nf\text{-}metric\text{-}space}(\operatorname{cc\text{-}normed\text{-}field}), u), \\ \quad \quad \quad \operatorname{nf\text{-}metric\text{-}space}(\operatorname{cc\text{-}normed\text{-}field}), f, y) \wedge \\ \quad (g \in (u \to \mathbb{C})) \wedge \\ \quad \forall y \in u. \\ \quad \quad \operatorname{is\text{-}continuous\text{-}at}(\operatorname{subspace\text{-}ms}(\operatorname{nf\text{-}metric\text{-}space}(\operatorname{cc\text{-}normed\text{-}field}), u), \\ \quad \quad \quad \operatorname{nf\text{-}metric\text{-}space}(\operatorname{cc\text{-}normed\text{-}field}), g, y) \wedge \\ \quad \forall z \in u.\; (h(z) \simeq (f(z) + g(z))) \wedge \\ \quad (\operatorname{conv3}(a, b, c) \subseteq u) \Rightarrow \\ \quad \quad \operatorname{tri\text{-}int}(h, a, b, c) \\ \quad \quad =\; (\operatorname{tri\text{-}int}(f, a, b, c) + \operatorname{tri\text{-}int}(g, a, b, c)) \end{array}}}
\vnbentry{\texttt{cc-\allowbreak{}open-\allowbreak{}subset-\allowbreak{}cc}}{\vnbstmt{\begin{array}{@{}l@{}} \forall u. \\ \quad \operatorname{is\text{-}open}(\operatorname{nf\text{-}metric\text{-}space}(\operatorname{cc\text{-}normed\text{-}field}), u) \Rightarrow \\ \quad \quad (u \subseteq \mathbb{C}) \end{array}}}
\vnbentry{\texttt{affine-\allowbreak{}primitive-\allowbreak{}on}}{\vnbstmt{\begin{array}{@{}l@{}} \forall u.; \forall p, q, w \in \mathbb{C}, f \in (u \to \mathbb{C}). \\ \quad \operatorname{is\text{-}open}(\operatorname{nf\text{-}metric\text{-}space}(\operatorname{cc\text{-}normed\text{-}field}), u) \wedge \\ \quad \forall z \in u. \\ \quad \quad (f(z) \simeq ((p \cdot z) + ((q \cdot ((z - w) \cdot (z - w))) \cdot \operatorname{recip}(2)))) \Rightarrow \\ \quad \quad \operatorname{holomorphic\text{-}on}(u, f) \wedge \\ \quad \quad \forall z \in u. \\ \quad \quad \quad \operatorname{is\text{-}diff\text{-}on}(\operatorname{cc\text{-}normed\text{-}field}, u, f, z, (p + (q \cdot (z - w)))) \end{array}}}
\vnbentry{\texttt{tri-\allowbreak{}int-\allowbreak{}affine-\allowbreak{}zero}}{\vnbstmt{\begin{array}{@{}l@{}} \forall u.; \forall p, q, w \in \mathbb{C}, g \in (u \to \mathbb{C}), a, b, c \in \mathbb{C}. \\ \quad \operatorname{is\text{-}open}(\operatorname{nf\text{-}metric\text{-}space}(\operatorname{cc\text{-}normed\text{-}field}), u) \wedge \\ \quad \forall z \in u.\; (g(z) \simeq (p + (q \cdot (z - w)))) \wedge \\ \quad (\operatorname{conv3}(a, b, c) \subseteq u) \Rightarrow \\ \quad \quad (\operatorname{tri\text{-}int}(g, a, b, c) = 0) \end{array}}}
\vnbentry{\texttt{holomorphic-\allowbreak{}on-\allowbreak{}continuous}}{\vnbstmt{\begin{array}{@{}l@{}} \forall u, f.; \forall y \in u. \\ \quad \operatorname{holomorphic\text{-}on}(u, f) \Rightarrow  \\ \quad \quad \operatorname{is\text{-}continuous\text{-}at}(\operatorname{subspace\text{-}ms}(\operatorname{nf\text{-}metric\text{-}space}(\operatorname{cc\text{-}normed\text{-}field}), u), \\ \quad \quad \quad \operatorname{nf\text{-}metric\text{-}space}(\operatorname{cc\text{-}normed\text{-}field}), f, y) \end{array}}}
\vnbentry{\texttt{cc-\allowbreak{}mid-\allowbreak{}coords}}{\vnbstmt{\begin{array}{@{}l@{}} \forall x, y \in \mathbb{C}. \\ \quad \operatorname{real\text{-}part}(\operatorname{cc\text{-}mid}(x, y)) \\ \quad =\; (1/2 \cdot (\operatorname{real\text{-}part}(x) + \operatorname{real\text{-}part}(y))) \wedge \\ \quad \operatorname{imag\text{-}part}(\operatorname{cc\text{-}mid}(x, y)) \\ \quad =\; (1/2 \cdot (\operatorname{imag\text{-}part}(x) + \operatorname{imag\text{-}part}(y))) \end{array}}}
\vnbentry{\texttt{std-\allowbreak{}mid}}{\vnbstmt{\begin{array}{@{}l@{}} \forall x, y \in \mathbb{C}. \\ \quad (0 \leq \operatorname{real\text{-}part}(x)) \wedge \\ \quad (0 \leq \operatorname{imag\text{-}part}(x)) \wedge \\ \quad ((\operatorname{real\text{-}part}(x) + \operatorname{imag\text{-}part}(x)) \leq 1) \wedge \\ \quad (0 \leq \operatorname{real\text{-}part}(y)) \wedge \\ \quad (0 \leq \operatorname{imag\text{-}part}(y)) \wedge \\ \quad ((\operatorname{real\text{-}part}(y) + \operatorname{imag\text{-}part}(y)) \leq 1) \Rightarrow \\ \quad \quad (0 \leq \operatorname{real\text{-}part}(\operatorname{cc\text{-}mid}(x, y))) \wedge \\ \quad \quad (0 \leq \operatorname{imag\text{-}part}(\operatorname{cc\text{-}mid}(x, y))) \wedge \\ \quad \quad (\operatorname{real\text{-}part}(\operatorname{cc\text{-}mid}(x, y)) + \operatorname{imag\text{-}part}(\operatorname{cc\text{-}mid}(x, y))) \\ \quad \quad \leq\; 1 \end{array}}}
\vnbentry{\texttt{std-\allowbreak{}dist}}{\vnbstmt{\begin{array}{@{}l@{}} \forall x, y \in \mathbb{C}. \\ \quad (0 \leq \operatorname{real\text{-}part}(x)) \wedge \\ \quad (0 \leq \operatorname{imag\text{-}part}(x)) \wedge \\ \quad ((\operatorname{real\text{-}part}(x) + \operatorname{imag\text{-}part}(x)) \leq 1) \wedge \\ \quad (0 \leq \operatorname{real\text{-}part}(y)) \wedge \\ \quad (0 \leq \operatorname{imag\text{-}part}(y)) \wedge \\ \quad ((\operatorname{real\text{-}part}(y) + \operatorname{imag\text{-}part}(y)) \leq 1) \Rightarrow \\ \quad \quad (\operatorname{magnitude}((x - y)) \leq 2) \end{array}}}
\vnbentry{\texttt{std-\allowbreak{}limit}}{\vnbstmt{\begin{array}{@{}l@{}} \forall s \in (\mathbb{N} \to \mathbb{C}), w \in \mathbb{C}. \\ \quad \forall n \in \mathbb{N}. \\ \quad \quad (0 \leq \operatorname{real\text{-}part}(s(n))) \wedge \\ \quad \quad (0 \leq \operatorname{imag\text{-}part}(s(n))) \wedge \\ \quad \quad ((\operatorname{real\text{-}part}(s(n)) + \operatorname{imag\text{-}part}(s(n))) \leq 1) \wedge \\ \quad \operatorname{converges\text{-}to}(\operatorname{cc\text{-}ms}, s, w) \Rightarrow \\ \quad \quad (0 \leq \operatorname{real\text{-}part}(w)) \wedge \\ \quad \quad (0 \leq \operatorname{imag\text{-}part}(w)) \wedge \\ \quad \quad ((\operatorname{real\text{-}part}(w) + \operatorname{imag\text{-}part}(w)) \leq 1) \end{array}}}
\vnbentry{\texttt{bary-\allowbreak{}mid}}{\vnbstmt{\begin{array}{@{}l@{}} \forall a, b, c, x, y \in \mathbb{C}. \\ \quad (((\operatorname{real\text{-}part}(\operatorname{cc\text{-}mid}(x, y)) \cdot a) + (\operatorname{imag\text{-}part}(\operatorname{cc\text{-}mid}(x, y)) \cdot b)) + (((1 - \operatorname{real\text{-}part}(\operatorname{cc\text{-}mid}(x, y))) - \operatorname{imag\text{-}part}(\operatorname{cc\text{-}mid}(x, y))) \cdot c)) \\ \quad =\; \operatorname{cc\text{-}mid}((((\operatorname{real\text{-}part}(x) \cdot a) + (\operatorname{imag\text{-}part}(x) \cdot b)) + (((1 - \operatorname{real\text{-}part}(x)) - \operatorname{imag\text{-}part}(x)) \cdot c)), \\ \quad \quad (((\operatorname{real\text{-}part}(y) \cdot a) + (\operatorname{imag\text{-}part}(y) \cdot b)) + (((1 - \operatorname{real\text{-}part}(y)) - \operatorname{imag\text{-}part}(y)) \cdot c))) \end{array}}}
\vnbentry{\texttt{bary-\allowbreak{}in-\allowbreak{}conv3}}{\vnbstmt{\begin{array}{@{}l@{}} \forall a, b, c, x \in \mathbb{C}. \\ \quad (0 \leq \operatorname{real\text{-}part}(x)) \wedge \\ \quad (0 \leq \operatorname{imag\text{-}part}(x)) \wedge \\ \quad ((\operatorname{real\text{-}part}(x) + \operatorname{imag\text{-}part}(x)) \leq 1) \Rightarrow \\ \quad \quad (((\operatorname{real\text{-}part}(x) \cdot a) + (\operatorname{imag\text{-}part}(x) \cdot b)) + (((1 - \operatorname{real\text{-}part}(x)) - \operatorname{imag\text{-}part}(x)) \cdot c)) \\ \quad \quad \in\; \operatorname{conv3}(a, b, c) \end{array}}}
\vnbentry{\texttt{bary-\allowbreak{}vertices}}{\vnbstmt{\begin{array}{@{}l@{}} \forall a, b, c \in \mathbb{C}. \\ \quad (((\operatorname{real\text{-}part}(1) \cdot a) + (\operatorname{imag\text{-}part}(1) \cdot b)) + (((1 - \operatorname{real\text{-}part}(1)) - \operatorname{imag\text{-}part}(1)) \cdot c)) \\ \quad =\; a \wedge \\ \quad (((\operatorname{real\text{-}part}(+i) \cdot a) + (\operatorname{imag\text{-}part}(+i) \cdot b)) + (((1 - \operatorname{real\text{-}part}(+i)) - \operatorname{imag\text{-}part}(+i)) \cdot c)) \\ \quad =\; b \wedge \\ \quad (((\operatorname{real\text{-}part}(0) \cdot a) + (\operatorname{imag\text{-}part}(0) \cdot b)) + (((1 - \operatorname{real\text{-}part}(0)) - \operatorname{imag\text{-}part}(0)) \cdot c)) \\ \quad =\; c \end{array}}}
\vnbentry{\texttt{bary-\allowbreak{}lip}}{\vnbstmt{\begin{array}{@{}l@{}} \forall a, b, c, x, y \in \mathbb{C}. \\ \quad \operatorname{magnitude}(((((\operatorname{real\text{-}part}(x) \cdot a) + (\operatorname{imag\text{-}part}(x) \cdot b)) + (((1 - \operatorname{real\text{-}part}(x)) - \operatorname{imag\text{-}part}(x)) \cdot c)) - (((\operatorname{real\text{-}part}(y) \cdot a) + (\operatorname{imag\text{-}part}(y) \cdot b)) + (((1 - \operatorname{real\text{-}part}(y)) - \operatorname{imag\text{-}part}(y)) \cdot c)))) \\ \quad \leq\; ((\operatorname{magnitude}((a - c)) + \operatorname{magnitude}((b - c))) \cdot \operatorname{magnitude}((x - y))) \end{array}}}
\vnbentry{\texttt{half-\allowbreak{}magnitude}}{\vnbstmt{\begin{array}{@{}l@{}} \forall x \in \mathbb{C}. \\ \quad (\operatorname{magnitude}((1/2 \cdot x)) = (1/2 \cdot \operatorname{magnitude}(x))) \end{array}}}
\vnbentry{\texttt{tri-\allowbreak{}child-\allowbreak{}geom-\allowbreak{}1}}{\vnbstmt{\begin{array}{@{}l@{}} \forall x, y, z \in \mathbb{C}. \\ \quad ((\operatorname{magnitude}((\operatorname{cc\text{-}mid}(x, y) - x)) + \operatorname{magnitude}((\operatorname{cc\text{-}mid}(x, z) - \operatorname{cc\text{-}mid}(x, y)))) + \operatorname{magnitude}((x - \operatorname{cc\text{-}mid}(x, z)))) \\ \quad =\; (1/2 \cdot ((\operatorname{magnitude}((y - x)) + \operatorname{magnitude}((z - y))) + \operatorname{magnitude}((x - z)))) \wedge \\ \quad \operatorname{magnitude}((x - x)) \\ \quad \leq\; ((\operatorname{magnitude}((y - x)) + \operatorname{magnitude}((z - y))) + \operatorname{magnitude}((x - z))) \end{array}}}
\vnbentry{\texttt{tri-\allowbreak{}child-\allowbreak{}geom-\allowbreak{}2}}{\vnbstmt{\begin{array}{@{}l@{}} \forall x, y, z \in \mathbb{C}. \\ \quad ((\operatorname{magnitude}((y - \operatorname{cc\text{-}mid}(x, y))) + \operatorname{magnitude}((\operatorname{cc\text{-}mid}(y, z) - y))) + \operatorname{magnitude}((\operatorname{cc\text{-}mid}(x, y) - \operatorname{cc\text{-}mid}(y, z)))) \\ \quad =\; (1/2 \cdot ((\operatorname{magnitude}((y - x)) + \operatorname{magnitude}((z - y))) + \operatorname{magnitude}((x - z)))) \wedge \\ \quad \operatorname{magnitude}((\operatorname{cc\text{-}mid}(x, y) - x)) \\ \quad \leq\; ((\operatorname{magnitude}((y - x)) + \operatorname{magnitude}((z - y))) + \operatorname{magnitude}((x - z))) \end{array}}}
\vnbentry{\texttt{tri-\allowbreak{}child-\allowbreak{}geom-\allowbreak{}3}}{\vnbstmt{\begin{array}{@{}l@{}} \forall x, y, z \in \mathbb{C}. \\ \quad ((\operatorname{magnitude}((\operatorname{cc\text{-}mid}(y, z) - \operatorname{cc\text{-}mid}(x, z))) + \operatorname{magnitude}((z - \operatorname{cc\text{-}mid}(y, z)))) + \operatorname{magnitude}((\operatorname{cc\text{-}mid}(x, z) - z))) \\ \quad =\; (1/2 \cdot ((\operatorname{magnitude}((y - x)) + \operatorname{magnitude}((z - y))) + \operatorname{magnitude}((x - z)))) \wedge \\ \quad \operatorname{magnitude}((\operatorname{cc\text{-}mid}(x, z) - x)) \\ \quad \leq\; ((\operatorname{magnitude}((y - x)) + \operatorname{magnitude}((z - y))) + \operatorname{magnitude}((x - z))) \end{array}}}
\vnbentry{\texttt{tri-\allowbreak{}child-\allowbreak{}geom-\allowbreak{}4}}{\vnbstmt{\begin{array}{@{}l@{}} \forall x, y, z \in \mathbb{C}. \\ \quad ((\operatorname{magnitude}((\operatorname{cc\text{-}mid}(y, z) - \operatorname{cc\text{-}mid}(x, y))) + \operatorname{magnitude}((\operatorname{cc\text{-}mid}(x, z) - \operatorname{cc\text{-}mid}(y, z)))) + \operatorname{magnitude}((\operatorname{cc\text{-}mid}(x, y) - \operatorname{cc\text{-}mid}(x, z)))) \\ \quad =\; (1/2 \cdot ((\operatorname{magnitude}((y - x)) + \operatorname{magnitude}((z - y))) + \operatorname{magnitude}((x - z)))) \wedge \\ \quad \operatorname{magnitude}((\operatorname{cc\text{-}mid}(x, y) - x)) \\ \quad \leq\; ((\operatorname{magnitude}((y - x)) + \operatorname{magnitude}((z - y))) + \operatorname{magnitude}((x - z))) \end{array}}}
\vnbentry{\texttt{tri-\allowbreak{}int-\allowbreak{}in-\allowbreak{}cc}}{\vnbstmt{\begin{array}{@{}l@{}} \forall u, f.; \forall a, b, c \in \mathbb{C}. \\ \quad (u \subseteq \mathbb{C}) \wedge \\ \quad (f \in (u \to \mathbb{C})) \wedge \\ \quad \forall y \in u. \\ \quad \quad \operatorname{is\text{-}continuous\text{-}at}(\operatorname{subspace\text{-}ms}(\operatorname{nf\text{-}metric\text{-}space}(\operatorname{cc\text{-}normed\text{-}field}), u), \\ \quad \quad \quad \operatorname{nf\text{-}metric\text{-}space}(\operatorname{cc\text{-}normed\text{-}field}), f, y) \wedge \\ \quad (\operatorname{conv3}(a, b, c) \subseteq u) \Rightarrow \\ \quad \quad (\operatorname{tri\text{-}int}(f, a, b, c) \in \mathbb{C}) \end{array}}}
\vnbentry{\texttt{goursat-\allowbreak{}nest}}{\vnbstmt{\begin{array}{@{}l@{}} \forall u, f.; \forall a, b, c \in \mathbb{C}. \\ \quad (u \subseteq \mathbb{C}) \wedge \\ \quad (f \in (u \to \mathbb{C})) \wedge \\ \quad \forall y \in u. \\ \quad \quad \operatorname{is\text{-}continuous\text{-}at}(\operatorname{subspace\text{-}ms}(\operatorname{nf\text{-}metric\text{-}space}(\operatorname{cc\text{-}normed\text{-}field}), u), \\ \quad \quad \quad \operatorname{nf\text{-}metric\text{-}space}(\operatorname{cc\text{-}normed\text{-}field}), f, y) \wedge \\ \quad (\operatorname{conv3}(a, b, c) \subseteq u) \Rightarrow \\ \quad \quad \exists w \in \mathbb{C}. \\ \quad \quad \quad (0 \leq \operatorname{real\text{-}part}(w)) \wedge \\ \quad \quad \quad (0 \leq \operatorname{imag\text{-}part}(w)) \wedge \\ \quad \quad \quad ((\operatorname{real\text{-}part}(w) + \operatorname{imag\text{-}part}(w)) \leq 1) \wedge \\ \quad \quad \quad \forall n \in \mathbb{N}. \exists x, \mathit{gouy}, z \in \mathbb{C}. \\ \quad \quad \quad \quad (0 \leq \operatorname{real\text{-}part}(x)) \wedge \\ \quad \quad \quad \quad (0 \leq \operatorname{imag\text{-}part}(x)) \wedge \\ \quad \quad \quad \quad ((\operatorname{real\text{-}part}(x) + \operatorname{imag\text{-}part}(x)) \leq 1) \wedge \\ \quad \quad \quad \quad (0 \leq \operatorname{real\text{-}part}(\mathit{gouy})) \wedge \\ \quad \quad \quad \quad (0 \leq \operatorname{imag\text{-}part}(\mathit{gouy})) \wedge \\ \quad \quad \quad \quad (\operatorname{real\text{-}part}(\mathit{gouy}) + \operatorname{imag\text{-}part}(\mathit{gouy})) \\ \quad \quad \quad \quad \leq\; 1 \wedge \\ \quad \quad \quad \quad (0 \leq \operatorname{real\text{-}part}(z)) \wedge \\ \quad \quad \quad \quad (0 \leq \operatorname{imag\text{-}part}(z)) \wedge \\ \quad \quad \quad \quad ((\operatorname{real\text{-}part}(z) + \operatorname{imag\text{-}part}(z)) \leq 1) \wedge \\ \quad \quad \quad \quad ((\operatorname{magnitude}((\mathit{gouy} - x)) + \operatorname{magnitude}((z - \mathit{gouy}))) + \operatorname{magnitude}((x - z))) \\ \quad \quad \quad \quad \leq\; (6 \cdot \operatorname{recip}(\operatorname{power}(2, n))) \wedge \\ \quad \quad \quad \quad (\operatorname{magnitude}((x - w)) \leq (12 \cdot \operatorname{recip}(\operatorname{power}(2, n)))) \wedge \\ \quad \quad \quad \quad \operatorname{magnitude}(\operatorname{tri\text{-}int}(f, a, b, c)) \\ \quad \quad \quad \quad \leq\; (\operatorname{power}(4, n) \cdot \operatorname{magnitude}(\operatorname{tri\text{-}int}(f, (((\operatorname{real\text{-}part}(x) \cdot a) + (\operatorname{imag\text{-}part}(x) \cdot b)) + (((1 - \operatorname{real\text{-}part}(x)) - \operatorname{imag\text{-}part}(x)) \cdot c)), (((\operatorname{real\text{-}part}(\mathit{gouy}) \cdot a) + (\operatorname{imag\text{-}part}(\mathit{gouy}) \cdot b)) + (((1 - \operatorname{real\text{-}part}(\mathit{gouy})) - \operatorname{imag\text{-}part}(\mathit{gouy})) \cdot c)), (((\operatorname{real\text{-}part}(z) \cdot a) + (\operatorname{imag\text{-}part}(z) \cdot b)) + (((1 - \operatorname{real\text{-}part}(z)) - \operatorname{imag\text{-}part}(z)) \cdot c))))) \end{array}}}
\vnbentry{\texttt{affine-\allowbreak{}holomorphic}}{\vnbstmt{\begin{array}{@{}l@{}} \forall u.; \forall p, q, w \in \mathbb{C}, g \in (u \to \mathbb{C}). \\ \quad \operatorname{is\text{-}open}(\operatorname{nf\text{-}metric\text{-}space}(\operatorname{cc\text{-}normed\text{-}field}), u) \wedge \\ \quad \forall z \in u.\; (g(z) \simeq (p + (q \cdot (z - w)))) \Rightarrow \\ \quad \quad \operatorname{holomorphic\text{-}on}(u, g) \end{array}}}
\vnbentry{\texttt{small-\allowbreak{}dyadic}}{\vnbstmt{\begin{array}{@{}l@{}} \forall m \in \mathbb{R}, r. \\ \quad ((0 \leq m) \wedge \operatorname{pos\text{-}rr}(r)) \Rightarrow  \\ \quad \quad \exists n \in \mathbb{N}. \\ \quad \quad \quad ((m \cdot \operatorname{recip}(\operatorname{power}(2, n))) < r) \end{array}}}
\vnbentry{\texttt{bary-\allowbreak{}side-\allowbreak{}near}}{\vnbstmt{\begin{array}{@{}l@{}} \forall a, b, c, p, q, x, w \in \mathbb{C}, r, m \in \mathbb{R}, z \in \operatorname{trace}(\operatorname{seg\text{-}path}((((\operatorname{real\text{-}part}(p) \cdot a) + (\operatorname{imag\text{-}part}(p) \cdot b)) + (((1 - \operatorname{real\text{-}part}(p)) - \operatorname{imag\text{-}part}(p)) \cdot c)), (((\operatorname{real\text{-}part}(q) \cdot a) + (\operatorname{imag\text{-}part}(q) \cdot b)) + (((1 - \operatorname{real\text{-}part}(q)) - \operatorname{imag\text{-}part}(q)) \cdot c))), 0, 1). \\ \quad (\operatorname{magnitude}((p - x)) \leq r) \wedge \\ \quad (\operatorname{magnitude}((q - p)) \leq r) \wedge \\ \quad (r \leq (6 \cdot m)) \wedge \\ \quad (\operatorname{magnitude}((x - w)) \leq (12 \cdot m)) \Rightarrow \\ \quad \quad \operatorname{magnitude}((z - (((\operatorname{real\text{-}part}(w) \cdot a) + (\operatorname{imag\text{-}part}(w) \cdot b)) + (((1 - \operatorname{real\text{-}part}(w)) - \operatorname{imag\text{-}part}(w)) \cdot c)))) \\ \quad \quad \leq\; ((24 \cdot (\operatorname{magnitude}((a - c)) + \operatorname{magnitude}((b - c)))) \cdot m) \end{array}}}
\vnbentry{\texttt{tri-\allowbreak{}int-\allowbreak{}abs-\allowbreak{}bound-\allowbreak{}sides}}{\vnbstmt{\begin{array}{@{}l@{}} \forall u, f.; \forall a, b, c \in \mathbb{C}, m \in \mathbb{R}. \\ \quad (u \subseteq \mathbb{C}) \wedge \\ \quad (f \in (u \to \mathbb{C})) \wedge \\ \quad \forall y \in u. \\ \quad \quad \operatorname{is\text{-}continuous\text{-}at}(\operatorname{subspace\text{-}ms}(\operatorname{nf\text{-}metric\text{-}space}(\operatorname{cc\text{-}normed\text{-}field}), u), \\ \quad \quad \quad \operatorname{nf\text{-}metric\text{-}space}(\operatorname{cc\text{-}normed\text{-}field}), f, y) \wedge \\ \quad (\operatorname{trace}(\operatorname{seg\text{-}path}(a, b), 0, 1) \subseteq u) \wedge \\ \quad (\operatorname{trace}(\operatorname{seg\text{-}path}(b, c), 0, 1) \subseteq u) \wedge \\ \quad (\operatorname{trace}(\operatorname{seg\text{-}path}(c, a), 0, 1) \subseteq u) \wedge \\ \quad \forall z \in \operatorname{trace}(\operatorname{seg\text{-}path}(a, b), 0, 1). \\ \quad \quad (\operatorname{magnitude}(f(z)) \leq m) \wedge \\ \quad \forall z \in \operatorname{trace}(\operatorname{seg\text{-}path}(b, c), 0, 1). \\ \quad \quad (\operatorname{magnitude}(f(z)) \leq m) \wedge \\ \quad \forall z \in \operatorname{trace}(\operatorname{seg\text{-}path}(c, a), 0, 1). \\ \quad \quad (\operatorname{magnitude}(f(z)) \leq m) \Rightarrow \\ \quad \quad \operatorname{magnitude}(\operatorname{tri\text{-}int}(f, a, b, c)) \\ \quad \quad \leq\; (((\operatorname{magnitude}((b - a)) + \operatorname{magnitude}((c - b))) + \operatorname{magnitude}((a - c))) \cdot m) \end{array}}}
\vnbentry{\texttt{goursat-\allowbreak{}triangle}}{\vnbstmt{\begin{array}{@{}l@{}} \forall u, f.; \forall a, b, c \in u. \\ \quad (\operatorname{holomorphic\text{-}on}(u, f) \wedge (\operatorname{conv3}(a, b, c) \subseteq u)) \Rightarrow  \\ \quad \quad (\operatorname{tri\text{-}int}(f, a, b, c) = 0) \end{array}}}
\vnbentry{\texttt{cc-\allowbreak{}diff-\allowbreak{}on-\allowbreak{}local-\allowbreak{}quadratic}}{\vnbstmt{\begin{array}{@{}l@{}} \forall s.; \forall f \in (s \to \mathbb{C}), a \in s, l \in \mathbb{C}, m \in \mathbb{R}, r. \\ \quad \operatorname{is\text{-}open}(\operatorname{nf\text{-}metric\text{-}space}(\operatorname{cc\text{-}normed\text{-}field}), s) \wedge \\ \quad (0 \leq m) \wedge \\ \quad \operatorname{pos\text{-}rr}(r) \wedge \\ \quad \forall x \in s. \\ \quad \quad (\operatorname{magnitude}((x - a)) < r) \Rightarrow  \\ \quad \quad \quad \operatorname{magnitude}(((f(x) - f(a)) - ((x - a) \cdot l))) \\ \quad \quad \quad \leq\; (m \cdot (\operatorname{magnitude}((x - a)) \cdot \operatorname{magnitude}((x - a)))) \Rightarrow \\ \quad \quad \operatorname{is\text{-}diff\text{-}on}(\operatorname{cc\text{-}normed\text{-}field}, s, f, a, l) \end{array}}}
\vnbentry{\texttt{cauchy-\allowbreak{}convex-\allowbreak{}segments}}{\vnbstmt{\begin{array}{@{}l@{}} \forall u, f.; \forall a, b, c \in u. \\ \quad (\operatorname{is\text{-}convex}(u) \wedge \operatorname{holomorphic\text{-}on}(u, f)) \Rightarrow  \\ \quad \quad \operatorname{seg\text{-}int}(f, a, c) \\ \quad \quad =\; (\operatorname{seg\text{-}int}(f, a, b) + \operatorname{seg\text{-}int}(f, b, c)) \end{array}}}
\vnbentry{\texttt{seg-\allowbreak{}int-\allowbreak{}sum}}{\vnbstmt{\begin{array}{@{}l@{}} \forall u, f, g.; \forall h \in (u \to \mathbb{C}), a, b \in \mathbb{C}. \\ \quad (u \subseteq \mathbb{C}) \wedge \\ \quad (f \in (u \to \mathbb{C})) \wedge \\ \quad \forall y \in u. \\ \quad \quad \operatorname{is\text{-}continuous\text{-}at}(\operatorname{subspace\text{-}ms}(\operatorname{nf\text{-}metric\text{-}space}(\operatorname{cc\text{-}normed\text{-}field}), u), \\ \quad \quad \quad \operatorname{nf\text{-}metric\text{-}space}(\operatorname{cc\text{-}normed\text{-}field}), f, y) \wedge \\ \quad (g \in (u \to \mathbb{C})) \wedge \\ \quad \forall y \in u. \\ \quad \quad \operatorname{is\text{-}continuous\text{-}at}(\operatorname{subspace\text{-}ms}(\operatorname{nf\text{-}metric\text{-}space}(\operatorname{cc\text{-}normed\text{-}field}), u), \\ \quad \quad \quad \operatorname{nf\text{-}metric\text{-}space}(\operatorname{cc\text{-}normed\text{-}field}), g, y) \wedge \\ \quad \forall z \in u.\; (h(z) \simeq (f(z) + g(z))) \wedge \\ \quad (\operatorname{trace}(\operatorname{seg\text{-}path}(a, b), 0, 1) \subseteq u) \Rightarrow \\ \quad \quad \operatorname{seg\text{-}int}(h, a, b) \\ \quad \quad =\; (\operatorname{seg\text{-}int}(f, a, b) + \operatorname{seg\text{-}int}(g, a, b)) \end{array}}}
\vnbentry{\texttt{seg-\allowbreak{}int-\allowbreak{}abs-\allowbreak{}bound}}{\vnbstmt{\begin{array}{@{}l@{}} \forall u, f.; \forall a, b \in \mathbb{C}, m \in \mathbb{R}. \\ \quad (u \subseteq \mathbb{C}) \wedge \\ \quad (f \in (u \to \mathbb{C})) \wedge \\ \quad \forall y \in u. \\ \quad \quad \operatorname{is\text{-}continuous\text{-}at}(\operatorname{subspace\text{-}ms}(\operatorname{nf\text{-}metric\text{-}space}(\operatorname{cc\text{-}normed\text{-}field}), u), \\ \quad \quad \quad \operatorname{nf\text{-}metric\text{-}space}(\operatorname{cc\text{-}normed\text{-}field}), f, y) \wedge \\ \quad (\operatorname{trace}(\operatorname{seg\text{-}path}(a, b), 0, 1) \subseteq u) \wedge \\ \quad \forall z \in \operatorname{trace}(\operatorname{seg\text{-}path}(a, b), 0, 1). \\ \quad \quad (\operatorname{magnitude}(f(z)) \leq m) \Rightarrow \\ \quad \quad \operatorname{magnitude}(\operatorname{seg\text{-}int}(f, a, b)) \\ \quad \quad \leq\; (\operatorname{magnitude}((b - a)) \cdot m) \end{array}}}
\vnbentry{\texttt{seg-\allowbreak{}int-\allowbreak{}const}}{\vnbstmt{\begin{array}{@{}l@{}} \forall u.; \forall p, w \in \mathbb{C}, g \in (u \to \mathbb{C}), a, b \in \mathbb{C}. \\ \quad \operatorname{is\text{-}open}(\operatorname{nf\text{-}metric\text{-}space}(\operatorname{cc\text{-}normed\text{-}field}), u) \wedge \\ \quad \forall z \in u.\; (g(z) \simeq (p + (0 \cdot (z - w)))) \wedge \\ \quad (\operatorname{trace}(\operatorname{seg\text{-}path}(a, b), 0, 1) \subseteq u) \Rightarrow \\ \quad \quad (\operatorname{seg\text{-}int}(g, a, b) = (p \cdot (b - a))) \end{array}}}
\vnbentry{\texttt{cauchy-\allowbreak{}convex-\allowbreak{}primitive}}{\vnbstmt{\begin{array}{@{}l@{}} \forall u, f.; \forall o \in u. \\ \quad (\operatorname{is\text{-}convex}(u) \wedge \operatorname{holomorphic\text{-}on}(u, f)) \Rightarrow  \\ \quad \quad \exists g \in (u \to \mathbb{C}). \\ \quad \quad \quad \operatorname{holomorphic\text{-}on}(u, g) \wedge \\ \quad \quad \quad \forall z \in u. \\ \quad \quad \quad \quad \operatorname{is\text{-}diff\text{-}on}(\operatorname{cc\text{-}normed\text{-}field}, u, g, z, f(z)) \end{array}}}
\vnbentry{\texttt{cauchy-\allowbreak{}convex-\allowbreak{}closed}}{\vnbstmt{\begin{array}{@{}l@{}} \forall u, f, m, c, a, b. \\ \quad \operatorname{is\text{-}convex}(u) \wedge \\ \quad \operatorname{holomorphic\text{-}on}(u, f) \wedge \\ \quad \operatorname{is\text{-}road}(m, c, a, b) \wedge \\ \quad (\operatorname{trace}(m, a, b) \subseteq u) \wedge \\ \quad (m(b) = m(a)) \Rightarrow \\ \quad \quad (\operatorname{line\text{-}int}(f, m, c, a, b) = 0) \end{array}}}

\vnbfile{theorem-library/goursat-exceptional.scm}{theorem-library/goursat-exceptional.scm}
\noindent{\small goursat-exceptional.scm -- CAUCHY-GOURSAT WITH ONE EXCEPTIONAL POINT (complex-analysis.pdf ch. 3, s. 3.2: Prop 3.6, Cor 3.7, Prop 3.8, Cor 3.9 in the notes' form: f CONTINUOUS on the open set U and HOLOMORPHIC on U \textbackslash{} \{x\} for some x in U).  Batch 35, 2026-09-26.}\par
\vnbentry{\texttt{seg-\allowbreak{}int-\allowbreak{}congruence}}{\vnbstmt{\begin{array}{@{}l@{}} \forall u, f, w.; \forall g \in (w \to \mathbb{C}), a, b \in \mathbb{C}. \\ \quad (u \subseteq \mathbb{C}) \wedge \\ \quad (f \in (u \to \mathbb{C})) \wedge \\ \quad \forall y \in u. \\ \quad \quad \operatorname{is\text{-}continuous\text{-}at}(\operatorname{subspace\text{-}ms}(\operatorname{nf\text{-}metric\text{-}space}(\operatorname{cc\text{-}normed\text{-}field}), u), \\ \quad \quad \quad \operatorname{nf\text{-}metric\text{-}space}(\operatorname{cc\text{-}normed\text{-}field}), f, y) \wedge \\ \quad (\operatorname{trace}(\operatorname{seg\text{-}path}(a, b), 0, 1) \subseteq u) \wedge \\ \quad (\operatorname{trace}(\operatorname{seg\text{-}path}(a, b), 0, 1) \subseteq w) \wedge \\ \quad \forall z \in \operatorname{trace}(\operatorname{seg\text{-}path}(a, b), 0, 1). \\ \quad \quad (g(z) \simeq f(z)) \Rightarrow \\ \quad \quad (\operatorname{seg\text{-}int}(g, a, b) = \operatorname{seg\text{-}int}(f, a, b)) \end{array}}}
\vnbentry{\texttt{tri-\allowbreak{}int-\allowbreak{}restrict}}{\vnbstmt{\begin{array}{@{}l@{}} \forall u, f, w.; \forall a, b, c \in \mathbb{C}. \\ \quad (u \subseteq \mathbb{C}) \wedge \\ \quad (f \in (u \to \mathbb{C})) \wedge \\ \quad \forall y \in u. \\ \quad \quad \operatorname{is\text{-}continuous\text{-}at}(\operatorname{subspace\text{-}ms}(\operatorname{nf\text{-}metric\text{-}space}(\operatorname{cc\text{-}normed\text{-}field}), u), \\ \quad \quad \quad \operatorname{nf\text{-}metric\text{-}space}(\operatorname{cc\text{-}normed\text{-}field}), f, y) \wedge \\ \quad (w \subseteq u) \wedge \\ \quad (\operatorname{conv3}(a, b, c) \subseteq w) \Rightarrow \\ \quad \quad \operatorname{tri\text{-}int}(\operatorname{restrict}(f, w), a, b, c) \\ \quad \quad =\; \operatorname{tri\text{-}int}(f, a, b, c) \end{array}}}
\vnbentry{\texttt{tri-\allowbreak{}int-\allowbreak{}fan}}{\vnbstmt{\begin{array}{@{}l@{}} \forall u, f.; \forall a, b, c \in \mathbb{C}, x \in \operatorname{conv3}(a, b, c). \\ \quad (u \subseteq \mathbb{C}) \wedge \\ \quad (f \in (u \to \mathbb{C})) \wedge \\ \quad \forall y \in u. \\ \quad \quad \operatorname{is\text{-}continuous\text{-}at}(\operatorname{subspace\text{-}ms}(\operatorname{nf\text{-}metric\text{-}space}(\operatorname{cc\text{-}normed\text{-}field}), u), \\ \quad \quad \quad \operatorname{nf\text{-}metric\text{-}space}(\operatorname{cc\text{-}normed\text{-}field}), f, y) \wedge \\ \quad (\operatorname{conv3}(a, b, c) \subseteq u) \Rightarrow \\ \quad \quad \operatorname{tri\text{-}int}(f, a, b, c) \\ \quad \quad =\; ((\operatorname{tri\text{-}int}(f, a, b, x) + \operatorname{tri\text{-}int}(f, b, c, x)) + \operatorname{tri\text{-}int}(f, c, a, x)) \end{array}}}
\vnbentry{\texttt{tri-\allowbreak{}int-\allowbreak{}collinear}}{\vnbstmt{\begin{array}{@{}l@{}} \forall u, f.; \forall a, b, c \in \mathbb{C}. \\ \quad (u \subseteq \mathbb{C}) \wedge \\ \quad (f \in (u \to \mathbb{C})) \wedge \\ \quad \forall y \in u. \\ \quad \quad \operatorname{is\text{-}continuous\text{-}at}(\operatorname{subspace\text{-}ms}(\operatorname{nf\text{-}metric\text{-}space}(\operatorname{cc\text{-}normed\text{-}field}), u), \\ \quad \quad \quad \operatorname{nf\text{-}metric\text{-}space}(\operatorname{cc\text{-}normed\text{-}field}), f, y) \wedge \\ \quad (\operatorname{conv3}(a, b, c) \subseteq u) \wedge \\ \quad (c \in \operatorname{trace}(\operatorname{seg\text{-}path}(a, b), 0, 1)) \vee \\ \quad (a \in \operatorname{trace}(\operatorname{seg\text{-}path}(b, c), 0, 1)) \vee \\ \quad (b \in \operatorname{trace}(\operatorname{seg\text{-}path}(c, a), 0, 1)) \Rightarrow \\ \quad \quad (\operatorname{tri\text{-}int}(f, a, b, c) = 0) \end{array}}}
\vnbentry{\texttt{cc-\allowbreak{}cross-\allowbreak{}coords}}{\vnbstmt{\begin{array}{@{}l@{}} \forall p, q \in \mathbb{C}. \\ \quad \operatorname{imag\text{-}part}((p \cdot \operatorname{conjugate}(q))) \\ \quad =\; ((\operatorname{imag\text{-}part}(p) \cdot \operatorname{real\text{-}part}(q)) - (\operatorname{real\text{-}part}(p) \cdot \operatorname{imag\text{-}part}(q))) \end{array}}}
\vnbentry{\texttt{cc-\allowbreak{}affinely-\allowbreak{}independent}}{\vnbstmt{\begin{array}{@{}l@{}} \forall a, b, c \in \mathbb{C}, p, q \in \mathbb{R}. \\ \quad \neg (\operatorname{imag\text{-}part}(((b - a) \cdot \operatorname{conjugate}((c - a)))) = 0) \wedge \\ \quad (((p \cdot (b - a)) + (q \cdot (c - a))) = 0) \Rightarrow \\ \quad \quad ((p = 0) \wedge (q = 0)) \end{array}}}
\vnbentry{\texttt{conv3-\allowbreak{}avoids-\allowbreak{}vertex}}{\vnbstmt{\begin{array}{@{}l@{}} \forall a, b, c \in \mathbb{C}, t \in \mathbb{R}. \\ \quad \neg (\operatorname{imag\text{-}part}(((b - a) \cdot \operatorname{conjugate}((c - a)))) = 0) \wedge \\ \quad (0 < t) \wedge \\ \quad (t \leq 1) \Rightarrow \\ \quad \quad \neg (a \in \operatorname{conv3}((a + (t \cdot (b - a))), b, c)) \wedge \\ \quad \quad \neg (a \in \operatorname{conv3}((a + (t \cdot (b - a))), c, (a + (t \cdot (c - a))))) \end{array}}}
\vnbentry{\texttt{conv3-\allowbreak{}near-\allowbreak{}vertex}}{\vnbstmt{\begin{array}{@{}l@{}} \forall a, b, c \in \mathbb{C}, z \in \operatorname{conv3}(a, b, c). \\ \quad \operatorname{magnitude}((z - a)) \\ \quad \leq\; (\operatorname{magnitude}((b - a)) + \operatorname{magnitude}((c - a))) \end{array}}}
\vnbentry{\texttt{collinear-\allowbreak{}cases}}{\vnbstmt{\begin{array}{@{}l@{}} \forall a, b, c \in \mathbb{C}. \\ \quad (\operatorname{imag\text{-}part}(((b - a) \cdot \operatorname{conjugate}((c - a)))) = 0) \Rightarrow  \\ \quad \quad (c \in \operatorname{trace}(\operatorname{seg\text{-}path}(a, b), 0, 1)) \vee \\ \quad \quad (a \in \operatorname{trace}(\operatorname{seg\text{-}path}(b, c), 0, 1)) \vee \\ \quad \quad (b \in \operatorname{trace}(\operatorname{seg\text{-}path}(c, a), 0, 1)) \end{array}}}
\vnbentry{\texttt{conv3-\allowbreak{}vertex-\allowbreak{}shrink}}{\vnbstmt{\begin{array}{@{}l@{}} \forall a, b, c \in \mathbb{C}, t \in \operatorname{ccint}(0, 1). \\ \quad \operatorname{conv3}(a, (a + (t \cdot (b - a))), (a + (t \cdot (c - a)))) \\ \quad \subseteq\; \operatorname{conv3}(a, b, c) \wedge \\ \quad (\operatorname{conv3}((a + (t \cdot (b - a))), b, c) \subseteq \operatorname{conv3}(a, b, c)) \wedge \\ \quad \operatorname{conv3}((a + (t \cdot (b - a))), c, (a + (t \cdot (c - a)))) \\ \quad \subseteq\; \operatorname{conv3}(a, b, c) \end{array}}}
\vnbentry{\texttt{tri-\allowbreak{}int-\allowbreak{}vertex-\allowbreak{}split}}{\vnbstmt{\begin{array}{@{}l@{}} \forall u, f.; \forall a, b, c \in \mathbb{C}, t \in \operatorname{ccint}(0, 1). \\ \quad (u \subseteq \mathbb{C}) \wedge \\ \quad (f \in (u \to \mathbb{C})) \wedge \\ \quad \forall y \in u. \\ \quad \quad \operatorname{is\text{-}continuous\text{-}at}(\operatorname{subspace\text{-}ms}(\operatorname{nf\text{-}metric\text{-}space}(\operatorname{cc\text{-}normed\text{-}field}), u), \\ \quad \quad \quad \operatorname{nf\text{-}metric\text{-}space}(\operatorname{cc\text{-}normed\text{-}field}), f, y) \wedge \\ \quad (\operatorname{conv3}(a, b, c) \subseteq u) \Rightarrow \\ \quad \quad \operatorname{tri\text{-}int}(f, a, b, c) \\ \quad \quad =\; ((\operatorname{tri\text{-}int}(f, a, (a + (t \cdot (b - a))), (a + (t \cdot (c - a)))) + \operatorname{tri\text{-}int}(f, (a + (t \cdot (b - a))), b, c)) + \operatorname{tri\text{-}int}(f, (a + (t \cdot (b - a))), c, (a + (t \cdot (c - a))))) \end{array}}}
\vnbentry{\texttt{goursat-\allowbreak{}vertex}}{\vnbstmt{\begin{array}{@{}l@{}} \forall u, f, x.; \forall b, c \in u. \\ \quad \operatorname{is\text{-}open}(\operatorname{nf\text{-}metric\text{-}space}(\operatorname{cc\text{-}normed\text{-}field}), u) \wedge \\ \quad (f \in (u \to \mathbb{C})) \wedge \\ \quad \forall y \in u. \\ \quad \quad \operatorname{is\text{-}continuous\text{-}at}(\operatorname{subspace\text{-}ms}(\operatorname{nf\text{-}metric\text{-}space}(\operatorname{cc\text{-}normed\text{-}field}), u), \\ \quad \quad \quad \operatorname{nf\text{-}metric\text{-}space}(\operatorname{cc\text{-}normed\text{-}field}), f, y) \wedge \\ \quad (x \in u) \wedge \\ \quad \operatorname{holomorphic\text{-}on}(\operatorname{difference}(u, \operatorname{singleton}(x)), \\ \quad \quad \operatorname{restrict}(f, \operatorname{difference}(u, \operatorname{singleton}(x)))) \wedge \\ \quad (\operatorname{conv3}(x, b, c) \subseteq u) \Rightarrow \\ \quad \quad (\operatorname{tri\text{-}int}(f, x, b, c) = 0) \end{array}}}
\vnbentry{\texttt{goursat-\allowbreak{}triangle-\allowbreak{}except}}{\vnbstmt{\begin{array}{@{}l@{}} \forall u, f, x.; \forall a, b, c \in u. \\ \quad \operatorname{is\text{-}open}(\operatorname{nf\text{-}metric\text{-}space}(\operatorname{cc\text{-}normed\text{-}field}), u) \wedge \\ \quad (f \in (u \to \mathbb{C})) \wedge \\ \quad \forall y \in u. \\ \quad \quad \operatorname{is\text{-}continuous\text{-}at}(\operatorname{subspace\text{-}ms}(\operatorname{nf\text{-}metric\text{-}space}(\operatorname{cc\text{-}normed\text{-}field}), u), \\ \quad \quad \quad \operatorname{nf\text{-}metric\text{-}space}(\operatorname{cc\text{-}normed\text{-}field}), f, y) \wedge \\ \quad (x \in u) \wedge \\ \quad \operatorname{holomorphic\text{-}on}(\operatorname{difference}(u, \operatorname{singleton}(x)), \\ \quad \quad \operatorname{restrict}(f, \operatorname{difference}(u, \operatorname{singleton}(x)))) \wedge \\ \quad (\operatorname{conv3}(a, b, c) \subseteq u) \Rightarrow \\ \quad \quad (\operatorname{tri\text{-}int}(f, a, b, c) = 0) \end{array}}}
\vnbentry{\texttt{cc-\allowbreak{}diff-\allowbreak{}on-\allowbreak{}of-\allowbreak{}eps-\allowbreak{}bound}}{\vnbstmt{\begin{array}{@{}l@{}} \forall s.; \forall g \in (s \to \mathbb{C}), a \in s, l \in \mathbb{C}. \\ \quad \operatorname{is\text{-}open}(\operatorname{nf\text{-}metric\text{-}space}(\operatorname{cc\text{-}normed\text{-}field}), s) \wedge \\ \quad \forall e. \\ \quad \quad \operatorname{pos\text{-}rr}(e) \Rightarrow  \\ \quad \quad \quad \exists d. \\ \quad \quad \quad \quad \operatorname{pos\text{-}rr}(d) \wedge \\ \quad \quad \quad \quad \forall y \in s. \\ \quad \quad \quad \quad \quad (\operatorname{magnitude}((y - a)) < d) \Rightarrow  \\ \quad \quad \quad \quad \quad \quad \operatorname{magnitude}(((g(y) - g(a)) - ((y - a) \cdot l))) \\ \quad \quad \quad \quad \quad \quad \leq\; (e \cdot \operatorname{magnitude}((y - a))) \Rightarrow \\ \quad \quad \operatorname{is\text{-}diff\text{-}on}(\operatorname{cc\text{-}normed\text{-}field}, s, g, a, l) \end{array}}}
\vnbentry{\texttt{cc-\allowbreak{}continuous-\allowbreak{}at-\allowbreak{}add-\allowbreak{}const}}{\vnbstmt{\begin{array}{@{}l@{}} \forall u.; \forall f, g \in (u \to \mathbb{C}), k \in \mathbb{C}, y \in u. \\ \quad (u \subseteq \mathbb{C}) \wedge \\ \quad \forall z \in u.\; (g(z) \simeq (f(z) + k)) \wedge \\ \quad \forall y \in u. \\ \quad \quad \operatorname{is\text{-}continuous\text{-}at}(\operatorname{subspace\text{-}ms}(\operatorname{nf\text{-}metric\text{-}space}(\operatorname{cc\text{-}normed\text{-}field}), u), \\ \quad \quad \quad \operatorname{nf\text{-}metric\text{-}space}(\operatorname{cc\text{-}normed\text{-}field}), f, y) \Rightarrow \\ \quad \quad \operatorname{is\text{-}continuous\text{-}at}(\operatorname{subspace\text{-}ms}(\operatorname{nf\text{-}metric\text{-}space}(\operatorname{cc\text{-}normed\text{-}field}), u), \\ \quad \quad \quad \operatorname{nf\text{-}metric\text{-}space}(\operatorname{cc\text{-}normed\text{-}field}), g, y) \end{array}}}
\vnbentry{\texttt{cauchy-\allowbreak{}convex-\allowbreak{}segments-\allowbreak{}except}}{\vnbstmt{\begin{array}{@{}l@{}} \forall u, f, x.; \forall a, b, c \in u. \\ \quad \operatorname{is\text{-}convex}(u) \wedge \\ \quad \operatorname{is\text{-}open}(\operatorname{nf\text{-}metric\text{-}space}(\operatorname{cc\text{-}normed\text{-}field}), u) \wedge \\ \quad (f \in (u \to \mathbb{C})) \wedge \\ \quad \forall y \in u. \\ \quad \quad \operatorname{is\text{-}continuous\text{-}at}(\operatorname{subspace\text{-}ms}(\operatorname{nf\text{-}metric\text{-}space}(\operatorname{cc\text{-}normed\text{-}field}), u), \\ \quad \quad \quad \operatorname{nf\text{-}metric\text{-}space}(\operatorname{cc\text{-}normed\text{-}field}), f, y) \wedge \\ \quad (x \in u) \wedge \\ \quad \operatorname{holomorphic\text{-}on}(\operatorname{difference}(u, \operatorname{singleton}(x)), \\ \quad \quad \operatorname{restrict}(f, \operatorname{difference}(u, \operatorname{singleton}(x)))) \Rightarrow \\ \quad \quad \operatorname{seg\text{-}int}(f, a, c) \\ \quad \quad =\; (\operatorname{seg\text{-}int}(f, a, b) + \operatorname{seg\text{-}int}(f, b, c)) \end{array}}}
\vnbentry{\texttt{cauchy-\allowbreak{}convex-\allowbreak{}primitive-\allowbreak{}except}}{\vnbstmt{\begin{array}{@{}l@{}} \forall u, f, x.; \forall o \in u. \\ \quad \operatorname{is\text{-}convex}(u) \wedge \\ \quad \operatorname{is\text{-}open}(\operatorname{nf\text{-}metric\text{-}space}(\operatorname{cc\text{-}normed\text{-}field}), u) \wedge \\ \quad (f \in (u \to \mathbb{C})) \wedge \\ \quad \forall y \in u. \\ \quad \quad \operatorname{is\text{-}continuous\text{-}at}(\operatorname{subspace\text{-}ms}(\operatorname{nf\text{-}metric\text{-}space}(\operatorname{cc\text{-}normed\text{-}field}), u), \\ \quad \quad \quad \operatorname{nf\text{-}metric\text{-}space}(\operatorname{cc\text{-}normed\text{-}field}), f, y) \wedge \\ \quad (x \in u) \wedge \\ \quad \operatorname{holomorphic\text{-}on}(\operatorname{difference}(u, \operatorname{singleton}(x)), \\ \quad \quad \operatorname{restrict}(f, \operatorname{difference}(u, \operatorname{singleton}(x)))) \Rightarrow \\ \quad \quad \exists g \in (u \to \mathbb{C}). \\ \quad \quad \quad \operatorname{holomorphic\text{-}on}(u, g) \wedge \\ \quad \quad \quad \forall z \in u. \\ \quad \quad \quad \quad \operatorname{is\text{-}diff\text{-}on}(\operatorname{cc\text{-}normed\text{-}field}, u, g, z, f(z)) \end{array}}}
\vnbentry{\texttt{cauchy-\allowbreak{}convex-\allowbreak{}closed-\allowbreak{}except}}{\vnbstmt{\begin{array}{@{}l@{}} \forall u, f, x, m, c, a, b. \\ \quad \operatorname{is\text{-}convex}(u) \wedge \\ \quad \operatorname{is\text{-}open}(\operatorname{nf\text{-}metric\text{-}space}(\operatorname{cc\text{-}normed\text{-}field}), u) \wedge \\ \quad (f \in (u \to \mathbb{C})) \wedge \\ \quad \forall y \in u. \\ \quad \quad \operatorname{is\text{-}continuous\text{-}at}(\operatorname{subspace\text{-}ms}(\operatorname{nf\text{-}metric\text{-}space}(\operatorname{cc\text{-}normed\text{-}field}), u), \\ \quad \quad \quad \operatorname{nf\text{-}metric\text{-}space}(\operatorname{cc\text{-}normed\text{-}field}), f, y) \wedge \\ \quad (x \in u) \wedge \\ \quad \operatorname{holomorphic\text{-}on}(\operatorname{difference}(u, \operatorname{singleton}(x)), \\ \quad \quad \operatorname{restrict}(f, \operatorname{difference}(u, \operatorname{singleton}(x)))) \wedge \\ \quad \operatorname{is\text{-}road}(m, c, a, b) \wedge \\ \quad (\operatorname{trace}(m, a, b) \subseteq u) \wedge \\ \quad (m(b) = m(a)) \Rightarrow \\ \quad \quad (\operatorname{line\text{-}int}(f, m, c, a, b) = 0) \end{array}}}

\vnbfile{theorem-library/road-juxtaposition-laws.scm}{theorem-library/road-juxtaposition-laws.scm}
\noindent{\small road-juxtaposition-laws.scm -- THE JUXTAPOSITION OF ROADS (Dieudonne 9.6, the notes' Proposition 3.1), THE TRIANGULAR ROAD, AND THE LENGTH OF A ROAD (the notes' (53)). Batch 37, 2026-09-26.  The vocabulary is structure-library/road-juxtaposition.scm (JUXTA, ROAD-LENGTH, TRI-ROAD, TRI-ROAD-DERIV).}\par
\vnbentry{\texttt{juxta-\allowbreak{}unfold}}{\vnbstmt{\begin{array}{@{}l@{}} \forall p, q, a, b, c. \\ \quad (\operatorname{juxta}(p, q, a, b, c) \simeq (\lambda t \in \operatorname{ccint}(a, c).\; \left(\text{if }(t \leq b)\text{ then }p(t)\text{ else }q(t)\right))) \end{array}}}
\vnbentry{\texttt{road-\allowbreak{}length-\allowbreak{}unfold}}{\vnbstmt{\begin{array}{@{}l@{}} \forall d, a, b. \\ \quad (\operatorname{road\text{-}length}(d, a, b) \simeq \int_{a}^{b} (\lambda t \in \operatorname{ccint}(a, b).\; \operatorname{magnitude}(d(t)))) \end{array}}}
\vnbentry{\texttt{tri-\allowbreak{}road-\allowbreak{}unfold}}{\vnbstmt{\begin{array}{@{}l@{}} \forall a, b, c. \\ \quad (\operatorname{tri\text{-}road}(a, b, c) \simeq \operatorname{juxta}(\operatorname{juxta}(\operatorname{seg\text{-}path}(a, b), (\lambda s \in \operatorname{ccint}(1, 2).\; \operatorname{seg\text{-}path}(b, c)((s - 1))), 0, 1, 2), (\lambda s \in \operatorname{ccint}(2, 3).\; \operatorname{seg\text{-}path}(c, a)((s - 2))), 0, 2, 3)) \end{array}}}
\vnbentry{\texttt{tri-\allowbreak{}road-\allowbreak{}deriv-\allowbreak{}unfold}}{\vnbstmt{\begin{array}{@{}l@{}} \forall a, b, c. \\ \quad (\operatorname{tri\text{-}road\text{-}deriv}(a, b, c) \simeq \operatorname{juxta}(\operatorname{juxta}(\operatorname{seg\text{-}deriv}(a, b), (\lambda s \in \operatorname{ccint}(1, 2).\; \operatorname{seg\text{-}deriv}(b, c)((s - 1))), 0, 1, 2), (\lambda s \in \operatorname{ccint}(2, 3).\; \operatorname{seg\text{-}deriv}(c, a)((s - 2))), 0, 2, 3)) \end{array}}}
\vnbentry{\texttt{juxta-\allowbreak{}apply-\allowbreak{}left}}{\vnbstmt{\begin{array}{@{}l@{}} \forall p, q.; \forall a, b, c \in \mathbb{R}, t \in \operatorname{ccint}(a, b). \\ \quad ((b \leq c) \Rightarrow (\operatorname{juxta}(p, q, a, b, c)(t) \simeq p(t))) \end{array}}}
\vnbentry{\texttt{juxta-\allowbreak{}apply-\allowbreak{}right}}{\vnbstmt{\begin{array}{@{}l@{}} \forall p, q.; \forall a, b, c \in \mathbb{R}, t \in \operatorname{ccint}(b, c). \\ \quad ((a \leq b) \wedge (b < t)) \Rightarrow  \\ \quad \quad (\operatorname{juxta}(p, q, a, b, c)(t) \simeq q(t)) \end{array}}}
\vnbentry{\texttt{juxta-\allowbreak{}apply-\allowbreak{}right-\allowbreak{}closed}}{\vnbstmt{\begin{array}{@{}l@{}} \forall p, q.; \forall a, b, c \in \mathbb{R}, t \in \operatorname{ccint}(b, c). \\ \quad ((a \leq b) \wedge (p(b) = q(b))) \Rightarrow  \\ \quad \quad (\operatorname{juxta}(p, q, a, b, c)(t) \simeq q(t)) \end{array}}}
\vnbentry{\texttt{juxta-\allowbreak{}in-\allowbreak{}fun}}{\vnbstmt{\begin{array}{@{}l@{}} \forall p, q.; \forall a, b, c \in \mathbb{R}, d. \\ \quad (a \leq b) \wedge \\ \quad (b \leq c) \wedge \\ \quad (p \in (\operatorname{ccint}(a, b) \to d)) \wedge \\ \quad (q \in (\operatorname{ccint}(b, c) \to d)) \Rightarrow \\ \quad \quad (\operatorname{juxta}(p, q, a, b, c) \in (\operatorname{ccint}(a, c) \to d)) \end{array}}}
\vnbentry{\texttt{juxta-\allowbreak{}trace-\allowbreak{}subset}}{\vnbstmt{\begin{array}{@{}l@{}} \forall p, q.; \forall a, b, c \in \mathbb{R}. \\ \quad (a \leq b) \wedge \\ \quad (b \leq c) \wedge \\ \quad (p \in (\operatorname{ccint}(a, b) \to \mathbb{C})) \wedge \\ \quad (q \in (\operatorname{ccint}(b, c) \to \mathbb{C})) \Rightarrow \\ \quad \quad \operatorname{trace}(\operatorname{juxta}(p, q, a, b, c), a, c) \\ \quad \quad \subseteq\; (\operatorname{trace}(p, a, b) \cup \operatorname{trace}(q, b, c)) \end{array}}}
\vnbentry{\texttt{juxta-\allowbreak{}trace}}{\vnbstmt{\begin{array}{@{}l@{}} \forall p, q.; \forall a, b, c \in \mathbb{R}. \\ \quad (a \leq b) \wedge \\ \quad (b \leq c) \wedge \\ \quad (p \in (\operatorname{ccint}(a, b) \to \mathbb{C})) \wedge \\ \quad (q \in (\operatorname{ccint}(b, c) \to \mathbb{C})) \wedge \\ \quad (p(b) = q(b)) \Rightarrow \\ \quad \quad \operatorname{trace}(\operatorname{juxta}(p, q, a, b, c), a, c) \\ \quad \quad =\; (\operatorname{trace}(p, a, b) \cup \operatorname{trace}(q, b, c)) \end{array}}}
\vnbentry{\texttt{right-\allowbreak{}limit-\allowbreak{}within-\allowbreak{}local}}{\vnbstmt{\begin{array}{@{}l@{}} \forall g, v, f, w, x, l, r. \\ \quad \operatorname{is\text{-}right\text{-}limit\text{-}within}(g, v, x, l) \wedge \\ \quad (w \subseteq \mathbb{R}) \wedge \\ \quad (f \in (w \to \mathbb{R})) \wedge \\ \quad \operatorname{pos\text{-}rr}(r) \wedge \\ \quad \forall t \in w.\; (((x < t) \wedge (t < (x + r))) \Rightarrow (t \in v)) \wedge \\ \quad \forall t \in w. \\ \quad \quad (((x < t) \wedge (t < (x + r))) \Rightarrow (f(t) \simeq g(t))) \Rightarrow \\ \quad \quad \operatorname{is\text{-}right\text{-}limit\text{-}within}(f, w, x, l) \end{array}}}
\vnbentry{\texttt{left-\allowbreak{}limit-\allowbreak{}within-\allowbreak{}local}}{\vnbstmt{\begin{array}{@{}l@{}} \forall g, v, f, w, x, l, r. \\ \quad \operatorname{is\text{-}left\text{-}limit\text{-}within}(g, v, x, l) \wedge \\ \quad (w \subseteq \mathbb{R}) \wedge \\ \quad (f \in (w \to \mathbb{R})) \wedge \\ \quad \operatorname{pos\text{-}rr}(r) \wedge \\ \quad \forall t \in w.\; ((((x - r) < t) \wedge (t < x)) \Rightarrow (t \in v)) \wedge \\ \quad \forall t \in w. \\ \quad \quad ((((x - r) < t) \wedge (t < x)) \Rightarrow (f(t) \simeq g(t))) \Rightarrow \\ \quad \quad \operatorname{is\text{-}left\text{-}limit\text{-}within}(f, w, x, l) \end{array}}}
\vnbentry{\texttt{regulated-\allowbreak{}on-\allowbreak{}glue}}{\vnbstmt{\begin{array}{@{}l@{}} \forall a, c, b \in \mathbb{R}, d, g, f. \\ \quad (a < c) \wedge \\ \quad (c < b) \wedge \\ \quad \operatorname{is\text{-}regulated\text{-}on}(d, a, c) \wedge \\ \quad \operatorname{is\text{-}regulated\text{-}on}(g, c, b) \wedge \\ \quad (f \in (\operatorname{ccint}(a, b) \to \mathbb{R})) \wedge \\ \quad \forall t \in \operatorname{ccint}(a, c). \\ \quad \quad ((t < c) \Rightarrow (f(t) \simeq d(t))) \wedge \\ \quad \forall t \in \operatorname{ccint}(c, b). \\ \quad \quad ((c < t) \Rightarrow (f(t) \simeq g(t))) \Rightarrow \\ \quad \quad \operatorname{is\text{-}regulated\text{-}on}(f, a, b) \end{array}}}
\vnbentry{\texttt{juxta-\allowbreak{}is-\allowbreak{}road}}{\vnbstmt{\begin{array}{@{}l@{}} \forall m_{1}, d, m_{2}, f, a, b, c. \\ \quad \operatorname{is\text{-}road}(m_{1}, d, a, b) \wedge \\ \quad \operatorname{is\text{-}road}(m_{2}, f, b, c) \wedge \\ \quad (\operatorname{m1}(b) = \operatorname{m2}(b)) \Rightarrow \\ \quad \quad \operatorname{is\text{-}road}(\operatorname{juxta}(m_{1}, m_{2}, a, b, c), \operatorname{juxta}(d, f, a, b, c), a, c) \end{array}}}
\vnbentry{\texttt{line-\allowbreak{}int-\allowbreak{}interior-\allowbreak{}congruence}}{\vnbstmt{\begin{array}{@{}l@{}} \forall u, f, m, c, a, b.; \forall d, g \in (\operatorname{ccint}(a, b) \to \mathbb{C}). \\ \quad (u \subseteq \mathbb{C}) \wedge \\ \quad (f \in (u \to \mathbb{C})) \wedge \\ \quad \forall y \in u. \\ \quad \quad \operatorname{is\text{-}continuous\text{-}at}(\operatorname{subspace\text{-}ms}(\operatorname{nf\text{-}metric\text{-}space}(\operatorname{cc\text{-}normed\text{-}field}), u), \\ \quad \quad \quad \operatorname{nf\text{-}metric\text{-}space}(\operatorname{cc\text{-}normed\text{-}field}), f, y) \wedge \\ \quad \operatorname{is\text{-}road}(m, c, a, b) \wedge \\ \quad (\operatorname{trace}(m, a, b) \subseteq u) \wedge \\ \quad \forall t \in \operatorname{ccint}(a, b).\; (d(t) \simeq m(t)) \wedge \\ \quad \forall t \in \operatorname{ooint}(a, b).\; (g(t) \simeq c(t)) \Rightarrow \\ \quad \quad (\operatorname{line\text{-}int}(f, d, g, a, b) = \operatorname{line\text{-}int}(f, m, c, a, b)) \end{array}}}
\vnbentry{\texttt{line-\allowbreak{}int-\allowbreak{}juxta}}{\vnbstmt{\begin{array}{@{}l@{}} \forall u, f, m_{1}, d, m_{2}, g, a, b, c. \\ \quad (u \subseteq \mathbb{C}) \wedge \\ \quad (f \in (u \to \mathbb{C})) \wedge \\ \quad \forall y \in u. \\ \quad \quad \operatorname{is\text{-}continuous\text{-}at}(\operatorname{subspace\text{-}ms}(\operatorname{nf\text{-}metric\text{-}space}(\operatorname{cc\text{-}normed\text{-}field}), u), \\ \quad \quad \quad \operatorname{nf\text{-}metric\text{-}space}(\operatorname{cc\text{-}normed\text{-}field}), f, y) \wedge \\ \quad \operatorname{is\text{-}road}(m_{1}, d, a, b) \wedge \\ \quad \operatorname{is\text{-}road}(m_{2}, g, b, c) \wedge \\ \quad (\operatorname{m1}(b) = \operatorname{m2}(b)) \wedge \\ \quad (\operatorname{trace}(m_{1}, a, b) \subseteq u) \wedge \\ \quad (\operatorname{trace}(m_{2}, b, c) \subseteq u) \Rightarrow \\ \quad \quad \operatorname{line\text{-}int}(f, \operatorname{juxta}(m_{1}, m_{2}, a, b, c), \operatorname{juxta}(d, g, a, b, c), a, c) \\ \quad \quad =\; (\operatorname{line\text{-}int}(f, m_{1}, d, a, b) + \operatorname{line\text{-}int}(f, m_{2}, g, b, c)) \end{array}}}
\vnbentry{\texttt{line-\allowbreak{}int-\allowbreak{}concat}}{\vnbstmt{\begin{array}{@{}l@{}} \forall u, f, m, g, o, h, a, b, c, d. \\ \quad (u \subseteq \mathbb{C}) \wedge \\ \quad (f \in (u \to \mathbb{C})) \wedge \\ \quad \forall y \in u. \\ \quad \quad \operatorname{is\text{-}continuous\text{-}at}(\operatorname{subspace\text{-}ms}(\operatorname{nf\text{-}metric\text{-}space}(\operatorname{cc\text{-}normed\text{-}field}), u), \\ \quad \quad \quad \operatorname{nf\text{-}metric\text{-}space}(\operatorname{cc\text{-}normed\text{-}field}), f, y) \wedge \\ \quad \operatorname{is\text{-}road}(m, g, a, b) \wedge \\ \quad \operatorname{is\text{-}road}(o, h, c, d) \wedge \\ \quad (m(b) = o(c)) \wedge \\ \quad (\operatorname{trace}(m, a, b) \subseteq u) \wedge \\ \quad (\operatorname{trace}(o, c, d) \subseteq u) \Rightarrow \\ \quad \quad \operatorname{line\text{-}int}(f, \\ \quad \quad \quad \operatorname{juxta}(m, (\lambda s \in \operatorname{ccint}(b, (b + (d - c))).\; o((c + (s - b)))), a, \\ \quad \quad \quad \quad b, (b + (d - c))), \\ \quad \quad \quad \operatorname{juxta}(g, (\lambda s \in \operatorname{ccint}(b, (b + (d - c))).\; h((c + (s - b)))), a, \\ \quad \quad \quad \quad b, (b + (d - c))), \\ \quad \quad \quad a, (b + (d - c))) \\ \quad \quad =\; (\operatorname{line\text{-}int}(f, m, g, a, b) + \operatorname{line\text{-}int}(f, o, h, c, d)) \end{array}}}
\vnbentry{\texttt{road-\allowbreak{}length-\allowbreak{}regulated}}{\vnbstmt{\begin{array}{@{}l@{}} \forall m, c, a, b. \\ \quad \operatorname{is\text{-}road}(m, c, a, b) \Rightarrow  \\ \quad \quad \operatorname{is\text{-}regulated\text{-}on}((\lambda t \in \operatorname{ccint}(a, b).\; \operatorname{magnitude}(c(t))), a, b) \end{array}}}
\vnbentry{\texttt{road-\allowbreak{}length-\allowbreak{}primitive}}{\vnbstmt{\begin{array}{@{}l@{}} \forall m, c, a, b. \\ \quad \operatorname{is\text{-}road}(m, c, a, b) \Rightarrow  \\ \quad \quad \exists g. \\ \quad \quad \quad \operatorname{is\text{-}primitive}(g, (\lambda t \in \operatorname{ccint}(a, b).\; \operatorname{magnitude}(c(t))), a, b) \end{array}}}
\vnbentry{\texttt{road-\allowbreak{}length-\allowbreak{}in-\allowbreak{}rr}}{\vnbstmt{\begin{array}{@{}l@{}} \forall m, c, a, b. \\ \quad \operatorname{is\text{-}road}(m, c, a, b) \Rightarrow  \\ \quad \quad (\operatorname{road\text{-}length}(c, a, b) \in \mathbb{R}) \end{array}}}
\vnbentry{\texttt{road-\allowbreak{}length-\allowbreak{}nonneg}}{\vnbstmt{\begin{array}{@{}l@{}} \forall m, c, a, b. \\ \quad \operatorname{is\text{-}road}(m, c, a, b) \Rightarrow  \\ \quad \quad (0 \leq \operatorname{road\text{-}length}(c, a, b)) \end{array}}}
\vnbentry{\texttt{pw-\allowbreak{}int-\allowbreak{}interior-\allowbreak{}congruence}}{\vnbstmt{\begin{array}{@{}l@{}} \forall f, i, \mathit{ppsi}, a, b. \\ \quad \operatorname{is\text{-}primitive}(f, i, a, b) \wedge \\ \quad (\mathit{ppsi} \in (\operatorname{ccint}(a, b) \to \mathbb{R})) \wedge \\ \quad \forall y \in \operatorname{ooint}(a, b).\; (i(y) \simeq \operatorname{ppsi}(y)) \Rightarrow \\ \quad \quad (\int_{a}^{b} i = \int_{a}^{b} \mathit{ppsi}) \end{array}}}
\vnbentry{\texttt{road-\allowbreak{}length-\allowbreak{}juxta}}{\vnbstmt{\begin{array}{@{}l@{}} \forall m_{1}, d, m_{2}, f, a, b, c. \\ \quad \operatorname{is\text{-}road}(m_{1}, d, a, b) \wedge \\ \quad \operatorname{is\text{-}road}(m_{2}, f, b, c) \wedge \\ \quad (\operatorname{m1}(b) = \operatorname{m2}(b)) \Rightarrow \\ \quad \quad \operatorname{road\text{-}length}(\operatorname{juxta}(d, f, a, b, c), a, c) \\ \quad \quad =\; (\operatorname{road\text{-}length}(d, a, b) + \operatorname{road\text{-}length}(f, b, c)) \end{array}}}
\vnbentry{\texttt{seg-\allowbreak{}shift-\allowbreak{}is-\allowbreak{}road}}{\vnbstmt{\begin{array}{@{}l@{}} \forall p, q \in \mathbb{C}, k, m \in \mathbb{R}. \\ \quad (m = (k + 1)) \Rightarrow  \\ \quad \quad \operatorname{is\text{-}road}((\lambda s \in \operatorname{ccint}(k, m).\; \operatorname{seg\text{-}path}(p, q)((s - k))), \\ \quad \quad \quad (\lambda s \in \operatorname{ccint}(k, m).\; \operatorname{seg\text{-}deriv}(p, q)((s - k))), k, m) \end{array}}}
\vnbentry{\texttt{seg-\allowbreak{}shift-\allowbreak{}trace}}{\vnbstmt{\begin{array}{@{}l@{}} \forall p, q \in \mathbb{C}, k, m \in \mathbb{R}. \\ \quad (m = (k + 1)) \Rightarrow  \\ \quad \quad \operatorname{trace}((\lambda s \in \operatorname{ccint}(k, m).\; \operatorname{seg\text{-}path}(p, q)((s - k))), k, m) \\ \quad \quad =\; \operatorname{trace}(\operatorname{seg\text{-}path}(p, q), 0, 1) \end{array}}}
\vnbentry{\texttt{road-\allowbreak{}length-\allowbreak{}segment}}{\vnbstmt{\begin{array}{@{}l@{}} \forall p, q \in \mathbb{C}. \\ \quad \operatorname{road\text{-}length}(\operatorname{seg\text{-}deriv}(p, q), 0, 1) \\ \quad =\; \operatorname{magnitude}((q - p)) \end{array}}}
\vnbentry{\texttt{seg-\allowbreak{}shift-\allowbreak{}length}}{\vnbstmt{\begin{array}{@{}l@{}} \forall p, q \in \mathbb{C}, k, m \in \mathbb{R}. \\ \quad (m = (k + 1)) \Rightarrow  \\ \quad \quad \operatorname{road\text{-}length}((\lambda s \in \operatorname{ccint}(k, m).\; \operatorname{seg\text{-}deriv}(p, q)((s - k))), k, m) \\ \quad \quad =\; \operatorname{magnitude}((q - p)) \end{array}}}
\vnbentry{\texttt{seg-\allowbreak{}shift-\allowbreak{}line-\allowbreak{}int}}{\vnbstmt{\begin{array}{@{}l@{}} \forall u, f.; \forall p, q \in \mathbb{C}, k, m \in \mathbb{R}. \\ \quad (u \subseteq \mathbb{C}) \wedge \\ \quad (f \in (u \to \mathbb{C})) \wedge \\ \quad \forall y \in u. \\ \quad \quad \operatorname{is\text{-}continuous\text{-}at}(\operatorname{subspace\text{-}ms}(\operatorname{nf\text{-}metric\text{-}space}(\operatorname{cc\text{-}normed\text{-}field}), u), \\ \quad \quad \quad \operatorname{nf\text{-}metric\text{-}space}(\operatorname{cc\text{-}normed\text{-}field}), f, y) \wedge \\ \quad (m = (k + 1)) \wedge \\ \quad (\operatorname{trace}(\operatorname{seg\text{-}path}(p, q), 0, 1) \subseteq u) \Rightarrow \\ \quad \quad \operatorname{line\text{-}int}(f, (\lambda s \in \operatorname{ccint}(k, m).\; \operatorname{seg\text{-}path}(p, q)((s - k))), \\ \quad \quad \quad (\lambda s \in \operatorname{ccint}(k, m).\; \operatorname{seg\text{-}deriv}(p, q)((s - k))), k, m) \\ \quad \quad =\; \operatorname{line\text{-}int}(f, \operatorname{seg\text{-}path}(p, q), \operatorname{seg\text{-}deriv}(p, q), 0, 1) \end{array}}}
\vnbentry{\texttt{tri-\allowbreak{}road-\allowbreak{}is-\allowbreak{}road}}{\vnbstmt{\begin{array}{@{}l@{}} \forall a, b, c \in \mathbb{C}. \\ \quad \operatorname{is\text{-}road}(\operatorname{tri\text{-}road}(a, b, c), \operatorname{tri\text{-}road\text{-}deriv}(a, b, c), 0, 3) \end{array}}}
\vnbentry{\texttt{tri-\allowbreak{}road-\allowbreak{}trace}}{\vnbstmt{\begin{array}{@{}l@{}} \forall a, b, c \in \mathbb{C}. \\ \quad \operatorname{trace}(\operatorname{tri\text{-}road}(a, b, c), 0, 3) \\ \quad =\; ((\operatorname{trace}(\operatorname{seg\text{-}path}(a, b), 0, 1) \cup \operatorname{trace}(\operatorname{seg\text{-}path}(b, c), 0, 1)) \cup \operatorname{trace}(\operatorname{seg\text{-}path}(c, a), 0, 1)) \end{array}}}
\vnbentry{\texttt{tri-\allowbreak{}int-\allowbreak{}is-\allowbreak{}line-\allowbreak{}int}}{\vnbstmt{\begin{array}{@{}l@{}} \forall u, f.; \forall a, b, c \in \mathbb{C}. \\ \quad (u \subseteq \mathbb{C}) \wedge \\ \quad (f \in (u \to \mathbb{C})) \wedge \\ \quad \forall y \in u. \\ \quad \quad \operatorname{is\text{-}continuous\text{-}at}(\operatorname{subspace\text{-}ms}(\operatorname{nf\text{-}metric\text{-}space}(\operatorname{cc\text{-}normed\text{-}field}), u), \\ \quad \quad \quad \operatorname{nf\text{-}metric\text{-}space}(\operatorname{cc\text{-}normed\text{-}field}), f, y) \wedge \\ \quad (\operatorname{trace}(\operatorname{seg\text{-}path}(a, b), 0, 1) \subseteq u) \wedge \\ \quad (\operatorname{trace}(\operatorname{seg\text{-}path}(b, c), 0, 1) \subseteq u) \wedge \\ \quad (\operatorname{trace}(\operatorname{seg\text{-}path}(c, a), 0, 1) \subseteq u) \Rightarrow \\ \quad \quad \operatorname{tri\text{-}int}(f, a, b, c) \\ \quad \quad =\; \operatorname{line\text{-}int}(f, \operatorname{tri\text{-}road}(a, b, c), \operatorname{tri\text{-}road\text{-}deriv}(a, b, c), 0, 3) \end{array}}}
\vnbentry{\texttt{road-\allowbreak{}length-\allowbreak{}tri}}{\vnbstmt{\begin{array}{@{}l@{}} \forall a, b, c \in \mathbb{C}. \\ \quad \operatorname{road\text{-}length}(\operatorname{tri\text{-}road\text{-}deriv}(a, b, c), 0, 3) \\ \quad =\; ((\operatorname{magnitude}((b - a)) + \operatorname{magnitude}((c - b))) + \operatorname{magnitude}((a - c))) \end{array}}}
\vnbentry{\texttt{line-\allowbreak{}int-\allowbreak{}length-\allowbreak{}bound}}{\vnbstmt{\begin{array}{@{}l@{}} \forall u, f, m, c, a, b.; \forall d \in \mathbb{R}. \\ \quad (u \subseteq \mathbb{C}) \wedge \\ \quad (f \in (u \to \mathbb{C})) \wedge \\ \quad \forall y \in u. \\ \quad \quad \operatorname{is\text{-}continuous\text{-}at}(\operatorname{subspace\text{-}ms}(\operatorname{nf\text{-}metric\text{-}space}(\operatorname{cc\text{-}normed\text{-}field}), u), \\ \quad \quad \quad \operatorname{nf\text{-}metric\text{-}space}(\operatorname{cc\text{-}normed\text{-}field}), f, y) \wedge \\ \quad \operatorname{is\text{-}road}(m, c, a, b) \wedge \\ \quad (\operatorname{trace}(m, a, b) \subseteq u) \wedge \\ \quad \forall \mathit{bny} \in \operatorname{trace}(m, a, b). \\ \quad \quad (\operatorname{magnitude}(f(\mathit{bny})) \leq d) \Rightarrow \\ \quad \quad \operatorname{magnitude}(\operatorname{line\text{-}int}(f, m, c, a, b)) \\ \quad \quad \leq\; (\operatorname{road\text{-}length}(c, a, b) \cdot d) \end{array}}}
\vnbentry{\texttt{road-\allowbreak{}length-\allowbreak{}affine-\allowbreak{}reparam}}{\vnbstmt{\begin{array}{@{}l@{}} \forall a, b, c, d, l, m \in \mathbb{R}, f, g, k. \\ \quad (0 < l) \wedge \\ \quad ((m + (l \cdot c)) = a) \wedge \\ \quad ((m + (l \cdot d)) = b) \wedge \\ \quad (c < d) \wedge \\ \quad \operatorname{is\text{-}road}(f, g, a, b) \wedge \\ \quad (k \in (\operatorname{ccint}(c, d) \to \mathbb{C})) \wedge \\ \quad \forall z \in \operatorname{ccint}(c, d). \\ \quad \quad (k(z) \simeq (l \cdot g((m + (l \cdot z))))) \Rightarrow \\ \quad \quad (\operatorname{road\text{-}length}(k, c, d) = \operatorname{road\text{-}length}(g, a, b)) \end{array}}}
\vnbentry{\texttt{tri-\allowbreak{}int-\allowbreak{}is-\allowbreak{}line-\allowbreak{}int-\allowbreak{}conv3}}{\vnbstmt{\begin{array}{@{}l@{}} \forall u, f.; \forall a, b, c \in \mathbb{C}. \\ \quad (u \subseteq \mathbb{C}) \wedge \\ \quad (f \in (u \to \mathbb{C})) \wedge \\ \quad \forall y \in u. \\ \quad \quad \operatorname{is\text{-}continuous\text{-}at}(\operatorname{subspace\text{-}ms}(\operatorname{nf\text{-}metric\text{-}space}(\operatorname{cc\text{-}normed\text{-}field}), u), \\ \quad \quad \quad \operatorname{nf\text{-}metric\text{-}space}(\operatorname{cc\text{-}normed\text{-}field}), f, y) \wedge \\ \quad (\operatorname{conv3}(a, b, c) \subseteq u) \Rightarrow \\ \quad \quad \operatorname{tri\text{-}int}(f, a, b, c) \\ \quad \quad =\; \operatorname{line\text{-}int}(f, \operatorname{tri\text{-}road}(a, b, c), \operatorname{tri\text{-}road\text{-}deriv}(a, b, c), 0, 3) \end{array}}}

\vnbfile{theorem-library/cc-exp.scm}{theorem-library/cc-exp.scm}
\noindent{\small cc-exp.scm -- THE ELEMENTARY FUNCTIONS exp, cos, sin ON CC (the notes 2.3.1, 2.3.2, everything that does not mention pi). Batch 31-A, 2026-09-25.  The definitions are structure-library/cc-elementary.scm (CC-EXP, CC-COS, CC-SIN by their power series, the user's decision of 2026-09-25).}\par
\vnbentry{\texttt{cc-\allowbreak{}exp-\allowbreak{}unfold}}{\vnbstmt{\begin{array}{@{}l@{}} \forall z. \\ \quad (\exp z \simeq \operatorname{cc\text{-}series\text{-}limit}((\lambda k \in \mathbb{N}.\; (\operatorname{recip}(\operatorname{factorial}(k)) \cdot \operatorname{power}(z, k))))) \end{array}}}
\vnbentry{\texttt{cc-\allowbreak{}cos-\allowbreak{}unfold}}{\vnbstmt{\begin{array}{@{}l@{}} \forall z. \\ \quad (\cos z \simeq \operatorname{cc\text{-}series\text{-}limit}((\lambda k \in \mathbb{N}.\; ((\operatorname{power}(-1, k) \cdot \operatorname{recip}(\operatorname{factorial}((2 \cdot k)))) \cdot \operatorname{power}(z, (2 \cdot k)))))) \end{array}}}
\vnbentry{\texttt{cc-\allowbreak{}sin-\allowbreak{}unfold}}{\vnbstmt{\begin{array}{@{}l@{}} \forall z. \\ \quad (\sin z \simeq \operatorname{cc\text{-}series\text{-}limit}((\lambda k \in \mathbb{N}.\; ((\operatorname{power}(-1, k) \cdot \operatorname{recip}(\operatorname{factorial}(((2 \cdot k) + 1)))) \cdot \operatorname{power}(z, ((2 \cdot k) + 1)))))) \end{array}}}
\vnbentry{\texttt{cps-\allowbreak{}radius-\allowbreak{}pos-\allowbreak{}inf}}{\vnbstmt{\begin{array}{@{}l@{}} \forall f \in (\mathbb{N} \to \mathbb{C}). \\ \quad \forall t \in \mathbb{R}. \\ \quad \quad (0 < t) \Rightarrow  \\ \quad \quad \quad \operatorname{series\text{-}converges}((\lambda k \in \mathbb{N}.\; (\operatorname{magnitude}(f(k)) \cdot \operatorname{power}(t, k)))) \Rightarrow \\ \quad \quad (\operatorname{cps\text{-}radius}(f) = \operatorname{pos\text{-}inf}) \end{array}}}
\vnbentry{\texttt{cex-\allowbreak{}exp-\allowbreak{}series-\allowbreak{}real-\allowbreak{}converges}}{\vnbstmt{\begin{array}{@{}l@{}} \forall t \in \mathbb{R}. \\ \quad (0 < t) \Rightarrow  \\ \quad \quad \operatorname{series\text{-}converges}((\lambda k \in \mathbb{N}.\; \\ \quad \quad \quad \quad (\operatorname{recip}(\operatorname{factorial}(k)) \cdot \operatorname{power}(t, k)))) \end{array}}}
\vnbentry{\texttt{cc-\allowbreak{}exp-\allowbreak{}radius}}{\vnbstmt{\begin{array}{@{}l@{}} \forall f \in (\mathbb{N} \to \mathbb{C}). \\ \quad \forall k \in \mathbb{N}.\; (f(k) = \operatorname{recip}(\operatorname{factorial}(k))) \Rightarrow  \\ \quad \quad (\operatorname{cps\text{-}radius}(f) = \operatorname{pos\text{-}inf}) \end{array}}}
\vnbentry{\texttt{cps-\allowbreak{}entire-\allowbreak{}converges}}{\vnbstmt{\begin{array}{@{}l@{}} \forall f \in (\mathbb{N} \to \mathbb{C}), z \in \mathbb{C}. \\ \quad (\operatorname{cps\text{-}radius}(f) = \operatorname{pos\text{-}inf}) \Rightarrow  \\ \quad \quad \operatorname{cc\text{-}series\text{-}converges}((\lambda k \in \mathbb{N}.\; (f(k) \cdot \operatorname{power}(z, k)))) \end{array}}}
\vnbentry{\texttt{cc-\allowbreak{}diff-\allowbreak{}on-\allowbreak{}local}}{\vnbstmt{\begin{array}{@{}l@{}} \forall u.; \forall f \in (u \to \mathbb{C}), v, g, a, l. \\ \quad \operatorname{is\text{-}open}(\operatorname{nf\text{-}metric\text{-}space}(\operatorname{cc\text{-}normed\text{-}field}), u) \wedge \\ \quad (v \subseteq u) \wedge \\ \quad \operatorname{is\text{-}diff\text{-}on}(\operatorname{cc\text{-}normed\text{-}field}, v, g, a, l) \wedge \\ \quad \forall x \in v.\; (f(x) = g(x)) \Rightarrow \\ \quad \quad \operatorname{is\text{-}diff\text{-}on}(\operatorname{cc\text{-}normed\text{-}field}, u, f, a, l) \end{array}}}
\vnbentry{\texttt{cps-\allowbreak{}entire-\allowbreak{}diff}}{\vnbstmt{\begin{array}{@{}l@{}} \forall f, a \in (\mathbb{N} \to \mathbb{C}), \mathit{cpda} \in \mathbb{C}. \\ \quad \forall j \in \mathbb{N}.\; (a(j) = ((j + 1) \cdot f((j + 1)))) \wedge \\ \quad (\operatorname{cps\text{-}radius}(f) = \operatorname{pos\text{-}inf}) \Rightarrow \\ \quad \quad \operatorname{is\text{-}diff\text{-}on}(\operatorname{cc\text{-}normed\text{-}field}, \mathbb{C}, \\ \quad \quad \quad (\lambda z \in \mathbb{C}.\; \\ \quad \quad \quad \quad \operatorname{cc\text{-}series\text{-}limit}((\lambda k \in \mathbb{N}.\; (f(k) \cdot \operatorname{power}(z, k))))), \\ \quad \quad \quad \mathit{cpda}, \\ \quad \quad \quad \operatorname{cc\text{-}series\text{-}limit}((\lambda k \in \mathbb{N}.\; (a(k) \cdot \operatorname{power}(\mathit{cpda}, k))))) \end{array}}}
\vnbentry{\texttt{cc-\allowbreak{}exp-\allowbreak{}as-\allowbreak{}power-\allowbreak{}series}}{\vnbstmt{\begin{array}{@{}l@{}} \forall f \in (\mathbb{N} \to \mathbb{C}), z \in \mathbb{C}. \\ \quad \forall k \in \mathbb{N}.\; (f(k) = \operatorname{recip}(\operatorname{factorial}(k))) \Rightarrow  \\ \quad \quad \operatorname{cc\text{-}series\text{-}converges}((\lambda k \in \mathbb{N}.\; \\ \quad \quad \quad \quad (\operatorname{recip}(\operatorname{factorial}(k)) \cdot \operatorname{power}(z, k)))) \wedge \\ \quad \quad \exp z \\ \quad \quad =\; \operatorname{cc\text{-}series\text{-}limit}((\lambda \mathit{cpsk} \in \mathbb{N}.\; \\ \quad \quad \quad \quad (f(\mathit{cpsk}) \cdot \operatorname{power}(z, \mathit{cpsk})))) \end{array}}}
\vnbentry{\texttt{cc-\allowbreak{}exp-\allowbreak{}converges}}{\vnbstmt{\begin{array}{@{}l@{}} \forall z \in \mathbb{C}. \\ \quad \operatorname{cc\text{-}series\text{-}converges}((\lambda k \in \mathbb{N}.\; \\ \quad \quad \quad (\operatorname{recip}(\operatorname{factorial}(k)) \cdot \operatorname{power}(z, k)))) \wedge \\ \quad (\exp z \in \mathbb{C}) \end{array}}}
\vnbentry{\texttt{cc-\allowbreak{}exp-\allowbreak{}deriv}}{\vnbstmt{\begin{array}{@{}l@{}} \forall z \in \mathbb{C}. \\ \quad \operatorname{is\text{-}diff\text{-}on}(\operatorname{cc\text{-}normed\text{-}field}, \mathbb{C}, (\lambda y \in \mathbb{C}.\; \exp y), z, \exp z) \end{array}}}
\vnbentry{\texttt{cc-\allowbreak{}exp-\allowbreak{}holomorphic}}{\vnbstmt{\operatorname{holomorphic\text{-}on}(\mathbb{C}, (\lambda y \in \mathbb{C}.\; \exp y))}}
\vnbentry{\texttt{cc-\allowbreak{}exp-\allowbreak{}zero}}{\vnbstmt{(\exp 0 = 1)}}
\vnbentry{\texttt{cex-\allowbreak{}sub-\allowbreak{}deriv}}{\vnbstmt{\begin{array}{@{}l@{}} \forall w, z \in \mathbb{C}. \\ \quad \operatorname{is\text{-}diff\text{-}on}(\operatorname{cc\text{-}normed\text{-}field}, \mathbb{C}, (\lambda y \in \mathbb{C}.\; (w - y)), z, -1) \end{array}}}
\vnbentry{\texttt{cex-\allowbreak{}shift-\allowbreak{}deriv-\allowbreak{}zero}}{\vnbstmt{\begin{array}{@{}l@{}} \forall w, z \in \mathbb{C}. \\ \quad \operatorname{is\text{-}diff\text{-}on}(\operatorname{cc\text{-}normed\text{-}field}, \mathbb{C}, \\ \quad \quad (\lambda y \in \mathbb{C}.\; (\exp y \cdot \exp (w - y))), z, 0) \end{array}}}
\vnbentry{\texttt{cc-\allowbreak{}exp-\allowbreak{}shift}}{\vnbstmt{\forall w, z \in \mathbb{C}.\; ((\exp z \cdot \exp (w - z)) = \exp w)}}
\vnbentry{\texttt{cc-\allowbreak{}exp-\allowbreak{}add}}{\vnbstmt{\forall z, w \in \mathbb{C}.\; (\exp (z + w) = (\exp z \cdot \exp w))}}
\vnbentry{\texttt{cc-\allowbreak{}exp-\allowbreak{}nonzero}}{\vnbstmt{\forall z \in \mathbb{C}.\; \neg (\exp z = 0)}}
\vnbentry{\texttt{cc-\allowbreak{}exp-\allowbreak{}maps-\allowbreak{}to-\allowbreak{}nonzero}}{\vnbstmt{\begin{array}{@{}l@{}} \forall z \in \mathbb{C}. \\ \quad (\exp z \in \operatorname{difference}(\mathbb{C}, \operatorname{singleton}(0))) \end{array}}}
\vnbentry{\texttt{cc-\allowbreak{}exp-\allowbreak{}neg}}{\vnbstmt{\forall z \in \mathbb{C}.\; (\exp -z = \operatorname{recip}(\exp z))}}
\vnbentry{\texttt{cc-\allowbreak{}exp-\allowbreak{}continuous}}{\vnbstmt{\begin{array}{@{}l@{}} \forall z \in \mathbb{C}. \\ \quad \operatorname{is\text{-}continuous\text{-}at}(\operatorname{subspace\text{-}ms}(\operatorname{nf\text{-}metric\text{-}space}(\operatorname{cc\text{-}normed\text{-}field}), \mathbb{C}), \\ \quad \quad \operatorname{nf\text{-}metric\text{-}space}(\operatorname{cc\text{-}normed\text{-}field}), (\lambda y \in \mathbb{C}.\; \exp y), z) \end{array}}}
\vnbentry{\texttt{cc-\allowbreak{}conjugate-\allowbreak{}sub}}{\vnbstmt{\begin{array}{@{}l@{}} \forall a, b \in \mathbb{C}. \\ \quad (\operatorname{conjugate}((a - b)) = (\operatorname{conjugate}(a) - \operatorname{conjugate}(b))) \end{array}}}
\vnbentry{\texttt{cc-\allowbreak{}magnitude-\allowbreak{}conjugate}}{\vnbstmt{\begin{array}{@{}l@{}} \forall z \in \mathbb{C}. \\ \quad (\operatorname{magnitude}(\operatorname{conjugate}(z)) = \operatorname{magnitude}(z)) \end{array}}}
\vnbentry{\texttt{cc-\allowbreak{}conjugate-\allowbreak{}power}}{\vnbstmt{\begin{array}{@{}l@{}} \forall k \in \mathbb{N}, z \in \mathbb{C}. \\ \quad \operatorname{conjugate}(\operatorname{power}(z, k)) \\ \quad =\; \operatorname{power}(\operatorname{conjugate}(z), k) \end{array}}}
\vnbentry{\texttt{cc-\allowbreak{}series-\allowbreak{}limit-\allowbreak{}conjugate}}{\vnbstmt{\begin{array}{@{}l@{}} \forall f, g \in (\mathbb{N} \to \mathbb{C}). \\ \quad \forall j \in \mathbb{N}.\; (g(j) = \operatorname{conjugate}(f(j))) \wedge \\ \quad \operatorname{cc\text{-}series\text{-}converges}(f) \Rightarrow \\ \quad \quad \operatorname{cc\text{-}series\text{-}converges}(g) \wedge \\ \quad \quad \operatorname{cc\text{-}series\text{-}limit}(g) \\ \quad \quad =\; \operatorname{conjugate}(\operatorname{cc\text{-}series\text{-}limit}(f)) \end{array}}}
\vnbentry{\texttt{cc-\allowbreak{}exp-\allowbreak{}conj}}{\vnbstmt{\begin{array}{@{}l@{}} \forall z \in \mathbb{C}. \\ \quad (\exp \operatorname{conjugate}(z) = \operatorname{conjugate}(\exp z)) \end{array}}}
\vnbentry{\texttt{cc-\allowbreak{}magnitude-\allowbreak{}sq-\allowbreak{}conjugate}}{\vnbstmt{\begin{array}{@{}l@{}} \forall z \in \mathbb{C}. \\ \quad (\operatorname{magnitude}(z) \cdot \operatorname{magnitude}(z)) \\ \quad =\; (z \cdot \operatorname{conjugate}(z)) \end{array}}}
\vnbentry{\texttt{cc-\allowbreak{}exp-\allowbreak{}real-\allowbreak{}valued}}{\vnbstmt{\forall t \in \mathbb{R}.\; (\exp t \in \mathbb{R})}}
\vnbentry{\texttt{cc-\allowbreak{}exp-\allowbreak{}real}}{\vnbstmt{\forall t \in \mathbb{R}.\; ((\exp t \in \mathbb{R}) \wedge (0 < \exp t))}}
\vnbentry{\texttt{cc-\allowbreak{}exp-\allowbreak{}magnitude-\allowbreak{}imaginary}}{\vnbstmt{\forall t \in \mathbb{R}.\; (\operatorname{magnitude}(\exp (+i \cdot t)) = 1)}}
\vnbentry{\texttt{cc-\allowbreak{}power-\allowbreak{}add}}{\vnbstmt{\begin{array}{@{}l@{}} \forall n, m \in \mathbb{N}, z \in \mathbb{C}. \\ \quad (\operatorname{power}(z, (m + n)) = (\operatorname{power}(z, m) \cdot \operatorname{power}(z, n))) \end{array}}}
\vnbentry{\texttt{cc-\allowbreak{}power-\allowbreak{}mul}}{\vnbstmt{\begin{array}{@{}l@{}} \forall k, m \in \mathbb{N}, z \in \mathbb{C}. \\ \quad (\operatorname{power}(z, (m \cdot k)) = \operatorname{power}(\operatorname{power}(z, m), k)) \end{array}}}
\vnbentry{\texttt{cc-\allowbreak{}power-\allowbreak{}mul-\allowbreak{}base}}{\vnbstmt{\begin{array}{@{}l@{}} \forall k \in \mathbb{N}, a, b \in \mathbb{C}. \\ \quad \operatorname{power}((a \cdot b), k) \\ \quad =\; (\operatorname{power}(a, k) \cdot \operatorname{power}(b, k)) \end{array}}}
\vnbentry{\texttt{factorial-\allowbreak{}mono}}{\vnbstmt{\begin{array}{@{}l@{}} \forall m, k \in \mathbb{N}. \\ \quad (\operatorname{factorial}(k) \leq \operatorname{factorial}((k + m))) \end{array}}}
\vnbentry{\texttt{cex-\allowbreak{}sign-\allowbreak{}magnitude}}{\vnbstmt{\begin{array}{@{}l@{}} \forall k \in \mathbb{N}, a \in \mathbb{R}. \\ \quad (0 \leq a) \Rightarrow  \\ \quad \quad (\operatorname{magnitude}((\operatorname{power}(-1, k) \cdot a)) = a) \end{array}}}
\vnbentry{\texttt{cc-\allowbreak{}cos-\allowbreak{}radius}}{\vnbstmt{\begin{array}{@{}l@{}} \forall f \in (\mathbb{N} \to \mathbb{C}). \\ \quad \forall k \in \mathbb{N}. \\ \quad \quad f(k) \\ \quad \quad =\; (\operatorname{power}(-1, k) \cdot \operatorname{recip}(\operatorname{factorial}((2 \cdot k)))) \Rightarrow \\ \quad \quad (\operatorname{cps\text{-}radius}(f) = \operatorname{pos\text{-}inf}) \end{array}}}
\vnbentry{\texttt{cc-\allowbreak{}sin-\allowbreak{}radius}}{\vnbstmt{\begin{array}{@{}l@{}} \forall f \in (\mathbb{N} \to \mathbb{C}). \\ \quad \forall k \in \mathbb{N}. \\ \quad \quad f(k) \\ \quad \quad =\; (\operatorname{power}(-1, k) \cdot \operatorname{recip}(\operatorname{factorial}(((2 \cdot k) + 1)))) \Rightarrow \\ \quad \quad (\operatorname{cps\text{-}radius}(f) = \operatorname{pos\text{-}inf}) \end{array}}}
\vnbentry{\texttt{cc-\allowbreak{}cos-\allowbreak{}as-\allowbreak{}power-\allowbreak{}series}}{\vnbstmt{\begin{array}{@{}l@{}} \forall f \in (\mathbb{N} \to \mathbb{C}), z \in \mathbb{C}. \\ \quad \forall k \in \mathbb{N}. \\ \quad \quad f(k) \\ \quad \quad =\; (\operatorname{power}(-1, k) \cdot \operatorname{recip}(\operatorname{factorial}((2 \cdot k)))) \Rightarrow \\ \quad \quad \operatorname{cc\text{-}series\text{-}converges}((\lambda k \in \mathbb{N}.\; \\ \quad \quad \quad \quad ((\operatorname{power}(-1, k) \cdot \operatorname{recip}(\operatorname{factorial}((2 \cdot k)))) \cdot \operatorname{power}(z, (2 \cdot k))))) \wedge \\ \quad \quad \cos z \\ \quad \quad =\; \operatorname{cc\text{-}series\text{-}limit}((\lambda \mathit{cpsk} \in \mathbb{N}.\; \\ \quad \quad \quad \quad (f(\mathit{cpsk}) \cdot \operatorname{power}(\operatorname{power}(z, 2), \mathit{cpsk})))) \end{array}}}
\vnbentry{\texttt{cc-\allowbreak{}sin-\allowbreak{}as-\allowbreak{}power-\allowbreak{}series}}{\vnbstmt{\begin{array}{@{}l@{}} \forall f \in (\mathbb{N} \to \mathbb{C}), z \in \mathbb{C}. \\ \quad \forall k \in \mathbb{N}. \\ \quad \quad f(k) \\ \quad \quad =\; (\operatorname{power}(-1, k) \cdot \operatorname{recip}(\operatorname{factorial}(((2 \cdot k) + 1)))) \Rightarrow \\ \quad \quad \operatorname{cc\text{-}series\text{-}converges}((\lambda k \in \mathbb{N}.\; \\ \quad \quad \quad \quad ((\operatorname{power}(-1, k) \cdot \operatorname{recip}(\operatorname{factorial}(((2 \cdot k) + 1)))) \cdot \operatorname{power}(z, ((2 \cdot k) + 1))))) \wedge \\ \quad \quad \sin z \\ \quad \quad =\; (z \cdot \operatorname{cc\text{-}series\text{-}limit}((\lambda \mathit{cpsk} \in \mathbb{N}.\; (f(\mathit{cpsk}) \cdot \operatorname{power}(\operatorname{power}(z, 2), \mathit{cpsk}))))) \end{array}}}
\vnbentry{\texttt{cc-\allowbreak{}cos-\allowbreak{}converges}}{\vnbstmt{\begin{array}{@{}l@{}} \forall z \in \mathbb{C}. \\ \quad \operatorname{cc\text{-}series\text{-}converges}((\lambda k \in \mathbb{N}.\; \\ \quad \quad \quad ((\operatorname{power}(-1, k) \cdot \operatorname{recip}(\operatorname{factorial}((2 \cdot k)))) \cdot \operatorname{power}(z, (2 \cdot k))))) \wedge \\ \quad (\cos z \in \mathbb{C}) \end{array}}}
\vnbentry{\texttt{cc-\allowbreak{}sin-\allowbreak{}converges}}{\vnbstmt{\begin{array}{@{}l@{}} \forall z \in \mathbb{C}. \\ \quad \operatorname{cc\text{-}series\text{-}converges}((\lambda k \in \mathbb{N}.\; \\ \quad \quad \quad ((\operatorname{power}(-1, k) \cdot \operatorname{recip}(\operatorname{factorial}(((2 \cdot k) + 1)))) \cdot \operatorname{power}(z, ((2 \cdot k) + 1))))) \wedge \\ \quad (\sin z \in \mathbb{C}) \end{array}}}
\vnbentry{\texttt{cc-\allowbreak{}series-\allowbreak{}even-\allowbreak{}odd}}{\vnbstmt{\begin{array}{@{}l@{}} \forall f, e, o \in (\mathbb{N} \to \mathbb{C}). \\ \quad \forall j \in \mathbb{N}.\; (e(j) = f((2 \cdot j))) \wedge \\ \quad \forall j \in \mathbb{N}.\; (o(j) = f(((2 \cdot j) + 1))) \wedge \\ \quad \operatorname{cc\text{-}series\text{-}converges}(f) \wedge \\ \quad \operatorname{cc\text{-}series\text{-}converges}(e) \wedge \\ \quad \operatorname{cc\text{-}series\text{-}converges}(o) \Rightarrow \\ \quad \quad \operatorname{cc\text{-}series\text{-}limit}(f) \\ \quad \quad =\; (\operatorname{cc\text{-}series\text{-}limit}(e) + \operatorname{cc\text{-}series\text{-}limit}(o)) \end{array}}}
\vnbentry{\texttt{cc-\allowbreak{}exp-\allowbreak{}i}}{\vnbstmt{\forall z \in \mathbb{C}.\; (\exp (+i \cdot z) = (\cos z + (+i \cdot \sin z)))}}
\vnbentry{\texttt{cc-\allowbreak{}cos-\allowbreak{}even}}{\vnbstmt{\forall z \in \mathbb{C}.\; (\cos -z = \cos z)}}
\vnbentry{\texttt{cc-\allowbreak{}sin-\allowbreak{}odd}}{\vnbstmt{\forall z \in \mathbb{C}.\; (\sin -z = -\sin z)}}
\vnbentry{\texttt{cc-\allowbreak{}sin-\allowbreak{}exp}}{\vnbstmt{\begin{array}{@{}l@{}} \forall z \in \mathbb{C}. \\ \quad (\sin z = ((\exp (+i \cdot z) - \exp -(+i \cdot z)) \cdot \operatorname{recip}((2 \cdot +i)))) \end{array}}}
\vnbentry{\texttt{cc-\allowbreak{}cos-\allowbreak{}exp}}{\vnbstmt{\begin{array}{@{}l@{}} \forall z \in \mathbb{C}. \\ \quad (\cos z = ((\exp (+i \cdot z) + \exp -(+i \cdot z)) \cdot \operatorname{recip}(2))) \end{array}}}
\vnbentry{\texttt{cex-\allowbreak{}lin-\allowbreak{}deriv}}{\vnbstmt{\begin{array}{@{}l@{}} \forall c, z \in \mathbb{C}. \\ \quad \operatorname{is\text{-}diff\text{-}on}(\operatorname{cc\text{-}normed\text{-}field}, \mathbb{C}, (\lambda y \in \mathbb{C}.\; (c \cdot y)), z, c) \end{array}}}
\vnbentry{\texttt{cc-\allowbreak{}exp-\allowbreak{}scaled-\allowbreak{}deriv}}{\vnbstmt{\begin{array}{@{}l@{}} \forall c, z \in \mathbb{C}. \\ \quad \operatorname{is\text{-}diff\text{-}on}(\operatorname{cc\text{-}normed\text{-}field}, \mathbb{C}, (\lambda y \in \mathbb{C}.\; \exp (c \cdot y)), z, \\ \quad \quad (c \cdot \exp (c \cdot z))) \end{array}}}
\vnbentry{\texttt{cc-\allowbreak{}diff-\allowbreak{}on-\allowbreak{}scale}}{\vnbstmt{\begin{array}{@{}l@{}} \forall u, f, a, l.; \forall c \in \mathbb{C}. \\ \quad \operatorname{is\text{-}diff\text{-}on}(\operatorname{cc\text{-}normed\text{-}field}, u, f, a, l) \Rightarrow  \\ \quad \quad \operatorname{is\text{-}diff\text{-}on}(\operatorname{cc\text{-}normed\text{-}field}, u, (\lambda y \in u.\; (c \cdot f(y))), a, (c \cdot l)) \end{array}}}
\vnbentry{\texttt{cex-\allowbreak{}exp-\allowbreak{}comb-\allowbreak{}deriv}}{\vnbstmt{\begin{array}{@{}l@{}} \forall a, b, c, d, z \in \mathbb{C}. \\ \quad \operatorname{is\text{-}diff\text{-}on}(\operatorname{cc\text{-}normed\text{-}field}, \mathbb{C}, \\ \quad \quad (\lambda y \in \mathbb{C}.\; ((a \cdot \exp (c \cdot y)) + (b \cdot \exp (d \cdot y)))), \\ \quad \quad z, ((a \cdot (c \cdot \exp (c \cdot z))) + (b \cdot (d \cdot \exp (d \cdot z))))) \end{array}}}
\vnbentry{\texttt{cc-\allowbreak{}cos-\allowbreak{}deriv}}{\vnbstmt{\begin{array}{@{}l@{}} \forall z \in \mathbb{C}. \\ \quad \operatorname{is\text{-}diff\text{-}on}(\operatorname{cc\text{-}normed\text{-}field}, \mathbb{C}, (\lambda y \in \mathbb{C}.\; \cos y), z, -\sin z) \end{array}}}
\vnbentry{\texttt{cc-\allowbreak{}sin-\allowbreak{}deriv}}{\vnbstmt{\begin{array}{@{}l@{}} \forall z \in \mathbb{C}. \\ \quad \operatorname{is\text{-}diff\text{-}on}(\operatorname{cc\text{-}normed\text{-}field}, \mathbb{C}, (\lambda y \in \mathbb{C}.\; \sin y), z, \cos z) \end{array}}}
\vnbentry{\texttt{cc-\allowbreak{}cos-\allowbreak{}holomorphic}}{\vnbstmt{\operatorname{holomorphic\text{-}on}(\mathbb{C}, (\lambda y \in \mathbb{C}.\; \cos y))}}
\vnbentry{\texttt{cc-\allowbreak{}sin-\allowbreak{}holomorphic}}{\vnbstmt{\operatorname{holomorphic\text{-}on}(\mathbb{C}, (\lambda y \in \mathbb{C}.\; \sin y))}}
\vnbentry{\texttt{cc-\allowbreak{}cos-\allowbreak{}add}}{\vnbstmt{\begin{array}{@{}l@{}} \forall z, w \in \mathbb{C}. \\ \quad (\cos (z + w) = ((\cos z \cdot \cos w) - (\sin z \cdot \sin w))) \end{array}}}
\vnbentry{\texttt{cc-\allowbreak{}sin-\allowbreak{}add}}{\vnbstmt{\begin{array}{@{}l@{}} \forall z, w \in \mathbb{C}. \\ \quad (\sin (z + w) = ((\cos w \cdot \sin z) + (\cos z \cdot \sin w))) \end{array}}}
\vnbentry{\texttt{cc-\allowbreak{}cos-\allowbreak{}sin-\allowbreak{}one}}{\vnbstmt{\begin{array}{@{}l@{}} \forall z \in \mathbb{C}. \\ \quad ((\operatorname{power}(\cos z, 2) + \operatorname{power}(\sin z, 2)) = 1) \end{array}}}

\vnbfile{theorem-library/cc-exp-kernel.scm}{theorem-library/cc-exp-kernel.scm}
\noindent{\small cc-exp-kernel.scm -- THE KERNEL OF exp AND THE CONSTANT pi (the notes 2.3.3: Lemma 2.18, Prop 2.19, and what they rest on). Batch 31-B, 2026-09-25.  exp, cos, sin are theorem-library/cc-exp.scm's (31-A, definitions structure-library/cc-elementary.scm); PI is defined in}\par
\vnbentry{\texttt{cek-\allowbreak{}limit-\allowbreak{}tail-\allowbreak{}le}}{\vnbstmt{\begin{array}{@{}l@{}} \forall s \in (\mathbb{N} \to \mathbb{C}), v, c \in \mathbb{C}, d \in \mathbb{R}, m \in \mathbb{N}. \\ \quad \operatorname{converges\text{-}to}(\operatorname{cc\text{-}ms}, s, v) \wedge \\ \quad \forall n \in \mathbb{N}. \\ \quad \quad ((m \leq n) \Rightarrow (\operatorname{magnitude}((s(n) - c)) \leq d)) \Rightarrow \\ \quad \quad (\operatorname{magnitude}((v - c)) \leq d) \end{array}}}
\vnbentry{\texttt{cek-\allowbreak{}half-\allowbreak{}ratio-\allowbreak{}partial}}{\vnbstmt{\begin{array}{@{}l@{}} \forall n \in \mathbb{N}, g \in (\mathbb{N} \to \mathbb{R}), m \in \mathbb{N}. \\ \quad \forall j \in \mathbb{N}.\; ((m \leq j) \Rightarrow ((2 \cdot g((j + 1))) \leq g(j))) \Rightarrow  \\ \quad \quad ((\operatorname{series\text{-}partial\text{-}sum}(g, (m + n)) - \operatorname{series\text{-}partial\text{-}sum}(g, m)) + (2 \cdot g((m + n)))) \\ \quad \quad \leq\; (2 \cdot g(m)) \end{array}}}
\vnbentry{\texttt{cc-\allowbreak{}series-\allowbreak{}ratio-\allowbreak{}tail}}{\vnbstmt{\begin{array}{@{}l@{}} \forall f \in (\mathbb{N} \to \mathbb{C}), m \in \mathbb{N}. \\ \quad \operatorname{cc\text{-}series\text{-}converges}(f) \wedge \\ \quad \forall j \in \mathbb{N}. \\ \quad \quad (m \leq j) \Rightarrow  \\ \quad \quad \quad ((2 \cdot \operatorname{magnitude}(f((j + 1)))) \leq \operatorname{magnitude}(f(j))) \Rightarrow \\ \quad \quad \operatorname{magnitude}((\operatorname{cc\text{-}series\text{-}limit}(f) - \operatorname{cc\text{-}series\text{-}partial\text{-}sum}(f, m))) \\ \quad \quad \leq\; (2 \cdot \operatorname{magnitude}(f(m))) \end{array}}}
\vnbentry{\texttt{cc-\allowbreak{}exp-\allowbreak{}tail}}{\vnbstmt{\begin{array}{@{}l@{}} \forall z \in \mathbb{C}, m \in \mathbb{N}. \\ \quad ((2 \cdot \operatorname{magnitude}(z)) \leq (m + 1)) \Rightarrow  \\ \quad \quad \operatorname{magnitude}((\exp z - \operatorname{cc\text{-}series\text{-}partial\text{-}sum}((\lambda k \in \mathbb{N}.\; (\operatorname{recip}(\operatorname{factorial}(k)) \cdot \operatorname{power}(z, k))), m))) \\ \quad \quad \leq\; (2 \cdot (\operatorname{recip}(\operatorname{factorial}(m)) \cdot \operatorname{power}(\operatorname{magnitude}(z), m))) \end{array}}}
\vnbentry{\texttt{cc-\allowbreak{}exp-\allowbreak{}near-\allowbreak{}one}}{\vnbstmt{\begin{array}{@{}l@{}} \forall z \in \mathbb{C}. \\ \quad (\operatorname{magnitude}(z) \leq 1) \Rightarrow  \\ \quad \quad \operatorname{magnitude}(((\exp z - 1) - z)) \\ \quad \quad \leq\; \operatorname{power}(\operatorname{magnitude}(z), 2) \end{array}}}
\vnbentry{\texttt{cc-\allowbreak{}exp-\allowbreak{}kernel-\allowbreak{}discrete}}{\vnbstmt{\begin{array}{@{}l@{}} \forall z \in \mathbb{C}. \\ \quad (((\exp z = 1) \wedge (\operatorname{magnitude}(z) < 1/2)) \Rightarrow (z = 0)) \end{array}}}
\vnbentry{\texttt{cek-\allowbreak{}exp-\allowbreak{}gt-\allowbreak{}one-\allowbreak{}small}}{\vnbstmt{\begin{array}{@{}l@{}} \forall t \in \mathbb{R}. \\ \quad (((0 < t) \wedge (t \leq 1/2)) \Rightarrow (1 < \exp t)) \end{array}}}
\vnbentry{\texttt{cek-\allowbreak{}exp-\allowbreak{}gt-\allowbreak{}one-\allowbreak{}n}}{\vnbstmt{\begin{array}{@{}l@{}} \forall n \in \mathbb{N}, t \in \mathbb{R}. \\ \quad (((0 < t) \wedge (t \leq (n \cdot 1/2))) \Rightarrow (1 < \exp t)) \end{array}}}
\vnbentry{\texttt{cc-\allowbreak{}exp-\allowbreak{}real-\allowbreak{}gt-\allowbreak{}one}}{\vnbstmt{\forall t \in \mathbb{R}.\; ((0 < t) \Rightarrow (1 < \exp t))}}
\vnbentry{\texttt{cc-\allowbreak{}exp-\allowbreak{}real-\allowbreak{}one-\allowbreak{}iff}}{\vnbstmt{\forall t \in \mathbb{R}.\; ((\exp t = 1) \Leftrightarrow (t = 0))}}
\vnbentry{\texttt{cc-\allowbreak{}exp-\allowbreak{}kernel-\allowbreak{}add}}{\vnbstmt{\begin{array}{@{}l@{}} \forall z, w \in \mathbb{C}. \\ \quad (((\exp z = 1) \wedge (\exp w = 1)) \Rightarrow (\exp (z + w) = 1)) \end{array}}}
\vnbentry{\texttt{cc-\allowbreak{}exp-\allowbreak{}kernel-\allowbreak{}neg}}{\vnbstmt{\forall z \in \mathbb{C}.\; ((\exp z = 1) \Rightarrow (\exp -z = 1))}}
\vnbentry{\texttt{cc-\allowbreak{}exp-\allowbreak{}kernel-\allowbreak{}sub}}{\vnbstmt{\begin{array}{@{}l@{}} \forall z, w \in \mathbb{C}. \\ \quad (((\exp z = 1) \wedge (\exp w = 1)) \Rightarrow (\exp (z - w) = 1)) \end{array}}}
\vnbentry{\texttt{cc-\allowbreak{}exp-\allowbreak{}kernel-\allowbreak{}imaginary}}{\vnbstmt{\begin{array}{@{}l@{}} \forall z \in \mathbb{C}. \\ \quad ((\exp z = 1) \Rightarrow (\operatorname{real\text{-}part}(z) = 0)) \end{array}}}
\vnbentry{\texttt{cc-\allowbreak{}exp-\allowbreak{}kernel-\allowbreak{}on-\allowbreak{}imaginary-\allowbreak{}axis}}{\vnbstmt{\begin{array}{@{}l@{}} \forall z \in \mathbb{C}. \\ \quad ((\exp z = 1) \Rightarrow (z = (+i \cdot \operatorname{imag\text{-}part}(z)))) \end{array}}}
\vnbentry{\texttt{cc-\allowbreak{}cos-\allowbreak{}real-\allowbreak{}part}}{\vnbstmt{\begin{array}{@{}l@{}} \forall t \in \mathbb{R}. \\ \quad (\cos t = \operatorname{real\text{-}part}(\exp (+i \cdot t))) \end{array}}}
\vnbentry{\texttt{cc-\allowbreak{}sin-\allowbreak{}real-\allowbreak{}part}}{\vnbstmt{\begin{array}{@{}l@{}} \forall t \in \mathbb{R}. \\ \quad (\sin t = \operatorname{imag\text{-}part}(\exp (+i \cdot t))) \end{array}}}
\vnbentry{\texttt{cc-\allowbreak{}cos-\allowbreak{}sin-\allowbreak{}real}}{\vnbstmt{\begin{array}{@{}l@{}} \forall t \in \mathbb{R}. \\ \quad ((\cos t \in \mathbb{R}) \wedge (\sin t \in \mathbb{R})) \end{array}}}
\vnbentry{\texttt{cc-\allowbreak{}cos-\allowbreak{}one-\allowbreak{}bounds}}{\vnbstmt{((5/12 \leq \cos 1) \wedge (\cos 1 \leq 7/12))}}
\vnbentry{\texttt{cc-\allowbreak{}cos-\allowbreak{}two-\allowbreak{}neg}}{\vnbstmt{(\cos 2 < 0)}}
\vnbentry{\texttt{cc-\allowbreak{}continuous-\allowbreak{}restrict-\allowbreak{}rr}}{\vnbstmt{\begin{array}{@{}l@{}} \forall f, g, x. \\ \quad \operatorname{is\text{-}continuous\text{-}at}(\operatorname{subspace\text{-}ms}(\operatorname{nf\text{-}metric\text{-}space}(\operatorname{cc\text{-}normed\text{-}field}), \mathbb{C}), \\ \quad \quad \operatorname{nf\text{-}metric\text{-}space}(\operatorname{cc\text{-}normed\text{-}field}), f, x) \wedge \\ \quad (g \in (\mathbb{R} \to \mathbb{R})) \wedge \\ \quad (x \in \mathbb{R}) \wedge \\ \quad \forall v \in \mathbb{R}.\; (g(v) = f(v)) \Rightarrow \\ \quad \quad \operatorname{is\text{-}continuous\text{-}at}(\operatorname{rr\text{-}ms}, \operatorname{rr\text{-}ms}, g, x) \end{array}}}
\vnbentry{\texttt{cc-\allowbreak{}neg-\allowbreak{}cos-\allowbreak{}continuous-\allowbreak{}rr}}{\vnbstmt{\begin{array}{@{}l@{}} ((\lambda v \in \mathbb{R}.\; (-1 \cdot \cos v)) \in (\mathbb{R} \to \mathbb{R})) \wedge \\ \forall x \in \mathbb{R}. \\ \quad \operatorname{is\text{-}continuous\text{-}at}(\operatorname{rr\text{-}ms}, \operatorname{rr\text{-}ms}, (\lambda v \in \mathbb{R}.\; (-1 \cdot \cos v)), x) \end{array}}}
\vnbentry{\texttt{cc-\allowbreak{}cos-\allowbreak{}real-\allowbreak{}continuous}}{\vnbstmt{\begin{array}{@{}l@{}} ((\lambda v \in \mathbb{R}.\; \cos v) \in (\mathbb{R} \to \mathbb{R})) \wedge \\ \forall x \in \mathbb{R}. \\ \quad \operatorname{is\text{-}continuous\text{-}at}(\operatorname{rr\text{-}ms}, \operatorname{rr\text{-}ms}, (\lambda v \in \mathbb{R}.\; \cos v), x) \end{array}}}
\vnbentry{\texttt{cc-\allowbreak{}cos-\allowbreak{}zero}}{\vnbstmt{(\cos 0 = 1)}}
\vnbentry{\texttt{cc-\allowbreak{}cos-\allowbreak{}has-\allowbreak{}zero}}{\vnbstmt{\exists t \in \mathbb{R}.\; ((0 < t) \wedge ((t < 2) \wedge (\cos t = 0)))}}
\vnbentry{\texttt{cc-\allowbreak{}exp-\allowbreak{}imaginary-\allowbreak{}period-\allowbreak{}exists}}{\vnbstmt{\begin{array}{@{}l@{}} \exists t \in \mathbb{R}. \\ \quad ((0 < t) \wedge ((t < 8) \wedge (\exp (+i \cdot t) = 1))) \end{array}}}
\vnbentry{\texttt{cc-\allowbreak{}exp-\allowbreak{}kernel-\allowbreak{}nontrivial}}{\vnbstmt{\exists z \in \mathbb{C}.\; ((\exp z = 1) \wedge \neg (z = 0))}}
\vnbentry{\texttt{cek-\allowbreak{}period-\allowbreak{}ge-\allowbreak{}half}}{\vnbstmt{\begin{array}{@{}l@{}} \forall t \in \mathbb{R}. \\ \quad (((0 < t) \wedge (\exp (+i \cdot t) = 1)) \Rightarrow (1/2 \leq t)) \end{array}}}
\vnbentry{\texttt{cc-\allowbreak{}exp-\allowbreak{}least-\allowbreak{}period}}{\vnbstmt{\begin{array}{@{}l@{}} \exists a \in \mathbb{R}. \\ \quad (0 < a) \wedge \\ \quad (\exp (+i \cdot a) = 1) \wedge \\ \quad \forall t \in \mathbb{R}. \\ \quad \quad (((0 < t) \wedge (\exp (+i \cdot t) = 1)) \Rightarrow (a \leq t)) \end{array}}}
\vnbentry{\texttt{cek-\allowbreak{}kernel-\allowbreak{}nn-\allowbreak{}multiple}}{\vnbstmt{\begin{array}{@{}l@{}} \forall n \in \mathbb{N}, w \in \mathbb{C}. \\ \quad ((\exp w = 1) \Rightarrow (\exp (n \cdot w) = 1)) \end{array}}}
\vnbentry{\texttt{cc-\allowbreak{}exp-\allowbreak{}kernel-\allowbreak{}zz-\allowbreak{}multiple}}{\vnbstmt{\begin{array}{@{}l@{}} \forall w \in \mathbb{C}, k \in \mathbb{Z}. \\ \quad ((\exp w = 1) \Rightarrow (\exp (k \cdot w) = 1)) \end{array}}}
\vnbentry{\texttt{cek-\allowbreak{}period-\allowbreak{}multiple-\allowbreak{}nn}}{\vnbstmt{\begin{array}{@{}l@{}} \forall n \in \mathbb{N}, a, r \in \mathbb{R}. \\ \quad (0 < a) \wedge \\ \quad (\exp (+i \cdot a) = 1) \wedge \\ \quad \forall t \in \mathbb{R}. \\ \quad \quad (((0 < t) \wedge (\exp (+i \cdot t) = 1)) \Rightarrow (a \leq t)) \wedge \\ \quad (0 \leq r) \wedge \\ \quad (r < (n \cdot a)) \wedge \\ \quad (\exp (+i \cdot r) = 1) \Rightarrow \\ \quad \quad \exists k \in \mathbb{N}.\; (r = (k \cdot a)) \end{array}}}
\vnbentry{\texttt{cek-\allowbreak{}period-\allowbreak{}multiple-\allowbreak{}zz}}{\vnbstmt{\begin{array}{@{}l@{}} \forall a, r \in \mathbb{R}. \\ \quad (0 < a) \wedge \\ \quad (\exp (+i \cdot a) = 1) \wedge \\ \quad \forall t \in \mathbb{R}. \\ \quad \quad (((0 < t) \wedge (\exp (+i \cdot t) = 1)) \Rightarrow (a \leq t)) \wedge \\ \quad (\exp (+i \cdot r) = 1) \Rightarrow \\ \quad \quad \exists k \in \mathbb{Z}.\; (r = (k \cdot a)) \end{array}}}
\vnbentry{\texttt{cc-\allowbreak{}pi-\allowbreak{}exists}}{\vnbstmt{\begin{array}{@{}l@{}} \exists p \in \mathbb{R}. \\ \quad (0 < p) \wedge \\ \quad \forall z \in \mathbb{C}. \\ \quad \quad ((\exp z = 1) \Leftrightarrow \exists k \in \mathbb{Z}.\; (z = (k \cdot ((2 \cdot p) \cdot +i)))) \end{array}}}
\vnbentry{\texttt{cek-\allowbreak{}pi-\allowbreak{}le}}{\vnbstmt{\begin{array}{@{}l@{}} \forall p, q. \\ \quad (p \in \mathbb{R}) \wedge \\ \quad (0 < p) \wedge \\ \quad \forall z \in \mathbb{C}. \\ \quad \quad ((\exp z = 1) \Leftrightarrow \exists k \in \mathbb{Z}.\; (z = (k \cdot ((2 \cdot p) \cdot +i)))) \wedge \\ \quad (q \in \mathbb{R}) \wedge \\ \quad (0 < q) \wedge \\ \quad \forall z \in \mathbb{C}. \\ \quad \quad ((\exp z = 1) \Leftrightarrow \exists k \in \mathbb{Z}.\; (z = (k \cdot ((2 \cdot q) \cdot +i)))) \Rightarrow \\ \quad \quad (q \leq p) \end{array}}}
\vnbentry{\texttt{cc-\allowbreak{}pi-\allowbreak{}unique}}{\vnbstmt{\begin{array}{@{}l@{}} \forall p, q. \\ \quad (p \in \mathbb{R}) \wedge \\ \quad (0 < p) \wedge \\ \quad \forall z \in \mathbb{C}. \\ \quad \quad ((\exp z = 1) \Leftrightarrow \exists k \in \mathbb{Z}.\; (z = (k \cdot ((2 \cdot p) \cdot +i)))) \wedge \\ \quad (q \in \mathbb{R}) \wedge \\ \quad (0 < q) \wedge \\ \quad \forall z \in \mathbb{C}. \\ \quad \quad ((\exp z = 1) \Leftrightarrow \exists k \in \mathbb{Z}.\; (z = (k \cdot ((2 \cdot q) \cdot +i)))) \Rightarrow \\ \quad \quad (p = q) \end{array}}}
\vnbentry{\texttt{pi-\allowbreak{}spec}}{\vnbstmt{\begin{array}{@{}l@{}} (\pi \in \mathbb{R}) \wedge \\ (0 < \pi) \wedge \\ \forall z \in \mathbb{C}. \\ \quad ((\exp z = 1) \Leftrightarrow \exists k \in \mathbb{Z}.\; (z = (k \cdot ((2 \cdot \pi) \cdot +i)))) \end{array}}}
\vnbentry{\texttt{pi-\allowbreak{}pos}}{\vnbstmt{((\pi \in \mathbb{R}) \wedge (0 < \pi))}}
\vnbentry{\texttt{cc-\allowbreak{}exp-\allowbreak{}one-\allowbreak{}iff}}{\vnbstmt{\begin{array}{@{}l@{}} \forall z \in \mathbb{C}. \\ \quad ((\exp z = 1) \Leftrightarrow \exists k \in \mathbb{Z}.\; (z = (k \cdot ((2 \cdot \pi) \cdot +i)))) \end{array}}}
\vnbentry{\texttt{cc-\allowbreak{}exp-\allowbreak{}two-\allowbreak{}pi-\allowbreak{}i}}{\vnbstmt{(\exp ((2 \cdot \pi) \cdot +i) = 1)}}
\vnbentry{\texttt{cc-\allowbreak{}cos-\allowbreak{}sin-\allowbreak{}two-\allowbreak{}pi}}{\vnbstmt{((\cos (2 \cdot \pi) = 1) \wedge (\sin (2 \cdot \pi) = 0))}}
\vnbentry{\texttt{cc-\allowbreak{}exp-\allowbreak{}i-\allowbreak{}pi}}{\vnbstmt{(\exp (+i \cdot \pi) = -1)}}
\vnbentry{\texttt{cc-\allowbreak{}cos-\allowbreak{}sin-\allowbreak{}pi}}{\vnbstmt{((\cos \pi = -1) \wedge (\sin \pi = 0))}}
\vnbentry{\texttt{cc-\allowbreak{}cos-\allowbreak{}sin-\allowbreak{}half-\allowbreak{}pi}}{\vnbstmt{((\cos (1/2 \cdot \pi) = 0) \wedge (\operatorname{power}(\sin (1/2 \cdot \pi), 2) = 1))}}
\vnbentry{\texttt{cc-\allowbreak{}exp-\allowbreak{}injective-\allowbreak{}strip}}{\vnbstmt{\begin{array}{@{}l@{}} \forall z, w \in \mathbb{C}. \\ \quad (-\pi \leq \operatorname{imag\text{-}part}(z)) \wedge \\ \quad (\operatorname{imag\text{-}part}(z) < \pi) \wedge \\ \quad (-\pi \leq \operatorname{imag\text{-}part}(w)) \wedge \\ \quad (\operatorname{imag\text{-}part}(w) < \pi) \wedge \\ \quad (\exp z = \exp w) \Rightarrow \\ \quad \quad (z = w) \end{array}}}
\vnbentry{\texttt{cc-\allowbreak{}exp-\allowbreak{}kernel-\allowbreak{}generator}}{\vnbstmt{\begin{array}{@{}l@{}} \exists g \in \mathbb{C}. \\ \quad (0 < \operatorname{imag\text{-}part}(g)) \wedge \\ \quad \forall z \in \mathbb{C}. \\ \quad \quad ((\exp z = 1) \Leftrightarrow \exists k \in \mathbb{Z}.\; (z = (k \cdot g))) \end{array}}}
\vnbentry{\texttt{cc-\allowbreak{}exp-\allowbreak{}kernel-\allowbreak{}generator-\allowbreak{}unique}}{\vnbstmt{\begin{array}{@{}l@{}} \forall g \in \mathbb{C}. \\ \quad (0 < \operatorname{imag\text{-}part}(g)) \wedge \\ \quad \forall z \in \mathbb{C}. \\ \quad \quad ((\exp z = 1) \Leftrightarrow \exists k \in \mathbb{Z}.\; (z = (k \cdot g))) \Rightarrow \\ \quad \quad (g = ((2 \cdot \pi) \cdot +i)) \end{array}}}
\vnbentry{\texttt{cc-\allowbreak{}pi-\allowbreak{}lt-\allowbreak{}four}}{\vnbstmt{(\pi < 4)}}
\vnbentry{\texttt{cc-\allowbreak{}sin-\allowbreak{}half-\allowbreak{}pi}}{\vnbstmt{(\sin (1/2 \cdot \pi) = 1)}}
\vnbentry{\texttt{cc-\allowbreak{}cos-\allowbreak{}two-\allowbreak{}pi}}{\vnbstmt{(\cos (2 \cdot \pi) = 1)}}
\vnbentry{\texttt{cc-\allowbreak{}sin-\allowbreak{}two-\allowbreak{}pi}}{\vnbstmt{(\sin (2 \cdot \pi) = 0)}}
\vnbentry{\texttt{cc-\allowbreak{}cos-\allowbreak{}pi}}{\vnbstmt{(\cos \pi = -1)}}
\vnbentry{\texttt{cc-\allowbreak{}sin-\allowbreak{}pi}}{\vnbstmt{(\sin \pi = 0)}}
\vnbentry{\texttt{cc-\allowbreak{}cos-\allowbreak{}half-\allowbreak{}pi}}{\vnbstmt{(\cos (1/2 \cdot \pi) = 0)}}

\vnbfile{theorem-library/cc-exp-surjective.scm}{theorem-library/cc-exp-surjective.scm}
\noindent{\small cc-exp-surjective.scm -- PROPOSITION 2.19 COMPLETED (exp is a bijection of each strip onto its image) AND THE BRIDGE cc-exp|RR = r-exp. Batch 34, 2026-09-26.  exp, cos, sin: theorem-library/cc-exp.scm (31-A); PI and the injectivity half of 2.19: theorem-library/cc-exp-kernel.scm (31-B).}\par
\vnbentry{\texttt{cc-\allowbreak{}exp-\allowbreak{}real-\allowbreak{}continuous-\allowbreak{}rr}}{\vnbstmt{\begin{array}{@{}l@{}} ((\lambda v \in \mathbb{R}.\; \exp v) \in (\mathbb{R} \to \mathbb{R})) \wedge \\ \forall x \in \mathbb{R}. \\ \quad \operatorname{is\text{-}continuous\text{-}at}(\operatorname{rr\text{-}ms}, \operatorname{rr\text{-}ms}, (\lambda v \in \mathbb{R}.\; \exp v), x) \end{array}}}
\vnbentry{\texttt{cc-\allowbreak{}diff-\allowbreak{}on-\allowbreak{}restrict-\allowbreak{}rr}}{\vnbstmt{\begin{array}{@{}l@{}} \forall f, g, a, l. \\ \quad \operatorname{is\text{-}diff\text{-}on}(\operatorname{cc\text{-}normed\text{-}field}, \mathbb{C}, f, a, l) \wedge \\ \quad (g \in (\mathbb{R} \to \mathbb{R})) \wedge \\ \quad (a \in \mathbb{R}) \wedge \\ \quad (l \in \mathbb{R}) \wedge \\ \quad \forall v \in \mathbb{R}.\; (g(v) = f(v)) \Rightarrow \\ \quad \quad \operatorname{is\text{-}diff\text{-}at}(g, a, l) \end{array}}}
\vnbentry{\texttt{cc-\allowbreak{}exp-\allowbreak{}real-\allowbreak{}deriv}}{\vnbstmt{\begin{array}{@{}l@{}} \forall t \in \mathbb{R}. \\ \quad \operatorname{is\text{-}diff\text{-}at}((\lambda v \in \mathbb{R}.\; \exp v), t, \exp t) \end{array}}}
\vnbentry{\texttt{cxs-\allowbreak{}log-\allowbreak{}exp-\allowbreak{}deriv}}{\vnbstmt{\begin{array}{@{}l@{}} \forall x \in \mathbb{R}. \\ \quad \operatorname{is\text{-}diff\text{-}at}((\lambda h \in \mathbb{R}.\; (\log\left(\exp h\right) - h)), x, 0) \end{array}}}
\vnbentry{\texttt{cc-\allowbreak{}exp-\allowbreak{}real-\allowbreak{}is-\allowbreak{}r-\allowbreak{}exp}}{\vnbstmt{\forall t \in \mathbb{R}.\; (\exp t = \exp\left(t\right))}}
\vnbentry{\texttt{cc-\allowbreak{}exp-\allowbreak{}real-\allowbreak{}surjective-\allowbreak{}pos}}{\vnbstmt{\begin{array}{@{}l@{}} \forall r \in \mathbb{R}. \\ \quad ((0 < r) \Rightarrow \exists s \in \mathbb{R}.\; (\exp s = r)) \end{array}}}
\vnbentry{\texttt{cc-\allowbreak{}circle-\allowbreak{}arg}}{\vnbstmt{\begin{array}{@{}l@{}} \forall a, b \in \mathbb{R}. \\ \quad ((\operatorname{power}(a, 2) + \operatorname{power}(b, 2)) = 1) \Rightarrow  \\ \quad \quad \exists h \in \mathbb{R}. \\ \quad \quad \quad ((-\pi \leq h) \wedge ((h < \pi) \wedge ((\cos h = a) \wedge (\sin h = b)))) \end{array}}}
\vnbentry{\texttt{cc-\allowbreak{}polar-\allowbreak{}form}}{\vnbstmt{\begin{array}{@{}l@{}} \forall w \in \mathbb{C}. \\ \quad \neg (w = 0) \Rightarrow  \\ \quad \quad \exists r, h \in \mathbb{R}. \\ \quad \quad \quad ((0 < r) \wedge ((-\pi \leq h) \wedge ((h < \pi) \wedge (w = (r \cdot \exp (+i \cdot h)))))) \end{array}}}
\vnbentry{\texttt{cc-\allowbreak{}exp-\allowbreak{}surjective-\allowbreak{}strip}}{\vnbstmt{\begin{array}{@{}l@{}} \forall w \in \mathbb{C}. \\ \quad \neg (w = 0) \Rightarrow  \\ \quad \quad \exists z \in \mathbb{C}. \\ \quad \quad \quad (-\pi \leq \operatorname{imag\text{-}part}(z)) \wedge \\ \quad \quad \quad (\operatorname{imag\text{-}part}(z) < \pi) \wedge \\ \quad \quad \quad (\exp z = w) \end{array}}}
\vnbentry{\texttt{cc-\allowbreak{}exp-\allowbreak{}bijection-\allowbreak{}strip}}{\vnbstmt{\begin{array}{@{}l@{}} (\lambda z \in \{q \in \mathbb{C} : ((-\pi \leq \operatorname{imag\text{-}part}(q)) \wedge (\operatorname{imag\text{-}part}(q) < \pi))\}.\; \\ \quad \exp z) \\ \in\; \mathrm{Bij}(\{q \in \mathbb{C} : ((-\pi \leq \operatorname{imag\text{-}part}(q)) \wedge (\operatorname{imag\text{-}part}(q) < \pi))\}, \operatorname{difference}(\mathbb{C}, \operatorname{singleton}(0))) \end{array}}}
\vnbentry{\texttt{cc-\allowbreak{}exp-\allowbreak{}open-\allowbreak{}strip-\allowbreak{}avoids-\allowbreak{}negatives}}{\vnbstmt{\begin{array}{@{}l@{}} \forall z \in \mathbb{C}. \\ \quad ((-\pi < \operatorname{imag\text{-}part}(z)) \wedge (\operatorname{imag\text{-}part}(z) < \pi)) \Rightarrow  \\ \quad \quad \neg ((\operatorname{imag\text{-}part}(\exp z) = 0) \wedge (\operatorname{real\text{-}part}(\exp z) \leq 0)) \end{array}}}
\vnbentry{\texttt{cc-\allowbreak{}exp-\allowbreak{}open-\allowbreak{}strip-\allowbreak{}onto}}{\vnbstmt{\begin{array}{@{}l@{}} \forall w \in \mathbb{C}. \\ \quad \neg ((\operatorname{imag\text{-}part}(w) = 0) \wedge (\operatorname{real\text{-}part}(w) \leq 0)) \Rightarrow  \\ \quad \quad \exists z \in \mathbb{C}. \\ \quad \quad \quad (-\pi < \operatorname{imag\text{-}part}(z)) \wedge \\ \quad \quad \quad (\operatorname{imag\text{-}part}(z) < \pi) \wedge \\ \quad \quad \quad (\exp z = w) \end{array}}}
\vnbentry{\texttt{cc-\allowbreak{}exp-\allowbreak{}bijection-\allowbreak{}open-\allowbreak{}strip}}{\vnbstmt{\begin{array}{@{}l@{}} (\lambda z \in \{q \in \mathbb{C} : ((-\pi < \operatorname{imag\text{-}part}(q)) \wedge (\operatorname{imag\text{-}part}(q) < \pi))\}.\; \\ \quad \exp z) \\ \in\; \mathrm{Bij}(\{q \in \mathbb{C} : ((-\pi < \operatorname{imag\text{-}part}(q)) \wedge (\operatorname{imag\text{-}part}(q) < \pi))\}, \{u \in \mathbb{C} : \neg ((\operatorname{imag\text{-}part}(u) = 0) \wedge (\operatorname{real\text{-}part}(u) \leq 0))\}) \end{array}}}

\vnbfile{theorem-library/analytic-log-laws.scm}{theorem-library/analytic-log-laws.scm}
\noindent{\small analytic-log-laws.scm -- ANALYTICITY AT A CENTRE (Def 2.11, Cor 2.14 at a), THE PRINCIPAL LOGARITHM, LEMMA 2.20 AND PROPOSITION 2.21. Batch 38, 2026-09-27.  Definitions: structure-library/analytic-log.scm (IS-ANALYTIC-AT, CC-LOG).  The source: scratchpad/triage/ca.txt 1142-1150}\par
\vnbentry{\texttt{alg-\allowbreak{}shift-\allowbreak{}in-\allowbreak{}ball}}{\vnbstmt{\begin{array}{@{}l@{}} \forall a \in \mathbb{C}, r \in \mathbb{R}, z \in \operatorname{ball}(\operatorname{cc\text{-}ms}, a, r). \\ \quad ((z \in \mathbb{C}) \wedge ((z - a) \in \operatorname{ball}(\operatorname{cc\text{-}ms}, 0, r))) \end{array}}}
\vnbentry{\texttt{alg-\allowbreak{}shift-\allowbreak{}diff}}{\vnbstmt{\begin{array}{@{}l@{}} \forall a \in \mathbb{C}, r \in \mathbb{R}, w \in \operatorname{ball}(\operatorname{cc\text{-}ms}, a, r). \\ \quad (0 < r) \Rightarrow  \\ \quad \quad \operatorname{is\text{-}diff\text{-}on}(\operatorname{cc\text{-}normed\text{-}field}, \operatorname{ball}(\operatorname{cc\text{-}ms}, a, r), \\ \quad \quad \quad (\lambda x \in \operatorname{ball}(\operatorname{cc\text{-}ms}, a, r).\; (x - a)), w, 1) \end{array}}}
\vnbentry{\texttt{cps-\allowbreak{}translate-\allowbreak{}diff}}{\vnbstmt{\begin{array}{@{}l@{}} \forall c, d \in (\mathbb{N} \to \mathbb{C}), a \in \mathbb{C}, r \in \mathbb{R}, w \in \operatorname{ball}(\operatorname{cc\text{-}ms}, a, r). \\ \quad \forall j \in \mathbb{N}.\; (d(j) = ((j + 1) \cdot c((j + 1)))) \wedge \\ \quad (0 < r) \wedge \\ \quad (r < \operatorname{cps\text{-}radius}(c)) \Rightarrow \\ \quad \quad \operatorname{is\text{-}diff\text{-}on}(\operatorname{cc\text{-}normed\text{-}field}, \operatorname{ball}(\operatorname{cc\text{-}ms}, a, r), \\ \quad \quad \quad (\lambda z \in \operatorname{ball}(\operatorname{cc\text{-}ms}, a, r).\; \\ \quad \quad \quad \quad \operatorname{cc\text{-}series\text{-}limit}((\lambda k \in \mathbb{N}.\; (c(k) \cdot \operatorname{power}((z - a), k))))), \\ \quad \quad \quad w, \\ \quad \quad \quad \operatorname{cc\text{-}series\text{-}limit}((\lambda k \in \mathbb{N}.\; (d(k) \cdot \operatorname{power}((w - a), k))))) \end{array}}}
\vnbentry{\texttt{cps-\allowbreak{}translate-\allowbreak{}holomorphic}}{\vnbstmt{\begin{array}{@{}l@{}} \forall c \in (\mathbb{N} \to \mathbb{C}), a \in \mathbb{C}, r \in \mathbb{R}. \\ \quad ((0 < r) \wedge (r < \operatorname{cps\text{-}radius}(c))) \Rightarrow  \\ \quad \quad \operatorname{holomorphic\text{-}on}(\operatorname{ball}(\operatorname{cc\text{-}ms}, a, r), \\ \quad \quad \quad (\lambda z \in \operatorname{ball}(\operatorname{cc\text{-}ms}, a, r).\; \\ \quad \quad \quad \quad \operatorname{cc\text{-}series\text{-}limit}((\lambda k \in \mathbb{N}.\; (c(k) \cdot \operatorname{power}((z - a), k)))))) \end{array}}}
\vnbentry{\texttt{cps-\allowbreak{}translate-\allowbreak{}nth-\allowbreak{}deriv}}{\vnbstmt{\begin{array}{@{}l@{}} \forall c \in (\mathbb{N} \to \mathbb{C}), a \in \mathbb{C}, r \in \mathbb{R}, k \in \mathbb{N}, g \in (\mathbb{N} \to \mathbb{C}). \\ \quad (0 < r) \wedge \\ \quad (r < \operatorname{cps\text{-}radius}(c)) \wedge \\ \quad \forall n \in \mathbb{N}. \\ \quad \quad g(n) \\ \quad \quad =\; ((\operatorname{factorial}((n + k)) \cdot \operatorname{recip}(\operatorname{factorial}(n))) \cdot c((n + k))) \Rightarrow \\ \quad \quad (\lambda z \in \operatorname{ball}(\operatorname{cc\text{-}ms}, a, r).\; \operatorname{cc\text{-}series\text{-}limit}((\lambda \mathit{cpsk} \in \mathbb{N}.\; (c(\mathit{cpsk}) \cdot \operatorname{power}((z - a), \mathit{cpsk})))))^{(k)} \\ \quad \quad =\; (\lambda z \in \operatorname{ball}(\operatorname{cc\text{-}ms}, a, r).\; \\ \quad \quad \quad \operatorname{cc\text{-}series\text{-}limit}((\lambda \mathit{cpsk} \in \mathbb{N}.\; \\ \quad \quad \quad \quad \quad (g(\mathit{cpsk}) \cdot \operatorname{power}((z - a), \mathit{cpsk}))))) \end{array}}}
\vnbentry{\texttt{taylor-\allowbreak{}coefficients-\allowbreak{}at}}{\vnbstmt{\begin{array}{@{}l@{}} \forall f, u.; \forall c \in (\mathbb{N} \to \mathbb{C}), a \in \mathbb{C}, r \in \mathbb{R}, k \in \mathbb{N}. \\ \quad (f \in (u \to \mathbb{C})) \wedge \\ \quad (0 < r) \wedge \\ \quad (r < \operatorname{cps\text{-}radius}(c)) \wedge \\ \quad (\operatorname{ball}(\operatorname{cc\text{-}ms}, a, r) \subseteq u) \wedge \\ \quad \forall z \in \operatorname{ball}(\operatorname{cc\text{-}ms}, a, r). \\ \quad \quad f(z) \\ \quad \quad =\; \operatorname{cc\text{-}series\text{-}limit}((\lambda \mathit{cpsk} \in \mathbb{N}.\; \\ \quad \quad \quad \quad (c(\mathit{cpsk}) \cdot \operatorname{power}((z - a), \mathit{cpsk})))) \Rightarrow \\ \quad \quad \operatorname{restrict}(f, \operatorname{ball}(\operatorname{cc\text{-}ms}, a, r))^{(k)}(a) \\ \quad \quad =\; (\operatorname{factorial}(k) \cdot c(k)) \end{array}}}
\vnbentry{\texttt{analytic-\allowbreak{}rep-\allowbreak{}holomorphic}}{\vnbstmt{\begin{array}{@{}l@{}} \forall f, u.; \forall c \in (\mathbb{N} \to \mathbb{C}), a \in \mathbb{C}, r \in \mathbb{R}. \\ \quad (f \in (u \to \mathbb{C})) \wedge \\ \quad (0 < r) \wedge \\ \quad (r < \operatorname{cps\text{-}radius}(c)) \wedge \\ \quad (\operatorname{ball}(\operatorname{cc\text{-}ms}, a, r) \subseteq u) \wedge \\ \quad \forall z \in \operatorname{ball}(\operatorname{cc\text{-}ms}, a, r). \\ \quad \quad f(z) \\ \quad \quad =\; \operatorname{cc\text{-}series\text{-}limit}((\lambda k \in \mathbb{N}.\; (c(k) \cdot \operatorname{power}((z - a), k)))) \Rightarrow \\ \quad \quad \operatorname{holomorphic\text{-}on}(\operatorname{ball}(\operatorname{cc\text{-}ms}, a, r), \operatorname{restrict}(f, \operatorname{ball}(\operatorname{cc\text{-}ms}, a, r))) \end{array}}}
\vnbentry{\texttt{analytic-\allowbreak{}deriv}}{\vnbstmt{\begin{array}{@{}l@{}} \forall f, u.; \forall c \in (\mathbb{N} \to \mathbb{C}), a \in \mathbb{C}, r \in \mathbb{R}. \\ \quad (f \in (u \to \mathbb{C})) \wedge \\ \quad (0 < r) \wedge \\ \quad (r < \operatorname{cps\text{-}radius}(c)) \wedge \\ \quad (\operatorname{ball}(\operatorname{cc\text{-}ms}, a, r) \subseteq u) \wedge \\ \quad \forall z \in \operatorname{ball}(\operatorname{cc\text{-}ms}, a, r). \\ \quad \quad f(z) \\ \quad \quad =\; \operatorname{cc\text{-}series\text{-}limit}((\lambda k \in \mathbb{N}.\; (c(k) \cdot \operatorname{power}((z - a), k)))) \wedge \\ \quad \operatorname{is\text{-}open}(\operatorname{cc\text{-}ms}, u) \Rightarrow \\ \quad \quad \operatorname{is\text{-}diff\text{-}on}(\operatorname{cc\text{-}normed\text{-}field}, u, f, a, c(1)) \end{array}}}
\vnbentry{\texttt{analytic-\allowbreak{}at-\allowbreak{}intro}}{\vnbstmt{\begin{array}{@{}l@{}} \forall f, u.; \forall c \in (\mathbb{N} \to \mathbb{C}), a \in \mathbb{C}, r \in \mathbb{R}. \\ \quad (f \in (u \to \mathbb{C})) \wedge \\ \quad (0 < r) \wedge \\ \quad (r < \operatorname{cps\text{-}radius}(c)) \wedge \\ \quad (\operatorname{ball}(\operatorname{cc\text{-}ms}, a, r) \subseteq u) \wedge \\ \quad \forall z \in \operatorname{ball}(\operatorname{cc\text{-}ms}, a, r). \\ \quad \quad f(z) \\ \quad \quad =\; \operatorname{cc\text{-}series\text{-}limit}((\lambda k \in \mathbb{N}.\; (c(k) \cdot \operatorname{power}((z - a), k)))) \wedge \\ \quad \operatorname{is\text{-}open}(\operatorname{cc\text{-}ms}, u) \Rightarrow \\ \quad \quad f \text{ is analytic at } a \end{array}}}
\vnbentry{\texttt{analytic-\allowbreak{}holomorphic-\allowbreak{}ball}}{\vnbstmt{\begin{array}{@{}l@{}} \forall f, u, a. \\ \quad f \text{ is analytic at } a \Rightarrow  \\ \quad \quad \exists r \in \mathbb{R}. \\ \quad \quad \quad (0 < r) \wedge \\ \quad \quad \quad (\operatorname{ball}(\operatorname{cc\text{-}ms}, a, r) \subseteq u) \wedge \\ \quad \quad \quad \operatorname{holomorphic\text{-}on}(\operatorname{ball}(\operatorname{cc\text{-}ms}, a, r), \operatorname{restrict}(f, \operatorname{ball}(\operatorname{cc\text{-}ms}, a, r))) \end{array}}}
\vnbentry{\texttt{cc-\allowbreak{}log-\allowbreak{}unique}}{\vnbstmt{\begin{array}{@{}l@{}} \forall a, b \in \mathbb{C}. \\ \quad (-\pi < \operatorname{imag\text{-}part}(a)) \wedge \\ \quad (\operatorname{imag\text{-}part}(a) < \pi) \wedge \\ \quad (-\pi < \operatorname{imag\text{-}part}(b)) \wedge \\ \quad (\operatorname{imag\text{-}part}(b) < \pi) \wedge \\ \quad (\exp a = \exp b) \Rightarrow \\ \quad \quad (a = b) \end{array}}}
\vnbentry{\texttt{cc-\allowbreak{}log-\allowbreak{}spec}}{\vnbstmt{\begin{array}{@{}l@{}} \forall w \in \mathbb{C}. \\ \quad \neg ((\operatorname{imag\text{-}part}(w) = 0) \wedge (\operatorname{real\text{-}part}(w) \leq 0)) \Rightarrow  \\ \quad \quad (\log\left(w\right) \in \mathbb{C}) \wedge \\ \quad \quad (-\pi < \operatorname{imag\text{-}part}(\log\left(w\right))) \wedge \\ \quad \quad (\operatorname{imag\text{-}part}(\log\left(w\right)) < \pi) \wedge \\ \quad \quad (\exp \log\left(w\right) = w) \end{array}}}
\vnbentry{\texttt{cc-\allowbreak{}log-\allowbreak{}in-\allowbreak{}slit}}{\vnbstmt{\begin{array}{@{}l@{}} \forall w \in \mathbb{C}. \\ \quad \neg ((\operatorname{imag\text{-}part}(w) = 0) \wedge (\operatorname{real\text{-}part}(w) \leq 0)) \Rightarrow  \\ \quad \quad (\log\left(w\right) \in \mathbb{C}) \wedge \\ \quad \quad (-\pi < \operatorname{imag\text{-}part}(\log\left(w\right))) \wedge \\ \quad \quad (\operatorname{imag\text{-}part}(\log\left(w\right)) < \pi) \end{array}}}
\vnbentry{\texttt{cc-\allowbreak{}exp-\allowbreak{}log}}{\vnbstmt{\begin{array}{@{}l@{}} \forall w \in \mathbb{C}. \\ \quad \neg ((\operatorname{imag\text{-}part}(w) = 0) \wedge (\operatorname{real\text{-}part}(w) \leq 0)) \Rightarrow  \\ \quad \quad (\exp \log\left(w\right) = w) \end{array}}}
\vnbentry{\texttt{cc-\allowbreak{}log-\allowbreak{}char}}{\vnbstmt{\begin{array}{@{}l@{}} \forall w, z \in \mathbb{C}. \\ \quad (-\pi < \operatorname{imag\text{-}part}(z)) \wedge \\ \quad (\operatorname{imag\text{-}part}(z) < \pi) \wedge \\ \quad (\exp z = w) \Rightarrow \\ \quad \quad (\log\left(w\right) = z) \end{array}}}
\vnbentry{\texttt{cc-\allowbreak{}log-\allowbreak{}exp}}{\vnbstmt{\begin{array}{@{}l@{}} \forall z \in \mathbb{C}. \\ \quad ((-\pi < \operatorname{imag\text{-}part}(z)) \wedge (\operatorname{imag\text{-}part}(z) < \pi)) \Rightarrow  \\ \quad \quad (\log\left(\exp z\right) = z) \end{array}}}
\vnbentry{\texttt{cc-\allowbreak{}log-\allowbreak{}real}}{\vnbstmt{\begin{array}{@{}l@{}} \forall t \in \mathbb{R}. \\ \quad ((0 < t) \Rightarrow (\log\left(t\right) = \log\left(t\right))) \end{array}}}
\vnbentry{\texttt{cc-\allowbreak{}log-\allowbreak{}one}}{\vnbstmt{(\log\left(1\right) = 0)}}
\vnbentry{\texttt{cc-\allowbreak{}log-\allowbreak{}polar}}{\vnbstmt{\begin{array}{@{}l@{}} \forall r, h \in \mathbb{R}. \\ \quad ((0 < r) \wedge ((-\pi < h) \wedge (h < \pi))) \Rightarrow  \\ \quad \quad (\log\left((r \cdot \exp (+i \cdot h))\right) = (\log\left(r\right) + (+i \cdot h))) \end{array}}}
\vnbentry{\texttt{cc-\allowbreak{}log-\allowbreak{}scale}}{\vnbstmt{\begin{array}{@{}l@{}} \forall s \in \mathbb{R}, w \in \mathbb{C}. \\ \quad (0 < s) \wedge \\ \quad \neg ((\operatorname{imag\text{-}part}(w) = 0) \wedge (\operatorname{real\text{-}part}(w) \leq 0)) \Rightarrow \\ \quad \quad (\log\left((s \cdot w)\right) = (\log\left(s\right) + \log\left(w\right))) \end{array}}}
\vnbentry{\texttt{cc-\allowbreak{}geometric-\allowbreak{}alternating}}{\vnbstmt{\begin{array}{@{}l@{}} \forall g \in (\mathbb{N} \to \mathbb{C}), z \in \mathbb{C}. \\ \quad \forall j \in \mathbb{N}.\; (g(j) = \operatorname{power}(-1, j)) \wedge \\ \quad (\operatorname{magnitude}(z) < 1) \Rightarrow \\ \quad \quad \operatorname{cc\text{-}series\text{-}converges}((\lambda k \in \mathbb{N}.\; (g(k) \cdot \operatorname{power}(z, k)))) \wedge \\ \quad \quad \operatorname{cc\text{-}series\text{-}limit}((\lambda k \in \mathbb{N}.\; (g(k) \cdot \operatorname{power}(z, k)))) \\ \quad \quad =\; \operatorname{recip}((1 + z)) \end{array}}}
\vnbentry{\texttt{cps-\allowbreak{}log-\allowbreak{}radius}}{\vnbstmt{\begin{array}{@{}l@{}} \forall c \in (\mathbb{N} \to \mathbb{C}), r \in \mathbb{R}. \\ \quad (c(0) = 0) \wedge \\ \quad \forall j \in \mathbb{N}. \\ \quad \quad (c((j + 1)) = (\operatorname{power}(-1, j) \cdot \operatorname{recip}((j + 1)))) \wedge \\ \quad (0 \leq r) \wedge \\ \quad (r < 1) \Rightarrow \\ \quad \quad (r < \operatorname{cps\text{-}radius}(c)) \end{array}}}
\vnbentry{\texttt{cps-\allowbreak{}log-\allowbreak{}series}}{\vnbstmt{\begin{array}{@{}l@{}} \forall c \in (\mathbb{N} \to \mathbb{C}), r \in \mathbb{R}, z \in \operatorname{ball}(\operatorname{cc\text{-}ms}, 0, r). \\ \quad (c(0) = 0) \wedge \\ \quad \forall j \in \mathbb{N}. \\ \quad \quad (c((j + 1)) = (\operatorname{power}(-1, j) \cdot \operatorname{recip}((j + 1)))) \wedge \\ \quad (0 < r) \wedge \\ \quad (r < 1) \Rightarrow \\ \quad \quad \operatorname{is\text{-}diff\text{-}on}(\operatorname{cc\text{-}normed\text{-}field}, \operatorname{ball}(\operatorname{cc\text{-}ms}, 0, r), \\ \quad \quad \quad (\lambda \mathit{cpdz} \in \operatorname{ball}(\operatorname{cc\text{-}ms}, 0, r).\; \\ \quad \quad \quad \quad \operatorname{cc\text{-}series\text{-}limit}((\lambda k \in \mathbb{N}.\; (c(k) \cdot \operatorname{power}(\mathit{cpdz}, k))))), \\ \quad \quad \quad z, \operatorname{recip}((1 + z))) \end{array}}}
\vnbentry{\texttt{alg-\allowbreak{}zero-\allowbreak{}deriv-\allowbreak{}ball}}{\vnbstmt{\begin{array}{@{}l@{}} \forall a \in \mathbb{C}, r \in \mathbb{R}, h \in (\operatorname{ball}(\operatorname{cc\text{-}ms}, a, r) \to \mathbb{C}), z \in \operatorname{ball}(\operatorname{cc\text{-}ms}, a, r). \\ \quad (0 < r) \wedge \\ \quad \forall \mathit{hbz} \in \operatorname{ball}(\operatorname{cc\text{-}ms}, a, r). \\ \quad \quad \operatorname{is\text{-}diff\text{-}on}(\operatorname{cc\text{-}normed\text{-}field}, \operatorname{ball}(\operatorname{cc\text{-}ms}, a, r), h, \mathit{hbz}, 0) \Rightarrow \\ \quad \quad (h(z) = h(a)) \end{array}}}
\vnbentry{\texttt{cps-\allowbreak{}log-\allowbreak{}deriv-\allowbreak{}at-\allowbreak{}one}}{\vnbstmt{\begin{array}{@{}l@{}} \forall c \in (\mathbb{N} \to \mathbb{C}), r \in \mathbb{R}, w \in \operatorname{ball}(\operatorname{cc\text{-}ms}, 1, r). \\ \quad (c(0) = 0) \wedge \\ \quad \forall j \in \mathbb{N}. \\ \quad \quad (c((j + 1)) = (\operatorname{power}(-1, j) \cdot \operatorname{recip}((j + 1)))) \wedge \\ \quad (0 < r) \wedge \\ \quad (r < 1) \Rightarrow \\ \quad \quad \operatorname{is\text{-}diff\text{-}on}(\operatorname{cc\text{-}normed\text{-}field}, \operatorname{ball}(\operatorname{cc\text{-}ms}, 1, r), \\ \quad \quad \quad (\lambda z \in \operatorname{ball}(\operatorname{cc\text{-}ms}, 1, r).\; \\ \quad \quad \quad \quad \operatorname{cc\text{-}series\text{-}limit}((\lambda k \in \mathbb{N}.\; (c(k) \cdot \operatorname{power}((z - 1), k))))), \\ \quad \quad \quad w, \operatorname{recip}(w)) \end{array}}}
\vnbentry{\texttt{log0-\allowbreak{}exp}}{\vnbstmt{\begin{array}{@{}l@{}} \forall c \in (\mathbb{N} \to \mathbb{C}), r \in \mathbb{R}, z \in \operatorname{ball}(\operatorname{cc\text{-}ms}, 1, r). \\ \quad (c(0) = 0) \wedge \\ \quad \forall j \in \mathbb{N}. \\ \quad \quad (c((j + 1)) = (\operatorname{power}(-1, j) \cdot \operatorname{recip}((j + 1)))) \wedge \\ \quad (0 < r) \wedge \\ \quad (r < 1) \Rightarrow \\ \quad \quad \exp \operatorname{cc\text{-}series\text{-}limit}((\lambda k \in \mathbb{N}.\; (c(k) \cdot \operatorname{power}((z - 1), k)))) \\ \quad \quad =\; z \end{array}}}
\vnbentry{\texttt{log0-\allowbreak{}bound}}{\vnbstmt{\begin{array}{@{}l@{}} \forall c \in (\mathbb{N} \to \mathbb{C}), r \in \mathbb{R}, z \in \operatorname{ball}(\operatorname{cc\text{-}ms}, 1, r). \\ \quad (c(0) = 0) \wedge \\ \quad \forall j \in \mathbb{N}. \\ \quad \quad (c((j + 1)) = (\operatorname{power}(-1, j) \cdot \operatorname{recip}((j + 1)))) \wedge \\ \quad (0 < r) \wedge \\ \quad (r < 1) \Rightarrow \\ \quad \quad \operatorname{magnitude}(\operatorname{cc\text{-}series\text{-}limit}((\lambda k \in \mathbb{N}.\; (c(k) \cdot \operatorname{power}((z - 1), k))))) \\ \quad \quad \leq\; (\operatorname{recip}((1 - \operatorname{magnitude}((z - 1)))) - 1) \end{array}}}
\vnbentry{\texttt{log0-\allowbreak{}is-\allowbreak{}log-\allowbreak{}on-\allowbreak{}ball}}{\vnbstmt{\begin{array}{@{}l@{}} \forall c \in (\mathbb{N} \to \mathbb{C}), r \in \mathbb{R}, z \in \operatorname{ball}(\operatorname{cc\text{-}ms}, 1, r). \\ \quad (c(0) = 0) \wedge \\ \quad \forall j \in \mathbb{N}. \\ \quad \quad (c((j + 1)) = (\operatorname{power}(-1, j) \cdot \operatorname{recip}((j + 1)))) \wedge \\ \quad (0 < r) \wedge \\ \quad ((r \cdot (1 + \pi)) \leq \pi) \Rightarrow \\ \quad \quad \log\left(z\right) \\ \quad \quad =\; \operatorname{cc\text{-}series\text{-}limit}((\lambda k \in \mathbb{N}.\; (c(k) \cdot \operatorname{power}((z - 1), k)))) \end{array}}}
\vnbentry{\texttt{cc-\allowbreak{}log-\allowbreak{}holomorphic-\allowbreak{}near-\allowbreak{}one}}{\vnbstmt{\begin{array}{@{}l@{}} \forall r \in \mathbb{R}. \\ \quad ((0 < r) \wedge ((r \cdot (1 + \pi)) \leq \pi)) \Rightarrow  \\ \quad \quad \operatorname{holomorphic\text{-}on}(\operatorname{ball}(\operatorname{cc\text{-}ms}, 1, r), (\lambda z \in \operatorname{ball}(\operatorname{cc\text{-}ms}, 1, r).\; \log\left(z\right))) \end{array}}}
\vnbentry{\texttt{cc-\allowbreak{}log-\allowbreak{}one-\allowbreak{}plus}}{\vnbstmt{\begin{array}{@{}l@{}} \forall c \in (\mathbb{N} \to \mathbb{C}), r \in \mathbb{R}, z \in \mathbb{C}. \\ \quad (c(0) = 0) \wedge \\ \quad \forall j \in \mathbb{N}. \\ \quad \quad (c((j + 1)) = (\operatorname{power}(-1, j) \cdot \operatorname{recip}((j + 1)))) \wedge \\ \quad (0 < r) \wedge \\ \quad ((r \cdot (1 + \pi)) \leq \pi) \wedge \\ \quad (\operatorname{magnitude}(z) < r) \Rightarrow \\ \quad \quad \log\left((1 + z)\right) \\ \quad \quad =\; \operatorname{cc\text{-}series\text{-}limit}((\lambda k \in \mathbb{N}.\; (c(k) \cdot \operatorname{power}(z, k)))) \end{array}}}
\vnbentry{\texttt{holomorphic-\allowbreak{}compose}}{\vnbstmt{\begin{array}{@{}l@{}} \forall u, v, f, g. \\ \quad \operatorname{holomorphic\text{-}on}(u, f) \wedge \\ \quad \operatorname{holomorphic\text{-}on}(v, g) \wedge \\ \quad \forall y \in u.\; (f(y) \in v) \Rightarrow \\ \quad \quad \operatorname{holomorphic\text{-}on}(u, (\lambda x \in u.\; g(f(x)))) \end{array}}}
\vnbentry{\texttt{cc-\allowbreak{}log-\allowbreak{}holomorphic-\allowbreak{}scale}}{\vnbstmt{\begin{array}{@{}l@{}} \forall r, s \in \mathbb{R}. \\ \quad ((0 < r) \wedge (((r \cdot (1 + \pi)) \leq \pi) \wedge (0 < s))) \Rightarrow  \\ \quad \quad \operatorname{holomorphic\text{-}on}(\operatorname{ball}(\operatorname{cc\text{-}ms}, s, (s \cdot r)), \\ \quad \quad \quad (\lambda z \in \operatorname{ball}(\operatorname{cc\text{-}ms}, s, (s \cdot r)).\; \\ \quad \quad \quad \quad \log\left(z\right))) \end{array}}}
\vnbentry{\texttt{log0-\allowbreak{}real}}{\vnbstmt{\begin{array}{@{}l@{}} \forall c \in (\mathbb{N} \to \mathbb{C}), r \in \mathbb{R}, z \in \operatorname{ball}(\operatorname{cc\text{-}ms}, 1, r). \\ \quad (c(0) = 0) \wedge \\ \quad \forall j \in \mathbb{N}. \\ \quad \quad (c((j + 1)) = (\operatorname{power}(-1, j) \cdot \operatorname{recip}((j + 1)))) \wedge \\ \quad (0 < r) \wedge \\ \quad (r < 1) \wedge \\ \quad (z \in \mathbb{R}) \Rightarrow \\ \quad \quad \operatorname{cc\text{-}series\text{-}limit}((\lambda k \in \mathbb{N}.\; (c(k) \cdot \operatorname{power}((z - 1), k)))) \\ \quad \quad \in\; \mathbb{C} \wedge \\ \quad \quad \operatorname{imag\text{-}part}(\operatorname{cc\text{-}series\text{-}limit}((\lambda k \in \mathbb{N}.\; (c(k) \cdot \operatorname{power}((z - 1), k))))) \\ \quad \quad =\; 0 \end{array}}}
\vnbentry{\texttt{log0-\allowbreak{}is-\allowbreak{}ln}}{\vnbstmt{\begin{array}{@{}l@{}} \forall c \in (\mathbb{N} \to \mathbb{C}), t \in \mathbb{R}. \\ \quad (c(0) = 0) \wedge \\ \quad \forall j \in \mathbb{N}. \\ \quad \quad (c((j + 1)) = (\operatorname{power}(-1, j) \cdot \operatorname{recip}((j + 1)))) \wedge \\ \quad (0 < t) \wedge \\ \quad (t < 2) \Rightarrow \\ \quad \quad \operatorname{cc\text{-}series\text{-}limit}((\lambda k \in \mathbb{N}.\; (c(k) \cdot \operatorname{power}((t - 1), k)))) \\ \quad \quad =\; \log\left(t\right) \end{array}}}

\vnbfile{theorem-library/cc-log-holomorphic.scm}{theorem-library/cc-log-holomorphic.scm}
\noindent{\small cc-log-holomorphic.scm -- THE SIGN OF cos ON [-pi/2, pi], pi $>$ 3, AND THE HOLOMORPHY OF THE PRINCIPAL LOGARITHM ON THE SLIT PLANE C \textbackslash{} (-inf, 0]; the two loose ends of Lemma 2.20 (the radius is exactly 1) and Remark 2.22 (log(1 - z) = - sum z\textasciicircum{}k/k).}\par
\vnbentry{\texttt{cc-\allowbreak{}cos-\allowbreak{}least-\allowbreak{}zero}}{\vnbstmt{\begin{array}{@{}l@{}} \forall t \in \mathbb{R}. \\ \quad (((0 < t) \wedge (t < (1/2 \cdot \pi))) \Rightarrow \neg (\cos t = 0)) \end{array}}}
\vnbentry{\texttt{cc-\allowbreak{}cos-\allowbreak{}pos-\allowbreak{}right}}{\vnbstmt{\begin{array}{@{}l@{}} \forall s \in \mathbb{R}. \\ \quad (((0 < s) \wedge (s < (1/2 \cdot \pi))) \Rightarrow (0 < \cos s)) \end{array}}}
\vnbentry{\texttt{cc-\allowbreak{}cos-\allowbreak{}pos-\allowbreak{}half}}{\vnbstmt{\begin{array}{@{}l@{}} \forall t \in \mathbb{R}. \\ \quad (((-(1/2 \cdot \pi) < t) \wedge (t < (1/2 \cdot \pi))) \Rightarrow (0 < \cos t)) \end{array}}}
\vnbentry{\texttt{cc-\allowbreak{}cos-\allowbreak{}nonpos-\allowbreak{}far}}{\vnbstmt{\begin{array}{@{}l@{}} \forall t \in \mathbb{R}. \\ \quad ((((1/2 \cdot \pi) \leq t) \wedge (t \leq \pi)) \Rightarrow (\cos t \leq 0)) \end{array}}}
\vnbentry{\texttt{cc-\allowbreak{}cos-\allowbreak{}pos-\allowbreak{}small}}{\vnbstmt{\forall t \in \mathbb{R}.\; (((0 < t) \wedge (t < 1)) \Rightarrow (0 < \cos t))}}
\vnbentry{\texttt{cc-\allowbreak{}pi-\allowbreak{}gt-\allowbreak{}two}}{\vnbstmt{(2 < \pi)}}
\vnbentry{\texttt{cc-\allowbreak{}cos-\allowbreak{}three-\allowbreak{}halves-\allowbreak{}pos}}{\vnbstmt{(0 < \cos 3/2)}}
\vnbentry{\texttt{cc-\allowbreak{}pi-\allowbreak{}gt-\allowbreak{}three}}{\vnbstmt{(3 < \pi)}}
\vnbentry{\texttt{cc-\allowbreak{}log-\allowbreak{}right-\allowbreak{}half}}{\vnbstmt{\begin{array}{@{}l@{}} \forall w \in \mathbb{C}. \\ \quad (0 < \operatorname{real\text{-}part}(w)) \Rightarrow  \\ \quad \quad (-(1/2 \cdot \pi) < \operatorname{imag\text{-}part}(\log\left(w\right))) \wedge \\ \quad \quad (\operatorname{imag\text{-}part}(\log\left(w\right)) < (1/2 \cdot \pi)) \end{array}}}
\vnbentry{\texttt{cc-\allowbreak{}log-\allowbreak{}mul}}{\vnbstmt{\begin{array}{@{}l@{}} \forall u, v \in \mathbb{C}. \\ \quad \neg ((\operatorname{imag\text{-}part}(u) = 0) \wedge (\operatorname{real\text{-}part}(u) \leq 0)) \wedge \\ \quad \neg ((\operatorname{imag\text{-}part}(v) = 0) \wedge (\operatorname{real\text{-}part}(v) \leq 0)) \wedge \\ \quad -\pi \\ \quad <\; (\operatorname{imag\text{-}part}(\log\left(u\right)) + \operatorname{imag\text{-}part}(\log\left(v\right))) \wedge \\ \quad (\operatorname{imag\text{-}part}(\log\left(u\right)) + \operatorname{imag\text{-}part}(\log\left(v\right))) \\ \quad <\; \pi \Rightarrow \\ \quad \quad (\log\left((u \cdot v)\right) = (\log\left(u\right) + \log\left(v\right))) \end{array}}}
\vnbentry{\texttt{cc-\allowbreak{}log-\allowbreak{}rotate}}{\vnbstmt{\begin{array}{@{}l@{}} \forall w \in \mathbb{C}. \\ \quad (0 < \operatorname{real\text{-}part}(w)) \Rightarrow  \\ \quad \quad (\log\left((+i \cdot w)\right) = (\log\left(w\right) + (+i \cdot (1/2 \cdot \pi)))) \end{array}}}
\vnbentry{\texttt{cc-\allowbreak{}log-\allowbreak{}rotate-\allowbreak{}neg}}{\vnbstmt{\begin{array}{@{}l@{}} \forall w \in \mathbb{C}. \\ \quad (0 < \operatorname{real\text{-}part}(w)) \Rightarrow  \\ \quad \quad (\log\left((-+i \cdot w)\right) = (\log\left(w\right) - (+i \cdot (1/2 \cdot \pi)))) \end{array}}}
\vnbentry{\texttt{cc-\allowbreak{}log-\allowbreak{}im-\allowbreak{}near-\allowbreak{}one}}{\vnbstmt{\begin{array}{@{}l@{}} \forall d, r \in \mathbb{R}, z \in \operatorname{ball}(\operatorname{cc\text{-}ms}, 1, r). \\ \quad ((0 < d) \wedge ((0 < r) \wedge (((r \cdot (1 + \pi)) \leq \pi) \wedge ((r \cdot (1 + d)) \leq d)))) \Rightarrow  \\ \quad \quad \neg ((\operatorname{imag\text{-}part}(z) = 0) \wedge (\operatorname{real\text{-}part}(z) \leq 0)) \wedge \\ \quad \quad (-d < \operatorname{imag\text{-}part}(\log\left(z\right))) \wedge \\ \quad \quad (\operatorname{imag\text{-}part}(\log\left(z\right)) < d) \end{array}}}
\vnbentry{\texttt{cc-\allowbreak{}log-\allowbreak{}holomorphic-\allowbreak{}at}}{\vnbstmt{\begin{array}{@{}l@{}} \forall a \in \mathbb{C}. \\ \quad \neg ((\operatorname{imag\text{-}part}(a) = 0) \wedge (\operatorname{real\text{-}part}(a) \leq 0)) \Rightarrow  \\ \quad \quad \exists p \in \mathbb{R}. \\ \quad \quad \quad (0 < p) \wedge \\ \quad \quad \quad \operatorname{ball}(\operatorname{cc\text{-}ms}, a, p) \\ \quad \quad \quad \subseteq\; \{u \in \mathbb{C} : \neg ((\operatorname{imag\text{-}part}(u) = 0) \wedge (\operatorname{real\text{-}part}(u) \leq 0))\} \wedge \\ \quad \quad \quad \operatorname{holomorphic\text{-}on}(\operatorname{ball}(\operatorname{cc\text{-}ms}, a, p), (\lambda x \in \operatorname{ball}(\operatorname{cc\text{-}ms}, a, p).\; \log\left(x\right))) \end{array}}}
\vnbentry{\texttt{cc-\allowbreak{}log-\allowbreak{}holomorphic}}{\vnbstmt{\begin{array}{@{}l@{}} \operatorname{holomorphic\text{-}on}(\{u \in \mathbb{C} : \neg ((\operatorname{imag\text{-}part}(u) = 0) \wedge (\operatorname{real\text{-}part}(u) \leq 0))\}, \\ \quad (\lambda w \in \{u \in \mathbb{C} : \neg ((\operatorname{imag\text{-}part}(u) = 0) \wedge (\operatorname{real\text{-}part}(u) \leq 0))\}.\; \\ \quad \quad \log\left(w\right))) \end{array}}}
\vnbentry{\texttt{cc-\allowbreak{}log-\allowbreak{}deriv}}{\vnbstmt{\begin{array}{@{}l@{}} \forall b \in \mathbb{C}. \\ \quad \neg ((\operatorname{imag\text{-}part}(b) = 0) \wedge (\operatorname{real\text{-}part}(b) \leq 0)) \Rightarrow  \\ \quad \quad \operatorname{is\text{-}diff\text{-}on}(\operatorname{cc\text{-}normed\text{-}field}, \\ \quad \quad \quad \{u \in \mathbb{C} : \neg ((\operatorname{imag\text{-}part}(u) = 0) \wedge (\operatorname{real\text{-}part}(u) \leq 0))\}, \\ \quad \quad \quad (\lambda w \in \{u \in \mathbb{C} : \neg ((\operatorname{imag\text{-}part}(u) = 0) \wedge (\operatorname{real\text{-}part}(u) \leq 0))\}.\; \\ \quad \quad \quad \quad \log\left(w\right)), \\ \quad \quad \quad b, \operatorname{recip}(b)) \end{array}}}
\vnbentry{\texttt{cc-\allowbreak{}log-\allowbreak{}one-\allowbreak{}minus}}{\vnbstmt{\begin{array}{@{}l@{}} \forall r \in \mathbb{R}, z \in \mathbb{C}. \\ \quad ((0 < r) \wedge (((r \cdot (1 + \pi)) \leq \pi) \wedge (\operatorname{magnitude}(z) < r))) \Rightarrow  \\ \quad \quad \log\left((1 - z)\right) \\ \quad \quad =\; -\operatorname{cc\text{-}series\text{-}limit}((\lambda k \in \mathbb{N}.\; ({k}^{\ast\!-1} \cdot \operatorname{power}(z, k)))) \end{array}}}
\vnbentry{\texttt{cc-\allowbreak{}series-\allowbreak{}terms-\allowbreak{}small}}{\vnbstmt{\begin{array}{@{}l@{}} \forall f \in (\mathbb{N} \to \mathbb{C}), e. \\ \quad (\operatorname{cc\text{-}series\text{-}converges}(f) \wedge \operatorname{pos\text{-}rr}(e)) \Rightarrow  \\ \quad \quad \exists n \in \mathbb{N}. \forall m \in \mathbb{N}. \\ \quad \quad \quad ((n \leq m) \Rightarrow (\operatorname{magnitude}(f(m)) \leq e)) \end{array}}}
\vnbentry{\texttt{cps-\allowbreak{}log-\allowbreak{}radius-\allowbreak{}le-\allowbreak{}one}}{\vnbstmt{\begin{array}{@{}l@{}} \forall c \in (\mathbb{N} \to \mathbb{C}), r \in \mathbb{R}. \\ \quad (c(0) = 0) \wedge \\ \quad \forall j \in \mathbb{N}. \\ \quad \quad (c((j + 1)) = (\operatorname{power}(-1, j) \cdot \operatorname{recip}((j + 1)))) \wedge \\ \quad (1 < r) \Rightarrow \\ \quad \quad \neg (r < \operatorname{cps\text{-}radius}(c)) \end{array}}}
\vnbentry{\texttt{cc-\allowbreak{}sin-\allowbreak{}pos-\allowbreak{}half}}{\vnbstmt{\begin{array}{@{}l@{}} \forall t \in \mathbb{R}. \\ \quad (((0 < t) \wedge (t < \pi)) \Rightarrow (0 < \sin t)) \end{array}}}
\vnbentry{\texttt{cps-\allowbreak{}log-\allowbreak{}radius-\allowbreak{}one}}{\vnbstmt{\begin{array}{@{}l@{}} \forall c \in (\mathbb{N} \to \mathbb{C}). \\ \quad (c(0) = 0) \wedge \\ \quad \forall j \in \mathbb{N}. \\ \quad \quad (c((j + 1)) = (\operatorname{power}(-1, j) \cdot \operatorname{recip}((j + 1)))) \Rightarrow \\ \quad \quad (\operatorname{cps\text{-}radius}(c) = 1) \end{array}}}

\vnbfile{theorem-library/winding-number-laws.scm}{theorem-library/winding-number-laws.scm}
\noindent{\small winding-number-laws.scm -- CAUCHY ON A DISC, THE WINDING NUMBER, THE CIRCLE AS A ROAD (complex-analysis.pdf ch. 3: Corollary 3.9 for a disc inside an open set; Definition 3.26 and the disc case of Corollary 3.30; Example 3.13 at the centre). Agent CA-1 of the October roadmap (docs/roadmap-2026-10-03.md, line (a), ...}\par
\vnbentry{\texttt{cc-\allowbreak{}ball-\allowbreak{}is-\allowbreak{}convex}}{\vnbstmt{\begin{array}{@{}l@{}} \forall c \in \mathbb{C}, r \in \mathbb{R}. \\ \quad \operatorname{is\text{-}convex}(\operatorname{ball}(\operatorname{cc\text{-}ms}, c, r)) \end{array}}}
\vnbentry{\texttt{line-\allowbreak{}int-\allowbreak{}domain-\allowbreak{}congruence}}{\vnbstmt{\begin{array}{@{}l@{}} \forall u, f, m, c, a, b, v.; \forall g \in (v \to \mathbb{C}). \\ \quad (u \subseteq \mathbb{C}) \wedge \\ \quad (f \in (u \to \mathbb{C})) \wedge \\ \quad \forall y \in u. \\ \quad \quad \operatorname{is\text{-}continuous\text{-}at}(\operatorname{subspace\text{-}ms}(\operatorname{nf\text{-}metric\text{-}space}(\operatorname{cc\text{-}normed\text{-}field}), u), \\ \quad \quad \quad \operatorname{nf\text{-}metric\text{-}space}(\operatorname{cc\text{-}normed\text{-}field}), f, y) \wedge \\ \quad \operatorname{is\text{-}road}(m, c, a, b) \wedge \\ \quad (\operatorname{trace}(m, a, b) \subseteq u) \wedge \\ \quad (\operatorname{trace}(m, a, b) \subseteq v) \wedge \\ \quad \forall x \in \operatorname{trace}(m, a, b).\; (g(x) \simeq f(x)) \Rightarrow \\ \quad \quad (\operatorname{line\text{-}int}(g, m, c, a, b) = \operatorname{line\text{-}int}(f, m, c, a, b)) \end{array}}}
\vnbentry{\texttt{holomorphic-\allowbreak{}on-\allowbreak{}restrict}}{\vnbstmt{\begin{array}{@{}l@{}} \forall u, f, v. \\ \quad \operatorname{holomorphic\text{-}on}(u, f) \wedge \\ \quad \operatorname{is\text{-}open}(\operatorname{nf\text{-}metric\text{-}space}(\operatorname{cc\text{-}normed\text{-}field}), v) \wedge \\ \quad (v \subseteq u) \Rightarrow \\ \quad \quad \operatorname{holomorphic\text{-}on}(v, \operatorname{restrict}(f, v)) \end{array}}}
\vnbentry{\texttt{cauchy-\allowbreak{}convex-\allowbreak{}sub}}{\vnbstmt{\begin{array}{@{}l@{}} \forall u, f, v, m, c, a, b. \\ \quad \operatorname{holomorphic\text{-}on}(u, f) \wedge \\ \quad \operatorname{is\text{-}open}(\operatorname{nf\text{-}metric\text{-}space}(\operatorname{cc\text{-}normed\text{-}field}), v) \wedge \\ \quad \operatorname{is\text{-}convex}(v) \wedge \\ \quad (v \subseteq u) \wedge \\ \quad \operatorname{is\text{-}road}(m, c, a, b) \wedge \\ \quad (\operatorname{trace}(m, a, b) \subseteq v) \wedge \\ \quad (m(b) = m(a)) \Rightarrow \\ \quad \quad (\operatorname{line\text{-}int}(f, m, c, a, b) = 0) \end{array}}}
\vnbentry{\texttt{cauchy-\allowbreak{}disc}}{\vnbstmt{\begin{array}{@{}l@{}} \forall u, f.; \forall c \in \mathbb{C}, r \in \mathbb{R}, m, d, a, b. \\ \quad \operatorname{holomorphic\text{-}on}(u, f) \wedge \\ \quad (0 < r) \wedge \\ \quad (\operatorname{ball}(\operatorname{cc\text{-}ms}, c, r) \subseteq u) \wedge \\ \quad \operatorname{is\text{-}road}(m, d, a, b) \wedge \\ \quad (\operatorname{trace}(m, a, b) \subseteq \operatorname{ball}(\operatorname{cc\text{-}ms}, c, r)) \wedge \\ \quad (m(b) = m(a)) \Rightarrow \\ \quad \quad (\operatorname{line\text{-}int}(f, m, d, a, b) = 0) \end{array}}}
\vnbentry{\texttt{cc-\allowbreak{}punctured-\allowbreak{}membership}}{\vnbstmt{\begin{array}{@{}l@{}} \forall z \in \mathbb{C}, y. \\ \quad ((y \in \operatorname{difference}(\mathbb{C}, \operatorname{singleton}(z))) \Leftrightarrow ((y \in \mathbb{C}) \wedge \neg (y = z))) \end{array}}}
\vnbentry{\texttt{cc-\allowbreak{}punctured-\allowbreak{}open}}{\vnbstmt{\begin{array}{@{}l@{}} \forall z \in \mathbb{C}. \\ \quad \operatorname{is\text{-}open}(\operatorname{cc\text{-}ms}, \operatorname{difference}(\mathbb{C}, \operatorname{singleton}(z))) \end{array}}}
\vnbentry{\texttt{winding-\allowbreak{}integrand-\allowbreak{}holomorphic}}{\vnbstmt{\begin{array}{@{}l@{}} \forall z \in \mathbb{C}. \\ \quad \operatorname{holomorphic\text{-}on}(\operatorname{difference}(\mathbb{C}, \operatorname{singleton}(z)), \\ \quad \quad (\lambda w \in \operatorname{difference}(\mathbb{C}, \operatorname{singleton}(z)).\; \\ \quad \quad \quad \operatorname{recip}((w - z)))) \end{array}}}
\vnbentry{\texttt{winding-\allowbreak{}in-\allowbreak{}cc}}{\vnbstmt{\begin{array}{@{}l@{}} \forall m, c, a, b.; \forall z \in \mathbb{C}. \\ \quad (\operatorname{is\text{-}road}(m, c, a, b) \wedge \neg (z \in \operatorname{trace}(m, a, b))) \Rightarrow  \\ \quad \quad (\operatorname{winding}(m, c, a, b, z) \in \mathbb{C}) \end{array}}}
\vnbentry{\texttt{winding-\allowbreak{}zero-\allowbreak{}outside-\allowbreak{}disc}}{\vnbstmt{\begin{array}{@{}l@{}} \forall m, d, a, b.; \forall c \in \mathbb{C}, r \in \mathbb{R}, z \in \mathbb{C}. \\ \quad \operatorname{is\text{-}road}(m, d, a, b) \wedge \\ \quad (m(b) = m(a)) \wedge \\ \quad (0 < r) \wedge \\ \quad (\operatorname{trace}(m, a, b) \subseteq \operatorname{ball}(\operatorname{cc\text{-}ms}, c, r)) \wedge \\ \quad \neg (z \in \operatorname{ball}(\operatorname{cc\text{-}ms}, c, r)) \Rightarrow \\ \quad \quad (\operatorname{winding}(m, d, a, b, z) = 0) \end{array}}}
\vnbentry{\texttt{winding-\allowbreak{}circle-\allowbreak{}center}}{\vnbstmt{\begin{array}{@{}l@{}} \forall c \in \mathbb{C}, r \in \mathbb{R}. \\ \quad (0 < r) \Rightarrow  \\ \quad \quad \operatorname{winding}(\operatorname{circle\text{-}path}(c, r), \operatorname{circle\text{-}deriv}(c, r), 0, (2 \cdot \pi), c) \\ \quad \quad =\; 1 \end{array}}}
\vnbentry{\texttt{road-\allowbreak{}holomorphic-\allowbreak{}image}}{\vnbstmt{\begin{array}{@{}l@{}} \forall u, f, c, m, d, a, b. \\ \quad \operatorname{holomorphic\text{-}on}(u, f) \wedge \\ \quad (c \in (u \to \mathbb{C})) \wedge \\ \quad \forall z \in u. \\ \quad \quad \operatorname{is\text{-}diff\text{-}on}(\operatorname{cc\text{-}normed\text{-}field}, u, f, z, c(z)) \wedge \\ \quad \forall y \in u. \\ \quad \quad \operatorname{is\text{-}continuous\text{-}at}(\operatorname{subspace\text{-}ms}(\operatorname{nf\text{-}metric\text{-}space}(\operatorname{cc\text{-}normed\text{-}field}), u), \\ \quad \quad \quad \operatorname{nf\text{-}metric\text{-}space}(\operatorname{cc\text{-}normed\text{-}field}), c, y) \wedge \\ \quad \operatorname{is\text{-}road}(m, d, a, b) \wedge \\ \quad (\operatorname{trace}(m, a, b) \subseteq u) \Rightarrow \\ \quad \quad \operatorname{is\text{-}road}((\lambda t \in \operatorname{ccint}(a, b).\; f(m(t))), \\ \quad \quad \quad (\lambda t \in \operatorname{ccint}(a, b).\; (c(m(t)) \cdot d(t))), a, b) \end{array}}}
\vnbentry{\texttt{circle-\allowbreak{}is-\allowbreak{}road}}{\vnbstmt{\begin{array}{@{}l@{}} \forall c \in \mathbb{C}, r \in \mathbb{R}. \\ \quad \operatorname{is\text{-}road}(\operatorname{circle\text{-}path}(c, r), \operatorname{circle\text{-}deriv}(c, r), 0, (2 \cdot \pi)) \end{array}}}

\vnbfile{theorem-library/analytic-from-integral.scm}{theorem-library/analytic-from-integral.scm}
\noindent{\small analytic-from-integral.scm -- THE CAUCHY-TYPE INTEGRAL IS A POWER SERIES (complex-analysis.pdf 3.5: Lemma 3.33 with Remark 3.34, the engine of Theorem 3.35; 3.7: the Cauchy estimate behind Prop 3.47 / Cor 3.48). Agent CA-2 of the October roadmap (docs/roadmap-2026-10-03.md 0(a)), 2026-10-03.}\par
\vnbentry{\texttt{aif-\allowbreak{}cont-\allowbreak{}product}}{\vnbstmt{\begin{array}{@{}l@{}} \forall u, f, g, h.; \forall y \in u. \\ \quad (u \subseteq \mathbb{C}) \wedge \\ \quad (f \in (u \to \mathbb{C})) \wedge \\ \quad \forall y \in u. \\ \quad \quad \operatorname{is\text{-}continuous\text{-}at}(\operatorname{subspace\text{-}ms}(\operatorname{nf\text{-}metric\text{-}space}(\operatorname{cc\text{-}normed\text{-}field}), u), \\ \quad \quad \quad \operatorname{nf\text{-}metric\text{-}space}(\operatorname{cc\text{-}normed\text{-}field}), f, y) \wedge \\ \quad (g \in (u \to \mathbb{C})) \wedge \\ \quad \forall y \in u. \\ \quad \quad \operatorname{is\text{-}continuous\text{-}at}(\operatorname{subspace\text{-}ms}(\operatorname{nf\text{-}metric\text{-}space}(\operatorname{cc\text{-}normed\text{-}field}), u), \\ \quad \quad \quad \operatorname{nf\text{-}metric\text{-}space}(\operatorname{cc\text{-}normed\text{-}field}), g, y) \wedge \\ \quad (h \in (u \to \mathbb{C})) \wedge \\ \quad \forall x \in u.\; (h(x) \simeq (f(x) \cdot g(x))) \Rightarrow \\ \quad \quad \operatorname{is\text{-}continuous\text{-}at}(\operatorname{subspace\text{-}ms}(\operatorname{nf\text{-}metric\text{-}space}(\operatorname{cc\text{-}normed\text{-}field}), u), \\ \quad \quad \quad \operatorname{nf\text{-}metric\text{-}space}(\operatorname{cc\text{-}normed\text{-}field}), h, y) \end{array}}}
\vnbentry{\texttt{aif-\allowbreak{}cont-\allowbreak{}sum}}{\vnbstmt{\begin{array}{@{}l@{}} \forall u, f, g, h.; \forall y \in u. \\ \quad (u \subseteq \mathbb{C}) \wedge \\ \quad (f \in (u \to \mathbb{C})) \wedge \\ \quad \forall y \in u. \\ \quad \quad \operatorname{is\text{-}continuous\text{-}at}(\operatorname{subspace\text{-}ms}(\operatorname{nf\text{-}metric\text{-}space}(\operatorname{cc\text{-}normed\text{-}field}), u), \\ \quad \quad \quad \operatorname{nf\text{-}metric\text{-}space}(\operatorname{cc\text{-}normed\text{-}field}), f, y) \wedge \\ \quad (g \in (u \to \mathbb{C})) \wedge \\ \quad \forall y \in u. \\ \quad \quad \operatorname{is\text{-}continuous\text{-}at}(\operatorname{subspace\text{-}ms}(\operatorname{nf\text{-}metric\text{-}space}(\operatorname{cc\text{-}normed\text{-}field}), u), \\ \quad \quad \quad \operatorname{nf\text{-}metric\text{-}space}(\operatorname{cc\text{-}normed\text{-}field}), g, y) \wedge \\ \quad (h \in (u \to \mathbb{C})) \wedge \\ \quad \forall x \in u.\; (h(x) \simeq (f(x) + g(x))) \Rightarrow \\ \quad \quad \operatorname{is\text{-}continuous\text{-}at}(\operatorname{subspace\text{-}ms}(\operatorname{nf\text{-}metric\text{-}space}(\operatorname{cc\text{-}normed\text{-}field}), u), \\ \quad \quad \quad \operatorname{nf\text{-}metric\text{-}space}(\operatorname{cc\text{-}normed\text{-}field}), h, y) \end{array}}}
\vnbentry{\texttt{aif-\allowbreak{}cont-\allowbreak{}holo-\allowbreak{}restrict}}{\vnbstmt{\begin{array}{@{}l@{}} \forall v, g, u.; \forall h \in (u \to \mathbb{C}), y \in u. \\ \quad \operatorname{holomorphic\text{-}on}(v, g) \wedge \\ \quad (u \subseteq v) \wedge \\ \quad \forall x \in u.\; (h(x) \simeq g(x)) \Rightarrow \\ \quad \quad \operatorname{is\text{-}continuous\text{-}at}(\operatorname{subspace\text{-}ms}(\operatorname{nf\text{-}metric\text{-}space}(\operatorname{cc\text{-}normed\text{-}field}), u), \\ \quad \quad \quad \operatorname{nf\text{-}metric\text{-}space}(\operatorname{cc\text{-}normed\text{-}field}), h, y) \end{array}}}
\vnbentry{\texttt{aif-\allowbreak{}cont-\allowbreak{}compose}}{\vnbstmt{\begin{array}{@{}l@{}} \forall u.; \forall g \in (u \to \mathbb{C}), v, \mathit{aiph}.; \forall h \in (u \to \mathbb{C}), y \in u. \\ \quad (u \subseteq \mathbb{C}) \wedge \\ \quad \forall y \in u. \\ \quad \quad \operatorname{is\text{-}continuous\text{-}at}(\operatorname{subspace\text{-}ms}(\operatorname{nf\text{-}metric\text{-}space}(\operatorname{cc\text{-}normed\text{-}field}), u), \\ \quad \quad \quad \operatorname{nf\text{-}metric\text{-}space}(\operatorname{cc\text{-}normed\text{-}field}), g, y) \wedge \\ \quad (v \subseteq \mathbb{C}) \wedge \\ \quad \forall x \in u.\; (g(x) \in v) \wedge \\ \quad \forall y \in v. \\ \quad \quad \operatorname{is\text{-}continuous\text{-}at}(\operatorname{subspace\text{-}ms}(\operatorname{nf\text{-}metric\text{-}space}(\operatorname{cc\text{-}normed\text{-}field}), v), \\ \quad \quad \quad \operatorname{nf\text{-}metric\text{-}space}(\operatorname{cc\text{-}normed\text{-}field}), \mathit{aiph}, y) \wedge \\ \quad \forall x \in u.\; (h(x) \simeq \operatorname{aiph}(g(x))) \Rightarrow \\ \quad \quad \operatorname{is\text{-}continuous\text{-}at}(\operatorname{subspace\text{-}ms}(\operatorname{nf\text{-}metric\text{-}space}(\operatorname{cc\text{-}normed\text{-}field}), u), \\ \quad \quad \quad \operatorname{nf\text{-}metric\text{-}space}(\operatorname{cc\text{-}normed\text{-}field}), h, y) \end{array}}}
\vnbentry{\texttt{aif-\allowbreak{}cont-\allowbreak{}affine}}{\vnbstmt{\begin{array}{@{}l@{}} \forall u.; \forall p, q, w \in \mathbb{C}, h \in (u \to \mathbb{C}), y \in u. \\ \quad ((u \subseteq \mathbb{C}) \wedge \forall x \in u.\; (h(x) \simeq (p + (q \cdot (x - w))))) \Rightarrow  \\ \quad \quad \operatorname{is\text{-}continuous\text{-}at}(\operatorname{subspace\text{-}ms}(\operatorname{nf\text{-}metric\text{-}space}(\operatorname{cc\text{-}normed\text{-}field}), u), \\ \quad \quad \quad \operatorname{nf\text{-}metric\text{-}space}(\operatorname{cc\text{-}normed\text{-}field}), h, y) \end{array}}}
\vnbentry{\texttt{aif-\allowbreak{}cont-\allowbreak{}shift}}{\vnbstmt{\begin{array}{@{}l@{}} \forall u.; \forall c \in \mathbb{C}, h \in (u \to \mathbb{C}), y \in u. \\ \quad ((u \subseteq \mathbb{C}) \wedge \forall x \in u.\; (h(x) \simeq (x - c))) \Rightarrow  \\ \quad \quad \operatorname{is\text{-}continuous\text{-}at}(\operatorname{subspace\text{-}ms}(\operatorname{nf\text{-}metric\text{-}space}(\operatorname{cc\text{-}normed\text{-}field}), u), \\ \quad \quad \quad \operatorname{nf\text{-}metric\text{-}space}(\operatorname{cc\text{-}normed\text{-}field}), h, y) \end{array}}}
\vnbentry{\texttt{aif-\allowbreak{}cont-\allowbreak{}const}}{\vnbstmt{\begin{array}{@{}l@{}} \forall u.; \forall k \in \mathbb{C}, h \in (u \to \mathbb{C}), y \in u. \\ \quad ((u \subseteq \mathbb{C}) \wedge \forall x \in u.\; (h(x) \simeq k)) \Rightarrow  \\ \quad \quad \operatorname{is\text{-}continuous\text{-}at}(\operatorname{subspace\text{-}ms}(\operatorname{nf\text{-}metric\text{-}space}(\operatorname{cc\text{-}normed\text{-}field}), u), \\ \quad \quad \quad \operatorname{nf\text{-}metric\text{-}space}(\operatorname{cc\text{-}normed\text{-}field}), h, y) \end{array}}}
\vnbentry{\texttt{aif-\allowbreak{}cont-\allowbreak{}recip-\allowbreak{}shift}}{\vnbstmt{\begin{array}{@{}l@{}} \forall u.; \forall c \in \mathbb{C}, h \in (u \to \mathbb{C}), y \in u. \\ \quad (u \subseteq \mathbb{C}) \wedge \\ \quad \forall x \in u.\; \neg (x = c) \wedge \\ \quad \forall x \in u.\; (h(x) \simeq \operatorname{recip}((x - c))) \Rightarrow \\ \quad \quad \operatorname{is\text{-}continuous\text{-}at}(\operatorname{subspace\text{-}ms}(\operatorname{nf\text{-}metric\text{-}space}(\operatorname{cc\text{-}normed\text{-}field}), u), \\ \quad \quad \quad \operatorname{nf\text{-}metric\text{-}space}(\operatorname{cc\text{-}normed\text{-}field}), h, y) \end{array}}}
\vnbentry{\texttt{aif-\allowbreak{}cont-\allowbreak{}power-\allowbreak{}of}}{\vnbstmt{\begin{array}{@{}l@{}} \forall u.; \forall g \in (u \to \mathbb{C}), n \in \mathbb{N}, h \in (u \to \mathbb{C}), y \in u. \\ \quad (u \subseteq \mathbb{C}) \wedge \\ \quad \forall y \in u. \\ \quad \quad \operatorname{is\text{-}continuous\text{-}at}(\operatorname{subspace\text{-}ms}(\operatorname{nf\text{-}metric\text{-}space}(\operatorname{cc\text{-}normed\text{-}field}), u), \\ \quad \quad \quad \operatorname{nf\text{-}metric\text{-}space}(\operatorname{cc\text{-}normed\text{-}field}), g, y) \wedge \\ \quad \forall x \in u.\; (h(x) \simeq \operatorname{power}(g(x), n)) \Rightarrow \\ \quad \quad \operatorname{is\text{-}continuous\text{-}at}(\operatorname{subspace\text{-}ms}(\operatorname{nf\text{-}metric\text{-}space}(\operatorname{cc\text{-}normed\text{-}field}), u), \\ \quad \quad \quad \operatorname{nf\text{-}metric\text{-}space}(\operatorname{cc\text{-}normed\text{-}field}), h, y) \end{array}}}
\vnbentry{\texttt{aif-\allowbreak{}magnitude-\allowbreak{}recip}}{\vnbstmt{\begin{array}{@{}l@{}} \forall d \in \mathbb{C}. \\ \quad \neg (d = 0) \Rightarrow  \\ \quad \quad \operatorname{magnitude}(\operatorname{recip}(d)) \\ \quad \quad =\; \operatorname{recip}(\operatorname{magnitude}(d)) \end{array}}}
\vnbentry{\texttt{aif-\allowbreak{}geometric-\allowbreak{}partial}}{\vnbstmt{\begin{array}{@{}l@{}} \forall n \in \mathbb{N}, p, v, \mathit{aiqv}, e. \\ \quad (p \in \mathbb{C}) \wedge \\ \quad (v \in \mathbb{C}) \wedge \\ \quad (\mathit{aiqv} \in \mathbb{C}) \wedge \\ \quad (e \in \mathbb{C}) \wedge \\ \quad ((e \cdot (1 - (\mathit{aiqv} \cdot v))) = v) \Rightarrow \\ \quad \quad \operatorname{cc\text{-}series\text{-}partial\text{-}sum}((\lambda j \in \mathbb{N}.\; \\ \quad \quad \quad \quad ((p \cdot \operatorname{power}(v, (j + 1))) \cdot \operatorname{power}(\mathit{aiqv}, j))), \\ \quad \quad \quad n) \\ \quad \quad =\; ((p \cdot e) - ((p \cdot \operatorname{power}((\mathit{aiqv} \cdot v), n)) \cdot e)) \end{array}}}
\vnbentry{\texttt{aif-\allowbreak{}trace-\allowbreak{}bounded}}{\vnbstmt{\begin{array}{@{}l@{}} \forall u.; \forall i \in (u \to \mathbb{C}), m, a, b. \\ \quad (u \subseteq \mathbb{C}) \wedge \\ \quad \forall y \in u. \\ \quad \quad \operatorname{is\text{-}continuous\text{-}at}(\operatorname{subspace\text{-}ms}(\operatorname{nf\text{-}metric\text{-}space}(\operatorname{cc\text{-}normed\text{-}field}), u), \\ \quad \quad \quad \operatorname{nf\text{-}metric\text{-}space}(\operatorname{cc\text{-}normed\text{-}field}), i, y) \wedge \\ \quad \operatorname{is\text{-}path}(m, a, b) \wedge \\ \quad (\operatorname{trace}(m, a, b) \subseteq u) \Rightarrow \\ \quad \quad \exists \mathit{aim} \in \mathbb{R}. \\ \quad \quad \quad (0 \leq \mathit{aim}) \wedge \\ \quad \quad \quad \forall \mathit{aiy} \in \operatorname{trace}(m, a, b). \\ \quad \quad \quad \quad (\operatorname{magnitude}(i(\mathit{aiy})) \leq \mathit{aim}) \end{array}}}
\vnbentry{\texttt{aif-\allowbreak{}scale-\allowbreak{}small}}{\vnbstmt{\begin{array}{@{}l@{}} \forall l \in \mathbb{R}, s. \\ \quad ((0 \leq l) \wedge \operatorname{pos\text{-}rr}(s)) \Rightarrow  \\ \quad \quad \operatorname{pos\text{-}rr}((s \cdot \operatorname{recip}((l + 1)))) \wedge \\ \quad \quad ((l \cdot (s \cdot \operatorname{recip}((l + 1)))) \leq s) \end{array}}}
\vnbentry{\texttt{aif-\allowbreak{}cont-\allowbreak{}scale}}{\vnbstmt{\begin{array}{@{}l@{}} \forall u.; \forall k \in \mathbb{C}, g, h \in (u \to \mathbb{C}), y \in u. \\ \quad (u \subseteq \mathbb{C}) \wedge \\ \quad \forall y \in u. \\ \quad \quad \operatorname{is\text{-}continuous\text{-}at}(\operatorname{subspace\text{-}ms}(\operatorname{nf\text{-}metric\text{-}space}(\operatorname{cc\text{-}normed\text{-}field}), u), \\ \quad \quad \quad \operatorname{nf\text{-}metric\text{-}space}(\operatorname{cc\text{-}normed\text{-}field}), g, y) \wedge \\ \quad \forall x \in u.\; (h(x) \simeq (k \cdot g(x))) \Rightarrow \\ \quad \quad \operatorname{is\text{-}continuous\text{-}at}(\operatorname{subspace\text{-}ms}(\operatorname{nf\text{-}metric\text{-}space}(\operatorname{cc\text{-}normed\text{-}field}), u), \\ \quad \quad \quad \operatorname{nf\text{-}metric\text{-}space}(\operatorname{cc\text{-}normed\text{-}field}), h, y) \end{array}}}
\vnbentry{\texttt{line-\allowbreak{}int-\allowbreak{}of-\allowbreak{}uniform-\allowbreak{}limit}}{\vnbstmt{\begin{array}{@{}l@{}} \forall u, m, c, a, b.; \forall g \in (u \to \mathbb{C}), d. \\ \quad (u \subseteq \mathbb{C}) \wedge \\ \quad \operatorname{is\text{-}road}(m, c, a, b) \wedge \\ \quad (\operatorname{trace}(m, a, b) \subseteq u) \wedge \\ \quad \forall y \in u. \\ \quad \quad \operatorname{is\text{-}continuous\text{-}at}(\operatorname{subspace\text{-}ms}(\operatorname{nf\text{-}metric\text{-}space}(\operatorname{cc\text{-}normed\text{-}field}), u), \\ \quad \quad \quad \operatorname{nf\text{-}metric\text{-}space}(\operatorname{cc\text{-}normed\text{-}field}), g, y) \wedge \\ \quad \forall k \in \mathbb{N}.\; (d(k) \in (u \to \mathbb{C})) \wedge \\ \quad \forall k \in \mathbb{N}, y \in u. \\ \quad \quad \operatorname{is\text{-}continuous\text{-}at}(\operatorname{subspace\text{-}ms}(\operatorname{nf\text{-}metric\text{-}space}(\operatorname{cc\text{-}normed\text{-}field}), u), \\ \quad \quad \quad \operatorname{nf\text{-}metric\text{-}space}(\operatorname{cc\text{-}normed\text{-}field}), d(k), y) \wedge \\ \quad \forall s. \\ \quad \quad \operatorname{pos\text{-}rr}(s) \Rightarrow  \\ \quad \quad \quad \exists p \in \mathbb{N}. \forall k \in \mathbb{N}, \mathit{aiy} \in \operatorname{trace}(m, a, b). \\ \quad \quad \quad \quad (p \leq k) \Rightarrow  \\ \quad \quad \quad \quad \quad (\operatorname{magnitude}((d(k)(\mathit{aiy}) - g(\mathit{aiy}))) \leq s) \Rightarrow \\ \quad \quad \operatorname{converges\text{-}to}(\operatorname{cc\text{-}ms}, (\lambda k \in \mathbb{N}.\; \operatorname{line\text{-}int}(d(k), m, c, a, b)), \\ \quad \quad \quad \operatorname{line\text{-}int}(g, m, c, a, b)) \end{array}}}
\vnbentry{\texttt{aif-\allowbreak{}geometric-\allowbreak{}relation}}{\vnbstmt{\begin{array}{@{}l@{}} \forall w, c, z \in \mathbb{C}. \\ \quad (\neg (w = c) \wedge \neg (w = z)) \Rightarrow  \\ \quad \quad (\operatorname{recip}((w - z)) \cdot (1 - ((z - c) \cdot \operatorname{recip}((w - c))))) \\ \quad \quad =\; \operatorname{recip}((w - c)) \end{array}}}
\vnbentry{\texttt{aif-\allowbreak{}geometric-\allowbreak{}remainder}}{\vnbstmt{\begin{array}{@{}l@{}} \forall c \in \mathbb{C}, r \in \mathbb{R}, w, z, p \in \mathbb{C}, m \in \mathbb{R}, k \in \mathbb{N}. \\ \quad (0 < r) \wedge \\ \quad (r \leq \operatorname{magnitude}((w - c))) \wedge \\ \quad (\operatorname{magnitude}((z - c)) < r) \wedge \\ \quad (\operatorname{magnitude}(p) \leq m) \Rightarrow \\ \quad \quad \operatorname{magnitude}((\operatorname{cc\text{-}series\text{-}partial\text{-}sum}((\lambda j \in \mathbb{N}.\; ((p \cdot \operatorname{power}(\operatorname{recip}((w - c)), (j + 1))) \cdot \operatorname{power}((z - c), j))), k) - (p \cdot \operatorname{recip}((w - z))))) \\ \quad \quad \in\; \mathbb{R} \wedge \\ \quad \quad \operatorname{magnitude}((\operatorname{cc\text{-}series\text{-}partial\text{-}sum}((\lambda j \in \mathbb{N}.\; ((p \cdot \operatorname{power}(\operatorname{recip}((w - c)), (j + 1))) \cdot \operatorname{power}((z - c), j))), k) - (p \cdot \operatorname{recip}((w - z))))) \\ \quad \quad \leq\; ((m \cdot \operatorname{power}((\operatorname{magnitude}((z - c)) \cdot \operatorname{recip}(r)), k)) \cdot \operatorname{recip}((r - \operatorname{magnitude}((z - c))))) \end{array}}}
\vnbentry{\texttt{geometric-\allowbreak{}expansion-\allowbreak{}uniform}}{\vnbstmt{\begin{array}{@{}l@{}} \forall u, c, r.; \forall z \in \mathbb{C}, i.; \forall m \in \mathbb{R}, s. \\ \quad (u \subseteq \mathbb{C}) \wedge \\ \quad (c \in \mathbb{C}) \wedge \\ \quad (r \in \mathbb{R}) \wedge \\ \quad (0 < r) \wedge \\ \quad \forall x \in u.\; (r \leq \operatorname{magnitude}((x - c))) \wedge \\ \quad (\operatorname{magnitude}((z - c)) < r) \wedge \\ \quad \forall x \in u.\; (i(x) \in \mathbb{C}) \wedge \\ \quad (0 \leq m) \wedge \\ \quad \forall x \in u.\; (\operatorname{magnitude}(i(x)) \leq m) \wedge \\ \quad \operatorname{pos\text{-}rr}(s) \Rightarrow \\ \quad \quad \exists p \in \mathbb{N}. \forall k \in \mathbb{N}, y \in u. \\ \quad \quad \quad (p \leq k) \Rightarrow  \\ \quad \quad \quad \quad \operatorname{magnitude}((\operatorname{cc\text{-}series\text{-}partial\text{-}sum}((\lambda j \in \mathbb{N}.\; ((i(y) \cdot \operatorname{power}(\operatorname{recip}((y - c)), (j + 1))) \cdot \operatorname{power}((z - c), j))), k) - (i(y) \cdot \operatorname{recip}((y - z))))) \\ \quad \quad \quad \quad \leq\; s \end{array}}}
\vnbentry{\texttt{aif-\allowbreak{}off-\allowbreak{}center}}{\vnbstmt{\begin{array}{@{}l@{}} \forall c \in \mathbb{C}, r \in \mathbb{R}, w, z \in \mathbb{C}. \\ \quad (0 < r) \wedge \\ \quad (r \leq \operatorname{magnitude}((w - c))) \wedge \\ \quad (\operatorname{magnitude}((z - c)) < r) \Rightarrow \\ \quad \quad (\neg (w = c) \wedge \neg (w = z)) \end{array}}}
\vnbentry{\texttt{aif-\allowbreak{}partial-\allowbreak{}integrals}}{\vnbstmt{\begin{array}{@{}l@{}} \forall k \in \mathbb{N}, u, c, r, i, m, d, a, b, z. \\ \quad (u \subseteq \mathbb{C}) \wedge \\ \quad (c \in \mathbb{C}) \wedge \\ \quad (r \in \mathbb{R}) \wedge \\ \quad (0 < r) \wedge \\ \quad \forall x \in u.\; (r \leq \operatorname{magnitude}((x - c))) \wedge \\ \quad (i \in (u \to \mathbb{C})) \wedge \\ \quad \forall y \in u. \\ \quad \quad \operatorname{is\text{-}continuous\text{-}at}(\operatorname{subspace\text{-}ms}(\operatorname{nf\text{-}metric\text{-}space}(\operatorname{cc\text{-}normed\text{-}field}), u), \\ \quad \quad \quad \operatorname{nf\text{-}metric\text{-}space}(\operatorname{cc\text{-}normed\text{-}field}), i, y) \wedge \\ \quad \operatorname{is\text{-}road}(m, d, a, b) \wedge \\ \quad (\operatorname{trace}(m, a, b) \subseteq u) \wedge \\ \quad (z \in \mathbb{C}) \wedge \\ \quad (\operatorname{magnitude}((z - c)) < r) \Rightarrow \\ \quad \quad (\lambda w \in u.\; \\ \quad \quad \quad \operatorname{cc\text{-}series\text{-}partial\text{-}sum}((\lambda j \in \mathbb{N}.\; \\ \quad \quad \quad \quad \quad ((i(w) \cdot \operatorname{power}(\operatorname{recip}((w - c)), (j + 1))) \cdot \operatorname{power}((z - c), j))), \\ \quad \quad \quad \quad k)) \\ \quad \quad \in\; (u \to \mathbb{C}) \wedge \\ \quad \quad \forall y \in u. \\ \quad \quad \quad \operatorname{is\text{-}continuous\text{-}at}(\operatorname{subspace\text{-}ms}(\operatorname{nf\text{-}metric\text{-}space}(\operatorname{cc\text{-}normed\text{-}field}), u), \\ \quad \quad \quad \quad \operatorname{nf\text{-}metric\text{-}space}(\operatorname{cc\text{-}normed\text{-}field}), \\ \quad \quad \quad \quad (\lambda w \in u.\; \\ \quad \quad \quad \quad \quad \operatorname{cc\text{-}series\text{-}partial\text{-}sum}((\lambda j \in \mathbb{N}.\; \\ \quad \quad \quad \quad \quad \quad \quad ((i(w) \cdot \operatorname{power}(\operatorname{recip}((w - c)), (j + 1))) \cdot \operatorname{power}((z - c), j))), \\ \quad \quad \quad \quad \quad \quad k)), \\ \quad \quad \quad \quad y) \wedge \\ \quad \quad \operatorname{line\text{-}int}((\lambda w \in u.\; \\ \quad \quad \quad \quad \operatorname{cc\text{-}series\text{-}partial\text{-}sum}((\lambda j \in \mathbb{N}.\; \\ \quad \quad \quad \quad \quad \quad ((i(w) \cdot \operatorname{power}(\operatorname{recip}((w - c)), (j + 1))) \cdot \operatorname{power}((z - c), j))), \\ \quad \quad \quad \quad \quad k)), \\ \quad \quad \quad m, d, a, b) \\ \quad \quad =\; \operatorname{cc\text{-}series\text{-}partial\text{-}sum}((\lambda \mathit{aitk} \in \mathbb{N}.\; \\ \quad \quad \quad \quad (\operatorname{line\text{-}int}((\lambda w \in u.\; (i(w) \cdot \operatorname{power}(\operatorname{recip}((w - c)), (\mathit{aitk} + 1)))), m, d, a, b) \cdot \operatorname{power}((z - c), \mathit{aitk}))), \\ \quad \quad \quad k) \end{array}}}
\vnbentry{\texttt{cauchy-\allowbreak{}kernel-\allowbreak{}integral-\allowbreak{}series}}{\vnbstmt{\begin{array}{@{}l@{}} \forall u, c, r, i, m, d, a, b.; \forall z \in \mathbb{C}. \\ \quad (u \subseteq \mathbb{C}) \wedge \\ \quad (c \in \mathbb{C}) \wedge \\ \quad (r \in \mathbb{R}) \wedge \\ \quad (0 < r) \wedge \\ \quad \forall x \in u.\; (r \leq \operatorname{magnitude}((x - c))) \wedge \\ \quad (i \in (u \to \mathbb{C})) \wedge \\ \quad \forall y \in u. \\ \quad \quad \operatorname{is\text{-}continuous\text{-}at}(\operatorname{subspace\text{-}ms}(\operatorname{nf\text{-}metric\text{-}space}(\operatorname{cc\text{-}normed\text{-}field}), u), \\ \quad \quad \quad \operatorname{nf\text{-}metric\text{-}space}(\operatorname{cc\text{-}normed\text{-}field}), i, y) \wedge \\ \quad \operatorname{is\text{-}road}(m, d, a, b) \wedge \\ \quad (\operatorname{trace}(m, a, b) \subseteq u) \wedge \\ \quad (\operatorname{magnitude}((z - c)) < r) \Rightarrow \\ \quad \quad \operatorname{cc\text{-}series\text{-}converges\text{-}to}((\lambda k \in \mathbb{N}.\; \\ \quad \quad \quad \quad (\operatorname{line\text{-}int}((\lambda w \in u.\; (i(w) \cdot \operatorname{power}(\operatorname{recip}((w - c)), (k + 1)))), m, d, a, b) \cdot \operatorname{power}((z - c), k))), \\ \quad \quad \quad \operatorname{line\text{-}int}((\lambda w \in u.\; (i(w) \cdot \operatorname{recip}((w - z)))), m, d, a, b)) \end{array}}}
\vnbentry{\texttt{analytic-\allowbreak{}from-\allowbreak{}integral-\allowbreak{}formula}}{\vnbstmt{\begin{array}{@{}l@{}} \forall u, c, r, i, m, d, a, b.; \forall p \in \mathbb{C}, f.; \forall z \in \mathbb{C}. \\ \quad (u \subseteq \mathbb{C}) \wedge \\ \quad (c \in \mathbb{C}) \wedge \\ \quad (r \in \mathbb{R}) \wedge \\ \quad (0 < r) \wedge \\ \quad \forall x \in u.\; (r \leq \operatorname{magnitude}((x - c))) \wedge \\ \quad (i \in (u \to \mathbb{C})) \wedge \\ \quad \forall y \in u. \\ \quad \quad \operatorname{is\text{-}continuous\text{-}at}(\operatorname{subspace\text{-}ms}(\operatorname{nf\text{-}metric\text{-}space}(\operatorname{cc\text{-}normed\text{-}field}), u), \\ \quad \quad \quad \operatorname{nf\text{-}metric\text{-}space}(\operatorname{cc\text{-}normed\text{-}field}), i, y) \wedge \\ \quad \operatorname{is\text{-}road}(m, d, a, b) \wedge \\ \quad (\operatorname{trace}(m, a, b) \subseteq u) \wedge \\ \quad \forall z \in \mathbb{C}. \\ \quad \quad (\operatorname{magnitude}((z - c)) < r) \Rightarrow  \\ \quad \quad \quad f(z) \\ \quad \quad \quad =\; (p \cdot \operatorname{line\text{-}int}((\lambda w \in u.\; (i(w) \cdot \operatorname{recip}((w - z)))), m, d, a, b)) \wedge \\ \quad (\operatorname{magnitude}((z - c)) < r) \Rightarrow \\ \quad \quad \operatorname{cc\text{-}series\text{-}converges\text{-}to}((\lambda k \in \mathbb{N}.\; \\ \quad \quad \quad \quad ((p \cdot \operatorname{line\text{-}int}((\lambda w \in u.\; (i(w) \cdot \operatorname{power}(\operatorname{recip}((w - c)), (k + 1)))), m, d, a, b)) \cdot \operatorname{power}((z - c), k))), \\ \quad \quad \quad f(z)) \end{array}}}
\vnbentry{\texttt{aif-\allowbreak{}coefficient-\allowbreak{}bound}}{\vnbstmt{\begin{array}{@{}l@{}} \forall u, c, r, i, m, d, a, b.; \forall \mathit{aim} \in \mathbb{R}, k \in \mathbb{N}. \\ \quad (u \subseteq \mathbb{C}) \wedge \\ \quad (c \in \mathbb{C}) \wedge \\ \quad (r \in \mathbb{R}) \wedge \\ \quad (0 < r) \wedge \\ \quad \forall x \in u.\; (r \leq \operatorname{magnitude}((x - c))) \wedge \\ \quad (i \in (u \to \mathbb{C})) \wedge \\ \quad \forall y \in u. \\ \quad \quad \operatorname{is\text{-}continuous\text{-}at}(\operatorname{subspace\text{-}ms}(\operatorname{nf\text{-}metric\text{-}space}(\operatorname{cc\text{-}normed\text{-}field}), u), \\ \quad \quad \quad \operatorname{nf\text{-}metric\text{-}space}(\operatorname{cc\text{-}normed\text{-}field}), i, y) \wedge \\ \quad \operatorname{is\text{-}road}(m, d, a, b) \wedge \\ \quad (\operatorname{trace}(m, a, b) \subseteq u) \wedge \\ \quad \forall \mathit{aiy} \in \operatorname{trace}(m, a, b). \\ \quad \quad (\operatorname{magnitude}(i(\mathit{aiy})) \leq \mathit{aim}) \Rightarrow \\ \quad \quad \operatorname{magnitude}(\operatorname{line\text{-}int}((\lambda w \in u.\; (i(w) \cdot \operatorname{power}(\operatorname{recip}((w - c)), (k + 1)))), m, d, \\ \quad \quad \quad \quad a, b)) \\ \quad \quad \leq\; (\operatorname{road\text{-}length}(d, a, b) \cdot (\mathit{aim} \cdot \operatorname{power}(\operatorname{recip}(r), (k + 1)))) \end{array}}}
\vnbentry{\texttt{cauchy-\allowbreak{}estimate-\allowbreak{}from-\allowbreak{}formula}}{\vnbstmt{\begin{array}{@{}l@{}} \forall u, c, r, i, m, d, a, b.; \forall \mathit{aim} \in \mathbb{R}, p \in \mathbb{C}, k \in \mathbb{N}. \\ \quad (u \subseteq \mathbb{C}) \wedge \\ \quad (c \in \mathbb{C}) \wedge \\ \quad (r \in \mathbb{R}) \wedge \\ \quad (0 < r) \wedge \\ \quad \forall x \in u.\; (r \leq \operatorname{magnitude}((x - c))) \wedge \\ \quad (i \in (u \to \mathbb{C})) \wedge \\ \quad \forall y \in u. \\ \quad \quad \operatorname{is\text{-}continuous\text{-}at}(\operatorname{subspace\text{-}ms}(\operatorname{nf\text{-}metric\text{-}space}(\operatorname{cc\text{-}normed\text{-}field}), u), \\ \quad \quad \quad \operatorname{nf\text{-}metric\text{-}space}(\operatorname{cc\text{-}normed\text{-}field}), i, y) \wedge \\ \quad \operatorname{is\text{-}road}(m, d, a, b) \wedge \\ \quad (\operatorname{trace}(m, a, b) \subseteq u) \wedge \\ \quad \forall \mathit{aiy} \in \operatorname{trace}(m, a, b). \\ \quad \quad (\operatorname{magnitude}(i(\mathit{aiy})) \leq \mathit{aim}) \Rightarrow \\ \quad \quad \operatorname{magnitude}((p \cdot \operatorname{line\text{-}int}((\lambda w \in u.\; (i(w) \cdot \operatorname{power}(\operatorname{recip}((w - c)), (k + 1)))), m, d, a, b))) \\ \quad \quad \leq\; (\operatorname{magnitude}(p) \cdot (\operatorname{road\text{-}length}(d, a, b) \cdot (\mathit{aim} \cdot \operatorname{power}(\operatorname{recip}(r), (k + 1))))) \end{array}}}
\vnbentry{\texttt{cauchy-\allowbreak{}kernel-\allowbreak{}series-\allowbreak{}absolute}}{\vnbstmt{\begin{array}{@{}l@{}} \forall u, c, r, i, m, d, a, b.; \forall z \in \mathbb{C}. \\ \quad (u \subseteq \mathbb{C}) \wedge \\ \quad (c \in \mathbb{C}) \wedge \\ \quad (r \in \mathbb{R}) \wedge \\ \quad (0 < r) \wedge \\ \quad \forall x \in u.\; (r \leq \operatorname{magnitude}((x - c))) \wedge \\ \quad (i \in (u \to \mathbb{C})) \wedge \\ \quad \forall y \in u. \\ \quad \quad \operatorname{is\text{-}continuous\text{-}at}(\operatorname{subspace\text{-}ms}(\operatorname{nf\text{-}metric\text{-}space}(\operatorname{cc\text{-}normed\text{-}field}), u), \\ \quad \quad \quad \operatorname{nf\text{-}metric\text{-}space}(\operatorname{cc\text{-}normed\text{-}field}), i, y) \wedge \\ \quad \operatorname{is\text{-}road}(m, d, a, b) \wedge \\ \quad (\operatorname{trace}(m, a, b) \subseteq u) \wedge \\ \quad (\operatorname{magnitude}((z - c)) < r) \Rightarrow \\ \quad \quad \operatorname{series\text{-}converges}((\lambda k \in \mathbb{N}.\; \\ \quad \quad \quad \quad \operatorname{magnitude}((\operatorname{line\text{-}int}((\lambda w \in u.\; (i(w) \cdot \operatorname{power}(\operatorname{recip}((w - c)), (k + 1)))), m, d, a, b) \cdot \operatorname{power}((z - c), k))))) \end{array}}}
\vnbentry{\texttt{analytic-\allowbreak{}at-\allowbreak{}from-\allowbreak{}integral-\allowbreak{}formula}}{\vnbstmt{\begin{array}{@{}l@{}} \forall u, c, r, i, m, d, a, b.; \forall p \in \mathbb{C}, \mathit{aiou}.; \forall f \in (\mathit{aiou} \to \mathbb{C}). \\ \quad (u \subseteq \mathbb{C}) \wedge \\ \quad (c \in \mathbb{C}) \wedge \\ \quad (r \in \mathbb{R}) \wedge \\ \quad (0 < r) \wedge \\ \quad \forall x \in u.\; (r \leq \operatorname{magnitude}((x - c))) \wedge \\ \quad (i \in (u \to \mathbb{C})) \wedge \\ \quad \forall y \in u. \\ \quad \quad \operatorname{is\text{-}continuous\text{-}at}(\operatorname{subspace\text{-}ms}(\operatorname{nf\text{-}metric\text{-}space}(\operatorname{cc\text{-}normed\text{-}field}), u), \\ \quad \quad \quad \operatorname{nf\text{-}metric\text{-}space}(\operatorname{cc\text{-}normed\text{-}field}), i, y) \wedge \\ \quad \operatorname{is\text{-}road}(m, d, a, b) \wedge \\ \quad (\operatorname{trace}(m, a, b) \subseteq u) \wedge \\ \quad \operatorname{is\text{-}open}(\operatorname{cc\text{-}ms}, \mathit{aiou}) \wedge \\ \quad (\operatorname{ball}(\operatorname{cc\text{-}ms}, c, r) \subseteq \mathit{aiou}) \wedge \\ \quad \forall z \in \mathbb{C}. \\ \quad \quad (\operatorname{magnitude}((z - c)) < r) \Rightarrow  \\ \quad \quad \quad f(z) \\ \quad \quad \quad =\; (p \cdot \operatorname{line\text{-}int}((\lambda w \in u.\; (i(w) \cdot \operatorname{recip}((w - z)))), m, d, a, b)) \Rightarrow \\ \quad \quad f \text{ is analytic at } c \end{array}}}

\vnbfile{structure-library/functoriality.scm}{structure-library/functoriality.scm}
\noindent{\small functoriality.scm -- every view-as is a FUNCTOR, and here is the proof.}\par
\vnbentry{\texttt{abelian-\allowbreak{}group-\allowbreak{}as-\allowbreak{}monoid-\allowbreak{}functorial}}{\vnbstmt{\begin{array}{@{}l@{}} \forall a, b, f. \\ \quad \operatorname{is\text{-}hom\text{-}abelian\text{-}group}(a, b, f) \Rightarrow  \\ \quad \quad \operatorname{is\text{-}hom\text{-}monoid}(\operatorname{abelian\text{-}group\text{-}as\text{-}monoid}(a), \operatorname{abelian\text{-}group\text{-}as\text{-}monoid}(b), f) \end{array}}}
\vnbentry{\texttt{commutative-\allowbreak{}ring-\allowbreak{}additive-\allowbreak{}ag-\allowbreak{}functorial}}{\vnbstmt{\begin{array}{@{}l@{}} \forall a, b, f. \\ \quad \operatorname{is\text{-}hom\text{-}commutative\text{-}ring}(a, b, f) \Rightarrow  \\ \quad \quad \operatorname{is\text{-}hom\text{-}abelian\text{-}group}(\operatorname{commutative\text{-}ring\text{-}additive\text{-}ag}(a), \operatorname{commutative\text{-}ring\text{-}additive\text{-}ag}(b), \\ \quad \quad \quad f) \end{array}}}
\vnbentry{\texttt{commutative-\allowbreak{}ring-\allowbreak{}multiplicative-\allowbreak{}cm-\allowbreak{}functorial}}{\vnbstmt{\begin{array}{@{}l@{}} \forall a, b, f. \\ \quad \operatorname{is\text{-}hom\text{-}commutative\text{-}ring}(a, b, f) \Rightarrow  \\ \quad \quad \operatorname{is\text{-}hom\text{-}comm\text{-}monoid}(\operatorname{commutative\text{-}ring\text{-}multiplicative\text{-}cm}(a), \\ \quad \quad \quad \operatorname{commutative\text{-}ring\text{-}multiplicative\text{-}cm}(b), f) \end{array}}}
\vnbentry{\texttt{complex-\allowbreak{}inner-\allowbreak{}product-\allowbreak{}space-\allowbreak{}as-\allowbreak{}module-\allowbreak{}functorial}}{\vnbstmt{\begin{array}{@{}l@{}} \forall a, b, f. \\ \quad \operatorname{is\text{-}hom\text{-}complex\text{-}inner\text{-}product\text{-}space}(a, b, f) \Rightarrow  \\ \quad \quad \operatorname{is\text{-}hom\text{-}module}(\operatorname{complex\text{-}inner\text{-}product\text{-}space\text{-}as\text{-}module}(a), \\ \quad \quad \quad \operatorname{complex\text{-}inner\text{-}product\text{-}space\text{-}as\text{-}module}(b), f) \end{array}}}
\vnbentry{\texttt{field-\allowbreak{}additive-\allowbreak{}ag-\allowbreak{}functorial}}{\vnbstmt{\begin{array}{@{}l@{}} \forall a, b, f. \\ \quad \operatorname{is\text{-}hom\text{-}field}(a, b, f) \Rightarrow  \\ \quad \quad \operatorname{is\text{-}hom\text{-}abelian\text{-}group}(\operatorname{field\text{-}additive\text{-}ag}(a), \operatorname{field\text{-}additive\text{-}ag}(b), f) \end{array}}}
\vnbentry{\texttt{field-\allowbreak{}as-\allowbreak{}euclidean-\allowbreak{}ring-\allowbreak{}functorial}}{\vnbstmt{\begin{array}{@{}l@{}} \forall a, b, f. \\ \quad \operatorname{is\text{-}hom\text{-}field}(a, b, f) \Rightarrow  \\ \quad \quad \operatorname{is\text{-}hom\text{-}euclidean\text{-}ring}(\operatorname{field\text{-}as\text{-}euclidean\text{-}ring}(a), \operatorname{field\text{-}as\text{-}euclidean\text{-}ring}(b), f) \end{array}}}
\vnbentry{\texttt{field-\allowbreak{}as-\allowbreak{}integral-\allowbreak{}domain-\allowbreak{}functorial}}{\vnbstmt{\begin{array}{@{}l@{}} \forall a, b, f. \\ \quad \operatorname{is\text{-}hom\text{-}field}(a, b, f) \Rightarrow  \\ \quad \quad \operatorname{is\text{-}hom\text{-}integral\text{-}domain}(\operatorname{field\text{-}as\text{-}integral\text{-}domain}(a), \operatorname{field\text{-}as\text{-}integral\text{-}domain}(b), f) \end{array}}}
\vnbentry{\texttt{module-\allowbreak{}vector-\allowbreak{}ag-\allowbreak{}functorial}}{\vnbstmt{\begin{array}{@{}l@{}} \forall a, b, f. \\ \quad \operatorname{is\text{-}hom\text{-}module}(a, b, f) \Rightarrow  \\ \quad \quad \operatorname{is\text{-}hom\text{-}abelian\text{-}group}(\operatorname{module\text{-}vector\text{-}ag}(a), \operatorname{module\text{-}vector\text{-}ag}(b), f) \end{array}}}
\vnbentry{\texttt{normed-\allowbreak{}ag-\allowbreak{}as-\allowbreak{}abelian-\allowbreak{}group-\allowbreak{}functorial}}{\vnbstmt{\begin{array}{@{}l@{}} \forall a, b, f. \\ \quad \operatorname{is\text{-}hom\text{-}normed\text{-}ag}(a, b, f) \Rightarrow  \\ \quad \quad \operatorname{is\text{-}hom\text{-}abelian\text{-}group}(\operatorname{normed\text{-}ag\text{-}as\text{-}abelian\text{-}group}(a), \operatorname{normed\text{-}ag\text{-}as\text{-}abelian\text{-}group}(b), f) \end{array}}}
\vnbentry{\texttt{normed-\allowbreak{}field-\allowbreak{}additive-\allowbreak{}ag-\allowbreak{}functorial}}{\vnbstmt{\begin{array}{@{}l@{}} \forall a, b, f. \\ \quad \operatorname{is\text{-}hom\text{-}normed\text{-}field}(a, b, f) \Rightarrow  \\ \quad \quad \operatorname{is\text{-}hom\text{-}abelian\text{-}group}(\operatorname{normed\text{-}field\text{-}additive\text{-}ag}(a), \operatorname{normed\text{-}field\text{-}additive\text{-}ag}(b), f) \end{array}}}
\vnbentry{\texttt{normed-\allowbreak{}field-\allowbreak{}as-\allowbreak{}commutative-\allowbreak{}ring-\allowbreak{}functorial}}{\vnbstmt{\begin{array}{@{}l@{}} \forall a, b, f. \\ \quad \operatorname{is\text{-}hom\text{-}normed\text{-}field}(a, b, f) \Rightarrow  \\ \quad \quad \operatorname{is\text{-}hom\text{-}commutative\text{-}ring}(\operatorname{normed\text{-}field\text{-}as\text{-}commutative\text{-}ring}(a), \\ \quad \quad \quad \operatorname{normed\text{-}field\text{-}as\text{-}commutative\text{-}ring}(b), f) \end{array}}}
\vnbentry{\texttt{normed-\allowbreak{}field-\allowbreak{}as-\allowbreak{}integral-\allowbreak{}domain-\allowbreak{}functorial}}{\vnbstmt{\begin{array}{@{}l@{}} \forall a, b, f. \\ \quad \operatorname{is\text{-}hom\text{-}normed\text{-}field}(a, b, f) \Rightarrow  \\ \quad \quad \operatorname{is\text{-}hom\text{-}integral\text{-}domain}(\operatorname{normed\text{-}field\text{-}as\text{-}integral\text{-}domain}(a), \\ \quad \quad \quad \operatorname{normed\text{-}field\text{-}as\text{-}integral\text{-}domain}(b), f) \end{array}}}
\vnbentry{\texttt{normed-\allowbreak{}vector-\allowbreak{}space-\allowbreak{}as-\allowbreak{}module-\allowbreak{}functorial}}{\vnbstmt{\begin{array}{@{}l@{}} \forall a, b, f. \\ \quad \operatorname{is\text{-}hom\text{-}normed\text{-}vector\text{-}space}(a, b, f) \Rightarrow  \\ \quad \quad \operatorname{is\text{-}hom\text{-}module}(\operatorname{normed\text{-}vector\text{-}space\text{-}as\text{-}module}(a), \operatorname{normed\text{-}vector\text{-}space\text{-}as\text{-}module}(b), \\ \quad \quad \quad f) \end{array}}}
\vnbentry{\texttt{normed-\allowbreak{}vector-\allowbreak{}space-\allowbreak{}as-\allowbreak{}normed-\allowbreak{}ag-\allowbreak{}functorial}}{\vnbstmt{\begin{array}{@{}l@{}} \forall a, b, f. \\ \quad \operatorname{is\text{-}hom\text{-}normed\text{-}vector\text{-}space}(a, b, f) \Rightarrow  \\ \quad \quad \operatorname{is\text{-}hom\text{-}normed\text{-}ag}(\operatorname{normed\text{-}vector\text{-}space\text{-}as\text{-}normed\text{-}ag}(a), \\ \quad \quad \quad \operatorname{normed\text{-}vector\text{-}space\text{-}as\text{-}normed\text{-}ag}(b), f) \end{array}}}
\vnbentry{\texttt{ring-\allowbreak{}additive-\allowbreak{}ag-\allowbreak{}functorial}}{\vnbstmt{\begin{array}{@{}l@{}} \forall a, b, f. \\ \quad \operatorname{is\text{-}hom\text{-}ring}(a, b, f) \Rightarrow  \\ \quad \quad \operatorname{is\text{-}hom\text{-}abelian\text{-}group}(\operatorname{ring\text{-}additive\text{-}ag}(a), \operatorname{ring\text{-}additive\text{-}ag}(b), f) \end{array}}}
\vnbentry{\texttt{ring-\allowbreak{}multiplicative-\allowbreak{}monoid-\allowbreak{}functorial}}{\vnbstmt{\begin{array}{@{}l@{}} \forall a, b, f. \\ \quad \operatorname{is\text{-}hom\text{-}ring}(a, b, f) \Rightarrow  \\ \quad \quad \operatorname{is\text{-}hom\text{-}monoid}(\operatorname{ring\text{-}multiplicative\text{-}monoid}(a), \operatorname{ring\text{-}multiplicative\text{-}monoid}(b), f) \end{array}}}



\end{document}
